% Generated body; only mechanical merging changes.



\begin{abstract}
We determine the thermal scaling dimension of the long-range Ising model in the
expansion parameter $\epsilon=2s-d$. The analytic result through $C_5(d)$ is
expressed in an explicitly defined master basis, with its graph weights,
subtraction prescription, regulator expansions and finite special-function
representations derived below. At $d=1$, the result uses only fixed
Gamma/polygamma values and ordinary generalized hypergeometric values;
the proof of the $\mathcal K$ reduction, the logarithmic-wheel reduction and
the five-value evaluation of $P_5(1)$ are included in full. At $d=2$ the
primitive periods factor in the same scalar values, and the nonprimitive
masters $\mathcal K_2$, $\mathcal L_2$ and $Z_2$ are evaluated in closed form, so that $C_5(2)$
is an explicit expression in those values, $\Gamma(\frac14)$, $\pi$, $\log2$, Catalan's
constant and $\zeta(3)$. At $d=3$ the remaining
rank-two structure is identified without claiming an unavailable scalar
summation. For generic $0<d<4$ the masters $\mathcal L_d$ and $Z_d$ are reduced
to at most three samples of finitely many hypergeometric functions with at most
two summation indices at unit argument and their first and second derivatives
in the rim exponent, using a proof that the top Fourier modes of products of
balanced kernels vanish; $\mathcal K_d$ is reduced, under stated analytic
inputs, to one Laurent coefficient of such a sample, and, by a second route,
to a polygamma part plus one two-fold Mellin--Barnes integral. A finite representation of a reduced analytic expression in nine
special-function values gives $C_6(1)=14331.68703275805$; exact sums over all
six-loop graphs at increasing quadrature degree, which do not use that expression,
give $14331.6870328(5)$, in agreement to $1.0\times10^{-8}$. Numerical finite-forest calculations
give $C_6(2)=362.366\pm0.004$, $C_6(3)=33.5754\pm0.0005$,
$C_7(1)=-322550\pm41$, $C_7(2)=-3876.1\pm2.1$,
$C_7(3)=-220.35\pm0.13$, $C_8(1)=(8.060\pm0.017)\times10^6$,
$C_8(2)=45990\pm33$ and $C_8(3)=1602.6\pm3.9$.
The integration kernels, regulator-dependent fixed-point extraction, local
subtractions, sampling designs, covariance propagation and convergence
checks supporting these estimates are retained. Numerical allowances are
empirical, not rigorous intervals or bounds on truncating the epsilon
expansion. Part~V records a second observable built from the same graphs and
master functions, the bulk one-point coefficient $a_{\phi^2}$ of the auxiliary
bulk field through $\epsilon^6$ at general $d$, from a calculation supplied by
the author whose checks are internal to this framework. This note records one
common graph/renormalization calculation;
it is not an independent field-theoretic certification of its higher-loop
results.
\end{abstract}

\clearpage
\SourcePart{Conclusions, conventions and renormalization}
\section{Scope and current coefficients}\label{gff:sec:guide}
The model is the long-range Ising model: spins $\sigma_i=\pm1$ on the sites of
$\mathbb Z^d$ with ferromagnetic energy
\[
 \mathcal H_{\rm lat}=-\sum_{i<j}J(|i-j|)\,\sigma_i\sigma_j,\qquad
 J(r)\sim r^{-(d+s)},
\]
in the critical long-range regime above $s=d/2$. Its continuum description is
a Gaussian field with covariance proportional to $|p|^{-s}$ in momentum space
and to $|x|^{-(d-s)}$ in position space, normalized to coefficient one, with a
quartic interaction of engineering dimension $2s-d=\epsilon$.
The observable is the scaling dimension of the thermal operator
$\mathcal E=[\phi^2]$:
\begin{equation}
 \epsilon=2s-d,\qquad
 \Delta_\phi=\frac{d-\epsilon}{4},\qquad
 \Delta_{\mathcal E}(d,\epsilon)=\sum_{n\ge0}C_n(d)\epsilon^n.
 \label{gff:cn:expansion}
\end{equation}
The general-dimensional long-range regime is fixed $0<d<4$ and sufficiently
small positive $\epsilon$; noninteger dimensions and intermediate meromorphic
integrals have the continuation prescriptions stated in the relevant
sections. This is a local expansion near $s=d/2$, not an unresummed
approximation over the entire interacting regime.

\begin{center}
\begin{tabular}{@{}c r r r@{}}
\toprule
Order & $d=1$ & $d=2$ & $d=3$\\\midrule
$C_0$ & $1/2$ & $1$ & $3/2$\\
$C_1$ & $-1/6$ & $-1/6$ & $-1/6$\\
$C_2$ & $0.657131263981$ & $0.308065413582$ & $0.181221785406$\\
$C_3$ & $-3.719856322902$ & $-0.874850012653$ & $-0.353581520433$\\
$C_4$ & $45.42488010356$ & $5.048868431$ & $1.228063545$\\
$C_5$ & $-731.13947662454$ & $-38.833645896$ & $-5.850351421$\\
$C_6$ & $14331.68703275805$ & $362.366\pm0.004$ & $33.5754\pm0.0005$\\
$C_7$ & $-322550\pm41$ & $-3876.1\pm2.1$ & $-220.35\pm0.13$\\
$C_8$ & $(8.060\pm0.017)10^6$ & $45990\pm33$ & $1602.6\pm3.9$\\
\bottomrule
\end{tabular}
\end{center}
The first six rows use the special-function evaluations, not the noisier
lower-order checks produced during high-loop integration; their closed forms
are collected in Section~\ref{gff:c1:closedForms}. The entry for
$C_6(1)$ is the finite special-function representation of
Section~\ref{gff:c1:c6analytic}. It evaluates a reduced analytic expression that
is taken as an input and is not derived in this note
(Section~\ref{gff:c1:status}). The independent deterministic sum over all graphs at
quadrature degree 22 (Section~\ref{gff:app:C6Numerical}) gives $14331.6870328(5)$. It has
no sampling error; its uncertainty is the empirical convergence-and-rounding
allowance of Section~\ref{gff:c6:uncertainty}, and the representation lies
$1.0\times10^{-8}$ below its unrounded value $14331.68703276838$. For the
remaining entries of the last three rows, the allowances above and the
sampling standard errors below are \emph{different} quantities:
\begin{center}
\begin{tabular}{@{}l r r@{}}
\toprule
Coefficient & Underlying estimator & Sampling standard error\\\midrule
$C_6(2)$ & $362.36647176375726$ & $0.0007997307$\\
$C_6(3)$ & $33.57542664131543$ & $0.0001185495235$\\
$C_7(1)$ & $-322549.8888129889$ & $10.08791295$\\
$C_7(2)$ & $-3876.1260927949925$ & $0.5244496691$\\
$C_7(3)$ & $-220.35448047392794$ & $0.0318188006$\\
$C_8(1)$ & $8059766.723129955$ & $4209.576151$\\
$C_8(2)$ & $45990.00263614446$ & $8.216820958$\\
$C_8(3)$ & $1602.6052229663408$ & $0.9516489064$\\
\bottomrule
\end{tabular}
\end{center}
Each allowance is the output of a stated rule applied to the sampling
standard error and, where they exist, to measured convergence steps. The
table gives the rule for each entry and when it was fixed: before the
production run, during it but before its result was known, or after the
result. ``Four SE'' means four
sampling standard errors; a half-split estimate, called a half-stream estimate
in Sections~\ref{gff:app:C7d2} and~\ref{gff:app:C8d2}, uses only the first or only the
second half of the replicates of every random stream.
\begin{center}\small
\begin{tabular}{@{}l l >{\raggedright\arraybackslash}p{0.44\textwidth} >{\raggedright\arraybackslash}p{0.19\textwidth}@{}}
\toprule
Entry & Allowance & Rule & Fixed\\\midrule
$C_6(1)$, graph sum & $\pm4.9\cdot10^{-7}$ & the last change with quadrature degree, plus the difference of the two cutoffs, a floating-point allowance and the lower-loop convergence (Section~\ref{gff:c6:uncertainty}) & before production\\
$C_6(2)$ & $\pm0.004$ & four SE, rounded up to a multiple of $0.001$ (Section~\ref{gff:prec:d2production2026}) & before the result was known; rounding step refined from $0.01$ afterwards\\
$C_6(3)$ & $\pm0.0005$ & four SE, rounded up to a multiple of $0.0001$ (Section~\ref{gff:prec:d3production2026}) & before production; rounding step refined from $0.001$ afterwards\\
$C_7(1)$ & $\pm41$ & the largest of four SE, twice the last refinement step plus four of its SE, twice the degree six to seven check step plus four of its SE, and the half-split displacement, rounded up to two significant figures (Section~\ref{gff:c7f:production}) & before production\\
$C_7(2)$ & $\pm2.1$ & the larger of four SE and the largest half-split deviation, rounded up to a multiple of $0.1$ (Section~\ref{gff:app:C7d2}) & after the result (the rule of an earlier, lower-precision $C_7(2)$ calculation)\\
$C_7(3)$ & $\pm0.13$ & four SE, rounded up to a multiple of $0.01$ (Section~\ref{gff:app:C7d3}) & before the result was known\\
$C_8(1)$ & $\pm17000$ & as for $C_7(1)$, with the degree eight to nine check step, and with twice the last degree step of the exact loop-five input as a fifth candidate (Section~\ref{gff:c8f:result}) & before production; batch count raised afterwards\\
$C_8(2)$ & $\pm33$ & as for $C_7(2)$, rounded up to a multiple of $1$ (Section~\ref{gff:app:C8d2}) & after the result\\
$C_8(3)$ & $\pm3.9$ & four SE, rounded up to a multiple of $0.1$ (Section~\ref{gff:app:C8d3}) & after the result\\
\bottomrule
\end{tabular}
\end{center}
For $C_7(1)$ the four candidates are $40.35$, $39.50$, $22.49$ and $12.58$,
so the allowance is four sampling standard errors. The radial quadrature of
$C_7(1)$ stops at degree six, and the rule adds no model of the truncation
beyond it. Adding the measured degree six to seven step and a geometric tail
at the last step ratio gives the post hoc estimate
$C_7(1)\approx-322524\pm38$ \eqref{gff:c7f:truncationAware}; its central value
lies inside the quoted interval. That tail model agrees with the
lower-order checks at fifth and sixth order but is not validated at seventh
order, where the step ratio is larger. The result is therefore
$C_7(1)=-322550\pm41$, and the truncation-aware value is a check
(Section~\ref{gff:c7f:checks}). For $C_8(1)$ the rule and the design were fixed
before production, but the number of production batches was not. Under the
stopping rule set first, nine batches give $8061480\pm5300$ with allowance
$\pm22000$; after interim results of five to eight batches had been seen, all
fourteen planned batches were completed (Section~\ref{gff:c8f:production}). At
degree eight the radial quadrature is converged well inside the allowance.
For $C_8(2)$ and $C_8(3)$ four sampling standard errors exceed the distances
of both half-split estimates from the full estimate. The coverage of the
allowance on $C_8(3)$ is conditional on a finite variance, for which
Section~\ref{gff:c8d3:tails} gives empirical evidence, not a bound.

Two eighth-order results are checked by a separately developed code. At
$d=2$ it shares the method and the per-graph input files of loops three to seven and gives
$C_8(2)=46006.1\pm13.4$ (sampling standard error), $+1.03$ combined standard
errors from the estimator (Section~\ref{gff:c8d2:independent}). At $d=3$ it
shares the method and the seventh-order integrator and gives
$C_8(3)=1600.98\pm3.97$, $-0.40$ combined standard errors from the estimator
(Section~\ref{gff:c8d3:independent}).
The additional estimator digits are not
certified physical digits. Numerical
errors do not cover shared mistakes in the graph/renormalization framework,
and are not bounds on the omitted orders in \eqref{gff:cn:expansion}.

Part~V (Section~\ref{gff:ope:section}) adds a second observable built from the same
graphs and master functions: the bulk one-point coefficient $a_{\phi^2}$ of the
composite $:\!\Phi^2\!:$ of the auxiliary bulk field, through $\epsilon^6$ at
fixed $0<d<4$, which is the horizontal coordinate of the bootstrap bounds of
\cite{gff:BLNV}. It comes from a calculation supplied by the author with its own
code. Its checks are internal to the master framework of this note, together
with comparisons against this note's own inputs (Section~\ref{gff:ope:checks});
Section~\ref{gff:ope:status} states what it establishes.

\subsection{Closed forms through fifth order}\label{gff:c1:closedForms}
The entries $C_0$ to $C_5$ of the table are values of exact expressions, not
statistical estimates. This subsection states the expressions together with
every definition they need; each formula below repeats one derived later in
the note, and symbols with more than one meaning in the note are listed in
Section~\ref{gff:c1:notation}. At $d=1$ every step is proved in this note
(Part~II). At $d=2$ the closed form of $C_5(2)$ is proved in this
note. At $d=3$ the inputs of $C_5(3)$, namely $Q_3$,
$P_4(3)$, $P_5(3)$, $\LL_3$, $Z_3$ and $\KK_3$, are evaluated numerically. The
values of $Q_3$, $P_4(3)$ and $P_5(3)$ from their fixed formulas agree to $10^{-74}$ or better
between evaluations at 90 and 100 working digits. Independent evaluations confirm $Q_3$ and $P_4(3)$
to about $10^{-79}$ and $2\cdot10^{-86}$ (a $100$-digit spectral evaluation), $P_5(3)$ to about
$2\cdot10^{-53}$ (a $60$-digit integration of the Gamma kernels), and $\LL_3$, $Z_3$ and $\KK_3$ to
$1.2\cdot10^{-93}$ or better. Propagated into $C_5(3)$ these differences are below $10^{-54}$, well
below one unit in the last of the 50 digits printed for it. These agreements are not interval
bounds. The spectral formula for $\KK_d$ below is proved at $d=2$ in
Section~\ref{gff:mgk:section}; at other
$d$ its derivation uses input~(M) of Section~\ref{gff:gdk:conventions} and the extraction of $[\tau^2]$ under the
trace, which is proved only at $d=2$ (Section~\ref{gff:gdk:status}). The value
of $\KK_3$ does not depend on (M): the two-fold Barnes representation
\eqref{gff:gdk:KMB}, derived without it, reproduces it to $4\cdot10^{-100}$.
Section~\ref{gff:c1:status} lists what is proved, what is conditional and what
is numerical.

\paragraph{Any dimension.} With $q=d/4$, the polygamma functions
$\psi_m(z)=d^{m+1}\log\Gamma(z)/dz^{m+1}$ and the combinations
\begin{equation*}\begin{aligned}
 A&=\psi_0(2q)+\psi_0(1)-2\psi_0(q),&
 R&=\psi_1(2q)-\psi_1(1),\\
 U&=\psi_2(q),& V&=\psi_2(1)+\psi_2(2q),&
 X&=\psi_3(2q)-\psi_3(1).
\end{aligned}\end{equation*}
the coefficients are
\begin{equation*}\boxed{\begin{aligned}
 C_0(d)&=\frac d2,&C_1(d)&=-\frac16,&C_2(d)&=\frac A9,\\
 C_3(d)&=\frac{A^2+R-8Q}{108},\\
 648C_4(d)&=16P_4-A^3-8AQ-3AR+\frac{53}{2}U-3V.
\end{aligned}}\end{equation*}
\begin{equation*}\boxed{\begin{aligned}
46656 C_5(d)={}&-1536\bigl(P_5+\mathcal K_d\bigr)
       +384AP_4-576\mathcal L_d-1728Q^2\\
 &+720A^2Q+1008QR+7152Z_d\\
 &-15A^4-90A^2R+414AU-108AV-117R^2-75X.
\end{aligned}}\end{equation*}
Here $Q_d$, $P_4(d)$ and $P_5(d)$, written $Q$, $P_4$ and $P_5$ in the boxed
formulas, are the primitive (zig-zag) periods with three, four and five loops, normalized by $h_d^{L-1}$ with $h_d=\Gamma(d/2)/\pi^{d/2}$,
and $\KK_d$, $\LL_d$ and $Z_d$ are nonprimitive master data. All of them are
spectral integrals. With
\begin{equation*}\begin{aligned}
 d_0(d)&=1,&
 d_\ell(d)&=\frac{(2\ell+d-2)\Gamma(\ell+d-2)}
 {\Gamma(d-1)\Gamma(\ell+1)}\quad(\ell\geq1),\\
 a_\ell&=\frac\ell2+\frac q2,&b_\ell&=a_\ell+q,\\
 t_\ell(\nu)&=\Gamma(2q)
 \frac{\Gamma(a_\ell+\ii\nu)\Gamma(a_\ell-\ii\nu)}
      {\Gamma(b_\ell+\ii\nu)\Gamma(b_\ell-\ii\nu)}.
\end{aligned}\end{equation*}
and the trace
\begin{equation*}
 \langle f\rangle_d=\frac1\pi\sum_{\ell=0}^{\infty}d_\ell(d)
                    \int_0^{\infty}f(\ell,\nu)\,d\nu.
\end{equation*}
the four-loop and three-loop periods are
\begin{equation*}
 3Q_d=\langle t_\ell^3\rangle_d,\qquad
 P_4(d)=\langle t_\ell^4\rangle_d.
\end{equation*}
and, with
\begin{equation*}\begin{aligned}
 w_1&=-\psi_0(q)+\frac12\Re\{\psi_0(a_\ell+\ii\nu)+\psi_0(b_\ell+\ii\nu)\},\\
 w_2&=\frac18\Re\{\psi_1(a_\ell+\ii\nu)-\psi_1(b_\ell+\ii\nu)\}.
\end{aligned}\end{equation*}
the two wheel master data are
\begin{equation*}\boxed{
 \mathcal L_d=\langle4t_\ell^4w_1\rangle_d,\qquad
 Z_d=\left\langle t_\ell^3\left(w_1^2+\frac{w_2}{3}\right)\right\rangle_d.
}\end{equation*}
For $\KK_d$ put $c_*=\psi_0(1)+\psi_0(2q)$ and, at $a=a_\ell$ and $b=b_\ell$,
\begin{equation*}\begin{aligned}
 p&=\Re\{\psi_0(a+\ii\nu)+\psi_0(b+\ii\nu)\},&
 h&=\Im\{\psi_0(a+\ii\nu)-\psi_0(b+\ii\nu)\},\\
 j&=\Re\{\psi_1(a+\ii\nu)-\psi_1(b+\ii\nu)\},&
 k&=\Im\{\psi_1(a+\ii\nu)+\psi_1(b+\ii\nu)\},\\
 r&=\Re\{\psi_2(a+\ii\nu)+\psi_2(b+\ii\nu)\},&
 E_2&=-\frac{h^2}{2}-\frac{5j}{8}.
\end{aligned}\end{equation*}
Then
\begin{equation*}\boxed{
 \mathcal K_d=-\frac{AU}{4}
 +\frac14\left\langle t_\ell^2
   \left[(c_*-p)E_2-\frac{hk}{2}-\frac{5r}{32}\right]\right\rangle_d.
}\end{equation*}
The master $\KK_d$ is defined as the finite second-order coefficient of a
subtracted two-loop master graph, and the spectral formula above is its
evaluation. The two-fold Mellin--Barnes integral \eqref{gff:gdk:KMB}, derived in
position space without that formula (Section~\ref{gff:gdk:positionSpace}), agrees
to about $10^{-100}$ with the proved values at $d=1$ \eqref{gff:eq:KprovedFinal}
and $d=2$ \eqref{gff:gc:K2closed} and with the residue series at $d=3$
(Section~\ref{gff:gdk:numerics}). The trace forms converge directly for $\frac23<d<4$ ($Q_d$, $Z_d$),
for $d>\frac12$ ($P_4$, $\LL_d$) and for $0<d<4$ ($\KK_d$); elsewhere they are
continued analytically.

The five-loop primitive needs the two-vector zig-zag spectrum. With
$p=d/2-1$ (not the $p$ of the formula for $\KK_d$),
\begin{align*}
 \tau_n^{(d)}(\nu)&=(-1)^n\pi^{d/2}
 \frac{\Gamma(n/2+d/8+\ii\nu)\Gamma(n/2+d/8-\ii\nu)}
      {\Gamma(n/2+3d/8+\ii\nu)\Gamma(n/2+3d/8-\ii\nu)},
 \\
 \rho_n^{(d)}(\nu)&=
 \frac{\Gamma(p)\Gamma(n+2p)(n+p)}
      {2\pi^{3d/2+1}\Gamma(2p)n!}
 \left[(n+p)^2+4\nu^2\right]
 \frac{\Gamma(p+2\ii\nu)\Gamma(p-2\ii\nu)}
      {\Gamma(2\ii\nu)\Gamma(-2\ii\nu)}.
\end{align*}
the periods with $L=3,4,5$ loops are, for $2<d<4$,
\begin{equation*}\boxed{
 P_L(d)=h_d^{L-1}\sum_{n=0}^{\infty}\int_0^{\infty}
        \rho_n^{(d)}(\nu)\,[\tau_n^{(d)}(\nu)]^{L+2}d\nu,
 \qquad L=3,4,5.
}\end{equation*}
and, for $0<d<2$, where the first pair of Gamma poles crosses the real contour,
\begin{equation*}\begin{aligned}
 P_L(d)=h_d^{L-1}\Bigg\{&
 \sum_{n=0}^{\infty}\int_0^{\infty}
       \rho_n^{(d)}(\nu)[\tau_n^{(d)}(\nu)]^{L+2}d\nu\\
 &+\sum_{n=1}^{\infty}
 \frac{\Gamma(n+2p)(n+p)n(n+2p)}
      {2\pi^{3d/2}n!\Gamma(-p)}
       [\tau_n^{(d)}(-\ii p/2)]^{L+2}\Bigg\}.
\end{aligned}\end{equation*}
At removable singular values such as $d=1,2$ the Gamma ratios are taken in their
joint analytic limit. The masters can be reduced further, to a few two-variable
hypergeometric samples and their low-order parameter jets. For generic $d$ this
is a finite representation, not a reduction to scalar values.

\paragraph{One dimension.} At $d=1$ every coefficient is a fixed value. With
$q=\frac14$, $\bar q=\frac34$ and
\begin{equation*}\boxed{\begin{aligned}
 \pA&=\psi_0(2q)+\psi_0(1)-2\psi_0(q),&
 \pD&=\psi_0(\bar q)-\psi_0(q),\\
 \pT&=\psi_1(q),&\overline{\pT}&=\psi_1(\bar q),\\
 \pU&=\psi_2(q),&\overline{\pU}&=\psi_2(\bar q).
\end{aligned}}\end{equation*}
whose values are, with Catalan's constant
$\Cat=\sum_{n\ge0}(-1)^n/(2n+1)^2$,
\begin{equation*}\begin{aligned}
 \pA&=A, &\pD&=\pi,\\
 \pT+\overline{\pT}&=2\pi^2,&
 \pT-\overline{\pT}&=16\Cat,\\
 \pU+\overline{\pU}&=-112\zeta(3),&
 \overline{\pU}-\pU&=4\pi^3,\\
 H&=\frac{5\pT-\overline{\pT}}6,&
 \log2&=\frac{\pA-\pD}{4}.
\end{aligned}\end{equation*}
and with $Q=\Gamma(\frac14)^4/\pi$, the coefficients are
\begin{equation*}\boxed{\begin{aligned}
 C_0&=\frac12,& C_1&=-\frac16,& C_2&=\frac{\pA}{9},\\[1mm]
 C_3&=\frac{\pA^2+(\pT+\overline{\pT})/6-8Q}{108}.
\end{aligned}}\end{equation*}
\begin{equation*}\boxed{\begin{aligned}
648C_4={}&16P_4-\pA^3-\frac{\pA}{2}(\pT+\overline{\pT})-8\pA Q\\
&+\frac{365\pU-6\overline{\pU}}{14}.
\end{aligned}}\end{equation*}
\begin{equation*}\boxed{\begin{aligned}
46656C_5={}&-1536P_5+96(7\pA-\pD)P_4+384\pD B_- -1728Q^2\\
&+16Q(45\pA^2+383\pT-64\overline{\pT})
 -768\pT(3\pT-\overline{\pT})\\
&-15\pA^4-15\pA^2(\pT+\overline{\pT})\\
&-\frac{6\pA}{7}(431\pU+18\overline{\pU})
 -\frac{83}{4}(\pT+\overline{\pT})^2.
\end{aligned}}\end{equation*}
The periods are ordinary hypergeometric values:
\[
 P_4=13B_--3B_+,\qquad
 P_5=16\pi^{3/2}\bigl[-35X_0+(22\sqrt2-13)X_++(22\sqrt2+13)X_-\bigr],
\]
with
\begin{equation*}\begin{aligned}
 B_-&=\Gamma(\tfrac14)^4 F_-,&
 B_+&=\frac{64\pi^4}{\Gamma(\tfrac14)^4}F_+,\\
 X_0&=\mathcal Y_2,&
 X_+&=c_8^7\mathcal Y_3-s_8^7\mathcal Y_5,&
 X_-&=c_8^7\mathcal Y_1+s_8^7\mathcal Y_7.
\end{aligned}\end{equation*}
\begin{equation*}\begin{aligned}
 F_-&={}_4F_3\!\left(\begin{matrix}\frac14,\frac14,\frac12,\frac12\\
          \frac34,\frac34,1\end{matrix};1\right),&
 F_+&={}_4F_3\!\left(\begin{matrix}\frac34,\frac34,\frac12,\frac12\\
          \frac54,\frac54,1\end{matrix};1\right),\\
 c_8&=\cos(\pi/8),&s_8&=\sin(\pi/8),
\end{aligned}\end{equation*}
\begin{equation*}\begin{aligned}
\mathcal Y_r={}&\frac r8
 \left[\frac{\Gamma((r+2)/8)}{\Gamma((r+6)/8)}\right]^7
 {}_9F_8\!\left(
 \begin{matrix}1,1+\frac r8,((r+2)/8)^{\times7}\\
 \frac r8,((r+6)/8)^{\times7}\end{matrix};-1\right),\\
&r\in\{1,2,3,5,7\}.
\end{aligned}\end{equation*}
Here $(a)^{\times n}$ stands for $n$ copies of the parameter $a$, and ${}_pF_q$
is the series of Section~\ref{gff:c1:notation}. At $r=2$ one upper and one lower
parameter $1$ cancel, so $\mathcal Y_2$ is an ${}_8F_7$ value. No elementary evaluation or minimality of these values is
asserted.

\paragraph{Two dimensions.} With $a=\log2$,
\begin{equation*}\boxed{
 Q_2=\frac16\left[\frac{\Gamma(1/4)}{\Gamma(3/4)}\right]^4
     =\frac{\Gamma(1/4)^8}{24\pi^4}.
}\end{equation*}
\begin{equation*}\boxed{\begin{aligned}
 C_0(2)&=1,&C_1(2)&=-\frac16,&C_2(2)&=\frac{4a}{9},\\
 C_3(2)&=\frac{4a^2}{27}-\frac{\Gamma(1/4)^8}{324\pi^4},\\
 648C_4(2)&=16P_4(2)-64a^3-32aQ_2-359\zeta(3),\\
 46656C_5(2)&=-1536\{P_5(2)+\mathcal K_2\}
 +1536aP_4(2)-576\mathcal L_2-1728Q_2^2\\
 &\quad +(11520a^2+28608\Cat)Q_2
 -3840a^4-21456a\zeta(3).
\end{aligned}}\end{equation*}
The two nonprimitive masters are $\LL_2=-2\log2\,P_4(2)$, proved in
Section~\ref{gff:gc:L2proof} for the series definition of $\LL_2$, and
$\KK_2=32\Cat^2-28\zeta(3)\log2$, proved in Section~\ref{gff:gc:K2proof} for its
master-path definition. The formula for $C_5(2)$ below also uses
$Z_2=4\Cat Q_2$ \eqref{gff:gs:Z2exact}. The primitive periods reuse the
one-dimensional values above (proved in Section~\ref{gff:fp:twoDimensional}):
\begin{equation*}\boxed{
 P_4(2)=\frac{B_-^2-B_+^2}{4\pi^3}.
}\end{equation*}
\begin{equation*}\boxed{\begin{aligned}
P_5(2)=\frac8\pi\bigl[&25X_0^2-28\sqrt2\,X_0(X_++X_-)\\
 &+18(X_+^2+X_-^2)+28X_+X_-\bigr].
\end{aligned}}\end{equation*}
Hence
\begin{equation*}\boxed{\begin{aligned}
 972\,C_5(2)={}&-32P_5(2)+56\log2\,P_4(2)-36Q_2^2
 +\bigl(240\log^22+596\Cat\bigr)Q_2\\
 &-1024\Cat^2+449\,\zeta(3)\log2-80\log^42 .
\end{aligned}}\end{equation*}
so $C_5(2)$ needs only $\Gamma(\frac14)$, $\pi$, $\log2$, $\Cat$, $\zeta(3)$,
$F_\pm$ and the five values $\mathcal Y_r$.

\paragraph{Three dimensions.} At $q=\frac34$ the polygamma combinations give
$A=2+4\log2-\pi$ and $R=\frac{\pi^2}3-4$, so $C_2(3)=(2+4\log2-\pi)/9$. The
higher coefficients follow from the formulas for any dimension with the $d=3$
values of the masters, which reduce to fixed two-variable hypergeometric
functions. No reduction of them to scalar ordinary ${}_pF_q$ values is known or
asserted. Evaluated this way,
$C_5(3)=-5.8503514207513272292179442304977443001109973463783\ldots$ to the $50$
digits shown.

\subsection{A finite special-function representation of $C_6(1)$}\label{gff:c1:c6analytic}
The value of $C_6(1)$ in the table is a finite formula in nine
special-function values. It comes from an analytic calculation in the package
\path{C6_research_note/}. That calculation takes a reduced analytic
expression for the six-loop coefficient as its input and converts each of
its sectors into one value at unit argument; Appendix~\ref{gff:app:C6derivation}
gives the conversion. With $A=\pi+4\log2$,
$Q=\Gamma(\frac14)^4/\pi$ and Catalan's constant $\Cat$,
\begin{equation}
\begin{aligned}
 2187\,C_6(1)={}&E-\frac{476}{3}\pi^3(32\Cat+9\pi^2)
 -2176\bigl[(32\Cat+9\pi^2)L_3+128\,C_{212}\bigr]\\
 &-4\pi^2\rho_4\,\mathcal G_4+192\pi\sqrt{290}\,\mathcal G_6
 +16\pi^3\rho_8\,\mathcal G_8-393\pi\sqrt{89}\,\mathcal H_W\\
 &+24\pi^{3/2}\bigl(\rho_7\,\mathcal G_7+196\pi\sqrt{185}\,\mathcal H_7\bigr)
 +7680\pi^3\,\mathcal H_*,
\end{aligned}
\label{gff:c1:c6nine}
\end{equation}
where $E$ is the polynomial
\begin{align*}
 E={}&\frac Q8\left[63A^3+120A\Cat+79A\pi^2+2013\pi^3+54204\zeta(3)\right]
 +Q^2(29\pi-249A)
 +\frac{21}{640}A^5+\frac7{64}A^3\pi^2\\
 &+\frac{1149}{64}A^2\pi^3+\frac{7863}{16}A^2\zeta(3)
 +192A\Cat(16\Cat+3\pi^2)+\frac{14853}{640}A\pi^4
 -\frac{544}{3}\Cat\pi^3\\
 &+\frac{37631}{128}\pi^5+\frac{722279}{48}\pi^2\zeta(3)
 -\frac{805959}{8}\zeta(5).
\end{align*}
The nine values are the following.
\begin{itemize}
\item $L_3=\operatorname{Ls}_3(\frac\pi2)$ and
 $C_{212}=\operatorname{Cl}_{2,1,2}(\frac\pi2)$, a log-sine integral and a
 multiple Clausen value:
 \[
  \operatorname{Ls}_3(\sigma)=-\int_0^\sigma\log^2\bigl|2\sin\tfrac\theta2\bigr|\dd\theta,\qquad
  \operatorname{Cl}_{2,1,2}(\sigma)=\sum_{n>m>k\ge1}\frac{\cos(n\sigma)}{n^2mk^2}.
 \]
\item Four Meijer $G$ values at unit argument, in the convention of
 \cite{gff:DLMF16}:
 \begin{gather*}
  \mathcal G_4=G^{2,3}_{5,5}\!\left(1\,\middle|\,\begin{matrix}(\frac12)^{\times4},\alpha_4\\
  0^{\times4},\alpha_4\end{matrix}\right),\qquad
  \mathcal G_6=G^{4,5}_{7,7}\!\left(1\,\middle|\,\begin{matrix}(\frac12)^{\times6},\alpha_6\\
  0^{\times6},\alpha_6\end{matrix}\right),\\
  \mathcal G_8=G^{3,8}_{11,11}\!\left(1\,\middle|\,\begin{matrix}-\frac14,(\frac12)^{\times8},
  \alpha_{8,+},\alpha_{8,-}\\0^{\times8},\frac34,\alpha_{8,+},\alpha_{8,-}\end{matrix}\right),\\
  \mathcal G_7=G^{3,8}_{10,10}\!\left(1\,\middle|\,\begin{matrix}
  -\frac14,(\frac12)^{\times7},\eta_{7,+},\eta_{7,-}\\
  0^{\times7},\frac34,\eta_{7,+},\eta_{7,-}\end{matrix}\right),
 \end{gather*}
 with the parameters
 \begin{gather*}
  C_d=96Q-42A^2-18A\pi-1227\pi^2,\qquad
  C_s=60Q-\tfrac{105}{4}A^2-9A\pi-\tfrac{2469}{4}\pi^2,\\
  \rho_4=\sqrt{C_d^2+C_s^2},\qquad
  \alpha_4=\frac1\pi\arctan\frac{C_s}{C_d},\qquad
  \alpha_6=\frac1\pi\arctan\frac{11}{13},\\
  \rho_8=\sqrt{2177203-1028322\sqrt2},\\
  \alpha_{8,\pm}=\frac1{2\pi}\Bigl[\arctan\frac{556}{1233-417\sqrt2}
  \pm\ii\operatorname{arcosh}\frac{1233-139\sqrt2}{\rho_8}\Bigr],\\
  q_{7,2}=259A-49\pi,\qquad q_{7,1}=481A-539\pi,\qquad q_{7,0}=407A-1127\pi,\\
  d_7=\sqrt{q_{7,1}^2-4q_{7,2}q_{7,0}},\qquad
  \rho_7=\sqrt{(q_{7,2}-q_{7,0})^2+q_{7,1}^2},\qquad
  \eta_{7,\pm}=\frac1\pi\arctan\frac{q_{7,1}\pm d_7}{2q_{7,2}},
 \end{gather*}
 with principal values throughout.
\item Three bivariate $H$ values at unit arguments, $\mathcal H_W$,
 $\mathcal H_7$ and $\mathcal H_*$, that is, two-fold Mellin--Barnes integrals.
 The first is
\begin{gather*}
 \mathcal H_W=\frac1{(2\pi\ii)^2}\int_{\frac14-\ii\infty}^{\frac14+\ii\infty}
 \int_{\frac18-\ii\infty}^{\frac18+\ii\infty}
 \frac{N(s)\,\Gamma(t)\Gamma(t+\frac12)\Gamma(\frac12-t)\Gamma(1-t)}
 {(\frac12-s-t)^2}\dd t\dd s,\\
 N(s)=\frac{\Gamma(s)^2\Gamma(\frac12-s)^3}
 {\Gamma(1-s)^2\Gamma(\frac12+s)\Gamma(s+\alpha_W)\Gamma(1-s-\alpha_W)},
\end{gather*}
 with $\alpha_W=\pi^{-1}\arctan\frac58$. The other two are
\begin{gather*}
 \mathcal H_7=\frac1{(2\pi\ii)^2}\int_{\frac14-\ii\infty}^{\frac14+\ii\infty}
 \int_{\frac1{16}-\ii\infty}^{\frac1{16}+\ii\infty}
 \frac{\mathcal A_7(s)\cot(2\pi t)}{(\frac12-t)(\frac12-s-t)}\dd t\dd s,\\
 \mathcal A_7(s)=\frac{\Gamma(s)^3\Gamma(\frac12-s)^7\Gamma(\frac54-s)}
 {\Gamma(1-s)^3\Gamma(-s)\Gamma(\frac14-s)\prod_{b\in\{\beta_7,\bar\beta_7\}}\Gamma(s+b)\Gamma(1-s-b)},\\
 \mathcal H_*=\frac1{(2\pi\ii)^2}\int_{\frac14-\ii\infty}^{\frac14+\ii\infty}
 \int_{\frac14-\ii\infty}^{\frac14+\ii\infty}
 \frac{\Psi_+(s)\Psi_-(t)}{\Gamma(1-s-t)}\dd t\dd s,\\
 \Psi_\sigma(s)=\frac{\Gamma(s)^3\Gamma(\frac12-s)^3}
 {\Gamma(\frac12+s)\prod_{b\in\{\beta_\sigma,\bar\beta_\sigma\}}\Gamma(s+b)\Gamma(1-s-b)},
\end{gather*}
 where $\beta_7=\pi^{-1}\arctan\bigl((13+\ii\sqrt{139})/14\bigr)$ and, for
 $\sigma=\pm1$,
 $\beta_\sigma=\pi^{-1}\arctan\bigl(B_\sigma/2+\ii\sqrt{C_\sigma-B_\sigma^2/4}\bigr)$
 with $B_\sigma=\frac{21}{20}+\sigma e_*$,
 $C_\sigma=\frac{\lambda_*}{10}+\sigma e_*\chi_*$, $e_*=\sqrt{40\lambda_*-279}/20$,
 $\chi_*=2(21\lambda_*-150)/(40\lambda_*-279)$, and $\lambda_*$ is the root in
 $(15,16)$ of $4\lambda^3-112\lambda^2+909\lambda-2250=0$. The bar denotes the
 complex conjugate; all inverse tangents are principal (complex for
 $\beta_7$ and $\beta_\sigma$).
\end{itemize}
The formula evaluates to
\[
 C_6(1)=14331.68703275805432044546325524621711831732\ldots
\]
The three bivariate values carry the three sectors of the reduced expression for
which no evaluation in standard constants is known: the wheel sum $W_0$, the
seventh-sector sum $M_7$ \eqref{gff:c6:residualSums} and the pairing $P_a$ of
primitive hypergeometric functions \eqref{gff:c6:primitivePairing}
(Section~\ref{gff:c6:analytic}). The
count of nine describes how the terms are packaged, not the number of
independent constants.

The reduced expression is an input: it is derived neither in this note nor
in the package. Two of its Laurent coefficients were missing from its first
version and were added by a correction; they contribute $+659.279$ to
$C_6(1)$. The deterministic sum over all six-loop
graphs at quadrature degree 22 does not use it and gives $14331.6870328(5)$.
The representation lies $1.03\times10^{-8}$ below its unrounded value,
$14331.68703276838$, which is $0.021$ of its allowance; the comparison tests
the reduced expression at the level of $3\times10^{-11}$ relative. This
confirms the reduced expression and its correction numerically;
it does not prove them. The package encloses an earlier form of the reduced
expression, \eqref{gff:c6d:intervalForm}, in an interval of width below
$1.2\times10^{-145}$ (Section~\ref{gff:c6d:interval}). That form reaches
\eqref{gff:c6:reduced} only through exact identities of the package that this
note does not reproduce, among them two inherited normalizations of the wheel
periods; the later forms, reassembled from the same intervals with $M_7$
enclosed through the tail bound \eqref{gff:c6d:M7tail}, have the same width. The
enclosure therefore holds for the reduced expression conditionally on those
identities, and it encloses the reduced expression, not $C_6(1)$. An
independent review (\path{LRI_C6_closed_form_review/}) recomputed the input
and the nine values with separate codes and found no error.

\subsection{Status of the analytic results}\label{gff:c1:status}
The table separates what this note proves from what it takes as an input or
confirms only numerically. ``Proved here'' means that the note contains the
complete argument; ``conditional'' names the unproved input; ``numerical''
means agreement between independent evaluations, without an interval bound
unless one is stated. The numerical coefficients $C_6(d)$ for $d=2,3$,
$C_7(d)$ and $C_8(d)$ are sampling estimates with the allowances of
Section~\ref{gff:sec:guide} and are not listed.
{\small
\begin{longtable}{@{}>{\raggedright\arraybackslash}p{0.33\textwidth}>{\raggedright\arraybackslash}p{0.22\textwidth}>{\raggedright\arraybackslash}p{0.39\textwidth}@{}}
\toprule Statement & Where & Status\\\midrule\endhead
\multicolumn{3}{@{}l}{\emph{One dimension}}\\
Graph weights, subtraction and RG assembly of $C_0$--$C_5$ & Secs.~\ref{gff:sec:RG}, \ref{gff:sec:graphs}, \ref{gff:sec:assembly}; App.~\ref{gff:app:certificate} & proved here; every four- and five-loop residue is listed exactly\\
$\KK$ \eqref{gff:eq:KprovedFinal} & Sec.~\ref{gff:app:KpsiProof} & proved here\\
Wheel reduction of $\LL$ & Sec.~\ref{gff:app:LreductionProof} & proved here\\
$P_5(1)$ in five fixed values \eqref{gff:eq:P5FixedAnswer} & Sec.~\ref{gff:sec:P5FixedProof} & proved here\\
Graph integration of $C_5(1)$ & Sec.~\ref{gff:sec:checks} & numerical: agrees with the closed form to $6\cdot10^{-13}$ with the span cutoff and $3\cdot10^{-12}$ with the maxgap cutoff, at degree 24; not a bound\\
$C_6(1)=14331.68703275805\ldots$ \eqref{gff:c1:c6nine} & Secs.~\ref{gff:c1:c6analytic}, \ref{gff:c6:analytic}; App.~\ref{gff:app:C6derivation} & conditional: the reduced expression \eqref{gff:c6:reduced} is taken as an input and derived neither here nor in \path{C6_research_note/}; its conversion to the nine values is derived in App.~\ref{gff:app:C6derivation}, from three inherited elliptic-moment evaluations \eqref{gff:c6d:Iseries}--\eqref{gff:c6d:contourM2} and published identities; the interval of Section~\ref{gff:c6d:interval} encloses an earlier form of the reduced expression, which reaches \eqref{gff:c6:reduced} only through package identities not reproduced here, among them two inherited wheel normalizations; confirmed numerically by the graph sum to $3\cdot10^{-11}$ relative\\
$C_6(1)=14331.6870328(5)$, graph sum & Sec.~\ref{gff:app:C6Numerical} & numerical: exact sum over all graphs at quadrature degree 22, with an empirical convergence allowance\\
\multicolumn{3}{@{}l}{\emph{Two dimensions}}\\
$Q_2=\Gamma(\frac14)^8/(24\pi^4)$ & \eqref{gff:gd:d2exact}, \eqref{gff:gc:Q2} & proved here (complex Selberg integral, gamma reflection)\\
$Z_2=4\Cat Q_2$ \eqref{gff:gs:Z2exact} & Sec.~\ref{gff:gs:Z2section} & proved here: a complex Selberg integral for the two-screening integral \eqref{gff:gs:W3d2}, with the wheel jets \eqref{gff:gs:wheelJets}\\
$P_4(2)=(B_-^2-B_+^2)/(4\pi^3)$ \eqref{gff:fp:P4two} & Sec.~\ref{gff:fp:twoDimensional} & proved here: finite-Fourier two-mode theorem and the elliptic moments of Section~\ref{gff:sec:oldmasters}; the first derivation of $b$ cites the mixed-moment integral and the Stieltjes identity of \cite{gff:Wan}; the second (Section~\ref{gff:mgp:P4second}) uses neither\\
$P_5(2)$, quadratic in $X_0$, $X_\pm$ \eqref{gff:fp:P5two} & Sec.~\ref{gff:fp:twoDimensional} & proved here: square-lattice factorization \eqref{gff:fp:primitiveSquare}, the one-dimensional three-mode function \eqref{gff:fx:eq:abc} and exact residue algebra\\
$\LL_2=-2\log2\,P_4(2)$ \eqref{gff:gc:L2closed}, for the series \eqref{gff:gs:Lseries} & Sec.~\ref{gff:gc:L2proof} & proved here: elliptic lemmas of Section~\ref{gff:sec:oldmasters} and one contour integral, with the Stieltjes identity of \cite{gff:Wan} cited. Three further proofs avoid that identity: rim--spoke reflection (Section~\ref{gff:gc:L2second}), Parseval pairings on the unit circle (Section~\ref{gff:mgp:parseval}) and the Fourier representation of the two modes (Section~\ref{gff:mgp:fourier}); the last two use the half-strip evaluation of the elliptic moments (Section~\ref{gff:mgp:halfStrip})\\
$\KK_2=32\Cat^2-28\zeta(3)\log2$ \eqref{gff:gc:K2closed}, for the master-path definition \eqref{gff:gs:Kdefinition} & Secs.~\ref{gff:gc:K2proof}, \ref{gff:mgk:riemann} & proved here: position-space proof (rotation lemma, sum of two squares, Laurent extraction) with the one-dimensional lemmas of Section~\ref{gff:app:KpsiProof}; a second position-space proof of \eqref{gff:gc:K2master}, by Riemann's bilinear relation (Section~\ref{gff:mgk:riemann})\\
$\KK_2$ also equals the series \eqref{gff:gs:Kseries} and the spectral form \eqref{gff:gs:Kspectral} & Secs.~\ref{gff:mgk:spectral}, \ref{gff:mgk:series} & proved here: regularization in the rim exponent, residues on the square lattice and periodization, with the one-dimensional lemmas of Section~\ref{gff:sec:Kdef}; for the series also the per-spin residue identity, proved at $d=2$ with explicit arc bounds\\
Mixed-trace step \eqref{gff:gs:RS}, for $0<\tau<\frac12$ & Sec.~\ref{gff:mgk:regulated} & proved here at $d=2$, in the full-Mellin-line reading: the closed-form trace \eqref{gff:mgk:regulatedEq} with \eqref{gff:gc:K2master}; at other $d$ it is input (M) of Section~\ref{gff:gdk:conventions}\\
Closed form of $C_5(2)$ \eqref{gff:gc:C5twoClosed} & Sec.~\ref{gff:app:generalClosed} & proved here: exact algebra from the rows above and \eqref{gff:gs:C5}\\
\multicolumn{3}{@{}l}{\emph{General dimension}}\\
$C_0$--$C_5(d)$ \eqref{gff:gs:C0C4}, \eqref{gff:gs:C5} & Sec.~\ref{gff:gs:patterns} & proved by exact computation: exact symbolic residues of the 26 four-loop and 129 five-loop classes (certificate and exact audit in \path{lri_general_special_functions_update/})\\
Spectral and series forms of $Q_d$, $P_4$, $\LL_d$, $Z_d$ \eqref{gff:gs:primitiveWheel}, \eqref{gff:gs:LZspectral}, \eqref{gff:gs:QPseries}--\eqref{gff:gs:Zseries} & Secs.~\ref{gff:gs:logWheel}, \ref{gff:gs:residueProof} & proved here\\
Spectral and series forms of $\KK_d$ \eqref{gff:gs:Kspectral}, \eqref{gff:gs:Kseries} & Secs.~\ref{gff:gs:singleKernel}, \ref{gff:gdk:status} & at $d=2$ proved here (Section~\ref{gff:mgk:section}); at other $d$ conditional on input (M) and on extracting $[\tau^2]$ under the trace\\
Finite-Fourier theorem of the primitive sector and finite evaluation rules for $Q_d$, $P_4(d)$, $P_5(d)$ & Sec.~\ref{gff:app:generalPeriodization} & proved here for $\frac12<q<1$, then by continuation\\
Double residue series of $P_5(d)$ \eqref{gff:gc:P5series} & Sec.~\ref{gff:gc:primitive}; checks in \ref{gff:mgc:gcChecks} & proved here for $\frac12<q<1$; the extension by uniform convergence to $\frac25<d<4$ uses convergence orders whose elementary estimates are stated but not written out; checked in double precision at $d=2,\frac{13}5,3$\\
Finite Horn representation of $\LL_d$, $Z_d$, $\KK_d$ \eqref{gff:gc:finiteClosure} & Sec.~\ref{gff:app:generalClosed} & proved here for $\frac23<d<2$, then by continuation of the complete combination\\
Finite Fourier theorem at a general rim exponent, top-mode theorem and its product form, residue functional & Secs.~\ref{gff:gdr:finite}, \ref{gff:gdr:top}, \ref{gff:gdr:residue}; status in \ref{gff:gdr:status} & proved here, with complete proofs in Sec.~\ref{gff:gdr:periodization}. Reviewed independently in the form before the merge (\path{LRI_C5_general_d/review/VERIFICATION.md}): the finite Fourier theorem, the top-mode theorem in its wheel-only form, the residue functional and the block decomposition; the product form separately (\path{review/lemmaG-proof}). Not reviewed: the explicit constants of \eqref{gff:mgf:gammaBounds}, \eqref{gff:mgf:multiplicityBound}, \eqref{gff:mgf:latticeSums} (Sec.~\ref{gff:gdk:growthProof}) and the contour of the residue-functional proof, changed at the merge\\
Sample formulas for $\LL_d$, $Z_d$ and the reduced $C_5(d)$ & Secs.~\ref{gff:gdr:samples}--\ref{gff:gdr:C5}; checks in \ref{gff:gdr:checks} & sample formulas proved (Sections~\ref{gff:mgl:section} and \ref{gff:mgz:proofs}); reduced $C_5(d)$ derived (exact algebra, rederived symbolically in a separate review), conditional on (M), and on (H) on the jet circle, unless $\KK_d$ comes from \eqref{gff:gdk:KMB}; confirmed numerically to about $10^{-48}$\\
Constant \eqref{gff:gdk:Sconstant} of the two-wheel sample and bilateral identity \eqref{gff:gdk:Sbilateral} & Sec.~\ref{gff:gdk:Sfamily}; status in \ref{gff:gdk:status} & proved here from the case $r=2$ of the top-mode theorem, the residue functional and the two-wheel convolution; not independently reviewed; constancy and value confirmed numerically\\
Growth theorem for the $R$ family of $\KK_d$ & Secs.~\ref{gff:gdk:Rfamily}, \ref{gff:gdk:growthProof}; status in \ref{gff:gdk:status} & proved here; parts (a) (sufficiency), (b) and (c) are the case $r=3$ of the proof of the top-mode theorem, which was reviewed independently (Section~\ref{gff:gdr:status}); the sharpness of $\mathbb H_q$ and the explicit constants not independently reviewed\\
Sample formulas for $\KK_d$ & Secs.~\ref{gff:gdk:threeNode}, \ref{gff:gdk:continuation}; status in \ref{gff:gdk:status} & given (M) and (J), proved on $\mathbb H_q$ under the conditions of Section~\ref{gff:gdk:status}; the residue functional and the three-node rule not independently reviewed; at $d=2$, (M) is proved for real $0<\tau<\frac12$ (Section~\ref{gff:mgk:regulated}; complex $\tau$ is not addressed there) and (J) in Section~\ref{gff:gc:K2laurent}; on the jet circle also conditional on (H), at every $d$\\
Barnes representation of $\KK_d$ \eqref{gff:gdk:KMB} & Sec.~\ref{gff:gdk:positionSpace} & derived here without (M), (J), (H); the statement re-derived by separate AI sessions; the printed derivation not reviewed\\
\multicolumn{3}{@{}l}{\emph{Three dimensions}}\\
Determinant and rank-two identification of $Q_3$, $P_4(3)$, $P_5(3)$, and the complex-node variant \eqref{gff:td:complexAnswers} & Secs.~\ref{gff:d3:determinantProof}, \ref{gff:td:complexVariant}, \ref{gff:app:d3RankTwo} & proved here\\
Specialization of $\LL_d$ and $Z_d$ to $d=3$ & Sec.~\ref{gff:gdr:special}; conditions in \ref{gff:gdr:status} & derived here: exact algebra on two identities proved in Section~\ref{gff:gdr:halfShift} (the half-shift proof at a general rim exponent, with Dixon's and Dougall's sums); the identities re-derived by separate AI sessions, the printed derivation not reviewed; confirmed numerically to $10^{-51}$\\
Values of $Q_3$, $P_4(3)$, $P_5(3)$ and $C_5(3)$ & Sec.~\ref{gff:d3:values} & numerical: independent evaluations agree to about $10^{-79}$ ($Q_3$), $2\cdot10^{-86}$ ($P_4(3)$) and $2\cdot10^{-53}$ ($P_5(3)$), which changes $C_5(3)$ by less than $10^{-54}$; not interval bounds. $C_5(3)$ does not depend on (M), through \eqref{gff:gdk:KMB}\\
No ${}_3F_2(1)$ or ${}_4F_3(1)$ relation for $\mathfrak H_3(0)$ of $Q_3$ \eqref{gff:td:Qanswer} in a bounded search; no relation for $\mathsf r_3(3)$, $\KK_3$ & Secs.~\ref{gff:d3:values}, \ref{gff:gdr:attempts}, \ref{gff:gdk:special} & numerical only: bounded catalogue and PSLQ searches in stated bases, not independence proofs; no search is stated for $P_4(3)$, $P_5(3)$\\
\multicolumn{3}{@{}l}{\emph{Bulk one-point coefficient (Part~V)}}\\
$a^{(2)}_{\phi^2}$, $a^{(3)}_{\phi^2}$ & Sec.~\ref{gff:ope:lowOrders} & derived there; agree with \cite{gff:BLNV}\\
$a^{(4)}_{\phi^2}$--$a^{(6)}_{\phi^2}$ and the Gamma-factor form \eqref{gff:ope:factorized} & Secs.~\ref{gff:ope:result}, \ref{gff:ope:graphs} & calculation supplied by the author (16 October 2026); checks internal to the master framework: the five- and six-vertex graphs are reconstructed from master jets and checked with other regulator directions, not integrated one by one; the crossed kite rests on the package's Mellin--Barnes representation\\
$b_5$ and $c_5$ separately \eqref{gff:ope:b5}, \eqref{gff:ope:c5} & Secs.~\ref{gff:ope:RG}, \ref{gff:ope:checks} & $3c_5-b_5$ equals \eqref{gff:gs:D5} exactly; the split agrees numerically with graph sums at $d=1$ and with the seventh- and eighth-order productions at $d=2,3$\\
Map to $C_j(d)$ \eqref{gff:ope:logmap} & Sec.~\ref{gff:ope:thermal} & exact algebra between two calculations with shared inputs; validates neither\\
Exact $a^{(n)}_{\phi^2}(1)$, $a^{(n)}_{\phi^2}(2)$ and the closed $d=2$ map \eqref{gff:ope:d2log} & Secs.~\ref{gff:ope:d1}, \ref{gff:ope:d2} & exact algebra on the proved $d=1$, $d=2$ masters, conditional on $\mathsf H_4$ of the supplied calculation\\
Comparison with the numerical bootstrap & Sec.~\ref{gff:ope:bootstrap} & exploratory: readings of published figures, no bootstrap problem solved; the spread of prescriptions is not an error bar\\
\bottomrule
\end{longtable}}

\subsection{Notation and the logical route}\label{gff:c1:notation}
In the one-dimensional coordinate prescription the covariance is
$\langle\phi(x)\phi(0)\rangle_0=|x|^{-(1-\epsilon)/2}$ and the
interaction is $g\phi^4/4!$. The graph weights and MS recursion are fixed
first. The one-dimensional master reductions then give the full analytic
coefficient; the residue certificate of Appendix~\ref{gff:app:certificate} lists
the residue of every four- and five-loop class.
The general-dimensional sections replace the one-dimensional integrals by
normalized Schwinger and spectral integrals, derive the master jets and
provide their finite-function representations; Sections~\ref{gff:gdr:section}
and \ref{gff:gdk:section} shorten those of the three nonprimitive masters to a
few two-variable hypergeometric samples and their low-order parameter jets.
The numerical sections
retain a separate, regulator-dependent finite-scheme formulation and its
complete error accounting. Part~V uses the same graphs, subtractions and
masters for a second observable, the bulk one-point coefficient. Appendix~A
lists the residue certificate, and Appendix~B derives the special-function
representation of $C_6(1)$ quoted in Section~\ref{gff:c1:c6analytic}.

\begin{longtable}{@{}p{0.23\textwidth}p{0.72\textwidth}@{}}
\toprule Symbol & Meaning\\\midrule\endhead
$C_n(d)$ & Physical coefficient of $\epsilon^n$ in $\Delta_{\mathcal E}$.\\
$b_L,c_L$ & $L$-loop beta and mass-anomalous-dimension coefficients; lowercase $c_L$ is not $C_L$.\\
$z$ & Analytic regulator, set equal to physical $\epsilon$ only after the required Laurent or Taylor data have been retained.\\
$d,q$ & Spatial dimension and $q=d/4$ in the general-dimensional sections. In the five-index one-dimensional master $\mathcal J(a,b,c,d,e)$, $d$ is only an edge exponent.\\
$\widehat b_L(z),\widehat c_L(z)$ & Finite-scheme RG functions, whose $z$ dependence cannot be dropped before fixed-point extraction.\\
$P_3(d),Q_d$ & Primitive normalization $P_3(d)=3Q_d$. At $d=1$, $Q=\Gamma(1/4)^4/\pi$.\\
$\mathcal K_d,\mathcal L_d,Z_d$ & Explicit nonprimitive master data; $\mathcal K_d$ is not an elliptic integral. Reduced at generic $d$ in Sections~\ref{gff:gdr:section} and \ref{gff:gdk:section}.\\
$K(m),\mathbf K(k)$ & Complete elliptic integrals in parameter and modulus conventions: $\mathbf K(k)=K(k^2)$. A complementary argument is not a derivative.\\
$\psi_m(z)$ & $d^{m+1}\log\Gamma(z)/dz^{m+1}$; $\psi_0$ is digamma.\\
$\Cat$ & Catalan's constant; distinct from graph symbols and matrix-argument hypergeometric functions.\\
$a_{\phi^2}$, $a^{(n)}_{\phi^2}(d)$, $a^{(n)}$ & The bulk one-point coefficient of $:\!\Phi^2\!:$ and its coefficient of $\epsilon^n$ (Part~V); not the squared defect OPE coefficient $\lambda^2_{\phi\phi\Eop}$.\\
(M) & Input of Section~\ref{gff:gdk:conventions}: the mixed trace $J_{2\tau}=C(\tau)\langle t_\ell^2R_\tau\rangle_d$ as an identity of meromorphic functions of $\tau$, with the trace \eqref{gff:gdk:fullTrace} read on the full Mellin line and continued from its strip of convergence; proved at $d=2$ for real $0<\tau<\frac12$ (Section~\ref{gff:mgk:regulated}).\\
(J) & Input of Section~\ref{gff:gdk:conventions}: the jet $J_t-J_t^\Gamma=\KK_dt^2+O(t^3)$ on the master path \eqref{gff:gs:Kdefinition}, whose orders $t^0$, $t^1$ are used explicitly; at $d=2$ its low orders are rederived in Section~\ref{gff:gc:K2laurent}.\\
(H) & Continuation hypothesis of Section~\ref{gff:gdk:continuation}: the continued three-node samples \eqref{gff:gdk:sample}, \eqref{gff:gdk:regularized} are jointly holomorphic in $(t,\tau)$ on a domain that contains the disk bounded by the jet circle and meets $\mathbb H_q$, off a closed discrete set.\\
$\mathbb H_q$ & The half-plane $\operatorname{Re}\tau>\tau_*(q)=\max(0,2q-1,\frac12-2q)$ \eqref{gff:gdk:domain}, on which the lattice sum of the $R$ family converges absolutely (Section~\ref{gff:gdk:Rfamily}).\\
jet circle & The circle $|\tau|=r_\tau$ around the pole $\tau=0$ on which the $\tau$-jet of \eqref{gff:gdk:Ksample} is taken as a Cauchy integral; it does not lie in $\mathbb H_q$ (Section~\ref{gff:gdk:continuation}).\\
$[\tau^2]$ under the trace & The interchange of the Taylor coefficient $[\tau^2]$ at $\tau=0$ with the trace, $[\tau^2]\langle t_\ell^2(R_\tau-S_\tau)\rangle_d=\langle t_\ell^2E_2\rangle_d$, used to derive \eqref{gff:gs:Kspectral} from \eqref{gff:gs:Kdefinition} (Sections~\ref{gff:gs:singleKernel}, \ref{gff:gdk:lowOrders}); justified only at $d=2$ (Section~\ref{gff:mgk:regularization}).\\
\midrule
\multicolumn{2}{@{}l}{\emph{Part~III families}}\\
$b$, $b_{\max}$, $b_{\rm spec}$ & The rim exponent of a balanced wheel with spoke exponent $d-b$ \eqref{gff:gs:balancedWheel}, complex in general; the critical wheel has $b=d/2$. Sections~\ref{gff:gdr:section} and \ref{gff:gdk:section} write $\sigma=b/2$ and $g=(d-b)/2$ ($g_b$ in Section~\ref{gff:gdk:section}). The spectral integral \eqref{gff:gdr:kernel} represents $W_r$ for $0<\operatorname{Re}b<b_{\rm spec}=d-\max(d,1)/r$, and the finite Fourier theorem of Section~\ref{gff:gdr:finite} holds for $0<\operatorname{Re}b<b_{\max}(r,d)$ \eqref{gff:gdr:bmax}, with $b_{\max}\le b_{\rm spec}$. Not $b_L$, $b_\ell$ or the mode $b$ of \eqref{gff:fp:abWheel}.\\
$W_r(d;b)$, $t_{\ell,b}$, $N_b$, $\Xi_r$ & The $r$-wheel $W_r(d;b)=\langle t_{\ell,b}^{\,r}\rangle_d$ \eqref{gff:gs:balancedWheel}, with $t_{\ell,b}(\nu)=N_b\,g_\ell(\ii\nu;b)$ for real $b$ and the normalization $N_b$ of \eqref{gff:gdr:kernel}; at $b=d/2$, $t_{\ell,b}=t_\ell$ and $N_{d/2}=\Gamma(2q)$. $\Xi_r(t)$ \eqref{gff:fp:periodization} and $\Xi_r(t;b)$ \eqref{gff:gdr:periodizationDef} are the periodized samples, and $\Xi_2$ (Sections~\ref{gff:gdr:checks}, \ref{gff:gdk:Sfamily}) is the case $r=2$; $\Xi_{\rm R}$ \eqref{gff:gdk:periodization} belongs to the $R$ family, and $\Xi_\Gamma$ (Section~\ref{gff:two:gammaSector}) is unrelated. The rim sample $\mathfrak t_3(b)$, first used in Section~\ref{gff:app:threeDimensionalReduction}, is defined in \eqref{gff:gdr:samples3}.\\
$T_{\ell n}$, $T_{mn}(x)$ & $T_{\ell n}$ \eqref{gff:gs:seriesBase}, with $S_{\ell n}$, $a_{\ell n}$, $b_{\ell n}$ there (not $S_\tau$, $a_\ell$, $b_\ell$), and $T_{\ell n}(b)$ at a rim exponent $b$ \eqref{gff:gdr:Tln}. $T_{mn}(x)=c_m(x)c_n(x)$ \eqref{gff:gc:L2rim} has the first index $m=\ell+n$, so that $T_{mn}(\frac12)=T_{\ell n}$ at $d=2$; $T_\ell(z)$, with one index, is a product kernel of Section~\ref{gff:gdr:top}.\\
$J_t$, $J_{2\tau}$, $C(\tau)$, $R_\tau$, $S_\tau$ & The master path $J_t=J_d(\frac d2,\frac d2,\frac d2,d-t,\frac d2)$ of the jet \eqref{gff:gs:Bjet}, its Gamma part $J_t^\Gamma$ and $\KK_d=[t^2](J_t-J_t^\Gamma)$ \eqref{gff:gs:Kdefinition}. With $t=2\tau$, $C(\tau)=\Gamma(2q)\Gamma(\tau)/\Gamma(2q-\tau)$ and the kernels $R_\tau(s)$, $S_\tau(s)$ \eqref{gff:gs:RS} form the mixed trace of input (M); the $R$ and $S$ families of $\KK_d$ are those of $R_\tau$ and $S_\tau$. Used in Sections~\ref{gff:gs:singleKernel}, \ref{gff:gc:K2proof}, \ref{gff:mgk:section} and \ref{gff:gdk:section}. $C(\tau)$ is not a coefficient $C_j$, $R_\tau$ is not $R$, and $S_\tau$ is not $S_{\ell n}$.\\
$\mu_1$, $\mu_2$, $\eta_m$, $\theta_m$, $\lambda_m$, $\gamma_C$ & Polygamma combinations at the arguments $a_{\ell n}$ of \eqref{gff:gs:seriesBase}: $\lambda_m$ \eqref{gff:gs:lambda} and $\eta_m$, $\theta_m$, $\mu_1$, $\mu_2$ \eqref{gff:gs:etaTheta}, used in the series of Section~\ref{gff:gs:seriesSection}. The $\lambda_j$ of \eqref{gff:gc:dictionary} and the $\eta_r$, $\theta_r$, $\mu_1$, $\mu_2$ of \eqref{gff:gc:wheelDictionary} are the same combinations, with $m=\ell+n$; Sections~\ref{gff:gc:L2rimDerivative}, \ref{gff:gdr:taylor} and \ref{gff:mgz:jetsSection} continue them in the rim exponent. $\gamma_C(a)=\Gamma(a)/\Gamma(1-a)$ (Sections~\ref{gff:app:generalDimension}, \ref{gff:gs:Z2section}, \ref{gff:gc:L2second}, \ref{gff:gdr:special}).\\
$\mathcal N_d$, $\mathcal N_x=N_x$ & $\mathcal N_d=\KK_d+\frac38\LL_d-\frac{149}{32}Z_d$ \eqref{gff:fp:remainingCombination}, the only combination in which the nonprimitive masters enter $C_5(d)$. $N_x=\Gamma(1-x)/\Gamma(x)$ \eqref{gff:gc:L2rim} is written $\mathcal N_x$ in Sections~\ref{gff:gdr:residue} and \ref{gff:gdr:special}: one object with two names, unrelated to $\mathcal N_d$; it equals $N_b$ at $d=2$, $x=\frac b2$.\\
\midrule
\multicolumn{2}{@{}l}{\emph{Symbols with more than one meaning; each is defined where it is used}}\\
$h$ & $h_d=\Gamma(d/2)/\pi^{d/2}$, and an $\Im\psi_0$ combination in the formula for $\KK_d$ (both Section~\ref{gff:c1:closedForms}); the propagator exponent $(1-\epsilon)/2$ or $(1-z)/2$ (Sections~\ref{gff:sec:RG}, \ref{gff:sec:graphs}); the Gamma ratios $h_b(z)$ (Section~\ref{gff:gdr:halfShift}) and $h_x(\sigma)$ (Section~\ref{gff:mgp:fourier}); further local meanings in Sections~\ref{gff:gc:primitive}, \ref{gff:mgk:series}, \ref{gff:app:threeDimensionalReduction}, \ref{gff:gdr:periodization}, \ref{gff:c6:analytic}, \ref{gff:cn:finiteKernels}, \ref{gff:cn:finiteAssembly}, \ref{gff:app:C7d2} and \ref{gff:c8d3:independent}.\\
$p$ & an $\Re\psi_0$ combination in the formula for $\KK_d$, and $p=d/2-1$ in the five-loop spectral formula (both Section~\ref{gff:c1:closedForms}); $p=D/2-1$ in Section~\ref{gff:sec:P5complete}; local regulators, momenta, edge powers, exponent sums, pole orders and auxiliary polynomials elsewhere, for example in Sections~\ref{gff:gc:L2second}, \ref{gff:mgk:singles}, \ref{gff:td:complexVariant}, \ref{gff:mgf:nodes} and \ref{gff:app:C7d2}.\\
$B$, $\BB$, $\mathcal B$, $\mathfrak B$ & Euler's $B(u,v)$; $\BB(u,v)=\int_{\RR}|x|^{-u}|1-x|^{-v}\dd x$ (Section~\ref{gff:sec:jets}); $\mathcal B(g)$ and $B(g)$, the same reduced beta function (Sections~\ref{gff:sec:RG}, \ref{gff:app:C6Numerical}, \ref{gff:mgc:c6assembly}); $\mathcal B_L(z)$, the mass-insertion amplitude (Section~\ref{gff:sec:graphs}); $B_\pm$, hypergeometric values (Section~\ref{gff:c1:closedForms}), not $B_\sigma$, $\sigma=\pm1$, of Section~\ref{gff:c1:c6analytic}; $B_d$ \eqref{gff:gs:convolution}; the residue weights $B_{r,k}(\sigma)$ (Section~\ref{gff:gdr:residue}); the Bernoulli polynomials $B_n(x)$ (Sections~\ref{gff:d3:values}, \ref{gff:gdk:continuation}); $\mathcal B(x)$, a mode amplitude \eqref{gff:gc:L2modes}, and $\mathcal B_\tau$ (Section~\ref{gff:mgk:regulated}); local $B_k$, $B_i(n)$ (Sections~\ref{gff:mgk:regularization}, \ref{gff:mgk:series}); $\mathfrak B_{\rm kn}$ (Section~\ref{gff:mgc:c6assembly}), unrelated to $\mathfrak B_c$ and the transforms $\mathfrak B_{1/4}$, $\mathfrak B_{1/8}$ of Appendix~\ref{gff:app:C6derivation}.\\
$H$, $H_7$, $\mathcal H_7$, $\mathfrak H_7$ & $H=2\pi^2/3+8\Cat$ (Sections~\ref{gff:c1:closedForms}, \ref{gff:cn:oneDresults}); $H_7$, a ${}_7F_6$ value (Section~\ref{gff:sec:oldmasters}); $\mathcal H_7$, a bivariate $H$ function, with $\mathcal H_W$, $\mathcal H_*$ (Section~\ref{gff:c1:c6analytic}); $\mathfrak H_7$, a two-variable Horn function (Section~\ref{gff:app:threeDimensionalReduction}), with the Horn series $\mathfrak H[\cdots]$ \eqref{gff:gc:HornDefinition} and the values $\mathfrak H_+$, $\mathfrak H_c$ (Section~\ref{gff:mgz:fiveSection}); (H), the continuation hypothesis (first block), and the half-plane $\mathbb H_q$; $H(a;v)$ \eqref{gff:gc:dictionary}; $\mathcal H_x$ \eqref{gff:gc:L2residueForm}; $H_6^\chi$ and $H_6(d)$, bare channels, not ${}_6F_5$ values (Sections~\ref{gff:app:C6Numerical}, \ref{gff:mgc:c6assembly}); the harmonic numbers $H_n$ (Section~\ref{gff:c6:analytic}, Appendix~\ref{gff:app:C6derivation}); the Hardy space $H^2$ (Section~\ref{gff:mgp:parseval}); $H(\tau)$ (Section~\ref{gff:mgk:series}); $H_{\rm R}$, $H_{\rm M}$, $H_{\rm L}$ (Section~\ref{gff:gdk:threeNode}); further local meanings in Sections~\ref{gff:sec:shifts}, \ref{gff:sec:checks}, \ref{gff:cn:finiteKernels} and \ref{gff:app:C7complete}.\\
$G_4$, $\mathcal G_4$, $\mathcal G_L$, $G_n$, $\Cat$ & $G_4=\Gamma(\frac14)^4$ (Section~\ref{gff:sec:oldmasters}); $\mathcal G_4$, $\mathcal G_6$, $\mathcal G_7$, $\mathcal G_8$, Meijer $G$ values (Section~\ref{gff:c1:c6analytic}); $\mathcal G_L(z)$, a Gamma prefactor (Section~\ref{gff:cn:finiteAssembly}); $\mathcal G_x(s)$, a Gamma product (Sections~\ref{gff:gc:L2second}, \ref{gff:mgp:fourier}), and $\mathcal G_\alpha(n)$ \eqref{gff:gdk:kernels}; $G_n(z)$, coupling counterterms (Section~\ref{gff:sec:graphs}); $G_0$, $G_1$, $G_2$ \eqref{gff:gs:Bproducts}; $\Cat=\mathsf G$, Catalan's constant; $\mathfrak G(\tau)$, with $\mathfrak F(\tau)$ \eqref{gff:eq:KproofFG}, not the graph set $\Gr_L$ (Section~\ref{gff:sec:graphs}); local $G_{\ell k}$, $G(\beta)$, $G_\ell(s)$ (Sections~\ref{gff:gc:primitive}, \ref{gff:mgp:halfStrip}, \ref{gff:mgk:residues}).\\
$Q$, $Q_d$, $\mathcal Q_3$ & $Q=\Gamma(\frac14)^4/\pi=Q_1$; $Q_d$, the three-loop period, written $Q$ in the general-dimensional formulas of Section~\ref{gff:c1:closedForms}; $Q=Q_3$ in Section~\ref{gff:app:C6d2Complete}; $\mathcal Q_3=[\tau^3]\mathcal Q(\tau)^2$ (Section~\ref{gff:gc:K2proof}).\\
$R$, $U$, $T$, $\mathcal T$ & polygamma combinations ($R$, $U$) and values ($\pT$, $\pU$) in Section~\ref{gff:c1:closedForms}; integrals or functions in Sections~\ref{gff:app:selberg}, \ref{gff:sec:jets} and \ref{gff:sec:oldmasters}, those of Section~\ref{gff:sec:oldmasters} reused in Section~\ref{gff:gc:L2proof}; $U$ is also the allowance of Section~\ref{gff:c6:uncertainty}; $\mathcal R$ has different meanings in Sections~\ref{gff:sec:graphs} and~\ref{gff:app:C6d2Complete}, and $\mathcal R_{\ell n}$, $\mathcal R_+$, $\mathcal R_{\rm m}$ are lattice sums (Sections~\ref{gff:gs:seriesSection}, \ref{gff:mgl:section}); the $R$ family of $\KK_d$ is that of $R_\tau$ (Part~III families); $R$, the number of replicates (Sections~\ref{gff:mgs:benchmarks}, \ref{gff:c8d2:independent}); $U_G$, $U_0$, $U_F$, Symanzik polynomials, and $R$, a parameter sum (Sections~\ref{gff:cn:finiteKernels}, \ref{gff:cn:finiteAssembly}, \ref{gff:app:C7d2}); $T_\tau(s)$, $T_{\mathsf h}(\eta)$ (Sections~\ref{gff:mgk:residues}, \ref{gff:mgk:singles}) and $T_{\rm R}$, $T_{\rm L}$ (Sections~\ref{gff:gdk:Rfamily}, \ref{gff:gdk:growthProof}); $\mathcal T_d$, the corrected tetrahedron (Section~\ref{gff:app:generalDimension}), against $\mathcal T=\sum_na_n^4h_n(8h_n^2+\delta_n)$ in the wheel identities of Section~\ref{gff:c6d:wheel}.\\
$E$ & $\mathcal E=[\phi^2]$; $E_2$ in the formula for $\KK_d$ (Section~\ref{gff:c1:closedForms}), unrelated to the polynomials $E_m$ of Section~\ref{gff:app:generalPeriodization}; $E$, a polynomial in constants (Section~\ref{gff:c1:c6analytic}); $E(a,b)$, the Euler integral of $K$ (Sections~\ref{gff:sec:Kdef}, \ref{gff:mgk:singles}); $E_{\mathfrak k}$, $E^\sharp$ (Section~\ref{gff:gc:L2second}); $E_0$, $E_0^*$, $E_*^*$ (Sections~\ref{gff:app:threeDimensionalReduction}, \ref{gff:td:complexVariant}); $E$, the number of edges of a graph (Section~\ref{gff:app:C6Precision}); the line sets $E_B$, $E(W)$ (Sections~\ref{gff:app:C7d2}, \ref{gff:c8d3:independent}).\\
$Z$ & $\mathcal Z_g$, $Z_r$ and the counterterm coefficients $Z_n(z)$ (Section~\ref{gff:sec:graphs}); $Z_d$, the wheel master, with $Z_2=4\Cat Q_2$ \eqref{gff:gs:Z2exact}.\\
$D$, $\pD$ & $\pD=\pi$ (Section~\ref{gff:c1:closedForms}); $D$, the dimension in Section~\ref{gff:sec:P5complete}; $D=\pi\cot\pi q$ \eqref{gff:gdr:D}, with $S$, $C$ there; $D_4$, $D_5$ and in general $D_L$, the combinations $3c_L-b_L$ (Sections~\ref{gff:sec:RG}, \ref{gff:app:generalDimension}), with $D_5^{\rm prim}$ (Section~\ref{gff:gs:patterns}). Not $3c_L-b_L$: $D_3$, $D_5$, polylogarithm values (Sections~\ref{gff:c6:analytic}, \ref{gff:c6d:polylog}), and $D_4$, $D_6$, wheel sums (Section~\ref{gff:c6d:wheel}); $D_2$, the two-loop kernel of Section~\ref{gff:two:conversion}; $\widetilde D_j$ \eqref{gff:gc:dictionary}; $D_s$, $D_r$ (Sections~\ref{gff:gs:patterns}, \ref{gff:gs:logWheel}); the sine products $D_q(t)$, $D_b(t)$, $D_T$, $D_{\rm R}$ (Sections~\ref{gff:app:generalPeriodization}, \ref{gff:gdr:periodization}, \ref{gff:gdr:top}, \ref{gff:gdk:Rfamily}); $\mathfrak D_L$ (Section~\ref{gff:gdr:samples}); $D(\ell)$ (Section~\ref{gff:mgk:series}); local $\mathcal D(\tau)$, $\mathcal D_2$, $\mathcal D_3$ (Section~\ref{gff:mgk:regularization}), $D_{\rm F}$ (Section~\ref{gff:mgp:fourier}) and $\mathcal D_4(y)$ (Section~\ref{gff:c6d:wheelcert}).\\
$\mathcal K$, $K$ & $\KK$ and $\KK_d$, the nonprimitive master; $\KP$ and, in Sections~\ref{gff:cn:finiteKernels}--\ref{gff:app:C6d2Complete}, $\mathcal K$, the pole-part projector; $K_G$, the forest kernels of Sections~\ref{gff:cn:finiteAssembly}--\ref{gff:app:C7d2}; $K(m)$, $\mathbf K(k)$, elliptic integrals; $\boldsymbol K_{\le4}$, the loop one-to-four vector of Section~\ref{gff:app:C7complete}; $K_\ell(s)$, the five-loop primitive kernel \eqref{gff:fp:primitiveStart}; $K_5$, the complete graph (Sections~\ref{gff:gs:patterns}, \ref{gff:mg7:primitive}); $K_n$, a Pochhammer ratio (Section~\ref{gff:gc:L2second}) and the sums $K_n$ (Section~\ref{gff:c6d:scope}); $\mathfrak K$, a compact set (Section~\ref{gff:gdk:growthProof}).\\
$G_{L,d}$ & the Gamma prefactor of Sections~\ref{gff:cn:finiteAssembly} and \ref{gff:app:C6Precision}; its two normalizations differ by a constant, so the log-jets $a_{L,j}$ and $\ell_j$ agree for $j\ge1$.\\
$r_j$, $\mathsf r_k$ & $r_j$, series coefficients (Section~\ref{gff:gs:seriesSection}); $\mathsf r_k$, $\mathsf s_k$, Laurent coefficients \eqref{gff:gdk:notation}, with $\mathsf r_{-1},\ldots,\mathsf r_2$ in \eqref{gff:gdk:rLow} and $\mathsf r_3=[\tau^3]\langle t_\ell^2R_\tau\rangle_d$ (Section~\ref{gff:gdk:section}); $r_3=2Q_2$, the primitive three-loop residue (Section~\ref{gff:cn:finiteKernels}); $\mathsf r_1$, $\mathsf r_2$ \eqref{gff:ope:r12}, the coefficients of the ratio $b_0^{\Phi,-}/b_0^{\Phi,+}$ of bulk-to-defect coefficients (Section~\ref{gff:ope:lowOrders}), not the $\mathsf r_1$, $\mathsf r_2$ of \eqref{gff:gdk:rLow}.\\
$w$, $\mathfrak w$ & $w_1$, $w_2$, spectral weights (Sections~\ref{gff:c1:closedForms}, \ref{gff:mgz:jetsSection}); the graph weights $w_4=w_A$ and $w_2=w_B$ (Sections~\ref{gff:sec:graphs}, \ref{gff:app:C6Numerical}); $w_H=3w_B-w_A$; $w(k)=\log\frac k{1-k^2}$ (Sections~\ref{gff:sec:oldmasters}, \ref{gff:mgp:halfStrip}); $w_{\KK}$, the weight of \eqref{gff:gs:Kspectral} at $d=2$ (Section~\ref{gff:mgk:conventions}); $w_j(\tau)$, the node weights \eqref{gff:gdk:threeNodeRule}; $\mathfrak w_{r,h}$, spin-tail weights (Section~\ref{gff:d3:values}), and $\mathfrak w_{r,j}(b)$, node weights \eqref{gff:mgf:Wsample}.\\
$z$, $z_1$, $z_2$ & $z$, the analytic regulator; $z_1$, $z_2$, deformation coordinates (Sections~\ref{gff:sec:jets}, \ref{gff:gs:patterns}) or Mellin--Barnes variables (Section~\ref{gff:gdk:positionSpace}); $z$ is also a local complex variable in proofs where the regulator does not enter or is called $\tau$ (Sections~\ref{gff:gc:L2second}, \ref{gff:mgp:parseval}, \ref{gff:mgk:regularization}, \ref{gff:gdr:periodization}), and $z_\ell=\frac\ell2+\ii\nu$ (Section~\ref{gff:mgf:nodes}).\\
$D_G$, $\varpi$ & local meanings in Sections~\ref{gff:cn:finiteKernels} and \ref{gff:cn:finiteAssembly} ($D_G$); Mellin--Barnes weights (Section~\ref{gff:gdk:positionSpace}), hypergeometric values (Section~\ref{gff:c6:analytic}) and Binet's function $\varpi(t)$ (Section~\ref{gff:gc:L2second}) ($\varpi$).\\
MOM & the finite zero-momentum scheme of Section~\ref{gff:cn:finiteAssembly}.\\
$a$, $\mathfrak a$ & $a=\log2$ and $a_\ell=\frac\ell2+\frac q2$ (Section~\ref{gff:c1:closedForms}), with $a_\ell=\frac{\ell+\sigma}2$ at a rim exponent \eqref{gff:gdr:kernel}; the finite Laurent coefficients $a_{L,k}$ (Sections~\ref{gff:sec:graphs}, \ref{gff:app:generalDimension}); the series terms $a_n(t)$ (Section~\ref{gff:sec:Kdef}); the log-jets $a_{L,j}$ (Sections~\ref{gff:cn:finiteAssembly}, \ref{gff:app:C6Precision}); $a_n$ in Section~\ref{gff:c6:analytic}; $a^{(n)}$ and the amplitudes $\mathsf a_n(\epsilon)$ (Part~V); $a=8N_{\rm e}/\pi$, a mode at $d=2$ \eqref{gff:fp:abWheel}; $a=a_\ell$ at fixed spin (Section~\ref{gff:mgk:conventions}); $a=\frac12+\frac z4$ (Sections~\ref{gff:cn:finiteKernels}, \ref{gff:app:C7d2}, \ref{gff:c8d2:independent}); the prefactors $a_r^0$, $a_r^+$, $a_r^c$ (Sections~\ref{gff:fp:arraysSection}, \ref{gff:mgz:fiveSection}); $a_G$ (Sections~\ref{gff:cn:finiteKernels}, \ref{gff:mg7:primitive}). $\mathfrak a_3$, $\mathfrak b_3$, $\mathfrak a_4$, $\mathfrak b_4$, mode amplitudes (Section~\ref{gff:app:threeDimensionalReduction}), continued in $b$ by \eqref{gff:gdr:halfShiftModes}; $\mathfrak a(x)$, $\mathfrak b(x)$ and $\mathfrak s(x)$ (Section~\ref{gff:gc:L2second}) are not these or the Horn values $\mathfrak s_0$, $\mathfrak s_{1/2}$ (Sections~\ref{gff:gdr:samples}, \ref{gff:mgl:section}).\\
$A$, $\mathcal A$ & $A$ and $A_d$ \eqref{gff:gd:blocks}, so $A_2=4\log2$ and $A_3=2+4\log2-\pi$; $A=\pA=A_1$ in Sections~\ref{gff:c1:closedForms} and \ref{gff:c1:c6analytic}; the Horn blocks $A_r^0$, $A_r^+$ \eqref{gff:fp:wheelBlocks} (Sections~\ref{gff:gdr:residue}, \ref{gff:mgz:fiveSection}); local $A_k$, $B_k$ and $A_i(m)$ (Sections~\ref{gff:mgk:regularization}, \ref{gff:mgk:series}), not $A_d$. $\mathcal A_L(z)$, $\mathcal A_L(d,z)$, the bare amplitudes (Sections~\ref{gff:sec:graphs}, \ref{gff:app:generalDimension}), and $\mathcal A_L^\chi$ (Section~\ref{gff:app:C6Numerical}); $\mathcal A_{\rm s}(t)$ (Section~\ref{gff:sec:P5complete}); $\mathcal A_{\ell k}$ (Sections~\ref{gff:app:generalPeriodization}, \ref{gff:gc:primitive}); $\mathcal A(x)$, a mode amplitude \eqref{gff:gc:L2modes}; $\mathcal A_\tau$ (Section~\ref{gff:mgk:regulated}); $\mathcal A(h,z)$ \eqref{gff:done:bareSeries} and $\mathcal A(x,z)$ (Section~\ref{gff:app:C7complete}); $\mathcal A_7(s)$ (Section~\ref{gff:c1:c6analytic}).\\
$b$, $c$ & $b_L$, $c_L$ (first block), and the rim exponent $b$ (Part~III families); $b=8M_{\rm e}/\pi$, a mode at $d=2$ \eqref{gff:fp:abWheel}; $b=b_\ell$ at fixed spin (Section~\ref{gff:mgk:conventions}); $b=\frac12+\frac{3z}4$ (Section~\ref{gff:c8d2:independent}); $b_0^{\Phi,\pm}$, bulk-to-defect coefficients (Section~\ref{gff:ope:lowOrders}); $c_j$ \eqref{gff:gc:dictionary}, so the $c_5$, $c_6$ of \eqref{gff:gc:kappa} are not loop coefficients; $c_n=(\frac12)_n/n!$ (Sections~\ref{gff:mgp:P4second}, \ref{gff:gc:L2proof}, \ref{gff:mgk:series}); $c_k(x)$ \eqref{gff:gc:L2rim}; $c_*$ \eqref{gff:gs:Kweights}; $c_\tau=\cos\pi\tau$ (Sections~\ref{gff:gc:K2proof}, \ref{gff:mgk:section}); the Fourier coefficients $c_{r,k}$ (Sections~\ref{gff:app:generalPeriodization}, \ref{gff:gdr:finite}) and $c_k(\tau)$ (Section~\ref{gff:gdk:Rfamily}); generic constants $c$, $c_0$ (Sections~\ref{gff:mgk:regularization}, \ref{gff:gdk:growthProof}).\\
$C$, $S$ & $C_n(d)$ (first block); $C(\tau)$ (Part~III families); $S=\sin\pi q$, $C=\cos\pi q$ \eqref{gff:gdr:D} (Sections~\ref{gff:gdr:samples}--\ref{gff:gdr:C5}); the Horn blocks $C_r(t)$ \eqref{gff:fp:wheelBlocks}, for example $C_3(t)$, not $C_3(d)$ (Sections~\ref{gff:gdr:residue}, \ref{gff:mgz:fiveSection}); (C), the condition \eqref{gff:gdr:BC}; generic constants $C$, $C_0$, $C_r$, $C_q$, $C_K$, $C_{\mathfrak K}$ in bounds (Sections~\ref{gff:gc:L2second}, \ref{gff:mgk:regularization}, \ref{gff:gdr:periodization}, \ref{gff:mgz:uniform}, \ref{gff:gdk:growthProof}); $C_d$, $C_s$, $C_\sigma$, $C_{212}$ of Section~\ref{gff:c1:c6analytic}; the zonal polynomials $C^{(2)}_\kappa$ (Sections~\ref{gff:app:d3RankTwo}, \ref{gff:d3:values}); $\mathsf C$ (Part~V); $S_{\ell n}$ \eqref{gff:gs:seriesBase}, the periodized sums $S_r(t)$, $S_{\rm R}$ (Sections~\ref{gff:app:generalPeriodization}, \ref{gff:gdr:periodization}, \ref{gff:gdk:Rfamily}) and the sectors $S_{\rm pol}$, $S_6$, $S_8$ (Section~\ref{gff:c6:analytic}).\\
$F$ & $F_4(v)$, $F_2(v)$, the bare amplitudes (Section~\ref{gff:sec:graphs}); $F_1(d,z)$, $F_2(d,z)$ of \eqref{gff:gd:twoLoopHyper} (Sections~\ref{gff:app:generalDimension}, \ref{gff:app:C7d2}, \ref{gff:c8d2:independent}); the polygamma polynomials $F_3$, $F_4$, $F_{3,1}$, $F_{4,1}$ \eqref{gff:gs:Fpolys} (Sections~\ref{gff:gs:seriesSection}, \ref{gff:mgl:section}, \ref{gff:mgz:jetsSection}); $F(\tau)$ and its coefficients $F_k$ \eqref{gff:gdk:notation}, so that $F_3=\mathsf r_3-\mathsf s_3$ in Section~\ref{gff:gdk:lowOrders} is not the polynomial $F_3$; $F(s)$ \eqref{gff:gs:Kholomorphic}; $F_{r,j}$ (Section~\ref{gff:app:generalClosed}); $F_\pm$, ${}_4F_3$ values (Section~\ref{gff:c1:closedForms}); local $F_{\rm L}$, $F_{\rm R}$ (Section~\ref{gff:gc:L2second}) and $F(\beta)$, $G(\beta)$, $L(\beta)$ (Section~\ref{gff:mgp:halfStrip}).\\
$g$ & the coupling (Section~\ref{gff:sec:RG}); the Barnes kernels $g_\ell(s)$, $g_\ell(s;b)$ (Sections~\ref{gff:app:generalSpecial}--\ref{gff:gdk:section}) and $\mathsf g_\ell(s)$ (Section~\ref{gff:gc:L2second}); $g=(d-b)/2$, the spoke half-exponent (Section~\ref{gff:gdr:section}), written $g_b$ in Section~\ref{gff:gdk:section}; the free exponents $g_i$ of $f_{g_i}$ (Section~\ref{gff:gdr:top}); $g(k)=\mathbf K(k')$ (Sections~\ref{gff:sec:oldmasters}, \ref{gff:mgp:halfStrip}); $g(m,n)$ (Section~\ref{gff:mgk:series}); $g_{\rm alt}$ (Section~\ref{gff:ope:RG}).\\
$I$ & $I_m(z)$ (Section~\ref{gff:sec:graphs}); the Euler moments $I_1,\ldots,I_6$ (Sections~\ref{gff:sec:Kdef}, \ref{gff:gc:K2proof}, \ref{gff:mgk:singles}, \ref{gff:mgk:riemann}); $I_1$, $I_3$, $I_5$, $I_7$ (Section~\ref{gff:sec:P5complete}); the pairings $I_{jk}$ (Section~\ref{gff:c6:analytic}, Appendix~\ref{gff:app:C6derivation}); identity matrices $I_2$, $I_L$ (Sections~\ref{gff:app:d3RankTwo}, \ref{gff:rt:kummerDiagnostic}, \ref{gff:cn:finiteKernels}); $I$, the number of model terms (Section~\ref{gff:gdk:continuation}); local $I_\ell$, $I_{\rm L}$, $I_{\rm R}$ and intervals $I$, $I_{\rm F}$ (Sections~\ref{gff:gdr:top}, \ref{gff:gc:L2second}, \ref{gff:gc:L2proof}, \ref{gff:mgp:fourier}).\\
$J$, $\mathsf J$ & $\JJ_m(z)$ (Section~\ref{gff:sec:graphs}) and the master $\mathcal J(a,b,c,d,e)$ (first block); the jet $J_d$ \eqref{gff:gs:Bjet} and the master path $J_t$, $J_{2\tau}$ (Part~III families); $J_0(q)$ (Section~\ref{gff:app:generalDimension}); $J_1$, $J_3$, $J_5$ (Section~\ref{gff:sec:P5complete}); $J_r(a;v)$ \eqref{gff:gc:dictionary}; the bare channels $J_5^\chi$, $J_5(d)$ (Sections~\ref{gff:app:C6Numerical}, \ref{gff:mgc:c6assembly}); the Jacobians $J$, $J_L$, $J_{f,L}$ (Sections~\ref{gff:meth:pooling}, \ref{gff:c7d2:production}); $J(|i-j|)$, the lattice coupling (Section~\ref{gff:sec:guide}); $\mathsf J_0(x)$, $\mathsf J_1(x)$, Barnes integrals (Section~\ref{gff:gc:L2second}), not $\mathsf J_d$ of Section~\ref{gff:ope:lowOrders} at $d=0,1$.\\
$P$ & the periods $P_3(d)=3Q_d$, $P_4(d)$, $P_5(d)$ (first block); $P_W$, a wheel period (Sections~\ref{gff:sec:wheelclassupdated}, \ref{gff:gs:logWheel}); the lattice weights $P_\pm$, $P_{\rm m}$ (Section~\ref{gff:mgl:section}) and $P_{\rm R}$, $P_{\rm M}$, $P_{\rm L}$ (Sections~\ref{gff:gdk:growthProof}, \ref{gff:gdk:threeNode}); $P_{10}$, the retained ten primitive six-loop kernels (Sections~\ref{gff:prec:d2production2026}, \ref{gff:meth:refinement}--\ref{gff:meth:crossdimension}); $P_a$, a pairing \eqref{gff:c6:primitivePairing}; local $P_f$, $P_{\mathsf h}$ (Section~\ref{gff:mgk:singles}).\\
$q$ & $q=d/4$ (first block); the nome $q=e^{-\pi\mathbf K(k')/\mathbf K(k)}$ (Section~\ref{gff:mgp:fourier}); the coefficients $q_n$ of $\mathsf q(z)$ (Section~\ref{gff:mgp:parseval}) and the sequence $q_n$ of Section~\ref{gff:c6d:wheel}; the parameters $q_{7,j}$ of Section~\ref{gff:c1:c6analytic}.\\
$s$, $\varsigma$ & the long-range exponent $s$ (Section~\ref{gff:sec:guide}); the Mellin variable $s$ (Sections~\ref{gff:gs:singleKernel}, \ref{gff:gdk:conventions}); $s_b=\sin\frac{\pi b}2$, $\varsigma_b=\cos\frac{\pi b}2$ (Section~\ref{gff:gdr:section}); $s_\tau=\sin\pi\tau$ (Sections~\ref{gff:gc:K2proof}, \ref{gff:mgk:section}); $\varsigma_f$, sign vectors (Section~\ref{gff:gc:K2proof}); $s_{\rm K}$, a parametric excess (Section~\ref{gff:rt:kummerDiagnostic}).\\
$s_G$ & $s_G(x)=\frac14\sum_e\log x_e-\frac L2\log S(x)$, the regulator logarithm (Section~\ref{gff:app:C6Precision}); $s_G=\sum_ie_i(e_i-1)$, the leg count of a graph (Section~\ref{gff:app:C7d2}).\\
$t$ & the path parameter of $J_t$ \eqref{gff:gs:Kdefinition}, not $t=d-e$ of \eqref{gff:gs:Acoords}; the periodization variable of $S_r(t)$, $\Xi_r(t)$ (Sections~\ref{gff:app:generalPeriodization}, \ref{gff:gdr:section}, \ref{gff:gdk:Rfamily}); the wheel eigenvalues $t_\ell$, $t_{\ell,b}$ (Section~\ref{gff:c1:closedForms}, Part~III families); a direction parameter in Section~\ref{gff:ope:lowOrders}.\\
$u$, $v$, $\sigma$ & $u_j$, the coefficients of $g_*$ (Section~\ref{gff:sec:RG}); $u$, $v$, $\sigma$ with $\tau$ (Section~\ref{gff:sec:Kdef}) or $\varrho$ (Sections~\ref{gff:gc:K2laurent}, \ref{gff:mgk:conventions}), Taylor data of $\mathfrak F$, $\mathfrak G$; $u(b)$, $v(b)$, the weights of \eqref{gff:gdr:twonode}, with $u_0$, $v_0$ (Section~\ref{gff:mgz:simplifiedSection}); local $u$, $v$ (Section~\ref{gff:gdr:halfShift}); $u_h$ \eqref{gff:fp:logCoefficients}; $v_m(x)$ (Section~\ref{gff:mgp:fourier}). $\sigma_i=\pm1$, the spins (Section~\ref{gff:sec:guide}), and $\sigma=\pm1$ (Section~\ref{gff:c1:c6analytic}); $\sigma=b/2$ (Sections~\ref{gff:gdr:section}, \ref{gff:gdk:section}), a free offset in Section~\ref{gff:gdr:top}; $\sigma=\operatorname{Im}t$ (Section~\ref{gff:td:complexVariant}); $\widehat\sigma$, standard errors (Part~IV); further local meanings in Sections~\ref{gff:sec:jets}, \ref{gff:sec:P5complete}, \ref{gff:app:generalDimension}, \ref{gff:gs:patterns}, \ref{gff:gs:Z2section} and \ref{gff:cn:finiteKernels}.\\
$V$, $X$, $\mathcal X$, $\mathsf X$ & $V$, $X$, polygamma constants \eqref{gff:gs:psiBasis}; $X_0$, $X_\pm$, the one-dimensional values of Section~\ref{gff:c1:closedForms}, against $X_0$, $X_h$ at $d=3$ \eqref{gff:td:P4data}, with $X_0=\Xi_4(0;\frac32)$ (Section~\ref{gff:gdr:special}); $\mathsf X_r(t)$ (Section~\ref{gff:app:threeDimensionalReduction}); $\mathcal X_d(\tau)$ \eqref{gff:gdk:notation}, with $\mathcal X_2(\tau)$ (Section~\ref{gff:mgk:conventions}); the samples $\mathcal X_0$, $\mathcal X_{1/2}$, $\mathcal X_T$ (Sections~\ref{gff:gdr:samples}, \ref{gff:mgl:section}), where the subscript labels a node, not a dimension; the model factors $X(m)$ (Section~\ref{gff:gdk:continuation}); $V$, $V_0$, elliptic moments (Sections~\ref{gff:sec:oldmasters}, \ref{gff:mgp:halfStrip}); $V_r(t)$, $V_r^{(0)}(t)$ (Sections~\ref{gff:app:threeDimensionalReduction}, \ref{gff:td:complexVariant}); the variances $V_C$, $V_{\rm pf}$, $V_\Delta$, $V_{f\Delta}$ (Section~\ref{gff:app:C7d2}) and $\widehat V_G$, $\widehat V_{\rm bench}$ (Section~\ref{gff:mgs:benchmarks}); further local meanings in Sections~\ref{gff:sec:shifts}, \ref{gff:gdr:periodization}, \ref{gff:cn:finiteKernels} and \ref{gff:ope:observable}.\\
$x_0$, $m_0$ & $x_0$, a vertex fixed at the origin (Sections~\ref{gff:sec:graphs}, \ref{gff:sec:oldmasters}, \ref{gff:sec:P5complete}, \ref{gff:ope:lowOrders}); $x_0$, a Gamma ratio \eqref{gff:gdk:modelCoefficients}; $m_0$, an integer cutoff (Section~\ref{gff:gdk:growthProof}), and $m_0$ of Section~\ref{gff:c6d:Hconvention}; $\mathfrak m_0$, a moment (Section~\ref{gff:mgp:fourier}).\\
$C_d$, $C_s$, $C_\sigma$, $d_7$, $q_{7,j}$, $\rho_4$, $\rho_7$, $\rho_8$; $W_0$, $M_7$, $P_a$ & parameters of the nine values of Section~\ref{gff:c1:c6analytic}, used again in Section~\ref{gff:c6:analytic} and Appendix~\ref{gff:app:C6derivation}: $C_d$, $C_s$, $C_\sigma$, $B_\sigma$, $d_7$, $q_{7,j}$, $\rho_4$, $\rho_7$, $\rho_8$, $\eta_{7,\pm}$ and $\lambda_*$ are not $C_n(d)$, $B_\pm$, $d_\ell$, $q$, $\rho_n^{(d)}$, $\eta_m$ or $\lambda_m$. Section~\ref{gff:c1:c6analytic} uses the wheel sum $W_0$, the seventh-sector sum $M_7$ and the pairing $P_a$ before their definitions \eqref{gff:c6:residualSums}, \eqref{gff:c6:primitivePairing}; $W_0$ is not a wheel $W_r$.\\
$\beta$ & $\beta_{\rm RG}(g)$ (Section~\ref{gff:sec:RG}); $\alpha$, $\beta$, $\gamma$, constants of the five-loop primitive (Sections~\ref{gff:sec:P5FixedProof}, \ref{gff:fp:twoDimensional}); $\beta_\ell=a_\ell+g$ \eqref{gff:gdr:kernel}; the free shifts $\beta_i$ of \eqref{gff:gdr:BC} (Sections~\ref{gff:gdr:top}, \ref{gff:gdk:topmode}); $\beta_k$, $\beta_n$, digamma sequences (Sections~\ref{gff:gc:L2proof}, \ref{gff:mgp:parseval}); $\beta_f$ (Section~\ref{gff:gc:K2proof}); $\beta_7$, $\beta_\sigma$ (Section~\ref{gff:c1:c6analytic}); the bubbles $\beta_{ij}$ (Section~\ref{gff:app:C7d2}).\\
$\gamma$ & $\gamma_r(g)$ (Section~\ref{gff:sec:RG}); $\alpha$, $\beta$, $\gamma$, constants of the five-loop primitive (Section~\ref{gff:sec:P5FixedProof}); Euler's constant $\gamma_E$ (Section~\ref{gff:sec:Kdef}), written $\gamma$ in Section~\ref{gff:mgp:fourier}; $\gamma_C$ (Part~III families); the Fourier coefficients $\gamma_m$ (Sections~\ref{gff:gdr:finite}, \ref{gff:gdr:top}, \ref{gff:gdk:growthProof}); $\gamma$, a divergent subgraph (Sections~\ref{gff:app:C7d2}, \ref{gff:c8d3:independent}).\\
$\eta$ & $\eta_m$, polygamma combinations (Part~III families); $\eta_r=1/(8\pi r)$, the imaginary offset of the nodes \eqref{gff:fp:nodes} (Sections~\ref{gff:gdr:residue}, \ref{gff:mgf:nodes}); $\eta$ and $\eta_{\rm R}$ \eqref{gff:gdk:domain}, decay exponents of the growth bounds (Sections~\ref{gff:app:generalPeriodization}, \ref{gff:gdr:periodization}, \ref{gff:gdr:finite}, \ref{gff:gdk:Rfamily}); $\eta_{7,\pm}$, parameters of $\mathcal G_7$ (Section~\ref{gff:c1:c6analytic}); $\eta$, the remainder of \eqref{gff:meth:functional}; a local real variable in Sections~\ref{gff:mgp:halfStrip}, \ref{gff:mgp:fourier}, \ref{gff:gc:K2proof}, \ref{gff:mgk:singles}, \ref{gff:mgk:riemann}, \ref{gff:gdk:Sfamily} and \ref{gff:two:gammaSector}.\\
$\theta$ & $\theta_m$ and the $\theta_r$ of \eqref{gff:gc:wheelDictionary}, polygamma combinations (Part~III families); $\theta=\theta_0=2q$ or $\theta_1=2q-1$, the index of the measure block $M_\theta$ (Section~\ref{gff:gc:primitive}); $\theta=-\frac14$, a parameter of a complex Selberg integral (Section~\ref{gff:app:generalDimension}); $\theta_i(z)$, Stirling remainders (Section~\ref{gff:gdk:growthProof}); $\theta_{\rm cv}$, a control coefficient (Section~\ref{gff:meth:crossdimension}); elsewhere an angle.\\
$\vartheta$ & the auxiliary coordinate of the edge exponents $\alpha_e$ (Section~\ref{gff:ope:method}); an angle, with $f_\vartheta=\mathbf K(\sin\vartheta)$ and $g_\vartheta=\mathbf K(\cos\vartheta)$ (Section~\ref{gff:mgp:parseval}), and a local angle in Section~\ref{gff:gc:K2proof}; $\vartheta_\tau(x)$, a Gamma ratio (Section~\ref{gff:mgk:regularization}); $\vartheta=2\psi_1(q)-\psi_1(2q)$ (Section~\ref{gff:gdk:positionSpace}).\\
$\kappa$ & $\kappa(d,z)$ \eqref{gff:gd:normalization}; $\kappa_{\theta,j}$, $\kappa^0_j$ \eqref{gff:gc:kappa}; $\kappa(S)$, a nominal endpoint power (Section~\ref{gff:mgs:benchmarks}); local exponents, partitions and constants in Sections~\ref{gff:gc:K2proof}, \ref{gff:gc:L2second}, \ref{gff:mgp:halfStrip}, \ref{gff:app:d3RankTwo}, \ref{gff:gdr:top}, \ref{gff:gdr:samples}, \ref{gff:gdr:taylor} and \ref{gff:c6:analytic}. The exponent $\varkappa$ (Sections~\ref{gff:gdr:top}, \ref{gff:gdk:growthProof}) is a different letter.\\
$\lambda$ & $\lambda_m$ and $\lambda_j$, the same polygamma combinations, continued as $\lambda_j(x)$ (Part~III families); $\lambda_j$, local to the wheel-tail paragraph of Section~\ref{gff:d3:values}; $\lambda_k$, the Stirling coefficients \eqref{gff:gdk:modelCoefficients}; $\lambda_e$, the completion exponents (Sections~\ref{gff:sec:jets}, \ref{gff:gs:patterns}); the bare couplings $\lambda_0$ (Section~\ref{gff:app:generalDimension}), $\lambda_B$ and $\lambda$ (Sections~\ref{gff:ope:observable}, \ref{gff:ope:RG}); $\lambda_*$, a cubic root (Sections~\ref{gff:c1:c6analytic}, \ref{gff:c6d:HPdef}); $\lambda^2_{\phi\phi\Eop}$, a defect OPE coefficient (Part~V); scale factors in Part~IV.\\
$\tau$ & the regulator of the master path, $t=2\tau$, and of $R_\tau$, $S_\tau$, $C(\tau)$ (Part~III families); $\tau_n(\nu)$, $\tau_n^{(d)}(\nu)$, the zig-zag spectra (Sections~\ref{gff:c1:closedForms}, \ref{gff:sec:P5complete}, \ref{gff:app:generalDimension}); $\tau$, a Taylor coefficient of $\mathfrak G$ (Section~\ref{gff:sec:Kdef}), written $\varrho$ in Section~\ref{gff:mgk:conventions}; $\tau_*(q)$ \eqref{gff:gdk:domain}; local parameters in Sections~\ref{gff:app:generalDimension}, \ref{gff:gs:Z2section} and \ref{gff:gc:L2second}, and the sums $\tau_n$ of Section~\ref{gff:c6d:wheel}.\\
$\phi$, $\varphi$, $\Phi$ & the field $\phi$ (Sections~\ref{gff:c1:notation}, \ref{gff:ope:observable}) and the bulk field $\Phi$ (Part~V); $\phi(x)=\Gamma(\frac14+x)/\Gamma(\frac34+x)$ (Section~\ref{gff:mgk:regularization}); the angle $\phi$ (Section~\ref{gff:mgp:fourier}); $\phi(R,R')$ (Section~\ref{gff:gc:K2proof}); $\varphi_x(s)$ \eqref{gff:gc:L2rim}, $\varphi_\tau(y)$ (Section~\ref{gff:mgk:riemann}), $\varphi(z)$, $\bar\varphi(z)$ \eqref{gff:mgl:phis} and $\varphi_m$ (Section~\ref{gff:gdr:special}); $\Phi_2$ (Section~\ref{gff:gc:d2closedForms}), $\Phi_t(m)$, $\widehat\Phi$ (Sections~\ref{gff:gdk:Rfamily}--\ref{gff:gdk:continuation}), $\Phi_-$, $\Phi_0$, $\Phi_+$ (Section~\ref{gff:c6:analytic}) and $\Phi_{\le4}$ (Section~\ref{gff:c8d3:independent}); local meanings in Sections~\ref{gff:app:selberg}, \ref{gff:gc:L2proof} and \ref{gff:mgf:nodes}.\\
$\psi$, $\Psi$ & $\psi_m$, the polygamma functions (first block), written $\psi=\psi_0$ in Sections~\ref{gff:app:C7d2}, \ref{gff:ope:lowOrders} and \ref{gff:c6d:seventh}; in the proof of the twofold-convolution lemma of Section~\ref{gff:mgp:fourier}, $\psi=\theta-\frac\phi2$ and $\psi_0=\pi-\frac\phi2$ are angles; $\Psi_\sigma(s)$ (Section~\ref{gff:c1:c6analytic}), $\Psi(\sigma)$ (Section~\ref{gff:mgp:fourier}), $\Psi_{r,m}(v;b)$ (Section~\ref{gff:gdr:halfShift}) and $\Psi_t(n)$, $\widehat\Psi$ (Sections~\ref{gff:gdk:Rfamily}--\ref{gff:gdk:continuation}).\\
$\omega$, $\Omega$ & $\omega(\rho)$, the constant of the Gamma-ratio bound \eqref{gff:mgf:gammaBounds} (Sections~\ref{gff:gdr:periodization}, \ref{gff:gdk:growthProof}), against $\omega_j$, the sine shifts of the top-mode proof (Section~\ref{gff:gdr:top}); $\omega=w(\kappa)+\log2$ (Section~\ref{gff:mgp:halfStrip}); $\omega_\tau$ (Section~\ref{gff:gc:K2proof}); $\omega=\pm1$ (Section~\ref{gff:app:threeDimensionalReduction}); $\omega_b$, $\omega_c$ (Appendix~\ref{gff:app:certificate}); $\Omega_r(t)$, $\Omega(t;x)$, $\Omega_r(t;x)$, $\Omega_{\rm R}$, periodized wheels (Sections~\ref{gff:fp:twoDimensional}, \ref{gff:gc:L2proof}, \ref{gff:gdr:special}, \ref{gff:gdk:special}); the simplex measure $d\Omega$ (Part~IV).\\
sans-serif letters & local to the section that defines them, except $\pA$, $\pD$, $\pT$, $\pU$ (Section~\ref{gff:c1:closedForms}), $\BB$ and $\Cat=\mathsf G$. Part~V uses $\mathsf H_j$, $\mathsf e_j$, $\mathsf f_j$, $\mathsf l_j$, $\mathsf S_n$, $\mathsf w_n$, $\mathsf L_n$, $\mathsf J_d$, $\mathsf Y$, $\mathsf C$, $\mathsf K_\times$, $\mathsf K_\parallel$, $\mathsf c_\times$, $\mathsf a_n$, $\mathsf r_1$, $\mathsf r_2$, $\mathsf F_{\rm sch}$, $\mathsf f^{\rm sch}_k$, to keep them apart from $h_d$, $B_\pm$, $\ell_j$, $W_r$, $\LL_d$, $\JJ$, $H$ and $C_n$. Part~III uses its own in Sections~\ref{gff:app:generalClosed}--\ref{gff:gdk:section}, for example $\mathsf c_j$ \eqref{gff:gc:L2rim}, $\mathsf X_r$, $\mathsf r_k$, $\mathsf s_k$, $\mathsf P_\tau$, $\mathsf F_i$, $\mathsf E$, $\mathsf N_q$ and the sequences $\mathsf x_m$, $\mathsf y_j$, $\mathsf z_n$ \eqref{gff:mgf:threeFamilies}; Section~\ref{gff:app:C7d2} uses $\mathsf P_2$, $\mathsf P_3$, $\mathsf p_2$, $\mathsf p_3$, and Appendix~\ref{gff:app:C6derivation} the matrices $\mathsf M$, $\mathsf N$. Clashes: the strip $\mathsf T$ (Section~\ref{gff:mgp:halfStrip}) and the transpose ${}^{\mathsf T}$ have the glyph of $\pT=\psi_1(q)$; the arrays $\mathsf A_{M,\alpha}$, $\mathsf B_{M,\alpha}$ \eqref{gff:gc:finiteClosure} and $\mathsf U_\pm$, $\mathsf V_\pm$ (Section~\ref{gff:fp:arraysSection}) are not $\pA$, $\BB$, $\pU$; $\mathsf J_0$, $\mathsf J_1$ (Section~\ref{gff:gc:L2second}) are not $\mathsf J_d$; $\mathsf a_j$ (Section~\ref{gff:gc:L2second}) is not $\mathsf a_n$; $\mathsf r_1$, $\mathsf r_2$ are not $\mathsf r_k$; and $\mathsf P_x$, $\mathsf Q_x$ (Section~\ref{gff:mgp:fourier}) are functions, not the pole sets $\mathsf P_T$, $\mathsf P_\tau$ (Sections~\ref{gff:gdr:top}, \ref{gff:gdk:Rfamily}).\\
\bottomrule
\end{longtable}
For reference,
\[
 K(m)=\int_0^{\pi/2}\frac{d\theta}{\sqrt{1-m\sin^2\theta}},\quad
 B(u,v)=\frac{\Gamma(u)\Gamma(v)}{\Gamma(u+v)},\quad
 (a)_n=\frac{\Gamma(a+n)}{\Gamma(a)}.
\]
Our ordinary hypergeometric convention is
\[
 {}_pF_q\!\left(\begin{matrix}a_1,\ldots,a_p\\b_1,\ldots,b_q\end{matrix};x\right)
 =\sum_{n=0}^\infty\frac{\prod_i(a_i)_n}{\prod_j(b_j)_n}\frac{x^n}{n!}.
\]
For $p=q+1$, positive real parametric excess ensures absolute convergence
at $x=1$ in the nonsingular cases used here \cite{gff:DLMF16}. Repeated
parameters may not be cancelled along a deformation unless they agree as
functions of that deformation, not merely at its evaluation point.
Contours, endpoint subtractions and continuation rules are part of the
master definitions, not optional numerical conventions. In particular,
$1/4\mapsto d/4$ alone does not generalize a one-dimensional formula:
$3d/4$ and $1-d/4$ coincide only at $d=1$.

\paragraph{Evidence and reproducibility.}
The derivations and definitions below specify the retained calculation.
Published lower-order results and general mathematical identities are
cited separately from this project's higher-loop claims. Algebraic or
master-level agreement does not independently validate every graph.
The public repository of this work contains the code, graph arrays,
numerical inputs and records behind every quoted number;
Section~\ref{gff:cn:reproduction} gives the package map, in which package paths are
relative to the folder \path{GFF/code/} of that repository, and the
provenance of every merged block. These files supply
reproducibility data, not unstated definitions in the formulas.

\section{Renormalization-group derivation}\label{gff:sec:RG}
Normalize the free covariance and perturbations by
\begin{equation}
 \langle\phi(x)\phi(0)\rangle_0=|x|^{-h},\qquad
 h=\frac{1-\eps}{2},\qquad
 S_{\rm int}=\int\dd x\left(\frac{g_0}{4!}\phi^4+\frac{r_0}{2}\phi^2\right).
\end{equation}
After subtraction of the critical local mass power divergence, the long-range kinetic term has no wave-function counterterm. This is the field-theory framework of Ref.~\cite{gff:BGHS}. In the mass-anomalous-dimension convention used here,
\begin{align}
 \beta_{\rm RG}(g)&=-\eps g+\sum_{L\geq1}b_Lg^{L+1},&
 \gamma_r(g)&=\sum_{L\geq1}c_Lg^L,\label{gff:eq:RGfunctions}\\
 \beta_{\rm RG}(g_*)&=0,&
 \Delta_{\Eop}&=1-\nu^{-1}=\frac{1-\eps}{2}+\gamma_r(g_*).
 \label{gff:eq:thermal}
\end{align}
Here $b_4,c_4$ are \emph{loop coefficients}, whereas $C_4,C_5$ are coefficients in the physical $\eps$ expansion. None is Euler's beta function
$B(u,v)=\Gamma(u)\Gamma(v)/\Gamma(u+v)$.

\subsection{Lagrange inversion}
Put
\begin{equation}
 \mathcal B(g)=b_1+b_2g+b_3g^2+b_4g^3+b_5g^4+\cdots,
 \qquad \eps=g_*\mathcal B(g_*).
\end{equation}
Then, for $n\geq1$,
\begin{equation}\boxed{
 C_n=-\tfrac12\delta_{n1}
 +\frac1n[g^{n-1}]\frac{\gamma_r'(g)}{\mathcal B(g)^n}.
}\label{gff:eq:Lagrange}\end{equation}
This follows directly by applying Lagrange--B\"urmann inversion to $\gamma_r(g_*(\eps))$. Equivalently, for $g_* = \sum_j u_j\eps^j$,
\begin{align}
 u_1&=b_1^{-1},&u_2&=-b_2b_1^{-3},\notag\\
 u_3&=2b_2^2b_1^{-5}-b_3b_1^{-4},\notag\\
 u_4&=-5b_2^3b_1^{-7}+5b_2b_3b_1^{-6}-b_4b_1^{-5},\notag\\
 u_5&=14b_2^4b_1^{-9}-21b_2^2b_3b_1^{-8}
 +(6b_2b_4+3b_3^2)b_1^{-7}-b_5b_1^{-6}.
\end{align}
The coefficients through three loops in this normalization are
\begin{align}
 b_1&=3,&c_1&=1,&b_2&=-6A,&c_2&=-A,\notag\\
 b_3&=6Q+\frac{87}{4}A^2-\frac34\pi^2-4\pi,&
 c_3&=\frac72A^2-\frac16\pi^2-\frac43\pi.
 \label{gff:eq:bc3}
\end{align}
Consequently,
\begin{equation}\boxed{
 C_4=\frac{60A^3-8Ab_3+18Ac_3+3c_4-b_4}{243},
}\label{gff:eq:C4RG}\end{equation}
and
\begin{equation}\boxed{\begin{aligned}
 C_5=\frac1{729}\bigl[&336A^4-60A^2b_3+108A^2c_3+b_3^2-3b_3c_3\\
 &-10Ab_4+24Ac_4+3c_5-b_5\bigr].
\end{aligned}}\label{gff:eq:C5RG}\end{equation}
At fifth order the new combination is $D_5=3c_5-b_5$. The four-loop combination $D_4=3c_4-b_4$, sufficient for $C_4$, is not sufficient by itself for $C_5$. Extending a four-loop fixed-point series without five-loop input would miss $D_5/729$.

\section{The complete graph and subtraction prescription}\label{gff:sec:graphs}
This section fixes the full four- and five-loop calculation, including nonprimitive graphs. Write $z$ for an auxiliary regulator, equal to $\eps$ on the physical regulator path, and set $h=(1-z)/2$.

\subsection{Ordered graphs and exact weights}
For loop order $L$, let $\Gr_L$ be the set of symmetric edge-multiplicity matrices on $0,\ldots,L$ satisfying
\begin{equation}\begin{gathered}
 m_{ii}=0,\quad m_{ij}\in\{0,1,2,3\},\quad
 \sum_{i<j}m_{ij}=2L,\quad d_i:=\sum_jm_{ij}\in\{2,3,4\},\\
 \sum_{\substack{i\in S\\j\notin S}}m_{ij}\geq2
 \quad\hbox{for every nonempty proper vertex subset }S.
\end{gathered}\label{gff:eq:graphs}\end{equation}
The cut condition expresses connectedness and one-particle irreducibility. Put $e_i=4-d_i$, so $\sum_i e_i=4$, and use the rational weights
\begin{equation}\boxed{
 w_4(m)=\frac{24}{\prod_i e_i!\prod_{i<j}m_{ij}!},\qquad
 w_2(m)=\frac{w_4(m)}{12}\sum_i e_i(e_i-1).
}\label{gff:eq:weights}\end{equation}
The mass insertion can be represented by choosing two external half-edges at the same vertex of a quartic graph, which yields the second weight. All ordered vertex labels are retained. The perturbative vertex factorial cancels the sum over orderings; there is no additional factor $(L+1)!$ in the following ordered integrals. The enumerated counts for $L=1,\ldots,5$ are
\begin{equation}
 |\Gr_L|=1,\;6,\;85,\;1890,\;59145.
\end{equation}
These are not counts of isomorphism classes.

\paragraph{A zero-dimensional benchmark of the weights.}
The weights \eqref{gff:eq:weights} can be checked without enumerating graphs. In
zero dimensions put
\[
 M_j(g)=\frac1{\sqrt{2\pi}}\int_{\RR}x^j
 \exp\Bigl[-\frac{x^2}{2}-\frac g{24}\bigl(x^4-6x^2+3\bigr)\Bigr]\dd x;
\]
the normal-ordered interaction removes tadpoles. With $r$ coupled to $x^2/2$,
the Legendre transformation gives the amputated one-particle-irreducible
four-point function and the mass insertion,
\begin{equation}
 \Gamma^{(4)}=3\frac{M_0^2}{M_2^2}-\frac{M_4M_0^3}{M_2^4},\qquad
 \partial_r\Gamma^{(2)}\big|_{r=0}=\frac{M_0M_4-M_2^2}{2M_2^2},
 \label{gff:eq:zeroDim}
\end{equation}
whose Gaussian moments give the exact series
\begin{align*}
 \Gamma^{(4)}={}&g-\tfrac32g^2+\tfrac{15}{4}g^3-\tfrac{105}{8}g^4+\tfrac{893}{16}g^5
 -\tfrac{8795}{32}g^6+\tfrac{97987}{64}g^7-\tfrac{10951025}{1152}g^8
 +\tfrac{49985915}{768}g^9-\cdots,\\
 \partial_r\Gamma^{(2)}={}&1-\tfrac12g+\tfrac34g^2-\tfrac{49}{24}g^3+\tfrac{343}{48}g^4
 -\tfrac{2875}{96}g^5+\tfrac{83711}{576}g^6-\tfrac{2763929}{3456}g^7
 +\tfrac{33994099}{6912}g^8-\cdots.
\end{align*}
At every loop order the weights obey
\begin{equation}
 \sum_{m\in\Gr_L}w_4(m)=(-1)^L(L+1)!\,[g^{L+1}]\Gamma^{(4)},\qquad
 \sum_{m\in\Gr_L}w_2(m)=(-1)^L(L+1)!\,[g^{L}]\,\partial_r\Gamma^{(2)}.
 \label{gff:eq:zeroDimSums}
\end{equation}
A class $G$ with automorphism group $\Aut(G)$ has $(L+1)!/|\Aut(G)|$ vertex
labelings, so the class weights \eqref{gff:c7f:weights} sum to
$(-1)^L[g^{L+1}]\Gamma^{(4)}$ and $(-1)^L[g^L]\partial_r\Gamma^{(2)}$. For
$L=1,\ldots,5$ the ordered sums are $3$, $\frac{45}{2}$, $315$,
$\frac{13395}{2}$, $\frac{395775}{2}$ and $1$, $\frac92$, $49$,
$\frac{1715}{2}$, $\frac{43125}{2}$; the five- and six-loop sums of the
enumeration in the table after \eqref{gff:c6:weights} are those of
\eqref{gff:eq:zeroDimSums}. A second enumerator, which builds the adjacency
matrices row by row and tests connectedness and the cut condition of
\eqref{gff:eq:graphs} by a bridge search, rather than building the graphs edge by
edge and testing every cut, reproduced every ordered-label graph and every
rational weight \eqref{gff:eq:weights}; its record is not part of the repository.
The class inventories of the eighth-order calculation
for $L=1,\ldots,8$ (Section~\ref{gff:c8f:changes}), whose loops one to seven
reproduce those of the seventh-order calculation, satisfy the class form
exactly (\path{LRI_C8_HPC/gen/zero_dim_check.py}, with its output in
\path{LRI_C8_HPC/inputs8/verification/logs/c8gen_fin_54591916.out}). The
zero-dimensional integral is a combinatorial test of the weights, not an
approximation to a critical exponent.

With $x_0=0,x_L=1$, define
\begin{equation}\begin{aligned}
 \JJ_m(z)&=\int_{0<x_1<\cdots<x_{L-1}<1}
 \prod_{i<j}(x_j-x_i)^{-hm_{ij}}\dd^{L-1}x,\\
 I_m(z)&=\frac{\JJ_m(z)}{Lz},\qquad
 \mathcal A_L(z)=\sum_m w_4(m)I_m(z),\qquad
 \mathcal B_L(z)=\sum_m w_2(m)I_m(z).
\end{aligned}\label{gff:eq:graphintegral}\end{equation}
The factor $1/(Lz)$ is the overall scale integral with the total vertex span cut off at one. Translation has been removed. The mass power divergence is treated by analytic continuation and critical mass subtraction, rather than by discarding the associated nontrivial graphs.

\subsection{Minimal subtraction}
The bare amplitudes are
\begin{equation}
 F_4(v)=v+\sum_{L\geq1}(-1)^L\mathcal A_L(z)v^{L+1},\qquad
 F_2(v)=1+\sum_{L\geq1}(-1)^L\mathcal B_L(z)v^L,
\end{equation}
with $\mathcal A_1=3/z$, $\mathcal B_1=1/z$. Write
\begin{equation}
 g_0=\mu^z\mathcal Z_g(g,z),\quad r_0=\mu^sr Z_r(g,z),\quad
 \mathcal Z_g=g+\sum_{n\geq2}G_n(z)g^n,\quad
 Z_r=1+\sum_{n\geq1}Z_n(z)g^n.
\end{equation}
The coefficients $G_n,Z_n$ contain only pole terms. If $\KP$ selects negative Laurent powers of $z$, finiteness of $F_4(\mathcal Z_g)$ and $Z_r F_2(\mathcal Z_g)$ gives
\begin{equation}
 G_n=-\KP[g^n]F_4(\mathcal Z_{g,<n}),\qquad
 Z_n=-\KP[g^n]\bigl[Z_{r,<n}F_2(\mathcal Z_g)\bigr].
 \label{gff:eq:MSrecursion}
\end{equation}
The RG coefficients follow from the simple poles:
\begin{equation}\boxed{
 b_L=L[z^{-1}]G_{L+1},\qquad c_L=L[z^{-1}]Z_L.
}\label{gff:eq:simplepoles}\end{equation}
Equivalently, the pole equations are
\begin{equation}
 z\mathcal Z_g+\beta_{\rm RG}\partial_g\mathcal Z_g=0,\qquad
 \beta_{\rm RG}\partial_g Z_r+\gamma_r Z_r=0.
 \label{gff:eq:poleeqs}
\end{equation}
They imply the useful recursions
\begin{equation}
 G_n=\frac{\sum_{\ell=1}^{n-1}b_\ell(n-\ell)G_{n-\ell}}{(n-1)z},\qquad
 Z_n=\frac{\sum_{\ell=1}^{n}[(n-\ell)b_\ell+c_\ell]Z_{n-\ell}}{nz},
\end{equation}
where $G_1=Z_0=1$.


\subsection{Primitive versus nonprimitive: the precise distinction}
A primitive diagram has an overall ultraviolet divergence and no divergent proper subdiagram. A nonprimitive diagram contains subdivergences and must be evaluated together with its counterterms. For example, a two-edge bubble has the local integral
\[
\int_0^\ell r^{-1+z}\,dr=\frac1z+\log\ell+O(z).
\]
If it occurs inside a larger graph, its pole subtraction multiplies the quotient graph. In particular,
\[
\left(-\frac pz\right)\left(\frac{q_{-1}}z+q_0+O(z)\right)
=-\frac{pq_{-1}}{z^2}-\frac{pq_0}{z}+O(1),
\]
so the larger graph's new simple pole depends on the lower-loop \emph{finite part} $q_0$. Primitive periods alone do not determine it. The forest formula implements these products systematically \cite{gff:Herzog}.

The phrase ``remaining nonprimitive'' used during the calculation also included the lower-loop fixed-point terms. Precisely,
\[
C_5=C_5^{\rm primitive}+\frac{\mathcal R-10Ab_4+24Ac_4+3c_5^{\rm nonprim}-b_5^{\rm nonprim}}{729},
\]
where $\mathcal R=336A^4-60A^2b_3+108A^2c_3+b_3^2-3b_3c_3$. The later separation into primitive, wheel-plus-assembly, and rest sectors is stated in Section~\ref{gff:sec:assembly}. Those pieces are bookkeeping within the chosen normalization; the universal object is their sum.

\subsection{Laurent orders needed for the fifth coefficient}
Let $a_{L,k}=[z^k]\mathcal A_L$ and $d_{L,k}=[z^k]\mathcal B_L$. Expanding \eqref{gff:eq:MSrecursion} gives
\begin{align}
 b_4&=36Aa_{2,0}-216a_{2,1}+48a_{3,0}-4a_{4,-1},\notag\\
 c_4&=26Ad_{2,0}-140d_{2,1}+40d_{3,0}-4d_{4,-1},\label{gff:eq:b4c4poles}\\
 b_5&=5a_{5,-1}-75a_{4,0}+450a_{3,1}-60Aa_{3,0}
       -1350a_{2,2}+585Aa_{2,1}-5b_3a_{2,0},\notag\\
 c_5&=5d_{5,-1}-65d_{4,0}+325d_{3,1}-\frac{95}{2}Ad_{3,0}
       -\frac{2275}{3}d_{2,2}\notag\\
    &\hspace{12mm}+\frac{725}{2}Ad_{2,1}
       -\frac{10b_3+5c_3}{3}d_{2,0}.\label{gff:eq:b5c5poles}
\end{align}
The lower-loop data must therefore be retained through
\begin{center}
\begin{tabular}{@{}cl@{}}
\toprule
Loop order & Required Laurent powers in $\mathcal A_L,\mathcal B_L$\\
\midrule
2 & $z^{-2},z^{-1},1,z,z^2$\\
3 & $z^{-3},z^{-2},z^{-1},1,z$\\
4 & $z^{-4},z^{-3},z^{-2},z^{-1},1$\\
5 & $z^{-5},z^{-4},z^{-3},z^{-2},z^{-1}$\\
\bottomrule
\end{tabular}
\end{center}
The finite four-loop terms $a_{4,0},d_{4,0}$ are new inputs relative to the computation of $C_4$. Simplifying $C_4$ alone cannot determine them.

\paragraph{What the higher-pole identities cannot detect.}\label{gff:mgc:higherPoles}
The pole equations \eqref{gff:eq:poleeqs} fix the higher poles of $G_n$ and $Z_n$ from the lower-order
simple poles. They do not fix a new simple pole. For example, the shift
\[
 b_5\longmapsto b_5+\delta b_5,\qquad
 \mathcal Z_g\longmapsto\mathcal Z_g+\frac{\delta b_5}{5z}\,g^6
\]
preserves $z\mathcal Z_g+\beta_{\rm RG}\partial_g\mathcal Z_g=0$ through $g^6$: the three new terms
are $\tfrac15\delta b_5g^6$, $-\tfrac65\delta b_5g^6$ and $\delta b_5g^6$. It changes the physical
coefficient by $\delta C_5=-\delta b_5/729$. More generally $\delta C_5=(3\delta c_5-\delta b_5)/729$
when the lower loops are held fixed. Agreement of the higher poles with their renormalization-group
values, used as a check in Sections~\ref{gff:sec:checks} and~\ref{gff:c6:checks} and at seventh and eighth
order, therefore tests the lower-order inputs and the pipeline but cannot detect missing new
simple-pole information. This caveat was written by OpenAI GPT-6 astra in the research archive behind
the sixth-order packages (package \path{LRI_C6_D2_HPC/reference/}).


\SourcePart{One-dimensional analytic coefficients through fifth order}

\section{Final one-dimensional formulas}\label{gff:cn:oneDresults}

We use
\[
\eps=2s-1,\qquad \Delta_\phi=\frac{1-\eps}{4},\qquad
\Delta_{\mathcal E}=\frac12-\frac\eps6+C_2\eps^2+C_3\eps^3+C_4\eps^4+C_5\eps^5+O(\eps^6).
\]
Here $\mathcal E=[\phi^2]$ is the thermal operator. The covariance is
$\langle\phi(x)\phi(0)\rangle_0=|x|^{-(1-\eps)/2}$, and the interaction is $g\phi^4/4!$.
Set
\begin{equation}
 A=\pi+4\log2,\qquad Q=\frac{\Gamma(1/4)^4}{\pi},\qquad
 H=\frac{2\pi^2}{3}+8\Cat,
 \qquad \Cat=\sum_{n=0}^{\infty}\frac{(-1)^n}{(2n+1)^2}.
\end{equation}
Thus $\Cat$ is Catalan's constant. It is not a graph amplitude.

The complete reduction obtained from the graph calculation is
\begin{equation}\boxed{\begin{aligned}
46656 C_5={}&-1536(P_5+\KK)+384AP_4-576\LL-1728Q^2\\
&+48Q(15A^2+149H+7\pi^2)-15A^4-30A^2\pi^2\\
&-828A\pi^3-21456A\zeta(3)-83\pi^4.
\end{aligned}}\label{gff:eq:full}
\end{equation}
The four constants $P_4,P_5,\LL,\KK$ are evaluated in Sections~\ref{gff:sec:Kdef}--\ref{gff:sec:P5complete}; their graph-normalized definitions and complete reductions are included. The preferred fixed-parameter presentation of $P_5$ is given next.

\subsection{\texorpdfstring{The five fixed hypergeometric values for $P_5$}{The five fixed hypergeometric values for P5}}
For $r\in\{1,2,3,5,7\}$ define
\begin{equation}\boxed{\begin{aligned}
\mathcal Y_r={}&\frac r8
 \left[\frac{\Gamma((r+2)/8)}{\Gamma((r+6)/8)}\right]^7\\
&\times{}_9F_8\!\left(
\begin{matrix}1,\,1+r/8,\,((r+2)/8)^{\times7}\\
r/8,\,((r+6)/8)^{\times7}\end{matrix};-1\right).
\end{aligned}}\label{gff:eq:P5FixedY}\end{equation}
Put $c_8=\cos(\pi/8)$ and $s_8=\sin(\pi/8)$. Then
\begin{equation}\boxed{\begin{aligned}
P_5=16\pi^{3/2}\Big\{&-35\mathcal Y_2\\
&+(22\sqrt2-13)\big(c_8^7\mathcal Y_3-s_8^7\mathcal Y_5\big)\\
&+(22\sqrt2+13)\big(c_8^7\mathcal Y_1+s_8^7\mathcal Y_7\big)\Big\}.
\end{aligned}}\label{gff:eq:P5FixedAnswer}\end{equation}
Every parameter is fixed and positive; all arguments are the ordinary
$-1$. Each series has excess $3/2$ and is absolutely convergent. For
$r=2$ one matching numerator/denominator parameter $1$ cancels, so
\begin{equation}
\mathcal Y_2=\frac{\pi^{7/2}}4\,
{}_8F_7\!\left(\begin{matrix}\frac54,(\frac12)^{\times7}\\
\frac14,1^{\times6}\end{matrix};-1\right).
\label{gff:eq:P5FixedY2}\end{equation}
Thus \eqref{gff:eq:P5FixedAnswer} uses four ${}_9F_8$ values and one
${}_8F_7$ value. It has no parameter derivative, Taylor-coefficient
extraction, coalescence limit, extra summation, multivariate function,
or unstated analytic-continuation sheet. Section~\ref{gff:sec:P5FixedProof}
proves the identity using a regularized bilateral sum with exactly
three Fourier harmonics. Inserting it in the polygamma form of $C_5$
makes the complete result an ordinary fixed-parameter
hypergeometric expression. No elementary evaluation or minimality of
this five-value basis is asserted.

\subsection{The reflected-quarter polygamma basis}
\label{gff:sec:psiAll}
The following is the derivative-free form of every coefficient $C_0,\ldots,C_5$.
The evaluation of $\KK$ used here is proved in Section~\ref{gff:app:KpsiProof},
and the wheel reduction in Section~\ref{gff:app:LreductionProof}.

\subsubsection{A basis adapted to \texorpdfstring{$1/4$ and $3/4$}{1/4 and 3/4}}
Put $q=1/4$ and $\bar q=1-q=3/4$. With the convention
$\psi_m(z)=d^{m+1}\log\Gamma(z)/dz^{m+1}$, define
\begin{equation}\boxed{\begin{aligned}
 \pA&=\psi_0(2q)+\psi_0(1)-2\psi_0(q),&
 \pD&=\psi_0(\bar q)-\psi_0(q),\\
 \pT&=\psi_1(q),&\overline{\pT}&=\psi_1(\bar q),\\
 \pU&=\psi_2(q),&\overline{\pU}&=\psi_2(\bar q).
\end{aligned}}\label{gff:eq:psiQuarterBasis}\end{equation}
An overline here means evaluation at the reflected argument, not complex
conjugation. There is no differentiation with respect to $q$ in the
coefficient formulas. At $q=1/4$, the exact dictionary is
\begin{equation}\begin{aligned}
 \pA&=A, &\pD&=\pi,\\
 \pT+\overline{\pT}&=2\pi^2,&
 \pT-\overline{\pT}&=16\Cat,\\
 \pU+\overline{\pU}&=-112\zeta(3),&
 \overline{\pU}-\pU&=4\pi^3,\\
 H&=\frac{5\pT-\overline{\pT}}6,&
 \log2&=\frac{\pA-\pD}{4}.
\end{aligned}\label{gff:eq:psiQuarterDictionary}\end{equation}
These follow from the polygamma reflection and multiplication identities
\cite{gff:DLMFPolygamma,gff:DLMFDigamma}. An entirely quarter-argument version of
$\pA$, apart from the reference value at $1$, is
\[
 \pA=\frac{4\psi_0(1)-5\psi_0(q)+\psi_0(\bar q)}{3}
 \qquad(q=1/4).
\]
We retain the gamma normalization $Q$ and the ordinary periods $P_4,P_5$.
For clarity, the wheel normalizations are
\[
 B_-=\Gamma(q)^4F_-,\qquad
 B_+=\frac{64\pi^4}{\Gamma(q)^4}F_+,\qquad P_4=13B_--3B_+.
\]
If desired as a uniform notation, $Q=\Gamma(q)^4/\pD$ and
$B_+=64\pD^4F_+/\Gamma(q)^4$ at $q=1/4$. This is not an evaluation of
$Q$ or the periods in polygamma values alone.

\subsubsection{The two master reductions}
The proved polygamma evaluation and the digamma-prefactor wheel identity are
\begin{equation}\boxed{\begin{aligned}
 \KK&=\frac12\left[\pT(3\pT-\overline{\pT})+\pA\pU\right],\\[1mm]
 \LL&=-\frac{3\pA-\pD}{6}P_4-\frac{2\pD}{3}B_-.
\end{aligned}}\label{gff:eq:KLpsiReflected}\end{equation}
Equivalently,
\[
 \LL=-\frac{13\pA-3\pD}{2}B_-
       +\frac{3\pA-\pD}{2}B_+.
\]
The first identity eliminates the entire six-${}_3F_2$ parameter derivative
in the definition of $\KK$. The second removes the parameter derivative in the
definition of the logarithmic wheel but does not eliminate the two wheel
periods $B_\pm$. Both are equalities to the original master
definitions, not merely equalities among proposed right-hand sides.

\subsubsection{\texorpdfstring{$C_0$ through $C_5$}{C0 through C5}}
Writing $\Delta_{\mathcal E}=\sum_{n=0}^5 C_n\eps^n+O(\eps^6)$, the
additional coefficient formulas are
\begin{equation}\boxed{\begin{aligned}
 C_0&=\frac12,& C_1&=-\frac16,& C_2&=\frac{\pA}{9},\\[1mm]
 C_3&=\frac{\pA^2+(\pT+\overline{\pT})/6-8Q}{108}.
\end{aligned}}\label{gff:eq:C0to3psi}\end{equation}
\begin{equation}\boxed{\begin{aligned}
648C_4={}&16P_4-\pA^3-\frac{\pA}{2}(\pT+\overline{\pT})-8\pA Q\\
&+\frac{365\pU-6\overline{\pU}}{14}.
\end{aligned}}\label{gff:eq:C4psiQuarter}\end{equation}
\begin{equation}\boxed{\begin{aligned}
46656C_5={}&-1536P_5+96(7\pA-\pD)P_4+384\pD B_- -1728Q^2\\
&+16Q(45\pA^2+383\pT-64\overline{\pT})
 -768\pT(3\pT-\overline{\pT})\\
&-15\pA^4-15\pA^2(\pT+\overline{\pT})\\
&-\frac{6\pA}{7}(431\pU+18\overline{\pU})
 -\frac{83}{4}(\pT+\overline{\pT})^2.
\end{aligned}}\label{gff:eq:C5psiQuarter}\end{equation}
Section~\ref{gff:app:KpsiProof} proves the master identity used in this substitution. All explicit logarithmic, zeta, Catalan, and power-of-$\pi$
coefficients in these expressions are written using fixed-argument
$\psi_0,\psi_1,\psi_2$. The non-polygamma input consists of the gamma
normalization $Q$, the two wheel periods, and $P_5$; their original
definitions and normalizations have not changed.

For checking \eqref{gff:eq:C5psiQuarter}, the wheel term is exactly
\[
384AP_4-576\LL=96(7\pA-\pD)P_4+384\pD B_-.
\]
An exact symbolic check
(\path{lri_psi_complete_update/verification/symbolic_checks.json}) compares each
$C_n$ with the formulas of Section~\ref{gff:cn:oneDresults}, in particular
\eqref{gff:eq:full}, after substituting \eqref{gff:eq:psiQuarterDictionary} and
\eqref{gff:eq:KLpsiReflected}. In particular, no arithmetic inference is based
on fitting decimal values of the periods.

\section{The full-line Selberg formula used in the expansions}\label{gff:app:selberg}
For
\[
 M(u,v,w)=\iint_{\RR^2}|xy|^{-u}|(1-x)(1-y)|^{-v}|x-y|^{-w}\dd x\dd y,
\]
let
\begin{align}
 U&=S_2(1-u,1-v,-w/2),\qquad c=\cos(\pi w/2),\\
 \theta_u&=\pi(u+w/2),\quad\theta_v=\pi(v+w/2),\\
 T&=\pi(u+v+w),\qquad \Phi=\pi(u+v+w/2),
\end{align}
where
\[
S_2(a,b,c)=\frac{\Gamma(a)\Gamma(b)\Gamma(a+c)\Gamma(b+c)\Gamma(1+2c)}
{\Gamma(a+b+c)\Gamma(a+b+2c)\Gamma(1+c)}.
\]
The real-line chamber sum is
\begin{equation}\begin{aligned}
 M=U\bigg[1&-\frac{2c(\sin\theta_u+\sin\theta_v)}{\sin T}\\
 &+\frac{\sin\theta_v\sin(\pi v)+\sin\theta_u\sin(\pi u)
 +2c\sin\theta_u\sin\theta_v}{\sin T\sin\Phi}\bigg].
\end{aligned}
\end{equation}
Apparent removable singularities are treated in the whole expression. The elementary Selberg evaluation is standard \cite{gff:Selberg}; the real-line chamber weighting follows by splitting each integration variable into $(-\infty,0),(0,1),(1,\infty)$ and mapping the resulting chambers to the Selberg square; the trigonometric factors are fixed by analytic continuation of those chamber integrals. The formula is stated in full so that no separate chamber table is an input. This formula supplies every symmetric two-variable expansion used here.

\section{Root and subtraction conventions}
For an $L$-loop graph the rooted bare integral has the overall factor $2/(Lz)$, with the two root positions fixed at zero and one and their separation integrated with a unit infrared cutoff. Roots must not both lie in a proper superficially divergent collision set. In the code, the root-safety test includes power-divergent and non-one-particle-irreducible collision sets as well as logarithmic ones. This avoids spurious infrared poles in a rooted integral. The counterterm forest itself still contains the proper logarithmic one-particle-irreducible subgraphs; induced contractions that create tadpoles vanish in the stated prescription.

In a quotient graph, the roots are inherited from the parent graph. Choosing arbitrary new quotient roots before subtracting can change finite terms and is not a legitimate shortcut. Counterterms of proper subgraphs can be computed with their own safe roots, since their local pole coefficients are independent of that infrared choice. These distinctions were enforced when producing the residue certificate and the numerical comparisons.

\section{Why the remaining diagrams need only three second-order expansions}\label{gff:sec:jets}
The nonprimitive graph calculation requires only the following three second-order master expansions.

A vertex connected to two other vertices can be integrated using
\begin{equation}
 \BB(u,v)=\int_\RR|x|^{-u}|1-x|^{-v}\dd x
 =\sqrt\pi\,
 \frac{\Gamma((1-u)/2)\Gamma((1-v)/2)\Gamma((u+v-1)/2)}
 {\Gamma(u/2)\Gamma(v/2)\Gamma((2-u-v)/2)}.
 \label{gff:eq:linebeta}
\end{equation}
It replaces its two edges by one edge of power $u+v-1$. A block with two internal vertices and two boundary vertices reduces to
\begin{equation}
 \JJ(a,b,c,d,e)=\iint_{\RR^2}
 \frac{\dd x\dd y}{|x|^a|1-x|^b|y|^c|1-y|^d|x-y|^e}.
 \label{gff:eq:master}
\end{equation}
A generic chamber decomposition of \eqref{gff:eq:master} is a finite sum of ordinary ${}_3F_2(1)$ functions. The crucial further simplification is that the full five-loop coefficient does \emph{not} require its arbitrary-parameter form: only its second-order regulator expansions at three critical index patterns occur.

Complete the five active indices to the ten edges on vertices $(0,1,x,y,\infty)$, ordered lexicographically:
\begin{equation}\begin{aligned}
\boldsymbol\lambda=(&3-a-b-c-d-e,\ a,\ c,\ b+d+e-1,\ b,\ d,\\
 &a+c+e-1,\ e,\ 2-a-b-e,\ 2-c-d-e).
\end{aligned}\label{gff:eq:completion}
\end{equation}
Every completed vertex has incident power sum two. Real-line conformal changes of variables permit permutations of the five completed vertices, followed by fixing three positions to $0,1,\infty$. This gives an exact $S_5$ action on the five active indices. Zero-edge representatives factor into \eqref{gff:eq:linebeta}.

After this action and the zero-edge factorizations, the critical representatives needed are
\begin{equation}
 T=(1/2,1/2,1/2,1/2,1/2),\quad
 A_0=(1/2,1/2,1/2,1/2,1),\quad
 B_0=(1/2,1/2,1/2,1,1/2).
\end{equation}
The exceptional higher-connectivity graphs are reduced to these same patterns by the regulator-corrected moves in Section~\ref{gff:sec:shifts}.

At five loops a nontrivial two-internal-vertex block is multiplied by at most two further simple-pole factors and the overall scale pole. An analytic remainder of total regulator degree at least three cannot then contribute to the simple pole. Consequently, the three second-order expansions below suffice. This order counting also covers the finite lower-loop terms required by the five-loop forests.

\subsection{Regular tetrahedron pattern: a gamma-valued Hessian}
At $T$, $\JJ=3Q$ and all first index derivatives vanish. On the balanced completed variations, the unique permutation-invariant quadratic form is proportional to $\sum_e(\delta\lambda_e)^2$. Evaluation along a symmetric Selberg path fixes it:
\begin{equation}\boxed{
 \JJ=3Q+\frac{QH}{4}\sum_{e=1}^{10}(\delta\lambda_e)^2+O(3).
}\label{gff:eq:Tjet}
\end{equation}
Here $O(3)$ denotes analytic terms of total regulator degree at least three.

For an explicit check, put
\[
 M(u,v,w)=\JJ(u,v,u,v,w).
\]
The full-line Selberg evaluation gives
\begin{equation}
 M(1/2,1/2,1/2)=3Q,\qquad \nabla M=0,\qquad
 \frac{\nabla^2 M}{3Q}=H
 \begin{pmatrix}2&1&1\\1&2&1\\1&1&1\end{pmatrix}.
 \label{gff:eq:Hessian}
\end{equation}
The appearance of Catalan's constant comes from
$\psi_1(1/4)=\pi^2+8\Cat$ and $\psi_1(3/4)=\pi^2-8\Cat$.
The Selberg formula used to establish these identities is recalled in Section~\ref{gff:app:selberg} \cite{gff:Selberg}.

\subsection{Pattern \texorpdfstring{$A_0$}{A0}: no new regular quadratic constant}
Write
\[
 a=u+p,\quad c=u-p,\quad b=v+q,\quad d=v-q,\quad e=w,
 \qquad \sigma_0=2-2u-w,\quad\sigma_1=2-2v-w.
\]
Subtracting the endpoint pole differences from the symmetric integral gives
\begin{equation}\boxed{
 \JJ=M(u,v,w)+R(w)\left(\frac{p^2}{\sigma_0}+\frac{q^2}{\sigma_1}\right)+O(3),
}\label{gff:eq:Ajet}
\end{equation}
where, through the needed order,
\begin{equation}
 R(w)=\pi^3+28\zeta(3)-\frac A2[\pi^3+28\zeta(3)](w-1)+O((w-1)^2).
\end{equation}
In \eqref{gff:eq:Ajet}, the remainder has total scaling degree at least three when all local regulators are scaled together; higher terms from the endpoint pole residues are included in this notation. The diagonal collision residue is independent of $p,q$ and cancels in the difference. The two endpoint residues yield $R(w)=\partial_x^2\BB(x,w)|_{x=(2-w)/2}$.

After the full endpoint pole parts have been removed, symmetry permits only two additional second-order constants, multiplying $p^2+q^2$ and $pq$. Both vanish. To determine them, set $u=v=1/2$, $w=1-t$, and use the two unique-star paths
\begin{align}
 p=t,\ q=0:&\quad
 \JJ=\BB(1/2+t,1/2)\BB(1-t,1-t),\\
 p=q=t/2:&\quad
 \JJ=\BB((1+t)/2,(1+t)/2)\BB(1-t,1-t).
\end{align}
Expansion of these exact gamma products through $t^2$ fixes both regular quadratic constants to zero. The symbolic verification code is not part of the public repository.

\subsection{Pattern \texorpdfstring{$B_0$}{B0}: exactly one new mixed derivative}\label{gff:sec:Bjet}
Use local coordinates
\begin{align}
 p_1&=1-d,&p_2&=2-b-d-e,\\
 z_1&=3-a-b-c-d-e,&z_2&=2-c-d-e.
\end{align}
Equivalently,
\[
 d=1-p_1,\quad b=1+p_1-e-p_2,\quad c=1+p_1-e-z_2,\quad
 a=e-p_1+p_2+z_2-z_1.
\]
The base point has $p_1=p_2=z_1=z_2=0$, $e=1/2$.
The only local pole hyperplanes are $p_1=0$ and $p_2=0$. A convenient complete pole part is
\begin{equation}\begin{aligned}
 \mathcal P={}&\frac2{p_1}\BB(e+p_2+z_2-z_1,1-p_2)
 +\frac2{p_2}\BB(e,1+p_1-e)-\frac4{p_1p_2}.
\end{aligned}
\end{equation}
After subtracting it, the function is jointly analytic in all five local coordinates.

On $z_2=0$, the $y$ integration is a unique three-point star. On $z_1=0$, the completed graph has a zero edge. These exact restrictions are
\begin{align}
 U_2(z_1)&=\BB(1+p_1-e,1-p_1)\BB(e+p_2-z_1,1-p_2),\\
 U_1(z_2)&=\BB(e-p_1+p_2+z_2,1+p_1-e-z_2)
           \BB(1+p_1-e-p_2,1-p_1),\\
 U_0&=U_2(0)=U_1(0).
\end{align}
Analytic interpolation between the two restrictions leaves only one second-order mixed term. Explicitly,
\begin{equation}\boxed{\begin{aligned}
 \JJ={}&U_2(z_1)+U_1(z_2)-U_0+z_1z_2\KK\\
 &+\frac2{p_1}\big[\BB(e+p_2+z_2-z_1,1-p_2)
 -\BB(e+p_2-z_1,1-p_2)\\
 &\hspace{16mm}-\BB(e+p_2+z_2,1-p_2)
 +\BB(e+p_2,1-p_2)\big]+O(3).
\end{aligned}}\label{gff:eq:Bjet}
\end{equation}
The subtracted analytic difference vanishes on both $z_1=0$ and $z_2=0$, hence is divisible by $z_1z_2$. Its coefficient at the base point is precisely the one constant $\KK$. Any dependence of that coefficient on another local index raises the total degree to at least three and is not needed.

This explains why the many apparent ordinary-hypergeometric parameter derivatives in the remaining graph classes collapse to a single new constant rather than a large independent basis.

\section{\texorpdfstring{Evaluation of $\mathcal K$}{Evaluation of K}}\label{gff:sec:Kdef}

\subsection{From the original master to six Euler integrals}
In \eqref{gff:eq:Bjet}, choose
\[
 p_1=p_2=z_1=z_2=t,\qquad e=1/2.
\]
Then
\[
 J_t=\JJ(1/2,1/2,1/2,1-t,1/2)
\]
and the gamma interpolation becomes
\begin{equation}
 J_t^\Gamma=\BB(1/2,1-t)^2+
 \frac2t\,[\BB(1/2,1-t)-\BB(1/2+2t,1-t)].
\end{equation}
Thus
\begin{equation}
 J_t-J_t^\Gamma=t^2\KK+O(t^3).\label{gff:eq:Kpath}
\end{equation}

Integrating $x$ in \eqref{gff:eq:master} gives the full-line elliptic star $F(y)$. Fold its three real-line intervals onto $(0,1)$, obtaining
\begin{equation}\begin{aligned}
 J_t=4\int_0^1[K(x)+K(1-x)]\big[&x^{-1/2}(1-x)^{t-1}\\
 &+x^{-t}(1-x)^{t-1}+x^{-1/2}(1-x)^{-t}\big]\dd x.
\end{aligned}\label{gff:eq:Jfold}
\end{equation}
This integral is initially convergent for small positive $t$; its regulated poles are handled explicitly by multiplication by $t^2$.

For the beta-weighted elliptic moment,
\begin{equation}
 E(a,b)=\int_0^1x^{a-1}(1-x)^{b-1}K(x)\dd x,
\end{equation}
Euler integration gives
\begin{equation}
 E(a,b)=\frac\pi2 B(a,b)
 {}_3F_2\!\left(\begin{matrix}1/2,1/2,a\\1,a+b\end{matrix};1\right).
\end{equation}
A Thomae transformation gives the alternative form
\begin{equation}
 E(a,b)=\frac\pi2\frac{\Gamma(b)^2}{\Gamma(b+1/2)^2}
 {}_3F_2\!\left(\begin{matrix}b,b,1-a\\b+1/2,b+1/2\end{matrix};1\right).
\end{equation}
These follow from the standard beta transform and unity-argument hypergeometric identities \cite{gff:DLMF16}.
Apply the appropriate form to each of the six terms in \eqref{gff:eq:Jfold}. Multiplication by $t^2$ gives exactly \eqref{gff:eq:Jhat}. Since
\[
 V_r(t)=t\BB(1/2+rt,1-t),
\]
\eqref{gff:eq:Kpath} proves the regular fourth-derivative formula \eqref{gff:eq:Kdef}.

For an additional check, the gamma interpolation alone has
\begin{equation}
 [t^2]J_t^\Gamma=\frac{A^4}{6}-\frac{A^2\pi^2}{6}
 -\frac{A\pi^3}{3}-12A\zeta(3)-\frac{\pi^4}{180}.
\end{equation}
The endpoint-subtracted logarithmic coefficient of \eqref{gff:eq:Jfold}, minus this expression, is a second numerical route to $\KK$.

\subsection{An analytic representation of the regulated master}
All derivatives in this section act on the displayed parameters; the argument of every hypergeometric function remains one. Define
\begin{align}
 f_1(t)&={}_3F_2\!\left(\begin{matrix}t,t,1/2\\1/2+t,1/2+t\end{matrix};1\right),&
 f_2(t)&={}_3F_2\!\left(\begin{matrix}t,t,t\\1/2+t,1/2+t\end{matrix};1\right),\\
 f_3(t)&={}_3F_2\!\left(\begin{matrix}1/2,1/2,t\\1,1/2+t\end{matrix};1\right),&
 f_4(t)&={}_3F_2\!\left(\begin{matrix}1/2,1/2,t\\1,1\end{matrix};1\right),\\
 f_5(t)&={}_3F_2\!\left(\begin{matrix}1/2,1/2,1/2\\1,3/2-t\end{matrix};1\right),&
 f_6(t)&={}_3F_2\!\left(\begin{matrix}1/2,1/2,1-t\\1,3/2-t\end{matrix};1\right).
\end{align}
Put
\begin{equation}\begin{aligned}
\widehat J(t)=2\pi\bigg\{&
 \frac{\Gamma(1+t)^2}{\Gamma(1/2+t)^2}[f_1(t)+f_2(t)]\\
&+\frac{t\sqrt\pi\,\Gamma(1+t)}{\Gamma(1/2+t)}f_3(t)
+t\Gamma(1+t)\Gamma(1-t)f_4(t)\\
&+\frac{t^2\sqrt\pi\,\Gamma(1-t)}{\Gamma(3/2-t)}[f_5(t)+f_6(t)]\bigg\},
\end{aligned}\label{gff:eq:Jhat}
\end{equation}
and, for $r=0,2$,
\begin{equation}
 V_r(t)=2\sqrt\pi\,
 \frac{\Gamma(1/4-rt/2)\Gamma(1+t/2)\Gamma(1/4+(r-1)t/2)}
 {\Gamma(1/4+rt/2)\Gamma(1/2-t/2)\Gamma(1/4+(1-r)t/2)}.
 \label{gff:eq:V}
\end{equation}
The new constant has the regular parameter-derivative evaluation
\begin{equation}\boxed{
 \KK=\frac1{24}\left.\frac{d^4}{dt^4}
 \left[\widehat J(t)-V_0(t)^2-2\big(V_0(t)-V_2(t)\big)\right]\right|_{t=0}.
}\label{gff:eq:Kdef}
\end{equation}
The bracket starts at order $t^4$. Every gamma factor is regular near zero, and the six hypergeometric series have parametric excess either $1/2$ or $1-t$. Their defining series and the required derivatives therefore converge near the evaluation point \cite{gff:DLMF16}. No singular defining series has to be assigned a value at $t=0$ in \eqref{gff:eq:Kdef}.

Two numerical evaluations that do not use the result of the proof below give
\begin{equation}
 \KK=39.33156912825903905750318715159786152304\ldots:
\end{equation}
a 32-node Cauchy evaluation of the fourth derivative in \eqref{gff:eq:Kdef} at 45
working digits, and the endpoint-subtracted logarithmic integral obtained
from \eqref{gff:eq:Jfold}, which agree to about $4.1\times10^{-33}$ (a 16-node
evaluation at 40 digits differs by $3.7\times10^{-17}$). The hypergeometric
evaluation uses gamma-ratio asymptotics and Hurwitz-zeta tail summation, the
integral explicitly subtracted endpoint expansions. At 110 working digits
the endpoint integral agrees with the proved value \eqref{gff:eq:KprovedFinal} to
$6\times10^{-109}$
(\path{lri_psi_complete_update/verification/numerical_checks.json}).
Equation~\eqref{gff:eq:Kdef} is an intermediate analytic representation; the following proof reduces this exact master to fixed polygamma values.

\paragraph{Conventions and special values.}
We use the derivative-order convention
\begin{equation}
 \psi_m(z)=\frac{d^{m+1}}{dz^{m+1}}\log\Gamma(z),\qquad m=0,1,2,\ldots.
 \label{gff:eq:polygammaConvention}
\end{equation}
Thus $\psi_0=\psi=\Gamma'/\Gamma$ is the digamma function, $\psi_1=\psi'$
is trigamma, and $\psi_2=\psi''$ is tetragamma. These are ordinary
one-variable polygamma functions, not the multivariate functions sometimes
denoted by uppercase $\Psi_1,\Psi_2$. In particular, a square such as
$\psi_1(1/4)^2$ means an ordinary product, not $\psi_2(1/4)$.
The standard series and rational-argument identities
\cite{gff:DLMFPolygamma,gff:DLMFDigamma} give
\begin{equation}\begin{aligned}
 \psi_0(1)&=-\gamma_E,&
 \psi_0(1/2)&=-\gamma_E-2\log2,\\
 \psi_0(1/4)&=-\gamma_E-\frac\pi2-3\log2,&
 \psi_0(3/4)&=-\gamma_E+\frac\pi2-3\log2,\\
 \psi_1(1/4)&=\pi^2+8\Cat,&
 \psi_1(3/4)&=\pi^2-8\Cat,\\
 \psi_2(1/4)&=-2\pi^3-56\zeta(3),&
 \psi_2(3/4)&=2\pi^3-56\zeta(3).
\end{aligned}\label{gff:eq:quarterPolygamma}
\end{equation}
For example, for $m\geq1$,
\[
 \psi_m(z)=(-1)^{m+1}m!\sum_{n=0}^{\infty}(n+z)^{-m-1}.
\]
Splitting the positive odd integers into the residue classes $1$ and $3$
modulo four yields the displayed quarter-argument values. At order two,
$\sum(4n+1)^{-3}+\sum(4n+3)^{-3}=7\zeta(3)/8$ and their difference is
$\pi^3/32$; the latter also follows from the reflection formula in
\cite{gff:DLMFPolygamma}. The project constants therefore obey the exact dictionary
\begin{equation}\boxed{\begin{aligned}
 A&=\psi_0(1/2)+\psi_0(1)-2\psi_0(1/4),\\
 H&=\psi_1(1/4)-\frac{\pi^2}{3}.
\end{aligned}}\label{gff:eq:AHpolygamma}
\end{equation}

\subsection{The compact evaluation and its proof}
\label{gff:app:KpsiProof}
The proof starts from the original master path
\eqref{gff:eq:Kpath} and does not assume a numerical integer relation, the
wheel reduction, or an unproved hypergeometric summation. The main step is
a bilinear identity for a pair of ordinary ${}_3F_2(1)$ functions. The
periodicity argument below includes the growth estimate needed to rule
out a nonconstant periodic factor.

\subsection{Statement and notation}
Write $\ell=\log2$, and define the parameter functions
\begin{equation}\begin{aligned}
 \mathfrak F(t)&={}_3F_2\!\left(\begin{matrix}1/2,1/2,t\\1,1\end{matrix};1\right),\\
 \mathfrak G(t)&={}_3F_2\!\left(\begin{matrix}t,t,t\\1/2+t,1/2+t\end{matrix};1\right).
\end{aligned}\label{gff:eq:KproofFG}\end{equation}
These are $f_4$ and $f_2$, respectively, in the six-function
formula. They are unrelated to the fixed wheel values $F_\pm$ or Catalan's
constant $\Cat$. We will prove
\begin{equation}\boxed{\begin{aligned}
 \KK&=128\Cat^2+24\pi^2\Cat-4\pi^3\ell-28(\pi+4\ell)\zeta(3)\\
 &=2\psi_1(1/4)^2-\pi^2\psi_1(1/4)+\frac A2\psi_2(1/4).
\end{aligned}}\label{gff:eq:KprovedFinal}\end{equation}
The two equalities have different roles: the first evaluates the original
master, and the second is the rational-argument polygamma dictionary.

\subsection{Reducing the six elliptic moments to two}
Keep the \emph{parameter} convention $K(x)$, and put
\[
 E(a,b)=\int_0^1 x^{a-1}(1-x)^{b-1}K(x)\,dx.
\]
For real $0<t<1$, define
\begin{align*}
 I_1&=E(1/2,t),&I_2&=E(t,1/2),&I_3&=E(1-t,t),\\
 I_4&=E(t,1-t),&I_5&=E(1/2,1-t),&I_6&=E(1-t,1/2).
\end{align*}
Then \eqref{gff:eq:Jfold} says $J_t=4\sum_{j=1}^6 I_j$.
Set $s_t=\sin(\pi t)$ and $c_t=\cos(\pi t)$.
The four exact moment relations are
\begin{equation}\boxed{\begin{aligned}
 I_1&=c_t I_3+s_t I_4,& I_5&=s_t I_3-c_t I_4,\\
 I_2&=s_t I_3,& I_6&=s_t I_4.
\end{aligned}}\label{gff:eq:KproofMomentRelations}\end{equation}
Here is a contour proof, including the branch signs. Integrate
$z^{-1/2}(1-z)^{t-1}K(z)$ on an upper-half-plane contour, using principal
powers continued from $0<z<1$. The boundary values for $x>1$ and $x<0$
are \cite{gff:DLMFEllipticConnection,gff:DLMF15}
\[
 K(x+\ii0)=x^{-1/2}\{K(1/x)+\ii K(1-1/x)\},\qquad
 K(x)=(1-x)^{-1/2}K\!\left(\frac{x}{x-1}\right).
\]
The three real intervals contribute
\[
 -\ii I_5+I_1-e^{-\ii\pi t}(I_3+\ii I_4)=0.
\]
Real and imaginary parts give the first two relations.
The analogous contour for $z^{-1/2}(1-z)^{t-1}K(1-z)$ gives
\[
 -I_5-\ii I_6+I_2-e^{-\ii\pi t}I_4=0,
\]
and hence the last two. At $z=0$ the worst singularity is
$z^{-1/2}\log z$; at $z=1$ it is $(1-z)^{t-1}\log(1-z)$.
The small arcs therefore vanish. The large arc is
$O(R^{t-1}\log R)$ and vanishes for $t<1$. These estimates justify the
contours throughout the indicated interval; the resulting identities
continue meromorphically in $t$.

Euler integration and the Thomae transformation already used in the main
text give \cite{gff:DLMF16}
\[
 I_3=\frac\pi2\frac{\Gamma(t)^2}{\Gamma(1/2+t)^2}\mathfrak G(t),\qquad
 I_4=\frac{\pi^2}{2s_t}\mathfrak F(t).
\]
Consequently,
\begin{equation}\boxed{\begin{aligned}
 J_t={}&2\pi(1+c_t+2s_t)
       \frac{\Gamma(t)^2}{\Gamma(1/2+t)^2}\mathfrak G(t)\\
 &+2\pi^2\left(2+\tan\frac{\pi t}{2}\right)\mathfrak F(t).
\end{aligned}}\label{gff:eq:KproofTwoFunctions}\end{equation}
This is an exact meromorphic reduction of the six-function master, not a
truncation in $t$.

\subsection{A contiguous relation with an explicit telescoping certificate}
Write
\[
 a_n(t)=\frac{(1/2)_n^2(t)_n}{(n!)^3},\qquad
 \mathfrak F(t)=\sum_{n=0}^{\infty}a_n(t).
\]
Direct rational algebra proves the termwise identity
\begin{equation}\begin{aligned}
&t^2a_n(t+1)-(2t^2-2t+3/4)a_n(t)+(t-1)^2a_n(t-1)\\
&\hspace{12mm}=
 \frac{(n+1)^3a_{n+1}(t)}{n+t}-\frac{n^3a_n(t)}{n+t-1}.
\end{aligned}\label{gff:eq:KproofTelescoper}\end{equation}
For generic $\Re t<0$ the right-hand side telescopes to zero:
$a_n(t)=O(n^{\Re t-2})$, so the upper boundary is $O(n^{\Re t})$,
and the lower boundary vanishes. Apparent exceptional parameter values
are handled by analytic continuation. Thus
\begin{equation}\boxed{
 t^2\mathfrak F(t+1)-(2t^2-2t+3/4)\mathfrak F(t)
 +(t-1)^2\mathfrak F(t-1)=0.
}\label{gff:eq:KproofRecurrence}\end{equation}
The defining series is holomorphic for $\Re t<1$. This recurrence extends
$\mathfrak F$ meromorphically to the plane, with possible poles only at
positive integers.

\subsection{The discrete Wronskian is constant, not merely periodic}
The function $\mathfrak F(1-t)$ solves the same recurrence. Its discrete
Wronskian with $\mathfrak F(t)$ is
\begin{equation}
 W(t)=t^2\{\mathfrak F(t+1)\mathfrak F(1-t)
                 -\mathfrak F(t)\mathfrak F(-t)\}.
 \label{gff:eq:KproofCasoratian}
\end{equation}
Equation~\eqref{gff:eq:KproofRecurrence} gives $W(t)=W(t-1)$.
To determine this periodic function, first use the Thomae identity
\begin{equation}
 \mathfrak F(1-t)=
 \frac{\Gamma(t)}{\Gamma(1-t)\Gamma(1/2+t)^2}\mathfrak G(t).
 \label{gff:eq:KproofThomae}
\end{equation}
Since $\mathfrak G$ is analytic near zero and $\mathfrak G(0)=1$, we have
$\mathfrak F(1-t)=1/(\pi t)+O(1)$. The apparent singularity of $W$ at
zero is removable and
\[
 W(0)=-\frac1{\pi^2}.
\]
All possible poles of $W$ are integers. Periodicity therefore makes each
one removable, and $W$ is entire.

A growth bound is essential here. The double Euler representation is
\begin{equation}
 \mathfrak F(t)=\frac1{\pi^2}\int_0^1\!\int_0^1
 \frac{(1-uv)^{-t}}{\sqrt{uv(1-u)(1-v)}}\,du\,dv,
 \qquad\Re t<1.
 \label{gff:eq:KproofDoubleEuler}
\end{equation}
It follows by expanding $(1-uv)^{-t}$ and applying two beta integrals;
absolute convergence at $u=v=1$ holds precisely for $\Re t<1$.
In particular, $|\mathfrak F(x+\ii y)|\leq\mathfrak F(x)$.
For $-1/2\leq\Re t\leq1/2$, the four functions
$\mathfrak F(t),\mathfrak F(-t),\mathfrak F(t-1),\mathfrak F(-t-1)$
are uniformly bounded by $\mathfrak F(1/2)$, independently of $\Im t$.
Solving \eqref{gff:eq:KproofRecurrence} at $t$ and $-t$ expresses
$\mathfrak F(1+t)$ and $\mathfrak F(1-t)$ in terms of those bounded
functions. Their rational coefficients are bounded for $|\Im t|\geq1$.
Hence
\[
 W(t)=O(1+|\Im t|^2)
 \quad\text{uniformly in a full vertical strip of width one}.
\]
An entire one-periodic function with this polynomial vertical growth is
constant. Explicitly, its $n$th Fourier coefficient can be integrated on
any horizontal translate of a period; its modulus is bounded by a
polynomial in $|y|$ times $e^{2\pi n y}$. Sending $y\to-\infty$ for
$n>0$ and $y\to+\infty$ for $n<0$ kills every nonzero mode.
We conclude that $W(t)\equiv-1/\pi^2$, or
\begin{equation}\boxed{
 \mathfrak F(t)\mathfrak F(-t)
 -\mathfrak F(1+t)\mathfrak F(1-t)=\frac1{\pi^2t^2}.
}\label{gff:eq:KproofBilinearF}\end{equation}
This argument supplies both the normalization and the reason that a
nonconstant periodic multiplier is impossible.

Apply \eqref{gff:eq:KproofThomae} at $t$ and $-t$, and use gamma reflection.
The equivalent identity, analytic near zero, is
\begin{equation}\boxed{
 \cos^2(\pi t)\,\mathfrak G(t)\mathfrak G(-t)
 +\pi^2t^2\,\mathfrak F(t)\mathfrak F(-t)=1.
}\label{gff:eq:KproofBilinearFG}\end{equation}
No numerical input has been used in this proof.

\subsection{Only one combination of higher Taylor coefficients is needed}
Write ordinary Taylor coefficients, without hidden factorials, as
\[
 \mathfrak F(t)=1+u t+v t^2+O(t^3),\qquad
 \mathfrak G(t)=1+\sigma t^3+\tau t^4+O(t^5).
\]
The absence of the first two terms in $\mathfrak G-1$ follows immediately
from its three upper parameters $t,t,t$. Extracting $[t^4]$ in
\eqref{gff:eq:KproofBilinearFG} gives
\[
 2\tau+\frac{\pi^4}{3}+\pi^2(2v-u^2)=0,
\]
hence the exact elimination rule
\begin{equation}\boxed{
 \tau+\pi^2v=\frac{\pi^2u^2}{2}-\frac{\pi^4}{6}.
}\label{gff:eq:KproofTaylorCancellation}\end{equation}
Neither $v$ nor $\tau$ has to be evaluated individually. The two lower
coefficients are elementary inverse-binomial sums:
\begin{equation}\begin{aligned}
 u&=\sum_{n\geq1}\frac{\binom{2n}{n}^2}{16^n n}
      =4\ell-\frac{8\Cat}{\pi},\\
 \sigma&=\sum_{n\geq1}\frac{16^n}{n^3\binom{2n}{n}^2}
      =8\pi\Cat-14\zeta(3).
\end{aligned}\label{gff:eq:KproofLowerCoefficients}\end{equation}
For completeness, both evaluations follow without an external summation
assumption. Differentiate \eqref{gff:eq:KproofDoubleEuler} at zero, set
$u=\sin^2\theta$, $v=\sin^2\phi$ in that integral, and use
\[
 \int_0^{\pi/2}\log(1-a\sin^2\phi)\,d\phi
 =\pi\log\frac{1+\sqrt{1-a}}2.
\]
This elementary identity follows by differentiating in $a$ and matching
at $a=0$. It gives
\[
 u=-\frac{16}{\pi}\int_0^{\pi/4}\log\cos x\,dx
   =4\ell-\frac{8\Cat}{\pi}.
\]
For $\sigma$, expand
\[
 (\arcsin x)^2=\frac12\sum_{n\geq1}
           \frac{4^n x^{2n}}{n^2\binom{2n}{n}}
\]
and integrate with $x=\sin\theta$, using
$\int_0^{\pi/2}\sin^{2n-1}\theta\,d\theta
=4^n/(2n\binom{2n}{n})$. The result is
\begin{align*}
 \sigma&=4\int_0^{\pi/2}\frac{\theta^2}{\sin\theta}\,d\theta
 =-8\int_0^{\pi/2}\theta\log\tan(\theta/2)\,d\theta\\
 &=8\pi\sum_{k\geq0}\frac{(-1)^k}{(2k+1)^2}
   -16\sum_{k\geq0}\frac1{(2k+1)^3}
 =8\pi\Cat-14\zeta(3).
\end{align*}
In the last step one may use an Abel factor in the cosine series
$\log\tan(\theta/2)=-2\sum_{k\geq0}\cos((2k+1)\theta)/(2k+1)$
and then take its limit. The endpoint logarithms are integrable, and
integration yields absolutely convergent sums.

\subsection{Extraction of the original master coefficient}
The exact gamma interpolation in the main text gives
\[
 \KK=[t^2]J_t-g_2,\qquad
 g_2=\frac{A^4}{6}-\frac{A^2\pi^2}{6}-\frac{A\pi^3}{3}
      -12A\zeta(3)-\frac{\pi^4}{180}.
\]
Here $[t^2]J_t$ is a Laurent coefficient. No ordinary derivative at a
pole is being taken. To expand \eqref{gff:eq:KproofTwoFunctions}, use
\[
 \frac{\pi\Gamma(1+t)^2}{\Gamma(1/2+t)^2}
 =\exp\left(4\ell t-\frac{\pi^2}{3}t^2+4\zeta(3)t^3
              -\frac{7\pi^4}{90}t^4+O(t^5)\right).
\]
Direct coefficient extraction now yields
\begin{equation}\begin{aligned}
 \KK={}&4A\sigma+4\tau+\pi^3u+4\pi^2v
 -32\ell^2\pi^2-8\ell\pi^3\\
 &+112\ell\zeta(3)+\frac{2\pi^4}{3}+28\pi\zeta(3).
\end{aligned}\label{gff:eq:KproofCoefficientAssembly}\end{equation}
Substituting \eqref{gff:eq:KproofTaylorCancellation} first reduces this to
\[
 \KK=4A\sigma+2\pi^2u^2+\pi^3u-32\ell^2\pi^2-8\ell\pi^3
       +28A\zeta(3).
\]
The elementary values \eqref{gff:eq:KproofLowerCoefficients} then give exactly
\[
 \KK=128\Cat^2+24\pi^2\Cat-4\pi^3\ell-28A\zeta(3).
\]
Finally \eqref{gff:eq:quarterPolygamma} proves \eqref{gff:eq:KprovedFinal} and
\eqref{gff:eq:KLpsiReflected}. Thus the compact evaluation is an analytic consequence of the original master definition. This completes the proof.

\paragraph{Scope of the proof.} The proof evaluates the one-dimensional
master constant defined by \eqref{gff:eq:Kdef}. It proves a local coefficient
identity by means of a global parameter identity; it does not claim that the
whole deformed master $J_t$ is a gamma quotient, and it does not
establish a corresponding identity in other dimensions.

\paragraph{Numerical checks of the identities.} The proof uses no numerical
input. Its identities were also tested numerically, separately from the
proof. These tests are consistency checks, not premises of the proof, and
none of them is an interval bound. Besides the 110-digit endpoint comparison
quoted after \eqref{gff:eq:Kdef}, the record
\path{lri_psi_complete_update/verification/numerical_checks.json} contains
the following tests. A direct evaluation of $\mathfrak F$ and $\mathfrak G$
at $t=0.2$ and 50 working digits satisfies the bilinear identity
\eqref{gff:eq:KproofBilinearFG} to $2.6728\times10^{-51}$. At 75 working digits,
with two settings of the asymptotic tail summation, the six-function master
$\widehat J(t)/t^2$ of \eqref{gff:eq:Jhat} and the two-function reduction
\eqref{gff:eq:KproofTwoFunctions} agree to better than $2\times10^{-68}$ at
$t=0.2$, $0.37$ and $0.12+0.07\ii$, and \eqref{gff:eq:KproofBilinearFG} holds to
better than $4\times10^{-72}$ at the two real points. At the complex point
the 75-digit bilinear residual is only $3\times10^{-59}$, and
$6\times10^{-48}$ with the longer tail, because the gamma-ratio tail cancels
there. Rerun at 130 working digits, the residual has modulus
$1.39\times10^{-102}$
(\path{lri_psi_complete_update/verification/complex_bilinear_130.json}); the
less accurate 75-digit values are kept in the record. At 110 digits the
integral forms of $u$ and $\sigma$ given after
\eqref{gff:eq:KproofLowerCoefficients} reproduce their closed forms to
$8\times10^{-112}$ and $7\times10^{-111}$. Finally, the coefficients
$C_0,\ldots,C_5$ of Section~\ref{gff:cn:oneDresults} and their reflected-quarter
forms agree in exact symbolic algebra (Section~\ref{gff:sec:psiAll}) and
numerically at 70 working digits, with
largest difference $9.3\times10^{-69}$, in $C_5$
(\path{lri_psi_complete_update/verification/coefficient_values.json}). The
numerical comparison reuses the 240-digit values of $F_\pm$
(Section~\ref{gff:sec:oldmasters}) from an earlier record, stored beside it as
\path{wheel_inputs.json}, and recomputes $P_5$ from its gamma moment. None of
these checks bears on the offset of the earlier full-graph control, which
Section~\ref{gff:sec:checks} explains.

\section{\texorpdfstring{The wheel periods $P_4$ and $\mathcal L$}{The wheel periods P4 and L}}\label{gff:sec:oldmasters}

\subsection{The four-loop wheel: two \texorpdfstring{${}_4F_3(1)$}{4F3(1)} values}
Let $G_4=\Gamma(1/4)^4$, and define
\[
 F_-={}_4F_3\!\left(\begin{matrix}1/4,1/4,1/2,1/2\\3/4,3/4,1\end{matrix};1\right),\qquad
 F_+={}_4F_3\!\left(\begin{matrix}3/4,3/4,1/2,1/2\\5/4,5/4,1\end{matrix};1\right).
\]
Then
\begin{equation}\boxed{P_4=13G_4F_--\frac{192\pi^4}{G_4}F_+.}\end{equation}
It evaluates to $2229.5220610709419642945251997485\ldots$.
It follows that the two-${}_4F_3$ expression of the fourth coefficient is
\begin{equation}\boxed{\begin{aligned}
648C_4={}&208G_4F_- -\frac{3072\pi^4}{G_4}F_+
-A^3-A\pi^2-53\pi^3-1436\zeta(3)-8AQ.
\end{aligned}}\label{gff:eq:C4compactUpdated}\end{equation}
For comparison with the earlier elliptic formulas, define
\[
H_4={}_4F_3\!\left(\begin{matrix}1/2,1/2,1/2,1/2\\1,1,1\end{matrix};1\right),\quad
H_7={}_7F_6\!\left(\begin{matrix}5/4,(1/2)^{\times6}\\1/4,1,1,1,1,1\end{matrix};1\right).
\]
The two elimination identities are
\begin{equation}\begin{aligned}
 H_4&=\frac{G_4}{4\pi^3}F_- -\frac{16\pi}{G_4}F_+,\\
 H_7&=\frac{G_4}{2\pi^4}F_- +\frac{32}{G_4}F_+.
\end{aligned}\end{equation}
Indeed, write $\mathbf K(k)=K(k^2)$, $k'=\sqrt{1-k^2}$ and
$N_\nu=\int_0^1 k^\nu\mathbf K(k)\mathbf K(k')\dd k$. Euler integration gives
\begin{equation}
 N_\nu=\frac{\pi^2}{8}\frac{\Gamma(a)^2}{\Gamma(a+1/2)^2}
 {}_4F_3\!\left(\begin{matrix}1/2,1/2,a,a\\1,a+1/2,a+1/2\end{matrix};1\right),
 \qquad a=\frac{\nu+1}{2}.
\end{equation}
The Landen substitution $k=(1-t)/(1+t)$, together with the mixed-to-square moment identity, gives
$N_{-1/2}-N_{1/2}=2N_0$ and
$N_{-1/2}+N_{1/2}=4\int_0^1\mathbf K(k)^2\dd k$.
Using $N_0=\pi^3H_4/8$ and
$\int_0^1\mathbf K(k)^2\dd k=\pi^4H_7/32$ proves the stated identities.
These elliptic-moment evaluations are discussed in \cite{gff:Wan}.



\subsubsection{From the wheel graph to the elliptic period}
For the wheel edges $02,03,04,12,13,14,24,34$, fix $x_0=0,x_1=1$. The two remaining outer-vertex integrations factor, giving
\[
P_4=\int_{\mathbb R}\frac{F(x)^2}{\sqrt{|x(1-x)|}}\,dx,\qquad
F(x)=\int_{\mathbb R}\frac{du}{\sqrt{|u(1-u)(u-x)|}}.
\]
Splitting at $u=0,x,1$ gives $F(x)=4[K(x)+K(1-x)]$ for $0<x<1$. The transformations
\[
F(x)=x^{-1/2}F(1/x)\quad(x>1),\qquad
F(x)=(1-x)^{-1/2}F\!\left(\frac{x}{x-1}\right)\quad(x<0)
\]
fold the tails. Hence
\[
P_4=16\int_0^1\frac{1+\sqrt x+\sqrt{1-x}}{\sqrt{x(1-x)}}
[K(x)+K(1-x)]^2\,dx.
\]
The ordered tetrahedron satisfies
\[
\int_{0<x<y<1}\frac{dx\,dy}{\sqrt{xy(1-x)(1-y)(y-x)}}
=\frac12S_2\!\left(\frac12,\frac12,-\frac14\right)=\frac{\Gamma(1/4)^4}{4\pi};
\]
its full-line period is $3Q$. The three-loop calculation of \cite{gff:BGHS}, with the tetrahedron value corrected in \cite{gff:BGH}, uses this normalization; the integral above is a Selberg integral \cite{gff:Selberg}.

\subsubsection{The moment identities and their contour proof}
Let $M=\int_0^1\mathbf K(k)^2\,dk$, $N=N_0$, and
$J=\int_0^1\mathbf K(k)\mathbf K(k')/k'\,dk$ with $k'=\sqrt{1-k^2}$.
Euler integration of the complementary moment gives
\[
\int_0^1 k^r\mathbf K(k')\,dk
=\frac\pi4\frac{\Gamma((r+1)/2)^2}{\Gamma((r+2)/2)^2},\quad r>-1.
\]
Expanding the other $\mathbf K$ proves the formula for $N_\nu$ above. The square-moment evaluation is $M=\pi^4H_7/32$ \cite{gff:Wan}. To see the mixed-to-square identity $J=2M$, integrate $\mathbf K(z)^2$ around the upper semicircle from $1$ to $-1$. Cauchy's theorem gives $-2M$ on the arc. The boundary-value Landen relation
\[
e^{i\theta}\mathbf K(e^{2i\theta})=\tfrac12[\mathbf K(\cos\theta)+i\mathbf K(\sin\theta)],
\quad 0<\theta<\pi/2,
\]
with the limit taken from inside the disk, makes the real part of that arc integral equal to $-J$; the imaginary part is odd under $\theta\mapsto\pi/2-\theta$. Integrable logarithmic endpoint singularities are handled by an interior limit. Thus $J=2M$.

The real Landen substitution $k=(1-t)/(1+t)$ preserves the mixed measure:
\[
\mathbf K(k)\mathbf K(k')|dk|=\mathbf K(t)\mathbf K(t')dt.
\]
The sum and difference of the half-integer weights transform to $2/t'$ and $2t/t'$, respectively. Therefore
\[
N_{-1/2}+N_{1/2}=2J=4M,\qquad N_{-1/2}-N_{1/2}=2N_0.
\]
Since
\[
N_{-1/2}=\frac{G_4}{16}F_-,\quad
N_{1/2}=\frac{4\pi^4}{G_4}F_+,\quad N_0=\frac{\pi^3}{8}H_4,
\]
these prove both elimination identities. Folding the wheel integral gives
$P_4=16(20M+16N)$, and consequently the stated two-${}_4F_3$ result.

\subsubsection{Four-loop counterterm combination}
The residue certificate, including the wheel $b_4^{\rm wheel}=-6P_4$, gives
\[
3c_4-b_4=6P_4+45AQ+\frac{405A^3-27A\pi^2-159\pi^3}{8}
-8\pi A-\frac{1077}{2}\zeta(3).
\]
Inserting this in the fourth-order RG formula eliminates the intermediate $\pi A$ terms and yields the compact $C_4$. This is an analytic reduction within the graph calculation, not an additional external benchmark.

\begin{equation}
 B_-=G_4F_-,\qquad B_+=\frac{64\pi^4}{G_4}F_+,
 \qquad P_4=13B_--3B_+.
 \label{gff:eq:LnormalizedBasis}
\end{equation}
The derivative-free formula is
\begin{equation}\boxed{\begin{aligned}
 \LL&=-(5\pi+26\log2)B_-+(\pi+6\log2)B_+\\
 &=-\frac{13A-3\pi}{2}B_-+\frac{3A-\pi}{2}B_+.
\end{aligned}}\label{gff:eq:LpolyWheel}\end{equation}

\begin{equation}\boxed{\begin{aligned}
\LL=-\frac23\Bigl\{&[\psi_0(1)-\psi_0(1/4)]P_4\\
&+[\psi_0(3/4)-\psi_0(1/4)]B_-\Bigr\}.
\end{aligned}}\label{gff:eq:Ldigamma}\end{equation}
The two digamma differences are $\pi/2+3\log2$ and $\pi$, respectively.
Numerically, $\LL=-5805.30162954880238014117987700138\ldots$.

\subsection{The logarithmic wheel: full contour proof}
\label{gff:app:LreductionProof}
We give the contour proof of the logarithmic-wheel reduction.
It proves \eqref{gff:eq:LpolyWheel} and \eqref{gff:eq:Ldigamma}, independently of
the evaluation of $\KK$. The symbols $M$, $N$, $U$, $V$, and $R$ below denote
local elliptic moments. They are distinct from the $V_r(t)$ of Section~\ref{gff:sec:Kdef}.

\subsection{Conventions and the folded integral}
In this section put $\ell=\log2$ and use $B_\pm$ from \eqref{gff:eq:LnormalizedBasis}.
The elliptic integral uses the \emph{modulus} convention:
\[
 \mathbf K(k)=\frac\pi2{}_2F_1(1/2,1/2;1;k^2),\qquad
 k'=\sqrt{1-k^2},\qquad f(k)=\mathbf K(k),\quad g(k)=\mathbf K(k').
\]
The prime in $k'$ denotes the complementary modulus, not differentiation.
In the parameter convention, $K(x)=\mathbf K(\sqrt{x})$; the full-line definition below uses it.
Define
\begin{align*}
 M&=\int_0^1 f(k)^2\dd k,&N&=\int_0^1 f(k)g(k)\dd k,\\
 U_0&=\int_0^1\frac{f(k)^2}{k'}\dd k,&
 V_0&=\int_0^1\frac{f(k)g(k)}{k'}\dd k,\\
 U&=\int_0^1\frac{f(k)^2\log(kk')}{k'}\dd k,&
 V&=\int_0^1\frac{f(k)g(k)\log(kk')}{k'}\dd k,\\
 w(k)&=\log\frac{k}{1-k^2},&
 R&=\int_0^1f(k)^2w(k)\dd k.
\end{align*}
All these integrals converge.

The wheel master is defined on the full line by
\[
 \LL=\int_{\mathbb R}\frac{F(x)^2\log|x(1-x)|}{\sqrt{|x(1-x)|}}\dd x,
 \qquad F(x)=\int_{\mathbb R}\frac{\dd u}{\sqrt{|u(1-u)(u-x)|}}.
\]
For $0<x<1$, $F(x)=4[K(x)+K(1-x)]$. Folding the two tails gives
\begin{align*}
\LL=16\int_0^1[K(x)+K(1-x)]^2\bigg[
&\frac{\log x+\log(1-x)}{\sqrt{x(1-x)}}\\
&+\frac{\log(1-x)-2\log x}{\sqrt{1-x}}
+\frac{\log x-2\log(1-x)}{\sqrt x}\bigg]\dd x.
\end{align*}
Substitution $x=k^2$ and reflection of the second tail term give the useful
identity
\begin{equation}
 \frac{\LL}{128}=U+V+\int_0^1(f+g)^2w\dd k.
 \label{gff:lp:eq:fold}
\end{equation}

\subsection{Real Landen identities}
Let $T(k)=(1-k)/(1+k)$. Standard Landen transformations \cite{gff:DLMFElliptic,gff:DLMFEllipticConnection} give
\[
 f(T(k))=\frac{1+k}{2}g(k),\qquad g(T(k))=(1+k)f(k).
\]
Consequently,
\begin{align}
 f(T)g(T)|T'|&=fg,&g(T)^2|T'|&=2f^2,&f(T)^2|T'|&=\tfrac12g^2.
 \label{gff:lp:eq:measures}
\end{align}
In particular, $\int_0^1g^2\dd k=2M$.
Since
\[
 w(T(k))=-w(k)-2\ell,
\]
we have
\begin{align}
 \int_0^1fgw\dd k&=-\ell N,&
 \int_0^1g^2w\dd k&=-2R-4\ell M.
 \label{gff:lp:eq:weighted}
\end{align}
Combining these with \eqref{gff:lp:eq:fold} yields
\begin{equation}
 \frac{\LL}{128}=U+V-R-4\ell M-2\ell N.
 \label{gff:lp:eq:reduction}
\end{equation}
It remains to evaluate just three logarithmic combinations $U,V,R$.

\subsection{Evaluation of \texorpdfstring{$U$}{U} by the Stieltjes transform}
Use the standard identity (Wan, Theorem 1, equation (15)) \cite{gff:Wan}:
\begin{equation}
 f(t)^2=\frac4\pi\int_0^1\frac{k f(k)g(k)}{1-t^2k^2}\dd k,
 \qquad 0\le t<1.
 \label{gff:lp:eq:stieltjes}
\end{equation}
Integration with respect to $\dd t$ gives
\[
 M=\frac4\pi\int_0^1fg\,\operatorname{artanh}k\dd k.
\]
Applying $T$ and $\operatorname{artanh}T(k)=-\tfrac12\log k$ therefore gives
\begin{equation}
 \int_0^1fg\log k\dd k=-\frac\pi2 M.
 \label{gff:lp:eq:mixedlog}
\end{equation}
Also $\log(1+T(k))=\ell-\log(1+k)$, so
\begin{equation}
 \int_0^1fg\log(1+k)\dd k=\frac\ell2N.
 \label{gff:lp:eq:pluslog}
\end{equation}
Integrating \eqref{gff:lp:eq:stieltjes} against $\dd t/\sqrt{1-t^2}$ proves
$U_0=2N$, using the elementary integral
$\int_0^{\pi/2}(1-k^2\sin^2\theta)^{-1}\dd\theta=\pi/(2k')$
and then complementing $k$.

For the logarithmic version, the tangent substitution gives
\[
\int_0^{\pi/2}\frac{\log(\sin\theta\cos\theta)}{1-k^2\sin^2\theta}\dd\theta
=-\frac\pi{2k'}\log\frac{(1+k')^2}{k'}.
\]
Insert this into \eqref{gff:lp:eq:stieltjes} and complement the outer variable:
\begin{align*}
 U&=-2\int_0^1\frac{kfg}{k'}\log\frac{(1+k')^2}{k'}\dd k\\
 &=-2\int_0^1fg\,[2\log(1+k)-\log k]\dd k.
\end{align*}
Equations \eqref{gff:lp:eq:mixedlog}--\eqref{gff:lp:eq:pluslog} now prove
\begin{equation}\boxed{U=-\pi M-2\ell N.}\label{gff:lp:eq:U}\end{equation}
The integral interchanges are justified by positivity for the nonlogarithmic
identities and by Tonelli applied to the negative logarithm for the last one;
the resulting moments are finite.

\subsection{Evaluation of \texorpdfstring{$R$}{R} by a first-quadrant contour}
Continue $\mathbf K(z)$ from its series at zero and put
\[
 W(z)=\operatorname{Log}z-\operatorname{Log}(1-z^2),
\]
with the logarithms analytic in the first quadrant and their indicated upper
boundary values on the real axis. Integrate $\mathbf K(z)^2W(z)$ around the first
quadrant. The large arc vanishes because the integrand is
$O(\log^3|z|/|z|^2)$. The small arcs at $0$ and $1$ also vanish; their
singularities are only integrable powers of logarithms. Thus
\begin{equation}
 \int_0^\infty\mathbf K(x+i0)^2W(x+i0)\dd x
 =i\int_0^\infty\mathbf K(iy)^2W(iy)\dd y.
 \label{gff:lp:eq:quadrant}
\end{equation}
The reciprocal and imaginary-modulus connection formulas \cite{gff:DLMFElliptic,gff:DLMFEllipticConnection} are
\[
 \mathbf K(x+i0)=\frac1x[f(1/x)+ig(1/x)]\quad(x>1),
 \qquad
 \mathbf K(iy)=\frac{1}{\sqrt{1+y^2}}f\!\left(\frac{y}{\sqrt{1+y^2}}\right).
\]
Furthermore $W(x+i0)=w(1/x)+i\pi$ for $x>1$ and
$\operatorname{Im}W(iy)=\pi/2$.
The real part of the left-hand side of \eqref{gff:lp:eq:quadrant} is therefore
\[
 2R-\int_0^1g^2w\dd k-2\pi N=4R+4\ell M-2\pi N.
\]
On the right-hand side, set $y=k/k'$; the real part is
\[
 -\frac\pi2\int_0^\infty\mathbf K(iy)^2\dd y
 =-\frac\pi2U_0=-\pi N.
\]
Equating these gives
\begin{equation}\boxed{R=\frac\pi4N-\ell M.}\label{gff:lp:eq:R}\end{equation}
In particular, the extra $\pi N$ term is a branch contribution; discarding the
imaginary part of the reciprocal-modulus continuation would lose it.

\subsection{Evaluation of \texorpdfstring{$V$}{V} by an upper-semicircle contour}
Integrate the same $\mathbf K(z)^2W(z)$ over the upper half of the unit disk. The
integral over the diameter from $-1$ to $1$ is
\[
 2R+i\pi M.
\]
On the semicircle, put $z=e^{2i\theta}$, $0<\theta<\pi/2$. The boundary-value
Landen identity, with continuation from the interior of the disk, reads
\[
 e^{i\theta}\mathbf K(e^{2i\theta})
 =\frac12[\mathbf K(\cos\theta)+i\mathbf K(\sin\theta)],
\]
This follows by applying
$\mathbf K(z)=(1+z)^{-1}\mathbf K(2\sqrt z/(1+z))$: as $z$ approaches
$e^{2i\theta}$ from inside the disk, its transformed modulus approaches
$\sec\theta$ from the upper side. The reciprocal-modulus formula then
fixes the displayed plus sign. Also,
\[
 W(e^{2i\theta})=-\log(4\sin\theta\cos\theta)+\frac{i\pi}{2}.
\]
Write $f_\theta=\mathbf K(\sin\theta)$, $g_\theta=\mathbf K(\cos\theta)$. Then
\[
 \mathbf K(z)^2\dd z=
 \left[-f_\theta g_\theta+\frac{i}{2}(g_\theta^2-f_\theta^2)\right]\dd\theta.
\]
Terms containing $g_\theta^2-f_\theta^2$ times a symmetric factor integrate to
zero. The complete arc integral is therefore
\[
 V+2\ell V_0-\frac{i\pi}{2}V_0.
\]
Cauchy's theorem implies, from its imaginary and real parts respectively,
\[
 V_0=2M,\qquad V=-2R-4\ell M.
\]
Together with \eqref{gff:lp:eq:R}, this proves
\begin{equation}\boxed{V=-\frac\pi2N-2\ell M.}\label{gff:lp:eq:V}\end{equation}
The endpoint limits are valid because the logarithmic singularities are
integrable and the small indentations contribute zero.

\subsection{Cancellation and reduction to the two \texorpdfstring{${}_4F_3$}{4F3} values}
Substituting \eqref{gff:lp:eq:U}, \eqref{gff:lp:eq:R}, and \eqref{gff:lp:eq:V} into
\eqref{gff:lp:eq:reduction} gives
\begin{equation}\boxed{
 \LL=-(128\pi+640\ell)M-(96\pi+512\ell)N.
}\label{gff:lp:eq:MN}\end{equation}
To eliminate $M,N$, introduce
\[
 N_\nu=\int_0^1 k^\nu f(k)g(k)\dd k,\qquad \nu>-1.
\]
Expanding $f$ and using the beta moment of $g$ yields
\[
 N_\nu=\frac{\pi^2}{8}
 \frac{\Gamma((\nu+1)/2)^2}{\Gamma((\nu+2)/2)^2}
 {}_4F_3\!\left(\begin{matrix}1/2,1/2,(\nu+1)/2,(\nu+1)/2\\
 1,(\nu+2)/2,(\nu+2)/2\end{matrix};1\right).
\]
Hence $N_{-1/2}=B_-/16$ and $N_{1/2}=B_+/16$.
Applying $T$ to their sum and difference gives
\[
 N_{-1/2}+N_{1/2}=2V_0=4M,
 \qquad N_{-1/2}-N_{1/2}=2N.
\]
Thus
\[
 M=\frac{B_-+B_+}{64},\qquad N=\frac{B_--B_+}{32}.
\]
Substitution into \eqref{gff:lp:eq:MN} proves \eqref{gff:eq:LpolyWheel}.
The nonlogarithmic folded integral gives $P_4=320M+256N=13B_--3B_+$,
which proves \eqref{gff:eq:Ldigamma}. No parameter-differentiation identity is required.

\section{The five-loop primitive period}\label{gff:sec:P5complete}

\subsection{Primitive graph and normalization}
A zig-zag decompletion has edges
\[
02,03,04,05,13,14,15,24,25,35.
\]
Fixing $x_0=0,x_5=1$ and performing two independent star integrations gives
\begin{equation}
P_5=\int_{\mathbb R^2}\frac{F(u)F(v)}{\sqrt{|v(1-u)(u-v)|}}\,du\,dv,
\quad F(x)=\int_{\mathbb R}\frac{dt}{\sqrt{|t(1-t)(t-x)|}}.
\label{gff:eq:P5graph}
\end{equation}
This integral is convergent and has no subdivergence subtractions. The zig-zag class has automorphism order $2$ and contributes $24P_5$ to $b_5$. The other two primitive classes share the product period $(3Q)^2=9Q^2$ and have automorphism orders $16$ and $12$. Their combined contribution is
\[
\left(\frac{48}{16}+\frac{48}{12}\right)9Q^2=63Q^2.
\]
All primitive mass weights vanish because the four external legs are on distinct vertices. This yields $C_5^{\rm primitive}=-(63Q^2+24P_5)/729$ in the present normalization.

\subsection{Spectral representation and continuation to one dimension}
The starting point is the conformal-triangle spectral representation of zig-zag graphs in \cite{gff:DIS}, in particular its eigenvalues, spectral measure and two-point reduction, equations (5.23), (5.27) and (6.2). In the present loop convention its eigenvalue power is $k=L+2$, where $L$ is the contribution's loop order in the quartic beta function.

Continuing the spectral formula to $D=1$ requires more than substituting $D=1$ into the density on the real spectral axis. A pole crossing contributes a discrete sum. Define
\begin{align}
 \tau_n(\nu)&=(-1)^n\sqrt\pi
 \left|\frac{\Gamma(n/2+1/8+i\nu)}
 {\Gamma(n/2+3/8+i\nu)}\right|^2,
 &&n=0,1,\label{gff:sp:eq:tau}\\
 \sigma_j&=(-1)^j\sqrt\pi\,
 \frac{\Gamma(j/2-1/8)\Gamma(j/2+3/8)}
 {\Gamma(j/2+1/8)\Gamma(j/2+5/8)},
 &&j=1,2,\ldots.\label{gff:sp:eq:sigma}
\end{align}
For the zig-zag family, the one-dimensional period is
\begin{equation}\label{gff:sp:eq:spec}
 \boxed{\displaystyle
 P_L=\frac1{\pi^2}\int_0^\infty
 \nu\tanh(2\pi\nu)
 \left[\tau_0(\nu)^k+\tau_1(\nu)^k\right]\dd\nu
 +\frac1{2\pi^2}\sum_{j=1}^\infty
 (j-\tfrac12)\sigma_j^k,
 \qquad k=L+2.
 }
\end{equation}
The first term alone is not the period.

\subsubsection{Origin of the discrete sum}
Here is the contour-continuation detail needed to fix its sign and normalization. Set $p=D/2-1$ and specialize the propagator parameter in \cite{gff:DIS} to $\beta=D/4$. The combined density is
\begin{align}\label{gff:sp:eq:rho}
 \rho_n^{(D)}(\nu)
 ={}&\frac{\Gamma(p)\Gamma(n+2p)(n+p)}
 {2\pi^{3D/2+1}\Gamma(2p)n!}
 \bigl[(n+p)^2+4\nu^2\bigr]\\[-2pt]
 &\qquad\times
 \frac{\Gamma(p+2i\nu)\Gamma(p-2i\nu)}
 {\Gamma(2i\nu)\Gamma(-2i\nu)}.\notag
\end{align}
The corresponding eigenvalue is
\[
 \tau_n^{(D)}(\nu)=(-1)^n\pi^{D/2}
 \frac{\Gamma(n/2+D/8+i\nu)\Gamma(n/2+D/8-i\nu)}
 {\Gamma(n/2+3D/8+i\nu)\Gamma(n/2+3D/8-i\nu)}.
\]
For $D>2$ the spectral formula uses the real $\nu$ contour. As $D$ passes below $2$, the poles $\nu=\pm ip/2$ exchange sides. Retaining the analytically continued contour gives a correction
\[
 -2\pi i\Res_{\nu=-ip/2}
 \rho_n^{(D)}(\nu)\tau_n^{(D)}(\nu)^k.
\]
Before taking $D\to1$, its coefficient is
\begin{equation}\label{gff:sp:eq:crossing}
 \frac{\Gamma(n+2p)(n+p)n(n+2p)}
 {2\pi^{3D/2}n!\Gamma(-p)}\,
 \tau_n^{(D)}(-ip/2)^k.
\end{equation}
The $n=0$ crossing cancels identically. For $n\geq1$, the limit of the prefactor in \eqref{gff:sp:eq:crossing} is $(n-1/2)/(2\pi^2)$, while the eigenvalue becomes \eqref{gff:sp:eq:sigma}. This gives the second term of \eqref{gff:sp:eq:spec}.

On the real axis, only $n=0,1$ remain in the pointwise $D\to1$ limit, each with density $\nu\tanh(2\pi\nu)/\pi^2$. The fact that higher-spin densities vanish on this axis does not remove their crossed-pole residues. Losing these residues would already spoil the known three-loop tetrahedron normalization.

The continuation and residue reduction here specialize the spectral construction of \cite{gff:DIS}; they are not an independently published evaluation of this Ising coefficient.

\subsection{Closing the contour and summing the gamma poles}
Define
\begin{equation}\label{gff:sp:eq:Ralpha}
 R_\alpha(x)=
 \frac{\Gamma(\alpha+x)\Gamma(\alpha-x)}
 {\Gamma(\alpha+1/4+x)\Gamma(\alpha+1/4-x)}.
\end{equation}
With $x=i\nu$, closing the continuous spectral integral to the right expresses it as
\begin{equation}\label{gff:sp:eq:contres}
 P_L^{\rm cont}=\pi^{k/2-1}
 \sum_{\Re x>0}\Res_x
 \left\{x\tan(2\pi x)
 \bigl[R_{1/8}(x)^k+(-1)^kR_{5/8}(x)^k\bigr]\right\}.
\end{equation}
One can use expanding rectangular contours avoiding the real poles. Gamma reflection and Stirling estimates give vanishing outer-contour contributions for $k>4$.

There are two pole families. The tangent poles are
$x_j=(2j-1)/4$, $j\geq1$, where
\[
 R_{5/8}(x_j)=-R_{1/8}(x_j),\qquad
 \sqrt\pi R_{1/8}(x_j)=\sigma_j.
\]
Their residues in \eqref{gff:sp:eq:contres} give exactly
\[
 -\frac1{2\pi^2}\sum_{j=1}^\infty(j-\tfrac12)\sigma_j^k.
\]
They cancel the discrete term of \eqref{gff:sp:eq:spec}. The complete period is therefore given just by the gamma poles at $x=\alpha+n$, $n\geq0$, for $\alpha=1/8,5/8$.

\subsection{A Barnes integral with no unwanted tangent poles}
Set
\[
s=2x-\frac14,\qquad
g(s)=\frac{\Gamma(s+\frac12)}{\Gamma(s+1)},\qquad
C(s)=\csc(\pi s),\qquad D(s)=\cot(\pi s).
\]
Reflection and duplication imply
\begin{align}
R_{1/8}(x)&=g(s)\big[1-\cot(\pi s/2)\big],\\
R_{5/8}(x)&=g(s)\big[1+\tan(\pi s/2)\big].
\end{align}
The two gamma-pole families have become $s=0,1,2,\ldots$. Define
\begin{equation}\begin{split}
S(s)=\tan(\pi s+\pi/4)\big\{
[1-\cot(\pi s/2)]^7-[1+\tan(\pi s/2)]^7\big\}.
\end{split}\end{equation}
Its unwanted tangent poles are at $s=1/4+j$. Their residues equal those of
$16\sqrt2\,\csc[\pi(s-1/4)]$. Consequently the kernel
\[
\mathcal R(s)=S(s)-16\sqrt2\,\csc[\pi(s-1/4)]
\]
has only integer poles on the right of $\Re s=-1/4$ and retains the desired residues. A rational trigonometric simplification gives
\begin{equation}\boxed{\begin{aligned}
\mathcal R(s)={}&-128C+160C^3+192C^5-128C^7\\
&-16DC-400DC^3+192DC^5.
\end{aligned}}\label{gff:fp:eq:kernel}\end{equation}
The SymPy check \path{lri_fixed_hypergeometric_update/public/kernel_identity.py} verifies this exactly after replacing
$\tan(\pi s/2)$ by an indeterminate.

Closing a vertical contour to the right is clockwise. Also,
$x\,\dd x=(s+1/4)\dd s/4$. Therefore
\begin{equation}
P_5=-\frac{\pi^{5/2}}4\frac1{2\pi\ii}
\int_{-1/4-\ii\infty}^{-1/4+\ii\infty}
(s+\tfrac14)g(s)^7\mathcal R(s)\dd s.
\label{gff:fp:eq:barnes}\end{equation}
No poles lie on this line. On rectangles avoiding the real poles, gamma-ratio asymptotics and the trigonometric factors justify the vanishing of the closing contributions. The remaining integer-pole sum converges absolutely.

\subsection{The reflection that removes four derivative orders}
For odd $j$, write
\begin{align}
I_j&=\frac1{2\pi\ii}\int_{-1/4-\ii\infty}^{-1/4+\ii\infty}
(s+\tfrac14)g(s)^7C(s)^j\dd s,\\
J_j&=\frac1{2\pi\ii}\int_{-1/4-\ii\infty}^{-1/4+\ii\infty}
(s+\tfrac14)g(s)^7D(s)C(s)^j\dd s.
\end{align}
The substitution $s\mapsto-1/2-s$ preserves the vertical line. Gamma reflection and elementary trigonometry give
\begin{equation}\boxed{
g(-\tfrac12-s)=-D(s)g(s),\qquad
C(-\tfrac12-s)=-\frac{C(s)}{D(s)},\qquad
D(-\tfrac12-s)=\frac1{D(s)}.
}\end{equation}
Together with $D^2=C^2-1$, these imply
\begin{equation}\boxed{
I_7=0,\qquad I_5=I_3,\qquad J_5=0,\qquad J_3=J_1.
}\label{gff:fp:eq:momentrelations}\end{equation}
For example, $I_7=-I_7$. For $I_3$ the reflected integrand gives
$I_3=-I_7+2I_5-I_3$, whence $I_3=I_5$. Similarly,
$J_5=-J_5$ and $J_1=-J_5+2J_3-J_1$.

Substitution into \eqref{gff:fp:eq:kernel}--\eqref{gff:fp:eq:barnes} yields
\begin{equation}\boxed{
P_5=\frac{\pi^{5/2}}4\big(128I_1-352I_3+416J_1\big).
}\label{gff:fp:eq:three}\end{equation}
The pole order is now at most three rather than seven. This is the precise mechanism behind the reduction from sixth to second parameter derivatives.

\subsubsection{Summing the three remaining pole families}
Define the analytic function
\begin{equation}
\mathcal A_{\rm s}(t)=\sum_{n=0}^\infty(-1)^n(n+\tfrac14+t)
\left[\frac{\Gamma(n+\frac12+t)}{\Gamma(n+1+t)}\right]^7.
\label{gff:fp:eq:A}\end{equation}
Termwise parameter derivatives through the needed order are absolutely convergent near zero. Using $\csc[\pi(n+t)]=(-1)^n\csc(\pi t)$ and
$\cot[\pi(n+t)]=\cot(\pi t)$, closing to the right gives
\begin{equation}\boxed{\begin{aligned}
I_1&=-\frac{\mathcal A_{\rm s}(0)}\pi,\\
I_3&=-\frac{\mathcal A_{\rm s}''(0)+\pi^2\mathcal A_{\rm s}(0)}{2\pi^3},\\
J_1&=-\frac{\mathcal A_{\rm s}'(0)}{\pi^2}.
\end{aligned}}\end{equation}
Here the minus signs are the clockwise-contour signs. For the third formula, $\csc(\pi t)\cot(\pi t)=(\pi t)^{-2}-1/6+O(t^2)$ has no simple-pole term. Substitution into \eqref{gff:fp:eq:three} gives

\begin{equation}
 P_5=\frac4{\sqrt\pi}
 \left[11\mathcal A_{\rm s}''(0)-26\pi\mathcal A_{\rm s}'(0)
                  +3\pi^2\mathcal A_{\rm s}(0)\right].
 \label{gff:cn:P5jet}
\end{equation}
This intermediate identity follows by substituting the three residues into
\eqref{gff:fp:eq:three}. The finite-periodization proof below eliminates all
three Taylor operations from the final evaluation.

\subsection{Derivation of the positive gamma integral}
On the same vertical line, put $s=-1/4+\ii y$ and
$G(y)=|\Gamma(1/4+\ii y)|^2$. Gamma reflection gives
\begin{equation}
g(s)=\frac{G(y)}{\sqrt2\,\pi}
\big[\cosh(\pi y)-\ii\sinh(\pi y)\big].
\end{equation}
Writing $z=\cosh(\pi y)-\ii\sinh(\pi y)$, one has
$C(s)=-\sqrt2/z$ and $D(s)=-\overline z/z$.
Taking conjugate halves of the contour integral yields
\begin{align}
I_1&=\frac1{8\pi^8}\int_0^\infty yG(y)^7
\big[\sinh^3(2\pi y)-3\sinh(2\pi y)\big]\dd y,\\
I_3&=-\frac1{2\pi^8}\int_0^\infty yG(y)^7\sinh(2\pi y)\dd y,\\
J_1&=\frac1{4\pi^8}\int_0^\infty yG(y)^7
\sinh(2\pi y)\cosh(2\pi y)\dd y.
\end{align}
Inserting these into \eqref{gff:fp:eq:three} and using
$4\sinh^3 v=\sinh(3v)-3\sinh v$ and
$2\sinh v\cosh v=\sinh(2v)$ proves \eqref{gff:fp:eq:gamma}.
This derivation fixes the coefficients $1,13,29$ and the factor $\pi^{-11/2}$ analytically; they were not obtained by fitting a numerical period.

\begin{equation}\boxed{
 P_5=\frac1{\pi^{11/2}}\int_0^\infty
 y|\Gamma(\tfrac14+\ii y)|^{14}
 [\sinh(6\pi y)+13\sinh(4\pi y)+29\sinh(2\pi y)]\,dy.
}\label{gff:fp:eq:gamma}\end{equation}
The integrand is $O(y^2)$ at zero and $O(y^{-5/2}e^{-\pi y})$ at infinity
by the Gamma-ratio asymptotics \cite{gff:DLMFStirling}; it is positive on the
integration interval. This provides an independent convergent numerical
representation of the same period.
\subsection{Finite periodization and the five fixed values}\label{gff:sec:P5FixedProof}

\subsubsection{The unilateral and bilateral sums}
Use the meromorphic function
\begin{equation}
f(s)=\left(s+\frac14\right)
\left[\frac{\Gamma(s+\frac12)}{\Gamma(s+1)}\right]^7,
\qquad
\mathcal A_{\rm s}(t)=\sum_{n=0}^{\infty}(-1)^n f(n+t).
\label{gff:fx:eq:fA}\end{equation}
This is the same unilateral sum as \eqref{gff:fp:eq:A}. By taking out
the $n=0$ factors and using
$(t+\frac14+n)/(t+\frac14)=(t+\frac54)_n/(t+\frac14)_n$,
\begin{equation}\begin{aligned}
\mathcal A_{\rm s}(t)={}&\left(t+\frac14\right)
 \left[\frac{\Gamma(t+\frac12)}{\Gamma(t+1)}\right]^7\\
&\times{}_9F_8\!\left(
\begin{matrix}1,t+\frac54,(t+\frac12)^{\times7}\\
t+\frac14,(t+1)^{\times7}\end{matrix};-1\right).
\end{aligned}\label{gff:fx:eq:Ahyper}\end{equation}
Consequently
\begin{equation}
\mathcal Y_r=\mathcal A_{\rm s}\!\left(\frac{r-2}{8}\right),
\qquad r=1,2,3,5,7.
\label{gff:fx:eq:Ysamples}\end{equation}
The parameter sum difference in \eqref{gff:fx:eq:Ahyper} is $3/2$;
gamma-ratio asymptotics give absolute and locally uniform convergence
away from poles. In particular all five evaluation points in
\eqref{gff:fx:eq:Ysamples} are nonsingular and have positive
hypergeometric parameters \cite{gff:DLMF16,gff:DLMFStirling}.

Now extend the index in \eqref{gff:fx:eq:fA} to all integers:
\begin{equation}
\Sigma(t)=\sum_{n\in\mathbb Z}(-1)^n f(n+t),
\qquad
\Xi(t)=\cos^7(\pi t)\,\Sigma(t).
\label{gff:fx:eq:bilateral}\end{equation}
These expressions are initially defined away from the half integers.
Gamma reflection gives the two useful identities
\begin{equation}
\frac{\Gamma(-s)}{\Gamma(\frac12-s)}
=-\cot(\pi s)\frac{\Gamma(s+\frac12)}{\Gamma(s+1)},
\qquad
f(-\tfrac12-s)=\cot^7(\pi s)f(s).
\label{gff:fx:eq:reflection}\end{equation}
The minus sign from the linear factor in $f$ cancels the seventh
power of the minus sign in the first identity. These are direct
consequences of the standard gamma reflection formula \cite{gff:DLMF5}.
Writing a negative summation index as $n=-m-1$ and using
\eqref{gff:fx:eq:reflection} shows, including the alternating sign, that
\begin{equation}\boxed{\begin{aligned}
\Sigma(t)&=\mathcal A_{\rm s}(t)
 -\tan^7(\pi t)\mathcal A_{\rm s}(\tfrac12-t),\\
\Xi(t)&=\cos^7(\pi t)\mathcal A_{\rm s}(t)
 -\sin^7(\pi t)\mathcal A_{\rm s}(\tfrac12-t).
\end{aligned}}\label{gff:fx:eq:twohalves}\end{equation}
Thus evaluating $\Xi$ requires only two ordinary unilateral
hypergeometric values, with no new special-function convention.

\subsubsection{Why only three harmonics survive}
We give the analytic growth argument explicitly, because periodicity
alone would not imply a finite Fourier polynomial.
For a compact set of $t$ avoiding half integers, Stirling's
ratio estimate and \eqref{gff:fx:eq:reflection} imply
\[
 f(n+t)=O(|n|^{-5/2})\qquad (n\to\pm\infty),
\]
locally uniformly. Indeed the positive-index estimate follows from
$\Gamma(s+1/2)/\Gamma(s+1)\sim s^{-1/2}$; reflection reduces the
negative-index estimate to the same right-half-plane bound. Hence
$\Sigma$ is meromorphic. Its only possible poles are at half integers
and have order at most seven. Multiplication by $\cos^7(\pi t)$
removes them. The same estimates after multiplication are locally
uniform at a half integer, so the removal applies to the entire
series, not just to individual summands. Therefore $\Xi$ is entire.

Reindexing the absolutely convergent bilateral sum gives
\[
\Sigma(t+1)=-\Sigma(t),\qquad \Xi(t+1)=\Xi(t).
\]
For $0\leq x\leq1$ and $|y|\geq1$, the reflection factors are uniformly
bounded. The right-half-plane gamma estimate therefore also gives
\begin{equation}
|\Sigma(x+\ii y)|\leq
 C\sum_{n\in\mathbb Z}(1+|n+\ii y|)^{-5/2}
=O(|y|^{-3/2}).
\label{gff:fx:eq:Sigmagrowth}\end{equation}
It follows that
\begin{equation}
|\Xi(x+\ii y)|=O\!\left(e^{7\pi|y|}|y|^{-3/2}\right).
\label{gff:fx:eq:Xigrowth}\end{equation}
Let $c_k=\int_0^1\Xi(x)e^{-2\pi\ii kx}\dd x$.
Entirety and periodicity allow the contour to be shifted vertically.
For $k\geq4$, shift down to $\operatorname{Im}t=-Y$ and use
\eqref{gff:fx:eq:Xigrowth} to obtain
\[
|c_k|\leq C Y^{-3/2}e^{(7-2k)\pi Y}\longrightarrow0.
\]
Shifting upward gives $c_k=0$ for $k\leq-4$. Thus only the seven
Fourier modes $-3,\ldots,3$ are possible.

There is one more symmetry. Reflection and reindexing imply
\[
\Sigma(-\tfrac12-t)=\cot^7(\pi t)\Sigma(t),
\qquad
\Xi(\tfrac12-t)=-\Xi(t).
\]
In the second equation periodicity has been used after multiplying by
$\cos^7(\pi t)$. If $\Xi(t)=\sum_{k=-3}^3c_ke^{2\pi\ii kt}$,
this gives $c_0=0$ and $c_{-k}=-(-1)^kc_k$. Consequently
\begin{equation}\boxed{
\Xi(t)=\alpha\cos(2\pi t)+\beta\sin(4\pi t)
       +\gamma\cos(6\pi t).
}\label{gff:fx:eq:threeharmonics}\end{equation}
The coefficients are real because \eqref{gff:fx:eq:twohalves} is real at
nonsingular real $t$. This establishes the three-dimensional space
analytically; the coefficients have not been fitted by assuming an
unproved trigonometric ansatz.

\subsubsection{Exact finite sampling and evaluation of the period}
Set
\begin{equation}\begin{aligned}
X_0&=\Xi(0)=\mathcal Y_2,\\
X_+&=\Xi(\tfrac18)=c_8^7\mathcal Y_3-s_8^7\mathcal Y_5,\\
X_-&=\Xi(-\tfrac18)=c_8^7\mathcal Y_1+s_8^7\mathcal Y_7.
\end{aligned}\label{gff:fx:eq:Xsamples}\end{equation}
Equation \eqref{gff:fx:eq:threeharmonics} gives exactly
\begin{equation}\begin{aligned}
\alpha&=\frac{X_0}{2}+\frac{X_++X_-}{2\sqrt2},\\
\beta&=\frac{X_+-X_-}{2},\\
\gamma&=\frac{X_0}{2}-\frac{X_++X_-}{2\sqrt2}.
\end{aligned}\label{gff:fx:eq:abc}\end{equation}
These are fixed, finite evaluation points; no spacing tends to zero
and there is no interpolation remainder.

For completeness, the link to the previously derived Barnes
integrals can be made without a numerical differentiation step.
The first equation in \eqref{gff:fx:eq:twohalves} gives
$\Sigma(t)=\mathcal A_{\rm s}(t)+O(t^7)$ at zero, since
$\mathcal A_{\rm s}(1/2-t)$ is analytic there. Hence
\begin{equation}\begin{aligned}
\mathcal A_{\rm s}(0)&=\alpha+\gamma,\\
\mathcal A_{\rm s}'(0)&=4\pi\beta,\\
\mathcal A_{\rm s}''(0)&=\pi^2(3\alpha-29\gamma).
\end{aligned}\label{gff:fx:eq:matchingjets}\end{equation}

Substitution of \eqref{gff:fx:eq:matchingjets} into \eqref{gff:cn:P5jet} gives
\[
 P_5=16\pi^{3/2}(9\alpha-26\beta-79\gamma).
\]
Using \eqref{gff:fx:eq:abc} therefore yields
\begin{equation}
 P_5=16\pi^{3/2}
 [-35X_0+(22\sqrt2-13)X_+ +(22\sqrt2+13)X_-],
 \label{gff:fx:eq:Px}
\end{equation}
which proves \eqref{gff:eq:P5FixedAnswer}. Differentiation is used only in
the derivation of \eqref{gff:cn:P5jet}, not in the final fixed-value formula.

\subsubsection{All five parameter arrays}
For clarity, the repeating parameters and the shifted pair in the
five functions are all displayed below. Each row also has the
numerator parameter $1$, as in \eqref{gff:eq:P5FixedY}.
\begin{center}
\begin{tabular}{ccccc}
\toprule
$r$ & shifted upper & seven equal uppers & shifted lower & seven equal lowers\\
\midrule
1 & $9/8$ & $3/8$ & $1/8$ & $7/8$\\
2 & $5/4$ & $1/2$ & $1/4$ & $1$\\
3 & $11/8$ & $5/8$ & $3/8$ & $9/8$\\
5 & $13/8$ & $7/8$ & $5/8$ & $11/8$\\
7 & $15/8$ & $9/8$ & $7/8$ & $13/8$\\
\bottomrule
\end{tabular}
\end{center}
The positive denominator parameters rule out singular defining
series. The only cancellation used is the valid, fixed-parameter
one at $r=2$. The eighths arise from the choice of the exact sampling
points $t=\pm1/8$, not from a change of physical dimension or from a
new regulator. There is no claim that the five individual values
are independent or that their appearance is minimal.

The eighths in the parameters arise from the fixed sample nodes
$t=0,\pm1/8$. They are not new regulators or a change in physical dimension.
At 125 working digits the independently evaluated Gamma integral and the
five-value expression differed by $3.63\times10^{-121}$; this is a numerical
consistency check rather than an interval enclosure. The period begins
\begin{equation*}\begin{split}
 P_5={}&28599.293176562103536646477021422057490321246031264\\
 &\hspace{11mm}16921356927935177776442703824905147018458419833882\ldots.
\end{split}\end{equation*}

\section{Restoring finite regulator corrections}\label{gff:sec:shifts}
The preceding expansions are inserted in the forest calculation, not merely in bare uniqueness integrals. The general forest framework is standard \cite{gff:Herzog}; the concrete deformations and order-by-order corrections below are part of this calculation.

Let the original edge powers be $m_e(1-z)/2$ and introduce
\[
 a_e(z,d)=m_e(1-z)/2+d\,v_e.
\]
The chosen $v_e$ have zero sum over the entire graph and over every proper logarithmically divergent subgraph. Therefore $d$ changes shapes of the integrals but not their logarithmic degrees.
At fixed $d$, write
\begin{equation}
 C_G(z,d)=\sum_{k\ge1}\frac{c_{G,k}(d)}{z^k},\qquad r_G(d)=c_{G,1}(d).
\end{equation}
The higher-pole identity is
\begin{equation}
 L_Gc_{G,k}(d)=\sum_{\gamma\subsetneq G}
 L_\gamma r_\gamma(d)c_{G/\gamma,k-1}(d),\qquad k\ge2,
 \label{gff:eq:higher}
\end{equation}
where the sum is over connected logarithmic one-particle-irreducible proper subgraphs. In particular, the simple residue belongs to the subgraph, not to the quotient.

Define the higher-pole part
\begin{equation}
 H_G(z,d)=\frac1{L_G}\sum_{\gamma\subsetneq G}
 \frac{L_\gamma r_\gamma(d)}z\,C_{G/\gamma}(z,d).
\end{equation}
Let $\mathfrak F'$ denote nonempty sets of disjoint proper logarithmic subgraphs; nested subgraphs are already included inside each $C_\gamma$. The physical residue is obtained from an integrable diagonal $d=z$ by
\begin{equation}\boxed{\begin{aligned}
 r_G(0)={}&-\Res_{z=0}\left[
 I_G(z,z)+\sum_{F\in\mathfrak F'}
 \Big(\prod_{\gamma\in F}C_\gamma(z,z)\Big)I_{G/F}(z,z)\right]\\
 &-\Res_{z=0}H_G(z,z).
\end{aligned}}\label{gff:eq:restore}
\end{equation}
Crucially, $C_\gamma(z,z)$ means first form the fixed-$d$ counterterm and then set $d=z$. It is not the counterterm obtained by setting $d=z$ before subtraction.

To prove \eqref{gff:eq:restore}, expand the fixed-$d$ forest-subtracted amplitude. Its poles are $-C_G(z,d)$. On the diagonal, the coefficient of $z^{-1}$ receives $-r_G(0)$ and also Taylor coefficients of the higher poles. Equation~\eqref{gff:eq:higher} identifies these additional terms as $-\Res H_G(z,z)$, giving the stated formula. This retains precisely the finite regulator corrections that were missing from bare uniqueness reductions.

\subsection{Seven regulator-corrected graph classes}
For six vertices use edge-word order
\[
(01,02,03,04,05,12,13,14,15,23,24,25,34,35,45).
\]
Class numbers are zero-based positions of lexicographically minimal edge words in the sorted 129-class list. The shifts used are:
\begin{center}\small
\begin{tabular}{@{}r l l@{}}\toprule
Class & Nonzero $v_e$ & Fixed roots\\\midrule
52 & $v_{13}=-2,\ v_{23}=2$ & $0,4$\\
56 & $v_{14}=-2,\ v_{24}=2$ & $0,4$\\
60 & $v_{14}=-2,\ v_{45}=2$ & $0,2$\\
79 & $v_{23}=2,\ v_{15}=-2$ & $0,4$\\
80 & $v_{12}=-2,\ v_{13}=2$ & $0,4$\\
91 & $v_{14}=2,\ v_{15}=-2,\ v_{25}=2,\ v_{34}=-2$ & $0,4$\\
127& $v_{05}=-4,\ v_{14}=2,\ v_{23}=2$ & $1,2$\\\bottomrule
\end{tabular}
\end{center}
The exceptional four-loop class with word \texttt{0011101211} uses $v_{24}=2,v_{34}=-2$, with roots $0,1$. Its shift is antisymmetric under a graph automorphism, so the first derivative of its simple residue vanishes. Its higher-pole derivatives are still retained through \eqref{gff:eq:higher}.

Ordinary degree-three unique stars reduce six of the five-loop classes. Class 79 uses a two-internal-vertex conformal three-point block. On its diagonal both internal vertices have incident power sum two, and conformal covariance fixes the external dependence. With $h=(1-z)/2$, its bare integral becomes
\begin{equation}
 I_{79}^{U}(z)=\frac2{5z}M(h,h,h+2z)
 \BB(h,1-3z)\BB(h,1-4z).
\end{equation}
It is gamma-valued. The regulator-restoration terms in \eqref{gff:eq:restore} use the tetrahedron Hessian \eqref{gff:eq:Hessian}. They give
\begin{equation}
 r_{79}=3Q\left[-\frac{13A^2}{20}+\frac{\pi^2}{12}+\frac{37H}{60}\right]
 =Q\left[-\frac{39A^2}{20}+\frac{89\pi^2}{60}+\frac{74\Cat}{5}\right].
\end{equation}
Dropping those restoration terms would not give this residue.

The complete resulting seven residues are recorded compactly below. Define
\begin{align}
 X={}&-\frac{25A^4}{12}+\frac{5A^2\pi^2}{6}
 +\frac{52A\pi^3}{15}+\frac{512A\zeta(3)}5-\frac{11\pi^4}{900},\\
 Y={}&Q\left[-\frac{39A^2}{20}+\frac{89\pi^2}{60}+\frac{74\Cat}{5}\right].
\end{align}
Then
\begin{align}
 r_{52}&=Q\left[-\frac{3A^2}{2}+\frac{41\pi^2}{30}+\frac{76\Cat}{5}\right],\\
 r_{56}&=X-\frac{16}5\KK,\\
 r_{60}&=-\frac{47A^4}{240}-\frac{11A^2\pi^2}{120}
 +\frac{19A\pi^3}{60}+\frac{43A\zeta(3)}5+\frac{3\pi^4}{400},\\
 r_{79}&=Y,\\
 r_{80}&=X+\frac45\KK,\\
 r_{91}&=-\frac{79A^4}{80}-\frac{43A^2\pi^2}{360}
 +\frac{17A\pi^3}{12}+\frac{117A\zeta(3)}5-\frac{623\pi^4}{3600},\\
 r_{127}&=Y.
\end{align}
These are unweighted counterterm residues. Their graph weights, words, and the residue expressions for all other required classes are in the residue certificate of Appendix~\ref{gff:app:certificate}, and in machine-readable form in \path{lri_numerical_resolution/verification/inputs/residue_certificate.json}.


\section{The six nonprimitive wheel-insertion topologies}\label{gff:sec:wheelclassupdated}
A primitive graph has no ultraviolet-divergent proper subgraph. A nonprimitive graph requires subtraction of such subgraphs in addition to its overall pole. The six classes in this section contain a one-loop bubble or a four-loop wheel subgraph. They were evaluated separately from the remaining master-integral classes and must be included exactly once.

Using the edge-word order of Section~\ref{gff:sec:shifts}, define
$d_G=-5(3w_2-w_4)/|\operatorname{Aut}G|$. The unweighted counterterm residues are
\begin{center}\small
\renewcommand{\arraystretch}{1.25}
\begin{tabular}{@{}r l r r l@{}}\toprule
Class & Edge word & $|\operatorname{Aut}G|$ & $d_G$ & $r_G$\\\midrule
51 & \texttt{000110111111110} &4&$15/2$&$-2AP_4/5-3\LL/10$\\
78 & \texttt{000111101110111} &2&$15$&$-2AP_4/5-\LL/10$\\
99 & \texttt{000120111111100} &2&$30$&$-2AP_4/5-\LL/10$\\
113& \texttt{000121101110110} &1&$60$&$-2AP_4/5+\LL/10$\\
124& \texttt{001110111112000} &8&$15/2$&$-2AP_4/5-\LL/2$\\
128& \texttt{001111011201100} &4&$15$&$-2AP_4/5-3\LL/10$\\\bottomrule
\end{tabular}
\end{center}
The table includes the subdivergence subtractions; these are not bare integrals. Weighted summation gives
\begin{equation}
 (3c_5-b_5)_{\rm wheel\ insertions}=-54AP_4-9\LL.
\end{equation}
Since $b_4^{\rm wheel}=-6P_4$ and $c_4^{\rm wheel}=0$, the lower-loop assembly supplies $60AP_4$. Thus
\begin{equation}
 C_5^{\rm wheel+assembly}=\frac{(-54AP_4-9\LL)+60AP_4}{729}
 =\frac{2AP_4-3\LL}{243}.
\end{equation}
The finite star expansion and the regulated wheel derivatives reduce the six residues to $P_4$ and $\LL$. The test programs are not part of the public repository; the detailed derivation is given below. This sector is not the entire nonprimitive contribution.

\subsection{Derivation of the residue table}
\subsubsection{First regulator derivative of the full-line star}
Introduce
\[
 F_h(x)=\int_\RR |u(1-u)(u-x)|^{-h}\dd u.
\]
For $0<x<1$, split the real line at $0,x,1$ to obtain
\begin{align}
F_h(x)={}&B(1-h,1-h)\Big[
 x^{1-2h}{}_2F_1(h,1-h;2-2h;x)\notag\\
&\hspace{32mm}+(1-x)^{1-2h}{}_2F_1(h,1-h;2-2h;1-x)\Big]\notag\\
&+B(1-h,3h-1)\Big[
{}_2F_1(h,3h-1;2h;x)
+{}_2F_1(h,3h-1;2h;1-x)\Big].\label{gff:wi:eq:fourchambers}
\end{align}
At $h=1/2$, the lower-parameter derivative terms cancel in the sum. Euler's transformation \cite{gff:DLMF15} gives
\[
 \left.\partial_a{}_2F_1(a,b;1;x)\right|_{a=b=1/2}
 =-\frac12\log(1-x){}_2F_1(1/2,1/2;1;x).
\]
Using this in \eqref{gff:wi:eq:fourchambers}, the beta-function derivative terms cancel as well, leaving
\begin{equation}\boxed{
 \left.\partial_hF_h(x)\right|_{h=1/2}
 =-\log|x(1-x)|F(x).
}\label{gff:wi:eq:Fprime}\end{equation}
The tails follow from $F_h(x)=x^{1-3h}F_h(1/x)$ for $x>1$ and the corresponding negative-tail transformation.

In the opposite-outer-root representation of the wheel, put
\[
 P_W(z)=\int_\RR |x(1-x)|^{-h}F_h(x)^2\dd x,
 \qquad h=(1-z)/2.
\]
Equation \eqref{gff:wi:eq:Fprime} implies
\begin{equation}\boxed{P_W'(0)=\frac32\LL.}\label{gff:wi:eq:Pprime}\end{equation}

\subsubsection{Edge derivatives and root changes}
In this rooted wheel, the fixed roots are opposite outer vertices. The remaining variables are the center $x$ and the two other outer vertices. Let $D_eP$ denote the derivative with respect to an independent edge power at the point where all eight powers are $1/2$.

The completed wheel is the octahedral graph. For conformally balanced index variations, averaging over its edge-transitive automorphism group makes the first variation vanish. In the opposite-root representation, the allowed deformations obey zero total edge shift and zero sum of spoke shifts. It follows that the index gradient is constant separately on the four rim edges and the four spokes. The two root-to-center spokes differentiate the explicit factor $|x(1-x)|^{-1/2}$, so each contributes $-\LL/2$. Equation \eqref{gff:wi:eq:Fprime} for one outer integrated vertex then fixes the rim derivative. Thus
\begin{equation}\boxed{
 D_{\mathrm{spoke}}P=-\frac12\LL,\qquad
 D_{\mathrm{rim}}P=-\frac14\LL.
}\label{gff:wi:eq:edgegrad}\end{equation}
Changing the fixed outer pair from opposite to adjacent gives, by the scale Jacobian,
\[
 P_W^{\mathrm{adj}\,\prime}(0)
 =P_W^{\mathrm{opp}\,\prime}(0)+4D_{\mathrm{rim}}P
 =\frac12\LL.
\]
Only these first-derivative statements are needed; they are not identities for arbitrary higher index variations.

\subsubsection{Bubble insertions on a rim or spoke}
The full-line convolution coefficient is
\[
\BB(u,v)=\sqrt\pi\,
\frac{\Gamma((1-u)/2)\Gamma((1-v)/2)\Gamma((u+v-1)/2)}
{\Gamma(u/2)\Gamma(v/2)\Gamma((2-u-v)/2)}.
\]
In particular, $\BB(h,2h)=2/z+A+O(z)$. For classes 99 and 113 a degree-three bubble endpoint can be integrated exactly. The rooted bare integral is
\[
 I_G(z)=\frac{2}{5z}\BB(h,2h)P_e(z),\qquad
 P_e'(0)=P_W'(0)-D_eP.
\]
The bubble counterterm contributes $-P_W(z)/z^2$ to the subtraction bracket. Its simple-pole extraction yields
\[
 r_G=-\frac25AP_4+\frac15P_W'(0)+\frac45D_eP.
\]
The modified quotient edge is a spoke for class 99 and a rim for class 113. Equations \eqref{gff:wi:eq:Pprime}--\eqref{gff:wi:eq:edgegrad} give their table entries.

\subsubsection{Splitting the degree-four center}
For classes 124 and 128 one needs the finite part of a star with one doubled edge. The local distributional expansion is
\begin{align}
&\int_\RR |u-y|^{-h}|u-z_1|^{-h}|u-x|^{-1+z}\dd u\notag\\
&\quad=|x-y|^{-h}|x-z_1|^{-h}
\left[\frac2z+A+2\log\frac{|x-y|\,|x-z_1|}{|y-z_1|}\right]+O(z).
\label{gff:wi:eq:starfinite}
\end{align}
One derivation shifts the two single-edge powers from $h$ to $h+z$, applies the unique star--triangle integral, and restores the finite logarithmic terms from the local bubble residue. Those logarithms cannot be dropped.

After subtraction,
\[
 r_G=-\frac25AP_4+\frac15P_W'(0)-\frac45\LL_\triangle,
\]
where $\LL_\triangle$ inserts the logarithm in \eqref{gff:wi:eq:starfinite} into the wheel. For opposite neighbors, $\LL_\triangle=\LL$. For adjacent neighbors, \eqref{gff:wi:eq:edgegrad} gives $\LL_\triangle=3\LL/4$. This proves classes 124 and 128.

\subsubsection{A wheel inserted into the outer one-loop graph}
Classes 51 and 78 have only the wheel as a proper ultraviolet-divergent subgraph. Invert about the degree-two root and fix the other root at one. The additional powers produced by inversion differ from their zero-regulator values by $O(z)$; their logarithms vanish at the only proper divergent collision, so they do not affect the constant term of the angular integral. One obtains
\[
 J_G(z)=\BB(h+2z,1-4z)P_W^{\mathrm{root\ pair}}(z)+O(z),
 \quad \BB(h+2z,1-4z)=\frac1{2z}+A+O(z).
\]
The wheel counterterm times the bubble quotient contributes only a double pole. Hence
\[
 r_G=-\frac25AP_4-\frac15P_W^{\mathrm{root\ pair}\,\prime}(0).
\]
The pair is opposite for class 51 and adjacent for class 78, proving the remaining entries.


\section{Assembly of the full coefficient}\label{gff:sec:assembly}
Write
\[
\beta(g)=-\eps g+\sum_{L\ge1}b_Lg^{L+1},\qquad
\gamma_r(g)=\sum_{L\ge1}c_Lg^L.
\]
The long-range kinetic term has no wave-function counterterm in this setup \cite{gff:BGHS}. At the nonzero fixed point,
\[
\Delta_{\mathcal E}=(1-\eps)/2+\gamma_r(g_*),\qquad\beta(g_*)=0.
\]
With $b_1=3,c_1=1,b_2=-6A,c_2=-A$ and
\[
 b_3=6Q+\frac{87}4A^2-\frac34\pi^2-4\pi,
 \qquad c_3=\frac72A^2-\frac16\pi^2-\frac43\pi,
\]
Lagrange inversion gives
\begin{equation}\begin{aligned}
729C_5={}&336A^4-60A^2b_3+108A^2c_3+b_3^2-3b_3c_3\\
&-10Ab_4+24Ac_4+3c_5-b_5.
\end{aligned}\label{gff:eq:RG}
\end{equation}

The primitive and wheel-insertion sectors are
\begin{equation}
 C_5^{\rm primitive}=-\frac{63Q^2+24P_5}{729},\qquad
 C_5^{\rm wheel+assembly}=\frac{2AP_4-3\LL}{243}.
\end{equation}
The result for all other subtracted terms, including their lower-loop assembly terms, is
\begin{equation}\boxed{\begin{aligned}
 C_5^{\rm rest}={}&-\frac{5A^4}{15552}
 +\frac{5A^2Q}{324}-\frac{5A^2\pi^2}{7776}
 -\frac{23A\pi^3}{1296}-\frac{149A\zeta(3)}{324}\\
 &+\frac{298\Cat Q}{243}-\frac{8\KK}{243}
 +\frac{4Q^2}{81}+\frac{319Q\pi^2}{2916}
 -\frac{83\pi^4}{46656}.
\end{aligned}}\label{gff:eq:rest}
\end{equation}
The coefficient of the only new master constant is $-8/243$. Adding the three sectors and using $H=2\pi^2/3+8\Cat$ proves \eqref{gff:eq:full} by exact algebra.

The enumeration and weighting are fixed by the following graph construction. With $e_i=4-\sum_jm_{ij}$,
\[
 w_4=\frac{24}{\prod_i e_i!\prod_{i<j}m_{ij}!},\qquad
 w_2=\frac{w_4}{12}\sum_i e_i(e_i-1),
\]
and an unlabeled $L$-loop class contributes
\[
 b_L=L(-1)^L\frac{w_4}{|\operatorname{Aut}G|}r_G,
 \qquad c_L=L(-1)^L\frac{w_2}{|\operatorname{Aut}G|}r_G.
\]
The 129 five-loop classes are exhausted by 92 explicitly evaluated remaining classes, six wheel-insertion classes, three primitive classes, and 28 classes whose $3w_2-w_4$ vanishes. The 26 non-wheel four-loop classes are also included. These disjoint sets cover the full enumeration; no unknown residue is hidden in \eqref{gff:eq:rest}.

\section{Independent full-graph numerical control}\label{gff:sec:checks}\label{gff:sec:evidenceledger}\label{gff:app:numericalResolution}
The following control checks the assembled coefficient independently of
special-function evaluation. Only the precision-preserving formulation and
its converged comparison are needed here. The code and records are in the
package \path{lri_numerical_resolution/} (Section~\ref{gff:cn:reproduction}).

The graph control uses the two coordinate cutoffs and meromorphic
Chebyshev quadrature defined in Section~\ref{gff:app:C6Numerical}, specialized
to five loops: the cutoff ``maxgap'' requires the largest gap between
consecutive ordered vertices to be below one, and the cutoff ``span''
requires the total vertex span to be below one. Below, $N_{\rm lower}$ is the
quadrature degree of the lower-loop inputs. In particular its integrand is \eqref{gff:c6:sectorIntegral}
with $L=5$, and its radial moments are \eqref{gff:c6:monomial}. The bare
amplitudes are renormalized by \eqref{gff:eq:MSrecursion} before applying
\eqref{gff:eq:C5RG}. Thus the special-function period values are not inputs
to this independent integration.

The combined bare simple pole is
\[
 H_{5,-1}=[z^{-1}](3\mathcal B_5-\mathcal A_5).
\]
Its exact sensitivity, with all lower-loop inputs held fixed, is
\[
 \delta C_5=\frac5{729}\,\delta H_{5,-1}.
\]
Indeed the new fifth-order contribution is $(3c_5-b_5)/729$, while
\eqref{gff:eq:b5c5poles} gives
$\delta(3c_5-b_5)=5\delta H_{5,-1}$.
The program therefore forms the exact integer channel weights before
integration; this removes 9,540 zero-weight ordered graphs from the
59,145-graph inventory, rather than subtracting two noisy channel sums
only at the end. Quadrature weights, smooth-kernel evaluation, within-group
accumulation and final Laurent assembly are all precision-controlled.
A same-degree comparison with the full two-channel calculation changes
$C_5$ by less than $1.7\times10^{-12}$. The combined channel has 335 rather
than 342 sector exponent groups, and exhaustive exact enumeration shows that
the groups with exponent $a=-2$ disappear from it: its most negative exponent
is $-\frac32$. At degree 18, replacing only the binary64 five-loop kernel by
the extended one, with the same lower-loop inputs, changes $C_5$ by
$-7.957\times10^{-9}$ (span) and $-7.965\times10^{-9}$ (maxgap), most of the
offset of the earlier version described after the table below. For an
independent pole check, the renormalization-group equations predict the
higher poles $H_{5,-p}$, $p\ge2$, from the lower loops; the largest residuals
are $1.96\times10^{-11}$ (span) and $1.74\times10^{-11}$ (maxgap). They are
diagnostics only and do not certify the simple pole, which carries the new
information.

With lower-loop degree $N_{\rm lower}=32$, the complete results are
\begin{center}
\begin{tabular}{@{}c l r r@{}}
\toprule
Degree & Cutoff & Graph value of $C_5$ & Graph value minus closed form\\\midrule
22 & span & $-731.1394766245375925$ & $+2.2334\times10^{-12}$\\
22 & maxgap & $-731.1394766245106192$ & $+2.9207\times10^{-11}$\\
24 & span & $-731.1394766245391892$ & $+6.3667\times10^{-13}$\\
24 & maxgap & $-731.1394766245368057$ & $+3.0201\times10^{-12}$\\
\bottomrule
\end{tabular}
\end{center}
The comparison target is the special-function evaluation
$-731.139476624539825853869126582527\ldots$.
These are observed residuals, not rigorous bounds on either representation.
The cutoff difference alone is not an accuracy certificate: a shared
regular contribution can have a common error that cancels in that
difference. The complete-channel and precision/degree comparisons are
therefore separate checks. An earlier version of this control, whose quadrature kernel used binary64
arithmetic, gave $-731.1394766152787$ with the span cutoff, $9.26\times10^{-9}$
above the closed form. The precision-preserving formulation removes that
offset; the analytic coefficient was not changed. The signed error ledger of
the earlier value is in
\path{lri_numerical_resolution/verification/records/resolution_audit.json}
(field \path{error_budgets}).

\subsection{Further checks of the analytic coefficient}\label{gff:sec:furtherChecks}
The following checks of the one-dimensional result were run with the
analytic calculation and do not use the full graph integration above.
\begin{itemize}
\item \emph{Exact algebra.} The two regular quadratic coefficients at pattern
$A_0$ vanish; the fourth-order weighted sum reproduces the analytic $3c_4-b_4$,
including all gamma and zeta terms; every five-loop class is accounted for
exactly once in the nonzero, primitive, wheel or zero-weight sets; and the
rationally weighted residue table sums to \eqref{gff:eq:rest}, whose addition to
the primitive and wheel sectors gives \eqref{gff:eq:full}. No coefficient is
inferred from a numerical integer-relation fit
(\path{lri_numerical_resolution/verification/records/symbolic_audit.json}).
\item \emph{Seven subtracted graphs.} With extended precision in the
quadrature weights, integrands and accumulation, seven classes were
integrated anew, with the analytic residues used only as comparison targets.
At interpolation degree 18 the numerical minus analytic residues are
$-1.02\times10^{-9}$ (class 52), $-4.95\times10^{-9}$ (56),
$+3.31\times10^{-10}$ (60), $-2.25\times10^{-9}$ (79), $-5.19\times10^{-10}$
(80), $-5.82\times10^{-11}$ (91) and $+2.27\times10^{-10}$ (127); at degree
22 those of classes 56 and 79 fall to $-1.02\times10^{-11}$ and
$-1.92\times10^{-11}$. The lower-loop inputs of these comparisons were
evaluated as bare span-cutoff integrals. The checks cover gamma-valued and
$\KK$-dependent residues; their record is not part of the repository.
\item \emph{Fresh evaluation of the masters.} Evaluating all four masters
without stored period values, at 65 working digits with 64 Cauchy nodes for
the parameter derivatives, reproduces
$C_5=-731.1394766245398258538691265825272306625\ldots$; a separate 28-digit
run agrees to its requested precision (record not part of the repository;
the replay of Section~\ref{gff:cn:reproduction} recomputes $C_5$ from the closed
forms).
\item \emph{Regenerated certificate.} In the code delivery that preceded
this note, the symbolic task regenerated all 26 non-wheel four-loop and 92
retained five-loop residues of Appendix~\ref{gff:app:certificate} from the
graph-class data, the master expansions and the forest rules, without reading
a saved residue; all 118 comparisons were exact. That record is not part of
the repository. Regenerating the table is stronger than adding the saved
rows, but it uses the same master identities and reduction code.
\end{itemize}
These tests, like the graph integration, do not establish the physical
five-loop coefficient independently of the shared graph and renormalization
framework.

\subsection{External lattice comparison: informative but inconclusive}\label{gff:cnm:latticeComparison}
The Monte Carlo data of Uzelac, Glumac and Ani\v ci\'c concern the same
classical long-range chain, with $J(r)\propto r^{-(1+s)}$ \cite{gff:cnm:UGA}.
Their Table~II gives $y_t=\nu^{-1}$, which here is $1-\Delta_{\mathcal E}$,
at three couplings. An earlier comparison resummed the series of $y_t$
through $\epsilon^5$ in six Borel--Pad\'e--Leroy forms: the approximants
$[3/2]$ and $[2/3]$ of the Borel--Leroy transform
$\sum_kc_kt^k/\Gamma(k+b+1)$, with $b=0,1,2$, followed by the Laplace
integral $\int_0^\infty e^{-t}t^b[L/M](\epsilon t)\,dt$. None of the six
approximants has a positive real pole at these couplings. No Monte Carlo
value was fitted.
\begin{center}
\begin{tabular}{@{}rrrr@{}}\toprule
$s$ & $\epsilon$ & $y_t$ of \cite{gff:cnm:UGA} & Resummation range\\\midrule
0.6&0.2&$0.50\pm0.02$&$0.51738$--$0.51774$\\
0.7&0.4&$0.51\pm0.02$&$0.51518$--$0.51886$\\
0.8&0.6&$0.47\pm0.02$&$0.50021$--$0.51231$\\\bottomrule
\end{tabular}
\end{center}
The third row shows a tension, and it must not be omitted from the
interpretation: at $s=0.8$ all six resummed values lie above the Monte Carlo
interval, the lowest by $0.010$. The resummation spread is not a confidence
interval. Changing $C_5$ by $20\%$ in either direction still gives pole-free
resummed values inside the Monte Carlo interval nearest to the expansion
point, at $s=0.6$, so these data do not isolate the fifth coefficient. With
$C_5$ lowered by $20\%$, however, the three $[3/2]$ approximants acquire
positive real Borel poles at all three couplings, and only the three $[2/3]$
values remain; with $C_5$ raised by $20\%$, all six stay pole-free. The
closed-form $C_0,\ldots,C_5$ of Section~\ref{gff:cn:oneDresults} reproduce the
three ranges to the digits shown:
\path{lri_numerical_resolution/public/lattice_comparison.py} recomputes the
ranges, the absence of positive Borel poles and the $20\%$ statement. No
simulation was run for this note.

\SourcePart{Arbitrary-dimensional analytic coefficients and specializations}

\section{General-dimensional normalization and lower-order derivation}\label{gff:app:generalDimension}

\subsection{Dimension, regulator and normalization}
In this section, $d$ always denotes the spatial dimension, not the dummy
exponent called $d$ in some earlier one-dimensional subintegrals. Take
\begin{equation}
 0<d<4,\qquad \epsilon=2s-d,\qquad
 \zeta=\frac s2=\frac{d+\epsilon}{4},\qquad
 \Delta_\phi=\frac{d-\epsilon}{4}.
 \label{gff:gd:conventions}
\end{equation}
The physical long-range kinetic term has $0<\zeta<1$; at fixed $d<4$ this
holds for sufficiently small positive $\epsilon$. Noninteger $d$ is defined
by analytic continuation of the dimensionally parameterized integrals.
Analytically continuing individual periods to $d=4$ is useful as a check,
but it is not the same as replacing this expansion by the short-range
Wilson--Fisher expansion: the absence of a wave-function counterterm, used
below, is a long-range statement \cite{gff:BGHS,gff:BGH}.

After tuning the relevant local mass term to the critical surface,
\begin{equation}
 \Delta_{\Eop}(d,\epsilon)=d-\nu^{-1}
 =\frac{d-\epsilon}{2}+\gamma_r(g_*),\qquad
 \beta(g)=-\epsilon g+\sum_{\ell\geq1}b_\ell(d)g^{\ell+1},\qquad
 \gamma_r(g)=\sum_{\ell\geq1}c_\ell(d)g^\ell.
 \label{gff:gd:RG}
\end{equation}
We use a minimal-subtraction convention with $b_1=3$, $c_1=1$ and an
exactly pole-only one-loop bubble. Write $z$ for the analytic regulator and
introduce the infrared covariance $(p^2+\mu^2)^{-(d+z)/4}$. The bare coupling
normalization is
\begin{equation}
 \lambda_0=\mu^z\kappa(d,z)\mathcal Z_g(g,z),\qquad
 \kappa(d,z)=(4\pi)^{d/2}
       \frac{\Gamma((d+z)/2)}{\Gamma(1+z/2)}.
 \label{gff:gd:normalization}
\end{equation}
The parameter $\mu$ is an infrared regulator, not the coefficient of the
physical $\phi^2$ perturbation. A direct momentum integration gives
\begin{equation}
 \kappa(d,z)\int\frac{d^dp}{(2\pi)^d}
       (p^2+1)^{-(d+z)/2}=\frac2z.
 \label{gff:gd:bubble}
\end{equation}
Hence the one-loop four-point and mass amplitudes are $3/z$ and $1/z$.
The extra $z$-dependent normalization in \eqref{gff:gd:normalization} must be
kept before extracting Laurent coefficients. Finite parts computed in this
massive-Schwinger prescription must not be mixed with finite parts from
the one-dimensional ordered-span prescription of Section~\ref{gff:sec:graphs}.

\subsection{The evaluated coefficients through the third order}
Let
\begin{align}
 q&=d/4,&
 A_d&=\psi_0(2q)+\psi_0(1)-2\psi_0(q),\notag\\
 R_d&=\psi_1(2q)-\psi_1(1),&
 M_d&=\frac{\Gamma(-q)\Gamma(2q)^2}{\Gamma(3q)}
     =-\frac{\Gamma(1-q)\Gamma(2q)^2}{q\Gamma(3q)},
 \label{gff:gd:blocks}
\end{align}
and let $\mathcal T_d$ be the corrected tetrahedron beta-function
coefficient $\alpha_{I_4}$ of \cite{gff:BGH}. Its independent integral definition
is given in \eqref{gff:gd:tetraMB}; put
\begin{equation}
 Q_d=\frac{\mathcal T_d}{6}.
 \label{gff:gd:Qdef}
\end{equation}
In particular $Q_1=Q$, but $Q_d$ is not defined by replacing the argument
of the Gamma functions in the one-dimensional value $Q$.
The coefficient formulas are
\begin{equation}\boxed{\begin{aligned}
 C_0(d)&=\frac d2,&C_1(d)&=-\frac16,&C_2(d)&=\frac{A_d}{9},\\
 C_3(d)&=\frac{A_d^2+R_d-8Q_d}{108}
        =\frac{A_d^2+R_d}{108}-\frac{\mathcal T_d}{81}.
\end{aligned}}\label{gff:gd:C03}\end{equation}
Here and below the polygamma convention is
$\psi_m(z)=d^{m+1}\log\Gamma(z)/dz^{m+1}$.
The corresponding three-loop MS coefficients in \eqref{gff:gd:normalization}
are
\begin{equation}\boxed{\begin{aligned}
 b_1&=3,&c_1&=1,&b_2&=-6A_d,&c_2&=-A_d,\\
 b_3&=6Q_d+\frac{87}{4}A_d^2-\frac94R_d+M_d,&
 c_3&=\frac72A_d^2-\frac12R_d+\frac13M_d.
\end{aligned}}\label{gff:gd:bc3}\end{equation}
The $M_d$ term cancels from the physical $C_3(d)$, but it cannot be dropped
from the RG functions used at fourth and fifth order.

At $d=1$,
\begin{equation}
 A_1=A,\qquad R_1=\frac{\pi^2}{3},\qquad M_1=-4\pi,
 \qquad Q_1=\frac{\Gamma(1/4)^4}{\pi}.
 \label{gff:gd:d1blocks}
\end{equation}
This recovers \eqref{gff:eq:bc3} and the one-dimensional $C_3$ of
Section~\ref{gff:cn:oneDresults}. Equation
\eqref{gff:gd:blocks} also displays a concrete obstruction to blind argument
substitution: $3q$ and $1-q$ cancel in $M_d$ only at $d=1$. In addition,
the reflected identity for $R_d$ is not obtained by keeping
$\psi_1(q)+\psi_1(1-q)=2\pi^2$ away from $q=1/4$.

\subsubsection{Derivation from the published three-loop calculation}
The three-loop beta functions of \cite{gff:BGHS}, Eqs.~(2.13) and~(2.14) there,
are sums over graphs with coefficients $\alpha_X$: $\alpha_D$ for the one-loop
graph, $\alpha_S$ for the two-loop graph, and $\alpha_U$, $\alpha_T$ and
$\alpha_{I_1},\ldots,\alpha_{I_4}$ for the seven three-loop graphs, each the
normalized graph amplitude of Eq.~(2.12) there. The tetrahedron coefficient
$\alpha_{I_4}$ given in \cite{gff:BGHS} was corrected in \cite{gff:BGH}; the corrected
value is required here. For clarity,
let $t_1=\psi_1(1)$, $t_q=\psi_1(q)$, $t_2=\psi_1(2q)$ and
$J=J_0(q)$. The coefficients quoted in \cite{gff:BGH}, specialized to one scalar,
are
\begin{align}
 \alpha_U=\alpha_{I_2}&=-t_1-t_q+2t_2+J,\notag\\
 \alpha_T&=\tfrac12A_d^2+\tfrac12t_1+t_q-\tfrac32t_2-J,\notag\\
 \alpha_{I_1}&=\tfrac32A_d^2+\tfrac12(t_1-t_2),&
 \alpha_{I_3}&=M_d/3,&\alpha_{I_4}&=6Q_d.
 \label{gff:gd:alphas}
\end{align}
An explicit definition, useful for reproducing the intermediate cancellation,
is
\begin{equation}\begin{aligned}
 J_0(q)=\frac1{\Gamma(q)^2}\sum_{n=1}^{\infty}
 &\frac{\Gamma(n+2q)\Gamma(n+q)^2}{n\,n!\,\Gamma(2q+2n)}\\
 &\times\left[2\psi_0(n+1)-\psi_0(n)-2\psi_0(n+q)
   -\psi_0(n+2q)+2\psi_0(2q+2n)\right].
\end{aligned}\label{gff:gd:J0}\end{equation}
After dividing the zero-momentum one-loop normalization out of the
coupling, write $\alpha_S/\alpha_D^2=-A_d+s_1z+O(z^2)$. The expansion
in \cite{gff:BGH} gives
\begin{equation}
 s_1=\frac14A_d^2+\frac34t_1+t_q-\frac74t_2-J.
 \label{gff:gd:s1}
\end{equation}
The two-loop bare amplitudes are then
\begin{equation}
 \mathcal A_2=\frac9{z^2}+\frac{3A_d}{z}-3s_1+O(z),\qquad
 \mathcal B_2=\frac2{z^2}+\frac{A_d}{2z}-\frac{s_1}{2}+O(z).
 \label{gff:gd:twoLoopPoles}
\end{equation}
Consequently, at $z=0$, the zero-momentum and MS variables obey
\begin{equation}
 g_{\rm ZM}=g_{\rm MS}-3s_1g_{\rm MS}^3+O(g^4),\qquad
 r_{\rm ZM}=r_{\rm MS}(1-\tfrac12s_1g_{\rm MS}^2+O(g^3)).
 \label{gff:gd:schemechange}
\end{equation}
These finite transformations imply
\begin{align}
 b_3^{\rm MS}&=3\alpha_T+12\alpha_U+12\alpha_{I_1}
                  +3\alpha_{I_3}+\mathcal T_d+9s_1,\notag\\
 c_3^{\rm MS}&=\alpha_T+2\alpha_{I_1}+\alpha_U+\alpha_{I_3}.
 \label{gff:gd:schemeb3}
\end{align}
Substitution yields \eqref{gff:gd:bc3}: all $J$ and $t_q$ terms cancel.
Alternatively, the physical coefficient can be checked directly before
choosing the MS scheme:
\begin{equation}\begin{aligned}
 C_3&=-\frac1{81}\left[-12A_d^2+9s_1+9\alpha_U
                            +6\alpha_{I_1}+\mathcal T_d\right]\\
 &=\frac{A_d^2+R_d}{108}-\frac{\mathcal T_d}{81}.
\end{aligned}\label{gff:gd:C3direct}\end{equation}
The cancellation is symbolic; it is not an integer-relation fit.

\subsection{A specified tetrahedron at every dimension}
The equal-index specialization of the corrected kite representation in
\cite{gff:BGH}, eqs.~(3.8)--(3.11), is
\begin{equation}\boxed{\begin{aligned}
\mathcal T_d={}&\frac{2\Gamma(2q)^2}{\Gamma(q)^4}
 \int_{-q/3-\ii\infty}^{-q/3+\ii\infty}\frac{du}{2\pi\ii}
 \int_{-q/3-\ii\infty}^{-q/3+\ii\infty}\frac{dv}{2\pi\ii}\\[-1mm]
&\quad\times
 \frac{\Gamma(-u)^2\Gamma(-v)^2\Gamma(q+u)\Gamma(q+v)
       \Gamma(q+u+v)^2\Gamma(-u-v)}
      {\Gamma(q-u)\Gamma(q-v)\Gamma(2q+u+v)}.
\end{aligned}}\label{gff:gd:tetraMB}\end{equation}
Both contours run upwards. Every numerator Gamma argument has positive
real part on these contours for $q>0$, so no pole separation is left
implicit. If $u=-q/3+\ii x$, $v=-q/3+\ii y$, the exponential part of
Stirling's bound is $\exp[-\pi(|x|+|y|+|x+y|)]$; the double integral is
absolutely convergent. There are no hidden regulator derivatives in this
representation. It is a twofold Mellin--Barnes integral, not a claim of
an ordinary single-variable hypergeometric evaluation for generic $d$.

For an independent geometric normalization, let $|e|=1$ and
\begin{equation}\begin{aligned}
 P_3^{\rm raw}(d)&=\int_{\mathbb R^d}\!d^dx\int_{\mathbb R^d}\!d^dy\,
 \left[x^2(x-e)^2y^2(y-e)^2(x-y)^2\right]^{-q},\\
 h_d&=\frac{\Gamma(d/2)}{\pi^{d/2}},\\
 \mathcal T_d&=2h_d^2P_3^{\rm raw}(d),\qquad
 Q_d=\frac{h_d^2}{3}P_3^{\rm raw}(d).
\end{aligned}\label{gff:gd:tetraCoordinate}\end{equation}
The spatial integral defines the same analytic family; the MB formula
provides a direct meaning when $d$ is not an integer.

\subsubsection{An exact two-dimensional Gamma evaluation}
For $d=2$ identify $\mathbb R^2$ with $\mathbb C$. Then
\begin{equation}
 P_3^{\rm raw}(2)=\int_{\mathbb C^2}
 \frac{d^2z\,d^2w}{|z|\,|1-z|\,|w|\,|1-w|\,|z-w|}.
 \label{gff:gd:complexTetra}
\end{equation}
Write $\gamma_C(a)=\Gamma(a)/\Gamma(1-a)$ for real equal holomorphic and
antiholomorphic parameters. The complex Selberg identity
\cite{gff:NeretinComplex}, with $n=2$, $\sigma=\tau=1/2$, $\theta=-1/4$, gives
\begin{align}
 P_3^{\rm raw}(2)
 &=2\pi^2\prod_{j=1}^{2}
 \frac{\gamma_C(1/2+(j-1)\theta)^2\gamma_C(j\theta)}
      {\gamma_C(1+j\theta)\gamma_C(\theta)}\notag\\
 &=\frac{\pi^2}{2}
       \left[\frac{\Gamma(1/4)}{\Gamma(3/4)}\right]^4.
 \label{gff:gd:complexSelberg}
\end{align}
The convergence inequalities hold strictly:
$\theta>-1/2$, $\sigma+\theta=\tau+\theta=1/4>0$,
$\sigma+\tau+\theta=3/4<1$, and
$\sigma+\tau+2\theta=1/2<1$.
Thus this is a convergent integral evaluation, not a formal continuation
of a divergent Selberg integral. Using $h_2=1/\pi$ yields
\begin{equation}\boxed{\begin{aligned}
 \mathcal T_2&=\left[\frac{\Gamma(1/4)}{\Gamma(3/4)}\right]^4,
 &Q_2&=\frac16\left[\frac{\Gamma(1/4)}{\Gamma(3/4)}\right]^4,\\
 C_2(2)&=\frac{4\log2}{9},
 &C_3(2)&=\frac{4\log^22}{27}
   -\frac1{81}\left[\frac{\Gamma(1/4)}{\Gamma(3/4)}\right]^4.
\end{aligned}}\label{gff:gd:d2exact}\end{equation}
The fact that quarter-arguments reappear at $d=2$ here is a consequence
of the complex Selberg reduction, not a justification for replacing
arguments in the $d=1$ formulas.

\subsection{Dimension-independent graph data and the d-dimensional kernel}
Use precisely the finite matrix set in \eqref{gff:eq:graphs}: $L+1$ vertices,
$2L$ internal edges, $m_{ij}\in\{0,1,2,3\}$, internal degrees between two
and four, and every nontrivial cut containing at least two edges. The
weights $w_4,w_2$ are those in \eqref{gff:eq:weights}. In a full $d$-dimensional
integral there is \emph{no ordering of vertices on a line}. Therefore the
labeled sum now includes the vertex factorial:
\begin{equation}\boxed{\begin{aligned}
 \mathcal A_L(d,z)&=\frac1{(L+1)!}\sum_{m\in\Gr_L}w_4(m)\mathcal I_m(d,z),\\
 \mathcal B_L(d,z)&=\frac1{(L+1)!}\sum_{m\in\Gr_L}w_2(m)\mathcal I_m(d,z).
\end{aligned}}\label{gff:gd:amplitudes}\end{equation}
Equivalently, sum one representative of each vertex-isomorphism class
with weights $w_4/|\Aut_v(m)|$, $w_2/|\Aut_v(m)|$. Parallel-edge
factorials are already present in $w_4,w_2$; $\Aut_v$ counts vertex
permutations, not a second set of edge permutations.

Let the $2L$ individual internal lines, including parallel lines, be
labeled $e=1,\ldots,2L$. The first Symanzik polynomial and simplex
measure are
\begin{equation}
 U_m(x)=\sum_{\text{spanning trees }T}\prod_{e\notin T}x_e,
 \qquad d\Omega=\delta\!\left(1-\sum_{e=1}^{2L}x_e\right)
                       \prod_{e=1}^{2L}dx_e.
 \label{gff:gd:U}
\end{equation}
Put $a=(d+z)/4$ and $r(d,z)=\Gamma(2a)/\Gamma(1+z/2)$. The complete
normalized graph kernel is
\begin{equation}\boxed{
 \mathcal I_m(d,z)=
 \frac{r(d,z)^L\Gamma(Lz/2)}{\Gamma(a)^{2L}}
 \int_{x_e\geq0}d\Omega\,
       \frac{\prod_{e=1}^{2L}x_e^{a-1}}{U_m(x)^{d/2}}.
}\label{gff:gd:kernel}\end{equation}
For a primitive graph it suffices to set $z=0$ in the simplex factor
when computing the residue. For nonprimitive graphs one must first
meromorphically continue the integral and apply the subtraction
recursion; the simplex integral at $z=0$ need not converge.

To derive \eqref{gff:gd:kernel}, apply a Schwinger parameter to every
$(p^2+1)^{-a}$, perform the $L$ Gaussian momentum integrals, and set the
Schwinger parameters equal to $\lambda x_e$. The powers of the common
scale are
\[
 \lambda^{2La-1-dL/2}=\lambda^{Lz/2-1},
 \qquad \int_0^\infty e^{-\lambda}\lambda^{Lz/2-1}d\lambda
       =\Gamma(Lz/2).
\]
The factors $(4\pi)^{-Ld/2}$ cancel against \eqref{gff:gd:normalization}.
For the one-loop bubble $U=x_1+x_2=1$ on the simplex,
\eqref{gff:gd:kernel} gives exactly $2/z$, confirming both the factor of two
and the vertex-factorial convention.

\paragraph{Power subgraphs and continuation.}
For fixed $0<d<4$ a proper unmarked two-point subgraph has a local
power divergence, not an independent wave-function counterterm. We use
the analytic critical-mass prescription of \cite{gff:BGHS}, as in the
one-dimensional calculation: its local mass power divergence is removed
by critical tuning and its analytic continuation is retained. In
particular, graphs with triple lines or vacuum-attached pieces are not
silently dropped. An ordinary domain from which to continue the raw
Schwinger integrals is real $z>d/3$: for an induced subgraph with $v$
vertices and $e$ edges, the ultraviolet condition is
$ez+d(2v-2-e)>0$, and the worst tadpole-free two-vertex case is $e=3$.
Continuing to $z=0$ implements the specified analytic regularization;
subsequent logarithmic subdivergences are removed by MS. No cutoff on
individual divergent simplex faces is to be substituted for this
prescription.

\begin{center}
\begin{tabular}{@{}rrrrr@{}}
\toprule
$L$ & Labeled matrices & Vertex classes & Primitive classes
 & Nonzero labeled $3w_2-w_4$\\\midrule
1&1&1&1&0\\
2&6&2&0&3\\
3&85&8&1&55\\
4&1890&27&1&1470\\
5&59145&129&3&49605\\\bottomrule
\end{tabular}
\end{center}
These counts are for the displayed matrix set, not for a minimal
nonprimitive master basis. Some classes can disappear or combine only
after analytic continuation and subtraction. The inventory file
\path{lri_general_special_functions_update/verification/graph_inventory.json},
written by \path{code/graph_inventory.py} of that package,
prints all 167 classes through five loops, their exact rational weights,
expanded edge lists and every monomial of $U_m$. The polynomials were
checked independently against the matrix-tree determinant at exact
integer edge parameters (record \path{verification/inventory_checks.json}). The identity
\begin{equation}
 \frac1{(L+1)!}\sum_m w_i(m)\mathcal I_m
 =\sum_{[m]}\frac{w_i(m)}{|\Aut_v(m)|}\mathcal I_m
 \label{gff:gd:orbitIdentity}
\end{equation}
follows from orbit--stabilizer and was also checked in the enumeration.

\subsection{An all-regulator two-loop master in ordinary hypergeometric functions}
There are only two two-loop vertex classes. The one-vertex-reducible
class has amplitude $(2/z)^2$; denote the other by $S_d(z)$. Their exact
weights give
\begin{equation}
 \mathcal A_2=\frac3{z^2}+3S_d(z),\qquad
 \mathcal B_2=\frac1{z^2}+\frac12S_d(z).
 \label{gff:gd:twoLoopAssembly}
\end{equation}
For the same $a=(d+z)/4$, put $b=(d+3z)/4$ and define
\begin{align}
 F_1(d,z)&={}_3F_2\!\left(
 \begin{matrix}d/2,z/2,a\\1-z/2,a+1/2\end{matrix};\frac14\right),\notag\\
 F_2(d,z)&={}_3F_2\!\left(
 \begin{matrix}z,2a,b\\1+z/2,b+1/2\end{matrix};\frac14\right),\notag\\
 W(d,z)&=\frac{r(d,z)^2\Gamma(1+z)\Gamma(1-z/2)\Gamma(b)^2}
 {\Gamma(a)^2\Gamma(d/2)\Gamma(2b)}.
 \label{gff:gd:twoLoopHyperDefs}
\end{align}
Then the full meromorphic two-loop amplitude is
\begin{equation}\boxed{
 S_d(z)=\frac{4F_1(d,z)-2W(d,z)F_2(d,z)}{z^2}.
}\label{gff:gd:twoLoopHyper}\end{equation}
The two defining ${}_3F_2$ series have argument $1/4$, strictly inside
the convergence disk. Near $z=0$ their parameters and the Gamma factors
in $W$ are regular. Thus \eqref{gff:gd:twoLoopHyper} determines all of the
two-loop Laurent data needed at fourth and fifth order, not just the
residue. The pole is explicit and the numerator is analytic.

\subsubsection{Derivation and a residue check}
Number the edges so that the Symanzik polynomial is
$U=(x_1+x_2)(x_3+x_4)+x_3x_4$. Set
$x_1+x_2=u$, $x_3=(1-u)v$, $x_4=(1-u)(1-v)$. The beta integral over
the split of $x_1,x_2$ and then over $u$ reduces the simplex to
\begin{equation}\begin{aligned}
 S_d(z)=&\frac{r(d,z)^2\Gamma(z)}{\Gamma(a)^4}
 B(a,a)B(2a,z/2)\\
 &\times\int_0^1[v(1-v)]^{a-1}
 {}_2F_1\!\left(\frac d2,\frac z2;\frac d2+z;
                       1-v+v^2\right)dv.
\end{aligned}\label{gff:gd:twoLoopEuler}\end{equation}
Apply the Gauss connection formula at $1-v(1-v)$ \cite{gff:DLMF15}.
Its two pieces have powers $[v(1-v)]^0$ and $[v(1-v)]^{z/2}$, respectively.
For either piece, termwise beta integration uses
\[
 \frac{B(c+n,c+n)}{B(c,c)}
  =\frac{(c)_n^2}{(2c)_{2n}}
  =\frac{(c)_n}{4^n(c+1/2)_n}.
\]
This gives exactly $F_1,F_2$. The first prefactor simplifies to $4/z^2$;
the second to $-2W/z^2$. The derivation starts where the beta integrals
converge and then continues meromorphically.

At $z=0$, $F_1=F_2=W=1$. The $n\geq1$ terms of the defining series
show $\partial_zF_2(d,0)=2\partial_zF_1(d,0)$, while
$\partial_z\log W(d,0)=-A_d/2$. It follows directly that
\begin{equation}
 4F_1-2WF_2=2+A_d z-s_1z^2+O(z^3),\qquad
 S_d(z)=\frac2{z^2}+\frac{A_d}{z}-s_1+O(z).
 \label{gff:gd:twoLoopResidue}
\end{equation}
Thus the known two-loop pole is recovered without evaluating the
intermediate hypergeometric derivatives separately. Derivatives appear
here only in the proof of a Taylor identity; the all-regulator formula
\eqref{gff:gd:twoLoopHyper} is an ordinary fixed-argument evaluation.

\subsection{Complete higher-order MS assembly}
For all the following quantities use the same prescription
\eqref{gff:gd:normalization}--\eqref{gff:gd:kernel}. Define
\begin{equation}
 a_{L,k}(d)=[z^k]\mathcal A_L(d,z),\qquad
 e_{L,k}(d)=[z^k]\mathcal B_L(d,z).
 \label{gff:gd:LaurentNotation}
\end{equation}
The letter $e$ here is a mass-amplitude Laurent coefficient, not an
external half-edge count $e_i$. The following finite recipe makes the
general-dimensional $C_4,C_5$ representation unambiguous.

Set
\begin{align}
 F_4(v)&=v+\sum_{L=1}^5(-1)^L\mathcal A_L v^{L+1},&
 F_2(v)&=1+\sum_{L=1}^5(-1)^L\mathcal B_Lv^L,\notag\\
 \mathcal Z_g&=g+\sum_{n=2}^{6}G_n(z)g^n,&
 Z_r&=1+\sum_{n=1}^{5}Z_n(z)g^n.
 \label{gff:gd:bareSeries}
\end{align}
All $G_n,Z_n$ are pure pole parts. Recursively impose
\begin{equation}\begin{aligned}
 G_n&=-\KP[g^n]F_4(\mathcal Z_{g,<n}),&
 Z_n&=-\KP[g^n]\{Z_{r,<n}F_2(\mathcal Z_g)\},\\
 b_L&=L[z^{-1}]G_{L+1},&c_L&=L[z^{-1}]Z_L.
\end{aligned}\label{gff:gd:MS}\end{equation}
Here $\KP$ retains negative powers of $z$. This is also an explicit
forest subtraction: multiplication by lower counterterms is carried
out before selecting a new simple pole. Keeping only primitive graphs
would not implement \eqref{gff:gd:MS}.

Write $A=A_d$, $R=R_d$, $M=M_d$, $Q=Q_d$ temporarily in the formulas
below. Expanding the recursion gives
\begin{align}
 b_4={}&36Aa_{2,0}-216a_{2,1}+48a_{3,0}-4a_{4,-1},\notag\\
 c_4={}&26Ae_{2,0}-140e_{2,1}+40e_{3,0}-4e_{4,-1},
 \label{gff:gd:bc4poles}\\
 b_5={}&5a_{5,-1}-75a_{4,0}+450a_{3,1}-60Aa_{3,0}
       -1350a_{2,2}+585Aa_{2,1}-5b_3a_{2,0},\notag\\
 c_5={}&5e_{5,-1}-65e_{4,0}+325e_{3,1}-\frac{95}{2}Ae_{3,0}
       -\frac{2275}{3}e_{2,2}\notag\\
 &\hspace{7mm}+\frac{725}{2}Ae_{2,1}
           -\frac{10b_3+5c_3}{3}e_{2,0}.
 \label{gff:gd:bc5poles}
\end{align}
Every $a_{L,k},e_{L,k}$ in these formulas is the coefficient of the
explicit finite integral sum \eqref{gff:gd:amplitudes}, with $U$ defined
in \eqref{gff:gd:U}; there are no unspecified diagram labels. The code
inventory is a convenience for enumerating this finite definition,
not a replacement for it.

For completeness, the required orders are
\begin{center}
\begin{tabular}{@{}p{0.12\textwidth}p{0.19\textwidth}p{0.61\textwidth}@{}}\toprule
Loop order & Highest Laurent power & Required analytic input\\\midrule
2 & $z^2$ & Full formula \eqref{gff:gd:twoLoopHyper}\\
3 & $z^1$ & Master expansions through the stated order\\
4 & $z^0$ & Master finite parts and the primitive period\\
5 & $z^{-1}$ & Subtracted simple poles and primitive periods\\\bottomrule
\end{tabular}
\end{center}
The combinations of these finite data required by the physical coefficients are reduced in Section~\ref{gff:app:generalSpecial}. They cannot be discarded merely because the primitive residues are known.

\subsubsection{The coefficient formulas through fifth order}
Lagrange inversion applied to
$\epsilon=g_*(3+b_2g_*+b_3g_*^2+b_4g_*^3+b_5g_*^4)$ gives
\begin{equation}
 C_n(d)=-\frac12\delta_{n1}
   +\frac1n[g^{n-1}]
      \frac{\gamma_r'(g)}{(3+b_2g+b_3g^2+b_4g^3+b_5g^4)^n}.
 \label{gff:gd:Lagrange}
\end{equation}
Define $D_4(d)=3c_4(d)-b_4(d)$ and $D_5(d)=3c_5(d)-b_5(d)$.
Substituting the evaluated \eqref{gff:gd:bc3} gives the collected general-$d$
formulas
\begin{equation}\boxed{
 243 C_4(d)=D_4-51A^3+9AR-48AQ-2AM,
}\label{gff:gd:C4}\end{equation}
\begin{equation}\boxed{\begin{aligned}
729 C_5(d)={}&D_5+10AD_4-6Ac_4
  -\frac{5541}{16}A^4-\frac{51}{4}A^2M-162A^2Q
  +\frac{315}{8}A^2R\\
 &+6MQ-\frac34MR+36Q^2-18QR+\frac{27}{16}R^2.
\end{aligned}}\label{gff:gd:C5}\end{equation}
Equations~\eqref{gff:gd:amplitudes}, \eqref{gff:gd:kernel}, and
\eqref{gff:gd:bc4poles}--\eqref{gff:gd:bc5poles} are part of these formulas:
$D_4,c_4,D_5$ are completely specified Laurent functionals, not
asserted new elementary constants. Nevertheless, leaving them as
integrals is a weaker endpoint than the finite hypergeometric
reduction obtained at $d=1$. In particular, a known value of $D_4$
would not by itself determine $c_4$ or the finite four-loop amplitudes
needed for $C_5$.

An equivalent form, identical to \eqref{gff:eq:C4RG} and \eqref{gff:eq:C5RG} and
convenient for checking signs, is
\begin{align}
243 C_4={}&60A^3-8Ab_3+18Ac_3+3c_4-b_4,\notag\\
729 C_5={}&336A^4-60A^2b_3+108A^2c_3+b_3^2-3b_3c_3
              -10Ab_4+24Ac_4+3c_5-b_5.
 \label{gff:gd:C45universal}
\end{align}
All terms in \eqref{gff:gd:C4}--\eqref{gff:gd:C5} have been checked by exact
symbolic algebra, including independent extraction of the simple-pole
coefficients in \eqref{gff:gd:bc4poles}--\eqref{gff:gd:bc5poles} from the RG pole
equations. Substituting the established one-dimensional RG data into
the universal assembly reproduces the retained $C_4(1),C_5(1)$ formulas.
That last check is an algebraic specialization, not a fresh numerical
evaluation of the new Schwinger integrals at $d=1$.

\subsubsection{The combined five-loop channel survives in arbitrary d}
The exact integrand weight used in the numerical resolution of the
one-dimensional problem was combinatorial:
\begin{equation}
 w_H=3w_2-w_4=\frac{w_4}{4}
                   \left[\sum_i e_i(e_i-1)-4\right].
 \label{gff:gd:combinedWeight}
\end{equation}
It therefore cancels the same 9540 of the 59145 labeled five-loop
matrices in every dimension. With lower-loop data held fixed,
\begin{equation}\boxed{
 \delta C_5(d)=\frac5{729}\,
  \delta[z^{-1}]\{3\mathcal B_5(d,z)-\mathcal A_5(d,z)\}.
}\label{gff:gd:sensitivity}\end{equation}
This does not remove the $a_{4,0},e_{4,0}$ terms. It identifies the
five-loop channel that should be integrated directly, before a
large numerical subtraction, when performing a general-$d$ evaluation.

\subsection{General-dimensional primitive periods}
The spectral construction extends the primitive periods, but
not by keeping the one-dimensional five hypergeometric parameter lists.
Let $P_L^{\rm raw}(d)$ be the conformal zig-zag decompletion with edge
exponents $q=d/4$, two fixed points at distance one and the remaining
$L-1$ vertices integrated over $\mathbb R^d$. Normalize
\begin{equation}
 P_L(d)=h_d^{L-1}P_L^{\rm raw}(d),\qquad h_d=\Gamma(d/2)/\pi^{d/2}.
 \label{gff:gd:primitiveNormalization}
\end{equation}
Thus $P_3(d)=3Q_d$, $P_4(1)=P_4$, and $P_5(1)=P_5$. The raw
coordinate integrals at integer dimension and their continued
spectral or parametric representations fix the graph definition.

Set $p=d/2-1$ and $k=L+2$. The eigenvalue and measure obtained from
\cite{gff:DIS}, eqs.~(5.23), (5.27), and (6.2), are
\begin{align}
 \tau_n^{(d)}(\nu)&=(-1)^n\pi^{d/2}
 \frac{\Gamma(n/2+d/8+\ii\nu)\Gamma(n/2+d/8-\ii\nu)}
      {\Gamma(n/2+3d/8+\ii\nu)\Gamma(n/2+3d/8-\ii\nu)},
 \label{gff:gd:tau}\\
 \rho_n^{(d)}(\nu)&=
 \frac{\Gamma(p)\Gamma(n+2p)(n+p)}
      {2\pi^{3d/2+1}\Gamma(2p)n!}
 \left[(n+p)^2+4\nu^2\right]
 \frac{\Gamma(p+2\ii\nu)\Gamma(p-2\ii\nu)}
      {\Gamma(2\ii\nu)\Gamma(-2\ii\nu)}.
 \label{gff:gd:rho}
\end{align}
For $2<d<4$ the ordinary real-axis representation is
\begin{equation}\boxed{
 P_L(d)=h_d^{L-1}\sum_{n=0}^{\infty}\int_0^{\infty}
        \rho_n^{(d)}(\nu)\,[\tau_n^{(d)}(\nu)]^{L+2}d\nu,
 \qquad L=3,4,5.
}\label{gff:gd:spectral}\end{equation}
The angular quantum number $n$ runs over an infinite tower. Its presence
is one reason the finite $d=1$ hypergeometric formula cannot be lifted
by changing quarter-arguments alone.

For $0<d<2$, $-1<p<0$ and the first pair of Gamma poles crosses the
real spectral contour. Continuing the contour as in the earlier
one-dimensional derivation adds
\begin{equation}\begin{aligned}
 P_L(d)=h_d^{L-1}\Bigg\{&
 \sum_{n=0}^{\infty}\int_0^{\infty}
       \rho_n^{(d)}(\nu)[\tau_n^{(d)}(\nu)]^{L+2}d\nu\\
 &+\sum_{n=1}^{\infty}
 \frac{\Gamma(n+2p)(n+p)n(n+2p)}
      {2\pi^{3d/2}n!\Gamma(-p)}
       [\tau_n^{(d)}(-\ii p/2)]^{L+2}\Bigg\}.
\end{aligned}\label{gff:gd:spectralContinued}\end{equation}
The $n=0$ crossing vanishes identically. At removable singular
parameter values such as $d=1,2$, take the joint analytic limit of the
Gamma ratios, not the separate factors. There are no additional
measure-pole crossings inside $0<d<2$; the eigenvalue Gamma arguments
for the $n\geq1$ residues stay away from their pole sets. The residue
coefficient is $-2\pi\ii\Res_{\nu=-\ii p/2}(\rho\tau^{L+2})$,
fixing its sign. At $d=1$, formula~\eqref{gff:gd:spectralContinued} reduces to the
continuous-plus-discrete formula \eqref{gff:sp:eq:spec} of
Section~\ref{gff:sec:P5complete}.
The numerical code
\path{lri_general_special_functions_update/code/primitive_spectral.py} (record
\path{verification/primitive_spectral_checks.json}) implements only \eqref{gff:gd:spectral} with
$d>2$; it does not silently substitute $d<2$ into that real-axis code.

As a normalization check, at $d=4$ one has
$\tau_n=(-1)^n4\pi^2/[(n+1)^2+4\nu^2]$. An elementary beta integration
then gives the graph-period limits
\begin{equation}
 P_3(4)=6\zeta(3),\qquad P_4(4)=20\zeta(5),\qquad
 P_5(4)=\frac{441}{8}\zeta(7).
 \label{gff:gd:period4D}
\end{equation}
Equivalently, $\mathcal T_4=12\zeta(3)$. These are checks of the
period normalization, not assertions that the complete long-range
RG functions have a regular short-range limit; $M_d$, for example,
has a pole at $d=4$.

\subsubsection{Primitive weights and the product periods}
The primitive classes at $L=3,4,5$ have vertex automorphism orders
$24$, $8$, and $(12,16,2)$ respectively. The four external legs are
on distinct vertices, so their mass weights vanish. At five loops,
the first two completed graphs have a three-vertex cut with two
components of two vertices each. Fix the cut vertices at $0,e,\infty$;
then the two remaining integrations decouple into two tetrahedron
integrals. Since every completed vertex has four propagators of
exponent $d/4$, the conformal normalization gives $P_3(d)^2=9Q_d^2$,
independently of $d$. The third class is the zig-zag $P_5(d)$.
Consequently the primitive pieces are
\begin{equation}\boxed{\begin{aligned}
 b_3^{\rm prim}(d)&=6Q_d,&
 b_4^{\rm prim}(d)&=-6P_4(d),\\
 b_5^{\rm prim}(d)&=63Q_d^2+24P_5(d),&
 c_L^{\rm prim}(d)&=0\quad(L=3,4,5).
\end{aligned}}\label{gff:gd:primitiveBeta}\end{equation}
Their direct new-loop contributions to $C_4,C_5$ are therefore
$2P_4(d)/81$ and $-[63Q_d^2+24P_5(d)]/729$. The lower-loop RG assembly
and nonprimitive graphs remain essential. The inventory file \path{graph_inventory.json} supplies
explicit completed edge lists and the cut vertices for the product
factorizations, rather than identifying them only by topology names.

\subsection{Numerical checks of the lower-order and primitive data}\label{gff:mgc:gdChecks}
These checks were run by OpenAI GPT-6 astra in the general-dimensional calculation (package
\path{lri_general_special_functions_update/}, scripts \path{tetrahedron.py},
\path{two_loop_master.py}, \path{primitive_spectral.py}, \path{symbolic_checks.py} and
\path{inventory_checks.py} in \path{code/}; records of the same names in \path{verification/}). They
are numerical consistency residuals, not interval enclosures.

\paragraph{Tetrahedron.}
A double-precision tensor trapezoidal quadrature of \eqref{gff:gd:tetraMB} used the contours
$\Re u=\Re v=-q/3$, spacings $0.015$ and $0.01$, and imaginary cutoffs $8$ and $10$. The finer run
differs from the exact values at $d=1,2,4$ (\eqref{gff:gd:d1blocks}, \eqref{gff:gd:d2exact} and
$\mathcal T_4=12\zeta(3)$) by about $3.65\times10^{-13}$, $6.43\times10^{-14}$ and
$5.63\times10^{-15}$. The resulting values, with $C_2(d)$ and $C_3(d)$ from \eqref{gff:gd:C03}, are
(record \path{C0_to_C3_values.csv})
\begin{center}
\begin{tabular}{@{}rrrr@{}}\toprule
$d$ & $\mathcal T_d$ & $C_2(d)$ & $C_3(d)$\\\midrule
1   &330.008919265&0.657131264&$-3.719856323$\\
1.5 &142.350694846&0.427213506&$-1.612227399$\\
2   &76.628287192 &0.308065414&$-0.874850013$\\
2.5 &46.354193318 &0.233416481&$-0.535556041$\\
3   &30.102615387 &0.181221785&$-0.353581520$\\
3.5 &20.493170714 &0.142032326&$-0.246028075$\\\bottomrule
\end{tabular}
\end{center}
The rows at $d=1,2,3$ agree with the table of Section~\ref{gff:sec:guide}. At $d=1.5$ and $2.5$ the record
values of $\mathcal T_d$ (the full double-precision values in \path{C0_to_C3_values.csv}) agree with
$6Q_d$ from the independent spectral code of \path{LRI_C5_general_d/references/reference_values.json}
to $1.0\times10^{-13}$ and $2.2\times10^{-13}$ absolute (comparison made for this note). The rounded numerical values
published with the corrected tetrahedron \cite{gff:BGH} are useful benchmarks, but their last printed
digits are not used as inputs.

\paragraph{Two-loop master.}
At 75 working digits the two-${}_3F_2$ formula \eqref{gff:gd:twoLoopHyper} was compared with the Euler
integral \eqref{gff:gd:twoLoopEuler} at $(d,z)=(1,0.4)$, $(2,0.3)$ and $(3,0.35)$. The differences were
$1.38\times10^{-44}$, below $2.2\times10^{-58}$ and below $4.5\times10^{-75}$. Cauchy extraction of the
analytic numerator of \eqref{gff:gd:twoLoopHyper}, with radii $0.025$ and $0.02$ and $72$ and $80$ nodes,
gives the Laurent data through $z^2$ at $d=1,2,3$; changing the settings changes the extracted
coefficients by less than $2.0\times10^{-69}$. At $d=1$ the finite parts in this massive-Schwinger
prescription are $[z^0]\mathcal A_2\simeq-55.190665$ and $[z^0]\mathcal B_2\simeq-9.198444$. They
differ from the ordered-span finite parts of Section~\ref{gff:sec:graphs}, as they must; the two sets are
not interchangeable (Section~\ref{gff:app:generalDimension}).

\paragraph{Primitive spectral formula.}
Gauss--Legendre quadrature of \eqref{gff:gd:spectral} after a map of the spectral half-line, with $192$
and $320$ nodes and Euler acceleration of the alternating angular sums, gives at $d=3$
$P_3(3)\simeq15.05130769366$, equal to $\mathcal T_3/2=3Q_3$ of the separate Mellin--Barnes
calculation, and $P_5(3)\simeq240.34031776$. The $P_4(3)$ run is an unextrapolated positive partial
sum: $62.0583263456$, $62.0583277437$, $62.0583279184$ and $62.0583279403$ at $128$, $256$, $512$ and
$1024$ angular modes, with no bound claimed for the omitted tail. All three agree with the
high-precision values of Section~\ref{gff:d3:values} to the accuracy of the quadrature. The same code at
$d=4$ reproduces the four-dimensional periods $P_3=6\zeta(3)$ and $P_5=\frac{441}8\zeta(7)$ within
$4.3\times10^{-13}$ and $3.4\times10^{-12}$. None of these is a determination of the complete $C_4(3)$
or $C_5(3)$.

\paragraph{Exact algebra and graph data.}
The symbolic check reconstructs the MS three-loop coefficients \eqref{gff:gd:bc3} from the zero-momentum
ones, shows the cancellation of $J_0$, expands the Lagrange inversion through fifth order, derives
\eqref{gff:gd:bc4poles}--\eqref{gff:gd:bc5poles} independently from the pole equations, and recovers the
one-dimensional coefficient formulas by substitution. The inventory check verifies all 167 Kirchhoff
polynomials by exact matrix-tree arithmetic, the orbit factors, the primitive mass weights, the two
product factorizations of \eqref{gff:gd:primitiveBeta} and the 9540 zero combined weights of
\eqref{gff:gd:combinedWeight}. No general-dimensional nonprimitive four- or five-loop integral is
evaluated numerically in these checks.


\section{General-dimensional master reduction and collected coefficients}\label{gff:app:generalSpecial}

\subsection{Conventions and the collected answer}
Here $d$ is the physical spatial dimension, not the exponent-shift
parameter $d$ of Section~\ref{gff:sec:shifts}. Set
\begin{equation}
 q=\frac d4,\qquad \epsilon=2s-d,\qquad
 \Delta_\phi=\frac{d-\epsilon}{4},\qquad
 \Delta_{\mathcal E}=\sum_{j=0}^5 C_j(d)\epsilon^j+O(\epsilon^6).
 \label{gff:gs:conventions}
\end{equation}
The physical perturbative range is fixed $0<d<4$ and sufficiently small
positive $\epsilon$. The endpoint $d=4$ is not substituted into
unrenormalized power-subgraph formulas. Define
\begin{equation}\begin{aligned}
 A&=\psi_0(2q)+\psi_0(1)-2\psi_0(q),&
 R&=\psi_1(2q)-\psi_1(1),\\
 U&=\psi_2(q),& V&=\psi_2(1)+\psi_2(2q),&
 X&=\psi_3(2q)-\psi_3(1).
\end{aligned}\label{gff:gs:psiBasis}\end{equation}
Here and below $Q=Q_d$, $P_4=P_4(d)$ and $P_5=P_5(d)$ have the
normalization of \eqref{gff:gd:primitiveNormalization}: $P_3=3Q$.
The three additional dimension-dependent functions are denoted by
$\mathcal K_d,\mathcal L_d,Z_d$. Explicit series definitions are
\eqref{gff:gs:Lseries}, \eqref{gff:gs:Zseries} and \eqref{gff:gs:Kseries}.
They are not free parameters. Sections~\ref{gff:gdr:section} and
\ref{gff:gdk:section} give shorter finite-sample evaluations of all three, and
\eqref{gff:gdr:C5reduced} is \eqref{gff:gs:C5} written in them.

The complete collected coefficients are
\begin{equation}\boxed{\begin{aligned}
 C_0(d)&=\frac d2,&C_1(d)&=-\frac16,&C_2(d)&=\frac A9,\\
 C_3(d)&=\frac{A^2+R-8Q}{108},\\
 648C_4(d)&=16P_4-A^3-8AQ-3AR+\frac{53}{2}U-3V.
\end{aligned}}\label{gff:gs:C0C4}\end{equation}
\begin{equation}\boxed{\begin{aligned}
46656 C_5(d)={}&-1536\bigl(P_5+\mathcal K_d\bigr)
       +384AP_4-576\mathcal L_d-1728Q^2\\
 &+720A^2Q+1008QR+7152Z_d\\
 &-15A^4-90A^2R+414AU-108AV-117R^2-75X.
\end{aligned}}\label{gff:gs:C5}\end{equation}
In particular, the extra Gamma factor associated with power-divergent
mass subgraphs cancels from the explicit coefficients of both physical
answers. It must nevertheless be retained until the graph and
fixed-point assembly is complete.

\paragraph{Exact one-dimensional specialization.}
At $d=1$, the dictionary is
\begin{equation}\begin{aligned}
 R&=\frac{\pi^2}{3},&U&=-2\pi^3-56\zeta(3),&V&=-16\zeta(3),\\
 X&=\frac{14\pi^4}{15},&Z_1&=QH,&
 H&=\frac{2\pi^2}{3}+8\Cat,\\
 \mathcal K_1&=\mathcal K,&\mathcal L_1&=\mathcal L.
\end{aligned}\label{gff:gs:d1Dictionary}\end{equation}
Substitution into \eqref{gff:gs:C0C4}--\eqref{gff:gs:C5} reproduces the
one-dimensional coefficients of Section~\ref{gff:cn:oneDresults} (for $C_5$,
\eqref{gff:eq:full}) exactly. In particular, the $7152Z_d$
term becomes $48\cdot149QH$. Neither $Z_d=Q_dH_d$ with a guessed
$H_d$, nor the one-dimensional polygamma formula for $\mathcal K$, is
assumed away from $d=1$.

\subsection{The evaluated RG combinations}
For comparison with the assembly formulas
\eqref{gff:gd:C4}--\eqref{gff:gd:C45universal}, define
\begin{equation}\begin{aligned}
 M&=\frac{\Gamma(-q)\Gamma(2q)^2}{\Gamma(3q)}
   =-\frac{\Gamma(1-q)\Gamma(2q)^2}{q\Gamma(3q)},\\
 N&=\psi_0(-q)+\psi_0(3q)-\psi_0(1)-\psi_0(2q),\qquad
 D_j=3c_j-b_j.
\end{aligned}\label{gff:gs:MN}\end{equation}
For $0<q<1$, the last Gamma representation has only positive arguments.
The recurrence $\psi_0(-q)=\psi_0(1-q)+1/q$ likewise avoids a negative
argument. The exact fourth-order results are
\begin{align}
 b_4={}&-\frac{393}{4}A^3-\frac{11}{2}AM-54AQ+\frac{117}{4}AR
       -\frac32MN-6P_4-\frac{165}{8}U+\frac94V,
 \label{gff:gs:b4}\\
 c_4={}&-\frac{127}{8}A^3-\frac76AM-3AQ+\frac{51}{8}AR
       -\frac12MN-\frac{57}{16}U+\frac38V,
 \label{gff:gs:c4}\\
 D_4={}&\frac{405}{8}A^3+2AM+45AQ-\frac{81}{8}AR
       +6P_4+\frac{159}{16}U-\frac98V.
 \label{gff:gs:D4}
\end{align}
Separate the fifth-order nonprimitive nonwheel part from its wheel and
primitive contributions:
\begin{equation}\begin{aligned}
 D_5^{\rm rem}={}&-\frac{16347}{64}A^4-\frac{57}{4}A^2M
 -\frac{1179}{4}A^2Q+\frac{3159}{32}A^2R-3AMN\\
 &-\frac{3657}{32}AU+\frac{189}{16}AV-24\mathcal K_d
 -6MQ+\frac34MR+\frac{135}{4}QR\\
 &-\frac{225}{64}R^2-\frac{75}{64}X+\frac{447}{4}Z_d,\\
 D_5={}&D_5^{\rm rem}-54AP_4-9\mathcal L_d-63Q^2-24P_5.
\end{aligned}\label{gff:gs:D5}\end{equation}
These expressions determine the three combinations $D_4$, $c_4$ and
$D_5$ left open in \eqref{gff:gd:C4}--\eqref{gff:gd:C5}. Inserting them into \eqref{gff:gd:C45universal}, together
with \eqref{gff:gd:bc3}, gives \eqref{gff:gs:C0C4}--\eqref{gff:gs:C5}; the $M$, $MN$
and $MR$ terms cancel. The script
\path{lri_general_special_functions_update/code/general_special/exact_audit.py}
checks this substitution in exact arithmetic. It computes $b_3$ and $c_3$
from the three-loop graph rows instead of taking \eqref{gff:gd:bc3} as input.

\subsection{Why the graph reduction extends to general dimension}\label{gff:gs:patterns}
Normalize every coordinate integration with
$h_d=\Gamma(d/2)/\pi^{d/2}$. The two-point convolution coefficient is
\begin{equation}
 B_d(a,b)=\Gamma(d/2)
 \frac{\Gamma((d-a)/2)\Gamma((d-b)/2)\Gamma((a+b-d)/2)}
 {\Gamma(a/2)\Gamma(b/2)\Gamma((2d-a-b)/2)}.
 \label{gff:gs:convolution}
\end{equation}
It multiplies $|x-y|^{-(a+b-d)}$. It is symmetric under permutations of
$a,b,2d-a-b$. The star--triangle relation applies when its three absolute-distance
exponents sum to $2d$. These are dimension-dependent Gamma identities,
not replacements of special one-dimensional constants \cite{gff:DIS}.

A reduced two-vertex master is
\begin{equation}
 J_d(a,b,c,f,e)=h_d^2\int_{\mathbb R^d}\!d^dx\int_{\mathbb R^d}\!d^dy\,
 \frac{1}{|x|^a|x-\hat e|^b|y|^c|y-\hat e|^f|x-y|^e},
 \qquad |\hat e|=1.
 \label{gff:gs:masterJ}
\end{equation}
Integrals outside their ordinary convergence domain are defined by
analytic continuation and the critical-mass prescription of \cite{gff:BGHS}
used in Section~\ref{gff:app:generalDimension}.
Its conformal completion is a balanced $K_5$ with ordered edges
$(01,02,03,04,12,13,14,23,24,34)$ and exponents
\begin{equation}\begin{split}
\boldsymbol\lambda={}&(3d-a-b-c-f-e,\ a,\ c,\ b+f+e-d,\ b,\ f,\\
 &a+c+e-d,\ e,\ 2d-a-b-e,\ 2d-c-f-e).
\end{split}\label{gff:gs:completion}\end{equation}
Each completed vertex has incident exponent sum $2d$; permutation of
its five vertices gives the exact completion covariance. The needed
Taylor data reduce to three critical patterns.

\subsubsection{Pattern A: the Gamma-determined jet}
For the critical point $(d/2,d/2,d/2,d/2,d)$ set
\begin{equation}\begin{aligned}
 t&=d-e,&p_0&=2d-a-c-e,&p_1&=2d-b-f-e,\\
 u&=(a-c)/2,&v&=(b-f)/2,&\sigma&=p_0+p_1.
\end{aligned}\label{gff:gs:Acoords}\end{equation}
An explicit pole subtraction is
\begin{equation}\begin{aligned}
 \mathcal P_A={}&\frac2{p_0}B_d\left(\frac{d+t}{2}-u,d-t\right)
 +\frac2{p_1}B_d\left(\frac{d+t}{2}-v,d-t\right)\\
 &+\frac2tB_d(d-p_0,d-p_1)-\frac4t\left(\frac1{p_0}+\frac1{p_1}\right).
\end{aligned}\label{gff:gs:Apole}\end{equation}
The analytic remainder through total degree two is
\begin{equation}\boxed{\begin{aligned}
J_d-\mathcal P_A={}&\sigma\left(\frac{AR}{2}+\frac U8\right)
+t\sigma\left(\frac{A^2R}{8}+\frac{AU}{16}-\frac{R^2}{8}\right)
-\frac{AV}{8}\sigma^2+O(3).
\end{aligned}}\label{gff:gs:Ajet}\end{equation}
This is fixed without evaluating a new master. On each of
$a+b+e=2d$, $c+f+e=2d$, and $a+b+c+f+e=3d$, a completion edge vanishes
and star--triangle reduction gives a product of two copies of
\eqref{gff:gs:convolution}. After \eqref{gff:gs:Apole}, the remainder is analytic.
An analytic difference vanishing on these three distinct hyperplanes is
divisible by their product. It consequently has no term of total degree
at most two. Expanding the Gamma products therefore determines
\eqref{gff:gs:Ajet} uniquely. The exact interpolation in
\path{lri_general_special_functions_update/code/general_special/derive_A.py}
uses rational directions only; it is finite polynomial linear algebra, not
numerical fitting.

\subsubsection{Pattern B: one genuinely independent quadratic coefficient}
At $(d/2,d/2,d/2,d,d/2)$ put the following local coordinates; $p_1$ here is
not the $p_1$ of \eqref{gff:gs:Acoords}, which equals $p_2$ here:
$p_1=d-f$, $p_2=2d-b-f-e$, $z_1=3d-a-b-c-f-e$,
$z_2=2d-c-f-e$. Define
\begin{align}
 G_2={}&B_d(d+p_1-e,d-p_1)B_d(e+p_2-z_1,d-p_2),\notag\\
 G_1={}&B_d(e-p_1+p_2+z_2,d+p_1-e-z_2)
              B_d(d+p_1-e-p_2,d-p_1),\notag\\
 G_0={}&B_d(d+p_1-e,d-p_1)B_d(e+p_2,d-p_2),\notag\\
 \Delta_B={}&B_d(e+p_2+z_2-z_1,d-p_2)-B_d(e+p_2-z_1,d-p_2)\notag\\
 &-B_d(e+p_2+z_2,d-p_2)+B_d(e+p_2,d-p_2).
 \label{gff:gs:Bproducts}
\end{align}
The needed jet is
\begin{equation}
 J_d=G_2+G_1-G_0+\frac2{p_1}\Delta_B+\mathcal K_d z_1z_2+O(3).
 \label{gff:gs:Bjet}
\end{equation}
The Gamma comparison agrees on $z_1=0$ and $z_2=0$. Completion symmetry
and analyticity make the degree-two difference proportional to $z_1z_2$.
Its normalization is fixed by the one-parameter path below. Its parameter
$t$ is not the $t=d-e$ of \eqref{gff:gs:Acoords}:
\begin{equation}\begin{aligned}
 J_t&=J_d(d/2,d/2,d/2,d-t,d/2),\\
 J_t^\Gamma&=B_d(d/2,d-t)^2+\frac2t
   \left[B_d(d/2,d-t)-B_d(d/2+2t,d-t)\right],\\
 \mathcal K_d&=[t^2](J_t-J_t^\Gamma).
\end{aligned}\label{gff:gs:Kdefinition}\end{equation}
This coefficient extraction is a \emph{proof and normalization definition};
the final evaluation \eqref{gff:gs:Kseries} has neither $t$ nor a derivative.
At $d=1$ it is exactly the one-dimensional master $\mathcal K$ of \eqref{gff:eq:Bjet}, not an unrelated
constant with a matching numerical value.

\subsubsection{Pattern T: tetrahedron symmetry}
At $a=b=c=f=e=d/2$, all ten completion edges equal $d/2$. Balanced
variations satisfy $\sum_{e\ni v}\delta\lambda_e=0$ at every vertex. Here
$e$ labels the ten completion edges; it is not the exponent $e$ of
\eqref{gff:gs:masterJ}.
There is no invariant linear form, and every invariant quadratic form
on this balanced subspace is proportional to $\sum_e(\delta\lambda_e)^2$.
Indeed the coefficients of a general permutation-invariant quadratic
form depend only on whether two edges coincide, meet, or are disjoint;
the balance equations reduce the latter two sums to the first.
Thus
\begin{equation}\boxed{
 J_d=3Q_d+\frac{Z_d}{4}\sum_{e=1}^{10}(\delta\lambda_e)^2+O(3).
}\label{gff:gs:Tjet}\end{equation}
This argument applies in every dimension; the coefficient $Z_d$ itself
is not fixed by symmetry. A central-edge-only change by $t$ has squared
completion variation $6t^2$. A balanced three-wheel path has $18t^2$,
which fixes the factors of three and nine used below.

\subsubsection{Forest sum and six wheel rows}
The graph set, exact Wick weights and subtraction prescription are those
already specified in \eqref{gff:gd:amplitudes}--\eqref{gff:gd:kernel}. Equivalently,
a rooted logarithmic $L$-loop coordinate graph has common radial factor
$2/(Lz)$ and edge exponents $m_{ij}(d-z)/2$. Its proper disjoint forests
are subtracted recursively; nested forests occur inside the subgraph
counterterms. The local counterterm is
\[
 Z_G=-\operatorname{PP}\left[I_G+
 \sum_{\varnothing\ne\mathcal F}\left(\prod_{\gamma\in\mathcal F}Z_\gamma\right)
 I_{G/\mathcal F}\right].
\]
The shift method of Section~\ref{gff:sec:shifts}, used on exceptional graphs, preserves the sum
of edge deformations in every divergent subgraph. Only convolution,
star--triangle, completion covariance and the jets above enter.
Consequently its combinatorics is dimension independent. Critical base
exponents scale with $d$; regulator slopes do \emph{not} scale with $d$.

The residue certificate for these classes consists of the files
\path{rows_4_0_27.json} and \path{rows_5_0_129.json} in
\path{lri_general_special_functions_update/code/general_special/}. They list
all 26 nonprimitive four-loop classes and the 92 nonzero five-loop classes
that are neither wheels nor primitive. At five loops six further classes
are wheels, three are primitive, and 28 have zero combined weight. The sum
is 129 classes. For each class the files give the canonical word, the
rational weights and the exact symbolic residue.
They give \eqref{gff:gs:b4}--\eqref{gff:gs:D5}; no period value is used to obtain
these polynomial coefficients.

For completeness the six wheel residues are
$-2AP_4/5+(\alpha D_s+\beta D_r)/5$, where $D_s=\frac14\partial_aP_W$ and
$D_r=\frac14\partial_bP_W$ are the spoke and rim derivatives at $a=b=d/2$ of the
rooted four-wheel $P_W$ of Section~\ref{gff:gs:logWheel}. The class index is the
zero-based position in the five-loop list of
\path{lri_general_special_functions_update/code/general_special/C5_wheel_nonprimitive/reduced_classes.json}.
The data are:
\begin{center}
\begin{tabular}{@{}rrrr@{}}\toprule
Class index & combined weight & $\alpha$ & $\beta$\\\midrule
51&$15/2$&2&2\\78&15&2&$-2$\\99&30&2&$-2$\\
113&60&$-2$&2\\124&$15/2$&6&$-2$\\128&15&6&$-6$\\\bottomrule
\end{tabular}
\end{center}
Their sum is $-54AP_4+24D_s-12D_r=-54AP_4-9\mathcal L_d$.
The completion/product argument behind \eqref{gff:gd:primitiveBeta} gives
$D_5^{\rm prim}=-63Q^2-24P_5$.

\subsection{A common spectral kernel for all nonprimitive master data}
\label{gff:gs:spectralSection}
Let
\begin{equation}\begin{aligned}
 d_0(d)&=1,&
 d_\ell(d)&=\frac{(2\ell+d-2)\Gamma(\ell+d-2)}
 {\Gamma(d-1)\Gamma(\ell+1)}\quad(\ell\geq1),\\
 a_\ell&=\frac\ell2+\frac q2,&b_\ell&=a_\ell+q,\\
 t_\ell(\nu)&=\Gamma(2q)
 \frac{\Gamma(a_\ell+\ii\nu)\Gamma(a_\ell-\ii\nu)}
      {\Gamma(b_\ell+\ii\nu)\Gamma(b_\ell-\ii\nu)}.
\end{aligned}\label{gff:gs:spectrum}\end{equation}
Here $a_\ell,b_\ell$ are spin-dependent Gamma arguments, not the RG
coefficients $b_j$ of \eqref{gff:gd:RG}. At $d=1$, $(d_0,d_1,d_2,\ldots)=(1,1,0,\ldots)$;
at $d=2$ the multiplicities are $(1,2,2,\ldots)$.
At noninteger $d<1$ they are continued Gamma ratios, not positive
integer multiplicities. For notational economy write
\begin{equation}
 \langle f\rangle_d=\frac1\pi\sum_{\ell=0}^{\infty}d_\ell(d)
                    \int_0^{\infty}f(\ell,\nu)\,d\nu.
 \label{gff:gs:trace}
\end{equation}
This is the \emph{one-vector wheel kernel}, not the two-vector zig-zag
measure in \eqref{gff:gd:rho}. The primitive $P_5(d)$ keeps the latter
representation and its stated continuation prescription.

For a balanced wheel with rim exponent $b$ and spoke exponent $d-b$,
the commuting $H_\alpha$ construction of \cite{gff:DIS} gives eigenvalue
\begin{equation}
 t_{\ell,b}(\nu)=\Gamma(d/2)\frac{\Gamma((d-b)/2)}{\Gamma(b/2)}
 \left|\frac{\Gamma(\ell/2+b/4+\ii\nu)}
 {\Gamma(\ell/2+d/2-b/4+\ii\nu)}\right|^2,
 \qquad W_r(d;b)=\langle t_{\ell,b}^{\,r}\rangle_d.
 \label{gff:gs:balancedWheel}
\end{equation}
One can also obtain it directly by applying the normalized convolution
operator $h_d|x-y|^{-b}|y|^{b-d}$ to homogeneous harmonic functions.
The spherical-harmonic addition theorem gives the factor
$d_\ell/(2\pi)$ on the full Mellin line, hence $d_\ell/\pi$ on the
positive half-line. In particular
\begin{equation}
 3Q_d=\langle t_\ell^3\rangle_d,\qquad
 P_4(d)=\langle t_\ell^4\rangle_d.
 \label{gff:gs:primitiveWheel}
\end{equation}
At the graph-period check point $d=4$, $t_\ell=(\nu^2+(\ell+1)^2/4)^{-1}$
and $d_\ell=(\ell+1)^2$, giving $W_3=6\zeta(3)$ and $W_4=20\zeta(5)$.

\subsubsection{The logarithmic wheel and the tetrahedron Hessian}
\label{gff:gs:logWheel}
Let a rooted four-wheel have four spokes of exponent $a$ and four rims
of exponent $b$, and let $P_W(a,b)$ be its period, normalized as in
\eqref{gff:gs:balancedWheel}, so that $P_W(d-b,b)=W_4(d;b)$. At $a=b=d/2$ its completion is the octahedron. The
completed exponents are four $a$'s, four $b$'s, two
$-2d+3a+2b$'s, two $2d-a-2b$'s, and the three absent-edge exponents
$2d-4a$, $4d-4a-4b$, $0$. Their vertex sums equal $2d$.
The octahedral symmetry has two edge orbits: 12 present and three absent.
The total variation is zero, while the sum over the present edges is
$8\,\delta a+4\,\delta b$. Its invariant first derivative therefore
obeys $\partial_aP_W=2\partial_bP_W$, or $D_s=2D_r$ when each derivative
is divided by four. It follows that
\begin{equation}
 \mathcal L_d=\frac43(D_r-2D_s)=-4D_r
   =\left.\partial_bW_4(d;b)\right|_{b=d/2},\qquad
 Z_d=\left.\frac19\partial_b^2W_3(d;b)\right|_{b=d/2}.
 \label{gff:gs:wheelJets}
\end{equation}
These derivatives are used only in the proof. For direct evaluation put
\begin{equation}\begin{aligned}
 w_1&=-\psi_0(q)+\frac12\Re\{\psi_0(a_\ell+\ii\nu)+\psi_0(b_\ell+\ii\nu)\},\\
 w_2&=\frac18\Re\{\psi_1(a_\ell+\ii\nu)-\psi_1(b_\ell+\ii\nu)\}.
\end{aligned}\label{gff:gs:w12}\end{equation}
Then
\begin{equation}\boxed{
 \mathcal L_d=\langle4t_\ell^4w_1\rangle_d,\qquad
 Z_d=\left\langle t_\ell^3\left(w_1^2+\frac{w_2}{3}\right)\right\rangle_d.
}\label{gff:gs:LZspectral}\end{equation}
The first derivative $\langle3t_\ell^3w_1\rangle_d$ vanishes identically;
it is a useful normalization and tail-control check.

\subsubsection{A convergent single-kernel formula for the subtracted master}\label{gff:gs:singleKernel}
Set $c_*=\psi_0(1)+\psi_0(2q)$ and, at $a=a_\ell,b=b_\ell$, define the
local abbreviations below. Here $h$ is not $h_d$, $j$ and $k$ are not
indices, and $r$ is not the wheel order:
\begin{equation}\begin{aligned}
 p&=\Re\{\psi_0(a+\ii\nu)+\psi_0(b+\ii\nu)\},&
 h&=\Im\{\psi_0(a+\ii\nu)-\psi_0(b+\ii\nu)\},\\
 j&=\Re\{\psi_1(a+\ii\nu)-\psi_1(b+\ii\nu)\},&
 k&=\Im\{\psi_1(a+\ii\nu)+\psi_1(b+\ii\nu)\},\\
 r&=\Re\{\psi_2(a+\ii\nu)+\psi_2(b+\ii\nu)\},&
 E_2&=-\frac{h^2}{2}-\frac{5j}{8}.
\end{aligned}\label{gff:gs:Kweights}\end{equation}
The complete subtraction collapses to
\begin{equation}\boxed{
 \mathcal K_d=-\frac{AU}{4}
 +\frac14\left\langle t_\ell^2
   \left[(c_*-p)E_2-\frac{hk}{2}-\frac{5r}{32}\right]\right\rangle_d.
}\label{gff:gs:Kspectral}\end{equation}
Unlike the unsubtracted $J_t$ of \eqref{gff:gs:Kdefinition}, this integral has no ultraviolet
pole. Its large-spin summand is $O(\ell^{-3}\log\ell)$.

Here is a derivation that fixes its finite part. Write $t=2\tau$ in
\eqref{gff:gs:Kdefinition} and $s=\ii\nu$. Here and below $s$ is the Mellin
variable, not the long-range exponent $s$ of \eqref{gff:gs:conventions}. The mixed trace of two critical
wheel operators and one operator of rim exponent $d-2\tau$ is
$J_{2\tau}=C(\tau)\langle t_\ell^2R_\tau\rangle_d$, where (the letters are
local: $C(\tau)$ is not a coefficient $C_j$, $R_\tau$ is not $R$ of
\eqref{gff:gs:psiBasis}, and $S_\tau$ is not $S_{\ell n}$ of \eqref{gff:gs:seriesBase})
\begin{equation}\begin{aligned}
 C(\tau)&=\frac{\Gamma(2q)\Gamma(\tau)}{\Gamma(2q-\tau)},\\
 R_\tau(s)&=\frac{\Gamma(a+s)\Gamma(b-\tau-s)}
                 {\Gamma(a+\tau+s)\Gamma(b-s)},\\
 S_\tau(s)&=\left[
 \frac{\Gamma(a-\tau/4+s)\Gamma(a-\tau/4-s)\Gamma(b+s)\Gamma(b-s)}
 {\Gamma(b+\tau/4+s)\Gamma(b+\tau/4-s)\Gamma(a+s)\Gamma(a-s)}
 \right]^2.
\end{aligned}\label{gff:gs:RS}\end{equation}
At $d=2$ this mixed-trace step is proved, for $0<\tau<\frac12$ and in its
full-Mellin-line reading, in Section~\ref{gff:mgk:regulated}. Its reading for
complex $\tau$ and its status at other $d$ are stated as input (M) in
Section~\ref{gff:gdk:conventions}.
The two-wheel convolution is elementary:
$W_2(d;b')=B_d(d-b',2b')$. Taking $b'=d/2-\tau$ gives the exact identity
\begin{equation}
 J_{2\tau}=C(\tau)^2
 \left[\frac{\Gamma(q-\tau/2)}{\Gamma(q+\tau/2)}\right]^4
 +C(\tau)\langle t_\ell^2(R_\tau-S_\tau)\rangle_d.
 \label{gff:gs:Ksubtraction}
\end{equation}
On the Mellin line, the real difference starts at order $\tau^2$:
\[
 \Re(R_\tau-S_\tau)
 =E_2\tau^2+\left[-pE_2-\frac{hk}{2}-\frac{5r}{32}\right]\tau^3+O(\tau^4).
\]
Using $C(\tau)=\tau^{-1}+c_*+O(\tau)$ and $t=2\tau$ yields the integral
in \eqref{gff:gs:Kspectral}. The coefficient of $t^2$ in the explicit Gamma
term of \eqref{gff:gs:Ksubtraction} is
\[
 \frac{A^4}{6}-\frac{A^2R}{2}+\frac{R^2}{8}
 +\frac{A(2V-U)}{12}-\frac X{48},
\]
whereas that of $J_t^\Gamma$ is
\[
 \frac{A^4}{6}-\frac{A^2R}{2}+\frac{R^2}{8}
 +\frac{A(V+U)}6-\frac X{48}.
\]
Their difference is exactly $-AU/4$. This derives
\eqref{gff:gs:Kspectral} from the master definition \eqref{gff:gs:Kdefinition}, rather than
from a numerical relation search, given the mixed-trace step and the
extraction of $[\tau^2]$ under the trace. Both are proved at $d=2$
(Sections~\ref{gff:mgk:regularization} and \ref{gff:mgk:regulated}). At other $d$ the first is input (M) of
Section~\ref{gff:gdk:conventions}, and the second is an unproved interchange of
$[\tau^2]$ with the trace (Section~\ref{gff:gdk:status}); both are confirmed
numerically in Section~\ref{gff:gdk:numerics}.

\subsection{Explicit polygamma--hypergeometric series: no integrals or parameter derivatives}
\label{gff:gs:seriesSection}
All quantities needed for the nonprimitive terms can now be written
using positive-argument Gamma and polygamma functions and integer sums.
For $\ell,n\geq0$ define
\begin{equation}\begin{aligned}
 T_{\ell n}&=\frac{(q)_{\ell+n}(1-q)_n}{(2q)_{\ell+n}\,n!},&
 S_{\ell n}&=\frac{(2q)_{\ell+n}\,n!}{(3q)_{\ell+n}(1+q)_n},\\
 a_{\ell n}&=\ell+n+q,& b_{\ell n}&=\ell+n+2q,\\
 \chi_m(n;q)&=(-1)^m\psi_m(n+1-q)+\psi_m(q)-(-1)^m\psi_m(1-q).
\end{aligned}\label{gff:gs:seriesBase}\end{equation}
In terms of \eqref{gff:gs:spectrum}, $a_{\ell n}=2a_\ell+n$ and
$b_{\ell n}=a_\ell+b_\ell+n$; they are not $a_\ell,b_\ell$.
Reflection gives $\chi_m(n;q)=\psi_m(q-n)$, but the positive-argument
expression is preferable for evaluating it. Thus the natural pair
$q=d/4$ and $1-q$ remains explicit; no trigonometric replacement is
required.

For $m=1,2,3,4$ set (these $\lambda_m$ are not the completion exponents
$\lambda_e$ of \eqref{gff:gs:completion})
\begin{equation}\begin{split}
 \lambda_m={}&\psi_{m-1}(a_{\ell n})-\psi_{m-1}(b_{\ell n})
 -(-1)^m\chi_{m-1}(n;q)\\
 &+(1+(-1)^m)\psi_{m-1}(1)-\psi_{m-1}(n+1).
\end{split}\label{gff:gs:lambda}\end{equation}
The following are explicit polygamma combinations, not differentiation
operators:
\begin{equation}\begin{aligned}
 \eta_m&=\frac12\{\psi_m(a_{\ell n})+(-1)^m\chi_m(n;q)\},
                       &&m=1,2,3,\\
 \theta_m&=\frac14\{\psi_{m+1}(a_{\ell n})-(-1)^m\chi_{m+1}(n;q)\},
                       &&m=1,2,\\
 \mu_1&=\frac12\{\psi_0(a_{\ell n})-\psi_0(q)
                       +\psi_0(n+1-q)-\psi_0(1-q)\},\\
 \mu_2&=\frac14\{\psi_1(a_{\ell n})-\psi_1(q)
                       +\psi_1(n+1-q)-\psi_1(1-q)\}.
\end{aligned}\label{gff:gs:etaTheta}\end{equation}
All subscripts $\ell,n$ on the right are suppressed in the following
polynomials:
\begin{equation}\begin{aligned}
 F_3&=\frac92\lambda_1^2+\frac32\lambda_2,\\
 F_{3,1}&=9\lambda_1\eta_1+\frac32\eta_2,\\
 F_{3,2}&=9\eta_1^2+9\lambda_1\theta_1+\frac32\theta_2,\\
 F_4&=\frac{32}{3}\lambda_1^3+8\lambda_1\lambda_2+\frac23\lambda_3,\\
 F_{4,1}&=32\lambda_1^2\eta_1+8\eta_1\lambda_2
                         +8\lambda_1\eta_2+\frac23\eta_3.
\end{aligned}\label{gff:gs:Fpolys}\end{equation}
The subscripts $1,2$ in these labels name the displayed polynomials;
they are not a request to differentiate a special function.

The primitive wheel periods and the two corresponding nonprimitive
master data are
\begin{equation}\boxed{\begin{aligned}
 Q_d&=\frac13\sum_{\ell=0}^{\infty}d_\ell(d)
                   \sum_{n=0}^{\infty}T_{\ell n}^3F_3,\\
 P_4(d)&=-\sum_{\ell=0}^{\infty}d_\ell(d)
                   \sum_{n=0}^{\infty}T_{\ell n}^4F_4.
\end{aligned}}\label{gff:gs:QPseries}\end{equation}
\begin{equation}\boxed{
 \mathcal L_d=-\sum_{\ell=0}^{\infty}d_\ell(d)\sum_{n=0}^{\infty}
 T_{\ell n}^4\bigl(4\mu_1F_4+F_{4,1}\bigr).
}\label{gff:gs:Lseries}\end{equation}
\begin{equation}\boxed{
 Z_d=\frac19\sum_{\ell=0}^{\infty}d_\ell(d)\sum_{n=0}^{\infty}T_{\ell n}^3
 \left[(9\mu_1^2+3\mu_2)F_3+6\mu_1F_{3,1}+F_{3,2}\right].
}\label{gff:gs:Zseries}\end{equation}
These are ordinary nonnegative-integer series at fixed dimension.
They contain only Pochhammer symbols and $\psi_0,\ldots,\psi_3$ at the
explicit arguments above.

For $\mathcal K_d$ define, for $j=0,1,2,3$,
\begin{equation}
 r_j=\frac12\lambda_{j+1}+\psi_j(b_{\ell n})+\psi_j(n+1)
                         -(1-(-1)^j)\psi_j(1),
 \qquad c_*=\psi_0(1)+\psi_0(2q),
 \label{gff:gs:rj}
\end{equation}
and the polynomial
\begin{equation}\begin{aligned}
 \mathcal R_{\ell n}={}&\frac{r_0-c_*}{2}
        (\lambda_1^3+3\lambda_1\lambda_2+\lambda_3)
   +\frac{3r_1}{4}(\lambda_1^2+\lambda_2)
   -\frac{\lambda_1r_2+r_3}{4}\\
 &-\frac{\lambda_1^4}{12}+\frac{\lambda_1^2\lambda_2}{8}
   +\frac{7\lambda_1\lambda_3}{24}
   +\frac{\lambda_2^2}{8}+\frac{\lambda_4}{12}.
\end{aligned}\label{gff:gs:Kpoly}\end{equation}
Then the explicit result is
\begin{equation}\boxed{
 \mathcal K_d=-\frac{AU}{4}
 +\sum_{\ell=0}^{\infty}d_\ell(d)
 \left\{\sum_{n=0}^{\infty}
 \left[\frac{M^2S_{\ell n}^2}{8}
       -\frac{T_{\ell n}^2\mathcal R_{\ell n}}4\right]\right\}.
}\label{gff:gs:Kseries}\end{equation}
The square brackets must be formed before summing: they are the two
pole families of the same convergent integral. The inner $n$ sum is
performed before the outer spin sum. Separating the two families or
rearranging this as an unrestricted double sum is not justified;
for example, absolute double summability fails in part of the physical
range even though the stated iterated sum converges.

\paragraph{Domains and continuation.}
For $2/3<d<4$, the $Q_d,Z_d$ series and their Barnes representations
converge directly; the corresponding lower threshold for $P_4$ and
$\mathcal L_d$ is $d>1/2$. The formulas cover $d=1,2,3$ without any
continuation. The subtracted $\mathcal K_d$ spin integral and its ordered
residue sum converge for $0<d<4$. Numerically, at $d=0.8$ the literal bracket
of \eqref{gff:gs:Kseries} loses its digits to cancellation far out in the radial
tail, so Section~\ref{gff:gdk:numerics} evaluates the canceled form
\eqref{gff:gc:Kcancelled} there instead. For smaller positive dimensions in
the other cases, use analytic continuation of the displayed special
functions from their common convergence strip, equivalently the
convergent normalized coordinate/Schwinger periods. Do not evaluate
individual divergent spin integrals and assign them finite numbers.
The analytically continued multiplicities at $d<1$ provide cancellations
that need not survive taking absolute values term by term.

\paragraph{Why these are hypergeometric-series data, rather than new graph names.}
The coefficient $d_\ell T_{\ell n}^{\,r}$ is a two-variable Horn-type
hypergeometric term: its ratios under $\ell\mapsto\ell+1$ and
$n\mapsto n+1$ are rational. Away from removable exceptional dimensions,
\[
 d_\ell(d)=\frac{(4q-2)_\ell(2q)_\ell}{(2q-1)_\ell\,\ell!}.
\]
Together with $(q)_{\ell+n}^r/(2q)_{\ell+n}^r$ and
$(1-q)_n^r/(n!)^r$, this prints the full parameter array of the underlying
Kamp\'e de F\'eriet-type series. The extra polygamma polynomials are
explicit resonant weights; the weighted series is not falsely labeled
as an unweighted ordinary ${}_pF_q$. Most importantly,
\eqref{gff:gs:seriesBase}--\eqref{gff:gs:Kseries} are executable definitions in
standard Gamma and polygamma functions, with no unsupplied master value.

\subsection{Residue proof of the series}\label{gff:gs:residueProof}
Consider the even Barnes moment for one fixed spin,
\begin{equation}
 I_r(a,g)=\int_{-\ii\infty}^{\ii\infty}\frac{ds}{2\pi\ii}
 \left[\frac{\Gamma(a+s)\Gamma(a-s)}
 {\Gamma(a+g+s)\Gamma(a+g-s)}\right]^r.
 \label{gff:gs:BarnesMoment}
\end{equation}
For the wheel at $b=d/2$, $a=a_\ell$ and $g=q$.
Closing an upward contour to the right is clockwise, so the integral
is the negative sum of residues at $s=a+n$. Put $s=a+n+u$. Then
\begin{equation}
 \frac{\Gamma(a+s)\Gamma(a-s)}
      {\Gamma(a+q+s)\Gamma(a+q-s)}
 =-\frac{T_{\ell n}}{\Gamma(2q)\,u}
       \exp\left(\sum_{m\geq1}\frac{\lambda_m}{m!}u^m\right).
 \label{gff:gs:localGamma}
\end{equation}
This follows from Gamma reflection and
\[
 \Gamma(-n-u)=\frac{(-1)^{n+1}\Gamma(1-u)}
 {n!\,u\prod_{j=1}^n(1+u/j)}.
\]
The regular logarithmic derivative of this last factor at order $m$
is $(1+(-1)^m)\psi_{m-1}(1)-\psi_{m-1}(n+1)$,
which is precisely the last line of \eqref{gff:gs:lambda}.
The second and third complete exponential Bell polynomials give
$F_3$ and $F_4$. Their signs are opposite because the pole orders are
three and four. This proves \eqref{gff:gs:QPseries}.

For a general balanced rim exponent put $x=b/2$, $g=d/2-x$.
The normalized residue prefactor is
\[
 -\frac{(x)_{\ell+n}(1-g)_n}{(d/2)_{\ell+n}\,n!}.
\]
Its logarithmic first and second $b$ derivatives at $b=d/2$ are
$\mu_1,\mu_2$. The corresponding first and second derivatives of
$\lambda_m$ are the already displayed $\eta_m,\theta_m$.
Differentiating the convergent residue identities in the proof and
then substituting these \emph{explicit} formulas yields
\eqref{gff:gs:Lseries}--\eqref{gff:gs:Zseries}. No differentiation remains in
those evaluation formulas.

For the subtracted master there are two right-hand pole families.
Write
\[
 g(s)=\frac{\Gamma(a+s)\Gamma(a-s)}{\Gamma(b+s)\Gamma(b-s)},\quad
 \ell(s)=\frac{g'(s)}{g(s)},\quad
 p(s)=\frac12\{\psi_0(a+s)+\psi_0(a-s)+\psi_0(b+s)+\psi_0(b-s)\}.
\]
Here $\ell(s)$ is the logarithmic derivative of $g(s)$, not the spin $\ell$. The holomorphic continuation
of the weight in \eqref{gff:gs:Kspectral} is
\begin{equation}
 F(s)=(c_*-p)\left(\frac{\ell(s)^2}{8}-\frac{5\ell'(s)}{16}\right)
       +\frac{\ell(s)p'(s)}4-\frac{5p''(s)}{32}.
 \label{gff:gs:Kholomorphic}
\end{equation}
Primes in this proof act on the elementary Gamma/polygamma kernel,
not on an unevaluated master. At $s=a+n+u$,
\[
 \ell(s)=-u^{-1}+\sum_{m\geq1}\lambda_m\frac{u^{m-1}}{(m-1)!},\qquad
 p(s)=\frac1{2u}+\sum_{j\geq0}r_j\frac{u^j}{j!}.
\]
Multiplying these finite jets through the required order gives the
negative residue $-T_{\ell n}^2\mathcal R_{\ell n}/4$ of
$\Gamma(2q)^2g(s)^2F(s)/4$. At the other family, $s=b+n+u$,
$g(s)^2$ has a double zero but $F(s)=-1/(2u^3)+O(u^{-2})$.
The resulting negative residue is $M^2S_{\ell n}^2/8$.
This proves both terms and the sign in \eqref{gff:gs:Kseries}.
The local residue polynomial was also checked independently by Cauchy
circles around both families at integer and noninteger dimensions.

\subsection{An additional exact two-dimensional master evaluation}\label{gff:gs:Z2section}
The dimension-two balanced three-wheel is a complex Selberg integral.
For $b$ close to one take, in the complex Selberg identity
\cite{gff:NeretinComplex}, the parameters $n=2$, $\sigma=b/2$, $\tau=1-b/2$ and
$\theta=-b/4$. These names are local to this paragraph; they are not the
$\sigma$ of \eqref{gff:gs:Acoords}, the $\tau$ of \eqref{gff:gs:RS} or the $\theta_m$ of
\eqref{gff:gs:etaTheta}. With $\gamma_C(x)=\Gamma(x)/\Gamma(1-x)$,
its normalized value simplifies to
\begin{equation}\boxed{
 W_3(2;b)=\frac12\frac{\gamma_C(b/4)^3}{\gamma_C(3b/4)}.
}\label{gff:gs:W3d2}\end{equation}
All convergence inequalities hold in a neighborhood of $b=1$:
$0<b<4/3$ suffices for the displayed two-screening integral.
At $b=1$ its logarithmic derivative vanishes, and its logarithmic
second derivative is
\[
 \frac34\{\psi_1(1/4)-\psi_1(3/4)\}=12\Cat.
\]
Combining this with \eqref{gff:gs:wheelJets} gives the exact check
\begin{equation}\boxed{
 Z_2=4\Cat\,Q_2,\qquad
 Q_2=\frac16\left[\frac{\Gamma(1/4)}{\Gamma(3/4)}\right]^4.
}\label{gff:gs:Z2exact}\end{equation}
This is a derivation from a convergent integral identity, not an
integer-relation fit to the numerical value of $Z_2$.

A second derivation of \eqref{gff:gs:W3d2} and of \eqref{gff:gs:Z2exact}, from Dixon's
theorem and the two-node formula of Section~\ref{gff:gdr:samples} without the complex
Selberg integral, is given in Section~\ref{gff:gdr:special}, \eqref{gff:mgz:R2}--\eqref{gff:mgz:Z2};
it holds for $0<\operatorname{Re}b<\frac43$.


\subsection{Special-function coefficient checks}
The one-dimensional specialization recovers the analytic formulas of
Part~II. The general master calculations give, at the precision relevant here,
\[
\begin{array}{c|rr}
 d&C_4(d)&C_5(d)\\\hline
 2&5.048868431&-38.833645896\\
 3&1.228063545&-5.850351421
\end{array}
\]
These are master/assembly evaluations, not a second full-graph calculation
in arbitrary dimension. The two-dimensional exact result
$Z_2=4\Cat Q_2$ supplies an independent master check. The closed forms
\eqref{gff:gc:L2closed}, \eqref{gff:gc:K2closed} make the $d=2$ entry of $C_5$ exact through
\eqref{gff:gc:C5twoClosed}.

\paragraph{Master values at $d=1,2,3$.}\label{gff:gs:masterTable}
The six masters of \eqref{gff:gs:C5} take the following values, correctly rounded to $30$ significant
digits:
\begin{center}\small
\begin{tabular}{@{}lrll@{}}\toprule
Master & $d$ & Value & Source\\\midrule
$Q_d$    & 1 & $55.0014865441629826199207662385$  & \eqref{gff:gd:d1blocks}\\
         & 2 & $12.7713811986469576577497636845$  & \eqref{gff:gc:Q2}\\
         & 3 & $5.01710256455385536959085484569$  & \eqref{gff:td:Qanswer}\\\addlinespace
$P_4(d)$ & 1 & $2229.52206107094196429452519975$  & \eqref{gff:eq:LnormalizedBasis}\\
         & 2 & $250.487315581867285531851756888$  & \eqref{gff:fp:P4two}\\
         & 3 & $62.0583279434066724664633885841$  & \eqref{gff:td:P4answer}\\\addlinespace
$P_5(d)$ & 1 & $28599.2931765621035366464770214$  & \eqref{gff:eq:P5FixedAnswer}\\
         & 2 & $1548.08264023312061513243907395$  & \eqref{gff:fp:P5two}\\
         & 3 & $240.340317759083549816366546388$  & \eqref{gff:td:P5answer}\\\addlinespace
$\mathcal K_d$ & 1 & $39.3315691282590390575031871516$ & \eqref{gff:eq:KprovedFinal}\\
         & 2 & $3.51810913859025224164909017031$  & \eqref{gff:gc:K2closed}\\
         & 3 & $1.14859464952717933655664631081$  & \eqref{gff:gs:Kseries}, record below\\\addlinespace
$\mathcal L_d$ & 1 & $-5805.30162954880238014117987700$ & \eqref{gff:eq:LpolyWheel}\\
         & 2 & $-347.249153123201131058476783075$ & \eqref{gff:gc:L2closed}\\
         & 3 & $-62.7105849683876693472961734335$ & records below\\\addlinespace
$Z_d$    & 1 & $764.931030199588577040161778769$  & \eqref{gff:gs:d1Dictionary}\\
         & 2 & $46.7925830723296966567388258780$  & \eqref{gff:gs:Z2exact}\\
         & 3 & $8.82092489306009248553686732817$  & records below\\\bottomrule
\end{tabular}
\end{center}
At $d=1$ and $d=2$ every entry is a closed form of this note, evaluated at $60$ digits; the entries
agree with all values of these masters printed elsewhere in the note. At $d=3$, $Q_3$, $P_4(3)$ and
$P_5(3)$ are the fixed rational formulas of Section~\ref{gff:app:threeDimensionalReduction}, evaluated at
$90$ digits with $160$ summed terms and a Hurwitz-zeta tail expansion of order $64$
(\path{LRI_C5_d3/inputs/general_d3/period_values.json}); Section~\ref{gff:d3:values} compares them with a less
converged $55$-digit run and with independent integrations of the Gamma kernels. $\mathcal K_3$ is the
$430$-digit ball-arithmetic evaluation of \eqref{gff:gs:Kseries} (\path{LRI_C5_d3/wf_K3/K3_value.json}); at
$d=3$ that series rests on input (M) of Section~\ref{gff:gdk:conventions} and on the extraction of
$[\tau^2]$ under the trace (Section~\ref{gff:gdk:status}). The two-fold Barnes integral \eqref{gff:gdk:KMB},
which uses neither, reproduces it to $4\cdot10^{-100}$ (Section~\ref{gff:gdk:numerics}).
$\mathcal L_3$ and $Z_3$ come from the $d=3$ formulas of Section~\ref{gff:gdr:special}, which rest on the
half-shift identities of Section~\ref{gff:gdr:halfShift} (records \path{LRI_C5_d3/wf_L3/L3_final_values.json},
about $158$ digits, and \path{LRI_C5_d3/wf_Z3/Z3_values.json}, about $250$ digits). An independent
$100$-digit evaluation of the
spectral traces of Section~\ref{gff:gs:spectralSection}, with the derivatives in $b$ of
\eqref{gff:gs:wheelJets} and the $\tau$-expansion of \eqref{gff:gs:Ksubtraction} taken on Cauchy circles
(\path{LRI_C5_d3/run_C_combined.json}), agrees with $\mathcal L_3$, $Z_3$ and $\mathcal K_3$ to
$1.2\times10^{-93}$, $2.9\times10^{-98}$ and $1.5\times10^{-96}$, and with $3Q_3$ and $P_4(3)$ to
$3.5\times10^{-79}$ and $2.1\times10^{-86}$. These agreements are numerical, not interval bounds.
Values at other dimensions (and $Z_d$ to $36$ digits) are in the validation tables of
Sections~\ref{gff:mgl:checks} and~\ref{gff:gdr:checks}; the first, lower-precision evaluations are described
in the paragraph on controls below.

\paragraph{An independent route to $P_5(2)$.}
The spectral formula \eqref{gff:gd:spectral} has a regular limit $d\to2$. Rename its angular index $n$ (of
\eqref{gff:gd:tau} and \eqref{gff:gd:rho}) to $\ell$. At $d=2$ one has $p=0$ and $h_2=1/\pi$, the joint limit
of the Gamma ratios in \eqref{gff:gd:rho} gives the multiplicities $d_\ell(2)=1,2,2,\ldots$, and
$\tau^{(2)}_\ell(\nu)=(-1)^\ell\pi\,t_\ell(\nu)$, with $t_\ell(\nu)$ of \eqref{gff:gs:spectrum} at $d=2$:
there $\Gamma(2q)=1$ and
$t_\ell(\nu)=|\Gamma(\ell/2+\frac14+\ii\nu)/\Gamma(\ell/2+\frac34+\ii\nu)|^2$. For $L=5$,
\begin{equation}
 P_5(2)=\frac1{2\pi}\sum_{\ell=0}^\infty(-1)^\ell d_\ell(2)
       \int_0^\infty(\ell^2+4\nu^2)\,t_\ell(\nu)^7\,d\nu .
 \label{gff:gs:P5d2}
\end{equation}
OpenAI GPT-6 astra evaluated \eqref{gff:gs:P5d2} in double precision, in the package
\path{lri_general_special_functions_update/} (script \path{code/general_special/spectral_evaluate.py},
record \path{verification/general_special/spectral_master_values.json}, entry \texttt{zigzag\_d2}). With
$192$, $384$ and $768$ quadrature nodes it agrees with the table to $6\times10^{-11}$, and the same
limit for $L=3$ reproduces $P_3(2)=3Q_2$. This route uses neither the square factorization of
Section~\ref{gff:fp:twoDimensional} nor the finite sampling of Section~\ref{gff:app:generalPeriodization}.

\paragraph{Controls of the master evaluations.}\label{gff:mgc:masterControls}
These controls belong to the same general-dimensional calculation (folder \path{code/general_special/}
of \path{lri_general_special_functions_update/}, records in \path{verification/general_special/}). The first values of $Q_d$, $P_4(d)$, $\mathcal L_d$, $Z_d$ and
$\mathcal K_d$ came from the common spectral kernel of Section~\ref{gff:gs:spectralSection}
(\path{spectral_evaluate.py}, record \path{spectral_master_values.json}). That evaluator is
independent of the two-contour kite representation, which was used separately to check $\mathcal K_d$
after extracting its pinch residue (\path{kappa_mb.py}, record \path{kappa_mb_checks.json}). The
common-kernel runs vary the spin cutoff ($128$, $256$, $512$), the number of quadrature nodes ($384$,
$512$, $768$) and the number of terms of an angular-tail expansion. That expansion follows from
centered Gamma ratios and even inverse powers of $\ell+(d-2)/2$; its fitted coefficients accelerate the
integration and are not an ansatz for a master identity. At $d=2$ and $d=3$ the two-contour and
common-kernel values of $\mathcal K_d$ agree to about $10^{-11}$ or better. The slower two-contour run
at $d=1$ is not assigned that accuracy; there the one-dimensional closed form is the comparison target.
Local Cauchy residues test both families of \eqref{gff:gs:Kseries}
(\path{local_residue_checks.json}), and raw finite sums of the series converge towards the spectral
values (\path{ordered_series_checks.json}); raw rectangular truncations are not assigned an error
bound, and the iterated convention of \eqref{gff:gs:Kseries} is kept. Exact SymPy calculations
(\path{exact_audit.py}) check the three-loop benchmark, both fixed-point assemblies, the cancellation of
$M$ and $N$, both one-dimensional specializations, the tetrahedron quadratic normalization and the local
$\mathcal K$ residue polynomial. The digits of these runs are superseded by the high-precision values of
Section~\ref{gff:gs:masterTable} (see also Sections~\ref{gff:gc:d2closedForms}, \ref{gff:d3:values}, \ref{gff:gdr:checks}
and~\ref{gff:gdk:numerics}); they remain
an independent lower-precision route. They are master-level checks, not a second general-dimensional
full-graph calculation, and not interval enclosures.

The one-dimensional polygamma form of $\mathcal K$ does not lift to general $d$. The combination
$2\psi_1(q)^2-3R\psi_1(q)+A\psi_2(q)/2$, with $A,R$ of \eqref{gff:gs:psiBasis}, equals $\mathcal K$ at
$d=1$ but gives $25.3748796\ldots$ at $d=2$, whereas $\mathcal K_2=3.5181091\ldots$ (record
\path{kappa_mb_checks.json}).


\section{Finite multivariable representations for the masters}\label{gff:app:generalClosed}

This section writes the masters of \eqref{gff:gs:C5} in standard function classes. For $\mathcal L_d$,
$Z_d$ and $\mathcal K_d$ it reaches finite combinations of unweighted multivariable Horn series at
unit arguments (Section~\ref{gff:gc:HornRules} and the finite arrays after it). This is a function-class
statement, not a reduction to a few ordinary ${}_pF_q$ values: the combinations are long, outside a
common convergence strip only the complete combination may be continued in $d$, and the weighted
double series of Section~\ref{gff:gs:seriesSection} remain preferable for numerical work. For the
primitive $P_5(d)$ the section gives a derivative-free double residue series
(Section~\ref{gff:gc:primitive}); the shorter finite sampling of the primitives is
Section~\ref{gff:app:generalPeriodization}, and at $d=2$ all masters have closed forms
(Section~\ref{gff:gc:d2closedForms}). Apart from those closed forms, the constructions of this section
were written by OpenAI GPT-6 astra (package \path{lri_general_closed_update/}).

\subsection{The two-dimensional normalization and its immediate consequences}
At $d=2$ the tetrahedron normalization \eqref{gff:gd:Qdef}, with $\mathcal T_2$ from
\eqref{gff:gd:d2exact}, gives
\begin{equation}\boxed{
 Q_2=\frac16\left[\frac{\Gamma(1/4)}{\Gamma(3/4)}\right]^4
     =\frac{\Gamma(1/4)^8}{24\pi^4}.
}\label{gff:gc:Q2}\end{equation}
Indeed, reflection gives $\Gamma(1/4)\Gamma(3/4)=\pi\sqrt2$, so the fourth
power of the ratio is $\Gamma(1/4)^8/(4\pi^4)$ \cite{gff:DLMF5}. The factor
$1/6$ is that of $Q_d=\mathcal T_d/6$ in \eqref{gff:gd:Qdef}; it must not be
dropped. Numerically,
\begin{equation}
 Q_2=12.7713811986469576577497636845240084220372849859347\ldots,
 \qquad Z_2=\frac{\Cat\,\Gamma(1/4)^8}{6\pi^4}.
 \label{gff:gc:Q2number}
\end{equation}
The second identity is $Z_2=4\Cat Q_2$ of \eqref{gff:gs:Z2exact}, written in the
same Gamma normalization.

For $a=\log2$, the dimension-two specialization of all coefficients is
\begin{equation}\boxed{\begin{aligned}
 C_0(2)&=1,&C_1(2)&=-\frac16,&C_2(2)&=\frac{4a}{9},\\
 C_3(2)&=\frac{4a^2}{27}-\frac{\Gamma(1/4)^8}{324\pi^4},\\
 648C_4(2)&=16P_4(2)-64a^3-32aQ_2-359\zeta(3),\\
 46656C_5(2)&=-1536\{P_5(2)+\mathcal K_2\}
 +1536aP_4(2)-576\mathcal L_2-1728Q_2^2\\
 &\quad +(11520a^2+28608\Cat)Q_2
 -3840a^4-21456a\zeta(3).
\end{aligned}}\label{gff:gc:d2coefficients}\end{equation}
The remaining dimension-two masters in this display are the fixed
hypergeometric functions specified below, not new undetermined constants; the two
nonprimitive ones, $\mathcal K_2$ and $\mathcal L_2$, are evaluated in closed form in
Section~\ref{gff:gc:d2closedForms}.
None of these identities is inferred from numerical values.

\subsection{Closed forms of the two-dimensional nonprimitive masters}
\label{gff:gc:d2closedForms}
The two masters of \eqref{gff:gc:d2coefficients} that are neither primitive periods nor
Gamma values have closed forms in constants that are already present:
\begin{equation}\boxed{
 \LL_2=-2\log2\,P_4(2)=-\frac A2\,P_4(2)
 =-\frac{\log2}{2\pi^3}\bigl(B_-^2-B_+^2\bigr),
}\label{gff:gc:L2closed}\end{equation}
\begin{equation}\boxed{
 \KK_2=32\Cat^2-28\zeta(3)\log2
 =\frac12\Bigl[\frac{\{\psi_1(\frac14)-\psi_1(\frac34)\}^2}{4}+AU\Bigr].
}\label{gff:gc:K2closed}\end{equation}
Here $A=4\log2$ and $U=\psi_2(\frac12)=-14\zeta(3)$ are the values of
\eqref{gff:gs:psiBasis} at $d=2$, $B_\pm$ are the retained one-dimensional values of
\eqref{gff:fp:inheritedValues}, and $\psi_1(\frac14)-\psi_1(\frac34)=16\Cat$ by
\eqref{gff:eq:quarterPolygamma}. The second form of \eqref{gff:gc:K2closed} reads
$\KK_2=\frac12[\Phi_2+AU]$ with $\Phi_2=64\Cat^2$. The one-dimensional result
\eqref{gff:eq:KprovedFinal} has the same shape, $\KK=\frac12[\Phi_1+AU]$ with
$\Phi_1=\psi_1(\frac14)\{3\psi_1(\frac14)-\psi_1(\frac34)\}$ and the $d=1$ values of
$A$ and $U$.

\paragraph{Status.}
Both identities are proved; no numerical identification enters either proof.
\begin{itemize}
\item \eqref{gff:gc:L2closed} is proved in Section~\ref{gff:gc:L2proof} for $\LL_2$ defined by
the series \eqref{gff:gs:Lseries} at $d=2$. The same proof establishes the two single-sum
relations \eqref{gff:gc:L2singles}, which were first found numerically. The
wheel-derivative reading $\partial_bW_4(2;b)|_{b=1}=-2\log2\,W_4(2;1)$ of
\eqref{gff:gs:wheelJets}, where $W_4(2;1)=P_4(2)$ by \eqref{gff:gs:primitiveWheel}, is the form
in which $\LL_d$ enters \eqref{gff:gs:D5}. It follows because the function $\mathcal W(x)$
of Section~\ref{gff:gc:L2proof} equals $W_4(2;2x)$ for $\frac25\leq x\leq\frac35$. This is
the per-spin residue evaluation of \eqref{gff:gs:BarnesMoment} at rim exponent $2x$,
written out at $d=2$ as \eqref{gff:mgp:Wresidue} in Section~\ref{gff:gc:L2second}; it is also the
case $d=2$ of \eqref{gff:gdr:functional} with \eqref{gff:gdr:d12lattice}. Since the two functions agree on an interval, so do their
derivatives. The wheel-derivative reading has also been proved directly, by a
rim--spoke reflection identity (Section~\ref{gff:gc:L2second}).
\item \eqref{gff:gc:K2closed} is proved in Section~\ref{gff:gc:K2proof} for $\KK_2$ defined by
the master path \eqref{gff:gs:Kdefinition}. The proof starts from the position-space
integral and uses neither \eqref{gff:gs:Kspectral} nor \eqref{gff:gs:Kseries}. The series
\eqref{gff:gs:Kseries} and the spectral form \eqref{gff:gs:Kspectral} are also proved to equal the right side
of \eqref{gff:gc:K2closed}, in Section~\ref{gff:mgk:section}, which also proves the mixed-trace step at $d=2$
and gives a second position-space proof of \eqref{gff:gc:K2master}.
\end{itemize}

Three further proofs were found independently. Two of them prove \eqref{gff:gc:L2closed}
through the rim derivative \eqref{gff:gc:L2derivative} and the residue functional
\eqref{gff:gc:L2residue} of Section~\ref{gff:gc:L2proof}, but obtain the derivatives of the two modes
differently: one evaluates the single sums of \eqref{gff:gc:L2reduction} by Hardy-space Parseval
pairings on the unit circle (Section~\ref{gff:mgp:parseval}), the other differentiates a Fourier
(Cauchy beta integral) representation of the two-mode coefficients (Section~\ref{gff:mgp:fourier}).
The logarithmic elliptic identities they need are proved in Section~\ref{gff:mgp:halfStrip} by
half-strip contours and the real Landen transformation, without Wan's Stieltjes identity. The third proves that the right side of
\eqref{gff:gs:Kspectral} at $d=2$ equals \eqref{gff:gc:K2closed}. It regularizes
\eqref{gff:gs:Ksubtraction} in the rim exponent, evaluates the regulated trace by residues on
the square lattice and periodizes the resulting single sums; it is written out in
Sections~\ref{gff:mgk:regularization}--\ref{gff:mgk:spectral}. Its statements for \eqref{gff:gs:Kseries} and
\eqref{gff:gs:Kdefinition} rest on the per-spin residue identity and on the mixed-trace step, which
Sections~\ref{gff:mgk:series} and \ref{gff:mgk:regulated} prove at $d=2$.

\paragraph{Numerical confirmation.}
These checks are not proof steps. A $70$-digit Euler--Maclaurin evaluation of the
square-lattice factorization of \eqref{gff:gs:Lseries}, and of \eqref{gff:gs:Kseries} through
the limit $\theta\to0$ of its factorized single sums at $z=e^{\ii\theta}$, gives the
residuals $4.5\times10^{-69}$ for \eqref{gff:gc:L2closed} and $2.0\times10^{-70}$ for
\eqref{gff:gc:K2closed}. The factorized single sums of \eqref{gff:gs:Kseries} diverge
individually, so the unrestricted double sum is not used; the limit $\theta\to0$ of the
lattice sum with weight $e^{\ii(m-n)\theta}$ preserves the spin order and equals the
ordered spin sum (part (c) of the proposition on the series and the strip limit,
Section~\ref{gff:mgk:series}). An independent evaluation from the spectral kernel, with
$\LL_2=\partial_bW_4(2;b)|_{b=1}$ as in \eqref{gff:gs:wheelJets} and $\KK_2$ from the spectral
form \eqref{gff:gs:Kspectral}, gives $1.7\times10^{-79}$ and $3.1\times10^{-62}$. To fifty
digits,
\[\begin{aligned}
 P_4(2)&=250.48731558186728553185175688791140716835571092358\ldots,\\
 \LL_2&=-347.24915312320113105847678307522316920816760199487\ldots,\\
 \KK_2&=3.5181091385902522416490901703077395226226728094710\ldots,\\
 P_5(2)&=1548.0826402331206151324390739457893403881102323229\ldots.
\end{aligned}\]

\paragraph{The five-loop coefficient.}
Inserting \eqref{gff:gc:L2closed} and \eqref{gff:gc:K2closed} into \eqref{gff:gc:d2coefficients}
and dividing by $48$ gives (exact algebra)
\begin{equation}\boxed{\begin{aligned}
 972\,C_5(2)={}&-32P_5(2)+56\log2\,P_4(2)-36Q_2^2
 +\bigl(240\log^22+596\Cat\bigr)Q_2\\
 &-1024\Cat^2+449\,\zeta(3)\log2-80\log^42 .
\end{aligned}}\label{gff:gc:C5twoClosed}\end{equation}
In the general-dimensional form \eqref{gff:gs:C5} the two substitutions are
$384AP_4-576\LL_2=672AP_4$ and $-1536\KK_2=-768(\Phi_2+AU)$ at $d=2$.
Thus $C_5(2)$ requires only $\Gamma(\frac14)$, $\pi$, $\log2$, $\Cat$, $\zeta(3)$,
the two ${}_4F_3(1)$ values $F_\pm$ of \eqref{gff:fp:FpmDictionary} and the five fixed
values \eqref{gff:fp:YDictionary} that enter \eqref{gff:fp:P5two}. All of them already occur in
the one-dimensional coefficients, and no multivariable value of
Section~\ref{gff:gc:HornRules} is needed at $d=2$. Numerically \eqref{gff:gc:C5twoClosed} gives
$C_5(2)=-38.83364589622536055293698\ldots$, in agreement with the table of
Section~\ref{gff:sec:guide}.

\paragraph{Remaining gaps.}
There is no gap in the proofs of \eqref{gff:gc:L2closed} and \eqref{gff:gc:K2closed} as stated.
The following related points are not settled in this note.
\begin{itemize}
\item Wan's Stieltjes identity \eqref{gff:lp:eq:stieltjes}, which enters the proof of
Section~\ref{gff:gc:L2proof} through $U_0=2N_{\rm e}$ and \eqref{gff:lp:eq:mixedlog}, is cited
rather than reproved there, exactly as in Section~\ref{gff:app:LreductionProof}. The proofs
of Sections~\ref{gff:gc:L2second}, \ref{gff:mgp:parseval} and \ref{gff:mgp:fourier} do not use it. The
half-strip contours of Section~\ref{gff:mgp:halfStrip} prove $U_0=2N_{\rm e}$ and reprove
\eqref{gff:lp:eq:U}, \eqref{gff:lp:eq:V} and \eqref{gff:lp:eq:R} without it, and Section~\ref{gff:mgp:P4second}
reproves \eqref{gff:fp:abWheel}, hence \eqref{gff:fp:P4two}, without it; among the proofs of
\eqref{gff:gc:L2closed} only the first still uses it, through \eqref{gff:lp:eq:mixedlog}.
\item The derivation of \eqref{gff:gs:Kspectral} in Section~\ref{gff:gs:spectralSection} uses the
mixed-trace representation $J_{2\tau}=C(\tau)\langle t_\ell^2R_\tau\rangle_d$ stated before
\eqref{gff:gs:Ksubtraction}. At $d=2$ it is proved for
$0<\tau<\frac12$, in its full-Mellin-line (real-part) reading only, in
Section~\ref{gff:mgk:regulated}; its reading for complex $\tau$ and its status at other $d$ are
input (M) of Section~\ref{gff:gdk:conventions}. The closed forms here do
not use it.
\item The four-exponent generalization of the rim--spoke reflection ratio of
Section~\ref{gff:gc:L2second} has been checked numerically but not proved; it is not used.
\end{itemize}



\subsection{A common residue dictionary}
For nonnegative integers $\ell,n$ put $m=\ell+n$, $N=\ell+2n$, and
\begin{equation}\begin{aligned}
 T_{\ell n}&=\frac{(q)_{\ell+n}(1-q)_n}{(2q)_{\ell+n}n!},\\
 d_0(4q)&=1,\qquad
 d_{j+1}(4q)=4q\frac{(2q+1)_j(4q-1)_j}{(2q)_j(2)_j},\quad j\geq0.
\end{aligned}\label{gff:gc:base}\end{equation}
The second line is a nonsingular way to implement angular multiplicities.
At $d=1$ the upper parameter $4q-1=0$ terminates the positive-spin series
at $j=0$, leaving just $\ell=0,1$. At $d=2$ every positive-spin
multiplicity is two. At noninteger $d<1$ this is an analytic multiplicity,
not a positive integer count of modes.

Define four useful classes of explicit scalar sequences:
\begin{equation}\begin{aligned}
 \Delta_r(a,b;v)&=\psi_r(v+a)-\psi_r(v+b),\qquad r\geq0,\\
 J_r(a;v)&=\psi_r(v+a),\qquad r\geq1,\\
 H(a;v)&=\psi_0(v+a)-\psi_0(a),\\
 c_j&=(-1)^{j+1}\psi_{j-1}(q)-\psi_{j-1}(1-q)
            +(1+(-1)^j)\psi_{j-1}(1),\\
 \widetilde D_j&=\Delta_{j-1}(q,2q;m)+\Delta_{j-1}(1-q,1;n),\qquad
 \lambda_j=c_j+\widetilde D_j.
\end{aligned}\label{gff:gc:dictionary}\end{equation}
Here $j\leq6$ suffices. All appearances of $\psi_r$ with a summation
index will be eliminated by finite Horn insertions below. Only $c_j$
and the other fixed-argument prefactors retain polygamma notation.

For the lower wheel functions, define, with the indicated finite ranges,
\begin{equation}\begin{aligned}
 \eta_r&=\frac{(-1)^r\psi_r(q)-\psi_r(1-q)}2
        +\frac{J_r(q;m)+J_r(1-q;n)}2,&1\leq r\leq3,\\
 \theta_r&=-\frac{(-1)^r\psi_{r+1}(q)+\psi_{r+1}(1-q)}4
        +\frac{J_{r+1}(q;m)+J_{r+1}(1-q;n)}4,&1\leq r\leq2,\\
 \mu_1&=\frac{H(q;m)+H(1-q;n)}2,\\
 \mu_2&=\frac{J_1(q;m)+J_1(1-q;n)-\psi_1(q)-\psi_1(1-q)}4.
\end{aligned}\label{gff:gc:wheelDictionary}\end{equation}
The ordinary finite polynomials needed are
\begin{equation}\begin{aligned}
 F_3&=\frac92\lambda_1^2+\frac32\lambda_2,&
 F_{3,1}&=9\lambda_1\eta_1+\frac32\eta_2,\\
 F_{3,2}&=9\eta_1^2+9\lambda_1\theta_1+\frac32\theta_2,\\
 F_4&=\frac{32}{3}\lambda_1^3+8\lambda_1\lambda_2+\frac23\lambda_3,\\
 F_{4,1}&=32\lambda_1^2\eta_1+8\eta_1\lambda_2
                      +8\lambda_1\eta_2+\frac23\eta_3.
\end{aligned}\label{gff:gc:wheelPolynomials}\end{equation}
Thus the already proved wheel definitions can be written as
\begin{equation}\boxed{\begin{aligned}
 Q_d&=\frac13\sum_{\ell,n\geq0}d_\ell(d)T_{\ell n}^{3}F_3,\\
 P_4(d)&=-\sum_{\ell,n\geq0}d_\ell(d)T_{\ell n}^{4}F_4,\\
 \mathcal L_d&=-\sum_{\ell,n\geq0}d_\ell(d)T_{\ell n}^{4}
                             (4\mu_1F_4+F_{4,1}),\\
 Z_d&=\frac19\sum_{\ell,n\geq0}d_\ell(d)T_{\ell n}^{3}
          \{(9\mu_1^2+3\mu_2)F_3+6\mu_1F_{3,1}+F_{3,2}\}.
\end{aligned}}\label{gff:gc:wheelSeries}\end{equation}
The subscripts on $F_{r,j}$ are labels of printed polynomials, not
unperformed derivatives. Equations~\eqref{gff:gc:base}--\eqref{gff:gc:wheelSeries}
are also the exact finite inputs for the unweighted Horn construction.

\subsection{A double residue series for \texorpdfstring{$P_5(d)$}{P5(d)}}\label{gff:gc:primitive}
This subsection completes the residue dictionary above for the remaining primitive (implementation
\path{code/general_closed/primitive_series.py} of \path{lri_general_closed_update/}). Like
\eqref{gff:gs:QPseries} for $Q_d$ and $P_4(d)$, the result is an explicit double series, with no Barnes
integral and no parameter derivative. At $d=1,2,3$ the coefficients use the fixed-value formulas
\eqref{gff:eq:P5FixedAnswer}, \eqref{gff:fp:P5two} and \eqref{gff:td:P5answer}; for other $d$ the stated
evaluator is \eqref{gff:fp:periodsFinite}. The series is an independent representation and a check of
them. The proof below uses objects defined in Section~\ref{gff:app:generalPeriodization}, which follows:
the Barnes form \eqref{gff:fp:primitiveStart} with the kernel \eqref{gff:fp:kernel}, the measure block
\eqref{gff:fp:measureBlock}, the array $\mathcal A_{\ell k}$ and the values $G_{\ell k}$ of
\eqref{gff:fp:measureAntisymmetry}, the two-channel identity \eqref{gff:fp:twoChannels} and the polynomials
$E_m$ of \eqref{gff:fp:BellFinite}.

In this subsection $\theta$ is $\theta_0=2q$ or $\theta_1=2q-1$, and $M_\theta$ is the measure block of
\eqref{gff:fp:measureBlock}; this $\theta$ is not a wheel sequence $\theta_r$ of \eqref{gff:gc:wheelDictionary}.
At the shifted arguments $q+N$, $N=0,1,2,\ldots$, the relation $\Gamma(x-N)=(-1)^N\Gamma(x)/(1-x)_N$
gives
\begin{equation}
 M_\theta(q+N)=\frac{\Gamma(\theta+q+N)\Gamma(\theta-q-N)}{\Gamma(q+N)\Gamma(-q-N)}
 =M_\theta(q)\,\frac{(\theta+q)_N(1+q)_N}{(q)_N(1+q-\theta)_N}.
 \label{gff:gc:Htheta}
\end{equation}
In the second channel below only $\ell\geq1$ contributes, so $N\geq1$ there, and one uses
\begin{equation}
 M_\theta(q+N)=M_\theta(q+1)\,
 \frac{(\theta+q+1)_{N-1}(2+q)_{N-1}}{(q+1)_{N-1}(2+q-\theta)_{N-1}},\qquad N\geq1,
 \label{gff:gc:HthetaShift}
\end{equation}
rather than the prefactor $M_\theta(q)$ of \eqref{gff:gc:Htheta}. For $\theta=2q-1$ that prefactor
contains $\Gamma(3q-1)$, which is singular at $q=\frac13$; there, for $N\geq1$, it multiplies the
vanishing $(3q-1)_N$. The value $M_{2q-1}(q)$ is itself singular at $q=\frac13$, so the $\ell=0$ terms
of the second channel are omitted rather than weighted by zero. In the block notation of
\eqref{gff:gc:positiveBlock}, $\ell=j+1$ and $N-1=j+2n$. The first channel uses both $\ell=0$ and
$\ell\geq1$.

With $c_j$, $\widetilde D_j$ and $\Delta_r$ of \eqref{gff:gc:dictionary} and $N=\ell+2n$ of
\eqref{gff:gc:base}, put, for $1\leq j\leq6$,
\begin{equation}\begin{aligned}
 \kappa_{\theta,j}&=\kappa^0_j+7\widetilde D_j
 +2^j\bigl\{\Delta_{j-1}(\theta+q,q;N)+\Delta_{j-1}(1+q,1+q-\theta;N)\bigr\},\\
 \kappa^0_j&=\begin{cases}(7-2^{j+1})\,c_j,&j\text{ odd},\\7c_j,&j\text{ even}.\end{cases}
\end{aligned}\label{gff:gc:kappa}\end{equation}
Thus $(\kappa^0_1,\ldots,\kappa^0_6)=(3c_1,7c_2,-9c_3,7c_4,-57c_5,7c_6)$, the same in both channels.
This follows from reflection and recurrence in the proof below, not from a fit. The
$\kappa_{\theta,j}$ are unrelated to the normalization $\kappa(d,z)$ of \eqref{gff:gd:normalization}.
Let $E_6$ be the polynomial of \eqref{gff:fp:BellFinite}, displayed in \eqref{gff:fp:BellExamples};
$E_6(x_1,\ldots,x_6)$ is the coefficient of $u^6$ in $\exp(\sum_{j=1}^6x_ju^j/j!)$. Then
\begin{equation}\boxed{\begin{aligned}
P_5(d)=\frac1{2\Gamma(2q)^2}
\sum_{\ell,n\geq0}(-1)^\ell d_\ell(d)\,T_{\ell n}^{7}
\Big[&M_{2q}(q+N)\,E_6(\kappa_{2q,1},\ldots,\kappa_{2q,6})\\
 &+\ell(\ell+4q-2)\,M_{2q-1}(q+N)\,
 E_6(\kappa_{2q-1,1},\ldots,\kappa_{2q-1,6})\Big].
\end{aligned}}\label{gff:gc:P5series}\end{equation}
Every ingredient is explicit and free of derivatives. The double series converges absolutely for
$\frac25<d<4$; smaller positive dimensions are reached by analytic continuation of the same master.

\emph{Proof.} Take $\frac12<q<1$ and start from the Barnes form \eqref{gff:fp:primitiveStart}, with
$p=2q-1$ and $a_\ell$, $b_\ell$, $g_\ell(s)$ of \eqref{gff:fp:kernel}; in this strip the contour
separates the pole families. The line points upwards and is closed to the right, so each spin
integral is minus the sum of its right residues. There are two families of right poles: $s=a_\ell+n$,
$n\geq0$, of order seven, from $g_\ell(s)^7$, and the measure poles $s=(p+k)/2$, $k\geq0$, from
$\Gamma(p-2s)$. The second family must not simply be discarded term by term.

\emph{Measure poles.} Since $\Res_{s=(p+k)/2}\Gamma(p-2s)=(-1)^{k+1}/(2\,k!)$ and
$\Gamma(p+k)\Gamma(-p-k)=(-1)^{k+1}\pi/[(p+k)\sin\pi p]$, minus the residue at $s=(p+k)/2$, times
the spin weight $(-1)^\ell d_\ell$, is $-\Gamma(2p+1)\sin(\pi p)\,\mathcal A_{\ell k}/(4\pi)$ before the
overall factor $\Gamma(2q)^5/2$, with $\mathcal A_{\ell k}$ and $G_{\ell k}=g_\ell((p+k)/2)$ of
\eqref{gff:fp:measureAntisymmetry}. The multiplicity $d_k$ in $\mathcal A_{\ell k}$ arises as
$\Gamma(2p+k)(p+k)/k!=\Gamma(2p+1)d_k/2$, by \eqref{gff:gs:spectrum} with $d-2=2p$. The factor is finite
in the initial strip and independent of $\ell,k$. Gamma reflection gives
$G_{k\ell}=(-1)^{\ell-k}G_{\ell k}$ for integer $\ell,k$ (Section~\ref{gff:app:generalPeriodization}), and
the power seven is odd, so
\begin{equation}
 \mathcal A_{k\ell}=-\mathcal A_{\ell k},\qquad \mathcal A_{\ell\ell}=0,\qquad
 \sum_{\ell,k\geq0}\mathcal A_{\ell k}=0.
 \label{gff:gc:antisymmetry}
\end{equation}
The last statement uses absolute convergence. For $\frac12<q<1$ and all $\ell,k\geq0$, with
$h=|\ell-k|$,
\begin{equation}
 |\mathcal A_{\ell k}|\leq\mathrm{const}\cdot(1+\ell+k)^{q-3}\,h\,(1+h)^{-7q}.
 \label{gff:gc:measureBound}
\end{equation}
Indeed, $d_\ell(d)=O((1+\ell)^{4q-2})$ by \eqref{gff:gs:spectrum} and $4q-2>0$, so
$d_\ell d_k\leq\mathrm{const}\cdot(1+\ell+k)^{8q-4}$; next,
$|(\ell+p)^2-(k+p)^2|=h(\ell+k+2p)\leq2h(1+\ell+k)$; finally, Gamma-ratio estimates
\cite{gff:DLMFStirling} bound the first ratio in $G_{\ell k}$ by a constant times $(1+\ell+k)^{-q}$ and,
for $\ell\geq k$, the second by a constant times $(1+h)^{-q}$, and $|G_{\ell k}|=|G_{k\ell}|$ covers
$\ell<k$. The right side of \eqref{gff:gc:measureBound} is summable over $\ell,k\geq0$: at fixed $\ell+k$
each value of $h$ occurs at most twice, so the factor $h(1+h)^{-7q}$ summed over these pairs is at most
$2\sum_{h\geq0}h(1+h)^{-7q}$, finite since $7q>2$, and $(1+\ell+k)^{q-3}$ is summable over $\ell+k$
since $q<2$. These estimates justify the exchange $\ell\leftrightarrow k$, so the measure
poles cancel by an exact involution; no divergent sum is discarded. (The periodized kernel \eqref{gff:fp:periodization} needs the
finer statement that canceling pairs lie in the same congruence class; it is shown after
\eqref{gff:fp:measureAntisymmetry}.)

\emph{The two channels.} Since $\Gamma(p+1\pm2s)=(p\pm2s)\Gamma(p\pm2s)$,
\[
 \bigl[(\ell+p)^2-4s^2\bigr]\Gamma(p+2s)\Gamma(p-2s)
 =\Gamma(p+1+2s)\Gamma(p+1-2s)+\ell(\ell+2p)\Gamma(p+2s)\Gamma(p-2s).
\]
After division by $\Gamma(2s)\Gamma(-2s)$ this is \eqref{gff:fp:twoChannels} at $V=2s$, because $p+1=2q$
and $2p=4q-2$. At $s=a_\ell+n+u$ one has $V=q+N+2u$, so the two channels are
$M_{2q}(q+N+2u)$ and $\ell(\ell+4q-2)M_{2q-1}(q+N+2u)$.

\emph{Local expansion.} At $s=a_\ell+n+u$, \eqref{gff:gs:localGamma} gives
$g_\ell(s)=-T_{\ell n}\exp(\sum_{j\geq1}\lambda_ju^j/j!)/(\Gamma(2q)u)$, with $\lambda_j=c_j+\widetilde
D_j$ of \eqref{gff:gc:dictionary}; these extend the $\lambda_m$ of \eqref{gff:gs:lambda} to $j\leq6$. For either
channel,
\[\begin{aligned}
 \log\frac{M_\theta(q+N+2u)}{M_\theta(q+N)}=\sum_{j\geq1}\frac{u^j}{j!}\Bigl\{
 &2^j\bigl[\psi_{j-1}(\theta+q+N)-\psi_{j-1}(q+N)\bigr]\\
 &+(-2)^j\bigl[\psi_{j-1}(\theta-q-N)-\psi_{j-1}(-q-N)\bigr]\Bigr\}.
\end{aligned}\]
Write $\mathrm{ct}_m(x)=\frac{d^m}{dx^m}\pi\cot(\pi x)$, so that reflection reads
$\psi_m(x)=(-1)^m\psi_m(1-x)-\mathrm{ct}_m(x)$ \cite{gff:DLMF5}. Applied to the two negative arguments, the
polygamma part of the second bracket becomes $2^j\Delta_{j-1}(1+q,1+q-\theta;N)$, and, since
$\mathrm{ct}_m$ has period one, its cotangent part is
$(-2)^j[\mathrm{ct}_{j-1}(-q)-\mathrm{ct}_{j-1}(\theta-q)]$. In both channels $\theta-q\equiv q\pmod1$,
and $\mathrm{ct}_m(-x)=(-1)^{m+1}\mathrm{ct}_m(x)$, so this equals
$(-2)^j[(-1)^j-1]\,\mathrm{ct}_{j-1}(q)$. It vanishes for even $j$. For odd $j$ it is
$2^{j+1}\mathrm{ct}_{j-1}(q)=-2^{j+1}c_j$, because then
$c_j=\psi_{j-1}(q)-\psi_{j-1}(1-q)=-\mathrm{ct}_{j-1}(q)$. Adding $7\lambda_j$ from $g_\ell(s)^7$ gives
exactly $\kappa_{\theta,j}$ of \eqref{gff:gc:kappa}. Hence $g_\ell(s)^7M_\theta(q+N+2u)$ equals
$-(T_{\ell n}/\Gamma(2q))^7\,M_\theta(q+N)\,u^{-7}\exp(\sum_{j\geq1}\kappa_{\theta,j}u^j/j!)$, and minus
its residue at $u=0$ is $(T_{\ell n}/\Gamma(2q))^7M_\theta(q+N)E_6(\kappa_{\theta,1},\ldots,
\kappa_{\theta,6})$, the sign coming from the seventh power. Summing over $n$ and $\ell$ with the
weights $(-1)^\ell d_\ell$ and multiplying by $\Gamma(2q)^5/2$ gives \eqref{gff:gc:P5series} for
$\frac12<q<1$. The local expansion is a device of the proof; the final polynomial needs no coefficient
extraction.

\emph{Convergence and continuation.} For fixed $\ell$ the terms of \eqref{gff:gc:P5series} are
$O(n^{-10q})$; for fixed $n$ each channel is $O(\ell^{q-2})$; in the joint large-index region they are
$O((\ell+n)^{-6q-2})$. These orders follow from $d_\ell(d)=O(\ell^{4q-2})$,
$T_{\ell n}=O((1+\ell+n)^{-q}(1+n)^{-q})$, $M_\theta(q+N)=O(N^{2\theta})$ and the boundedness of
$\kappa_{\theta,j}$; the elementary estimates are not written out. Hence the series converges
absolutely for $\frac1{10}<q<1$, locally uniformly away from the removable parameter singularities
treated by \eqref{gff:gc:HthetaShift}. Uniform convergence extends the identity from $\frac12<q<1$ to this
range, and smaller positive dimensions follow by analytic continuation of the same master; no finite
subtraction is added. $\square$

The numerical checks of \eqref{gff:gc:P5series} are collected with the other checks of this section in
Section~\ref{gff:mgc:gcChecks}.

\subsection{Making the \texorpdfstring{$\mathcal K_d$}{Kd} cancellation explicit}
The paired convention of \eqref{gff:gs:Kseries} is essential. We expose its common non-summable part before
passing to separate Horn terms. Put
\begin{equation}\begin{aligned}
 x&=\pi\cot(\pi q),\qquad P=\pi^2,\\
 R_\infty&=\frac{(P+x^2)^2}{2}=\frac{\pi^4}{2\sin^4(\pi q)},\\
 J_j&=J_j(2q;m)+J_j(1;n),\quad 1\leq j\leq3,\\
 H&=H(2q;m)+H(1;n).
\end{aligned}\label{gff:gc:Knotation}\end{equation}
Here $P$ is just $\pi^2$ and is not a graph period. The first four fixed
constants in \eqref{gff:gc:dictionary} simplify to
\begin{equation}
(c_1,c_2,c_3,c_4)=
\left(-x,-x^2-\frac23P,-2x^3-2Px,
 -6x^4-8Px^2-\frac{28}{15}P^2\right).
\label{gff:gc:cFirstFour}\end{equation}
In particular $c_1^3+3c_1c_2+c_3=0$. This identity removes an otherwise
standalone logarithmic large-index term.

The full remaining polynomial is
\begin{equation}\begin{split}
\mathcal P={}&\frac{\widetilde D_1^4}{6}-\frac{2x\widetilde D_1^3}{3}
 +\frac54\widetilde D_1^2\widetilde D_2+\frac34\widetilde D_1^2J_1
 -\left(\frac{13P}{12}+\frac{x^2}{4}\right)\widetilde D_1^2\\
&-\frac{5x}{2}\widetilde D_1\widetilde D_2+\frac5{12}\widetilde D_1\widetilde D_3
 -\frac{3x}{2}\widetilde D_1J_1-\frac14\widetilde D_1J_2
 +\left(\frac43Px+x^3\right)\widetilde D_1\\
&+\frac12\widetilde D_2^2+\frac34\widetilde D_2J_1
 +\left(-\frac{11P}{12}+\frac{x^2}{4}\right)\widetilde D_2
 -\frac{5x}{12}\widetilde D_3-\frac1{24}\widetilde D_4\\
&+H\left(\frac12\widetilde D_1^3-\frac32x\widetilde D_1^2+\frac32\widetilde D_1\widetilde D_2
       -P\widetilde D_1-\frac32x\widetilde D_2+\frac12\widetilde D_3\right)\\
&-\frac P2J_1+\frac x4J_2-\frac14J_3.
\end{split}\label{gff:gc:KregularPolynomial}\end{equation}
Every monomial contains at least one $\widetilde D_j$ or $J_j$; the possible $H$
factor does not occur alone. Direct exact algebra in the residue
polynomial gives $\mathcal R_{\ell n}=R_\infty+\mathcal P$.

Gauss summation \cite{gff:DLMFGaussUnity} gives a second useful identity:
\begin{equation}\begin{aligned}
 r(c)&=\frac{\Gamma(c)^2}{\Gamma(c-q)\Gamma(c+q)}
      ={}_2F_1\!\left(\begin{matrix}q,-q\\c\end{matrix};1\right),\\
 \delta(c)&=r(c)-1
 =-\frac{q^2}{c}{}_3F_2\!\left(\begin{matrix}1,1-q,1+q\\2,c+1\end{matrix};1\right).
\end{aligned}\label{gff:gc:GaussDelta}\end{equation}
The second line follows by isolating the zeroth term and shifting the
remaining index. Both series converge for the positive $c$ used here.
Using $S_{\ell n}$ and $M$ of \eqref{gff:gs:Kseries}, gamma reflection shows
\begin{equation}
 \frac{M^2S_{\ell n}^2}{T_{\ell n}^2}
 =2R_\infty\,r(m+2q)^2r(n+1)^2.
 \label{gff:gc:Kratio}
\end{equation}
Therefore the exact canceled form is
\begin{equation}\boxed{\begin{aligned}
\mathcal K_d=-\frac{AU}{4}
+\frac14\sum_{\ell,n\geq0}d_\ell(d)T_{\ell n}^2
\Big\{R_\infty\big[& (1+\delta(m+2q))^2(1+\delta(n+1))^2-1\big]
-\mathcal P\Big\}.
\end{aligned}}\label{gff:gc:Kcancelled}\end{equation}
The bracket involving $\delta$ has exactly eight nonconstant monomials,
with coefficients $\binom2i\binom2j$, $0\leq i,j\leq2$, $(i,j)\ne(0,0)$.
Together with \eqref{gff:gc:KregularPolynomial}, this is a finite polynomial
input to Horn functions; the two original divergent family sums have
not been assigned independent values.

For $0<q<1/2$, the monomial sums obtained this way are absolutely
convergent. In the joint large-index region the uncancelled base scales
as $O((\ell+n)^{-2})$, and every printed monomial supplies at least one
additional inverse power, up to harmless logarithms. At fixed $n$,
the remaining angular bound is $O(\ell^{2q-2}\log\ell)$, which explains
the upper endpoint $q=1/2$. At fixed $\ell$, the radial bound is summable
for positive $q$. The endpoint and larger dimensions must be treated
by continuation of the \emph{complete finite combination}, not by
separately summing divergent monomial series at argument one. The paired iterated series remains a direct alternative there.

\subsection{The standard hypergeometric function and the finite parameter rules}
\label{gff:gc:HornRules}
For finite lists of parameter--vector pairs define the ordinary
multivariable Horn-type generalized hypergeometric series
\begin{equation}\begin{split}
\mathfrak H\!\left[
 \begin{matrix}(a_\alpha,\boldsymbol u_\alpha)_\alpha\\
               (b_\beta,\boldsymbol v_\beta)_\beta\end{matrix};
 \boldsymbol z\right]
=\sum_{\boldsymbol n\in\mathbb N_0^r}
\frac{\prod_\alpha(a_\alpha)_{\boldsymbol u_\alpha\cdot\boldsymbol n}}
     {\prod_\beta(b_\beta)_{\boldsymbol v_\beta\cdot\boldsymbol n}}
\prod_{i=1}^r\frac{z_i^{n_i}}{n_i!}.
\end{split}\label{gff:gc:HornDefinition}\end{equation}
All vectors used below have nonnegative integer entries. Repeated pairs
are allowed. Thus each coefficient ratio on increasing a single index
is an explicit rational function of the indices. This is the standard
Horn/Gamma-ratio class, not a new LRI-specific function whose evaluation
has been left unspecified. It includes generalized Lauricella-type
arrays; the $\ell+2n$ vectors in the primitive make the more general
Horn notation convenient \cite{gff:BytevKniehlMoch,gff:MikhailovVolchanskiy}.

A finite set of identities is sufficient to print all its parameters.
For $a,b>0$, nonnegative integer $v$ and the indicated orders,
\begin{equation}\begin{aligned}
 J_r(a;v)&=(-1)^{r+1}r!\sum_{k\geq0}\frac1{(v+a+k)^{r+1}},&&r\geq1,\\
 \Delta_r(a,b;v)&=(-1)^{r+1}r!(b-a)
 \sum_{j=0}^{r}\sum_{k\geq0}
 \frac1{(v+a+k)^{r+1-j}(v+b+k)^{j+1}},&&r\geq0,\\
 H(a;v)&=v\sum_{k\geq0}\frac1{(a+k)(a+v+k)},\\
 \delta(a+v)&=-q^2\sum_{k\geq0}
 \frac{(1-q)_k(1+q)_k}{(2)_k(a+v)_{k+1}}.
\end{aligned}\label{gff:gc:insertionIdentities}\end{equation}
The first three follow from polygamma series and a finite partial-fraction
identity \cite{gff:DLMFPolygamma}; the last is \eqref{gff:gc:GaussDelta}.
For a product, assign a \emph{different} index $k$ to every factor.
The choices $0\leq j\leq r$ are finite, so expanding any of the printed
master polynomials produces finitely many unweighted series.
There is no harmonic-number or polygamma factor left in their coefficients.

Every reciprocal affine power is encoded by one pair of parameter rows:
\begin{equation}
 (a+\boldsymbol v\cdot\boldsymbol n)^{-s}
 =a^{-s}\frac{(a)_{\boldsymbol v\cdot\boldsymbol n}^{s}}
               {(a+1)_{\boldsymbol v\cdot\boldsymbol n}^{s}}.
 \label{gff:gc:reciprocalEncoding}
\end{equation}
For the last identity in \eqref{gff:gc:insertionIdentities}, use
\begin{equation}
 \frac1{(a+v)_{k+1}}=\frac1a\frac{(a)_v}{(a+1)_{v+k}}.
 \label{gff:gc:GaussEncoding}
\end{equation}
An otherwise absent $k!$ is supplied by an upper pair $(1,\boldsymbol e_k)$
in \eqref{gff:gc:HornDefinition}, which cancels its standard denominator.
This fixes all factorial conventions explicitly.

The polynomial factor $v$ in a harmonic insertion, and the primitive's
$\ell(\ell+4q-2)$, also require no Euler derivative. Expand any resulting
power of an index into falling factorials:
\begin{equation}
 n^s=\sum_{j=0}^s\left\{\begin{matrix}s\\j\end{matrix}\right\}
 n^{\underline j},\qquad
 (m+j)^{\underline j}=j!\frac{(j+1)_m}{(1)_m}.
 \label{gff:gc:falling}
\end{equation}
Shift $n=m+j$ in each term. Each old parameter pair
$(a,\boldsymbol v)$ is shifted to $(a+v_nj,\boldsymbol v)$, with a
prefactor $(a)_{v_nj}$ in the numerator or denominator as appropriate.
Multiply by $j!z_n^j$ and add the upper/lower pair
$(j+1,\boldsymbol e_n)/(1,\boldsymbol e_n)$. The braces in
\eqref{gff:gc:falling} are Stirling numbers of the second kind, so this
operation is a finite algebraic rule.

\subsubsection{Explicit starting arrays}
For any of the powers $r=2,3,4,7$, split the original double sum into
$\ell=0$ and $\ell=j+1$. Before inserting weights, the zero-spin block
has scalar prefactor one and rows
\begin{equation}
 \begin{matrix}
 (1;1),\ (q;1)^{\times r},\ (1-q;1)^{\times r}\\
 (2q;1)^{\times r},\ (1;1)^{\times r}
 \end{matrix},\qquad z_n=1.
 \label{gff:gc:zeroBlock}
\end{equation}
The positive-spin block has scalar prefactor $4q/2^r$ and rows
\begin{equation}
 \begin{matrix}
 (1;1,0),(1;0,1),(2q+1;1,0),(4q-1;1,0),
 (q+1;1,1)^{\times r},(1-q;0,1)^{\times r}\\
 (2q;1,0),(2;1,0),(2q+1;1,1)^{\times r},(1;0,1)^{\times r}
 \end{matrix},\qquad (z_j,z_n)=(1,1).
 \label{gff:gc:positiveBlock}
\end{equation}
Here semicolons separate a scalar parameter from its index vector;
these are precisely the pairs in \eqref{gff:gc:HornDefinition}.
\subsubsection{An explicit parameter-array example}
For illustration, it is enough to print one unambiguous insertion
instead of thousands of nearly identical arrays. The zero-spin contribution
of $T_{0n}^3\Delta_0(q,2q;n)$ alone is
\begin{equation}\begin{split}
 -\frac1{2q}\mathfrak H\!\left[
 \begin{matrix}
 (1;1,0),(1;0,1),(q;1,0)^{\times3},(1-q;1,0)^{\times3},
 (q;1,1),(2q;1,1)\\
 (2q;1,0)^{\times3},(1;1,0)^{\times3},
 (q+1;1,1),(2q+1;1,1)
 \end{matrix};(1,1)\right].
\end{split}\label{gff:gc:example}\end{equation}
Its coefficient is obtained from
$\Delta_0(q,2q;n)=-q\sum_k[(n+q+k)(n+2q+k)]^{-1}$.
There is no remaining derivative or finite-index harmonic sum in
\eqref{gff:gc:example}. All other entries follow the same printed rules.
The two-variable example is deliberately elementary; products insert
additional variables, up to the bounds in the following table.

\subsection{Finite arrays, scope of continuation, and what is now closed}
The scalar polynomial inputs for the array construction are precisely
\eqref{gff:gc:wheelPolynomials}, \eqref{gff:gc:wheelSeries},
\eqref{gff:gc:KregularPolynomial}--\eqref{gff:gc:Kcancelled}.
They have finite degree. Applying
\eqref{gff:gc:insertionIdentities}--\eqref{gff:gc:falling} therefore gives a
finite list of conventional parameter arrays for every master:
\begin{equation}
 M(q)=M_{\rm scalar}(q)+\sum_{\alpha=1}^{N_M}
 p_{M,\alpha}(q)\,
 \mathfrak H[\mathsf A_{M,\alpha}(q);
             \mathsf B_{M,\alpha}(q);\boldsymbol z_{M,\alpha}],
 \label{gff:gc:finiteClosure}
\end{equation}
where $M$ is one of $\mathcal L,Z,\mathcal K$. The primitive masters use the smaller finite-sample representation of Section~\ref{gff:app:generalPeriodization}; their residue series (the $Q_d$ and $P_4(d)$ lines of \eqref{gff:gc:wheelSeries}, and \eqref{gff:gc:P5series}) are not converted into arrays here, and Section~\ref{gff:mgc:fpChecks} compares the counts such a conversion gives.
Only $M=\mathcal K$ has a separate scalar term, $-AU/4$.
The coefficients $p_{M,\alpha}$ use rational functions of $q$, Gamma
values and \emph{fixed-argument} polygamma values through $\psi_3$.
All parameter rows are explicit affine functions of $q$ with rational
coefficients, and the nonprimitive argument vectors contain only $1$. No unspecified generic-index kite function or Feynman
integral is included in \eqref{gff:gc:finiteClosure}.

The supplied direct compiler cancels identical parameter pairs and
merges literally identical arrays but does not attempt contiguous
relations or permutations of all auxiliary axes. Its nonminimal counts are
\begin{center}\small
\begin{tabular}{@{}lrr@{}}\toprule
Master & Number of fixed Horn values & Maximum variables per value\\\midrule
$\mathcal L_d$ & 244 & 6\\
$Z_d$ & 318 & 6\\
$\mathcal K_d$ & 386 & 6\\
\bottomrule
\end{tabular}
\end{center}
These counts quantify an important tradeoff: a finite unweighted
function-class expression has been reached, but the mechanical expansion
is \emph{not shorter} than the weighted double series. It gives finite function-class closure, not a short scalar-hypergeometric evaluation.
The counts assert neither minimality nor independence. A much shorter
representation is obtained by a different route in Sections~\ref{gff:gdr:section}
and \ref{gff:gdk:section}: periodization in the rim exponent and the top-mode
theorem \eqref{gff:gdr:topless} reduce $\mathcal L_d$ and $Z_d$ to the values and
first and second rim-exponent derivatives of eight functions with at most two
summation indices (the inventory after \eqref{gff:gdr:C5reduced}), and, under the
inputs stated in Section~\ref{gff:gdk:conventions}, $\mathcal K_d$ to a third-order
$\tau$-jet of three-node samples of three two-variable Horn functions
\eqref{gff:gdk:Ksample}. These forms keep parameter derivatives, so they lie
outside the unweighted class of \eqref{gff:gc:finiteClosure}.

The array files, their scalar-parameter dictionary, and the compiler
are included in \path{lri_general_closed_update/} under
\path{code/general_closed/} and
\path{verification/general_closed/horn_arrays/}.
Each JSON row records the scalar prefactor, the full upper and lower
parameter--vector lists with multiplicities, and the argument vector.
The equations above specify a finite way of generating those same lists
without needing the program; the files are a compact machine-readable
certificate, not a replacement for a missing mathematical definition.

\paragraph{Continuation is part of the formula.}
There is a nonempty common absolute-convergence strip
\begin{equation}
 \frac16<q<\frac12,\qquad\text{equivalently }\frac23<d<2,
 \label{gff:gc:commonStrip}
\end{equation}
for this finite construction. This follows from the wheel bounds and the canceled $\mathcal K$ bounds. The tail insertions have fixed sign or are finite signed sums
of positive denominators there, and their absolute sums have the same
bounds as the weights they replace. This justifies each finite expansion
and the added summation indices in that strip.

For other $0<d<4$, use the unique dimensional analytic continuation of
\emph{the complete finite combination} from \eqref{gff:gc:commonStrip}.
At a removable exceptional value, evaluate that continued combination,
not a separate pole part of each summand. In particular some individual
$\mathcal K$ Horn series diverge at unity for $d\geq2$, even though their
continued combination is the finite physical master. This is not an
extra free regulator, nor does it introduce an adjustable finite part;
it is the continuation already fixed by the original subtracted master.
It does mean that the result must not be advertised as separately
convergent unity-argument series for every physical dimension. Direct
paired summation or the original independent master integral is often
numerically preferable at these points.

Thus \eqref{gff:gs:C0C4}--\eqref{gff:gs:C5}, together with the finite parameter
rules, give all coefficients through $C_5(d)$ in a specified standard
multivariable hypergeometric class. Sections~\ref{gff:gdr:section} and
\ref{gff:gdk:section} shorten this considerably, at the price of low-order
parameter derivatives; \eqref{gff:gdr:C5reduced} is $C_5(d)$ in that form. A
finite reduction to ordinary ${}_pF_q$ values for generic $d$ remains a
different, stronger question. It is not asserted by these results, and the
searches of Section~\ref{gff:gdr:attempts}, which found no such reduction of the
remaining constants, are numerical and rule nothing out. At $d=2$ such a reduction
holds for the whole coefficient: Section~\ref{gff:gc:d2closedForms} expresses $C_5(2)$
through ${}_4F_3(1)$, ${}_9F_8(-1)$ and ${}_8F_7(-1)$ values and fixed constants.

\paragraph{Checks of the finite multivariable construction.}\label{gff:mgc:gcChecks}
These checks belong to the package \path{lri_general_closed_update/} (scripts
\path{check_new_identities.py}, \path{check_primitive_sums.py}, \path{check_horn_compiler.py} and
\path{check_insertion_prefactors.py} in \path{code/general_closed/}; records \path{identity_checks.json},
\path{primitive_sum_checks.json}, \path{horn_compiler_checks.json} and
\path{insertion_prefactor_checks.json} in \path{verification/general_closed/}). They are consistency
checks, not interval bounds. The Gamma-ratio and
eighth-power forms of $Q_2$ in \eqref{gff:gc:Q2}, evaluated separately at 60 working digits, differ by about
$-4.98\times10^{-60}$. The same run checks the canceled polynomial \eqref{gff:gc:KregularPolynomial}, the
normalization \eqref{gff:gc:Kratio} and the shifted Gauss identity \eqref{gff:gc:GaussDelta} at
$q=1/4,3/10,1/2,3/4$ and several nonzero indices, with residuals below $2.0\times10^{-58}$. The Horn
compiler is checked by 1038 exact rational comparisons of normalized coefficients at $q=3/10$, over 519
sampled rows and including inserted tail indices, and by 12 falling-factorial shift identities. A further
336 exact rational checks cover the scalar prefactor of every type of tail insertion, and exact algebra
verifies the polynomial $E_6$ and both dimension-two substitutions in \eqref{gff:gc:d2coefficients}. These
checks verify the parameter encoding; they do not sum every continued multivariable endpoint function
numerically.

For the double residue series \eqref{gff:gc:P5series} of Section~\ref{gff:gc:primitive}: at $60$
digits, Cauchy circles around poles of both channels reproduce the local residues, with largest
normalized residual below $8.8\times10^{-55}$; the antisymmetry \eqref{gff:gc:antisymmetry} was checked
for unequal spins at $q=\frac35$ and $\frac34$; and the constants $\kappa^0_j$ agree with the
unreduced reflection expressions. Double-precision sums of \eqref{gff:gc:P5series}, with $112$ spins, an
Euler transform of the alternating spin sum and radial cutoffs $128$ to $1024$, were compared with
spectral quadrature of \eqref{gff:gd:spectral} at $224$ and $352$ nodes:
\begin{center}\small
\begin{tabular}{@{}rrr@{}}\toprule
$d$ & residue-sum value of $P_5(d)$ & residue sum minus $352$-node control\\\midrule
2 & $1548.0826402331213$ & $-4.43\times10^{-11}$\\
$13/5$ & $473.7956970211414$ & $-6.37\times10^{-12}$\\
3 & $240.3403177590616$ & $-2.77\times10^{-11}$\\\bottomrule
\end{tabular}
\end{center}
The displayed digits are outputs, not certified digits. Against the high-precision values of
Section~\ref{gff:gs:masterTable} the residue sums differ by $+6.8\times10^{-13}$ at $d=2$ and
$-2.2\times10^{-11}$ at $d=3$. At $d=1$ the raw radial series converges much more slowly: against a
$60$-digit evaluation of the Gamma integral its residual falls from $1.05\times10^{-5}$ at cutoff
$128$ to $5.86\times10^{-8}$ at cutoff $1024$. No high-precision agreement is claimed there; the
five-value formula \eqref{gff:eq:P5FixedAnswer} remains the evaluator at $d=1$.

The rigorous parts of the construction are the algebraic normalization, the local residue
identities, the involutive cancellation of the measure poles of $P_5(d)$ (Section~\ref{gff:gc:primitive})
and the finite Horn conversion in its stated convergence domain, followed by the prescribed
continuation; the numerical residuals corroborate these steps but are not interval bounds. The
physical identification of the coefficients rests on the graph reduction and forest assembly of
Sections~\ref{gff:app:generalDimension} and~\ref{gff:app:generalSpecial}; no independent
general-dimensional five-loop graph calculation is claimed.


\section{Finite periodization of the primitive sector}\label{gff:app:generalPeriodization}

\subsection{The three Barnes kernels}
Put $q=d/4$, $p=2q-1$ and
\begin{equation}\begin{aligned}
 a_\ell&=(\ell+q)/2,& b_\ell&=a_\ell+q,\\
 g_\ell(s)&=\frac{\Gamma(a_\ell+s)\Gamma(a_\ell-s)}
 {\Gamma(b_\ell+s)\Gamma(b_\ell-s)},&
 D_q(t)&=\sin(\pi t)\sin\bigl(\pi(t+q)\bigr).
\end{aligned}\label{gff:fp:kernel}\end{equation}
Use $d_\ell=d_\ell(4q)$ from \eqref{gff:gs:spectrum}. A version with no
spurious $0/0$ at $d=1,2$ is
\begin{equation}
 d_0=1,\qquad d_{k+1}=4q\,
 \frac{(2q+1)_k(4q-1)_k}{(2q)_k(2)_k},\qquad k\geq0.
 \label{gff:fp:degeneracy}
\end{equation}
The normalized wheel periods are
\begin{equation}
 W_r=\Gamma(2q)^r\sum_{\ell\geq0}d_\ell
 \int_{-\ii\infty}^{\ii\infty}\frac{ds}{2\pi\ii}\,g_\ell(s)^r,
 \qquad W_3=3Q_d,\quad W_4=P_4(d).
 \label{gff:fp:wheelStart}
\end{equation}
For the five-loop primitive define
\begin{equation}\begin{aligned}
 K_\ell(s)&=\bigl[(\ell+p)^2-4s^2\bigr]
 \frac{\Gamma(p+2s)\Gamma(p-2s)}{\Gamma(2s)\Gamma(-2s)}
 g_\ell(s)^7,\\
 P_5(d)&=\frac{\Gamma(2q)^5}{2}\sum_{\ell\geq0}(-1)^\ell d_\ell
 \int_{-\ii\infty}^{\ii\infty}\frac{ds}{2\pi\ii}\,K_\ell(s).
\end{aligned}\label{gff:fp:primitiveStart}\end{equation}
Initially take $1/2<q<1$ so that the contours separate the pole families
without crossed-pole terms. The result below is then continued with the
same prescription as the original primitive master. The kernels and
normalization are the wheel and conformal-triangle kernels already
specified in this note, based on the operator framework of \cite{gff:DIS}.
No new renormalization convention is introduced.

For the proof only, periodize the kernels:
\begin{equation}\begin{aligned}
 S_r(t)&=\sum_{\ell\geq0}d_\ell\sum_{n\in\mathbb Z}
             g_\ell(a_\ell+n+t)^r,&&r=3,4,\\
 S_7(t)&=\sum_{\ell\geq0}(-1)^\ell d_\ell\sum_{n\in\mathbb Z}
             K_\ell(a_\ell+n+t),\\
 \Xi_r(t)&=D_q(t)^rS_r(t),&&r=3,4,7.
\end{aligned}\label{gff:fp:periodization}\end{equation}
The infinite sums in \eqref{gff:fp:periodization} are proof devices, not
additional operations in the finite evaluation rule below.

\subsection{Why finitely many Fourier modes suffice}
\subsubsection{Wheel poles, reflection and growth}
For $r=3,4$ the only pole classes of $S_r$ modulo one are $t=0$ and
$t=-q$, each of order at most $r$. Gamma reflection gives, for $n\geq0$,
\begin{equation}
 g_\ell(a_\ell+n+t)=
 -\frac{\sin\pi(q-t)}{\sin\pi t}
 \frac{\Gamma(\ell+n+q+t)\Gamma(n+1-q+t)}
 {\Gamma(\ell+n+2q+t)\Gamma(n+1+t)}.
 \label{gff:fp:reflectedRight}
\end{equation}
Reflection, recurrence and the Gamma-ratio asymptotics used here are
standard identities \cite{gff:DLMF5,gff:DLMFStirling}; the periodization argument
is given explicitly rather than inferred from those identities alone.

On compact sets avoiding the pole classes, split the lattice into
$n\geq0$, $n\leq-\ell-1$, and $-\ell\leq n\leq-1$. In the middle part
put $n=-j-1$, $\ell=j+k+1$; then
\begin{equation}
 g_\ell(a_\ell+n+t)=
 \frac{\Gamma(k+q+t)\Gamma(j+1-t)}
 {\Gamma(k+2q+t)\Gamma(j+q+1-t)}.
 \label{gff:fp:middleFactor}
\end{equation}
These bounds give locally uniform convergence, including after the
possible pole factors are cleared. Consequently $\Xi_r$ is entire.
Evenness of $g_\ell(s)$ and the integer reindexing $n\mapsto-\ell-n$
give
\begin{equation}
 \Xi_r(t+1)=\Xi_r(t),\qquad \Xi_r(-q-t)=\Xi_r(t).
 \label{gff:fp:symmetry}
\end{equation}

For $t=x+\ii y$, $0\leq x\leq1$, $|y|\geq1$, the right-hand sum is
bounded by a constant times
\[
 \sum_{\ell,n\geq0}(1+\ell)^{4q-2}
 (1+|\ell+n+\ii y|)^{-rq}(1+|n+\ii y|)^{-rq}.
\]
The middle expression \eqref{gff:fp:middleFactor} has the analogous bound
with weight $(1+j+k)^{4q-2}$. Integral comparison, initially for
$1/2<q<1$, gives
\begin{equation}
 |S_r(x+\ii y)|=O\bigl(|y|^{-2q(r-2)}\bigr),\qquad
 |\Xi_r(x+\ii y)|=O\bigl(e^{2\pi r|y|}|y|^{-2q(r-2)}\bigr).
 \label{gff:fp:wheelGrowth}
\end{equation}
For example the endpoint $\ell\gg n,|y|$ has power
$\ell^{q(4-r)-2}$, which is summable in this strip. This also controls
the endpoints of the scaled integral, not just simultaneous large
indices.

\subsubsection{Five-loop measure poles cancel class by class}
The extra measure poles of $K_\ell$ occur at $2s=\pm(p+k)$,
$k\geq0$. Their \emph{weighted} right residues have the form
\begin{equation}\begin{aligned}
 \mathcal A_{\ell k}&=(-1)^\ell d_\ell d_k
 \bigl[(\ell+p)^2-(k+p)^2\bigr]G_{\ell k}^{\,7},\\
 G_{\ell k}&=
 \frac{\Gamma((\ell+k+3q-1)/2)}{\Gamma((\ell+k+5q-1)/2)}
 \frac{\Gamma((\ell-k+1-q)/2)}{\Gamma((\ell-k+1+q)/2)}.
\end{aligned}\label{gff:fp:measureAntisymmetry}\end{equation}
Gamma reflection \cite{gff:DLMF5}, applied to the second ratio, gives
$G_{k\ell}=(-1)^{\ell-k}G_{\ell k}$ and hence
$\mathcal A_{k\ell}=-\mathcal A_{\ell k}$; the diagonal residue is zero.
For the present purpose cancellation of the \emph{total} measure-residue
sum is not by itself sufficient. One must check the congruence classes.
A right measure pole is at
\[
 t=\frac{q-1+k-\ell}{2}\pmod1.
\]
Exchanging $k,\ell$ changes this by an integer, so the canceling
residues are in the same class. A left measure pole is at
$t=(1-3q-k-\ell)/2\pmod1$, also unchanged by the exchange. Thus the
extra poles cancel locally in the spin-summed periodization, not only
after the final contour integration. Absolute convergence in the initial
strip, which the bound \eqref{gff:gc:measureBound} provides, permits the pairing. Coincident pole parameters are handled by
continuation of the complete expression, not by keeping a separate
singular residue.

After this cancellation only the two order-seven Gamma-pole classes
remain; $\Xi_7$ is entire and satisfies \eqref{gff:fp:symmetry}. The measure
contributes at most a factor of order $(1+\ell+n+|y|)^{4q}$ to the
bound behind \eqref{gff:fp:wheelGrowth}, so
\begin{equation}
 |S_7(x+\ii y)|=O(|y|^{-6q}),\qquad
 |\Xi_7(x+\ii y)|=O(e^{14\pi|y|}|y|^{-6q}).
 \label{gff:fp:primitiveGrowth}
\end{equation}
The endpoint $\ell\gg n,|y|$ now has power $\ell^{q-2}$, again
summable for $q<1$.

\subsubsection{Finite-dimensionality theorem}
For an entire one-periodic function the Fourier coefficient is unchanged
when its integration contour is translated vertically. Apply this to
\eqref{gff:fp:wheelGrowth} or \eqref{gff:fp:primitiveGrowth}. For mode $k\geq r$,
translation down by $\ii Y$ bounds its absolute value by
$C e^{2\pi(r-k)Y}Y^{-\eta}$ with $\eta>0$; it therefore vanishes,
including the boundary mode $k=r$. Upward translation deals with
$k\leq-r$. Reflection \eqref{gff:fp:symmetry} pairs the remaining modes.
Thus, exactly,
\begin{equation}\boxed{
 \Xi_r(t)=\sum_{k=0}^{r-1}c_{r,k}(q)
                 \cos\bigl(2\pi k(t+q/2)\bigr),\qquad r=3,4,7.
}\label{gff:fp:finiteFourier}\end{equation}
For real $q$ the coefficients are real, since the meromorphic kernel is
real on nonsingular real arguments. This is a proved finite Fourier
space, not a fitted trigonometric ansatz. For the wheels the top coefficient
also vanishes: $c_{3,2}(q)=c_{4,3}(q)=0$ for $\frac2r<d<4$, by
\eqref{gff:gdr:topless}, proved in Section~\ref{gff:gdr:top} for a general rim
exponent. The sums in \eqref{gff:fp:finiteFunctional} may therefore stop at
$k=r-2$ for $r=3,4$; the rules below remain exact as written. No such
statement is made for $r=7$, whose kernel is not a product of balanced
factors.

\subsection{A finite evaluation rule, with no derivatives}
A stable fixed sampling grid is
\begin{equation}\begin{aligned}
 N_r&=2r-1,& \eta_r&=\frac1{8\pi r},&
 t_j&=-\frac q2+\frac{j}{N_r}+\ii\eta_r,
 &&j=0,\ldots,r-1,\\
 Y_j&=\Re\Xi_r(t_j).
\end{aligned}\label{gff:fp:nodes}\end{equation}
The nonzero imaginary part avoids every real pole class. It is a fixed
number, not a regulator to be sent to zero. Cosine orthogonality gives
\begin{equation}\boxed{\begin{aligned}
 c_{r,0}&=\frac1{N_r}\left(Y_0+2\sum_{j=1}^{r-1}Y_j\right),\\
 c_{r,k}&=\frac{2}{N_r\cosh(k/(4r))}
 \left[Y_0+2\sum_{j=1}^{r-1}Y_j
             \cos\frac{2\pi kj}{N_r}\right],\quad 1\leq k<r.
\end{aligned}}\label{gff:fp:cosineTransform}\end{equation}
Only finite sums appear here. For complex $q$, use all $2r-1$ complex
samples and the ordinary discrete Fourier transform, rather than using
complex conjugation to reconstruct half of them.

The residue functional can also be printed without differentiating
anything. For $1\leq h<r$ set
\begin{equation}\begin{aligned}
 L_h(q)&=(-1)^{h+1}\psi_{h-1}(1-q)-\psi_{h-1}(q),\\
 u_h(q)&=r\left[(1+(-1)^h)\psi_{h-1}(1)-L_h(q)\right].
\end{aligned}\label{gff:fp:logCoefficients}\end{equation}
These particular polygamma combinations are elementary reflection
values. If $c=\cot(\pi q)$, an implementation requiring no higher
polygamma calls may use
\begin{equation}\begin{aligned}
 \frac{u_1}{r\pi}&=-c,&
 \frac{u_2}{r\pi^2}&=c^2+\frac43,&
 \frac{u_3}{r\pi^3}&=-2c(1+c^2),\\
 \frac{u_4}{r\pi^4}&=\frac{32}{15}+8c^2+6c^4,&
 \frac{u_5}{r\pi^5}&=-8c(1+c^2)(2+3c^2),\\
 \frac{u_6}{r\pi^6}&=\frac{1024}{63}+136c^2+240c^4+120c^6.
\end{aligned}\label{gff:fp:elementaryWeights}\end{equation}
Only $u_1,\ldots,u_{r-1}$ are used. These follow from reflection and the
integer polygamma values, and introduce no additional period constants.

Define finite polynomials $E_m$, only for $0\leq m\leq6$, by
\begin{equation}
 E_0=1,\qquad E_m(x_1,\ldots,x_m)
 =\frac1m\sum_{h=1}^m\frac{x_h}{(h-1)!}
                         E_{m-h}(x_1,\ldots,x_{m-h}).
 \label{gff:fp:BellFinite}
\end{equation}
This is a finite polynomial definition, not a parameter derivative.
For orientation,
\begin{equation}\begin{aligned}
 E_2&=(x_1^2+x_2)/2,& E_3&=(x_1^3+3x_1x_2+x_3)/6,\\
720E_6={}&x_1^6+15x_1^4x_2+20x_1^3x_3+45x_1^2x_2^2+15x_1^2x_4\\
 &+60x_1x_2x_3+6x_1x_5+15x_2^3+15x_2x_4+10x_3^2+x_6.
\end{aligned}\label{gff:fp:BellExamples}\end{equation}
Write
\begin{equation}\begin{aligned}
 B_{r,k}(q)=\frac12\bigl[&e^{\ii\pi kq}
       E_{r-1}(u_1+2\pi\ii k,u_2,\ldots,u_{r-1})\\
 &+e^{-\ii\pi kq}
       E_{r-1}(u_1-2\pi\ii k,u_2,\ldots,u_{r-1})\bigr],\\
 \mathcal E_r(q;\boldsymbol Y)&=
 -\frac{\sum_{k=0}^{r-1}B_{r,k}(q)c_{r,k}}
             {(\pi\sin\pi q)^r}.
\end{aligned}\label{gff:fp:finiteFunctional}\end{equation}
Then the finite primitive formulas are
\begin{equation}\boxed{
 Q_d=\frac{\Gamma(2q)^3}{3}\mathcal E_3,
 \qquad P_4(d)=\Gamma(2q)^4\mathcal E_4,
 \qquad P_5(d)=\frac{\Gamma(2q)^5}{2}\mathcal E_7.
}\label{gff:fp:periodsFinite}\end{equation}
To check the sign and factors, close each Barnes contour to the right.
Its orientation is clockwise, so its value is minus the sum of right
residues. That sum is $\Res_{t=0}S_r(t)$. In
$D_q(t)^{-r}$ the leading coefficient is $(\pi\sin\pi q)^{-r}t^{-r}$;
the logarithmic coefficients of the analytic factor are exactly
\eqref{gff:fp:logCoefficients}. Expanding a cosine then gives
\eqref{gff:fp:finiteFunctional}. Residues are used in this proof only; the
final formula \eqref{gff:fp:periodsFinite} contains finite arithmetic and
fixed special-function values.

\subsection{Explicit one- and two-variable functions for every sample}
\label{gff:fp:arraysSection}
For finite upper/lower lists of parameter--integer-vector pairs use
\begin{equation}
 \mathfrak H[\mathsf U;\mathsf V;\boldsymbol z]
 =\sum_{\boldsymbol n\geq0}
 \frac{\prod_{(a,v)\in\mathsf U}(a)_{v\cdot\boldsymbol n}}
      {\prod_{(b,w)\in\mathsf V}(b)_{w\cdot\boldsymbol n}}
 \prod_i\frac{z_i^{n_i}}{n_i!}.
 \label{gff:fp:HornDefinition}
\end{equation}
This is the unweighted Horn convention already used in the note.
Signed integer-vector entries are permitted; $(a)_{-m}$ means
$\Gamma(a-m)/\Gamma(a)$, not a truncated product. Consecutive coefficient
ratios are rational in the indices. There are no index-dependent
polygammas in the following arrays. The function definition contains its
usual series, just as the definition of ${}_pF_q$ does; the master formula
has no additional outer infinite sum. This distinction is not a claim
that a two-variable function has become a one-variable function.

For two variables put $\boldsymbol e=(1,0)$,
$\boldsymbol f=(0,1)$, $\boldsymbol c=\boldsymbol e+\boldsymbol f$.
Repeated rows are denoted by a superscript $\times r$. Define
\begin{equation}\begin{aligned}
 A_r^0(t)&=a_r^0(t)\,
 {}_{2r+1}F_{2r}\!\left(
 \begin{matrix}1,(q+t)^{\times r},(1-q+t)^{\times r}\\
 (2q+t)^{\times r},(1+t)^{\times r}\end{matrix};1\right),\\
 a_r^0(t)&=\left[
 \frac{\Gamma(q+t)\Gamma(-t)}{\Gamma(2q+t)\Gamma(q-t)}\right]^r,\\
 A_r^+(t)&=a_r^+(t)\mathfrak H[\mathsf U_+;\mathsf V_+;(1,1)],\\
 a_r^+(t)&=4q\left[
 \frac{\Gamma(q+1+t)\Gamma(-t)}{\Gamma(2q+1+t)\Gamma(q-t)}\right]^r,\\
 C_r(t)&=a_r^c(t)\mathfrak H[\mathsf U_c;\mathsf V_c;(1,1)],\\
 a_r^c(t)&=4q\left[
 \frac{\Gamma(q+t)\Gamma(1-t)}{\Gamma(2q+t)\Gamma(1+q-t)}\right]^r.
\end{aligned}\label{gff:fp:wheelBlocks}\end{equation}
The complete two-variable arrays are
\begin{equation}\begin{aligned}
\mathsf U_+={}&\{(q+1+t,\boldsymbol c)^{\times r},
 (1,\boldsymbol e),(2q+1,\boldsymbol e),(4q-1,\boldsymbol e),
 (1,\boldsymbol f),(1-q+t,\boldsymbol f)^{\times r}\},\\
\mathsf V_+={}&\{(2q+1+t,\boldsymbol c)^{\times r},
 (2q,\boldsymbol e),(2,\boldsymbol e),(1+t,\boldsymbol f)^{\times r}\},\\
\mathsf U_c={}&\{(2q+1,\boldsymbol c),(4q-1,\boldsymbol c),
 (1,\boldsymbol e),(1-t,\boldsymbol e)^{\times r},
 (1,\boldsymbol f),(q+t,\boldsymbol f)^{\times r}\},\\
\mathsf V_c={}&\{(2q,\boldsymbol c),(2,\boldsymbol c),
 (1+q-t,\boldsymbol e)^{\times r},(2q+t,\boldsymbol f)^{\times r}\}.
\end{aligned}\label{gff:fp:wheelArrays}\end{equation}
For $r=3,4$ every sample is therefore
\begin{equation}\boxed{
 \Xi_r(t)=D_q(t)^r\bigl[
 A_r^0(t)+A_r^+(t)+A_r^0(1-q-t)+A_r^+(1-q-t)+C_r(t)\bigr].
}\label{gff:fp:wheelSample}\end{equation}
The right tail gives the first two functions, evenness maps the left
tail to $1-q-t$, and \eqref{gff:fp:middleFactor} gives the last function.
Equation \eqref{gff:fp:degeneracy} then proves every Pochhammer row in
\eqref{gff:fp:wheelArrays}. This is also a direct coefficient-by-coefficient
proof of the sample formula.

\subsubsection{The eight primitive sample blocks}
The $r=7$ arrays are obtained from the same printed blocks by finite row
insertions, as follows. Put
\begin{equation}
 v=q+2t,\quad \theta_0=2q,\quad\theta_1=2q-1,\qquad
 M_\theta(v)=\frac{\Gamma(\theta+v)\Gamma(\theta-v)}
                  {\Gamma(v)\Gamma(-v)}.
 \label{gff:fp:measureBlock}
\end{equation}
The identity underlying the two channels is
\begin{equation}
 \bigl[(\ell+p)^2-V^2\bigr]
 \frac{\Gamma(p+V)\Gamma(p-V)}{\Gamma(V)\Gamma(-V)}
 =M_{2q}(V)+\ell(\ell+4q-2)M_{2q-1}(V).
 \label{gff:fp:twoChannels}
\end{equation}
For clarity, a ``block'' below includes its scalar prefactor, both lists
and its argument vector.

Define a right insertion $\mathcal R_{\theta,w,h}$ on any block by
multiplying its prefactor by $M_\theta(w)$ and adjoining the rows
\begin{equation}
 \text{upper: }(\theta+w,h),(1+w,h),\qquad
 \text{lower: }(w,h),(1+w-\theta,h).
 \label{gff:fp:rightInsertion}
\end{equation}
It multiplies its coefficient by $M_\theta(w+h\cdot\boldsymbol n)$.
This follows by reflecting the negative Gamma argument and using that
$h\cdot\boldsymbol n$ is an integer. A middle insertion
$\mathcal C_{\theta,w,h}$ multiplies the prefactor by $M_\theta(w)$
and adjoins instead
\begin{equation}
 \text{upper: }(\theta+w,h),(\theta-w,-h),\qquad
 \text{lower: }(w,h),(-w,-h).
 \label{gff:fp:middleInsertion}
\end{equation}
This is the same identity before reflecting the index difference.

Start with $A_7^0,A_7^+,C_7$. Let $A^-$ be $A_7^+$ with a minus sign
in its prefactor and arguments $(-1,1)$; let $C^-$ be $C_7$ with a minus
sign and arguments $(-1,-1)$. Let $\widehat A^-$ be obtained from $A^-$
by multiplying the prefactor by $4q-1$, replacing its upper row
$(4q-1,\boldsymbol e)$ by $(4q,\boldsymbol e)$ and its lower row
$(2,\boldsymbol e)$ by $(1,\boldsymbol e)$. Define $\widehat C^-$
with the same replacements along $\boldsymbol c$. These replacements
supply the factor $\ell(\ell+4q-2)$ in the second channel of
\eqref{gff:fp:twoChannels}.

There are three right blocks and two middle blocks:
\begin{equation}\begin{aligned}
 Z(t)&=\mathcal R_{\theta_0,v,2} A_7^0(t),\\
 P_0(t)&=\mathcal R_{\theta_0,v+1,\boldsymbol e+2\boldsymbol f}A^-(t),&
 P_1(t)&=\mathcal R_{\theta_1,v+1,\boldsymbol e+2\boldsymbol f}\widehat A^-(t),\\
 C_0(t)&=\mathcal C_{\theta_0,v-1,\boldsymbol f-\boldsymbol e}C^-(t),&
 C_1(t)&=\mathcal C_{\theta_1,v-1,\boldsymbol f-\boldsymbol e}\widehat C^-(t).
\end{aligned}\label{gff:fp:primitiveBlocks}\end{equation}
The single-variable vector in the first line is the integer $2$.
Duplication of those step-two Pochhammers turns $Z(t)$ into an ordinary
${}_{19}F_{18}(1)$. All other blocks have only two indices.
The complete fixed sample is
\begin{equation}\boxed{\begin{aligned}
 \Xi_7(t)=D_q(t)^7\bigl[&Z(t)+P_0(t)+P_1(t)\\
 &+Z(1-q-t)+P_0(1-q-t)+P_1(1-q-t)+C_0(t)+C_1(t)\bigr].
\end{aligned}}\label{gff:fp:primitiveSample}\end{equation}
Equations \eqref{gff:fp:nodes}--\eqref{gff:fp:primitiveSample} specify every
parameter, argument and prefactor using finite lists; the JSON templates
in \path{lri_general_periodization_update/verification/general_periodization/arrays/}
are machine-readable copies, not missing definitions.

\paragraph{Convergence and continuation.}
The proofs above start at $1/2<q<1$. The right and middle Gamma-ratio
estimates give direct absolute convergence of the displayed blocks for
$q>1/(2r)$ in a wheel and $q>1/10$ in the primitive, always with $q<1$.
The endpoint in a positive-spin right block imposes the additional
condition $q(4-r)<1$ for a wheel and $q<1$ for the primitive; these are
automatic in the stated ranges. Thus the new $Q,P_4,P_5$ formulas are
direct endpoint-series formulas for $d>2/3,1/2,2/5$, respectively.
For smaller positive dimensions use their continued complete
combinations. Exceptional real parameter coincidences are dealt with
by the same continuation; the off-real sampling points themselves are
nonsingular. In particular $d=1,2,3$ are valid specializations, though the
retained five ordinary values remain a better presentation for $P_5(1)$.

\subsection{A stronger result at two dimensions: reuse the one-dimensional values}
\label{gff:fp:twoDimensional}
In this subsection $B_-,B_+$ and $X_0,X_+,X_-$ refer to the
\emph{retained one-dimensional} values, not to new functions of $d$:
\begin{equation}\begin{aligned}
 B_-&=\Gamma(\tfrac14)^4 F_-,&
 B_+&=\frac{64\pi^4}{\Gamma(\tfrac14)^4}F_+,\\
 X_0&=\mathcal Y_2,&
 X_+&=c_8^7\mathcal Y_3-s_8^7\mathcal Y_5,&
 X_-&=c_8^7\mathcal Y_1+s_8^7\mathcal Y_7.
\end{aligned}\label{gff:fp:inheritedValues}\end{equation}
For a self-contained dimension-two parameter dictionary, these are
\begin{equation}\begin{aligned}
 F_-&={}_4F_3\!\left(\begin{matrix}\frac14,\frac14,\frac12,\frac12\\
          \frac34,\frac34,1\end{matrix};1\right),&
 F_+&={}_4F_3\!\left(\begin{matrix}\frac34,\frac34,\frac12,\frac12\\
          \frac54,\frac54,1\end{matrix};1\right),\\
 c_8&=\cos(\pi/8),&s_8&=\sin(\pi/8),
\end{aligned}\label{gff:fp:FpmDictionary}\end{equation}
\begin{equation}\begin{aligned}
\mathcal Y_r={}&\frac r8
 \left[\frac{\Gamma((r+2)/8)}{\Gamma((r+6)/8)}\right]^7
 {}_9F_8\!\left(
 \begin{matrix}1,1+\frac r8,((r+2)/8)^{\times7}\\
 \frac r8,((r+6)/8)^{\times7}\end{matrix};-1\right),\\
&r\in\{1,2,3,5,7\}.
\end{aligned}\label{gff:fp:YDictionary}\end{equation}
At $r=2$ one upper and lower parameter $1$ cancel, giving an ordinary
${}_8F_7$; all parameter lists are fixed and positive rational.
The new identities are
\begin{equation}\boxed{
 P_4(2)=\frac{B_-^2-B_+^2}{4\pi^3}.
}\label{gff:fp:P4two}\end{equation}
\begin{equation}\boxed{\begin{aligned}
P_5(2)=\frac8\pi\bigl[&25X_0^2-28\sqrt2\,X_0(X_++X_-)\\
 &+18(X_+^2+X_-^2)+28X_+X_-\bigr].
\end{aligned}}\label{gff:fp:P5two}\end{equation}
Thus no new hypergeometric value is needed for either two-dimensional
primitive. The four-loop result uses the same two ${}_4F_3(1)$ values;
the five-loop result uses exactly the preferred four ${}_9F_8(-1)$ and
one ${}_8F_7(-1)$ with positive rational parameters. The formulas are
quadratic rather than linear. By themselves they do not evaluate the
two-dimensional nonprimitive sector; that evaluation is completed in
Section~\ref{gff:gc:d2closedForms}, with the proofs in Sections~\ref{gff:gc:L2proof}
and \ref{gff:gc:K2proof}.

\subsubsection{Factorization of the wheel lattice}
At $q=1/2$ the multiplicities are $d_0=1$, $d_{\ell>0}=2$.
In a right half-lattice put $k=\ell+n$, $j=n$. Gamma reflection gives
\[
 g_\ell(a_\ell+n+t)=
 -\cot(\pi t)\,
 \frac{\Gamma(k+\frac12+t)}{\Gamma(k+1+t)}
 \frac{\Gamma(j+\frac12+t)}{\Gamma(j+1+t)}.
\]
The diagonal has multiplicity one and the two off-diagonal half-lattices
have multiplicity two. They therefore form a full square lattice.
Define the meromorphic bilateral series and its regularization
\begin{equation}\begin{aligned}
 B_r(t)&=\sum_{n\in\mathbb Z}
      \left[\frac{\Gamma(n+\frac12+t)}{\Gamma(n+1+t)}\right]^r,\\
 \Omega_r(t)&=\cos^r(\pi t)B_r(t).
\end{aligned}\label{gff:fp:OmegaWheel}\end{equation}
The index sums are absolute for $r>2$. Direct reflection or the square
lattice just described gives
\begin{equation}
 \Xi_r(t)=(-1)^r\Omega_r(t)^2\qquad(d=2).
 \label{gff:fp:wheelSquare}\end{equation}
For $r=4$, the growth and reflection argument behind \eqref{gff:fp:finiteFourier} now yields
only two modes:
\begin{equation}
 \Omega_4(t)=a+b\cos\bigl(2\pi(t+\tfrac14)\bigr).
 \label{gff:fp:Omega4Modes}\end{equation}
There is no fitting in this statement: $\Omega_4$ is entire, one-periodic,
even about $-1/4$, and has growth $O(e^{4\pi|\Im t|}|\Im t|^{-1})$,
which excludes the boundary second mode.

Let $M_{\rm e}=\int_0^1\mathbf K(k)^2\,dk$ and
$N_{\rm e}=\int_0^1\mathbf K(k)\mathbf K(k')\,dk$, with the modulus
convention of Section~\ref{gff:app:LreductionProof}. The established elliptic
identities give
\begin{equation}
 M_{\rm e}=\frac{B_-+B_+}{64},\qquad
 N_{\rm e}=\frac{B_--B_+}{32}.
 \label{gff:fp:MN}\end{equation}
Here is a short analytic determination of both coefficients in
\eqref{gff:fp:Omega4Modes}. At $t=0$ the reflected half of $\Omega_4$ vanishes
to fourth order. Hence
\[
 a=\pi^2\,{}_4F_3\left(\begin{matrix}\frac12,\frac12,\frac12,\frac12\\1,1,1\end{matrix};1\right)
   =\frac{8N_{\rm e}}\pi.
\]
For $N(\nu)=\int_0^1 k^\nu\mathbf K(k)\mathbf K(k')\,dk$, the mixed
moment beta integral of \cite{gff:Wan} gives
\begin{equation}
 N(\nu)=\frac{\pi^2}{8}\sum_{n\geq0}
 \frac{(\frac12)_n^2}{(n!)^2}
 \frac{\Gamma(n+(\nu+1)/2)^2}{\Gamma(n+(\nu+2)/2)^2}.
 \label{gff:fp:mixedMomentProof}
\end{equation}
Differentiation in this convergent \emph{proof} representation gives
$\Omega_4'(0)=32N'(0)/\pi$. The already proved Stieltjes--Landen
identity \eqref{gff:lp:eq:mixedlog} says $N'(0)=-\pi M_{\rm e}/2$.
As $\Omega_4'(0)=-2\pi b$, it follows that
\begin{equation}
 a=\frac{8N_{\rm e}}\pi,\qquad b=\frac{8M_{\rm e}}\pi.
 \label{gff:fp:abWheel}\end{equation}
No derivative remains in these coefficients.
Substitute $\Xi_4=(a+b\cos(2\pi(t+1/4)))^2$ into
\eqref{gff:fp:finiteFunctional}. At $q=1/2$ the only contributing coefficient
is $c_{4,1}=2ab$, and $P_4(2)=8ab/\pi=512M_{\rm e}N_{\rm e}/\pi^3$.
Equation \eqref{gff:fp:P4two} follows immediately.

\subsubsection{A second determination of \texorpdfstring{$b$}{b}, without the Stieltjes identity}
\label{gff:mgp:P4second}
The value $a=8N_{\rm e}/\pi$ above uses only Euler integration. The value of $b$ above
rests on the mixed-moment series \eqref{gff:fp:mixedMomentProof} and on \eqref{gff:lp:eq:mixedlog},
whose proof uses the Stieltjes identity \eqref{gff:lp:eq:stieltjes} of \cite{gff:Wan}. The
following derivation uses neither. It was written by the Claude Code session that wrote the
third and fourth proofs of \eqref{gff:gc:L2closed} in Section~\ref{gff:gc:L2proof} (package
\path{LRI_C5_d2_simplify/}, folder \path{proof/}).

Put $c_n=(\frac12)_n/n!$, so that $\Gamma(n+\frac12)/\Gamma(n+1)=\sqrt\pi\,c_n$. At $t=0$
the reflected half of $\Omega_4$ vanishes to fourth order and $\cos^4\pi t$ has zero
derivative, so
\[
 -2\pi b=\Omega_4'(0)=4\pi^2\sum_{n\geq0}c_n^4\bigl[\psi_0(n+\tfrac12)-\psi_0(n+1)\bigr].
\]
Item (ii) of the pairing proposition \eqref{gff:mgp:pairings} below evaluates this single
sum by Parseval's identity on the unit circle. It uses only the generating function
\eqref{gff:mgp:gfDisk} (first line) and the circle values \eqref{gff:mgp:circle}, and gives
$b=4V_0/\pi$, with $V_0=\int_0^1\mathbf K(k)\mathbf K(k')\,dk/k'$ as defined before
\eqref{gff:lp:eq:fold}. The imaginary part of the upper-semicircle contour that proves
\eqref{gff:lp:eq:V} gives $V_0=2M_{\rm e}$; that contour uses the boundary-value Landen relation
and not \eqref{gff:lp:eq:stieltjes}. Hence
$b=8M_{\rm e}/\pi$, which is \eqref{gff:fp:abWheel}. The Fourier representation of
Section~\ref{gff:mgp:fourier} gives the same value independently, $b=v_1(\frac12)/(4\pi)=4V_0/\pi$.
For $a$ the same two routes give $a=4U_0/\pi$, and $U_0=2N_{\rm e}$ is proved without
\eqref{gff:lp:eq:stieltjes} in Section~\ref{gff:mgp:halfStrip}. Since \eqref{gff:fp:MN} uses only
$V_0=2M_{\rm e}$ and the real Landen substitution (the reduction to the two ${}_4F_3$
values after \eqref{gff:lp:eq:MN}), \eqref{gff:fp:P4two} is proved without the Stieltjes
identity and without the mixed-moment series. As a check, not a proof step, the
$70$-digit single sums give $8ab/\pi-(B_-^2-B_+^2)/(4\pi^3)=-1.1\times10^{-67}$
(\path{LRI_C5_d2_simplify/proof/synth/synth_checks.log}).


\subsubsection{The five-loop square and the retained five values}
Put
\begin{equation}\begin{aligned}
 G(s)&=(2s+\tfrac12)
           \left[\frac{\Gamma(s+\frac12)}{\Gamma(s+1)}\right]^7,\\
 \Omega_{\rm a}(t)&=\cos^7(\pi t)
                      \sum_{n\in\mathbb Z}(-1)^nG(n+t).
\end{aligned}\label{gff:fp:OmegaAlternating}\end{equation}
At $q=1/2$ the measure in \eqref{gff:fp:primitiveStart} is
$\ell^2-4s^2$. On the square lattice it becomes
$-(2k+\frac12+2t)(2j+\frac12+2t)$ and
$(-1)^\ell=(-1)^{k-j}$. Consequently
\begin{equation}
 \Xi_7(t)=\Omega_{\rm a}(t)^2\qquad(d=2).
 \label{gff:fp:primitiveSquare}\end{equation}
The sign is plus: the negative quadratic measure and the odd seventh
power of the Gamma-reflection sign cancel.

Most usefully, $G(s)=2f(s)$ for the \emph{one-dimensional} function in
\eqref{gff:fx:eq:fA}. Thus
\begin{equation}
 \Omega_{\rm a}(t)=2\Xi^{(1)}(t)
 =2\bigl[\alpha\cos2\pi t+\beta\sin4\pi t+\gamma\cos6\pi t\bigr],
 \label{gff:fp:oneDimensionalSquare}
\end{equation}
where $\alpha,\beta,\gamma$ are exactly those in
\eqref{gff:fx:eq:abc}. Applying the finite residue polynomial gives
\begin{equation}
 P_5(2)=\frac8\pi
 \left(\alpha^2-14\alpha\gamma+8\beta^2+113\gamma^2\right).
 \label{gff:fp:P5abc}
\end{equation}
One convenient exact algebra check is that at $q=1/2$
\[
 \bigl(B_{7,0},B_{7,1},\ldots,B_{7,6}\bigr)
       =\pi^6(20,0,28,0,84,0,924).
\]
Expanding the square in \eqref{gff:fp:oneDimensionalSquare} and using the
normalization $P_5=\mathcal E_7/2$ gives \eqref{gff:fp:P5abc}; inserting
\eqref{gff:fx:eq:abc} gives \eqref{gff:fp:P5two}. This proves the quadratic
factorization without a numerical integer-relation search.

\subsection{Numerical checks of the finite periodization}\label{gff:mgc:fpChecks}
These checks were run by OpenAI GPT-6 astra in the calculation that introduced the periodization
(package \path{lri_general_periodization_update/}; scripts in \path{code/general_periodization/},
records \path{wheel_periodization.json}, \path{horn_template_checks.json},
\path{p5_periodization.json}, \path{d2_factorization.json}, \path{d2_control.json},
\path{d2_inherited_periods.json} and \path{exact_algebra.json} in
\path{verification/general_periodization/}). They are numerical consistency checks, not interval
bounds.

\paragraph{Function counts.}
The finite evaluation \eqref{gff:fp:periodsFinite} with the arrays of Section~\ref{gff:fp:arraysSection}
needs far fewer special-function values than the unweighted Horn arrays obtained by applying the rules
of Section~\ref{gff:gc:HornRules} to the wheel series \eqref{gff:gc:wheelSeries} and to the derivative-free
double-residue series \eqref{gff:gc:P5series} of $P_5(d)$:
\begin{center}\small
\begin{tabular}{@{}lrrrr@{}}\toprule
 & Horn values & variables & periodization values, at most & variables\\\midrule
$Q_d$ &20&4&15&2\\
$P_4(d)$&56&5&20&2\\
$P_5(d)$&15060&8&56&2\\\bottomrule
\end{tabular}
\end{center}
For $P_5(d)$ the Horn count is $6024$ for the first channel of \eqref{gff:gc:P5series} and $9036$ for the
second. The sampling counts are the five blocks of \eqref{gff:fp:wheelSample} at each of the $r$ nodes of
\eqref{gff:fp:nodes} ($r=3,4$) and the eight blocks of \eqref{gff:fp:primitiveSample} at the seven nodes for $r=7$.
The counts include the one-variable subblocks, and taking the real part of a complex value counts as one
evaluation. They are upper bounds and assert neither minimality nor independence. The gain is the
elimination of the auxiliary indices that encoded residue polygammas, not a proof that the
multivariable functions reduce to ordinary ones.

\paragraph{Fixed-spin controls.}
At 35 working digits, fixed-spin wheel samples were evaluated as ordinary hypergeometric functions and
compared with real-axis Barnes integrals. The cases $(q,r,\ell)=(0.6,3,0)$, $(0.7,3,2)$, $(0.6,4,1)$ and
$(0.25,4,0)$ differ by less than $3\times10^{-30}$, and further off-grid samples check the finite Fourier
polynomial \eqref{gff:fp:finiteFourier}. These test building blocks of the theorem, not a rederived graph
sum.

\paragraph{Parameter arrays and finite transforms.}
At 65 working digits, 44 coefficients of the Horn blocks, including the signed-index middle blocks,
were compared with the Gamma kernels at $q=1/4,1/2,0.63,3/4$; the largest normalized residual is
$3.81\times10^{-65}$. The cosine transforms \eqref{gff:fp:cosineTransform} were checked on every basis mode
for $r=3,4,7$, with residuals below $4.75\times10^{-66}$. The exact derivation of the arrays is
\eqref{gff:fp:reflectedRight}--\eqref{gff:fp:middleFactor} and the finite row rules; these are coefficient
checks.

\paragraph{Noninteger dimensions.}
The spin-summed finite-sampling expression for $P_5(d)$ was compared with spectral quadrature of
\eqref{gff:gd:spectral}, with 640 nodes and 160 spin terms; the sampling computation used 128 spin terms,
Euler acceleration and doubled radial cutoffs. The difference is $-1.39\times10^{-9}$ at $d=2.4$ and
$+7.38\times10^{-8}$ at $d=3$, and the off-grid periodization residuals are below $2.3\times10^{-11}$.
These runs use double-precision Gamma kernels inside wider accumulations, so their residuals are not
promoted to high precision because the finite weights were computed at 50 digits. The five-loop residue
functional is sensitive to cancellation. Imaginary-axis Chebyshev nodes far from the real axis are
particularly ill conditioned; the grid \eqref{gff:fp:nodes} controls this amplification much better, but it
is not an error certificate.

\paragraph{Two dimensions.}
The square factorizations of Section~\ref{gff:fp:twoDimensional} were evaluated first at 55 digits, using
Gamma-ratio series with a Hurwitz-zeta asymptotic tail and a finite Euler transform for the alternating
kernel. A 70-digit rerun changed the direct cutoff from 96 to 144, the tail order from 44 to 58 and the
Euler settings from $(20,260)$ to $(32,340)$, and used the grid \eqref{gff:fp:nodes}. The $Q_2$ value agrees
with $\Gamma(1/4)^8/(24\pi^4)$ to about 67 decimal places, and $P_4(2)$ and $P_5(2)$ are the values
quoted in Section~\ref{gff:gc:d2closedForms}.
The retained values $F_\pm$ and the five $\mathcal Y_r$ were then recomputed with literal ordinary
hypergeometric calls at 60 working digits. The formulas \eqref{gff:fp:P4two} and \eqref{gff:fp:P5two} differ
from the 70-digit square-factorization values by $-5.97\times10^{-59}$ and $+1.59\times10^{-58}$. The
quadratic change of basis from \eqref{gff:fp:P5abc} to \eqref{gff:fp:P5two} was checked by exact symbolic
algebra. No stored period value was an input to either route.

\paragraph{What the finite Fourier theorem does not give.}
The theorem \eqref{gff:fp:finiteFourier} concerns the Barnes-shift variable $t$. It does not imply finite
dependence on an independent rim exponent $b$, so it cannot replace the first and second $b$
derivatives that define $\mathcal L_d$ and $Z_d$ in \eqref{gff:gs:wheelJets} by a finite sampling grid in
$b$; that would be an unjustified interpolation. Sections~\ref{gff:gdr:section} and~\ref{gff:gdk:section} treat
the rim exponent by a separate periodization (Section~\ref{gff:gdr:periodization}) and keep low-order
parameter derivatives.

\subsection{\texorpdfstring{Proof of the closed form of $\mathcal L_2$}{Proof of the closed form of L2}}
\label{gff:gc:L2proof}
We prove \eqref{gff:gc:L2closed} for the series \eqref{gff:gs:Lseries}. The rim exponent of the
wheel lattice of Section~\ref{gff:fp:twoDimensional} is deformed away from one, the
two-mode statement \eqref{gff:fp:Omega4Modes} is proved for every nearby rim exponent, and
the two derivative sums that appear are evaluated with the elliptic lemmas of
Section~\ref{gff:app:LreductionProof} and one further contour integral. This is the
one-dimensional contour proof of Section~\ref{gff:app:LreductionProof}, transferred to the
square lattice. Throughout, $d=2$, $q=\frac12$, $c_k=(\frac12)_k/k!$ (not the constants
$c_j$ of \eqref{gff:gc:dictionary}), and
\[
 \alpha_k=\psi_0(k+\tfrac12)-\psi_0(k+1),\qquad
 \beta_k=\psi_0(k+\tfrac12)-\psi_0(\tfrac12),
\]
that is, $\alpha_k=\Delta_0(\frac12,1;k)$ and $\beta_k=H(\frac12;k)$ in the notation of
\eqref{gff:gc:dictionary}. The letters $a,b$ denote the coefficients of
\eqref{gff:fp:Omega4Modes}. The functions $f,g$, the weight $w$ and the moments $U_0,R$
are those of Section~\ref{gff:app:LreductionProof}, whose $M,N$ are $M_{\rm e},N_{\rm e}$
of \eqref{gff:fp:MN}.

The first proof (Sections~\ref{gff:gc:L2rimDerivative}--\ref{gff:gc:L2conclusion}) and the second
(Section~\ref{gff:gc:L2second}) were written by a Claude Code session with prover and reviewer
subagents (package
\path{LRI_C5_d2_proof/}, folders \path{L2_singlesums/} and \path{L2_structural/}). Another Claude
Code session, with prover, verifier and synthesizer subagents, wrote the half-strip evaluation
of Section~\ref{gff:mgp:halfStrip}, the third and fourth proofs (Sections~\ref{gff:mgp:parseval}
and~\ref{gff:mgp:fourier}) and the remark on the integer-relation form of the collapse inside the
first proof (package \path{LRI_C5_d2_simplify/}, folder \path{proof/}; the third proof is in its
subfolder \path{L2-elliptic/}, the half-strip evaluation and the fourth proof in
\path{L2-fourier/}). That session had first found \eqref{gff:gc:L2closed} and the relations
\eqref{gff:gc:L2singles} numerically, in the top level of the same package.

\subsubsection{The series as a rim derivative}\label{gff:gc:L2rimDerivative}
Let $I=[\frac25,\frac35]$. For $x\in I$ and $m,n,k\geq0$ put
\begin{equation}\begin{aligned}
 \varphi_x(s)&=\frac{\Gamma(s+x)}{\Gamma(s+1)},\qquad
 N_x=\frac{\Gamma(1-x)}{\Gamma(x)},\qquad
 c_k(x)=\frac{(x)_k}{k!},\qquad T_{mn}(x)=c_m(x)c_n(x),\\
 \lambda_j(x)&=\Delta_{j-1}(x,1;m)+\Delta_{j-1}(x,1;n)+\mathsf c_j(x),\qquad
 \sum_{j\geq1}\mathsf c_j(x)\frac{u^j}{j!}=\log\frac{\pi u}{\sin\pi u}
 +\log\frac{\sin\pi(x+u)}{\sin\pi x},
\end{aligned}\label{gff:gc:L2rim}\end{equation}
so that $\mathsf c_1=\pi\cot\pi x$, $\mathsf c_2=\frac{\pi^2}3-\frac{\pi^2}{\sin^2\pi x}$ and
$\mathsf c_3=\frac{2\pi^3\cos\pi x}{\sin^3\pi x}$; at $x=\frac12$ these are the constants
$c_j$ of \eqref{gff:gc:dictionary}. Let $F_4(x)$ be the polynomial $F_4$ of
\eqref{gff:gs:Fpolys} evaluated at these $\lambda_j(x)$.

\paragraph{One summand as a residue.}
Gamma reflection \cite{gff:DLMF5} gives $N_x\sin(\pi x)\varphi_x(m)\varphi_x(n)/\pi=T_{mn}(x)$,
and $\partial_u^j\log\varphi_x(k+u)|_{u=0}=\Delta_{j-1}(x,1;k)$. Hence, as Laurent
expansions at $u=0$,
\begin{equation}
 N_x\frac{\sin\pi(x+u)}{\sin\pi u}\,\varphi_x(m+u)\varphi_x(n+u)
 =\frac{T_{mn}(x)}{u}\exp\Bigl(\sum_{j\geq1}\lambda_j(x)\frac{u^j}{j!}\Bigr).
 \label{gff:gc:L2local}
\end{equation}
Since $F_4$ is the complete Bell coefficient $[u^3]\exp(4\sum_j\lambda_ju^j/j!)$, the
residue at $u=0$ of the fourth power of the left side of \eqref{gff:gc:L2local} is
$T_{mn}^4F_4$. Equation \eqref{gff:gc:L2local} is the two-dimensional form of
\eqref{gff:gs:localGamma} for rim exponent $2x$.

\paragraph{Identification at $x=\frac12$.}
Put $m=\ell+n$. Then $T_{mn}(\frac12)=T_{\ell n}$ of \eqref{gff:gs:seriesBase}, and the
$\lambda_j(\frac12)$, $j\leq3$, are the $\lambda_j$ of \eqref{gff:gs:lambda}, since
$\mathsf c_1(\frac12)=\mathsf c_3(\frac12)=0$ and
$\mathsf c_2(\frac12)=-\frac{2\pi^2}3=2[\psi_1(1)-\psi_1(\frac12)]$.
Moreover $\frac12\partial_x\log T_{mn}|_{x=1/2}=\mu_1$ and
$\frac12\partial_x\lambda_j|_{x=1/2}=\eta_j$ of \eqref{gff:gs:etaTheta}, because
$\partial_x\Delta_{j-1}(x,1;k)=\psi_j(k+x)$ and
$\frac12\mathsf c_j'(\frac12)=(-\frac{\pi^2}2,0,-\pi^4)=\frac12((-1)^j-1)\psi_j(\frac12)$. As
$F_{4,1}=\sum_j(\partial F_4/\partial\lambda_j)\,\eta_j$, the summand of \eqref{gff:gs:Lseries}
is
\[
 T_{\ell n}^4\bigl(4\mu_1F_4+F_{4,1}\bigr)
 =\tfrac12\,\partial_x\bigl[T_{mn}(x)^4F_4(x)\bigr]_{x=1/2}.
\]
These identifications are exact algebra in independent polygamma symbols (machine
checked in \path{LRI_C5_d2_proof/L2_singlesums/verify_symbolic.py},
\S\S1--3 and 5).

\paragraph{The square lattice.}
For $x\in I$ and $k\geq0$, Stirling's formula \cite{gff:DLMFStirling} gives
$0<c_k(x)\leq C(k+1)^{x-1}$. The quantities $\Delta_{j-1}(x,1;k)$, $\psi_j(k+x)$ ($j\geq1$),
$\mathsf c_j$ and $\mathsf c_j'$ are bounded, and $|\psi_0(k+x)-\psi_0(x)|\leq C\log(k+2)$. Hence
$T_{mn}^4F_4$ and its $x$-derivative are bounded by
$Cc_m(x)^4c_n(x)^4\log(m+2)\log(n+2)$, which is summable over $\mathbb N_0^2$ uniformly on
$I$ because $4(x-1)\leq-\frac85$. The summands are symmetric in $m,n$, and at $d=2$ the
multiplicities of \eqref{gff:gs:spectrum} are $d_0=1$ and $d_{\ell\geq1}=2$. Therefore
$\sum_\ell d_\ell\sum_ny_{\ell+n,n}=\sum_{m,n\geq0}y_{mn}$ for every symmetric absolutely
summable array $(y_{mn})$. Define
\[
 \mathcal W(x)=-\sum_{m,n\geq0}T_{mn}(x)^4F_4(x),\qquad x\in I .
\]
By \eqref{gff:gs:QPseries}, $\mathcal W(\frac12)=P_4(2)$. By \eqref{gff:gs:Lseries} and
termwise differentiation (Weierstrass M-test),
\begin{equation}
 \LL_2=\tfrac12\,\mathcal W'(\tfrac12).
 \label{gff:gc:L2derivative}
\end{equation}
The function $\mathcal W(x)$ equals the wheel $W_4(2;2x)$ of \eqref{gff:gs:balancedWheel} (see
the Status paragraph of Section~\ref{gff:gc:d2closedForms}); the proof below does not use
this.

\paragraph{Residue form.}
Put $\mathcal H_x(u)=\sum_{k\geq0}\varphi_x(k+u)^4$. On the circle $|u|=\frac1{10}$ we have
$|\varphi_x(k+u)|\leq C(k+1)^{x-1}$ uniformly, so the double series of fourth powers of
the left side of \eqref{gff:gc:L2local} converges uniformly there, to
$N_x^4\sin^4\pi(x+u)\,\mathcal H_x(u)^2/\sin^4\pi u$. Integrating term by term over the
circle,
\begin{equation}
 \mathcal W(x)=-N_x^4\Res_{u=0}
 \frac{\sin^4\pi(x+u)\,\mathcal H_x(u)^2}{\sin^4\pi u}.
 \label{gff:gc:L2residueForm}
\end{equation}

\subsubsection{Two Fourier modes for every rim exponent}
For $x\in I$ define
\[
 \Omega(t;x)=\sin^4\pi(t+x)\,\mathcal H_x(t)
 +\sin^4(\pi t)\sum_{j\geq0}\Bigl[\frac{\Gamma(j+1-t)}{\Gamma(j+2-x-t)}\Bigr]^4 .
\]
At $x=\frac12$ this is $\Omega_4(t)$ of \eqref{gff:fp:OmegaWheel}. We claim
\begin{equation}
 N_x^2\,\Omega(t;x)=\mathcal A(x)+\mathcal B(x)\cos2\pi\bigl(t+\tfrac x2\bigr),
 \qquad \mathcal A(\tfrac12)=a,\quad\mathcal B(\tfrac12)=b,
 \label{gff:gc:L2modes}
\end{equation}
with $\mathcal A,\mathcal B$ real analytic on the interior of $I$.

Gamma reflection gives, for $j\geq0$ and $m\in\mathbb Z$,
\begin{gather*}
 \sin^4(\pi t)\Bigl[\frac{\Gamma(j+1-t)}{\Gamma(j+2-x-t)}\Bigr]^4
 =\sin^4\pi(t+x)\,\varphi_x(t-j-1)^4,\\
 \sin^4(\pi t)\,\varphi_x(-m-x-t)^4=\sin^4\pi(t+x)\,\varphi_x(m+t)^4 .
\end{gather*}
By the first identity, $\Omega(t;x)=\sin^4\pi(t+x)\sum_{n\in\mathbb Z}\varphi_x(n+t)^4$
away from the poles, which is one-periodic. By the second, with $n\mapsto-n$,
$\Omega(-x-t;x)=\Omega(t;x)$. The two series in the definition converge locally
uniformly away from $t\in-x-\mathbb N_0$ and $t\in1+\mathbb N_0$, respectively. Their
poles there have order at most four and are canceled by the zeros of
$\sin^4\pi(t+x)$ and $\sin^4\pi t$, so $\Omega$ is entire. For $\Re z\geq0$, $|z|\geq1$
and $x\in I$, Stirling's formula with its error bound in the closed right half-plane
\cite{gff:DLMFStirling} gives, uniformly in $x$,
\[
 \Bigl|\frac{\Gamma(z+x)}{\Gamma(z+1)}\Bigr|\leq C|z|^{x-1},\qquad
 \Bigl|\frac{\Gamma(z)}{\Gamma(z+1-x)}\Bigr|\leq C|z|^{x-1}.
\]
With $z=n+t$ and $z=j+1-t$, for $t=\sigma+\ii y$, $0\leq\sigma\leq1$, $|y|\geq1$, this
gives
$|\Omega(t;x)|\leq Ce^{4\pi|y|}\sum_{n\geq0}(n^2+y^2)^{2x-2}\leq C'e^{4\pi|y|}|y|^{4x-3}$,
and $4x-3<0$. The argument of \eqref{gff:fp:finiteFourier} now applies: shifting the
contour of the $k$-th Fourier coefficient to $\Im t=\mp Y$ removes every mode with
$|k|\geq2$, including $|k|=2$ because of the decaying factor $|y|^{4x-3}$. The
reflection $\Omega(-x-t)=\Omega(t)$ gives $\omega_1=e^{2\pi\ii x}\omega_{-1}$ for the two
remaining nonconstant coefficients, which yields the form \eqref{gff:gc:L2modes}, and
$\mathcal A,\mathcal B$ are real because $\Omega$ is real on the real axis. Finally
$\mathcal A(x)=N_x^2\int_0^1\Omega(\sigma+\ii;x)\,d\sigma$ and
$\mathcal B(x)=2e^{\ii\pi x}N_x^2\int_0^1\Omega(\sigma+\ii;x)e^{2\pi\ii(\sigma+\ii)}d\sigma$.
On $\Im t=1$ both series converge locally uniformly for complex $x$ near $I$, so
$\mathcal A$ and $\mathcal B$ are real analytic. At $x=\frac12$, $N_{1/2}=1$ and
\eqref{gff:gc:L2modes} is \eqref{gff:fp:Omega4Modes}.

\subsubsection{The residue functional and an exact reduction}
By the definition of $\Omega$,
$\Omega(u;x)-\sin^4\pi(x+u)\mathcal H_x(u)$ equals $\sin^4(\pi u)$ times a function
analytic for $|u|<1$, and $\sin\pi(x+u)\neq0$ near $u=0$. In \eqref{gff:gc:L2residueForm}
the factor $N_x^2\sin^4\pi(x+u)\mathcal H_x(u)$ may therefore be replaced by
\eqref{gff:gc:L2modes}, since the difference contributes only terms regular at $u=0$. Hence
\begin{equation}\begin{aligned}
 \mathcal W(x)&=-\bigl[\rho_0\mathcal A^2+2\rho_1\mathcal A\mathcal B+\rho_2\mathcal B^2\bigr],
 \qquad
 \rho_p(x)=\Res_{u=0}\frac{\cos^p(2\pi u+\pi x)}{\sin^4\pi u\,\sin^4\pi(u+x)},\\
 \rho_0&=-\frac{4(\cos2\pi x+4)\cos\pi x}{\pi\sin^7\pi x},\qquad
 \rho_1=\frac{4(4\sin^2\pi x-5)}{\pi\sin^7\pi x},\qquad
 \rho_2=\frac{4(\sin^2\pi x-5)\cos\pi x}{\pi\sin^7\pi x}.
\end{aligned}\label{gff:gc:L2residue}\end{equation}
At $x=\frac12$ one has $\rho_0=\rho_2=0$, $\rho_1=-\frac4\pi$, $\rho_0'=12$, $\rho_1'=0$ and
$\rho_2'=16$. So $\mathcal W(\frac12)=8ab/\pi$, the value of $P_4(2)$ found after
\eqref{gff:fp:abWheel}, and by
\eqref{gff:gc:L2derivative}
\[
 \LL_2=-6a^2-8b^2+\frac4\pi\bigl[\mathcal A'(\tfrac12)\,b+a\,\mathcal B'(\tfrac12)\bigr].
\]

At $t=0$ the second series in $\Omega$ carries the factor $\sin^4\pi t$. Evaluating
\eqref{gff:gc:L2modes} and its $t$-derivative at $t=0$ gives, for $x\in I$,
\begin{gather*}
 N_x^2\sin^4(\pi x)\,\mathcal H_x(0)=\mathcal A+\mathcal B\cos\pi x,\\
 N_x^2\bigl[4\pi\sin^3\pi x\cos\pi x\,\mathcal H_x(0)+\sin^4\pi x\,\partial_u\mathcal H_x(0)\bigr]
 =-2\pi\mathcal B\sin\pi x .
\end{gather*}
The series $\mathcal H$ is jointly analytic in $(u,x)$ near $(0,\frac12)$ and may be
differentiated termwise. At $x=\frac12$ we have $N_{1/2}=1$,
$(N_x^2)'=-4\psi_0(\frac12)$, $\varphi_{1/2}(k)=\sqrt\pi\,c_k$ and
\[\begin{aligned}
 \mathcal H&=\pi^2\sum_kc_k^4,&
 \partial_u\mathcal H&=4\pi^2\sum_kc_k^4\alpha_k,\\
 \partial_x\mathcal H&=4\pi^2\sum_kc_k^4\psi_0(k+\tfrac12),&
 \partial_x\partial_u\mathcal H&=4\pi^2\sum_kc_k^4\bigl[4\psi_0(k+\tfrac12)\alpha_k+\psi_1(k+\tfrac12)\bigr].
\end{aligned}\]
The two relations and their $x$-derivatives at $\frac12$ then give
$a=\pi^2\sum_kc_k^4$ and $b=-2\pi\sum_kc_k^4\alpha_k$; the terms in $\psi_0(\frac12)$
cancel, and
\begin{equation}\begin{gathered}
 \mathcal A'(\tfrac12)=\pi b+4\pi^2\mathcal S_1,\qquad
 \mathcal B'(\tfrac12)=2\pi a-2\pi\mathcal S_2,\\
 \mathcal S_1=\sum_{k\geq0}c_k^4\beta_k,\qquad
 \mathcal S_2=\sum_{k\geq0}c_k^4\bigl[4\alpha_k\beta_k+\psi_1(k+\tfrac12)\bigr].
\end{gathered}\label{gff:gc:L2taylor}\end{equation}
All four sums converge absolutely, since $c_k^4\sim(\pi k)^{-2}$. Consequently
\begin{equation}\boxed{
 \LL_2=2a^2-4b^2+16\pi b\,\mathcal S_1-8a\,\mathcal S_2,\qquad P_4(2)=\frac{8ab}\pi .
}\label{gff:gc:L2reduction}\end{equation}
This is an exact identity; it uses no relation among the single sums.

\paragraph{Remark (the integer-relation form of the collapse).}
The integer-relation search that first found \eqref{gff:gc:L2closed} returned, in the present
normalization, $\LL_2=3a^2-2b^2+8\pi b\,\mathcal S_1-12a\,\mathcal S_2$. It differs from
\eqref{gff:gc:L2reduction} by $a^2+2b^2-8\pi b\,\mathcal S_1-4a\,\mathcal S_2$, which vanishes
only by the single-sum relations \eqref{gff:gc:L2singles} below. The four numbers $a^2$, $b^2$,
$\pi b\,\mathcal S_1$ and $a\,\mathcal S_2$ therefore satisfy a $\mathbb Z$-linear relation,
and the searched form is one member of a one-parameter family of representations. Only
\eqref{gff:gc:L2reduction} follows from the structure without those relations, and it is the
form used here; both give the same value. The algebra is checked exactly by
\path{LRI_C5_d2_simplify/proof/synth/synth_checks.py}.


\subsubsection{Generating functions and elliptic moments}
For $0<k<1$,
\begin{equation}\begin{aligned}
 \sum_{n\geq0}c_n^2\alpha_nk^{2n}&=-g(k)-\frac2\pi f(k)\log k,\\
 \sum_{n\geq0}c_n^2\beta_nk^{2n}&=-\frac2\pi f(k)\log k',\\
 \sum_{n\geq0}c_n^2\bigl[4\alpha_n\beta_n+2\psi_1(n+\tfrac12)\bigr]k^{2n}
 &=4g(k)\log k'+2\pi f(k)+\frac8\pi f(k)\log k\,\log k'.
\end{aligned}\label{gff:gc:L2generating}\end{equation}
To prove these, put $z=k^2$, $\Phi(x,z)={}_2F_1(x,x;1;z)$ and
\[
 \Lambda_x(z)=\sum_{n\geq0}\frac{(x)_n^2}{n!^2}\bigl[2\psi_0(n+1)-2\psi_0(n+x)-\log z\bigr]z^n
 =\frac{\Gamma(x)^2}{\Gamma(2x)}\,{}_2F_1(x,x;2x;1-z),
\]
the logarithmic case $m=0$ of the connection formula 15.8.10 of \cite{gff:DLMF15}. At
$x=\frac12$ it gives $\Lambda_{1/2}(z)=\pi\,{}_2F_1(\frac12,\frac12;1;1-z)=2g(k)$, which
is the first line, since $\sum_nc_n^2k^{2n}=\frac2\pi f(k)$. Euler's transformation $\Phi(x,z)=(1-z)^{1-2x}\Phi(1-x,z)$ (15.8.1 of
\cite{gff:DLMF15}), differentiated at $x=\frac12$, gives
$\partial_x\Phi(\frac12,z)=-\log(1-z)\,\Phi(\frac12,z)$. Termwise,
$\partial_x\Phi(\frac12,z)=2\sum_nc_n^2\beta_nz^n$, which is the second line. The
connection formula 15.8.4 of \cite{gff:DLMF15} for the parameters $(x,x;1)$, with its Gamma
coefficients simplified by reflection, reads
\[
 \Lambda_x(z)=(1-z)^{1-2x}\Lambda_{1-x}(z)+2\pi\cot(\pi x)\,\Phi(x,z),
 \qquad 0<x<1,\ x\neq\tfrac12 .
\]
For fixed $|z|<1$ the series converge locally uniformly in $x$, so both sides are
analytic on $0<x<1$ and the identity also holds at $x=\frac12$. Its $x$-derivative
there is $\partial_x\Lambda|_{1/2}=-\log(1-z)\,\Lambda_{1/2}(z)-\pi^2\Phi(\frac12,z)$.
Termwise,
$\partial_x\Lambda|_{1/2}=-\sum_nc_n^2[4\alpha_n\beta_n+2\psi_1(n+\frac12)]z^n
-2\log z\sum_nc_n^2\beta_nz^n$, and with the second line this is the third.

Write $\mathrm{FG}[\phi]=\int_0^1f(k)g(k)\phi(k)\,dk$ and
$\mathrm{GG}[\phi]=\int_0^1g(k)^2\phi(k)\,dk$. Section~\ref{gff:app:LreductionProof} gives the
following values: $\mathrm{GG}[1]=2M_{\rm e}$ by \eqref{gff:lp:eq:measures};
$\mathrm{FG}[w]=-\log2\,N_{\rm e}$ and $\mathrm{GG}[w]=-2R-4\log2\,M_{\rm e}$ by
\eqref{gff:lp:eq:weighted}; $\mathrm{FG}[\log k]=-\frac\pi2M_{\rm e}$ by \eqref{gff:lp:eq:mixedlog};
$R=\frac\pi4N_{\rm e}-\log2\,M_{\rm e}$ by \eqref{gff:lp:eq:R}; and $U_0=2N_{\rm e}$.
One further input is needed:
\begin{equation}
 \mathrm{GG}[\log k]=\int_0^1\mathbf K(k')^2\log k\,dk=-\pi N_{\rm e}.
 \label{gff:gc:L2contour}
\end{equation}
It is the real part of the first-quadrant identity \eqref{gff:lp:eq:quadrant} with $W(z)$
replaced by $\operatorname{Log}z$. The arc estimates are unchanged: the integrand is
$O(\log^3|z|/|z|^2)$ at infinity and has integrable logarithmic singularities at $0$ and
$1$. After $x=1/k$ on $(1,\infty)$ the boundary value $\operatorname{Log}(x+\ii0)=-\log k$
multiplies $(f+\ii g)^2$, so the two contributions $\pm\int_0^1f^2\log k\,dk$ cancel and
the real part of the left side is $\mathrm{GG}[\log k]$. On the imaginary axis
$\operatorname{Log}(\ii y)=\log y+\frac{\ii\pi}2$, so after $y=k/k'$ the real part of the
right side is $-\frac\pi2U_0=-\pi N_{\rm e}$. Since $\log k'=\frac12(\log k-w)$, it
follows that
\begin{equation}
 \mathrm{FG}[\log k']=-\frac\pi4M_{\rm e}+\frac{\log2}2N_{\rm e},\qquad
 \mathrm{GG}[\log k']=\log2\,M_{\rm e}-\frac\pi4N_{\rm e}.
 \label{gff:gc:L2moments}
\end{equation}

\subsubsection{The single sums and the conclusion}\label{gff:gc:L2conclusion}
The complementary moment of Section~\ref{gff:sec:oldmasters},
$\int_0^1k^rg(k)\,dk=\frac\pi4\Gamma(\frac{r+1}2)^2/\Gamma(\frac{r+2}2)^2$ for $r>-1$, and
its $r$-derivative at $r=2n$ (dominated convergence) give
\[
 c_n^2=\frac4{\pi^2}\int_0^1k^{2n}g(k)\,dk,\qquad
 c_n^2\alpha_n=\frac4{\pi^2}\int_0^1k^{2n}\log k\;g(k)\,dk .
\]
In each single sum write one factor $c_n^2$, or $c_n^2\alpha_n$, in this form, and sum
the remaining factor with \eqref{gff:gc:L2generating}. For $\mathcal S_2$ use
\[
 c_n^4\bigl[4\alpha_n\beta_n+\psi_1(n+\tfrac12)\bigr]
 =\tfrac12\bigl\{c_n^2\bigl[4\alpha_n\beta_n+2\psi_1(n+\tfrac12)\bigr]\cdot c_n^2
 +2c_n^2\alpha_n\cdot2c_n^2\beta_n\bigr\}.
\]
The interchanges of sum and integral are justified by Tonelli's theorem, since
$c_n^2>0$, $\alpha_n<0$ and $\beta_n\geq0$ have fixed signs, and, for the first product
above, by Fubini's theorem, since that product is absolutely summable. The two
contributions $\pm\frac{16}{\pi^3}\int_0^1fg\log k\log k'\,dk$ cancel, and
\[
 \sum_nc_n^4=\frac{8N_{\rm e}}{\pi^3},\qquad
 \sum_nc_n^4\alpha_n=-\frac{4M_{\rm e}}{\pi^2},\qquad
 \mathcal S_1=-\frac8{\pi^3}\,\mathrm{FG}[\log k'],\qquad
 \mathcal S_2=\frac8{\pi^2}\,\mathrm{GG}[\log k']+\frac4\pi N_{\rm e}.
\]
Hence $a=8N_{\rm e}/\pi$ and $b=8M_{\rm e}/\pi$, in agreement with \eqref{gff:fp:abWheel},
and \eqref{gff:gc:L2moments} gives the two single-sum relations
\begin{equation}\boxed{
 \mathcal S_1=\frac{\pi b-2a\log2}{4\pi^2},\qquad
 \mathcal S_2=\frac a4+\frac{b\log2}\pi .
}\label{gff:gc:L2singles}\end{equation}
With $\hat a=a/\pi^2=\sum_kc_k^4$ and $\hat b=b/\pi^2=-\frac2\pi\sum_kc_k^4\alpha_k$
they read $\mathcal S_1=\frac14(\pi\hat b-2\hat a\log2)$ and
$\mathcal S_2=\frac{\pi^2}4\hat a+\pi\hat b\log2$, the two relations first found
numerically. Inserting \eqref{gff:gc:L2singles} into
\eqref{gff:gc:L2reduction}, the $a^2$ and $b^2$ terms cancel:
\[
 \LL_2=2a^2-4b^2+\frac{4b}{\pi}(\pi b-2a\log2)-8a\Bigl(\frac a4+\frac{b\log2}\pi\Bigr)
 =-\frac{16\log2}{\pi}\,ab=-2\log2\,P_4(2).
\]
With $P_4(2)=(B_-^2-B_+^2)/(4\pi^3)$ from \eqref{gff:fp:P4two}, this proves
\eqref{gff:gc:L2closed}. Equivalently, \eqref{gff:gc:L2taylor} and \eqref{gff:gc:L2singles} say
$\mathcal A'(\frac12)=2\pi b-2\log2\,a$ and $\mathcal B'(\frac12)=\frac{3\pi}2a-2\log2\,b$,
that is, $(4^x\mathcal A)'=2\pi\,4^x\mathcal B$ and
$(4^x\mathcal B)'=\frac{3\pi}2\,4^x\mathcal A$ at $x=\frac12$.

\paragraph{Checks.}
These are not proof steps. The closed forms of $\sum c_k^4$, $\sum c_k^4\alpha_k$,
$\mathcal S_1$ and $\mathcal S_2$ agree with $70$-digit Euler--Maclaurin values of the
single sums to $3\times10^{-70}$. With $\mathcal A'$ and $\mathcal B'$ obtained from the termwise
$x$-derivative of $\Omega(t;x)$ at generic $t$, \eqref{gff:gc:L2taylor} and the chain
$\frac12\mathcal W'(\frac12)=\LL_2$ hold to $10^{-82}$. The generating functions
\eqref{gff:gc:L2generating} and the moments \eqref{gff:gc:L2contour}, \eqref{gff:gc:L2moments} hold to
about $35$ digits.

\subsubsection{A second proof by rim--spoke reflection}
\label{gff:gc:L2second}
An independent proof works directly with the wheel derivative of
\eqref{gff:gs:wheelJets}. This subsection states the result and its consequences and then proves them in full. For $\frac14<x<\frac34$ it
shows that
\[
 W_4(2;2x)=\frac{2^{4-4x}\Gamma(1-x)^4\Gamma(x)^6}{\pi^{11/2}\Gamma(2x-\frac12)}
 \bigl[\mathfrak a(x)\mathfrak b(1-x)+\mathfrak b(x)\mathfrak a(1-x)\bigr],
\]
where $\mathfrak a(x)+\mathfrak b(x)\cos2\pi\tau$ is the two-mode form of the entire
function
\[
 \pi^4\sum_{k\in\mathbb Z}
 \bigl[\Gamma(1-\tfrac x2+k+\tau)\,\Gamma(1-\tfrac x2-k-\tau)\bigr]^{-4},
\]
with $\mathfrak a(\frac12)=a$ and $\mathfrak b(\frac12)=b$. The inputs are the
square-lattice factorization of the wheel for every rim exponent, the
growth argument of \eqref{gff:fp:finiteFourier}, and a Barnes integral of well-poised type.
That integral is proved in its terminating case from the Pfaff--Saalsch\"utz sum (16.4.3 of \cite{gff:DLMF16}),
Whipple's transformations (16.4.14 and 16.4.15 of \cite{gff:DLMF16}) and Barnes' first lemma
(DLMF 5.13.3), and is extended by Carlson's theorem. The bracket is
invariant under $x\to1-x$, which exchanges rim and spoke exponents. Hence, with
$\gamma_C(z)=\Gamma(z)/\Gamma(1-z)$ as in \eqref{gff:gs:W3d2},
\[
 \frac{W_4(2;b)}{W_4(2;2-b)}=2^{4-4b}\frac{\gamma_C(b/2)^2}{\gamma_C(b-\frac12)},
 \qquad \frac12<b<\frac32 ,
\]
and at the fixed point $b=1$ only the Gamma prefactor contributes to the logarithmic
derivative. This gives
$\partial_b\log W_4(2;b)|_{b=1}=\frac12(-4\log2-4\psi_0(\frac12)+6\psi_0(\frac12)-2\psi_0(\frac12))
=-2\log2$. As a check, the factorized form agrees with a direct spectral evaluation of
$W_4(2;2x)$ at $x=\frac25,\frac12,\frac35$ to $2\times10^{-70}$.

The machine checks of the proof below are
\path{LRI_C5_d2_proof/L2_structural/checks_all.wl} (log \path{checks_all.log}); the
$70$-digit spectral check just quoted is recorded in
\path{LRI_C5_d2_proof/review_L2_structural_numerics_r0/}.

\paragraph{Notation.}
Throughout this subsubsection $d=2$, the rim exponent is $b=2x$, and
$N_x=\Gamma(1-x)/\Gamma(x)=\gamma_C(1-x)$ and $\varphi_x(s)=\Gamma(s+x)/\Gamma(s+1)$ are as in
\eqref{gff:gc:L2rim}. Since $\Gamma(d/2)=1$, the eigenvalue \eqref{gff:gs:balancedWheel} is
$t_{\ell,2x}(\nu)=N_x\,\mathsf g_\ell(\ii\nu)$ with
\[
 \mathsf g_\ell(s)=\frac{\Gamma(\frac\ell2+\frac x2+s)\,\Gamma(\frac\ell2+\frac x2-s)}
 {\Gamma(\frac\ell2+1-\frac x2+s)\,\Gamma(\frac\ell2+1-\frac x2-s)},
\]
the kernel $g_\ell(s;2x)$ of \eqref{gff:gdr:kernel} at $d=2$, and $W_4(2;2x)=\langle t_{\ell,2x}^4\rangle_2$
with $d_0=1$, $d_{\ell\geq1}=2$. For $0<x<\frac34$ let
\[
 \Omega(t;x)=\pi^4\sum_{k\in\mathbb Z}\bigl[\Gamma(1+k+t)\,\Gamma(1-x-k-t)\bigr]^{-4}
 =\sin^4\pi(t+x)\sum_{k\in\mathbb Z}\varphi_x(k+t)^4 ,
\]
where the second form follows termwise from Gamma reflection \cite{gff:DLMF5}. For $x\in I$
this is the function $\Omega(t;x)$ of \eqref{gff:gc:L2modes}, and $\Omega(t;\frac12)=\Omega_4(t)$
of \eqref{gff:fp:OmegaWheel}. The function displayed above in the variable $\tau$ is
$\Omega(\tau-\frac x2;x)$, and its two modes are written $\mathfrak a(x)$, $\mathfrak b(x)$. The
statement to be proved is
\begin{equation}
 W_4(2;2x)=\mathfrak p(x)\bigl[\mathfrak a(x)\mathfrak b(1-x)+\mathfrak b(x)\mathfrak a(1-x)\bigr],\qquad
 \mathfrak p(x)=\frac{2^{4-4x}\Gamma(1-x)^4\Gamma(x)^6}{\pi^{11/2}\Gamma(2x-\frac12)},\qquad
 \tfrac14<x<\tfrac34 .
 \label{gff:mgp:rimSpoke}
\end{equation}
At $x=\frac12$, $\mathfrak p(\frac12)=4/\pi$ and \eqref{gff:mgp:rimSpoke} is $P_4(2)=8ab/\pi$.

\paragraph{Lemma (two modes of a reciprocal-Gamma kernel).}
\emph{Let $0<\kappa_1,\dots,\kappa_4<1$ with $\kappa=\kappa_1+\dots+\kappa_4>\frac52$, and let
$\mathfrak k(s)=\prod_{j=1}^4[\Gamma(\kappa_j+s)\Gamma(\kappa_j-s)]^{-1}$, an entire even
function. Then}
\begin{itemize}
\item[(i)] \emph{$|\mathfrak k(s)|\le C\,e^{4\pi|\Im s|}(1+|s|)^{4-2\kappa}$ for all $s\in\mathbb C$;}
\item[(ii)] \emph{$E_{\mathfrak k}(\tau)=\sum_{k\in\mathbb Z}\mathfrak k(k+\tau)$ converges absolutely
and locally uniformly, and $E_{\mathfrak k}$ is entire, even and one-periodic;}
\item[(iii)] \emph{$E_{\mathfrak k}(\tau)=\int_{\mathbb R}\mathfrak k(s)\dd s
+2\cos(2\pi\tau)\int_{\mathbb R}\mathfrak k(s)\cos(2\pi s)\dd s$.}
\end{itemize}
\emph{Proof.} (i) By reflection,
$[\Gamma(\kappa_j+s)\Gamma(\kappa_j-s)]^{-1}=\pi^{-1}\sin\pi(\kappa_j-s)\,\Gamma(s+1-\kappa_j)/\Gamma(s+\kappa_j)$.
For $\Re s\geq0$ the Gamma ratio has no poles, because $1-\kappa_j>0$, and it is
$O((1+|s|)^{1-2\kappa_j})$ by the Gamma-ratio asymptotics \cite{gff:DLMFStirling} for $|s|\geq1$
and by continuity on the compact rest; moreover $|\sin\pi(\kappa_j-s)|\le e^{\pi|\Im s|}$. For
$\Re s\le0$ use $\mathfrak k(-s)=\mathfrak k(s)$. (ii) On a compact set $K$, (i) gives
$|\mathfrak k(k+\tau)|\le C_K(1+|k|)^{4-2\kappa}$, and $2\kappa-4>1$; the Weierstrass test applies.
Evenness and periodicity follow by re-indexing. (iii) For $\tau=u+\ii y$ with $0\le u\le1$, (i)
and comparison with $\int_{\mathbb R}(1+|v+\ii y|)^{4-2\kappa}\dd v$ give
$|E_{\mathfrak k}(u+\ii y)|\le C'e^{4\pi|y|}(1+|y|)^{5-2\kappa}$, with $5-2\kappa<0$. For an entire
one-periodic function the coefficient
$e_m=\int_0^1E_{\mathfrak k}(u+\ii y)e^{-2\pi\ii m(u+\ii y)}\dd u$ does not depend on $y$, so
$|e_m|\le C'e^{2\pi my+4\pi|y|}(1+|y|)^{5-2\kappa}$. Letting $y\to-\infty$ for $m\geq2$ and
$y\to+\infty$ for $m\le-2$ gives $e_m=0$; for $|m|=2$ this uses $(1+|y|)^{5-2\kappa}\to0$.
Evenness gives $e_{-1}=e_1$. Finally $\int_{\mathbb R}|\mathfrak k|<\infty$, so by Fubini
$e_0=\sum_k\int_0^1\mathfrak k(k+u)\dd u=\int_{\mathbb R}\mathfrak k$ and
$e_1=\int_{\mathbb R}\mathfrak k(s)e^{-2\pi\ii s}\dd s=\int_{\mathbb R}\mathfrak k(s)\cos(2\pi s)\dd s$.
$\square$

This is the growth argument of \eqref{gff:fp:finiteFourier} applied to a reciprocal-Gamma
kernel. Two cases are used.

(a) $\kappa_j=1-\frac x2$ for all $j$, with $0<x<\frac34$, so that $\kappa=4-2x>\frac52$. Then
$\pi^4E_{\mathfrak k}(\tau)=\Omega(\tau-\frac x2;x)$, so that for $0<x<\frac34$
\begin{equation}
 \Omega(t;x)=\mathfrak a(x)+\mathfrak b(x)\cos2\pi\bigl(t+\tfrac x2\bigr),\qquad
 \mathfrak a(x)=\pi^4\!\int_{\mathbb R}\frac{\dd s}{\mathcal G_x(s)^4},\qquad
 \mathfrak b(x)=2\pi^4\!\int_{\mathbb R}\frac{\cos(2\pi s)\dd s}{\mathcal G_x(s)^4},
 \label{gff:mgp:twoModes}
\end{equation}
with $\mathcal G_x(s)=\Gamma(1-\frac x2+s)\Gamma(1-\frac x2-s)$. On $I$, \eqref{gff:gc:L2modes} gives
$\mathfrak a=\mathcal A/N_x^2$ and $\mathfrak b=\mathcal B/N_x^2$, so $\mathfrak a(\frac12)=a$ and
$\mathfrak b(\frac12)=b$. Evaluating at $t=-\frac x2$ and $t=\frac12-\frac x2$, the two
combinations $\mathfrak a\pm\mathfrak b$ are well-poised bilateral series:
$\mathfrak a(x)+\mathfrak b(x)=\sin^4\frac{\pi x}2\sum_{k\in\mathbb Z}[\Gamma(k+\frac x2)/\Gamma(k+1-\frac x2)]^4$
and
$\mathfrak a(x)-\mathfrak b(x)=\cos^4\frac{\pi x}2\sum_{k\in\mathbb Z}[\Gamma(k+\frac12+\frac x2)/\Gamma(k+\frac32-\frac x2)]^4$.

(b) $\kappa_j=\frac12+p_j$ with $0<p_j<\frac12$ and $p=p_1+\dots+p_4>\frac12$, so that
$\kappa=2+p>\frac52$. Then
\[
 E^\sharp(\tau)=\sum_{k\in\mathbb Z}\prod_{j=1}^4
 \bigl[\Gamma(\tfrac12+p_j+k+\tau)\,\Gamma(\tfrac12+p_j-k-\tau)\bigr]^{-1}
\]
is a constant plus a multiple of $\cos2\pi\tau$. If $p_j=\frac x2$ for all $j$, then
$\pi^4E^\sharp(\tau)=\Omega(\tau-\frac{1-x}2;1-x)$, because $1-\frac{1-x}2=\frac12+\frac x2$.

\paragraph{Proposition (the lattice step).}
\emph{For $0<x<\frac34$ put $\mathcal F_x(s)=[\Gamma(\frac x2+s)\Gamma(\frac x2-s)]^4$ and}
\[
 \mathsf J_0(x)=\int_{-\ii\infty}^{\ii\infty}\frac{\dd s}{2\pi\ii}\,\mathcal F_x(s),\qquad
 \mathsf J_1(x)=\int_{-\ii\infty}^{\ii\infty}\frac{\dd s}{2\pi\ii}\,\cos(2\pi s)\,\mathcal F_x(s),
\]
\emph{which converge absolutely because $|\Gamma(\frac x2+\ii y)|^2\sim2\pi|y|^{x-1}e^{-\pi|y|}$. Then}
\begin{equation}\begin{aligned}
 W_4(2;2x)&=-N_x^4\Res_{u=0}\frac{\Omega(u;x)^2}{\sin^4\pi u\,\sin^4\pi(u+x)}
 =\frac{N_x^4}{\pi^4}\int_{-\ii\infty}^{\ii\infty}\frac{\dd s}{2\pi\ii}\,
 \Omega\bigl(s-\tfrac x2;x\bigr)\,\mathcal F_x(s)\\
 &=\frac{N_x^4}{\pi^4}\bigl[\mathfrak a(x)\mathsf J_0(x)+\mathfrak b(x)\mathsf J_1(x)\bigr].
\end{aligned}\label{gff:mgp:lattice}\end{equation}
\emph{In particular, with $\mathcal W$ of \eqref{gff:gc:L2residueForm},}
\begin{equation}
 W_4(2;2x)=\mathcal W(x)\qquad(x\in I).
 \label{gff:mgp:Wresidue}
\end{equation}

\emph{Proof.} (i) \emph{Spectral form.} Since $t_{\ell,2x}(\nu)$ is even in $\nu$,
\eqref{gff:gs:trace} gives
\[
 W_4(2;2x)=N_x^4\sum_{\ell\geq0}d_\ell\int_{-\ii\infty}^{\ii\infty}\frac{\dd s}{2\pi\ii}\,\mathsf g_\ell(s)^4 .
\]
The double sum and integral converge absolutely for $x<\frac34$: at fixed $\ell$,
$t_{\ell,2x}(\nu)^4=O(\nu^{8x-8})$, and $\int_0^\infty t_{\ell,2x}^4\dd\nu=O(\ell^{8x-7})$.

(ii) \emph{Lattice identity.} For $\ell\geq0$ and $n\in\mathbb Z$, Gamma reflection applied to
$\Gamma(-n-u)$ and to $1/\Gamma(1-x-n-u)$ gives
\[
 \mathsf g_\ell\bigl(\tfrac\ell2+\tfrac x2+n+u\bigr)
 =-\frac{\sin\pi(x+u)}{\sin\pi u}\,\varphi_x(\ell+n+u)\,\varphi_x(n+u),
\]
so $-N_x\,\mathsf g_\ell(\frac\ell2+\frac x2+n+u)$ is the left side of \eqref{gff:gc:L2local} with
$m=\ell+n$. With $(k,j)=(\ell+n,n)$ the
weights $d_0=1$, $d_{\ell\geq1}=2$ turn $\sum_{\ell\geq0}d_\ell\sum_{n\in\mathbb Z}$ into the sum
over all $(k,j)\in\mathbb Z^2$, because the summand $\varphi_x(k+u)^4\varphi_x(j+u)^4$ is
symmetric in $k,j$. Hence
\[
 \mathsf S(u)=\sum_{\ell\geq0}d_\ell\sum_{n\in\mathbb Z}\mathsf g_\ell\bigl(\tfrac\ell2+\tfrac x2+n+u\bigr)^4
 =\frac{\sin^4\pi(x+u)}{\sin^4\pi u}\Bigl[\sum_{k\in\mathbb Z}\varphi_x(k+u)^4\Bigr]^2
 =\frac{\Omega(u;x)^2}{\sin^4\pi u\,\sin^4\pi(u+x)} .
\]
On a circle $|u|=r$ with $0<r<\min(x,1-x)$ one has $|\varphi_x(k+u)|\le C(1+|k|)^{x-1}$
uniformly, and $4x-4<-1$, so the double series converges absolutely and uniformly there.
This is \eqref{gff:gdr:d12lattice} at $d=2$.

(iii) \emph{Residues of the two-dimensional integral.} For each $\ell$ the poles of
$\mathsf g_\ell$ in $\Re s>0$ are at $s=\frac\ell2+\frac x2+n$, $n\geq0$, from
$\Gamma(\frac\ell2+\frac x2-s)$; those of $\Gamma(\frac\ell2+\frac x2+s)$ lie in $\Re s<0$. By
reflection,
\[
 \mathsf g_\ell(s)=\frac{\Gamma(s+\frac\ell2+\frac x2)}{\Gamma(s+\frac\ell2+1-\frac x2)}\cdot
 \frac{\Gamma(s-\frac\ell2+\frac x2)}{\Gamma(s-\frac\ell2+1-\frac x2)}\cdot
 \frac{\sin\pi(\frac x2+s-\frac\ell2)}{\sin\pi(\frac x2-s+\frac\ell2)} .
\]
On the arcs $|s-\frac\ell2-\frac x2|=N+\frac12$ in $\Re s\geq0$, which pass midway between
consecutive poles, the sine ratio is bounded uniformly in $N$, and for large $N$ both Gamma
ratios are $O(|s|^{x-1})$. So $\mathsf g_\ell^4=O(|s|^{8x-8})$ there, and since the arcs
have length $O(N)$ and $8x-7<0$ their contributions vanish. Closing the contour to the
right, which is clockwise,
$\int_{-\ii\infty}^{\ii\infty}\frac{\dd s}{2\pi\ii}\mathsf g_\ell^4
=-\sum_{n\geq0}\Res_{u=0}\mathsf g_\ell(\frac\ell2+\frac x2+n+u)^4$. For $n<0$ the function
$\mathsf g_\ell(\frac\ell2+\frac x2+n+u)$ is regular at $u=0$: there $\Gamma(-n-u)$ is regular,
and so is $\Gamma(\ell+n+x+u)$, since $x\notin\mathbb Z$. Summing over $\ell$ with the weights
$d_\ell$, and using the uniform convergence of (ii) on $|u|=r$ to take the sums inside one
contour integral, gives $W_4(2;2x)=-N_x^4\Res_{u=0}\mathsf S(u)$, the first equality in
\eqref{gff:mgp:lattice}.

For $k\le-1$, $1/\Gamma(k+1+u)$ has a simple zero at $u=0$, so $\varphi_x(k+u)^4=O(u^4)$, and
all terms of $\mathsf S$ that contain such a factor are regular at $u=0$. Hence
$\Res_{u=0}\mathsf S(u)=\Res_{u=0}[\sin^4\pi(x+u)\,\mathcal H_x(u)^2/\sin^4\pi u]$ with $\mathcal H_x$
of \eqref{gff:gc:L2residueForm}, and for $x\in I$ that equation gives \eqref{gff:mgp:Wresidue}. This is
the per-spin residue evaluation of \eqref{gff:gs:BarnesMoment} at rim exponent $2x$; it is also the
case $d=2$ of \eqref{gff:gdr:functional} with \eqref{gff:gdr:d12lattice}.

(iv) \emph{Reflection.} For $k\in\mathbb Z$, $1/\Gamma(k+1+u)=(-1)^{k+1}\sin(\pi u)\Gamma(-k-u)/\pi$,
so $\varphi_x(k+u)^4=\pi^{-4}\sin^4(\pi u)\,\mathcal F_x(\frac x2+k+u)$. With
$\mathcal P_x(u)=\sum_{k\in\mathbb Z}\mathcal F_x(\frac x2+k+u)$ this gives
$\Omega(u;x)=\pi^{-4}\sin^4\pi(x+u)\sin^4(\pi u)\,\mathcal P_x(u)$, hence
\[
 \frac{\Omega(u;x)^2}{\sin^4\pi u\,\sin^4\pi(u+x)}=\pi^{-4}\,\Omega(u;x)\,\mathcal P_x(u).
\]

(v) \emph{Residues of the one-dimensional integral.} The function
$\Omega(s-\frac x2;x)=\mathfrak a+\mathfrak b\cos2\pi s$ is entire, with
$|\Omega(s-\frac x2;x)|\le(|\mathfrak a|+|\mathfrak b|)e^{2\pi|\Im s|}$. Writing
$\Gamma(\frac x2-s)=\pi/[\sin\pi(\frac x2-s)\Gamma(1-\frac x2+s)]$ gives
$|\mathcal F_x(s)|\le C|s|^{4x-4}e^{-4\pi|\Im s|}$ on the arcs $|s-\frac x2|=N+\frac12$ in
$\Re s\geq0$. So the integrand is $O(|s|^{4x-4})$ there, and the arc contributions vanish
because $4x-3<0$. The poles of $\mathcal F_x$ in $\Re s>0$ are at $s=\frac x2+n$, $n\geq0$, and
therefore
\[
 \int_{-\ii\infty}^{\ii\infty}\frac{\dd s}{2\pi\ii}\,\Omega\bigl(s-\tfrac x2;x\bigr)\mathcal F_x(s)
 =-\sum_{n\geq0}\Res_{u=0}\bigl[\Omega(n+u;x)\,\mathcal F_x(\tfrac x2+n+u)\bigr]
 =-\Res_{u=0}\bigl[\Omega(u;x)\,\mathcal P_x(u)\bigr],
\]
by the periodicity of $\Omega$ and the uniform convergence of $\mathcal P_x$ near $u=0$; the
terms with $n<0$ are regular at $u=0$. With (iv) this is the second equality of
\eqref{gff:mgp:lattice}, and \eqref{gff:mgp:twoModes} gives the third. $\square$

\paragraph{Lemma (a well-poised Barnes integral).}
\emph{Let $\Re\mathsf a_j>0$ ($j=1,2,3$), $\mathsf a=\mathsf a_1+\mathsf a_2+\mathsf a_3$ and
$\Re(\mathsf a+\frac12-\mathsf b)>0$. Then}
\begin{equation}\begin{aligned}
 &\int_{-\ii\infty}^{\ii\infty}\frac{\dd s}{2\pi\ii}\,
 \frac{\prod_{j=1}^3\Gamma(\mathsf a_j+s)\Gamma(\mathsf a_j-s)}{\Gamma(\mathsf b+s)\Gamma(\mathsf b-s)}\\
 &\qquad=\frac{2\sqrt\pi\,2^{2\mathsf b-2\mathsf a}\prod_{i<j}\Gamma(\mathsf a_i+\mathsf a_j)\prod_i\Gamma(2\mathsf a_i)}
 {\Gamma(\mathsf a+\frac12-\mathsf b)\prod_i\Gamma(\mathsf a_i+\mathsf b)}\;
 {}_4F_3\!\left(\begin{matrix}2\mathsf b-1,\ \mathsf b-\mathsf a_1,\ \mathsf b-\mathsf a_2,\ \mathsf b-\mathsf a_3\\
 \mathsf b+\mathsf a_1,\ \mathsf b+\mathsf a_2,\ \mathsf b+\mathsf a_3\end{matrix};1\right).
\end{aligned}\label{gff:mgp:BarnesLemma}\end{equation}
\emph{The ${}_4F_3$ is well-poised and converges, since its parameter excess is
$1+2\mathsf a-2\mathsf b$.}

At $\mathsf b=\mathsf a_3$ this is Barnes' first lemma (DLMF 5.13.3). At $\mathsf b=\frac12$
it is the de Branges--Wilson integral (DLMF 5.13.5) with one parameter equal to $0$, since
$\Gamma(s)\Gamma(-s)/[\Gamma(2s)\Gamma(-2s)]=4\cos\pi s=4\pi/[\Gamma(\frac12+s)\Gamma(\frac12-s)]$.
Write $I_{\rm L}$ and $I_{\rm R}$ for the two sides of \eqref{gff:mgp:BarnesLemma}.

\emph{Proof.} (a) \emph{Terminating case $\mathsf a_3=\mathsf b+n$, $n\in\mathbb N_0$.} Take
$\mathsf a_1,\mathsf a_2>0$ and $\mathsf b>\frac12$ real and generic
($\mathsf b-\mathsf a_1,\mathsf b-\mathsf a_2\notin\mathbb Z$). Then
$\Gamma(\mathsf a_3\pm s)/\Gamma(\mathsf b\pm s)=(\mathsf b+s)_n(\mathsf b-s)_n$. The
Pfaff--Saalsch\"utz sum (16.4.3 of \cite{gff:DLMF16}), with its parameters taken as
$\mathsf a_2+s$, $\mathsf a_2-s$, $\mathsf a_2+\mathsf b$ and $1+\mathsf a_2-\mathsf b-n$, gives
\[
 (\mathsf b+s)_n(\mathsf b-s)_n=(\mathsf a_2+\mathsf b)_n(\mathsf b-\mathsf a_2)_n
 \sum_{m=0}^n\frac{(-n)_m(\mathsf a_2+s)_m(\mathsf a_2-s)_m}
 {m!\,(\mathsf a_2+\mathsf b)_m(1+\mathsf a_2-\mathsf b-n)_m}.
\]
Since $(\mathsf a_2\pm s)_m\Gamma(\mathsf a_2\pm s)=\Gamma(\mathsf a_2+m\pm s)$, Barnes' first lemma
evaluates each term:
$\int\frac{\dd s}{2\pi\ii}\Gamma(\mathsf a_1+s)\Gamma(\mathsf a_1-s)\Gamma(\mathsf a_2+m+s)\Gamma(\mathsf a_2+m-s)
=\Gamma(2\mathsf a_1)\Gamma(\mathsf a_1+\mathsf a_2+m)^2\Gamma(2\mathsf a_2+2m)/\Gamma(2\mathsf a_1+2\mathsf a_2+2m)$.
With $(2c)_{2m}=4^m(c)_m(c+\frac12)_m$ this gives
\[\begin{gathered}
 I_{\rm L}=(\mathsf a_2+\mathsf b)_n(\mathsf b-\mathsf a_2)_n\,
 \frac{\Gamma(2\mathsf a_1)\Gamma(2\mathsf a_2)\Gamma(\mathsf a_1+\mathsf a_2)^2}{\Gamma(2\mathsf a_1+2\mathsf a_2)}\,F_{\rm L},\\
 F_{\rm L}={}_4F_3\!\left(\begin{matrix}-n,\ \mathsf a_1+\mathsf a_2,\ \mathsf a_2,\ \mathsf a_2+\frac12\\
 \mathsf a_2+\mathsf b,\ 1+\mathsf a_2-\mathsf b-n,\ \mathsf a_1+\mathsf a_2+\frac12\end{matrix};1\right).
\end{gathered}\]
On the right,
$F_{\rm R}={}_4F_3(2\mathsf b-1,\mathsf b-\mathsf a_1,\mathsf b-\mathsf a_2,-n;\mathsf b+\mathsf a_1,\mathsf b+\mathsf a_2,2\mathsf b+n;1)$
is the terminating very-well-poised ${}_7F_6$ with first parameter $2\mathsf b-1$ and further
parameters $\mathsf b-\mathsf a_1$, $\mathsf b-\mathsf a_2$, $\mathsf b-\frac12$, $\mathsf b$, $-n$: the
pair $\mathsf b-\frac12$, $\mathsf b$ cancels the very-well-poised factor and itself. Whipple's
transformation (16.4.15 of \cite{gff:DLMF16}, with $f=-n$, $(b,c)\to(\mathsf b-\mathsf a_1,\mathsf b-\mathsf a_2)$
and $(d,e)\to(\mathsf b-\frac12,\mathsf b)$) gives
\[
 F_{\rm R}=\frac{(2\mathsf b)_n(\frac12)_n}{(\mathsf b+\frac12)_n(\mathsf b)_n}\,
 {}_4F_3\!\left(\begin{matrix}\mathsf a_1+\mathsf a_2,\ \mathsf b-\frac12,\ \mathsf b,\ -n\\
 \mathsf b+\mathsf a_1,\ \mathsf b+\mathsf a_2,\ \frac12-n\end{matrix};1\right),
\]
and Whipple's transformation of a balanced ${}_4F_3$ (16.4.14 of \cite{gff:DLMF16}, with
$(a;b,c;d,e,f)\to(\mathsf a_1+\mathsf a_2;\mathsf b-\frac12,\mathsf b;\mathsf b+\mathsf a_2,\mathsf b+\mathsf a_1,\frac12-n)$)
turns the last ${}_4F_3$ into
$\frac{(\mathsf b-\mathsf a_2)_n(\frac12-n-\mathsf a_1-\mathsf a_2)_n}{(\mathsf b+\mathsf a_1)_n(\frac12-n)_n}F_{\rm L}$.
Hence $F_{\rm R}=K_nF_{\rm L}$ with
\[
 K_n=\frac{(2\mathsf b)_n(\frac12)_n(\mathsf b-\mathsf a_2)_n(\frac12-n-\mathsf a_1-\mathsf a_2)_n}
 {(\mathsf b+\frac12)_n(\mathsf b)_n(\mathsf b+\mathsf a_1)_n(\frac12-n)_n}.
\]
The prefactor of $I_{\rm R}$ times $K_n$ equals the prefactor of $I_{\rm L}$. Indeed, with
$\mathsf a_3=\mathsf b+n$: $\Gamma(\mathsf a_1+\mathsf b+n)=\Gamma(\mathsf a_1+\mathsf b)(\mathsf b+\mathsf a_1)_n$,
$\Gamma(\mathsf a_2+\mathsf b+n)=\Gamma(\mathsf a_2+\mathsf b)(\mathsf a_2+\mathsf b)_n$,
$\Gamma(2\mathsf b+2n)(2\mathsf b)_n/\Gamma(2\mathsf b+n)=(2\mathsf b)_{2n}=4^n(\mathsf b)_n(\mathsf b+\frac12)_n$,
and $(\frac12)_n(\frac12-n-c)_n/(\frac12-n)_n=(\frac12+c)_n$ for $c=\mathsf a_1+\mathsf a_2$, which
cancels against $\Gamma(\mathsf a+\frac12-\mathsf b)=\Gamma(c+\frac12)(c+\frac12)_n$. What remains is
$2\sqrt\pi\,2^{-2c}\Gamma(2c)/[\Gamma(c)\Gamma(c+\frac12)]=1$, by duplication \cite{gff:DLMF5}. So
$I_{\rm L}=I_{\rm R}$ when $\mathsf a_3=\mathsf b+n$.

(b) \emph{Carlson's theorem.} Keep $\mathsf a_1,\mathsf a_2>0$ and a generic
$\mathsf b\in(\frac12,\frac12+\mathsf a_1+\mathsf a_2)$, put $\mathsf a_3=\mathsf b+z$ with $\Re z\geq0$, and
let $f(z)=(I_{\rm L}-I_{\rm R})/\Gamma(\mathsf b+z)^2$. Then $f$ is analytic for $\Re z>-\varepsilon$,
and $f(n)=0$ for $n\in\mathbb N_0$ by (a). It is bounded on $\Re z\geq0$. First, by the Gamma
bound below with $w=\mathsf b+z$,
\[
 \Bigl|\frac{I_{\rm L}}{\Gamma(\mathsf b+z)^2}\Bigr|\le\frac{e^{2/3}}{2\pi}\int_{\mathbb R}
 \frac{|\Gamma(\mathsf a_1+\ii y)|^2|\Gamma(\mathsf a_2+\ii y)|^2}{|\Gamma(\mathsf b+\ii y)|^2}\dd y<\infty .
\]
Second, by duplication, $I_{\rm R}/\Gamma(\mathsf b+z)^2$ is a constant times
$\frac{\Gamma(\mathsf a_1+\mathsf b+z)\Gamma(\mathsf a_2+\mathsf b+z)\Gamma(\mathsf b+\frac12+z)}
{\Gamma(\mathsf a_1+\mathsf a_2+\frac12+z)\Gamma(2\mathsf b+z)\Gamma(\mathsf b+z)}\,F_{\rm R}(z)$, where
$F_{\rm R}(z)$ is the ${}_4F_3$ of \eqref{gff:mgp:BarnesLemma} at $\mathsf a_3=\mathsf b+z$. The Gamma ratio has
equal total weights in numerator and denominator and no poles for $\Re z\geq0$, so it is
bounded there \cite{gff:DLMFStirling}. For $\Re z\geq0$ and $k\geq0$,
$|k-z|\le|k+z|\le|k+2\mathsf b+z|$, so $|(-z)_n/(2\mathsf b+z)_n|\le1$ and
$|F_{\rm R}(z)|\le\sum_n\bigl|\frac{(2\mathsf b-1)_n(\mathsf b-\mathsf a_1)_n(\mathsf b-\mathsf a_2)_n}
{n!\,(\mathsf b+\mathsf a_1)_n(\mathsf b+\mathsf a_2)_n}\bigr|<\infty$, the terms being
$O(n^{2\mathsf b-2-2\mathsf a_1-2\mathsf a_2})$ with exponent below $-1$. A function that is analytic and
bounded on $\Re z\geq0$ and vanishes at every nonnegative integer is identically zero
(Carlson's theorem \cite{gff:TitchmarshCarlson}). So $f\equiv0$.

(c) \emph{Continuation.} Both sides of \eqref{gff:mgp:BarnesLemma} are analytic on the convex domain
$\{\Re\mathsf a_j>0,\ \Re(\mathsf a+\frac12-\mathsf b)>0\}\subset\mathbb C^4$. On the left, the
imaginary axis separates the poles of $\Gamma(\mathsf a_j+s)$ from those of $\Gamma(\mathsf a_j-s)$ and
the integrand decays like $e^{-2\pi|\Im s|}$, locally uniformly. On the right,
$1/[\Gamma(\mathsf b+\mathsf a_j)(\mathsf b+\mathsf a_j)_n]=1/\Gamma(\mathsf b+\mathsf a_j+n)$ is entire and the
series converges locally uniformly. By (b) the two sides agree for real
$(\mathsf a_1,\mathsf a_2,\mathsf b)$ in an open subset of $\mathbb R^3$ and $\Re\mathsf a_3\geq\mathsf b$.
Continuation in $\mathsf a_3$, and then in $(\mathsf a_1,\mathsf a_2,\mathsf b)$ (an analytic function that
vanishes on an open subset of $\mathbb R^3$ vanishes identically), gives equality on the whole
domain. $\square$

\paragraph{Lemma (a uniform Gamma bound).}
\emph{For $\Re w\geq\frac12$ and real $y$, $|\Gamma(w+\ii y)\Gamma(w-\ii y)|\le e^{2/3}|\Gamma(w)|^2$.}

\emph{Proof.} Binet's first formula \cite{gff:WhittakerWatsonBinet} gives, for $\Re\zeta>0$,
$\log\Gamma(\zeta)=(\zeta-\frac12)\log\zeta-\zeta+\frac12\log2\pi+\int_0^\infty\varpi(t)e^{-t\zeta}\,\frac{\dd t}t$
with $\varpi(t)=\frac12-\frac1t+\frac1{e^t-1}=\frac12[\coth\frac t2-\frac2t]$. The expansion
$\coth v=\frac1v+2v\sum_{n\geq1}(v^2+n^2\pi^2)^{-1}$ (DLMF 4.36.3) gives
$0<\coth v-\frac1v<\frac v3$ for $v>0$, hence $0<\varpi(t)/t<\frac1{12}$, and the remainder
integral has modulus at most $1/(12\Re\zeta)\le\frac16$. In
$\log|\Gamma(w+\ii y)\Gamma(w-\ii y)/\Gamma(w)^2|$ the terms $-\zeta$ and $\frac12\log2\pi$ cancel. The
main terms give $\mathfrak g(y)+\mathfrak g(-y)-2\mathfrak g(0)$ with
$\mathfrak g(t)=\Re[(w+\ii t-\frac12)\log(w+\ii t)]$, and, with $\zeta=w+\ii t$,
$\mathfrak g''(t)=-\Re[\zeta^{-1}+\frac12\zeta^{-2}]\le(\frac12-\Re w)/|\zeta|^2\le0$. So $\mathfrak g$ is
concave and $\mathfrak g(y)+\mathfrak g(-y)-2\mathfrak g(0)\le0$. The four remainders contribute at
most $\frac23$. $\square$

\paragraph{Proposition (Barnes--reciprocal duality).}
\emph{Let $0<p_j<\frac12$ ($j=1,\dots,4$), $p=p_1+\dots+p_4>\frac12$, and}
\[
 \mathfrak l(p_1,\dots,p_4)=\frac{2^{2-2p}\prod_{i<j}\Gamma(p_i+p_j)\prod_i\Gamma(2p_i)}
 {\pi^{9/2}\,\Gamma(p-\frac12)} .
\]
\emph{Then, with $E^\sharp$ of case (b) above, for all $\tau\in\mathbb C$}
\begin{equation}
 \int_{-\ii\infty}^{\ii\infty}\frac{\dd s}{2\pi\ii}\bigl[\cos2\pi s+\cos2\pi\tau\bigr]
 \prod_{j=1}^4\Gamma(p_j+s)\Gamma(p_j-s)
 =4\pi^7\,\mathfrak l(p_1,\dots,p_4)\,E^\sharp(\tau).
 \label{gff:mgp:duality}
\end{equation}

\emph{Proof.} The left side is a constant plus a multiple of $\cos2\pi\tau$, and so is the
right side, by case (b). Both sides are symmetric in the $p_j$. It therefore suffices to
prove equality at $\tau_4=\frac12+p_4$ when $p_3\neq p_4$: by symmetry it then holds at
$\tau_3=\frac12+p_3$ as well, and $\cos2\pi\tau_j=-\cos2\pi p_j$ takes different values at
$j=3,4$, because $\cos2\pi(\cdot)$ is injective on $(0,\frac12)$; a constant plus a multiple of
$\cos2\pi\tau$ is fixed by its values at two different values of $\cos2\pi\tau$. The case $p_3=p_4$ follows by continuity of both sides in
the $p_j$ (dominated convergence on the left, uniform convergence of $E^\sharp$ on the right).

\emph{Left side at $\tau_4$.} By reflection and
$\cos2\pi s-\cos2\pi p_4=2\sin\pi(p_4+s)\sin\pi(p_4-s)$,
$(\cos2\pi s-\cos2\pi p_4)\Gamma(p_4+s)\Gamma(p_4-s)=2\pi^2/[\Gamma(1-p_4+s)\Gamma(1-p_4-s)]$. So the left
side is $2\pi^2$ times the left side of \eqref{gff:mgp:BarnesLemma} with
$(\mathsf a_1,\mathsf a_2,\mathsf a_3,\mathsf b)=(p_1,p_2,p_3,1-p_4)$; the condition of that lemma is
$p>\frac12$.

\emph{Right side at $\tau_4$.} For $k\geq0$ the factor $j=4$ contains $1/\Gamma(-k)=0$. For
$k=-1-n$ with $n\geq0$ the summand is
$\prod_j[\Gamma(p_j+p_4-n)\Gamma(1+p_j-p_4+n)]^{-1}$, and $1/\Gamma(z-n)=(-1)^n(1-z)_n/\Gamma(z)$ gives
\[
 E^\sharp(\tau_4)=\prod_{j=1}^4\frac1{\Gamma(p_j+p_4)\Gamma(1+p_j-p_4)}\;
 {}_4F_3\!\left(\begin{matrix}1-2p_4,\ 1-p_1-p_4,\ 1-p_2-p_4,\ 1-p_3-p_4\\
 1+p_1-p_4,\ 1+p_2-p_4,\ 1+p_3-p_4\end{matrix};1\right),
\]
the ${}_4F_3$ of \eqref{gff:mgp:BarnesLemma} at $\mathsf b=1-p_4$. The prefactors agree: after
cancelling $\prod_{j\le3}\Gamma(p_j+p_4)$ and $\Gamma(2p_4)$,
$4\pi^7\mathfrak l/\prod_j[\Gamma(p_j+p_4)\Gamma(1+p_j-p_4)]$ reduces to
$4\pi^{5/2}2^{2-2p}\prod_{i<j\le3}\Gamma(p_i+p_j)\prod_{i\le3}\Gamma(2p_i)/[\Gamma(p-\frac12)\prod_{j\le3}\Gamma(1+p_j-p_4)]$,
which is $2\pi^2$ times the prefactor of \eqref{gff:mgp:BarnesLemma} at $\mathsf b=1-p_4$. $\square$

\paragraph{Corollary (the symmetric case).}
\emph{For $\frac14<x<\frac34$ and
$\mathfrak l_x=\mathfrak l(\frac x2,\frac x2,\frac x2,\frac x2)$, that is
$\mathfrak l_x=2^{2-4x}\Gamma(x)^{10}/(\pi^{9/2}\Gamma(2x-\frac12))$,}
\[
 \mathsf J_0(x)=4\pi^3\,\mathfrak l_x\,\mathfrak b(1-x),\qquad
 \mathsf J_1(x)=4\pi^3\,\mathfrak l_x\,\mathfrak a(1-x).
\]
\emph{Proof.} Take all $p_j=\frac x2$ in \eqref{gff:mgp:duality}, so that $p=2x>\frac12$. The left
side is $\mathsf J_1(x)+\mathsf J_0(x)\cos2\pi\tau$. By case (b) and \eqref{gff:mgp:twoModes} at $1-x$,
which lies in $(0,\frac34)$, the right side is
$4\pi^7\mathfrak l_x\pi^{-4}\Omega(\tau-\frac{1-x}2;1-x)=4\pi^3\mathfrak l_x[\mathfrak a(1-x)+\mathfrak b(1-x)\cos2\pi\tau]$.
Compare the two modes. $\square$

\paragraph{Proof of \eqref{gff:mgp:rimSpoke}.}
By \eqref{gff:mgp:lattice} and the corollary,
$W_4(2;2x)=N_x^4\pi^{-4}\,4\pi^3\mathfrak l_x[\mathfrak a(x)\mathfrak b(1-x)+\mathfrak b(x)\mathfrak a(1-x)]$, and
$\frac4\pi N_x^4\mathfrak l_x=\mathfrak p(x)$. $\square$

\paragraph{Corollary (rim--spoke reflection).}
\emph{For $\frac12<b<\frac32$,}
\[
 \frac{W_4(2;b)}{W_4(2;2-b)}=\frac{\mathfrak p(b/2)}{\mathfrak p(1-b/2)}
 =2^{4-4b}\,\frac{\gamma_C(b/2)^2}{\gamma_C(b-\frac12)} .
\]
\emph{Proof.} The bracket
$\mathfrak s(x)=\mathfrak a(x)\mathfrak b(1-x)+\mathfrak b(x)\mathfrak a(1-x)$ is invariant under $x\to1-x$,
and $W_4(2;2x)>0$ because its spectral integrand is positive, so $\mathfrak s\neq0$. Hence
$W_4(2;2x)/W_4(2;2-2x)=\mathfrak p(x)/\mathfrak p(1-x)
=2^{4-8x}\Gamma(x)^2\Gamma(\frac32-2x)/[\Gamma(1-x)^2\Gamma(2x-\frac12)]
=2^{4-8x}\gamma_C(x)^2/\gamma_C(2x-\frac12)$; put $b=2x$. $\square$

\paragraph{Corollary (the logarithmic wheel).}
\emph{$\partial_bW_4(2;b)|_{b=1}=-2\log2\,P_4(2)$, and $\LL_2$ of \eqref{gff:gs:Lseries} equals
$-2\log2\,P_4(2)$.}

\emph{Proof.} The kernel $\mathcal G_x(s)^{-4}$ of \eqref{gff:mgp:twoModes} is entire in $(x,s)$. Part
(i) of the two-mode lemma holds verbatim for complex $\kappa_j$ with $0<\Re\kappa_j<1$; for
$\kappa_j=1-\frac{x'}2$ it gives $|\mathcal G_{x'}(s)|^{-4}\le C(1+|s|)^{4\Re x'-4}$ for real $s$,
locally uniformly in $x'$. Cauchy's estimate on a circle $|x'-x|=r$ then gives
$|\partial_x\mathcal G_x(s)^{-4}|\le C(1+|s|)^{4(x+r)-4}$, which is integrable when $x+r<\frac34$. By
dominated convergence $\mathfrak a,\mathfrak b\in C^1(0,\frac34)$, so $\mathfrak s\in C^1(\frac14,\frac34)$, and
$\mathfrak s(1-x)=\mathfrak s(x)$ gives $\mathfrak s'(\frac12)=0$. Therefore
\[\begin{gathered}
 \frac{d}{dx}W_4(2;2x)\Big|_{x=\frac12}=\mathfrak p'(\tfrac12)\,\mathfrak s(\tfrac12)
 =(\log\mathfrak p)'(\tfrac12)\,W_4(2;1),\\
 (\log\mathfrak p)'(x)=-4\log2-4\psi_0(1-x)+6\psi_0(x)-2\psi_0(2x-\tfrac12),
\end{gathered}\]
and $(\log\mathfrak p)'(\frac12)=-4\log2$, the three $\psi_0(\frac12)$ terms cancelling. With $b=2x$,
$\partial_bW_4(2;b)|_{b=1}=-2\log2\,W_4(2;1)$, and $W_4(2;1)=P_4(2)$ by \eqref{gff:gs:primitiveWheel}.
This is the wheel-derivative reading \eqref{gff:gs:wheelJets} of $\LL_2$. For the series reading,
\eqref{gff:mgp:Wresidue} and \eqref{gff:gc:L2derivative} give
$\LL_2=\frac12\mathcal W'(\frac12)=\partial_bW_4(2;b)|_{b=1}$. $\square$

\paragraph{Remarks.}
(1) The reflection $b\to2-b$ exchanges the rim and spoke exponents of the two-dimensional
wheel. It is not a graph symmetry, since planar duality maps the wheel to itself with the
same exponents. It is a hypergeometric identity, \eqref{gff:mgp:BarnesLemma} and
\eqref{gff:mgp:duality}, between a Barnes integral of Gamma products and a periodization of
reciprocal Gamma products. $\LL_2$ is the derivative at the fixed point $b=1$, where only the
Gamma prefactor $\mathfrak p$ contributes.
(2) This route gives only the combination of single sums that $\LL_2$ needs, not the two
relations \eqref{gff:gc:L2singles} separately.
(3) The duality \eqref{gff:mgp:duality} is proved for four different parameters, but the lattice
step with four different rim exponents was not written out; the resulting four-exponent
ratio is the numerical statement listed under the remaining gaps of
Section~\ref{gff:gc:d2closedForms}.
(4) Checks, not proof steps: \path{checks_all.log} records the residue functional
\eqref{gff:gc:L2residue} in exact algebra; \eqref{gff:mgp:rimSpoke} at $x=\frac12,\frac25,\frac13,\frac3{10},
0.27,0.6,0.7$ to all $35$ working digits; $(\log\mathfrak p)'(\frac12)=-4\log2$ and the reflection
ratio exactly; \eqref{gff:mgp:duality} at complex $\tau$ and \eqref{gff:mgp:BarnesLemma} at complex
parameters to about $20$ digits; and the identities $F_{\rm R}=K_nF_{\rm L}$ and the prefactor
ratio of step (a) exactly for $n=0,\dots,6$. The script \path{extra_checks.wl} of the same folder
tests the Gamma bound on random samples.


\subsubsection{The elliptic moments by half-strip contours}
\label{gff:mgp:halfStrip}
The two further proofs below use the elliptic moments $U_0$, $V_0$, $U$, $V$, $R$ defined
before \eqref{gff:lp:eq:fold}. This subsubsection evaluates all of them by one method, which does
not use the Stieltjes identity \eqref{gff:lp:eq:stieltjes}. The notation is the one introduced with
these moments: the modulus convention for $\mathbf K$, $f(k)=\mathbf K(k)$, $g(k)=\mathbf K(k')$,
the weight $w(k)=\log\frac k{1-k^2}$, and the moments $M,N$, which are $M_{\rm e},N_{\rm e}$ of
\eqref{gff:fp:MN}; $T(k)=\frac{1-k}{1+k}$ is the Landen map of \eqref{gff:lp:eq:measures}.

\paragraph{Proposition (elliptic moments).}
\begin{equation}\begin{gathered}
 U_0=2N_{\rm e},\qquad V_0=2M_{\rm e},\qquad R=\frac\pi4N_{\rm e}-\log2\,M_{\rm e},\\
 U+\log2\,U_0=-\pi M_{\rm e},\qquad V+\log2\,V_0=-\frac\pi2N_{\rm e}.
\end{gathered}\label{gff:mgp:ellipticMoments}\end{equation}
\emph{Equivalently \eqref{gff:lp:eq:U}, \eqref{gff:lp:eq:V}, \eqref{gff:lp:eq:R}, $U_0=2N_{\rm e}$ and
$V_0=2M_{\rm e}$ hold, and they are proved here without \eqref{gff:lp:eq:stieltjes}.}

\emph{Proof.} \emph{Angle form.} With $k=\sin\beta$ one has $\dd k/k'=\dd\beta$ and
$kk'=\frac12\sin2\beta$. So, with $F(\beta)=\mathbf K(\sin\beta)$, $G(\beta)=\mathbf K(\cos\beta)$ and
$L(\beta)=\log\sin2\beta$,
\[\begin{aligned}
 U_0&=\int_0^{\pi/2}F^2\dd\beta,& V_0&=\int_0^{\pi/2}FG\dd\beta,\\
 U+\log2\,U_0&=\int_0^{\pi/2}F^2L\dd\beta,& V+\log2\,V_0&=\int_0^{\pi/2}FGL\dd\beta .
\end{aligned}\]

(1) \emph{Holomorphy.} On the strip $\mathsf T=\{0<\Re\beta<\frac\pi2\}$ the functions $F$, $G$ and
$L$ (principal branch) are holomorphic. For $\beta=\alpha+\ii\eta$ with $0<\alpha<\frac\pi2$,
$\sin\beta=\sin\alpha\cosh\eta+\ii\cos\alpha\sinh\eta$ has positive real part and is real only for
$\eta=0$, where it lies in $(0,1)$; hence $\sin^2\beta$ is real only there and avoids the cut
$[1,\infty)$ of $\mathbf K$ as a function of $k^2$. The same holds for
$\cos\beta=\cos\alpha\cosh\eta-\ii\sin\alpha\sinh\eta$. Next, $\sin2\beta=\sin2\alpha\cosh2\eta+\ii\cos2\alpha\sinh2\eta$ is real only
if $\eta=0$ or $\alpha=\frac\pi4$, and then it is positive.

(2) \emph{Boundary values.} For $y>0$ put $\kappa=\tanh y$ and $\kappa'=\operatorname{sech}y$. The
imaginary-modulus transformation $\mathbf K(\ii k/k')=k'\mathbf K(k)$ and the reciprocal-modulus
transformation
$\mathbf K(X+\ii0)=X^{-1}[\mathbf K(X^{-1})+\ii\mathbf K(\sqrt{1-X^{-2}})]$ for $X>1$
\cite{gff:DLMFEllipticConnection} give, as limits from inside $\mathsf T$,
\begin{equation}\begin{aligned}
 F(\ii y)&=\kappa'f(\kappa),& F(\tfrac\pi2+\ii y)&=\kappa'[g(\kappa)+\ii f(\kappa)],&
 L(\ii y)&=w(\kappa)+\log2+\tfrac{\ii\pi}2,\\
 G(\ii y)&=\kappa'[g(\kappa)-\ii f(\kappa)],& G(\tfrac\pi2+\ii y)&=\kappa'f(\kappa),&
 L(\tfrac\pi2+\ii y)&=w(\kappa)+\log2-\tfrac{\ii\pi}2 .
\end{aligned}\label{gff:mgp:strip}\end{equation}
Indeed $\sin(\epsilon+\ii y)\to\ii\sinh y$ with $\ii\sinh y=\ii\kappa/\kappa'$;
$\sin(\frac\pi2-\epsilon+\ii y)\to\cosh y+\ii0$; $\cos(\epsilon+\ii y)\to\cosh y-\ii0$;
$\cos(\frac\pi2-\epsilon+\ii y)\to-\ii\sinh y$, and $\mathbf K$ is even; $\sin2(\epsilon+\ii y)\to\ii\sinh2y$ and
$\sin2(\frac\pi2-\epsilon+\ii y)\to-\ii\sinh2y$, both from the right half-plane; finally
$\log\sinh2y=\log(2\kappa/\kappa'^2)=w(\kappa)+\log2$.

(3) \emph{Bounds.} Let $\Theta\in\{F^2,FG,F^2L,FGL\}$. With the boundary values included,
$|\Theta(\beta)|\le C(1+|\log|\beta||^3+|\log|\beta-\frac\pi2||^3)$ for $0\le\Im\beta\le1$, because
$\mathbf K(k)=O(1+|\log(1-k^2)|)$ uniformly as $k^2\to1$ in $\mathbb C\setminus[1,\infty)$ (from the
expansion in powers of $k'^2$, DLMF 19.12.1) and $L(\beta)=O(1+|\log|\beta||+|\log|\beta-\frac\pi2||)$.
For $\Im\beta\geq1$, $|\Theta(\beta)|\le C(1+\Im\beta)^3e^{-2\Im\beta}$, because
$|\sin\beta|,|\cos\beta|\geq\sinh\Im\beta$, $\mathbf K(k)=O(|k|^{-1}\log|k|)$ as $|k|\to\infty$
\cite{gff:DLMFEllipticConnection}, and $L=O(1+\Im\beta)$. These two uniform estimates, which include
the one-sided boundary values of (2), are standard and are stated here without further detail.

(4) \emph{Half-strip identity.} Cauchy's theorem on the rectangle
$[\epsilon,\frac\pi2-\epsilon]\times[0,Y]$, followed by $\epsilon\to0$ (dominated convergence with the
bounds of (3) and locally uniform convergence to the boundary values) and $Y\to\infty$ (the top
edge is $O(Y^3e^{-2Y})$), gives
\begin{equation}
 \int_0^{\pi/2}\Theta(\beta)\dd\beta=-\ii\int_0^\infty\bigl[\Theta(\tfrac\pi2+\ii y)-\Theta(\ii y)\bigr]\dd y .
 \label{gff:mgp:halfStripId}
\end{equation}

(5) \emph{Evaluation.} Write $f,g$ for $f(\kappa),g(\kappa)$ and $\omega=w(\kappa)+\log2$; all integrals
below are over $\kappa\in(0,1)$ and converge, since the integrands have at most $\log^3$ endpoint
singularities. With $\dd y=\dd\kappa/\kappa'^2$ and \eqref{gff:mgp:strip}, \eqref{gff:mgp:halfStripId} gives
\[\begin{aligned}
 U_0&=2\!\int fg-\ii\!\int(g^2-2f^2),&\qquad V_0&=2\!\int f^2,\\
 U+\log2\,U_0&=\int\bigl[2fg\,\omega-\tfrac\pi2g^2\bigr]-\ii\int\bigl[(g^2-2f^2)\omega+\pi fg\bigr],&\qquad
 V+\log2\,V_0&=\int\bigl[2f^2\omega-\pi fg\bigr].
\end{aligned}\]
The left sides are real, so the imaginary parts vanish:
$U_0=2N_{\rm e}$, $\int g^2=2M_{\rm e}$, $V_0=2M_{\rm e}$, $\int(g^2-2f^2)\omega=-\pi N_{\rm e}$, and
$U+\log2\,U_0=2\int fg\,\omega-\pi M_{\rm e}$, $V+\log2\,V_0=2\int f^2\omega-\pi N_{\rm e}$.

(6) \emph{Landen.} By \eqref{gff:lp:eq:measures}, $f(T)g(T)|T'|=fg$ and $g(T)^2|T'|=2f^2$, and
$w(T(\kappa))=-w(\kappa)-2\log2$ gives $\omega(T(\kappa))=-\omega(\kappa)$. Substituting $\kappa\to T(\kappa)$
gives $\int fg\,\omega=-\int fg\,\omega=0$ and $\int g^2\omega=-2\int f^2\omega$. With (5),
$-4\int f^2\omega=-\pi N_{\rm e}$, that is $\int f^2\omega=\frac\pi4N_{\rm e}$. Therefore
$U+\log2\,U_0=-\pi M_{\rm e}$, $V+\log2\,V_0=\frac\pi2N_{\rm e}-\pi N_{\rm e}=-\frac\pi2N_{\rm e}$ and
$R=\int f^2w=\int f^2\omega-\log2\,M_{\rm e}=\frac\pi4N_{\rm e}-\log2\,M_{\rm e}$. With $U_0=2N_{\rm e}$
and $V_0=2M_{\rm e}$ these are $U=-\pi M_{\rm e}-2\log2\,N_{\rm e}$, $V=-\frac\pi2N_{\rm e}-2\log2\,M_{\rm e}$
and $R=\frac\pi4N_{\rm e}-\log2\,M_{\rm e}$, that is \eqref{gff:lp:eq:U}, \eqref{gff:lp:eq:V} and
\eqref{gff:lp:eq:R}. $\square$

The half-strip contours thus reprove the three logarithmic evaluations of the
one-dimensional wheel by a single method. Conversely, \eqref{gff:mgp:ellipticMoments} follows from
\eqref{gff:lp:eq:U}, \eqref{gff:lp:eq:V}, \eqref{gff:lp:eq:R} and $U_0=2N_{\rm e}$, $V_0=2M_{\rm e}$ as proved in
Section~\ref{gff:app:LreductionProof}; that second proof uses \eqref{gff:lp:eq:stieltjes}. As a check,
not a proof step, a $60$-digit run (symmetry, Slurm job 567960;
\path{LRI_C5_d2_simplify/proof/L2-fourier/check_strip_dps60.log}) confirms the boundary
values \eqref{gff:mgp:strip} to $4\times10^{-62}$, the four strip identities and the Landen steps to
$5\times10^{-63}$, and the identities for $U+\log2\,U_0$ and $V+\log2\,V_0$ to $2\times10^{-63}$ and
$2\times10^{-65}$.


\subsubsection{A third proof: Parseval pairings on the unit circle}
\label{gff:mgp:parseval}
This proof keeps the rim derivative \eqref{gff:gc:L2derivative}, the residue functional
\eqref{gff:gc:L2residue}, the jets \eqref{gff:gc:L2taylor} and the reduction \eqref{gff:gc:L2reduction} of
Section~\ref{gff:gc:L2proof}. It evaluates the four single sums $a=\pi^2\sum_kc_k^4$,
$b=-2\pi\sum_kc_k^4\alpha_k$, $\mathcal S_1$ and $\mathcal S_2$ differently: each is a Parseval
pairing of two power series on the unit circle, and the circle values turn the pairings into the
elliptic moments $U_0$, $V_0$, $U$, $V$ defined before \eqref{gff:lp:eq:fold}. With
\eqref{gff:mgp:ellipticMoments} no step uses \eqref{gff:lp:eq:stieltjes}. Here $c_n$, $\alpha_n$,
$\beta_n$ are as in Section~\ref{gff:gc:L2proof},
\[
 y_1(z)={}_2F_1(\tfrac12,\tfrac12;1;z)=\sum_{n\geq0}c_n^2z^n,\qquad y_3(z)={}_2F_1(\tfrac12,\tfrac12;1;1-z),
\]
and, as in the upper-semicircle contour that proves \eqref{gff:lp:eq:V},
$f_\vartheta=\mathbf K(\sin\vartheta)$ and $g_\vartheta=\mathbf K(\cos\vartheta)=f_{\pi/2-\vartheta}$ for
$0<\vartheta<\frac\pi2$.

\paragraph{Generating functions on the slit disk.}
For $|z|<1$, $z\notin(-1,0]$, with principal logarithms,
\begin{equation}\begin{aligned}
 \sum_{n\geq0}c_n^2\alpha_nz^n&=-\frac\pi2\,y_3(z)-\frac12\log z\;y_1(z),\\
 \sum_{n\geq0}c_n^2\beta_nz^n&=-\frac12\log(1-z)\,y_1(z),\\
 \sum_{n\geq0}c_n^2\bigl[4\alpha_n\beta_n+2\psi_1(n+\tfrac12)\bigr]z^n
 &=\pi\log(1-z)\,y_3(z)+\pi^2y_1(z)+\log z\,\log(1-z)\,y_1(z).
\end{aligned}\label{gff:mgp:gfDisk}\end{equation}
At $z=k^2$ with $0<k<1$, where $y_1=\frac2\pi f(k)$ and $y_3=\frac2\pi g(k)$, these are the three
lines of \eqref{gff:gc:L2generating}. The left sides are power series with polynomially bounded
coefficients, hence analytic for $|z|<1$. The right sides are analytic on the slit disk, since
$y_1$ is analytic off $[1,\infty)$, $y_3$ off $(-\infty,0]$ and $\log z$ off $(-\infty,0]$. The slit
disk is connected and the two sides agree on $(0,1)$, so they agree on the whole slit disk.

\paragraph{Lemma (values on the circle).}
\emph{For $0<\vartheta<\frac\pi4$,}
\begin{equation}
 y_1(e^{4\ii\vartheta})=\frac1\pi\,e^{-\ii\vartheta}\bigl[g_\vartheta+\ii f_\vartheta\bigr],\qquad
 y_3(e^{-4\ii\vartheta})=\frac2\pi\,e^{\ii\vartheta}f_\vartheta .
 \label{gff:mgp:circle}
\end{equation}
\emph{Proof.} The function $y_1$ is analytic on $\mathbb C\setminus[1,\infty)$, and
$y_1(e^{4\ii\vartheta})=\frac2\pi\mathbf K(e^{2\ii\vartheta})$. The boundary-value Landen relation
$e^{\ii\vartheta}\mathbf K(e^{2\ii\vartheta})=\frac12[\mathbf K(\cos\vartheta)+\ii\mathbf K(\sin\vartheta)]$, proved with the
upper-semicircle contour behind \eqref{gff:lp:eq:V}, gives the first identity. (The
coefficients $1$ and $\ii$ of the reciprocal-modulus value used there can be confirmed at the
endpoint. For $s>1$, $y_1(s+\ii0)$ solves the hypergeometric equation on $(1,\infty)$, whose
solutions are spanned by $s^{-1/2}y_1(1/s)$ and $s^{-1/2}y_1(1-1/s)$. As $s\downarrow1$ the first
behaves like $\pi^{-1}\log\frac{16}{s-1}$ and the second tends to $1$, while
$y_1(z)=\pi^{-1}[4\log2-\log(1-z)]+o(1)$ as $z\to1$ \cite{gff:DLMF15} gives
$y_1(s+\ii0)=\pi^{-1}[\log\frac{16}{s-1}+\ii\pi]+o(1)$. So the coefficients are $1$ and $\ii$.)
For the second identity, both sides of
${}_2F_1(\frac12,\frac12;1;1-w^2)=\frac2{1+w}\,{}_2F_1(\frac12,\frac12;1;(\frac{1-w}{1+w})^2)$ are
analytic for $\Re w>0$, where $w^2\notin(-\infty,0]$ and $|\frac{1-w}{1+w}|<1$, and they agree on
$(0,1)$, where this is the real Landen identity $f(T(k))=\frac{1+k}2g(k)$ behind
\eqref{gff:lp:eq:measures}. Take $w=e^{-2\ii\vartheta}$. Then $1-w^2=1-e^{-4\ii\vartheta}$,
$(1-w)/(1+w)=\ii\tan\vartheta$ and $2/(1+w)=e^{\ii\vartheta}/\cos\vartheta$, and Pfaff's transformation
\cite{gff:DLMF15} gives
${}_2F_1(\frac12,\frac12;1;-\tan^2\vartheta)=\cos\vartheta\,{}_2F_1(\frac12,\frac12;1;\sin^2\vartheta)
=\frac2\pi\cos\vartheta\,f_\vartheta$. $\square$

\paragraph{Proposition (pairings).}
\begin{equation}
 a=\frac{4U_0}\pi,\qquad b=\frac{4V_0}\pi,\qquad
 \mathcal S_1=-\frac2{\pi^3}\bigl(U+2\log2\,U_0\bigr),\qquad
 \mathcal S_2=\frac a2+\frac4{\pi^2}\bigl(V+2\log2\,V_0\bigr).
 \label{gff:mgp:pairings}
\end{equation}
\emph{Proof.} If $\mathsf p(z)=\sum_np_nz^n$ and $\mathsf q(z)=\sum_nq_nz^n$ have real
square-summable coefficients, they lie in the Hardy space $H^2$, have radial limits almost
everywhere, and
$\sum_np_nq_n=\frac1{2\pi}\int_{-\pi}^{\pi}\mathsf p(e^{\ii\theta})\,\mathsf q(e^{-\ii\theta})\dd\theta$. All
coefficient sequences used below are square summable: $c_n^2=O(1/n)$,
$c_n^2\beta_n=O(\log n/n)$, $c_n^2\alpha_n=O(n^{-2})$ and
$c_n^2[4\alpha_n\beta_n+2\psi_1(n+\frac12)]=O(\log n/n^2)$. For $\theta\in(-\pi,\pi)\setminus\{0\}$ the
radius to $e^{\ii\theta}$ lies in the slit disk, and the right sides of \eqref{gff:mgp:gfDisk} are
continuous up to $e^{\ii\theta}$, so the radial limits are the values of the closed forms. On the
circle $z=e^{\ii\theta}$ one has $\log\bar z=-\ii\theta$ and
$\log(1-z)+\log(1-\bar z)=2\log|2\sin\frac\theta2|$; $\Im\log(1-\bar z)$ is odd in $\theta$;
$|y_1(z)|^2$ is even; and $y_1(e^{-\ii\theta})y_3(e^{\ii\theta})=\overline{y_1(e^{\ii\theta})y_3(e^{-\ii\theta})}$,
since $y_1,y_3$ have real Taylor coefficients. All integrands below have only integrable
logarithmic singularities. For $0<\theta<\pi$ and $\theta=4\vartheta$, \eqref{gff:mgp:circle} gives
$|y_1(e^{\ii\theta})|^2=\pi^{-2}(f_\vartheta^2+g_\vartheta^2)$ and
$y_1(e^{\ii\theta})y_3(e^{-\ii\theta})=2\pi^{-2}(f_\vartheta g_\vartheta+\ii f_\vartheta^2)$. The fold
$\vartheta\mapsto\frac\pi2-\vartheta$ exchanges $f$ and $g$ and fixes $\sin2\vartheta$; so for any
weight $\chi(\sin2\vartheta)$ it turns $\int_0^{\pi/4}(f^2+g^2)\chi\dd\vartheta$ into
$\int_0^{\pi/2}f^2\chi\dd\vartheta$ and $\int_0^{\pi/4}fg\,\chi\dd\vartheta$ into
$\frac12\int_0^{\pi/2}fg\,\chi\dd\vartheta$. Finally, by the angle form in the proof of
\eqref{gff:mgp:ellipticMoments} (with $\beta=\vartheta$),
$\int_0^{\pi/2}f_\vartheta^2\log(2\sin2\vartheta)\dd\vartheta=U+2\log2\,U_0$ and
$\int_0^{\pi/2}f_\vartheta g_\vartheta\log(2\sin2\vartheta)\dd\vartheta=V+2\log2\,V_0$.

(i) $a/\pi^2=\sum c_n^2\cdot c_n^2=\frac1{2\pi}\int|y_1|^2\dd\theta
=\frac4{\pi^3}\int_0^{\pi/4}(f_\vartheta^2+g_\vartheta^2)\dd\vartheta=\frac{4U_0}{\pi^3}$.

(ii) $-b/(2\pi)=\sum c_n^2\cdot c_n^2\alpha_n$ is the pairing of $y_1$ with the first line of
\eqref{gff:mgp:gfDisk}. The integrand is $-\frac\pi2y_1(z)y_3(\bar z)+\frac{\ii\theta}2|y_1(z)|^2$, and the
second term is odd. So
$b=\pi\Re\int_0^\pi y_1(e^{\ii\theta})y_3(e^{-\ii\theta})\dd\theta=\frac8\pi\int_0^{\pi/4}f_\vartheta g_\vartheta\dd\vartheta
=\frac{4V_0}\pi$.

(iii) $\mathcal S_1=\sum c_n^2\cdot c_n^2\beta_n=-\frac1{4\pi}\int\log(1-\bar z)\,|y_1(z)|^2\dd\theta$, by
the second line of \eqref{gff:mgp:gfDisk}. Dropping the odd part and setting $\theta=4\vartheta$,
\[
 \mathcal S_1=-\frac2{\pi^3}\int_0^{\pi/4}\log(2\sin2\vartheta)(f_\vartheta^2+g_\vartheta^2)\dd\vartheta
 =-\frac2{\pi^3}\int_0^{\pi/2}\log(2\sin2\vartheta)f_\vartheta^2\dd\vartheta
 =-\frac2{\pi^3}\bigl(U+2\log2\,U_0\bigr).
\]

(iv) Since
$c_n^4[4\alpha_n\beta_n+\psi_1(n+\frac12)]=(2c_n^2\beta_n)(c_n^2\alpha_n)+c_n^2\cdot\frac12c_n^2[4\alpha_n\beta_n+2\psi_1(n+\frac12)]$,
$\mathcal S_2$ is the sum of two pairings. By \eqref{gff:mgp:gfDisk} the combined integrand is
\[
 \pi\log|2\sin\tfrac\theta2|\;y_1(z)y_3(\bar z)+\frac{\pi^2}2|y_1(z)|^2
 -\ii\theta\log|2\sin\tfrac\theta2|\,|y_1(z)|^2 .
\]
The last term is odd. The middle one contributes $\frac{\pi^2}2\cdot\frac a{\pi^2}=\frac a2$. The first
one contributes
$\frac2{\pi^2}\int_0^\pi\log(2\sin\frac\theta2)f_{\theta/4}g_{\theta/4}\dd\theta
=\frac4{\pi^2}\int_0^{\pi/2}\log(2\sin2\vartheta)f_\vartheta g_\vartheta\dd\vartheta=\frac4{\pi^2}(V+2\log2\,V_0)$.
$\square$

\paragraph{Conclusion.}
Inserting \eqref{gff:mgp:ellipticMoments} into \eqref{gff:mgp:pairings} gives $a=8N_{\rm e}/\pi$ and
$b=8M_{\rm e}/\pi$, which reproves \eqref{gff:fp:abWheel}, and
\[
 \mathcal S_1=\frac{2M_{\rm e}}{\pi^2}-\frac{4\log2\,N_{\rm e}}{\pi^3}=\frac{\pi b-2a\log2}{4\pi^2},\qquad
 \mathcal S_2=\frac{2N_{\rm e}}\pi+\frac{8\log2\,M_{\rm e}}{\pi^2}=\frac a4+\frac{b\log2}\pi ,
\]
which are \eqref{gff:gc:L2singles}. With \eqref{gff:gc:L2reduction} this gives $\LL_2=-2\log2\,P_4(2)$
exactly as at the end of Section~\ref{gff:gc:L2proof}, and with \eqref{gff:fp:P4two} it proves
\eqref{gff:gc:L2closed}. The inputs are Euler's and Pfaff's transformations, the connection formulas
of \cite{gff:DLMF15} behind \eqref{gff:gc:L2generating}, the Landen transformations, Parseval's identity in
$H^2$, the boundary-value Landen relation and the half-strip contours of
Section~\ref{gff:mgp:halfStrip}; the Stieltjes identity is not used.

\paragraph{Remarks.}
(1) The pairings \eqref{gff:mgp:pairings} use only the circle values. Given $U_0=2N_{\rm e}$ and
$V_0=2M_{\rm e}$, the first relation of \eqref{gff:gc:L2singles} is therefore equivalent to
$U+\log2\,U_0=-\frac\pi2V_0$, that is to \eqref{gff:lp:eq:U}, and the second to
$V+\log2\,V_0=-\frac\pi4U_0$, that is to \eqref{gff:lp:eq:V}. The two single-sum relations are the
unit-circle images of the logarithmic evaluations of the one-dimensional wheel.
(2) Two relations among second-order logarithmic moments were found numerically while this route
was developed; they are not used. With $\langle\chi\rangle=\int_0^1f(k)g(k)\chi(k)\dd k$,
$2\langle\log^2k\rangle-\langle\log k\log(1-k)\rangle-3\langle\log k\log(1+k)\rangle
=\frac{\pi^2}2N_{\rm e}+\pi\log2\,M_{\rm e}$, which is equivalent to the second relation of
\eqref{gff:gc:L2singles} through a second-order moment representation that has only been sketched,
and
$32\langle\log^2(1+k)\rangle=22\langle\log^2k\rangle-7\pi^2N_{\rm e}+8\log^22\,N_{\rm e}$. Both come
from an integer-relation search at $40$ digits
(\path{LRI_C5_d2_simplify/proof/L2-elliptic/probe_moments.py}) and are numerical statements only.
(3) Checks, not proof steps: the scripts \path{check_parseval.py}, \path{collapse_sympy.py},
\path{check_collapse.py} and \path{check_rimderivative.py} of the same folder test the generating
functions, the circle values and the pairings at $30$ digits, the rim derivative
\eqref{gff:gc:L2derivative} with the $70$-digit single sums, and the collapse algebra exactly (no logs
of these runs are shipped); the synthesis check \path{LRI_C5_d2_simplify/proof/synth/synth_checks.log}
gives, with the $70$-digit single sums, residuals $9\times10^{-72}$ and $1.1\times10^{-69}$ for the
two relations \eqref{gff:gc:L2singles}, $7\times10^{-68}$ for \eqref{gff:gc:L2reduction}, and
$1.7\times10^{-67}$ for $\LL_2+16\log2\,ab/\pi$.


\subsubsection{A fourth proof: the Fourier representation of the two modes}
\label{gff:mgp:fourier}
This proof keeps the rim derivative \eqref{gff:gc:L2derivative} and the residue functional
\eqref{gff:gc:L2residue}, but obtains the modes $\mathcal A$, $\mathcal B$ and their $x$-derivatives at
$x=\frac12$ without single sums and without growth estimates. Cauchy's beta integral writes the
Fourier coefficients of $\Omega(t;x)$ as convolutions of $(2\cos\frac\theta2)^{-x}$ on $(-\pi,\pi)$,
which have compact support, and at $x=\frac12$ these convolutions are complete elliptic
integrals. Throughout,
$I_{\rm F}=(\frac38,\frac58)$, which contains $I$, and $D_{\rm F}=\{x\in\mathbb C:|x-\frac12|<\frac18\}$. For
real $x\in I_{\rm F}$ put
\[\begin{gathered}
 h_x(\sigma)=\frac1{\Gamma(1-\frac x2+\sigma)\,\Gamma(1-\frac x2-\sigma)},\qquad
 \mathsf P_x(\theta)=(2\cos\tfrac\theta2)^{-x}\,\mathbf 1_{(-\pi,\pi)}(\theta),\\
 \mathsf Q_x=\mathsf P_x*\mathsf P_x,\qquad v_m(x)=(\mathsf Q_x*\mathsf Q_x)(2\pi m).
\end{gathered}\]
Thus $h_x=\mathcal G_x^{-1}$ in the notation of \eqref{gff:mgp:twoModes}. By reflection,
$h_x(\sigma)=\sin\pi(1-\frac x2-\sigma)\,\Gamma(\frac x2+\sigma)/[\pi\Gamma(1-\frac x2+\sigma)]$, so $h_x$ is
real, even and continuous on $\mathbb R$ with $|h_x(\sigma)|\le C(1+|\sigma|)^{x-1}$; hence
$h_x^2\in L^2(\mathbb R)$ and $h_x^4\in L^1(\mathbb R)$. The terms of $\Omega(t;x)$ satisfy
$\pi^4[\Gamma(1+s)\Gamma(1-x-s)]^{-4}=\pi^4h_x(\sigma)^4$ at $s=\sigma-\frac x2$.

\paragraph{Lemma (Cauchy's beta integral).}
\emph{For real $c>-1$ and real $\sigma$,}
\[
 \int_{-\pi}^{\pi}(2\cos\tfrac\theta2)^{c}\,e^{\ii\sigma\theta}\dd\theta
 =\frac{2\pi\,\Gamma(c+1)}{\Gamma(1+\frac c2+\sigma)\,\Gamma(1+\frac c2-\sigma)} .
\]
\emph{Proof.} Put $s=\sigma-\frac c2$ and $w=e^{\ii\theta}$. For $|\theta|<\pi$,
$1+w=2\cos\frac\theta2\,e^{\ii\theta/2}$, so with principal powers
$(2\cos\frac\theta2)^ce^{\ii\sigma\theta}=(1+w)^cw^s$, and the integral equals
$-\ii\oint_{|w|=1}(1+w)^cw^{s-1}\dd w$. Let $\Re s>0$. The integrand is holomorphic in the disk
slit along $(-1,0]$ and integrable at $w=0$. For $0<\rho<1$, collapsing the circle $|w|=\rho$
onto the two sides of the slit $[-\rho,0]$, where $w^{s-1}=r^{s-1}e^{\pm\ii\pi(s-1)}$ at $w=-r$,
gives $\oint_{|w|=\rho}(1+w)^cw^{s-1}\dd w=2\ii\sin(\pi s)\int_0^\rho(1-r)^cr^{s-1}\dd r$. As
$\rho\to1$ the right side tends to $2\ii\sin(\pi s)\,\mathrm B(s,c+1)$; the left side tends to the
integral over the unit circle by dominated convergence, with the majorant
$|1+\rho e^{\ii\theta}|^c\le(4\rho)^{c/2}|\cos\frac\theta2|^c$ for $-1<c<0$ (since
$|1+\rho e^{\ii\theta}|^2-4\rho\cos^2\frac\theta2=(1-\rho)^2$) and $|1+\rho e^{\ii\theta}|^c\le2^c$ for
$c\geq0$. Hence the integral is
$2\sin(\pi s)\Gamma(s)\Gamma(c+1)/\Gamma(s+c+1)=2\pi\Gamma(c+1)/[\Gamma(1-s)\Gamma(1+c+s)]$, which is the
claim. Both sides are entire in $\sigma$, so the restriction $\Re s>0$ is removed by continuation.
$\square$

With $c=-x$ the lemma reads
$h_x(\sigma)=(2\pi\Gamma(1-x))^{-1}\int_{\mathbb R}\mathsf P_x(\theta)e^{\ii\sigma\theta}\dd\theta$.

\paragraph{Lemma (the twofold convolution).}
\emph{For $0<|\phi|<2\pi$ put $k=\cos\frac\phi4$ and $k'=|\sin\frac\phi4|$. For real $x\in I_{\rm F}$,}
\begin{equation}
 \mathsf Q_x(\phi)=2\int_0^{\pi/2}\frac{(2k\cos\eta)^{1-2x}}{\sqrt{1-k^2\sin^2\eta}}\dd\eta
 =(2k)^{1-2x}\,\mathrm B(\tfrac12,1-x)\,{}_2F_1(\tfrac12,\tfrac12;\tfrac32-x;k^2).
 \label{gff:mgp:Qform}
\end{equation}
\emph{Both expressions define holomorphic functions of $x\in D_{\rm F}$ that agree there, and,
uniformly for $x\in D_{\rm F}$,}
\begin{equation}
 |\mathsf Q_x(\phi)|\le C\,|\sin\tfrac\phi2|^{-1/4},\qquad
 |\partial_x\mathsf Q_x(\phi)|\le C\,|\sin\tfrac\phi2|^{-1/3}.
 \label{gff:mgp:Qbounds}
\end{equation}
\emph{At $x=\frac12$,}
\begin{equation}
 \mathsf Q_{1/2}(\phi)=2\mathbf K(k),\qquad
 \partial_x\mathsf Q_x(\phi)\big|_{x=1/2}=\pi\mathbf K(k')-2\mathbf K(k)\log(4kk').
 \label{gff:mgp:Qhalf}
\end{equation}
\emph{Proof.} By evenness take $0<\phi<2\pi$. The integrand of $\mathsf P_x*\mathsf P_x$ is supported
on $\theta\in(\phi-\pi,\pi)$, and $4\cos\frac\theta2\cos\frac{\phi-\theta}2=2[\cos\psi-\cos\psi_0]$ with
$\psi=\theta-\frac\phi2$ and $\psi_0=\pi-\frac\phi2$. So
$\mathsf Q_x(\phi)=\int_{-\psi_0}^{\psi_0}[2(\cos\psi-\cos\psi_0)]^{-x}\dd\psi$. The substitution
$\sin\frac\psi2=k\sin\eta$ with $k=\sin\frac{\psi_0}2=\cos\frac\phi4$ gives
$2(\cos\psi-\cos\psi_0)=4k^2\cos^2\eta$ and $\dd\psi=2k\cos\eta\,\dd\eta/\sqrt{1-k^2\sin^2\eta}$, which
is the $\eta$-integral. Expanding $(1-k^2\sin^2\eta)^{-1/2}$ and integrating termwise with
$\int_0^{\pi/2}\sin^{2n}\eta\cos^{1-2x}\eta\dd\eta=\Gamma(n+\frac12)\Gamma(1-x)/[2\Gamma(n+\frac32-x)]$ gives
the ${}_2F_1$. For fixed $\phi$ both expressions are holomorphic in $x\in D_{\rm F}$ and agree for
real $x$, hence on $D_{\rm F}$.

Bounds: for $x\in D_{\rm F}$, $\xi=\Re(1-2x)\in(-\frac14,\frac14)$, and since $k\cos\eta\le1$,
$|(2k\cos\eta)^{1-2x}|=(2k\cos\eta)^\xi\le2(k\cos\eta)^{-1/4}$ and
$|\partial_x(2k\cos\eta)^{1-2x}|\le4|\log(2k\cos\eta)|(k\cos\eta)^{-1/4}\le C(k\cos\eta)^{-1/3}$. For
$0<s'<1$, $\int_0^{\pi/2}(\cos\eta)^{-s'}(1-k^2\sin^2\eta)^{-1/2}\dd\eta\le C_{s'}k'^{-s'}$: on
$[0,\frac\pi4]$ the integrand is bounded, and on $[\frac\pi4,\frac\pi2]$ one uses
$(1-k^2\sin^2\eta)^{1/2}\geq\max(\cos\eta,k'/\sqrt2)$ and the variable $\cos\eta$. Since
$kk'=\frac12|\sin\frac\phi2|$, \eqref{gff:mgp:Qbounds} follows, and differentiation under the
$\eta$-integral is justified by the same majorants.

At $x=\frac12$ the $\eta$-integral is $2\mathbf K(k)$, and
$\partial_x\mathsf Q_x|_{x=1/2}=-4\log(2k)\mathbf K(k)-4\int_0^{\pi/2}\log\cos\eta\,(1-k^2\sin^2\eta)^{-1/2}\dd\eta$.
The next lemma gives $-4\log(2k)\mathbf K(k)-2\mathbf K(k)\log\frac{k'}k+\pi\mathbf K(k')$, which is
\eqref{gff:mgp:Qhalf}. $\square$

\paragraph{Lemma (a logarithmic moment).}
\emph{For $0<k<1$,}
\[
 \int_0^{\pi/2}\frac{\log\cos\eta}{\sqrt{1-k^2\sin^2\eta}}\dd\eta
 =\frac12\mathbf K(k)\log\frac{k'}{k}-\frac\pi4\mathbf K(k') .
\]
\emph{Proof.} With $\eta=\operatorname{am}(u,k)$ the left side is $\int_0^{\mathbf K}\log\operatorname{cn}u\,\dd u$,
$\mathbf K=\mathbf K(k)$. The theta-quotient form of the Jacobian functions \cite{gff:DLMFTheta} gives
$\operatorname{cn}u=\sqrt{k'/k}\,\theta_2(\zeta,q)/\theta_4(\zeta,q)$ with $\zeta=\pi u/(2\mathbf K)$ and
$q=e^{-\pi\mathbf K(k')/\mathbf K(k)}$, and the product formulas \cite{gff:DLMFTheta} give
\[
 \log\frac{\theta_2(\zeta,q)}{\theta_4(\zeta,q)}=\log(2q^{1/4})+\log\cos\zeta
 +\sum_{n\geq1}\bigl[\log(1+2q^{2n}\cos2\zeta+q^{4n})-\log(1-2q^{2n-1}\cos2\zeta+q^{4n-2})\bigr].
\]
For $0\le r<1$, $\log(1\pm2r\cos2\zeta+r^2)=2\Re\log(1\pm re^{2\ii\zeta})$ has mean zero over
$[0,\pi]$, hence over $[0,\frac\pi2]$ by the symmetry $\zeta\mapsto\pi-\zeta$, and the series
converges uniformly. The mean of $\log\cos\zeta$ over $[0,\frac\pi2]$ is $-\log2$. So the mean of
$\log\operatorname{cn}$ over $[0,\mathbf K]$ is $\frac12\log\frac{k'}k+\frac14\log q$, and
$\int_0^{\mathbf K}\log\operatorname{cn}u\,\dd u=\frac{\mathbf K}2\log\frac{k'}k-\frac\pi4\mathbf K(k')$. $\square$

\paragraph{Proposition (two modes from compact Fourier support).}
\emph{For real $x\in I_{\rm F}$ and all $t\in\mathbb C$,
$\Omega(t;x)=\mathfrak a(x)+\mathfrak b(x)\cos2\pi(t+\frac x2)$, with}
\begin{equation}
 \mathfrak a(x)=\pi^4\!\int_{\mathbb R}h_x^4\dd\sigma=\frac{\pi\,v_0(x)}{8\,\Gamma(1-x)^4},\qquad
 \mathfrak b(x)=2\pi^4\!\int_{\mathbb R}h_x^4\cos(2\pi\sigma)\dd\sigma=\frac{\pi\,v_1(x)}{4\,\Gamma(1-x)^4}.
 \label{gff:mgp:fourierModes}
\end{equation}
\emph{Consequently, on $I$, $\mathcal A(x)=\sin^2(\pi x)\,v_0(x)/(8\pi)$ and
$\mathcal B(x)=\sin^2(\pi x)\,v_1(x)/(4\pi)$.}

\emph{Proof.} Let $\omega_m=\int_0^1\Omega(t;x)e^{-2\pi\ii mt}\dd t$. The terms of $\Omega$ are
integrable on the real line, so unfolding (Fubini) and $s=\sigma-\frac x2$ give
$\omega_m=\pi^4e^{\ii\pi mx}\int_{\mathbb R}h_x(\sigma)^4e^{-2\pi\ii m\sigma}\dd\sigma$. By the Cauchy beta
integral with $c=-x$ and Fubini ($\mathsf P_x\in L^1$),
$h_x(\sigma)^2=(2\pi\Gamma(1-x))^{-2}\int_{\mathbb R}\mathsf Q_x(\phi)e^{\ii\sigma\phi}\dd\phi$. Since
$\mathsf Q_x\in L^1\cap L^2$ by \eqref{gff:mgp:Qbounds}, Plancherel's theorem gives, with
$\widehat\chi(\theta)=\int\chi(\sigma)e^{-\ii\theta\sigma}\dd\sigma$, the identity
$\widehat{h_x^2}=\mathsf Q_x/(2\pi\Gamma(1-x)^2)$ in $L^2$. Parseval's identity for the $L^2$ functions
$h_x^2$ and $h_x^2e^{\ii\xi\sigma}$ then gives
\[
 \int_{\mathbb R}h_x^4e^{-\ii\xi\sigma}\dd\sigma
 =\frac1{2\pi}\int_{\mathbb R}\widehat{h_x^2}(\theta)\,\overline{\widehat{h_x^2}(\theta-\xi)}\dd\theta
 =\frac{(\mathsf Q_x*\mathsf Q_x)(\xi)}{8\pi^3\Gamma(1-x)^4}.
\]
Thus $\omega_m=\pi e^{\ii\pi mx}v_m(x)/(8\Gamma(1-x)^4)$. Now $\mathsf Q_x$ is even, nonnegative and
supported in $[-2\pi,2\pi]$, so $\mathsf Q_x*\mathsf Q_x$ is even and supported in $[-4\pi,4\pi]$, and at
$\xi=\pm4\pi$ the convolution integral runs over a single point. Hence $v_{-1}=v_1$ and $v_m=0$
for $|m|\geq2$. A continuous one-periodic function whose Fourier coefficients vanish for
$|m|\geq2$ equals the corresponding trigonometric polynomial (Fej\'er), so for real $t$,
$\Omega(t;x)=\frac\pi{8\Gamma(1-x)^4}[v_0(x)+2v_1(x)\cos(2\pi t+\pi x)]$, and by analyticity for all
$t$. The integral forms follow from $\omega_0$ and $\omega_1$. Finally
$N_x^2/\Gamma(1-x)^4=[\Gamma(x)\Gamma(1-x)]^{-2}=\sin^2(\pi x)/\pi^2$ and \eqref{gff:gc:L2modes} give the
expressions for $\mathcal A$ and $\mathcal B$. $\square$

For $x\in I$ this is a second proof of \eqref{gff:gc:L2modes}: the modes with $|m|\geq2$ are removed
by the compact support of $\mathsf Q_x*\mathsf Q_x$, without growth estimates. The same argument
should work for every real $x<\frac34$, where $\mathsf Q_x$ behaves like $|\phi|^{1-2x}$ near $0$ and is
still square integrable; this extension is a sketch, is not proved here and is not used (the
general case is \eqref{gff:mgp:twoModes}).

\paragraph{Lemma (the convolutions at $x=\frac12$).}
\emph{$v_0$ and $v_1$ extend holomorphically to $D_{\rm F}$, with
$v_m'(x)=\int_{\mathbb R}\partial_x[\mathsf Q_x(\phi)\mathsf Q_x(2\pi m-\phi)]\dd\phi$. At $x=\frac12$,}
\begin{equation}
 v_0=32\,U_0,\qquad v_1=16\,V_0,\qquad
 v_0'=32\bigl[\pi V_0-2U-4\log2\,U_0\bigr],\qquad
 v_1'=16\bigl[\pi U_0-2V-4\log2\,V_0\bigr].
 \label{gff:mgp:vHalf}
\end{equation}
\emph{Proof.} Use the holomorphic extension \eqref{gff:mgp:Qform}. By \eqref{gff:mgp:Qbounds},
$|\mathsf Q_x(\phi)\mathsf Q_x(2\pi m-\phi)|$ and its $x$-derivative are bounded, uniformly for
$x\in D_{\rm F}$, by $C|\sin\frac\phi2|^{-7/12}$ on the support ($|\phi|<2\pi$ for $m=0$,
$0<\phi<2\pi$ for $m=1$), which is integrable. Holomorphy and differentiation under the integral
follow. By evenness, $v_0'=4\int_0^{2\pi}(\partial_x\mathsf Q_x)\mathsf Q_x\dd\phi$ and
$v_1'=2\int_0^{2\pi}(\partial_x\mathsf Q_x)(\phi)\,\mathsf Q_x(2\pi-\phi)\dd\phi$. Substitute
$\phi=2\pi-4\beta$ with $\beta\in(0,\frac\pi2)$: then $\cos\frac\phi4=\sin\beta$,
$\cos\frac{2\pi-\phi}4=\cos\beta$ and $4kk'=2\sin2\beta$, so \eqref{gff:mgp:Qhalf} gives
$\mathsf Q_{1/2}(\phi)=2F(\beta)$, $\mathsf Q_{1/2}(2\pi-\phi)=2G(\beta)$ and
$\partial_x\mathsf Q_x(\phi)|_{1/2}=\pi G(\beta)-2F(\beta)[\log2+\log\sin2\beta]$, with $F,G$ as in the
proof of \eqref{gff:mgp:ellipticMoments}. With $\dd\phi=4\dd\beta$, the angle forms given there and
$\int_0^{\pi/2}G^2\dd\beta=U_0$, this gives \eqref{gff:mgp:vHalf}. $\square$

\paragraph{Conclusion.}
By \eqref{gff:mgp:fourierModes}, since $\sin^2\pi x$ has zero derivative at $x=\frac12$, and by
\eqref{gff:mgp:vHalf},
\[
 a=\frac{v_0(\frac12)}{8\pi}=\frac{4U_0}\pi,\qquad b=\frac{v_1(\frac12)}{4\pi}=\frac{4V_0}\pi,\qquad
 \mathcal A'(\tfrac12)=\frac{v_0'(\frac12)}{8\pi},\qquad \mathcal B'(\tfrac12)=\frac{v_1'(\frac12)}{4\pi}.
\]
With \eqref{gff:mgp:ellipticMoments}, $a=8N_{\rm e}/\pi$, $b=8M_{\rm e}/\pi$,
$-2U-4\log2\,U_0=2\pi M_{\rm e}-4\log2\,N_{\rm e}$ and $-2V-4\log2\,V_0=\pi N_{\rm e}-4\log2\,M_{\rm e}$,
so that
\begin{equation}
 \mathcal A'(\tfrac12)=16M_{\rm e}-\frac{16\log2}\pi N_{\rm e}=2\pi b-2\log2\,a,\qquad
 \mathcal B'(\tfrac12)=12N_{\rm e}-\frac{16\log2}\pi M_{\rm e}=\frac{3\pi}2a-2\log2\,b .
 \label{gff:mgp:AprimeBprime}
\end{equation}
Inserted into the expression for $\LL_2$ displayed after \eqref{gff:gc:L2residue},
\[
 \LL_2=-6a^2-8b^2+\frac4\pi\Bigl[(2\pi b-2\log2\,a)\,b+a\Bigl(\frac{3\pi}2a-2\log2\,b\Bigr)\Bigr]
 =-\frac{16\log2}\pi\,ab=-2\log2\,P_4(2).
\]
The rotation part $(2\pi b,\frac{3\pi}2a)$ of \eqref{gff:mgp:AprimeBprime} cancels $-6a^2-8b^2$ exactly,
and the diagonal part $-2\log2$ leaves $-2\log2$ times the coefficient $8/\pi$ of $ab$ in
$P_4(2)=8ab/\pi$. With \eqref{gff:gc:L2taylor}, \eqref{gff:mgp:AprimeBprime} is equivalent to
\eqref{gff:gc:L2singles}, which this route therefore proves once more; \eqref{gff:mgp:AprimeBprime} are
the relations quoted after \eqref{gff:gc:L2singles} and in \eqref{gff:gdr:landen}. The inputs are Cauchy's
beta integral, Plancherel's theorem, Stirling's formula, the theta-function form of
$\operatorname{cn}$ and the half-strip contours of Section~\ref{gff:mgp:halfStrip}; the Stieltjes
identity is not used.

\paragraph{Remark (a real-line form).}
Put $\mathfrak m_j=\int_{\mathbb R}h_{1/2}(\sigma)^4\cos(2\pi j\sigma)\dd\sigma$, so that $\mathfrak m_0=a/\pi^4$
and $\mathfrak m_1=b/(2\pi^4)$, and $\Psi(\sigma)=\psi_0(\frac34+\sigma)+\psi_0(\frac34-\sigma)$. Then
\[
 \int_{\mathbb R}h_{1/2}^4\,\Psi\dd\sigma=2\pi\mathfrak m_1-(5\log2+2\gamma)\mathfrak m_0,\qquad
 \int_{\mathbb R}h_{1/2}^4\,\Psi\cos(2\pi\sigma)\dd\sigma=\frac{3\pi}8\mathfrak m_0-(5\log2+2\gamma)\mathfrak m_1 .
\]
To see this, differentiate the integral forms \eqref{gff:mgp:fourierModes} under the integral. One has
$\partial_xh_x^4=2h_x^4\Psi_x$ with $\Psi_x(\sigma)=\psi_0(1-\frac x2+\sigma)+\psi_0(1-\frac x2-\sigma)$, which
is entire in $\sigma$ because the poles of $\psi_0(1-\frac x2-\sigma)$ sit at the fourth-order zeros of
$h_x^4$; the reflected form of $h_x$ gives the majorant
$|\partial_xh_x^4|\le C\log(2+|\sigma|)(1+|\sigma|)^{4x-4}$, uniformly for $x$ near $\frac12$. So
$\mathfrak a'(\frac12)=2\pi^4\int h_{1/2}^4\Psi$ and $\mathfrak b'(\frac12)=4\pi^4\int h_{1/2}^4\Psi\cos2\pi\sigma$.
Since $\mathfrak a=\mathcal A/N_x^2$, $\mathfrak b=\mathcal B/N_x^2$ and $(N_x^{-2})'(\frac12)=4\psi_0(\frac12)$,
\eqref{gff:mgp:AprimeBprime} gives $\mathfrak a'(\frac12)=2\pi b-(10\log2+4\gamma)a$ and
$\mathfrak b'(\frac12)=\frac{3\pi}2a-(10\log2+4\gamma)b$, which are the displayed identities. By
duplication and reflection,
$\Psi(\sigma)=\psi_0(\frac12+2\sigma)+\psi_0(\frac12-2\sigma)-2\log2+\pi\sec2\pi\sigma$, where the combination
is regular; so the first identity also reads
$\int_{\mathbb R}h_{1/2}^4[\psi_0(\frac12+2\sigma)+\psi_0(\frac12-2\sigma)+\pi\sec2\pi\sigma]\dd\sigma
=2\pi\mathfrak m_1-(3\log2+2\gamma)\mathfrak m_0$. This form is not used.

\paragraph{Checks.}
These are not proof steps. A $60$-digit run (symmetry, Slurm job 567960; logs
\path{LRI_C5_d2_simplify/proof/L2-fourier/check_reduction_dps60.log} and
\path{check_strip_dps60.log}, script \path{check_reduction.py}) gives: the Cauchy beta integral to
$1.4\times10^{-42}$; $\mathsf Q_{1/2}=2\mathbf K$ and \eqref{gff:mgp:Qhalf} to $6\times10^{-65}$; $a=4U_0/\pi$
and $b=4V_0/\pi$ to $4\times10^{-65}$; the derivatives from \eqref{gff:mgp:vHalf} against \eqref{gff:gc:L2taylor}
with the $70$-digit single sums to $6\times10^{-63}$; \eqref{gff:mgp:AprimeBprime} to $5\times10^{-63}$ and
$3\times10^{-64}$; and the collapse with these elliptic values of $a$, $b$, $\mathcal A'$, $\mathcal B'$
against the $70$-digit $\LL_2$ to $7\times10^{-62}$. The direct quadrature of $\mathsf Q_x$ at
$x=0.6$ is limited by its endpoint singularity ($2\times10^{-26}$). The symbolic collapse is checked
exactly by \path{collapse.py} (log \path{collapse.log}).


\subsection{\texorpdfstring{Proof of the closed form of $\mathcal K_2$}{Proof of the closed form of K2}}
\label{gff:gc:K2proof}
We prove \eqref{gff:gc:K2closed} for $\KK_2=[t^2](J_t-J_t^\Gamma)$ of \eqref{gff:gs:Kdefinition}
at $d=2$. Write $t=2\tau$. Since $h_2=1/\pi$, the path master is the kite
\eqref{gff:done:kite} at $\boldsymbol p=(1,1,1,2-2\tau,1)$:
\[
 J_{2\tau}=\frac1{\pi^2}\int_{\mathbb C^2}
 \frac{d^2x\,d^2y}{|x|\,|1-x|\,|y|\,|1-y|^{2-2\tau}\,|x-y|},\qquad0<\tau<1 .
\]
We use the parameter convention $K(m)$ of Section~\ref{gff:c1:notation}, the moments
\[
 E(a,b)=\int_0^1x^{a-1}(1-x)^{b-1}K(x)\,dx,
\]
the six moments $I_1,\ldots,I_6$ listed before
\eqref{gff:eq:KproofMomentRelations}, evaluated at $t=\tau$, and the functions
$\mathfrak F,\mathfrak G$ of \eqref{gff:eq:KproofFG}. Put $s_\tau=\sin\pi\tau$ and
$c_\tau=\cos\pi\tau$. At $d=1$ the master is a linear functional of the six moments,
\eqref{gff:eq:Jfold}; at $d=2$ it is a quadratic form in them. This proof uses none of the
periodization machinery of this Section; it is placed next to the proof for $\LL_2$ for
convenience. It was written by the Claude Code session that wrote the first proof of
Section~\ref{gff:gc:L2proof} (package \path{LRI_C5_d2_proof/}, folder \path{K2_selberg/});
the closed form \eqref{gff:gc:K2closed} itself was first found numerically by the other session named
there.

\subsubsection{Real cells and elliptic moments}
For $\xi=(\xi_1,\xi_2)\in\RR^2$ let
\[
 (L_1,\ldots,L_5)=(\xi_1,\,1-\xi_1,\,\xi_2,\,1-\xi_2,\,\xi_2-\xi_1),\qquad
 (\kappa_1,\ldots,\kappa_5)=(-\tfrac12,-\tfrac12,-\tfrac12,\tau-1,-\tfrac12),
\]
and $\Upsilon(\xi)=\prod_f|L_f(\xi)|^{\kappa_f}$. Then
$\pi^2J_{2\tau}=\int_{\mathbb C^2}\prod_f|L_f(x,y)|^{2\kappa_f}d^2x\,d^2y$, where $L_f(x,y)$
is the same affine form at complex arguments. The five lines $L_f=0$ cut $\RR^2$ into
twelve open cells $R$: for $\xi_2$ in one of $(-\infty,0)$, $(0,1)$, $(1,\infty)$, the
points $0,1,\xi_2$ cut the $\xi_1$-line into four intervals. On each cell the sign vector
$\varsigma(R)=(\operatorname{sgn}L_f)_f$ is constant, and distinct cells have distinct sign
vectors. For $0<\tau<1$,
\begin{equation}
 \Upsilon_R:=\int_R\Upsilon(\xi)\,d\xi=2I_{j(R)},
 \label{gff:gc:K2cells}
\end{equation}
with $j(R)$ given by
\begin{center}
\begin{tabular}{@{}lc@{\qquad}lc@{\qquad}lc@{}}\toprule
cell & $j$ & cell & $j$ & cell & $j$\\\midrule
$\xi_1<\xi_2<0$ & 6 & $0<\xi_1<\xi_2<1$ & 1 & $\xi_2<0,\ 1<\xi_1$ & 5\\
$\xi_2<\xi_1<0$ & 5 & $0<\xi_2<\xi_1<1$ & 2 & $0<\xi_2<1<\xi_1$ & 1\\
$\xi_1<0<\xi_2<1$ & 2 & $0<\xi_1<1<\xi_2$ & 3 & $1<\xi_1<\xi_2$ & 4\\
$\xi_1<0,\ 1<\xi_2$ & 4 & $\xi_2<0<\xi_1<1$ & 6 & $1<\xi_2<\xi_1$ & 3\\\bottomrule
\end{tabular}
\end{center}
Indeed $\Upsilon\geq0$, so by Tonelli we may integrate over $\xi_1$ first. For
$e_1<e_2<e_3$ and $P(\xi_1)=|(\xi_1-e_1)(\xi_1-e_2)(\xi_1-e_3)|^{-1/2}$, the substitutions
$\xi_1=e_3-(e_3-e_1)/\sin^2\theta$, $e_1+(e_2-e_1)\sin^2\theta$,
$e_3-(e_3-e_2)\sin^2\theta$ and $e_1+(e_3-e_1)/\sin^2\theta$ on the four intervals give
\begin{gather*}
 \int_{-\infty}^{e_1}P=\int_{e_2}^{e_3}P=2rK(m),\qquad
 \int_{e_1}^{e_2}P=\int_{e_3}^{\infty}P=2rK(1-m),\\
 m=\frac{e_3-e_2}{e_3-e_1},\qquad r=(e_3-e_1)^{-1/2}.
\end{gather*}
With $\{e_i\}=\{0,1,\xi_2\}$, the outer weight $|\xi_2|^{-1/2}|1-\xi_2|^{\tau-1}r$ becomes
the Euler weight $v^{-1/2}(1-v)^{\tau-1}$, $v^{-1/2}(1-v)^{-\tau}$ or
$v^{-\tau}(1-v)^{\tau-1}$ after $v=\xi_2$ on $(0,1)$, $v=-\xi_2/(1-\xi_2)$ on
$(-\infty,0)$ or $v=1/\xi_2$ on $(1,\infty)$, respectively. The arguments of $K$ become
$v$ or $1-v$, and the table follows. All twelve integrals are finite for $0<\tau<1$.

\subsubsection{The rotation lemma}
For each form write $|L_f(x,y)|^2=\alpha_f+\beta_f^2$, where $x=x_1+\ii x_2$,
$y=y_1+\ii y_2$, $\alpha_f=L_f(x_1,y_1)^2$, and $\beta_f$ is the linear part of $L_f$
evaluated at $(x_2,y_2)$; this holds because the $L_f$ have real coefficients. Put
$\xi=(x_1-x_2,y_1-y_2)$ and $\eta=(x_1+x_2,y_1+y_2)$. Then
$\alpha_f-\beta_f^2=L_f(\xi)L_f(\eta)$ and $dx_1\,dx_2\,dy_1\,dy_2=\frac14\,d\xi\,d\eta$.
We claim
\begin{equation}
 \pi^2J_{2\tau}=-\frac14\sum_{R,R'}e^{\ii\pi\phi(R,R')}\,\Upsilon_R\Upsilon_{R'},\qquad
 \phi(R,R')=\sum_{f:\ \varsigma_f(R)\neq\varsigma_f(R')}\kappa_f ,
 \label{gff:gc:K2rotation}
\end{equation}
the sum running over all $144$ ordered pairs of cells.

(a) \emph{Convergence.} Since $\alpha+\beta^2\geq|\alpha-\beta^2|$ and all $\kappa_f<0$,
the integrand $\prod_f(\alpha_f+\beta_f^2)^{\kappa_f}$ is at most $\Upsilon(\xi)\Upsilon(\eta)$,
whose integral is $\frac14(\int_{\RR^2}\Upsilon)^2=4(\sum_jI_j)^2<\infty$ by
\eqref{gff:gc:K2cells}. In particular $J_{2\tau}$ converges for $0<\tau<1$.

(b) \emph{An analytic family.} For $\Re z>0$ put
$\Psi(z)=z^2\int_{\RR^4}\prod_f(\alpha_f+z^2\beta_f^2)^{\kappa_f}\,dx_1\,dx_2\,dy_1\,dy_2$, with
principal powers; off a null set $\alpha_f+z^2\beta_f^2\notin(-\infty,0]$. If
$|\arg z|\leq\theta_0<\frac\pi2$ and $r_1\leq|z|\leq r_2$, the inequality
$|\alpha+Be^{\ii\vartheta}|\geq(\alpha+B)\cos\frac\vartheta2$ for $\alpha,B\geq0$ gives
$|\alpha_f+z^2\beta_f^2|\geq\cos\theta_0\,\min(1,r_1^2)\,(\alpha_f+\beta_f^2)$. The
integrand is therefore dominated by a constant multiple of the integrable function of
(a), and $\Psi$ is holomorphic on $\Re z>0$ by Morera's theorem, Fubini's theorem and
dominated convergence. For real $z>0$ the substitution $(x_2,y_2)\to(x_2,y_2)/z$ gives
$\Psi(z)=\pi^2J_{2\tau}$, so $\Psi$ is constant by the identity theorem.

(c) \emph{The boundary value.} Let $z=e^{\ii\theta}$ with $\theta\uparrow\frac\pi2$. Since
$|\alpha_f+e^{2\ii\theta}\beta_f^2|\geq|\alpha_f-\beta_f^2|$, the integrand is dominated by
$\Upsilon(\xi)\Upsilon(\eta)$ for every $\theta$. Since
$\Im(\alpha_f+e^{2\ii\theta}\beta_f^2)>0$ when $\beta_f\neq0$, almost everywhere
$(\alpha_f+e^{2\ii\theta}\beta_f^2)^{\kappa_f}\to|L_f(\xi)L_f(\eta)|^{\kappa_f}
e^{\ii\pi\kappa_f[L_f(\xi)L_f(\eta)<0]}$. Dominated convergence and $z^2\to-1$ give
$\pi^2J_{2\tau}=-\frac14\iint\Upsilon(\xi)\Upsilon(\eta)
\exp(\ii\pi\sum_f\kappa_f[\varsigma_f(\xi)\neq\varsigma_f(\eta)])\,d\xi\,d\eta$. The phase is
constant on each product of two cells, which proves \eqref{gff:gc:K2rotation}.

This is the Kawai--Lewellen--Tye (KLT) factorization \cite{gff:C6KLT} behind \eqref{gff:done:kiteFactorization}, proved
here directly at the point needed.

\subsubsection{The master path as a sum of two squares}
With $\omega_\tau=e^{\ii\pi\tau}$ the phase in \eqref{gff:gc:K2rotation} is
$(-\ii)^n(-\omega_\tau)^m$, where $n$ counts the forms with exponent $-\frac12$ whose signs
differ on $R$ and $R'$, and $m=1$ if the sign of $1-\xi_2$ differs and $m=0$ otherwise.
Collecting the $144$ terms with the table of \eqref{gff:gc:K2cells}, a finite exact
computation (machine checked in \path{LRI_C5_d2_proof/K2_selberg/region_sum_matrix.py}),
gives, with $\omega=\omega_\tau$,
\[\begin{aligned}
 \pi^2J_{2\tau}=4\bigl[&\ii I_1I_2-\ii\omega I_1I_4-\ii I_1I_5+I_1I_6-\ii\omega I_2I_3
 -\omega I_2I_4+I_2I_5\\
 &+\ii I_3I_4+\ii\omega I_3I_5-\omega I_3I_6+\ii\omega I_4I_6+\ii I_5I_6\bigr].
\end{aligned}\]
Now insert the moment relations \eqref{gff:eq:KproofMomentRelations} at $t=\tau\in(0,1)$ and
$\omega_\tau=c_\tau+\ii s_\tau$. (The second contour in the proof of
\eqref{gff:eq:KproofMomentRelations} also needs, on $x<0$, the boundary value
$K(1-x-\ii0)=(1-x)^{-1/2}[K(\frac1{1-x})-\ii K(\frac{-x}{1-x})]$. It follows from the
displayed formula for $K(x+\ii0)$, $x>1$, by Schwarz reflection, and that proof uses it
without displaying it.) All terms except $s_\tau^2(I_3^2+I_4^2)$ cancel; for instance, the coefficient of
$I_3I_4$ becomes $2\ii-2\omega_\tau(s_\tau+\ii c_\tau)=0$. Euler's integral gives
$I_3=\pi^2\mathfrak F(1-\tau)/(2s_\tau)$ and $I_4=\pi^2\mathfrak F(\tau)/(2s_\tau)$, because
$B(\tau,1-\tau)=\pi/s_\tau$. Therefore, for $0<\tau<1$,
\begin{equation}\boxed{
 J_{2\tau}=\pi^2\bigl[\mathfrak F(\tau)^2+\mathfrak F(1-\tau)^2\bigr]
 =\frac4{\pi^2}s_\tau^2\bigl(I_3^2+I_4^2\bigr)
 =\frac4{\pi^2}\bigl[s_\tau I_1I_4+I_2I_5\bigr].
}\label{gff:gc:K2master}\end{equation}
The last form is \eqref{gff:done:kiteFactorization} at this point: with $a=b=A'=c=-\frac12$
and $B'=\tau-1$, the four factors $\mathcal O(a,b,A',B',c)$, $\mathcal M(a,b,A',B',c)$,
$\mathcal O(A',B',a,b,c)$ and $\mathcal M(A',B',a,b,c)$ take the values $2I_1$, $2I_4$,
$2s_\tau I_3=2I_2$ and $2I_5$. That identity is thereby
proved here; it is not an input. Equation \eqref{gff:gc:K2master} is the two-dimensional
partner of \eqref{gff:eq:KproofTwoFunctions}, with no cross term. It is symmetric under
$t\to2-t$, as completion covariance requires. At $\tau=\frac12$ Dixon's sum
\cite{gff:DLMF16} gives $\mathfrak F(\frac12)=\Gamma(\frac14)^4/(4\pi^3)$, so
$J_1=\Gamma(\frac14)^8/(8\pi^4)=3Q_2$, in agreement with \eqref{gff:gc:Q2}.

\subsubsection{Laurent extraction}
\label{gff:gc:K2laurent}
Write $\mathfrak F(\tau)=1+u\tau+v\tau^2+O(\tau^3)$ and
$\mathfrak G(\tau)=1+\sigma\tau^3+\varrho\tau^4+O(\tau^5)$, where $\varrho$ is the
coefficient called $\tau$ in \eqref{gff:eq:KproofTaylorCancellation} and $\mathfrak G-1=O(\tau^3)$
because every term $n\geq1$ of $\mathfrak G$ carries $(\tau)_n^3$. By
\eqref{gff:eq:KproofThomae}, $\pi\mathfrak F(1-\tau)=\tau^{-1}\mathcal Q(\tau)\mathfrak G(\tau)$
for $0<\tau<1$, with
$\mathcal Q(\tau)=\pi\Gamma(1+\tau)/[\Gamma(1-\tau)\Gamma(\frac12+\tau)^2]$. The series of
$\mathfrak F$ converges locally uniformly on $\Re\tau<1$, its terms being
$O(n^{\Re\tau-2})$; $\mathfrak G$ is holomorphic on $|\tau|<\frac12$; and $\mathcal Q$ is
holomorphic and zero-free there, with $\mathcal Q(0)=1$. Hence
\[
 J_{2\tau}=\pi^2\mathfrak F(\tau)^2+\tau^{-2}\mathcal Q(\tau)^2\mathfrak G(\tau)^2
\]
is meromorphic on $|\tau|<\frac12$, with a double pole at $\tau=0$ only. By the identity
theorem it is the continuation in which \eqref{gff:gs:Kdefinition} takes its Laurent
coefficient. At $d=2$, \eqref{gff:gs:convolution} gives
\[
 J^\Gamma_{2\tau}=\mathcal V(\tau)^2+\frac{\mathcal V(\tau)-\widetilde{\mathcal V}(\tau)}\tau,\qquad
 \mathcal V=\frac{\Gamma(\tau)\Gamma(\frac12-\tau)}{\Gamma(1-\tau)\Gamma(\frac12+\tau)},\quad
 \widetilde{\mathcal V}=\frac{\Gamma(\tau)\Gamma(\frac12+\tau)\Gamma(\frac12-2\tau)}
 {\Gamma(1-\tau)\Gamma(\frac12-\tau)\Gamma(\frac12+2\tau)},
\]
which are $B_2(1,2-2\tau)$ and $B_2(1+4\tau,2-2\tau)$. The series
$\log\Gamma(1+\tau)=-\gamma_E\tau+\sum_{k\geq2}(-1)^k\zeta(k)\tau^k/k$ and
$\log\Gamma(\frac12+\tau)=\frac12\log\pi+\sum_{k\geq1}\psi_{k-1}(\frac12)\tau^k/k!$, with
$\psi_0(\frac12)=-\gamma_E-2\log2$, $\psi_1(\frac12)=\frac{\pi^2}2$,
$\psi_2(\frac12)=-14\zeta(3)$ and $\psi_3(\frac12)=\pi^4$ \cite{gff:DLMFPolygamma}, give
\begin{gather*}
 \log\mathcal Q(\tau)^2=8\log2\,\tau-\pi^2\tau^2+8\zeta(3)\tau^3-\frac{\pi^4}6\tau^4+O(\tau^5),\\
 J^\Gamma_{2\tau}=\frac{e^{8\tau\log2}}{\tau^2}-20\zeta(3)\tau-48\zeta(3)\log2\,\tau^2+O(\tau^3).
\end{gather*}
The last coefficient is four times the coefficient of $t^2$ in $J^\Gamma_t$ stated after
\eqref{gff:gs:Ksubtraction}, evaluated at $q=\frac12$. Write $\mathcal Q_k=[\tau^k]\mathcal Q(\tau)^2$;
in particular $\mathcal Q_3=\frac{256}3\log^32-8\pi^2\log2+8\zeta(3)$ and
$\mathcal Q_4=\frac{512}3\log^42-32\pi^2\log^22+64\zeta(3)\log2+\frac{\pi^4}3$. The Laurent
coefficients are
\begin{center}
\begin{tabular}{@{}cll@{}}\toprule
order & $J_{2\tau}$ & $J^\Gamma_{2\tau}$\\\midrule
$\tau^{-2}$ & $1$ & $1$\\
$\tau^{-1}$ & $8\log2$ & $8\log2$\\
$\tau^{0}$ & $\pi^2+\mathcal Q_2=32\log^22$ & $32\log^22$\\
$\tau^{1}$ & $2\pi^2u+2\sigma+\mathcal Q_3$ & $\frac{256}3\log^32-20\zeta(3)$\\
$\tau^{2}$ & $\pi^2(u^2+2v)+2\varrho+16\sigma\log2+\mathcal Q_4$ & $\frac{512}3\log^42-48\zeta(3)\log2$\\
\bottomrule
\end{tabular}
\end{center}
The orders $\tau^{-2},\tau^{-1},\tau^0$ cancel identically. At order $\tau$ the
difference is $2[\sigma+\pi^2u-4\pi^2\log2+14\zeta(3)]$, which vanishes by
\eqref{gff:eq:KproofLowerCoefficients}; this is the two-dimensional instance of
$J_t-J_t^\Gamma=O(t^2)$ in \eqref{gff:gs:Bjet}, here derived rather than assumed. At order
$\tau^2$ the coefficients $v$ and $\varrho$ enter only through
$2(\varrho+\pi^2v)=\pi^2u^2-\frac{\pi^4}3$, which is \eqref{gff:eq:KproofTaylorCancellation}.
Since $t=2\tau$, $[t^2]=\frac14[\tau^2]$, and
\[
 \KK_2=\frac{\pi^2}2u^2+4\sigma\log2-8\pi^2\log^22+28\zeta(3)\log2 .
\]
With $u=4\log2-8\Cat/\pi$ and $\sigma=8\pi\Cat-14\zeta(3)$ from
\eqref{gff:eq:KproofLowerCoefficients}, this is $\KK_2=32\Cat^2-28\zeta(3)\log2$.
Eliminating $\sigma$ with the order-$\tau$ relation instead gives
$\KK_2=\frac{\pi^2}2(u-A)^2+\frac12AU$ with $\pi(A-u)=8\Cat$, the second form of
\eqref{gff:gc:K2closed}. All coefficient algebra here is exact and machine checked
(\path{LRI_C5_d2_proof/K2_selberg/laurent_K2_sympy.py} and
\path{LRI_C5_d2_proof/K2_selberg/laurent_K2.wl}).

The one-dimensional inputs \eqref{gff:eq:KproofThomae}, \eqref{gff:eq:KproofTaylorCancellation} and
\eqref{gff:eq:KproofLowerCoefficients} concern $\mathfrak F$ and $\mathfrak G$ only and are
proved in Section~\ref{gff:app:KpsiProof}. In that proof \eqref{gff:eq:KproofThomae} is also used
at $-t$, with $\mathfrak F(1+t)$ continued by \eqref{gff:eq:KproofRecurrence}. This is
legitimate: both sides of \eqref{gff:eq:KproofThomae} are meromorphic in a punctured disk
around $t=0$ and agree for $0<t<\frac12$, so they agree there identically.

\subsubsection{Further proofs and checks}
\label{gff:gc:K2further}
Section~\ref{gff:mgk:section} gives a second position-space proof of \eqref{gff:gc:K2master}, by
Riemann's bilinear relation (Section~\ref{gff:mgk:riemann}). It also proves at $d=2$ that the
spectral form \eqref{gff:gs:Kspectral} and the series \eqref{gff:gs:Kseries} equal \eqref{gff:gc:K2closed},
and it proves the mixed-trace step for $0<\tau<\frac12$; its items (i)--(v) say where each
statement is proved. Together with Section~\ref{gff:gc:K2laurent}, these results show that at $d=2$
the definition \eqref{gff:gs:Kdefinition}, the spectral form and the series agree; the agreement of
the definition with the series does not use the mixed-trace step.

\paragraph{Checks.}
These are not proof steps. The closed form \eqref{gff:gc:K2master} agrees with a
two-dimensional quadrature of the $y$-integral, with the inner star \eqref{gff:mgk:star}
inserted, to between $38$ and $68$ digits at seven values of $\tau$. Direct
nested quadratures of the twelve cell integrals agree with $2I_j$ to between $18$ and
$24$ digits, and \eqref{gff:gs:Ksubtraction}, evaluated with an independent spectral code,
reproduces \eqref{gff:gc:K2master} to $54$ digits at $\tau=0.1$ and $\tau=0.23$. At
$\tau=\frac3{10}$ the twelve-term form above, with the moments computed by direct
quadrature, agrees with
$\pi^2[\mathfrak F(\tau)^2+\mathfrak F(1-\tau)^2]=47.790785413226103671\ldots$ to $37$
digits.

\subsection{\texorpdfstring{Spectral and series forms of $\KK_2$, the mixed-trace step, and a second position-space proof}{Spectral and series forms of K2, the mixed-trace step, and a second position-space proof}}
\label{gff:mgk:section}
Section~\ref{gff:gc:K2proof} proves \eqref{gff:gc:K2closed} for the master-path definition
\eqref{gff:gs:Kdefinition}. This subsection proves, at $d=2$, the statements that connect that
definition with the spectral form \eqref{gff:gs:Kspectral} and the series \eqref{gff:gs:Kseries}.
\begin{enumerate}
\item[(i)] For $0<\tau<\frac12$ the two regulated traces of \eqref{gff:gs:Ksubtraction} are
evaluated in closed form:
$C(\tau)\langle t_\ell^2R_\tau\rangle_2=\pi^2[\mathfrak F(\tau)^2+\mathfrak F(1-\tau)^2]$ and
$\langle t_\ell^2S_\tau\rangle_2=C(\tau)\mathcal X_2(\tau)^2$ (Section~\ref{gff:mgk:regulated}).
\item[(ii)] With \eqref{gff:gc:K2master} this proves the mixed-trace representation
$J_{2\tau}=C(\tau)\langle t_\ell^2R_\tau\rangle_2$ stated before \eqref{gff:gs:Ksubtraction}, at
$d=2$, for $0<\tau<\frac12$ and in the full-Mellin-line reading (Section~\ref{gff:mgk:regulated}).
\item[(iii)] The coefficient $[\tau^2]$ may be taken under the trace, and the right side of
\eqref{gff:gs:Kspectral} at $d=2$ equals $32\Cat^2-28\zeta(3)\log2$; as a by-product,
$\langle t_\ell^2E_2\rangle_2=-21\zeta(3)$ (Sections~\ref{gff:mgk:regularization} and \ref{gff:mgk:spectral}).
\item[(iv)] Each inner $n$ sum of \eqref{gff:gs:Kseries} at $d=2$ equals the corresponding spin term
of \eqref{gff:gs:Kspectral}, and the spin sum converges absolutely. Hence \eqref{gff:gs:Kseries} equals
\eqref{gff:gc:K2closed}, and the strip limit used for the $70$-digit evaluation of
Section~\ref{gff:gc:d2closedForms} computes exactly this ordered sum (Section~\ref{gff:mgk:series}).
\item[(v)] A second proof of \eqref{gff:gc:K2master} from the position-space integral evaluates the
inner star by Riemann's bilinear relation and the outer integral by a Stokes pairing on the
slit plane (Section~\ref{gff:mgk:riemann}).
\end{enumerate}
All five statements are proved here. Besides standard theorems (Stirling's formula, Gamma
reflection, Gauss's and Thomae's unity-argument identities, the residue theorem, Liouville's
theorem, Green's theorem, dominated convergence and the identity theorem), the proofs use only
one-dimensional lemmas of Section~\ref{gff:sec:Kdef}: the Euler integral for $E(a,b)$, the moment
relations \eqref{gff:eq:KproofMomentRelations}, the Thomae
relation \eqref{gff:eq:KproofThomae}, and the Taylor data \eqref{gff:eq:KproofLowerCoefficients} and
\eqref{gff:eq:KproofTaylorCancellation}. These concern the elliptic moments and the two functions
$\mathfrak F,\mathfrak G$ only, and do not depend on the dimension. No numerical identification
enters. The results are specific to $d=2$: at other dimensions the mixed-trace step remains
input~(M) of Section~\ref{gff:gdk:conventions}, and the interchange of $[\tau^2]$ with the trace is
not proved.

The material comes from two packages, both written by Claude Code sessions:
\begin{itemize}
\item package \path{LRI_C5_d2_simplify/} (folders \path{proof/} and \path{proof/K2-parseval/}),
by the session that wrote the third and fourth proofs of \eqref{gff:gc:L2closed}: the regularized
spectral route, Sections~\ref{gff:mgk:regularization}--\ref{gff:mgk:spectral}, apart from the parts
listed in the next item;
\item package \path{LRI_C5_d2_proof/}, by the session that wrote the first proof of
Section~\ref{gff:gc:L2proof} and the proof of Section~\ref{gff:gc:K2proof}: in folder \path{K2_gauss/},
the sharper estimate at the end of Section~\ref{gff:mgk:regularization}, the second evaluation of the
single sums (Section~\ref{gff:mgk:singles}), a second derivation of the Laurent extraction of
Section~\ref{gff:mgk:spectral} and the per-spin identity of Section~\ref{gff:mgk:series} with its arc
bounds; in folders \path{K2_d1template/} and \path{K2_d1template_fix1/}, the revised contour
bounds used in Section~\ref{gff:mgk:residues}, the third evaluation of the single sums, the
mixed-trace corollary of Section~\ref{gff:mgk:regulated} and the Riemann bilinear proof of
Section~\ref{gff:mgk:riemann}.
\end{itemize}
Both packages contain the absolute convergence of the series and the strip limit of
Section~\ref{gff:mgk:series}.

\subsubsection{Conventions at \texorpdfstring{$d=2$}{d=2}}\label{gff:mgk:conventions}
Throughout, $d=2$ and $q=\frac12$. Then $a_\ell=\frac\ell2+\frac14$, $b_\ell=\frac\ell2+\frac34$,
$d_0=1$, $d_\ell=2$ for $\ell\geq1$, and $\Gamma(2q)=1$, so $t_\ell(\nu)=g_\ell(\ii\nu)$ with the
kernel $g_\ell$ of \eqref{gff:fp:kernel}. At fixed spin we write $a=a_\ell$, $b=b_\ell$; then
$b-a=\frac12$ and $a+b=\ell+1$. The quantities of \eqref{gff:gs:RS} are
\[
 C(\tau)=\frac{\Gamma(\tau)}{\Gamma(1-\tau)}=\frac1\tau+c_*+O(\tau),\qquad c_*=2\psi_0(1),\qquad
 \mathcal X_2(\tau)=\Bigl[\frac{\Gamma(\frac12-\frac\tau2)}{\Gamma(\frac12+\frac\tau2)}\Bigr]^2,
\]
where $\mathcal X_2$ is $\mathcal X_d$ of \eqref{gff:gdk:notation} at $d=2$. The constants of
\eqref{gff:gs:psiBasis} are $A=4\log2$ and $U=\psi_2(\frac12)=-14\zeta(3)$, so
$-AU/4=14\zeta(3)\log2$. We write
\[
 w_{\KK}=(c_*-p)E_2-\frac{hk}2-\frac{5r}{32}
\]
for the weight in \eqref{gff:gs:Kspectral}, with $p,h,j,k,r,E_2$ of \eqref{gff:gs:Kweights}. Since $h,k$
are odd in $\nu$ and $p,j,r$ are even, $E_2$ and $w_{\KK}$ are even in $\nu$.

The kernel $R_\tau(s)$ is not even in $s$. As in \eqref{gff:gdk:fullTrace}, every trace in this
subsection is read on the full Mellin line,
\[
 \langle f\rangle_2=\frac1{2\pi}\sum_{\ell\geq0}d_\ell\int_{\RR}f(\ell,\nu)\dd\nu .
\]
For $f$ even in $\nu$ this is \eqref{gff:gs:trace}. For real $\tau$,
$R_\tau(-\ii\nu)=\overline{R_\tau(\ii\nu)}$, so $\langle t_\ell^2R_\tau\rangle_2$ is
$\frac1\pi\sum_\ell d_\ell\int_0^\infty\Re\{t_\ell^2R_\tau\}\dd\nu$, the real-part reading used in
\eqref{gff:gs:Ksubtraction}. We put $s_\tau=\sin\pi\tau$ and $c_\tau=\cos\pi\tau$, as in
Section~\ref{gff:gc:K2proof}, and $\xi=\frac\ell2+\ii\nu$, $\bar\xi=\frac\ell2-\ii\nu$ (not the
real points $\xi$ of Section~\ref{gff:gc:K2proof}), so that $\Re\xi=\ell/2\geq0$,
$a\pm\ii\nu=\frac14+\xi$ or $\frac14+\bar\xi$, and $b\pm\ii\nu=\frac34+\xi$ or $\frac34+\bar\xi$.
The functions $\mathfrak F,\mathfrak G$ are those of \eqref{gff:eq:KproofFG}, with the Taylor data
of Section~\ref{gff:gc:K2laurent}:
\begin{equation}\begin{gathered}
 \mathfrak F(\tau)=1+u\tau+v\tau^2+O(\tau^3),\qquad
 \mathfrak G(\tau)=1+\sigma\tau^3+\varrho\tau^4+O(\tau^5),\\
 u=4\log2-\frac{8\Cat}\pi,\qquad \sigma=8\pi\Cat-14\zeta(3),\qquad
 \varrho+\pi^2v=\frac{\pi^2u^2}2-\frac{\pi^4}6 .
\end{gathered}\label{gff:mgk:Taylor}\end{equation}
Here $\varrho$ is the coefficient called $\tau$ in \eqref{gff:eq:KproofTaylorCancellation}; the
letter $\tau$ is the regulator throughout this subsection.

\subsubsection{Analytic regularization in the rim exponent}\label{gff:mgk:regularization}

\paragraph{Lemma (uniform Gamma ratios).}
\emph{Let $\Lambda_0\subset\mathbb C$ be compact and $c_0>0$. There is $C_0$ such that for
all $\alpha,\beta\in\Lambda_0$ and all $z$ with $\Re z\geq0$ and $\Re(z+\alpha)\geq c_0$,}
\[
 \Bigl|\frac{\Gamma(z+\alpha)}{\Gamma(z+\beta)}-(1+z)^{\alpha-\beta}\Bigr|
 \leq\frac{C_0\,|(1+z)^{\alpha-\beta}|}{1+|z|},\qquad
 \Bigl|\frac{\Gamma(z+\alpha)}{\Gamma(z+\beta)}\Bigr|\leq C_0(1+|z|)^{\Re(\alpha-\beta)} .
\]
\emph{The same holds for $z$ in a closed sector $|\arg z|\leq\pi-\delta$, $|z|\geq R$, after
enlarging $C_0$ and $R$.}

\emph{Proof.} For $|z|$ large, $z+\alpha$ and $z+\beta$ lie in a sector $|\arg w|\leq\pi-\delta'$,
where $\log\Gamma(w)=(w-\frac12)\log w-w+\frac12\log2\pi+O(|w|^{-1})$ holds uniformly
\cite{gff:DLMFStirling}. Expanding gives
$\log\Gamma(z+\alpha)-\log\Gamma(z+\beta)=(\alpha-\beta)\log(1+z)+O(|z|^{-1})$ uniformly in
$\alpha,\beta\in\Lambda_0$. On the remaining compact set of $z$ both sides are continuous,
because $1/\Gamma$ is entire and $\Re(z+\alpha)\geq c_0$. \hfill$\square$

\paragraph{Lemma (majorant).}
\emph{Fix $0<r<\frac12$. There is $C_r$ such that for all $\ell\geq0$, $\nu\in\RR$ and complex
$|\tau|\leq r$,}
\begin{equation}
 \bigl|t_\ell(\nu)^2\,[R_\tau(\ii\nu)-S_\tau(\ii\nu)]\bigr|\leq C_r(1+|\xi|)^{-3+2r}.
 \label{gff:mgk:majorant}
\end{equation}
\emph{Hence $\mathcal D(\tau)=\langle t_\ell^2(R_\tau-S_\tau)\rangle_2$ converges absolutely and
uniformly on $|\tau|\leq r$, and $\mathcal D$ is holomorphic in $|\tau|<\frac12$.} The function
$\mathcal D$ is $F(\tau)$ of \eqref{gff:gdk:notation} at $d=2$.

\emph{Proof.} With $\phi(x)=\Gamma(\frac14+x)/\Gamma(\frac34+x)$ one has
$t_\ell=\phi(\xi)\phi(\bar\xi)$,
\[
 \begin{gathered}
 R_\tau(\ii\nu)=\frac{\Gamma(\frac14+\xi)}{\Gamma(\frac14+\tau+\xi)}\cdot
 \frac{\Gamma(\frac34-\tau+\bar\xi)}{\Gamma(\frac34+\bar\xi)},\qquad
 S_\tau(\ii\nu)=\vartheta_\tau(\xi)\vartheta_\tau(\bar\xi),\\
 \vartheta_\tau(x)=\Bigl[\frac{\Gamma(\frac14-\frac\tau4+x)\,\Gamma(\frac34+x)}
 {\Gamma(\frac14+x)\,\Gamma(\frac34+\frac\tau4+x)}\Bigr]^2 .
 \end{gathered}
\]
Every Gamma function in a numerator has an argument with real part at least
$\frac14-\frac r4>0$. The lemma on uniform Gamma ratios gives $|t_\ell|^2\leq C(1+|\xi|)^{-2}$ and
\[
 R_\tau(\ii\nu),\ S_\tau(\ii\nu)=(1+\xi)^{-\tau}(1+\bar\xi)^{-\tau}\Bigl(1+O\Bigl(\frac1{1+|\xi|}\Bigr)\Bigr)
\]
uniformly in $|\tau|\leq r$. Since $\arg(1+\bar\xi)=-\arg(1+\xi)$, the product
$(1+\xi)^{-\tau}(1+\bar\xi)^{-\tau}$ equals $|1+\xi|^{-2\tau}$, whose modulus is at most
$(1+|\xi|)^{2r}$. Hence
\begin{equation}
 R_\tau(\ii\nu)-S_\tau(\ii\nu)=|1+\xi|^{-2\tau}\,O\bigl((1+|\xi|)^{-1}\bigr)\qquad
 \text{uniformly in }|\tau|\leq r,
 \label{gff:mgk:RminusS}
\end{equation}
and \eqref{gff:mgk:majorant} follows. The majorant is summable:
$\int_{\RR}(1+|\frac\ell2+\ii\nu|)^{-3+2r}\dd\nu\leq C(1+\ell)^{-2+2r}$ and
$\sum_\ell(1+\ell)^{-2+2r}<\infty$. The integrands are holomorphic in $\tau$, so dominated
convergence and Morera's theorem give holomorphy. \hfill$\square$

\paragraph{Corollary (regularized form).}
\emph{(i) $\mathcal D(\tau)=\mathcal D_2\tau^2+\mathcal D_3\tau^3+O(\tau^4)$, with
$\mathcal D_2=\langle t_\ell^2E_2\rangle_2$ and
$\mathcal D_3=\langle t_\ell^2[-pE_2-\frac{hk}2-\frac{5r}{32}]\rangle_2$.
(ii) $C(\tau)\mathcal D(\tau)$ is holomorphic in $|\tau|<\frac12$, and
$[\tau^2]\{C(\tau)\mathcal D(\tau)\}=\langle t_\ell^2w_{\KK}\rangle_2$. Hence the right side of
\eqref{gff:gs:Kspectral} at $d=2$ is}
\begin{equation}
 -\frac{AU}4+\frac14\langle t_\ell^2w_{\KK}\rangle_2
 =14\zeta(3)\log2+\frac14[\tau^2]\bigl\{C(\tau)\mathcal D(\tau)\bigr\}.
 \label{gff:mgk:Kreg}
\end{equation}
\emph{(iii) For real $0<\tau<\frac12$,
$\mathcal D(\tau)=\langle t_\ell^2R_\tau\rangle_2-\langle t_\ell^2S_\tau\rangle_2$, and both
traces converge absolutely.}

\emph{Proof.} (i) By Cauchy's formula on $|\tau|=r$ and Fubini's theorem, justified by
\eqref{gff:mgk:majorant}, $[\tau^n]\mathcal D=\langle t_\ell^2\,[\tau^n](R_\tau-S_\tau)\rangle_2$.
Taylor expansion of the Gamma functions gives
\[\begin{aligned}
 \log R_\tau(s)&=\sum_{k\geq1}\frac{(-\tau)^k\psi_{k-1}(b-s)-\tau^k\psi_{k-1}(a+s)}{k!},\\
 \log S_\tau(s)&=2\sum_{k\geq1}\frac{(-\frac\tau4)^k\{\psi_{k-1}(a+s)+\psi_{k-1}(a-s)\}
 -(\frac\tau4)^k\{\psi_{k-1}(b+s)+\psi_{k-1}(b-s)\}}{k!}.
\end{aligned}\]
Expanding the exponentials to third order and symmetrizing under $s\to-s$ gives the finite
polynomial identity
\[
 \tfrac12\bigl[(R_\tau-S_\tau)(s)+(R_\tau-S_\tau)(-s)\bigr]
 =E_2\tau^2+\Bigl[-pE_2-\frac{hk}2-\frac{5r}{32}\Bigr]\tau^3+O(\tau^4),
\]
with $p,h,j,k,r$ of \eqref{gff:gs:Kweights} written in the polygamma values at $a\pm s$, $b\pm s$.
This is the expansion displayed after \eqref{gff:gs:Ksubtraction}. It is checked in exact computer
algebra, at $q=\frac12$ in \path{LRI_C5_d2_simplify/proof/K2-parseval/symbolic_checks.py} (part A)
and for arbitrary real $a,b$ in \path{LRI_C5_d2_proof/K2_gauss/check_expansion.wl}. The parts
odd in $\nu$ integrate to zero in the full-line trace. (ii) $C(\tau)=\tau^{-1}\Gamma(1+\tau)/\Gamma(1-\tau)$
has a simple pole with residue $1$, so
$[\tau^2](C\mathcal D)=\mathcal D_3+c_*\mathcal D_2=\langle t_\ell^2w_{\KK}\rangle_2$.
(iii) For real $\tau\in(0,\frac12)$ the lemma on uniform Gamma ratios gives
$|t_\ell^2R_\tau|,|t_\ell^2S_\tau|\leq C(1+|\xi|)^{-2-2\tau}$, which is summable. \hfill$\square$

Part (i) supplies, at $d=2$, the interchange of $[\tau^2]$ with the trace that the derivation of
\eqref{gff:gs:Kspectral} after \eqref{gff:gs:Ksubtraction} uses.

\paragraph{A sharper estimate for the even part.}
The bound \eqref{gff:mgk:majorant} does not use the cancellation in the even part
$D^{\rm e}_\tau(\nu)=\frac12[R_\tau(\ii\nu)+R_\tau(-\ii\nu)]-S_\tau(\ii\nu)$. Fix $\rho_0=\frac1{32}$ and
write $z=a+\ii\nu$. There is $C$ such that, for all $\ell\geq0$, $\nu\in\RR$ and $|\tau|\leq\rho_0$,
\begin{equation}
 |t_\ell(\nu)|\leq C(1+|z|)^{-1},\qquad
 |D^{\rm e}_\tau(\nu)|\leq C|\tau|^2(1+|z|)^{-2+2\rho_0}.
 \label{gff:mgk:sharp}
\end{equation}
\emph{Proof.} The first bound is $|\Gamma(w)/\Gamma(w+\frac12)|\leq C|w|^{-1/2}$ for
$\Re w\geq\frac14$ \cite{gff:DLMFStirling}, with compactness. For the second, all arguments
$w\in\{a\pm s,b\pm s\}$ have $\Re w\geq\frac14>\rho_0$, so the series
$\log\Gamma(w+\alpha\tau)-\log\Gamma(w)=\sum_{k\geq1}\psi_{k-1}(w)(\alpha\tau)^k/k!$ with
$|\alpha|\leq1$ converge. The polygamma functions obey
$|\psi_{k-1}(w)|\leq(k-1)!\sum_{n\geq0}|w+n|^{-k}\leq(k-1)!\,4^kC_0(1+|w|)^{-2}$ for $k\geq3$, and
$\psi_1(w)-\psi_1(w+\frac12)=O(|w|^{-2})$, $\psi_{k-1}(w)=O(|w|^{1-k})$,
$\psi_0(w)=\log w+O(|w|^{-1})$. Put $A_k=\psi_{k-1}(a+s)+\psi_{k-1}(a-s)$,
$B_k=\psi_{k-1}(b+s)+\psi_{k-1}(b-s)$,
$\Delta_{\rm e}=\frac12[\log R_\tau(s)+\log R_\tau(-s)]-\log S_\tau(s)$ and
$\varepsilon_{\rm o}=\frac12[\log R_\tau(s)-\log R_\tau(-s)]$. The two expansions above give
\[
 [\tau^k]\Delta_{\rm e}=\frac1{k!}\Bigl[-\bigl(\tfrac12+2(-4)^{-k}\bigr)A_k
 +\bigl(\tfrac{(-1)^k}2+2\cdot4^{-k}\bigr)B_k\Bigr].
\]
The coefficient vanishes for $k=1$; it is $-\frac5{16}(A_2-B_2)=O(|z|^{-2})$ for $k=2$ and
$-\frac{15}{192}(A_3+B_3)=O(|z|^{-2})$ for $k=3$; and it is $O((8^k/k)|z|^{-2})$ for $k\geq4$.
Hence $|\Delta_{\rm e}|\leq C|\tau|^2|z|^{-2}$. Similarly $|\varepsilon_{\rm o}|\leq C|\tau||z|^{-1}$;
its $\tau^1$ coefficient is $\frac12[(\psi_0(b+s)-\psi_0(a+s))-(\psi_0(b-s)-\psi_0(a-s))]$.
Finally $|\log S_\tau|\leq2|\tau|\log(2+|z|)+C|\tau|$. Since
$D^{\rm e}_\tau=S_\tau(e^{\Delta_{\rm e}}\cosh\varepsilon_{\rm o}-1)$, it is bounded by
$C(2+|z|)^{2\rho_0}(|\Delta_{\rm e}|+|\varepsilon_{\rm o}|^2)$, which proves \eqref{gff:mgk:sharp} for
$|z|\geq z_0$. For $|z|\leq z_0$ (finitely many spins and bounded $\nu$) compactness and
$D^{\rm e}_\tau=O(\tau^2)$ suffice. \hfill$\square$

Consequently $C(\tau)D^{\rm e}_\tau$, whose pole at $\tau=0$ is cancelled, is holomorphic on
$|\tau|\leq\rho_0$, and by the maximum principle $|t_\ell^2C(\tau)D^{\rm e}_\tau|\leq C(1+|z|)^{-4+2\rho_0}$
there. Since $[\tau^2]\{C(\tau)D^{\rm e}_\tau\}=w_{\KK}$ (by the expansion in the proof of part (i) of
the corollary and $C(\tau)=\tau^{-1}+c_*+O(\tau)$), Cauchy's estimate
gives $|t_\ell^2w_{\KK}|\leq C(1+|z|)^{-4+2\rho_0}$, so the spin term of \eqref{gff:gs:Kspectral} at $d=2$ is
$O(\ell^{-3+2\rho_0})$, consistent with the decay $O(\ell^{-3}\log\ell)$ stated after
\eqref{gff:gs:Kspectral}.

\subsubsection{Residue families on the square lattice}\label{gff:mgk:residues}
For $0<\tau<\frac12$ define the lattice
functions (the sans-serif letters are local to this subsection)
\begin{equation}\begin{aligned}
 \mathsf h(x)&=\frac{\Gamma(x+\frac12)^3}{\Gamma(x+1)^2\,\Gamma(x+\frac12+\tau)},&
 \mathsf k(x)&=\frac{\Gamma(x+1-\tau)^3}{\Gamma(x+\frac32-\tau)^2\,\Gamma(x+1)},\\
 \mathsf m(x)&=\Bigl[\frac{\Gamma(x+\frac{1-\tau}2)}{\Gamma(x+1)}\Bigr]^2,&
 \mathsf n(x)&=\frac{\Gamma(x+\frac12)^2\,\Gamma(x+1-\tau)}{\Gamma(x+1)^3}.
\end{aligned}\label{gff:mgk:lattice}\end{equation}
For $x\geq0$, Stirling's formula gives $\mathsf h,\mathsf k,\mathsf m,\mathsf n=O((1+x)^{-1-\tau})$.
Since $\mathsf h'/\mathsf h$ and $\mathsf m'/\mathsf m$ are combinations of digamma differences,
which are $O(1/(1+x))$, also $\mathsf h',\mathsf m'=O((1+x)^{-2-\tau})$. Note the shift relation
$\mathsf h(x+\frac12-\tau)=\mathsf k(x)$. The single sums
\[
 \begin{gathered}
 \Sigma_{\mathsf h}=\sum_{m\geq0}\mathsf h(m),\quad \Sigma'_{\mathsf h}=\sum_{m\geq0}\mathsf h'(m),\quad
 \Sigma_{\mathsf k}=\sum_{m\geq0}\mathsf k(m),\\
 \Sigma_{\mathsf m}=\sum_{m\geq0}\mathsf m(m),\quad
 \Sigma'_{\mathsf m}=\sum_{m\geq0}\mathsf m'(m),\quad \Sigma_{\mathsf n}=\sum_{k\geq0}\mathsf n(k)
 \end{gathered}
\]
therefore converge absolutely for $\tau>0$. At fixed spin the two integrands are
\[
 \Theta^{\rm R}_\ell(s)=g_\ell(s)^2R_\tau(s)
 =\frac{\Gamma(a+s)^3\,\Gamma(a-s)^2\,\Gamma(b-\tau-s)}{\Gamma(b+s)^2\,\Gamma(b-s)^3\,\Gamma(a+\tau+s)},\qquad
 \Theta^{\rm S}_\ell(s)=g_\ell(s)^2S_\tau(s)=g_\ell(s;1-\tau)^2,
\]
where $g_\ell(s;1-\tau)$ is the balanced kernel of rim exponent $1-\tau$ of \eqref{gff:gdk:Skernel},
that is $\Gamma(a'+s)\Gamma(a'-s)/[\Gamma(b'+s)\Gamma(b'-s)]$ with $a'=a-\frac\tau4$ and
$b'=b+\frac\tau4$. Thus
$\langle t_\ell^2R_\tau\rangle_2=\sum_\ell d_\ell\int_{-\ii\infty}^{\ii\infty}\frac{ds}{2\pi\ii}\Theta^{\rm R}_\ell(s)$,
and similarly for $S_\tau$.

\paragraph{Lemma (per-spin residues).}
\emph{Let $\ell\geq0$, $0<\tau<\frac12$ and $m=\ell+n$. Then, with absolutely convergent series,}
\begin{equation}
 \int_{-\ii\infty}^{\ii\infty}\frac{ds}{2\pi\ii}\,\Theta^{\rm R}_\ell(s)
 =\sum_{n\geq0}\bigl[\rho_{\mathsf h}(m,n)+\rho_{\mathsf k}(m,n)\bigr],\qquad
 \int_{-\ii\infty}^{\ii\infty}\frac{ds}{2\pi\ii}\,\Theta^{\rm S}_\ell(s)=\sum_{n\geq0}\rho_{\mathsf m}(m,n),
 \label{gff:mgk:perSpin}
\end{equation}
\[\begin{aligned}
 \rho_{\mathsf h}(m,n)&=-\frac1{\pi^2c_\tau}\Bigl[\mathsf h'(m)\mathsf h(n)+\mathsf h(m)\mathsf h'(n)
 +\pi\tan(\pi\tau)\,\mathsf h(m)\mathsf h(n)\Bigr],\qquad
 \rho_{\mathsf k}(m,n)=\frac{s_\tau^3}{\pi c_\tau^2}\,\mathsf k(m)\mathsf k(n),\\
 \rho_{\mathsf m}(m,n)&=-\frac1{\pi^2}\Bigl[\cos^2(\tfrac{\pi\tau}2)\bigl(\mathsf m'(m)\mathsf m(n)
 +\mathsf m(m)\mathsf m'(n)\bigr)+\pi s_\tau\,\mathsf m(m)\mathsf m(n)\Bigr].
\end{aligned}\]

\emph{Proof.} \emph{Poles.} In $\Re s>0$ the poles of $\Theta^{\rm R}_\ell$ are the double poles
$s=a+n$, from $\Gamma(a-s)^2$, and the simple poles $s=b-\tau+n$, from $\Gamma(b-\tau-s)$,
$n\geq0$; the other factors are finite and nonzero there because $\tau\notin\frac12+\mathbb Z$ and
$\tau\notin\mathbb Z$. The two families are disjoint because $b-\tau-a=\frac12-\tau\notin\mathbb Z$,
and none lies on $\Re s=0$, since $a>0$ and $b-\tau\geq\frac34-\tau>0$. The factors with left
poles are holomorphic on $\Re s\geq0$.

\emph{Arcs.} Gamma reflection of the three factors with right poles, together with
$\sin\pi(b-s)=\cos\pi(a-s)$, gives the identity $\Theta^{\rm R}_\ell=G_\ell\,T_\tau$, with
\[
 \begin{gathered}
 G_\ell(s)=\Bigl[\frac{\Gamma(a+s)\Gamma(1-b+s)}{\Gamma(b+s)\Gamma(1-a+s)}\Bigr]^2
 \frac{\Gamma(a+s)\Gamma(1-b+s)}{\Gamma(a+\tau+s)\Gamma(1-b+\tau+s)},\\
 T_\tau(s)=\cot^2\pi(a-s)\,\frac{\cos\pi(a-s)}{\sin\pi(b-\tau-s)} .
 \end{gathered}
\]
For fixed $\ell$, $G_\ell$ is holomorphic on $\{\Re s\geq0,\ |s|\geq b\}$; there every Gamma
argument stays in a sector $|\arg|\leq\frac34\pi$, and Stirling's formula gives
$|G_\ell(s)|\leq C_\ell|s|^{-2-2\tau}$. Inside $|s|<b$ the factors have poles and zeros that cancel
pairwise in $\Theta^{\rm R}_\ell$: the triple poles of $\Gamma(1-b+s)^3$ at $s=b-1-k$ against the
triple zeros of $\cot^2\pi(a-s)\cos\pi(a-s)$; the double poles of $\cot^2\pi(a-s)$ at $s=a-k$,
$k\geq1$, against the double zeros of $1/\Gamma(1-a+s)^2$; and the poles of
$1/\sin\pi(b-\tau-s)$ at $s=b-\tau-k$, $k\geq1$, against the zeros of $1/\Gamma(1-b+\tau+s)$. The
residue theorem is therefore applied to $\Theta^{\rm R}_\ell$ itself. The factor $T_\tau$ is
$1$-periodic and bounded on $|\Im s|\geq1$, and it is bounded on the circles
$|s|=\rho_N:=a+N+\frac34-\frac\tau2$, which stay at distance at least $\frac14-O(\rho_N^{-1})$
from both pole families. So $|\Theta^{\rm R}_\ell|\leq C\rho_N^{-2-2\tau}$ on the right
half-circles, and the arcs, together with the parts of the Mellin line beyond $\rho_N$
(where $|\Theta^{\rm R}_\ell(\ii\nu)|\leq C(1+|\nu|)^{-2-2\tau}$), contribute
$O(\rho_N^{-1-2\tau})$. The right half-discs are traversed clockwise, so
\[
 \int_{-\ii\infty}^{\ii\infty}\frac{ds}{2\pi\ii}\,\Theta^{\rm R}_\ell
 =-\sum_{n\geq0}\Bigl[\Res_{s=a+n}+\Res_{s=b-\tau+n}\Bigr]\Theta^{\rm R}_\ell ,
\]
the residue series converging absolutely by the bounds below. (The number of poles of each family
inside a given arc is $N$ or $N+1$; this does not matter, since the series converges absolutely.)

\emph{Residues.} At $s=a+n+u$ one has $a+s=m+\frac12+u$, $b+s=m+1+u$ and
$a+\tau+s=m+\frac12+\tau+u$, while reflection gives
\[
 \frac{\Gamma(a-s)^2\,\Gamma(b-\tau-s)}{\Gamma(b-s)^3}
 =\frac{\cos^3\pi u}{\sin^2\pi u\,\cos\pi(\tau+u)}\,\mathsf h(n+u),
\]
the signs $(-1)^{2n}$, $(-1)^{n}$ and $(-1)^{3n}$ of the three reflected factors cancelling. Hence
$\Theta^{\rm R}_\ell=(\pi u)^{-2}\Psi_\tau(u)\,\mathsf h(m+u)\mathsf h(n+u)$ with
$\Psi_\tau(u)=(\pi u)^2\cos^3\pi u/[\sin^2\pi u\cos\pi(\tau+u)]=c_\tau^{-1}(1+\pi u\tan\pi\tau+O(u^2))$,
and the residue is
$\pi^{-2}[c_\tau^{-1}(\mathsf h'(m)\mathsf h(n)+\mathsf h(m)\mathsf h'(n))+\pi s_\tau c_\tau^{-2}\mathsf h(m)\mathsf h(n)]
=-\rho_{\mathsf h}(m,n)$. At $s=b-\tau+n+u$ only $\Gamma(b-\tau-s)=\Gamma(-n-u)$ is singular, with
residue $-(-1)^n/n!$ in $s$. At $u=0$ the remaining factors are
$\Gamma(a+s)^3/[\Gamma(b+s)^2\Gamma(a+\tau+s)]=\mathsf k(m)$ and
\[
 \frac{\Gamma(a-s)^2}{\Gamma(b-s)^3}=\frac{\Gamma(\tau-\frac12-n)^2}{\Gamma(\tau-n)^3}
 =(-1)^n\frac{s_\tau^3}{\pi c_\tau^2}\,\frac{\Gamma(n+1-\tau)^3}{\Gamma(n+\frac32-\tau)^2},
\]
so the residue is $-(s_\tau^3/\pi c_\tau^2)\mathsf k(m)\mathsf k(n)=-\rho_{\mathsf k}(m,n)$.

\emph{The $S$ family.} The function $\Theta^{\rm S}_\ell=g_\ell(s;1-\tau)^2$ has in $\Re s>0$ only
the double poles $s=a'+n$. Put $x_\tau=\frac{1-\tau}2$; then $2a'=\ell+x_\tau$, $a'+b'=\ell+1$ and
$b'-a'=1-x_\tau$. At $s=a'+n+u$, reflection gives
\[
 \begin{gathered}
 g_\ell(s;1-\tau)=-\frac{\sin\pi(x_\tau+u)}{\sin\pi u}\,
 \frac{\Gamma(m+x_\tau+u)\,\Gamma(n+x_\tau+u)}{\Gamma(m+1+u)\,\Gamma(n+1+u)},\\
 \Theta^{\rm S}_\ell=\frac{\sin^2\pi(x_\tau+u)}{\sin^2\pi u}\,\mathsf m(m+u)\mathsf m(n+u).
 \end{gathered}
\]
Since $\sin^2\pi(x_\tau+u)/\sin^2\pi u=(\pi u)^{-2}[\sin^2\pi x_\tau+\pi u\sin2\pi x_\tau+O(u^2)]$,
$\sin\pi x_\tau=\cos\frac{\pi\tau}2$ and $\sin2\pi x_\tau=s_\tau$, the residue is
$\pi^{-2}[\pi s_\tau\mathsf m(m)\mathsf m(n)+\cos^2(\frac{\pi\tau}2)(\mathsf m'(m)\mathsf m(n)+\mathsf m(m)\mathsf m'(n))]
=-\rho_{\mathsf m}(m,n)$. On the circles $|s|=a'+N+\frac12$, at distance $\frac12$ from the poles,
reflection and Stirling's formula give $|g_\ell(s;1-\tau)|\leq C|s|^{-1-\tau}$, so the arcs vanish
and the second identity of \eqref{gff:mgk:perSpin} follows as the first. Absolute convergence follows
from $|\rho_{\mathsf h}|,|\rho_{\mathsf k}|,|\rho_{\mathsf m}|\leq C(1+m)^{-1-\tau}(1+n)^{-1-\tau}$.
\hfill$\square$

\paragraph{Proposition (factorized traces).}
\emph{For $0<\tau<\frac12$,}
\begin{equation}\begin{aligned}
 \langle t_\ell^2R_\tau\rangle_2&=-\frac1{\pi^2c_\tau}\Bigl[2\Sigma_{\mathsf h}\Sigma'_{\mathsf h}
 +\pi\tan(\pi\tau)\,\Sigma_{\mathsf h}^2\Bigr]+\frac{s_\tau^3}{\pi c_\tau^2}\,\Sigma_{\mathsf k}^2,\\
 \langle t_\ell^2S_\tau\rangle_2&=-\frac1{\pi^2}\Bigl[2\cos^2(\tfrac{\pi\tau}2)\,\Sigma_{\mathsf m}\Sigma'_{\mathsf m}
 +\pi s_\tau\,\Sigma_{\mathsf m}^2\Bigr].
\end{aligned}\label{gff:mgk:factorized}\end{equation}
\emph{Proof.} If $\rho$ is symmetric on $\mathbb Z_{\geq0}^2$ and absolutely summable, then
$\sum_{\ell\geq0}d_\ell\sum_{n\geq0}\rho(\ell+n,n)=\sum_{m,n\geq0}\rho(m,n)$: the spin $\ell$ labels the
diagonal $m-n=\ell$, $d_0=1$ counts the main diagonal, and $d_\ell=2$ counts the two diagonals
$m-n=\pm\ell$, which carry equal sums. Apply this to $\rho_{\mathsf h}+\rho_{\mathsf k}$ and to
$\rho_{\mathsf m}$, which are symmetric and absolutely summable for $\tau>0$, and factor the
absolutely convergent double sums. \hfill$\square$

This is the spectral form, at $\tau>0$, of the square-lattice factorization \eqref{gff:fp:wheelSquare}.
With $\hat{\mathsf h}(z)=\sum_m\mathsf h(m)z^m$ and similarly for the other sequences, the Fourier
series $\sum_{m,n}(\rho_{\mathsf h}+\rho_{\mathsf k})(m,n)e^{\ii(m-n)\theta}$ is the same bilinear
combination of $\hat{\mathsf h},\hat{\mathsf h}',\hat{\mathsf k}$ and their conjugates at
$z=e^{\ii\theta}$. For $\tau>0$ all coefficient sequences are absolutely summable, these functions
are continuous on the closed disk, and the strip limit $\theta\to0$ is evaluation at $z=1$. At
$\tau=0$ the same limit needs the logarithmic singular parts of the single sums
(Section~\ref{gff:mgk:series}).

\subsubsection{The single sums in closed form}\label{gff:mgk:singles}

\paragraph{Lemma (two-mode periodization).}
\emph{For $0<\tau<\frac12$,}
\begin{equation}
 \Sigma'_{\mathsf h}=-\pi s_\tau\,\Sigma_{\mathsf n},\qquad
 \Sigma_{\mathsf h}=c_\tau\,\Sigma_{\mathsf n}+s_\tau^2\,\Sigma_{\mathsf k}.
 \label{gff:mgk:periodTwo}
\end{equation}
\emph{Proof.} \emph{Bilateral sum.} For $k\geq0$, reflection gives identically in $\eta$
\[
 \mathsf h(-k-1+\eta)=T_{\mathsf h}(\eta)\,\mathsf n(k+\tfrac12-\eta),\qquad
 T_{\mathsf h}(\eta)=\frac{\sin^2(\pi\eta)\,\cos\pi(\eta+\tau)}{\cos^3(\pi\eta)},
\]
since the three factors $\Gamma(-k-\frac12+\eta)^3$, $\Gamma(-k+\eta)^{-2}$ and
$\Gamma(-k-\frac12+\tau+\eta)^{-1}$ produce the signs $(-1)^{3k+3}$, $(-1)^{2k}$ and $(-1)^{k+1}$,
which cancel. The series $\sum_{m\geq0}\mathsf h(m+\eta)$ and $\sum_{k\geq0}\mathsf n(k+\frac12-\eta)$
converge locally uniformly and are holomorphic for $\Re\eta>-\frac12$ and $\Re\eta<1$
respectively, because $\mathsf h,\mathsf n=O(|x|^{-1-\tau})$ and the poles of $\mathsf n$ lie at
$-\frac12-\mathbb Z_{\geq0}$ and $-1+\tau-\mathbb Z_{\geq0}$. Hence the bilateral sum
$\Sigma^{\mathbb Z}_{\mathsf h}(\eta)=\sum_{m\in\mathbb Z}\mathsf h(m+\eta)$ satisfies, on the strip
$-\frac12<\Re\eta<1$,
\[
 \Omega_{\mathsf h}(\eta):=\cos^3(\pi\eta)\,\Sigma^{\mathbb Z}_{\mathsf h}(\eta)
 =\cos^3(\pi\eta)\sum_{m\geq0}\mathsf h(m+\eta)
 +\sin^2(\pi\eta)\cos\pi(\eta+\tau)\sum_{k\geq0}\mathsf n(k+\tfrac12-\eta),
\]
and the right side is holomorphic there. \emph{Entire and antiperiodic.} $\Sigma^{\mathbb Z}_{\mathsf h}$
is $1$-periodic and $\cos^3(\pi\eta)$ changes sign under $\eta\mapsto\eta+1$. Since the strip has
width $\frac32>1$, $\Omega_{\mathsf h}$ extends to an entire function with
$\Omega_{\mathsf h}(\eta+1)=-\Omega_{\mathsf h}(\eta)$. \emph{Growth.} For $0\leq\Re\eta\leq1$ and
$|\Im\eta|=|y|\geq1$, the lemma on uniform Gamma ratios in the sector $|\arg|\leq\pi-\delta$ for
$|y|\geq R$, and compactness for $1\leq|y|\leq R$ (there are no poles off the real axis), give
$|\sum_{m\geq0}\mathsf h(m+\eta)|,|\sum_{k\geq0}\mathsf n(k+\frac12-\eta)|\leq C\sum_{k\geq0}|k+\ii y|^{-1-\tau}
\leq C'|y|^{-\tau}$; for the term with real part in $[-\frac12,0)$ the argument stays in such a
sector because $|y|\geq1$. Hence $|\Omega_{\mathsf h}(\eta)|\leq Ce^{3\pi|y|}|y|^{-\tau}$.
\emph{Two modes.} An entire function of period $2$ that changes sign under $\eta\mapsto\eta+1$ has
an expansion $\sum_{k\ \rm odd}\hat\omega_ke^{\ii\pi k\eta}$, with
$\hat\omega_k=\frac12\int_0^2\Omega_{\mathsf h}(x+\ii y)e^{-\ii\pi k(x+\ii y)}\dd x$ for every $y$.
Thus $|\hat\omega_k|\leq Ce^{3\pi|y|}|y|^{-\tau}e^{\pi ky}$. Letting $y\to-\infty$ for $k\geq3$ and
$y\to+\infty$ for $k\leq-3$ gives $\hat\omega_k=0$ for $|k|\geq3$; for $|k|=3$ this uses the factor
$|y|^{-\tau}$, that is $\tau>0$. Therefore
$\Omega_{\mathsf h}(\eta)=\alpha_0\cos\pi\eta+\beta_0\sin\pi\eta$. \emph{Three evaluations.} At
$\eta=0$ (value and derivative; the second term of $\Omega_{\mathsf h}$ has a double zero there,
and termwise differentiation is legitimate), $\alpha_0=\Sigma_{\mathsf h}$ and
$\pi\beta_0=\Sigma'_{\mathsf h}$. At $\eta=\frac12$, $\beta_0=\Omega_{\mathsf h}(\frac12)
=\cos(\frac\pi2+\pi\tau)\Sigma_{\mathsf n}=-s_\tau\Sigma_{\mathsf n}$. At $\eta=\frac12-\tau$ the
second term vanishes, $\cos^3(\frac\pi2-\pi\tau)=s_\tau^3$, and
$\sum_m\mathsf h(m+\frac12-\tau)=\Sigma_{\mathsf k}$ by the shift relation, so
$\alpha_0s_\tau+\beta_0c_\tau=s_\tau^3\Sigma_{\mathsf k}$. These three relations give
\eqref{gff:mgk:periodTwo}. \hfill$\square$

\paragraph{Lemma (one-mode periodization).}
\emph{For $0<\tau<\frac12$ and $x_\tau=\frac{1-\tau}2$,}
\begin{equation}
 \Sigma_{\mathsf m}=\Gamma(\tau)\,\mathcal X_2(\tau),\qquad
 \Sigma'_{\mathsf m}=-2\pi\tan(\tfrac{\pi\tau}2)\,\Sigma_{\mathsf m},\qquad
 \sum_{m\in\mathbb Z}\mathsf m(m+\eta)=\Sigma_{\mathsf m}\,\frac{\sin^2\pi x_\tau}{\sin^2\pi(x_\tau+\eta)} .
 \label{gff:mgk:periodOne}
\end{equation}
\emph{Proof.} Gauss's summation \cite{gff:DLMFGaussUnity} gives
$\Sigma_{\mathsf m}=\Gamma(x_\tau)^2{}_2F_1(x_\tau,x_\tau;1;1)=\Gamma(x_\tau)^2\Gamma(1-2x_\tau)/\Gamma(1-x_\tau)^2
=\Gamma(\tau)\mathcal X_2(\tau)$, since $1-2x_\tau=\tau>0$. Reflection gives
$\mathsf m(-k-1+\eta)=\frac{\sin^2\pi\eta}{\sin^2\pi(x_\tau+\eta)}\tilde{\mathsf m}(k-\eta)$ with
$\tilde{\mathsf m}(x)=[\Gamma(x+1)/\Gamma(x+2-x_\tau)]^2$. Hence
\[
 \sin^2\pi(x_\tau+\eta)\sum_{m\in\mathbb Z}\mathsf m(m+\eta)
 =\sin^2\pi(x_\tau+\eta)\sum_{m\geq0}\mathsf m(m+\eta)+\sin^2(\pi\eta)\sum_{k\geq0}\tilde{\mathsf m}(k-\eta)
\]
is holomorphic on $-x_\tau<\Re\eta<1$, a strip of width larger than $1$; it is $1$-periodic, hence
entire, and it is $O(e^{2\pi|y|}|y|^{-\tau})$. Its Fourier modes $e^{2\pi\ii k\eta}$ with $k\neq0$
vanish by the argument of the previous proof (for $|k|=1$ through the factor $|y|^{-\tau}$), so it
is a constant $c_{\mathsf m}$. The reflected part has a double zero at $\eta=0$, so
$c_{\mathsf m}=\sin^2(\pi x_\tau)\Sigma_{\mathsf m}$, which is the last identity, and
$\Sigma'_{\mathsf m}/\Sigma_{\mathsf m}=\frac{d}{d\eta}\log\sin^{-2}\pi(x_\tau+\eta)\big|_0=-2\pi\cot\pi x_\tau
=-2\pi\tan\frac{\pi\tau}2$. \hfill$\square$

The last identity in \eqref{gff:mgk:periodOne} is an instance of Dougall's bilateral ${}_2H_2$
summation; the proof above does not assume it.

\paragraph{Lemma (Euler and Thomae).}
\emph{For $0<\tau<\frac12$,}
\begin{equation}
 \Sigma_{\mathsf n}=\pi\Gamma(1-\tau)\,\mathfrak F(1-\tau)=\frac{\pi\Gamma(\tau)}{\Gamma(\frac12+\tau)^2}\,\mathfrak G(\tau),\qquad
 \Sigma_{\mathsf k}=\frac{\Gamma(1-\tau)^3}{\Gamma(\frac32-\tau)^2}\,\mathfrak G(1-\tau)=\Gamma(\tau)\Gamma(1-\tau)^2\,\mathfrak F(\tau).
 \label{gff:mgk:singleValues}
\end{equation}
\emph{Proof.} The first form of each is the definition of the series:
$\Sigma_{\mathsf n}=\sum_k\Gamma(k+\frac12)^2\Gamma(k+1-\tau)/k!^3$ and
$\Sigma_{\mathsf k}=\sum_k\Gamma(k+1-\tau)^3/[\Gamma(k+\frac32-\tau)^2k!]$. The second forms follow from
\eqref{gff:eq:KproofThomae} at $t=\tau$ and at $t=1-\tau$; the latter reads
$\mathfrak G(1-\tau)=\Gamma(\frac32-\tau)^2\Gamma(\tau)\Gamma(1-\tau)^{-1}\mathfrak F(\tau)$, and its
${}_3F_2(1)$ has parametric excess $\tau>0$. \hfill$\square$

In particular $s_\tau\Sigma_{\mathsf k}=\pi\Gamma(1-\tau)\mathfrak F(\tau)$, since
$\Gamma(\tau)\Gamma(1-\tau)=\pi/s_\tau$.

\paragraph{Second evaluation: a bilateral Laurent lemma.}
The derivative sums can also be obtained without the growth argument for $\Omega_{\mathsf h}$. Let
$f\in\{\mathsf h,\mathsf m\}$, with $\beta_f=\frac12$ and pole order $p_f=3$ for $\mathsf h$, and
$\beta_f=x_\tau$, $p_f=2$ for $\mathsf m$.
\begin{enumerate}
\item[(a)] For each $\eta_0>0$ there is $C_{\eta_0}$ with $|f(K)|\leq C_{\eta_0}(1+|K|)^{-1-\tau}$
whenever $\operatorname{dist}(K,-\beta_f-\mathbb Z_{\geq0})\geq\eta_0$. For $\Re K\geq0$ this is
Stirling's formula. For $\Re K<0$ use the reflection identities
\[
 \begin{gathered}
 \mathsf h(K)=\frac{\sin^2(\pi K)\,\sin\pi(K+\frac12+\tau)}{\sin^3\pi(K+\frac12)}\,
 \frac{\Gamma(-K)^2\,\Gamma(\frac12-\tau-K)}{\Gamma(\frac12-K)^3},\\
 \mathsf m(K)=\frac{\sin^2(\pi K)}{\sin^2\pi(K+x_\tau)}\,\frac{\Gamma(-K)^2}{\Gamma(1-x_\tau-K)^2}:
 \end{gathered}
\]
the trigonometric factors are $1$-periodic, bounded away from the poles, and have finite limits as
$\Im K\to\pm\infty$, and the Gamma ratios are $O(|K|^{-1-\tau})$ by Stirling's formula in
$\Re(-K)>0$.
\item[(b)] $P_f(u)=\sum_{k\in\mathbb Z}f(k+u)$ is meromorphic and $1$-periodic, with poles only in
$-\beta_f+\mathbb Z$, of order at most $p_f$; its Laurent coefficients $a_{-j}$ at $u=-\beta_f$ are
the sums of the termwise Laurent coefficients (by (a), the convergence is absolute and locally
uniform off the poles and uniform on small circles around $-\beta_f$). Moreover
$\sup_{0\leq\Re u\leq1}|P_f(u)|\leq C\sum_k|k+u|^{-1-\tau}\leq C'|\Im u|^{-\tau}\to0$ as
$|\Im u|\to\infty$.
\item[(c)] $P_f(u)=a_{-2}\,\pi^2/\sin^2\pi(u+\beta_f)+a_{-3}\,\pi^3\cos\pi(u+\beta_f)/\sin^3\pi(u+\beta_f)$,
with $a_{-3}=0$ for $\mathsf m$. Indeed, since $\pi^2/\sin^2\pi\epsilon=\epsilon^{-2}+O(1)$,
$\pi^3\cos\pi\epsilon/\sin^3\pi\epsilon=\epsilon^{-3}+O(\epsilon)$ and
$\pi\cot\pi\epsilon=\epsilon^{-1}+O(\epsilon)$, the function
$P_f-a_{-2}\pi^2\sin^{-2}\pi(u+\beta_f)-a_{-3}\pi^3\cos\pi(u+\beta_f)\sin^{-3}\pi(u+\beta_f)-a_{-1}\pi\cot\pi(u+\beta_f)$
is entire, $1$-periodic and bounded, hence constant by Liouville's theorem. Its limits at
$\pm\ii\infty$ are $\pm\ii\pi a_{-1}$, so $a_{-1}=0$ and the constant vanishes.
\item[(d)] $\sum_{k\geq0}f'(k)=P_f'(0)$: for $k\geq1$, $f(u-k)$ has a double zero at $u=0$, from
$1/\Gamma(K+1)^2$, while the other factors are finite at $K=-k$ because $\frac12-k$,
$\frac12+\tau-k$ and $x_\tau-k$ are not poles; termwise differentiation of the two halves of $P_f$
is justified by local uniform convergence.
\end{enumerate}
For $\mathsf m$, only $\mathsf m(-j-x_\tau+\epsilon)\simeq\epsilon^{-2}/[j!^2\Gamma(1-x_\tau-j)^2]$
contributes at $u=-x_\tau$, and $1/\Gamma(1-x_\tau-j)=(-1)^j\sin(\pi x_\tau)\Gamma(x_\tau+j)/\pi$, so
$a_{-2}=\pi^{-2}\sin^2(\pi x_\tau)\Sigma_{\mathsf m}$. Then (c) is the last identity of
\eqref{gff:mgk:periodOne}, and (d) gives $\Sigma'_{\mathsf m}=-2\pi\cot(\pi x_\tau)\Sigma_{\mathsf m}$. For
$\mathsf h$, only the terms $\mathsf h(-\frac12-j+\epsilon)$ have a pole at $u=-\frac12$. With
$\Gamma(-j+\epsilon)^3=(-1)^j\epsilon^{-3}/j!^3+\dots$, $1/\Gamma(\frac12-j)^2=\Gamma(\frac12+j)^2/\pi^2$ and
$1/\Gamma(\tau-j)=(-1)^js_\tau\Gamma(1+j-\tau)/\pi$,
\[
 a_{-3}=\frac{s_\tau}{\pi^3}\sum_{j\geq0}\frac{\Gamma(\frac12+j)^2\,\Gamma(1+j-\tau)}{j!^3}
 =\frac{s_\tau}{\pi^3}\,\Sigma_{\mathsf n},
\]
which converges since $\tau>0$. With $\beta_f=\frac12$, (c) reads
$P_{\mathsf h}(u)=a_{-2}\pi^2\sec^2\pi u-a_{-3}\pi^3\tan(\pi u)\sec^2(\pi u)$, so (d) gives
$\Sigma'_{\mathsf h}=P_{\mathsf h}'(0)=-\pi^4a_{-3}=-\pi s_\tau\Sigma_{\mathsf n}$, the first identity
of \eqref{gff:mgk:periodTwo}. (As a by-product, $P_{\mathsf h}(0)=a_{-2}\pi^2=\Sigma_{\mathsf h}$,
because $\mathsf h(-k)=0$ for $k\geq1$.) The undifferentiated sum follows from the Euler integral
of Section~\ref{gff:sec:Kdef}: with the moments $I_j$ listed before \eqref{gff:eq:KproofMomentRelations}, at $t=\tau$,
termwise integration with $c_n^2B(n+\frac12,\tau)=\Gamma(\tau)\mathsf h(n)/\pi$, where
$c_n=(\frac12)_n/n!$, gives
\begin{equation}
 \Sigma_{\mathsf h}=\frac{\pi^{3/2}}{\Gamma(\frac12+\tau)}\,
 {}_3F_2\!\left(\begin{matrix}\frac12,\frac12,\frac12\\1,\frac12+\tau\end{matrix};1\right)
 =\frac{2I_1}{\Gamma(\tau)}
 =\frac{\pi c_\tau\Gamma(\tau)}{\Gamma(\frac12+\tau)^2}\,\mathfrak G(\tau)+\frac{\pi^2}{\Gamma(\tau)}\,\mathfrak F(\tau),
 \label{gff:mgk:Sh}
\end{equation}
the last step by $I_1=c_\tau I_3+s_\tau I_4$ of \eqref{gff:eq:KproofMomentRelations} and the values of
$I_3,I_4$ stated after it. By \eqref{gff:mgk:singleValues} the two terms are $c_\tau\Sigma_{\mathsf n}$
and $s_\tau^2\Sigma_{\mathsf k}$, which reproves the second identity of \eqref{gff:mgk:periodTwo}.

\paragraph{Third evaluation: the logarithmic expansion of $K$.}
The derivative sum is also a single elliptic moment. For $0<x<1$ the logarithmic connection
formula \cite{gff:DLMF15} gives
\[
 K(1-x)=\frac12\sum_{n\geq0}c_n^2x^n\bigl[2\psi_0(n+1)-2\psi_0(n+\tfrac12)-\log x\bigr],
\]
with positive terms. Since $\int_0^1x^{n-1/2}(1-x)^{\tau-1}\log x\dd x
=B(n+\frac12,\tau)[\psi_0(n+\frac12)-\psi_0(n+\frac12+\tau)]$ and
$\mathsf h'(n)/\mathsf h(n)=3\psi_0(n+\frac12)-2\psi_0(n+1)-\psi_0(n+\frac12+\tau)$, Tonelli's theorem
gives
\[
 I_2=E(\tau,\tfrac12)=\int_0^1x^{-1/2}(1-x)^{\tau-1}K(1-x)\dd x
 =-\frac{\Gamma(\tau)}{2\pi}\,\Sigma'_{\mathsf h},
\]
so $\Sigma'_{\mathsf h}=-2\pi I_2/\Gamma(\tau)$. By \eqref{gff:eq:KproofMomentRelations},
$I_2=s_\tau I_3=\frac{\pi^2}2\mathfrak F(1-\tau)$ and
$I_1=c_\tau I_3+s_\tau I_4=\frac{\pi^2}2[\mathfrak F(\tau)+\cot(\pi\tau)\mathfrak F(1-\tau)]$. Hence
\[
 \Sigma_{\mathsf h}=\frac{\pi^2}{\Gamma(\tau)}\bigl[\mathfrak F(\tau)+\cot(\pi\tau)\,\mathfrak F(1-\tau)\bigr],\qquad
 \Sigma'_{\mathsf h}=-\frac{\pi^3}{\Gamma(\tau)}\,\mathfrak F(1-\tau),\qquad
 \Sigma_{\mathsf k}=\Gamma(\tau)\Gamma(1-\tau)^2\,\mathfrak F(\tau),
\]
in agreement with \eqref{gff:mgk:periodTwo} and \eqref{gff:mgk:singleValues}. This route uses the elliptic
moments and the moment relations in place of periodization.

\subsubsection{The regulated traces and the mixed-trace step}\label{gff:mgk:regulated}

\paragraph{Theorem (regulated traces).}
\emph{For $0<\tau<\frac12$ both traces converge absolutely, and}
\begin{equation}\boxed{
 C(\tau)\langle t_\ell^2R_\tau\rangle_2=\pi^2\bigl[\mathfrak F(\tau)^2+\mathfrak F(1-\tau)^2\bigr],\qquad
 \langle t_\ell^2S_\tau\rangle_2=C(\tau)\,\mathcal X_2(\tau)^2 .}
 \label{gff:mgk:regulatedEq}
\end{equation}
\emph{Proof.} Insert \eqref{gff:mgk:periodTwo} into \eqref{gff:mgk:factorized}. Since
$-2\Sigma_{\mathsf h}\Sigma'_{\mathsf h}=2\pi s_\tau\Sigma_{\mathsf n}\Sigma_{\mathsf h}$ and
$\Sigma_{\mathsf h}=c_\tau\Sigma_{\mathsf n}+s_\tau^2\Sigma_{\mathsf k}$,
\[
 -\frac1{\pi^2c_\tau}\Bigl[2\Sigma_{\mathsf h}\Sigma'_{\mathsf h}+\pi\frac{s_\tau}{c_\tau}\Sigma_{\mathsf h}^2\Bigr]
 =\frac{s_\tau}{\pi c_\tau^2}\,\Sigma_{\mathsf h}\bigl(2c_\tau\Sigma_{\mathsf n}-\Sigma_{\mathsf h}\bigr)
 =\frac{s_\tau}{\pi c_\tau^2}\bigl(c_\tau^2\Sigma_{\mathsf n}^2-s_\tau^4\Sigma_{\mathsf k}^2\bigr).
\]
Adding the $\mathsf k$ family, the factors $c_\tau^{-2}$ cancel:
\[
 \langle t_\ell^2R_\tau\rangle_2=\frac{s_\tau}\pi\bigl[\Sigma_{\mathsf n}^2+s_\tau^2\Sigma_{\mathsf k}^2\bigr]
 =\pi s_\tau\Gamma(1-\tau)^2\bigl[\mathfrak F(1-\tau)^2+\mathfrak F(\tau)^2\bigr],
\]
by \eqref{gff:mgk:singleValues}. Multiplication by $C(\tau)$, with
$C(\tau)\,\pi s_\tau\Gamma(1-\tau)^2=\pi s_\tau\Gamma(\tau)\Gamma(1-\tau)=\pi^2$, gives the first
identity. For the second, \eqref{gff:mgk:periodOne} in \eqref{gff:mgk:factorized} gives
\[
 \langle t_\ell^2S_\tau\rangle_2=\frac{\Sigma_{\mathsf m}^2}\pi\Bigl[4\cos^2(\tfrac{\pi\tau}2)\tan(\tfrac{\pi\tau}2)-s_\tau\Bigr]
 =\frac{s_\tau\Gamma(\tau)^2\mathcal X_2(\tau)^2}\pi=C(\tau)\,\mathcal X_2(\tau)^2 .\qquad\square
\]

The second identity is \eqref{gff:gdk:Sclosed} at $d=2$, which Section~\ref{gff:gdk:Sfamily} derives at
every dimension from the two-wheel convolution $W_2(d;b')=B_d(d-b',2b')$; here it is obtained from
Gauss's sum and the periodization (Dougall) identity of \eqref{gff:mgk:periodOne}. The second and third
evaluations of the single sums give the same result: with
$\mathcal A_\tau=\Gamma(\tau)\mathfrak G(\tau)/\Gamma(\frac12+\tau)^2$ and
$\mathcal B_\tau=\pi\mathfrak F(\tau)/\Gamma(\tau)$ one has $\Sigma_{\mathsf h}=\pi(c_\tau\mathcal A_\tau+\mathcal B_\tau)$,
$\Sigma'_{\mathsf h}=-\pi^2s_\tau\mathcal A_\tau$ and $\Sigma_{\mathsf k}=\pi\mathcal B_\tau/s_\tau^2$, and
\eqref{gff:mgk:factorized} becomes
$\langle t_\ell^2R_\tau\rangle_2=\pi s_\tau\mathcal A_\tau^2+\pi\mathcal B_\tau^2/s_\tau$, that is
\[
 C(\tau)\langle t_\ell^2R_\tau\rangle_2=\pi^2\mathfrak F(\tau)^2
 +s_\tau^2\frac{\Gamma(\tau)^4}{\Gamma(\frac12+\tau)^4}\,\mathfrak G(\tau)^2,
\]
which is \eqref{gff:mgk:regulatedEq} by \eqref{gff:eq:KproofThomae}. The assembly with the third evaluation
is checked in exact algebra in \path{LRI_C5_d2_proof/K2_d1template_fix1/exact_algebra_v2.wl} (part S2),
and the periodization route numerically (Section~\ref{gff:mgk:status}).
At $d=1$ the regulated master is linear in $\mathfrak F$ and $\mathfrak G$
\eqref{gff:eq:KproofTwoFunctions}; at $d=2$ it is a sum of two squares, which is where $\Cat^2$ comes
from: $u=4\log2-8\Cat/\pi$ enters the extraction below through $2\pi^2u^2$. The two-mode lemma is
the second-power analogue of \eqref{gff:fp:Omega4Modes}; here the decay $|y|^{-\tau}$ supplied by the
regulator is what removes the third harmonics. The periodization route uses neither elliptic
integrals nor the moment relations, only residues, Gauss's sum, Thomae's relation and
periodization, whereas the position-space proofs of \eqref{gff:gc:K2master} relate the same single sums
to the elliptic moments. At $\tau>0$ the only single sums are $\Sigma_{\mathsf h},\Sigma'_{\mathsf h},
\Sigma_{\mathsf k},\Sigma_{\mathsf m},\Sigma'_{\mathsf m}$ and $\Sigma_{\mathsf n}$, all absolutely convergent;
the divergent single sums of the factorization at $\tau=0$, and their finite parts, are never
evaluated.

\paragraph{Corollary (the mixed-trace step at $d=2$).}
\emph{For $0<\tau<\frac12$,}
\begin{equation}
 J_{2\tau}=C(\tau)\langle t_\ell^2R_\tau\rangle_2 ,
 \label{gff:mgk:mixedTrace}
\end{equation}
\emph{with the trace read on the full Mellin line, equivalently as the real part of the half-line
integral \eqref{gff:gs:trace}. This is input (M) of Section~\ref{gff:gdk:conventions} at $d=2$ on this
interval.}

\emph{Proof.} By \eqref{gff:gc:K2master}, $J_{2\tau}=\pi^2[\mathfrak F(\tau)^2+\mathfrak F(1-\tau)^2]$ for
$0<\tau<1$, and by \eqref{gff:mgk:regulatedEq} the right side of \eqref{gff:mgk:mixedTrace} has the same
value for $0<\tau<\frac12$. \hfill$\square$

The proofs of \eqref{gff:gc:K2master} in Section~\ref{gff:gc:K2proof} and in Section~\ref{gff:mgk:riemann} use
neither the trace nor (M), and the proof of \eqref{gff:mgk:regulatedEq} does not use (M); so
\eqref{gff:mgk:mixedTrace} is proved, not assumed. With the second identity of
\eqref{gff:mgk:regulatedEq}, which is the two-wheel identity at $d=2$, it gives \eqref{gff:gs:Ksubtraction}
at $d=2$ for $0<\tau<\frac12$, and with part (i) of the corollary of
Section~\ref{gff:mgk:regularization} the derivation of \eqref{gff:gs:Kspectral} from
\eqref{gff:gs:Kdefinition} given after \eqref{gff:gs:Ksubtraction} is complete at $d=2$. As a
normalization check, at $\tau=\frac12$ one has $R_{1/2}=t_\ell$ and $C(\frac12)=1$, so the right
side of \eqref{gff:mgk:mixedTrace} is $\langle t_\ell^3\rangle_2=3Q_2$ by \eqref{gff:gs:primitiveWheel},
while \eqref{gff:gc:K2master} gives $J_1=3Q_2$ by Dixon's sum. The point $\tau=\frac12$ lies outside
the interval of the corollary, because the two pole families of $\Theta^{\rm R}_\ell$ merge there,
so this is a consistency check only. The reading of (M) for complex $\tau$, and (M) at other
dimensions, are not addressed here.

\subsubsection{The spectral form at \texorpdfstring{$d=2$}{d=2}}\label{gff:mgk:spectral}

\paragraph{Theorem.}
\emph{The right side of \eqref{gff:gs:Kspectral} at $d=2$ equals $32\Cat^2-28\zeta(3)\log2$.}

\emph{Proof.} By part (iii) of the corollary of Section~\ref{gff:mgk:regularization} and by
\eqref{gff:mgk:regulatedEq}, for $0<\tau<\frac12$,
\begin{equation}
 C(\tau)\mathcal D(\tau)=\pi^2\mathfrak F(\tau)^2+\pi^2\mathfrak F(1-\tau)^2-C(\tau)^2\mathcal X_2(\tau)^2 .
 \label{gff:mgk:CD}
\end{equation}
As in Section~\ref{gff:gc:K2laurent}, \eqref{gff:eq:KproofThomae} gives
$\pi\mathfrak F(1-\tau)=\tau^{-1}\mathcal Q(\tau)\mathfrak G(\tau)$ with
$\mathcal Q(\tau)=\pi\Gamma(1+\tau)/[\Gamma(1-\tau)\Gamma(\frac12+\tau)^2]$, holomorphic and zero-free on
$|\tau|<\frac12$ with $\mathcal Q(0)=1$; and $\tau C(\tau)\mathcal X_2(\tau)$ is holomorphic there with
value $1$ at $0$. So the right side of \eqref{gff:mgk:CD} is meromorphic on $|\tau|<\frac12$. The left
side is holomorphic there by the corollary of Section~\ref{gff:mgk:regularization}. Two functions
meromorphic on a punctured disk that agree on $(0,\frac12)$ agree identically, so their Laurent
coefficients at $0$ coincide. From $\log\Gamma(c+x)-\log\Gamma(c)=\sum_k\psi_{k-1}(c)x^k/k!$,
$\psi_0(\frac12)-\psi_0(1)=-2\log2$, $\psi_k(1)=(-1)^{k+1}k!\,\zeta(k+1)$ and
$\psi_k(\frac12)=(2^{k+1}-1)\psi_k(1)$ \cite{gff:DLMFPolygamma},
\[
 \begin{gathered}
 \log\mathcal Q(\tau)^2=8\log2\,\tau-\pi^2\tau^2+8\zeta(3)\tau^3-\frac{\pi^4}6\tau^4+O(\tau^5),\\
 \log\bigl[\tau C(\tau)\mathcal X_2(\tau)\bigr]^2=8\log2\,\tau+\zeta(3)\tau^3+O(\tau^5).
 \end{gathered}
\]
The first is the series of Section~\ref{gff:gc:K2laurent}. Hence
$C\mathcal D=\pi^2\mathfrak F(\tau)^2+\tau^{-2}\{\mathcal Q(\tau)^2\mathfrak G(\tau)^2-[\tau C(\tau)\mathcal X_2(\tau)]^2\}$,
and with \eqref{gff:mgk:Taylor} and the coefficients $\mathcal Q_k=[\tau^k]\mathcal Q(\tau)^2$ of
Section~\ref{gff:gc:K2laurent},
\begin{equation}\begin{aligned}
 [\tau^{-2}]\,C\mathcal D&=1-1=0,\qquad [\tau^{-1}]\,C\mathcal D=8\log2-8\log2=0,\\
 [\tau^{0}]\,C\mathcal D&=\pi^2+\mathcal Q_2-32\log^22=0,\\
 [\tau^{1}]\,C\mathcal D&=2\pi^2u+2\sigma+7\zeta(3)-8\pi^2\log2,\\
 [\tau^{2}]\,C\mathcal D&=2(\varrho+\pi^2v)+\pi^2u^2+\frac{\pi^4}3+56\zeta(3)\log2\\
 &\quad-32\pi^2\log^22+16\sigma\log2 .
\end{aligned}\label{gff:mgk:laurent}\end{equation}
In the last line the $\tau^2$ coefficient of $\pi^2\mathfrak F^2$ is $\pi^2(u^2+2v)$, that of
$\tau^{-2}\mathcal Q^2\mathfrak G^2$ is $\mathcal Q_4+2\sigma\mathcal Q_1+2\varrho$ with $\mathcal Q_1=8\log2$
and $\mathcal Q_4$ of Section~\ref{gff:gc:K2laurent}, and that of $\tau^{-2}[\tau C\mathcal X_2]^2$ is
$\frac{512}3\log^42+8\zeta(3)\log2$; in the line before, $\mathcal Q_3$ of Section~\ref{gff:gc:K2laurent}
and $[\tau^3][\tau C\mathcal X_2]^2=\frac{256}3\log^32+\zeta(3)$ enter. The vanishing
of the first three coefficients is the consistency required by $\mathcal D=O(\tau^2)$. The
cancellation rule in \eqref{gff:mgk:Taylor} removes $v$, $\varrho$ and $\pi^4$:
$[\tau^2]\,C\mathcal D=2\pi^2u^2-32\pi^2\log^22+16\sigma\log2+56\zeta(3)\log2$. With $u$ and $\sigma$
from \eqref{gff:mgk:Taylor}, $2\pi^2u^2=32\pi^2\log^22-128\pi\Cat\log2+128\Cat^2$ and
$16\sigma\log2=128\pi\Cat\log2-224\zeta(3)\log2$, so $[\tau^2]\,C\mathcal D=128\Cat^2-168\zeta(3)\log2$.
Finally \eqref{gff:mgk:Kreg} gives
\[
 14\zeta(3)\log2+\tfrac14\bigl(128\Cat^2-168\zeta(3)\log2\bigr)=32\Cat^2-28\zeta(3)\log2 .\qquad\square
\]

\paragraph{Corollary (a spectral moment).}
\begin{equation}
 \langle t_\ell^2E_2\rangle_2=\mathcal D_2=[\tau^1]\,C\mathcal D
 =2\pi^2u+2\sigma+7\zeta(3)-8\pi^2\log2=-21\zeta(3).
 \label{gff:mgk:E2}
\end{equation}
This is \eqref{gff:gdk:E2} at $d=2$, $\frac32\psi_2(\frac12)=-21\zeta(3)$, proved here without (M) or
(J), and with the interchange of $[\tau^2]$ and the trace justified by the corollary of
Section~\ref{gff:mgk:regularization}.

\paragraph{Relation to the definition.}
With \eqref{gff:mgk:mixedTrace}, $J_{2\tau}-J^\Gamma_{2\tau}=C(\tau)\mathcal D(\tau)-[J^\Gamma_{2\tau}-C(\tau)^2\mathcal X_2(\tau)^2]$,
and by \eqref{gff:gdk:Delta} at $d=2$ the bracket is $\frac32U\tau+AU\tau^2+O(\tau^3)$. So the order-$\tau$
coefficient of $J_{2\tau}-J^\Gamma_{2\tau}$ is $[\tau^1]C\mathcal D-\frac32U=-21\zeta(3)+21\zeta(3)=0$,
which is the vanishing found in Section~\ref{gff:gc:K2laurent}, and
$\frac14[\tau^2](J_{2\tau}-J^\Gamma_{2\tau})=\frac14[\tau^2]C\mathcal D-\frac{AU}4$ is the right side of
\eqref{gff:mgk:Kreg}. At $d=2$ the definition \eqref{gff:gs:Kdefinition}, the spectral form
\eqref{gff:gs:Kspectral} and the regulated trace therefore give the same expression identically, not
only the same number. The exact algebra of both routes is checked in
\path{LRI_C5_d2_simplify/proof/K2-parseval/symbolic_checks.py} (part B) and
\path{LRI_C5_d2_simplify/proof/synth/synth_checks_K2.py}.

\subsubsection{The series at \texorpdfstring{$d=2$}{d=2}}\label{gff:mgk:series}
At $q=\frac12$
the data of \eqref{gff:gs:seriesBase} and \eqref{gff:gs:Kseries} are $M^2=16$, $T_{\ell n}=c_mc_n$ and
$S_{\ell n}=e_me_n$, with $m=\ell+n$, $c_k=(\frac12)_k/k!=\Gamma(k+\frac12)/[\sqrt\pi\,\Gamma(k+1)]$ and
$e_k=k!/(\frac32)_k$. Write
\[
 D(\ell)=\sum_{n\geq0}\Bigl[\frac{M^2S_{\ell n}^2}8-\frac{T_{\ell n}^2\mathcal R_{\ell n}}4\Bigr],
\]
so that \eqref{gff:gs:Kseries} at $d=2$ reads $-\frac{AU}4+\sum_\ell d_\ell D(\ell)$.

\paragraph{Lemma (local forms).}
\emph{Let $u\to0$. At $s=a+n+u$,}
\[
 g_\ell(s)=-\pi\cot(\pi u)\,c(m+u)\,c(n+u),\qquad c(x)=\frac{\Gamma(x+\frac12)}{\sqrt\pi\,\Gamma(x+1)},
\]
\emph{with $\psi_0(a+s)=\psi_0(m+\frac12+u)$, $\psi_0(b+s)=\psi_0(m+1+u)$,
$\psi_0(a-s)=\psi_0(n+1+u)+\pi\cot\pi u$ and $\psi_0(b-s)=\psi_0(n+\frac12+u)-\pi\tan\pi u$. At
$s=b+n+u$,}
\[
 g_\ell(s)=\tan(\pi u)\,\frac{\Gamma(m+1+u)\,\Gamma(n+1+u)}{\Gamma(m+\frac32+u)\,\Gamma(n+\frac32+u)},
\]
\emph{with $\psi_0(a+s)=\psi_0(m+1+u)$, $\psi_0(b+s)=\psi_0(m+\frac32+u)$,
$\psi_0(a-s)=\psi_0(n+\frac32+u)-\pi\tan\pi u$ and $\psi_0(b-s)=\psi_0(n+1+u)+\pi\cot\pi u$.}

\emph{Proof.} Use $\Gamma(z)\Gamma(1-z)=\pi/\sin\pi z$ and $\psi_0(1-z)-\psi_0(z)=\pi\cot\pi z$
\cite{gff:DLMF5,gff:DLMFPolygamma}; for example
$\Gamma(-n-u)/\Gamma(\frac12-n-u)=[\sin\pi(\frac12-n-u)/\sin\pi(-n-u)]\,\Gamma(n+\frac12+u)/\Gamma(n+1+u)
=-\cot(\pi u)\,\Gamma(n+\frac12+u)/\Gamma(n+1+u)$. \hfill$\square$

The products $c(m)c(n)=T_{\ell n}$ and $\Gamma(k+1)/\Gamma(k+\frac32)=2e_k/\sqrt\pi$ connect these
local forms with \eqref{gff:gs:seriesBase}. Inserting them into the holomorphic weight $F(s)$ of
\eqref{gff:gs:Kholomorphic} and expanding in $u$ gives the $q=\frac12$ case of the residues computed in
Section~\ref{gff:gs:residueProof}:
\begin{equation}
 \Res_{s=a+n}\frac{g_\ell^2F}4=\frac{T_{\ell n}^2\mathcal R_{\ell n}}4,\qquad
 \Res_{s=b+n}\frac{g_\ell^2F}4=-2e_m^2e_n^2=-\frac{M^2S_{\ell n}^2}8 .
 \label{gff:mgk:seriesResidues}
\end{equation}
Both are identities between Laurent coefficients, polynomial in the independent symbols
$\psi_j(m+\frac12),\psi_j(m+1),\psi_j(n+\frac12),\psi_j(n+1)$ (or their shifts by $\frac12$) with
exact $\pi$, $\psi_j(1)$ and $\psi_j(\frac12)$; they are verified exactly in
\path{LRI_C5_d2_proof/K2_gauss/check_residue_series.wl}.

\paragraph{Proposition (the per-spin identity).}
\emph{For every $\ell\geq0$ the inner series $D(\ell)$ converges, and}
\begin{equation}
 D(\ell)=\int_{-\ii\infty}^{\ii\infty}\frac{ds}{2\pi\ii}\,\frac{g_\ell(s)^2F(s)}4
 =\frac1{4\pi}\int_0^\infty t_\ell(\nu)^2\,w_{\KK}\dd\nu .
 \label{gff:mgk:innerSpin}
\end{equation}
\emph{Proof.} On $s=\ii\nu$ one has $g_\ell'/g_\ell=\psi_0(a+s)-\psi_0(a-s)-\psi_0(b+s)+\psi_0(b-s)=2\ii h$,
its derivative is $2j$, and $p(\ii\nu)=p$, $p'(\ii\nu)=\ii k$, $p''(\ii\nu)=r$ for the function $p(s)$
of Section~\ref{gff:gs:residueProof}. Hence $F(\ii\nu)=(c_*-p)(-\frac{h^2}2-\frac{5j}8)-\frac{hk}2-\frac{5r}{32}=w_{\KK}$,
which is real and even in $\nu$; so is $g_\ell(\ii\nu)=t_\ell(\nu)$. This gives the second equality.
For the first, close the contour to the right along $|s|=\rho_N:=a+N+\frac34$; these arcs stay at
distance at least $\frac14$ from both pole families $a+\mathbb Z_{\geq0}$ and $b+\mathbb Z_{\geq0}$. On
the arcs $|g_\ell(s)|\leq C|s|^{-1}$: indeed $\Gamma(a+s)/\Gamma(b+s)=O(|s|^{-1/2})$ by Stirling's
formula for $\Re s\geq0$, and
\[
 \frac{\Gamma(a-s)}{\Gamma(b-s)}=\frac{\sin\pi(b-s)}{\sin\pi(a-s)}\cdot\frac{\Gamma(1-b+s)}{\Gamma(1-a+s)},
\]
where the trigonometric ratio is bounded at distance at least $\frac14$ from $a+\mathbb Z$ and has
finite limits as $\Im s\to\pm\infty$, and the Gamma ratio is $O(|s|^{-1/2})$. In the same way
$\psi_0(a-s)=\psi_0(1-a+s)-\pi\cot\pi(a-s)$, and the analogous formula for $\psi_0(b-s)$, give
$|F|\leq C\log^2|s|$ on the arcs. The arc integrals are therefore $O(\rho_N^{-1}\log^2\rho_N)\to0$. At
fixed $\ell$ each residue family converges absolutely, with terms $O(n^{-2}\log^4n)$, and the
residues are summed in increasing $n$. Closing upward to the right is clockwise, so
\eqref{gff:mgk:seriesResidues} gives
$\int\frac{ds}{2\pi\ii}\frac{g_\ell^2F}4=-\sum_n[\Res_{a+n}+\Res_{b+n}]=D(\ell)$. \hfill$\square$

This is the per-spin identity behind \eqref{gff:gs:Kseries}, at $d=2$, with the contour estimates that
Section~\ref{gff:gs:residueProof} does not write out.

\paragraph{Proposition (the series and the strip limit).}
\emph{(a) For each $r\in(0,\frac12)$ there is $C_r'$ with
$|t_\ell(\nu)^2w_{\KK}(\ell,\nu)|\leq C_r'(1+|\xi|)^{-3+2r}$; in particular
$\langle|t_\ell^2w_{\KK}|\rangle_2<\infty$. (b) $\sum_\ell d_\ell|D(\ell)|<\infty$, and \eqref{gff:gs:Kseries} at
$d=2$ equals $-\frac{AU}4+\frac14\langle t_\ell^2w_{\KK}\rangle_2=32\Cat^2-28\zeta(3)\log2$.
(c) Let $g(m,n)$ be the symmetric square-lattice summand with $g(\ell+n,n)$ the summand of $D(\ell)$,
written as a finite sum $\sum_i\alpha_iA_i(m)B_i(n)$ of products of real sequences of bounded
variation that are $O(\log^3(m+2)/(m+1))$. Put $\hat A_i(z)=\sum_mA_i(m)z^m$. Then
$h(\theta)=\sum_i\alpha_i\hat A_i(e^{\ii\theta})\overline{\hat B_i(e^{\ii\theta})}$ is continuous on
$(0,2\pi)$, its Fourier coefficients are $D(|\ell|)$, and
$\lim_{\theta\to0}h(\theta)=\sum_{\ell\geq0}d_\ell D(\ell)$.} Such a factorization, with $27$ distinct
single sequences, is the exact symbolic one behind the $70$-digit evaluation
(\path{LRI_C5_d2_simplify/factor_K2.py}, output \path{K2_factorized.json}); each sequence is $c_k^2$
or $e_k^2$ times polygamma differences with at most one factor $\psi_0(k+1)$.

\emph{Proof.} (a) Let $H(\tau)=\frac12C(\tau)[(R_\tau-S_\tau)(\ii\nu)+(R_\tau-S_\tau)(-\ii\nu)]$. It is
holomorphic in $|\tau|<\frac12$, since the pole of $C$ is cancelled by the double zero of the even
part, and $[\tau^2]H=w_{\KK}$ by the corollary of Section~\ref{gff:mgk:regularization}. By
\eqref{gff:mgk:RminusS}, $|H|\leq C(1+|\xi|)^{-1+2r}$ on $|\tau|=r$, and Cauchy's estimate gives
$|w_{\KK}|\leq r^{-2}C(1+|\xi|)^{-1+2r}$; multiply by $t_\ell^2\leq C(1+|\xi|)^{-2}$. (b) By
\eqref{gff:mgk:innerSpin}, $|D(\ell)|\leq\frac1{8\pi}\int_{\RR}|t_\ell^2w_{\KK}|\dd\nu$, so (a) gives absolute
convergence of the spin sum, and the value follows from the theorem of Section~\ref{gff:mgk:spectral}.
(c) By Dirichlet's test (bounded variation and $A_i(m)\to0$) the series $\hat A_i(e^{\ii\theta})$
converge locally uniformly on $(0,2\pi)$, so $h$ is continuous there. Since $A_i,B_i\in\ell^2$, the
product $\hat A_i\overline{\hat B_i}$ is integrable on the circle, and its $\ell$th Fourier
coefficient is the correlation $\sum_nA_i(n+\ell)B_i(n)$, with $A_i(k)=0$ for $k<0$. Summing over
$i$, the coefficient is $\sum_ng(n+\ell,n)=D(\ell)$ for $\ell\geq0$, and $D(|\ell|)$ for $\ell<0$ by
the symmetry of $g$. By (b) these coefficients are absolutely summable, so the Fourier series
converges uniformly to a continuous function equal to $h$ almost everywhere, hence everywhere on
$(0,2\pi)$. Letting $\theta\to0$ gives $\sum_{\ell\in\mathbb Z}D(|\ell|)=\sum_{\ell\geq0}d_\ell D(\ell)$.
\hfill$\square$

At $d=2$ the inner series of \eqref{gff:gs:Kseries} therefore sum in an absolutely convergent spin order,
although the unrestricted double series is not absolutely convergent: its square-lattice single sums
diverge like $\sum_k\log^pk/k$. Part (c) shows that the strip limit $\theta\to0$ of the lattice sum
with weight $e^{\ii(m-n)\theta}$, used for the $70$-digit evaluation of
Section~\ref{gff:gc:d2closedForms} (\path{LRI_C5_d2_simplify/eval_K2.py}), computes exactly this
ordered sum. The sharper bound \eqref{gff:mgk:sharp} gives $D(\ell)=O(\ell^{-3+2\rho_0})$.

\paragraph{Remark (why the single sums at $\tau=0$ are avoided).}
This is a numerical observation and is not used. At $\tau=0$ the $\mathcal R$ part of the bilinear
form $g(m,n)$ has rank $9$, and the $S_{\ell n}$ part adds one more. In the variables
$\psi_{j-1}(k+\frac12)-\psi_{j-1}(\frac12)$ and $\psi_{j-1}(k+1)-\psi_{j-1}(1)$, the span of all parameter
derivatives of ${}_2F_1(\frac12+\alpha,\frac12+\beta;1+\gamma;z)$ through total order four has dimension
$22$, and the row space of the form has three directions outside it, at weights $2$, $3$ and $4$
(\path{LRI_C5_d2_proof/K2_gauss/rowspace.py}). An integer-relation search on the $70$-digit single sums
(\path{LRI_C5_d2_proof/K2_gauss/pslq_singles.py}) indicates that individual sums contain constants
absent from the final answer, for example
\[
 \sum_kc_k^2\bigl[\psi_2(k+\tfrac12)-\psi_2(k+1)\bigr]
 \overset{\rm num}{=}3\pi^3+4\pi\log^22-64\Cat\log2-128\,\Im\operatorname{Li}_3\bigl(\tfrac{1+\ii}2\bigr),
\]
while the weight-five sums are not identified in a small basis. The regulator $\tau>0$ makes every
single sum convergent, and the only non-Gamma functions left are $\mathfrak F$ and $\mathfrak G$,
whose unknown coefficients $v,\varrho$ cancel by \eqref{gff:eq:KproofTaylorCancellation}.

\subsubsection{A second position-space proof: Riemann's bilinear relation}\label{gff:mgk:riemann}
This subsubsection proves \eqref{gff:gc:K2master} from the integral
\[
 J_{2\tau}=\frac1{\pi^2}\int_{\mathbb C^2}\frac{d^2x\,d^2y}{|x|\,|1-x|\,|y|\,|1-y|^{2-2\tau}\,|x-y|},
 \qquad0<\tau<1,
\]
of Section~\ref{gff:gc:K2proof}, without the cell decomposition and rotation of that section. Here
$d^2x$ is Lebesgue measure on $\mathbb C\cong\RR^2$, so $dx\wedge d\bar x=-2\ii\,d^2x$; $K(m)$ is the
parameter convention, holomorphic on $\mathbb C\setminus[1,\infty)$; $I_1,\dots,I_6$ are the moments
listed before \eqref{gff:eq:KproofMomentRelations}, at $t=\tau$; $\mathbb C_{\rm cut}=\mathbb C\setminus((-\infty,0]\cup[1,\infty))$;
all powers are principal; and the letters $\omega,\Omega_y,\Pi_j$ below are local (they are not the
$\omega_\tau$ of Section~\ref{gff:gc:K2proof} or the $\Omega_r$ of \eqref{gff:fp:OmegaWheel}).

\paragraph{Lemma (properties of $K$).}
(i) For $m\in\mathbb C\setminus[1,\infty)$, $K(m)=\frac12\int_0^1ds/\sqrt{s(1-s)(1-ms)}$, by $s=\sin^2\theta$
in the definition of Section~\ref{gff:c1:notation}. (ii) For $m\in\mathbb C\setminus[1,\infty)$,
$K(m)=(1-m)^{-1/2}K(\frac m{m-1})$ \cite{gff:DLMF15}; in particular $K(x)=(1-x)^{-1/2}K(\frac x{x-1})$ for
$x<0$, with $\frac x{x-1}\in(0,1)$, and $K(1-x)=x^{-1/2}K(1-\frac1x)$ for $x>1$. (iii) For $x>1$,
$K(x\pm\ii0)=x^{-1/2}[K(\frac1x)\pm\ii K(1-\frac1x)]$. (iv) For $|1-m|<1$, $m\notin[1,\infty)$,
$K(m)=\frac12\sum_nc_n^2(1-m)^n[2\psi_0(n+1)-2\psi_0(n+\frac12)-\log(1-m)]$ \cite{gff:DLMF15}. (v) There
is $C$ such that, for all $m\in\mathbb C\setminus[1,\infty)$ and for the boundary values (iii),
$|K(m)|\leq C(1+|\log|1-m||)$ for $|m|\leq3$ and $|K(m)|\leq C|m|^{-1/2}(1+\log|m|)$ for $|m|\geq3$.

\emph{Proof.} (iii) In (i) put $m=x\pm\ii\eta$ with $\eta\downarrow0$. For $s>1/x$,
$1-ms\to(1-xs)\mp\ii0$ with $1-xs<0$, so $(1-ms)^{-1/2}\to\pm\ii(xs-1)^{-1/2}$. Since $|1-ms|\geq|1-xs|$
and $(s(1-s)|1-xs|)^{-1/2}\in L^1(0,1)$, dominated convergence applies. On $(0,1/x)$ substitute
$s=\sigma/x$, and on $(1/x,1)$ substitute $s=1-(1-\frac1x)\sigma$; the two pieces become
$x^{-1/2}K(1/x)$ and $\pm\ii x^{-1/2}K(1-1/x)$. The restriction of $K$ to the upper (lower) half-plane
continues analytically across $(1,\infty)$, because the hypergeometric equation has analytic
coefficients there, so $K$ extends continuously from each side with these values (the $+$ case is the
formula used in the proof of \eqref{gff:eq:KproofMomentRelations}). (v) If $|1-m|\leq\frac12$, use (iv)
and $|\log(1-m)|\leq|\log|1-m||+\pi$. If $|m|\leq3$ and $|1-m|\geq\frac12$, use (i) and
$|1-ms|^2=(1-s\Re m)^2+s^2(\Im m)^2$: for $\Re m\leq\frac34$, $|1-ms|\geq\frac14$; for $\Re m\geq\frac54$,
$|1-ms|\geq|1-s\Re m|$ and $\int_0^1(s(1-s)|1-xs|)^{-1/2}ds$ is bounded for $x\in[\frac54,3]$; for
$\frac34<\Re m<\frac54$, $(\Im m)^2\geq\frac3{16}$, so $|1-ms|\geq\frac38$ for $s\leq\frac12$ and
$|1-ms|\geq\frac{\sqrt3}8$ for $s\geq\frac12$. If $|m|\geq3$, (ii) gives
$K(m)=(1-m)^{-1/2}K(\frac m{m-1})$ with $|1-\frac m{m-1}|=|m-1|^{-1}\leq\frac12$, which reduces to the
first case. The boundary values obey the same bounds by continuity. \hfill$\square$

\paragraph{Green's theorem on slit domains.}
Below, Green's theorem $\int_Vd\alpha=\oint_{\partial V}\alpha$, or Cauchy's theorem, is applied to a
domain $V=\{r<|z-e|$ for the marked points $e$, $|z|<R\}\setminus\Sigma$, where $\Sigma$ is a finite
union of segments (slits) and $\alpha=\mathsf g\,d\bar z$ or $\mathsf g\,dz$. In every case $\mathsf g$ is
real-analytic and bounded on $V$ and extends continuously to each open slit from each side, and
$d\alpha$ is integrable. Delete the $\eta$-neighborhood of $\Sigma$; on the closure of what remains
$\alpha$ is $C^1$, so Green's theorem applies. As $\eta\to0$, the area integral converges by
dominated convergence; the integrals over the two curves parallel to a slit converge to the integrals
of the one-sided boundary values, by uniform convergence on compact parts of the open slits and
boundedness; and the remaining boundary pieces have length $O(\eta)$. So the slit edges enter with
the positive orientation of $V$: each edge is traversed with $V$ on its left.

\paragraph{Reduction to the inner star.}
The integrand is positive, so by Tonelli's theorem
\begin{equation}
 J_{2\tau}=\frac1{\pi^2}\int_{\mathbb C}d^2y\,|y|^{-1}|1-y|^{2\tau-2}\,\mathcal I(y),\qquad
 \mathcal I(y)=\int_{\mathbb C}\frac{d^2x}{|x|\,|1-x|\,|x-y|}\in(0,\infty],
 \label{gff:mgk:starReduction}
\end{equation}
with both sides in $[0,\infty]$. For $y\notin\{0,1\}$, $\mathcal I(y)<\infty$: the three factors are
locally integrable in $\RR^2$, and the integrand is $O(|x|^{-3})$ at infinity.

\paragraph{Lemma (the inner star is a Riemann bilinear relation).}
\emph{For $y\in\mathbb C_{\rm cut}$,}
\begin{equation}\boxed{
 \mathcal I(y)=8\,\Re\bigl[K(y)\,\overline{K(1-y)}\bigr].}
 \label{gff:mgk:star}
\end{equation}
\emph{Proof.} Fix $y\in\mathbb C_{\rm cut}$ and put $\Pi_y(x)=x(x-1)(x-y)$.

\emph{(a) The cut.} Let $\Gamma_y=[0,y]\cup[y,1]\cup[1,+\infty)$, oriented from $0$ to $\infty$. Since
$y\notin(-\infty,0]\cup[1,\infty)$, the segment $[0,y]$ does not meet $[1,\infty)$; $[y,1]$ meets
$[1,\infty)$ only at $1$, and $[0,y]$ meets $[y,1]$ only at $y$ (for $y\in(0,1)$ the three pieces are
collinear). So $\Gamma_y$ is a simple arc from $0$ to $\infty$, and $\mathbb C\setminus\Gamma_y$ is
simply connected. $\Pi_y$ has no zero there, so there is a holomorphic $\omega$ on
$\mathbb C\setminus\Gamma_y$ with $\omega^2=1/\Pi_y$, and $|\omega|^2=1/(|x||1-x||x-y|)$.

\emph{(b) Edge values and jumps.} Near any point of an open piece of $\Gamma_y$, $\Pi_y$ has two
holomorphic square roots on a small disk, and $\omega$ equals $\pm1/(\text{that root})$ on each side.
Hence $\omega$ has continuous boundary values $\omega_L$ (left of $\Gamma_y$) and $\omega_R$ (right of
$\Gamma_y$) on each open piece, and the ratio $\omega_R/\omega_L\in\{\pm1\}$ is a constant
$\varepsilon_1,\varepsilon_2,\varepsilon_3$ on $(0,y)$, $(y,1)$, $(1,\infty)$. A small circle around $0$
crosses $\Gamma_y$ once, on $(0,y)$; since $\Pi_y$ has a simple zero inside, $\Pi_y^{-1/2}$ changes sign
along the circle, so $\varepsilon_1=-1$. A small circle around $y$ crosses $(0,y)$ and $(y,1)$ once each,
so $\varepsilon_1\varepsilon_2=-1$ and $\varepsilon_2=+1$. A small circle around $1$ crosses $(y,1)$ and
$(1,\infty)$, so $\varepsilon_3=-1$. The direction of crossing does not matter, because
$\varepsilon_k^{-1}=\varepsilon_k$.

\emph{(c) Primitive.} $|\omega|\leq C|x-e|^{-1/2}$ near $e\in\{0,y,1\}$, and $|\omega|\leq C|x|^{-3/2}$ for
large $|x|$. Let $\Omega_y$ be the primitive of $\omega$ on $\mathbb C\setminus\Gamma_y$ with
$\Omega_y(x)\to0$ as $x\to0$; the limit exists by the Cauchy criterion, since radial and circular paths
near $0$ give $|\Omega_y(x')-\Omega_y(x'')|\leq C(|x'|^{1/2}+|x''|^{1/2})$. $\Omega_y$ is bounded: it
is continuous up to the edges from each side, has limits at $0,y,1$ within each adjacent sector, and
for $|x|\geq R_0:=2+|y|$ it differs from its value on the circle $|x|=R_0$ by the integral along a
radial segment, which does not meet $\Gamma_y$ (for $\arg x\neq0$) and is $O(R_0^{-1/2})$. The edge
values satisfy $d\Omega_{y,L}=\omega_L\,dx$ and $d\Omega_{y,R}=\omega_R\,dx$ along $\Gamma_y$. At the
corners $y$ and $1$ the left edges of the two adjacent pieces border the same sector, and so do the
right edges; $|\omega|$ is integrable there, so the edge values are continuous through the corners.
Put
\[
 \Pi_1=\int_0^y\omega_L\,dx,\qquad \Pi_2=\int_y^1\omega_L\,dx,\qquad \Pi_3=\int_1^\infty\omega_L\,dx
 \qquad(\text{along }\Gamma_y).
\]
By (b) the edge values are: on $(0,y)$, $\Omega_{y,L}=\int_0^x\omega_L$ and $\Omega_{y,R}=-\Omega_{y,L}$;
at $y$, $\pm\Pi_1$; on $(y,1)$, $\Omega_{y,L}=\Pi_1+\int_y^x\omega_L$ and
$\Omega_{y,R}=-\Pi_1+\int_y^x\omega_L$; at $1$, $\Pi_1+\Pi_2$ and $-\Pi_1+\Pi_2$; on $(1,\infty)$,
$\Omega_{y,L}=\Pi_1+\Pi_2+\int_1^x\omega_L$ and $\Omega_{y,R}=-\Pi_1+\Pi_2-\int_1^x\omega_L$. Thus
$\Omega_{y,L}+\Omega_{y,R}=0$ on $(0,y)$, $\Omega_{y,L}-\Omega_{y,R}=2\Pi_1$ on $(y,1)$, and
$\Omega_{y,L}+\Omega_{y,R}=2\Pi_2$ on $(1,\infty)$.

\emph{(d) Stokes.} On $\mathbb C\setminus\Gamma_y$,
$d(\Omega_y\bar\omega\,d\bar x)=\omega\bar\omega\,dx\wedge d\bar x=-2\ii|\omega|^2d^2x$. Apply Green's
theorem on slit domains to $\{x\notin\Gamma_y:\ |x|<R,\ |x-e|>r,\ e=0,y,1\}$. The positively oriented
boundary runs along the left edge of $\Gamma_y$ in its direction and along the right edge against it,
so
\[
 -2\ii\int|\omega|^2d^2x=\int_{\Gamma_y}\bigl(\Omega_{y,L}\bar\omega_L-\Omega_{y,R}\bar\omega_R\bigr)d\bar x
 +(\text{circle terms}).
\]
The circle terms vanish in the limit: on $|x|=r$, $|\Omega_y|\leq Cr^{1/2}$ and $|\omega|\leq Cr^{-1/2}$,
so the term is $O(r)$; on $|x-y|=r$ and $|x-1|=r$, $\Omega_y$ is bounded and $|\omega|\leq Cr^{-1/2}$, so
the term is $O(r^{1/2})$; on $|x|=R$, $\Omega_y$ is bounded and $|\omega|\leq CR^{-3/2}$, so the term is
$O(R^{-1/2})$. The edge integrands are integrable. Letting $r\to0$ and $R\to\infty$ (monotone
convergence on the left) gives $-2\ii\,\mathcal I(y)=\int_{\Gamma_y}(\Omega_{y,L}\bar\omega_L-\Omega_{y,R}\bar\omega_R)\,d\bar x$.
By (b) and (c) the integrand is $0$ on $(0,y)$, $2\Pi_1\bar\omega_L$ on $(y,1)$ and $2\Pi_2\bar\omega_L$
on $(1,\infty)$. Hence $-2\ii\,\mathcal I(y)=2\Pi_1\bar\Pi_2+2\Pi_2\bar\Pi_3$.

\emph{(e) Cauchy.} The same exhaustion applied to $\omega\,dx$ gives
$\int_{\Gamma_y}(\omega_L-\omega_R)\,dx=0$, that is $2\Pi_1+0+2\Pi_3=0$, so $\Pi_3=-\Pi_1$. Therefore
$-2\ii\,\mathcal I(y)=2(\Pi_1\bar\Pi_2-\bar\Pi_1\Pi_2)=4\ii\,\Im(\Pi_1\bar\Pi_2)$, that is
$\mathcal I(y)=2\,\Im(\bar\Pi_1\Pi_2)$. This is Riemann's bilinear relation for the holomorphic
differential $dx/w$ on the elliptic curve $w^2=\Pi_y(x)$, whose periods are $2\Pi_1$ and $2\Pi_2$; the
factor $\frac12$ in $\mathcal I=\frac12\Im(\overline{2\Pi_1}\,2\Pi_2)$ is the inverse of the degree $2$ of
the covering $(x,w)\mapsto x$.

\emph{(f) The periods.} On $[0,y]$ put $x=ys$. Then $\Pi_y(ys)=y^2s(1-s)(1-ys)$, and
$s(1-s)(1-ys)\notin(-\infty,0]$ for $s\in(0,1)$ because $ys\notin[1,\infty)$. Hence
$\omega_L(ys)=\eta_1/(y\sqrt{s(1-s)(1-ys)})$ with a constant $\eta_1\in\{\pm1\}$, and by (i) of the
lemma on $K$, $\Pi_1=2\eta_1K(y)$. On $[y,1]$ put $x=1-(1-y)\sigma'$. Then
$\Pi_y=-(1-y)^2\sigma'(1-\sigma')(1-(1-y)\sigma')$, with $\sigma'(1-\sigma')(1-(1-y)\sigma')\notin(-\infty,0]$,
so $\omega_L=\eta_2/(\ii(1-y)\sqrt{\sigma'(1-\sigma')(1-(1-y)\sigma')})$ and, with $dx=-(1-y)d\sigma'$,
$\Pi_2=-2\ii\eta_2K(1-y)$. So (e) gives $\mathcal I(y)=8\epsilon_y\Re[K(y)\overline{K(1-y)}]$ with
$\epsilon_y=-\eta_1\eta_2\in\{\pm1\}$.

\emph{(g) The sign.} $\mathcal I(y)>0$, so $\Re[K(y)\overline{K(1-y)}]\neq0$ on $\mathbb C_{\rm cut}$. This
function is continuous on the connected set $\mathbb C_{\rm cut}$ and equals $K(\frac12)^2>0$ at $y=\frac12$,
so it is positive everywhere and $\epsilon_y\equiv1$. \hfill$\square$

For real $y\in(0,1)$ the sign can also be read off directly: continuation along the upper edge gives
$\omega_L=|\Pi_y|^{-1/2}$ on $(0,y)$, $\ii|\Pi_y|^{-1/2}$ on $(y,1)$ and $-|\Pi_y|^{-1/2}$ on $(1,\infty)$,
hence $\Pi_1=2K(y)$, $\Pi_2=2\ii K(1-y)$ and $\Pi_3=-2K(y)$ (using $\int_1^\infty dx/\sqrt{\Pi_y}=2K(y)$
via $x=1/s$), which confirms $\Pi_3=-\Pi_1$ and $\mathcal I(y)=8K(y)K(1-y)$. The algebra of (d)--(f) is
checked exactly in \path{LRI_C5_d2_proof/K2_d1template_fix1/exact_algebra_v2.wl} (part S0a).

\paragraph{Edge values of the outer integrand.}
Let $\varphi_\tau(y)=y^{-1/2}(1-y)^{\tau-1}$, $\mathsf u=\varphi_\tau K(y)$ and $\mathsf v=\varphi_\tau K(1-y)$.
These are holomorphic on $\mathbb C_{\rm cut}$, and $|\varphi_\tau|^2=|y|^{-1}|1-y|^{2\tau-2}$, so by
\eqref{gff:mgk:star}
\begin{equation}
 |y|^{-1}|1-y|^{2\tau-2}\,\mathcal I(y)=4(\mathsf u\bar{\mathsf v}+\bar{\mathsf u}\mathsf v)=8\Re(\mathsf u\bar{\mathsf v}),
 \qquad y\in\mathbb C_{\rm cut}.
 \label{gff:mgk:uv}
\end{equation}
Define the positive functions
\[\begin{aligned}
 f_-(x)&=|x|^{-1/2}(1-x)^{\tau-1}K(x),& g_-(x)&=|x|^{-1/2}(1-x)^{\tau-3/2}K\bigl(\tfrac1{1-x}\bigr)&&(x<0),\\
 f_+(x)&=x^{-1}(x-1)^{\tau-1}K(\tfrac1x),& g_+(x)&=x^{-1}(x-1)^{\tau-1}K(1-\tfrac1x)&&(x>1).
\end{aligned}\]
The principal powers have the one-sided values $(x\pm\ii0)^{-1/2}=\mp\ii|x|^{-1/2}$ for $x<0$ and
$(1-(x\pm\ii0))^{\tau-1}=-(x-1)^{\tau-1}e^{\mp\ii\pi\tau}$ for $x>1$. With (ii) and (iii) of the lemma
on $K$ the edge values are
\begin{equation}\begin{aligned}
 x<0:&\qquad \mathsf u(x\pm\ii0)=\mp\ii f_-(x),\qquad \mathsf v(x\pm\ii0)=-f_-(x)\mp\ii g_-(x),\\
 x>1:&\qquad \mathsf u(x\pm\ii0)=-e^{\mp\ii\pi\tau}\bigl(f_+(x)\pm\ii g_+(x)\bigr),\qquad
 \mathsf v(x\pm\ii0)=-e^{\mp\ii\pi\tau}g_+(x).
\end{aligned}\label{gff:mgk:edgeValues}\end{equation}
For the $\mathsf v$ entry at $x<0$: $1-y=(1-x)\mp\ii0$ with $1-x>1$, so (iii) gives
$K(1-y)=(1-x)^{-1/2}[K(\frac1{1-x})\mp\ii K(\frac x{x-1})]$, and $K(\frac x{x-1})=(1-x)^{1/2}K(x)$ by (ii).
For $x>1$, $K(1-x)=x^{-1/2}K(1-\frac1x)$ by (ii), so $\mathsf v$ jumps only through $\varphi_\tau$. The
substitutions $w=\frac x{x-1}$, $w=\frac1{1-x}$, $w=\frac1x$ and $w=\frac1x$ give (exact check S0f)
\begin{equation}
 \int_{-\infty}^0f_-=I_5,\quad \int_{-\infty}^0g_-=I_6,\quad \int_1^\infty f_+=I_3,\quad
 \int_1^\infty g_+=I_4,\quad \int_0^1\mathsf u=I_1,\quad \int_0^1\mathsf v=I_2 .
 \label{gff:mgk:edgeIntegrals}
\end{equation}
By (v) of the lemma on $K$: near $0$, $|\mathsf u|\leq C|y|^{-1/2}$ and
$|\mathsf v|\leq C|y|^{-1/2}(1+|\log|y||)$; near $1$, $|\mathsf u|\leq C|1-y|^{\tau-1}(1+|\log|1-y||)$ and
$|\mathsf v|\leq C|1-y|^{\tau-1}$; for $|y|\geq4$, $|\mathsf u|,|\mathsf v|\leq C|y|^{\tau-2}(1+\log|y|)$; and
on compact subsets of $\mathbb C\setminus\{0,1\}$ both are bounded, uniformly up to the edges. Hence
$\mathsf u\bar{\mathsf v}\in L^1(\mathbb C)$ for $0<\tau<1$, and $J_{2\tau}<\infty$ by
\eqref{gff:mgk:starReduction} and \eqref{gff:mgk:uv}. The moment relations \eqref{gff:eq:KproofMomentRelations} at
$t=\tau$ are the real and imaginary parts of Cauchy's theorem for $\mathsf u$ and $\mathsf v$ on the upper
half-plane, $\int_{\RR}\mathsf u(x+\ii0)\dd x=\int_{\RR}\mathsf v(x+\ii0)\dd x=0$, read with
\eqref{gff:mgk:edgeValues} and \eqref{gff:mgk:edgeIntegrals}; this is the contour proof given after
\eqref{gff:eq:KproofMomentRelations}.

\paragraph{Lemma (the outer integral).}
\emph{For $0<\tau<1$,}
\begin{equation}
 \int_{\mathbb C}|y|^{-1}|1-y|^{2\tau-2}\,\mathcal I(y)\,d^2y=4\bigl[I_5^2+2s_\tau I_1I_4-I_4^2\bigr].
 \label{gff:mgk:outer}
\end{equation}
\emph{Proof.} \emph{(a) Primitive.} $\mathbb C_{\rm cut}$ is star-shaped with respect to every point of
$(0,1)$: a segment from $c\in(0,1)$ to $y\in\mathbb C_{\rm cut}\setminus\RR$ meets $\RR$ only at $c$. Put
$\mathsf u_1(y)=\int_{[0,y]}\mathsf u(z)\,dz$. The integral converges since $|\mathsf u|\leq C|z|^{-1/2}$
near $0$, and it differs from $\int_{[1/2,y]}\mathsf u$ by a constant (Cauchy's theorem for triangles), so
$\mathsf u_1$ is a holomorphic primitive of $\mathsf u$ on $\mathbb C_{\rm cut}$. Along $[0,y]$,
$|1-\rho e^{\ii\theta}|\geq|1-\rho|$; with the bounds above,
$|\mathsf u(\rho e^{\ii\theta})|\leq C\rho^{-1/2}|1-\rho|^{\tau-1}(1+|\log|1-\rho||)$ for $\rho\leq3$ and
$\leq C\rho^{\tau-2}(1+\log\rho)$ for $\rho\geq3$, so $\mathsf u_1$ is bounded on $\mathbb C_{\rm cut}$.
Moreover $|\mathsf u_1(y)|\leq C|y|^{1/2}$ near $0$, and $\mathsf u_1(y)=I_1+\int_{[1,y]}\mathsf u$ (Cauchy's
theorem for the triangle $0,1,y$), so $\mathsf u_1(y)\to I_1$ as $y\to1$.

\emph{(b) Edge values of $\mathsf u_1$.} With $\hat f_-(x)=\int_x^0f_-$, $\hat f_+(x)=\int_1^xf_+$ and
$\hat g_+(x)=\int_1^xg_+$,
\[
 \mathsf u_1(x\pm\ii0)=\pm\ii\hat f_-(x)\quad(x<0),\qquad
 \mathsf u_1(x\pm\ii0)=I_1-e^{\mp\ii\pi\tau}\bigl(\hat f_+(x)\pm\ii\hat g_+(x)\bigr)\quad(x>1).
\]

\emph{(c) Stokes.} On $\mathbb C_{\rm cut}$,
$d(\mathsf u_1\bar{\mathsf v}\,d\bar y)=\mathsf u\bar{\mathsf v}\,dy\wedge d\bar y=-2\ii\,\mathsf u\bar{\mathsf v}\,d^2y$.
Apply Green's theorem on slit domains to $\{y\in\mathbb C_{\rm cut}:\ r<|y|<R,\ |y-1|>r\}$. The upper
edges of both slits are traversed in the $+x$ direction and the lower edges in the $-x$ direction. With
$[\mathsf g]^+_-=\mathsf g(x+\ii0)-\mathsf g(x-\ii0)$,
\[
 -2\ii\int\mathsf u\bar{\mathsf v}\,d^2y=\int_{-R}^{-r}[\mathsf u_1\bar{\mathsf v}]^+_-\dd x
 +\int_{1+r}^{R}[\mathsf u_1\bar{\mathsf v}]^+_-\dd x+(\text{circle terms}).
\]

\emph{(d) Circle terms.} On $|y|=r$: $|\mathsf u_1|\leq Cr^{1/2}$ and $|\mathsf v|\leq Cr^{-1/2}(1+|\log r|)$,
so the term is $O(r(1+|\log r|))$. On $|y-1|=r$: $|\mathsf u_1|\leq C$ and $|\mathsf v|\leq Cr^{\tau-1}$,
so the term is $O(r^\tau)$. On $|y|=R$: $|\mathsf u_1|\leq C$ and $|\mathsf v|\leq CR^{\tau-2}(1+\log R)$, so
the term is $O(R^{\tau-1}(1+\log R))$. All three vanish in the limit, because $0<\tau<1$.

\emph{(e) Jumps.} From \eqref{gff:mgk:edgeValues} and (b),
\[
 \begin{aligned}
 x<0:&\quad [\mathsf u_1\bar{\mathsf v}]^+_-=(\ii\hat f_-)(-f_-+\ii g_-)-(-\ii\hat f_-)(-f_--\ii g_-)=-2\ii\,\hat f_-f_-,\\
 x>1:&\quad [\mathsf u_1\bar{\mathsf v}]^+_-=-2\ii s_\tau I_1\,g_++2\ii\,\hat g_+g_+ .
 \end{aligned}
\]

\emph{(f) Limit.} Since $\mathsf u\bar{\mathsf v}\in L^1$, the area integral converges. On $(-\infty,0)$,
$\hat f_-'=-f_-$, $\hat f_-(0)=0$ and $\hat f_-(-\infty)=I_5$, so $\int_{-\infty}^0(-2\ii\hat f_-f_-)\dd x=-\ii I_5^2$.
On $(1,\infty)$, $\hat g_+'=g_+$, $\hat g_+(1)=0$ and $\hat g_+(\infty)=I_4$, so
$\int_1^\infty(-2\ii s_\tau I_1g_++2\ii\hat g_+g_+)\dd x=-2\ii s_\tau I_1I_4+\ii I_4^2$. Hence
\[
 \int_{\mathbb C_{\rm cut}}\mathsf u\bar{\mathsf v}\,d^2y=\frac\ii2\bigl(-\ii I_5^2-2\ii s_\tau I_1I_4+\ii I_4^2\bigr)
 =\frac12\bigl(I_5^2+2s_\tau I_1I_4-I_4^2\bigr).
\]

\emph{(g)} By \eqref{gff:mgk:uv}, the left side of \eqref{gff:mgk:outer} is
$8\Re\int\mathsf u\bar{\mathsf v}\,d^2y=4(I_5^2+2s_\tau I_1I_4-I_4^2)$; the complement of
$\mathbb C_{\rm cut}$ has measure zero. \hfill$\square$

\emph{Second evaluation.} Take instead the primitive $\mathsf v_1(y)=\int_{[1,y]}\mathsf v$ of $\mathsf v$,
based at $1$. On $x<0$ it has edge values $\mathsf v_1(x\pm\ii0)=-I_2+\hat f_-(x)\pm\ii\hat g_-(x)$, with
$\hat g_-(x)=\int_x^0g_-$, and the same steps give
$\int\mathsf v\bar{\mathsf u}\,d^2y=\frac12(I_4^2+2I_2I_5-I_5^2)$ (exact check S0c). Since
$\int\mathsf v\bar{\mathsf u}=\overline{\int\mathsf u\bar{\mathsf v}}$ and both values are real, the two
evaluations must agree; given \eqref{gff:eq:KproofMomentRelations} they do, the difference of the two
brackets being $2I_4^2(1-s_\tau^2-c_\tau^2)=0$. Their average,
$\int\mathsf u\bar{\mathsf v}\,d^2y=\frac12(s_\tau I_1I_4+I_2I_5)$, is an identity of the bracket polynomials
that needs no relation among the $I_j$; it is the bilinear form of \eqref{gff:done:kiteFactorization} at
this point, the last form of \eqref{gff:gc:K2master}.

\paragraph{Assembly.}
By \eqref{gff:mgk:starReduction}, \eqref{gff:mgk:outer}, \eqref{gff:eq:KproofMomentRelations} and the values
$I_3=\pi^2\mathfrak F(1-\tau)/(2s_\tau)$, $I_4=\pi^2\mathfrak F(\tau)/(2s_\tau)$ of
Section~\ref{gff:gc:K2proof},
\[
 J_{2\tau}=\frac4{\pi^2}\bigl(I_5^2+2s_\tau I_1I_4-I_4^2\bigr)=\frac4{\pi^2}s_\tau^2\bigl(I_3^2+I_4^2\bigr)
 =\pi^2\bigl[\mathfrak F(1-\tau)^2+\mathfrak F(\tau)^2\bigr]<\infty ,
\]
which is \eqref{gff:gc:K2master}; the other forms there follow from $I_2=s_\tau I_3$, $I_6=s_\tau I_4$ and the
average above. All of this algebra is checked exactly in
\path{LRI_C5_d2_proof/K2_d1template_fix1/exact_algebra_v2.wl} (parts S0a--S0f). This proof uses neither
the cell table \eqref{gff:gc:K2cells}, nor the rotation lemma \eqref{gff:gc:K2rotation}, nor
\eqref{gff:done:kiteFactorization}.

\paragraph{Remarks.}
(1) \emph{Power counting.} Convergence already follows from the bounds above; for completeness, the
integrand of $J_{2\tau}$ in $\RR^4$ has the following strict margins: $x\to0$, $x\to1$, $x\to y$ and
$y\to0$, margin $1$; $y\to1$, margin $2\tau$; $x,y\to0$ jointly, $4-3=1$; $x,y\to1$ jointly,
$4-(4-2\tau)=2\tau$; $x\to\infty$ alone, $2-3=-1$; $y\to\infty$ alone, $2-(4-2\tau)=2\tau-2$; $x,y\to\infty$
jointly, $4-(6-2\tau)=2\tau-2$; $x,y\to\infty$ with $x-y$ bounded, total $2\tau-3$, with $|x-y|^{-1}$
integrable. Mixed configurations factorize into these cases. For $0<\tau<1$ every margin has the
required sign, so the point $(1,1,1,2-2\tau,1)$ lies in the interior of the convergence domain in all
five exponents. (2) \emph{Relation to $d=1$.} At $d=1$, integrating $x$ over $\RR$ gives the full-line
elliptic star, which is linear in $K$, and folding gives the linear functional \eqref{gff:eq:Jfold}. At
$d=2$, integrating $x$ over $\mathbb C$ gives the bilinear star \eqref{gff:mgk:star}, and the $y$ integration
becomes the bilinear Stokes pairing \eqref{gff:mgk:outer}. The six Euler moments, Euler integration and the
moment relations are common to both dimensions. The higher Taylor coefficients enter the extraction of
Section~\ref{gff:gc:K2laurent} as $2(\varrho+\pi^2v)$, against $4(\varrho+\pi^2v)$ in
\eqref{gff:eq:KproofCoefficientAssembly} at $d=1$.

\subsubsection{Status and checks}\label{gff:mgk:status}
\begin{center}\small
\begin{tabular}{@{}p{0.47\textwidth}p{0.12\textwidth}p{0.35\textwidth}@{}}\toprule
Statement & Status & Inputs\\\midrule
Majorant \eqref{gff:mgk:majorant}, regularized form \eqref{gff:mgk:Kreg}, $[\tau^2]$ under the trace; sharper bound \eqref{gff:mgk:sharp}
 & proved & Stirling \cite{gff:DLMFStirling}; the expansion after \eqref{gff:gs:Ksubtraction}, checked exactly\\
Per-spin residues \eqref{gff:mgk:perSpin}, factorized traces \eqref{gff:mgk:factorized} & proved & reflection, Stirling\\
Single sums \eqref{gff:mgk:periodTwo}, \eqref{gff:mgk:periodOne}, \eqref{gff:mgk:singleValues}, \eqref{gff:mgk:Sh} (three evaluations)
 & proved & Gauss \cite{gff:DLMFGaussUnity}, \eqref{gff:eq:KproofThomae}, \eqref{gff:eq:KproofMomentRelations}, Euler integral, \cite{gff:DLMF15}\\
Regulated traces \eqref{gff:mgk:regulatedEq} & proved & the rows above\\
Mixed-trace step \eqref{gff:mgk:mixedTrace} at $d=2$, $0<\tau<\frac12$, full-line reading & proved & \eqref{gff:mgk:regulatedEq}, \eqref{gff:gc:K2master}\\
Right side of \eqref{gff:gs:Kspectral} at $d=2$ equals $32\Cat^2-28\zeta(3)\log2$; $\langle t_\ell^2E_2\rangle_2=-21\zeta(3)$
 & proved & \eqref{gff:eq:KproofLowerCoefficients}, \eqref{gff:eq:KproofTaylorCancellation}\\
Per-spin identity \eqref{gff:mgk:innerSpin}; \eqref{gff:gs:Kseries} at $d=2$ equals the same; strip limit
 & proved & \eqref{gff:mgk:seriesResidues} (exact algebra), the factorization of \path{factor_K2.py} for part (c)\\
Inner star \eqref{gff:mgk:star}, outer integral \eqref{gff:mgk:outer}, second proof of \eqref{gff:gc:K2master} & proved & \eqref{gff:eq:KproofMomentRelations}, Euler integral, \cite{gff:DLMF15}\\
(M) for complex $\tau$; (M) and $[\tau^2]$ under the trace at $d\neq2$ & not addressed & input (M) of Section~\ref{gff:gdk:conventions}\\\bottomrule
\end{tabular}
\end{center}

\paragraph{Checks.}
None of the steps above uses numerics; the following are checks only, each with its record in the
package named.
\begin{itemize}
\item \emph{Periodization and traces} (\path{check_periodization.py} in
\path{LRI_C5_d2_simplify/proof/K2-parseval/}):
at $\tau=0.13$, $0.3$ and $0.41$ both identities \eqref{gff:mgk:periodTwo}, the first two of
\eqref{gff:mgk:periodOne}, the Euler--Thomae step \eqref{gff:mgk:singleValues} and the first identity of
\eqref{gff:mgk:regulatedEq} hold to $10^{-34}$ or better; for example
$\langle t_\ell^2R_{0.3}\rangle_2=20.73663823644757695029316\ldots$ from both sides.
\item \emph{Residues} (\path{check_perspin.py} in the same folder): \eqref{gff:mgk:perSpin} at $\tau=0.3$,
$\ell=0,3$, Barnes integral against residue sum, to $6\times10^{-19}$; \eqref{gff:mgk:innerSpin} at
$\ell=0,1$ to $10^{-21}$; the factorized summand of \path{factor_K2.py} against \eqref{gff:gs:Kpoly} at four
lattice points to $4\times10^{-20}$. The residues of \eqref{gff:mgk:perSpin} by Cauchy circles agree with
the formulas to $6\times10^{-54}$, and the identity $\Theta^{\rm R}_\ell=G_\ell T_\tau$ holds at a test
point to $10^{-78}$ (\path{LRI_C5_d2_proof/review_K2_d1template_numerics_r0/logs/a01_9_20.568115.log}).
\item \emph{Single sums}: the third evaluation, and \eqref{gff:mgk:factorized} computed from direct sums,
agree with the closed forms and with $\pi^2[\mathfrak F(\tau)^2+\mathfrak F(1-\tau)^2]$ to
$3.3\times10^{-67}$ at $\tau=\frac17,\frac3{10},\frac9{20}$
(\path{LRI_C5_d2_proof/review_K2_d1template_numerics_r0/logs/a4b_*.log}). In the second evaluation,
Gauss's sum, the Thomae step and the undifferentiated single sums hold to about $10^{-50}$, the
derivative identities to $10^{-42}$ (limited by finite differences), and the formula of part (c) for
$P_{\mathsf h}$ at six points $(\tau,u)$ to about $24$ digits
(\path{LRI_C5_d2_proof/K2_gauss/logs/check_singles.log}, \path{check_dLambda.log},
\path{check_bilateral.log}).
\item \emph{Algebra}: the expansion after \eqref{gff:gs:Ksubtraction}, both Laurent extractions and
\eqref{gff:mgk:seriesResidues} are exact (\path{K2-parseval/symbolic_checks.out},
\path{synth/synth_checks_K2.log}, \path{K2_gauss/logs/check_expansion.log}, \path{check_extraction.log},
\path{check_residue_series.log}); the Taylor inputs \eqref{gff:mgk:Taylor} hold to $10^{-23}$
(\path{check_taylor_inputs.out}).
\item \emph{Spectral moment}: a direct spectral evaluation of $\langle t_\ell^2E_2\rangle_2$ on symmetry
(spins $\ell<4096$, Richardson extrapolation in $\ell$; job 567967, \path{K2-parseval/e2check-567967.out})
gives $-25.243194966351479993$, equal to $-21\zeta(3)$ within $7.3\times10^{-24}$. A direct spectral
integration of $C(\tau)\langle t_\ell^2(R_\tau-S_\tau)\rangle_2$ at $\tau=\frac1{10},\frac15$ agrees with the
right side of \eqref{gff:mgk:CD} to about $14$ digits (\path{K2_gauss/logs/check_endtoend.log}).
\item \emph{Extraction without the cancellation rule}: a $128$-point Cauchy integral on $|\tau|=\frac18$,
at $90$ digits, of the true ${}_3F_2$ functions in \eqref{gff:mgk:CD}, which does not use
\eqref{gff:eq:KproofTaylorCancellation}, gives $-\frac{AU}4+\frac14[\tau^2]C\mathcal D$ equal to
$32\Cat^2-28\zeta(3)\log2$ to $88$ digits, $[\tau^1]C\mathcal D=-21\zeta(3)$ to $20$ digits, and agreement
with the $70$-digit series value to $2\times10^{-65}$, the truncation of that value
(\path{K2_gauss/logs/check_K2_cauchy.log}).
\item \emph{Position space}: the edge values \eqref{gff:mgk:edgeValues} hold to $7\times10^{-30}$, and the
Stokes expressions and moment relations to $2\times10^{-28}$, at $\tau=0.3,0.45,0.6$
(\path{LRI_C5_d2_proof/K2_d1template_fix1/check_boundary_values_mp.out}); \eqref{gff:mgk:star} agrees with an
independent Feynman-parameter representation to $56$--$58$ digits at fourteen points
(\path{check_lemma1_star.out}); a fresh composite run at $\tau=\frac{37}{100}$, inserting \eqref{gff:mgk:star}
into \eqref{gff:mgk:starReduction} and integrating over $y$, gives
$J_{0.74}=41.7880979669900734758237249270162127\ldots$, equal to
$\pi^2[\mathfrak F(\tau)^2+\mathfrak F(1-\tau)^2]$ to $38$ digits (\path{composite_tau037.out}).
\item \emph{The tetrahedron point}: at $\tau=\frac12$, \eqref{gff:gc:K2master} gives
$J_1=2\pi^2\mathfrak F(\frac12)^2=3Q_2$ (Dixon's sum, checked to $578$ digits), and its second derivative,
$\frac{\pi^2}2[\mathfrak F'(\frac12)^2+\mathfrak F(\frac12)\mathfrak F''(\frac12)]$, equals $Z_2=4\Cat Q_2$ of
\eqref{gff:gs:Z2exact} to $77$ digits (\path{review_K2_d1template_numerics_r0/logs/c6b.568116.log}). By
\eqref{gff:gs:Tjet}, a change of the single exponent $f$ by $\delta$ has squared completion variation
$4\delta^2$, so this is a check of \eqref{gff:gc:K2master} against the complex Selberg evaluation of $Z_2$.
\end{itemize}


\subsection{Scope of the primitive reduction}
The finite sample formulas use at most two-variable hypergeometric functions
for generic dimension; at $d=2$ the factorizations of Section~\ref{gff:fp:twoDimensional}
use the same scalar values as at $d=1$. For generic dimension the nonprimitive
functions are the explicit finite multivariable combinations of
Section~\ref{gff:app:generalClosed}; Sections~\ref{gff:gdr:section} and \ref{gff:gdk:section}
extend the periodization of this section to a general rim exponent and to the
family $R_\tau$ of \eqref{gff:gs:RS}, and reduce them to two-variable samples and
their low-order parameter jets. At $d=2$ all three are known in closed form:
$Z_2=4\Cat Q_2$ by \eqref{gff:gs:Z2exact}, and $\mathcal L_2$, $\mathcal K_2$ by
\eqref{gff:gc:L2closed} and \eqref{gff:gc:K2closed} (Sections~\ref{gff:gc:L2proof} and
\ref{gff:gc:K2proof}).
Their physical occurrence can be collected as
\begin{equation}
 \mathcal N_d=\mathcal K_d+\frac38\mathcal L_d-\frac{149}{32}Z_d.
 \label{gff:fp:remainingCombination}
\end{equation}
This is a change of basis, not an additional evaluation. In this combination \eqref{gff:gs:C5} reads
\begin{equation}\boxed{\begin{aligned}
46656C_5(d)={}&-1536\bigl(P_5(d)+\mathcal N_d\bigr)+384AP_4(d)
 -1728Q_d^2+720A^2Q_d+1008Q_dR\\
 &-15A^4-90A^2R+414AU-108AV-117R^2-75X,
\end{aligned}}\label{gff:fp:C5target}\end{equation}
since $1536\cdot\frac38=576$ and $1536\cdot\frac{149}{32}=7152$. Thus $\mathcal K_d$, $\mathcal L_d$
and $Z_d$ enter $C_5(d)$ only through $\mathcal N_d$, which appears together with $P_5(d)$. This is a
remark on the structure of \eqref{gff:gs:C5}; Sections~\ref{gff:gdr:section} and~\ref{gff:gdk:section} reduce the
three constituents separately. The finite Fourier
spectra used in Sections~\ref{gff:gdr:section} and \ref{gff:gdk:section} belong to the
one-parameter families through which the masters are defined, the rim exponent
in \eqref{gff:gs:wheelJets} and $\tau$ in \eqref{gff:gs:RS}; for $\mathcal K_d$ this
route also rests on the inputs of Section~\ref{gff:gdk:conventions}.

\section{Three dimensions: half-shift determinants and rank-two functions}\label{gff:app:threeDimensionalReduction}

The half-shift reduction and the rank-two identification of this section, with their checks, were
written by OpenAI GPT-6 astra, in the packages \path{lri_d3_reduction_update/} and
\path{lri_d3_matrix_update/} (Section~\ref{gff:cn:reproduction}).
The integer dimension $d=3$ admits a further reduction of the primitive periods $Q_3$, $P_4(3)$ and
$P_5(3)$, but not the square factorization of $d=2$ (Section~\ref{gff:fp:twoDimensional}). Its
multiplicity $d_\ell=2\ell+1$ produces an antisymmetric determinant on two lattices offset by one
half (Section~\ref{gff:d3:determinantProof}). This eliminates one Fourier-parity sector of the samples
\eqref{gff:fp:finiteFourier} in favor of ordinary hypergeometric values. The complementary sector is not
fixed by that identity. The formulas below evaluate it by a small, explicitly specified set of
two-variable Horn values, which are identified, from Section~\ref{gff:app:d3RankTwo} on, with real
rank-two matrix-argument hypergeometric functions. The formulas contain no parameter
derivatives, Barnes integrals or external infinite sums. They do not prove that the remaining
two-variable functions are finite combinations of ordinary ${}_pF_q$ values.

The tetrahedron takes the form
\begin{equation}\boxed{
 Q_3=4\pi^2\mathfrak H_3(0)
       -\frac{\Gamma(\tfrac14)^{10}}{48\pi^{7/2}}.
}\label{gff:td:Qanswer}\end{equation}
Here $Q_3$ means $Q_{d=3}$, not the one-dimensional three-loop coefficient.
The fixed \emph{two-variable} function $\mathfrak H_3(0)$ is fully specified
in \eqref{gff:td:Hbase} below.  In particular \eqref{gff:td:Qanswer} is not a
Gamma-only evaluation of the tetrahedron.

For the real formulas with positive rational parameters proved below, the numbers of distinct
function evaluations are bounded as in the following table. For comparison, its second column
counts the values of functions of at most two variables in the generic-dimensional samples
\eqref{gff:fp:wheelSample} and \eqref{gff:fp:primitiveSample} at the nodes \eqref{gff:fp:nodes}: five blocks
at each of the $r$ nodes for $r=3,4$, and eight blocks at each of the seven nodes for $r=7$,
one-variable blocks included and the real part of a complex value counted once.
\begin{center}\small
\begin{tabular}{@{}lrrrr@{}}\toprule
Period & Generic-$d$ samples, at most & Two-variable values & Ordinary values & Total\\\midrule
$Q_3$ & 15 & 1 & 0 & 1\\
$P_4(3)$ & 20 & 2 & 3 & 5\\
$P_5(3)$ & 56 & 5 & 18 & 23\\\bottomrule
\end{tabular}
\end{center}
Gamma factors and finite arithmetic are not counted. The counts are upper bounds, not independence
or minimality statements. The fixed complex nodes of Section~\ref{gff:td:complexVariant} reduce the
two-variable counts further, to one for $P_4(3)$ and two for $P_5(3)$, at the price of fixed complex
parameters.

\subsection{Ordinary hypergeometric moments}
Write $h(x)=\Gamma(x+\tfrac18)/\Gamma(x+\tfrac78)$.  For
$\omega=\pm1$ and nonnegative integer $m$ define the ordinary function
\begin{equation}\begin{aligned}
 \mathfrak U_{r,m}^{(\omega)}(u)={}&
 u^m h(u)^r\,
 {}_{r+m+1}F_{r+m}\!\left(
 \begin{matrix}1,(u+\tfrac18)^{\times r},(u+1)^{\times m}\\
 (u+\tfrac78)^{\times r},u^{\times m}\end{matrix};\omega\right).
\end{aligned}\label{gff:td:Udefinition}\end{equation}
Repeated parameters with multiplicity zero are absent.  Identical fixed
upper and lower parameters may be canceled.  Its series is
$\sum_{n\ge0}\omega^n(n+u)^m h(n+u)^r$.  The required pairs are
$(r,m)=(3,0),(3,1),(4,0),(4,1),(7,1),(7,3)$; all satisfy
$3r/4-m>1$, so the displayed ordinary series converge absolutely at both
arguments $\omega=1$ and $-1$ \cite{gff:DLMF16}.

Use the following finite combination, also entirely in ordinary functions:
\begin{equation}\boxed{\begin{aligned}
 \Psi_{r,m}^{(\omega)}(u)={}&
 \sin^r\!\bigl(\pi(u+\tfrac18)\bigr)
       \mathfrak U_{r,m}^{(\omega)}(u)\\
 &+\omega(-1)^m\sin^r\!\bigl(\pi(\tfrac18-u)\bigr)
       \mathfrak U_{r,m}^{(\omega)}(1-u).
\end{aligned}}\label{gff:td:Psidefinition}\end{equation}
In all real formulas below $0<u<1$; all upper and lower parameters of the
uncancelled ordinary functions are positive rational numbers.  At
$u=1/8$ the second term vanishes exactly, so its function need not be
computed.  The symbol $\omega$ is a hypergeometric argument, not the
physical expansion parameter $\epsilon$.

The finite-mode identities needed here are
\begin{equation}\boxed{\begin{aligned}
 \Psi_{3,0}^{(+)}(u)&=\mathfrak a_3\cos\pi u,&
 \Psi_{3,1}^{(+)}(u)&=\mathfrak b_3\sin\pi u,\\
 \Psi_{4,0}^{(+)}(u)&=\mathfrak a_4+\mathfrak c_4\cos2\pi u,&
 \Psi_{4,1}^{(+)}(u)&=\mathfrak b_4\sin2\pi u,\\
 \Psi_{7,m}^{(-)}(u)&=\sum_{j=1}^3 h_{m,j}\sin(2\pi j u),
 &&m=1,3.
\end{aligned}}\label{gff:td:momentModes}\end{equation}
The coefficient $\mathfrak c_4$ will cancel from every answer.  The useful
amplitudes have the fixed-value definitions
\begin{equation}\begin{aligned}
 \mathfrak a_3&=\frac{\Psi_{3,0}^{(+)}(1/8)}{\cos(\pi/8)},&
 \mathfrak b_3&=\frac{\Psi_{3,1}^{(+)}(1/8)}{\sin(\pi/8)},\\
 \mathfrak a_4&=\Psi_{4,0}^{(+)}(1/4),&
 \mathfrak b_4&=\frac{\Psi_{4,1}^{(+)}(1/8)}{\sin(\pi/4)}.
\end{aligned}\label{gff:td:wheelAmplitudes}\end{equation}
For $m=1,3$ let $Y_{m,-}=\Psi_{7,m}^{(-)}(1/8)$,
$Y_{m,0}=\Psi_{7,m}^{(-)}(1/4)$ and
$Y_{m,+}=\Psi_{7,m}^{(-)}(3/8)$. The three coefficients $h_{m,j}$, unrelated to
the function $h(x)$ above, are explicitly
\begin{equation}\begin{aligned}
 h_{m,1}&=\frac{Y_{m,0}}2+\frac{Y_{m,-}+Y_{m,+}}{2\sqrt2},\\
 h_{m,2}&=\frac{Y_{m,-}-Y_{m,+}}2,\\
 h_{m,3}&=-\frac{Y_{m,0}}2+\frac{Y_{m,-}+Y_{m,+}}{2\sqrt2}.
\end{aligned}\label{gff:td:hcoefficients}\end{equation}
This is a finite sine transform on a proved three-mode space, not
numerical interpolation of an arbitrary function.

\subsubsection{Two Gamma evaluations within the ordinary sector}
Dixon's sum and the Rogers--Dougall sum \cite[eqs.~16.4.4 and
16.4.9]{gff:DLMF16} give
\begin{equation}\boxed{
 \mathfrak a_3\mathfrak b_3
   =\frac{\Gamma(\tfrac14)^{10}}{16\sqrt2\,\pi^4},
 \qquad
 \mathfrak b_4=\frac{\Gamma(\tfrac14)^6}{16\pi^{5/2}}.
}\label{gff:td:gammaAmplitudes}\end{equation}
For clarity the particular sums used in this step are
\begin{equation}\begin{aligned}
 {}_3F_2\!\left(\begin{matrix}(\tfrac14)^{\times3}\\1,1\end{matrix};1\right)
 &=\frac{\Gamma(\tfrac98)\Gamma(\tfrac58)}
 {\Gamma(\tfrac54)\Gamma(\tfrac78)^2\Gamma(\tfrac34)},\\
 {}_4F_3\!\left(\begin{matrix}\tfrac98,(\tfrac14)^{\times3}\\
 \tfrac18,1,1\end{matrix};1\right)
 &=\frac{\Gamma(\tfrac58)\Gamma(\tfrac18)}
 {\Gamma(\tfrac54)\Gamma(\tfrac34)\Gamma(\tfrac38)^2},\\
 {}_5F_4\!\left(\begin{matrix}\tfrac98,(\tfrac14)^{\times4}\\
 \tfrac18,1,1,1\end{matrix};1\right)
 &=\frac{\sqrt\pi}{\Gamma(\tfrac54)\Gamma(\tfrac34)^3}.
\end{aligned}\label{gff:td:DixonDougall}\end{equation}
For the second identity insert a canceled pair $5/8$ into the
Rogers--Dougall ${}_5F_4$, with parameters $(a,b,c,d)=(1/4,1/4,1/4,5/8)$ in
the notation of \cite[eq.~16.4.9]{gff:DLMF16}; this $d$ is not the dimension.
For the third take $(a,b,c,d)=(1/4,1/4,1/4,1/4)$.  Their convergence conditions hold.
Substitution into \eqref{gff:td:wheelAmplitudes}, followed by reflection and
duplication \cite{gff:DLMF5}, proves \eqref{gff:td:gammaAmplitudes}.  Thus the
Gamma powers in \eqref{gff:td:Qanswer} are analytic evaluations, not PSLQ fits.

\subsection{The remaining two-variable functions and complete sample formulas}
Use the Horn convention \eqref{gff:fp:HornDefinition}.  Put
$\boldsymbol e=(1,0)$, $\boldsymbol f=(0,1)$,
$\boldsymbol c=(1,1)$ and $\boldsymbol v=(1,2)$, and set
\begin{equation}\begin{aligned}
 \mathsf U_r(t)=\{&(1,\boldsymbol e),(1,\boldsymbol f),
 (\tfrac32,\boldsymbol e),
 (\tfrac34+t,\boldsymbol c)^{\times r},
 (\tfrac14+t,\boldsymbol f)^{\times r}\},\\
 \mathsf V_r(t)=\{&(\tfrac12,\boldsymbol e),
 (\tfrac32+t,\boldsymbol c)^{\times r},
 (1+t,\boldsymbol f)^{\times r}\}.
\end{aligned}\label{gff:td:Harrays}\end{equation}
For $r=3,4$ define the fixed two-variable function
\begin{equation}\boxed{
 \mathfrak H_r(t)=\mathfrak H[
       \mathsf U_r(t);\mathsf V_r(t);(1,1)].
}\label{gff:td:Hbase}\end{equation}
For $r=7$ adjoin three upper/lower row pairs:
\begin{equation}\begin{aligned}
 x_t&=t+\tfrac58,& y_t&=t+\tfrac18,& v_t&=2t+\tfrac34,\\
 \mathfrak H_7(t)=\mathfrak H\Big[&
 \mathsf U_7(t)\cup\{(1+x_t,\boldsymbol c),
       (1+y_t,\boldsymbol f),(1+v_t,\boldsymbol v)\};\\
 &\mathsf V_7(t)\cup\{(x_t,\boldsymbol c),
       (y_t,\boldsymbol f),(v_t,\boldsymbol v)\};(-1,1)\Big].
\end{aligned}\label{gff:td:Hseven}\end{equation}
Every row is printed here; there is no unevaluated coefficient-extraction
instruction.  The two upper rows with parameter $1$ cancel the two
factorials of the Horn convention.  The ratio
$(3/2)_\ell/(1/2)_\ell=2\ell+1$ supplies the three-dimensional angular
multiplicity.  The three additional pairs supply $x$, $y$ and $x+y$.
All real nodes used below have positive rational rows.

For compact scalar prefactors write
\begin{equation}\begin{aligned}
 G_r(t)&=\left[
 \frac{\Gamma(\tfrac34+t)\Gamma(\tfrac14+t)}
      {\Gamma(\tfrac32+t)\Gamma(1+t)}\right]^r,\\
 V_r(t)&=G_r(t)\mathfrak H_r(t),&&r=3,4,\\
 V_7(t)&=4x_ty_tv_tG_7(t)\mathfrak H_7(t).
\end{aligned}\label{gff:td:Vfinite}\end{equation}
The symbol $V_r(t)$ here is local to this section; it is not the
polygamma constant $V$ in the general-dimensional coefficient formula.

Let $\mathsf X_r(t)=\Xi_r(t)|_{q=3/4}$, with $\Xi_r$ as in
\eqref{gff:fp:periodization}, and define
\begin{equation}\begin{aligned}
 t^\dagger&=\tfrac14-t,& f(t)&=-\tfrac12\cos2\pi t,\\
 D(t)&=\sin\pi t\sin\pi(t+\tfrac34),&
 T(t)&=\tan\pi(\tfrac34+2t),& z_t&=\tfrac78-t.
\end{aligned}\label{gff:td:sampleScalars}\end{equation}
A sample in a wheel is now just
\begin{equation}\boxed{\begin{aligned}
 \mathsf X_r(t)={}&f(t)^rV_r(t)+f(t^\dagger)^rV_r(t^\dagger)\\
 &+2D(t)^r\bigl[
 \mathfrak U_{r,1}^{(+)}(x_t)\mathfrak U_{r,0}^{(+)}(z_t)
 +\mathfrak U_{r,0}^{(+)}(x_t)\mathfrak U_{r,1}^{(+)}(z_t)\bigr],
 \qquad r=3,4.
\end{aligned}}\label{gff:td:wheelSample}\end{equation}
The corresponding five-loop sample is
\begin{equation}\boxed{\begin{aligned}
 \mathsf X_7(t)=T(t)\Big\{&f(t)^7V_7(t)-f(t^\dagger)^7V_7(t^\dagger)\\
 &+8D(t)^7\bigl[
 \mathfrak U_{7,3}^{(-)}(x_t)\mathfrak U_{7,1}^{(-)}(z_t)
 -\mathfrak U_{7,1}^{(-)}(x_t)\mathfrak U_{7,3}^{(-)}(z_t)\bigr]\Big\}.
\end{aligned}}\label{gff:td:sevenSample}\end{equation}
Only fixed ordinary and two-variable functions appear.  In particular
the middle triangular part of the $(\ell,n)$ lattice of
\eqref{gff:fp:periodization} factorizes into ordinary products.  At $t=0$ the left and middle pieces vanish,
giving the especially simple nonsingular formulas
\begin{equation}\boxed{\begin{aligned}
 \mathsf X_3(0)&=-2\sqrt2\,\pi^{3/2}\mathfrak H_3(0),\\
 \mathsf X_4(0)&=4\pi^2\mathfrak H_4(0),\\
 \mathsf X_7(0)&=\frac{15\sqrt2}{8}\pi^{7/2}\mathfrak H_7(0).
\end{aligned}}\label{gff:td:zeroSamples}\end{equation}
These are values of already regularized functions, not singular
parameter limits to be taken by the reader.

\subsection{Primitive formulas with only real rational parameters}
Let $c_8=\cos(\pi/8)$ and $s_8=\sin(\pi/8)$, as in \eqref{gff:fp:FpmDictionary}.
For the four-loop answer use
\begin{equation}\begin{aligned}
 X_0&=4\pi^2\mathfrak H_4(0),\\
 X_h&=\frac1{32}\left[
 \frac{\Gamma(\tfrac78)\Gamma(\tfrac38)}
      {\Gamma(\tfrac{13}8)\Gamma(\tfrac98)}\right]^4
       \mathfrak H_4(\tfrac18)
       +4s_8^8\mathfrak U_{4,0}^{(+)}(\tfrac34)
                \mathfrak U_{4,1}^{(+)}(\tfrac34),\\
 \mathfrak a_4&=c_8^4\mathfrak U_{4,0}^{(+)}(\tfrac14)
               +s_8^4\mathfrak U_{4,0}^{(+)}(\tfrac34),&
 \mathfrak b_4&=\frac{\Gamma(\tfrac14)^6}{16\pi^{5/2}}.
\end{aligned}\label{gff:td:P4data}\end{equation}
Here $X_0=\mathsf X_4(0)$ and $X_h=\mathsf X_4(1/8)$; they are not the
one-dimensional $X_0,X_+,X_-$ of Section~\ref{gff:app:generalPeriodization}.
The answer is
\begin{equation}\boxed{
 P_4(3)=2\pi X_0-10\pi X_h
            +(14\sqrt2-20)\pi\mathfrak a_4\mathfrak b_4.
}\label{gff:td:P4answer}\end{equation}
This uses two $\mathfrak H_4$ values, two ordinary ${}_5F_4(1)$ values,
and one ordinary ${}_6F_5(1)$ value, plus Gamma factors.  No claim that
the two $\mathfrak H_4$ values have become ordinary functions is intended.

For $P_5(3)$ form the two ordinary-product constants
\begin{equation}\begin{aligned}
 \alpha&=h_{3,2}h_{1,1}-h_{3,1}h_{1,2},\\
 \beta&=h_{3,2}h_{1,3}-h_{3,3}h_{1,2},\\
 \theta(t)&=2\pi(t+\tfrac38),\qquad w(t)=\cos\theta(t),\\
 \mathcal O(t)&=16w(t)\bigl(1-w(t)^2\bigr)
             \bigl[-\alpha-3\beta+4\beta w(t)^2\bigr].
\end{aligned}\label{gff:td:oddSeven}\end{equation}
All $h_{m,j}$ were given by the finite ordinary evaluations
\eqref{gff:td:hcoefficients}.  Use three regularized samples, with $u_j=\cos(4\pi(t_j+3/8))$ the variable of
\eqref{gff:td:evenPolynomial} below:
\begin{equation}\begin{aligned}
 t_0&=0,&t_+&=\tfrac1{16},&t_-&=-\tfrac1{16},\\
 u_0&=0,&u_+&=\tfrac1{\sqrt2},&u_-&=-\tfrac1{\sqrt2},\\
 E_j&=\frac{\mathsf X_7(t_j)-\mathcal O(t_j)}{1-u_j},
 &&j=0,+,-.
\end{aligned}\label{gff:td:sevenE}\end{equation}
Equations \eqref{gff:td:sevenSample} and \eqref{gff:td:zeroSamples} specify each
sample without derivatives.  The result is
\begin{equation}\boxed{\begin{aligned}
 P_5(3)={}&4\sqrt2\,\pi^{3/2}
 \left[-35E_0+
 \left(22+\frac{13}{\sqrt2}\right)E_+
 +\left(22-\frac{13}{\sqrt2}\right)E_-\right]\\
 &-448\pi^{3/2}(2\beta-\alpha).
\end{aligned}}\label{gff:td:P5answer}\end{equation}
The only two-variable values are
\begin{equation}
 \mathfrak H_7(0),\quad\mathfrak H_7(\tfrac1{16}),\quad
 \mathfrak H_7(-\tfrac1{16}),\quad
 \mathfrak H_7(\tfrac3{16}),\quad
 \mathfrak H_7(\tfrac5{16}).
 \label{gff:td:fiveHornNodes}
\end{equation}
All their uncancelled parameters are positive rational, and their
arguments are $(-1,1)$.  The ten ordinary values for $h_{1,j},h_{3,j}$
use $u=1/8,1/4,3/4,3/8,5/8$ and $m=1,3$.  Eight further ordinary values
in the two middle determinants use $u=9/16,11/16,13/16,15/16$ and $m=1,3$.
In total $P_5(3)$ therefore uses eighteen ordinary values.
The excess of every one of these ordinary series is $17/4-m>0$.

\subsection{Proof of the dimension-three determinant}\label{gff:d3:determinantProof}
\subsubsection{The half-offset lattice, not a dimension-two square}
At $q=3/4$ set $k=\ell+n$, $j=n$ in the bilateral sums
\eqref{gff:fp:periodization}, with $g_\ell$ and $a_\ell$ as in \eqref{gff:fp:kernel}, and put
\[
 x=k+t+5/8,\qquad y=j+t+1/8,\qquad u=t+1/8.
\]
Gamma reflection gives, as a meromorphic identity,
\begin{equation}
 g_\ell(a_\ell+n+t)
   =-\frac{\sin\pi(3/4-t)}{\sin\pi t}\,h(x)h(y),
 \qquad d_\ell=2\ell+1=2(x-y).
 \label{gff:td:latticeKernel}
\end{equation}
Here $k\ge j$, with both indices otherwise bilateral.  The two sequences
are $h(k+u+1/2)$ and $h(j+u)$; unlike $d=2$, they are not the same
sequence on the same integer lattice.

The five-loop Gamma measure becomes elementary:
\begin{equation}\boxed{
 \frac{\Gamma(\tfrac12+v)\Gamma(\tfrac12-v)}
      {\Gamma(v)\Gamma(-v)}=-v\tan\pi v.
}\label{gff:td:elementaryMeasure}\end{equation}
This follows directly from the two reflection formulas in
\cite{gff:DLMF5}.  At a sampled pole-family shift $v=2s=x+y$,
\[
 (\ell+\tfrac12)^2-v^2=-4xy,
 \qquad (-1)^\ell=(-1)^{k-j}.
\]
Multiplying by $d_\ell$ gives the polynomial
$8xy(x^2-y^2)$ times the common tangent $T(t)$ and the alternating
sign.  Splitting the bilateral lattice into its two infinite tails and
its middle part now proves \eqref{gff:td:wheelSample}--\eqref{gff:td:sevenSample}.
For example the middle part has $\ell=j+k+1$, $j,k\ge0$, and
$x=k+t+5/8$, $y=j+7/8-t$.  Its weight is $2(x+y)$ in a wheel and the
rank-two polynomial $8(x^3y-xy^3)T(t)$ in the primitive.  These separate
into the displayed ordinary products.

\subsubsection{Completing only the antisymmetric half}
A shift $t\mapsto t+1/2$ interchanges the two lattices.  In a wheel the
factor $x-y$ changes sign.  In the primitive the extra alternating sign
also changes sign.  After reindexing the two complementary half-planes,
the full-square sums factor into determinants.  The result is
\begin{equation}\begin{aligned}
 &\mathsf X_r(t)-(-1)^r\mathsf X_r(t+\tfrac12)\\
 &\quad=2(-1)^r\bigl[
 \Psi_{r,1}^{(+)}(u+\tfrac12)\Psi_{r,0}^{(+)}(u)
 -\Psi_{r,0}^{(+)}(u+\tfrac12)\Psi_{r,1}^{(+)}(u)\bigr],
 \quad r=3,4,\\[1mm]
 &\mathsf X_7(t)-\mathsf X_7(t+\tfrac12)\\
 &\quad=-8T(t)\bigl[
 \Psi_{7,3}^{(-)}(u+\tfrac12)\Psi_{7,1}^{(-)}(u)
 -\Psi_{7,1}^{(-)}(u+\tfrac12)\Psi_{7,3}^{(-)}(u)\bigr].
\end{aligned}\label{gff:td:determinants}\end{equation}
The factors are obtained from
$\sin\pi(t+1/4)\sin\pi(t+3/4)=\cos(2\pi t)/2$.
All moments involved obey $3r/4-m>1$, so the product sums used in the
reindexing are absolutely convergent.  The formulas extend across
removable trigonometric singularities by equality of the complete
meromorphic expressions, not by discarding a pole family.

For the proof of the finite modes, reflection gives
\[
 h(-x)=\frac{\sin\pi(x+1/8)}{\sin\pi(1/8-x)}h(x).
\]
It turns \eqref{gff:td:Psidefinition} into the regularized bilateral moment
$\sin^r\pi(u+1/8)\sum_{n\in\mathbb Z}\omega^n(n+u)^m h(n+u)^r$.
This proof representation has all poles canceled, and satisfies
\begin{equation}\begin{aligned}
 \Psi(u+1)&=(-1)^r\omega\Psi(u),&
 \Psi(-u)&=(-1)^m\Psi(u),\\
 |\Psi(x+\ii y)|&=O\!\left(e^{r\pi|y|}|y|^{m+1-3r/4}\right)
 &&(|y|\to\infty).
\end{aligned}\label{gff:td:momentGrowth}\end{equation}
The growth estimate follows from centered Gamma-ratio asymptotics and
absolute integral comparison \cite{gff:DLMFStirling}.  Periodicity or
antiperiodicity supplies integer or half-integer Fourier frequencies.
Vertical contour translation excludes their boundary frequencies when
$m+1-3r/4<0$.  Reflection then leaves exactly the modes in
\eqref{gff:td:momentModes}.  This proves finite-dimensionality, rather than
assuming it from a few samples.

Substituting those modes in \eqref{gff:td:determinants} yields
\begin{equation}\boxed{\begin{aligned}
 \mathsf X_3(t)+\mathsf X_3(t+\tfrac12)
       &=-2\mathfrak a_3\mathfrak b_3,\\
 \mathsf X_4(t)-\mathsf X_4(t+\tfrac12)
       &=4\mathfrak a_4\mathfrak b_4\cos\theta(t),\\
 \mathsf X_7(t)-\mathsf X_7(t+\tfrac12)
       &=2\mathcal O(t).
\end{aligned}}\label{gff:td:projectedAnswer}\end{equation}
In the cosine convention \eqref{gff:fp:finiteFourier}, the first two lines
prove $c_{3,2}=0$, $c_{3,0}=-\mathfrak a_3\mathfrak b_3$,
$c_{4,3}=0$, and $c_{4,1}=2\mathfrak a_4\mathfrak b_4$.
The last line fixes
\begin{equation}
 c_{7,1}=-4\alpha-4\beta,\qquad
 c_{7,3}=4\alpha+8\beta,\qquad c_{7,5}=-4\beta.
 \label{gff:td:sevenOddCoefficients}
\end{equation}
Of these, $c_{3,2}=0$ and $c_{4,3}=0$ hold in every dimension, by the top-mode
theorem \eqref{gff:gdr:topless} of Section~\ref{gff:gdr:top}; the identifications
$c_{3,0}=-\mathfrak a_3\mathfrak b_3$ and $c_{4,1}=2\mathfrak a_4\mathfrak b_4$
and the coefficients \eqref{gff:td:sevenOddCoefficients} are the additional
restrictions special to three dimensions. The identifications of $c_{3,0}$ and $c_{4,1}$ also
hold at a general rim exponent $b$, by the same proof (Section~\ref{gff:gdr:halfShift}); the $d=3$ forms
of $\LL_d$ and $Z_d$ in Section~\ref{gff:gdr:special} use them.

\subsubsection{Finite residue algebra and the remaining even polynomial}
For a self-contained normalization check, the weights
$B_{r,k}/\pi^{r-1}$ in \eqref{gff:fp:finiteFunctional} at $q=3/4$ are
\begin{equation}\begin{aligned}
 r=3:&\quad(8,-6\sqrt2,12),\\
 r=4:&\quad(32,-24\sqrt2,40,-40\sqrt2),\\
 r=7:&\quad(2624,-1904\sqrt2,2912,-2352\sqrt2,
             4032,-3696\sqrt2,7392).
\end{aligned}\label{gff:td:residueWeights}\end{equation}
Together with \eqref{gff:fp:periodsFinite} at $q=3/4$, that is
$Q_3=\Gamma(3/2)^3\mathcal E_3/3$,
$P_4(3)=\Gamma(3/2)^4\mathcal E_4$, and
$P_5(3)=\Gamma(3/2)^5\mathcal E_7/2$, these weights give
\eqref{gff:td:Qanswer} and \eqref{gff:td:P4answer} immediately.  The clockwise
right closure supplies the minus sign in $\mathcal E_r$.

For $r=7$ there are also exact zeros
\begin{equation}
 \mathsf X_7(\tfrac18)=\mathsf X_7(\tfrac58)=0.
 \label{gff:td:sevenZeros}
\end{equation}
Indeed $v=\ell+3/4+2n+2t$ is an integer at either point, and the
measure in \eqref{gff:td:elementaryMeasure} vanishes.  No Gamma pole of
$g_\ell$ occurs there.  The termwise vanishing is legitimate in the
absolutely convergent regularized sum.  Hence the even part is exactly
\begin{equation}
 \mathsf X_7(t)-\mathcal O(t)
  =(1-u)\bigl(p_0+p_1u+p_2u^2\bigr),\qquad
 u=\cos(4\pi(t+3/8)).
 \label{gff:td:evenPolynomial}
\end{equation}
The even-mode contribution to the $r=7$ residue weight is
$-32(9p_0+13p_1+22p_2)$; the odd contribution is
$1792\sqrt2(2\beta-\alpha)$.  For every quadratic polynomial $p$,
\begin{equation}
 9p_0+13p_1+22p_2
 =-35p(0)+\left(22+\frac{13}{\sqrt2}\right)p(1/\sqrt2)
           +\left(22-\frac{13}{\sqrt2}\right)p(-1/\sqrt2).
 \label{gff:td:quadraticRule}
\end{equation}
This is verified by comparing three coefficients.  It proves
\eqref{gff:td:P5answer} without a derivative, infinitesimal step, or numerical
interpolation error.

Crucially, \eqref{gff:td:projectedAnswer} does not determine the complementary
parity sector.  It leaves one coefficient for $r=3$, two for $r=4$, and
three for the even $r=7$ polynomial.  Equations \eqref{gff:td:Hbase} and
\eqref{gff:td:Hseven} evaluate that remaining information.  This explains
why the dimension-two square proof of Section~\ref{gff:fp:twoDimensional} cannot simply be repeated with the
replacement $d_\ell=2\ell+1$.  It does not prove that some different
identity could never reduce the residual Horn values further.

\subsection{A variant with fixed complex nodes}\label{gff:td:complexVariant}
The real rational formulas above are the preferred presentation when parameter simplicity
matters. Two further exact finite-node identities give smaller two-variable counts. Define
\begin{equation}\begin{aligned}
 t_4^*&=-\tfrac38+\frac{\ii\log2}{4\pi},\qquad
 t_7^*=-\tfrac38+\frac{\ii}{4\pi}\operatorname{arcosh}\frac{22}{13},\\
 E_0^*&=\mathsf X_7(0)-\mathcal O(0),\qquad
 E_*^*=\mathsf X_7(t_7^*)-\mathcal O(t_7^*),
\end{aligned}\label{gff:td:complexNodes}\end{equation}
with $\mathcal O$ as in \eqref{gff:td:oddSeven}; $E_0^*$ is the $E_0$ of \eqref{gff:td:sevenE}. Then
\begin{equation}\boxed{\begin{aligned}
 P_4(3)&=8\pi\left[3\sqrt2\,\mathfrak a_4\mathfrak b_4-\mathsf X_4(t_4^*)\right],\\
 P_5(3)&=\frac{2\sqrt2\,\pi^{3/2}}{11}\left[29E_0^*-\frac{2197}{9}E_*^*\right]
        -448\pi^{3/2}(2\beta-\alpha).
\end{aligned}}\label{gff:td:complexAnswers}\end{equation}
Both nodes lie on the line $\operatorname{Re}t=-\frac38$. There every $\cos\bigl(2\pi k(t+\frac38)\bigr)$ is real,
so the samples \eqref{gff:fp:finiteFourier}, whose coefficients are real, are real.

\emph{Proof.} Since $\Xi_r$ is entire, \eqref{gff:fp:finiteFourier} and \eqref{gff:td:evenPolynomial} are
identities of entire functions whose right sides are trigonometric polynomials, so they hold at
complex $t$. By \eqref{gff:fp:finiteFourier} at $q=\frac34$, with $c_{4,3}=0$ and
$c_{4,1}=2\mathfrak a_4\mathfrak b_4$ (Section~\ref{gff:d3:determinantProof}),
$\mathsf X_4(t)=c_{4,0}+2\mathfrak a_4\mathfrak b_4\cos\theta(t)+c_{4,2}\cos2\theta(t)$ with $\theta(t)$
as in \eqref{gff:td:oddSeven}. The weights \eqref{gff:td:residueWeights}, inserted in
\eqref{gff:fp:finiteFunctional} and \eqref{gff:fp:periodsFinite}, give
$P_4(3)=-\pi\bigl(8c_{4,0}-6\sqrt2\,c_{4,1}+10c_{4,2}\bigr)$. At $t_4^*$ one has
$\cos\theta=\cosh(\frac12\log2)=\frac3{2\sqrt2}$ and $\cos2\theta=\cos\bigl(4\pi(t_4^*+\frac38)\bigr)=\frac54$,
which gives the first line. For the second, \eqref{gff:td:evenPolynomial} gives $E_0^*=p(0)$ and, since
$u=\cos\bigl(4\pi(t_7^*+\frac38)\bigr)=\frac{22}{13}$ and $1-\frac{22}{13}=-\frac9{13}$,
$E_*^*=-\frac9{13}\,p(\frac{22}{13})$. The rule
\[
 9p_0+13p_1+22p_2=\frac{29}{22}\,p(0)+\frac{169}{22}\,p\bigl(\tfrac{22}{13}\bigr),
\]
exact for every quadratic $p$ (compare the coefficients of $1$, $u$ and $u^2$), replaces
\eqref{gff:td:quadraticRule} in the proof of \eqref{gff:td:P5answer}; $169\cdot13=2197$. $\square$

A sample on this line needs only one Horn value. By \eqref{gff:td:Harrays} and \eqref{gff:td:Vfinite},
with $x=x_t+k$, $y=y_t+n$ and $k=\ell+n$ as in \eqref{gff:td:latticeKernel},
$V_r(t)=\sum_{k\ge n\ge0}2(x-y)h(x)^rh(y)^r$ for $r=3,4$ and
$V_7(t)=8\sum_{k\ge n\ge0}(-1)^{k-n}xy(x^2-y^2)h(x)^7h(y)^7$. A shift $t\mapsto t+1$ moves $x$ and $y$
by one and so removes the boundary row $n=0$; the sum of that row is denoted $V_r^{(0)}(t)$:
\begin{equation}\begin{aligned}
 V_r(t+1)&=V_r(t)-V_r^{(0)}(t),\\
 V_r^{(0)}(t)&=2h(y_t)^r\left[\mathfrak U_{r,1}^{(+)}(x_t)-y_t\,\mathfrak U_{r,0}^{(+)}(x_t)\right],
 \qquad r=3,4,\\
 V_7^{(0)}(t)&=8h(y_t)^7\left[y_t\,\mathfrak U_{7,3}^{(-)}(x_t)-y_t^3\,\mathfrak U_{7,1}^{(-)}(x_t)\right].
\end{aligned}\label{gff:td:boundaryShift}\end{equation}
Let $t=-\frac38+\ii\sigma$ with $\sigma>0$. Then $t^\dagger=1+\bar t$, and since $V_r$ is real on the
real axis, $V_r(t^\dagger)=\overline{V_r(t)-V_r^{(0)}(t)}$. The other ordinary factors reduce by
$\mathfrak U_{r,m}^{(\omega)}(u+1)=\omega\bigl[\mathfrak U_{r,m}^{(\omega)}(u)-u^mh(u)^r\bigr]$.
This proves the counts of one $\mathfrak H_4$ value for $P_4(3)$ and two $\mathfrak H_7$ values for
$P_5(3)$, with no second complex Horn value hidden in the samples. The nodes are fixed constants,
not regulators. They do introduce fixed complex parameters, some with negative real part (at
$t_4^*$ the row $\frac14+t$ of \eqref{gff:td:Harrays} has real part $-\frac18$); there are no
denominator poles at the stated nodes. No minimality claim is made for this variant. Its numerical
agreement with \eqref{gff:td:P4answer} and \eqref{gff:td:P5answer} is given in Section~\ref{gff:d3:values}.

\subsection{The real rank-two hypergeometric normalization}\label{gff:app:d3RankTwo}
This subsection and the next two identify the residual functions $\mathfrak H_r$ with standard
real matrix-argument hypergeometric functions of rank two; Section~\ref{gff:rt:kummerDiagnostic} then
records a matrix transformation that maps the tetrahedron to itself and so does not reduce it. The identification is exact and
coefficientwise, not a numerical fit. It does not reduce them to ordinary one-variable ${}_pF_q$
values: the ordinary-hypergeometric endpoint reached for the primitive periods at $d=2$
(Section~\ref{gff:fp:twoDimensional}) is not reached at $d=3$. What it adds is a classification: the
residual objects are low-order functions on a rank-two symmetric cone, not unclassified Horn
arrays. The rational-node formulas above remain the preferred explicit specification of the
periods.

For a partition $\kappa=(m,n)$ with $m\geq n\geq0$ (here $m$ is a part of
$\kappa$, not the moment order $m$ of \eqref{gff:td:Udefinition}), write
\begin{equation}
 [a]_\kappa=(a)_m(a-\tfrac12)_n,
 \qquad |\kappa|=m+n.
 \label{gff:rt:partitionPochhammer}
\end{equation}
From here to the end of this section the Jack parameter is $\alpha=2$, the
matrix size is two, and the physical dimension is $d=3$. This $\alpha$ is
unrelated to the constant $\alpha$ of \eqref{gff:td:oddSeven}.  These three numbers
refer to different notions.  Let $C_\kappa^{(2)}$ denote the
$C$-normalized real zonal polynomial, so that
$\sum_{|\kappa|=k}C_\kappa^{(2)}(X)=(\operatorname{tr}X)^k$.
We use the standard definition \cite{gff:MacdonaldHyperI,gff:DLMFMatrixHyper}
\begin{equation}
 {}_pF_q^{(\alpha=2)}\!\left(
 \begin{matrix}\boldsymbol a\\\boldsymbol b\end{matrix};zI_2\right)
 =\sum_{m\geq n\geq0}
 \frac{\prod_i[a_i]_{(m,n)}}{\prod_j[b_j]_{(m,n)}}
 \frac{C_{(m,n)}^{(2)}(I_2)}{(m+n)!}\,z^{m+n}.
 \label{gff:rt:matrixDefinition}
\end{equation}
This definition fixes the conventional special function; it is not a
claim that writing a summation as a function symbol has removed its
rank-two structure.

The identity that matches the angular multiplicity is
\begin{equation}\boxed{
 \frac{C_{(m,n)}^{(2)}(I_2)}{(m+n)!}
 =\frac{2(m-n)+1}{(\tfrac32)_m n!}
 =\frac{2(m-n)+1}{[\tfrac32]_{(m,n)}}.
}\label{gff:rt:hookWeight}\end{equation}
Here is an explicit hook-product verification, including normalization.
If $a(s)$ and $l(s)$ are the arm and leg lengths of a box of $\kappa$, the
Jack hook product is
\begin{equation}\begin{aligned}
 j_{(m,n)}(2)
 &=\prod_{s\in(m,n)}(2a(s)+l(s)+1)(2a(s)+l(s)+2)\\
 &=2^{2(m+n)}
 \frac{(\tfrac32)_m\,m!\,n!\,(\tfrac12)_n}{2(m-n)+1}.
\end{aligned}\label{gff:rt:hooks}\end{equation}
The specialization formula
$C_\kappa^{(2)}(I_2)/|\kappa|!
=2^{2|\kappa|}[1]_\kappa/j_\kappa(2)$ then gives
\eqref{gff:rt:hookWeight}.  In particular the factor $2\ell+1$ with
$\ell=m-n$ is exactly the rank-two real-zonal weight after canceling
one $[3/2]_\kappa$.  This is why a real matrix-argument function appears.

\subsection{A standard \texorpdfstring{${}_3F_2$}{3F2} of matrix argument for the tetrahedron}
Retain the residual functions $\mathfrak H_r(t)$ of \eqref{gff:td:Hbase} and
\eqref{gff:td:Hseven} and put
\begin{equation}
 a=t+\frac34,\qquad b=t+\frac32.
 \label{gff:rt:ab}
\end{equation}
For $r=3,4$ their exact identification is
\begin{equation}\boxed{
 \mathfrak H_r(t)=
 {}_{r+1}F_r^{(\alpha=2)}\!\left(
 \begin{matrix}a^{\times r},\ \frac32\\ b^{\times r}\end{matrix};I_2\right).
}\label{gff:rt:wheelMatrix}\end{equation}
To verify it, use $m=\ell+n$ in the explicit Horn arrays.  The summand
of either side is
\begin{equation}
 (2\ell+1)
 \left[\frac{(a)_{\ell+n}(a-\frac12)_n}
 {(b)_{\ell+n}(b-\frac12)_n}\right]^r.
 \label{gff:rt:wheelCoefficient}
\end{equation}
On the matrix side this follows from canceling
$[3/2]_{(m,n)}$ against \eqref{gff:rt:hookWeight}.  Thus no summation theorem
or exchange of a conditionally convergent sum is being assumed.
The endpoint sums here are absolutely convergent: along $m\to\infty$
with $n$ fixed the majorant is $m^{1-3r/4}$, and the remaining $n$
factors are summable for $r=3,4$.

At $t=0$ one upper parameter $3/2$ cancels a lower parameter. The
tetrahedron formula \eqref{gff:td:Qanswer} therefore becomes
\begin{equation}\boxed{
 Q_3=4\pi^2
 {}_3F_2^{(\alpha=2)}\!\left(
 \begin{matrix}\frac34,\frac34,\frac34\\\frac32,\frac32\end{matrix};I_2\right)
 -\frac{\Gamma(\frac14)^{10}}{48\pi^{7/2}}.
}\label{gff:rt:Q3}\end{equation}
The same value is the rim sample $\mathfrak t_3(b)$ of \eqref{gff:gdr:samples3} at $d=3$ and
$b=\frac32$: there $\sigma=g=\frac34$, and the summand $d_\ell T_{\ell n}(\frac32)^3$ of
\eqref{gff:gdr:Tln} is \eqref{gff:rt:wheelCoefficient} with $r=3$, $t=0$. Its second $b$-derivative enters
$Z_3$ (Section~\ref{gff:gdr:special}).
Both remaining four-loop values likewise have short conventional arrays:
\begin{equation}\boxed{\begin{aligned}
 \mathfrak H_4(0)
 &={}_4F_3^{(\alpha=2)}\!\left(
 \begin{matrix}(\frac34)^{\times4}\\(\frac32)^{\times3}\end{matrix};I_2\right),\\
 \mathfrak H_4(\tfrac18)
 &={}_5F_4^{(\alpha=2)}\!\left(
 \begin{matrix}(\frac78)^{\times4},\frac32\\(\frac{13}{8})^{\times4}\end{matrix};I_2\right).
\end{aligned}}\label{gff:rt:P4blocks}\end{equation}
Substitution in \eqref{gff:td:P4answer} gives an equivalent
matrix-argument presentation of $P_4(3)$ without changing its ordinary
moment factors.  Neither of these matrix functions is to be evaluated
by an ordinary scalar \texttt{hyper} call.

A local coefficient check makes the distinction unambiguous.  With the
same upper and lower parameters as in \eqref{gff:rt:Q3},
\begin{equation}\begin{aligned}
 {}_3F_2^{(\alpha=2)}(\tfrac34,\tfrac34,\tfrac34;
                         \tfrac32,\tfrac32;zI_2)
 &=1+\frac38z+\frac{2769}{12800}z^2
        +\frac{15113}{102400}z^3+O(z^4),\\
 {}_3F_2(\tfrac34,\tfrac34,\tfrac34;
                         \tfrac32,\tfrac32;z)
 &=1+\frac{3}{16}z+\frac{1029}{12800}z^2
        +\frac{9317}{204800}z^3+O(z^4).
\end{aligned}\label{gff:rt:scalarDistinction}\end{equation}
The first coefficients already differ.  This checks the normalization;
it does not preclude an identity with some \emph{other} scalar
hypergeometric functions or parameters.

\subsection{Five loops: an exact three-term standard-function dictionary}
The polynomial insertion in $\mathfrak H_7$ can also be expressed using
ordinary parameter shifts of standard rank-two functions, rather than
a bespoke extra Horn row. Set (this $v$ is local; it is not $v_t$ of \eqref{gff:td:Hseven})
\begin{equation}
 v=t+\frac58,\qquad \delta_v=v(v-\tfrac12).
 \label{gff:rt:vdelta}
\end{equation}
Define three standard matrix-argument functions, all at $-I_2$:
\begin{equation}\begin{aligned}
 \mathcal J_A(t)&={}_9F_8^{(\alpha=2)}\!\left(
 \begin{matrix}a^{\times7},\frac32,v+2\\b^{\times7},v\end{matrix};-I_2\right),\\
 \mathcal J_B(t)&={}_{10}F_9^{(\alpha=2)}\!\left(
 \begin{matrix}a^{\times7},\frac32,v+1,v+1\\b^{\times7},v,v\end{matrix};-I_2\right),\\
 \mathcal J_C(t)&={}_9F_8^{(\alpha=2)}\!\left(
 \begin{matrix}a^{\times7},\frac32,v+1\\b^{\times7},v\end{matrix};-I_2\right).
\end{aligned}\label{gff:rt:JABC}\end{equation}
Then the exact dictionary is
\begin{equation}\boxed{
 \mathfrak H_7(t)=
 \frac{\delta_{v+1}\mathcal J_A(t)
       -\delta_v\mathcal J_B(t)-\mathcal J_C(t)}{2v-\frac12}.
}\label{gff:rt:H7matrix}\end{equation}
The five arguments required by \eqref{gff:td:P5answer}, listed in
\eqref{gff:td:fiveHornNodes}, are
$t=0,\pm1/16,3/16,5/16$.  They all have $v>1/2$, and all generalized
Pochhammer factors in these arrays have positive real starting
arguments.  At $t=0$, canceling one common $3/2$ changes the three
orders to ${}_8F_7$, ${}_9F_8$ and ${}_8F_7$, still \emph{of matrix
argument}.  No coalescence limit or parameter differentiation is used.

Here is the coefficientwise proof, which also fixes the sign of the
matrix argument.  For $\kappa=(m,n)$ put
$x=m+v$, $y=n+v-1/2$.  Direct Pochhammer cancellation gives
\begin{equation}\begin{aligned}
 \frac{[v+1]_\kappa}{[v]_\kappa}&=\frac{xy}{\delta_v},\\
 \frac{[v+2]_\kappa}{[v]_\kappa}
 &=\frac{xy(x+1)(y+1)}{\delta_v\delta_{v+1}}.
\end{aligned}\label{gff:rt:insertionRatios}\end{equation}
The only polynomial identity needed is
\begin{equation}
 xy(x+y)=xy(x+1)(y+1)-(xy)^2-xy.
 \label{gff:rt:insertionPolynomial}
\end{equation}
Consequently the right-hand side of \eqref{gff:rt:H7matrix} multiplies the
common coefficient \eqref{gff:rt:wheelCoefficient}, with $r=7$, by
$xy(x+y)/(\delta_v(2v-1/2))$, exactly the three polynomial ratios in the
Horn array \eqref{gff:td:Hseven} of $\mathfrak H_7$.  Finally
$(-1)^\ell=(-1)^{m-n}=(-1)^{m+n}$, so the matrix argument is $-I_2$,
\emph{not} $\operatorname{diag}(-1,1)$.
Each of the three endpoint series separately converges absolutely.
For example, the largest polynomial insertion has degree two in each
of $m,n$; its absolute summand is bounded by a constant times
$(m+1)^{-9/4}(n+1)^{-13/4}$ on $m\geq n$, which is summable.

This dictionary standardizes the function class, but increases the
number of displayed functions per Horn block from one to three.
It is not advertised as a more economical numerical formula.  In
particular, it does not eliminate the complementary ordered-lattice
information isolated by the half-shift determinant.

\subsection{A checked transformation that does not finish the tetrahedron}\label{gff:rt:kummerDiagnostic}
The matrix Kummer transformation \cite[Eq.~(35.8.5)]{gff:DLMFMatrixHyper} is a useful diagnostic. For
$s_{\rm K}=b_1+b_2-a_1-a_2-a_3$ it writes the ${}_3F_2^{(\alpha=2)}$ at $I_2$ with the parameter array
$(a_1,a_2,a_3;b_1,b_2)$ as the one with the array
\begin{equation}
 (a_1,a_2,a_3;b_1,b_2)\longmapsto(b_1-a_1,b_1-a_2,a_3;b_1,s_{\rm K}+a_3),
 \label{gff:rt:kummerMap}
\end{equation}
times the prefactor $\Gamma_2(b_2)\Gamma_2(s_{\rm K})/[\Gamma_2(b_2-a_3)\Gamma_2(s_{\rm K}+a_3)]$. Here
$\Gamma_2(a)=\sqrt\pi\,\Gamma(a)\Gamma(a-\frac12)$ is the multivariate Gamma function of order two,
and the conditions $\operatorname{Re}b_2,\operatorname{Re}s_{\rm K}>\frac12$ hold below. For the array of \eqref{gff:rt:Q3},
$s_{\rm K}=\frac34$: the transformed array is identical, and the prefactor equals one. Parameter
permutations add nothing, because the three upper parameters and the two lower parameters are
separately equal. Transformations generated by this map and these permutations therefore expose no
canceling parameter and no terminating series. This is a statement about this explicit route, not
a proof that a different reduction to ordinary ${}_pF_q$ values cannot exist. The script
\path{code/rank_two_checks.py} of \path{lri_d3_matrix_update/} checks the fixed point exactly.

The determinantal formulas for \emph{Hermitian} matrix-argument functions have Jack parameter
$\alpha=1$, whereas \eqref{gff:rt:hookWeight} fixes $\alpha=2$ here. Transplanting a Hermitian formula
would change the coefficients. The distinction between real and Hermitian matrix-argument
transformations is made explicitly in \cite{gff:RichardsZhengMatrix}; it explains why the familiar
determinant machinery is not itself the missing scalar identity.

\subsection{Three-dimensional values and limits of the reduction}\label{gff:d3:values}
The fixed rational formulas give
\[
\begin{aligned}
 Q_3&=5.01710256455385536959085484569018\ldots,\\
 P_4(3)&=62.0583279434066724664633885840967\ldots,\\
 P_5(3)&=240.340317759083549816366546388004\ldots.
\end{aligned}
\]
The digits shown are correctly rounded; $\ldots$ marks omitted digits.
Fresh integrations of the Gamma kernels on the real axis, which do not use these formulas, agree
with these values to about $10^{-51}$, $10^{-52}$ and $10^{-53}$, respectively. For $Q_3$ and $P_4(3)$
there is a stronger independent check: the $100$-digit spectral evaluation of
Section~\ref{gff:gs:masterTable} agrees with $3Q_3$ and $P_4(3)$ to $3.5\times10^{-79}$ and
$2.1\times10^{-86}$. For $P_5(3)$ the kernel integration is the only independent check. The formulas
themselves, evaluated at 55 and at 90 working digits with changed truncation and tail order, differ
by $2.0\times10^{-53}$, $2.9\times10^{-53}$ and $1.6\times10^{-49}$; the 55-digit run is the less
accurate one, since the 90-digit values agree with a further 100-digit evaluation to better than
$10^{-74}$. These are numerical residuals, not interval bounds. The rank-two identity is a
coefficientwise identity with Jack parameter $\alpha=2$: neither its
superscript nor its matrix argument may be dropped. No complete reduction
of these remaining values, or of $\mathcal N_3$ (\eqref{gff:fp:remainingCombination}
at $d=3$), to scalar ordinary
${}_pF_q$ values is asserted.

\paragraph{How the values were obtained and checked.}
The formula side sums the ordinary and two-variable series of the formulas above to a truncation
point and adds an asymptotic tail built from the Stirling expansion of the Gamma ratios and Hurwitz
zeta values. The runs at 55 and 90 working digits used truncation $80$ with tail order $36$ and
truncation $160$ with tail order $64$ (\path{code/general_d3/verify_d3.py}, record
\path{verification/general_d3/period_values.json} of \path{lri_d3_reduction_update/}). The 100-digit
run used truncation $220$ and tail order $80$ (\path{code/recompute_d3_periods.py}, record
\path{verification/fresh_primitive_values.json} of \path{lri_d3_matrix_update/}); it repeats the
formula side and is not an independent integral. The tails are numerical accelerations, not
interval bounds and not fits.

The independent check (\path{code/general_d3/check_spectral.py}, record
\path{verification/general_d3/fresh_spectral.json}) integrates the real-axis Gamma kernels afresh,
without reading stored period values, at 45 and 60 working digits with different quadrature
breakpoints and spin cutoffs. Its wheel spin tail comes from analytically integrated
centered-Stirling coefficients, not from a fit to the answer. With the symbols $\lambda_j$, $e_j$
and $\mathfrak w_{r,h}$ local to this paragraph, put $e_0=1$ and
\begin{equation}\begin{aligned}
 \lambda_j&=-\frac{rB_{2j+1}(\frac18)}{j(2j+1)},\qquad
 e_j=\frac1j\sum_{k=1}^jk\lambda_ke_{j-k},\\
 \mathfrak w_{r,h}&=2\Gamma(\tfrac32)^r\,\Gamma(\tfrac{3r}2+2h-1)
 \sum_{j=0}^h\frac{e_je_{h-j}}{\Gamma(\frac{3r}4+2j)\,\Gamma(\frac{3r}4+2h-2j)},
\end{aligned}\label{gff:td:analyticWheelTail}\end{equation}
where $B_n(x)$ are the Bernoulli polynomials. In $\langle t_\ell^r\rangle_3$ of
Section~\ref{gff:c1:closedForms}, which is $3Q_3$ for $r=3$ and $P_4(3)$ for $r=4$, the spins
$\ell<\ell_{\rm c}$ are integrated directly, and the sum over the spins $\ell\ge\ell_{\rm c}$ is
approximated by the truncated asymptotic series
$\sum_{h=0}^{h_{\max}}\mathfrak w_{r,h}\,\zeta(\frac{3r}2-2+2h,\ell_{\rm c}+\frac12)$, with the
Hurwitz zeta function; the 60-digit run used $\ell_{\rm c}=72$ and $h_{\max}=32$. This is a
numerical acceleration, not an extra operation in \eqref{gff:td:Qanswer} or \eqref{gff:td:P4answer}. The
five-loop kernel uses an alternating spin sum, accelerated separately by an Euler transformation
(232 spins in the 60-digit run). In that run the formula values minus the kernel values are
$-1.47\times10^{-51}$ for $Q_3$, $-1.53\times10^{-52}$ for $P_4(3)$ and $-2.41\times10^{-53}$ for
$P_5(3)$. These residuals are the differences of \path{period_values.json} (90 digits) and
\path{fresh_spectral.json} (60 digits), also listed in the delivered record
\path{verification/general_d3/comparison_summary.json}.

At 75 working digits, evaluations at the complex points $t=0.063+0.041\ii$ and $t=-0.082+0.027\ii$
confirm the half-shift determinants \eqref{gff:td:determinants}, with largest residual
$2.42\times10^{-67}$, and evaluations at two complex points $u$ confirm the finite modes
\eqref{gff:td:momentModes} to $5\times10^{-70}$ (\path{verification/general_d3/identities.json}). In the
same run the complex-node formulas \eqref{gff:td:complexAnswers} differ from \eqref{gff:td:P4answer} and
\eqref{gff:td:P5answer} by $7.0\times10^{-71}$ and $2.7\times10^{-65}$. Twenty-five exact symbolic checks
of the finite residue algebra, the sample rules and the polynomial identities, and 288 comparisons of
rational Horn coefficients, pass (\path{code/general_d3/check_exact.py}, record
\path{verification/general_d3/exact_algebra.json}). Six literal ordinary-hypergeometric evaluations
at 40 working digits agree with the tail evaluator to $2.06\times10^{-38}$
(\path{verification/general_d3/literal_hypergeometric.json}). All these residuals are observed values,
not bounds. They are master-level checks, not a new three-dimensional full-graph calculation or a
rigorous enclosure.

\paragraph{Checks of the rank-two identification.}
The script \path{code/rank_two_checks.py} of \path{lri_d3_matrix_update/} (record
\path{verification/rank_two_checks.json}) builds the Jack hook products independently and compares
them with \eqref{gff:rt:hookWeight} on all $91$ partitions $\kappa=(m,n)$ with $m\le12$. On the same
partitions it compares the coefficients of \eqref{gff:rt:wheelMatrix} for $r=3,4$ at $t=0,\frac18,\frac27$,
and those of \eqref{gff:rt:H7matrix} at the five nodes \eqref{gff:td:fiveHornNodes}, with the Horn
coefficients. With the polynomial identity \eqref{gff:rt:insertionPolynomial}, the Kummer fixed point of
Section~\ref{gff:rt:kummerDiagnostic} and the normalization
$\sum_{|\kappa|=k}C^{(2)}_\kappa(I_2)/k!=2^k/k!$ for $k\le12$, these are 1107 exact checks, all passing.
Direct expansions of the matrix functions and of the Horn series agree through total degree $90$ at
$z=\frac14$ (for $\mathfrak H_3(0)$ and $\mathfrak H_4(\frac18)$) and $z=-\frac14$ (for
$\mathfrak H_7(\pm\frac1{16})$), with $z$ as in \eqref{gff:rt:matrixDefinition}, at 90 working digits.
These are coefficient and interior-series checks, not interval bounds on the endpoint values.

\paragraph{A bounded search for a scalar form of the tetrahedron.}
Consider the 6479 convergent ordinary values $F={}_3F_2(\boldsymbol a;\boldsymbol b;1)$ and
$F={}_4F_3(\boldsymbol a;\boldsymbol b;1)$ whose upper parameters lie in
$\{\frac18,\frac28,\ldots,1\}$ and lower parameters in $\{\frac12,\frac34,1,\frac54,\frac32\}$, with
positive excess and no upper parameter equal to a lower one. For each of them PSLQ was run on the
vector $\bigl(\mathfrak H_3(0),1,\Gamma(\frac14)^{10}/\pi^{11/2},F\bigr)$ at 85 working digits, with
tolerance $10^{-57}$, coefficients up to $10^6$ and at most 800 iterations. No relation with a
nonzero coefficient of $\mathfrak H_3(0)$ was returned (\path{code/scalar_catalog_search.py} of
\path{lri_d3_matrix_update/}, record \path{verification/scalar_catalog_search.json}). This finite
search is not an independence result and does not rule out other parameter arrays, arguments,
prefactors or nonlinear identities. It concerns the tetrahedron only.

\paragraph{The open primitive problem.}
The smallest open primitive question is to evaluate the single rank-two
${}_3F_2^{(\alpha=2)}$ of \eqref{gff:rt:Q3} by finitely many ordinary scalar hypergeometric values. The
four- and five-loop complements need the related rank-two functions \eqref{gff:rt:P4blocks} and
\eqref{gff:rt:H7matrix}. No such scalarization is proved here. Reducing these primitive functions would
still leave $\mathcal N_3$ for $C_5(3)$.

\paragraph{Coefficients through fourth order.}
At $d=3$, that is $q=\frac34$, the constants \eqref{gff:gs:psiBasis} are
\begin{equation}\begin{aligned}
 A&=2+4\log2-\pi,& R&=\frac{\pi^2}3-4,& U&=2\pi^3-56\zeta(3),\\
 V&=16-16\zeta(3),& X&=\frac{14\pi^4}{15}-96,
\end{aligned}\label{gff:td:physicalDictionary}\end{equation}
by reflection, recurrence and the integer and half-integer polygamma values
\cite{gff:DLMFPolygamma,gff:DLMFDigamma}. Here $V$ and $X$ are these constants, not $V_r(t)$ of
\eqref{gff:td:Vfinite}, $\mathsf X_r(t)$ or $X_0$, $X_h$ of \eqref{gff:td:P4data}. Then \eqref{gff:gs:C0C4} reads
\begin{equation}\boxed{\begin{gathered}
 C_0(3)=\tfrac32,\qquad C_1(3)=-\tfrac16,\qquad C_2(3)=\tfrac A9,\\
 108\,C_3(3)=A^2+\tfrac{\pi^2}3-4-8Q_3,\\
 648\,C_4(3)=16P_4(3)-A^3-8AQ_3-\pi^2A+12A+53\pi^3-1436\zeta(3)-48,
\end{gathered}}\label{gff:td:C0C4}\end{equation}
with $Q_3$ from \eqref{gff:td:Qanswer} and $P_4(3)$ from \eqref{gff:td:P4answer}. This gives
$C_4(3)=1.22806354532762307997\ldots$; it is a reevaluation of the coefficient formula, not a new
four-loop graph integration.

With these values and those of the $d=3$ nonprimitive masters $\LL_3$, $Z_3$ and $\KK_3$
(Section~\ref{gff:gs:masterTable}), \eqref{gff:gs:C5}
gives
\[
 C_5(3)=-5.8503514207513272292179442304977443001109973463783\ldots
\]
to the $50$ digits shown (\path{LRI_C5_d3/c5_value.json}). The formula-side values of $Q_3$,
$P_4(3)$ and $P_5(3)$ are stable to $10^{-74}$ or better, and the independent checks above and in
Section~\ref{gff:gs:masterTable} confirm $Q_3$, $P_4(3)$, $P_5(3)$ and the nonprimitive masters to about
$10^{-79}$, $2\cdot10^{-86}$, $2\cdot10^{-53}$ and $1.2\cdot10^{-93}$ or better. Propagated through
\eqref{gff:gs:C5}, these differences change $C_5(3)$ by less than $10^{-54}$: the largest part, about
$8\times10^{-55}$, comes from $P_5(3)$, since $\partial C_5/\partial P_5=-1536/46656\approx-0.033$.
$\KK_3$ is the $430$-digit value of
\eqref{gff:gs:Kseries}; the two-fold Barnes integral \eqref{gff:gdk:KMB}, which uses
none of the inputs of Section~\ref{gff:gdk:conventions}, reproduces it to
$4\cdot10^{-100}$ (Section~\ref{gff:gdk:numerics}), and the generic-dimensional
route of Sections~\ref{gff:gdr:section} and \ref{gff:gdk:section} reproduces
$C_5(3)$ to $5\cdot10^{-45}$ (Section~\ref{gff:gdr:checks}).

\section{Generic-dimensional reduction of \texorpdfstring{$\LL_d$}{Ld} and \texorpdfstring{$Z_d$}{Zd}}\label{gff:gdr:section}

The nonprimitive masters $\LL_d$ and $Z_d$ of \eqref{gff:gs:C5} are rim jets of the balanced wheels,
$\LL_d=\partial_bW_4(d;b)|_{b=d/2}$ and $Z_d=\frac19\partial_b^2W_3(d;b)|_{b=d/2}$ \eqref{gff:gs:wheelJets}.
Section~\ref{gff:app:generalPeriodization} periodizes the wheels only at $b=d/2$. This section extends the finite
periodization to a general rim exponent $b$ (Section~\ref{gff:gdr:periodization}) and proves that the top Fourier modes of
every product of balanced kernels satisfying the convergence condition (C) of \eqref{gff:gdr:BC} vanish in every
dimension, in particular for every balanced wheel (Section~\ref{gff:gdr:top}). With these two statements it gives, for
$\LL_d$ and $Z_d$, in place of the six-variable Horn representation of \eqref{gff:gc:finiteClosure} (244 and 318 values),
samples at three nodes, real for real $b$: five two-variable Horn functions at unit argument for $\LL_d$, and two
${}_7F_6(1)$ values and three two-variable Horn values for $Z_d$, together with their first and second
$b$-derivatives (Sections~\ref{gff:gdr:samples}--\ref{gff:gdr:C5}; the unweighted, derivative-free form
\eqref{gff:gc:finiteClosure} remains valid). $\KK_d$ is treated in Section~\ref{gff:gdk:section}; $P_5(d)$ is not treated
here. Section~\ref{gff:gdr:samples} states the sample formulas of the two wheels, and Sections~\ref{gff:mgl:section}
and~\ref{gff:mgz:proofs} prove them for the four-wheel and the three-wheel and write them out in Horn values. The
package \path{LRI_C5_general_d/} holds the code and records: \path{foundation/} for Section~\ref{gff:gdr:periodization},
\path{L/} for Section~\ref{gff:mgl:section}, \path{Z/} for Section~\ref{gff:mgz:proofs} and \path{specialize/} for the rest;
the review records are in \path{review/}.

\paragraph{Provenance.}
This section, with the code and records in \path{foundation/}, \path{L/}, \path{Z/} and \path{specialize/}, was
written by a Claude Code session. Its complete proofs, with the detailed checks, were added on 16 October 2026 by the
same session, from that session's separate notes, which are not part of the repository. A second Claude Code session,
which wrote
Section~\ref{gff:gdk:section}, supplied two inputs: the explicit constants of \eqref{gff:mgf:gammaBounds},
\eqref{gff:mgf:multiplicityBound} and \eqref{gff:mgf:latticeSums}, taken from its growth proof
(Section~\ref{gff:gdk:growthProof}; package \path{LRI_C5_general_d/K_general/}), and the spectral reference values
\path{LRI_C5_general_d/references/reference_values.json}, computed with its code
\path{LRI_C5_general_d/refvals/hp_general_d.py}. The $d=3$ identities of Section~\ref{gff:gdr:halfShift} state their own
provenance. Section~\ref{gff:gdr:status} says which parts the reviews covered.

\paragraph{Status.}
The general-$b$ statements of Section~\ref{gff:gdr:periodization}, the lattice identities \eqref{gff:gdr:d12lattice} and the
sample formulas of Section~\ref{gff:gdr:samples}, with the two-node formula \eqref{gff:gdr:twonode}, have written proofs here
(Sections~\ref{gff:gdr:periodization}, \ref{gff:mgl:section}, \ref{gff:mgz:proofs} and \ref{gff:gdr:special}). So do the $d=3$
identities \eqref{gff:gdr:halfShiftC} and \eqref{gff:gdr:halfShiftGamma} of Section~\ref{gff:gdr:halfShift}, whose printed
derivation has not been reviewed. The formulas
\eqref{gff:gdr:P4}--\eqref{gff:gdr:C5reduced} are \emph{derived}: exact algebra on top of these theorems, confirmed
numerically (Section~\ref{gff:gdr:checks}). The reviews and checks named in this section and in
Section~\ref{gff:gdk:section} were made by separate AI sessions with their own codes; there was no human referee.
Statements that a further reduction does not exist are \emph{numerical only}. Section~\ref{gff:gdr:status} lists the
status of every statement, the conditional inputs and what each review covered. The evaluation of $\KK_d$ quoted in
Section~\ref{gff:gdr:C5} is that of Section~\ref{gff:gdk:section}, which has its own status list (Section~\ref{gff:gdk:status})
and was not part of the review \path{LRI_C5_general_d/review/VERIFICATION.md}.

\paragraph{Notation.}
In this section $q=d/4$. Except in Section~\ref{gff:gdr:top}, where $\sigma$ is a free offset and $g$ in $f_g$ a free
exponent, $\sigma=b/2$ and $g=(d-b)/2=2q-\sigma$ is the spoke half-exponent; the Barnes kernel $g_\ell$ carries its
spin index. Further $s_b=\sin\frac{\pi b}2$, $\varsigma_b=\cos\frac{\pi b}2$, and in
Sections~\ref{gff:gdr:samples}--\ref{gff:gdr:C5}, at $b=d/2$,
\begin{equation}
 S=\sin\pi q,\qquad C=\cos\pi q,\qquad D=\pi\cot\pi q=\psi_0(1-q)-\psi_0(q).
 \label{gff:gdr:D}
\end{equation}
In Section~\ref{gff:gdr:periodization} and in the bounds of the proofs of Sections~\ref{gff:mgl:section}
and~\ref{gff:mgz:uniform}, $C$ (also $C_K$) denotes a generic constant; (C) is the condition \eqref{gff:gdr:BC}, and $C_r(t)$ is a block of Section~\ref{gff:gdr:residue}. This $D$ is not the dimension $D$ of Section~\ref{gff:sec:P5complete} and not the functions $D_b(t)$, $D_T(t)$ of \eqref{gff:gdr:periodizationDef}, \eqref{gff:gdr:TG}. $A,R,U,V,X$ are the polygamma values
\eqref{gff:gs:psiBasis}; at $q=\frac14$, $A=\pA$ and $U=\pU$ of \eqref{gff:eq:psiQuarterBasis}, and $D=\pD=\pi$. The
calligraphic $\mathcal A,\mathcal B$ are the mode amplitudes of \eqref{gff:gc:L2modes}, not polygamma values. Likewise the samples $\mathcal X_T,\mathcal X_0,\mathcal X_{1/2}$ of \eqref{gff:gdr:samples4} are unrelated to $X$, the $R$ family of $\KK_d$ is unrelated to $R$, and the simplex variable $V$ of Section~\ref{gff:gdr:top} is unrelated to $V$. A prime
denotes $\partial_b$. Functions of $b$ such as $\mathfrak t_3(b)$ of \eqref{gff:gdr:samples3} also depend on $d$; where $d$ is fixed it is
written first, as in $\mathfrak t_3(1;b)$. Some letters have meanings here that are local to this section: $g$ is the
spoke half-exponent, not the coupling; $\sigma=b/2$ is a rim offset, not a spin; $s_b$ is a sine, not the exponent $s$;
$S$ and $C$ are the trigonometric values \eqref{gff:gdr:D}, not $C_n$, $C(\tau)$ or the condition (C); and $\omega(\rho)$
of \eqref{gff:mgf:gammaBounds} is a constant, unrelated to the numbers $\omega_j$ of Step~1 of the proof of
\eqref{gff:gdr:topG}.

\subsection{Finite periodization at a general rim exponent}\label{gff:gdr:periodization}
\paragraph{Kernel.}
For $0<d<4$ and a rim exponent $b$ (complex in general) put
\begin{equation}\begin{gathered}
 N_b=\frac{\Gamma(2q)\Gamma(g)}{\Gamma(\sigma)},\qquad
 a_\ell=\frac{\ell+\sigma}2,\qquad \beta_\ell=a_\ell+g,\qquad
 g_\ell(s;b)=\frac{\Gamma(a_\ell+s)\Gamma(a_\ell-s)}{\Gamma(\beta_\ell+s)\Gamma(\beta_\ell-s)},\\
 W_r(d;b)=\langle t_{\ell,b}^{\,r}\rangle_d
 =N_b^{\,r}\sum_{\ell\geq0}d_\ell\int_{-\ii\infty}^{\ii\infty}\frac{ds}{2\pi\ii}\,g_\ell(s;b)^r .
\end{gathered}\label{gff:gdr:kernel}\end{equation}
The eigenvalue of \eqref{gff:gs:balancedWheel} is $t_{\ell,b}(\nu)=N_b\,g_\ell(\ii\nu;b)$ for real $b$; for complex $b$,
$W_r$ and $\langle t_{\ell,b}^{\,r}\rangle_d$ denote the analytic continuation \eqref{gff:gdr:kernel}, with $|\cdot|^2$
replaced by $\Gamma(a_\ell+\ii\nu)\Gamma(a_\ell-\ii\nu)/(\Gamma(\beta_\ell+\ii\nu)\Gamma(\beta_\ell-\ii\nu))$. Here $2a_\ell=\ell+\sigma$ and
$a_\ell+\beta_\ell=\ell+2q$. At $b=d/2$ one has $\sigma=g=q$, $N_{d/2}=\Gamma(2q)$, $\beta_\ell$ is the $b_\ell$ of
\eqref{gff:fp:kernel}, and \eqref{gff:gdr:kernel} is \eqref{gff:fp:wheelStart}. The multiplicities \eqref{gff:fp:degeneracy} obey
$|d_\ell|\le C(1+\ell)^{d-2}$ and
\begin{equation}
 \sum_{\ell\geq0}d_\ell z^\ell=(1+z)(1-z)^{1-d},\qquad|z|<1 .
 \label{gff:gdr:Dgen}
\end{equation}
For $\operatorname{Re}b>0$ the imaginary axis separates the two pole families of every kernel, and the integral in
\eqref{gff:gdr:kernel} converges absolutely for $2r\operatorname{Re}g>\max(d,1)$. So \eqref{gff:gdr:kernel} represents $W_r$
on $0<\operatorname{Re}b<b_{\rm spec}=d-\max(d,1)/r$.

On the contour $|g_\ell(\ii\nu;b)|$ is comparable to $(1+\ell+|\nu|)^{-2\operatorname{Re}g}$. The
$\nu$-integral therefore converges absolutely for $2r\operatorname{Re}g>1$ and the spin sum for
$2r\operatorname{Re}g>d$; together these are the condition just stated. The integral also converges
at some $\operatorname{Re}b\le0$, but there it no longer represents $W_r$. The function $W_r(d;b)$ is
singular at $b=0$, where the $\ell=0$ contour is pinched (a pole of order $r-1$), and at
$b=b_{\rm spec}$, where the spin or the $\nu$ integration diverges. For $b$ near $\frac d2$ the disc
$|b-\frac d2|<\min(\frac d2,\,b_{\rm spec}-\frac d2)$ lies in this strip of analyticity. The bound
$|d_\ell|\le C(1+\ell)^{d-2}$ holds for every $0<d<4$ and is proved below, \eqref{gff:mgf:multiplicityBound};
for $d<1$ the $d_\ell$ with $\ell\geq2$ are negative.


\paragraph{Periodization and strip.}
Define
\begin{equation}
 S_r(t;b)=\sum_{\ell\geq0}d_\ell\sum_{n\in\mathbb Z}g_\ell(a_\ell+n+t;b)^r,\qquad
 D_b(t)=\sin(\pi t)\sin\bigl(\pi(t+\sigma)\bigr),\qquad
 \Xi_r(t;b)=D_b(t)^rS_r(t;b).
 \label{gff:gdr:periodizationDef}
\end{equation}
At $b=d/2$ these are $S_r,D_q,\Xi_r$ of \eqref{gff:fp:periodization}; the subscript of $D_q$ there is $\sigma=q$, so $D_q$ is $D_b$ at $b=2q$, not at $b=q$. Put $\alpha=r\operatorname{Re}g$ and
\begin{equation}
 b_{\max}(r,d)=\min\Bigl\{d-\frac dr,\;d-\frac{2(d-1)}r,\;d-\frac1r\Bigr\},
 \label{gff:gdr:bmax}
\end{equation}
that is $d-\frac1r$ for $d\le1$, $\frac{(r-1)d}r$ for $1\le d\le2$ and $d-\frac{2(d-1)}r$ for $d\ge2$. The condition
$\operatorname{Re}b<b_{\max}$ is the conjunction of (i) $2\alpha>d$, (ii) $\alpha>d-1$ and (iii) $2\alpha>1$. One has
$b_{\max}\le b_{\rm spec}$, with equality for $d\le2$, and $b=d/2$ lies in $0<b<b_{\max}(r,d)$ iff $\frac2r<d<4$:
$d>\frac12$ for $W_4$ and $d>\frac23$ for $W_3$.

\paragraph{Lattice split and estimates.}
Let $f(z)=\Gamma(z)/\Gamma(z+g)$. For $\ell\geq0$, $n\in\mathbb Z$, $j,k\geq0$ and $t$ off the classes $t\equiv0$,
$t\equiv-\sigma\pmod1$, Gamma reflection \cite{gff:DLMF5} gives
\begin{equation}\begin{aligned}
 g_\ell(a_\ell+n+t)&=f(\ell+\sigma+n+t)\,f(-n-t)\\
 &=-\frac{\sin\pi(g-t)}{\sin\pi t}\,
 \frac{\Gamma(\ell+n+\sigma+t)\,\Gamma(n+1-g+t)}{\Gamma(\ell+n+2q+t)\,\Gamma(n+1+t)},\\
 g_{j+k+1}(a_{j+k+1}-j-1+t)&=\frac{\Gamma(k+\sigma+t)\,\Gamma(j+1-t)}{\Gamma(k+2q+t)\,\Gamma(j+1+g-t)},\\
 g_\ell(a_\ell+n+t)&=g_\ell\bigl(a_\ell+(-\ell-n-1)+(1-\sigma-t)\bigr).
\end{aligned}\label{gff:gdr:reflected}\end{equation}
At $b=d/2$ the first two are \eqref{gff:fp:reflectedRight} and \eqref{gff:fp:middleFactor}. The lattice splits into the
right part $n\geq0$, the middle part $-\ell\le n\le-1$ and the left part $n\le-\ell-1$; by the last identity the left
part is the right part at the mirror point $1-\sigma-t$. The Gamma-ratio bound
$|\Gamma(z+u)/\Gamma(z+v)|\le C|z|^{\operatorname{Re}(u-v)}$ \cite{gff:DLMFStirling} gives, on compact sets avoiding the
pole classes, $|g_\ell(a_\ell+n+t)|^r\le C(1+\ell+n)^{-\alpha}(1+n)^{-\alpha}$ in the right part and
$C(1+k)^{-\alpha}(1+j)^{-\alpha}$ in the middle part. For $t=x+\ii y$ with $x$ bounded and
$|y|\geq y_0>|\operatorname{Im}\sigma|+1$ the same holds with $1+\ell+n,1+n,1+k,1+j$ replaced by
$\ell+n+|y|,\dots,j+|y|$. The constants are uniform for $b$ in compact subsets of $\operatorname{Re}b>0$. Finally
$\sum_{\ell,n\geq0}(1+\ell)^{d-2}(1+\ell+n)^{-\alpha}(1+n)^{-\alpha}<\infty$ iff $\alpha>\max(\frac d2,d-1,\frac12)$,
that is iff $\operatorname{Re}b<b_{\max}$; under this condition the version with $Y$ in place of the $1$'s is
$O(Y^{\max(d,1)-2\alpha}(1+\log Y))$, and the same holds for the middle sum. The paragraph \emph{Proofs of the estimates} below proves them in full, with explicit constants in \eqref{gff:mgf:gammaBounds} and \eqref{gff:mgf:latticeSums}.

\paragraph{Proofs of the estimates.}
\emph{The identities \eqref{gff:gdr:reflected}.} For $s=a_\ell+n+t$ one has $a_\ell+s=\ell+\sigma+n+t$,
$a_\ell-s=-n-t$, $\beta_\ell+s=\ell+\sigma+n+t+g$ and $\beta_\ell-s=g-n-t$, which is the factorized
form. Gamma reflection \cite{gff:DLMF5} gives, for every integer $n$,
\[
 \frac{\Gamma(-n-t)}{\Gamma(g-n-t)}=-\frac{\sin\pi(g-t)}{\sin\pi t}\,
 \frac{\Gamma(n+1-g+t)}{\Gamma(n+1+t)},
\]
and $\ell+\sigma+n+t+g=\ell+2q+n+t$; this is the second line. The substitution $\ell=j+k+1$, $n=-j-1$
in the factorized form gives the third line. Finally $g_\ell$ is even in $s$, and
$-(a_\ell+n+t)=a_\ell+(-\ell-n-1)+(1-\sigma-t)$ because $2a_\ell=\ell+\sigma$.

\emph{The Gamma-ratio bound.} Stirling's formula with its error bound \cite{gff:DLMFStirling} states that
for $|\arg w|<\pi$, with principal branches,
$\operatorname{Ln}\Gamma(w)=(w-\frac12)\operatorname{Ln}w-w+\frac12\ln2\pi+\epsilon_\Gamma(w)$ with
$|\epsilon_\Gamma(w)|\le\sec^2(\frac12\arg w)/(12|w|)$. From it: let $\rho\geq0$ and
$\omega(\rho)=\rho+\frac32\rho^2+\frac{11}3\rho^3+\frac43$. If $\operatorname{Re}z\geq0$,
$|z|\geq\max(1,2\rho)$ and $|u|,|v|\le\rho$, then $\Gamma(z+u)$ and $\Gamma(z+v)$ are finite and
nonzero, and
\begin{equation}\begin{gathered}
 \frac{\Gamma(z+u)}{\Gamma(z+v)}=z^{u-v}e^{\theta},\qquad|\theta|\le\frac{\omega(\rho)}{|z|},\\
 e^{-\omega(\rho)/|z|}e^{-\frac\pi2|\operatorname{Im}(u-v)|}|z|^{\operatorname{Re}(u-v)}
 \le\Bigl|\frac{\Gamma(z+u)}{\Gamma(z+v)}\Bigr|\le
 e^{\omega(\rho)/|z|}e^{\frac\pi2|\operatorname{Im}(u-v)|}|z|^{\operatorname{Re}(u-v)} .
\end{gathered}\label{gff:mgf:gammaBounds}\end{equation}
\emph{Proof.} Put $\xi=u/z$, so $|\xi|\le\frac12$, and $w=z+u=z(1+\xi)$. Then
$|w|\geq|z|/2\geq\frac12$ and $|\arg w|\le|\arg z|+|\arg(1+\xi)|\le\frac\pi2+\frac\pi6$, so $w$ is
neither a pole nor a zero of $\Gamma$, and $|\epsilon_\Gamma(w)|\le\sec^2(\pi/3)/(12|w|)\le2/(3|z|)$.
The imaginary part of $\operatorname{Ln}z+\operatorname{Ln}(1+\xi)$ lies in
$[-\frac{2\pi}3,\frac{2\pi}3]$, so $\operatorname{Ln}w=\operatorname{Ln}z+\operatorname{Ln}(1+\xi)$.
Write $\operatorname{Ln}(1+\xi)=\xi-\xi^2/2+\xi_3$ with
$|\xi_3|\le|\xi|^3/(3(1-|\xi|))\le\frac23|\xi|^3$. Multiplying out with $z\xi=u$,
\[
 \bigl(z+u-\tfrac12\bigr)\operatorname{Ln}(1+\xi)=u+h(z,u),\qquad
 h(z,u)=\frac{u^2}{2z}-\frac{u}{2z}-\frac{u^3}{2z^2}+\frac{u^2}{4z^2}+\bigl(z+u-\tfrac12\bigr)\xi_3 .
\]
With $|u|\le\rho\le|z|/2$, $|z|\geq1$ and $|z+u-\frac12|\le2|z|$, each term is bounded, and
$|h(z,u)|\le(\frac\rho2+\frac34\rho^2+\frac{11}6\rho^3)/|z|$. Hence
$\operatorname{Ln}\Gamma(z+u)=(z+u-\frac12)\operatorname{Ln}z-z+\frac12\ln2\pi+h(z,u)+\epsilon_\Gamma(z+u)$.
Subtracting the same identity for $v$ gives
$\operatorname{Ln}\Gamma(z+u)-\operatorname{Ln}\Gamma(z+v)=(u-v)\operatorname{Ln}z+\theta$ with
$|\theta|\le2(\frac\rho2+\frac34\rho^2+\frac{11}6\rho^3)/|z|+\frac4{3|z|}=\omega(\rho)/|z|$.
Exponentiating gives the first line of \eqref{gff:mgf:gammaBounds}. The second follows from
$|z^{u-v}|=|z|^{\operatorname{Re}(u-v)}e^{-\operatorname{Im}(u-v)\arg z}$ and $|\arg z|\le\frac\pi2$.
$\square$ The bound is a version of the large-$z$ asymptotics of a Gamma ratio \cite{gff:DLMFStirling}
that is uniform in $u,v$ and in the closed right half-plane, with an explicit constant. It is used
below for arguments $m+c+t$ with an integer $m\geq0$ and a bounded constant $c$: one puts
$z=m+t+X$, where $X\geq0$ bounds $|\operatorname{Re}t|$, so that $\operatorname{Re}z\geq m$, and
absorbs $c-X$ into $u$ or $v$; arguments $m+c-t$ are treated in the same way.

\emph{The multiplicities.} For $q\neq\frac14$, \eqref{gff:fp:degeneracy} gives
$d_{k+1}=\frac2{\Gamma(4q-1)}(2q+k)\frac{\Gamma(k+4q-1)}{\Gamma(k+2)}$. By \eqref{gff:mgf:gammaBounds}
with $z=k\geq6$ and $\rho=3$, $\Gamma(k+4q-1)/\Gamma(k+2)=k^{4q-3}e^\theta$ with
$|\theta|\le\omega(3)/k$, uniformly in $q\in(0,1)$. The finitely many $k<6$ are continuous in $q$, and
so is $2/\Gamma(4q-1)$, also at $q=\frac14$, where it vanishes; near $q=\frac14$ one uses the form
\eqref{gff:fp:degeneracy}, with $(4q-1)_k=(4q-1)(4q)_{k-1}$, and at $q=\frac14$ the multiplicities are
$(1,1,0,0,\dots)$. Hence
\begin{equation}
 |d_\ell|\le C_q\,(1+\ell)^{d-2},\qquad d_\ell\,(1+\ell)^{2-d}\to\frac2{\Gamma(d-1)}\quad(\ell\to\infty),
 \label{gff:mgf:multiplicityBound}
\end{equation}
with $C_q=\sup_\ell|d_\ell|(1+\ell)^{2-d}$ finite and bounded on compact subsets of $0<q<1$, across
$d=1$ included (the limit is $0$ at $d=1$).

\emph{The right and middle bounds.} Let $\operatorname{Re}b>0$ and let $K$ be a compact set off the
classes $t\equiv0$ and $t\equiv-\sigma$. In the second line of \eqref{gff:gdr:reflected},
\eqref{gff:mgf:gammaBounds} bounds $\Gamma(\ell+n+\sigma+t)/\Gamma(\ell+n+2q+t)$ by
$C(1+\ell+n)^{-\operatorname{Re}g}$ and $\Gamma(n+1-g+t)/\Gamma(n+1+t)$ by $C(1+n)^{-\operatorname{Re}g}$,
since $\sigma-2q=-g$, for all indices with $\ell+n$, respectively $n$, above a bound fixed by $K$ and
$b$. In the third line it bounds $\Gamma(k+\sigma+t)/\Gamma(k+2q+t)$ and
$\Gamma(j+1-t)/\Gamma(j+1+g-t)$ by $C(1+k)^{-\operatorname{Re}g}$ and $C(1+j)^{-\operatorname{Re}g}$. The
prefactor $\sin\pi(g-t)/\sin\pi t$ is bounded on $K$. Each ratio depends on a single combination of
the summation indices ($\ell+n$, $n$, $k$ or $j$). A ratio with a small index therefore occurs for
infinitely many lattice points but takes only finitely many values; these are continuous on $K$ and
bounded there one by one. Raising to the power $r$ gives the bounds stated above,
$(1+\ell+n)^{-\alpha}(1+n)^{-\alpha}$ and $(1+k)^{-\alpha}(1+j)^{-\alpha}$ with
$\alpha=r\operatorname{Re}g$. Now let $t=x+\ii y$ with $x$ in a fixed bounded interval and
$|y|\geq y_0>|\operatorname{Im}\sigma|+1$. The numerator arguments $\ell+n+\sigma+t$, $n+1-g+t$ and
$k+\sigma+t$ have imaginary parts of modulus at least $|y|-|\operatorname{Im}\sigma|>1$ (note
$\operatorname{Im}g=-\operatorname{Im}\sigma$), and $j+1-t$ has $|\operatorname{Im}|=|y|>1$; so no
Gamma function in the numerators meets a pole, and the prefactor is bounded. With $z$ chosen as above,
$|z|$ is comparable to the index plus $|y|$. For $|y|$ or the index above $2\rho+1$, where $\rho$
bounds the shifts, \eqref{gff:mgf:gammaBounds} applies directly; the remaining factor values, with small
index and bounded $|y|\geq y_0$, are finitely many continuous functions on a compact set that meets no
pole. This gives the bounds with $1+\ell+n$, $1+n$, $1+k$, $1+j$ replaced by $\ell+n+|y|$, $n+|y|$,
$k+|y|$, $j+|y|$. For complex $b$ this is why $y_0$ must exceed $|\operatorname{Im}\sigma|+1$: the class
$-\sigma+\mathbb Z$ lies on $\operatorname{Im}t=-\operatorname{Im}\sigma$. All quantities are continuous
in $b$, so the constants are uniform for $b$ in a compact set $B$ of $\operatorname{Re}b>0$, provided
$K$ avoids the pole classes of every $b\in B$, $y_0>1+\max_B|\operatorname{Im}\sigma|$, and $\alpha$ is
replaced by $r\min_B\operatorname{Re}g$.

\emph{The lattice sums.} Let $0<d<4$ and $\alpha>\max(\frac d2,d-1,\frac12)$, and put
$\bar\alpha=2\alpha-\max(d-1,0)$, so that $\bar\alpha>1$ and $1-\bar\alpha=\max(d,1)-2\alpha$. Then for
$Y\geq1$
\begin{equation}\begin{gathered}
 \Sigma^{\rm r}(Y)=\sum_{\ell,n\geq0}(1+\ell)^{d-2}(\ell+n+Y)^{-\alpha}(n+Y)^{-\alpha}
 \le C^{\rm r}\,Y^{\max(d,1)-2\alpha}(1+\log Y),\\
 \Sigma^{\rm m}(Y)=\sum_{j,k\geq0}(1+j+k)^{d-2}(k+Y)^{-\alpha}(j+Y)^{-\alpha}\le C^{\rm m}\,Y^{d-2\alpha},
\end{gathered}\label{gff:mgf:latticeSums}\end{equation}
where, with $C_1=1+2^{\alpha-d+1}/(\alpha-d+1)$,
\[
 C^{\rm r}=C_1\Bigl[1+\frac1{\bar\alpha-1}+\frac1{(\bar\alpha-1)^2}\Bigr],\qquad
 C^{\rm m}=\begin{cases}2^{d-1}C_1\bigl(1+\frac1{\alpha-1}\bigr),&d\geq2,\\[2pt]
 \Bigl[\frac2d+\frac{2^{\alpha-d/2}}{\alpha-d/2}\Bigr]^2,&d<2.\end{cases}
\]
The two expressions for $C^{\rm m}$ are each continuous on their closed ranges, and $C_1$ and $C^{\rm r}$
are continuous on the open set of $(d,\alpha)$ defined by the condition; so all the constants are
bounded on its compact subsets, across $d=1$ and $d=2$ included. Since $Y^{d-2\alpha}\le
Y^{\max(d,1)-2\alpha}(1+\log Y)$ for $Y\geq1$, the middle sum obeys the bound of the right sum as
well. The left part of a lattice sum, with $k$ in place of $n$, has the form of the right one.
\emph{Proof.} (i) For $M\geq1$ and an exponent $e\in(-2,2)$ with $\alpha>e+1$ put
$\Sigma_e(M)=\sum_{\ell\geq0}(1+\ell)^e(\ell+M)^{-\alpha}$ and $u=1+\ell$. Since
$\ell+M\geq\max(u,M)$, $\Sigma_e(M)\le M^{-\alpha}\sum_{1\le u\le M}u^e+\sum_{u>M}u^{e-\alpha}$. Here
$\sum_{1\le u\le M}u^e\le M^{\max(e+1,0)}(1+\log M)$: for $e>0$ each term is at most $M^e$; for
$-1\le e\le0$, $u^e\le M^{e+1}/u$; for $e<-1$, $u^e\le1/u$; and $\sum_{u\le M}1/u\le1+\log M$. Since
$u^{e-\alpha}$ decreases, $\sum_{u>M}u^{e-\alpha}\le\int_{\lfloor M\rfloor}^\infty v^{e-\alpha}dv
\le2^{\alpha-e-1}M^{e+1-\alpha}/(\alpha-e-1)$, using $\lfloor M\rfloor\geq M/2$. With $e=d-2$ this
gives $\Sigma_{d-2}(M)\le C_1M^{\max(d-1,0)-\alpha}(1+\log M)$, and for $d\geq2$ the same without the
logarithm. (ii) The right sum is $\Sigma^{\rm r}(Y)=\sum_{n\geq0}(n+Y)^{-\alpha}\Sigma_{d-2}(n+Y)\le
C_1\sum_{n\geq0}(n+Y)^{-\bar\alpha}(1+\log(n+Y))$. The function $v\mapsto v^{-\bar\alpha}(1+\log v)$
decreases on $[1,\infty)$ because $\bar\alpha\geq1$, so the sum is at most
\[
 Y^{-\bar\alpha}(1+\log Y)+\int_Y^\infty v^{-\bar\alpha}(1+\log v)\,dv
 =Y^{-\bar\alpha}(1+\log Y)+Y^{1-\bar\alpha}\Bigl[\frac{1+\log Y}{\bar\alpha-1}+\frac1{(\bar\alpha-1)^2}\Bigr],
\]
which is at most $[1+\frac1{\bar\alpha-1}+\frac1{(\bar\alpha-1)^2}]\,Y^{1-\bar\alpha}(1+\log Y)$; times $C_1$ this is the bound with $C^{\rm r}$. (iii) The middle sum for $d\geq2$: by
$(1+j+k)^{d-2}\le2^{d-2}[(1+j)^{d-2}+(1+k)^{d-2}]$,
$\Sigma^{\rm m}(Y)\le2^{d-1}\Sigma_{d-2}(Y)\sum_k(k+Y)^{-\alpha}$, with $\Sigma_{d-2}(Y)\le
C_1Y^{d-1-\alpha}$ and $\sum_k(k+Y)^{-\alpha}\le(1+\frac1{\alpha-1})Y^{1-\alpha}$, because
$\alpha>d-1\geq1$. (iv) The middle sum for $d<2$: here $1+j+k\geq\sqrt{(1+j)(1+k)}$, so
$(1+j+k)^{d-2}\le[(1+j)(1+k)]^{(d-2)/2}$ and $\Sigma^{\rm m}(Y)\le\Sigma_{(d-2)/2}(Y)^2$. Since
$(d-2)/2\in(-1,0)$, $\sum_{1\le u\le Y}u^{(d-2)/2}\le Y^{d/2}/(d/2)$, and the tail needs
$\alpha>d/2$; this gives the stated $C^{\rm m}$. $\square$

\emph{The condition is also necessary.} At $Y=1$, \eqref{gff:mgf:latticeSums} gives
$\sum_{\ell,n\geq0}(1+\ell)^{d-2}(1+\ell+n)^{-\alpha}(1+n)^{-\alpha}<\infty$ for
$\alpha>\max(\frac d2,d-1,\frac12)$; the converse also holds. Split the sum into $n\le\ell$ and
$n>\ell$. For $n\le\ell$, $1+\ell+n$ is comparable to $1+\ell$, and $\sum_{n\le\ell}(1+n)^{-\alpha}$ to
$1$, $1+\log(1+\ell)$ or $(1+\ell)^{1-\alpha}$ for $\alpha>1$, $\alpha=1$, $\alpha<1$; this part
converges iff $\alpha>d-1$ (when $\alpha>1$), $d<2$ (when $\alpha=1$) or $2\alpha>d$ (when $\alpha<1$).
For $n>\ell$, $1+\ell+n$ is comparable to $1+n$, and $\sum_{\ell<n}(1+\ell)^{d-2}$ to $n^{d-1}$,
$1+\log n$ or $1$ for $d>1$, $d=1$, $d<1$; this part converges iff $2\alpha>\max(d,1)$. The
conjunction of the two conditions is $\alpha>\max(\frac d2,d-1,\frac12)$: the second gives the first
and last entries, and if it holds, the first part gives $\alpha>d-1$ in each of its three cases (for
$\alpha<1$ the condition $2\alpha>d$ forces $d<2$, so $\alpha>\frac d2>d-1$).

\emph{The entries of $b_{\max}$.} With $\alpha=r\operatorname{Re}g=r(d-\operatorname{Re}b)/2$ the
condition reads $\operatorname{Re}b<b_{\max}(r,d)$. Its three entries come from three regions of the
lattice: (i) $2\alpha>d$ from the region where $\ell$ and $n$ are both large, (ii) $\alpha>d-1$ from
the end $\ell\gg n$, and (iii) $2\alpha>1$ from the end $n\gg\ell$. For $r=4$, \eqref{gff:gdr:bmax} is
$d-\frac14$, $\frac34d$, $\frac{d+1}2$ on the ranges $d\le1$, $1\le d\le2$, $d\geq2$; for $r=3$ it is
$d-\frac13$, $\frac23d$, $\frac{d+2}3$. Some values, with the spectral bound of \eqref{gff:gdr:kernel}:
\begin{center}\small
\begin{tabular}{@{}lcccccccc@{}}\toprule
$d$ & 0.8 & 1 & 1.25 & 1.5 & 2 & 2.5 & 2.75 & 3\\\midrule
$b_{\max}(3,d)$ & 0.4667 & 0.6667 & 0.8333 & 1 & 1.3333 & 1.5 & 1.5833 & 1.6667\\
$b_{\rm spec}(3,d)$ & 0.4667 & 0.6667 & 0.8333 & 1 & 1.3333 & 1.6667 & 1.8333 & 2\\
$b_{\max}(4,d)$ & 0.55 & 0.75 & 0.9375 & 1.125 & 1.5 & 1.75 & 1.875 & 2\\
$b_{\rm spec}(4,d)$ & 0.55 & 0.75 & 0.9375 & 1.125 & 1.5 & 1.875 & 2.0625 & 2.25\\\bottomrule
\end{tabular}
\end{center}
The two bounds differ only for $d>2$ (paragraph \emph{Beyond $b_{\max}$} of Section~\ref{gff:mgf:nodes}).
At $d=2$ and $r=4$, $b_{\max}=\frac32$ is $x=\frac b2<\frac34$ for the one-variable sum
$\Omega(t;x)$ of \eqref{gff:gc:L2modes}, a range that contains the interval used in
Section~\ref{gff:gc:L2proof}; at $d=3$, $b_{\max}(4,3)=2$ is the range ``every $b<2$'' of the $d=3$
paragraph of Section~\ref{gff:gdr:special}. For $r=3$, $b_{\max}(3,2)=\frac43$ is the range of
\eqref{gff:gs:W3d2}, and $b_{\max}(3,3)=\frac53$. The exponent $\eta$ of part (c) of
Section~\ref{gff:gdr:finite} is $2\alpha-\max(d,1)>0$ throughout the strip, by (i) and (iii), so the
exclusion of the boundary mode needs no extra condition.


\subsubsection{Finite Fourier space}\label{gff:gdr:finite}
Let $r\geq2$, $0<d<4$ and $0<\operatorname{Re}b<b_{\max}(r,d)$, $\sigma\notin\mathbb Z$. Then:
\begin{enumerate}
\item[(a)] $S_r(t;b)$ converges absolutely and locally uniformly off the classes $t\equiv0$ and $t\equiv-\sigma$; it
is one-periodic and meromorphic, with poles of order at most $r$ in these two classes only.
\item[(b)] $\Xi_r(t;b)$ is entire, $\Xi_r(t+1;b)=\Xi_r(t;b)$ and $\Xi_r(-\sigma-t;b)=\Xi_r(t;b)$.
\item[(c)] For $0\le x\le1$ and $|y|\geq2+|\operatorname{Im}b|$,
$|\Xi_r(x+\ii y;b)|\le C\,e^{2\pi r|y|}|y|^{-\eta}(1+\log|y|)$ with $\eta=r(d-\operatorname{Re}b)-\max(d,1)>0$ (a decay exponent, unrelated to the weights $\eta_m$ of \eqref{gff:gs:etaTheta}).
\item[(d)] There are coefficients $c_{r,k}(q,b)$, analytic in $b$ on the strip and real for real $b$, with
\begin{equation}
 \Xi_r(t;b)=\sum_{k=0}^{r-1}c_{r,k}(q,b)\cos\bigl(2\pi k(t+\tfrac b4)\bigr).
 \label{gff:gdr:finiteFourier}
\end{equation}
$\Xi_r(t;b)$ is jointly analytic in $(t,b)$, so every sample at a fixed node is analytic in $b$.
\end{enumerate}
\emph{Proof.} (a) By the right and
middle bounds and \eqref{gff:mgf:latticeSums}, the right and middle parts converge absolutely and
uniformly on compact sets off the pole classes, and so does the left part, which is the right part at
the mirror point. Each term is meromorphic in $t$. By \eqref{gff:gdr:reflected} poles come only from
$\Gamma(-n-t)$ ($t\in\mathbb Z$) and $\Gamma(\ell+\sigma+n+t)$ ($t\in-\sigma+\mathbb Z$), each to order
at most $r$. If $t_0$ is such a point, $(t-t_0)^r$ times each term is analytic in a small disc around
$t_0$; the series of these functions converges uniformly on the boundary circle, hence, by the maximum
principle, in the disc, so $(t-t_0)^rS_r$ is analytic near $t_0$. Periodicity follows from $n\mapsto
n+1$ in an absolutely convergent sum. (b) $D_b^r$ vanishes to order $r$ on both classes, so $\Xi_r$ is
entire, and $D_b(t+1)=D_b(t)$. Evenness of $g_\ell$ gives
$g_\ell(a_\ell+n-\sigma-t)=g_\ell(a_\ell+(-\ell-n)+t)$, and $n\mapsto-\ell-n$ is a bijection of
$\mathbb Z$, so $S_r(-\sigma-t)=S_r(t)$; also $D_b(-\sigma-t)=\sin\pi(\sigma+t)\sin\pi t=D_b(t)$. (c)
$|D_b(x+\ii y)|\le Ce^{2\pi|y|}$. For $|y|\geq2+|\operatorname{Im}b|=2+2|\operatorname{Im}\sigma|$ the
large-$|y|$ bounds apply at $t$ and, for the left part, at the mirror point $1-\sigma-t$, whose
imaginary part has modulus at least $|y|-|\operatorname{Im}\sigma|\geq\max(|y|/2,\,2+|\operatorname{Im}\sigma|)$.
With \eqref{gff:mgf:latticeSums} at $Y=|y|$ (or $|y|/2$), each of the three parts of $S_r$ is bounded by
$C|y|^{\max(d,1)-2\alpha}(1+\log|y|)$, and $\eta=2\alpha-\max(d,1)$. (d) As for \eqref{gff:fp:finiteFourier}, the coefficient
$\gamma_m=\int_0^1\Xi_r(x+\ii y)e^{-2\pi\ii m(x+\ii y)}dx$ does not depend on $y$ (Cauchy's theorem on
a rectangle, the vertical sides cancelling by periodicity). For $m\geq r$ take $y=-Y$ with
$Y\geq2+|\operatorname{Im}b|$; by (c), $|\gamma_m|\le Ce^{2\pi(r-m)Y}Y^{-\eta}(1+\log Y)\to0$, also for
$m=r$ since $\eta>0$. For $m\le-r$ take $y=+Y$. So $\Xi_r=\sum_{|m|\le r-1}\gamma_me^{2\pi\ii mt}$. The
reflection in (b), $\Xi_r(-\sigma-t)=\sum_m\gamma_{-m}e^{2\pi\ii m\sigma}e^{2\pi\ii mt}$, gives
$\gamma_m=\gamma_{-m}e^{2\pi\ii m\sigma}$. With $c_{r,0}=\gamma_0$ and
$c_{r,k}=2\gamma_ke^{-\ii\pi k\sigma}=2\gamma_{-k}e^{\ii\pi k\sigma}$ one has
$\gamma_ke^{2\pi\ii kt}+\gamma_{-k}e^{-2\pi\ii kt}=c_{r,k}\cos(2\pi k(t+\frac b4))$, which is
\eqref{gff:gdr:finiteFourier}. For the analyticity in $b$ let $B$ be a compact subset of the strip and
$y_B=2+\max_B|\operatorname{Im}b|$. On the segment $t=x+\ii y_B$, $0\le x\le1$, the series for $S_r$
converges uniformly in $(x,b)\in[0,1]\times B$, by the large-$|y|$ bounds with constants uniform on $B$
and \eqref{gff:mgf:latticeSums} with $\alpha=r\min_B\operatorname{Re}g>\max(\frac d2,d-1,\frac12)$. Its
terms are analytic in $b$ there, since the segment stays at distance at least one from both pole
classes, and $D_b(t)^r$ is entire in $(t,b)$. Hence $\Xi_r(x+\ii y_B;b)$ is analytic in $b$, and so is
$\gamma_m(b)=\int_0^1\Xi_r(x+\ii y_B;b)e^{-2\pi\ii m(x+\ii y_B)}dx$ on the interior of $B$. The same
argument, with the Cauchy formula in $t$ on small circles around the pole classes, shows that
$\Xi_r(t;b)$ is jointly analytic in $(t,b)$. For real $b$, $\Xi_r$ is real on the real axis, so
$\gamma_{-m}=\overline{\gamma_m}$ and the $c_{r,k}$ are real. $\square$ (The point $\sigma\in\mathbb Z$ lies in
the strip only at $b=2$ when $r\geq3$ and $d>(2r-2)/(r-2)$, never for $r=3$ or for $r=4$, $d\le3$; it is treated in
the paragraph \emph{Integer $\sigma$} below. The masters use only $b$ near $\frac d2$, where $0<\sigma<1$.)

\paragraph{Integer $\sigma$.}
In the strip, $\sigma\in\mathbb Z$ can occur only at the real point $b=2$, because $b_{\max}\le
d-\frac dr<d<4$. The point $b=2$ lies in the strip iff $b_{\max}(r,d)>2$, that is iff $r\geq3$ and
$d>(2r-2)/(r-2)$. This never happens for $r=3$; it happens for $d>3$ when $r=4$, for $d>\frac83$ when
$r=5$, and for $d>(2r-2)/(r-2)$, which decreases to $2$, for larger $r$. At $b=2$ the two pole classes
merge. $\Xi_r$ is still entire: the poles of $S_r$ in $\mathbb Z$ have order at most $2r$, and
$D_2(t)^r=(-1)^r\sin^{2r}\pi t$. The analyticity argument in (d) does not use $\sigma\notin\mathbb Z$,
so the $c_{r,k}$ are analytic across $b=2$, and the finite Fourier statement holds there as well. The
top-mode statement \eqref{gff:gdr:topless} holds there by continuity, and also directly, since
\eqref{gff:gdr:topG} allows every $\sigma\in\mathbb C$. Only the weights of \eqref{gff:gdr:functional} and of
\eqref{gff:mgf:Wsample}, which carry $(\sin\pi\sigma)^{-r}$, are singular at $b=2$. $W_r$ is analytic
there, so these poles cancel in the sum (paragraph \emph{Radius of the sample jets} of
Section~\ref{gff:mgf:nodes}). None of the cases used in this note ($r=3,4$ with $d\le3$) contains $b=2$ in
its strip.


\subsubsection{The top modes vanish}\label{gff:gdr:top}
The top-mode statement holds for every product of balanced kernels that satisfies the convergence condition (C)
below, with different spoke exponents and independent shifts of the left and right Gamma arguments; the wheels are the
case of equal factors. Write
$f_g(w)=\Gamma(w)/\Gamma(w+g)$, so that $f=f_g$ in \eqref{gff:gdr:reflected}. In this subsection $g_i,\alpha_i,\beta_i$
($i=1,\dots,r$) are numbers attached to the $i$-th factor: the kernel $g_\ell(s;b)$ is always written with its arguments,
and $\alpha_i,\beta_i$ are unrelated to $\alpha=r\operatorname{Re}g$ and to $\beta_\ell$ in \eqref{gff:gdr:kernel}. Let $0<d<4$,
$r\geq2$ and $\sigma\in\mathbb C$, and let $g_i,\alpha_i,\beta_i\in\mathbb C$ satisfy
\begin{equation}
 \text{(B)}\ \ \sigma+\alpha_i+\beta_i+g_i=\frac d2\ \ (i=1,\dots,r),\qquad
 \text{(C)}\ \ \varkappa:=\sum_{i=1}^r\operatorname{Re}g_i>\max\Bigl(\frac d2,\,d-1,\,\frac12\Bigr).
 \label{gff:gdr:BC}
\end{equation}
Put
\begin{equation}\begin{gathered}
 T_\ell(z)=\prod_{i=1}^rf_{g_i}(\ell+\sigma+\alpha_i+z)\,f_{g_i}(\beta_i-z),\qquad
 S_T(t)=\sum_{\ell\geq0}d_\ell\sum_{n\in\mathbb Z}T_\ell(n+t),\\
 D_T(t)=\prod_{i=1}^r\sin\pi(t+\sigma+\alpha_i)\,\sin\pi(t-\beta_i),\qquad \Xi_T=D_TS_T .
\end{gathered}\label{gff:gdr:TG}\end{equation}
Here $T_\ell(z)$, with one index, is unrelated to $T_{\ell n}(b)$ of \eqref{gff:gdr:Tln}. Then $S_T$ converges absolutely and locally uniformly off the $2r$ pole classes $-\sigma-\alpha_i+\mathbb Z$ and
$\beta_i+\mathbb Z$; $\Xi_T$ is entire and one-periodic, $\Xi_T=\sum_{|m|\le r-1}\gamma_me^{2\pi\ii mt}$, with $\gamma_m$ analytic in
$(\sigma,\alpha_i,g_i)$ on the convex set defined by (C) ($\beta_i$ fixed by (B)); and
\begin{equation}\boxed{
 \gamma_{1-r}=\gamma_{r-1}=0,\qquad \Xi_T(t)=\sum_{|m|\le r-2}\gamma_m\,e^{2\pi\ii mt},}
 \label{gff:gdr:topG}
\end{equation}
equivalently $S_T(t)=O(e^{-4\pi|\operatorname{Im}t|})$ as $|\operatorname{Im}t|\to\infty$. Individual $\operatorname{Re}g_i$ may
be zero or negative; only their sum is constrained. It is proved below, and Section~\ref{gff:gdk:topmode} applies it to the $R$ family of $\KK_d$.

\paragraph{Mixed wheels and wheels.}
A product $\prod_ig_\ell(s+\delta_i;b_i)$ of kernels \eqref{gff:gdr:kernel} with rim exponents $b_i$ and shifts $\delta_i$,
sampled on the lattice $s=\frac{\ell+\sigma}2+n+t$ with any offset $\sigma$, is $T_\ell(n+t)$ with
\begin{equation}
 g_i=\frac{d-b_i}2,\qquad \alpha_i=\frac12\Bigl(\frac{b_i}2-\sigma\Bigr)+\delta_i,\qquad
 \beta_i=\frac12\Bigl(\frac{b_i}2-\sigma\Bigr)-\delta_i ,
 \label{gff:gdr:mixed}
\end{equation}
since $g_\ell(s;b)=f_{(d-b)/2}(\frac\ell2+\frac b4+s)\,f_{(d-b)/2}(\frac\ell2+\frac b4-s)$. Then (B) holds automatically,
(C) says that the mean rim exponent satisfies $\frac1r\sum_i\operatorname{Re}b_i<b_{\max}(r,d)$, and the shifts are free.
For the wheels, $b_i=b$, $\delta_i=0$ and $\sigma=b/2$, one has $\alpha_i=\beta_i=0$, $g_i=g$,
$T_\ell(n+t)=g_\ell(a_\ell+n+t;b)^r$ and $D_T=D_b^r$, so $\Xi_T=\Xi_r$, $\varkappa=\alpha$, and (C) is
$\operatorname{Re}b<b_{\max}$. With $c_{r,r-1}=2\gamma_{r-1}e^{-\ii\pi(r-1)\sigma}$ this gives, under the hypotheses of
Section~\ref{gff:gdr:finite},
\begin{equation}\boxed{
 c_{r,r-1}(q,b)=0,\qquad
 \Xi_r(t;b)=\sum_{k=0}^{r-2}c_{r,k}(q,b)\cos\bigl(2\pi k(t+\tfrac b4)\bigr),}
 \label{gff:gdr:topless}
\end{equation}
and equivalently $S_r(t;b)=O(e^{-4\pi|\operatorname{Im}t|})$. This wheel statement is the one used in the rest of this
section; Section~\ref{gff:gdk:section} uses the product case $r=3$ and the case $r=2$ of \eqref{gff:gdr:topG}.

\paragraph{The $R$ family of $\KK_d$.}
Section~\ref{gff:gdk:topmode} applies \eqref{gff:gdr:topG} to the $R$ family of $\KK_d$, a product of three balanced kernels
that is not a wheel. On the $\tau$-circle used for $\KK_d$ that application rests in addition on the continuation
hypothesis of Section~\ref{gff:gdk:continuation}; the separate lattice-sum code \path{LRI_C5_general_d/review/lemmaG-numeric} of Section~\ref{gff:gdr:checks}
confirms it at two points of that kind (top modes below $2\cdot10^{-47}$ at 54 digits).

\emph{Proof of \eqref{gff:gdr:topG}.} Its steps are as follows; the paragraph \emph{Details of the top-mode proof} after the proof gives in full the estimates and computations that the steps state in short form. Let $Y_0=1+\max_i\max(|\operatorname{Im}(\sigma+\alpha_i)|,|\operatorname{Im}\beta_i|)$, so that all pole
classes lie in $|\operatorname{Im}t|\le Y_0-1$.

\emph{Step 0 (estimates and finite modes).} Split the lattice into the right part $n\geq0$, the middle part $\ell=j+k+1$,
$n=-j-1$, and the left part $n=-\ell-1-k$ ($j,k\geq0$); there is no mirror map in general. Gamma reflection writes the
factors with large negative argument as
$f_g(\beta-n-t)=\frac{\sin\pi(\beta+g-t)}{\sin\pi(\beta-t)}\frac{\Gamma(n+1-\beta-g+t)}{\Gamma(n+1-\beta+t)}$ in the right part and
$f_g(\sigma+\alpha-1-k+t)=\frac{\sin\pi(\sigma+\alpha+g+t)}{\sin\pi(\sigma+\alpha+t)}\frac{\Gamma(k+2-\sigma-\alpha-g-t)}{\Gamma(k+2-\sigma-\alpha-t)}$
in the left part. The Gamma-ratio bound applied to every factor gives a power $-\operatorname{Re}g_i$ of a base common to
all factors, so the exponents add up to $-\varkappa$, and the lattice sums of Section~\ref{gff:gdr:periodization} apply with
$\alpha=\varkappa$. Step~0 of the paragraph \emph{Details of the top-mode proof} carries this out: $S_T$ converges
absolutely and locally uniformly off the pole classes; for $0\le x\le1$ and $|y|\geq\max(Y_0,2)$,
$|S_T(x+\ii y)|\le C|y|^{\max(d,1)-2\varkappa}(1+\log|y|)$; $\Xi_T$ is entire and one-periodic, its modes $|m|\geq r$
vanish, and the $\gamma_m$ are analytic in the parameters.

\emph{Step 1 (two half-planes).} On $\operatorname{Im}t>Y_0$, $S_T$ is analytic, one-periodic and tends to zero, so
$S_T=\sum_{m\geq1}s_m^+w^m$ with $w=e^{2\pi\ii t}$; on $\operatorname{Im}t<-Y_0$, $S_T=\sum_{m\geq1}s_m^-w^{-m}$. Let $\omega_j$ run over the $2r$
numbers $\sigma+\alpha_i$ and $-\beta_i$. Since
$\sin\pi(t+\omega)=\frac\ii2e^{-\ii\pi(t+\omega)}(1-e^{2\pi\ii\omega}w)=-\frac\ii2e^{\ii\pi(t+\omega)}(1-e^{-2\pi\ii\omega}w^{-1})$, one has
$D_T=\kappa_+w^{-r}\prod_j(1-e^{2\pi\ii\omega_j}w)=\kappa_-w^r\prod_j(1-e^{-2\pi\ii\omega_j}w^{-1})$ with
$\kappa_\pm=(-\frac14)^re^{\mp\ii\pi\sum_j\omega_j}$, and comparing coefficients gives
\[
 \gamma_{1-r}=\kappa_+s_1^+,\qquad\gamma_{r-1}=\kappa_-s_1^- .
\]
For the wheels the reflection in Section~\ref{gff:gdr:finite}(b) relates the two. In general there is no such symmetry at
complex parameters, and the two coefficients are treated separately.

\emph{Step 2 (unfolding).} For $Y>Y_0$ the double series converges absolutely on the segment $x+\ii Y$, $0\le x\le1$ (this
uses (C)), so
\[
 s_1^+=\int_0^1S_T(x+\ii Y)\,e^{-2\pi\ii(x+\ii Y)}\,dx=\sum_{\ell\geq0}d_\ell I_\ell,\qquad
 I_\ell=\int_{\mathbb R+\ii Y}T_\ell(z)\,e^{-2\pi\ii z}\,dz,\qquad \sum_\ell|d_\ell||I_\ell|<\infty .
\]

\emph{Step 3 (Laplace form and $I_\ell$).} Assume first $\operatorname{Re}g_i>0$ for all $i$ (removed in Step~7). Rotating the
ray of Euler's integral $\Gamma(g)f_g(w)=\int_0^\infty e^{-wu}(1-e^{-u})^{g-1}du$, with $u=\mp\ii v$ through the first or the
fourth quadrant, gives, first for $\operatorname{Re}w>0$ and then by analytic continuation,
\begin{equation}\begin{aligned}
 f_g(w)&=\frac{-\ii}{\Gamma(g)}\int_0^\infty e^{\ii wv}\bigl(1-e^{\ii v}\bigr)^{g-1}dv&&(\operatorname{Im}w>0),\\
 f_g(w)&=\frac{\ii}{\Gamma(g)}\int_0^\infty e^{-\ii wv}\bigl(1-e^{-\ii v}\bigr)^{g-1}dv&&(\operatorname{Im}w<0),
\end{aligned}\label{gff:gdr:laplace}\end{equation}
with principal powers. The singularities at $v\in2\pi\mathbb N$ are integrable, since
$|(1-e^{\pm\ii v})^{g-1}|\le e^{\pi|\operatorname{Im}g|/2}|2\sin\frac v2|^{\operatorname{Re}g-1}$. On $\operatorname{Im}z=Y$ every left argument
$\ell+\sigma+\alpha_i+z$ has imaginary part at least $1$ and every right argument $\beta_i-z$ at most $-1$. Let $a_\pm$ be the
convolutions on $[0,\infty)$ of $\mathsf h_i^+(v)=e^{\ii\alpha_iv}(1-e^{\ii v})^{g_i-1}$ and of $\mathsf h_i^-(v)=e^{-\ii\beta_iv}(1-e^{-\ii v})^{g_i-1}$,
that is $a_\pm(V)=\int_{\Delta_V}\prod_i\mathsf h_i^\pm(v_i)\,d\mu$ with $\Delta_V=\{v_i\geq0,\sum_iv_i=V\}$. By Fubini,
$\prod_if_{g_i}(\ell+\sigma+\alpha_i+z)=c_+\int_0^\infty e^{\ii(\ell+\sigma+z)V}a_+(V)\,dV$ and
$\prod_if_{g_i}(\beta_i-z)=c_-\int_0^\infty e^{\ii zV}a_-(V)\,dV$, where $c_\pm=\prod_i(\mp\ii)/\Gamma(g_i)$ and
$c_+c_-=\kappa_\Gamma:=\prod_i\Gamma(g_i)^{-2}$.

\emph{Endpoint behavior.} Fix $0<\delta<\frac\pi2$, and let $\lambda(V)=1+|\log|2\pi-V||$ if $\varkappa=1$ and $\lambda=1$
otherwise (this $\lambda$ is local to Step~3 and unrelated to the weights $\lambda_m$ of Section~\ref{gff:gdr:taylor}). Then
\begin{equation}
 |a_\pm(V)|\le\begin{cases}C\,V^{\varkappa-1},&0<V\le\pi,\\
 C\,|2\pi-V|^{\min(\varkappa-1,0)}\lambda(V),&\pi\le V\le2\pi+\delta.\end{cases}
 \label{gff:gdr:endpoint}
\end{equation}
The paragraph \emph{Details of the top-mode proof} proves \eqref{gff:gdr:endpoint} from
$|\mathsf h_i^\pm(v)|\le C|2\sin\frac v2|^{\operatorname{Re}g_i-1}$ and Dirichlet's integral. For $\varkappa>1$ the
bound at $2\pi$ is attained in general, and at $\varkappa=1$ it is at most logarithmic. Hence, near $V=0$,
\begin{equation}
 a_+(V)\,a_-(2\pi-V)=O\bigl(V^{\,\varkappa-1+\min(\varkappa-1,0)}\bigr)\qquad(\text{times at most }1+|\log V|\text{ at }\varkappa=1),
 \label{gff:gdr:endprod}
\end{equation}
and symmetrically near $V=2\pi$. The exponent exceeds $-1$ iff $2\varkappa>1$.

Along $z=x+\ii Y$ the two products are $c_\pm$ times the Fourier transforms $\widehat{a^Y_\pm}(x)=\int e^{\ii xV}a^Y_\pm(V)\,dV$ of
$a^Y_+(V)=e^{\ii(\ell+\sigma)V-YV}a_+(V)$ and $a^Y_-(V)=e^{-YV}a_-(V)$, extended by zero to $V\le0$; both lie in $L^1(\mathbb R)$.
So $T_\ell(z)=\kappa_\Gamma\widehat{a^Y_+*a^Y_-}(x)$. Since $|f_{g_i}(w)|\le C(1+|w|)^{-\operatorname{Re}g_i}$ on horizontal lines at
distance at least one from the real axis, $\widehat{a^Y_+*a^Y_-}=O(|x|^{-2\varkappa})$ is integrable by $2\varkappa>1$. On
$[0,2\pi+\delta]$ the functions $a^Y_\pm$ are square integrable by \eqref{gff:gdr:endpoint} and $2\varkappa>1$, so $a^Y_+*a^Y_-$ is
continuous at $V=2\pi$. Fourier inversion at this point gives
\begin{equation}
 I_\ell=2\pi\kappa_\Gamma\int_0^{2\pi}e^{\ii(\ell+\sigma)V}a_+(V)\,a_-(2\pi-V)\,dV
 =2\pi\kappa_\Gamma\int_{\widetilde\Delta}e^{\ii(\ell+\sigma)V}\prod_{i=1}^r\mathsf h_i^+(v_i)\,\mathsf h_i^-(v_i')\,d\mathrm{vol},
 \label{gff:gdr:Il}
\end{equation}
where $\widetilde\Delta=\{(v,v')\in[0,\infty)^{2r}:V+V'=2\pi\}$, $V=\sum_iv_i$, $V'=\sum_iv_i'$, and
$d\mathrm{vol}=dv_1\cdots dv_r\,dv'_1\cdots dv'_{r-1}$ is its Lebesgue measure, invariant under permutations of the $2r$
coordinates.

\emph{Step 4 (spin sum).} By Step~2 and Abel's theorem, $s_1^+=\lim_{\rho\uparrow1}\sum_\ell d_\ell\rho^\ell I_\ell$. For $\rho<1$
the sum may be taken inside \eqref{gff:gdr:Il}, since $a_+(V)a_-(2\pi-V)\in L^1(0,2\pi)$ by \eqref{gff:gdr:endprod}, and it
produces \eqref{gff:gdr:Dgen} at $z=\rho e^{\ii V}$. For $\frac12\le\rho<1$, $|1-\rho e^{\ii V}|\geq\sqrt2|\sin\frac V2|$, so the integrand
is dominated by $C|\sin\frac V2|^{-\max(d-1,0)}|a_+(V)a_-(2\pi-V)|$. This function is bounded on compact subsets of
$(0,2\pi)$ and is $O(V^{\,\varkappa-1+\min(\varkappa-1,0)-\max(d-1,0)})$ at both ends (with $2\pi-V$ at the upper end and at
most a logarithm at $\varkappa=1$). The bound is integrable iff $\varkappa+\min(\varkappa-1,0)>\max(d-1,0)$, which is exactly (C): for
$\varkappa\geq1$ it reads $\varkappa>d-1$, and for $\varkappa<1$ it reads $2\varkappa>\max(d,1)$, which forces $d<2$ and hence
$\varkappa>\frac d2>d-1$. For the wheels this is the full hypothesis $b<b_{\max}$, not only $b<b_{\rm spec}$ (Step~2 uses
it as well). Dominated convergence gives
\[
 s_1^+=2\pi\kappa_\Gamma\int_{\widetilde\Delta}\Psi\,d\mathrm{vol},\qquad
 \Psi=e^{\ii\sigma V}(1+e^{\ii V})(1-e^{\ii V})^{1-d}\prod_{i=1}^r\mathsf h_i^+(v_i)\,\mathsf h_i^-(v_i') .
\]

\emph{Step 5 (balance and exchange).} For $0<v<2\pi$, $1-e^{\pm\ii v}=2\sin\frac v2\,e^{\pm\ii(v-\pi)/2}$, so that
$(1-e^{\ii v})^{g-1}=e^{\ii(g-1)(v-\pi)}(1-e^{-\ii v})^{g-1}$ for principal powers, and by (B)
\begin{equation}
 \mathsf h_i^+(v)=e^{-\ii\pi(g_i-1)}\,e^{\ii(\frac d2-1-\sigma)v}\,\mathsf h_i^-(v)\qquad(0<v<2\pi),
 \label{gff:gdr:kernelBalance}
\end{equation}
with the same rate $\frac d2-1-\sigma$ for every factor. In the interior of $\widetilde\Delta$ all coordinates lie in $(0,2\pi)$.
With $1+e^{\ii V}=2\cos\frac V2\,e^{\ii V/2}$ and $(1-e^{\ii V})^{1-d}=(2\sin\frac V2)^{1-d}e^{\ii(1-d)(V-\pi)/2}$ the rates in $V$
cancel, and
\[
 \Psi=e^{\ii\pi(d-1)/2-\ii\pi\sum_i(g_i-1)}\,2\cos\frac V2\Bigl(2\sin\frac V2\Bigr)^{1-d}\prod_{i=1}^r\mathsf h_i^-(v_i)\,\mathsf h_i^-(v_i').
\]
The exchange $\iota(v,v')=(v',v)$ maps $\widetilde\Delta$ onto itself, preserves $d\mathrm{vol}$ and sends $V$ to $2\pi-V$. Every
factor of $\Psi$ is $\iota$-invariant except $\cos\frac V2$, which changes sign. So $\Psi\circ\iota=-\Psi$, $s_1^+=0$ and
$\gamma_{1-r}=0$. Integrated over the fibers of $V$, this is the odd one-variable integral
$s_1^+=2\pi\kappa_\Gamma e^{\ii\pi(d-1)/2-\ii\pi\sum_i(g_i-1)}\int_0^{2\pi}a_-(V)a_-(2\pi-V)(2\sin\frac V2)^{1-d}2\cos\frac V2\,dV$.
For the wheels, $a_-(V)=e^{-\ii(g-1)(V-r\pi)/2}\mathsf p(V)$ with $\mathsf p(V)=\int_{\Delta_V}\prod_i(2\sin\frac{v_i}2)^{g-1}d\mu$, positive for
real $b$, and the prefactor becomes $2\pi e^{\ii\pi(\sigma+1/2)}/\Gamma(g)^{2r}$: the integrand is real up to a constant
phase. Equivalently, writing $\Psi=2\cos\frac V2(2\sin\frac V2)^{1-d}\,\Pi\,e^{\ii\Theta}$ with
$\Pi=\prod_i(2\sin\frac{v_i}2\cdot2\sin\frac{v_i'}2)^{g_i-1}$, the phase obeys
$\Theta-\Theta\circ\iota=\sum_i(\sigma+\alpha_i+\beta_i+g_i-\frac d2)(v_i-v_i')$ for all $(v,v')$. So per-factor balance, not
total balance, makes it $\iota$-invariant.

\emph{Step 6 (the other top mode).} Since $d$ and the $d_\ell$ are real, $\overline{\Xi_T(\bar t;p)}=\Xi_T(t;\bar p)$ for the
parameter point $p=(\sigma,\alpha_i,g_i)$, and $\bar p$ again satisfies (B) and (C). Hence
$\gamma_{r-1}(p)=\overline{\gamma_{1-r}(\bar p)}=0$ by Steps~1--5 at $\bar p$. The direct computation on
$\operatorname{Im}z=-Y$ gives the same: the two forms of \eqref{gff:gdr:laplace} exchange their roles, every exponential has $\ii$
replaced by $-\ii$, and the integrand is again odd under $\iota$.

\emph{Step 7 (continuation).} $\operatorname{Re}g_i>0$ was used only in Steps~3--5. The $\gamma_m$ are analytic on the
convex, hence connected, set (C), and $\gamma_{\pm(r-1)}$ vanish on its nonempty open part where all
$\operatorname{Re}g_i>0$; so they vanish on all of it. $\square$

\paragraph{Details of the top-mode proof.}
The steps above state some estimates and computations in short form; this paragraph gives them in
full, and with it the proof of \eqref{gff:gdr:topG} is complete. Throughout, $\beta_i$ is given by (B). The set (C) is an open half-space of
$\mathbb C^{2r+1}\cong\mathbb R^{4r+2}$, hence convex and connected. For a compact set $B$ in it, let
$Y_B$ be the maximum of $Y_0$ and $\varkappa_B$ the minimum of $\varkappa$ over $B$; then
$\varkappa_B>\max(\frac d2,d-1,\frac12)$. Let $\mathsf P_T$ be the union of the $2r$ classes
$-\sigma-\alpha_i+\mathbb Z$ and $\beta_i+\mathbb Z$.

\emph{Step 0.} In the right part the Gamma-ratio bound \eqref{gff:mgf:gammaBounds} is applied to the left
factors $f_{g_i}(\ell+\sigma+\alpha_i+n+t)$ as they stand (base $1+\ell+n$) and to the right factors in
the first reflected form (base $1+n$); in the middle part to all factors as they stand,
$f_{g_i}(k+\sigma+\alpha_i+t)$ and $f_{g_i}(j+1+\beta_i-t)$ (bases $1+k$ and $1+j$); in the left part to
the left factors in the second reflected form (base $1+k$) and to the right factors as they stand,
$f_{g_i}(\ell+k+1+\beta_i-t)$ (base $1+\ell+k$). Factor $i$ contributes the power
$-\operatorname{Re}g_i$ of a base comparable, uniformly in $i$, to one of these, so the exponents add
up to $-\varkappa$ whatever the signs of the individual $\operatorname{Re}g_i$. On a compact set $K$
off $\mathsf P_T$ the sine ratios of the reflected forms are bounded, and, as in the paragraph
\emph{Proofs of the estimates}, the finitely many factor values with small arguments are bounded on
$K$ one by one, although each occurs for infinitely many lattice points. This gives
$|T_\ell(n+t)|\le C_K(1+\ell+n)^{-\varkappa}(1+n)^{-\varkappa}$, $C_K(1+k)^{-\varkappa}(1+j)^{-\varkappa}$
and $C_K(1+k)^{-\varkappa}(1+\ell+k)^{-\varkappa}$ in the three parts. For $|y|\geq Y_0$: by (B),
$\operatorname{Im}(\beta_i+g_i)=-\operatorname{Im}(\sigma+\alpha_i)$ and
$\operatorname{Im}(\sigma+\alpha_i+g_i)=-\operatorname{Im}\beta_i$. So every Gamma argument in a
numerator, namely $\ell+n+\sigma+\alpha_i+t$, $n+1-\beta_i-g_i+t$, $k+\sigma+\alpha_i+t$,
$j+1+\beta_i-t$, $k+2-\sigma-\alpha_i-g_i-t$ and $\ell+k+1+\beta_i-t$, and every sine argument in a
denominator, has imaginary part of modulus at least $|y|-Y_0+1\geq1$. No numerator Gamma function then
meets a pole, the sine ratios are bounded, and the bounds hold with each base $1+\cdots$ replaced by
$\cdots+|y|$. The constants are uniform for parameters in $B$, with $\varkappa_B$ in place of
$\varkappa$ and $|y|\geq Y_B$, provided $K$ avoids $\mathsf P_T$ for every parameter point of $B$. For
the wheels the right and middle bounds are those of the paragraph \emph{Proofs of the estimates}. With
$|d_\ell|\le C(1+\ell)^{d-2}$ and \eqref{gff:mgf:latticeSums} at $\alpha=\varkappa$ (the left sum has the
form of the right one, with $k$ in place of $n$) the three parts converge absolutely and uniformly on
compact subsets off $\mathsf P_T$. Each term is meromorphic in $t$; its poles come from
$\Gamma(\ell+\sigma+\alpha_i+n+t)$ and $\Gamma(\beta_i-n-t)$, are simple, and lie on the classes
$-\sigma-\alpha_i+\mathbb Z$ and $\beta_i+\mathbb Z$. At a point $t_0$ of $\mathsf P_T$ every term
therefore has a pole of order at most the number of the $2r$ classes that contain $t_0$; multiplied by
$(t-t_0)$ to that power, the series converges uniformly on a small circle around $t_0$, so the product
with $S_T$ is analytic near $t_0$, as in part (a) above. Periodicity follows from $n\mapsto n+1$. Each
factor of $D_T$ has simple zeros exactly on one class, so $D_T$ vanishes at $t_0$ to the same order,
and $\Xi_T$ is entire and one-periodic. For $0\le x\le1$ and $|y|\geq\max(Y_0,2)$, the large-$|y|$
bounds and \eqref{gff:mgf:latticeSums} give $|S_T(x+\ii y)|\le C|y|^{\max(d,1)-2\varkappa}(1+\log|y|)$, and
$|D_T(x+\ii y)|\le Ce^{2\pi r|y|}$. Taking $y=-Y$ for $m\geq r$ and $y=Y$ for $m\le-r$ in the Fourier
coefficient gives $|\gamma_m|\le Ce^{2\pi(r-|m|)Y}Y^{\max(d,1)-2\varkappa}(1+\log Y)\to0$, because
$2\varkappa>\max(d,1)$ by (C). For the analyticity, on the segment $t=x+\ii Y_B$, $0\le x\le1$, which
stays at distance at least one from $\mathsf P_T$ for every parameter point of $B$, the series of $S_T$
converges uniformly in $x$ and in the parameters, its terms are analytic in the parameters, and $D_T$
is entire in $t$ and the parameters. Hence $\gamma_m=\int_0^1\Xi_T(x+\ii Y_B)e^{-2\pi\ii m(x+\ii Y_B)}dx$
is analytic on the interior of $B$.

\emph{Step 1.} The coefficients of the expansion in the upper half-plane with $m\le0$ vanish: the
coefficient integral $\int_0^1S_T(x+\ii y)e^{-2\pi\ii m(x+\ii y)}dx$ does not depend on $y>Y_0$, its
exponential is bounded by one for $m\le0$, and $S_T\to0$ uniformly as $y\to\infty$. The relation
$\gamma_{1-r}=\kappa_+s_1^+$ compares the coefficients of the lowest power $w^{1-r}$ in $\Xi_T=D_TS_T$
with the expansion of Step~0; the lower half-plane is the same.

\emph{Step 3: the Laplace forms \eqref{gff:gdr:laplace}.} For $\operatorname{Re}w>0$,
$\Gamma(g)f_g(w)=B(w,g)=\int_0^\infty e^{-wu}(1-e^{-u})^{g-1}du$ (put $x=e^{-u}$ in Euler's integral).
Substitute $u=-\ii v$, with $v$ on the positive imaginary axis, and rotate the ray to $[0,\infty)$
through the first quadrant. If also $\operatorname{Im}w>0$, then $|e^{\ii wv}|=e^{-\operatorname{Im}(wv)}$
decays exponentially in the closed quadrant, since $\arg w+\arg v\in(0,\pi)$ there. Moreover
$1-e^{\ii v}$ has nonnegative real part in the quadrant and vanishes only at $v\in2\pi\mathbb N_0$,
where the integrand is $O(|v-2\pi k|^{\operatorname{Re}g-1})$, which is integrable. These points lie on
the boundary of the quadrant: $v=0$ at its corner and $v=2\pi k$, $k\geq1$, on its real side. A small
quarter circle of radius $\epsilon$ at $0$ and small semicircles at $2\pi k$, $k\geq1$, contribute
$O(\epsilon^{\operatorname{Re}g})$ each and nothing in the limit. This gives the first formula for
$\operatorname{Re}w>0$, $\operatorname{Im}w>0$. The substitution $u=\ii v$ and the fourth quadrant give
the second formula for $\operatorname{Re}w>0$, $\operatorname{Im}w<0$. Both right sides are analytic on
the respective half-planes, because
$|(1-e^{\pm\ii v})^{g-1}|\le e^{\pi|\operatorname{Im}g|/2}|2\sin\frac v2|^{\operatorname{Re}g-1}$ is
$2\pi$-periodic and locally integrable; so \eqref{gff:gdr:laplace} holds on the whole half-planes. For real
$g$ the second formula is the complex conjugate of the first; for complex $g$ it is the separate
rotation just described. On $\operatorname{Im}z=Y>Y_0$ one has
$|e^{\ii(\ell+\sigma+\alpha_i+z)v}|=e^{-(Y+\operatorname{Im}(\sigma+\alpha_i))v}\le e^{-v}$ and
$|e^{-\ii(\beta_i-z)v}|=e^{-(Y-\operatorname{Im}\beta_i)v}\le e^{-v}$, so the two products of Step~3 are
absolutely convergent integrals and Fubini applies.

\emph{Step 3: the endpoint bound \eqref{gff:gdr:endpoint}.} On $[0,3\pi]$ the exponentials in
$\mathsf h_i^\pm$ and the factors $e^{-\operatorname{Im}g_i\arg(1-e^{\pm\ii v})}$ are bounded, so
$|\mathsf h_i^\pm(v)|\le C|2\sin\frac v2|^{\operatorname{Re}g_i-1}$ there. The modulus
$|2\sin\frac v2|$ lies between $cv$ and $v$ on $(0,\pi+\delta]$ and between $c|2\pi-v|$ and $|2\pi-v|$
on $[\frac\pi2,3\pi)$. Hence, whatever the sign of $\operatorname{Re}g_i-1$,
$|\mathsf h_i^\pm(v)|\le Cv^{\operatorname{Re}g_i-1}$ on the first interval and
$|\mathsf h_i^\pm(v)|\le C|2\pi-v|^{\operatorname{Re}g_i-1}$ on the second. For $V\le\pi$ every
coordinate of $\Delta_V$ is at most $\pi$, and Dirichlet's integral gives
$|a_\pm(V)|\le C\int_{\Delta_V}\prod_iv_i^{\operatorname{Re}g_i-1}d\mu
=CV^{\varkappa-1}\prod_i\Gamma(\operatorname{Re}g_i)/\Gamma(\varkappa)$. For $\pi\le V\le2\pi+\delta$ at
most one coordinate of a point of $\Delta_V$ exceeds $V/2$. Where none does, all coordinates are below
$\pi+\delta$, and the same Dirichlet bound gives $O(1)$. Where $v_j>V/2$, write $s=V-v_j<V/2$ and
$\epsilon_V=2\pi-V\in[-\delta,\pi]$, so that $2\pi-v_j=\epsilon_V+s$. Bounding factor $j$ by
$C|\epsilon_V+s|^{\operatorname{Re}g_j-1}$ and integrating the other $r-1$ factors over their simplex
of size $s$ (Dirichlet again, with exponent $\varkappa-\operatorname{Re}g_j>0$) leaves
\[
 C\int_0^{V/2}|\epsilon_V+s|^{\operatorname{Re}g_j-1}s^{\varkappa-\operatorname{Re}g_j-1}ds .
\]
For $\operatorname{Re}g_j\geq1$ this is bounded. For $\operatorname{Re}g_j<1$, the part
$s\le2|\epsilon_V|$ is $|\epsilon_V|^{\varkappa-1}$ times the convergent integral
$\int_0^2|1\pm u|^{\operatorname{Re}g_j-1}u^{\varkappa-\operatorname{Re}g_j-1}du$ (substitute
$s=|\epsilon_V|u$), and on $s>2|\epsilon_V|$ one has $|\epsilon_V+s|\geq s/2$, which leaves
$C\int_{2|\epsilon_V|}^{2\pi}s^{\varkappa-2}ds$. Both are
$O(|\epsilon_V|^{\min(\varkappa-1,0)}\lambda)$. Summing over $j$ proves \eqref{gff:gdr:endpoint}. At
$\varkappa=1$ the bound is at most logarithmic: when $\operatorname{Im}\sum_ig_i\neq0$,
$a_\pm(2\pi-\epsilon_V)$ stays bounded and oscillates like $\epsilon_V^{\ii\operatorname{Im}\sum_ig_i}$
(checked numerically in the review \path{LRI_C5_general_d/review/lemmaG-proof} for $r=2$,
$g=(0.45+0.3\ii,\,0.55)$), and the logarithm can occur
only when $\operatorname{Im}\sum_ig_i=0$. Only the upper bound is used. For the wheels at real $b$, by
dominated convergence, $a_-(2\pi^-)=e^{-\ii(g-1)(2-r)\pi/2}\mathsf p(2\pi)$ is finite and nonzero, so
for $\varkappa>1$ the exponent in \eqref{gff:gdr:endprod} is $\varkappa-1$, not $2\varkappa-2$.

\emph{Step 3: the formula \eqref{gff:gdr:Il}.} The functions $a^Y_\pm$ lie in $L^1(\mathbb R)$:
$|a^Y_+(V)|$ is at most the convolution of the functions
$Ce^{-(Y+\operatorname{Im}(\sigma+\alpha_i))v}|2\sin\frac v2|^{\operatorname{Re}g_i-1}$, each integrable
on $[0,\infty)$, and likewise for $a^Y_-$. The bound $|f_{g_i}(w)|\le C(1+|w|)^{-\operatorname{Re}g_i}$
on horizontal lines at distance at least one from the real axis follows from Gamma reflection and
\eqref{gff:mgf:gammaBounds}. For $U<2\pi+\delta$ the convolution
$(a^Y_+*a^Y_-)(U)=\int_0^Ua^Y_+(V)a^Y_-(U-V)dV$ involves $a^Y_\pm$ only on $[0,2\pi+\delta]$, where
they are square integrable by \eqref{gff:gdr:endpoint} and $2\varkappa>1$, and a convolution of two
square-integrable functions is continuous. Since the convolution and its Fourier transform are both
integrable, Fourier inversion holds at this point of continuity,
$\int_{\mathbb R}\widehat{a^Y_+*a^Y_-}(x)e^{-2\pi\ii x}dx=2\pi(a^Y_+*a^Y_-)(2\pi)$. Together with
$e^{-2\pi\ii z}=e^{2\pi Y}e^{-2\pi\ii x}$, and the factor $e^{-2\pi Y}$ in $(a^Y_+*a^Y_-)(2\pi)$, this
is the first form of \eqref{gff:gdr:Il}. The second form is Fubini; the absolute integrand obeys the same
bounds. The measure $d\mathrm{vol}$ is $(2r)^{-1/2}$ times the surface measure of the hyperplane
$V+V'=2\pi$, hence invariant under every permutation of the $2r$ coordinates.

\emph{Step 4.} For $\frac12\le\rho<1$, $|1-\rho e^{\ii V}|^2=(1-\rho)^2+4\rho\sin^2\frac V2\geq
2\sin^2\frac V2$, so $|(1+\rho e^{\ii V})(1-\rho e^{\ii V})^{1-d}|\le C|\sin\frac V2|^{-\max(d-1,0)}$.
Dominated convergence gives first the one-variable form
$s_1^+=2\pi\kappa_\Gamma\int_0^{2\pi}e^{\ii\sigma V}(1+e^{\ii V})(1-e^{\ii V})^{1-d}a_+(V)a_-(2\pi-V)\,dV$,
and the form on $\widetilde\Delta$ is Fubini, since after integration over the fibers of $V$ the
function $|\Psi|$ is dominated by the same integrable function. For $d>2$ the entry $\varkappa>d-1$ of
(C) is stronger than the spectral condition $2\varkappa>d$, and by \eqref{gff:gdr:endprod} the argument does
not extend past $b_{\max}$.

\emph{Step 6.} Conjugating the expansion of Step~1 at $\bar p$ gives the expansion in the lower
half-plane at $p$, so $s_m^-(p)=\overline{s_m^+(\bar p)}$ and
$\gamma_{r-1}(p)=\kappa_-s_1^-(p)=\kappa_-\overline{s_1^+(\bar p)}=0$; this is equivalent to the
relation $\gamma_m(\bar p)=\overline{\gamma_{-m}(p)}$ used above. At real parameters it is the reality
$\gamma_{-m}=\overline{\gamma_m}$; at complex parameters it relates two different parameter points and
is not a symmetry of a single $\Xi_T$.

\emph{Step 7.} The part of (C) where all $\operatorname{Re}g_i>0$ is open and nonempty; it contains the
points with all $\operatorname{Re}g_i=\varkappa/r$.

\emph{The equivalent form.} On $\operatorname{Im}t>Y_0$, $D_T=\kappa_+w^{-r}(1+O(w))$ with
$\kappa_+\neq0$, and $\Xi_T=\sum_{|m|\le r-2}\gamma_mw^m=O(e^{2\pi(r-2)\operatorname{Im}t})$; so
$S_T=\Xi_T/D_T=O(e^{-4\pi\operatorname{Im}t})$ as $\operatorname{Im}t\to+\infty$. Conversely, this decay
forces $s_1^+=0$: let $Y\to\infty$ in its coefficient integral. The lower half-plane is the same with
$\kappa_-$. For the wheels, with $c_{r,r-1}=2\gamma_{r-1}e^{-\ii\pi(r-1)\sigma}$, this is the
equivalent form stated after \eqref{gff:gdr:topless}.

\emph{The wheel case and the earlier proof.} For the wheels this proof is the earlier, wheel-only proof
that the review \path{LRI_C5_general_d/review/VERIFICATION.md} checked (proof review in
\path{LRI_C5_general_d/review/foundation-rigor}; Section~\ref{gff:gdr:status}), with Steps~3--4 as corrected after that
review, up to three changes: the parameters are complex throughout (the earlier version reduced to real
$b$ by analyticity in $b$), the second top mode comes from Step~6 instead of the reflection of part (b)
of Section~\ref{gff:gdr:finite}, and Step~5 is written through \eqref{gff:gdr:kernelBalance}.


Two ingredients enter: the balance of every single factor (Step~5) and the factor $1+z$ in \eqref{gff:gdr:Dgen}, which
is odd about $V=\pi$. Both are essential: controls that break either one have nonzero top modes. The numerical tests
of the theorem and of its proof, with these controls, are collected in Section~\ref{gff:gdr:checks} (item \emph{Top-mode
theorem and its product form}).

\paragraph{Relation to the periodization at $b=d/2$ and to the half-shift reduction.}
At $b=d/2$ and $\frac2r<d<4$, \eqref{gff:gdr:topless} sharpens \eqref{gff:fp:finiteFourier} for the wheels: $c_{3,2}(q)=0$ and
$c_{4,3}(q)=0$ in every such dimension, and the sum in \eqref{gff:fp:finiteFunctional} stops at $k=r-2$. It says nothing
about $\Xi_7$: the kernel $K_\ell$ of \eqref{gff:fp:primitiveStart} carries the extra factor
$[(\ell+p)^2-4s^2]\Gamma(p+2s)\Gamma(p-2s)/(\Gamma(2s)\Gamma(-2s))$ and the multiplicities $(-1)^\ell d_\ell$, so it is not of the
form \eqref{gff:gdr:TG}. In Section~\ref{gff:app:threeDimensionalReduction} the relations $c_{3,2}=0$ and $c_{4,3}=0$ come
out of the half-shift determinant \eqref{gff:td:projectedAnswer}; they are the $d=3$ cases of \eqref{gff:gdr:topless}, which
holds for every $d$ and every $b$ in the strip. What is special to three dimensions is the identification of the
surviving modes with products of ordinary moment amplitudes, $c_{3,0}=-\mathfrak a_3\mathfrak b_3$ and
$c_{4,1}=2\mathfrak a_4\mathfrak b_4$, and the odd $r=7$ coefficients \eqref{gff:td:sevenOddCoefficients}. At $d=1$ and
$d=2$, \eqref{gff:gdr:topless} also follows from the one-variable lattices \eqref{gff:gdr:d12lattice} below (at $d=2$, $b=1$
and $r=4$ it is \eqref{gff:fp:Omega4Modes}); at generic $d$ only the proof of Section~\ref{gff:gdr:top} gives it.

\subsubsection{Residue functional, explicit weights and blocks}\label{gff:gdr:residue}
For $0<\operatorname{Re}b<b_{\max}(r,d)$ and $\sigma\notin\mathbb Z$,
\begin{equation}
 W_r(d;b)=-N_b^{\,r}\Res_{t=0}\frac{\Xi_r(t;b)}{D_b(t)^r}
 =-\Bigl[\frac{N_b}{\pi\sin\pi\sigma}\Bigr]^r\sum_{k=0}^{r-2}B_{r,k}(\sigma)\,c_{r,k}(q,b),
 \label{gff:gdr:functional}
\end{equation}
where $B_{r,k}$ is \eqref{gff:fp:finiteFunctional} with the logarithmic coefficients \eqref{gff:fp:logCoefficients}, both
with $q$ replaced by $\sigma$. The proof follows the one after \eqref{gff:fp:periodsFinite}; it is given in full in the
paragraph \emph{The residue functional in full} below. The weights depend on $b$ only, not on $q$. For $r=3,4$,
\begin{equation}\boxed{\begin{aligned}
 W_3(d;b)&=-\frac{2N_b^3}{\pi s_b^5}\Bigl[(1+2\varsigma_b^2)\,c_{3,0}+3\varsigma_b\,c_{3,1}\Bigr],\\
 W_4(d;b)&=\frac{4N_b^4}{\pi s_b^7}\Bigl[\varsigma_b(3+2\varsigma_b^2)\,c_{4,0}+(1+4\varsigma_b^2)\,c_{4,1}
 +5\varsigma_b\,c_{4,2}\Bigr].
\end{aligned}}\label{gff:gdr:explicit}\end{equation}
The weights of the omitted top modes, $B_{3,2}=6\pi^2/s_b^2$ and $B_{4,3}=-20\pi^3/s_b^3$, are nonzero, so the
formulas below depend on \eqref{gff:gdr:topless}.

\paragraph{The residue functional in full.}
\emph{Proof of \eqref{gff:gdr:functional}.} Let
$0<\operatorname{Re}b<b_{\max}(r,d)$ and $\sigma\notin\mathbb Z$; then $b_{\max}\le b_{\rm spec}$, so
\eqref{gff:gdr:kernel} represents $W_r$. For each spin close the Barnes contour to the right by the arcs of
the circles $|s-a_\ell|=n_0+\frac12$, $n_0\in\mathbb N$, that lie in $\operatorname{Re}s\geq0$. The part of the disc
$|s-a_\ell|<n_0+\frac12$ in $\operatorname{Re}s>0$ contains the right poles $a_\ell+n$, $0\le n\le n_0$, and no left
pole, since the left poles $-a_\ell-n$ have negative real part. The arcs stay at distance at least $\frac12$ from
$a_\ell+\mathbb Z$, because $|s-a_\ell-m|\geq|n_0+\frac12-|m||\geq\frac12$ for every integer $m$, so the sine ratio
produced by reflecting $\Gamma(a_\ell-s)/\Gamma(\beta_\ell-s)$ is bounded on them, and $g_\ell(s)^r
=O(|s|^{-2r\operatorname{Re}g})$ there; since $2r\operatorname{Re}g>1$, the arcs contribute nothing in
the limit. (The package's own proof closes the contour on the circles $|s|=a_\ell+N+\frac12$, whose radius is real
only for real $b$. The circles centred at $a_\ell$ keep the distance $\frac12$ from $a_\ell+\mathbb Z$ for every
complex $b$. This change was made when the proof was merged into this note, on 16 October 2026, and has not been
reviewed.) The closure is clockwise, so the $\ell$-th integral is
$-\sum_{n\geq0}\Res_{s=a_\ell+n}g_\ell(s)^r$, and
$\Res_{s=a_\ell+n}g_\ell(s)^r=\Res_{t=0}g_\ell(a_\ell+n+t)^r$. In $S_r$ the point $t=0$ is a pole only
of the terms with $n\geq0$: for $n<0$, $\Gamma(-n-t)$ is regular at $t=0$, and $t=0\equiv-\sigma$ would
need $\sigma\in\mathbb Z$. By the uniform convergence on a small circle $|t|=\epsilon$ (paragraph
\emph{Proofs of the estimates}),
$\sum_\ell d_\ell\sum_{n\geq0}\Res_{s=a_\ell+n}g_\ell^r=\Res_{t=0}S_r=\Res_{t=0}\Xi_r/D_b^r$, which is
the first equality of \eqref{gff:gdr:functional}. Next,
\[
 \log\frac{(\pi t\sin\pi\sigma)^r}{D_b(t)^r}
 =r\log\frac{\pi t}{\sin\pi t}-r\log\frac{\sin\pi(t+\sigma)}{\sin\pi\sigma}
 =\sum_{h\geq1}u_h\frac{t^h}{h!},
\]
with $u_h$ the coefficients \eqref{gff:fp:logCoefficients} at $q\to\sigma$,
$u_h(\sigma)=r[(1+(-1)^h)\psi_{h-1}(1)-(-1)^{h+1}\psi_{h-1}(1-\sigma)+\psi_{h-1}(\sigma)]$. Indeed
$\log(\pi t/\sin\pi t)=\sum_{m\geq1}\zeta(2m)t^{2m}/m$, with $2(h-1)!\,\zeta(h)=(1+(-1)^h)\psi_{h-1}(1)$
for even $h$, and $\partial_\sigma^h\log\sin\pi\sigma=-\psi_{h-1}(\sigma)-(-1)^h\psi_{h-1}(1-\sigma)$
\cite{gff:DLMFPolygamma}. With the polynomials $E_m$ of \eqref{gff:fp:BellFinite}, which are the Taylor
coefficients of $\exp\sum_hx_ht^h/h!$, therefore
$\Res_{t=0}e^{2\pi\ii mt}D_b(t)^{-r}=(\pi\sin\pi\sigma)^{-r}E_{r-1}(u_1+2\pi\ii m,u_2,\dots,u_{r-1})$.
Inserting $\Xi_r=\sum_{|m|\le r-2}\gamma_me^{2\pi\ii mt}$, which is \eqref{gff:gdr:topless}, gives the
exponential form
\begin{equation}
 W_r(d;b)=-\Bigl[\frac{N_b}{\pi\sin\pi\sigma}\Bigr]^r\sum_{|m|\le r-2}\gamma_m\,
 E_{r-1}(u_1+2\pi\ii m,u_2,\dots,u_{r-1}),
 \label{gff:mgf:expForm}
\end{equation}
and $\gamma_{\pm k}=\frac12c_{r,k}e^{\pm\ii\pi k\sigma}$ ($k\geq1$), $\gamma_0=c_{r,0}$ turn it into
\eqref{gff:gdr:functional} with
$B_{r,k}(\sigma)=\frac12[e^{\ii\pi k\sigma}E_{r-1}(u_1+2\pi\ii k,u_2,\dots)+e^{-\ii\pi k\sigma}
E_{r-1}(u_1-2\pi\ii k,u_2,\dots)]$, which is \eqref{gff:fp:finiteFunctional} at $q\to\sigma$. The
functional depends on $b$ only, through $\sigma$; the dimension enters only through $N_b$ and the
$c_{r,k}(q,b)$. By \eqref{gff:fp:elementaryWeights} at $q\to\sigma$ the $u_h$ are elementary,
$u_1/(r\pi)=-\varsigma_b/s_b$, $u_2/(r\pi^2)=\varsigma_b^2/s_b^2+\frac43$ and
$u_3/(r\pi^3)=-2(\varsigma_b/s_b)(1+\varsigma_b^2/s_b^2)$, and for the $b$-derivatives of the weights
$\partial_\sigma u_h=u_{h+1}-r(1-(-1)^h)\psi_h(1)$. For $r=3,4$,
\begin{equation}\begin{gathered}
 B_{3,0}=\frac{2\pi^2(1+2\varsigma_b^2)}{s_b^2},\qquad B_{3,1}=\frac{6\pi^2\varsigma_b}{s_b^2},\qquad
 B_{3,2}=\frac{6\pi^2}{s_b^2},\\
 B_{4,0}=-\frac{4\pi^3\varsigma_b(3+2\varsigma_b^2)}{s_b^3},\qquad
 B_{4,1}=-\frac{4\pi^3(1+4\varsigma_b^2)}{s_b^3},\qquad
 B_{4,2}=-\frac{20\pi^3\varsigma_b}{s_b^3},\qquad B_{4,3}=-\frac{20\pi^3}{s_b^3},
\end{gathered}\label{gff:mgf:Bweights}\end{equation}
and \eqref{gff:gdr:explicit} follows. $\square$ At $d=2$ and $r=4$, with $x=\frac b2$, $N_b$ equals $\mathcal N_x$ of
Section~\ref{gff:gdr:special}, and inserting $c_{4,0},c_{4,1},c_{4,2}$ of the $d=2$ paragraph there into
\eqref{gff:gdr:explicit} reproduces the weights $\rho_0,\rho_1,\rho_2$ of \eqref{gff:gc:L2residue}. At $b=1$
($\varsigma_b=0$, $s_b=1$, $N_1=1$) it reduces to $P_4(2)=4c_{4,1}/\pi=8\mathcal A(\frac12)\mathcal
B(\frac12)/\pi$.


Each sample is a sum of five blocks, \eqref{gff:mgf:blocks} below. The blocks are
\eqref{gff:fp:wheelBlocks} with the arrays \eqref{gff:fp:wheelArrays}, where the rim parameter $q$ is replaced by $\sigma$ in
$q+t$, $q+1+t$ and in the mirror point, and the spoke parameter $q$ by $g$ in $1-q+t$, $1+q-t$ and $\Gamma(q-t)$. Every
parameter is affine in $(q,b)$, and all five blocks converge absolutely exactly for $\operatorname{Re}b<b_{\max}$
(paragraph \emph{Proof of the block decomposition} below). The sections below do not use generic nodes. They sample at the pole node
and at the reflection centers, where the blocks collapse to fewer functions.

\paragraph{Proof of the block decomposition.}
Write $A_r^0,A_r^+,C_r$ for the blocks
\eqref{gff:fp:wheelBlocks} with the substitutions just stated; at $b=\frac d2$ they are the
blocks of \eqref{gff:fp:wheelSample}. The claim is
\begin{equation}
 \Xi_r(t;b)=D_b(t)^r\bigl[A_r^0(t)+A_r^+(t)+A_r^0(1-\sigma-t)+A_r^+(1-\sigma-t)+C_r(t)\bigr]
 \qquad(t\notin\mathbb Z\cup(-\sigma+\mathbb Z)).
 \label{gff:mgf:blocks}
\end{equation}
In the right part with $\ell=0$, the factorized form in \eqref{gff:gdr:reflected} and
\[
 f(\sigma+n+t)=\frac{\Gamma(\sigma+t)}{\Gamma(2q+t)}\frac{(\sigma+t)_n}{(2q+t)_n},\qquad
 f(-n-t)=\frac{\Gamma(-t)}{\Gamma(g-t)}\frac{(1-g+t)_n}{(1+t)_n},
\]
which use $\sigma+g=2q$, give the series of $A_r^0(t)$ term by term; the upper parameter $1$ of the
${}_{2r+1}F_{2r}$ cancels $n!$. For $\ell=j+1$,
$f(j+1+\sigma+n+t)=[\Gamma(1+\sigma+t)/\Gamma(1+2q+t)]\,(1+\sigma+t)_{j+n}/(1+2q+t)_{j+n}$, and the
multiplicity $d_{j+1}=4q\,(2q+1)_j(4q-1)_j/((2q)_j(2)_j)$ of \eqref{gff:fp:degeneracy} supplies the prefactor
$4q$ and the rows $(2q+1,\boldsymbol e),(4q-1,\boldsymbol e);(2q,\boldsymbol e),(2,\boldsymbol e)$;
the rows $(1,\boldsymbol e)$ and $(1,\boldsymbol f)$ cancel the factorials of \eqref{gff:fp:HornDefinition}.
This gives $A_r^+(t)$. In the middle part, the third line of \eqref{gff:gdr:reflected},
\[
 \frac{\Gamma(k+\sigma+t)}{\Gamma(k+2q+t)}\,\frac{\Gamma(j+1-t)}{\Gamma(j+1+g-t)}
 =\frac{\Gamma(\sigma+t)\Gamma(1-t)}{\Gamma(2q+t)\Gamma(1+g-t)}\,
 \frac{(\sigma+t)_k(1-t)_j}{(2q+t)_k(1+g-t)_j},
\]
with the multiplicity $d_{j+k+1}=4q\,(2q+1)_{j+k}(4q-1)_{j+k}/((2q)_{j+k}(2)_{j+k})$, gives $C_r(t)$,
with the multiplicity rows along $\boldsymbol c$. By the last line of \eqref{gff:gdr:reflected} the left
part is the right part at $1-\sigma-t$. All the series converge absolutely by
\eqref{gff:mgf:latticeSums}, so the rearrangement is justified, and multiplying by $D_b(t)^r$ gives
\eqref{gff:mgf:blocks}. Each block also converges directly under part of the condition: the
${}_{2r+1}F_{2r}(1)$ has parameter excess $r(2q+t)+r(1+t)-1-r(\sigma+t)-r(1-g+t)=2rg-1$, so it needs
condition (iii); $A_r^+$ needs (i)--(iii), and $C_r$ needs (i) and (ii). Together this is exactly
$\operatorname{Re}b<b_{\max}$. At integer $d$ the multiplicity rows terminate: $d_{j+1}/(4q)$ is
$\delta_{j0}$ at $d=1$, $1$ at $d=2$ and $1+\frac23j$ at $d=3$.

\paragraph{Removable points of individual blocks.}
As functions of $t$ the five blocks are sub-sums of $S_r$, so each is analytic off $\mathbb
Z\cup(-\sigma+\mathbb Z)$. In the factorized form, a Gamma prefactor times a Horn value, an individual
Horn value has poles where one of its lower parameters is a non-positive integer, cancelled by a zero of
the same order of its prefactor. This happens for $A_r^0(t)$ at $t\in-2q-\mathbb N_0$, for
$A_r^0(1-\sigma-t)$ at $t\in1+g+\mathbb N_0$, for $A_r^+(t)$ at $t\in-1-2q-\mathbb N_0$, for
$A_r^+(1-\sigma-t)$ at $t\in2+g+\mathbb N_0$, and for $C_r(t)$ at $t\in-2q-\mathbb N_0$ and
$t\in1+g+\mathbb N_0$. At these points \eqref{gff:mgf:blocks} holds by continuity. For real $b$ the points
are real, and the nodes of \eqref{gff:mgf:Wsample} below, which have $\operatorname{Im}t_j=\eta_r>0$, never
meet them.

\subsubsection{Fixed nodes, jets, evaluation of the samples, and the continuation beyond \texorpdfstring{$b_{\max}$}{bmax}}\label{gff:mgf:nodes}
The formulas of this subsection are exact algebra on
Sections~\ref{gff:gdr:finite}--\ref{gff:gdr:residue}; the continuation beyond $b_{\max}$ is derived, not proved
(Section~\ref{gff:gdr:status}). Sections~\ref{gff:gdr:samples}--\ref{gff:gdr:C5} do not use this subsection. It
describes the form implemented by the general-$b$ code \path{LRI_C5_general_d/foundation/wheel_general_b.py}
(function \texttt{W(r,d,b,dps)}), which Section~\ref{gff:gdr:checks} tests.

\paragraph{Nodes that do not move with $b$.}
Let $N$ be either $N_r=2r-1$ of \eqref{gff:fp:nodes} or the reduced value $N_r'=2r-3$ permitted by
\eqref{gff:gdr:topless}, and $M=(N-1)/2$. Put
\[
 t_j=-\frac q2+\frac jN+\ii\eta_r,\qquad j=0,\dots,N-1,
\]
with $\eta_r=1/(8\pi r)$ of \eqref{gff:fp:nodes}; for $N=N_r$ the first $r$ of them are the nodes of
\eqref{gff:fp:nodes}. Since $\Xi_r$ is a trigonometric polynomial with exponential modes $|m|\le M$, the
discrete Fourier transform of the $N$ complex samples recovers all modes without aliasing,
$\gamma_m=N^{-1}\sum_j\Xi_r(t_j;b)e^{-2\pi\ii mt_j}$. Inserted into \eqref{gff:mgf:expForm} this gives, for
$0<\operatorname{Re}b<b_{\max}(r,d)$ and $\sigma\notin\mathbb Z$,
\begin{equation}\boxed{\begin{aligned}
 W_r(d;b)&=N_b^{\,r}\sum_{j=0}^{N-1}\mathfrak w_{r,j}(b)\,\Xi_r(t_j;b),\\
 \mathfrak w_{r,j}(b)&=-\frac{1}{N(\pi\sin\pi\sigma)^r}\sum_{|m|\le M}
 E_{r-1}(u_1+2\pi\ii m,u_2,\dots,u_{r-1})\,e^{-2\pi\ii mt_j}.
\end{aligned}}\label{gff:mgf:Wsample}\end{equation}
The weights are elementary: by \eqref{gff:fp:elementaryWeights} at $q\to\sigma$ they are polynomials in
$\pi$, $\cot\frac{\pi b}2$ and $1/\sin\frac{\pi b}2$, times $e^{-2\pi\ii mt_j}$. The nodes do not
depend on $b$, and every factor is analytic in $b$ on the strip (the samples by part (d) of
Section~\ref{gff:gdr:finite}); so the product rule and \eqref{gff:gs:wheelJets} give, with
$\partial_b\log N_b=-\frac12[\psi_0(g)+\psi_0(\sigma)]$,
\begin{equation}\begin{aligned}
 \LL_d&=\sum_j\Bigl[\partial_b\bigl(N_b^4\mathfrak w_{4,j}\bigr)\,\Xi_4(t_j;b)
 +N_b^4\mathfrak w_{4,j}\,\partial_b\Xi_4(t_j;b)\Bigr]_{b=d/2},\\
 Z_d&=\frac19\sum_j\Bigl[\partial_b^2\bigl(N_b^3\mathfrak w_{3,j}\bigr)\,\Xi_3
 +2\partial_b\bigl(N_b^3\mathfrak w_{3,j}\bigr)\,\partial_b\Xi_3+N_b^3\mathfrak w_{3,j}\,\partial_b^2\Xi_3\Bigr]_{b=d/2}.
\end{aligned}\label{gff:mgf:jets}\end{equation}
Here $\partial_b^k\Xi_r(t_j;b)=\sum_{p\le k}\binom kp\partial_b^{k-p}[D_b(t_j)^r]\,\partial_b^p[\cdots]$,
with $\partial_bD_b(t)^r=\frac{r\pi}2\cot\pi(t+\sigma)\,D_b(t)^r$, and the bracket is the sum of the
five blocks of \eqref{gff:mgf:blocks}. The $b$-derivatives of the blocks at fixed $t_j$ are first and second
derivatives of the ${}_{2r+1}F_{2r}(1)$ values and of the two-variable Horn values with respect to their
parameters, which are all affine in $b$ with slopes $\pm\frac12$, together with polygamma values from
the Gamma prefactors; in the mirrored blocks the argument $1-\sigma-t_j$ moves with $b$, which is again
a parameter derivative. In this form $\LL_d$ and $Z_d$ are finite combinations, with elementary
coefficients, of one- and two-variable hypergeometric values at unit argument and their first and
second parameter derivatives. With the reduced nodes $W_3$ needs $3$ complex samples and $W_4$ needs
$5$. The $b$-dependent symmetric nodes $-\frac\sigma2+\frac j{N_r}+\ii\eta_r$ with the real cosine
transform \eqref{gff:fp:cosineTransform} are equally valid, and need $2$ and $3$ real samples; their
$b$-derivatives also contain $t$-derivatives of the samples, which is why the fixed nodes are used for
the jets. The code implements the symmetric nodes as a cross-check (option \texttt{nodes='note'}).

\paragraph{Radius of the sample jets.}
$W_r(d;b)$ is analytic up to $b_{\rm spec}$, but the individual samples $\Xi_r(t_j;b)$ and the
coefficients $c_{r,k}(q,b)$ with $k\le r-2$ have a simple pole at $b=b_{\max}$ when
$b_{\max}<b_{\rm spec}$, that is for $d>2$ (nonzero residue in every case computed; paragraph
\emph{Beyond $b_{\max}$} below). Moreover, at a fixed node the individual blocks have poles in $b$
where $\sigma+t_j\in\mathbb Z$ (at complex $b$, since $\operatorname{Im}t_j=\eta_r\neq0$); they are
removable in $\Xi_r(t_j;b)$ because of the zero of $D_b(t_j)^r$. The derivatives in \eqref{gff:mgf:jets} at
$b=\frac d2$ are unaffected. When they are computed from Taylor series or Cauchy circles of the
samples, however, the radius must stay below $b_{\max}-\frac d2$; for $d\geq2$ this is $\frac12$ for
$r=4$ and $\frac{4-d}6$ for $r=3$. A second restriction concerns the weights: $\mathfrak w_{r,j}(b)$
carries $(\sin\pi\sigma)^{-r}$ and has a pole of order $r$ at $b=2$. The disc
$|b-\frac d2|<b_{\max}-\frac d2$ contains $b=2$ exactly when $b_{\max}>2$ (paragraph \emph{Integer
$\sigma$} of Section~\ref{gff:gdr:finite}): for $r=4$ when $d>3$, and for $r=5$ when $d>\frac83$. There the
samples and $W_r$ are analytic, and the poles of the weights cancel only in the sum \eqref{gff:mgf:Wsample}.
Cauchy circles of the weighted products $N_b^r\mathfrak w_{r,j}\Xi_r(t_j;b)$, or of the weights alone,
must then stay inside $|b-\frac d2|<|2-\frac d2|$. The jets \eqref{gff:mgf:jets} at $b=\frac d2\neq2$ are
unaffected, and the cases used in this note ($r=3,4$ with $d\le3$) do not meet $b=2$.

\paragraph{Evaluation of the Horn values.}
The module \path{LRI_C5_general_d/foundation/hypsums.py} evaluates the sums of the five blocks as an
exact finite part plus an asymptotic tail, and \path{wheel_general_b.py} assembles
\eqref{gff:mgf:blocks} and \eqref{gff:mgf:Wsample}. Every sequence that occurs is a Gamma-ratio sequence
$\mathsf y_n$, a constant times a product of integer powers of Pochhammer ratios $(e)_n/(e')_n$. Its
large-$n$ expansion $\mathsf y_n\sim\sum_k\upsilon_kn^{-p_{\mathsf y}-k}$ follows from the
Bernoulli-polynomial expansion of $\log\Gamma(n+e)-\log\Gamma(n+e')$ \cite{gff:DLMFStirling}.
\emph{Single sums:}
$\sum_{n\geq0}\mathsf y_n=\sum_{n<N_{\rm H}}\mathsf y_n+\sum_k\upsilon_k\,\zeta(p_{\mathsf y}+k,N_{\rm H})$,
with Hurwitz zeta values \cite{gff:DLMFHurwitz} computed by Euler--Maclaurin summation in which every term
carries the factor $N_{\rm H}^{-p_{\mathsf y}-k}$, so that they have relative precision for complex
exponents. (The routine \texttt{mpmath.zeta(s,a)} has an absolute error floor of about $10^{-P}$ at $P$
working digits for complex $s$ and is not used.) The same formula, with the Hurwitz values continued in
the exponent, is the analytic continuation (the zeta-regularized value) when the series diverges.
\emph{Two-variable sums:} $A_r^+$ and $C_r$ have the form $\sum_{m\geq0}\mathsf x_m\mathsf q_m$ with a
Cauchy convolution $\mathsf q_m=\sum_{n=0}^m\mathsf y_{m-n}\mathsf z_n$. For $A_r^+$,
$\mathsf x_m=[(1+\sigma+t)_m/(1+2q+t)_m]^r$, $\mathsf y_j=d_{j+1}/(4q)$ and
$\mathsf z_n=[(1-g+t)_n/(1+t)_n]^r$; for $C_r$, $\mathsf x_m=d_{m+1}/(4q)$,
$\mathsf y_j=[(1-t)_j/(1+g-t)_j]^r$ and $\mathsf z_k=[(\sigma+t)_k/(2q+t)_k]^r$. The code computes
$\mathsf q_m$ exactly for $m<M_{\rm H}$ and uses its large-$m$ expansion for $m\geq M_{\rm H}$. If
$\mathsf y_j\sim\sum_a\upsilon_aj^{-p_{\mathsf y}-a}$ and $\mathsf z_n\sim\sum_c\upsilon'_cn^{-p_{\mathsf z}-c}$,
the generating function of $\mathsf y$ at $x=e^{-w}\to1$ has the singular part
$\sum_a\upsilon_a\Gamma(1-p_{\mathsf y}-a)w^{p_{\mathsf y}+a-1}$ and the regular part
$\sum_k\mathsf m^{\mathsf y}_k(-w)^k/k!$, where $\mathsf m^{\mathsf y}_k$ is the zeta-regularized moment
``$\sum_jj^k\mathsf y_j$'' (Mellin asymptotics of harmonic sums \cite{gff:FlajoletGourdonDumas}; transfer at
the singularity $x=1$ \cite{gff:FlajoletOdlyzko}), and likewise for $\mathsf z$. Multiplying and
transferring back gives three families,
\begin{equation}\begin{aligned}
 \mathsf q_m&\sim\sum_i\bigl[\mathfrak e^{(1)}_im^{-p_{\mathsf y}-i}+\mathfrak e^{(2)}_im^{-p_{\mathsf z}-i}
 +\mathfrak e^{(3)}_im^{1-p_{\mathsf y}-p_{\mathsf z}-i}\bigr],\\
 \mathfrak e^{(1)}_i&=\sum_{a+k=i}\upsilon_a\frac{(p_{\mathsf y}+a)_k}{k!}\mathsf m^{\mathsf z}_k,\qquad
 \mathfrak e^{(2)}_i=\sum_{c+k=i}\upsilon'_c\frac{(p_{\mathsf z}+c)_k}{k!}\mathsf m^{\mathsf y}_k,\\
 \mathfrak e^{(3)}_i&=\sum_{a+c=i}\upsilon_a\upsilon'_c\frac{\Gamma(1-p_{\mathsf y}-a)\Gamma(1-p_{\mathsf z}-c)}
 {\Gamma(2-p_{\mathsf y}-p_{\mathsf z}-a-c)} .
\end{aligned}\label{gff:mgf:threeFamilies}\end{equation}
The tail $\sum_{m\geq M_{\rm H}}\mathsf x_m\mathsf q_m$ is then a finite combination of three batches of
Hurwitz values. For the multiplicity sequence at non-integer $d$ the regularized moments are known in
closed form. By \eqref{gff:gdr:Dgen},
$\sum_j(d_{j+1}/4q)e^{-jw}=[(1+e^{-w})(1-e^{-w})^{1-d}-1]e^w/(4q)$, whose first term has only
non-integer powers of $w$; hence
\begin{equation}
 \mathsf m^{\mathsf y}_k=\frac{(-1)^{k+1}}{4q}\qquad\bigl(\mathsf y_j=d_{j+1}/(4q),\ d\notin\mathbb Z\bigr),
 \label{gff:mgf:muD}
\end{equation}
which the numerical moments reproduce. The other moments are computed with a head $N_{\rm H}\le
M_{\rm H}/2$, so that the cancellation in $\sum_{n<N_{\rm H}}n^k\mathsf z_n$ is suppressed by
$M_{\rm H}^{-k}$ in the tail. \emph{Exceptional parameters:} for non-integer $d$ the coefficients in
\eqref{gff:mgf:threeFamilies} have removable poles when $r(d-b)/2\in\mathbb Z$ (the $\log m$ resonance of
two families). At such points \texttt{W} returns the mean over $8$ points of a circle of radius about
$10^{-(P+10)/8}$ around $b$, which is exact for analytic $W_r$ up to the eighth power of the radius,
with extra working digits. At $d=1,2,3$ the multiplicities are polynomial in $j$ (or $\delta_{j0}$);
then only convergent moments enter, the coefficient of the third family is
$\upsilon_a\upsilon'_c(-p_{\mathsf y}-a)!/(1-p_{\mathsf z}-c)_{1-p_{\mathsf y}-a}$, the middle block is a finite sum of
products of single moments, and no resonance occurs inside the strip. \emph{Parameters and cost:} the
defaults are $M_{\rm H}=\max(80,1.6P')$, $K_{\rm H}=\max(30,0.8P'+10)$ asymptotic terms and
$N_{\rm H}=M_{\rm H}/2$, with $P'=P+15$ guard digits at $P$ requested digits. One $W_4$ at $50$ digits
takes about $45$\,s on one Intel core of the symmetry cluster, and all $846$ evaluations of the checks
of the general-$b$ code in Section~\ref{gff:gdr:checks} took $6$ minutes on one $64$-core AMD node. The
accuracy of this method is established by convergence in $(M_{\rm H},K_{\rm H},N_{\rm H})$ and by the
independent checks, not by interval bounds. \emph{Independent spectral code:}
\path{LRI_C5_general_d/foundation/spectral.py} evaluates \eqref{gff:gdr:kernel} on the Mellin line, by
tanh--sinh quadrature for the spins $\ell<40$ and, for $\ell\geq40$, by the Stirling expansion in
$z_\ell=\frac\ell2+\ii\nu$ with
$\frac1{2\pi}\int_{\mathbb R}z_\ell^{-e}\bar z_\ell^{-e'}d\nu=\ell^{1-e-e'}\Gamma(e+e'-1)/(\Gamma(e)\Gamma(e'))$
and a Hurwitz-zeta spin tail. It shares no code with the finite-sample formula except the Hurwitz
routine.

\paragraph{Beyond $b_{\max}$.}
For $d>2$ one has $b_{\max}<b_{\rm spec}$; the $\ell$-endpoint condition (ii) binds. The continuation
meant here is the continuation in $b$ of each block value at a fixed node $t$. Each block of
\eqref{gff:mgf:blocks} is a finite part, plus Hurwitz zeta values with exponents affine in $b$, plus
convergent remainders, and the evaluation above computes exactly this continuation (zeta regularization
of the tails). So the samples $\Xi_r(t_j;b)$ and the $c_{r,k}$ continue meromorphically in $b$. On
$b_{\max}\le b<b_{\rm spec}$ their only pole is at $b_{\max}$: the next candidates $b_{\max}+2k/r$
exceed $b_{\rm spec}$, since $b_{\rm spec}-b_{\max}=(d-2)/r<2/r$. Both sides of \eqref{gff:mgf:Wsample}
are meromorphic and agree on the strip, so the identity persists on $b_{\max}\le b<b_{\rm spec}$ with
continued block values. The pole at $b_{\max}$ was derived in the review
(\path{LRI_C5_general_d/review/foundation-rigor}); the derivation is an endpoint asymptotic sketch, not a
written proof. The divergence at $b_{\max}$ comes from the ends $\ell\to\infty$ of the lattice at fixed
$n$ and at fixed $\ell+n$. There $d_\ell\sim2\ell^{d-2}/\Gamma(d-1)$ by \eqref{gff:mgf:multiplicityBound},
$f(\ell+c)^r\sim\ell^{-rg}$, and $\sum_\ell\ell^{d-2-rg}$ has the pole
$1/(rg-d+1)=-2/(r(b-b_{\max}))$. Collecting these ends gives
\begin{equation}\begin{gathered}
 \Res_{b=b_{\max}}\Xi_r(t;b)=-\frac{4}{r\,\Gamma(d-1)}\,D_b(t)^r\bigl[\Phi_{\max}(-t)+\Phi_{\max}(\sigma+t)\bigr],\\
 \Phi_{\max}(u)=\sum_{m\in\mathbb Z}\Bigl[\frac{\Gamma(m+u)}{\Gamma(m+u+g)}\Bigr]^r,
\end{gathered}\label{gff:mgf:bmaxres}\end{equation}
at $\sigma=b_{\max}/2$ and $g=(d-b_{\max})/2$, where $rg=d-1>1$. The residue has three properties.
(1) Its top mode is proportional to $\cos(\pi(rg+2-d)/2)$, which vanishes at $b_{\max}$; this follows
from Steps~1, 3 and 5 of the proof of \eqref{gff:gdr:topG} applied to $\Phi_{\max}$, with $a_\pm(2\pi)$
finite because $rg>1$, and it is consistent with \eqref{gff:gdr:topless} holding up to the boundary. (2)
Its lower modes are nonzero in every case computed, so the continued $c_{r,k}$ with $k\le r-2$ have a
simple pole. (3) The functional annihilates it: $\Phi_{\max}(\sigma+t)$ is regular at $t=0$, and
$\Res_{t=0}\Phi_{\max}(-t)=-\sum_{n\geq0}\Res_{z=-n}f(z)^r=0$, because $f^r$ has poles only in
$-\mathbb N_0$ and is $O(|z|^{-rg})$ on the circles $|z|=n_0+\frac12$, with $rg>1$. Hence $W_r$ has no
pole at $b_{\max}$, as the spectral representation requires. The same review checked this numerically at
$20$ digits for $(d,r)=(2.5,4)$, $(2.75,3)$, $(3.5,4)$ and $(3.5,3)$, with its scripts
\path{bmax_residue.py}, \path{bmax_normalization.py} and \path{bmax_compare.py} (symmetry job 569252 for the first
two): the top mode of
\eqref{gff:mgf:bmaxres} is $5\cdot10^{-21}$ to $2\cdot10^{-18}$ of the largest mode, and the lower modes are
of order one; the functional applied to \eqref{gff:mgf:bmaxres} vanishes to $2\cdot10^{-21}$--$10^{-18}$
relative; at $d=2.5$, $r=4$ the predicted residues of $c_{4,0},\dots,c_{4,3}$,
$(-9.109258,\,-11.819680,\,-2.719024,\,0)$, agree with the estimate
$\frac\delta2[c_{4,k}(b_{\max}+\delta)-c_{4,k}(b_{\max}-\delta)]$ from the continued code to
$O(\delta^2)$, at most $4\cdot10^{-5}$ at $\delta=10^{-3}$ and $4\cdot10^{-7}$ at $\delta=10^{-4}$. This
also confirms that the zeta-regularized evaluation is the analytic continuation. The continued formula
agrees with the spectral values at $(d,b)=(2.5,1.8)$ and $(2.5,1.76)$ for $r=4$ and $(2.5,1.55)$ for
$r=3$ to $10^{-14}$ at $15$ digits, and at $50$ digits as in Section~\ref{gff:gdr:checks}. For $d<1$
condition (iii) binds and coincides with the spectral bound. Nothing downstream uses the continuation:
every use of these formulas in this note is at $b=\frac d2$, inside the strip.


\subsection{Samples, modes and the two masters}\label{gff:gdr:samples}
By \eqref{gff:gdr:topless} and \eqref{gff:gdr:explicit}, $W_4$ needs three modes and $W_3$ two. In the variable
$\tilde t=t+\frac b4$ the trigonometric polynomial is $\sum_kc_{r,k}\cos2\pi k\tilde t$, even about $\tilde t=0$ and
$\tilde t=\frac12$. We sample it at the pole node $t=0$ ($\tilde t=\frac b4$) and at these reflection centers.

\paragraph{Four-wheel.}
Three real samples fix the three modes:
\begin{equation}
 \mathcal X_T(b)=N_b^4\Xi_4(0;b),\qquad
 \mathcal X_0(b)=N_b^4\Xi_4(-\tfrac b4;b),\qquad
 \mathcal X_{1/2}(b)=N_b^4\Xi_4(\tfrac12-\tfrac b4;b).
 \label{gff:gdr:samples4}
\end{equation}
With $x=b/4$, $y=d/2-b/4$ (in the definitions of $\mathfrak h_u,\mathfrak s_u$ only; elsewhere in this section $t=x+\ii y$, and $x$ is the argument of $\Omega_r(t;x)$ in Section~\ref{gff:gdr:special}) they are five two-variable Horn functions at argument $(1,1)$:
$\mathcal X_T=(\pi s_b)^4\mathfrak h_x$,
$\mathcal X_0=\kappa_0^4(\mathfrak s_0+2\xi_0^4\mathfrak h_1)$,
$\mathcal X_{1/2}=\kappa_h^4(\mathfrak s_{1/2}+2\xi_h^4\mathfrak h_{1/2})$, where
\[
 \mathfrak h_u=\sum_{\ell,n\ge0}d_\ell\Bigl[\frac{(u+x)_{\ell+n}}{(u+y)_{\ell+n}}\Bigr]^4
 \Bigl[\frac{(u+1-y)_n}{(u+1-x)_n}\Bigr]^4,\qquad
 \mathfrak s_u=\sum_{j,k\ge0}d_{j+k+2u}\Bigl[\frac{(u+x)_j(u+x)_k}{(u+y)_j(u+y)_k}\Bigr]^4,
\]
$\kappa_0=N_b\sin^2(\pi x)\Gamma(x)^2/\Gamma(y)^2$, $\kappa_h=N_b\cos^2(\pi x)\Gamma(\frac12+x)^2/\Gamma(\frac12+y)^2$,
$\xi_0=\frac{x(1-y)}{y(1-x)}$, $\xi_h=\frac{1-2y}{1-2x}$. The index $u$ of $\mathfrak h_u$ and $\mathfrak s_u$ is unrelated to the weight $u(b)$ of \eqref{gff:gdr:twonode}. In particular
$\mathfrak h_x(d;b)=\sum d_\ell T_{\ell n}(b)^4$ with
\begin{equation}
 T_{\ell n}(b)=\frac{(\sigma)_{\ell+n}(1-g)_n}{(2q)_{\ell+n}\,n!},
 \label{gff:gdr:Tln}
\end{equation}
which at $b=d/2$ is the $T_{\ell n}$ of \eqref{gff:gs:seriesBase}.

Section~\ref{gff:mgl:section} proves these sample formulas for every $r$, \eqref{gff:mgl:sampT}--\eqref{gff:mgl:samph}.

By \eqref{gff:gdr:topless} the samples are $N_b^{-4}\mathcal X_0=c_{4,0}+c_{4,1}+c_{4,2}$,
$N_b^{-4}\mathcal X_{1/2}=c_{4,0}-c_{4,1}+c_{4,2}$ and $N_b^{-4}\mathcal X_T=c_{4,0}+\varsigma_bc_{4,1}+c_{4,2}\cos\pi b$.
This system has determinant $4\sin^2\pi\sigma\ne0$. Solving it and inserting into \eqref{gff:gdr:explicit} gives, for
every $b$ in the strip,
\begin{equation}
 W_4(d;b)=\frac{2}{\pi s_b^7}\Bigl[(1+\varsigma_b)^2(1+2\varsigma_b)\mathcal X_0-(1-\varsigma_b)^2(1-2\varsigma_b)
 \mathcal X_{1/2}-2\varsigma_b\mathcal X_T\Bigr].
 \label{gff:gdr:threenode}
\end{equation}
At $b=d/2$, and differentiating once (the product rule done once; all samples at $b=d/2$),
\begin{gather}
 P_4(d)=\frac{2}{\pi S^7}\Bigl[(1+C)^2(1+2C)\mathcal X_0-(1-C)^2(1-2C)\mathcal X_{1/2}-2C\mathcal X_T\Bigr],
 \label{gff:gdr:P4}\\
 \begin{aligned}
 \LL_d={}&-\frac1{S^8}\Bigl[(1+C)^2(4+9C+8C^2)\mathcal X_0+(1-C)^2(4-9C+8C^2)\mathcal X_{1/2}
 -2(1+6C^2)\mathcal X_T\Bigr]\\&+\frac{2}{\pi S^7}\,\mathfrak D_L,
 \end{aligned}
 \label{gff:gdr:L}\\
 \mathfrak D_L=(1+C)^2(1+2C)\,\mathcal X_0'-(1-C)^2(1-2C)\,\mathcal X_{1/2}'-2C\,\mathcal X_T'.
 \label{gff:gdr:frozen}
\end{gather}
Written out in the five Horn values and their first $b$-derivatives, \eqref{gff:gdr:L} is \eqref{gff:mgl:Lfinal}, derived in Section~\ref{gff:mgl:section}. The derivative
data enter only through the single frozen-weight combination $\mathfrak D_L$, the $b$-derivative of the bracket of
\eqref{gff:gdr:P4} with its coefficients held at $b=d/2$. The first line of \eqref{gff:gdr:L} is the derivative of the
coefficients of \eqref{gff:gdr:threenode}.

\paragraph{Three-wheel.}
Two real samples suffice (proofs in Section~\ref{gff:mgz:nodes}): the pole node
$t=0$ and the self-mirror node $t_b=\frac12-\frac b4$, where $1-\sigma-t_b=t_b$ and $\tilde t=\frac12$. At $t=0$ only the leading
coefficients of the triple poles survive, and at $t_b$ the kernel is sampled on the midpoint lattice, because
$a_\ell+n+t_b=n+\frac{\ell+1}2$. With
\begin{equation}
 \mathfrak t_3(b)=\sum_{\ell,n\ge0}d_\ell T_{\ell n}(b)^3,\qquad
 \widehat M_3(b)=\sum_{\ell\ge0}d_\ell\!\!\sum_{s\in\frac{\ell+1}2+\mathbb Z}\!\![N_bg_\ell(s;b)]^3 ,
 \label{gff:gdr:samples3}
\end{equation}
one has $N_b^3\Xi_3(0;b)=-(\pi s_b)^3\mathfrak t_3$ and
$N_b^3\Xi_3(t_b;b)=(\frac{1+\varsigma_b}2)^3\widehat M_3$. Note that $\widehat M_3$ is the bare lattice sum: the
factor $D_b(t_b)^3=(1+\varsigma_b)^3/8$ carried by the sample is written separately. The two samples are
$c_{3,0}+\varsigma_bc_{3,1}$ and $c_{3,0}-c_{3,1}$, and \eqref{gff:gdr:explicit} gives
\begin{equation}
 W_3(d;b)=u(b)\mathfrak t_3(b)+v(b)\widehat M_3(b),\qquad
 u(b)=\frac{2\pi^2(1+2\varsigma_b)}{s_b^2},\qquad v(b)=\frac{\varsigma_b(1+\varsigma_b)^2}{2\pi s_b^3}.
 \label{gff:gdr:twonode}
\end{equation}
$\mathfrak t_3$ is a ${}_7F_6(1)$ plus one two-variable Horn value, $\widehat M_3$ a ${}_7F_6(1)$ plus two. In the
block decomposition of Section~\ref{gff:gdr:residue}, $\widehat M_3$ is the sum of a right-part sum (the ${}_7F_6(1)$ and one
Horn value) and a middle-part sum (one Horn value); the count of Section~\ref{gff:gdr:C5} treats these two as separate
functions. In terms of
the normalized zero mode
$\hat c_0=N_b^3c_{3,0}=\frac{\varsigma_b(1+\varsigma_b)^2}8\widehat M_3-\frac{\pi^3s_b^3}{1+\varsigma_b}\mathfrak t_3$,
\eqref{gff:gdr:twonode} reads $W_3=\frac{6\pi^2}{s_b^2}\mathfrak t_3+\frac{4}{\pi s_b^3}\hat c_0$. Then, at $b=d/2$,
\begin{gather}
 3Q_d=\frac{6\pi^2}{S^2}\mathfrak t_3+\frac{4}{\pi S^3}\hat c_0,
 \label{gff:gdr:Q}\\
 Z_d=\frac{2\pi^2}{3S^2}\mathfrak t_3''-\frac{\pi^2\cos(\pi d/2)}{6S^2}Q_d
 +\frac19\Bigl[\frac{\pi(1+C^2)}{S^5}\hat c_0-\frac{4C}{S^4}\hat c_0'+\frac{4}{\pi S^3}\hat c_0''\Bigr].
 \label{gff:gdr:Z}
\end{gather}
Here $\partial_bW_3(d;d/2)=0$ has been used to remove $\mathfrak t_3'$. Its proof is Pattern~T of Section~\ref{gff:gs:patterns},
\eqref{gff:gs:Tjet}, together with \eqref{gff:gs:wheelJets}: on the balanced $K_5$
subspace there is no linear form invariant under the permutation group $S_5$ (not the lattice sum $S_r$ of \eqref{gff:gdr:periodizationDef} at $r=5$), and the balanced three-wheel path is a balanced variation. (\eqref{gff:gs:LZspectral} records the vanishing
only as a check.) The formula \eqref{gff:gdr:Z}, and the $Z_d$ part of \eqref{gff:gdr:C5reduced}, depend on this input; the
full jet formula $9Z_d=(u\mathfrak t_3+v\widehat M_3)''$ and \eqref{gff:gdr:Q} do not. At $d=1,2$ the input also follows
from the closed forms of Section~\ref{gff:gdr:special}. The second derivatives enter only through
$\frac19(u\mathfrak t_3''+v\widehat M_3'')$ with $u=\frac{2\pi^2(1+2C)}{S^2}$, $v=\frac{C(1+C)^2}{2\pi S^3}$, and
the first derivatives only through $\hat c_0'$.

\paragraph{Analyticity and termwise derivatives.}
The five Horn functions, $\mathfrak t_3$ and $\widehat M_3$ are analytic on the strip, and their $b$-derivatives
are the termwise derivatives. The bounds of Section~\ref{gff:gdr:periodization} hold uniformly for $b$ in compact
subsets of the strip, with $\alpha$ above the threshold $\max(\frac d2,d-1,\frac12)$ by a fixed margin; the
polygamma weights produced by differentiation grow only logarithmically and cost part of that margin
(for the five Horn functions, the proof of \eqref{gff:mgl:dseries} in Section~\ref{gff:mgl:section}; for the three-wheel, Section~\ref{gff:mgz:uniform}).


\subsection{Proofs of the four-wheel samples and the explicit formula for \texorpdfstring{$\LL_d$}{Ld}}\label{gff:mgl:section}
The code \path{LRI_C5_general_d/L/L_general.py} evaluates the formulas of this subsection. It proves the
sample formulas stated after \eqref{gff:gdr:samples4} for every $r\geq2$, writes \eqref{gff:gdr:threenode} and \eqref{gff:gdr:L}
out in the five Horn values and their first $b$-derivatives, and collects the checks of that form
(Section~\ref{gff:mgl:checks}).

\paragraph{Setting and status.}
Throughout, $r\geq2$, $0<d<4$, $0<\operatorname{Re}b<b_{\max}(r,d)$ of \eqref{gff:gdr:bmax} and $\sigma\notin\mathbb Z$, the
hypotheses of Section~\ref{gff:gdr:finite}. As in Section~\ref{gff:gdr:samples},
\begin{equation}
 x=\frac b4,\qquad y=\frac d2-\frac b4,\qquad \sigma=2x,\qquad g=y-x,\qquad x+y=2q,\qquad \tilde t=t+x,
 \label{gff:mgl:xy}
\end{equation}
and in this subsection $x,y$ always have this meaning. For $d\neq2$ the multiplicities \eqref{gff:gs:spectrum} have the
very-well-poised form
\begin{equation}
 d_\ell=\frac{(4q-2)_\ell\,(2q)_\ell}{(2q-1)_\ell\,\ell!},
 \label{gff:mgl:mult}
\end{equation}
since $\Gamma(d-1)=(d-2)\Gamma(d-2)$, $(2\ell+d-2)/(d-2)=(2q)_\ell/(2q-1)_\ell$ and
$\Gamma(\ell+d-2)/(\Gamma(d-2)\,\ell!)=(4q-2)_\ell/\ell!$; here $(2q-1)_\ell\neq0$ because $2q-1\in(-1,1)$ and
$2q\neq1$. The form \eqref{gff:fp:degeneracy} has no $0/0$ at $d=2$, where $d_\ell=2-\delta_{\ell0}$.
Status: the lattice form \eqref{gff:mgl:tauform}--\eqref{gff:mgl:evenform}, the samples
\eqref{gff:mgl:sampT}--\eqref{gff:mgl:samph} and the derivative series \eqref{gff:mgl:dseries} are proved here, from the
estimates of Section~\ref{gff:gdr:periodization} and the finite Fourier theorem of Section~\ref{gff:gdr:finite}. The formulas
\eqref{gff:mgl:W4five}, \eqref{gff:mgl:W3three}, \eqref{gff:mgl:P4five} and \eqref{gff:mgl:Lfinal} are derived: exact algebra on these
statements, on the top-mode theorem \eqref{gff:gdr:topless} and on the residue functional \eqref{gff:gdr:explicit}, checked
symbolically. All numbers in Section~\ref{gff:mgl:checks} are numerical. A separate review with its own codes
(\path{LRI_C5_general_d/review/L-rigor/}) rederived every step of the package's own derivation by hand and
symbolically and found no wrong formula (for the steps added in this note, see Section~\ref{gff:gdr:status}); its
corrections (the continuity step at $t=0$ in the proof of \eqref{gff:mgl:sampT}, the restriction of the three-dimensional
moment statement to $\mathfrak s_0,\mathfrak s_{1/2}$, and the list of digamma arguments) are incorporated below.

\paragraph{Lattice sums in the variable $\tilde t$.}
Let
\begin{equation}
 \varphi(z)=\frac{\Gamma(z+x)}{\Gamma(z+y)},\qquad \bar\varphi(z)=\frac{\Gamma(z+1-y)}{\Gamma(z+1-x)},
 \label{gff:mgl:phis}
\end{equation}
\begin{equation}\begin{gathered}
 \mathcal R_+(u)=\sum_{\ell,n\geq0}d_\ell\bigl[\varphi(\ell+n+u)\,\bar\varphi(n+u)\bigr]^r,\\
 \mathcal R_{\rm m}(\tilde t)=\sum_{j,k\geq0}d_{j+k+1}\bigl[\varphi(k+\tilde t)\,\varphi(j+1-\tilde t)\bigr]^r,\qquad
 \mathcal R_{\rm s}(\tilde t)=\sum_{j,k\geq0}d_{j+k}\bigl[\varphi(k+\tilde t)\,\varphi(j-\tilde t)\bigr]^r,
\end{gathered}\label{gff:mgl:sums}\end{equation}
\begin{equation}\begin{gathered}
 P_+(\tilde t)=\bigl[\sin\pi(\tilde t+x)\sin\pi(\tilde t-y)\bigr]^r,\qquad
 P_-(\tilde t)=\bigl[\sin\pi(\tilde t-x)\sin\pi(\tilde t+y)\bigr]^r,\\
 P_{\rm m}(\tilde t)=\bigl[\sin\pi(\tilde t-x)\sin\pi(\tilde t+x)\bigr]^r=D_b(t)^r .
\end{gathered}\label{gff:mgl:weights}\end{equation}
For
$t=\tilde t-x$ off the pole classes $t\equiv0$ and $t\equiv-\sigma$,
\begin{align}
 \Xi_r(t;b)&=P_+(\tilde t)\,\mathcal R_+(\tilde t)+P_-(\tilde t)\,\mathcal R_+(1-\tilde t)
 +P_{\rm m}(\tilde t)\,\mathcal R_{\rm m}(\tilde t)
 \label{gff:mgl:tauform}\\
 &=P_+(\tilde t)\,\mathcal R_+(1+\tilde t)+P_-(\tilde t)\,\mathcal R_+(1-\tilde t)
 +P_{\rm m}(\tilde t)\,\mathcal R_{\rm s}(\tilde t),
 \label{gff:mgl:evenform}
\end{align}
with absolutely convergent series. The second form is manifestly even in $\tilde t$: $P_+(-\tilde t)=P_-(\tilde t)$,
$P_{\rm m}$ is even, and $\mathcal R_{\rm s}$ is even (exchange $j$ and $k$).

\emph{Proof.} Split the lattice of \eqref{gff:gdr:periodizationDef} into the right part $n\geq0$, the middle part
$\ell=j+k+1$, $n=-j-1$ ($j,k\geq0$) and the left part $n\le-\ell-1$, as in Section~\ref{gff:gdr:periodization}. By the
estimates there and Section~\ref{gff:gdr:finite}(a) the three parts converge absolutely off the pole classes, so they may
be summed separately and rearranged.

\emph{Right part.} By the first line of \eqref{gff:gdr:reflected}, $g_\ell(a_\ell+n+t)=f(\ell+\sigma+n+t)\,f(-n-t)$ with
$f(z)=\Gamma(z)/\Gamma(z+g)$. Since $\sigma+t=\tilde t+x$ and $x+g=y$, $f(\ell+\sigma+n+t)=\varphi(\ell+n+\tilde t)$.
Gamma reflection in the form $\Gamma(z-n)\Gamma(1-z+n)=(-1)^n\pi/\sin\pi z$ \cite{gff:DLMF5}, used for $z=x-\tilde t$ and
$z=y-\tilde t$, gives
\[
 f(-n-t)=\frac{\Gamma(x-\tilde t-n)}{\Gamma(y-\tilde t-n)}
 =\frac{\sin\pi(y-\tilde t)}{\sin\pi(x-\tilde t)}\,\bar\varphi(n+\tilde t).
\]
Since $D_b(t)=\sin\pi(\tilde t-x)\sin\pi(\tilde t+x)$, the prefactor is
$D_b(t)\sin\pi(y-\tilde t)/\sin\pi(x-\tilde t)=\sin\pi(\tilde t+x)\sin\pi(\tilde t-y)$, and the right part is
$P_+(\tilde t)\mathcal R_+(\tilde t)$.

\emph{Left part.} By the last line of \eqref{gff:gdr:reflected}, $n\mapsto n'=-\ell-n-1\geq0$ maps the left part
bijectively onto the right part at the mirror point $t'=1-\sigma-t$, that is at $\tilde t'=1-\tilde t$. Since
$D_b(1-\sigma-t)=\sin\pi(\sigma+t)\sin\pi t=D_b(t)$, the left part is $P_+(1-\tilde t)\mathcal R_+(1-\tilde t)$, and
$P_+(1-\tilde t)=P_-(\tilde t)$.

\emph{Middle part.} In the second line of \eqref{gff:gdr:reflected}, $k+\sigma+t=k+\tilde t+x$, $k+2q+t=k+\tilde t+y$,
$j+1-t=j+1-\tilde t+x$ and $j+1+g-t=j+1-\tilde t+y$, so the summand is
$[\varphi(k+\tilde t)\varphi(j+1-\tilde t)]^r$ with multiplicity $d_{j+k+1}$, and the prefactor is
$D_b(t)^r=P_{\rm m}(\tilde t)$. This proves \eqref{gff:mgl:tauform}.

\emph{Second form.} $\mathcal R_+(\tilde t)-\mathcal R_+(1+\tilde t)=\sum_\ell d_\ell[\varphi(\ell+\tilde t)\,
\bar\varphi(\tilde t)]^r$ is the column $n=0$. From $\sin\pi(\tilde t-y)\,\Gamma(1-y+\tilde t)=-\pi/\Gamma(y-\tilde t)$
and $1/\Gamma(1-x+\tilde t)=\sin\pi(x-\tilde t)\,\Gamma(x-\tilde t)/\pi$,
\begin{equation}
 \sin\pi(\tilde t+x)\sin\pi(\tilde t-y)\,\bar\varphi(\tilde t)
 =\sin\pi(\tilde t+x)\sin\pi(\tilde t-x)\,\varphi(-\tilde t),
 \label{gff:mgl:colid}
\end{equation}
so $P_+(\tilde t)[\mathcal R_+(\tilde t)-\mathcal R_+(1+\tilde t)]
=P_{\rm m}(\tilde t)\sum_\ell d_\ell[\varphi(\ell+\tilde t)\varphi(-\tilde t)]^r$. This is the row $j=0$ of
$\mathcal R_{\rm s}(\tilde t)$, and the rows $j\geq1$ of $\mathcal R_{\rm s}(\tilde t)$ are $\mathcal R_{\rm m}(\tilde t)$
after $j\to j+1$. All rearrangements are of absolutely convergent series. $\square$

\paragraph{The three samples.}
Normalize the sums by their first terms: for $u\in\{x,\frac12,1\}$ and $u\in\{0,\frac12\}$ respectively,
\begin{equation}
 \mathfrak h^{(r)}_u=\sum_{\ell,n\geq0}d_\ell\Bigl[\frac{(u+x)_{\ell+n}}{(u+y)_{\ell+n}}\Bigr]^r
 \Bigl[\frac{(u+1-y)_n}{(u+1-x)_n}\Bigr]^r,\qquad
 \mathfrak s^{(r)}_u=\sum_{j,k\geq0}d_{j+k+2u}\Bigl[\frac{(u+x)_j\,(u+x)_k}{(u+y)_j\,(u+y)_k}\Bigr]^r,
 \label{gff:mgl:five}
\end{equation}
so that $\mathcal R_+(u)=[\varphi(u)\bar\varphi(u)]^r\mathfrak h^{(r)}_u$,
$\mathcal R_{\rm s}(0)=\varphi(0)^{2r}\mathfrak s^{(r)}_0$ and
$\mathcal R_{\rm m}(\frac12)=\varphi(\frac12)^{2r}\mathfrak s^{(r)}_{1/2}$. For $r=4$ these are the $\mathfrak h_u$,
$\mathfrak s_u$ of Section~\ref{gff:gdr:samples}, and the superscript is dropped; for $r=3$, $\mathfrak h^{(3)}_x$ is
$\mathfrak t_3$ of \eqref{gff:gdr:samples3}. At $u=x$ one has $u+x=\sigma$, $u+y=2q$, $u+1-y=1-g$ and $u+1-x=1$, so
$\mathfrak h^{(r)}_x=\sum_{\ell,n}d_\ell T_{\ell n}(b)^r$ with $T_{\ell n}(b)$ of \eqref{gff:gdr:Tln}. In the notation
\eqref{gff:fp:HornDefinition}, with $\boldsymbol e=(1,0)$, $\boldsymbol f=(0,1)$, $\boldsymbol c=\boldsymbol e+\boldsymbol f$,
the first summation index $\ell$ (or $j$), the second $n$ (or $k$), and argument $(1,1)$,
\begin{equation}\begin{aligned}
 \mathfrak h^{(r)}_u={}&\mathfrak H\bigl[\{(u+x,\boldsymbol c)^{\times r},(4q-2,\boldsymbol e),(2q,\boldsymbol e),
 (1,\boldsymbol f),(u+1-y,\boldsymbol f)^{\times r}\};\\
 &\qquad\{(u+y,\boldsymbol c)^{\times r},(2q-1,\boldsymbol e),(u+1-x,\boldsymbol f)^{\times r}\};(1,1)\bigr],\\
 \mathfrak s^{(r)}_0={}&\mathfrak H\bigl[\{(4q-2,\boldsymbol c),(2q,\boldsymbol c),(1,\boldsymbol e),(1,\boldsymbol f),
 (x,\boldsymbol e)^{\times r},(x,\boldsymbol f)^{\times r}\};\\
 &\qquad\{(2q-1,\boldsymbol c),(1,\boldsymbol c),(y,\boldsymbol e)^{\times r},(y,\boldsymbol f)^{\times r}\};(1,1)\bigr],\\
 \mathfrak s^{(r)}_{1/2}={}&4q\,\mathfrak H\bigl[\{(2q+1,\boldsymbol c),(4q-1,\boldsymbol c),(1,\boldsymbol e),
 (1,\boldsymbol f),(\tfrac12+x,\boldsymbol e)^{\times r},(\tfrac12+x,\boldsymbol f)^{\times r}\};\\
 &\qquad\{(2q,\boldsymbol c),(2,\boldsymbol c),(\tfrac12+y,\boldsymbol e)^{\times r},
 (\tfrac12+y,\boldsymbol f)^{\times r}\};(1,1)\bigr].
\end{aligned}\label{gff:mgl:arrays}\end{equation}
The first two use \eqref{gff:mgl:mult} and hold for $d\neq2$; the third uses \eqref{gff:fp:degeneracy} and holds for every
$d$. The upper pairs $(1,\boldsymbol e)$, $(1,\boldsymbol f)$ cancel the factorials of \eqref{gff:fp:HornDefinition}
where the summand has none. At $b=d/2$ ($x=\frac q2$, $y=\frac{3q}2$) the parameters are
\begin{equation}\begin{aligned}
 \mathfrak h_x&=\sum_{\ell,n}d_\ell\Bigl[\tfrac{(q)_{\ell+n}}{(2q)_{\ell+n}}\Bigr]^4\Bigl[\tfrac{(1-q)_n}{n!}\Bigr]^4
 =\sum_{\ell,n}d_\ell T_{\ell n}^4,&
 \mathfrak h_1&=\sum_{\ell,n}d_\ell\Bigl[\tfrac{(1+\frac q2)_{\ell+n}}{(1+\frac{3q}2)_{\ell+n}}\Bigr]^4
 \Bigl[\tfrac{(2-\frac{3q}2)_n}{(2-\frac q2)_n}\Bigr]^4,\\
 \mathfrak h_{1/2}&=\sum_{\ell,n}d_\ell\Bigl[\tfrac{(\frac{1+q}2)_{\ell+n}}{(\frac{1+3q}2)_{\ell+n}}\Bigr]^4
 \Bigl[\tfrac{(\frac{3-3q}2)_n}{(\frac{3-q}2)_n}\Bigr]^4,&
 \mathfrak s_0&=\sum_{j,k}d_{j+k}\Bigl[\tfrac{(\frac q2)_j(\frac q2)_k}{(\frac{3q}2)_j(\frac{3q}2)_k}\Bigr]^4,\\
 \mathfrak s_{1/2}&=\sum_{j,k}d_{j+k+1}\Bigl[\tfrac{(\frac{1+q}2)_j(\frac{1+q}2)_k}{(\frac{1+3q}2)_j(\frac{1+3q}2)_k}\Bigr]^4,
\end{aligned}\label{gff:mgl:fiveCenter}\end{equation}
with $T_{\ell n}$ of \eqref{gff:gs:seriesBase}. Thus $\mathfrak h_x$ is the unweighted version of the residue series
\eqref{gff:gs:QPseries}. With $\kappa_0,\kappa_h,\xi_0,\xi_h$ of Section~\ref{gff:gdr:samples}, that is
$\kappa_0=N_b\sin^2(\pi x)\varphi(0)^2$, $\kappa_h=N_b\cos^2(\pi x)\varphi(\frac12)^2$,
$\xi_0=x(1-y)/(y(1-x))$ and $\xi_h=(1-2y)/(1-2x)$, the samples are
\begin{align}
 N_b^r\,\Xi_r(0;b)&=(-1)^r(\pi s_b)^r\,\mathfrak h^{(r)}_x,\label{gff:mgl:sampT}\\
 N_b^r\,\Xi_r(-\tfrac b4;b)&=(-1)^r\kappa_0^{\,r}\bigl[\mathfrak s^{(r)}_0+2\xi_0^{\,r}\mathfrak h^{(r)}_1\bigr],
 \label{gff:mgl:samp0}\\
 N_b^r\,\Xi_r(\tfrac12-\tfrac b4;b)&=\kappa_h^{\,r}\bigl[\mathfrak s^{(r)}_{1/2}+2\xi_h^{\,r}\mathfrak h^{(r)}_{1/2}\bigr].
 \label{gff:mgl:samph}
\end{align}
For $r=4$ these are the expressions for $\mathcal X_T,\mathcal X_0,\mathcal X_{1/2}$ given after \eqref{gff:gdr:samples4}.
For $r=3$, \eqref{gff:mgl:sampT} is $N_b^3\Xi_3(0;b)=-(\pi s_b)^3\mathfrak t_3$, which Section~\ref{gff:mgz:nodes} also
proves directly, \eqref{gff:mgz:node0}, and comparison of \eqref{gff:mgl:samph} with the midpoint sample \eqref{gff:mgz:nodeb}
gives
$(\frac{1+\varsigma_b}2)^3\widehat M_3=\kappa_h^3[\mathfrak s^{(3)}_{1/2}+2\xi_h^3\mathfrak h^{(3)}_{1/2}]$: the right
and left parts, equal at $\tilde t=\frac12$, give $2\kappa_h^3\xi_h^3\mathfrak h^{(3)}_{1/2}$ (whose row $\ell=0$ is a
${}_7F_6(1)$), and the middle part gives $\kappa_h^3\mathfrak s^{(3)}_{1/2}$.

\emph{Proof.} (i) The node $t=0$, $\tilde t=x$. It lies in the pole class $t\equiv0$, so \eqref{gff:mgl:tauform} is first
extended to it. Assume $g\neq1$. Written in $t$, the terms of $\mathcal R_+(\tilde t)$ are
$[\Gamma(\ell+n+\sigma+t)\Gamma(n+1-g+t)/(\Gamma(\ell+n+2q+t)\Gamma(n+1+t))]^r$, the Gamma ratios of the first line of
\eqref{gff:gdr:reflected} without their prefactor $-\sin\pi(g-t)/\sin\pi t$, which has moved into $P_+$. The terms of
$\mathcal R_+(1-\tilde t)$ are the same ratios at $1-\sigma-t$,
$[\Gamma(\ell+n+1-t)\Gamma(n+2-2q-t)/(\Gamma(\ell+n+1+2q-\sigma-t)\Gamma(n+2-\sigma-t))]^r$, and the terms of
$\mathcal R_{\rm m}(\tilde t)$ are the ratios of the second line of \eqref{gff:gdr:reflected}. All their poles lie in
$-\sigma-\mathbb N_0$, $g-1-\mathbb N_0$, $\mathbb N=\{1,2,\dots\}$ and $2-2q+\mathbb N_0$. None of these sets contains
$0$, since $\operatorname{Re}\sigma>0$, $q<1$, $g\neq1$ and $\operatorname{Re}g<2$. Choose a closed disc $|t|\le\delta$
that avoids these sets and the class $-\sigma+\mathbb Z$. On it the shifts range over a compact set on which no Gamma
argument meets a pole, so the Gamma-ratio bound of Section~\ref{gff:gdr:periodization} holds uniformly, and the terms,
multiplicities included, are bounded by $C(1+\ell)^{d-2}(1+\ell+n)^{-\alpha}(1+n)^{-\alpha}$ (the first two sums) and
$C(1+j+k)^{d-2}(1+k)^{-\alpha}(1+j)^{-\alpha}$ (the middle sum), $\alpha=r\operatorname{Re}g$. By the lattice-sum criterion of Section~\ref{gff:gdr:periodization}, which holds
since $\operatorname{Re}b<b_{\max}$, the three series converge uniformly on the disc. Hence both sides of
\eqref{gff:mgl:tauform} are analytic on the disc ($\Xi_r$ is entire by Section~\ref{gff:gdr:finite}(b), the weights are
entire); they agree on the punctured disc, so they agree at $t=0$. At $\tilde t=x$ the weights $P_-$ and $P_{\rm m}$
vanish, so only the right sum survives, with $P_+(x)=[\sin\pi\sigma\,\sin\pi(x-y)]^r=(-1)^r[\sin\pi\sigma\sin\pi g]^r$.
With $\varphi(x)=\Gamma(\sigma)/\Gamma(2q)$, $\bar\varphi(x)=\Gamma(1-g)$ and $\Gamma(g)\Gamma(1-g)\sin\pi g=\pi$,
\[
 N_b\sin\pi\sigma\sin\pi g\;\varphi(x)\bar\varphi(x)=\pi\sin\pi\sigma=\pi s_b,
\]
which together with $\mathcal R_+(x)=[\varphi(x)\bar\varphi(x)]^r\mathfrak h^{(r)}_x$ gives \eqref{gff:mgl:sampT}. At the
isolated point $g=1$, that is $b=d-2$ (in the strip only for $d>2$), $\bar\varphi(x)$ has a pole. There
\eqref{gff:mgl:sampT} follows by continuity in $b$, since both sides are analytic in $b$: the left side by
Section~\ref{gff:gdr:finite}(d), the right side because the terms of $\mathfrak h^{(r)}_x=\sum d_\ell T_{\ell n}(b)^r$ are
polynomials in $b$ and obey the same majorant uniformly near $b=d-2$. The masters use only $b$ near $\frac d2$, where
$g=q<1$.

(ii) The node $\tilde t=0$, $t=-\frac b4$. It is off the pole classes, since $t\equiv0$ or $t\equiv-\sigma$ would need
$\frac b4\in\mathbb Z$. In \eqref{gff:mgl:evenform}, $P_+(0)=P_-(0)=(-1)^r[\sin\pi x\sin\pi y]^r$ and
$P_{\rm m}(0)=(-1)^r\sin^{2r}\pi x$, so
$\Xi_r=(-1)^r\{2[\sin\pi x\sin\pi y\,\varphi(1)\bar\varphi(1)]^r\mathfrak h^{(r)}_1
+[\sin^2(\pi x)\varphi(0)^2]^r\mathfrak s^{(r)}_0\}$.

(iii) The node $\tilde t=\frac12$, $t=\frac12-\frac b4$. It is off the pole classes, since either class would need
$b\equiv2\pmod4$, that is $\sigma\in\mathbb Z$. In \eqref{gff:mgl:tauform}, $P_+(\frac12)=P_-(\frac12)=[\cos\pi x\cos\pi y]^r$
and $P_{\rm m}(\frac12)=\cos^{2r}\pi x$, so
$\Xi_r=2[\cos\pi x\cos\pi y\,\varphi(\frac12)\bar\varphi(\frac12)]^r\mathfrak h^{(r)}_{1/2}
+[\cos^2(\pi x)\varphi(\frac12)^2]^r\mathfrak s^{(r)}_{1/2}$.

Finally, $\varphi(1)=\frac xy\varphi(0)$ and $\bar\varphi(1)=\frac{1-y}{1-x}\Gamma(1-y)/\Gamma(1-x)$ with
$\Gamma(1-y)/\Gamma(1-x)=(\sin\pi x/\sin\pi y)\varphi(0)$ by reflection, and
$\bar\varphi(\frac12)=\Gamma(\frac32-y)/\Gamma(\frac32-x)=\frac{1-2y}{1-2x}(\cos\pi x/\cos\pi y)\varphi(\frac12)$ by
$\Gamma(\frac12+z)\Gamma(\frac12-z)=\pi/\cos\pi z$. Hence
\begin{equation}
 \sin\pi x\sin\pi y\;\varphi(1)\bar\varphi(1)=\xi_0\sin^2(\pi x)\,\varphi(0)^2,\qquad
 \cos\pi x\cos\pi y\;\varphi(\tfrac12)\bar\varphi(\tfrac12)=\xi_h\cos^2(\pi x)\,\varphi(\tfrac12)^2,
 \label{gff:mgl:rhoid}
\end{equation}
and multiplying (ii) and (iii) by $N_b^r$ gives \eqref{gff:mgl:samp0} and \eqref{gff:mgl:samph}. $\square$

\paragraph{The three-node formula in Horn values.}
By \eqref{gff:gdr:topless} the three samples are, in units of $N_b^4$, $c_{4,0}+c_{4,1}+c_{4,2}$ at $\tilde t=0$,
$c_{4,0}-c_{4,1}+c_{4,2}$ at $\tilde t=\frac12$ and $c_{4,0}+\varsigma_bc_{4,1}+\cos(\pi b)\,c_{4,2}$ at $t=0$. The
determinant is $4\sin^2\pi\sigma=4s_b^2\neq0$, and
\[\begin{gathered}
 N_b^4c_{4,1}=\frac{\mathcal X_0-\mathcal X_{1/2}}2,\qquad
 N_b^4c_{4,0}=\frac{\mathcal X_0+\mathcal X_{1/2}}2-N_b^4c_{4,2},\\
 N_b^4c_{4,2}=\frac1{2s_b^2}\Bigl[\frac{\mathcal X_0+\mathcal X_{1/2}}2+\varsigma_bN_b^4c_{4,1}-\mathcal X_T\Bigr].
\end{gathered}\]
Inserting these into \eqref{gff:gdr:explicit} gives \eqref{gff:gdr:threenode}, and with \eqref{gff:mgl:sampT}--\eqref{gff:mgl:samph} at
$r=4$, for every $b$ in the strip,
\begin{equation}\boxed{\begin{aligned}
 W_4(d;b)={}&-\frac{4\pi^3\varsigma_b}{s_b^3}\,\mathfrak h_x
 +\frac{2(1+\varsigma_b)^2(1+2\varsigma_b)\kappa_0^4}{\pi s_b^7}\bigl[\mathfrak s_0+2\xi_0^4\mathfrak h_1\bigr]\\
 &-\frac{2(1-\varsigma_b)^2(1-2\varsigma_b)\kappa_h^4}{\pi s_b^7}\bigl[\mathfrak s_{1/2}+2\xi_h^4\mathfrak h_{1/2}\bigr].
\end{aligned}}\label{gff:mgl:W4five}\end{equation}
For $r=3$ and $0<\operatorname{Re}b<b_{\max}(3,d)$ the solve of the two samples at $t=0$ and $\tilde t=\frac12$ is
the proof of \eqref{gff:gdr:twonode} in Section~\ref{gff:mgz:nodes}; it gives
$W_3(d;b)=-\frac{2N_b^3}{\pi s_b^5}[(1+2\varsigma_b)\Xi_3(0;b)-2\varsigma_b(1-\varsigma_b)\Xi_3(\frac12-\frac b4;b)]$,
and with \eqref{gff:mgl:sampT} and \eqref{gff:mgl:samph} at $r=3$ this is
\begin{equation}
 W_3(d;b)=\frac{2\pi^2(1+2\varsigma_b)}{s_b^2}\,\mathfrak h^{(3)}_x
 +\frac{4\varsigma_b(1-\varsigma_b)\kappa_h^3}{\pi s_b^5}\bigl[\mathfrak s^{(3)}_{1/2}+2\xi_h^3\mathfrak h^{(3)}_{1/2}\bigr],
 \label{gff:mgl:W3three}
\end{equation}
which is \eqref{gff:gdr:twonode} with its midpoint sample written in Horn values. Neither formula holds without the
top-mode theorem: a mode $c_{4,3}$ would enter $W_4$ with the nonzero weight $B_{4,3}=-20\pi^3/s_b^3$ of
Section~\ref{gff:gdr:residue} and would not be fixed by three samples. The agreement of \eqref{gff:mgl:W4five} with independent
values at $b\neq d/2$ (Section~\ref{gff:mgl:checks}) is therefore also a numerical test of \eqref{gff:gdr:topless}. The exact
solve and the insertion into \eqref{gff:gdr:explicit} are checked symbolically by
\path{LRI_C5_general_d/L/sym_functional.py}, which also rederives \eqref{gff:gdr:explicit} from \eqref{gff:fp:finiteFunctional}
with $q$ replaced by $\sigma$.

\paragraph{The four-wheel period.}
At $b=d/2$ one has $x=\frac q2$, $y=\frac{3q}2$, $\sigma=g=q$, $N_{d/2}=\Gamma(2q)$, $s_b=S$, $\varsigma_b=C$, so
$\kappa_0=\Gamma(2q)\sin^2(\frac{\pi q}2)\,[\Gamma(\frac q2)/\Gamma(\frac{3q}2)]^2$,
$\kappa_h=\Gamma(2q)\cos^2(\frac{\pi q}2)\,[\Gamma(\frac{1+q}2)/\Gamma(\frac{1+3q}2)]^2$,
$\xi_0=\frac{2-3q}{3(2-q)}$ and $\xi_h=\frac{1-3q}{1-q}$. Put, for $b$ near $d/2$,
\begin{equation}\begin{gathered}
 \mathsf k_0(b)=\frac{2(1+\varsigma_b)^2\kappa_0^4}{\pi s_b^7},\qquad
 \mathsf k_h(b)=\frac{2(1-\varsigma_b)^2\kappa_h^4}{\pi s_b^7},\\
 \mathfrak S_0=\mathfrak s_0+2\xi_0^4\mathfrak h_1=\frac{\mathcal X_0}{\kappa_0^4},\qquad
 \mathfrak S_h=\mathfrak s_{1/2}+2\xi_h^4\mathfrak h_{1/2}=\frac{\mathcal X_{1/2}}{\kappa_h^4},
\end{gathered}\label{gff:mgl:kS}\end{equation}
so that \eqref{gff:mgl:W4five} reads
$W_4=-\frac{4\pi^3\varsigma_b}{s_b^3}\mathfrak h_x+(1+2\varsigma_b)\mathsf k_0(b)\mathfrak S_0
+(2\varsigma_b-1)\mathsf k_h(b)\mathfrak S_h$. Write $\mathsf k_0,\mathsf k_h$ for the values at $b=d/2$; with
$1+C=2\cos^2\frac{\pi q}2$, $1-C=2\sin^2\frac{\pi q}2$ and $S=2\sin\frac{\pi q}2\cos\frac{\pi q}2$,
\[
 \mathsf k_0=\frac{\Gamma(2q)^4}{16\pi}\Bigl[\frac{\Gamma(\frac q2)}{\Gamma(\frac{3q}2)}\Bigr]^8
 \frac{\tan\frac{\pi q}2}{\cos^2\frac{\pi q}2},\qquad
 \mathsf k_h=\frac{\Gamma(2q)^4}{16\pi}\Bigl[\frac{\Gamma(\frac{1+q}2)}{\Gamma(\frac{1+3q}2)}\Bigr]^8
 \frac{\cot\frac{\pi q}2}{\sin^2\frac{\pi q}2}.
\]
Then, for $\frac12<d<4$,
\begin{equation}\boxed{
 P_4(d)=-\frac{4\pi^3C}{S^3}\,\mathfrak h_x+(1+2C)\,\mathsf k_0\,\mathfrak S_0+(2C-1)\,\mathsf k_h\,\mathfrak S_h,}
 \label{gff:mgl:P4five}
\end{equation}
which is \eqref{gff:gdr:P4} with the samples written in Horn values. The symbols $\mathsf k_0,\mathsf k_h,\mathfrak S_0,
\mathfrak S_h$ and $\mathsf v_0,\mathsf v_h$ below are local to this subsection.

\paragraph{Derivative series.}
A prime denotes the total derivative $d/db$ at $b=d/2$; in $\mathfrak h_x$ the argument $u=x$ moves with $b$, while
$u=\frac12,1$ in $\mathfrak h_u$ and $u=0,\frac12$ in $\mathfrak s_u$ are fixed. Let $H(a;m)=\psi_0(a+m)-\psi_0(a)$ as
in \eqref{gff:gc:dictionary}, and write $\mathfrak h_u[\ell,n]$ and $\mathfrak s_u[j,k]$ for the summands of
\eqref{gff:mgl:fiveCenter}. For $\frac12<d<4$ the series
\begin{equation}\begin{aligned}
 \mathfrak h_x'&=2\sum_{\ell,n\geq0}d_\ell T_{\ell n}^4\bigl[H(q;\ell+n)+H(1-q;n)\bigr]
 =4\sum_{\ell,n\geq0}d_\ell T_{\ell n}^4\mu_1,\\
 \mathfrak h_u'&=\sum_{\ell,n\geq0}\mathfrak h_u[\ell,n]\bigl[H(u+\tfrac q2;\ell+n)+H(u+\tfrac{3q}2;\ell+n)\\
 &\hphantom{{}=\sum_{\ell,n\geq0}\mathfrak h_u[\ell,n]\bigl[}+H(u+1-\tfrac{3q}2;n)+H(u+1-\tfrac q2;n)\bigr],\\
 \mathfrak s_u'&=2\sum_{j,k\geq0}\mathfrak s_u[j,k]\bigl[H(u+\tfrac q2;j)+H(u+\tfrac{3q}2;j)\bigr],
\end{aligned}\label{gff:mgl:dseries}\end{equation}
with $u=\frac12,1$ in the second line, $u=0,\frac12$ in the third and $\mu_1$ of \eqref{gff:gs:etaTheta}, converge absolutely and are the $b$-derivatives at $b=d/2$. For general $r$ the
prefactors $2,1,2$ become $\frac r2,\frac r4,\frac r2$.

\emph{Proof.} $x'=\frac14$, $y'=-\frac14$ and $\partial_a\log(a)_m=H(a;m)$. In $\mathfrak h_u$ with $u$ fixed, the
parameters $u+x$ and $u+1-y$ stand in numerators and move with slope $\frac14$, while $u+y$ and $u+1-x$ stand in
denominators and move with slope $-\frac14$, so all four contributions carry the sign $+$, each with weight
$\frac r4$. In $\mathfrak h_x$ the numerator parameters $u+x=\sigma$ and $u+1-y=1-g$ move with slope $\frac12$, and
$u+y=2q$, $u+1-x=1$ are fixed; this gives weight $\frac r2$, and at $b=d/2$ the bracket is $2\mu_1$. In
$\mathfrak s_u$ the contributions are $\frac r4[H(u+x;j)+H(u+y;j)+H(u+x;k)+H(u+y;k)]$, and the summand is symmetric in
$j\leftrightarrow k$, so the $k$-terms sum to the $j$-terms. For convergence, let $B$ be a closed disc around $d/2$ inside
the strip $0<\operatorname{Re}b<b_{\max}(4,d)$ on which $\sigma\notin\mathbb Z$; it exists because $d/2$ lies in the
strip for $\frac12<d<4$. On $B$ the summands of the five functions are analytic in $b$ (no denominator Pochhammer symbol
vanishes, since $\operatorname{Re}y>0$ and $\operatorname{Re}(1-x)>0$ there), and by the Gamma-ratio bound, whose
constants are uniform for parameters in compact sets, they are bounded by
$C(1+\ell)^{d-2}(1+\ell+n)^{-\alpha_B}(1+n)^{-\alpha_B}$ for $\mathfrak h_u$ and
$C(1+j+k)^{d-2}(1+j)^{-\alpha_B}(1+k)^{-\alpha_B}$ for $\mathfrak s_u$, with
$\alpha_B=\min_B4\operatorname{Re}g>\max(\frac d2,d-1,\frac12)$. By the lattice-sum criterion of
Section~\ref{gff:gdr:periodization} the five series converge uniformly on $B$. By Weierstrass' theorem their sums are
analytic inside $B$ and may be differentiated termwise, and Cauchy's estimate on a circle around $d/2$ inside $B$
bounds the differentiated terms by a constant times the same majorant, so the derivative series converge absolutely.
The weights $H(a;m)$ grow like $\log(1+m)$, which is the logarithmic loss mentioned in Section~\ref{gff:gdr:samples}.
$\square$

\paragraph{The rim jet in Horn values.}
For $\frac12<d<4$,
\begin{equation}\boxed{\begin{aligned}
 \LL_d={}&\frac{2\pi^4(1+2C^2)}{S^4}\,\mathfrak h_x-\frac{4\pi^3C}{S^3}\,\mathfrak h_x'\\
 &+\mathsf k_0\Bigl\{\bigl[(1+2C)\mathsf v_0-\pi S\bigr]\mathfrak S_0
 +(1+2C)\bigl[\mathfrak s_0'+2\xi_0^4\mathfrak h_1'+2(\xi_0^4)'\mathfrak h_1\bigr]\Bigr\}\\
 &+\mathsf k_h\Bigl\{\bigl[(2C-1)\mathsf v_h-\pi S\bigr]\mathfrak S_h
 +(2C-1)\bigl[\mathfrak s_{1/2}'+2\xi_h^4\mathfrak h_{1/2}'+2(\xi_h^4)'\mathfrak h_{1/2}\bigr]\Bigr\},
\end{aligned}}\label{gff:mgl:Lfinal}\end{equation}
with the five derivatives \eqref{gff:mgl:dseries},
\[\begin{aligned}
 \mathsf v_0&=\partial_b\log\mathsf k_0(b)\big|_{b=d/2}
 =-4\psi_0(q)+2\psi_0(\tfrac q2)+2\psi_0(\tfrac{3q}2)+\pi\cot\tfrac{\pi q}2-\tfrac{3\pi}2\cot\pi q,\\
 \mathsf v_h&=\partial_b\log\mathsf k_h(b)\big|_{b=d/2}
 =-4\psi_0(q)+2\psi_0(\tfrac{1+q}2)+2\psi_0(\tfrac{1+3q}2)-\pi\tan\tfrac{\pi q}2-\tfrac{3\pi}2\cot\pi q,\\
 (\xi_0^4)'&=2\xi_0^4\Bigl[\frac4{3q}+\frac1{2-q}\Bigr]+\frac{2\xi_0^3}{3(2-q)},\qquad
 (\xi_h^4)'=\frac{4(1-2q)\,\xi_h^3}{(1-q)^2}.
\end{aligned}\]
Every coefficient is finite for $\frac12<d<4$. Written with the samples, \eqref{gff:mgl:Lfinal} is \eqref{gff:gdr:L}: both are
the $b$-derivative at $d/2$ of \eqref{gff:gdr:threenode}$\,=\,$\eqref{gff:mgl:W4five}, and
$\mathcal X_0'=(\kappa_0^4)'\mathfrak S_0+\kappa_0^4\mathfrak S_0'$, $\mathcal X_{1/2}'=(\kappa_h^4)'\mathfrak S_h
+\kappa_h^4\mathfrak S_h'$, $\mathcal X_T'=(\pi^4s_b^4)'\mathfrak h_x+\pi^4S^4\mathfrak h_x'$.

\emph{Proof.} Near $b=d/2$ one has $0<\sigma<1$, so \eqref{gff:mgl:W4five} holds on a neighborhood of $d/2$. Every factor
is analytic there, the five Horn functions by the proof of \eqref{gff:mgl:dseries}, so the product rule applies to
$W_4=-\frac{4\pi^3\varsigma_b}{s_b^3}\mathfrak h_x+(1+2\varsigma_b)\mathsf k_0(b)\mathfrak S_0
+(2\varsigma_b-1)\mathsf k_h(b)\mathfrak S_h$. With $\partial_b\varsigma_b=-\frac\pi2s_b$ and
$\partial_bs_b=\frac\pi2\varsigma_b$ one finds $\partial_b(-4\pi^3\varsigma_b/s_b^3)=2\pi^4(1+2\varsigma_b^2)/s_b^4$ and
$\partial_b(1+2\varsigma_b)=\partial_b(2\varsigma_b-1)=-\pi s_b$, which give the first line of \eqref{gff:mgl:Lfinal} and the
terms $-\pi S\,\mathsf k_0\mathfrak S_0$, $-\pi S\,\mathsf k_h\mathfrak S_h$. Next, with
$\partial_b\log N_b=-\frac12[\psi_0(g)+\psi_0(\sigma)]$, $\partial_b\log\varphi(0)^2=\frac12[\psi_0(x)+\psi_0(y)]$,
$\partial_b\log\varphi(\frac12)^2=\frac12[\psi_0(\frac12+x)+\psi_0(\frac12+y)]$,
$\partial_b\log\sin^2\pi x=\frac\pi2\cot\pi x$ and $\partial_b\log\cos^2\pi x=-\frac\pi2\tan\pi x$,
\[\begin{aligned}
 \partial_b\log\mathsf k_0(b)&=-\frac{\pi s_b}{1+\varsigma_b}-\frac{7\pi\varsigma_b}{2s_b}
 -2\psi_0(g)-2\psi_0(\sigma)+2\psi_0(x)+2\psi_0(y)+2\pi\cot\pi x,\\
 \partial_b\log\mathsf k_h(b)&=\frac{\pi s_b}{1-\varsigma_b}-\frac{7\pi\varsigma_b}{2s_b}
 -2\psi_0(g)-2\psi_0(\sigma)+2\psi_0(\tfrac12+x)+2\psi_0(\tfrac12+y)-2\pi\tan\pi x .
\end{aligned}\]
At $b=d/2$, $\pi S/(1+C)=\pi\tan\frac{\pi q}2$, $\pi S/(1-C)=\pi\cot\frac{\pi q}2$ and
$7\pi C/(2S)=\frac{7\pi}2\cot\pi q$; with $\pi\cot\frac{\pi q}2-\pi\tan\frac{\pi q}2=2\pi\cot\pi q$ this gives the
printed $\mathsf v_0$ and $\mathsf v_h$, and the terms $(1+2C)\mathsf v_0\mathsf k_0\mathfrak S_0$ and
$(2C-1)\mathsf v_h\mathsf k_h\mathfrak S_h$. The remaining terms are $(1+2C)\mathsf k_0\mathfrak S_0'$ and
$(2C-1)\mathsf k_h\mathfrak S_h'$, with $\mathfrak S_0'=\mathfrak s_0'+2\xi_0^4\mathfrak h_1'+2(\xi_0^4)'\mathfrak h_1$
and likewise for $\mathfrak S_h'$. Finally $(\xi_0^4)'=4\xi_0^4\partial_b\log\xi_0$ with
$\partial_b\log\xi_0=\frac14[\frac1x+\frac1y+\frac1{1-x}+\frac1{1-y}]=\frac12[\frac4{3q}+\frac1{2-3q}+\frac1{2-q}]$
at $b=d/2$, and $(\xi_h^4)'=4\xi_h^4\partial_b\log\xi_h$ with
$\partial_b\log\xi_h=\frac12[\frac1{1-2x}+\frac1{1-2y}]=\frac12[\frac1{1-q}+\frac1{1-3q}]$; the poles at
$q=\frac23$ and $q=\frac13$ cancel against $\xi_0\propto2-3q$ and $\xi_h\propto1-3q$, which gives the printed forms.
The coefficients are finite for $0<q<1$: $S\neq0$, the trigonometric functions of $\frac{\pi q}2$ and $\cot\pi q$ are
finite, the digamma arguments are positive, and $2-q,1-q\neq0$. $\square$

The closed coefficients were compared with direct symbolic differentiation of \eqref{gff:mgl:W4five} at eleven values of
$d$, including $\frac43$, $2$ and $\frac83$, to $10^{-40}$ (\path{LRI_C5_general_d/L/check_closed_coeffs.py}), and the
review of \path{LRI_C5_general_d/review/L-rigor/} repeated this with its own coefficient functions (Section~\ref{gff:gdr:checks}).

\paragraph{Moving and fixed nodes.}
The nodes $\tilde t=0,\frac12$ are $t=-\frac b4$ and $t=\frac12-\frac b4$, which move with $b$. Since
$\Xi_4(t;b)=\sum_kc_{4,k}\cos2\pi k\tilde t$ is even about both points, $\partial_t\Xi_4=0$ there, and the total
$b$-derivative of these samples equals their $b$-derivative at fixed $t$. In the fixed-$t$ form the right half-lattice
is differentiated only in its numerator parameters ($\sigma+t$ and $1-g+t$ move with slope $\frac12$, while $2q+t$
and $1+t$ are fixed) and the left half-lattice only in its denominator parameters; their average is the symmetric
derivative of \eqref{gff:mgl:dseries}, in which every parameter moves with slope $\pm\frac14$. The node $t=0$ is fixed by
construction. No $t$-derivative of any sample enters \eqref{gff:mgl:Lfinal}.

\paragraph{Relation to the residue series.}
The series \eqref{gff:gs:Lseries} is a residue at $t=0$: it needs the Laurent data of the right half-lattice at $t=0$ up
to third order, through the polynomials $F_4$ and $F_{4,1}$ of \eqref{gff:gs:Fpolys} (with $\lambda_3$ and $\eta_3$).
Formula \eqref{gff:mgl:Lfinal} keeps only the zeroth-order data at $t=0$, namely $\mathfrak h_x=\sum d_\ell T_{\ell n}^4$ and
$\mathfrak h_x'=4\sum d_\ell T_{\ell n}^4\mu_1$, and replaces the higher Laurent coefficients by the values at the two
reflection centers. Section~\ref{gff:gdr:taylor} is the intermediate form with the pole node alone, one Laurent order
below \eqref{gff:gs:Lseries}.

\paragraph{Special dimensions.}
At $b=d/2$ one has $y=3x$, and the triple-angle formulas give $\sin\pi y=(1+2C)\sin\pi x$ and
$\cos\pi y=(2C-1)\cos\pi x$, so $P_\pm(0)=(1+2C)^rP_{\rm m}(0)$ and $P_\pm(\frac12)=(2C-1)^rP_{\rm m}(\frac12)$. In
\eqref{gff:mgl:P4five} and \eqref{gff:mgl:Lfinal} the sample at $\tilde t=0$ carries the factor $1+2C$, and $\mathfrak h_1$
carries in addition $\xi_0^4$ (or $\xi_0^3$ in $(\xi_0^4)'$), with $\xi_0\propto2-3q$; the sample at
$\tilde t=\frac12$ carries $2C-1$, and $\mathfrak h_{1/2}$ carries in addition $\xi_h^4$ or $\xi_h^3$, with
$\xi_h\propto1-3q$. Hence at $d=\frac83$ ($q=\frac23$, $C=-\frac12$, $y=1$) $\mathfrak h_1$ drops out of both $P_4$ and
$\LL_d$: its coefficient in $W_4(d;b)$ has a zero of order five at $b=d/2$, a simple zero of $1+2\varsigma_b$ times the
fourth power of the simple zero of $\xi_0$. There $\mathfrak s_0$ enters $\LL_d$ only through its value, with
coefficient $-\pi S\mathsf k_0$. At $d=\frac43$ ($q=\frac13$, $C=\frac12$, $y=\frac12$) the same holds for
$\mathfrak h_{1/2}$ and $\mathfrak s_{1/2}$. At $d=2$ ($C=0$) the coefficient of $\mathfrak h_x'$ vanishes, and
$\mathfrak h_x$ drops out of $P_4(2)$. Thus $P_4(\frac43)$ and $P_4(\frac83)$ need three Horn values, $\LL_{4/3}$ and
$\LL_{8/3}$ four values and three derivatives, and $\LL_2$ five values and four derivatives. In the code the vanishing
coefficients come out zero to working precision.

\paragraph{The five functions at $d=1,2,3$.}
At $d=1$, $d_\ell=\delta_{\ell0}+\delta_{\ell1}$, so each $\mathfrak h_u$ is a single sum, a ${}_9F_8(1)$ (an
${}_8F_7(1)$ for $u=x$) plus the analogous single sum of the row $\ell=1$; $\mathfrak s_0=1+2(x/y)^4$, which is $\frac{83}{81}$ at
$b=\frac12$ with $\mathfrak s_0'=\frac{64}{243}$, and $\mathfrak s_{1/2}=1$. The samples at $\tilde t=0$ and
$\tilde t=\frac12$ coincide, since $c_{4,1}\equiv0$ there (Section~\ref{gff:gdr:special}).
At $d=2$, $x+y=1$, so $\bar\varphi=\varphi$, and $d_\ell=2-\delta_{\ell0}$. Put $\chi_m=\varphi(m+u)^r$ and
$\mathsf m(u)=\sum_{n\geq0}[(u+x)_n/(u+y)_n]^r$, so that $\sum_{m\geq0}\chi_m=\varphi(u)^r\mathsf m(u)$. In
$\mathcal R_+(u)$ set $m=\ell+n$; the weight $d_{m-n}$ is $2$ for $m>n$ and $1$ for $m=n$, so
$\mathcal R_+(u)=\sum_{m\geq n\geq0}d_{m-n}\chi_m\chi_n=(\sum_m\chi_m)^2$, the symmetrization of the proof of
\eqref{gff:gdr:d12lattice}. Likewise $d_{j+k+1}=2$ and $d_{j+k}=2-\delta_{j0}\delta_{k0}$ give
$\mathcal R_{\rm m}(\frac12)=2\varphi(\frac12)^{2r}\mathsf m(\frac12)^2$ and
$\mathcal R_{\rm s}(0)=2[\varphi(0)^r\mathsf m(0)]^2-\varphi(0)^{2r}$, and $\mathsf m(0)=1+(x/y)^r\mathsf m(1)$. Hence
\[
 \mathfrak h^{(r)}_u=\mathsf m(u)^2,\qquad \mathfrak s^{(r)}_{1/2}=2\,\mathfrak h^{(r)}_{1/2},\qquad
 \mathfrak s^{(r)}_0=2\bigl[1+(x/y)^r\mathsf m(1)\bigr]^2-1\qquad(d=2),
\]
with $\mathsf m(x)={}_rF_{r-1}(\sigma,\dots,\sigma;1,\dots,1;1)$; for real $b$, $\mathsf m(u)>0$ is the positive square
root of $\mathfrak h^{(r)}_u$. The $d=2$ rows of the test data in Section~\ref{gff:mgl:checks} satisfy the last two
relations to the printed precision.
At $d=3$, $d_\ell=2\ell+1$ is linear, so $\mathfrak s_0$ and $\mathfrak s_{1/2}$ are products of one-variable moments:
with $\mathsf n_i(u)=\sum_{j\geq0}j^i[(u+x)_j/(u+y)_j]^4$,
$\mathfrak s_0=4\mathsf n_0(0)\mathsf n_1(0)+\mathsf n_0(0)^2$ and
$\mathfrak s_{1/2}=4\mathsf n_0(\frac12)\mathsf n_1(\frac12)+3\mathsf n_0(\frac12)^2$. This does not hold for
$\mathfrak h_x,\mathfrak h_{1/2},\mathfrak h_1$: with $m=\ell+n$ they are nested sums over $m\geq n\geq0$ with weight
$2(m-n)+1$, genuine two-variable values (rank two). The reduction of $\LL_3$ in Section~\ref{gff:gdr:special} works instead
with the half-shift structure of the full samples. Its Dougall closed form of the amplitude $\mathfrak b_4$ rests on the
linear multiplicity $2\ell+1$, which turns the half-shifted lattice sum into a very-well-poised bilateral series. For
non-integer $d$ the multiplicity \eqref{gff:mgl:mult} is not polynomial; no Gamma-product amplitude was found there, and
none was searched for systematically. The closed forms of $\LL_1$, $\LL_2$ and the $d=3$ reduction are obtained in
Section~\ref{gff:gdr:special} from the mode form, not from \eqref{gff:mgl:Lfinal}; \eqref{gff:mgl:Lfinal} reproduces all three
values numerically (Section~\ref{gff:mgl:checks}).

\paragraph{Function count.}
At generic $\frac12<d<4$, \eqref{gff:mgl:P4five} needs the five two-variable Horn values
$\mathfrak h_x$, $\mathfrak h_{1/2}$, $\mathfrak h_1$, $\mathfrak s_0$, $\mathfrak s_{1/2}$ at argument $(1,1)$, and
\eqref{gff:mgl:Lfinal} the same five values and their five first parameter derivatives \eqref{gff:mgl:dseries}. The coefficients
need trigonometric functions of $\pi q$ and $\frac{\pi q}2$; $\Gamma(2q)$, $\Gamma(\frac q2)$, $\Gamma(\frac{3q}2)$,
$\Gamma(\frac{1+q}2)$, $\Gamma(\frac{1+3q}2)$; $\psi_0$ at $q,\frac q2,\frac{3q}2,\frac{1+q}2,\frac{1+3q}2$ (by the
duplication formula, equivalently $\psi_0(q),\psi_0(3q),\psi_0(\frac q2),\psi_0(\frac{3q}2)$ and $\log2$); and rational
functions of $q$. In the unweighted Horn class \eqref{gff:fp:HornDefinition} the derivatives cost $19$ three-variable
values: each weight $H(a;m)=m\sum_{k\geq0}[(a+k)(a+m+k)]^{-1}$ of \eqref{gff:gc:insertionIdentities} adds one summation
index, encoded by \eqref{gff:gc:reciprocalEncoding}, and the factor $m$ costs one array for $m=n$ or $m=j$ and two for
$m=\ell+n$ \eqref{gff:gc:falling}. This gives $3$ arrays for $\mathfrak h_x'$, $6$ each for $\mathfrak h_{1/2}'$ and
$\mathfrak h_1'$, and $2$ each for $\mathfrak s_0'$ and $\mathfrak s_{1/2}'$. So $\LL_d$ needs $5$ two-variable and $19$
three-variable unweighted Horn values, in place of the $244$ values with up to six variables of \eqref{gff:gc:finiteClosure}.
In the split form used by the code (next paragraph) the five functions are four one-variable values, an
${}_8F_7(1)$ in $\mathfrak h_x$, two ${}_9F_8(1)$ in $\mathfrak h_{1/2},\mathfrak h_1$ and a ${}_7F_6(1)$ in
$\mathfrak s_0$, and five two-variable values, each with its first derivative. For $P_4(d)$, the sampling
\eqref{gff:fp:wheelSample} at the nodes \eqref{gff:fp:nodes} uses $20$ complex block values ($A_4^0,A_4^+$ at eight complex
points and $C_4$ at four; $15$ with \eqref{gff:gdr:topless}); \eqref{gff:mgl:P4five} uses five real Horn values. The five
derivative values are the new constants of the generic-dimensional problem. Whether they are independent is not
decided here; at $d=1,2$ they reduce, and at $d=3$ they reduce, with the identities of Section~\ref{gff:gdr:halfShift}, to
the two constants $\mathfrak d_H,\mathfrak d_a$ (Section~\ref{gff:gdr:special}). Contiguous relations do not connect the five functions, since their parameters differ by
non-integer shifts ($x$ and $\frac12$; Section~\ref{gff:gdr:attempts}).

\paragraph{Evaluation.}
\path{LRI_C5_general_d/L/L_general.py} evaluates \eqref{gff:mgl:W4five}, and \eqref{gff:mgl:W3three}, for any complex $b$ in the
strip, in a split form. Let $\hat d_m=d_{m+1}/(4q)=(2q+1)_m(4q-1)_m/((2q)_m(2)_m)$ be the ratio of
\eqref{gff:fp:degeneracy}. The row $\ell=0$ of $\mathfrak h_u$ is a one-variable sum, and the rows $\ell=j+1\geq1$ are
$4q[(u+x)/(u+y)]^r\sum_{j,n\geq0}\hat d_j\mathsf x_{j+n}\mathsf z_n$ with
$\mathsf x_m=[(u+1+x)_m/(u+1+y)_m]^r$ and $\mathsf z_n=[(u+1-y)_n/(u+1-x)_n]^r$, sequences of the kind
$\mathsf x,\mathsf z$ of \eqref{gff:mgf:threeFamilies}. Further $\mathfrak s_{1/2}=4q\sum_{j,k}\hat d_{j+k}a_ja_k$ with
$a_j=[(\frac12+x)_j/(\frac12+y)_j]^r$, and $\mathfrak s_0=1+4q(x/y)^r[\sum_k\hat d_k\tilde a_k
+\sum_{j,k}\hat d_{j+k}\tilde a_ja_k]$ with $a_j=[(x)_j/(y)_j]^r$ and $\tilde a_k=[(1+x)_k/(1+y)_k]^r$. Each sum is
evaluated as an exact finite part plus an asymptotic tail by the method of the paragraph \emph{Evaluation of the Horn
values} of Section~\ref{gff:mgf:nodes}, with the routines of \path{LRI_C5_general_d/foundation/hypsums.py} and the
default truncation parameters and guard digits stated there; at $d=1,2,3$ the multiplicities are used as exact
polynomials. The
$b$-derivatives are exact: \path{LRI_C5_general_d/L/jet.py} implements first-order jets (a value and its derivative,
combined by the product rule), and \path{LRI_C5_general_d/L/hypsums_jet.py} is the same algorithm with every
operation jet-valued, including the Hurwitz zeta derivative in its first argument and the Bernoulli-polynomial tail
coefficients. Evaluating at the jet based at $b=\frac d2$ returns the value and the derivative of the evaluated
expression, that is the series \eqref{gff:mgl:dseries} with their tails (now with $\log m$ terms) accelerated in the same
way. No finite difference and no Cauchy circle in $b$ enters $\LL_d$. One pair $(P_4(d),\LL_d)$ at 50 digits takes
about 15\,s on one AMD core of the symmetry cluster (5\,s at 30 digits on a laptop); the earlier route of
\path{LRI_C5_general_d/foundation/}, a 96-point Cauchy circle in $b$ of the finite-sample $W_4$, took about 45\,s per
point at 50 digits. The function \texttt{L\_explicit} evaluates \eqref{gff:mgl:Lfinal} as printed, with the closed-form
coefficients of \texttt{coeffs\_closed} times the five values and the five derivative series; it agrees with the jet
derivative of the whole of \eqref{gff:mgl:W4five} to working precision, and \path{LRI_C5_general_d/L/sym_L.py} does the same
with coefficients obtained by symbolic differentiation. Usage: \texttt{python3 L\_general.py d [dps] [b]}, with $d$
given exactly (for example \texttt{4/3}) and $b=d/2$ by default; the script prints $W_4$, $W_3$ and their
$b$-derivatives and, at $b=d/2$, the value of \eqref{gff:mgl:Lfinal}. Resonant parameters, $r(d-b)/2\in\mathbb Z$ at
non-integer $d$, are refused; at $b=d/2$ they occur for $r=3$ at $d=\frac43,\frac83$ and never for $r=4$.

\subsubsection{Checks of the five-Horn formula}\label{gff:mgl:checks}
All runs were made on the symmetry cluster. The records are in \path{LRI_C5_general_d/L/results/} (\path{summary.json},
\path{val50.json}, \path{val70.json}, \path{cfd40.json} and \path{ref_d*.json}), produced by \path{run_validation.py},
\path{refcheck_r4.py}, \path{check_func_derivs.py}, \path{analyze.py} and \path{make_tables.py} in
\path{LRI_C5_general_d/L/}. References: at $d=1$ the proved closed forms $P_4=13B_--3B_+$ (Section~\ref{gff:cn:oneDresults})
and \eqref{gff:eq:KLpsiReflected}; at $d=2$ \eqref{gff:fp:P4two} and $\LL_2=-2\log2\,P_4(2)$ \eqref{gff:gc:L2closed}; at $d=3$ the
160-digit values of \path{LRI_C5_d3}; at $d=0.8,1.25,1.5,2.5,2.75$ the 41- to 50-digit values of the independent spectral
code in \path{LRI_C5_general_d/references/reference_values.json}; at $d=0.6,\frac43,\frac83,3.5$ the same spectral code,
run at 32 digits for $r=4$ only (\path{LRI_C5_general_d/L/refcheck_r4.py}; at $d=0.6$ the point $b=d/2$ lies outside the
$r=3$ strip and the $r=3$ spectral sum does not converge). The spectral code shares no code with the five-Horn formula.

\paragraph{Values of $\LL_d$.}
$\LL_d$ from \eqref{gff:mgl:Lfinal} with exact jet derivatives, at the working precision of the run (50 digits; 70 digits
at $d=1,1.5,2,2.5,3$). ``dev.'' is the relative deviation from the reference, ``conv.'' the relative change of $\LL_d$
when all truncation parameters are enlarged by the factor $1.4$.
\begin{center}\footnotesize
\begin{tabular}{@{}l l l l l l@{}}\toprule
$d$ & $\LL_d$ (five-Horn formula) & reference & dev.\ $P_4$ & dev.\ $\LL_d$ & conv.\ $\LL_d$\\\midrule
0.6 & $-45086.56690064459550422$ & spectral, own run (32 d.) & $3.6\cdot10^{-32}$ & $7.3\cdot10^{-33}$ & $1.9\cdot10^{-44}$\\
0.8 & $-14228.69045883150419322995007$ & spectral reference (47 d.) & $1.0\cdot10^{-45}$ & $2.7\cdot10^{-45}$ & $3.2\cdot10^{-45}$\\
1 & $-5805.301629548802380141$ & closed form (proved) & $1.0\cdot10^{-80}$ & $3.4\cdot10^{-81}$ & $<10^{-80}$\\
1.25 & $-2361.293891754120963656812701$ & spectral reference (49 d.) & $3.3\cdot10^{-45}$ & $8.7\cdot10^{-46}$ & $1.4\cdot10^{-47}$\\
$4/3$ & $-1818.786953676526007105$ & spectral, own run (32 d.) & $3.7\cdot10^{-24}$ & $5.5\cdot10^{-24}$ & $2.8\cdot10^{-49}$\\
1.5 & $-1127.932175171310391303895184$ & spectral reference (50 d.) & $9.7\cdot10^{-47}$ & $3.6\cdot10^{-45}$ & $1.2\cdot10^{-58}$\\
2 & $-347.2491531232011310585$ & $-2\log2\,P_4(2)$ (proved) & $4.7\cdot10^{-81}$ & $8.9\cdot10^{-81}$ & $<10^{-80}$\\
2.5 & $-136.796077287838337518533388$ & spectral reference (49 d.) & $3.8\cdot10^{-45}$ & $2.5\cdot10^{-45}$ & $2.4\cdot10^{-58}$\\
$8/3$ & $-104.035830813482931953$ & spectral, own run (32 d.) & $2.6\cdot10^{-26}$ & $3.5\cdot10^{-26}$ & $2.2\cdot10^{-47}$\\
2.75 & $-91.22798596406115896228950228$ & spectral reference (50 d.) & $5.3\cdot10^{-46}$ & $1.9\cdot10^{-46}$ & $2.5\cdot10^{-47}$\\
3 & $-62.7105849683876693473$ & \path{LRI_C5_d3} (160 d.) & $6.0\cdot10^{-81}$ & $1.7\cdot10^{-81}$ & $<10^{-80}$\\
3.5 & $-31.81370266461412493531$ & spectral, own run (32 d.) & $2.8\cdot10^{-27}$ & $3.5\cdot10^{-27}$ & $9.5\cdot10^{-47}$\\
\bottomrule\end{tabular}
\end{center}

\paragraph{Away from the center.}
Relative deviation of \eqref{gff:mgl:W4five}, \eqref{gff:mgl:W3three} and the two-node formula \eqref{gff:gdr:twonode} from the
56-digit values of the independent spectral code, from the exact $W_3(2;b)$ of \eqref{gff:gs:W3d2}, and, at $d=2$, of the
ratio $W_4(2;1.05)/W_4(2;0.95)$ from the rim--spoke reflection
$W_4(2;b)/W_4(2;2-b)=2^{4-4b}\gamma_C(b/2)^2/\gamma_C(b-\frac12)$, which follows from the formula for $W_4(2;2x)$ in
Section~\ref{gff:gc:L2second}. The columns \eqref{gff:mgl:W4five} and \eqref{gff:mgl:W3three} come from
\path{LRI_C5_general_d/L/L_general.py}, at 70 digits for $d=1.5,2,2.5$ and at 50 digits otherwise. The column
\eqref{gff:gdr:twonode} comes from \path{LRI_C5_general_d/Z/offdiag_check.py} (record
\path{LRI_C5_general_d/Z/results/offdiag60.json}), at 50 digits against the spectral values and at 60 digits against
the finite-sample formula \eqref{gff:mgf:Wsample} with five fixed nodes (``fixed nodes''), which shares the Horn summation
code and is therefore not independent; that code also gives the value column. The two $W_3$ columns differ only at
$d=1.5$ and $2.5$, where the working precisions differ.
\begin{center}\scriptsize
\begin{tabular}{@{}l l l l l l l@{}}\toprule
$d$ & $b$ & $\operatorname{Re}W_3(d;b)$ (24 digits) & reference & dev.\ \eqref{gff:mgl:W4five}
 & dev.\ \eqref{gff:mgl:W3three} & dev.\ \eqref{gff:gdr:twonode}\\\midrule
$0.8$ & $0.37$ & \texttt{266.761267027041907531165} & spectral & $3.5\cdot10^{-51}$ & $2.9\cdot10^{-49}$ & $2.9\cdot10^{-49}$\\
$0.8$ & $0.43$ & \texttt{268.642244930670511767657} & spectral & $3.5\cdot10^{-50}$ & $1.5\cdot10^{-48}$ & $1.5\cdot10^{-48}$\\
$1.25$ & $0.575$ & \texttt{107.42914710548489909118} & spectral & $5.8\cdot10^{-53}$ & $2.7\cdot10^{-50}$ & $2.7\cdot10^{-50}$\\
$1.25$ & $0.675$ & \texttt{108.364244882288165222595} & spectral & $6.8\cdot10^{-53}$ & $1.1\cdot10^{-49}$ & $1.1\cdot10^{-49}$\\
$1.5$ & $0.7$ & \texttt{72.7303680788392574835365} & spectral & $2.0\cdot10^{-54}$ & $1.5\cdot10^{-52}$ & $1.5\cdot10^{-50}$\\
$1.5$ & $0.8$ & \texttt{73.0949788822424244641612} & spectral & $2.3\cdot10^{-53}$ & $8.1\cdot10^{-52}$ & $1.2\cdot10^{-50}$\\
$2.5$ & $1.2$ & \texttt{23.3776070284027512618263} & spectral & $8.1\cdot10^{-57}$ & $2.4\cdot10^{-54}$ & $2.8\cdot10^{-51}$\\
$2.5$ & $1.3$ & \texttt{23.403771830554642453335} & spectral & $7.9\cdot10^{-56}$ & $9.9\cdot10^{-54}$ & $3.2\cdot10^{-51}$\\
$2.75$ & $1.325$ & \texttt{18.6817874548469398300141} & spectral & $4.8\cdot10^{-55}$ & $1.6\cdot10^{-51}$ & $1.6\cdot10^{-51}$\\
$2.75$ & $1.425$ & \texttt{18.6977204798280561759503} & spectral & $2.3\cdot10^{-52}$ & $7.5\cdot10^{-51}$ & $7.5\cdot10^{-51}$\\\addlinespace
$2$ & $0.95$ & & exact \eqref{gff:gs:W3d2} & (ratio, below) & $5.8\cdot10^{-81}$ & \\
$2$ & $1.05$ & & exact \eqref{gff:gs:W3d2} & (ratio, below) & $5.6\cdot10^{-81}$ & \\
$2$ & $1.05/0.95$ & & reflection & $9.8\cdot10^{-81}$ (ratio) & & \\\addlinespace
$1$ & $0.37$ & \texttt{216.134210851145764012993} & fixed nodes & & & $3.4\cdot10^{-62}$\\
$1$ & $0.3+0.1\,\ii$ & \texttt{223.841209588900262170043} & fixed nodes & & & $4.1\cdot10^{-62}$\\
$2$ & $0.8$ & \texttt{45.6639529602144152055429} & fixed nodes & & & $4.5\cdot10^{-63}$\\
$2$ & $1.2+0.05\,\ii$ & \texttt{49.6299926202374115388838} & fixed nodes & & & $1.9\cdot10^{-62}$\\
$3$ & $1.3$ & \texttt{16.4598664890623652308798} & fixed nodes & & & $5.6\cdot10^{-62}$\\
$3$ & $1.6$ & \texttt{15.5022646572795884994581} & fixed nodes & & & $1.2\cdot10^{-62}$\\
$3$ & $1.45+0.05\,\ii$ & \texttt{15.0586428466284212402012} & fixed nodes & & & $3.5\cdot10^{-62}$\\
$1.5$ & $0.9+0.05\,\ii$ & \texttt{86.1127570666953915845079} & fixed nodes & & & $4.6\cdot10^{-56}$\\
$2.5$ & $1.4$ & \texttt{25.6947810198617030782933} & fixed nodes & & & $1.8\cdot10^{-55}$\\
$3.5$ & $1.7$ & \texttt{10.2964255108971126897525} & fixed nodes & & & $6.3\cdot10^{-55}$\\
$0.8$ & $0.42$ & \texttt{263.52442188748393591694} & fixed nodes & & & $3.5\cdot10^{-55}$\\
\bottomrule\end{tabular}
\end{center}

\needspace{14\baselineskip}
\paragraph{Test data.}
The five Horn values and their total first $b$-derivatives at $b=d/2$ ($u=x$ moving in $\mathfrak h_x$), 32 of the at
least 55 digits computed in the 70-digit runs. At $d=1$, $\mathfrak s_0=\frac{83}{81}$, $\mathfrak s_0'=\frac{64}{243}$
and $\mathfrak s_{1/2}=1$ exactly.
\begin{center}\scriptsize
\begin{tabular}{@{}l l l l@{}}\toprule
$d$ & function & value at $b=d/2$ & $\partial_b$ at $b=d/2$\\\midrule
1 & $\mathfrak h_x$ & $1.1186363871641870683496192575256$ & $1.285246154178478427414416657749$\\
 & $\mathfrak h_1$ & $2.5987822357262795715169656086966$ & $7.9424211032565657723245715768249$\\
 & $\mathfrak h_{1/2}$ & $1.6764843387459078434810926719847$ & $3.9229772806249080806880221308364$\\
 & $\mathfrak s_0$ & $1.0246913580246913580246913580247$ & $0.26337448559670781893004115226337$\\
 & $\mathfrak s_{1/2}$ & $1.0$ & $0.0$\\
\addlinespace
1.5 & $\mathfrak h_x$ & $1.177392591060703875855976293685$ & $1.2933702591309377564756247705028$\\
 & $\mathfrak h_1$ & $2.7285923368887621906692408263826$ & $7.0118518320221630889168328700506$\\
 & $\mathfrak h_{1/2}$ & $1.6390896218805694790051124408507$ & $2.9123315039974903478917414402412$\\
 & $\mathfrak s_0$ & $1.0548507244990532977247501057129$ & $0.44524402644527226197972497348997$\\
 & $\mathfrak s_{1/2}$ & $2.2461861367840794529465463415088$ & $3.3311337225326803025747929195969$\\
\addlinespace
2 & $\mathfrak h_x$ & $1.2513473666877450268884376715415$ & $1.4377231145268788593624919295969$\\
 & $\mathfrak h_1$ & $2.8105997380603038789311930383846$ & $6.5768099722236687733793558671552$\\
 & $\mathfrak h_{1/2}$ & $1.6274006603148409180018539998063$ & $2.4813324518913317498491423929508$\\
 & $\mathfrak s_0$ & $1.0836461096219775566294370388355$ & $0.5495504249104147258401786138706$\\
 & $\mathfrak s_{1/2}$ & $3.2548013206296818360037079996126$ & $4.9626649037826634996982847859016$\\
\addlinespace
2.5 & $\mathfrak h_x$ & $1.3376359713522616550547394154586$ & $1.6441674626487503441319240236304$\\
 & $\mathfrak h_1$ & $2.8658700646971037650444991547941$ & $6.312896218106930957543704713237$\\
 & $\mathfrak h_{1/2}$ & $1.6288526582621043253519113017864$ & $2.2658474252647133894689489695958$\\
 & $\mathfrak s_0$ & $1.1103501379648233369119496859969$ & $0.61260865668373670979419911003516$\\
 & $\mathfrak s_{1/2}$ & $4.1368224751542518108391408828496$ & $5.8497720890012283607483552162793$\\
\addlinespace
3 & $\mathfrak h_x$ & $1.4335280971237659036402062676254$ & $1.8841954627648666804462998053832$\\
 & $\mathfrak h_1$ & $2.9047690244187603057569160926227$ & $6.1245414235188740303623918912455$\\
 & $\mathfrak h_{1/2}$ & $1.6357377373331371059669409314044$ & $2.1424411510929036014226183493683$\\
 & $\mathfrak s_0$ & $1.1349627686583730036642106378911$ & $0.65028123499659424582133045301897$\\
 & $\mathfrak s_{1/2}$ & $4.945802884283162665221750648488$ & $6.3455500640456084878526885322906$\\
\bottomrule\end{tabular}
\end{center}

\paragraph{Per-function derivative check.}
For each of the five functions the jet value and derivative at $b=d/2$ were compared, at 40 digits, with the mean and
the first Fourier coefficient of its plain (non-jet) values on a 48-point circle of radius
$0.05\min(\frac d2,b_{\max}-\frac d2)$ around $d/2$ (\path{LRI_C5_general_d/L/check_func_derivs.py},
\path{LRI_C5_general_d/L/results/cfd40.json}). This checks the jet implementation for every function separately; it is
not an independent evaluation. Each cell gives the relative deviation of the value / of the derivative, the latter
relative to $\max(|f'|,1)$ (at $d=1$, $\mathfrak s_{1/2}=1$ has $f'=0$).
\begin{center}\scriptsize
\begin{tabular}{@{}l l l l l l@{}}\toprule
$d$ & $\mathfrak h_x$ & $\mathfrak h_1$ & $\mathfrak h_{1/2}$ & $\mathfrak s_0$ & $\mathfrak s_{1/2}$\\\midrule
1 & $1{\cdot}10^{-45}$ / $3{\cdot}10^{-45}$ & $6{\cdot}10^{-46}$ / $4{\cdot}10^{-46}$ & $2{\cdot}10^{-45}$ / $1{\cdot}10^{-44}$ & $3{\cdot}10^{-45}$ / $2{\cdot}10^{-44}$ & $<10^{-60}$ / $7{\cdot}10^{-56}$\\
1.5 & $4{\cdot}10^{-45}$ / $2{\cdot}10^{-44}$ & $4{\cdot}10^{-45}$ / $1{\cdot}10^{-44}$ & $3{\cdot}10^{-45}$ / $2{\cdot}10^{-44}$ & $4{\cdot}10^{-45}$ / $2{\cdot}10^{-44}$ & $5{\cdot}10^{-44}$ / $3{\cdot}10^{-43}$\\
2 & $2{\cdot}10^{-45}$ / $9{\cdot}10^{-45}$ & $6{\cdot}10^{-46}$ / $2{\cdot}10^{-45}$ & $1{\cdot}10^{-45}$ / $5{\cdot}10^{-45}$ & $4{\cdot}10^{-46}$ / $1{\cdot}10^{-44}$ & $1{\cdot}10^{-45}$ / $1{\cdot}10^{-45}$\\
2.5 & $2{\cdot}10^{-45}$ / $2{\cdot}10^{-45}$ & $1{\cdot}10^{-45}$ / $2{\cdot}10^{-45}$ & $8{\cdot}10^{-46}$ / $9{\cdot}10^{-45}$ & $2{\cdot}10^{-48}$ / $7{\cdot}10^{-45}$ & $3{\cdot}10^{-44}$ / $3{\cdot}10^{-43}$\\
3 & $1{\cdot}10^{-45}$ / $8{\cdot}10^{-45}$ & $2{\cdot}10^{-45}$ / $9{\cdot}10^{-46}$ & $1{\cdot}10^{-45}$ / $5{\cdot}10^{-45}$ & $2{\cdot}10^{-45}$ / $8{\cdot}10^{-45}$ & $3{\cdot}10^{-46}$ / $4{\cdot}10^{-45}$\\
\bottomrule\end{tabular}
\end{center}

\paragraph{Summary.}
At $d=1,2,3$, $\LL_d$ and $P_4(d)$ agree with the proved or 160-digit values to $10^{-60}$ or better at 50 digits and
to $10^{-80}$ or better at 70 digits. At the five non-integer dimensions with 41- to 50-digit spectral references the
agreement is $10^{-45}$ to $10^{-47}$, that is the accuracy of the references (the internal convergence is $10^{-45}$
to $10^{-49}$ at 50 digits and about $10^{-58}$ at 70 digits); at $d=0.6,\frac43,\frac83,3.5$ it is $10^{-24}$ to
$10^{-33}$, limited by the 32-digit spectral runs. Away from the center, $W_4(d;b)$ and $W_3(d;b)$ at the ten points
$b=d/2\pm0.05$ ($\pm0.03$ at $d=0.8$) agree with the 56-digit spectral values to $1.5\cdot10^{-48}$ or better ($W_4$:
$3.5\cdot10^{-50}$ or better), and at $d=2$ with the exact $W_3(2;b)$ and the rim--spoke reflection to $10^{-80}$. The
$r=3$ by-product $3Q_d$ agrees with the references at the eight dimensions where they exist to $10^{-45}$ to
$10^{-81}$, and $\partial_bW_3(d;b)|_{b=d/2}$, which vanishes by \eqref{gff:gs:Tjet} with \eqref{gff:gs:wheelJets}, comes out
$10^{-47}$ to $10^{-86}$ relative to $W_3$; this tests the jets of the $r=3$ functions. Each jet derivative agrees with
the 48-point Cauchy-circle derivative to $3\cdot10^{-43}$ or better at 40 digits (previous table). The accelerated
evaluation is an asymptotic-series method; its accuracy rests on the convergence in the truncation parameters and on
the independent comparisons above, not on interval bounds.


\subsection{The three-wheel at a general rim exponent: proofs}\label{gff:mgz:proofs}
The code of this subsection is in \path{LRI_C5_general_d/Z/}. A separate review with
its own codes rederived the derivations of the package's own text
(\path{LRI_C5_general_d/review/Z-rigor/}; for the steps added in this note, see
Section~\ref{gff:gdr:status}).

This subsection proves the three-wheel statements of Section~\ref{gff:gdr:samples}:
the two samples, the two-node formula \eqref{gff:gdr:twonode} with its zero-mode
form, \eqref{gff:gdr:Q} and \eqref{gff:gdr:Z}. It also gives the five functions behind
$\mathfrak t_3$ and $\widehat M_3$, the jets of the weights $u,v$, series forms of
the jets and the first-derivative identity. Throughout $r=3$, $0<d<4$, and $b$
lies in the strip
\[
 0<\operatorname{Re}b<b_{\max}(3,d)=\min\Bigl\{\frac{2d}3,\ \frac{d+2}3,\ d-\frac13\Bigr\}<2
\]
of \eqref{gff:gdr:bmax}, which contains $b=\frac d2$ if and only if $d>\frac23$. On
this strip $0<\operatorname{Re}\sigma<1$, so $\sigma\notin\mathbb Z$. The inputs are
the finite Fourier theorem of Section~\ref{gff:gdr:finite}, the top-mode theorem
\eqref{gff:gdr:topless}, the residue functional \eqref{gff:gdr:explicit}, the reflected
form \eqref{gff:gdr:reflected} and the block decomposition of
Section~\ref{gff:gdr:residue}. By \eqref{gff:gdr:topless} the three-wheel sample has two
modes,
\begin{equation}
 \Xi_3(t;b)=c_{3,0}(q,b)+c_{3,1}(q,b)\cos2\pi\Bigl(t+\frac b4\Bigr).
 \label{gff:mgz:twomodes}
\end{equation}

\subsubsection{Analyticity in the rim exponent}\label{gff:mgz:uniform}
\emph{Statement.} $\mathfrak t_3$ and $\widehat M_3$ of \eqref{gff:gdr:samples3}, the
five functions of \eqref{gff:mgz:five} below and the weighted series of
\eqref{gff:mgz:series} are analytic in $b$ on the strip. Every $b$-derivative is the
termwise derivative, and the differentiated series again converge locally
uniformly.

\emph{Proof.} Let $K$ be a compact subset of the strip and
$\alpha_K=\min_{b\in K}3\operatorname{Re}g$. Since $\operatorname{Re}b<b_{\max}(3,d)$
on $K$, the conditions (i)--(iii) stated after \eqref{gff:gdr:bmax} hold at every
point of $K$ with $\alpha=3\operatorname{Re}g$, and by compactness
$\alpha_K>\max(\frac d2,d-1,\frac12)$. The estimates of
Section~\ref{gff:gdr:periodization} (paragraph ``Lattice split and estimates'') come
from the Gamma-ratio bound $|\Gamma(z+u)/\Gamma(z+v)|\le C|z|^{\operatorname{Re}(u-v)}$,
which is uniform for $u,v$ in compact sets. Here it is applied to Gamma ratios
whose parameters $\sigma$, $g$, $2q$ are affine in $b$, and only its part for $t$
in a compact set is used, not the large-$|\operatorname{Im}t|$ part. Write
$T_{\ell n}=[(\sigma)_{\ell+n}/(2q)_{\ell+n}]\,[(1-g)_n/n!]$ as in \eqref{gff:gdr:Tln}.
It is a product of two polynomials in $b$, bounded on $K$ by
$C_K(1+\ell+n)^{-\operatorname{Re}g}$ and $C_K(1+n)^{-\operatorname{Re}g}$. With
$|d_\ell|\le C(1+\ell)^{d-2}$ there is a constant $C_K$ with
\[
 |d_\ell\,T_{\ell n}(b)^3|\le C_K(1+\ell)^{d-2}(1+\ell+n)^{-\alpha_K}(1+n)^{-\alpha_K}
 \qquad(\ell,n\ge0,\ b\in K).
\]
The terms of $\widehat M_3$ are $N_b^3g_\ell(a_\ell+n+t_b;b)^3$ with
$t_b=\frac{1-\sigma}2$, since $a_\ell+n+t_b=n+\frac{\ell+1}2$. For $b\in K$ the
node $t_b$ stays at a positive distance from both pole classes $t\equiv0$ and
$t\equiv-\sigma\pmod 1$, because $\operatorname{Re}t_b$ and
$\operatorname{Re}(t_b+\sigma)=\operatorname{Re}\frac{1+\sigma}2$ lie in compact
subintervals of $(0,\frac12)$ and $(\frac12,1)$. So the right part ($n\ge0$, by the
first line of \eqref{gff:gdr:reflected}) and the middle part (by its second line) obey
the two bounds of Section~\ref{gff:gdr:periodization} with a constant independent of
$b\in K$; the finitely many terms with small arguments are continuous in $b$ on
$K$. The left part equals the right part, because $1-\sigma-t_b=t_b$ and the last
identity of \eqref{gff:gdr:reflected} holds. These majorants do not depend on $b\in K$,
and they are summable by the criterion at the end of the same paragraph, since
$\alpha_K>\max(\frac d2,d-1,\frac12)$. Every term is analytic in $b$ on the strip:
$T_{\ell n}$ is a polynomial in $b$, and at the fixed lattice points
$s=n+\frac{\ell+1}2$ the numerator Gamma functions of $N_bg_\ell(s;b)$ have the
arguments $\ell+n+\frac{1+\sigma}2$, $\frac{\sigma-1}2-n$ and $g$, none of which
is a pole for $0<\operatorname{Re}\sigma<1$ and $\operatorname{Re}g>0$. By the
Weierstrass $M$-test and Weierstrass's theorem on locally uniform limits of
analytic functions, $\mathfrak t_3$ and $\widehat M_3$ are analytic on the strip,
every $b$-derivative is the termwise derivative, and the differentiated series
again converge locally uniformly. The same holds for the subseries of
\eqref{gff:mgz:five}. For the weighted series of \eqref{gff:mgz:series} it can also be seen
directly: the weights grow at most logarithmically in the indices, and replacing
$\alpha_K$ by a slightly smaller exponent, still above $\max(\frac d2,d-1,\frac12)$,
gives summable majorants that are uniform on $K$. $\square$

\subsubsection{The two samples and the two-node formula}\label{gff:mgz:nodes}
\emph{Pole node.} For $b$ in the strip,
\begin{equation}
 \Xi_3(0;b)=-\Bigl(\frac{\pi s_b}{N_b}\Bigr)^3\mathfrak t_3(b).
 \label{gff:mgz:node0}
\end{equation}
This is the case $r=3$ of \eqref{gff:mgl:sampT}, which Section~\ref{gff:mgl:section} proves for every $r$ by extending
the lattice form \eqref{gff:mgl:tauform} to $t=0$; the proof below takes the limit term by term instead.

\emph{Proof.} By Section~\ref{gff:gdr:finite}(a) and its proof,
$t^3S_3(t;b)=\sum_\ell d_\ell\sum_nt^3g_\ell(a_\ell+n+t;b)^3$ converges uniformly
near $t=0$, and only the terms with $n\ge0$ are singular at $t=0$: the poles of the
middle part in the class $t\equiv0$ lie at $t\ge1$, by the second line of
\eqref{gff:gdr:reflected}, the same holds for the left part by the mirror identity,
and the class $t\equiv-\sigma$ does not contain $0$ since $\sigma\notin\mathbb Z$.
Since $D_b(t)/t\to\pi s_b$, the value $\Xi_3(0)=\lim_{t\to0}D_b(t)^3S_3(t)$ is
the sum over $\ell\ge0$, $n\ge0$ of $d_\ell\lim_{t\to0}[D_b(t)g_\ell(a_\ell+n+t)]^3$.
By the first line of \eqref{gff:gdr:reflected},
\[\begin{aligned}
 D_b(t)\,g_\ell(a_\ell+n+t)\ \to\ {}&-s_b\sin(\pi g)\,
 \frac{\Gamma(\ell+n+\sigma)\Gamma(n+1-g)}{\Gamma(\ell+n+2q)\,n!}\\
 ={}&-s_b\sin(\pi g)\Gamma(1-g)\frac{\Gamma(\sigma)}{\Gamma(2q)}\,T_{\ell n}
 =-\frac{\pi s_b}{N_b}\,T_{\ell n},
\end{aligned}\]
using $\sin(\pi g)\Gamma(1-g)=\pi/\Gamma(g)$. Cubing and summing gives
\eqref{gff:mgz:node0}. $\square$

The function $T_{\ell n}$ is the normalized residue prefactor at a general rim
exponent: $N_bg_\ell(a_\ell+n+u;b)=-T_{\ell n}(b)/u+O(1)$ for $n\ge0$ (as in
Section~\ref{gff:gdr:taylor}), so
$\mathfrak t_3=-\sum_{\ell,n}d_\ell\lim_{u\to0}u^3[N_bg_\ell(a_\ell+n+u;b)]^3$ collects
only the \emph{leading} coefficients of the triple poles. The residues of the cube
need the full Laurent data, which is the origin of the polygamma polynomial $F_3$
of \eqref{gff:gs:QPseries}.

\emph{Self-mirror node.} Let $t_b=\frac{1-\sigma}2=\frac12-\frac b4$. For $b$ in
the strip, $1-\sigma-t_b=t_b$, $\cos2\pi(t_b+\frac b4)=-1$,
$D_b(t_b)=\cos^2\frac{\pi\sigma}2=\frac{1+\varsigma_b}2$, and
\begin{equation}
 \Xi_3(t_b;b)=\Bigl(\frac{1+\varsigma_b}2\Bigr)^3\frac{\widehat M_3(b)}{N_b^3},\qquad
 \widehat M_3(b)=N_b^3\bigl[2A_3^0(t_b)+2A_3^+(t_b)+C_3(t_b)\bigr],
 \label{gff:mgz:nodeb}
\end{equation}
with the blocks $A_3^0$, $A_3^+$, $C_3$ of Section~\ref{gff:gdr:residue}, written out
in \eqref{gff:mgz:blocks3} below. In the Horn values of Section~\ref{gff:mgl:section} the
same sample is \eqref{gff:mgl:samph} at $r=3$.

\emph{Proof.} The first three statements are direct:
$D_b(t_b)=\sin\frac{\pi(1-\sigma)}2\sin\frac{\pi(1+\sigma)}2=\cos^2\frac{\pi\sigma}2$
and $2\cos^2\frac{\pi\sigma}2=1+\cos\pi\sigma=1+\varsigma_b$. Since
$0<\operatorname{Re}\sigma<1$, $t_b$ lies in neither pole class, so
$\Xi_3(t_b)=D_b(t_b)^3S_3(t_b)$ with an absolutely convergent $S_3(t_b)$. Moreover
$a_\ell+n+t_b=n+\frac{\ell+1}2$; as $n$ runs over $\mathbb Z$ this runs over the
midpoint lattice $\frac{\ell+1}2+\mathbb Z$ of \eqref{gff:gdr:samples3}, hence
$N_b^3S_3(t_b;b)=\widehat M_3(b)$. Finally the block decomposition of
Section~\ref{gff:gdr:residue} at a point equal to its own mirror point $1-\sigma-t$
gives the block form. $\square$

The midpoint lattice is invariant under $s\mapsto-s$, as is $g_\ell(s;b)$.

\emph{Proof of \eqref{gff:gdr:twonode}.} By \eqref{gff:mgz:twomodes},
$y_0:=\Xi_3(0;b)=c_{3,0}+\varsigma_bc_{3,1}$ and
$y_b:=\Xi_3(t_b;b)=c_{3,0}-c_{3,1}$. On the strip $1+\varsigma_b\neq0$
($\varsigma_b=-1$ only for $b\in2+4\mathbb Z$), so
$c_{3,0}=(y_0+\varsigma_by_b)/(1+\varsigma_b)$ and
$c_{3,1}=(y_0-y_b)/(1+\varsigma_b)$. Then
\[
 (1+2\varsigma_b^2)c_{3,0}+3\varsigma_bc_{3,1}
 =\frac{(1+\varsigma_b)(1+2\varsigma_b)\,y_0-2\varsigma_bs_b^2\,y_b}{1+\varsigma_b},
\]
and the first line of \eqref{gff:gdr:explicit} gives
\[
 W_3(d;b)=-\frac{2N_b^3(1+2\varsigma_b)}{\pi s_b^5}\,y_0
 +\frac{4N_b^3\varsigma_b}{\pi s_b^3(1+\varsigma_b)}\,y_b .
\]
Inserting \eqref{gff:mgz:node0} and \eqref{gff:mgz:nodeb} gives \eqref{gff:gdr:twonode} for
every $b$ in the strip and every $0<d<4$. $\square$ (With the midpoint sample in
Horn values the same formula is \eqref{gff:mgl:W3three}.)

\emph{Zero-mode form.} With the normalized zero mode
$\hat c_0(b)=N_b^3c_{3,0}(q,b)=N_b^3\int_0^1\Xi_3(t;b)\dd t$,
\begin{equation}
 W_3(d;b)=\frac{6\pi^2}{s_b^2}\mathfrak t_3(b)+\frac{4}{\pi s_b^3}\hat c_0(b),\qquad
 \hat c_0(b)=\frac{\varsigma_b(1+\varsigma_b)^2}8\widehat M_3(b)
 -\frac{\pi^3s_b^3}{1+\varsigma_b}\mathfrak t_3(b).
 \label{gff:mgz:zeromode}
\end{equation}
\emph{Proof.} $\hat c_0=N_b^3(y_0+\varsigma_by_b)/(1+\varsigma_b)$ with
\eqref{gff:mgz:node0} and \eqref{gff:mgz:nodeb} gives the second relation. Since
$y_0=c_{3,0}+\varsigma_bc_{3,1}$, one has
$(1+2\varsigma_b^2)c_{3,0}+3\varsigma_bc_{3,1}=3y_0-2s_b^2c_{3,0}$ identically (no
division by $\varsigma_b$ occurs, so $\varsigma_b=0$ is allowed), and the first line
of \eqref{gff:gdr:explicit} with \eqref{gff:mgz:node0} gives the first relation. $\square$
This is the form quoted after \eqref{gff:gdr:twonode}; at $b=\frac d2$, where
$W_3(d;\frac d2)=3Q_d$ by \eqref{gff:fp:wheelStart}, it is \eqref{gff:gdr:Q}.

At $b=\frac d2$, \eqref{gff:gdr:twonode} is a finite formula for the primitive,
$3Q_d=u(\frac d2)\mathfrak t_3(\frac d2)+v(\frac d2)\widehat M_3(\frac d2)$, with five
real hypergeometric values at the two real nodes $0$ and $\frac{1-q}2$, whereas
\eqref{gff:fp:periodsFinite} uses three complex samples with five blocks each. The extra
input is the top-mode theorem \eqref{gff:gdr:topless}. For complex $b$ the node $t_b$ is
complex; both nodes are real for real $b$.

\subsubsection{The five functions}\label{gff:mgz:fiveSection}
At $r=3$ and a general rim exponent the blocks of Section~\ref{gff:gdr:residue} are
\begin{equation}\begin{aligned}
 A_3^0(t)&=a_3^0(t)\,{}_7F_6\!\left(\begin{matrix}1,(\sigma+t)^{\times3},(1-g+t)^{\times3}\\
 (2q+t)^{\times3},(1+t)^{\times3}\end{matrix};1\right),&
 a_3^0(t)&=\Bigl[\frac{\Gamma(\sigma+t)\Gamma(-t)}{\Gamma(2q+t)\Gamma(g-t)}\Bigr]^3,\\
 A_3^+(t)&=a_3^+(t)\,\mathfrak H_+(t),&
 a_3^+(t)&=4q\Bigl[\frac{\Gamma(1+\sigma+t)\Gamma(-t)}{\Gamma(1+2q+t)\Gamma(g-t)}\Bigr]^3,\\
 C_3(t)&=a_3^c(t)\,\mathfrak H_c(t),&
 a_3^c(t)&=4q\Bigl[\frac{\Gamma(\sigma+t)\Gamma(1-t)}{\Gamma(2q+t)\Gamma(1+g-t)}\Bigr]^3,
\end{aligned}\label{gff:mgz:blocks3}\end{equation}
with the two Horn values
\begin{equation}\begin{aligned}
 \mathfrak H_+(t)&=\sum_{j,n\ge0}\frac{d_{j+1}}{4q}
 \Bigl[\frac{(1+\sigma+t)_{j+n}}{(1+2q+t)_{j+n}}\Bigr]^3\Bigl[\frac{(1-g+t)_n}{(1+t)_n}\Bigr]^3,\\
 \mathfrak H_c(t)&=\sum_{j,k\ge0}\frac{d_{j+k+1}}{4q}
 \Bigl[\frac{(1-t)_j}{(1+g-t)_j}\Bigr]^3\Bigl[\frac{(\sigma+t)_k}{(2q+t)_k}\Bigr]^3 .
\end{aligned}\label{gff:mgz:horn3}\end{equation}
These are \eqref{gff:fp:wheelBlocks} and $\mathfrak H[\mathsf U_+;\mathsf V_+;(1,1)]$,
$\mathfrak H[\mathsf U_c;\mathsf V_c;(1,1)]$ of \eqref{gff:fp:wheelArrays} at $r=3$ with
the rim parameter $q\to\sigma$ and the spoke parameter $q\to g$, as stated in
Section~\ref{gff:gdr:residue}. Here $d_{j+1}/(4q)=(2q+1)_j(4q-1)_j/[(2q)_j(2)_j]$ by
\eqref{gff:fp:degeneracy}; it equals $\delta_{j0}$, $1$ and $1+\frac23j$ at $d=1,2,3$.

Split $\mathfrak t_3=\mathfrak t_{3,0}+\mathfrak t_{3,+}$ into the spin $\ell=0$ and
$\ell\ge1$, and $\widehat M_3=\widehat M_{3,0}+\widehat M_{3,+}+\widehat M_{3,c}$ by
\eqref{gff:mgz:nodeb}:
\begin{equation}\begin{aligned}
 \mathfrak t_{3,0}(b)&={}_7F_6\!\left(\begin{matrix}1,\sigma,\sigma,\sigma,1-g,1-g,1-g\\
 2q,2q,2q,1,1,1\end{matrix};1\right),&
 \mathfrak t_{3,+}(b)&=4q\Bigl(\frac{\sigma}{2q}\Bigr)^3\mathfrak H_+(0),\\
 \widehat M_{3,0}(b)&=2N_b^3A_3^0(t_b),\qquad \widehat M_{3,+}(b)=2N_b^3A_3^+(t_b),&
 \widehat M_{3,c}(b)&=N_b^3C_3(t_b).
\end{aligned}\label{gff:mgz:five}\end{equation}
The spin-zero row is the ${}_7F_6(1)$ because
$T_{0n}=(\sigma)_n(1-g)_n/[(2q)_n\,n!]$, the upper parameter $1$ cancelling one
$n!$; $\mathfrak t_{3,+}$ follows from
$T_{j+1,n}=\frac{\sigma}{2q}\frac{(1+\sigma)_{j+n}(1-g)_n}{(1+2q)_{j+n}\,n!}$ and
$d_{j+1}=4q\cdot d_{j+1}/(4q)$. (The subscript $0$ of $\mathfrak t_{3,0}$ and
$\widehat M_{3,0}$ names the spin-zero block; these are not Fourier coefficients
$c_{3,0}$.) At the self-mirror node, $\sigma+t_b=1-t_b=\frac{1+\sigma}2$ and
$2q+t_b=1+g-t_b=2q+\frac{1-\sigma}2$, so the two sequences in $\mathfrak H_c(t_b)$
coincide and
\[
 \mathfrak H_c(t_b)=\sum_{j,k\ge0}\frac{d_{j+k+1}}{4q}\,\mathfrak q_j\mathfrak q_k,
 \qquad
 \mathfrak q_k=\Bigl[\frac{(\frac{1+\sigma}2)_k}{(2q+\frac{1-\sigma}2)_k}\Bigr]^3 .
\]
The remaining parameters at $t_b$ are $1-g+t_b=\frac{3+\sigma}2-2q$ and
$1+t_b=\frac{3-\sigma}2$. Every parameter is affine in $b$, so the $b$-derivatives
of the five functions are directional parameter derivatives
(table below). The prefactors $a_3^0,a_3^+,a_3^c$ at $t_b$ and $N_b$
contribute only polygamma values at fixed arguments. So $\mathfrak t_3$ is one
${}_7F_6(1)$ and one two-variable Horn value, and $\widehat M_3$ is one ${}_7F_6(1)$
and two two-variable Horn values, as stated in Section~\ref{gff:gdr:samples}.

\begin{center}\footnotesize
\begin{tabular}{@{}llll@{}}\toprule
function & upper parameters (slope in $b$) & lower parameters (slope in $b$) & at $b=d/2$\\\midrule
$\mathfrak t_{3,0}$: ${}_7F_6$ at $t=0$ & $\sigma^{\times3}$ ($\frac12$), $(1-g)^{\times3}$ ($\frac12$) & $(2q)^{\times3}$, $1^{\times3}$ (0)
& $q,\ 1-q;\ 2q,\ 1$\\
$\mathfrak H_+(0)$ & $1+\sigma$ ($\frac12$), $1-g$ ($\frac12$) & $1+2q$, $1$ (0) & $1+q,\ 1-q;\ 1+2q,\ 1$\\
$A_3^0(t_b)$: ${}_7F_6$ & $\frac{1+\sigma}2$ ($\frac14$), $\frac{3+\sigma}2-2q$ ($\frac14$) &
$2q+\frac{1-\sigma}2$ ($-\frac14$), $\frac{3-\sigma}2$ ($-\frac14$) & $\frac{1+q}2,\ \frac{3-3q}2;\ \frac{1+3q}2,\ \frac{3-q}2$\\
$\mathfrak H_+(t_b)$ & $\frac{3+\sigma}2$ ($\frac14$), $\frac{3+\sigma}2-2q$ ($\frac14$) &
$1+2q+\frac{1-\sigma}2$ ($-\frac14$), $\frac{3-\sigma}2$ ($-\frac14$) & $\frac{3+q}2,\ \frac{3-3q}2;\ \frac{3+3q}2,\ \frac{3-q}2$\\
$\mathfrak H_c(t_b)$ & $\frac{1+\sigma}2$ ($\frac14$), twice & $2q+\frac{1-\sigma}2$ ($-\frac14$), twice &
$\frac{1+q}2;\ \frac{1+3q}2$\\\bottomrule
\end{tabular}
\end{center}
The table lists the parameters of the five hypergeometric functions of
\eqref{gff:mgz:five} (each parameter repeated three times, $r=3$) and their slopes in
$b$. The multiplicity factor $d_{j+1}/(4q)$ (or $d_{j+k+1}/(4q)$) of the Horn values
and the prefactor $4q$ do not depend on $b$.

\subsubsection{Rim jets and their series}\label{gff:mgz:jetsSection}
Primes denote $\partial_b$; by Section~\ref{gff:mgz:uniform} every function above is
analytic in $b$ on the strip. For $\frac23<d<4$, at $b=\frac d2$,
\begin{equation}
 9Z_d=u''\mathfrak t_3+2u'\mathfrak t_3'+u\mathfrak t_3''
 +v''\widehat M_3+2v'\widehat M_3'+v\widehat M_3'' ,
 \label{gff:mgz:Zjets}
\end{equation}
with $\mathfrak t_3^{(k)}=\mathfrak t_{3,0}^{(k)}+\mathfrak t_{3,+}^{(k)}$ and
$\widehat M_3^{(k)}=\widehat M_{3,0}^{(k)}+\widehat M_{3,+}^{(k)}+\widehat M_{3,c}^{(k)}$
the jets of the five functions \eqref{gff:mgz:five}, and, with $S,C$ of \eqref{gff:gdr:D},
\begin{equation}\begin{aligned}
 u&=\frac{2\pi^2(1+2C)}{S^2},& u'&=-\frac{2\pi^3(1+C+C^2)}{S^3},&
 u''&=\frac{\pi^4(1+5C+2C^2+C^3)}{S^4},\\
 v&=\frac{C(1+C)^2}{2\pi S^3},& v'&=-\frac{(1+C)^2(1+2C)}{4S^4},&
 v''&=\frac{\pi(1+C)^3(2+C)}{4S^5}.
\end{aligned}\label{gff:mgz:uvjets}\end{equation}
\emph{Proof.} By \eqref{gff:gs:wheelJets}, $9Z_d=\partial_b^2W_3(d;b)$ at $b=\frac d2$,
which lies in the strip for $\frac23<d<4$. Differentiate \eqref{gff:gdr:twonode} twice,
termwise by Section~\ref{gff:mgz:uniform}, using $\partial_bs_b=\frac\pi2\varsigma_b$ and
$\partial_b\varsigma_b=-\frac\pi2s_b$; for example
$u'=-2\pi^3[s_b^2+\varsigma_b(1+2\varsigma_b)]/s_b^3$ and
$s_b^2+\varsigma_b+2\varsigma_b^2=1+\varsigma_b+\varsigma_b^2$. $\square$
Formula \eqref{gff:mgz:Zjets} is exact differentiation and does not use the input
$\partial_bW_3(d;\frac d2)=0$.

\emph{Series form of the jets.} On the strip,
\begin{equation}\begin{aligned}
 \mathfrak t_3'&=3\sum_{\ell,n\ge0}d_\ell T_{\ell n}^3\mu_1,&
 \mathfrak t_3''&=\sum_{\ell,n\ge0}d_\ell T_{\ell n}^3\bigl(9\mu_1^2+3\mu_2\bigr),\\
 \widehat M_3'&=3\sum_{\ell\ge0}d_\ell\sum_{s}[N_bg_\ell(s;b)]^3w_1,&
 \widehat M_3''&=\sum_{\ell\ge0}d_\ell\sum_{s}[N_bg_\ell(s;b)]^3\bigl(9w_1^2+3w_2\bigr),
\end{aligned}\label{gff:mgz:series}\end{equation}
with $s\in\frac{\ell+1}2+\mathbb Z$ and the explicit weights (arguments
$(\ell,n;b)$ of $\mu_m$ and $(\ell,s;b)$ of $w_m$ suppressed)
\begin{equation}\begin{aligned}
 \mu_1&=\tfrac12\bigl[\psi_0(\sigma+\ell+n)-\psi_0(\sigma)+\psi_0(1-g+n)-\psi_0(1-g)\bigr],\\
 \mu_2&=\tfrac14\bigl[\psi_1(\sigma+\ell+n)-\psi_1(\sigma)+\psi_1(1-g+n)-\psi_1(1-g)\bigr],\\
 w_1&=-\tfrac12\bigl[\psi_0(g)+\psi_0(\sigma)\bigr]
 +\tfrac14\textstyle\sum_\pm\bigl[\psi_0(a_\ell\pm s)+\psi_0(\beta_\ell\pm s)\bigr],\\
 w_2&=\tfrac14\bigl[\psi_1(g)-\psi_1(\sigma)\bigr]
 +\tfrac1{16}\textstyle\sum_\pm\bigl[\psi_1(a_\ell\pm s)-\psi_1(\beta_\ell\pm s)\bigr].
\end{aligned}\label{gff:mgz:weights}\end{equation}
At $b=\frac d2$, $\mu_1,\mu_2$ are those of \eqref{gff:gs:etaTheta}, and $w_1,w_2$ are
the spectral weights of \eqref{gff:gs:w12} continued to $s=\ii\nu$. Consequently
\begin{equation}
 9Z_d=\sum_{\ell,n\ge0}d_\ell T_{\ell n}^3\bigl[(9\mu_1^2+3\mu_2)\,u+6\mu_1\,u'+u''\bigr]
 +\sum_{\ell\ge0}d_\ell\sum_{s\in\frac{\ell+1}2+\mathbb Z}[N_bg_\ell(s;b)]^3
 \bigl[(9w_1^2+3w_2)\,v+6w_1\,v'+v''\bigr],
 \label{gff:mgz:Zseries2}
\end{equation}
at $b=\frac d2$, with $u,\dots,v''$ of \eqref{gff:mgz:uvjets}.

\emph{Proof.} $\mu_1=\partial_b\log T_{\ell n}$ and $\mu_2=\partial_b\mu_1$, since
$\partial_b\sigma=\frac12=\partial_b(1-g)$. At fixed $s$,
$w_1=\partial_b\log[N_bg_\ell(s;b)]$ and $w_2=\partial_bw_1$, since
$\partial_ba_\ell=\frac14=-\partial_b\beta_\ell$ and
$\partial_b\log N_b=-\frac12[\psi_0(g)+\psi_0(\sigma)]$. The midpoint lattice does not
move with $b$, so \eqref{gff:mgz:series} are the termwise derivatives of
\eqref{gff:gdr:samples3}, which Section~\ref{gff:mgz:uniform} justifies; \eqref{gff:mgz:Zseries2}
is \eqref{gff:mgz:Zjets} with these series. The weights were also checked term by term
against numerical differentiation, to $10^{-30}$
(\path{LRI_C5_general_d/Z/check_weights.py}). $\square$

\emph{Comparison with \eqref{gff:gs:Zseries}.} The series \eqref{gff:gs:Zseries} is
$9Z_d=\sum d_\ell T_{\ell n}^3[(9\mu_1^2+3\mu_2)F_3+6\mu_1F_{3,1}+F_{3,2}]$, with the
index-dependent triple-pole weight $F_3$ of \eqref{gff:gs:Fpolys} and its rim
derivatives $F_{3,1},F_{3,2}$. In \eqref{gff:mgz:Zseries2} the same structure appears
with $F_3,F_{3,1},F_{3,2}$ replaced by the index-independent elementary numbers
$u,u',u''$, at the price of the midpoint-lattice term. In that term
$\frac19\widehat M_3''$ is the midpoint-lattice sum of the spectral integrand
$t_\ell^3(w_1^2+w_2/3)$ of \eqref{gff:gs:LZspectral}, whose integral is $Z_d$ itself.

\subsubsection{\texorpdfstring{The first-derivative identity and the simplified formula for $Z_d$}{The first-derivative identity and the simplified formula for Zd}}\label{gff:mgz:simplifiedSection}
\emph{First rim derivative.} $\partial_bW_3(d;\frac d2)=0$ is equivalent to
$u'\mathfrak t_3+u\mathfrak t_3'+v'\widehat M_3+v\widehat M_3'=0$, that is, at
$b=\frac d2$,
\begin{equation}
 8\pi^3S^2(1+2C)\,\mathfrak t_3'+2SC(1+C)^2\,\widehat M_3'
 =8\pi^4S(1+C+C^2)\,\mathfrak t_3+\pi(1+C)^2(1+2C)\,\widehat M_3 ,
 \label{gff:mgz:first}
\end{equation}
and, in the zero-mode form \eqref{gff:mgz:zeromode}, to
\begin{equation}
 3\pi^3S^2\,\mathfrak t_3'+2S\,\hat c_0'=3\pi C\,\bigl(\pi^3S\,\mathfrak t_3+\hat c_0\bigr).
 \label{gff:mgz:firstzm}
\end{equation}
\emph{Proof.} The first is $4\pi S^4\,\partial_bW_3$ with \eqref{gff:mgz:uvjets}. For the
second write $W_3=u_0\mathfrak t_3+v_0\hat c_0$ with $u_0=6\pi^2/s_b^2$ and
$v_0=4/(\pi s_b^3)$, by \eqref{gff:mgz:zeromode}; at $b=\frac d2$,
$u_0'=-6\pi^3C/S^3$ and $v_0'=-6C/S^4$, and \eqref{gff:mgz:firstzm} is
$\frac{\pi S^4}2\,\partial_bW_3=0$. $\square$

Given the input $\partial_bW_3(d;\frac d2)=0$, which is proved by Pattern~T
\eqref{gff:gs:Tjet} with \eqref{gff:gs:wheelJets} (Section~\ref{gff:gdr:samples}), the first
derivative of the residue sum is not independent:
\begin{equation}
 \mathfrak t_3'=\frac{\pi C}S\,\mathfrak t_3+\frac{C}{\pi^2S^2}\,\hat c_0
 -\frac2{3\pi^3S}\,\hat c_0'\qquad\bigl(b=\tfrac d2\bigr)
 \label{gff:mgz:Rprime}
\end{equation}
for every $d$. In the basis $(\mathfrak t_3,\widehat M_3)$, \eqref{gff:mgz:first}
eliminates $\mathfrak t_3'$ for $d\neq\frac83$ (where $1+2C=0$), or $\widehat M_3'$ for
$d\neq2$ (where $C=0$). At $d=1$, \eqref{gff:mgz:firstzm} reads
$\mathfrak t_3'(1;\frac12)=\pi\,\mathfrak t_3(1;\frac12)$, which Dixon's closed form
\eqref{gff:mgz:Dixon1} confirms. At $d=3$, with the Gamma form of $\hat c_0(3;b)$
(Section~\ref{gff:gdr:halfShift}), it reads
$\mathfrak t_3'(3;\frac32)=-\pi\,\mathfrak t_3(3;\frac32)+\Gamma(\frac14)^{10}/(64\pi^{9/2})$,
an identity also found in the $d=3$ computation of \path{LRI_C5_d3/wf_Z3/}.

\emph{Proof of \eqref{gff:gdr:Z}.} Assume $\partial_bW_3(d;\frac d2)=0$. At $b=\frac d2$,
with $\hat c_0^{(k)}=\partial_b^k\hat c_0|_{b=d/2}$,
\begin{equation}
 Z_d=\frac{2\pi^2}{3S^2}\,\mathfrak t_3''-\frac{\pi^4\cos(\pi d/2)}{3S^4}\,\mathfrak t_3
 +\frac19\Bigl[\frac{3\pi}{S^3}\,\hat c_0-\frac{4C}{S^4}\,\hat c_0'
 +\frac4{\pi S^3}\,\hat c_0''\Bigr],
 \label{gff:mgz:simplified}
\end{equation}
and, with \eqref{gff:gdr:Q}, equivalently \eqref{gff:gdr:Z}. Indeed, at $b=\frac d2$,
\[
 u_0'=-\frac{6\pi^3C}{S^3},\qquad u_0''=\frac{3\pi^4(1+2C^2)}{S^4},\qquad
 v_0'=-\frac{6C}{S^4},\qquad v_0''=\frac{3\pi(1+3C^2)}{S^5}.
\]
Insert \eqref{gff:mgz:Rprime} into
$9Z_d=W_3''=u_0''\mathfrak t_3+2u_0'\mathfrak t_3'+u_0\mathfrak t_3''+v_0''\hat c_0
+2v_0'\hat c_0'+v_0\hat c_0''$. The coefficient of $\mathfrak t_3$ becomes
$u_0''+2u_0'\pi C/S=3\pi^4(1-2C^2)/S^4=-3\pi^4\cos(\pi d/2)/S^4$; that of $\hat c_0$
becomes $v_0''+2u_0'C/(\pi^2S^2)=3\pi(1-C^2)/S^5=3\pi/S^3$; that of $\hat c_0'$ becomes
$2v_0'-4u_0'/(3\pi^3S)=-4C/S^4$. This is \eqref{gff:mgz:simplified}. Replacing
$\mathfrak t_3$ by $\frac{S^2}{6\pi^2}(3Q_d-\frac4{\pi S^3}\hat c_0)$ from
\eqref{gff:gdr:Q} turns the $\mathfrak t_3$ term into
$-\frac{\pi^2\cos(\pi d/2)}{6S^2}Q_d+\frac{2\pi\cos(\pi d/2)}{9S^5}\hat c_0$, and
$3S^2+2\cos2\pi q=1+C^2$ gives \eqref{gff:gdr:Z}. $\square$

The same computation shows that the full jet formula \eqref{gff:mgz:Zjets} and
\eqref{gff:mgz:simplified} differ exactly by
$-\frac{2\pi C}{9S}\,\partial_bW_3(d;\frac d2)$ (the review \path{LRI_C5_general_d/review/Z-rigor/} confirmed this
symbolically). So \eqref{gff:gdr:Z} depends on the input only through this term, which
is absent at $d=2$ where $C=0$. In \eqref{gff:mgz:simplified} the second derivatives
enter only through $\frac{2\pi^2}{3S^2}\mathfrak t_3''+\frac4{9\pi S^3}\hat c_0''$,
whose second-derivative part is $\frac19(u\mathfrak t_3''+v\widehat M_3'')$, exactly as
in \eqref{gff:mgz:Zjets}. The data needed in \eqref{gff:mgz:simplified} are $\mathfrak t_3$,
$\mathfrak t_3''$ and the jets of the zero mode, which is computed from the two nodes
by \eqref{gff:mgz:zeromode}; so the complete set of inputs remains
$\{\mathfrak t_3^{(k)},\widehat M_3^{(k)}\}_{k\le2}$, that is, at most 15
special-function numbers (two ${}_7F_6(1)$ values and three two-variable Horn values,
each with its first and second $b$-derivatives), subject to the two exact relations
$W_3(d;\frac d2)=3Q_d$ and \eqref{gff:mgz:first}.

\emph{The count.} The 15 numbers enter only through the six combinations
$\mathfrak t_3^{(k)},\widehat M_3^{(k)}$, $k\le2$, and the second derivatives only
through the single combination $\frac19(u\mathfrak t_3''+v\widehat M_3'')$. After the
two exact relations, $Z_d$ is an elementary multiple of $Q_d$ plus an elementary
combination of three further numbers, for instance $\mathfrak t_3$ (or $\hat c_0$),
$\hat c_0'$ and $\frac19(u\mathfrak t_3''+v\widehat M_3'')$ at $b=\frac d2$; this is
the count of Section~\ref{gff:gdr:C5}. The choice works at every $d$, because the
coefficients $6\pi^2/S^2$ and $4/(\pi S^3)$ of $\mathfrak t_3$ and $\hat c_0$ in
\eqref{gff:gdr:Q}, and $3\pi^3S^2$ and $2S$ of $\mathfrak t_3'$ and $\hat c_0'$ in
\eqref{gff:mgz:firstzm}, never vanish for $0<d<4$. At $d=1$ ($\hat c_0\equiv0$) and at
$d=2$ ($\widehat M_3\equiv0$) only the second-derivative combination remains, and
Section~\ref{gff:gdr:special} evaluates it in closed form. All functions used here belong
to the classes of \eqref{gff:fp:wheelBlocks} and \eqref{gff:fp:wheelArrays}, with the rim
parameter $\sigma$ and the spoke parameter $g$ separated.

\subsection{The pole node alone: residue series of one lower Laurent order}\label{gff:gdr:taylor}
At $t=0$ the left and middle half-lattices carry $[\sin\pi t]^r$, so $\Xi_r^{(k)}(0;b)$ for $k\le r-1$ comes from the
right half-lattice alone. With $r-1$ modes, $k\le r-2$ suffices. Let $\lambda_m(\ell,n;b)$ be \eqref{gff:gs:lambda} with
the rim and spoke separated,
\[
 \lambda_m=\psi_{m-1}(\ell+n+\sigma)-\psi_{m-1}(\ell+n+2q)-(-1)^m\psi_{m-1}(g-n)
 +(1+(-1)^m)\psi_{m-1}(1)-\psi_{m-1}(n+1),
\]
so that $N_bg_\ell(a_\ell+n+u;b)=-\frac{T_{\ell n}(b)}u\exp\sum_m\lambda_m\frac{u^m}{m!}$ as in \eqref{gff:gs:localGamma}.
Then, with $\kappa_b=\pi\cot\frac{\pi b}2$,
\begin{equation}\boxed{
 W_3(d;b)=\sum_{\ell,n\ge0}d_\ell T_{\ell n}^3\bigl[2\pi^2-6\kappa_b\lambda_1\bigr],\qquad
 W_4(d;b)=\sum_{\ell,n\ge0}d_\ell T_{\ell n}^4\Bigl[4\kappa_b\bigl(4\lambda_1^2+\lambda_2-\tfrac{\pi^2}3\bigr)-8\pi^2\lambda_1\Bigr].}
 \label{gff:gdr:taylorW}
\end{equation}
At $b=d/2$, $\kappa_b=D$ and the $\lambda_m$ are those of \eqref{gff:gs:lambda}:
\begin{equation}
 3Q_d=\sum d_\ell T_{\ell n}^3\bigl(2\pi^2-6D\lambda_1\bigr),\qquad
 P_4(d)=\sum d_\ell T_{\ell n}^4\Bigl[4D\bigl(4\lambda_1^2+\lambda_2-\tfrac{\pi^2}3\bigr)-8\pi^2\lambda_1\Bigr],
 \label{gff:gdr:taylorQP}
\end{equation}
to be compared with $F_3$ (quadratic, $\lambda_2$) and $F_4$ (cubic, $\lambda_3$) of \eqref{gff:gs:QPseries}.
Differentiating \eqref{gff:gdr:taylorW} in $b$ (termwise, as for \eqref{gff:gs:Lseries}; $\partial_b\lambda_m=\eta_m$,
$\partial_b\eta_m=\theta_m$, $\partial_b\log T_{\ell n}=\mu_1$, $\partial_b\mu_1=\mu_2$ of \eqref{gff:gs:etaTheta};
$\kappa_b'=-\frac{\pi^2}{2S^2}$, $\kappa_b''=\frac{\pi^3C}{2S^3}$):
\begin{equation}\begin{aligned}
 \LL_d&=\sum d_\ell T_{\ell n}^4\bigl[4\mu_1\widehat F_4+\widehat F_4'\bigr],&
 9Z_d&=\sum d_\ell T_{\ell n}^3\bigl[(9\mu_1^2+3\mu_2)\widehat F_3+6\mu_1\widehat F_3'+\widehat F_3''\bigr],\\
 \widehat F_4&=4D\bigl(4\lambda_1^2+\lambda_2-\tfrac{\pi^2}3\bigr)-8\pi^2\lambda_1,&
 \widehat F_4'&=-\tfrac{2\pi^2}{S^2}\bigl(4\lambda_1^2+\lambda_2-\tfrac{\pi^2}3\bigr)+4D(8\lambda_1\eta_1+\eta_2)-8\pi^2\eta_1,\\
 \widehat F_3&=2\pi^2-6D\lambda_1,&
 \widehat F_3'&=\tfrac{3\pi^2}{S^2}\lambda_1-6D\eta_1,\qquad
 \widehat F_3''=-\tfrac{3\pi^3C}{S^3}\lambda_1+\tfrac{6\pi^2}{S^2}\eta_1-6D\theta_1 .
\end{aligned}\label{gff:gdr:taylorLZ}\end{equation}
Here $\widehat F_3$, $\widehat F_4$ are the weights of \eqref{gff:gdr:taylorQP}, the analogues of $F_3$, $F_4$ of
\eqref{gff:gs:Fpolys}, and primes are their $b$-derivatives.
These replace \eqref{gff:gs:Lseries} and \eqref{gff:gs:Zseries}: the polygamma weights drop by one order ($\lambda_3,\eta_3$
no longer occur in $\LL_d$, and $\lambda_2,\eta_2,\theta_2$ no longer occur in $Z_d$). In this form $P_4,\LL_d$ need
one Horn family and its parameter derivatives of order $\le2$ (shift) and $\le1$ (rim); $Q_d,Z_d$ need one family
with orders $\le1$ and $\le2$. Status: derived (exact algebra, \path{LRI_C5_general_d/specialize/sym_taylor.py}); \eqref{gff:gdr:taylorW}
is confirmed numerically, and \eqref{gff:gdr:taylorLZ} through $b$-derivatives of \eqref{gff:gdr:taylorW}.

\subsection{Specialization to \texorpdfstring{$d=1,2,3$}{d=1,2,3}}\label{gff:gdr:special}
\paragraph{The one-variable function behind $d=1$ and $d=2$.}
Let $\Omega_r(t;x)=\sin^r\pi(t+x)\sum_{n\in\mathbb Z}[\Gamma(n+t+x)/\Gamma(n+t+1)]^r$. At $x=\frac12$ this is the
$\Omega_r(t)$ of \eqref{gff:fp:OmegaWheel}, and $\Omega_4(t;x)$ is the $\Omega(t;x)$ of \eqref{gff:gc:L2modes}. For every $b$
in the strip,
\begin{equation}
 d=1:\ \Xi_r(t;b)=(-1)^r2^{-rb}\,\Omega_r(2t;b),\qquad
 d=2:\ \Xi_r(t;b)=(-1)^r\,\Omega_r(t;\tfrac b2)^2 .
 \label{gff:gdr:d12lattice}
\end{equation}
\emph{Proof at $d=1$.} Here $d_0=d_1=1$, $d_\ell=0$ for $\ell\geq2$, $2q=\frac12$ and $1-g=\frac12+\sigma$. In the
reflected form \eqref{gff:gdr:reflected}, valid for every $n\in\mathbb Z$, Legendre's duplication formula
$\Gamma(w)\Gamma(w+\frac12)=2^{1-2w}\sqrt\pi\,\Gamma(2w)$ \cite{gff:DLMF5} turns the $\ell=0$ and $\ell=1$ Gamma ratios at
$z=n+t$ into $2^{1-2\sigma}\Gamma(2z+b)/\Gamma(2z+1)$ and $2^{1-2\sigma}\Gamma(2z+1+b)/\Gamma(2z+2)$. These are the
terms $m=2n$ and $m=2n+1$ of the bilateral sum of $\Omega_r(2t;b)$, so the two rows merge into one lattice. The
prefactor is $D_b(t)\sin\pi(g-t)/\sin\pi t=\sin\pi(t+\sigma)\cos\pi(t+\sigma)=\frac12\sin\pi(2t+b)$. The sign of
\eqref{gff:gdr:reflected} gives $(-1)^r$, and the powers of two combine to $2^{-r}2^{r(1-2\sigma)}=2^{-rb}$.

\emph{Proof at $d=2$.} Here $d_0=1$, $d_\ell=2$ for $\ell\geq1$ (by \eqref{gff:gdr:Dgen}), $2q=1$ and $g=1-\sigma$. With
$\phi(z)=\Gamma(z+\sigma)/\Gamma(z+1)$ (local to this proof; not the $\varphi$ of \eqref{gff:mgl:phis}) the factorized form in \eqref{gff:gdr:reflected} and Gamma reflection give
$f(\ell+\sigma+n+t)=\phi(\ell+n+t)$ and $f(-n-t)=-[\sin\pi(\sigma+t)/\sin\pi t]\,\phi(n+t)$, so
$D_b(t)\,g_\ell(a_\ell+n+t)=-\sin^2\pi(t+\sigma)\,\phi(n+t)\phi(\ell+n+t)$. With $\chi_m=\phi(m+t)^r$ the
multiplicities turn $\sum_{n\in\mathbb Z}\sum_{\ell\geq0}d_\ell\chi_n\chi_{n+\ell}$ into the symmetrization
$(\sum_{m\in\mathbb Z}\chi_m)^2$. All sums converge absolutely, since $|\chi_m|\asymp|m|^{-r\operatorname{Re}g}$ and
$r\operatorname{Re}g>1$ is $\operatorname{Re}b<b_{\max}(r,2)$. This is the second identity. Both identities are
first obtained off the pole classes and then hold everywhere, since both sides are entire. $\square$

Since $\Omega_r(t+1;x)=(-1)^r\Omega_r(t;x)$ and $\Omega_r$ grows like $e^{r\pi|\operatorname{Im}t|}$ times a negative
power, the frequencies of $\Omega_r$ lie in $\mathbb Z+\frac r2$ and have modulus at most $\frac r2-1$. So at $d=1,2$
\eqref{gff:gdr:d12lattice} gives \eqref{gff:gdr:topless} directly, and at $d=1$ more: only the modes $k\equiv r\pmod2$ occur,
so $c_{4,1}\equiv0$ and $c_{3,0}\equiv0$.

\paragraph{Two amplitudes and their first-derivative relations.}
By \eqref{gff:gc:L2modes}, $\mathcal N_x^2\Omega_4(t;x)=\mathcal A(x)+\mathcal B(x)\cos2\pi(t+\frac x2)$ with
$\mathcal N_x=\Gamma(1-x)/\Gamma(x)$, the $N_x$ of \eqref{gff:gc:L2rim}, and Section~\ref{gff:gc:L2proof} proves, at the single
point $x=\frac12$,
\begin{equation}
 \mathcal A'=2\pi\mathcal B-2\log2\,\mathcal A,\qquad \mathcal B'=\tfrac{3\pi}2\mathcal A-2\log2\,\mathcal B,\qquad
 \mathcal A(\tfrac12)=\frac{8N_{\rm e}}\pi,\quad \mathcal B(\tfrac12)=\frac{8M_{\rm e}}\pi,
 \label{gff:gdr:landen}
\end{equation}
with $M_{\rm e}=(B_-+B_+)/64$, $N_{\rm e}=(B_--B_+)/32$ of \eqref{gff:fp:MN} and $x$-derivatives. These are
first-derivative relations at $x=\frac12$ (from \eqref{gff:gc:L2taylor} and the elliptic and Landen identities
\eqref{gff:gc:L2singles}), not a differential system in $x$; equivalently $(4^x\mathcal A)'=2\pi\,4^x\mathcal B$ and
$(4^x\mathcal B)'=\frac{3\pi}2\,4^x\mathcal A$ there. They are what makes the rim derivative collapse at $d=1$ and
$d=2$: both dimensions sample the \emph{same} $\Omega_4$ at the \emph{same} point $x=\frac12$ ($b=\frac12$ at $d=1$,
$b/2=\frac12$ at $d=2$).

\paragraph{$d=1$.}
By \eqref{gff:gdr:d12lattice} with $x=b$: $c_{4,0}=2^{-4b}\mathcal A(x)/\mathcal N_x^2$, $c_{4,1}\equiv0$,
$c_{4,2}=2^{-4b}\mathcal B(x)/\mathcal N_x^2$, so $\mathcal X_0\equiv\mathcal X_{1/2}$
and only two samples are independent. At $b=\frac12$ ($N_b=\sqrt\pi$, $\varsigma_b=s_b=2^{-1/2}$,
$\partial_b\log N_b=-\psi_0(\frac14)$), \eqref{gff:gdr:landen} gives
\begin{equation}\begin{aligned}
 P_4(1)&=32\pi\mathcal A+40\pi\mathcal B,\\
 \LL_1&=-2(\pi+\log2)\,P_4(1)+52\pi^2\mathcal A+64\pi^2\mathcal B
 =-(96\pi+512\log2)N_{\rm e}-(128\pi+640\log2)M_{\rm e},
\end{aligned}\label{gff:gdr:L1}\end{equation}
which is \eqref{gff:lp:eq:MN}, that is \eqref{gff:eq:KLpsiReflected}:
$\LL_1=-\frac{3\pA-\pD}6P_4-\frac{2\pD}3B_-$ with $\pA=\pi+4\log2$, $\pD=\pi$ and
$B_-=8\pi(c_{4,0}+2c_{4,2})$. The frozen derivative is elementary in the two amplitudes,
$\mathfrak D_L=\sqrt2\,(5\mathcal X_0'-\mathcal X_T')=\frac{\sqrt2\pi^2}4\bigl[(\frac{31\pi}2-8\log2)\mathcal A
+(19\pi-10\log2)\mathcal B\bigr]$, and the single Horn derivative is
$\mathfrak h_x'=4\sum_kc_k(\frac12)^4[\psi_0(k+\frac12)-\psi_0(\frac12)]=4\mathcal S_1$ of \eqref{gff:gc:L2taylor}, with
$c_k(x)$ of \eqref{gff:gc:L2rim} and $\mathfrak h_x(1;b)={}_4F_3(b,b,b,b;1,1,1;1)$.

\paragraph{$d=2$.}
By \eqref{gff:gdr:d12lattice} with $x=b/2$: $c_{4,0}=(\mathcal A^2+\frac12\mathcal B^2)/\mathcal N_x^4$,
$c_{4,1}=2\mathcal A\mathcal B/\mathcal N_x^4$, $c_{4,2}=\mathcal B^2/(2\mathcal N_x^4)$, that is
$\mathcal X_T=\mathcal A^2$, $\mathcal X_0=(\mathcal A+\mathcal B)^2$, $\mathcal X_{1/2}=(\mathcal A-\mathcal B)^2$
at $b=1$. Here $C=0$, so $D=0$ and $\varsigma_{b}$ passes through zero: the weight of $\mathcal X_T'$ in
\eqref{gff:gdr:frozen} vanishes, and $\partial_b\log(N_b/\mathcal N_x)=0$. Then
\[
 \LL_2=-(6\mathcal A^2+8\mathcal B^2)+\frac2\pi(\mathcal X_0'-\mathcal X_{1/2}')
 =-(6\mathcal A^2+8\mathcal B^2)+\frac4\pi\bigl[\mathcal A'\mathcal B+\mathcal A\mathcal B'\bigr]
 =-\frac{16\log2}\pi\mathcal A\mathcal B=-2\log2\,P_4(2),
\]
with $x$-derivatives on $\mathcal A,\mathcal B$ (so $\partial_b=\frac12\partial_x$), which is \eqref{gff:gc:L2residue}
and \eqref{gff:gc:L2closed}. The motion of the trigonometric weights ($\partial_b\varsigma_b=-\frac\pi2$) produces
$-(6\mathcal A^2+8\mathcal B^2)$, which the rotation terms $2\pi\mathcal B$, $\frac{3\pi}2\mathcal A$ of
\eqref{gff:gdr:landen} cancel exactly; only the normalization $4^x$ survives. At $d=1$ the same relations leave
$-\frac{2\pD}3B_-$. If $\pA$ and $\pD$ in \eqref{gff:eq:KLpsiReflected} are replaced by their general-$d$ values $A$ of
\eqref{gff:gs:psiBasis} and $D$ of \eqref{gff:gdr:D}, the $B_-$ term vanishes at $d=2$ and the formula becomes exactly
$\LL_2=-\frac A2P_4(2)$ with $A=4\log2$, which is \eqref{gff:gc:L2closed}. At $d=3$ no analogue of $B_-$ with elementary
coefficients in the samples was found (an integer-relation search by lattice reduction at 155 digits, repeated at 125 digits, in \path{LRI_C5_d3}, found no relation of $\LL_3$ with the $P_4(3)$ blocks; a numerical statement). In
\eqref{gff:gdr:taylorQP} the same $D=0$ leaves $3Q_2=2\pi^2\sum d_\ell T_{\ell n}^3$ (one unweighted Horn value, the
complex Selberg value of \eqref{gff:gs:W3d2}) and $P_4(2)=-8\pi^2\sum d_\ell T_{\ell n}^4\lambda_1$.

\paragraph{The half-shift identities at $d=3$ for a general rim exponent.}\label{gff:gdr:halfShift}
The $d=3$ forms of $\LL_d$ and $Z_d$ below need two identities at rim exponents $b\neq\frac32$:
$c_{4,1}(b)=2\mathfrak a_4(b)\mathfrak b_4(b)$, and a Gamma form of $c_{3,0}(b)$. Both follow from the
half-shift proof of Section~\ref{gff:d3:determinantProof}, repeated at a general $b$. The derivation is
this project's own work: a Claude Code session and its subagents made it in the package
\path{LRI_C5_d3/} (2 October 2026), where it was recorded as a sketch that follows that proof step by
step, and three separate AI sessions per statement tried to refute the statements, re-deriving steps by
hand or with codes of their own, without finding an error. The derivation is written out here; as
printed it has not been reviewed. At $d=3$ one has $q=\frac34$, $\sigma=\frac b2$ and $g=\frac32-\sigma=\frac{3-b}2$. In this
paragraph $x$, $y$, $u$ and $v$ are local, as in Section~\ref{gff:d3:determinantProof}, and
\[
 h_b(z)=\frac{\Gamma(z+\frac{1-g}2)}{\Gamma(z+\frac{1+g}2)},\qquad u=t+\frac{1-g}2 .
\]
At $b=\frac32$, $\frac{1-g}2=\frac18$, $h_b$ is the function $h$ of \eqref{gff:td:Udefinition} and $u=t+\frac18$.

\emph{Lattice.} Put $k=\ell+n$, $j=n$, $x=k+u+\frac12$ and $y=j+u$. The first line of \eqref{gff:gdr:reflected}
gives
\begin{equation}
 g_\ell(a_\ell+n+t;b)=-\frac{\sin\pi(g-t)}{\sin\pi t}\,h_b(x)\,h_b(y),\qquad d_\ell=2\ell+1=2(x-y),
 \label{gff:gdr:halfShiftKernel}
\end{equation}
which is \eqref{gff:td:latticeKernel} at $b=\frac32$. Since $g=\frac32-\sigma$, the prefactor of $\Xi_r$ in
\eqref{gff:gdr:periodizationDef} combines as $D_b(t)\cdot(-\sin\pi(g-t)/\sin\pi t)=\frac12\sin2\pi(t+\sigma)
=-\sin\pi(u+\frac{1-g}2)\sin\pi(u+\frac12+\frac{1-g}2)$, so that
\[
 \Xi_r(t;b)=(-1)^r\Bigl[\sin\pi\bigl(u+\tfrac{1-g}2\bigr)\sin\pi\bigl(u+\tfrac12+\tfrac{1-g}2\bigr)\Bigr]^r
 \sum_{k\ge j}2(x-y)\,h_b(x)^r\,h_b(y)^r ,
\]
summed over all integers $k\ge j$. At $b=\frac32$ the bracket is $\sin\pi(t+\frac14)\sin\pi(t+\frac34)=\frac12\cos2\pi t$,
the factor used in Section~\ref{gff:d3:determinantProof}.

\emph{Half shift.} For $m=0,1$ define the regularized bilateral moments
\[
 \Psi_{r,m}(v;b)=\sin^r\pi\bigl(v+\tfrac{1-g}2\bigr)\sum_{n\in\mathbb Z}(n+v)^m\,h_b(n+v)^r .
\]
By the reflection $h_b(-z)=h_b(z)\sin\pi(z+\frac{1-g}2)/\sin\pi(\frac{1-g}2-z)$ they are the functions
\eqref{gff:td:Psidefinition} with $\omega=+1$ and $\frac18,\frac78$ replaced by $\frac{1-g}2,\frac{1+g}2$; at $b=\frac32$
they are $\Psi^{(+)}_{r,m}$. Their series converge absolutely for $r\operatorname{Re}g>m+1$, since
$|h_b(z)|$ is comparable to $|z|^{-\operatorname{Re}g}$. The shift $t\mapsto t+\frac12$ moves $u$ by $\frac12$,
exchanges the lattices $\mathbb Z+u+\frac12$ of $x$ and $\mathbb Z+u$ of $y$, and changes the sign of the
trigonometric bracket. If $r\operatorname{Re}g>2$, all sums involved converge absolutely, and the reindexing
of Section~\ref{gff:d3:determinantProof} gives \eqref{gff:td:determinants} at general $b$:
\[
 \Xi_r(t;b)-(-1)^r\,\Xi_r(t+\tfrac12;b)=2(-1)^r\bigl[\Psi_{r,1}(u+\tfrac12;b)\Psi_{r,0}(u;b)
 -\Psi_{r,0}(u+\tfrac12;b)\Psi_{r,1}(u;b)\bigr].
\]

\emph{Finite modes.} As in \eqref{gff:td:momentGrowth}, $\Psi_{r,m}(v+1;b)=(-1)^r\Psi_{r,m}(v;b)$,
$\Psi_{r,m}(-v;b)=(-1)^m\Psi_{r,m}(v;b)$, and $|\Psi_{r,m}(v;b)|=O(e^{r\pi|\operatorname{Im}v|}
|\operatorname{Im}v|^{m+1-r\operatorname{Re}g})$ as $|\operatorname{Im}v|\to\infty$, now with the exponent
$r\operatorname{Re}g$ in place of $\frac{3r}4$. Contour translation excludes the boundary frequencies when
$m+1<r\operatorname{Re}g$. For $m\le1$ this is $\operatorname{Re}g>\frac12$, that is $\operatorname{Re}b<2=b_{\max}(4,3)$,
for $r=4$, and $\operatorname{Re}g>\frac23$, that is $\operatorname{Re}b<\frac53=b_{\max}(3,3)$, for $r=3$,
by \eqref{gff:gdr:bmax}; these are also the conditions of the half shift. Reflection then leaves the modes of
\eqref{gff:td:momentModes},
\begin{equation}
 \Psi_{4,0}=\mathfrak a_4(b)+\mathfrak c_4(b)\cos2\pi v,\quad \Psi_{4,1}=\mathfrak b_4(b)\sin2\pi v,\quad
 \Psi_{3,0}=\mathfrak a_3(b)\cos\pi v,\quad \Psi_{3,1}=\mathfrak b_3(b)\sin\pi v,
 \label{gff:gdr:halfShiftModes}
\end{equation}
which define the amplitudes at general $b$; at $b=\frac32$ they are those of \eqref{gff:td:momentModes}.

\emph{Fourier coefficients.} By \eqref{gff:gdr:finiteFourier}, $\Xi_r(t;b)=\sum_kc_{r,k}(\frac34,b)\cos2\pi k(u+\frac14)$,
since $t+\frac b4=u+\frac14$. Inserting \eqref{gff:gdr:halfShiftModes} in the half-shift identity gives
$\Xi_4(t;b)-\Xi_4(t+\frac12;b)=-4\mathfrak a_4\mathfrak b_4\sin2\pi u$ and
$\Xi_3(t;b)+\Xi_3(t+\frac12;b)=-2\mathfrak a_3\mathfrak b_3$, while the left sides are
$-2c_{4,1}\sin2\pi u+2c_{4,3}\sin6\pi u$ and $2c_{3,0}-2c_{3,2}\cos4\pi u$. Hence, for $0<\operatorname{Re}b<2$ and
$0<\operatorname{Re}b<\frac53$ respectively,
\begin{equation}
 c_{4,1}(\tfrac34,b)=2\mathfrak a_4(b)\mathfrak b_4(b),\qquad c_{3,0}(\tfrac34,b)=-\mathfrak a_3(b)\mathfrak b_3(b),
 \label{gff:gdr:halfShiftC}
\end{equation}
and $c_{4,3}=c_{3,2}=0$, as \eqref{gff:gdr:topless} also gives. At $b=\frac32$ this is \eqref{gff:td:projectedAnswer}.

\emph{Gamma evaluations.} Let $1<b<\frac53$ for $r=3$ and $1<b<2$ for $r=4$, so that $\frac12<g<1$. At the node
$v=\frac{1-g}2$ the second term of \eqref{gff:td:Psidefinition} carries the factor $\sin^r\pi(\frac{1-g}2-v)=0$, and
$h_b(n+\frac{1-g}2)=\Gamma(n+1-g)/\Gamma(n+1)$. With
$\mathfrak U_{r,m}=\sum_{n\ge0}(n+\frac{1-g}2)^mh_b(n+\frac{1-g}2)^r$, which is $\mathfrak U^{(+)}_{r,m}(\frac{1-g}2)$ of
\eqref{gff:td:Udefinition} at general $b$,
\[
 \mathfrak a_3(b)=\frac{\sin^3(\pi g)}{\sin(\pi g/2)}\,\mathfrak U_{3,0},\qquad
 \mathfrak b_3(b)=\frac{\sin^3(\pi g)}{\cos(\pi g/2)}\,\mathfrak U_{3,1},\qquad
 \mathfrak b_4(b)=\sin^3(\pi g)\,\mathfrak U_{4,1}.
\]
Dixon's sum \cite[eq.~16.4.4]{gff:DLMF16} with $a=b=c=1-g$, and the Rogers--Dougall sum \cite[eq.~16.4.9]{gff:DLMF16}
with $(a,b,c,d)=(1-g,1-g,1-g,1-\frac g2)$, after inserting the canceled pair $1-\frac g2$, and with
$(a,b,c,d)=(1-g,1-g,1-g,1-g)$ (in the notation of \cite{gff:DLMF16}; these $b,c,d$ are not the rim exponent or the
dimension) give
\begin{gather*}
 \mathfrak U_{3,0}=\frac{\Gamma(1-g)^3\,\Gamma(\frac{3-g}2)\,\Gamma(\frac{3g-1}2)}
 {\Gamma(2-g)\,\Gamma(\frac{1+g}2)^2\,\Gamma(g)},\qquad
 \mathfrak U_{3,1}=\frac{1-g}2\,\frac{\Gamma(1-g)^3\,\Gamma(\frac{2-g}2)\,\Gamma(\frac{3g-2}2)}
 {\Gamma(2-g)\,\Gamma(g)\,\Gamma(\frac g2)^2},\\
 \mathfrak U_{4,1}=\frac{1-g}2\,\frac{\Gamma(1-g)^4\,\Gamma(2g-1)}{\Gamma(2-g)\,\Gamma(g)^3}.
\end{gather*}
Their convergence conditions hold: $g>\frac13$ for Dixon's sum, $g>\frac23$ and $g>\frac12$ for the two
Rogers--Dougall sums. Duplication and reflection \cite{gff:DLMF5} then give
\begin{equation}
 \mathfrak a_3(b)\mathfrak b_3(b)=\frac{\pi^5\,\Gamma(3g-2)}{\Gamma(g)^9},\qquad
 \mathfrak b_4(b)=\frac{\pi^3\,\Gamma(2g-1)}{2\,\Gamma(g)^6}.
 \label{gff:gdr:halfShiftGamma}
\end{equation}
Both sides are analytic in $b$ on the strips of \eqref{gff:gdr:halfShiftC}, where the poles of $\Gamma(3g-2)$ and
$\Gamma(2g-1)$ ($\operatorname{Re}b\ge\frac53$ and $\operatorname{Re}b\ge2$) do not occur, so \eqref{gff:gdr:halfShiftGamma}
holds there by analytic continuation. At $b=\frac32$ it is \eqref{gff:td:gammaAmplitudes}. Consequently
$\hat c_0(3;b)=N_b^3c_{3,0}=-N_b^3\pi^5\Gamma(3g-2)/\Gamma(g)^9$, and, with $\nu_b$ of the next paragraph,
$\nu_b^{1/2}\mathfrak b_4=\pi^3\Gamma(2g-1)/[2\Gamma(\sigma)^2\Gamma(g)^4]$, whose logarithmic $b$-derivative at
$b=\frac32$ is $\psi_0(\frac34)-\psi_0(\frac12)=\frac\pi2-\log2$.

\emph{Checks} (records in \path{LRI_C5_d3/}; numerical, not interval bounds). A code of one of the verifiers
(\path{wf_vL3num/}, records \path{modes_b*.json}), which shares no summation routine with the evaluation code
\path{wf_L3/gen3.py}, fits $\Psi_{4,0}$, $\Psi_{4,1}$ and $\Xi_4$ at $90$ digits at $b=0.4$, $0.8$, $1.2$, $1.5$, $1.6$,
$1.8$ and $1.9$: the modes absent from \eqref{gff:gdr:halfShiftModes}, and $c_{4,3}$, are below $3\cdot10^{-69}$,
$c_{4,1}-2\mathfrak a_4\mathfrak b_4$ is below $2\cdot10^{-68}$, and at the six points other than $b=0.4$ the
function $W_4(3;b)$ assembled from the fitted coefficients agrees with an independent spectral code to
$5.1\cdot10^{-50}$ or better. The product $\mathfrak a_3\mathfrak b_3$ of \eqref{gff:gdr:halfShiftGamma} holds at
$b=1.1$, $1.25$, $1.4$, $1.5$, $1.6$ and $1.3+0.2\ii$ to $2\cdot10^{-60}$ at $60$ digits, and, with a second, literal
hypergeometric evaluator of the moments (not run at $b=1.6$), to $6\cdot10^{-29}$ at $30$ digits
(\path{wf_Z3_gamma_verify/check_amp.py}, logs \path{check_dps30.log}, \path{check_dps60.log}). The checks of the
separate review at $b=0.8$, outside the range of the Gamma evaluation, and at other points are quoted after
\eqref{gff:mgz:W3d3}, and the $d=3$ dictionary below is checked against the $160$-digit values in
Section~\ref{gff:gdr:checks}.

\paragraph{$d=3$.}
By \eqref{gff:gdr:halfShiftC}, $c_{4,1}(b)=2\mathfrak a_4(b)\mathfrak b_4(b)$ for $0<b<2$, with the moment amplitudes of
\eqref{gff:gdr:halfShiftModes}, and \eqref{gff:gdr:halfShiftGamma} gives $\mathfrak b_4(b)$ in Gamma form, so that
$\partial_b[\nu_b^{1/2}\mathfrak b_4]=(\frac\pi2-\log2)\mathfrak b_4$ at $b=\frac32$, where
$\nu_b=[\Gamma(\frac32-\frac b2)/\Gamma(\frac b2)]^4=N_b^4/\Gamma(\frac32)^4$ equals one there. Let $X_0=\Xi_4(0;\frac32)$; this is the $X_0$ of \eqref{gff:td:P4data}, not the sample $\mathcal X_0$ of \eqref{gff:gdr:samples4}. The dictionary to
\eqref{gff:gdr:samples4} is
\begin{gather*}
 \mathcal X_T=N_b^4X_0,\qquad \mathcal X_0-\mathcal X_{1/2}=4N_b^4\mathfrak a_4\mathfrak b_4,\qquad
 \mathfrak d_H=\partial_b\bigl[\nu_b(4c_{4,0}+5c_{4,2})\bigr]=\frac{16}{\pi^2}\partial_b\bigl[N_b^4(4c_{4,0}+5c_{4,2})\bigr],\\
 \mathfrak d_a=\partial_b\bigl[\nu_b^{1/2}\mathfrak a_4\bigr]
 =\frac{8}{\pi^2}\frac{\partial_b(N_b^4c_{4,1})}{\mathfrak b_4}-\Bigl(\frac\pi2-\log2\Bigr)\mathfrak a_4,
\end{gather*}
where the first equalities for $\mathfrak d_H$ and $\mathfrak d_a$ are their definitions. Inserting the dictionary into
\eqref{gff:gdr:L} gives
$\LL_3=4\pi P_4(3)-2\pi^2X_0+6\sqrt2\pi(\pi-2\log2)\mathfrak a_4\mathfrak b_4-2\pi\mathfrak d_H+12\sqrt2\pi\mathfrak b_4\mathfrak d_a$.
The derivative constants occur only in the combination $-2\pi\mathfrak d_H+12\sqrt2\pi\mathfrak b_4\mathfrak d_a$.
It has the same derivative content as $\widehat P'$, the $b$-derivative at $b=\frac32$ of
$\widehat P(b)=-2\pi\nu_b(4c_{4,0}+5c_{4,2})+6\sqrt2\pi\nu_bc_{4,1}$, in terms of which the $160$-digit
evaluation (\path{LRI_C5_d3/wf_L3/}) writes
$\LL_3=4\pi P_4(3)-2\pi^2X_0+\widehat P'$. The two differ by value terms:
\[
 \frac{2}{\pi S^7}\mathfrak D_L=\widehat P'-8\sqrt2\,N_b^4c_{4,1}+32\,N_b^4c_{4,2}\qquad(b=\tfrac32),
\]
numerically $-1332.48$ against $\widehat P'=274.55$; the identity holds to $10^{-57}$ with the 160-digit $d=3$ values stored in
\path{LRI_C5_d3/wf_L3/L3_final_values.json}. So $\mathfrak D_L$ is not literally $\widehat P'$. The split into a Horn-sector and an
ordinary-sector constant follows the rank-two structure. The collapse fails at $d=3$ because the lattice no longer
reduces to a one-variable bilateral sum: the multiplicity $2\ell+1$ gives rank two, the Dougall factor
$\mathfrak b_4$ closes but $\mathfrak a_4'$ and the Horn-sector derivative do not (lattice reduction at 155 digits,
\path{LRI_C5_d3}; numerical), and $D=-\pi\ne0$.

\paragraph{$Z_d$ at $d=1,2,3$.}
In \eqref{gff:gdr:Z} the inputs are $\mathfrak t_3(b)$ and $\hat c_0(b)$. At $d=1$, \eqref{gff:gdr:d12lattice} makes
$\Xi_3$ odd under $t\to t+\frac12$, so $\hat c_0\equiv0$; with $\cos\frac\pi2=0$ and
$\mathfrak t_3(1;b)={}_3F_2(b,b,b;1,1;1)$, which is Dixon's sum \cite{gff:DLMF16}, this gives
$Z_1=\frac{4\pi^2}3\mathfrak t_3''(\frac12)=QH$ of \eqref{gff:gs:d1Dictionary}. At $d=2$ the same symmetrization as in
the proof of \eqref{gff:gdr:d12lattice} gives $\widehat M_3\equiv0$ (the midpoint sum is the square of a bilateral sum
that is odd under $m\to-1-m$) and $\mathfrak t_3(2;b)=[{}_3F_2(\sigma,\sigma,\sigma;1,1;1)]^2$, Dixon's sum squared,
because $T_{\ell n}=\chi_{\ell+n}\chi_n$ with $\chi_m=(\sigma)_m/m!$. Hence $\hat c_0=-\pi^3s_b^3\mathfrak t_3/(1+\varsigma_b)$
and $u\mathfrak t_3=\frac12\gamma_C(b/4)^3/\gamma_C(3b/4)$, the complex Selberg value \eqref{gff:gs:W3d2}, here obtained
from Dixon's theorem; $Z_2=4\Cat Q_2$ follows from it alone (proofs in the next paragraph, \eqref{gff:mgz:R2}--\eqref{gff:mgz:Z2}). At $d=3$, $\hat c_0(b)=-N_b^3\pi^5\Gamma(3g-2)/\Gamma(g)^9$, by \eqref{gff:gdr:halfShiftC} and
\eqref{gff:gdr:halfShiftGamma}. Then $\hat c_0,\hat c_0',\hat c_0''$ are elementary, the zero-mode part of
\eqref{gff:gdr:Z} is $\sqrt2\,(4\pi+32\Cat/\pi)\,\hat c_0(\frac32)$, and
$Z_3=\frac{4\pi^2}3\mathfrak t_3''(\frac32)-(4\Cat+\frac{\pi^2}2)\frac{\Gamma(1/4)^{10}}{16\pi^{7/2}}$. The collapse
stops at $\mathfrak t_3''$, because $\mathfrak t_3(3;b)$ is a genuine two-variable Horn function.

\paragraph{Proofs of the three-wheel statements at $d=1,2,3$.}
\emph{$d=1$.} For $0<b<\frac23=b_{\max}(3,1)$,
\begin{equation}
 \mathfrak t_3(1;b)={}_3F_2\!\left(\begin{matrix}b,b,b\\1,1\end{matrix};1\right)
 =\frac{\Gamma(1+\frac b2)\,\Gamma(1-\frac{3b}2)}{\Gamma(1+b)\,\Gamma(1-\frac b2)^2\,\Gamma(1-b)},\qquad
 W_3(1;b)=\frac{6\pi^2}{s_b^2}\,\mathfrak t_3(1;b).
 \label{gff:mgz:Dixon1}
\end{equation}
\emph{Proof.} At $d=1$ only the modes $k\equiv r\pmod2$ occur (after
\eqref{gff:gdr:d12lattice}), so $c_{3,0}\equiv0$, $\hat c_0\equiv0$, and
\eqref{gff:mgz:zeromode} gives $W_3=6\pi^2\mathfrak t_3/s_b^2$. At $d=1$,
$d_0=d_1=1$ and $d_\ell=0$ for $\ell\ge2$ (the factor $(4q-1)_j$ in
\eqref{gff:fp:degeneracy} vanishes for $j\ge1$), $2q=\frac12$ and $1-g=\frac12+\sigma$.
Legendre's duplication $(x)_n(x+\frac12)_n=4^{-n}(2x)_{2n}$ \cite{gff:DLMF5} gives
$T_{0n}=(b)_{2n}/(2n)!$ and $T_{1n}=(b)_{2n+1}/(2n+1)!$. Hence
$\mathfrak t_3(1;b)=\sum_{m\ge0}[(b)_m/m!]^3$, which is Dixon's well-poised
${}_3F_2$ with $a=b=c$ \cite{gff:DLMF16}; it converges for
$\operatorname{Re}b<\frac23=b_{\max}(3,1)$, and Dixon's theorem gives the Gamma
product. $\square$ Equivalently, the vanishing zero mode is the relation
$\widehat M_3(1;b)=8\pi^3s_b^3\,\mathfrak t_3(1;b)/[\varsigma_b(1+\varsigma_b)^3]$
between the two lattices.

\emph{$Z_1$ from Dixon's theorem.} Since $\hat c_0\equiv0$ for all $b$,
$W_3(1;b)=6\pi^2\mathfrak t_3(1;b)/s_b^2$ on $0<b<\frac23$. At $b=\frac12$ one has
$s_b=\varsigma_b=2^{-1/2}$, and the logarithmic derivatives of the Gamma product
\eqref{gff:mgz:Dixon1}, with the recurrences $\psi_m(z+1)=\psi_m(z)+(-1)^mm!\,z^{-m-1}$
and $\psi_1(\frac12)=\frac{\pi^2}2$, are
\[
 (\log\mathfrak t_3)'(\tfrac12)=\psi_0(\tfrac34)-\psi_0(\tfrac14)=\pi,\qquad
 (\log\mathfrak t_3)''(\tfrac12)=\tfrac52\psi_1(\tfrac14)-\tfrac12\psi_1(\tfrac34)-\pi^2
 =\pi^2+24\Cat ,
\]
using $\psi_1(\frac14)=\pi^2+8\Cat$ and $\psi_1(\frac34)=\pi^2-8\Cat$. Further
$\mathfrak t_3(1;\frac12)=\Gamma(\frac54)\Gamma(\frac14)/[\Gamma(\frac32)\Gamma(\frac34)^2\Gamma(\frac12)]
=\Gamma(\frac14)^4/(4\pi^3)$ by \eqref{gff:mgz:Dixon1} and
$\Gamma(\frac34)=\sqrt2\,\pi/\Gamma(\frac14)$, so that \eqref{gff:gdr:Q} gives
$Q_1=4\pi^2\mathfrak t_3(1;\frac12)=\Gamma(\frac14)^4/\pi$. The first derivative of
$W_3(1;b)$ at $b=\frac12$ is
$12\pi^2\mathfrak t_3[(\log\mathfrak t_3)'-\pi]=0$, so the input
$\partial_bW_3=0$ holds here by Dixon's theorem; the second derivative is
$24\pi^4\mathfrak t_3-24\pi^3\mathfrak t_3'+12\pi^2\mathfrak t_3''=12\pi^2\mathfrak t_3''$.
Hence, with $\mathfrak t_3''/\mathfrak t_3=(\log\mathfrak t_3)''+[(\log\mathfrak t_3)']^2
=2\pi^2+24\Cat$,
\begin{equation}
 Z_1=\frac{4\pi^2}3\,\mathfrak t_3''(1;\tfrac12)=Q_1\Bigl(\frac{2\pi^2}3+8\Cat\Bigr),
 \label{gff:mgz:Z1}
\end{equation}
the value $Q_1H$ of \eqref{gff:gs:d1Dictionary}, here obtained from Dixon's theorem
without Pattern~T.

\emph{$d=2$.} For $0<\operatorname{Re}b<\frac43=b_{\max}(3,2)$, $\widehat M_3(2;b)=0$, and
$\mathfrak t_3(2;b)$ is the square of Dixon's sum at $\sigma=\frac b2$:
\begin{equation}\begin{aligned}
 \mathfrak t_3(2;b)&=\Bigl[{}_3F_2\!\left(\begin{matrix}\sigma,\sigma,\sigma\\1,1\end{matrix};1\right)\Bigr]^2
 =\biggl[\frac{\Gamma(1+\frac b4)\,\Gamma(1-\frac{3b}4)}{\Gamma(1+\frac b2)\,\Gamma(1-\frac b4)^2\,
 \Gamma(1-\frac b2)}\biggr]^2
 =\frac{s_b^2}{4\pi^2(1+2\varsigma_b)}\,\frac{\gamma_C(b/4)^3}{\gamma_C(3b/4)} ,
\end{aligned}\label{gff:mgz:R2}\end{equation}
with $\gamma_C(x)=\Gamma(x)/\Gamma(1-x)$ as in \eqref{gff:gs:W3d2}. Hence
$c_{3,0}=c_{3,1}$ and
\begin{equation}
 W_3(2;b)=u(b)\,\mathfrak t_3(2;b)=\frac12\,\frac{\gamma_C(b/4)^3}{\gamma_C(3b/4)},
 \label{gff:mgz:W3d2}
\end{equation}
which is the complex Selberg value \eqref{gff:gs:W3d2}. Note also
$\mathfrak t_3(2;b)=\mathfrak t_3(1;\frac b2)^2$.

\emph{Proof.} At $d=2$: $g=1-\sigma$, $d_0=1$, $d_\ell=2$ for $\ell\ge1$. Put
$\varrho=\frac{1+\sigma}2$ and $\varphi_m=\Gamma(m+\varrho)/\Gamma(m+2-\varrho)$,
$m\in\mathbb Z$ (both local to this proof). For $s=n+\frac{\ell+1}2$ one has
$a_\ell+s=\ell+n+\varrho$, $a_\ell-s=\varrho-1-n$ and
$\beta_\ell\pm s=a_\ell\pm s+2-2\varrho$. Gamma reflection \cite{gff:DLMF5} gives
$\Gamma(\varrho-1-n)/\Gamma(1-\varrho-n)=-\varphi_n$, because
$\sin\pi(1-\varrho-n)/\sin\pi(\varrho-1-n)=-1$. So
$g_\ell(n+\frac{\ell+1}2)=-\varphi_{\ell+n}\varphi_n$ and
\[
 \frac{\widehat M_3}{N_b^3}=-\sum_{\ell\ge0}d_\ell\sum_{n\in\mathbb Z}\varphi_{\ell+n}^3\varphi_n^3
 =-\sum_{m\ge n}(2-\delta_{mn})\,\varphi_m^3\varphi_n^3
 =-\Bigl(\sum_{m\in\mathbb Z}\varphi_m^3\Bigr)^2 .
\]
The same reflection gives $\varphi_{-1-m}=-\varphi_m$, so the bilateral sum vanishes.
All sums converge absolutely: $|\varphi_m|\asymp|m|^{\operatorname{Re}\sigma-1}$, and
$3(1-\operatorname{Re}\sigma)>1$ is $\operatorname{Re}b<\frac43$. Then
$y_b=c_{3,0}-c_{3,1}=0$ by \eqref{gff:mgz:nodeb}.

At the pole node the same rearrangement applies. At $d=2$ one has $2q=1$ and
$1-g=\sigma$, so $T_{\ell n}=\chi_{\ell+n}\chi_n$ with $\chi_m=(\sigma)_m/m!$, and
$\mathfrak t_3=\sum_{m\ge n\ge0}(2-\delta_{mn})\chi_m^3\chi_n^3=(\sum_{m\ge0}\chi_m^3)^2$.
The inner sum is Dixon's ${}_3F_2(\sigma,\sigma,\sigma;1,1;1)$, as at $d=1$ with
$b\to\sigma$ \cite{gff:DLMF16}; it converges absolutely for
$\operatorname{Re}(2-3\sigma)>0$, again $\operatorname{Re}b<\frac43$. This is the first
equality of \eqref{gff:mgz:R2}, and Dixon's theorem gives the second. For the third put
$x=\frac b4$. Gamma reflection and $\Gamma(1+x)=x\Gamma(x)$ give
\[
 \Bigl[\frac{\Gamma(1+x)\Gamma(1-3x)}{\Gamma(1+2x)\Gamma(1-x)^2\Gamma(1-2x)}\Bigr]^2
 \frac{\gamma_C(3x)}{\gamma_C(x)^3}
 =\frac{x^2\,\Gamma(3x)\Gamma(1-3x)}{[\Gamma(1+2x)\Gamma(1-2x)]^2\,\Gamma(x)\Gamma(1-x)}
 =\frac{\sin\pi x\,\sin^22\pi x}{4\pi^2\sin3\pi x},
\]
and $\sin3\pi x=\sin\pi x\,(1+2\cos2\pi x)$ turns the right side into
$s_b^2/(4\pi^2(1+2\varsigma_b))$. Finally \eqref{gff:gdr:twonode} with $\widehat M_3=0$
gives $W_3(2;b)=u(b)\mathfrak t_3(2;b)$, and
$u(b)\,s_b^2/(4\pi^2(1+2\varsigma_b))=\frac12$. $\square$

This is the square structure of the two-dimensional lattice \eqref{gff:gdr:d12lattice}
at the two nodes: at the self-mirror node the square vanishes, and at the pole node
it is the square of Dixon's sum. It derives the complex Selberg value
\eqref{gff:gs:W3d2} from \eqref{gff:gdr:twonode} and Dixon's theorem, without the complex
Selberg integral. Its agreement with \eqref{gff:gs:W3d2}, which
Section~\ref{gff:gs:Z2section} obtains from that integral, is an exact check of
\eqref{gff:gdr:twonode} and of the theorems it rests on (Sections~\ref{gff:gdr:finite},
\ref{gff:gdr:top}, \ref{gff:gdr:residue}) at $d=2$ on the range $0<b<\frac43$ where both
derivations hold.

\emph{$Z_2$ from Dixon's theorem.} Since $\widehat M_3\equiv0$, \eqref{gff:mgz:Zjets}
gives $9Z_2=(u\mathfrak t_3)''(1)$ with $u\mathfrak t_3$ of \eqref{gff:mgz:W3d2}. Its
logarithmic derivative vanishes at $b=1$ and its logarithmic second derivative there
is $\frac34[\psi_1(\frac14)-\psi_1(\frac34)]=12\Cat$, as in
Section~\ref{gff:gs:Z2section}, so $(u\mathfrak t_3)''(1)=12\Cat\,(u\mathfrak t_3)(1)$; with
$3Q_2=u(1)\mathfrak t_3(2;1)=2\pi^2\,\mathfrak t_3(1;\frac12)^2=\Gamma(\frac14)^8/(8\pi^4)$
this gives
\begin{equation}
 Z_2=4\Cat\,Q_2 ,
 \label{gff:mgz:Z2}
\end{equation}
which is \eqref{gff:gs:Z2exact}. Here it follows from Dixon's theorem alone, without the
complex Selberg integral and without the input $\partial_bW_3=0$ (which holds here
because the logarithmic derivative vanishes).

\emph{$d=3$.} At $d=3$, $\Gamma(2q)=\Gamma(\frac32)=\frac{\sqrt\pi}2$ and
$N_b^3=\frac{\pi^{3/2}}8\Gamma(g)^3/\Gamma(\sigma)^3$. With
$c_{3,0}(b)=-\mathfrak a_3(b)\mathfrak b_3(b)=-\pi^5\Gamma(3g-2)/\Gamma(g)^9$, the Gamma
form \eqref{gff:gdr:halfShiftC}, \eqref{gff:gdr:halfShiftGamma}, one has
$\frac4{\pi s_b^3}\hat c_0=-\frac{\pi^{11/2}\Gamma(3g-2)}{2s_b^3\Gamma(\sigma)^3\Gamma(g)^6}$
with $3g-2=\frac{5-3b}2$, and \eqref{gff:mgz:zeromode} becomes
\begin{equation}
 W_3(3;b)=\frac{6\pi^2\,\mathfrak t_3(3;b)}{\sin^2(\pi b/2)}
 -\frac{\pi^{11/2}\,\Gamma(\frac{5-3b}2)}{2\sin^3(\pi b/2)\,\Gamma(\frac{3-b}2)^6\,\Gamma(\frac b2)^3},
 \label{gff:mgz:W3d3}
\end{equation}
the general-$b$ formula of the $d=3$ computation (\path{LRI_C5_d3/wf_Z3/}). By
Section~\ref{gff:gdr:halfShift} the Gamma form holds on the whole strip
$0<\operatorname{Re}b<\frac53=b_{\max}(3,3)$. A run of the separate review with the block code of
\path{LRI_C5_general_d/Z/} confirms it numerically at $b=0.8$, below the range
$1<b<\frac53$ of the original Gamma evaluation in \path{LRI_C5_d3/}, at
$b=1.2$, $1.35$, $1.6$ and at the complex point $b=1.45+0.05\ii$
(\path{LRI_C5_general_d/review/Z-rigor/d3_gamma_offcenter.log}); an own code of that
review confirms it at $b=1.2$ and $b=1.45+0.05\ii$ (Section~\ref{gff:gdr:checks}).

\emph{$Z_3$.} Here $S=\frac1{\sqrt2}$, $C=-\frac1{\sqrt2}$ and
$\cos\frac{3\pi}2=0$. With the Gamma form,
$\log(-\hat c_0)=\mathrm{const}+\log\Gamma(3g-2)-3\log\Gamma(\sigma)-6\log\Gamma(g)$,
so at $b=\frac32$ ($g=\sigma=\frac34$)
\[
 (\log(-\hat c_0))'=\tfrac32\bigl[\psi_0(\tfrac34)-\psi_0(\tfrac14)\bigr]=\tfrac{3\pi}2,\qquad
 (\log(-\hat c_0))''=\tfrac94\bigl[\psi_1(\tfrac14)-\psi_1(\tfrac34)\bigr]=36\Cat,
\]
and $\hat c_0(\frac32)=-\Gamma(\frac14)^{10}/(2^{15/2}\pi^{5/2})$. In
\eqref{gff:mgz:simplified} the coefficients are $3\pi/S^3=6\sqrt2\pi$,
$-4C/S^4=8\sqrt2$ and $4/(\pi S^3)=8\sqrt2/\pi$, so the zero-mode part is
$\frac{\sqrt2}9[6\pi+12\pi+8(\frac{9\pi}4+\frac{36\Cat}\pi)]\hat c_0
=\sqrt2\,(4\pi+\frac{32\Cat}\pi)\,\hat c_0(\frac32)$ and
\begin{equation}
 Z_3=\frac{4\pi^2}3\,\mathfrak t_3''(3;\tfrac32)+\sqrt2\Bigl(4\pi+\frac{32\Cat}{\pi}\Bigr)\hat c_0(\tfrac32)
 =\frac{4\pi^2}3\,\mathfrak t_3''(3;\tfrac32)-\Bigl(4\Cat+\frac{\pi^2}2\Bigr)\frac{\Gamma(1/4)^{10}}{16\pi^{7/2}},
 \label{gff:mgz:Z3}
\end{equation}
the formula of the paragraph above. It uses the Gamma form \eqref{gff:gdr:halfShiftGamma} and the input
$\partial_bW_3=0$ (through \eqref{gff:mgz:simplified}). At $d=1,2$ the whole function
$\mathfrak t_3(d;b)$ is a Gamma product (Dixon's sum and its square); at $d=3$,
$\mathfrak t_3''(3;\frac32)$ remains as the one new constant.


\begin{center}\small
\begin{tabular}{@{}l>{\raggedright\arraybackslash}p{0.5\textwidth}>{\raggedright\arraybackslash}p{0.36\textwidth}@{}}\toprule
 & $\LL_d$ (four-wheel) & $Z_d$ (three-wheel)\\\midrule
$d=1$ & duplication to $\Omega_4(\cdot;\frac12)$; relations \eqref{gff:gdr:landen}; $B_-$ term $\propto D$ remains
 & $\hat c_0\equiv0$; $\mathfrak t_3$ Dixon\\
$d=2$ & square of $\Omega_4(\cdot;\frac12)$; relations \eqref{gff:gdr:landen}; $D=0$ leaves only $-2\log2$
 & $\widehat M_3\equiv0$; $\mathfrak t_3$ Dixon squared ($=$ Selberg)\\
$d=3$ & rank two; only $\mathfrak b_4$ (Dougall) closes; $\mathfrak D_L$ not reduced & $\hat c_0$ Gamma; $\mathfrak t_3''$ not reduced\\
generic & $\mathfrak D_L$ not reduced & $\hat c_0$, $\hat c_0'$, $\frac19(u\mathfrak t_3''+v\widehat M_3'')$ not reduced\\\bottomrule
\end{tabular}
\end{center}

\subsection{The five-loop coefficient}\label{gff:gdr:C5}
Inserting \eqref{gff:gdr:P4}, \eqref{gff:gdr:L}, \eqref{gff:gdr:Q}, \eqref{gff:gdr:Z} into \eqref{gff:gs:C5} (exact algebra; $A,R,U,V,X$
of \eqref{gff:gs:psiBasis}) gives
\begin{equation}\boxed{\begin{aligned}
 46656\,C_5(d)={}&-1536\bigl[P_5(d)+\KK_d\bigr]-1728Q_d^2
 +\Bigl[720A^2+1008R-\frac{1192\pi^2\cos(\pi d/2)}{S^2}\Bigr]Q_d\\
 &+k_0\mathcal X_0+k_{1/2}\mathcal X_{1/2}+k_T\mathcal X_T-\frac{1152}{\pi S^7}\mathfrak D_L
 +\frac{4768\pi^2}{S^2}\mathfrak t_3''\\
 &+\frac{2384}{3}\Bigl[\frac{\pi(1+C^2)}{S^5}\hat c_0-\frac{4C}{S^4}\hat c_0'+\frac{4}{\pi S^3}\hat c_0''\Bigr]\\
 &-15A^4-90A^2R+414AU-108AV-117R^2-75X,
\end{aligned}}\label{gff:gdr:C5reduced}\end{equation}
\[\begin{aligned}
 k_0&=(1+C)^2\Bigl[\frac{768A(1+2C)}{\pi S^7}+\frac{576(4+9C+8C^2)}{S^8}\Bigr],\\
 k_{1/2}&=(1-C)^2\Bigl[-\frac{768A(1-2C)}{\pi S^7}+\frac{576(4-9C+8C^2)}{S^8}\Bigr],\\
 k_T&=-\frac{1536AC}{\pi S^7}-\frac{1152(1+6C^2)}{S^8},\qquad
 Q_d=\frac{2\pi^2}{S^2}\mathfrak t_3+\frac{4}{3\pi S^3}\hat c_0 .
\end{aligned}\]
The term $384AP_4-576\LL_d$ of \eqref{gff:gs:C5} is $\sum_ik_i\mathcal X_i-\frac{1152}{\pi S^7}\mathfrak D_L$, and
$7152Z_d$ contributes the $\mathfrak t_3''$, $\hat c_0$ and $\cos(\pi d/2)Q_d$ terms. The masters $P_5(d)$ and $\KK_d$
enter \eqref{gff:gdr:C5reduced} only through $-1536[P_5(d)+\KK_d]$, as in \eqref{gff:gs:C5}. $P_5(d)$ is
\eqref{gff:fp:periodsFinite}. $\KK_d$ is defined by \eqref{gff:gs:Kdefinition}. It is evaluated by \eqref{gff:gs:Kseries}, whose
derivation uses the mixed-trace step (input (M) of Section~\ref{gff:gdk:conventions}) and the extraction of $[\tau^2]$
under the trace, by the Mellin--Barnes form \eqref{gff:gdk:KMB}, which uses none of the inputs of
Section~\ref{gff:gdk:conventions}, or by \eqref{gff:gdk:Kr3}; any of these may be substituted in this term without changing
the rest of \eqref{gff:gdr:C5reduced}. With \eqref{gff:gdk:Kr3},
\[
 -1536\,\KK_d=384AU-576c_*U-384(\mathsf r_3-\mathsf s_3),
\]
where $c_*=\psi_0(1)+\psi_0(2q)$ is that of \eqref{gff:gs:Kweights} and $\mathsf r_3,\mathsf s_3$ are the Laurent
coefficients of \eqref{gff:gdk:notation}, with $\mathsf r_3=[\tau^3]\langle t_\ell^2R_\tau\rangle_d$ (sans-serif; not the $r_3$ of \eqref{gff:gs:Kpoly} or of \eqref{gff:loc:rawAggregate}). This is exact given (M) and (J); (J) is the path form of \eqref{gff:gs:Bjet}, on
which \eqref{gff:gs:C5} itself rests. When $\mathsf r_3$ is taken from three-node samples on the jet circle, as in
\eqref{gff:gdk:Ksample}, hypothesis (H) of Section~\ref{gff:gdk:continuation} is needed as well. On the range $\frac23<d<4$
of \eqref{gff:gdr:C5reduced} the caveat for $d\le\frac47$ of Section~\ref{gff:gdk:conventions} does not arise.

\begin{center}\small
\begin{tabular}{@{}>{\raggedright\arraybackslash}p{0.13\textwidth}>{\raggedright\arraybackslash}p{0.33\textwidth}>{\raggedright\arraybackslash}p{0.48\textwidth}@{}}\toprule
master & before & with \eqref{gff:gdr:C5reduced}\\\midrule
$P_4(d)$ & \eqref{gff:fp:wheelSample}: 20 complex block values & 3 samples: 5 two-variable Horn values
 $\mathfrak h_x,\mathfrak h_1,\mathfrak h_{1/2},\mathfrak s_0,\mathfrak s_{1/2}$\\
$\LL_d$ & \eqref{gff:gc:finiteClosure}: 244 values, $\le6$ variables & the same 5 functions and their first
 $b$-derivatives, in the one combination $\mathfrak D_L$\\
$Q_d$ & \eqref{gff:fp:periodsFinite}: 3 complex samples & $\mathfrak t_3$, $\widehat M_3$: 2 one-variable ${}_7F_6(1)$ and
 3 two-variable Horn values\\
$Z_d$ & \eqref{gff:gc:finiteClosure}: 318 values, $\le6$ variables & the 5 functions of the $Q_d$ row with first and second
 $b$-derivatives, through $\hat c_0'$ and $\frac19(u\mathfrak t_3''+v\widehat M_3'')$\\
$\KK_d$ & \eqref{gff:gc:finiteClosure}: 386 values, $\le6$ variables & Section~\ref{gff:gdk:section}: a polygamma part
 plus $\mathsf r_3$, a third-order $\tau$-jet of three-node samples of three two-variable Horn functions
 \eqref{gff:gdk:Ksample}, under the inputs of Section~\ref{gff:gdk:conventions}; or one two-fold Barnes integral
 \eqref{gff:gdk:KMB}\\
$P_5(d)$ & \eqref{gff:fp:primitiveSample} & not treated in this section\\\bottomrule
\end{tabular}
\end{center}
Apart from $P_5$ and $\KK_d$, $C_5(d)$ needs eight hypergeometric functions with at most two summation indices at
unit argument: five for $r=4$ ($\mathfrak h_x,\mathfrak h_1,\mathfrak h_{1/2},\mathfrak s_0,\mathfrak s_{1/2}$) and three for $r=3$ ($\mathfrak t_3$ and the right-part and middle-part sums of $\widehat M_3$, Section~\ref{gff:gdr:samples}). Split into one-variable rows and two-variable remainders, they give six one-variable rows, namely ${}_8F_7$, two
${}_9F_8$ and ${}_7F_6$ for $r=4$ and two ${}_7F_6$ for $r=3$, plus eight two-variable Horn values. $C_5(d)$ also needs the first $b$-derivatives of all eight functions, the second $b$-derivatives of the three functions for $r=3$, and elementary numbers ($\pi$; $\sin,\cos$ of $\pi q,\pi q/2$; $\Gamma$,
$\psi_0$, $\psi_1$ at $q/2,q,3q/2,2q,(1\pm q)/2,(1+3q)/2,(3\pm q)/2$). After the exact relations they enter as eight
numbers per dimension: $\mathcal X_T,\mathcal X_0,\mathcal X_{1/2},\mathfrak D_L$ and
$\mathfrak t_3,\hat c_0,\hat c_0',\frac19(u\mathfrak t_3''+v\widehat M_3'')$. For $Q_d$ and $Z_d$ together these are
four numbers. Since $Q_d$ is one combination of $\mathfrak t_3$ and $\hat c_0$, \eqref{gff:gdr:Q}, $Z_d$ needs three numbers beyond the value of $Q_d$: one of $\mathfrak t_3,\hat c_0$, then $\hat c_0'$ and $\frac19(u\mathfrak t_3''+v\widehat M_3'')$.

\subsection{Further reductions tried}\label{gff:gdr:attempts}
\begin{itemize}
\item \emph{Frozen weights} (exact): one derivative constant for $\LL_d$, one first- and one second-derivative
 constant for $Z_d$ (above).
\item \emph{Pole node only} (exact): \eqref{gff:gdr:taylorW}--\eqref{gff:gdr:taylorLZ}, one Horn family per wheel, one Laurent
 order below \eqref{gff:gs:QPseries}--\eqref{gff:gs:Zseries}. It trades functions for parameter derivatives, so in the
 unweighted Horn class it is not shorter than \eqref{gff:gdr:C5reduced}.
\item \emph{Node choices.} $t=0$ is the only node with a single block; the reflection centers need two; a general
 node three. A real two-node form of $W_4$ using $t=0$ exists only for $|\cos\frac{\pi b}2|\le\frac12$
 (second node $\cos2\pi\tilde t_2=3\varsigma_b/(1+2\varsigma_b^2)$, a general point moving with $b$), so it does not
 shorten $\LL_d$.
\item \emph{Summation theorems for the zero-spin rows.} The row of $\mathfrak h_x$ (and of $\mathfrak t_3$) is a
 well-poised ${}_{2r}F_{2r-1}(1)$ with $a=\sigma$, pairs $(\sigma,1)$, $(1-g,2q)$, not very-well-poised; the rows at
 the reflection centers are not well-poised for generic $b$. Dixon and Dougall do not apply for $r\ge3$; at $d=1$
 only the sum of the $\ell=0,1$ rows is summable. Contiguous relations do not connect the functions (shifts
 $x,\frac12,1$ are non-integer apart). Negative.
\item \emph{Gamma form of $\hat c_0$} (true at $d=1,2,3$): PSLQ at 80 digits at $d=\frac45,\frac85,\frac{12}5,\frac{16}5$
 of $\log|\hat c_0|$ against $\log$ of $\pi,2,5,\frac{1+\sqrt5}2,\Gamma(\frac15),\Gamma(\frac25)$ (all Gamma and
 trigonometric values at multiples of $\frac1{10}$): no relation with coefficients below $10^{11}$. Negative,
 numerical.
\item \emph{Elementary closure of $\mathfrak D_L$} (true at $d=1,2$): PSLQ at $d=\frac43,\frac83$ of $\LL_d$
 against $\{\mathcal X_T,\mathcal X_0,\mathcal X_{1/2}\}\times\{1,\sqrt3,\frac{\log2}\pi,\frac{\log3}\pi,
 \frac{\sqrt3\log2}\pi,\frac{\sqrt3\log3}\pi\}$ (95 digits): no relation stable under a change of precision;
 none with coefficients below about $10^4$. Negative, numerical.

\item \emph{Polygamma closure of $Z_d/Q_d$ and $\mathfrak t_3''/\mathfrak t_3$} (true
 at $d=1,2$, \eqref{gff:mgz:Z1}, \eqref{gff:mgz:Z2}). A bounded PSLQ probe at 70 digits, on
 the 90-digit values of the two-node item of Section~\ref{gff:gdr:checks}, at $d=1.25,1.5,2.5$, against the 15
 constants $1$, $\pi^2$, $\psi_1(q)$, $\psi_1(2q)$, $\psi_1(\frac q2)$,
 $\psi_1(1-\frac q2)$, $\psi_1(\frac{3q}2)$, $\psi_1(3q)$, $\psi_0(q)^2$,
 $\psi_0(q)\psi_0(2q)$, $\psi_0(2q)^2$, $\pi\psi_0(q)$, $\pi\psi_0(2q)$,
 $\pi^2/\sin^2\pi q$, $\pi\cot\pi q$ (coefficient bound $10^4$): at $d=1.5$ only a
 relation internal to the basis (target coefficient zero), at $d=1.25$ and $2.5$ only
 relations with coefficients of size $10^3$--$10^4$, the size expected of spurious
 relations at this precision (\path{LRI_C5_general_d/Z/pslq_quick.py},
 \path{LRI_C5_general_d/Z/results/pslq_quick.log}). Negative, numerical; not an
 independence proof. No closed form of $\hat c_0$ is known at generic $d$: it vanishes
 at $d=1$, equals $-\pi^3s_b^3\mathfrak t_3/(1+\varsigma_b)$ with $\mathfrak t_3$ the
 square of Dixon's sum at $d=2$, and is a Gamma product in $b$ at $d=3$. The $d=2$ and
 $d=3$ Gamma forms have different shapes, so a uniform Gamma product in $(d,b)$ is not
 expected.

\end{itemize}

\subsection{Checks}\label{gff:gdr:checks}
Computations on the symmetry computing cluster of the Perimeter Institute
(\path{LRI_C5_general_d/specialize/results/}) and the separate review
(\path{LRI_C5_general_d/review/}). The package replay, \texttt{make -f public.mk replay}, compares with the records
only the numbers listed in \path{GFF/results/closed_forms_general/reduction_d/expected.json}. These do not include the
numbers of Sections~\ref{gff:mgf:nodes}, \ref{gff:mgl:section} and \ref{gff:mgz:proofs}, nor those of the items below on
the general-$b$ periodization, on the review probes and the product form in more detail, and on the two-node
formula; those numbers are read from the records named there, and no replay target compares them with the text.
\begin{itemize}
\item \eqref{gff:gdr:d12lattice} against the five-Horn samples: $d=1$, $b=\frac12,0.45$ and $d=2$, $b=1,1.1$, $r=3,4$:
 $10^{-44}$--$10^{-46}$. Relations \eqref{gff:gdr:landen} and $\mathcal A,\mathcal B$ against $B_\pm$: $5\cdot10^{-46}$.
 $P_4(1),\LL_1,P_4(2),\LL_2$ from the modes: $\le2\cdot10^{-45}$.
\item $d=3$ dictionary against the 160-digit values of \path{LRI_C5_d3/wf_L3}: $c_{4,k}$, $X_0$, $\mathfrak a_4$,
 $\mathfrak b_4$, $\mathfrak d_H$, $\mathfrak d_a$, $P_4(3)$, $\LL_3$ all to $1.2\cdot10^{-51}$ (50 digits);
 $\LL_3$ formula $5\cdot10^{-64}$.
\item \eqref{gff:gdr:taylorW} against the multi-node values ($d=1.5,2.5$; $b=d/2$, $d/2+0.04$; $r=3,4$):
 $5\cdot10^{-21}$--$3\cdot10^{-23}$ (Cauchy circle in the shift).
\item \eqref{gff:gdr:P4}, \eqref{gff:gdr:L}, \eqref{gff:gdr:Q}, \eqref{gff:gdr:Z} against the five-Horn and two-node codes at
 $d=1$, $2$, $3$, $0.8$, $1.25$, $1.5$, $2.5$, $2.75$ (50 digits): $\le2.1\cdot10^{-48}$; the block
 $384AP_4-576\LL-1728Q^2+(720A^2+1008R)Q+7152Z$ and its reduced form \eqref{gff:gdr:C5reduced}: $\le1.2\cdot10^{-48}$;
 against closed forms ($d=1,2$) $\le4\cdot10^{-50}$, against the independent spectral references $4\cdot10^{-50}$
 (50-digit reference values, 60-digit samples; \path{LRI_C5_general_d/review/joint/numerics-end-to-end}), at
 $d=3$ $6\cdot10^{-34}$ (limited by the 33-digit $Q_3$ used).

\item The explicit five-Horn forms \eqref{gff:mgl:W4five} and \eqref{gff:mgl:Lfinal} against the proved values at $d=1,2$, the
 160-digit values at $d=3$, the independent spectral references at nine non-integer dimensions and off-center spectral
 values at ten points, with per-function checks of the jet derivatives: Section~\ref{gff:mgl:checks}.

\item End to end, with $P_5(d)$ from \eqref{gff:eq:P5FixedAnswer}, \eqref{gff:fp:P5two} and \eqref{gff:td:P5answer} (values in
 Section~\ref{gff:gs:masterTable}) and $\KK_d$ from \eqref{gff:gdk:Kr3} ($\mathsf r_3$ from the jets of
 Section~\ref{gff:gdk:numerics}):
 \[\begin{aligned}
 C_5(1)&=-731.139476624539825853869126582527231,\\
 C_5(2)&=-38.8336458962253605529369824675329232,\\
 C_5(3)&=-5.850351420751327229217944230497744300111.
 \end{aligned}\]
 They agree with $C_5(1)$ of \eqref{gff:eq:C5psiQuarter}, printed to 33 decimals in Section~\ref{gff:sec:checks}, to $3\cdot10^{-34}$ (all
 printed digits), with \eqref{gff:gc:C5twoClosed} to $2\cdot10^{-50}$ and with \path{LRI_C5_d3} to $5\cdot10^{-45}$
 (\path{LRI_C5_general_d/review/joint/numerics-end-to-end}).
\item \emph{Top-mode theorem and its product form} (Section~\ref{gff:gdr:top}). The balance of every single factor and
 the factor $1+z$ in \eqref{gff:gdr:Dgen} are both essential. An unbalanced wheel with the same pole classes ($d=1$,
 $r=4$, $\sigma=0.2$, spoke
half-exponent $0.25$ instead of $0.3$) has $|\gamma_{\pm3}|$ equal to $1.7\%$ of the largest mode, against $10^{-41}$ for
the balanced kernel at 40 digits; an independent 50-digit code (\path{LRI_C5_general_d/review/topmode-numeric}) finds unbalanced top
modes of $6.7\cdot10^{-4}$ to $4.6\cdot10^{-2}$, linear in the imbalance, and $1.9\cdot10^{-3}$ when the factor $1+z$ is
dropped. For balanced wheels the same code finds $|c_{r,r-1}|/\max_k|c_{r,k}|$ between $2.4\cdot10^{-51}$ and
$1.6\cdot10^{-49}$ at 22 configurations with $r=2,\dots,6$, $0.8\le d\le3.2$ and complex $b$, against surviving modes of
$1.6\cdot10^{-3}$ to $1$. For products of different factors, a separate 48-digit lattice-sum code
 (\path{LRI_C5_general_d/review/lemmaG-numeric}; reruns at 54 digits on other compute nodes) finds top modes
 between $2.8\cdot10^{-48}$ and $4.7\cdot10^{-44}$ of the largest mode, at the level of the vanishing boundary
 modes $\gamma_{\pm r}$, at 12 mixed wheels ($r=3,4$; $d=1.3$, $2.4$, $3.1$; unequal $b_i$, real and complex shifts
 and offsets, one with $\operatorname{Re}g_1<0$),
 $6.5\cdot10^{-47}$--$5.2\cdot10^{-45}$ at four $R$ kernels with real $\tau$ and $3.0\cdot10^{-46}$ at one wheel;
 its reruns give at most $9.3\cdot10^{-54}$. The two top modes vanish separately, also where $|\gamma_{r-2}|$ and
 $|\gamma_{2-r}|$ differ by up to $79\%$ of the largest mode, and they stay
 below $2\cdot10^{-47}$ (54 digits) at four analytically continued points on the boundary of (C) or outside it, two
 of them in the $R$ family. A per-factor imbalance $\delta_{\rm bal}$ at fixed total balance switches them on
 linearly, with relative top modes
 $0.047$ down to $0.0013$ for $\delta_{\rm bal}=0.15$ down to $0.0015$. The per-spin code of
 Section~\ref{gff:gdk:numerics}
 (\path{LRI_C5_general_d/K_general/note_fragment/topmode/perspin_topmode.py}; its modes agree with those of the 48-digit lattice-sum code to $10^{-15}$) finds top modes of at
most $6.5\cdot10^{-16}$ of the largest mode, at the level of its extrapolation, for five $R$ kernels (real and complex
$\tau$) and two random mixed kernels ($r=3,4$), against $3\cdot10^{-3}$ to $6\cdot10^{-2}$ for controls that keep the total balance but break (B), use
multiplicities of another dimension, or shift one right argument. The short check
 \path{LRI_C5_general_d/foundation/check_topG.py} (30 digits, about 20\,s on a laptop) finds $10^{-24}$ to $9\cdot10^{-16}$, falling like the truncation error, at
four mixed points with $r=3,4$ and $d=0.6$, $1.7$, $3.3$, three of them with one $\operatorname{Re}g_i<0$ (Step~7),
against $9\cdot10^{-10}$ to $2\cdot10^{-5}$ for controls with right shifts $\pm0.1$ at fixed total balance; it also
 confirms \eqref{gff:gdr:kernelBalance} and \eqref{gff:gdr:mixed} to $10^{-29}$. No further mode
vanishes in general: at the points tested, $c_{5,3}$ and $c_{6,4}$ are $1.6\cdot10^{-3}$ to $0.16$ of the largest mode.
 The proof review \path{LRI_C5_general_d/review/lemmaG-proof} tested Steps~1--6 numerically in unbalanced cases,
 and the endpoint exponents of \eqref{gff:gdr:endpoint} (item on the product form below).
\item Independent review, own codes: top modes as quoted in the item above; the printed formulas of
 Sections~\ref{gff:gdr:samples}--\ref{gff:gdr:special} and \eqref{gff:gdr:C5reduced} rederived symbolically with difference $0$
 (\eqref{gff:gdr:taylorLZ} by hand);
 off-center $W_3,W_4$ from independent samples (exact big-integer convolutions) to $44$--$53$ digits; $P_4$ and
 $\LL_d$ against spectral references to $31$--$39$ and $31$--$38$ digits; the coefficient functions of the
 five-Horn formula for $\LL_d$ at 11 values of $d$ to $1.2\cdot10^{-38}$; $Z_1,Z_2,Z_3$ to $30.7$, $35.2$, $29.5$
 digits; the reduced $C_5$ block against the reference-master block to $30.4$--$35.6$ digits at 8 values of $d$.

\item \emph{The general-$b$ periodization} (Section~\ref{gff:gdr:periodization}). The checks use the code
 \path{wheel_general_b.py} with \path{hypsums.py} of \path{LRI_C5_general_d/foundation/}, driven by \path{run_validation.py} on the symmetry
 cluster (job 568940 on one 64-core AMD node, $50$ requested digits and $15$ guard digits;
 supplementary job 568959); records \path{foundation/results/val50.json} and \path{valx.json}, with the
 comparison \path{val50_analysis.txt} and \path{val50_summary.json} written by \path{analyze_val.py}.
 References: the closed forms of Section~\ref{gff:c1:closedForms} at $d=1,2$ (for $P_4$, $Q$, $\LL$ and $Z$),
 the archived $d=3$ values ($\geq90$ digits), and at non-integer $d$ the independent spectral values of
 \path{LRI_C5_general_d/references/reference_values.json} (41--50 digits at $b=\frac d2$ and 56 digits
 off it when these runs were made; the file was later rewritten with 50 significant digits).
 \begin{itemize}
 \item At $b=\frac d2$, with the fixed nodes of \eqref{gff:mgf:Wsample} ($N=N_r$, complex discrete Fourier
 transform) and with the $b$-dependent nodes and the real cosine transform \eqref{gff:fp:cosineTransform}
 (which at $b=\frac d2$ is literally the rule of Section~\ref{gff:app:generalPeriodization}),
 $W_4=P_4(d)$ and $W_3=3Q_d$ agree with the references to $2.2\cdot10^{-66}$--$1.5\cdot10^{-63}$ at
 $d=1,2,3$ (the working precision) and to $9.7\cdot10^{-47}$--$3.8\cdot10^{-45}$ at
 $d=0.8,1.25,1.5,2.5,2.75$, which is the accuracy of the references there (42--50 digits).
 \item At $b\neq\frac d2$ ($b=\frac d2\pm0.05$, at $d=0.8$ $b=\frac d2\pm0.03$; $r=3,4$; the five
 non-integer $d$, 20 points), the fixed-node values agree with the spectral references to
 $3.1\cdot10^{-55}$--$3.0\cdot10^{-49}$, and the $b$-dependent nodes (not run at $d=0.8$) to
 $3.3\cdot10^{-54}$--$7.7\cdot10^{-50}$. A second spectral code, \path{foundation/spectral.py}, which
 shares only the Hurwitz routine, agrees at six of these points
 ($d=1.5$, $b=0.7,0.8$, $r=3,4$; $d=2.5$, $b=1.2$, $r=3,4$) to $8.2\cdot10^{-54}$--$1.9\cdot10^{-50}$. At
 $d=2$ the exact $W_3(2;b)$ of \eqref{gff:gs:W3d2} is reproduced at $b=0.8$ and $1.2$ to $1.3\cdot10^{-65}$
 and $7.0\cdot10^{-65}$.
 \item Structure: $|c_{r,r-1}|/\max_k|c_{r,k}|$ lies between $4\cdot10^{-66}$ and $9\cdot10^{-49}$ at
 every point; at an extra point $t_*=0.2371+0.0917\ii$ the trigonometric polynomial reproduces
 $\Xi_r(t_*;b)$ to at most $1.0\cdot10^{-62}$ at $d=1,2,3$ and $1.2\cdot10^{-49}$ at non-integer $d$; at
 $d=1$ also $c_{4,1}=c_{4,3}=0$ and $c_{3,0}=c_{3,2}=0$, and at $d=2$ $c_{3,0}=c_{3,1}$, to working
 precision, as \eqref{gff:gdr:d12lattice} requires. The reduced nodes $N=N_r'$ agree with $N=N_r$ to
 $6.4\cdot10^{-65}$ ($d=2$, $b=1$, $r=4$), $4.1\cdot10^{-63}$ ($d=3$, $b=1.5$, $r=3$) and
 $1.2\cdot10^{-53}$--$2.5\cdot10^{-50}$ at four non-integer points.
 \item The $b$-dependence: Cauchy integrals of the finite-sample $W_r(d;b)$ over $96$ points of the
 circle $|b-\frac d2|=\rho$ (phases $2\pi(k+\frac12)/96$, conjugate symmetry used), with $\rho$ one third
 of the distance to the nearest singularity of $W_r$ or to $b_{\max}$, so that the truncation error is
 below about $3^{-96}$, give $\LL_d=\partial_bW_4$ and $Z_d=\frac19\partial_b^2W_3$ at $b=\frac d2$ to
 $3.0\cdot10^{-63}$--$1.0\cdot10^{-46}$ relative at $d=1,2,3$ and to
 $1.2\cdot10^{-46}$--$3.5\cdot10^{-45}$ at $d=0.8,1.25,1.5,2.5,2.75$ (references of 41--50 digits). The
 first derivative $\partial_bW_3(d;\frac d2)$, which vanishes by \eqref{gff:gs:Tjet} with
 \eqref{gff:gs:wheelJets}, comes out between $10^{-61}$ and $10^{-44}$. The sentence after the table that
 follows this list compares these values with the other evaluations.
 \item The jet formula \eqref{gff:mgf:jets} evaluated literally at $30$ digits (\path{foundation/sample_jets.py},
 job 568967; the per-block values and $b$-derivatives at the nodes are stored in
 \path{foundation/results/sample_jets.json}), with the sample derivatives at the reduced fixed nodes
 taken from Cauchy circles of $48$ points of radius one sixth of the distance to the nearest
 singularity of the samples: $\LL_d$ and $Z_d$ at $d=1,1.5,2,2.5,3$ to
 $2.4\cdot10^{-31}$--$2.8\cdot10^{-30}$, and $W_r(d;\frac d2)$ to $5.3\cdot10^{-32}$--$2.7\cdot10^{-30}$.
 \item Beyond $b_{\max}$ (Section~\ref{gff:mgf:nodes}): the continued Horn values against
 \path{spectral.py} at $(d,b,r)=(2.5,1.8,4)$, $(2.5,1.55,3)$, $(2.75,1.95,4)$ and $(2.75,1.65,3)$ to
 $4.7\cdot10^{-45}$, $4.9\cdot10^{-48}$, $2.1\cdot10^{-41}$ and $1.7\cdot10^{-49}$, limited by the
 cancellation near the pole of the $c_{r,k}$ at $b_{\max}$. At the resonant points $(d,b,r)=(\frac43,\frac23,3)$
 and $(2.5,1.5,4)$, where $r(d-b)/2\in\mathbb Z$ and the circle mean is used, to $3.5\cdot10^{-61}$ and
 $1.5\cdot10^{-61}$.
 \end{itemize}
 \begin{center}\small
 \begin{tabular}{@{}lll@{}}\toprule
 $d$ & $P_4(d)=W_4(d;\frac d2)$ & $3Q_d=W_3(d;\frac d2)$\\\midrule
 0.8 & \texttt{4403.5659705473683681} & \texttt{259.83164202536385403}\\
 1.25 & \texttt{1120.1530930090638933} & \texttt{104.22818434903994884}\\
 1.5 & \texttt{632.7128141805458393} & \texttt{71.175347423099135692}\\
 2.5 & \texttt{118.19407871190995074} & \texttt{23.177096659039552764}\\
 2.75 & \texttt{84.737926139821720537} & \texttt{18.545946016773332171}\\\bottomrule
 \end{tabular}
 \end{center}
 The table gives the finite-sample $P_4(d)$ and $3Q_d$ (fixed nodes), within $3.8\cdot10^{-45}$ of the
 spectral references. The values of $\LL_d$ and $Z_d$ from the $b$-circles at these five dimensions
 (\path{foundation/results/val50_summary.json}) are within $3.5\cdot10^{-45}$ of the references and agree in
 every printed digit with the table of $\LL_d$ in Section~\ref{gff:mgl:checks} and the table of $Z_d$ in the
 two-node item below.
\item \emph{Caution on the structural checks.} The periodicity, reflection and off-grid checks test the
 structure of Section~\ref{gff:gdr:finite}, not the accuracy of the samples. They cannot see a sample error
 that is itself a reflection-symmetric trigonometric polynomial of degree at most $r$; a $t$-independent
 error in a tail amplitude, multiplied by $D_b(t)^r$ and the Gamma prefactors, is of this kind. Only the
 top and boundary modes and the comparisons with independent references test accuracy. Errors of this
 kind occurred in the reviewer's own code and were caught only by the top and boundary modes and the
 references.
\item \emph{The review probes of Sections~\ref{gff:gdr:finite}--\ref{gff:gdr:residue}, in more detail.}
 \path{LRI_C5_general_d/review/topmode-numeric} (own lattice sums by two-dimensional Euler--Maclaurin
 summation with tanh--sinh quadrature, $50$ digits, symmetry jobs 569215 and 569260--569313): besides
 the top modes quoted above, the boundary modes are $4\cdot10^{-51}$ to
 $2.1\cdot10^{-49}$; at $r=2$, $d=1.3$, $b=0.6+0.1\ii$, where $\alpha=0.7$ barely exceeds the bound
 $0.65$, $\Xi_2$ is constant to $2.7\cdot10^{-51}$; the weights $B_{r,k}$ of \eqref{gff:gdr:functional} agree
 with a direct series residue to $1.7\cdot10^{-51}$--$7.8\cdot10^{-49}$; $W_r$ agrees with the spectral
 references to 47--50 digits; the stored $c_{r,k}$ of \path{foundation/results/val50.json} agree to
 $1.4\cdot10^{-50}$--$4.2\cdot10^{-50}$; the unbalanced controls quoted above stand against a
 floor of $10^{-30}$. \path{LRI_C5_general_d/review/foundation-rigor}
 (double-precision brute-force lattice sums, \path{brute_xi.py}, 17--25 nodes, symmetry jobs 569242 and
 569252): the five blocks of \eqref{gff:mgf:blocks} agree with the brute-force sums to
 $2.5\cdot10^{-16}$--$1.4\cdot10^{-14}$, and the top modes are $1.9\cdot10^{-15}$ to $2.7\cdot10^{-14}$
 in the fast cases ($d=3.5$; $d=0.9$, with negative multiplicities; complex $b$); in the slow cases the
 residuals halve when the cutoffs are doubled, so they are truncation error. These checks also cover
 $\sigma>1$ at $d=3.9$.
\item \emph{The checks of the product form, in more detail.} The proof check
 \path{LRI_C5_general_d/review/lemmaG-proof} (record \path{review.json}; the test scripts
 \path{t1_laplace.py} to \path{t5b_r4.py} saved no logs) tested: the Laplace forms \eqref{gff:gdr:laplace}
 for complex $g$, also at $\operatorname{Re}w<0$, to $10^{-20}$--$10^{-26}$; the relation
 $\gamma_{1-r}=\kappa_+s_1^+$ with $s_1^+$ from Step~4 and all its constants, against lattice sums in
 unbalanced cases, where the top mode is nonzero, so that the derivation and not only the final
 cancellation is tested: $7\cdot10^{-11}$ ($r=2$, real), $2\cdot10^{-8}$ ($r=2$, complex; limited by the
 lattice truncation), $2\cdot10^{-13}$ ($r=2$, total balance only), $10^{-12}$ ($r=3$, $g=(1,2,3)$,
 complex) and $3\cdot10^{-12}$ ($r=4$, $g=(1,1,2,2)$, $d=2.9$, total balance only), and
 $10^{-22}$--$2\cdot10^{-27}$ in balanced cases; the lower half-plane of Step~6, by the direct formula and
 by the conjugation, to $3\cdot10^{-12}$ at a complex unbalanced point with $r=2$, where
 $|\gamma_1|\neq|\gamma_{-1}|$; and the endpoint exponents of \eqref{gff:gdr:endpoint}, with fitted slopes
 $-0.300$ at $\varkappa=0.7$ and $-0.600$ at $\varkappa=0.4$ on both sides of $2\pi$, bounded at
 $\varkappa=1.2$, and bounded and oscillating at $\varkappa=1$ with complex $g_i$. As a check that the
 mechanism does not prove too much, its formal extension to $r=1$ at $d=2$ predicts
 $\sum_{m\in\mathbb Z}\Gamma(m+a)/\Gamma(m+a+g)=0$ for $\operatorname{Re}g>1$, a known corollary of
 Dougall's bilateral sum. The 48-digit code \path{LRI_C5_general_d/review/lemmaG-numeric} (record
 \path{probe_report.json}) shares nothing with the general-$b$ code, with \path{check_topG.py} or with the
 $\KK_d$ code. It samples $\Xi_T$ at $2r+1$ points on a horizontal line and Fourier transforms; the
 boundary modes $\gamma_{\pm r}$, which vanish by Step~0, set the noise floor. It splits $S_T$ into its
 right, middle and left parts, each a sum over a quadrant, evaluated by exact head sums, exact
 Abel--Plana tails and the corner double integral (tanh--sinh rules matched to the endpoint exponents,
 Gauss--Legendre panels); Gamma ratios come from its own asymptotic series, and \texttt{mpmath.zeta} is
 not used. Main runs are at $48$ digits; every case inside (C) and every continued point was rerun at
 $54$ digits with finer settings and a disjoint set of nodes (symmetry jobs 569336--569347 and 569461).
 Validation: closed-form quadrant sums to $1.1\cdot10^{-52}$; main against fine runs to
 $5\cdot10^{-44}$; the 50-digit modes of \path{topmode-numeric} at one wheel reproduced to all 40
 printed digits, and the modes of the per-spin code of Section~\ref{gff:gdk:numerics} at three of its points
 to its accuracy of $10^{-15}$. Controls: a per-factor imbalance at fixed total balance
 ($\beta_1\to\beta_1+\delta_{\rm bal}$, $\beta_2\to\beta_2-\delta_{\rm bal}$) at a real mixed wheel with
 $r=3$, $d=1.3$ gives relative top modes $0.047$, $0.030$, $0.0116$, $0.0041$, $0.00128$ for
 $\delta_{\rm bal}=0.15$, $0.05$, $0.015$, $0.005$, $0.0015$, with $\gamma_{-2}/\delta_{\rm bal}$
 converging to about $1.72+1.30\ii$ with an $O(\delta_{\rm bal})$ correction, so the onset is linear,
 while $\gamma_{\pm3}$ stay at $1.4$--$1.6\cdot10^{-44}$; at the complex version of that point and complex
 $\delta_{\rm bal}$ the two top modes are nonzero and unrelated ($0.025$ and $0.015$ at
 $\delta_{\rm bal}=0.05+0.03\ii$, $0.0032$ and $0.0018$ at $\delta_{\rm bal}=0.005+0.003\ii$); an $R$
 kernel with right shifts $\pm0.15$, a control of the per-spin code, gives $0.0135$ and $0.0029$, as that
 code does; outside (C), at two mixed wheels with $2\varkappa-d=-0.06$ and (B) exact, the plain lattice
 sum diverges (largest $|\gamma_m|$ about $10^{102}$ at truncation depth $26$ and $10^{134}$ at depth
 $36$), while the analytically continued sum (the leading power of the divergent corner integral
 subtracted and restored analytically) stays a trigonometric polynomial of degree $r-1$ with vanishing
 top modes ($6.8\cdot10^{-42}$--$3.7\cdot10^{-41}$ at $48$ digits, at most $6.1\cdot10^{-50}$ at $54$).
 The continuation reproduces a synthetic closed form outside its convergence region to
 $7.6\cdot10^{-17}$, and continued and plain evaluations agree to $8\cdot10^{-44}$ at a point inside
 (C). So (C) acts as a convergence hypothesis, as Step~7 suggests. One complex case was also evaluated
 on the lines $\operatorname{Im}t=0.2$ and $-0.25$, with modes agreeing to $1.2\cdot10^{-42}$. All error
 bars of this probe are empirical (convergence ladders, independent reruns, must-vanish modes), with no
 interval bounds. Endpoint exponents near $0.1$ converge slowly in the quadrature step; the cases avoid
 that regime or were cross-checked. The continued points test only failures of $2\varkappa>d$ by less
 than one; the $\tau$-circle of Section~\ref{gff:gdk:continuation} at $d=2.5$, where also
 $\varkappa>d-1$ fails, was not tested.


\item \emph{Two-node formula for the three-wheel} (Section~\ref{gff:mgz:proofs}). Code in
 \path{LRI_C5_general_d/Z/}, records in \path{LRI_C5_general_d/Z/results/}.
 \path{z_blocks.py} evaluates the five functions of \eqref{gff:mgz:five} at any complex
 $b$ with the accelerated Horn sums of \path{LRI_C5_general_d/foundation/hypsums.py}
 and \path{wheel_general_b.py} (exact finite part, Stirling tails,
 Euler--Maclaurin Hurwitz zeta values with relative precision, three-family
 convolution asymptotics, exact regularized multiplicity moments) and assembles
 \eqref{gff:mgz:Zjets}, its reduced forms and \eqref{gff:mgz:simplified}. The main engine
 \path{z_jets.py} reruns the same algorithms with every parameter in the ring
 $\mathbb C[\varepsilon]/(\varepsilon^3)$, $b=\frac d2+\varepsilon$, so the
 $b$-derivatives of the five functions are exact derivatives of the accelerated
 evaluation, with no numerical differentiation; at jet order $0$ it reproduces the
 block values of the general-$b$ code to working precision. \path{z_cauchy.py}
 computes the same jets from Cauchy circles in $b$; it is a cross-check, and it is
 used at the resonant dimensions $\frac{3d}4\in\mathbb Z$ ($d=\frac43,\frac83$), where
 the accelerated sums have removable singularities at $b=\frac d2$.
 \path{Z_general.py} is the entry point ($Z_d$ to 25 digits in about 5\,s on one
 laptop core); the full 90-digit set of the table below took 12\,s on one
 64-core node of the symmetry cluster.

 References: at $d=1,2$ the closed forms \eqref{gff:gs:d1Dictionary} and
 \eqref{gff:gs:Z2exact}; at $d=3$ the 250-digit value of \path{LRI_C5_d3/wf_Z3/}; at
 $d=0.8$, $1.25$, $1.5$, $2.5$, $2.75$ the independent spectral values of
 \path{LRI_C5_general_d/references/reference_values.json} (spectral code
 \path{LRI_C5_general_d/refvals/hp_general_d.py}); at $d=\frac43$, $1.75$, $\frac83$, $3.5$, $3.75$ new runs of the same
 spectral code, copied unchanged to \path{LRI_C5_general_d/Z/spectral_ref/} (45
 digits, 64-point $b$-circle). The spectral code shares no code with the
 finite-sample evaluation. At $d=0.7$, just above the threshold $d=\frac23$, the
 spectral run did not converge: its own $\partial_bW_3$ came out as $-83$, and its
 $b$-circle passes within $0.013$ of the spectral convergence bound
 $b_{\rm spec}=d-\frac13=0.367$ of Section~\ref{gff:gdr:periodization}. There only a
 consistency check is available: a $b$-circle of the general-$b$ finite-sample
 formula with five fixed complex nodes and all five blocks per node
 (\path{LRI_C5_general_d/Z/fd_circle_check.py}), which shares the Horn summation code
 and is therefore not independent.

\begin{center}\scriptsize
\begin{tabular}{@{}l l l l l l@{}}\toprule
$d$ & $Z_d$ (two-node formula, 36 digits shown) & reference & rel.\ dev. & $|\partial_bW_3|$ & engine, digits\\\midrule
$1$ & \texttt{764.931030199588577040161778769488397} & \eqref{gff:gs:d1Dictionary} & $4.0{\cdot}10^{-91}$ & $1.1{\cdot}10^{-97}$ & jets, 90\\
$2$ & \texttt{46.7925830723296966567388258779810842} & \eqref{gff:gs:Z2exact} & $3.9{\cdot}10^{-91}$ & $3.5{\cdot}10^{-98}$ & jets, 90\\
$3$ & \texttt{8.82092489306009248553686732817408087} & $d=3$, 250 digits & $4.4{\cdot}10^{-91}$ & $5.3{\cdot}10^{-98}$ & jets, 90\\
$0.8$ & \texttt{1870.91240775271737784994106229072511} & spectral, 46 digits & $2.0{\cdot}10^{-45}$ & $4.7{\cdot}10^{-70}$ & jets, 90\\
$1.25$ & \texttt{312.281876942721713720205721011622893} & spectral, 41 digits & $1.9{\cdot}10^{-46}$ & $1.7{\cdot}10^{-73}$ & jets, 90\\
$1.5$ & \texttt{149.908309291859708760442421585801485} & spectral, 44 digits & $2.9{\cdot}10^{-45}$ & $4.3{\cdot}10^{-73}$ & jets, 90\\
$2.5$ & \texttt{18.78702819646921677529557752590324} & spectral, 44 digits & $1.2{\cdot}10^{-45}$ & $2.1{\cdot}10^{-74}$ & jets, 90\\
$2.75$ & \texttt{12.6721261451666548401131859193957078} & spectral, 44 digits & $1.2{\cdot}10^{-46}$ & $1.8{\cdot}10^{-74}$ & jets, 90\\
$4/3$ & \texttt{240.897293321182080101413442492954567} & new spectral run & $1.0{\cdot}10^{-39}$ & $4.3{\cdot}10^{-44}$ & Cauchy, 60\\
$1.75$ & \texttt{80.4310856344043874032911580057932935} & new spectral run & $1.4{\cdot}10^{-40}$ & $1.7{\cdot}10^{-73}$ & jets, 90\\
$8/3$ & \texttt{14.3946259191595698392288065462545566} & new spectral run & $1.0{\cdot}10^{-39}$ & $4.3{\cdot}10^{-43}$ & Cauchy, 60\\
$3.5$ & \texttt{4.60407732413756930916907954325772311} & new spectral run & $1.4{\cdot}10^{-40}$ & $1.4{\cdot}10^{-74}$ & jets, 90\\
$3.75$ & \texttt{3.42691133073098483556250678138596949} & new spectral run & $1.4{\cdot}10^{-40}$ & $1.5{\cdot}10^{-74}$ & jets, 90\\
$0.7$ & \texttt{3193.82811885501679042755334997746735} & fixed-node circle$^\ast$ & $7.9{\cdot}10^{-40}$ & $7.8{\cdot}10^{-70}$ & jets, 90\\
\bottomrule\end{tabular}
\end{center}
 The table gives $Z_d$ from \eqref{gff:mgz:Zjets} against the references (records
 \path{LRI_C5_general_d/Z/results/Z_values.json}, \path{zj90.json},
 \path{zc60x.json}, \path{compare_refs.log}). ``rel.\ dev.'' is relative to the
 reference. The spectral references are good to the number of digits stated, or to
 about $10^{-39}$ for the new runs. The deviations of the five rows marked
 ``spectral'' were computed with the reference file as it stood before it was rebuilt
 with 50 printed digits (\path{compare_refs.log} predates the rebuild); the stated
 accuracies of the references did not change. $|\partial_bW_3|$ is the first rim derivative at
 $b=\frac d2$ computed from the same jets; it must vanish. $^\ast$At $d=0.7$ the
 reference is the fixed-node $b$-circle described above, a consistency check only.

 The two-node identity \eqref{gff:gdr:twonode} itself was tested away from $b=\frac d2$,
 including complex $b$ (\path{LRI_C5_general_d/Z/offdiag_check.py}, record
 \path{LRI_C5_general_d/Z/results/offdiag60.json}), against the 56-digit spectral
 values of \path{LRI_C5_general_d/references/reference_values.json} (evaluated at 50
 digits) and against the five-fixed-node finite-sample formula (60 digits; not
 independent, as above). The values and deviations are in the last column of the table in the paragraph
 \emph{Away from the center} of Section~\ref{gff:mgl:checks}, next to those of the five-Horn formulas.

 Block-level checks with the 90-digit exact jets (value, first and second
 $b$-derivative at $b=\frac d2$; \path{LRI_C5_general_d/Z/check_blocks.py},
 \path{check_formulas.py}, record \path{results/check_formulas.log}): at $d=1$,
 $\mathfrak t_3$ against \eqref{gff:mgz:Dixon1} to $1.0\cdot10^{-91}$, $1.5\cdot10^{-91}$,
 $7.3\cdot10^{-91}$, and $\hat c_0$ (which must vanish) at most $4.1\cdot10^{-90}$,
 $5.9\cdot10^{-89}$, $1.7\cdot10^{-87}$ in absolute value; at $d=2$, $\mathfrak t_3$
 against \eqref{gff:mgz:R2} to $1.3\cdot10^{-92}$, $7.0\cdot10^{-91}$, $8.5\cdot10^{-91}$,
 and $\widehat M_3$ (which must vanish) at most $7.0\cdot10^{-105}$, $9.0\cdot10^{-98}$,
 $8.0\cdot10^{-97}$; at $d=3$, $\mathfrak t_3$ against the values of
 $\mathfrak t_3(3;b)$ and its first two derivatives at $b=\frac32$ from the
 independent $d=3$ code (\path{LRI_C5_d3/wf_Z3/}) to $1.3\cdot10^{-90}$,
 $1.1\cdot10^{-90}$, $1.1\cdot10^{-90}$, and $\hat c_0$ against
 $-N_b^3\pi^5\Gamma(3g-2)/\Gamma(g)^9$ to $1.2\cdot10^{-90}$, $1.8\cdot10^{-90}$,
 $1.4\cdot10^{-90}$; at all $d$, the jets \eqref{gff:mgz:uvjets} against numerical
 differentiation to $5.2\cdot10^{-110}$. Also, at all twelve dimensions of the 90-digit runs, \eqref{gff:mgz:first},
 \eqref{gff:mgz:firstzm}, \eqref{gff:mgz:Rprime}, \eqref{gff:mgz:simplified} and \eqref{gff:gdr:Z} hold
 to about $10^{-73}$ and \eqref{gff:gdr:Q} to $3\cdot10^{-87}$;
 the $d=1$ relation between the two lattices holds to $4.6\cdot10^{-91}$, and the
 $d=3$ identity for $\mathfrak t_3'(3;\frac32)$ to $4.5\cdot10^{-91}$.

 Summary. (i) The two-node identity holds away from $b=\frac d2$, also for complex
 $b$, to $10^{-48}$--$10^{-51}$ against the spectral values and to
 $10^{-55}$--$10^{-63}$ against the fixed-node formula. (ii) $Z_1,Z_2,Z_3$ agree with
 their references to $4\cdot10^{-91}$ with the exact-jet engine at 90 digits, and to
 about $10^{-60}$ at 60 digits. (iii) At the five dimensions with spectral references
 the agreement is $10^{-45}$--$10^{-46}$, the stated precision of those references;
 at the five dimensions of the new spectral runs it is $10^{-39}$--$10^{-40}$, the
 precision of those runs. (iv) The first rim derivative computed from the same jets
 vanishes to $10^{-70}$--$10^{-98}$ (exact jets, 90 digits) and to $4\cdot10^{-43}$
 (Cauchy engine). (v) The exact-jet values at 60 and 90 working digits agree to
 $10^{-53}$--$10^{-61}$ at 12 dimensions, and the Cauchy-circle engine (64 points;
 its error is the aliasing of the circle) agrees with the 90-digit jets to
 $10^{-42}$--$10^{-55}$ at 10 dimensions. (vi) The block-level identities above hold
 to about $10^{-90}$. The Cauchy circles in $b$ use the same blocks and are a
 consistency check, not an independent evaluation. The separate review
 (\path{LRI_C5_general_d/review/Z-rigor/}, own codes) found in addition: the direct
 midpoint-lattice sum at $d=1$ against
 $8\pi^3s_b^3\mathfrak t_3/[\varsigma_b(1+\varsigma_b)^3]$ to
 $10^{-37}$--$10^{-42}$ at $b=0.2,0.3,0.5$; $\mathfrak t_3(2;b)$ against
 \eqref{gff:mgz:R2} to about $10^{-30}$ at five points including complex $b$, and the
 reflection identities behind $\widehat M_3(2;b)=0$ to about $10^{-30}$;
 at $d=3$ an own evaluation of $\mathfrak t_3$ and $\widehat M_3$ reproducing the Gamma
 form of $\hat c_0$ to $3\cdot10^{-30}$ ($b=1.2$) and $3\cdot10^{-28}$
 ($b=1.45+0.05\ii$). That review did not check the resonant dimensions
 $d=\frac43,\frac83$, where the only reference is the spectral code. All these
 numbers are numerical evidence from asymptotic-series evaluations, not interval
 bounds; together with the top-mode tests above they also support
 \eqref{gff:gdr:topless} at $r=3$, and at $d=2$ the identity \eqref{gff:mgz:W3d2} is an exact
 check of it.

\end{itemize}

\subsection{Status, limits and open points}\label{gff:gdr:status}\label{gff:gdr:open}
\begin{itemize}
\item \emph{Proved, with written proofs:} the finite Fourier theorem (Section~\ref{gff:gdr:finite}), the top-mode theorem
 \eqref{gff:gdr:topless} and its product form \eqref{gff:gdr:topG}, the residue functional \eqref{gff:gdr:functional} and the block
 decomposition; the identities
 \eqref{gff:gdr:d12lattice}; the sample formulas of Section~\ref{gff:gdr:samples} (Section~\ref{gff:mgl:section}, \eqref{gff:mgl:tauform}--\eqref{gff:mgl:samph}, with the derivative series \eqref{gff:mgl:dseries}; for the three-wheel, \eqref{gff:mgz:node0} and \eqref{gff:mgz:nodeb} of Section~\ref{gff:mgz:nodes}) and the analyticity of the samples in $b$.
 Checks: a separate review with its own codes (\path{LRI_C5_general_d/review/VERIFICATION.md}, proof review in
 \path{LRI_C5_general_d/review/foundation-rigor}) rederived by hand the proofs of Section~\ref{gff:gdr:periodization} in
 the form they had before they were merged into this note: the reflected forms and the estimates with unspecified
 constants, the finite Fourier theorem, the top-mode proof in its earlier wheel-only form, the residue functional,
 the block decomposition and the fixed-node formula. It also rederived the identities \eqref{gff:gdr:d12lattice}; the
 proof of the $d=2$ identity printed here was written after it. It did not cover the explicit constants of
 \eqref{gff:mgf:gammaBounds}, \eqref{gff:mgf:multiplicityBound} and \eqref{gff:mgf:latticeSums}, which come from
 Section~\ref{gff:gdk:growthProof} and have had no independent review, nor the contour in the proof of
 \eqref{gff:gdr:functional}, which was changed at the merge (paragraph \emph{The residue functional in full} of
 Section~\ref{gff:gdr:residue}), nor the derivation of the $d=3$ identities in Section~\ref{gff:gdr:halfShift}. Within the
 review a second reader checked parts of the top-mode proof, and a 50-digit
 code of its own (\path{LRI_C5_general_d/review/topmode-numeric}) tested the statements. The review found no printed
 formula wrong and one wrong
 intermediate estimate (the endpoint exponent in Steps~3--4 of the top-mode proof); Steps~3--4 as printed
 here use exactly the hypothesis $b<b_{\max}$. A second separate check of the proof in product form, which contains the wheel case,
 is in \path{LRI_C5_general_d/review/lemmaG-proof}; it also tested the endpoint exponents of \eqref{gff:gdr:endpoint}
 numerically.
 Beyond the wheel case, the product form needs the product estimates with $\varkappa=\sum_i\operatorname{Re}g_i$, the endpoint bound in product
 form, the kernel identity \eqref{gff:gdr:kernelBalance} with the exchange on $\widetilde\Delta$, the conjugation step for
 the second top mode and the continuation in the exponents. The check in
 \path{LRI_C5_general_d/review/lemmaG-proof} rederived every step by hand and tested Steps~1--6 numerically in
 unbalanced cases; it found no fatal or major issue and seven minor wording and documentation items, which are applied here. A separate 48-digit lattice-sum code
 (\path{LRI_C5_general_d/review/lemmaG-numeric}; Section~\ref{gff:gdr:checks}) confirms \eqref{gff:gdr:topG} numerically. The check in \path{LRI_C5_general_d/review/lemmaG-proof} covered the statements
 \eqref{gff:gdr:topG} and \eqref{gff:gdr:topless}, including the $R$-family parameters used in
 Section~\ref{gff:gdk:topmode} and the equivalence of (C) with $\tau\in\mathbb H_q$ there; Section~\ref{gff:gdk:topmode} uses these statements. An earlier separate proof of the case $r=3$,
 written for the $\KK_d$ calculation, is superseded by the proof printed here; where the two differ, that review found
 the earlier threshold $|\operatorname{Im}t|\ge2$ insufficient for complex parameters, which is why the proof here uses
 $Y_0$. In the case $r=3$, Step~0 of the proof gives parts (a) (sufficiency), (b) and (c) of the growth theorem of
 Section~\ref{gff:gdk:Rfamily}; the sharpness of $\mathbb H_q$ and the explicit constants are proved in
 Section~\ref{gff:gdk:growthProof}, which that review did not cover. The use of \eqref{gff:gdr:topG} for the $R$ family of $\KK_d$ at continued points also needs the continuation of
 the regularized samples in $\tau$ (hypothesis (H) of Section~\ref{gff:gdk:continuation}). The reviews
 \path{LRI_C5_general_d/review/L-rigor/} and \path{LRI_C5_general_d/review/Z-rigor/} read the derivations of
 Sections~\ref{gff:mgl:section} and \ref{gff:mgz:proofs} in the package's own text. Details that this note writes out where
 that text only states them, for instance the majorants in the proof of \eqref{gff:mgl:dseries} and the regularity of the
 middle and left parts at $t=0$ in the proof of \eqref{gff:mgz:node0}, were added at the merge and have not been reviewed.
\item \emph{Derived}, meaning exact algebra on the above: \eqref{gff:gdr:threenode}, \eqref{gff:gdr:P4}--\eqref{gff:gdr:Z},
 \eqref{gff:gdr:taylorW}--\eqref{gff:gdr:taylorLZ}, the specializations of Section~\ref{gff:gdr:special} at $d=1,2,3$ (at $d=3$ with the
 identities \eqref{gff:gdr:halfShiftC} and \eqref{gff:gdr:halfShiftGamma} of Section~\ref{gff:gdr:halfShift}), and
 \eqref{gff:gdr:C5reduced}. All were rederived symbolically by the review of \path{LRI_C5_general_d/review/VERIFICATION.md},
 the $d=3$ specializations with the two identities as inputs, and are confirmed numerically (Section~\ref{gff:gdr:checks});
 \eqref{gff:gdr:taylorLZ} only through $b$-derivatives of \eqref{gff:gdr:taylorW}, not by literal summation.
 The explicit forms \eqref{gff:mgl:W4five}, \eqref{gff:mgl:W3three}, \eqref{gff:mgl:P4five} and \eqref{gff:mgl:Lfinal} of
 Section~\ref{gff:mgl:section} are derived in the same sense; their coefficients were also compared with direct symbolic
 differentiation (Section~\ref{gff:mgl:section}).
\item \emph{Fixed-node form and continuation.} The fixed-node formula \eqref{gff:mgf:Wsample} and the jets
 \eqref{gff:mgf:jets} are exact algebra on \eqref{gff:gdr:functional}, confirmed numerically
 (Section~\ref{gff:gdr:checks}); Sections~\ref{gff:gdr:samples}--\ref{gff:gdr:C5} do not use them. The continuation
 of the samples beyond $b_{\max}$ (Section~\ref{gff:mgf:nodes}) is derived as the continuation of each block
 value at a fixed node; the residue \eqref{gff:mgf:bmaxres} is an endpoint asymptotic derivation made in the
 review \path{LRI_C5_general_d/review/foundation-rigor}, not a written proof, and it is confirmed numerically and not
 used downstream.
\item \emph{Three-wheel, Section~\ref{gff:mgz:proofs}.} Proved here, given the theorems of
 Sections~\ref{gff:gdr:finite}--\ref{gff:gdr:residue}: the analyticity statement of
 Section~\ref{gff:mgz:uniform}, the two samples \eqref{gff:mgz:node0}, \eqref{gff:mgz:nodeb},
 the two-node formula \eqref{gff:gdr:twonode} for every $b$ in the strip and every
 $0<d<4$, and its zero-mode form \eqref{gff:mgz:zeromode}. Derived (exact
 differentiation and algebra, rederived symbolically by the review
 \path{LRI_C5_general_d/review/Z-rigor/}):
 \eqref{gff:mgz:Zjets}, \eqref{gff:mgz:series}, \eqref{gff:mgz:Zseries2}, the equivalence of
 $\partial_bW_3(d;\frac d2)=0$ with \eqref{gff:mgz:first} and with \eqref{gff:mgz:firstzm},
 \eqref{gff:mgz:Rprime} and \eqref{gff:mgz:simplified}; of these only \eqref{gff:mgz:Rprime} and
 \eqref{gff:mgz:simplified}, hence \eqref{gff:gdr:Z}, use $\partial_bW_3(d;\frac d2)=0$.
 Proved here at $d=1,2$, from Dixon's theorem and without Pattern~T:
 \eqref{gff:mgz:Dixon1}, \eqref{gff:mgz:Z1}, $\widehat M_3(2;b)=0$, \eqref{gff:mgz:R2},
 \eqref{gff:mgz:W3d2} and \eqref{gff:mgz:Z2}. At $d=3$,
 \eqref{gff:mgz:W3d3} and \eqref{gff:mgz:Z3} follow with the Gamma form of
 $\hat c_0(3;b)$, proved in Section~\ref{gff:gdr:halfShift}, whose printed derivation has not been reviewed. Open: whether $\mathfrak t_3''$, $\widehat M_3'$, $\widehat M_3''$
 (equivalently the zero-mode jets) reduce further at generic $d$.

\item \emph{Conditional inputs.} \eqref{gff:gdr:Z} and the $Z_d$ part of \eqref{gff:gdr:C5reduced} use
 $\partial_bW_3(d;d/2)=0$. This input is proved, by Pattern~T \eqref{gff:gs:Tjet} with \eqref{gff:gs:wheelJets}; it is listed here only to record that \eqref{gff:gdr:Z} depends on it, while \eqref{gff:gdr:Q} and the full jet formula $9Z_d=(u\mathfrak t_3+v\widehat M_3)''$ do not. The $d=3$ specializations
 of $Z_d$ and $\LL_d$ use the Gamma form of $\hat c_0$ and the identity $c_{4,1}(b)=2\mathfrak a_4(b)\mathfrak b_4(b)$.
 Both are proved in Section~\ref{gff:gdr:halfShift}, which repeats the half-shift proof of
 Section~\ref{gff:d3:determinantProof} at a general rim exponent; separate AI sessions re-derived both statements, and
 the printed derivation has not been reviewed. They are confirmed numerically: the Gamma form also at $b=0.8$ and complex $b$,
 and the $d=3$ dictionary to $1.2\cdot10^{-51}$ (Section~\ref{gff:gdr:checks}). If $\KK_d$ is taken from
 \eqref{gff:gdk:Kr3}, \eqref{gff:gdr:C5reduced} is conditional on (M), and on (H) when $\mathsf r_3$ is sampled on the jet
 circle (Section~\ref{gff:gdk:status}); with \eqref{gff:gdk:KMB} it is not.
\item \emph{Numerical:} all values of Section~\ref{gff:gdr:checks}, the evaluation of the Horn values and their
 derivatives (asymptotic series, no interval bounds), and every statement that a further reduction does not exist
 (Section~\ref{gff:gdr:attempts}; PSLQ bounds in stated bases, not independence proofs).
\item \emph{Ranges.} $\frac12<d<4$ for $P_4,\LL_d$ and $\frac23<d<4$ for $Q_d,Z_d$, where $b=d/2$ lies in the strip;
 the continuation below these thresholds was not tested. At $d=\frac43,\frac83$ the accelerated evaluation of the
 $Z_d$ jets has removable resonances and uses Cauchy circles in $b$. The continuation of the samples beyond
 $b_{\max}$ for $d>2$ (a simple pole at $b_{\max}$ whose residue the functional annihilates; Section~\ref{gff:mgf:nodes}) is not used anywhere. For $C_5(d)$ the range is $\frac23<d<4$.
\item \emph{Outside this section:} $\KK_d$ (Section~\ref{gff:gdk:section}, status in Section~\ref{gff:gdk:status}) and
 $P_5(d)$. Since $P_5(d)$ was evaluated to high precision only at $d=1,2,3$ (in double precision also at $d=2.4$ and
 $\frac{13}5$, Sections~\ref{gff:mgc:fpChecks} and \ref{gff:mgc:gcChecks}), $C_5(d)$ itself was evaluated only at $d=1,2,3$
 ($C_5(1)=-731.139476624\ldots$, $C_5(2)=-38.8336458962\ldots$, $C_5(3)=-5.85035142075\ldots$; full values in
 Section~\ref{gff:gdr:checks}); the reduced block of \eqref{gff:gdr:C5reduced} without $P_5,\KK_d$ was evaluated at
 eight dimensions.
\item \emph{Open:} whether $\mathfrak D_L$, $\hat c_0$, $\hat c_0'$ and $\frac19(u\mathfrak t_3''+v\widehat M_3'')$
 reduce further at generic $d$, as they do at $d=1,2$ and partly at $d=3$.
\end{itemize}

\section{The nonprimitive master \texorpdfstring{$\KK_d$}{Kd} at generic dimension}\label{gff:gdk:section}

This section evaluates the nonprimitive master $\KK_d$ of \eqref{gff:gs:C5} for generic $0<d<4$ (for $d\le\frac47$
conditional on the separating-contour construction, Section~\ref{gff:gdk:conventions}). For $\KK_d$ it gives, in place of
the 386 Horn values of \eqref{gff:gc:finiteClosure} with up to six variables each, a low-order parameter jet of
two-variable functions, obtained by periodizing the family $R_\tau$ of \eqref{gff:gs:RS} as in
Sections~\ref{gff:app:generalPeriodization} and \ref{gff:gdr:periodization}. There are three steps.
\begin{enumerate}
\item The $S$ family of \eqref{gff:gs:RS} is elementary in every dimension (Section~\ref{gff:gdk:Sfamily}).
\item Given the inputs (M) and (J) of Section~\ref{gff:gdk:conventions}, the Gamma subtraction of \eqref{gff:gs:Kdefinition}
fixes the $R$ family through order $\tau^2$. The only
non-elementary number in $\KK_d$ is then the Laurent coefficient $\mathsf r_3=[\tau^3]\langle t_\ell^2R_\tau\rangle_d$
(Section~\ref{gff:gdk:lowOrders}).
\item $\langle t_\ell^2R_\tau\rangle_d$ is a three-node sample of an entire periodic function with the Fourier
modes $|k|\le1$, on the half-plane $\mathbb H_q$ of \eqref{gff:gdk:domain} with $\operatorname{Re}\tau<\frac32q$, $\tau\ne q$, and, by the continuation
hypothesis (H) of Section~\ref{gff:gdk:continuation}, on the circle that gives the jet. Each sample consists of three
two-variable Horn functions at argument $(1,1)$ (Sections~\ref{gff:gdk:Rfamily}--\ref{gff:gdk:continuation}).
\end{enumerate}
By a second route, which uses none of these inputs, $\KK_d$ is a polygamma part plus one two-fold Barnes integral
(Section~\ref{gff:gdk:positionSpace}). The specializations to $d=1,2,3$ reproduce \eqref{gff:eq:KprovedFinal},
\eqref{gff:gc:K2closed} and the $430$-digit \eqref{gff:gs:Kseries} value of $\KK_3$ (Section~\ref{gff:gdk:special}).
Section~\ref{gff:gdk:numerics} collects the numerical checks, and Section~\ref{gff:gdk:status} states what is proved, what is
conditional and what is numerical. Codes and logs are in the package folder \path{LRI_C5_general_d/K_general/}.
The material of this section is this project's own work, done in Claude Code sessions in that folder; the proof of
the growth theorem in Section~\ref{gff:gdk:growthProof} was written there as a separate proof file and is merged here in
full.

\subsection{Conventions and two inputs from the master path}\label{gff:gdk:conventions}
Put $q=d/4$. We use $A,R,U,V,X$ of \eqref{gff:gs:psiBasis}, $c_*=\psi_0(1)+\psi_0(2q)$ of \eqref{gff:gs:Kweights},
$C(\tau)$, $R_\tau$, $S_\tau$ of \eqref{gff:gs:RS}, and $d_\ell,a_\ell,b_\ell$ of \eqref{gff:gs:spectrum}. For a rim exponent
$b$ (not to be confused with $b_\ell$), $g_\ell(s;b)$ and $N_b$ are those of \eqref{gff:gdr:kernel}; we write
$a_\ell(b)=\ell/2+b/4$ and $g_b=(d-b)/2$ for the $a_\ell$ and $g$ there (so $a_\ell(b)+g_b$ is its $\beta_\ell$), and keep
$a_\ell,b_\ell$ for their values at $b=d/2$, as in \eqref{gff:gs:spectrum}. For real $b$, $t_{\ell,b}(\nu)=N_b\,g_\ell(\ii\nu;b)$
in \eqref{gff:gs:balancedWheel}; for complex $b$, $t_{\ell,b}$ and $\langle t_{\ell,b}^r\rangle_d$ denote the analytic
continuation \eqref{gff:gdr:kernel}, in which $|\cdot|^2$ of \eqref{gff:gs:balancedWheel} is replaced by
$\Gamma(a+\ii\nu)\Gamma(a-\ii\nu)/(\Gamma(\beta+\ii\nu)\Gamma(\beta-\ii\nu))$ with $a=a_\ell(b)$, $\beta=a_\ell(b)+g_b$.
Further, $g_\ell(s;d/2)=g_\ell(s)$ is the kernel of \eqref{gff:fp:kernel}. We write $t_\ell(s)=\Gamma(2q)g_\ell(s)$, so that
$t_\ell(\ii\nu)$ is \eqref{gff:gs:spectrum}. The
kernel $R_\tau$ is not even in $s$, so every trace is read on the full Mellin line,
\begin{equation}
 \langle t_\ell^2R_\tau\rangle_d=\Gamma(2q)^2\sum_{\ell\ge0}d_\ell\int_{-\ii\infty}^{\ii\infty}\frac{ds}{2\pi\ii}\,
 g_\ell(s)^2R_\tau(s).
 \label{gff:gdk:fullTrace}
\end{equation}
For real $\tau$ the integrand of \eqref{gff:gdk:fullTrace} takes complex-conjugate values at $s=\pm\ii\nu$, so this is
the real-part reading of \eqref{gff:gs:trace} used in \eqref{gff:gs:Kspectral} and in
\eqref{gff:mgk:mixedTrace} (Section~\ref{gff:mgk:regulated}); for complex $\tau$ it is the continuation of that trace, in the sense fixed in (M) below.
For even kernels such as
$t_\ell^2S_\tau$ it coincides with \eqref{gff:gs:trace}. Put
\begin{equation}\begin{gathered}
 \mathcal X_d(\tau)=\Bigl[\frac{\Gamma(q-\tau/2)}{\Gamma(q+\tau/2)}\Bigr]^2,\qquad
 F(\tau)=\langle t_\ell^2(R_\tau-S_\tau)\rangle_d=\sum_kF_k\tau^k,\\
 \mathsf r_k=[\tau^k]\langle t_\ell^2R_\tau\rangle_d,\qquad \mathsf s_k=[\tau^k]\langle t_\ell^2S_\tau\rangle_d .
\end{gathered}\label{gff:gdk:notation}\end{equation}
The Laurent coefficients $\mathsf r_k$ are not the polygamma combinations $r_j$ of \eqref{gff:gs:rj}. In this section
$C(\tau)$ is that of \eqref{gff:gs:RS} and is written with its argument; a bare $C$, also with a subscript such as $C_\ell$, is a generic
constant; $C_5(d)$ is the coefficient of \eqref{gff:gs:C5}. The constants $S$, $C$, $D$ of \eqref{gff:gdr:D} are not used, and $\mathcal X_d(\tau)$ is unrelated to the samples
$\mathcal X_0,\mathcal X_{1/2},\mathcal X_T$ of Section~\ref{gff:gdr:samples}.

Two statements of Section~\ref{gff:app:generalSpecial} about the master path $J_t$ of \eqref{gff:gs:Kdefinition} enter
below.
\begin{itemize}
\item[(M)] \emph{Mixed trace.} $J_{2\tau}=C(\tau)\langle t_\ell^2R_\tau\rangle_d$, as stated before
\eqref{gff:gs:Ksubtraction}, as an identity of meromorphic functions of $\tau$. The expression \eqref{gff:gdk:fullTrace}
converges absolutely, with the Mellin line separating the left and right pole families of every spin, only for
$\max(0,\frac12-2q)<\operatorname{Re}\tau<\frac32q$. Outside this strip, and in particular near $\tau=0$,
$\langle t_\ell^2R_\tau\rangle_d$ denotes its meromorphic continuation, which by (M) is $J_{2\tau}/C(\tau)$; the
$\mathsf r_k$ of \eqref{gff:gdk:notation} are its Laurent coefficients at $\tau=0$. The strip is nonempty iff
$d>\frac47$. For $d\le\frac47$ there is no strip to continue from: the trace must then be defined with contours that
separate the pole families for finitely many spins, which is only sketched (Section~\ref{gff:gdk:threeNode}), and the
statements below that involve $\langle t_\ell^2R_\tau\rangle_d$ or $\mathsf r_k$ are conditional on that
construction. (M) is proved at $d=2$ for
$0<\tau<\frac12$ (\eqref{gff:mgk:mixedTrace} in Section~\ref{gff:mgk:regulated}). At other $d$ it is confirmed numerically (Section~\ref{gff:gdk:numerics}).
\item[(J)] \emph{Jet.} $J_t-J_t^\Gamma=\KK_dt^2+O(t^3)$ on the path of \eqref{gff:gs:Kdefinition}. This is
\eqref{gff:gs:Bjet} restricted to that path, and at $d=1$ it is \eqref{gff:eq:Kpath}. Section~\ref{gff:app:generalSpecial} derives \eqref{gff:gs:Bjet}
from completion symmetry and analyticity; it is listed here as an input because its
lower orders $t^0,t^1$ on the path are used explicitly, in \eqref{gff:gdk:rLow} and \eqref{gff:gdk:MBidentities}. At $d=2$ its low
orders are rederived in Section~\ref{gff:gc:K2laurent}, and at non-integer $d$ it is confirmed numerically
(Section~\ref{gff:gdk:numerics}).
\end{itemize}

\subsection{\texorpdfstring{The $S$ family is elementary}{The S family is elementary}}\label{gff:gdk:Sfamily}
\paragraph{Closed form of the $S$ family.}
Let $0<d<4$ and $\max(0,\frac12-2q)<\operatorname{Re}\tau<2q$; this strip is nonempty for $d>\frac12$. Then
\begin{equation}\boxed{
 \langle t_\ell^2S_\tau\rangle_d=C(\tau)\,\mathcal X_d(\tau)^2
 =\frac{\Gamma(2q)\,\Gamma(\tau)\,\Gamma(q-\tau/2)^4}{\Gamma(2q-\tau)\,\Gamma(q+\tau/2)^4},}
 \label{gff:gdk:Sclosed}
\end{equation}
and everywhere else $\langle t_\ell^2S_\tau\rangle_d$ denotes the meromorphic continuation of the right side in
$\tau$ and $d$. Equivalently $C(\tau)\langle t_\ell^2S_\tau\rangle_d=C(\tau)^2\mathcal X_d(\tau)^2$, the first term of
\eqref{gff:gs:Ksubtraction}; \eqref{gff:gdk:Sclosed} makes that term an explicit statement about the $S$ family.

\emph{Proof.} Put $b'=d/2-\tau$. Then $a_\ell(b')=a_\ell-\tau/4$, $a_\ell(b')+g_{b'}=b_\ell+\tau/4$ and
$N_{b'}=\Gamma(2q)\Gamma(q+\tau/2)/\Gamma(q-\tau/2)$, so \eqref{gff:gs:RS} gives the kernel identity
\begin{equation}
 t_\ell(s)^2S_\tau(s)=\Gamma(2q)^2g_\ell(s;b')^2=\mathcal X_d(\tau)\,\bigl[N_{b'}\,g_\ell(s;b')\bigr]^2 .
 \label{gff:gdk:Skernel}
\end{equation}
Hence $\langle t_\ell^2S_\tau\rangle_d=\mathcal X_d(\tau)\langle t_{\ell,b'}^2\rangle_d=\mathcal X_d(\tau)W_2(d;b')$,
the balanced two-wheel of \eqref{gff:gs:balancedWheel}. A two-wheel is a single convolution, so
$W_2(d;b')=B_d(d-b',2b')$ by \eqref{gff:gs:convolution}; this is the identity used before \eqref{gff:gs:Ksubtraction}, and
spectrally it is Parseval's identity for the dilatation-invariant operator with eigenvalue $t_{\ell,b'}$. For complex
$b'$ it extends because both sides are holomorphic in $\tau$ on the strip, the trace by locally uniform absolute
convergence. With
$(d-b',2b')=(2q+\tau,4q-2\tau)$, \eqref{gff:gs:convolution} gives $B_d=C(\tau)\mathcal X_d(\tau)$. Finally,
$|g_\ell(\ii\nu;b')|\asymp(1+\ell+|\nu|)^{-2\operatorname{Re}g_{b'}}$ with $g_{b'}=q+\tau/2$, so the $\nu$-integral and
the spin sum converge absolutely iff $4\operatorname{Re}g_{b'}>\max(d,1)$; together with $\operatorname{Re}b'>0$, which
keeps the $\ell=0$ contour between its two pole families, this is the stated strip. $\square$

\paragraph{Laurent form.}
Expanding $\log\Gamma$ gives $\log[\tau\langle t_\ell^2S_\tau\rangle_d]=\sum_{h\ge1}x_h\tau^h/h!$ with
\begin{equation}\begin{gathered}
 x_h=\psi_{h-1}(1)-(-1)^h\psi_{h-1}(2q)-[h\ \text{odd}]\,2^{3-h}\psi_{h-1}(q),\\
 (x_1,x_2,x_3,x_4)=\bigl(c_*-4\psi_0(q),\,-R,\,V-U,\,-X\bigr).
\end{gathered}\label{gff:gdk:xh}\end{equation}
Thus $\mathsf s_{m-1}=[u^m]\exp\sum_hx_hu^h/h!$, which is the polynomial $E_m(x_1,\dots,x_m)$ of
\eqref{gff:fp:BellFinite}:
\begin{equation}\begin{gathered}
 \mathsf s_{-1}=1,\qquad \mathsf s_0=x_1,\qquad \mathsf s_1=\tfrac12(x_1^2+x_2),\qquad
 \mathsf s_2=\tfrac16(x_1^3+3x_1x_2+x_3),\\
 \mathsf s_3=\tfrac1{24}(x_1^4+6x_1^2x_2+4x_1x_3+3x_2^2+x_4).
\end{gathered}\label{gff:gdk:sk}\end{equation}

\paragraph{Remarks.}
(i) At $d=2$, \eqref{gff:gdk:Sclosed} is the value of $\langle t_\ell^2S_\tau\rangle_2$ obtained in
\eqref{gff:mgk:regulatedEq} (Section~\ref{gff:mgk:regulated}) from Gauss's sum and the bilateral (Dougall) identity of
\eqref{gff:mgk:periodOne} on the square lattice. At generic $d$ no Dixon, Dougall or
Thomae evaluation is needed. (ii) By \eqref{gff:gdk:Skernel} the spin-summed periodization of $t_\ell^2S_\tau$ is a
balanced two-wheel of rim exponent $d/2-\tau$. For $\tau\in\mathbb H_q$ of \eqref{gff:gdk:domain} its $\Xi$ is constant
in $t$, by the case $r=2$ of \eqref{gff:gdr:topG} (here $r$ counts the factors of the product; it is unrelated to the Laurent coefficients $\mathsf r_k$) with $\alpha_i=\beta_i=-\tau/4$ and $g_i=q+\tau/2$, that is
\eqref{gff:gdr:mixed} with $b_i=d/2-\tau$; \eqref{gff:gdk:Sclosed} does not use this. (iii) For a rim exponent $b$ with
$0<\operatorname{Re}b<b_{\max}(2,d)$ put $\sigma=b/2$, as in Section~\ref{gff:gdr:section}, so that $g_b=2q-\sigma$. By the
case $r=2$ of \eqref{gff:gdr:topless}, $\Xi_2(t;b)=c_{2,0}(q,b)$ does not depend on $t$. The residue functional
\eqref{gff:gdr:functional} at $r=2$, where $B_{2,0}(\sigma)=-2\pi\cot\pi\sigma$, together with $W_2(d;b)=B_d(d-b,2b)$ from
\eqref{gff:gs:convolution} and $N_b$ of \eqref{gff:gdr:kernel}, gives the constant:
\begin{equation}
 \Xi_2(t;b)=c_{2,0}(q,b)=\frac{\pi\sin^3(\pi\sigma)\,W_2(d;b)}{2\cos(\pi\sigma)\,N_b^2}
 =\frac{\pi\sin^3(\pi\sigma)\,\Gamma(\sigma)^4\,\Gamma(2q-2\sigma)}{2\cos(\pi\sigma)\,\Gamma(2q)\,\Gamma(g_b)^4\,\Gamma(2\sigma)} .
 \label{gff:gdk:Sconstant}
\end{equation}
At $b=d/2-\tau$, that is $\sigma=q-\frac\tau2$, with $\tau\in\mathbb H_q$ and $\operatorname{Re}\tau<2q$, this reads
$\Gamma(2q)^2\Xi_2=\pi\sin^3(\pi\sigma)\langle t_\ell^2S_\tau\rangle_d/(2\cos\pi\sigma)$, by \eqref{gff:gdk:Sclosed}. With the
Gamma-reflected form of $g_\ell(a_\ell+n+t)$ in \eqref{gff:gdr:reflected}, $D_b^2S_2=c_{2,0}$ becomes a bilateral identity,
for $t\notin\mathbb Z\cup(-\sigma+\mathbb Z)\cup(g_b+\mathbb Z)$:
\begin{equation}
 \sum_{\ell\ge0}d_\ell\sum_{n\in\mathbb Z}
 \Bigl[\frac{\Gamma(\ell+n+\sigma+t)}{\Gamma(\ell+n+2q+t)}\Bigr]^2\Bigl[\frac{\Gamma(n+1-g_b+t)}{\Gamma(n+1+t)}\Bigr]^2
 =\frac{c_{2,0}(q,b)}{\sin^2\pi(t+\sigma)\,\sin^2\pi(g_b-t)} .
 \label{gff:gdk:Sbilateral}
\end{equation}
At $d=2$ the left side is the square of a single bilateral sum, and \eqref{gff:gdk:Sbilateral} is the square of Dougall's
${}_2H_2$ sum, an instance of which is the last identity of \eqref{gff:mgk:periodOne} (with $x_\tau=\sigma$ and $\eta=t$);
at generic $d$ it is its analogue for the multiplicities $d_\ell$. The package record
\path{LRI_C5_general_d/K_general/S_part/onemode.jsonl} (code \path{onemode_check.py} there) confirms the constancy of
$\Xi_2$ at $b=d/2-\tau$ and the value \eqref{gff:gdk:Sconstant}, at four nodes each, to $4\cdot10^{-17}$ at seven points
$0.8\le d\le3.5$ and to $2\cdot10^{-39}$ at two points at $d=1$.

\subsection{\texorpdfstring{Low orders of the $R$ family and the single new number}{Low orders of the R family and the single new number}}\label{gff:gdk:lowOrders}
At $t=2\tau$, \eqref{gff:gs:convolution} gives $B_d(2q,4q-2\tau)=C(\tau)\Gamma(q-\tau)/\Gamma(q+\tau)$ and
$B_d(2q+4\tau,4q-2\tau)=C(\tau)\Gamma(q-2\tau)\Gamma(q+\tau)/[\Gamma(q+2\tau)\Gamma(q-\tau)]$, so the Gamma
interpolation of \eqref{gff:gs:Kdefinition} is
\begin{equation}
 J^\Gamma_{2\tau}=C(\tau)^2\frac{\Gamma(q-\tau)^2}{\Gamma(q+\tau)^2}
 +\frac{C(\tau)}\tau\Bigl[\frac{\Gamma(q-\tau)}{\Gamma(q+\tau)}
 -\frac{\Gamma(q-2\tau)\Gamma(q+\tau)}{\Gamma(q+2\tau)\Gamma(q-\tau)}\Bigr].
 \label{gff:gdk:JGamma}
\end{equation}
Expanding $\log\Gamma$ gives the exact Laurent series
\begin{equation}
 \Delta(\tau):=J^\Gamma_{2\tau}-C(\tau)^2\mathcal X_d(\tau)^2=\frac32\,U\tau+AU\tau^2+O(\tau^3).
 \label{gff:gdk:Delta}
\end{equation}
The orders $\tau^{-2},\tau^{-1},\tau^0$ cancel identically, and $[\tau^2]\Delta=AU$ is $-4$ times the difference
$-AU/4$ of the $t^2$ coefficients stated after \eqref{gff:gs:Ksubtraction}. A symbolic check for general $q$ is
\path{LRI_C5_general_d/K_general/note_fragment/fragment_checks/check_delta_symbolic.wls}.

\paragraph{Low orders.}
Assume (M) and (J). Then $F(\tau)=\frac32U\tau^2+F_3\tau^3+O(\tau^4)$; that is,
\begin{equation}
 \mathsf r_{-1}=1,\qquad \mathsf r_0=\mathsf s_0,\qquad \mathsf r_1=\mathsf s_1,\qquad
 \mathsf r_2=\mathsf s_2+\tfrac32U,
 \label{gff:gdk:rLow}
\end{equation}
and
\begin{equation}\boxed{
 \KK_d=\frac U8\bigl[\psi_0(1)+\psi_0(2q)+4\psi_0(q)\bigr]+\frac{\mathsf r_3-\mathsf s_3}4 ,}
 \label{gff:gdk:Kr3}
\end{equation}
with $\mathsf s_3$ the polygamma polynomial \eqref{gff:gdk:sk}. Equivalently,
\begin{equation}
 \KK_d=-\frac{AU}4+\frac14[\tau^2]\Bigl\{C(\tau)\langle t_\ell^2R_\tau\rangle_d-C(\tau)^2\mathcal X_d(\tau)^2\Bigr\}.
 \label{gff:gdk:Ktau}
\end{equation}
Equation \eqref{gff:gdk:Ktau} uses only (M), \eqref{gff:gdk:Delta} and the definition \eqref{gff:gs:Kdefinition}, since
$[\tau^2](J_{2\tau}-J^\Gamma_{2\tau})=4\KK_d$ for the Laurent coefficients whatever the lower orders are; here (J)
enters only \eqref{gff:gdk:rLow} and \eqref{gff:gdk:Kr3}.

\emph{Proof.} By (M) and \eqref{gff:gdk:Sclosed}, $J_{2\tau}=C(\tau)^2\mathcal X_d(\tau)^2+C(\tau)F(\tau)$. By (J) at $t=2\tau$,
$J_{2\tau}-J^\Gamma_{2\tau}=4\KK_d\tau^2+O(\tau^3)$. Subtracting and using \eqref{gff:gdk:Delta},
\[
 C(\tau)F(\tau)=\tfrac32U\tau+(AU+4\KK_d)\tau^2+O(\tau^3).
\]
This is \eqref{gff:gdk:Ktau}. Since $C(\tau)=\tau^{-1}+c_*+O(\tau)$, it follows that
$F(\tau)=\frac32U\tau^2+[AU+4\KK_d-\frac32c_*U]\tau^3+O(\tau^4)$, which is \eqref{gff:gdk:rLow}. With
$F_3=\mathsf r_3-\mathsf s_3$ and $A=c_*-2\psi_0(q)$ this gives \eqref{gff:gdk:Kr3}. $\square$

\paragraph{The spectral moment.}
On the Mellin line $\operatorname{Re}(R_\tau-S_\tau)=E_2\tau^2+O(\tau^3)$, with the weight $E_2=-h^2/2-5j/8$ of
\eqref{gff:gs:Kweights}, as displayed after \eqref{gff:gs:Ksubtraction} ($h$, $j$ the weights of \eqref{gff:gs:Kweights}, not the
index $h$ of \eqref{gff:gdk:xh} or the node index $j$ of \eqref{gff:gdk:threeNodeRule}; this $E_2$ is not the polynomial $E_2$ of \eqref{gff:fp:BellFinite}, which gives $\mathsf s_1$ in \eqref{gff:gdk:sk}). If $[\tau^2]$ may be taken under the trace (the
step used in the derivation of \eqref{gff:gs:Kspectral}, justified only at $d=2$, by the corollary of
Section~\ref{gff:mgk:regularization}), then
$F_2=\langle t_\ell^2E_2\rangle_d$, and \eqref{gff:gdk:rLow}, which assumes (M) and (J), gives
\begin{equation}\boxed{
 \langle t_\ell^2E_2\rangle_d=\frac32\,\psi_2\Bigl(\frac d4\Bigr).}
 \label{gff:gdk:E2}
\end{equation}
At $d=2$ this is the value $-21\zeta(3)$ of \eqref{gff:mgk:E2}; by \eqref{gff:eq:quarterPolygamma} it is
$-3\pi^3-84\zeta(3)$ at $d=1$ and $3\pi^3-84\zeta(3)$ at $d=3$.

\paragraph{Consequence.}
Equation \eqref{gff:gdk:Kr3} puts all non-elementary content of $\KK_d$ into the single number $\mathsf r_3$. The
following sections evaluate $\langle t_\ell^2R_\tau\rangle_d$, and hence $\mathsf r_3$, by finite sampling.

\subsection{\texorpdfstring{The $R$ family: lattice form, growth and finite Fourier space}{The R family: lattice form, growth and finite Fourier space}}\label{gff:gdk:Rfamily}
\paragraph{Lattice form.}
Let $f^{\rm R}_\ell(s)=\Gamma(2q)^2g_\ell(s)^2R_\tau(s)$ be the summand of \eqref{gff:gdk:fullTrace}. In Sections~\ref{gff:gdk:Rfamily}--\ref{gff:gdk:special}, $t$ is the periodization variable, not the path parameter of $J_t$. Direct substitution
gives, for $\ell\ge0$ and $n\in\mathbb Z$,
\begin{equation}\begin{gathered}
 f^{\rm R}_\ell(a_\ell+n+t)=\Gamma(2q)^2\,\Phi_t(\ell+n)\,\Psi_t(n),\\
 \Phi_t(m)=\frac{\Gamma(m+q+t)^3}{\Gamma(m+2q+t)^2\Gamma(m+q+\tau+t)},\qquad
 \Psi_t(n)=\frac{\Gamma(-n-t)^2\Gamma(q-\tau-n-t)}{\Gamma(q-n-t)^3}.
\end{gathered}\label{gff:gdk:latticeFactor}\end{equation}
Reflection gives $\Psi_t(n)=T_{\rm R}(t)\Phi_{t+1-2q}(n)$ with
$T_{\rm R}(t)=\sin^3\pi(q-t)/[\sin^2(\pi t)\sin\pi(q-\tau-t)]$. The summand is therefore one Gamma ratio
$\widetilde\Phi(w)=\Gamma(w+q)^3/[\Gamma(w+2q)^2\Gamma(w+q+\tau)]$ on the two lattices $w=m+t$ and $w=n+t+1-d/2$,
with the Toeplitz weight $d_{m-n}$, where $\sum_\ell d_\ell z^\ell=(1+z)(1-z)^{1-d}$. A half-shift as in
\eqref{gff:td:determinants} maps this structure to itself only if $2(1-d/2)\in\mathbb Z$, that is at $d=1,2,3$; at
generic $d$ the half-shift factorization of \eqref{gff:td:determinants} is not available.

\paragraph{Product structure.}
In the notation \eqref{gff:gdr:kernel} the kernel is a product of three balanced factors:
\begin{equation}
 t_\ell(s)^2R_\tau(s)=\Gamma(2q)^2\,g_\ell\bigl(s;\tfrac d2\bigr)^2\,
 g_\ell\Bigl(s+\frac{\tau-q}2;\,d-2\tau\Bigr).
 \label{gff:gdk:Rproduct}
\end{equation}
Indeed $a_\ell(d-2\tau)\pm(s+\frac{\tau-q}2)$ equal $a_\ell+s$ and $b_\ell-\tau-s$, and
$a_\ell(d-2\tau)+\tau\pm(s+\frac{\tau-q}2)$ equal $a_\ell+\tau+s$ and $b_\ell-s$. Each factor has rim plus spoke
exponent $d$, but the third has rim exponent $d-2\tau$ and its own shift. Equation \eqref{gff:gdk:Rproduct} is the
kernel form of the mixed trace of two critical wheel operators and one operator of rim exponent $d-2\tau$
(Section~\ref{gff:gs:spectralSection}), and $C(\tau)=N_{d-2\tau}$.

\paragraph{Periodization.}
Define
\begin{equation}\begin{gathered}
 S_{\rm R}(t;\tau)=\sum_{\ell\ge0}d_\ell\sum_{n\in\mathbb Z}f^{\rm R}_\ell(a_\ell+n+t),\qquad
 \Xi_{\rm R}=D_{\rm R}S_{\rm R},\\
 D_{\rm R}(t)=\sin^2(\pi t)\sin^3\bigl(\pi(t+q)\bigr)\sin\bigl(\pi(t-q+\tau)\bigr).
\end{gathered}\label{gff:gdk:periodization}\end{equation}
The pole classes are $\mathsf P_\tau=\mathbb Z\cup(-q+\mathbb Z)\cup(q-\tau+\mathbb Z)$, of orders at most $2,3,1$,
matching the zeros of $D_{\rm R}$. Put
\begin{equation}\begin{gathered}
 \tau_*(q)=\max\bigl(0,\,2q-1,\,\tfrac12-2q\bigr),\qquad
 \mathbb H_q=\{\tau\in\mathbb C:\operatorname{Re}\tau>\tau_*(q)\},\\
 \eta_{\rm R}=2\operatorname{Re}\tau-\max(0,1-4q).
\end{gathered}\label{gff:gdk:domain}\end{equation}
On $\mathbb H_q$ one has $\eta_{\rm R}>0$.

\paragraph{Growth and finite modes.}
Let $0<d<4$.
\begin{itemize}
\item[(a)] For $t\notin\mathsf P_\tau\cup(q+\mathbb Z)$ the double series $S_{\rm R}$ converges absolutely iff
$\tau\in\mathbb H_q$. For $\tau\in\mathbb H_q$ it converges locally uniformly off $\mathsf P_\tau$, and $S_{\rm R}$ is
meromorphic and one-periodic in $t$, with poles only in $\mathsf P_\tau$ of orders at most those of the zeros of
$D_{\rm R}$.
\item[(b)] For every compact set of pairs $(q,\tau)$ with $\tau\in\mathbb H_q$ there are constants $Y$ and $C$ with
$|S_{\rm R}(x+\ii y)|\le C|y|^{-\eta_{\rm R}}(1+\log|y|)$ for all real $x$ and $|y|\ge Y$. The right part
($n\ge0$), the middle part ($-\ell\le n\le-1$) and the left part ($n\le-\ell-1$) obey this bound separately for
$0\le x\le1$.
\item[(c)] $\Xi_{\rm R}$ is entire and one-periodic in $t$, jointly holomorphic on $\mathbb C\times\mathbb H_q$, and
$\Xi_{\rm R}(t;\tau)=\sum_{|k|\le2}c_k(\tau)e^{2\pi\ii kt}$ with $c_k$ holomorphic on $\mathbb H_q$.
\end{itemize}

\emph{Proof.} Section~\ref{gff:gdk:growthProof} gives the complete proof, with explicit constants. Parts (a)
(sufficiency), (b) and (c) are also the case $r=3$ of Step~0 of the proof of \eqref{gff:gdr:topG}, with the parameters of
Section~\ref{gff:gdk:topmode}, for which $\varkappa=2q+\operatorname{Re}\tau$ in \eqref{gff:gdr:BC} and
$\max(d,1)-2\varkappa=-\eta_{\rm R}$. In outline: Gamma reflection writes each summand of the right, middle and left
parts as a trigonometric factor times two Gamma ratios, each of total exponent $-(2q+\tau)$; a Stirling bound that is uniform
in the closed right half-plane and two lattice-sum estimates give (b) and the convergence in (a); matching lower bounds
give the divergence outside $\mathbb H_q$; and once $D_{\rm R}$ has cleared the poles, (b) removes the modes
$|k|\ge3$. $\square$

The exponent $\eta_{\rm R}$ plays the role of $2q(r-2)$ in \eqref{gff:fp:wheelGrowth}. The estimate holds on the whole
half-plane $\mathbb H_q$; only the residue functional of Section~\ref{gff:gdk:threeNode} needs
$\operatorname{Re}\tau<\frac32q$. Periodicity and analyticity improve (b) to
$S_{\rm R}=O(e^{-2\pi|\operatorname{Im}t|})$ on $\mathbb H_q$ (Section~\ref{gff:gdk:growthProof}), while each part
decays only like a power.

\subsection{The top modes vanish}\label{gff:gdk:topmode}
The kernel \eqref{gff:gdk:Rproduct} is a product of three balanced kernels with different rim exponents and shifts, so the
product form \eqref{gff:gdr:topG} of the top-mode theorem of Section~\ref{gff:gdr:top} applies to it; its wheel case
\eqref{gff:gdr:topless} alone does not.

\paragraph{The $R$ family has the modes $|k|\le1$.}
For $0<d<4$ and $\tau\in\mathbb H_q$, the coefficients $c_{\pm2}$ in part (c) of the growth theorem of
Section~\ref{gff:gdk:Rfamily} vanish:
\begin{equation}\boxed{
 \Xi_{\rm R}(t;\tau)=c_0(\tau)+c_1(\tau)e^{2\pi\ii t}+c_{-1}(\tau)e^{-2\pi\ii t},}
 \label{gff:gdk:threeModes}
\end{equation}
equivalently $S_{\rm R}=O(e^{-4\pi|\operatorname{Im}t|})$.

\emph{Proof.} By \eqref{gff:gdk:Rproduct}, $f^{\rm R}_\ell(a_\ell+n+t)=\Gamma(2q)^2T_\ell(n+t)$ with $T_\ell$ of
\eqref{gff:gdr:TG} for $r=3$, $\sigma=q$, $(g_1,g_2,g_3)=(q,q,\tau)$, $\alpha_i=0$ and
$(\beta_1,\beta_2,\beta_3)=(0,0,q-\tau)$; this is \eqref{gff:gdr:mixed} with $b=(\frac d2,\frac d2,d-2\tau)$ and
$\delta=(0,0,\frac{\tau-q}2)$. Then $D_T$, $S_T$ and $\Xi_T$ of \eqref{gff:gdr:TG} are $D_{\rm R}$,
$\Gamma(2q)^{-2}S_{\rm R}$ and $\Gamma(2q)^{-2}\Xi_{\rm R}$, so $\gamma_m=\Gamma(2q)^{-2}c_m$. Condition (B) of
\eqref{gff:gdr:BC} holds factor by factor, and (C) reads $2q+\operatorname{Re}\tau>\max(2q,4q-1,\frac12)$, that is
$\tau\in\mathbb H_q$. $\square$

A three-wheel has the same exponential modes; there reflection pairs $k=\pm1$ into the single cosine mode $c_{3,1}$,
while $R_\tau$ is not even and $c_{\pm1}$ are independent. For complex $\tau$ no reflection relates $c_2$ and
$c_{-2}$; Step~6 of the proof of \eqref{gff:gdr:topG} treats them separately.

\paragraph{Remarks.}
(i) Read as a wheel with the single pair $(\sigma,g_3)=(q,\tau)$, the third factor is off balance by $\tau-q$; only the
per-factor balance (B) holds, which is why the product form is needed; with its own shift $\beta_3=q-\tau$ the factor
is balanced and \eqref{gff:gdr:kernelBalance} applies. (ii) Per-factor balance is needed: total balance does not suffice,
and the multiplicities must be those of the dimension $d$ in (B), as the following two controls show numerically. At
$d=1.7$, $\tau=0.55+0.3\ii$, a control kernel with right shifts $\beta=(0.15,-0.15,q-\tau)$, which keeps the total
balance, has $|c_{-2}|/\max_k|c_k|=1.4\cdot10^{-2}$, and the $R$ kernel with the multiplicities of $d'=1.9$, for which
the phase $\Theta$ of Step~5 of the proof of \eqref{gff:gdr:topG} obeys $\Theta-\Theta\circ\iota=(d-d')(V-\pi)$, has
$6\cdot10^{-2}$; for the balanced kernel the ratio is below $10^{-15}$, within the spread of the spin-sum
extrapolation. (iii) \eqref{gff:gdr:topG} also contains the $S$ family ($r=2$, rim exponent $d/2-\tau$, by
\eqref{gff:gdk:Skernel}). The consequences of \eqref{gff:gdr:topless} for Section~\ref{gff:app:threeDimensionalReduction}
are stated in Section~\ref{gff:gdr:top}, paragraph ``Relation to the periodization at $b=d/2$ and to the half-shift
reduction''.

\subsection{Proof of the growth theorem}\label{gff:gdk:growthProof}
This subsection proves part (b) of the growth theorem of Section~\ref{gff:gdk:Rfamily} with explicit constants, derives (a)
and (c) from Step~0 of the proof of \eqref{gff:gdr:topG}, and proves the converse in (a), that is the sharpness of
$\mathbb H_q$, and the exponential decay of $S_{\rm R}$. It uses three estimates proved in
Section~\ref{gff:gdr:periodization}: the Gamma-ratio bound \eqref{gff:mgf:gammaBounds}, the multiplicity bound
\eqref{gff:mgf:multiplicityBound} and the lattice sums \eqref{gff:mgf:latticeSums}. Throughout, $0<q<1$, $d=4q$, $t=x+\ii y$
and $\varkappa=2q+\operatorname{Re}\tau$; $\log$, $\arg\in(-\pi,\pi]$ and complex powers are principal; and, besides
the generic constants $C$ of Section~\ref{gff:gdk:conventions}, a bare lowercase $c$ is a positive generic constant. In
terms of $\varkappa$, $\tau\in\mathbb H_q$ means
\begin{equation}
 \varkappa>2q\ \ (\text{that is, }\operatorname{Re}\tau>0),\qquad \varkappa>4q-1,\qquad 2\varkappa>1 ,
 \label{gff:gdk:threeConditions}
\end{equation}
which is the condition of \eqref{gff:mgf:latticeSums} with $\alpha=\varkappa$; its three entries come from the region where
both lattice indices are large, from the end $\ell\to\infty$ and from the end $n\to\infty$ (paragraph \emph{The
entries of $b_{\max}$} of Section~\ref{gff:gdr:periodization}). So $\tau_*=\frac12-2q$ for $q\le\frac14$, $\tau_*=0$ for
$\frac14\le q\le\frac12$ and $\tau_*=2q-1$ for $q\ge\frac12$, and $\eta_{\rm R}=2\varkappa-\max(d,1)$.

\paragraph{Reflected form of the three parts.}
Put
\begin{equation}\begin{aligned}
 \mathsf F_1(z)&=\frac{\Gamma(z+q)^3}{\Gamma(z+2q)^2\,\Gamma(z+q+\tau)},&
 \mathsf F_2(z)&=\frac{\Gamma(z+1-q)^3}{\Gamma(z+1)^2\,\Gamma(z+1-q+\tau)},\\
 \mathsf F_3(z)&=\frac{\Gamma(z)^2\,\Gamma(z+q-\tau)}{\Gamma(z+q)^3},&
 \mathsf F_4(z)&=\frac{\Gamma(z+1-2q)^2\,\Gamma(z+1-q-\tau)}{\Gamma(z+1-q)^3},
\end{aligned}\label{gff:gdk:Fratios}\end{equation}
so that $\mathsf F_1=\widetilde\Phi$ and $\Phi_t(m)=\mathsf F_1(m+t)$ in \eqref{gff:gdk:latticeFactor} (the
$\mathsf F_i$ are not the Laurent coefficients $F_k$ of \eqref{gff:gdk:notation}). The reflection
formula $\Gamma(w)\Gamma(1-w)=\pi/\sin\pi w$ \cite{gff:DLMF5} and $\sin\pi(c-n)=(-1)^n\sin\pi c$ give, for integers
$n,j,p\ge0$,
\[
 \Psi_t(n)=T_{\rm R}(t)\,\mathsf F_2(n+t),\qquad \Psi_t(-j-1)=\mathsf F_3(j+1-t),\qquad
 \Phi_t(-1-p)=T_{\rm L}(t)\,\mathsf F_4(p+1-t),
\]
with $T_{\rm R}$ of Section~\ref{gff:gdk:Rfamily} and
$T_{\rm L}(t)=\sin^2\pi(2q-1+t)\,\sin\pi(q+\tau-1+t)/\sin^3\pi(q-1+t)$; the reflection signs combine to
$(-1)^{4n}=(-1)^{6p}=1$. Hence the right part ($n\ge0$), the middle part ($n=-j-1$, $\ell=j+k+1$) and the left part
($n=-\ell-1-p$) of $S_{\rm R}$ are
\begin{equation}\begin{aligned}
 S^{\rightarrow}(t)&=\Gamma(2q)^2\,T_{\rm R}(t)\sum_{\ell,n\ge0}d_\ell\,\mathsf F_1(\ell+n+t)\,\mathsf F_2(n+t),\\
 S^{\mid}(t)&=\Gamma(2q)^2\sum_{j,k\ge0}d_{j+k+1}\,\mathsf F_1(k+t)\,\mathsf F_3(j+1-t),\\
 S^{\leftarrow}(t)&=\Gamma(2q)^2\,T_{\rm L}(t)\sum_{\ell,p\ge0}d_\ell\,\mathsf F_4(p+1-t)\,\mathsf F_3(\ell+p+1-t),
\end{aligned}\label{gff:gdk:parts}\end{equation}
that is $\Gamma(2q)^2P_{\rm R}H_{\rm R}$, $\Gamma(2q)^2P_{\rm M}H_{\rm M}$ and $\Gamma(2q)^2P_{\rm L}H_{\rm L}$ of
\eqref{gff:gdk:sample}, and $S_{\rm R}=S^{\rightarrow}+S^{\mid}+S^{\leftarrow}$ wherever the double series converges
absolutely. Since $R_\tau(-s)\neq R_\tau(s)$, the left part is not the right part at a reflected point, unlike the
wheels of Section~\ref{gff:gdr:periodization}, and it is estimated separately. For $0\le x\le1$ every argument $z$ in
\eqref{gff:gdk:parts} is $m+t$ or $m+1-t$ with an integer $m\ge0$ (one of $\ell+n$, $n$, $k$, $j$, $p$, $\ell+p$), so
$\operatorname{Re}z\ge m$ and $|z|\ge\max(m,|y|)$.

\paragraph{Gamma ratios and multiplicities.}
Let $\rho\ge2+|\tau|$. Every parameter in \eqref{gff:gdk:Fratios} then has modulus at most $\rho$, and each
$\mathsf F_i$ is a product of three ratios of \eqref{gff:mgf:gammaBounds},
\[\begin{aligned}
 \mathsf F_1&=\Bigl[\frac{\Gamma(z+q)}{\Gamma(z+2q)}\Bigr]^2\frac{\Gamma(z+q)}{\Gamma(z+q+\tau)},&
 \mathsf F_2&=\Bigl[\frac{\Gamma(z+1-q)}{\Gamma(z+1)}\Bigr]^2\frac{\Gamma(z+1-q)}{\Gamma(z+1-q+\tau)},\\
 \mathsf F_3&=\Bigl[\frac{\Gamma(z)}{\Gamma(z+q)}\Bigr]^2\frac{\Gamma(z+q-\tau)}{\Gamma(z+q)},&
 \mathsf F_4&=\Bigl[\frac{\Gamma(z+1-2q)}{\Gamma(z+1-q)}\Bigr]^2\frac{\Gamma(z+1-q-\tau)}{\Gamma(z+1-q)},
\end{aligned}\]
whose exponents $u-v$ add up to $-2q-\tau$. Since $z^az^b=z^{a+b}$ for principal powers, \eqref{gff:mgf:gammaBounds} gives,
for $\operatorname{Re}z\ge0$, $|z|\ge2\rho$ and $i=1,\dots,4$,
\begin{equation}\begin{gathered}
 \mathsf F_i(z)=z^{-2q-\tau}e^{\theta_i(z)},\qquad |\theta_i(z)|\le\frac{3\omega(\rho)}{|z|},\\
 e^{-3\omega(\rho)/|z|-\frac\pi2|\operatorname{Im}\tau|}\,|z|^{-\varkappa}\le|\mathsf F_i(z)|\le
 e^{3\omega(\rho)/|z|+\frac\pi2|\operatorname{Im}\tau|}\,|z|^{-\varkappa}.
\end{gathered}\label{gff:gdk:Fbound}\end{equation}
By \eqref{gff:mgf:multiplicityBound}, $|d_\ell|\le C_q(1+\ell)^{d-2}$ with $C_q$ bounded for $q$ in compact subsets of
$(0,1)$, across $d=1$, and $d_\ell(1+\ell)^{2-d}\to2/\Gamma(d-1)$; the limit vanishes only at $q=\frac14$, so for
$q\neq\frac14$ also $|d_\ell|\ge c\,\ell^{d-2}$ for large $\ell$. By \eqref{gff:fp:degeneracy}, with
$(4q-1)_k=(4q-1)(4q)_{k-1}$, all $d_\ell\ge0$ for $q\ge\frac14$, while for $q<\frac14$ ($d<1$) $d_\ell<0$ for every
$\ell\ge2$.

\paragraph{Trigonometric factors.}
\emph{Lemma.} Let $c_{\rm T}=8/(1-e^{-2\pi})^3\approx8.045$. If $|y|\ge1$ and $|y+\operatorname{Im}\tau|\ge1$, then
$|T_{\rm R}(t)|,|T_{\rm L}(t)|\le c_{\rm T}e^{\pi|\operatorname{Im}\tau|}$. For all $t$,
$|D_{\rm R}(t)|\le e^{6\pi|y|+\pi|\operatorname{Im}\tau|}$, and if $|y|\ge1+|\operatorname{Im}\tau|$, then
$|D_{\rm R}(t)|\ge\bigl(\frac{1-e^{-2\pi}}2\bigr)^6e^{6\pi|y|-\pi|\operatorname{Im}\tau|}$.

\emph{Proof.} From $|\sin\pi w|^2=\sin^2(\pi\operatorname{Re}w)+\sinh^2(\pi\operatorname{Im}w)$ one has
$\sinh(\pi|\operatorname{Im}w|)\le|\sin\pi w|\le e^{\pi|\operatorname{Im}w|}$, and
$\sinh(\pi b)\ge\frac12(1-e^{-2\pi})e^{\pi b}$ for $b\ge1$. The sines of $T_{\rm R}$ have imaginary parts $-y$
(numerator, three times) and $y$, $y$, $-(y+\operatorname{Im}\tau)$ (denominator); those of $T_{\rm L}$ have $y$, $y$,
$y+\operatorname{Im}\tau$ (numerator) and $y$ three times (denominator); those of $D_{\rm R}$ have $y$ five times and
$y+\operatorname{Im}\tau$ once. Insert these bounds and use
$|y|-|\operatorname{Im}\tau|\le|y+\operatorname{Im}\tau|\le|y|+|\operatorname{Im}\tau|$. $\square$

\paragraph{Proof of (b), with constants.}
With $\alpha=\varkappa$, the sums $\Sigma^{\rm r}(Y)$ and $\Sigma^{\rm m}(Y)$ of \eqref{gff:mgf:latticeSums} obey
$\Sigma^{\rm r}(Y)\le C^{\rm r}Y^{-\eta_{\rm R}}(1+\log Y)$ and $\Sigma^{\rm m}(Y)\le C^{\rm m}Y^{-2\operatorname{Re}\tau}$
for $Y\ge1$, since $\max(d,1)-2\varkappa=-\eta_{\rm R}$ and $d-2\varkappa=-2\operatorname{Re}\tau$. Let $\mathfrak K$
be a compact set of pairs $(q,\tau)$ with $0<q<1$ and $\tau\in\mathbb H_q$, and take maxima over $\mathfrak K$. Put
$\rho_{\mathfrak K}=2+\max|\tau|$ and $Y_{\mathfrak K}=2\rho_{\mathfrak K}$, which exceeds
$1+\max|\operatorname{Im}\tau|$; let
$c_{\mathfrak K}=2^\varkappa\exp\bigl(3\omega(\rho_{\mathfrak K})/Y_{\mathfrak K}+\frac\pi2|\operatorname{Im}\tau|\bigr)$,
a function of $(q,\tau)$, and
\[
 C_{\mathfrak K}=\max\,\Gamma(2q)^2\,C_q\,c_{\mathfrak K}^2\,\bigl[c_{\rm T}e^{\pi|\operatorname{Im}\tau|}C^{\rm r}
 +2^{\max(d-2,0)}C^{\rm m}\bigr],
\]
with $C^{\rm r}$, $C^{\rm m}$ at $\alpha=\varkappa$; it is finite because $C_q$, $C^{\rm r}$ and $C^{\rm m}$ are
bounded on compact sets. \emph{For $(q,\tau)\in\mathfrak K$ and $|y|\ge Y_{\mathfrak K}$ the three series
\eqref{gff:gdk:parts} converge absolutely for $0\le x\le1$, $|S^{\rightarrow}|$ and $|S^{\leftarrow}|$ are at most
$C_{\mathfrak K}|y|^{-\eta_{\rm R}}(1+\log|y|)$, $|S^{\mid}|$ is at most $C_{\mathfrak K}|y|^{-2\operatorname{Re}\tau}$,
and for all real $x$}
\begin{equation}\begin{aligned}
 |S_{\rm R}(x+\ii y)|&\le3C_{\mathfrak K}\,|y|^{-\eta_{\rm R}}(1+\log|y|),\\
 |\Xi_{\rm R}(x+\ii y)|&\le3C_{\mathfrak K}\,e^{\pi|\operatorname{Im}\tau|}e^{6\pi|y|}\,|y|^{-\eta_{\rm R}}(1+\log|y|).
\end{aligned}\label{gff:gdk:growthBound}\end{equation}
This is (b) with $Y=Y_{\mathfrak K}$ and $C=3C_{\mathfrak K}$, the bound on the parts holding for $0\le x\le1$.

\emph{Proof.} Fix $(q,\tau)\in\mathfrak K$, $0\le x\le1$ and $|y|\ge Y_{\mathfrak K}$. Every argument $z$ in
\eqref{gff:gdk:parts} has $\operatorname{Re}z\ge0$ and $|z|\ge\max(m,|y|)\ge\frac12(m+|y|)$, in particular
$|z|\ge Y_{\mathfrak K}=2\rho_{\mathfrak K}$, so \eqref{gff:gdk:Fbound} with $\rho=\rho_{\mathfrak K}$ gives
$|\mathsf F_i(z)|\le c_{\mathfrak K}(m+|y|)^{-\varkappa}$. The lemma on trigonometric factors applies, since
$|y|\ge1+|\operatorname{Im}\tau|$. With $|d_\ell|\le C_q(1+\ell)^{d-2}$ and
$|d_{j+k+1}|\le C_q(2+j+k)^{d-2}\le2^{\max(d-2,0)}C_q(1+j+k)^{d-2}$, the series \eqref{gff:gdk:parts} are dominated
termwise:
\[
 |S^{\rightarrow}|,\ |S^{\leftarrow}|\le\Gamma(2q)^2c_{\rm T}e^{\pi|\operatorname{Im}\tau|}C_q\,c_{\mathfrak K}^2\,
 \Sigma^{\rm r}(|y|),\qquad
 |S^{\mid}|\le\Gamma(2q)^2\,2^{\max(d-2,0)}C_q\,c_{\mathfrak K}^2\,\Sigma^{\rm m}(|y|);
\]
for the left part the dominating sum $\sum_{\ell,p\ge0}(1+\ell)^{d-2}(\ell+p+|y|)^{-\varkappa}(p+|y|)^{-\varkappa}$ is
again $\Sigma^{\rm r}(|y|)$. This gives the bounds on the parts. Their absolute convergence also justifies
$S_{\rm R}=S^{\rightarrow}+S^{\mid}+S^{\leftarrow}$ for $0\le x\le1$; for other $x$ one uses
$S_{\rm R}(t+1)=S_{\rm R}(t)$, the shift $n\mapsto n+1$ of the absolutely convergent double series. Finally
$|\Xi_{\rm R}|\le|D_{\rm R}||S_{\rm R}|$ and the bound on $D_{\rm R}$. $\square$

The constant is explicit but crude, because $\omega(\rho)$ grows like $\frac{11}3\rho^3$. For $|y|\ge Y\ge Y_{\mathfrak K}$
one may replace $Y_{\mathfrak K}$ by $Y$ in $c_{\mathfrak K}$; as $Y\to\infty$ the constant decreases to the same
expression with $c_{\mathfrak K}$ replaced by $2^\varkappa e^{\pi|\operatorname{Im}\tau|/2}$ (Section~\ref{gff:gdk:numerics},
item Growth constants).

\paragraph{Proof of (a) and (c).}
Sufficiency in (a), and (c), are the case $r=3$ of Step~0 of the proof of \eqref{gff:gdr:topG}, carried out in full in the
paragraph \emph{Details of the top-mode proof} of Section~\ref{gff:gdr:top}, with the parameters of
Section~\ref{gff:gdk:topmode}: $\sigma=q$, $(g_1,g_2,g_3)=(q,q,\tau)$, $\alpha_i=0$ and
$(\beta_1,\beta_2,\beta_3)=(0,0,q-\tau)$. Then $f^{\rm R}_\ell(a_\ell+n+t)=\Gamma(2q)^2T_\ell(n+t)$,
$D_T=D_{\rm R}$, $\varkappa=2q+\operatorname{Re}\tau$, and condition (C) of \eqref{gff:gdr:BC} is $\tau\in\mathbb H_q$. The
six pole classes of Step~0 are $-q+\mathbb Z$ three times, $\mathbb Z$ twice and $q-\tau+\mathbb Z$ once, so
$\mathsf P_T=\mathsf P_\tau$: through $\Gamma(\ell+n+q+t)^3$, $\Gamma(-n-t)^2$ and $\Gamma(q-\tau-n-t)$ a term of
$S_{\rm R}$ has poles of orders at most $3$ on $-q+\mathbb Z$, $2$ on $\mathbb Z$ and $1$ on $q-\tau+\mathbb Z$, added
where classes meet, which are the zero orders of $D_{\rm R}$. Step~0 gives absolute and locally uniform convergence
off $\mathsf P_\tau$; uniform convergence of $(t-t_0)^o$ times the series on a small circle around each
$t_0\in\mathsf P_\tau$, with $o$ the order of the zero of $D_{\rm R}$ at $t_0$, hence the pole orders in (a);
periodicity by $n\mapsto n+1$; an entire one-periodic $\Xi_{\rm R}$; the vanishing of the modes $|k|\ge3$, the
boundary modes included because $\eta_{\rm R}=2\varkappa-\max(d,1)>0$; and coefficients $c_k=\Gamma(2q)^2\gamma_k$
that are holomorphic in the parameters, hence in $\tau\in\mathbb H_q$. As a trigonometric polynomial in $t$ with
holomorphic coefficients, $\Xi_{\rm R}$ is jointly holomorphic on $\mathbb C\times\mathbb H_q$. The converse in (a) is
proved in the paragraph \emph{Sharpness of $\mathbb H_q$} below. $\square$

\paragraph{Exponential decay of the total.}
For $y>Y_{\mathfrak K}$ there are no poles, since $\mathsf P_\tau$ lies on the
lines $\operatorname{Im}t=0$ and $\operatorname{Im}t=-\operatorname{Im}\tau$; so $S_{\rm R}(t)=\widehat S(e^{2\pi\ii t})$
with $\widehat S$ holomorphic on $0<|w|<e^{-2\pi Y_{\mathfrak K}}$. By \eqref{gff:gdk:growthBound}, $\widehat S(w)\to0$ as
$w\to0$, so $\widehat S$ extends holomorphically with $\widehat S(0)=0$ (Riemann's theorem on removable
singularities). The maximum principle for $\widehat S(w)/w$ on $|w|\le e^{-2\pi(Y_{\mathfrak K}+1)}$ then gives
$|S_{\rm R}(x+\ii y)|\le e^{2\pi(Y_{\mathfrak K}+1)}M_{\mathfrak K}\,e^{-2\pi y}$ for $y\ge Y_{\mathfrak K}+1$, where
$M_{\mathfrak K}=3C_{\mathfrak K}(Y_{\mathfrak K}+1)^{-\eta_{\rm R}}(1+\log(Y_{\mathfrak K}+1))$ bounds $|S_{\rm R}|$
on $\operatorname{Im}t=Y_{\mathfrak K}+1$; for $y\le-Y_{\mathfrak K}-1$ one uses $e^{-2\pi\ii t}$ in place of $w$. So
$S_{\rm R}=O(e^{-2\pi|\operatorname{Im}t|})$ and $\Xi_{\rm R}=O(e^{4\pi|\operatorname{Im}t|})$, uniformly in $x$ and on
$\mathfrak K$, while numerically the parts decay only like powers (Section~\ref{gff:gdk:numerics}, item Growth): the
exponential decay of the total is a cancellation between the parts. Equation \eqref{gff:gdk:threeModes} improves this to
$S_{\rm R}=O(e^{-4\pi|\operatorname{Im}t|})$.

\paragraph{Sharpness of $\mathbb H_q$.}
\emph{Let $0<q<1$ and $\tau\notin\mathbb H_q$. For every $t\notin\mathsf P_\tau\cup(q+\mathbb Z)$ the double series
$S_{\rm R}(t)$ does not converge absolutely.} So the lattice sum, and with it any estimate of the type (b), exists
exactly on $\mathbb H_q$, up to the discrete set of exceptional $t$; this is the converse in (a).

\emph{Proof.} Fix $t$, and put $\rho=2+|\tau|$ and $m_0=\lceil2|\operatorname{Re}t|+2+2\rho\rceil$. For an integer
$m\ge m_0$ the arguments $z=m+t$ and $z=m+1-t$ of \eqref{gff:gdk:parts} have
$\operatorname{Re}z\ge m-|\operatorname{Re}t|\ge\max(2\rho,\frac12(1+m))$, so \eqref{gff:gdk:Fbound} gives
$c\,|z|^{-\varkappa}\le|\mathsf F_i(z)|\le C|z|^{-\varkappa}$, and $\mathsf F_i(z)\neq0$ there;
$T_{\rm R}(t)\neq0$ because $t\notin q+\mathbb Z$; and for $q\neq\frac14$, $|d_\ell|\ge c\,\ell^{d-2}$ for large
$\ell$. One of the conditions \eqref{gff:gdk:threeConditions} fails.
\begin{itemize}
\item $2\varkappa\le1$. The terms of the right part with $\ell=0$,
$d_0\Gamma(2q)^2T_{\rm R}(t)\mathsf F_1(n+t)\mathsf F_2(n+t)$, have modulus at least $c\,n^{-2\varkappa}$ for large
$n$ and are not summable. This case contains $q=\frac14$, where $\tau_*=0=\frac12-2q$.
\item $2\varkappa>1$ and $\varkappa\le4q-1$, so $q>\frac38$. Fix $j=j_0\ge m_0$ in the middle part. Its terms
$d_{j_0+k+1}\Gamma(2q)^2\mathsf F_1(k+t)\mathsf F_3(j_0+1-t)$, with $\mathsf F_3(j_0+1-t)\neq0$ fixed, have modulus
at least $c\,k^{4q-2-\varkappa}$ for large $k$, and $4q-2-\varkappa\ge-1$.
\item $2\varkappa>1$, $\varkappa>4q-1$ and $\varkappa\le2q$, so $q>\frac14$ and $\operatorname{Re}\tau\le0$. In the
middle part the block $N\le j,k<2N$, for large $N$, has $N^2$ terms, each of modulus at least
$c\,N^{4q-2}N^{-2\varkappa}$, because $j+k+1$ lies between $N$ and $4N$ and $|k+t|,|j+1-t|\le3N$. The block
contributes at least $c\,N^{-2\operatorname{Re}\tau}\ge c$, and the sum over the blocks $N=2^i$ diverges. $\square$
\end{itemize}

\paragraph{Is the exponent $\eta_{\rm R}$ attained?}
For $q\ge\frac14$, where all $d_\ell\ge0$ and $\eta_{\rm R}=2\operatorname{Re}\tau$, the bound on the right and left
parts is attained numerically: $\max_x|S^{\rightarrow}||y|^{\eta_{\rm R}}$ tends to a constant at $d=1.25$, $1.5$,
$2.5$ and $3$ (records \path{n5_case0.log} and \path{n5_case2.log} to \path{n5_case5.log} in
\path{LRI_C5_general_d/K_general/note_fragment/growth/checks/}). Numerically it is not attained at $d=0.8$
($q<\frac14$), the only such point tested: there the right and left parts decay like $|y|^{-2\operatorname{Re}\tau}$,
faster than $|y|^{-\eta_{\rm R}}$ (\path{n5_case1.log}). The absolute values taken in the proof lose the cancellation
$\sum_\ell d_\ell=\lim_{z\to1}(1+z)(1-z)^{1-4q}=0$ of the negative multiplicities.

\subsection{Residue functional, three-node rule and the Horn samples}\label{gff:gdk:threeNode}
\paragraph{Three-node rule.}
Let $\tau\in\mathbb H_q$ with $\operatorname{Re}\tau<\frac32q$ and $\tau\neq q$. Then
\begin{equation}
 \langle t_\ell^2R_\tau\rangle_d=-\Res_{t=0}\frac{\Xi_{\rm R}}{D_{\rm R}}-\Res_{t=q-\tau}\frac{\Xi_{\rm R}}{D_{\rm R}}
 =\sum_{|k|\le1}\rho_kc_k,
 \label{gff:gdk:functional}
\end{equation}
\begin{equation}
 \rho_k=-\frac{2\pi\ii k-\pi\bigl[3\cot\pi q+\cot\pi(\tau-q)\bigr]}{\pi^2\sin^3(\pi q)\sin\pi(\tau-q)}
 -\frac{e^{2\pi\ii k(q-\tau)}}{\pi\sin^2\pi(q-\tau)\,\sin^3\pi(2q-\tau)},
 \label{gff:gdk:rho}
\end{equation}
and for every real $t_0$ with $t_0+\frac j3\notin\mathsf P_\tau$,
\begin{equation}\boxed{
 \langle t_\ell^2R_\tau\rangle_d=\sum_{j=0}^2w_j(\tau)\,\Xi_{\rm R}\bigl(t_0+\tfrac j3;\tau\bigr),\qquad
 w_j(\tau)=\frac13\sum_{|k|\le1}\rho_k(\tau)\,e^{-2\pi\ii k(t_0+j/3)}.}
 \label{gff:gdk:threeNodeRule}
\end{equation}

\emph{Proof.} The left poles of $f^{\rm R}_\ell$ are $s=-a_\ell-n$, and the right poles are $s=a_\ell+n$ and
$s=b_\ell-\tau+n$, $n\ge0$, that is $t=0$ and $t=q-\tau$ on the lattice $s=a_\ell+n+t$. The imaginary axis separates
the two families because $a_\ell\ge\frac q2>0$ and $\operatorname{Re}(b_\ell-\tau)\ge\frac32q-\operatorname{Re}\tau>0$.
The two right families are disjoint because $q-\tau\notin\mathbb Z$, which on the strip excludes only $\tau=q$.
Moreover $2q-\tau\notin\mathbb Z$ holds automatically: $2q-\tau=m\in\mathbb Z$ would give
$\operatorname{Re}\tau=2q-m$, which is either $\ge2q>\frac32q$ or $\le2q-1\le\tau_*(q)$.

The trace converges absolutely. Put $\varkappa=2q+\operatorname{Re}\tau$ (the $\varkappa$ of \eqref{gff:gdr:BC} for this
kernel). On the axis, the Gamma-ratio bound \eqref{gff:mgf:gammaBounds}, which gives
$|\Gamma(z+u)/\Gamma(z+v)|\le C|z|^{\operatorname{Re}(u-v)}$ for $\operatorname{Re}z\ge0$ and $|z|$ large, applied with
$z=a_\ell\pm\ii\nu$, gives $|f^{\rm R}_\ell(\ii\nu)|\le C(a_\ell+|\nu|)^{-2\varkappa}$
for $a_\ell+|\nu|$ large. So each $\nu$-integral converges ($2\varkappa>1$), and with $|d_\ell|\le C(1+\ell)^{d-2}$,
$\sum_\ell|d_\ell|\int|f^{\rm R}_\ell(\ii\nu)|\,d\nu\le C\sum_{\ell\ge1}\ell^{4q-2}\ell^{1-2\varkappa}<\infty$, since
$4q-1-2\varkappa=-1-2\operatorname{Re}\tau$ and $\operatorname{Re}\tau>0$.

Close each Barnes contour of \eqref{gff:gdk:fullTrace} to the right, on arcs $|s|=R_N\to\infty$, $\operatorname{Re}s\ge0$,
with radii chosen so that the arcs keep a distance at least $\frac18$ from both right lattices. On the arcs, reflection
of $\Gamma(a_\ell-s)^2\Gamma(b_\ell-\tau-s)/\Gamma(b_\ell-s)^3$ \cite{gff:DLMF5} gives a ratio of three sines over three
sines, bounded away from the lattices and for large $|\operatorname{Im}s|$, times two Gamma ratios of the type above.
Hence $|f^{\rm R}_\ell(s)|\le C_\ell|s|^{-4q-2\operatorname{Re}\tau}$ with $4q+2\operatorname{Re}\tau>1$, so the arcs do
not contribute, and the integral is minus the sum of the right residues. The terms of $S_{\rm R}$ with $n<0$ are
holomorphic at $t=0$ and $t=q-\tau$, since a pole there would need $q\in\mathbb Z$, $q-\tau\in\mathbb Z$ or
$2q-\tau\in\mathbb Z$. Uniform convergence on small circles (Section~\ref{gff:gdk:growthProof}) therefore identifies the sums
of residues with the residues of $S_{\rm R}=\Xi_{\rm R}/D_{\rm R}$, and \eqref{gff:gdk:threeModes} gives
$\sum_{|k|\le1}\rho_kc_k$. The three-point discrete Fourier transform is exact for modes $|k|\le1$. $\square$

Part (c) of the growth theorem alone gives a five-node rule, with the nodes $t_0+j/5$ and the modes $|k|\le2$. By \eqref{gff:gdk:threeModes}, three nodes give the same $\KK_d$ at three fifths of the sampling cost
(Section~\ref{gff:gdk:numerics}). On the jet circle at $d=2$, $\rho_k$ has a triple pole at $\tau=0$, from
$\sin^3\pi(2q-\tau)$, which cancels in the sum; there $\tau=0$ is also a spurious pole of individual pieces of the
continued samples (Section~\ref{gff:gdk:continuation}), and the exceptional set $\mathsf E$ of hypothesis (H), stated in Section~\ref{gff:gdk:continuation}, still meets the
closed disk $|\tau|\le r_\tau$ only in $\{0\}$.

\paragraph{Separating contours (sketch).}
The strip of \eqref{gff:gdk:functional}, $\tau\in\mathbb H_q$ with $\operatorname{Re}\tau<\frac32q$, is nonempty iff
$d>\frac47$. For $\operatorname{Re}\tau\ge\frac32q$ the imaginary axis no longer separates the pole families for the
finitely many spins with $\operatorname{Re}(b_\ell-\tau)\le0$. For $d\le\frac47$ this happens at every
$\tau\in\mathbb H_q$, because then $\tau_*(q)=\frac12-2q\ge\frac32q$. For those spins one replaces the axis by a
contour that separates the left pole family from the two right ones. The proof above then gives \eqref{gff:gdk:functional},
for $\tau\notin q+\mathbb Z$, for the resulting separating-contour trace. That trace is jointly meromorphic in
$(q,\tau)$. Its poles lie at $\tau\in2q+\mathbb N_0$, where the left family $s=-a_\ell-n$ and the right family
$s=b_\ell-\tau+n'$ pinch the contour ($\tau=\ell+2q+n+n'$). There the factor $1/\sin^3\pi(2q-\tau)$ of $\rho_k$ in
\eqref{gff:gdk:rho} is singular as well, and $\tau=2q$, which has $\operatorname{Re}\tau\ge\frac32q$, lies in
$\mathbb H_q$ for $q>\frac18$. The trace agrees with
the axis trace wherever the latter is defined, so it is the meromorphic continuation of the trace in $\tau$ and, for
$d\le\frac47$, also in $q$. This is a sketch: the construction is not carried out in full in this note
(Section~\ref{gff:gdk:status}).

\paragraph{The samples as Horn functions.}
Put
\[\begin{aligned}
 \widehat\Phi(m)&=\frac{(q+t)_m^3}{(2q+t)_m^2(q+\tau+t)_m},&
 \widehat\Psi(n)&=\frac{(1-q+t)_n^3}{(1+t)_n^2(1-q+\tau+t)_n},\\
 \widehat\Psi_{\rm M}(m)&=\frac{(1-t)_m^2(1+q-\tau-t)_m}{(1+q-t)_m^3},&
 \widehat\Phi_{\rm L}(p)&=\frac{(2-2q-t)_p^2(2-q-\tau-t)_p}{(2-q-t)_p^3}.
\end{aligned}\]
For $\tau\in\mathbb H_q$ and $t\notin\mathsf P_\tau$,
\begin{equation}\boxed{
 \Xi_{\rm R}(t;\tau)=D_{\rm R}(t)\,\Gamma(2q)^2\bigl[P_{\rm R}H_{\rm R}+P_{\rm M}H_{\rm M}+P_{\rm L}H_{\rm L}\bigr],}
 \label{gff:gdk:sample}
\end{equation}
with $P_{\rm R}=\Phi_t(0)\Psi_t(0)$, $P_{\rm M}=\Phi_t(0)\Psi_t(-1)$, $P_{\rm L}=\Phi_t(-1)\Psi_t(-1)$ from
\eqref{gff:gdk:latticeFactor}, and
\[\begin{gathered}
 H_{\rm R}=\sum_{\ell,n\ge0}d_\ell\widehat\Phi(\ell+n)\widehat\Psi(n),\qquad
 H_{\rm M}=\sum_{j,k\ge0}d_{j+k+1}\widehat\Psi_{\rm M}(j)\widehat\Phi(k),\\
 H_{\rm L}=\sum_{\ell,p\ge0}d_\ell\widehat\Psi_{\rm M}(\ell+p)\widehat\Phi_{\rm L}(p).
\end{gathered}\]
In the convention of \eqref{gff:fp:HornDefinition}, with $\boldsymbol e=(1,0)$, $\boldsymbol f=(0,1)$ and
$\boldsymbol c=\boldsymbol e+\boldsymbol f$, these are $H_{\rm R}=\mathfrak H[\mathsf U_{\rm R};\mathsf V_{\rm R};(1,1)]$,
$H_{\rm M}=4q\,\mathfrak H[\mathsf U_{\rm M};\mathsf V_{\rm M};(1,1)]$ and
$H_{\rm L}=\mathfrak H[\mathsf U_{\rm L};\mathsf V_{\rm L};(1,1)]$, where
\begin{equation}\begin{aligned}
 \mathsf U_{\rm R}={}&\{(4q-2,\boldsymbol e),(2q,\boldsymbol e),(q+t,\boldsymbol c)^{\times3},(1,\boldsymbol f),
 (1-q+t,\boldsymbol f)^{\times3}\},\\
 \mathsf V_{\rm R}={}&\{(2q-1,\boldsymbol e),(2q+t,\boldsymbol c)^{\times2},(q+\tau+t,\boldsymbol c),
 (1+t,\boldsymbol f)^{\times2},(1-q+\tau+t,\boldsymbol f)\},\\
 \mathsf U_{\rm M}={}&\{(4q-1,\boldsymbol c),(2q+1,\boldsymbol c),(1,\boldsymbol e),(1-t,\boldsymbol e)^{\times2},
 (1+q-\tau-t,\boldsymbol e),\\
 &\qquad(1,\boldsymbol f),(q+t,\boldsymbol f)^{\times3}\},\\
 \mathsf V_{\rm M}={}&\{(2q,\boldsymbol c),(2,\boldsymbol c),(1+q-t,\boldsymbol e)^{\times3},(2q+t,\boldsymbol f)^{\times2},
 (q+\tau+t,\boldsymbol f)\},\\
 \mathsf U_{\rm L}={}&\{(4q-2,\boldsymbol e),(2q,\boldsymbol e),(1-t,\boldsymbol c)^{\times2},(1+q-\tau-t,\boldsymbol c),\\
 &\qquad(1,\boldsymbol f),(2-2q-t,\boldsymbol f)^{\times2},(2-q-\tau-t,\boldsymbol f)\},\\
 \mathsf V_{\rm L}={}&\{(2q-1,\boldsymbol e),(1+q-t,\boldsymbol c)^{\times3},(2-q-t,\boldsymbol f)^{\times3}\}.
\end{aligned}\label{gff:gdk:arrays}\end{equation}

\emph{Proof.} On the right part ($n\ge0$) the Pochhammer form of \eqref{gff:gdk:latticeFactor} gives
$\Phi_t(\ell+n)\Psi_t(n)=P_{\rm R}\widehat\Phi(\ell+n)\widehat\Psi(n)$; the reflection signs combine to
$(-1)^{6n}=1$. On the middle part put $n=-j-1$ and $\ell=j+k+1$, and on the left part $n=-\ell-1-p$. The
multiplicities are $d_\ell=(4q-2)_\ell(2q)_\ell/[(2q-1)_\ell\,\ell!]$ (Section~\ref{gff:gs:seriesSection}) for the
$\boldsymbol e$ rows and \eqref{gff:fp:degeneracy} for the $\boldsymbol c$ rows of $H_{\rm M}$, and the rows
$(1,\boldsymbol e)$, $(1,\boldsymbol f)$ cancel the factorials of \eqref{gff:fp:HornDefinition}. Absolute convergence on
$\mathbb H_q$ (Section~\ref{gff:gdk:Rfamily}) justifies the regrouping. $\square$

The subscripts ${\rm R},{\rm M},{\rm L}$ of $H,P,\mathsf U,\mathsf V$ and of $T_{\rm R},T_{\rm L}$ name the right,
middle and left parts of the lattice; on $S,D,\Xi,\eta,\Omega$ the subscript ${\rm R}$ names the $R$ family.

All parameters are affine in $(q,t,\tau)$. The series $H_{\rm R}$ and $H_{\rm L}$ converge absolutely on
$\mathbb H_q$, and $H_{\rm M}$ on $\operatorname{Re}\tau>\max(0,2q-1)$. At $d=2$ the rows
$(4q-2,\boldsymbol e)/(2q-1,\boldsymbol e)$ are $0/0$; there one splits off $\ell=0$ as in \eqref{gff:fp:wheelBlocks}. At
$d=3$ and $\tau=q$, where $R_q=g_\ell$, the rows of $H_{\rm R}$ reduce to those of $\mathfrak H_3(t)$ in
\eqref{gff:td:Harrays}.

\subsection{\texorpdfstring{Continuation in $\tau$ and the jet}{Continuation in tau and the jet}}\label{gff:gdk:continuation}
The jet of \eqref{gff:gdk:Kr3} is taken on a circle $|\tau|=r_\tau$ around the pole $\tau=0$ of
$\langle t_\ell^2R_\tau\rangle_d$. Since $\operatorname{Re}\tau>0$ is one of the conditions defining $\mathbb H_q$,
no such circle lies in $\mathbb H_q$; for $d=0.8$, $2.5$, $2.75$ and $3$ the whole circle used for the values of Section~\ref{gff:gdk:numerics}
($r_\tau=0.04$ or $0.08$) lies outside it. The samples are continued as follows.

Each of $\widehat\Phi,\widehat\Psi,\widehat\Psi_{\rm M},\widehat\Phi_{\rm L}$ behaves as a constant times $m^{a-1}$ with
$a=1-2q-\tau$. With $\mathsf b_\alpha(m)=\Gamma(m+\alpha)/\Gamma(m+1)$, subtract an $I$-term model,
$X(m)=\sum_{i<I}x_i\mathsf b_{a-i}(m)+\delta X(m)$ with $\delta X=O(m^{a-1-I})$. Here $X$, $Y$ stand for any of the four factors above and $x_i$, $y_i$ for their model coefficients; they are not the constant $X$ of \eqref{gff:gs:psiBasis} or the $x_h$ of \eqref{gff:gdk:xh}. The coefficients $x_i$ follow from the
Bernoulli-polynomial expansion of $\log X(m)$: each is the Gamma ratio $x_0$ below times a polynomial in Bernoulli
polynomials of the parameters. Each factor has the form $X(m)=\prod_{i=1}^3(u_i)_m/(v_i)_m$, with upper and lower
parameters $u_i,v_i$ read off above (local names) and $\sum_i(u_i-v_i)=a-1$. As $m\to\infty$ \cite{gff:DLMFStirling},
\begin{equation}\begin{gathered}
 X(m)\sim x_0\,m^{a-1}\exp\sum_{k\ge1}\lambda_km^{-k},\qquad
 x_0=\prod_{i=1}^3\frac{\Gamma(v_i)}{\Gamma(u_i)},\\
 \lambda_k=\frac{(-1)^{k+1}}{k(k+1)}\sum_{i=1}^3\bigl[B_{k+1}(u_i)-B_{k+1}(v_i)\bigr],
\end{gathered}\label{gff:gdk:modelCoefficients}\end{equation}
with the Bernoulli polynomials $B_{k+1}$ (not the $B_d$ of \eqref{gff:gs:convolution} or the weights $B_{r,k}$ of
\eqref{gff:gdr:functional}). The basis function $\mathsf b_{a-i}$ has the same expansion with the single pair
$(u,v)=(a-i,1)$ and the constant $1$. Matching the coefficients of $m^{a-1-j}$, $j<I$, gives $x_1,\dots,x_{I-1}$ by
a triangular solve with unit diagonal, starting from $x_0$. The model sums are closed Gamma
forms, by Gauss's sum \cite{gff:DLMFGaussUnity} after $(2q)_\ell/(2q-1)_\ell=1+\ell/(2q-1)$, and by Chu--Vandermonde:
\begin{equation}\begin{aligned}
 \mathcal G_\alpha(n)&:=\sum_{\ell\ge0}d_\ell\,\mathsf b_\alpha(\ell+n)
 =\frac{\Gamma(n+\alpha)\Gamma(2-4q-\alpha)\,(2n+2-4q+\alpha)}{\Gamma(n+3-4q)\Gamma(1-\alpha)},\\
 \mathcal S(\alpha,\beta)&:=\sum_{n\ge0}\mathcal G_\alpha(n)\,\mathsf b_\beta(n)
 =\frac{\Gamma(\alpha)\Gamma(\beta)\Gamma(2-4q-\alpha-\beta)\,(2-4q+\alpha-\beta)}{\Gamma(1-\alpha)\Gamma(3-4q-\beta)},\\
 \widetilde{\mathcal G}_\alpha(j)&:=\sum_{k\ge0}d_{j+k+1}\mathsf b_\alpha(k)
 =\frac{d_{j+1}\Gamma(\alpha)\Gamma(j+2)\Gamma(2-4q-\alpha)}{\Gamma(3-4q)\Gamma(j+2-\alpha)}
 \Bigl[2-4q-\alpha+\frac{(4q-1+j)\alpha}{2q+j}\Bigr],\\
 \mathcal V(\alpha,\beta)&:=\sum_{j,k\ge0}d_{j+k+1}\mathsf b_\beta(j)\,\mathsf b_\alpha(k)
 =\frac{2(1-2q)\Gamma(\alpha)\Gamma(\beta)\Gamma(2-4q-\alpha-\beta)\,(4q-\alpha-\beta)}{\Gamma(3-4q)\Gamma(2-\alpha-\beta)},
\end{aligned}\label{gff:gdk:kernels}\end{equation}
valid where the sums converge and by continuation elsewhere. For the Toeplitz blocks, $H_{\rm R}$ with
$(X,Y)=(\widehat\Phi,\widehat\Psi)$ and $H_{\rm L}$ with $(X,Y)=(\widehat\Psi_{\rm M},\widehat\Phi_{\rm L})$, and for the middle
block $H_{\rm M}$ with $(X,Y)=(\widehat\Phi,\widehat\Psi_{\rm M})$, this gives
\begin{equation}\begin{aligned}
 \sum_{\ell,n\ge0}d_\ell X(\ell+n)Y(n)={}&\sum_{\ell,n\ge0}d_\ell\,\delta X(\ell+n)Y(n)
 +\sum_{n\ge0}\Bigl[\sum_ix_i\mathcal G_{a-i}(n)\Bigr]\delta Y(n)\\
 &+\sum_{i,i'}x_iy_{i'}\mathcal S(a-i,a-i'),\\
 \sum_{j,k\ge0}d_{j+k+1}Y(j)X(k)={}&\sum_{j,k\ge0}d_{j+k+1}\,\delta Y(j)\,\delta X(k)
 +\sum_{j\ge0}\delta Y(j)\sum_ix_i\widetilde{\mathcal G}_{a-i}(j)\\
 &+\sum_{k\ge0}\delta X(k)\sum_iy_i\widetilde{\mathcal G}_{a-i}(k)+\sum_{i,i'}x_iy_{i'}\mathcal V(a-i,a-i').
\end{aligned}\label{gff:gdk:regularized}\end{equation}
For $I\ge1$ the two-variable remainders converge absolutely in a neighborhood of $\tau=0$ for every $0<q<1$, and so
do the one-variable remainders. For $d\neq2$ the only singularity at $\tau=0$ is $\Gamma(2\tau)$, in
$\mathcal S(a,a)$ and $\mathcal V(a,a)$; at $d=2$ the two spurious families below also contain $\tau=0$, and these
poles of individual pieces cancel in the sum. Spurious poles of individual pieces, at $\tau\in1-2q+\mathbb Z$, at $\tau\in2q-1-\mathbb Z$ and at
$\tau=q$, bound the admissible radius of the circle. On $\mathbb H_q$ the right sides of \eqref{gff:gdk:regularized} equal
the lattice sums, by an absolutely convergent rearrangement.

\paragraph{Hypothesis (H) and transfer.}
(H) states that there are a bounded connected open set $\mathcal U$ containing the closed disk $|\tau|\le r_\tau$ and
meeting $\mathbb H_q$, and a closed discrete set $\mathsf E\subset\mathcal U$ with $\mathsf E\cap\{|\tau|\le r_\tau\}\subseteq\{0\}$, such
that the continued samples \eqref{gff:gdk:sample}, \eqref{gff:gdk:regularized} are jointly holomorphic in $(t,\tau)$ for
$\tau\in\mathcal U\setminus\mathsf E$ and $t\notin\mathsf P_\tau$, locally uniformly. The construction above gives this, except
that the uniform remainder bounds of the inverse-factorial expansions \cite{gff:DLMFStirling} are not written out. The radius
$r_\tau$ is also taken so small that the disk contains no zero of $\sin\pi(\tau-q)$ or $\sin\pi(2q-\tau)$ other than
$\tau=0$, so that the weights $w_j$ of \eqref{gff:gdk:threeNodeRule} have no other singularity there. Given (H), the
Fourier coefficients $c_k(\tau)=\int_0^1\Xi_{\rm R}(x+\ii y_1;\tau)e^{-2\pi\ii k(x+\ii y_1)}dx$, with
$y_1>1+\sup_{\mathcal U}|\operatorname{Im}\tau|$, so that the segment avoids $\mathsf P_\tau$, are holomorphic on
$\mathcal U\setminus\mathsf E$: the integrand is jointly continuous and holomorphic in $\tau$, so by Fubini's theorem
its integral over the boundary of every small triangle in $\mathcal U\setminus\mathsf E$ vanishes, and Morera's
theorem applies. For $\tau\in(\mathcal U\cap\mathbb H_q)\setminus\mathsf E$, $c_k(\tau)$ is the $k$-th Fourier
coefficient of the entire one-periodic function $\Xi_{\rm R}(\cdot\,;\tau)$, which does not depend on the line
$\operatorname{Im}t=y_1$. The set $\mathcal U\setminus\mathsf E$ is connected, since a closed discrete subset does not
disconnect a planar domain, and $(\mathcal U\cap\mathbb H_q)\setminus\mathsf E$ is a nonempty open subset of it. The
coefficients with $|k|\ge2$ vanish on that subset by \eqref{gff:gdk:threeModes}, hence identically by the identity
theorem. Both sides of
$\Xi_{\rm R}(t;\tau)=\sum_{|k|\le1}c_k(\tau)e^{2\pi\ii kt}$ are holomorphic on the set of $(t,\tau)$ with
$\tau\in\mathcal U\setminus\mathsf E$ and $t\notin\mathsf P_\tau$, which is connected because its base $\mathcal U\setminus\mathsf E$
and its fibers $\mathbb C\setminus\mathsf P_\tau$ are; they agree for $\tau\in(\mathcal U\cap\mathbb H_q)\setminus\mathsf E$, so
they agree on all of it, at the real nodes included. The right side of \eqref{gff:gdk:threeNodeRule} is then holomorphic on
the connected set $(\mathcal U\setminus\mathsf E)\cup\mathbb H_q$ off the poles of the $\rho_k$, and equals the trace on the
residue strip, so it is the continuation of $\langle t_\ell^2R_\tau\rangle_d$, and
\begin{equation}\boxed{\begin{aligned}
 \KK_d&=-\frac{AU}4+\frac14[\tau^2]\Bigl\{C(\tau)\sum_{j=0}^2w_j(\tau)\,\Xi_{\rm R}\bigl(t_0+\tfrac j3;\tau\bigr)
 -C(\tau)^2\mathcal X_d(\tau)^2\Bigr\},\\
 \mathsf r_3&=[\tau^3]\sum_{j=0}^2w_j(\tau)\,\Xi_{\rm R}\bigl(t_0+\tfrac j3;\tau\bigr).
\end{aligned}}\label{gff:gdk:Ksample}\end{equation}
Thus, for $d\neq2$, $\mathsf r_3$ is a finite combination, with elementary coefficients, of the $\tau$-derivatives
of order at most three at $\tau=0$ of nine two-variable remainder sums (three blocks at three nodes), of one-variable
remainder sums, and of Gamma values. At $d=2$ the triple pole of $\rho_k$ and the spurious poles of individual pieces
at $\tau=0$ raise the orders of the derivatives that enter; see the open item on a form adapted to $d=2$
(Section~\ref{gff:gdk:status}). In the computations below the jet is taken with a Cauchy circle; explicit polygamma-weighted
two-variable series for these derivatives have not been written out.

\paragraph{No part-by-part estimate on the jet circle.}
The transfer cannot be replaced by a direct estimate of the lattice parts on the circle. Heuristically, for large
$|\operatorname{Im}t|$ the right part behaves like a $t$-dependent factor times
$\sum_{n\ge0}(n+t)^{-1-2\tau}\approx t^{-2\tau}/(2\tau)$: on $\mathbb H_q$ it decays like
$|\operatorname{Im}t|^{-2\operatorname{Re}\tau}$ (item Growth of Section~\ref{gff:gdk:numerics}, for $q\ge\frac14$), and for
$\operatorname{Re}\tau<0$ it grows like $|\operatorname{Im}t|^{2|\operatorname{Re}\tau|}$. The total decays fast:
given (H), $|\Xi_{\rm R}(x+\ii y;\tau)|\le\sum_{|k|\le1}|c_k(\tau)|\,e^{2\pi|y|}$, and since
$|\sin\pi w|\ge\sinh\pi|\operatorname{Im}w|$ for each of the six sine factors of $D_{\rm R}$,
$S_{\rm R}=\Xi_{\rm R}/D_{\rm R}=O(e^{-4\pi|y|})$ for $|y|\ge1+|\operatorname{Im}\tau|$. At six points of the jet
circles ($d=0.8,1.25,1.5,2.5,3$; five of the six outside $\mathbb H_q$) the largest of the continued right, middle and
left parts, times $|y|^{2\operatorname{Re}\tau}$, is constant to a few percent for $y=1,2,3$, while
$|\Xi_{\rm R}|e^{-2\pi|y|}$ is constant; at $y=3$ the total is $1.2\cdot10^{15}$ to $1.6\cdot10^{16}$ times smaller than
the largest part (records \path{n6_case0.log} to \path{n6_case5.log}, with the \path{.json} files of the same names, in
\path{LRI_C5_general_d/K_general/note_fragment/growth/checks/}). On the circle the decay
of the total is thus a cancellation between the parts, which no part-by-part bound can give; this is why (H) is a
continuation statement.

\paragraph{Status of (H) and a route to a proof.}
The construction \eqref{gff:gdk:regularized} is expected to satisfy (H). Its kernels \eqref{gff:gdk:kernels} depend on
$(q,\tau)$ only and are meromorphic in $\tau$; their poles, and the points $\tau\in q+\mathbb Z$ where pole classes merge, belong to
$\mathsf E$. By \eqref{gff:gdk:modelCoefficients} its model coefficients are the Gamma ratios $x_0$ times polynomials in
Bernoulli polynomials of the parameters, hence holomorphic in $(t,\tau)$ wherever the $x_0$ are. Each model term gains one power
in the remainder sums (the package uses $I=26$ to $30$). Only the locally uniform remainder bound \cite{gff:DLMFStirling}
is missing. An independent proof of (H), not carried out, could proceed as follows for the right block, with $t$ in a
compact strip. (i) Subtract from each of the two Gamma ratios of the right part, $\widetilde\Phi(z)$ at $z=\ell+n+t$
and $\widetilde\Phi(z+1-2q)$ at $z=n+t$ (Section~\ref{gff:gdk:Rfamily}), the single ratio
$\mathfrak m_c(z)=\Gamma(z+c)/\Gamma(z+c+2q+\tau)$ with a constant $c$. The Gamma-ratio bound \eqref{gff:mgf:gammaBounds}
makes the differences $O(|z|^{-\varkappa-1})$, $\varkappa=2q+\operatorname{Re}\tau$, and the remainder double sums then
converge absolutely for $\operatorname{Re}\tau>\max(-\frac12,2q-2,-2q)$. (ii) The spin sum of $\mathfrak m_c$ is a Beta
integral against $\sum_\ell d_\ell x^\ell=(1+x)(1-x)^{1-d}$ of \eqref{gff:gdr:Dgen}, which gives two Beta functions.
(iii) The remaining $n$-sum of $\mathfrak m_c$ times these Beta functions is a ${}_3F_2(1)$ with one upper parameter
equal to $1$ and parameter excess $2\tau$. A two-term Thomae relation, \cite[eq.~16.4.11]{gff:DLMF16} applied twice, maps a
${}_3F_2(1)$ with upper parameters $(1,u,u')$, lower parameters $(v,v')$ and excess $s=v+v'-1-u-u'$ (local names) to
$\Gamma(v)\Gamma(v')\Gamma(s)/[\Gamma(s+u)\Gamma(s+u')]$ times the ${}_3F_2(1)$ with upper parameters
$(v-1,v'-1,s)$ and lower parameters $(s+u,s+u')$, whose excess is $1$. With $s=2\tau$ this continues the $n$-sum as
$\Gamma(2\tau)$ times a ${}_3F_2(1)$ that converges for all $\tau$. The left block is a Toeplitz sum of the same type;
the middle block needs the two-variable kernels
$\widetilde{\mathcal G}$ and $\mathcal V$ of \eqref{gff:gdk:kernels}.

\subsection{\texorpdfstring{Specializations to $d=1,2,3$}{Specializations to d=1,2,3}}\label{gff:gdk:special}
\paragraph{$d=1$.}
Here $d_\ell=(1,1,0,\dots)$ and the offset is $1-d/2=\frac12$. The two spins merge into one sum over consecutive
pairs $(w,w+\frac12)$ of the lattice $t+\frac12\mathbb Z$, which has period $\frac12$ in $t$, while
$D_{\rm R}T_{\rm R}=-\frac18\cos^3(2\pi t)$ is antiperiodic under $t\to t+\frac12$. Hence
$\Xi_{\rm R}(t+\frac12)=-\Xi_{\rm R}(t)$: only $c_{\pm1}$ survive, and $c_0=c_{\pm2}=0$. Numerically, $C(\tau)$ times
the sampled $\langle t_\ell^2R_\tau\rangle_1$ reproduces \eqref{gff:eq:KproofTwoFunctions} with the path parameter of $J_t$ set to $2\tau$. With
$\KK_1=\KK$ of \eqref{gff:eq:KprovedFinal} (see \eqref{gff:gs:d1Dictionary}), \eqref{gff:gdk:Kr3} gives, conditionally on (M) at $d=1$ (confirmed numerically above), the
closed form
\[\begin{aligned}
 \mathsf r_3(1)&=\mathsf s_3(1)+4\KK_1-2\bigl(\pi^3+28\zeta(3)\bigr)\bigl(3\gamma_E+7\log2+\pi\bigr)\\
 &=767.2452094446567860904798157376224\ldots,
\end{aligned}\]
because $\frac U8[\psi_0(1)+\psi_0(\frac12)+4\psi_0(\frac14)]=\frac12(\pi^3+28\zeta(3))(3\gamma_E+7\log2+\pi)$ at
$q=\frac14$. No further condition enters: (J) at $d=1$ is \eqref{gff:eq:Kpath}, on which \eqref{gff:eq:KprovedFinal} already
rests.

\paragraph{$d=2$.}
Here the offset vanishes and $d_\ell=(1,2,2,\dots)$. Then
\[
 \Psi_t(n)=\frac{\cos^3(\pi t)}{\sin^2(\pi t)\cos\pi(t+\tau)}\,\Phi_t(n),\qquad
 \Xi_{\rm R}(t)=-\Omega_{\rm R}(t)^2,\qquad \Omega_{\rm R}(t)=\cos^3(\pi t)\sum_{m\in\mathbb Z}\Phi_t(m),
\]
because the lattice sum becomes the square of $\sum_m\Phi_t(m)$. The function $\Omega_{\rm R}$ is antiperiodic, and,
being entire with $|\Omega_{\rm R}|^2=|\Xi_{\rm R}|=O(e^{2\pi|\operatorname{Im}t|})$ by \eqref{gff:gdk:threeModes}, it has
only the two modes $e^{\pm\ii\pi t}$. This is the square-lattice structure of
Sections~\ref{gff:mgk:residues}--\ref{gff:mgk:regulated}, where the trace is evaluated in closed form,
$C(\tau)\langle t_\ell^2R_\tau\rangle_2=\pi^2[\mathfrak F(\tau)^2+\mathfrak F(1-\tau)^2]$, that is \eqref{gff:gc:K2master}; the
generic samples reproduce it numerically. At $d=2$ both (M) and (J) are proved,
(M) for $0<\tau<\frac12$ in Section~\ref{gff:mgk:regulated} and (J) in Section~\ref{gff:gc:K2laurent}, so \eqref{gff:gdk:Kr3} and \eqref{gff:gc:K2closed} give the exact value
\[\begin{aligned}
 \mathsf r_3(2)&=\mathsf s_3(2)+128\Cat^2-168\,\zeta(3)\log2-42\gamma_E\zeta(3)\\
 &=35.63714944308140680727039972168179\ldots .
\end{aligned}\]

\paragraph{$d=3$.}
Here the offset is $\frac12$ and $d_\ell=2\ell+1$. The half-shift argument of \eqref{gff:td:determinants}, repeated for
this kernel, gives the following for $\frac12<\tau<\frac98$; the derivation is not written out here, and the identity
is checked numerically at six points to $1.1\cdot10^{-24}$ (Section~\ref{gff:gdk:numerics}).
\[
 \Xi_{\rm R}(t)+\Xi_{\rm R}(t+\tfrac12)=-\frac\pi{16}\cos^3(2\pi t)\,
 \bigl[\mathcal M^1_{\rm o}\mathcal M^0_{\rm e}-\mathcal M^0_{\rm o}\mathcal M^1_{\rm e}\bigr],\qquad
 \mathcal M^k_{\rm e,o}=\sum_{n\in\mathbb Z}(n+u)^k\,\widetilde\Phi\bigl(n+u-\tfrac58\bigr)\big|_{q=3/4},
\]
with $u=t+\frac18$ (e) or $u=t+\frac58$ (o) and $\widetilde\Phi$ of Section~\ref{gff:gdk:Rfamily}, so that
$\widetilde\Phi(w-\frac58)|_{q=3/4}=\Gamma(w+\frac18)^3/[\Gamma(w+\frac78)^2\Gamma(w+\frac18+\tau)]$; at $\tau=q=\frac34$
this is the cube of the function $\Gamma(x+\frac18)/\Gamma(x+\frac78)$ of \eqref{gff:td:Udefinition}.
By \eqref{gff:gdk:threeModes} the left side is the constant $2c_0$. So $c_0$ is an ordinary $2\times2$ determinant of
one-variable bilateral moments, while this identity does not reduce $c_{\pm1}$; this is the $R$ analogue of
\eqref{gff:td:projectedAnswer}. No closed form of $\KK_3$ is known. The residue series \eqref{gff:gs:Kseries}, whose
derivation away from $d=2$ uses (M) and the extraction of $[\tau^2]$ under the trace (Section~\ref{gff:gdk:status}), summed
in ball arithmetic, gives $430$ digits (\path{LRI_C5_d3/wf_K3/K3_value.json}),
$\KK_3=1.14859464952717933655664631081067318763\ldots$; the position-space value of
Section~\ref{gff:gdk:positionSpace}, which uses neither, agrees with it to $3.6\cdot10^{-100}$ (Section~\ref{gff:gdk:numerics}).
With this value, \eqref{gff:gdk:Kr3}, conditionally on (M) and (J) at $d=3$, gives
$\mathsf r_3(3)=4.511759358236005780256620897961575\ldots$. A PSLQ search for $\KK_3$ at $420$ digits found no integer relation in the bases listed in
\path{LRI_C5_d3/wf_K3/pslq_K3_results_420.json}.

\paragraph{Consistency.}
At $d=1,2,3$ the five-node version of \eqref{gff:gdk:Ksample} (nodes $t_0+j/5$, modes $|k|\le2$) agrees with
\eqref{gff:eq:KprovedFinal}, \eqref{gff:gc:K2closed} and the $430$-digit \eqref{gff:gs:Kseries} value of $\KK_3$ to
$5\cdot10^{-37}$, $5\cdot10^{-48}$ and
$1.4\cdot10^{-43}$
(relative), and the position-space value of Section~\ref{gff:gdk:positionSpace} agrees with them to about $10^{-100}$
(Section~\ref{gff:gdk:numerics}). The Laurent coefficients $\mathsf r_3(1)$ and $\mathsf r_3(2)$ of the sampled jets agree
with the closed forms above to $10^{-32}$.

\subsection{Position-space representation}\label{gff:gdk:positionSpace}
By completion symmetry $J_t=J_d(d/2,d/2,d/2,d-t,d/2)=J_d(d-t,d/2,d/2,d/2,d/2)$. The $y$-integral is then the
one-loop star with indices $(q,q,q)$, which has a two-fold Mellin--Barnes representation, and the $x$-integral is a
bubble. With $\mathsf N_q=\Gamma(2q)^2/\Gamma(q)^4$ and $z_{12}=z_1+z_2$,
\begin{equation}
 J_t=\mathsf N_q\int\!\!\int\frac{dz_1\,dz_2}{(2\pi\ii)^2}\,
 \Gamma(-z_1)^2\Gamma(-z_2)^2\Gamma(q+z_{12})^2\frac{\Gamma(q+z_2)}{\Gamma(q-z_2)}\,
 \frac{\Gamma(\frac t2+z_1)\Gamma(q-\frac t2-z_{12})}{\Gamma(2q-\frac t2-z_1)\Gamma(q+\frac t2+z_{12})},
 \label{gff:gdk:MB}
\end{equation}
on $-\frac t2<\operatorname{Re}z_1<0$, $-q<\operatorname{Re}z_2<0$, $\operatorname{Re}z_{12}>-q$. Moving
$\operatorname{Re}z_1$ to the left of $-\frac t2$ crosses one simple pole. By Barnes' first lemma its residue is
$J^{\rm pole}_t=\Gamma(2q)\Gamma(\frac t2)^2\Gamma(q-\frac t2)^4/[\Gamma(q)^4\Gamma(2q-t)]$, and the remaining
integral is holomorphic at $t=0$. The coefficient $[t^2]$ of $J_t-J_t^\Gamma$ is therefore the sum of the polygamma part $\KK_d^\Gamma:=[t^2](J^{\rm pole}_t-J^\Gamma_t)$ and the coefficient $\mathcal I_d$ of $t^2$ in the remaining integral:
\begin{equation}\boxed{
 \KK_d=\KK_d^\Gamma+\mathcal I_d,}
 \label{gff:gdk:KMB}
\end{equation}
\[\begin{gathered}
 \KK_d^\Gamma=\frac{A^2\vartheta}2+\frac A2\bigl[\psi_2(2q)-\psi_2(q)\bigr]-\frac{R\vartheta}4+\frac{\vartheta^2}8
 +\frac{2\psi_3(q)-7\psi_3(2q)}{48},\qquad \vartheta=2\psi_1(q)-\psi_1(2q),\\
 \mathcal I_d=\mathsf N_q\int\!\!\int\frac{dz_1\,dz_2}{(2\pi\ii)^2}\,\mathfrak M_0\,\frac{\varpi_1^2+\varpi_2}2,\\
 \mathfrak M_0=\frac{\Gamma(-z_1)^2\Gamma(z_1)}{\Gamma(2q-z_1)}\,\frac{\Gamma(-z_2)^2\Gamma(q+z_2)}{\Gamma(q-z_2)}\,
 \Gamma(q+z_{12})\Gamma(q-z_{12}),\\
 \varpi_1=\tfrac12\bigl[\psi_0(z_1)+\psi_0(2q-z_1)-\psi_0(q-z_{12})-\psi_0(q+z_{12})\bigr],\\
 \varpi_2=\tfrac14\bigl[\psi_1(z_1)-\psi_1(2q-z_1)+\psi_1(q-z_{12})-\psi_1(q+z_{12})\bigr],
\end{gathered}\]
on $-q<\operatorname{Re}z_1,\operatorname{Re}z_2<0$, $\operatorname{Re}z_{12}>-q$ (for example
$\operatorname{Re}z_1=\operatorname{Re}z_2=-q/3$). By (J) the orders $t^0$ and $t^1$ of $J_t-J_t^\Gamma$ vanish,
which gives two Barnes identities:
\begin{equation}
 \mathsf N_q\int\!\!\int\frac{dz_1\,dz_2}{(2\pi\ii)^2}\,\mathfrak M_0=\psi_1(2q)-2\psi_1(q),\qquad
 \mathsf N_q\int\!\!\int\frac{dz_1\,dz_2}{(2\pi\ii)^2}\,\mathfrak M_0\,\varpi_1=-\Bigl[A\vartheta+\frac{\psi_2(2q)}2-\psi_2(q)\Bigr].
 \label{gff:gdk:MBidentities}
\end{equation}
This route uses only the definition \eqref{gff:gs:Kdefinition}, which defines $\KK_d$ as the coefficient $[t^2]$; it uses
none of (M), (J), (H), and not the spectral kernel. Its agreement with
\eqref{gff:gdk:Ksample} at the level of $10^{-41}$ to $10^{-47}$ (Section~\ref{gff:gdk:numerics}) is therefore a numerical
test of (M), and of (H) on the jet circle, at non-integer $d$. The completed graph of $J_t$ can be put into this star-and-bubble form in $60$ ways (choice of the vertex sent to infinity, the two external vertices, the star vertex and the bubble vertex). A computer enumeration shows that each of the $60$ two-fold Mellin--Barnes integrands contains a squared Gamma function, so this framework has no derivative-free residue expansion of $J_t$. The
kite with general indices can also be written as a sum of four Kamp\'e de F\'eriet series $F^{2:2;2}_{2:1;1}(1,1)$ and two Horn double series. This form is not derived or used in this note; it is recorded, with its convergence conditions, in
\path{LRI_C5_general_d/K_general/position_space/position_space_summary.json}.

\subsection{Numerical confirmation}\label{gff:gdk:numerics}
\paragraph{Values.}
The first table gives $\KK_d$ from \eqref{gff:gdk:KMB}, evaluated to $100$ digits by the trapezoid rule in the ball arithmetic of the Arb library~\cite{gff:Arb} on $\operatorname{Re}z_1=\operatorname{Re}z_2=-q/3$
(\path{LRI_C5_general_d/K_general/position_space/K_d_values_mb_D100.json}), and $\mathsf r_3$ obtained from it through
\eqref{gff:gdk:Kr3}. That column assumes (M) and (J); it agrees with the $\tau^3$ coefficients of the $34$-digit Laurent
tables of the directly sampled five-node jets to $3\cdot10^{-34}$ (relative). It is printed to $34$ digits, so that
\eqref{gff:gdk:Kr3} with the printed $\mathsf r_3$ reproduces the printed $\KK_d$ to $3\cdot10^{-32}$ (relative;
$10^{-33}$ against the $100$-digit values). The digits shown are correctly rounded; $\ldots$ marks omitted digits.
\begin{center}\small
\begin{tabular}{@{}lll@{}}\toprule
 $d$ & $\KK_d$, position space \eqref{gff:gdk:KMB} & $\mathsf r_3=[\tau^3]\langle t_\ell^2R_\tau\rangle_d$ by \eqref{gff:gdk:Kr3}\\\midrule
$0.8$ & $92.372770239922771632645391899036\ldots$ & $1894.360867513991151448939858484147\ldots$\\
$1$ & $39.331569128259039057503187151598\ldots$ & $767.2452094446567860904798157376224\ldots$\\
$1.25$ & $17.203345729921688327085720106549\ldots$ & $302.1379775430212297584696879029900\ldots$\\
$1.5$ & $9.0182974015948205470847431046793\ldots$ & $136.5685443212260369206603722998423\ldots$\\
$2$ & $3.5181091385902522416490901703077\ldots$ & $35.63714944308140680727039972168179\ldots$\\
$2.5$ & $1.8405012954285880307301818602723\ldots$ & $11.46731776044183878576775168950130\ldots$\\
$2.75$ & $1.4292491886603120000991755368839\ldots$ & $6.984719816576732523926795215898772\ldots$\\
$3$ & $1.1485946495271793365566463108107\ldots$ & $4.511759358236005780256620897961575\ldots$\\\bottomrule
\end{tabular}
\end{center}

\paragraph{Routes.}
The second table gives relative deviations from the position-space values. ``Spectral reference'' is the
$\tau$-circle evaluation of \eqref{gff:gs:Kspectral} in \path{LRI_C5_general_d/references/reference_values.json}
(50 printed digits, of which that file states 45 to 47 to be correct). ``Abel--Plana'' is \eqref{gff:gs:Kspectral} with an Abel--Plana spin
tail at 60 digits, equal to the residue series \eqref{gff:gc:Kcancelled} to better than $10^{-53}$ for $d\le2.5$ and to
$5\cdot10^{-49}$ at $d=2.75$
(\path{LRI_C5_general_d/K_general/refcheck/}). The two jets are the first line of \eqref{gff:gdk:Ksample}, that is
\eqref{gff:gdk:Ktau}, evaluated with the five nodes $t_0+j/5$ (modes $|k|\le2$) and with the three nodes of
\eqref{gff:gdk:threeNodeRule}, on a Cauchy circle of $48$ points and radius $0.08$ ($0.04$ at $d=0.8$), at $72$ digits.
The three-node circles at $d=1.5$, $2.5$ and $3$ were computed by the separate check code of an independent AI review of Sections~\ref{gff:gdr:section} and \ref{gff:gdk:section}
(\path{LRI_C5_general_d/review/joint/numerics-end-to-end/k3_jets.json}); through \eqref{gff:gdk:Kr3} instead of
\eqref{gff:gdk:Ktau} they give $3.8\cdot10^{-45}$, $5.0\cdot10^{-44}$ and $-8.1\cdot10^{-43}$
(\path{LRI_C5_general_d/K_general/note_fragment/test_build_joint/check_k3_routes.log}). The circles at $d=0.8$,
$2.5$, $2.75$ and $3$ lie entirely outside $\mathbb H_q$. The last column has the opposite sign convention: it gives the
position-space value minus the reference, relative to the reference, for the references \eqref{gff:eq:KprovedFinal},
\eqref{gff:gc:K2closed} and the $430$-digit \eqref{gff:gs:Kseries} value of $\KK_3$, the last conditional on (M) and on the
extraction of $[\tau^2]$ under the trace.
\begin{center}\small
\begin{tabular}{@{}lccccc@{}}\toprule
 $d$ & spectral reference & Abel--Plana & five-node jet & three-node jet &
 \begin{tabular}[b]{@{}c@{}}position space vs\\ \eqref{gff:eq:KprovedFinal}, \eqref{gff:gc:K2closed}, \eqref{gff:gs:Kseries}\end{tabular}\\\midrule
$0.8$ & $-1.9\cdot10^{-49}$ & $8.2\cdot10^{-57}$ & $-8.0\cdot10^{-47}$ & $-8.1\cdot10^{-47}$ & --\\
$1$ & -- & -- & $-4.8\cdot10^{-37}$ & -- & $5.2\cdot10^{-101}$\\
$1.25$ & $1.8\cdot10^{-49}$ & $-2.4\cdot10^{-55}$ & $-1.0\cdot10^{-41}$ & $-1.0\cdot10^{-41}$ & --\\
$1.5$ & $1.9\cdot10^{-49}$ & $2.1\cdot10^{-56}$ & $-9.6\cdot10^{-46}$ & $-1.4\cdot10^{-45}$ & --\\
$2$ & -- & -- & $5.2\cdot10^{-48}$ & -- & $-1.5\cdot10^{-100}$\\
$2.5$ & $5.2\cdot10^{-51}$ & $-1.2\cdot10^{-55}$ & $5.4\cdot10^{-44}$ & $4.1\cdot10^{-44}$ & --\\
$2.75$ & $2.1\cdot10^{-50}$ & $-8.3\cdot10^{-56}$ & $6.0\cdot10^{-44}$ & $-3.5\cdot10^{-44}$ & --\\
$3$ & -- & -- & $-1.4\cdot10^{-43}$ & $-9.5\cdot10^{-43}$ & $3.6\cdot10^{-100}$\\\bottomrule
\end{tabular}
\end{center}

\paragraph{Further checks.}
All are numerical, with empirical error control and no interval enclosures.
\begin{itemize}
\item \emph{$S$ family.} \eqref{gff:gdk:Sclosed} agrees with the spectral trace at $36$ points ($0.8\le d\le3.5$, real
and complex $\tau$) to $2\cdot10^{-38}$ or better, with an independent real-axis lattice route at $18$ points to
$4\cdot10^{-37}$ or better, and with the continued spectral trace on $\tau$-circles to $8\cdot10^{-29}$.
\item \emph{Low orders.} With the $34$-digit Laurent tables of the five-node jets,
$|\mathsf r_0-\mathsf s_0|$, $|\mathsf r_1-\mathsf s_1|$ and $|\mathsf r_2-\mathsf s_2-\frac32U|$ are at most
$5\cdot10^{-32}$ at all eight $d$, and \eqref{gff:gdk:Kr3} reproduces $\KK_d$ to $10^{-33}$. In the $72$-digit jets
$|F_2-\frac32U|$ lies between $10^{-47}$ and $4\cdot10^{-35}$.
\item \emph{Inputs (M) and (J).} The mixed trace (M) holds pointwise at three complex $\tau$ on the $\tau$-circles of the spectral reference (paragraph Routes)
at $d=0.8,1.25,1.5,2.5,2.75$ to $4\cdot10^{-32}$. For (J), at $d=1.25$ the full quadratic jet of
$J_d-(G_2+G_1-G_0+\frac2{p_1}\Delta_B)$ of \eqref{gff:gs:Bjet} at its base point $(d/2,d/2,d/2,d,d/2)$, reconstructed from its restrictions to $18$ random lines through that point, has only the $z_1z_2$ coefficient, equal to
$\KK_d$ to $2\cdot10^{-12}$; the other $14$ coefficients are below $3\cdot10^{-12}$. The orders $t^0,t^1$, that is
\eqref{gff:gdk:MBidentities}, hold to $10^{-101}$ at all eight $d$.
\item \emph{Samples.} \eqref{gff:gdk:sample} with \eqref{gff:gdk:regularized} agrees with the spectral trace at eight
$(d,\tau)$ points, three of them outside $\mathbb H_q$, to $1.1\cdot10^{-19}$ or better (relative; at $d=0.8,1.25$
against the spectral trace with an analytic $\nu$-tail). The termwise identities behind \eqref{gff:gdk:latticeFactor}
and \eqref{gff:gdk:arrays}, the kernels \eqref{gff:gdk:kernels} and $\rho_k$ of \eqref{gff:gdk:rho} hold to $6\cdot10^{-39}$
(\path{LRI_C5_general_d/K_general/note_fragment/fragment_checks/check_blocks.py}). At $d=1,2$ the generic samples
reproduce \eqref{gff:eq:KproofTwoFunctions} and \eqref{gff:gc:K2master} at two values of $\tau$ each, to $10^{-21}$ and
$2\cdot10^{-19}$ (absolute), and at $d=3$ the determinant form of Section~\ref{gff:gdk:special} holds to $1.1\cdot10^{-24}$
at six points.
\item \emph{Growth.} At six points of $\mathbb H_q$ ($0.8\le d\le3$, real and complex $\tau$) the three parts obey
the bound of Section~\ref{gff:gdk:Rfamily}, which is attained for $q\ge\frac14$, while the total $S_{\rm R}$ is up to
$10^{25}$ times smaller and decays like $e^{-4\pi|\operatorname{Im}t|}$.
\item \emph{Growth constants.} These test the estimates used in Section~\ref{gff:gdk:growthProof}. Stirling's error
bound, the input of \eqref{gff:mgf:gammaBounds}, holds at $4000$ random points with $|\arg w|\le0.97\pi$ and
$0.05\le|w|\le10^4$, the largest ratio of $|\epsilon_\Gamma(w)|$ to the bound being $0.99999973$; in
\eqref{gff:mgf:gammaBounds}, $|\theta||z|/\omega(\rho)\le0.29$ on random samples for $\rho=0.25,\dots,3$, the corner
$|z|=\max(1,2\rho)$, $\arg z\to\pm\frac\pi2$ included. For $0.02\le q\le0.99$ the constant $C_q$ of
\eqref{gff:mgf:multiplicityBound}, evaluated over $\ell\le2\cdot10^4$, lies between $1.04$ and $7.9$, and
$d_\ell(1+\ell)^{2-d}$ approaches $2/\Gamma(d-1)$. At ten pairs $(q,\alpha)$, some close to the thresholds
\eqref{gff:gdk:threeConditions}, and $Y=2,\dots,512$, the sums $\Sigma^{\rm r}(Y)$ and $\Sigma^{\rm m}(Y)$ of
\eqref{gff:mgf:latticeSums} (FFT convolutions of $2^{22}$ terms with a fitted tail) stay below their bounds, the largest
ratio being $0.81$, and the exact limits of $Y^{2\alpha-\max(d,1)}\Sigma^{\rm r}(Y)$ (for $d\neq1$) and
$Y^{2\alpha-d}\Sigma^{\rm m}(Y)$ as $Y\to\infty$ lie below $C^{\rm r}$ and $C^{\rm m}$. The Horn-block code of
\eqref{gff:gdk:sample}, used for the item Growth, reproduces the three parts \eqref{gff:gdk:parts}, computed as
Richardson-extrapolated direct lattice sums in ball arithmetic, at twelve complex $t$ with $|\operatorname{Im}t|\le4$
($d=0.8,1.5,2.5,3.4$, $\operatorname{Re}\tau=2.37$ or $2.6$) to between $2\cdot10^{-17}$ and $5\cdot10^{-9}$
(relative), within the error of the direct sums. At the six points of the item Growth, $C_{\mathfrak K}$ of
Section~\ref{gff:gdk:growthProof} for the single point $\mathfrak K=\{(q,\tau)\}$ is $4\cdot10^{31}$ to $4\cdot10^{44}$;
its limit for large $|y|$ is $2.6\cdot10^3$ to $2.7\cdot10^4$, which is $200$ to $1000$ times the observed maximum of
$|S^{\rightarrow}|$, $|S^{\mid}|$, $|S^{\leftarrow}|$ times $|y|^{\eta_{\rm R}}/(1+\log|y|)$. Records:
\path{n1_stirling.log}, \path{n2_mult.log}, \path{n3_lattice.log}, \path{n4_crossval.log} and
\path{n8_constants.log} in \path{LRI_C5_general_d/K_general/note_fragment/growth/checks/}.
\item \emph{Top modes.} With the Horn-block code at eleven points, $0.3\le d\le3.9$, including two continued points,
$|c_{\pm2}|/\max_k|c_k|$ lies between $6\cdot10^{-33}$ and $3\cdot10^{-27}$ and falls with the truncation error. An
independent per-spin code gives ratios below $10^{-15}$ at five points of $\mathbb H_q$, where single spins contribute
$O(1)$, each below the spread of its spin-sum extrapolation (at most $8.5\cdot10^{-16}$); three unbalanced controls give between $3\cdot10^{-3}$ and $6\cdot10^{-2}$. At six points of
the jet circles the continued samples have $|c_{\pm2}|/\max_k|c_k|\le1.3\cdot10^{-37}$.
\item \emph{Three nodes.} At $d=0.8,1.25,2.75$ the three-node $\KK_d$ equals the five-node value to
$3\cdot10^{-49}$, $3.5\cdot10^{-47}$ and $9.5\cdot10^{-44}$, at $0.5$ to $0.65$ of the wall time per circle point.
At the two points of the item Top modes nearest to $d=4$, $d=3.7$ and $3.9$, the three-node and five-node traces
differ by $8\cdot10^{-24}$ and $1.5\cdot10^{-21}$ (relative), more than the top modes, which are about $10^{-27}$ of
the largest mode. The cause is the conditioning of the residue functional, not the rule: at $d=3.9$,
$\tau=1.2+0.3\ii$, the weights \eqref{gff:gdk:rho} have $|\rho_k|\approx2\cdot10^4$, through the factor
$\cot(\pi q)/\sin^3(\pi q)$,
against a trace of modulus about $9$, so \eqref{gff:gdk:functional} cancels about five digits and amplifies the truncation
error of the samples, for both rules alike; at $d=3.9$ the difference is $5\cdot10^{-14}$ with the coarser of the
two truncations (\path{modes_d3.7_t0.95+0.15j.json} and \path{modes_d3.9_t1.2+0.3j.json} in
\path{LRI_C5_general_d/K_general/note_fragment/topmode/results/}).
\item \emph{Position space.} Two contours and step sizes agree to $3\cdot10^{-55}$ (relative) at
$d=0.8,1.25,1.5,2.5$. An independent double-precision
evaluation of \eqref{gff:gdk:KMB} and \eqref{gff:gdk:MBidentities}, written from the formulas as printed here, reproduces
$\KK_d$ to $7\cdot10^{-14}$ (relative) and the two identities to $10^{-12}$ absolute ($3\cdot10^{-15}$ relative) at all eight $d$
(\path{LRI_C5_general_d/K_general/note_fragment/fragment_checks/check_barnes.py}).
\item \emph{A caution.} The literal bracket of \eqref{gff:gs:Kseries}, summed with the Abel--Plana tails of
\path{LRI_C5_general_d/K_general/refcheck/routeB_residue.py}, fails numerically at $d=0.8$, the only $d<1$ tested: it
gives $70.3$ at $40$ digits and $2539.5$ at $60$ digits in place of $\KK_{0.8}=92.37\ldots$, while the canceled form
\eqref{gff:gc:Kcancelled}, summed by the same code, works. For $1.25\le d\le3$ the two brackets agree to the working
precision. The cause is numerical. The two brackets agree term by term, also at non-integer and complex indices, but
the literal one is a difference of two nearly equal terms of order $n^{-d}$ in the radial index $n$, and that code
fixes its constants at $20$ digits beyond its working precision, so the difference keeps no correct digit beyond
$n\approx10^{80}$ at $60$ digits ($10^{60}$ at $40$ digits). At $d=0.8$ the radial tail of that code reaches
$n\approx10^{394}$ ($3.5\cdot10^{279}$), where $n$ times the corrupted term is not small; at $d=1.25$ it stops near
$n\approx10^{93}$ ($10^{67}$), and $n$ times the error stays below $10^{-70}$ ($10^{-50}$), far below the
working precision (\path{check_kbracket_d08.log} and \path{check_kbracket_d08_huge.log} in
\path{LRI_C5_general_d/K_general/note_fragment/test_build_joint/}). This does not bear on the convergence statement
after \eqref{gff:gs:Kseries}.
\end{itemize}

\subsection{Status}\label{gff:gdk:status}
\paragraph{Proved.}
The following have written proofs, in this section or in Section~\ref{gff:gdr:section}. Here an independent review is a separate AI session that re-derived the statement and tried to refute it, and a numerical probe is a separate code. Unless stated otherwise, a statement has had no independent review.
\begin{itemize}
\item The closed form \eqref{gff:gdk:Sclosed} and its Laurent form \eqref{gff:gdk:sk}: Gamma algebra, the kernel identity
\eqref{gff:gdk:Skernel}, and the two-wheel convolution $W_2(d;b')=B_d(d-b',2b')$ of \eqref{gff:gs:convolution}, also
used for \eqref{gff:gs:Ksubtraction}. Three separate AI sessions re-derived it independently and tried to refute it; none
found an error.
\item The constant \eqref{gff:gdk:Sconstant} and the bilateral identity \eqref{gff:gdk:Sbilateral} (Section~\ref{gff:gdk:Sfamily},
remark (iii)): consequences of the case $r=2$ of \eqref{gff:gdr:topless}, of \eqref{gff:gdr:functional} and of
\eqref{gff:gs:convolution}; the constancy and the value at $b=d/2-\tau$ are confirmed numerically there. No independent
review.
\item The Laurent series \eqref{gff:gdk:Delta}: exact Gamma algebra, checked symbolically for general $q$.
\item The lattice and product forms \eqref{gff:gdk:latticeFactor} and \eqref{gff:gdk:Rproduct}, the sample formula
\eqref{gff:gdk:sample} with the arrays \eqref{gff:gdk:arrays}, and the kernels \eqref{gff:gdk:kernels}: exact derivations. Separate AI sessions re-derived the lattice and product forms, the sample formula and the
arrays and tried to refute them, without finding an error; all of them, the kernels included, are rechecked
numerically to $6\cdot10^{-39}$ (Section~\ref{gff:gdk:numerics}, item Samples).
\item The growth theorem of Section~\ref{gff:gdk:Rfamily}, proved in Section~\ref{gff:gdk:growthProof}: growth bound with
explicit constants, finite Fourier space with $|k|\le2$, and sharpness of $\mathbb H_q$. Parts (a) (sufficiency), (b) and (c) are the case $r=3$ of
\eqref{gff:gdr:topG}; an independent AI review re-derived every step of that proof, Step~0 included, by hand (Section~\ref{gff:gdr:status}). The sharpness of $\mathbb H_q$ and the explicit constants of Section~\ref{gff:gdk:growthProof} have had no independent
review.
\item The theorem \eqref{gff:gdk:threeModes} ($R$ family, modes $|k|\le1$) on $\mathbb H_q$. It is the case $r=3$ of the
product form \eqref{gff:gdr:topG} of the top-mode theorem, which has a written proof (Section~\ref{gff:gdr:top}). An independent AI review re-derived every step of that proof by hand and found no fatal or major issue (record \path{LRI_C5_general_d/review/lemmaG-proof}); it also checked
the $R$-family parameters of Section~\ref{gff:gdk:topmode} and the equivalence of condition (C) of \eqref{gff:gdr:BC}
with $\tau\in\mathbb H_q$. A separate $48$-digit lattice-sum code finds the top modes of four $R$ kernels below
$5.2\cdot10^{-45}$ of the largest mode (Section~\ref{gff:gdr:checks}; \path{LRI_C5_general_d/review/lemmaG-numeric}).
\item The residue functional \eqref{gff:gdk:functional} and the three-node rule \eqref{gff:gdk:threeNodeRule} for
$\tau\in\mathbb H_q$, $\operatorname{Re}\tau<\frac32q$, $\tau\neq q$.
\item The representation \eqref{gff:gdk:MB} and \eqref{gff:gdk:KMB}: exact derivation in Section~\ref{gff:gdk:positionSpace}
from \eqref{gff:gs:Kdefinition} (completion symmetry, the two-fold Mellin--Barnes form of the one-loop star, Barnes' first
lemma). It uses none of (M), (J), (H), since \eqref{gff:gs:Kdefinition} defines $\KK_d$ as the coefficient $[t^2]$. Separate AI sessions re-derived the statement and tried to refute it, without finding an error; the
derivation printed in Section~\ref{gff:gdk:positionSpace} has not been reviewed. A separate
code, written from the formulas as printed here, re-evaluates them (Section~\ref{gff:gdk:numerics}, item Position space;
\path{LRI_C5_general_d/K_general/note_fragment/fragment_checks/check_barnes.py}).
\end{itemize}

\paragraph{Conditional.}
\begin{itemize}
\item \eqref{gff:gdk:rLow} and \eqref{gff:gdk:Kr3} are exact given (M) and (J), and \eqref{gff:gdk:Ktau} given (M) alone. The
note proves (M) only at $d=2$; at other $d$ it is numerical. Its strongest test is the agreement of
\eqref{gff:gdk:Ksample} with the position-space route, which uses none of (M), (J), (H); on the jet circle this agreement
tests (M) and (H) together. (J) is the path form of \eqref{gff:gs:Bjet}, derived in Section~\ref{gff:app:generalSpecial}, on which
\eqref{gff:gs:C5} itself rests. At $d=2$ both inputs are proved, (M) for $0<\tau<\frac12$ in
Section~\ref{gff:mgk:regulated} and (J) in Section~\ref{gff:gc:K2laurent}, so there these statements and the value of $\mathsf r_3(2)$ in
Section~\ref{gff:gdk:special} are proved. Elsewhere, including the values of $\mathsf r_3(1)$ and $\mathsf r_3(3)$ in
Section~\ref{gff:gdk:special} and the $\mathsf r_3$ column of the first table of Section~\ref{gff:gdk:numerics}, they are
conditional on (M) and (J). At $d=1$, (J) is \eqref{gff:eq:Kpath}, on which \eqref{gff:eq:KprovedFinal} itself rests, so
there $\mathsf r_3(1)$ is conditional on (M) alone.
\item The two Barnes identities \eqref{gff:gdk:MBidentities} hold given (J).
\item \eqref{gff:gdk:E2} needs, in addition, $[\tau^2]$ under the trace. The derivation of \eqref{gff:gs:Kspectral} uses the
same step, which is justified only at $d=2$ (corollary of Section~\ref{gff:mgk:regularization}); there
\eqref{gff:gdk:E2} is \eqref{gff:mgk:E2}.
\item On the jet circle, \eqref{gff:gdk:threeModes}, \eqref{gff:gdk:threeNodeRule} and \eqref{gff:gdk:Ksample} need (H), the
analyticity in $\tau$ of the continued samples; the first line of \eqref{gff:gdk:Ksample} needs in addition only (M),
through \eqref{gff:gdk:Ktau}. None of the circles used in Section~\ref{gff:gdk:numerics} lies in $\mathbb H_q$.
\item For $d\le\frac47$ the strip of (M) is empty; there the residue functional, and with it the definition of
$\langle t_\ell^2R_\tau\rangle_d$ and of the $\mathsf r_k$, is only sketched. For $\operatorname{Re}\tau\ge\frac32q$
the residue functional is only sketched (Section~\ref{gff:gdk:threeNode}, paragraph ``Separating contours'').
\item For $C_5(d)$ these inputs reduce to (M), plus (H) when $\mathsf r_3$ is sampled on the jet circle as in
\eqref{gff:gdk:Ksample}: (J) adds no condition, since \eqref{gff:gs:C5} already rests on \eqref{gff:gs:Bjet}, and with $\KK_d$
from \eqref{gff:gdk:KMB} none of (M), (J), (H) enters. The range of \eqref{gff:gdr:C5reduced} is $\frac23<d<4$, so the caveat
for $d\le\frac47$ never affects $C_5(d)$.
\end{itemize}

\paragraph{Numerical only.}
All digit agreements of Section~\ref{gff:gdk:numerics}, and the $d=3$ determinant form of $c_0$ in
Section~\ref{gff:gdk:special}. The $100$-digit values are not interval enclosures: ball
arithmetic bounds the rounding error, while discretization and truncation are controlled empirically. The tested
range is $0.8\le d\le3$ for $\KK_d$ and $0.3\le d\le3.9$ for the top modes. A reduction of $\mathsf r_3$ to known
constants was not attempted at generic $d$; at $d=3$ the PSLQ search was negative. The half-shift factorization of
\eqref{gff:td:determinants} is not available at generic $d$, and there is no derivative-free Mellin--Barnes residue
expansion of $J_t$ (Section~\ref{gff:gdk:positionSpace}).

\paragraph{Open.}
\begin{itemize}
\item An independent review of the sharpness of $\mathbb H_q$ and of the explicit constants of the growth theorem
(Section~\ref{gff:gdk:growthProof}).
\item A proof of (M) at generic $d$, and a written proof of (H) with uniform remainder bounds.
\item Explicit polygamma-weighted two-variable series for $\mathsf r_3$ (the $\tau$-derivatives of the remainders in
\eqref{gff:gdk:regularized}), and a form adapted to $d=2$, where individual pieces have removable singularities.
\item A PSLQ search on $\mathsf r_3$ or $\KK_d$ with the $100$-digit values, keeping $n\log_{10}(\max|c|)$ at least
ten digits below the working precision and confirming any relation at a second precision.
\item Dimensions $d<0.8$, and the five-loop primitive kernel, which \eqref{gff:gdr:topG} does not cover.
\item A confluent series for $J_t$ on the path of \eqref{gff:gs:Kdefinition} for $d\ge2$ at small $t$. The confluent
form recorded in \path{LRI_C5_general_d/K_general/position_space/position_space_summary.json} holds only for
$0<d<2+2\operatorname{Re}t$; the confluent limit of the formula with four Kamp\'e de F\'eriet and two Horn series
mentioned at the end of Section~\ref{gff:gdk:positionSpace}, which would cover $d\ge2$, has not been written.
\end{itemize}

\SourcePart{Numerical sixth- to eighth-order coefficients}

\section{The sixth-order coefficient in one dimension}
\label{gff:app:C6Numerical}

\subsection{Result, convention, and scope}
This section extends the one-dimensional calculation by one perturbative
order. It does not extend the arbitrary-dimensional formula to sixth
order and does not alter any fifth-order identity of the analytic parts of
this note.
The regulator and normalization remain
\begin{equation}
 \epsilon=2s-1,\qquad
 \Delta_{\mathcal E}=\frac12-\frac\epsilon6+
 \sum_{n=2}^{6}C_n(1)\epsilon^n+O(\epsilon^7).
 \label{gff:c6:convention}
\end{equation}
The result of the graph-based numerical calculation is
\begin{equation}\boxed{C_6(1)=14331.68703277\pm0.00000049,}
 \label{gff:c6:answer}
\end{equation}
that is, $C_6(1)=14331.6870328(5)$. The value is the mean of the two
infrared-cutoff evaluations at quadrature degree $N_*=22$. Both are
computed by exact summation over every nonzero six-loop graph and every
gap ordering, as described in Section~\ref{gff:c6:exactSums}. The quoted
uncertainty is the pre-registered empirical allowance of
Section~\ref{gff:c6:uncertainty}. It adds the last observed change with
degree, the residual difference between the two cutoffs, a measured
floating-point allowance, and the convergence of the lower-loop inputs.
It is not a rigorous interval enclosure or a probability guarantee, and
it is not a bound on the omitted orders of the $\epsilon$ series.

This is a six-loop graph integration. It is not a Pad\'e prediction,
large-order fit, interpolation in dimension, or extrapolation from the
five previously known physical coefficients, which do not determine
$C_6$. No Monte Carlo sampling enters the final value. An earlier
evaluation combined exhaustive degree-six sums, deterministic hard-sector
refinements and Monte Carlo corrections for the higher degrees. It gave
$14331.77075$ with sampling standard error $0.14344$, quoted as
$14331.8\pm0.5$. The exact sums show that this estimate was high by
$0.084$, or $0.58$ of its standard error; its sampled level differences
are consistent with the exact ones (Section~\ref{gff:c6:checks}).

Section~\ref{gff:c6:analytic} compares \eqref{gff:c6:answer} with a finite
special-function representation of a reduced analytic expression for
$C_6(1)$, derived separately; the two agree to $1.0\times10^{-8}$.
Section~\ref{gff:c1:c6analytic} quotes that representation,
$14331.68703275805$, as the value of $C_6(1)$, and \eqref{gff:c6:answer} is its
independent numerical confirmation.

\subsection{Exactly which new amplitudes are required}
Write the beta function and mass anomalous dimension in the
minimal-subtraction conventions of Sections~\ref{gff:sec:RG}
and~\ref{gff:sec:graphs}:
\begin{equation}
 \beta(g)=-\epsilon g+g^2 B(g),\quad
 B(g)=3+b_2g+\cdots+b_6g^5,\qquad
 \gamma_r(g)=g+c_2g^2+\cdots+c_6g^6.
 \label{gff:c6:RGconventions}
\end{equation}
Since $\epsilon=g_*B(g_*)$, Lagrange inversion gives, for $n\ge2$,
\begin{equation}
 C_n=\frac1n[g^{n-1}]\gamma_r'(g)B(g)^{-n}.
 \label{gff:c6:Lagrange}
\end{equation}
With $A=\pi+4\log2$, $b_2=-6A$, and $c_2=-A$, the exact sixth-order
assembly is
\begin{equation}
\begin{aligned}
 2187C_6={}&3c_6-b_6+2016A^5-448A^3b_3+672A^3c_3
 -84A^2b_4+168A^2c_4\\
 &+\frac{49}{3}Ab_3^2-42Ab_3c_3-12Ab_5+30Ac_5
 +\frac73b_3b_4-4b_3c_4-3b_4c_3.
\end{aligned}\label{gff:c6:RGassembly}
\end{equation}
The script \path{LRI_C6_1D_HPC/reference/verify_algebra.py} is an independent
exact-rational polynomial check of \eqref{gff:c6:RGassembly}. It also checks the sensitivities
$\partial C_6/\partial c_6=1/729$ and
$\partial C_6/\partial b_6=-1/2187$.

Let $\mathcal A_L^\chi(\epsilon)$ and
$\mathcal B_L^\chi(\epsilon)$ denote the bare coupling and mass
amplitudes for the cutoff $\chi$. These are amplitudes, not the scalar
$A$ or the function $B(g)$ in \eqref{gff:c6:RGconventions}. Define
\begin{equation}
 H_6^\chi=[\epsilon^{-1}]
 (3\mathcal B_6^\chi-\mathcal A_6^\chi),\qquad
 J_5^\chi=[\epsilon^0]
 (8\mathcal B_5^\chi-3\mathcal A_5^\chi).
 \label{gff:c6:newInputs}
\end{equation}
The only new five- and six-loop inputs to the physical assembly are
these two numbers. Specifically,
\begin{equation}\boxed{
 C_6=R_6^\chi+\frac4{243}J_5^\chi-\frac2{729}H_6^\chi.}
 \label{gff:c6:bareAssembly}
\end{equation}
Here $R_6^\chi$ is obtained from the same Laurent minimal-subtraction
recursion, setting the five-loop finite terms and six-loop simple poles
to zero, while retaining all previously required poles. Two-loop
amplitudes are now needed through $\epsilon^3$, three-loop amplitudes
through $\epsilon^2$, and four-loop amplitudes through $\epsilon^1$.
Thus $R_6^\chi$ contains newly calculated higher Laurent data even
though it contains no new six-loop master.

For clarity, the new contributions to the renormalization coefficients
are $b_6=108a_{5,0}-6a_{6,-1}+\cdots$ and
$c_6=96e_{5,0}-6e_{6,-1}+\cdots$, where $a,e$ are the coefficients of
$\mathcal A,\mathcal B$. The combination
$3c_6-b_6=36J_5-6H_6+\cdots$ proves the factors and signs in
\eqref{gff:c6:bareAssembly}. Unit shifts of $J_5$ and $H_6$ were also
checked through the full high-precision recursion.

\subsection{Exhaustive graph enumeration and exact weights}
At $L$ loops the ordered-coordinate representation has $L+1$ vertices
and $2L$ internal edges. An integer multiplicity $m_{ij}\in\{0,1,2,3\}$
is assigned to each pair. Tadpoles are omitted, as in the five-loop graph set
\eqref{gff:eq:graphs}; every vertex has degree between two and four, and every
nontrivial cut has at least two crossing edges. The latter requirement
enforces connectivity and one-particle irreducibility.

If $e_i=4-\deg(i)$, the rational Wick weights \eqref{gff:eq:weights} are, with
$w_A=w_4$ and $w_B=w_2$,
\begin{equation}
 w_A(G)=\frac{24}{\prod_i e_i!\prod_{i<j}m_{ij}!},\qquad
 w_B(G)=\frac{w_A(G)}{12}\sum_i e_i(e_i-1),\qquad
 w_H=3w_B-w_A.
 \label{gff:c6:weights}
\end{equation}
All weights are stored exactly as integers in units of $1/1728$.
The fresh enumeration gives
\begin{center}
\begin{tabular}{lrr}
\toprule
 & Five loops & Six loops\\\midrule
Labeled 1PI graphs & 59,145 & 2,427,005\\
Zero $w_H$ graphs & 9,540 & 289,590\\
Nonzero $w_H$ graphs & 49,605 & 2,137,415\\
$\sum_G w_A$ & 197,887.5 & 7,716,476.25\\
$\sum_G w_B$ & 21,562.5 & 732,471.25\\
$\sum_G w_H$ & $-133,200$ & $-5,519,062.5$\\
\bottomrule
\end{tabular}
\end{center}
These are labeled graph counts, not counts of topological isomorphism
classes or master integrals. Every surviving six-loop $w_H$ is negative;
zero-weight graphs are removed by exact integer arithmetic before
integration. The 720 gap orderings reduce to 360 after spatial reversal;
this is an exact symmetry reduction, not a sampled subset of orderings.

\subsection{Meromorphic quadrature and the two cutoffs}
Let $D=L-1$ and order the $L$ consecutive gaps by magnitude. If $R$ is
the largest gap, the smaller ones are
$Rt_0,Rt_0t_1,\ldots,Rt_0\cdots t_{D-1}$, with $0\le t_j\le1$.
For an edge $e$, factor its distance as
\begin{equation}
 |x_i-x_j|=R\left(\prod_{j<r_e}t_j\right)Q_e(t),\qquad Q_e(t)>0,
 \quad
 b_j=\sum_{e:r_e>j}m_e,\qquad
 a_j=L-2-j-\frac{b_j}{2}.
 \label{gff:c6:sectorGeometry}
\end{equation}
Here $r_e$ is the rank of the largest gap crossed by that edge. Each
$Q_e$ is a sum of products of consecutive $t_j$, including a constant
term $1$. Define
\begin{equation}
 H(t)=1+t_0+t_0t_1+\cdots+t_0\cdots t_{D-1}.
 \label{gff:c6:spanFunction}
\end{equation}
The two infrared prescriptions are $R<1$ (maxgap) and $RH(t)<1$
(span). Integrating the common scale exactly, a weighted graph-sector
contributes
\begin{equation}
\begin{aligned}
 \frac{w_G}{L\epsilon}
 \int_{[0,1]^D}\!\prod_jdt_j\,t_j^{a_j+\epsilon b_j/2}
 \prod_e Q_e(t)^{-m_e/2}
 \exp\left\{\epsilon\left[
 \frac12\sum_e m_e\log Q_e(t)
 -L\,\mathbf1_{\chi=\mathrm{span}}\log H(t)\right]\right\}.
\end{aligned}\label{gff:c6:sectorIntegral}
\end{equation}
The integral is meromorphically continued in $\epsilon$ before taking
its Laurent coefficients. It is not an ordinary convergent integral at
$\epsilon=0$ in every sector.

Each smooth factor is tensor-interpolated on the degree-$N$ Chebyshev
grid. Integration of the interpolation polynomial is then exact as a
meromorphic function. For one monomial,
\begin{equation}
 \int_0^1 t^{a+\epsilon b/2+k}\,dt
 =\frac1{a+k+1+\epsilon b/2}.
 \label{gff:c6:monomial}
\end{equation}
If $a+k+1=0$, its contribution is the explicit pole $2/(b\epsilon)$;
otherwise the usual geometric expansion gives every needed Laurent
term. This implements endpoint subtraction without estimating a
regulated integral by a small nonzero $\epsilon$.

The quadrature weights are generated with 100-decimal-digit arithmetic
and stored both in binary64 and as pairs of binary64 numbers whose sum
is accurate to about $10^{-32}$. The binary64 tables agree bitwise with an
independent 80-digit generation, apart from entries that vanish
identically, and they reproduce the exact Laurent coefficients of
polynomial test integrands at every degree used. The arithmetic of the
sums themselves is described in Section~\ref{gff:c6:exactSums}.

\subsection{Exact summation at increasing degree}\label{gff:c6:exactSums}
Let $\boldsymbol H^{(N)}$ collect the twelve six-loop Laurent
coefficients at quadrature degree $N$: powers $-6$ through $-1$ for
maxgap, followed by the same powers for span. Each entry is a finite sum
over all 2,137,415 nonzero six-loop graphs, all 360 reversal classes, and
all $(N+1)^5$ tensor nodes of the interpolated integrand
\eqref{gff:c6:sectorIntegral}. These sums are evaluated in full for
$N=6,8,\ldots,22$; no graph, sector or node is sampled. Two exact
reorganizations make this affordable. Graphs with the same bucket vector
$(b_j)$ in a given sector share their Laurent weight tensor, and the edge
factors $\prod_eQ_e^{-m_e/2}$ and the logarithms in
\eqref{gff:c6:sectorIntegral} factor into two halves of the non-adjacent
edges. The complete sequence through $N=22$ used about 45 node-hours on
64-core processors.

Each graph-sector term is assigned to one of two strata according to its
exponents $a_j$ in \eqref{gff:c6:sectorGeometry}. Let $p$ be the number of
axes with integral $a_j\le-1$, which carry a Laurent pole in
\eqref{gff:c6:monomial}. The hard stratum is
\begin{equation}
 \mathcal S=\{p\ge2\}\ \cup\ \{\text{some }a_j<-1\},
 \label{gff:c6:strata}
\end{equation}
and the light stratum is its complement, in which every axis satisfies
$a_j\ge-1$ and at most one axis produces a pole. The strata are summed in
different arithmetic.
\begin{itemize}
\item The light stratum uses binary64 vector arithmetic on blocks of
eight nodes. Within each block the contributions of all weight groups are
accumulated with error-free two-sum compensation; the blocks are then
reduced in extended precision in a fixed order. Each sector value is
therefore independent of the thread count, the division of sectors
among jobs, and restarts. Its sensitivity to random one-unit perturbations
in the last place of every node datum and weight is at most $7\times10^{-10}$ in
$C_6$ and does not grow with $N$.
\item The hard stratum uses 80-bit extended precision (64-bit
significand) throughout, with the weights read from the double-length
tables. In binary64 the hard-stratum error grows with the degree because
of cancellation within the stratum: the difference between binary64 and
extended precision is $5\times10^{-9}$, $7.6\times10^{-8}$ and
$1.5\times10^{-7}$ in $C_6$ at $N=6$, $8$ and $10$.
\end{itemize}
Plain binary64 accumulation of the sector sums, as in the earlier
deterministic degree-six evaluation, shifts $C_6$ by about $10^{-5}$ at
$N=6$ and $2\times10^{-5}$ at $N=8$. The earlier degree-six value differs
from the present one by about $10^{-5}$ for this reason.

Before production, the implementation was required to reproduce the
earlier deterministic degree-six sums and the earlier hard-stratum sums at
$N=6,8,10$ within those records' own summation errors, to give the
complete sums as light plus hard stratum, and to agree between binary64
and extended precision within the measured errors. Identical executables
give bitwise identical binary64 results on AMD and Intel processors;
extended-precision results agree to $1.7\times10^{-16}$ relative, the
difference coming from vendor-specific x87 transcendental instructions.

\subsection{Five-loop and lower-loop inputs}
The five-loop data in \eqref{gff:c6:bareAssembly} are obtained by the same
exact summation over the five-loop graphs, with the 120 gap orderings reduced
to 60 by spatial reversal; a graph is removed exactly when its weight in the
combination being summed vanishes, $w_H=3w_B-w_A$ for $\mathcal H_5$ (9,540
of the 59,145 graphs) or $8w_B-3w_A$ for $\mathcal J_5$, in extended precision, at degrees $M=12,16,\ldots,40$. The
coupling and mass amplitudes are separated through
$\mathcal A_5=8\mathcal H_5-3\mathcal J_5$ and
$\mathcal B_5=3\mathcal H_5-\mathcal J_5$, where $\mathcal H_5$ and
$\mathcal J_5$ are the combinations $3\mathcal B_5-\mathcal A_5$ and
$8\mathcal B_5-3\mathcal A_5$. Their poles enter $R_6^\chi$, and the
finite part of $\mathcal J_5$ is $J_5^\chi$. From $M=28$ onward,
successive changes of $C_6$ are at most $4.4\times10^{-9}$, below the measured
five-loop rounding sensitivity of $1.0\times10^{-8}$. The two-, three- and
four-loop amplitudes, through the orders listed after
\eqref{gff:c6:bareAssembly}, are evaluated at degree 28; the change from
degree 22 is $-1.54\times10^{-9}$ (maxgap) and $+4.14\times10^{-10}$ (span)
in $C_6$. Recomputing them in binary64 instead of extended precision
changes $C_6$ by $6.0\times10^{-8}$, which implies an extended-precision
rounding error of about $3\times10^{-11}$. The inputs used in the final
assembly are
\begin{center}
\begin{tabular}{lrr}
\toprule
 & maxgap & span\\\midrule
$R_6^\chi$, lower degree 28, five-loop poles at $M=40$ & 39814.47028245011 & 81364.48627171610\\
$J_5^\chi$, $M=40$ & $-231603.7671545272$ & $-1836182.1032583138$\\
$K_\chi=R_6^\chi+\frac4{243}J_5^\chi$ & 36002.06259266366 & 51139.26646499489\\
\bottomrule
\end{tabular}
\end{center}
so that $C_6^\chi(N)=K_\chi-\frac2{729}H_6^\chi(N)$ for each cutoff.

\subsection{Degree sequence, stopping rule and uncertainty}\label{gff:c6:uncertainty}
\begin{center}
\begin{tabular}{rrrrr}
\toprule
$N$ & $C_6^{\mathrm{maxgap}}(N)$ & $C_6^{\mathrm{span}}(N)$ & maxgap$-$span & $\rho(N)$\\\midrule
6 & 14334.549146861 & 14332.554509427 & $1.99$ & $5.5\times10^{-10}$\\
8 & 14331.816869671 & 14331.694925294 & $0.122$ & $1.9\times10^{-10}$\\
10 & 14331.694513177 & 14331.686134511 & $8.4\times10^{-3}$ & $7.4\times10^{-10}$\\
12 & 14331.687501161 & 14331.686917319 & $5.8\times10^{-4}$ & $1.9\times10^{-9}$\\
14 & 14331.687061994 & 14331.687022076 & $4.0\times10^{-5}$ & $4.1\times10^{-9}$\\
16 & 14331.687034432 & 14331.687031833 & $2.6\times10^{-6}$ & $2.9\times10^{-9}$\\
18 & 14331.687032845 & 14331.687032694 & $1.5\times10^{-7}$ & $3.3\times10^{-8}$\\
20 & 14331.687032736 & 14331.687032727 & $9.1\times10^{-9}$ & $4.5\times10^{-9}$\\
22 & 14331.687032770 & 14331.687032767 & $2.5\times10^{-9}$ & $8.8\times10^{-8}$\\
\bottomrule
\end{tabular}
\end{center}
The convergence is geometric. Each step of two degrees reduces the change
by a factor of about $0.06$ for maxgap and $0.1$ for span, and the two
cutoffs, which differ by $2.0$ at $N=6$, approach a common value. The
estimator, the uncertainty and the stopping rule were fixed before any
degree beyond $N=8$ was computed. Degrees $6$ to $16$ were unconditional.
From $N=16$ on, the next degree was computed unless both last changes
$\Delta^\chi(N)=C_6^\chi(N)-C_6^\chi(N-2)$ and the cutoff difference were
below $2\rho(N)$, where $\rho(N)$ is the floating-point estimate described
below. At $N=22$ the changes are $+3.3\times10^{-8}$ and $+4.0\times10^{-8}$
and the cutoff difference is $2.5\times10^{-9}$, all below
$2\rho(22)=1.8\times10^{-7}$, so the sequence stopped at $N_*=22$.

The headline value is $\frac12[C_6^{\mathrm{maxgap}}(N_*)+C_6^{\mathrm{span}}(N_*)]$,
with the allowance
\begin{equation}
 U=\mathrm{ceil}_2\Bigl(2\max_\chi|\Delta^\chi(N_*)|
 +|C_6^{\mathrm{maxgap}}-C_6^{\mathrm{span}}|
 +4\rho(N_*)+U_5+U_{5,\mathrm{round}}+U_{\mathrm{low}}\Bigr),
 \label{gff:c6:allowance}
\end{equation}
where $\mathrm{ceil}_2$ rounds up to two significant figures. The terms
are $8.0\times10^{-8}$, $2.5\times10^{-9}$, $3.5\times10^{-7}$,
$8.7\times10^{-9}$, $4.1\times10^{-8}$ and $3.1\times10^{-9}$, so that
$U=4.9\times10^{-7}$.
\begin{itemize}
\item $\rho(N)$ is a rounding estimate. A fixed subset of 36 sectors is
recomputed with every node datum and weight moved by one unit in the
last place, in random directions, separately for the two strata. The
resulting changes are scaled to all 360 sectors and multiplied by a
calibration factor of $1.03$. That factor is the largest observed ratio
of the true binary64 error, measured against extended precision, to the
same perturbation estimate in binary64 at $N=6,8,10$.
\item $U_5$ is twice the change of $C_6$ between $M=36$ and $M=40$.
$U_{5,\mathrm{round}}$ is four times the calibrated perturbation estimate
for the $M=40$ five-loop data. $U_{\mathrm{low}}$ is twice the lower-degree
change quoted above.
\end{itemize}
The allowance is dominated by the rounding term. Increasing the degree
further would not reduce it. The separately measured rounding of the
lower-loop amplitudes, about $3\times10^{-11}$, does not change $U$.

\subsection{Independent checks}\label{gff:c6:checks}
The poles $\epsilon^{-6}$ through $\epsilon^{-2}$ of $3\mathcal B_6-\mathcal A_6$
are fixed by lower-order renormalization, so their errors test the whole
six-loop pipeline at every degree. For maxgap, the double-pole error is
$9.2\times10^{-2}$, $8.0\times10^{-4}$, $5.2\times10^{-6}$ and
$6.9\times10^{-8}$ at $N=10,14,18,22$, and the triple-pole error at $N=22$
is $2.7\times10^{-9}$; span behaves in the same way. The two cutoffs,
whose six-loop integrands and subtraction constants differ, agree to
$2.5\times10^{-9}$ at $N_*$.

The earlier Monte Carlo evaluation sampled the light-stratum differences
between degrees. Comparing its level estimates with the exact differences,
in units of its standard errors, gives $-0.57$ and $-0.60$ for $8-6$,
$+0.86$ and $+1.33$ for $10-8$, and $+0.09$ and $+1.26$ for the all-graph
$12-10$ diagnostic (maxgap and span). That estimator was therefore
consistent within its errors; its final value was $0.58$ standard
errors above the present one.

The seventh-order calculation of Section~\ref{gff:app:C7complete} produces
$C_6$ as a by-product, in a different representation: coordinate
forests, randomized integration, and radial quadrature through degree
six. Its degree-six value $14332.43593\pm0.30500$ lies $+2.46$ of its
standard errors above \eqref{gff:c6:answer}. Adding its measured degree-six
to degree-seven check step gives $14331.98210$, at $+0.97$; adding further
a geometric tail model gives $14331.62055$, at $-0.22$. The degree-six
tension is the radial truncation of that representation, discussed in
Section~\ref{gff:app:C7complete}; agreement beyond the measured step is
conditional on the tail model.

\subsection{A finite special-function representation}\label{gff:c6:analytic}
A separate analytic calculation, recorded in the package
\path{C6_research_note/}, starts from an analytic reduction of the
six-loop coefficient to explicitly defined periods and series. It converts that expression into a finite formula in nine
special-function values. The reduction itself is derived neither in this note
nor in the package. We state it as the input \eqref{gff:c6:reduced} and compare
its value with \eqref{gff:c6:answer}. Two of its leading Laurent coefficients,
$R_{147,-3}=R_{148,-3}=\frac25$ in the package's numbering, were missing from
the first version of the input and were added by a correction; their
contribution is given below.

In this subsection, for $n\ge0$,
\[
 a_n=\binom{2n}{n}4^{-n},\qquad
 h_n=\log2+H_n-H_{2n},\qquad
 \delta_n=\psi_1(n+\tfrac12)-\psi_1(n+1),\qquad
 y_n=n+\tfrac14,
\]
with $H_n^{(p)}=\sum_{j=1}^nj^{-p}$ and $H_n=H_n^{(1)}$; $\delta_n>0$ is a
trigamma difference, not a Kronecker delta. With $A=\pi+4\log2$,
$Q=\Gamma(\frac14)^4/\pi$ and Catalan's constant $\Cat$, the input reads
\begin{equation}
 2187\,C_6(1)=E+S_{\mathrm{pol}}+S_{\mathrm{low}}+S_6+S_8
 -393\pi^3W_0+24\pi^3M_7+96P_a,
 \label{gff:c6:reduced}
\end{equation}
with the known polynomial
\begin{align}
 E={}&\frac Q8\left[63A^3+120A\Cat+79A\pi^2+2013\pi^3+54204\zeta(3)\right]
 +Q^2(29\pi-249A)
 +\frac{21}{640}A^5+\frac7{64}A^3\pi^2\notag\\
 &+\frac{1149}{64}A^2\pi^3+\frac{7863}{16}A^2\zeta(3)
 +192A\Cat(16\Cat+3\pi^2)+\frac{14853}{640}A\pi^4
 -\frac{544}{3}\Cat\pi^3\notag\\
 &+\frac{37631}{128}\pi^5+\frac{722279}{48}\pi^2\zeta(3)
 -\frac{805959}{8}\zeta(5),
 \label{gff:c6:Epoly}
\end{align}
which already contains the restored coefficients, and the polylogarithmic
sector
\begin{equation}
 S_{\mathrm{pol}}=4352\bigl[(32\Cat+9\pi^2)D_3-64D_5\bigr],\qquad
 D_3=\operatorname{Im}\operatorname{Li}_{2,1}(\ii,1),\quad
 D_5=\operatorname{Re}\operatorname{Li}_{2,1,2}(\ii,1,1),
\end{equation}
where
$\operatorname{Li}_{s_1,\ldots,s_k}(z_1,\ldots,z_k)
=\sum_{n_1>\cdots>n_k\ge1}\prod_jz_j^{n_j}n_j^{-s_j}$.
The three ordinary sectors are
\begin{align}
 S_{\mathrm{low}}={}&4\pi^3C_d\varpi_4+2\pi^4C_s\varpi_6,\notag\\
 C_d={}&96Q-42A^2-18A\pi-1227\pi^2,\qquad
 C_s=60Q-\tfrac{105}{4}A^2-9A\pi-\tfrac{2469}{4}\pi^2,\notag\\
 S_6={}&24\Bigl[\frac{6Q^3}{\pi}(2\Phi_--1)-104\pi^5\Phi_0
 +\frac{2^{18}\pi^5}{Q^3}\Phi_+\Bigr],\notag\\
 S_8={}&\pi^6(4932-5004\sqrt2)\varpi_8
 +556\pi^2\bigl[(17+12\sqrt2)Y_1-(17-12\sqrt2)Y_7\bigr],
 \label{gff:c6:ordinarySectors}
\end{align}
with the hypergeometric values at unit argument (all parametric excesses
are positive)
\begin{gather*}
 \varpi_4={}_4F_3\bigl((\tfrac12)^{\times4};1^{\times3};1\bigr),\quad
 \varpi_6={}_7F_6\bigl(\tfrac54,(\tfrac12)^{\times6};\tfrac14,1^{\times5};1\bigr),\quad
 \varpi_8={}_9F_8\bigl(\tfrac54,(\tfrac12)^{\times8};\tfrac14,1^{\times7};1\bigr),\\
 \Phi_-={}_7F_6\bigl(1,(\tfrac14)^{\times6};(\tfrac34)^{\times6};1\bigr),\quad
 \Phi_0={}_6F_5\bigl((\tfrac12)^{\times6};1^{\times5};1\bigr),\quad
 \Phi_+={}_7F_6\bigl(1,(\tfrac34)^{\times6};(\tfrac54)^{\times6};1\bigr),\\
 Y_r=t\Bigl[\frac{\Gamma(t+\frac14)}{\Gamma(t+\frac34)}\Bigr]^8
 {}_{10}F_9\bigl(1,1+t,(t+\tfrac14)^{\times8};t,(t+\tfrac34)^{\times8};1\bigr),
 \qquad t=\frac r8.
\end{gather*}
The three remaining sectors are sums with no known evaluation in standard
constants:
\begin{align}
 W_0={}&\sum_{n\ge0}a_n^4\bigl[8+5\pi(4n+1)a_n^2\bigr]\delta_n,\qquad
 M_7=\sum_{n\ge0}(-1)^na_n^7F_n,\notag\\
 F_n={}&(98H_{2n}+74\log2-6\pi)T_n-98\pi\kappa_n
 -98(11u_n-13\pi y_n)\Bigl(\frac{\pi^2}3+4H_{2n}^{(2)}\Bigr)\notag\\
 &-8624y_n\bigl[\zeta(3)-H_{2n}^{(3)}\bigr],\notag\\
 T_n={}&80\pi^2y_n-26\pi u_n+11v_n,\qquad
 \kappa_n=28\pi^2y_n-10\pi u_n+3v_n,\notag\\
 u_n={}&1-14y_nh_n,\qquad
 v_n=196y_nh_n^2-28h_n+7y_n(\delta_n-\pi^2),
 \label{gff:c6:residualSums}
\end{align}
and
\begin{equation}
 P_a=\pi^4\bigl(15I_{22}+24I_{11}-4I_{00}\bigr),\qquad
 I_{ij}=\int_0^1\frac{p_i(x)\,p_j(1-x)}{\pi^{i+j}\sqrt{x(1-x)}}\dd x,
 \label{gff:c6:primitivePairing}
\end{equation}
with $p_0=v(x)^2$, $p_1=\sqrt Q\,v(x)\,w(1-x)$, $p_2=Q\,w(1-x)^2$,
$v(x)={}_2F_1(\frac14,\frac14;1;x)$ and
$w(x)={}_2F_1(\frac14,\frac14;\frac12;x)$. The sum $W_0$ includes $n=0$,
whose term is $(8+5\pi)\pi^2/3$. Its terms decay like $14/(\pi^2n^4)$; those
of $M_7$ alternate in sign and are $O(n^{-5/2}\log n)$; the pairings have
integrable endpoint singularities, so all three converge absolutely.

\paragraph{Special-function form.}
Appendix~\ref{gff:app:C6derivation} converts every sector of \eqref{gff:c6:reduced}
into one value at unit argument, using gamma reflection, sine addition, beta
integration and Mellin--Barnes transforms. The result is the nine-value
formula \eqref{gff:c1:c6nine}, whose values and parameters
Section~\ref{gff:c1:c6analytic} defines. In particular
$D_3=-\frac12\operatorname{Ls}_3(\frac\pi2)-\frac{7\pi^3}{192}$,
$D_5=\operatorname{Cl}_{2,1,2}(\frac\pi2)$,
$W_0=\sqrt{89}\,\mathcal H_W/\pi^2$,
$\pi^{3/2}M_7=\rho_7\,\mathcal G_7+196\pi\sqrt{185}\,\mathcal H_7$ and
$P_a=80\pi^3\mathcal H_*$. The count of nine values describes how the terms
are packaged, not the number of independent constants: sine addition merges
two $G$ values into one with an angle parameter, and $\alpha_4$ depends on
$Q$, $A$ and $\pi$.

\paragraph{Value and comparison with the graph sum.}
The representation evaluates to
\[
 C_6^{\mathrm{rep}}(1)=14331.68703275805432044546325524621711831732\ldots,
\]
with the following sector contributions, in units of $C_6$:
\begin{center}
\begin{tabular}{@{}lr@{\qquad}lr@{}}
\toprule
$E/2187$ & $-1334.069267305$ & $S_8/2187$ & $4446.133424647$\\
$S_{\mathrm{pol}}/2187$ & $-30.841591973$ & $-393\pi^3W_0/2187$ & $-437.775914461$\\
$S_{\mathrm{low}}/2187$ & $-944.671764724$ & $24\pi^3M_7/2187$ & $-5217.069054363$\\
$S_6/2187$ & $3148.152054246$ & $96P_a/2187$ & $14701.829146691$\\
\bottomrule
\end{tabular}
\end{center}
It lies $1.03\times10^{-8}$ below the unrounded value $14331.68703276838$ of
\eqref{gff:c6:answer}, which is $0.021$ of the allowance $U$. The two added
Laurent coefficients contribute $+659.278983$ to $C_6$; without them
\eqref{gff:c6:reduced} would give $13672.408$. Because \eqref{gff:c6:answer} is
computed from the graphs directly, without \eqref{gff:c6:reduced}, the agreement
tests the reduction and its correction numerically at the level of
$3\times10^{-11}$ relative. It is not a proof of the reduction.

The package encloses $C_6^{\mathrm{rep}}(1)$ in an outward-rounded interval of
width below $1.2\times10^{-145}$, using explicit bounds on the truncated tails
of the series for $W_0$, $M_7$ and $P_a$. That enclosure is rigorous for \eqref{gff:c6:reduced} as stated,
conditional on the identities the package uses (Section~\ref{gff:c6d:interval});
the tail of $M_7$ is bounded by \eqref{gff:c6d:M7tail}. An independent review
(\path{LRI_C6_closed_form_review/}) re-evaluated \eqref{gff:c6:reduced} from the definitions above to
about 65 digits, recomputed the nine values with separate codes (the
two-fold integrals directly, to 9--14 digits), and checked the algebra of
the conversions. It found no error. Some conversions rest on further
identities. The coefficient series for $S_6$ and $S_8$ and the two
periodization identities for the wheel are derived in
Appendix~\ref{gff:app:C6derivation}, the latter from three elliptic-moment
evaluations that the package takes over without derivation,
\eqref{gff:c6d:Iseries}--\eqref{gff:c6d:contourM2}; the Barnes representation of
Bailey type \cite{gff:Mishev} and a hypergeometric symmetry are published
identities. All were confirmed numerically, the series and periodization
identities to better than $10^{-63}$.
For $W_0$, $M_7$ and $P_a$ no reduction to standard constants is known, and
a PSLQ search in the package found no relation in the three bases it tested.

\subsection{Result in the physical series}
Combining the retained analytic lower-order coefficients with the
sixth-order value of Section~\ref{gff:c6:analytic} gives
\begin{equation}
\begin{aligned}
 \Delta_{\mathcal E}={}&\frac12-\frac\epsilon6
 +0.657131263981\,\epsilon^2
 -3.719856322902\,\epsilon^3\\
 &+45.42488010356\,\epsilon^4
 -731.13947662454\,\epsilon^5
 +14331.68703275805\,\epsilon^6+O(\epsilon^7).
\end{aligned}\label{gff:c6:physicalSeries}
\end{equation}
The graph sum of this section gives the sixth-order coefficient as
$14331.6870328(5)$, in agreement to $1.0\times10^{-8}$. Neither number is a bound on
the omitted $O(\epsilon^7)$ term.

\section{Finite Schwinger kernels and local subtraction}\label{gff:cn:finiteKernels}
Unless stated otherwise, this section fixes $d=2$, $A=4\log2$,
$a=(2+z)/4$ and $r(2,z)=1$. The kernel normalization is exactly
\eqref{gff:gd:kernel}. For a primitive $L$-loop graph define
\begin{equation}
 a_G=[z^{-1}]\mathcal I_G(2,z),\qquad
 \Pi_G=\frac L2a_G
 =\frac1{\pi^L}\int d\Omega\,\frac{\prod_e x_e^{-1/2}}{U_G(x)}.
 \label{gff:c6d2p:periodNorm}
\end{equation}
The perturbative sign is not included in $a_G$. At six loops every
primitive mass weight vanishes, so
\[
 C_{6,\mathrm{prim}}(2)=\frac2{729}
 \sum_{G\ {\rm primitive}}w_A(G)a_G.
\]
The following local one- and two-subgraph constructions define the finite
kernels that the previous two-dimensional estimator
(Section~\ref{gff:meth:pooling}) retained as numerical inputs,
\eqref{gff:meth:retained}. The 2026 two-dimensional production integrates
all six-loop kernels in every replicate (Section~\ref{gff:prec:d2production2026}).
Their MS conversions are given separately so that their finite-scheme
values are not confused with partial MS contributions. The numerical
values quoted in this section are archived records of that superseded
estimator. They and the programs that produced them are kept in the
research archive \texttt{LRI\_complete\_research\_bundle.zip} (packages
\path{lri_C6_d2_primitive_update/}, \path{lri_C6_d2_subtraction_update/}
and \path{lri_C6_d2_two_forest_update/}), which is not part of this
repository; no current coefficient depends on them.

\subsection{The sixth-order assembly at general dimension}\label{gff:mgc:c6assembly}
The material of this subsection was written by OpenAI GPT-6 astra in the earlier sixth-order
calculations at $d=2$ and $d=3$ (package \path{LRI_C6_D2_HPC/reference/}), with one exception: the
value of $\mathfrak B_{\rm kn}/2187$ below was computed from the closed forms by a Claude Code session on
16 October 2026, when this subsection was merged into the note; the earlier calculation had evaluated it
from numerical master values. The same session rechecked the polynomial \eqref{gff:mgc:knownPolynomial} and
the six-loop totals.

\paragraph{Assembly.}
The sixth-order formula \eqref{gff:c6:RGassembly} is the case $n=6$ of the Lagrange formula
\eqref{gff:c6:Lagrange}, an identity between polynomials in the renormalization-group coefficients. In the
bubble-normalized MS prescription \eqref{gff:gd:normalization}, with
\[
 \beta(g)=-\epsilon g+g^2B(g),\quad B(g)=3-6A_dg+b_3g^2+\cdots+b_6g^5,\qquad
 \gamma_r(g)=g-A_dg^2+\sum_{j=3}^{6}c_jg^j,
\]
it therefore holds at every $d$ with $A$ replaced by $A_d$ of \eqref{gff:gd:blocks}; in particular
$A_2=4\log2$ and $A_3=2+4\log2-\pi$. All $b_j,c_j$ in one assembly must belong to the same
prescription. The sensitivities to $c_6$ and $b_6$ quoted after \eqref{gff:c6:RGassembly} hold at
every $d$ as well, and the physical $C_0,\ldots,C_5$ at a fixed dimension cannot supply the independent
six-loop combination $3c_6-b_6$. For the bare amplitudes \eqref{gff:gd:amplitudes} put
\begin{equation}
 J_5(d)=[z^0]\bigl(8\mathcal B_5(d,z)-3\mathcal A_5(d,z)\bigr),\qquad
 H_6(d)=[z^{-1}]\bigl(3\mathcal B_6(d,z)-\mathcal A_6(d,z)\bigr).
 \label{gff:mgc:bareChannels}
\end{equation}
The Laurent recursion \eqref{gff:gd:MS} then gives, exactly as \eqref{gff:c6:bareAssembly} at $d=1$,
\begin{equation}\boxed{
 C_6(d)=R_6(d)+\frac4{243}J_5(d)-\frac2{729}H_6(d),
}\label{gff:mgc:bareAssembly}\end{equation}
where $R_6(d)$ is the result of the same recursion with the five-loop finite parts and the six-loop
simple poles set to zero and all other inputs retained. The new bare terms are
$b_6=108a_{5,0}-6a_{6,-1}+\cdots$ and $c_6=96e_{5,0}-6e_{6,-1}+\cdots$ in the notation
\eqref{gff:gd:LaurentNotation}, and their combination in \eqref{gff:c6:RGassembly} is exactly
$(36J_5-6H_6)/2187$. The two-, three-, four- and five-loop amplitudes are needed through $z^3$, $z^2$,
$z^1$ and $z^0$, and the six-loop ones through $z^{-1}$, each with all its higher poles; the
one-dimensional span and maxgap finite parts
belong to a different prescription and are not inputs at $d\neq1$. These identities were checked by
exact rational polynomial arithmetic; the checker is in the research archive
\texttt{LRI\_complete\_research\_bundle.zip}, which is not part of this repository.

\paragraph{A fifth-order channel that reenters.}
At fifth order 28 vertex-isomorphism classes have $w_4=3w_2$ and drop out of $3\mathcal B_5-\mathcal A_5$,
whose weight $w_H$ is \eqref{gff:gd:combinedWeight}. In $J_5$ their weight is $8w_2-3w_4=-w_2\neq0$, so all 28 reenter at sixth
order. Reusing only the fifth-order combined-channel data would omit them. This is consistent with the
fifth-order result: these terms cancel there, but not here.

\paragraph{Six-loop inventory.}
The 2,427,005 labeled six-loop 1PI graphs of the table after \eqref{gff:c6:weights} fall into 660
vertex-isomorphism classes. With the weights divided by $|\Aut|$, as in the right side of
\eqref{gff:gd:orbitIdentity},
\[
 \sum_Gw_A=\frac{97987}{64},\qquad \sum_Gw_B=\frac{83711}{576},\qquad
 \sum_G\bigl(3w_B-w_A\bigr)=-\frac{105125}{96}.
\]
Times $7!$ these are the six-loop totals of the table after \eqref{gff:c6:weights}, and the first two are the
zero-dimensional coefficients of \eqref{gff:eq:zeroDimSums}. Every class has twelve separately indexed
internal edges, and its first Symanzik polynomial is stored as a list of complementary spanning-tree
masks; the six-loop inventory contains 140,781 such monomials.

\paragraph{MS bookkeeping at $d=2$.}
The MS-converted sector values below are quoted together with the five-loop mass terms
$-6Ac_5/2187$. The reason is the following exact identity. Use the MS values of $b_3,c_3$
\eqref{gff:gd:bc3}, $b_4,c_4$ \eqref{gff:gs:b4}--\eqref{gff:gs:c4} and $D_5=3c_5-b_5$ \eqref{gff:gs:D5}, and
eliminate $b_5=3c_5-D_5$ in \eqref{gff:c6:RGassembly}; since $-12Ab_5+30Ac_5=12AD_5-6Ac_5$, at $d=2$
\begin{equation}\boxed{
 2187\,C_6(2)=\mathfrak B_{\rm kn}+3c_6(2)-b_6(2)-6A\,c_5(2),\qquad A=4\log2 .
}\label{gff:mgc:MStarget}\end{equation}
With the dictionary $R=0$, $M=-4$, $N=4-A$, $U=-14\zeta(3)$, $V=-4\zeta(3)$ and $X=0$ of
\eqref{gff:gs:psiBasis} and \eqref{gff:gs:MN} at $q=1/2$, writing $Q=Q_2$ and $P_j=P_j(2)$,
\begin{equation}\begin{aligned}
 \mathfrak B_{\rm kn}={}&-\frac{14385}{16}A^5-1719A^3Q+186A^3
 -\frac{771}{2}A^2P_4+\frac{164985}{16}A^2\zeta(3)+174A^2\\
 &-288A\KK_2-108A\LL_2-288AP_5-852AQ^2+240AQ+1341AZ_2-\frac{16}{3}A\\
 &-84P_4Q+32P_4+\frac{5511}{2}Q\zeta(3)+144Q-718\zeta(3).
\end{aligned}\label{gff:mgc:knownPolynomial}\end{equation}
Every coefficient follows by exact expansion of $\mathcal R(b_3,c_3,b_4,c_4)+12AD_5$, with
$\mathcal R$ of \eqref{gff:done:lowerR}. With $Z_2=4\Cat Q_2$ \eqref{gff:gs:Z2exact}, the closed forms
\eqref{gff:gc:L2closed}, \eqref{gff:gc:K2closed}, and $P_4(2)$, $P_5(2)$ of Section~\ref{gff:gc:d2closedForms},
\[
 \frac{\mathfrak B_{\rm kn}}{2187}=-1285.866857775584\ldots,\qquad
 C_6(2)=-1285.866857775584\ldots+\frac{3c_6-b_6-24\log2\,c_5}{2187}.
\]
The fifth-order physical coefficient fixes $D_5$ but not $c_5$ separately, so $c_5$ is a genuine input
of \eqref{gff:mgc:MStarget}. This decomposition was used to organize the partial MS sectors below; the
current estimate does not use it, since it computes the physical combination in the finite scheme of
Section~\ref{gff:cn:finiteAssembly}.

\paragraph{Why a different representation was needed at $d=2$.}
The same 660 six-loop classes occur in every dimension, so the graph count is not what makes $d=2$
harder. On the line, seven vertices can be ordered and replaced by six positive gaps; removing the
common scale leaves five dimensionless gap variables, every distance is a sum of consecutive gaps, and
the collision sectors \eqref{gff:c6:sectorGeometry} have an efficient implementation. The
two-dimensional Schwinger representation has 12 edge parameters and an 11-dimensional projective
simplex. (This describes the representation, not a minimal dimension: in coordinate space seven points
have 14 real coordinates, of which translation, rotation and scale remove four.) Without a linear
ordering and with many more relative configurations, naive tensor quadrature and untuned sampling are
much less effective. An earlier attempt at $d=2,3$ sampled the unsubtracted Schwinger sector integrals
with Taylor endpoint subtraction, both uniformly and with a subset dynamic-programming proposal. The
separate large endpoint subtraction terms caused severe variance and loss of precision: the error
propagated to $C_6(2)$ from $H_6$ alone remained of order $10^5$, and no value was accepted. More samples
do not repair a poor subtraction representation. The common finite kernels, the exact root-scale
coefficients and the regulator-aware fixed-point assembly of this section and of
Sections~\ref{gff:app:C6d2Complete} and~\ref{gff:cn:finiteAssembly} address that obstacle directly. Two
dimensions also have advantages: complex-plane factorization, the complex Selberg integral
\eqref{gff:gd:complexSelberg} and the bilinear ${}_3F_2$ form \eqref{gff:done:kiteFactorization}, which allow
many counterterms to be evaluated analytically.


\subsection{Exact primitive inventory and a product reduction}
Use the complete 660-class six-loop inventory. A graph is declared
primitive only if every proper induced vertex subset $S$ with $|S|\ge2$
satisfies $e(S)<2(|S|-1)$. This selects ten classes. All edge labels,
parallel-line labels, rational Wick weights and spanning-tree polynomials
come from the unchanged input inventory. The class indices in this section
are zero-based indices in the project's 660-class six-loop inventory
(archived as \path{data/inventory6.json} in the research archive), not
indices in any external graph catalog; the same indices identify the
six-loop kernels in \path{LRI_C6_D2_HPC/inputs/inventory.json}.

Completing each selected graph to a four-regular graph on eight vertices
partitions the ten classes into five isomorphism classes:
\begin{center}
\begin{tabular}{@{}c l r l@{}}\toprule
Group & Uncompleted class indices & $\sum w_A$ & Period\\\midrule
0 & $639,654,657$ & $16$ & $P_3(2)P_4(2)$\\
1 & $648$ & $12$ & $\Pi_1$\\
2 & $649,651,655$ & $48$ & $\Pi_2$\\
3 & $653,658$ & $12$ & $\Pi_3$\\
4 & $659$ & $1/6$ & $\Pi_4$\\\bottomrule
\end{tabular}
\end{center}
Each equality within a group follows from conformal completion, not from
numerical comparison. In the present two-dimensional normalization each
critical coordinate edge has absolute-distance power one, and each
completed vertex has incident power four. Inversion of an integration
variable therefore cancels its Jacobian, proving invariance under the
choice of the deleted vertex. This is the weighted completion argument;
the general completion and product identities are reviewed in
\cite{gff:SchnetzD2Focus}.

For the representative $639$, the three-vertex separator is $\{4,5,6\}$.
Removing it leaves components $\{0,1\}$ and $\{2,3,7\}$.
Adding the uniquely required triangles on the separator completes the
components to $K_5$ and the octahedral graph. Fixing the separator at
$0,1,\infty$ makes their integrations independent. Thus their normalized
periods multiply, with no additional measure factor:
\begin{equation}
 \Pi_0=P_3(2)P_4(2)=3Q_2P_4(2).
 \label{gff:c6d2p:product}
\end{equation}
The explicit separator and child-graph isomorphisms are checked in the
archived record
\par\noindent{\small\path{lri_C6_d2_primitive_update/verification/primitive_reduction_checks.json}}\par
\noindent of the research archive.
The complete primitive contribution can therefore be written
\begin{equation}\boxed{
 C_{6,\rm prim}(2)=\frac2{2187}
 \left[48Q_2P_4(2)+12\Pi_1+48\Pi_2+12\Pi_3+\frac16\Pi_4\right].
}\label{gff:c6d2p:fivePeriods}\end{equation}
Only the last four periods require new numerical integration. The first
term is evaluated from the retained ordinary-hypergeometric formula for
$P_4(2)$ and the exact Gamma value of $Q_2$:
\begin{equation}\begin{aligned}
 Q_2&=\frac{\Gamma(1/4)^8}{24\pi^4},\\
 P_3(2)P_4(2)&=9597.2069781654209014729791042940933\ldots,\\
 \frac{96}{2187}Q_2P_4(2)
 &=140.42552505774735658305227770343438\ldots.
\end{aligned}\label{gff:c6d2p:productValue}\end{equation}
The unreduced ten-class integration is retained as an independent check of
this reduction; the product value is not fitted to its output.

\subsection{Positive tropical sampling and a stable cycle determinant}
At $d=2$, the primitive residue of the normalized Schwinger integral is
\begin{equation}
 a_G=\frac{2}{L\pi^L}
 \int_{x_e\ge0}\!d\Omega\;
 \frac{\prod_{e=1}^{2L}x_e^{-1/2}}{U_G(x)}.
 \label{gff:c6d2p:integral}
\end{equation}
There is no endpoint analytic continuation in this primitive integral.
We use tropical importance sampling of the positive projective integrand,
in the framework of \cite{gff:BorinskyD2Focus}. Here its specialized construction
is spelled out to fix the normalization.

In an edge ordering $\sigma$, set the largest edge parameter to one and use
$x_{\sigma(j)}=\prod_{k=j}^{2L-1}t_k$ for $j<2L$.
If $h(S)$ is the cycle rank of the first $|S|$ edges, write
\begin{equation}
 \alpha(S)=\frac{|S|-2h(S)}2,
 \qquad U_G=\left(\prod_{j=1}^{2L-1}t_j^{h(S_j)}\right)\widetilde U_G.
 \label{gff:c6d2p:hepp}
\end{equation}
All proper-prefix $\alpha(S_j)$ are strictly positive for the selected graphs.
The integral over a sector has the form
$\int\prod_jt_j^{\alpha(S_j)-1}\,dt_j/\widetilde U_G$.
The exact sector-normalization recursion is
\begin{equation}
 Z_G(\varnothing)=1,\qquad
 Z_G(S)=\frac{\sum_{e\in S}Z_G(S\setminus\{e\})}{\alpha(S)}
 \quad(S\ne E),\qquad
 Z_G(E)=\sum_{e\in E}Z_G(E\setminus\{e\}).
 \label{gff:c6d2p:dp}
\end{equation}
At each backward step choose the removed edge with probability proportional
to $Z_G(S\setminus\{e\})$. Independently sample
$t_j=U_j^{1/\alpha(S_j)}$, with $U_j$ uniform on $(0,1)$.
For graph weights $w_G$, sample graphs proportionally to $w_GZ_G(E)$.
If $W=\sum_Gw_GZ_G(E)$, each sample of the weighted bare residue is simply
\begin{equation}
 Y=\frac{2W}{L\pi^L}\frac1{\widetilde U_G}.
 \label{gff:c6d2p:estimator}
\end{equation}
Thus the estimator is positive and bounded; it does not sample a difference
of divergent contributions. This positivity does not extend to the omitted
nonprimitive subtraction problem.

A direct evaluation of $\widetilde U_G$ sums all spanning-tree monomials.
For the production runs we instead construct a minimum spanning tree in
the sampled edge order. Let $f_i$ be its $L$ chords and let $e$ range over
its $L$ tree edges. Define
\begin{equation}
 V_{ie}=\begin{cases}
 \sigma_{ie}\sqrt{x_e/x_{f_i}},&e\text{ lies in the fundamental cycle of }f_i,\\
 0,&\text{otherwise},
 \end{cases}
 \qquad \sigma_{ie}\in\{+1,-1\}.
 \label{gff:c6d2p:cycleV}
\end{equation}
The matrix-tree identity gives the scaled form
\begin{equation}\boxed{
 \widetilde U_G=\det(I_L+VV^{\mathsf T}).
}\label{gff:c6d2p:det}
\end{equation}
To see the normalization, the fundamental-cycle matrix has one unit entry
on its own chord and zero entries on all other chords. Factoring
$\prod_i x_{f_i}$ from its Gram determinant leaves
$\det(I+VV^{\mathsf T})$. The minimum-spanning-tree property makes that
product the leading Hepp monomial and ensures $x_e\le x_{f_i}$ on each
fundamental cycle. Hence $|V_{ie}|\le1$ and every eigenvalue of
$I+VV^{\mathsf T}$ is at least one. Cholesky evaluation avoids the poor
conditioning of an unscaled Laplacian near a sector boundary.

An independent 60-digit check compares the monomial sum with this determinant
on 180 configurations, including exponentially small sector variables. The
largest relative difference is $4.49\times10^{-61}$. Production runs also
perform pointwise checks before accumulating the batch estimates. The
graph input and both C++ implementations are in the research archive
(package \path{lri_C6_d2_primitive_update/}), which is not part of this
repository.

\subsection{Normalization and the exact set of graphs}
Work at $d=2$, $\epsilon=2s-2$. With the regulator $z$ used for Laurent
coefficients, $a=\tfrac12+z/4$ and the loop normalization \eqref{gff:gd:kernel} becomes
\begin{equation}
 r(2,z)=1,\qquad
 I_G(z)=\frac{\Gamma(Lz/2)}{\Gamma(\frac12+z/4)^{2L}}
  \int_{\Delta_{2L-1}}\frac{\prod_e x_e^{a-1}}{U_G(x)}\,d\Omega,
 \quad U_G=\sum_T\prod_{e\notin T}x_e.
 \label{gff:loc:normalization}
\end{equation}
Here every internal line is a separate Schwinger variable, including
parallel lines. The weights $w_A,w_B$ are the rational vertex-isomorphism
weights of the six-loop inventory used throughout this section; set
$w_H=3w_B-w_A$. No additional
vertex factorial is introduced in this section.

Let $\mathcal S_{26}$ consist of the classes with exactly one connected,
proper, one-particle-irreducible logarithmic subgraph $\gamma$, and no
power-divergent subgraph. Both $\gamma$ and $G/\gamma$ are primitive.
The class IDs and aggregate weights are as follows; a number $j$ in the
second column denotes the unchanged inventory label \texttt{L6G}$j$,
not its zero-based array index.
\begin{center}\small
\begin{tabular}{@{}cp{0.64\textwidth}rr@{}}
\toprule
$h=\ell(\gamma)$ & Class numbers & Number & $\sum w_H$\\\midrule
1 & 524, 529, 539, 563, 575, 587, 588, 596, 599, 607, 613, 614,
    617, 619, 620, 621, 627 &17&$-279/2$\\
3 & 642 &1&$-4$\\
5 & 638, 639, 641, 647, 648, 651, 653, 657 &8&$-93/4$\\\bottomrule
\end{tabular}
\end{center}
The total weight is $-667/4$. The archived machine-readable record gives every
edge list, the subgraph edge mask, the rational weights, and the exact
isomorphism to the primitive subgraph and quotient. A separate checker
also reconstructs their spanning-tree polynomials; its 158 exact checks
include both aggregate conversion identities below.

\subsection{A common-integrand subtraction with no large endpoint difference}
Let $h=\ell(\gamma)$, $k=L-h$, and use $\rho=G/\gamma$ for the quotient.
Define
\begin{equation}
 X_\gamma=\sum_{e\in\gamma}x_e,\quad
 X_\rho=\sum_{e\notin\gamma}x_e,\quad
 S=X_\gamma+X_\rho,\quad
 U_0=U_\gamma U_\rho,\quad U_\times=U_G-U_0.
 \label{gff:loc:U0}
\end{equation}
The edge variables of the quotient are the original edges outside
$\gamma$. The monomials of $U_0$ are a subset of those of $U_G$;
therefore $U_\times$ has nonnegative coefficients. This is checked
exactly, not by evaluating the polynomials at selected points.

The auxiliary amplitude is
\begin{equation}
 \boxed{D_G(z)=I_G(z)-\frac{k}{L}I_\gamma(z)I_\rho(z).}
 \label{gff:loc:Ddefinition}
\end{equation}
Its common projective kernel is
\begin{equation}
 \boxed{\frac1{U_G}-\frac{X_\rho}{S\,U_0}
 =\frac{X_\gamma U_0-X_\rho U_\times}
 {S\,U_0\,U_G}.}
 \label{gff:loc:cancelledKernel}
\end{equation}
In simplex gauge $S=1$; retaining $S$ makes the homogeneity correct in
any Hepp gauge. This is not a replacement of the product by a local
approximation: its integral is exactly the product term in
\eqref{gff:loc:Ddefinition}.

\paragraph{Proof of the factor $k/L$.}
In the product integral, separate the two sets of variables by
$x_\gamma=\lambda u$ and $x_\rho=(1-\lambda)v$, with $u,v$ each on
a unit simplex. Their scale factors are
$\lambda^{hz/2-1}(1-\lambda)^{kz/2-1}$.
Without a scale weight the beta integral exactly combines the two
Gamma prefactors into $\Gamma(Lz/2)$. Inserting $X_\rho/S=1-\lambda$
multiplies the beta integral by
\begin{equation}
 \frac{B(hz/2,1+kz/2)}{B(hz/2,kz/2)}=\frac{k}{L}.
 \label{gff:loc:betaProof}
\end{equation}
The identity holds in the ordinary convergence region and then by
meromorphic continuation. Hence the subtraction has the correct
normalization at nonzero $z$, not just at leading pole order.

\paragraph{Endpoint cancellation.}
When $\gamma$ shrinks, $U_\times$ has at least one extra power of the
subgraph scale compared with $U_0$. The first numerator term in
\eqref{gff:loc:cancelledKernel} has this power through $X_\gamma$.
The artificial singularity when the complementary factor shrinks
is removed similarly, using the fact that it is not a proper
subdivergence of the original $G$. Since the two factors are primitive,
these are the only zero exponents introduced by the product denominator.
The cancellation is applied to the monomial exponents before floating
point evaluation. The implementation never forms two nearly equal,
individually singular reciprocal polynomials.

More explicitly, order the $2L$ edges, write the ordered Schwinger
variables as products of $t_j\in(0,1)$, and let $E_j$ be the first $j$
edges. If
\begin{equation}
 h_0(E_j)=\min_{m\in\operatorname{supp} U_0}|m\cap E_j|,\qquad
 \alpha_j=\frac j2-h_0(E_j)
       +\mathbf1_{E_j=E_\gamma\ \mathrm{or}\ E_j=E_\rho},
 \label{gff:loc:alpha}
\end{equation}
then every proper-sector $\alpha_j$ is positive for these graphs.
The indicated powers are stripped from the numerator and $t_j$ is
sampled with density $\alpha_jt_j^{\alpha_j-1}$. Sector probabilities
are proportional to the corresponding product of $1/\alpha_j$ and
are generated by a subset dynamic program. The normalization for the
simple residue is $2/(L\pi^L)$. The remaining rational kernel is bounded
on each closed sector after these structural powers have been removed;
this distinguishes it from an earlier estimator of the same graphs, which
sampled the unsubtracted integrands with importance weights taken from their
logarithmic moments and was much noisier.

Define the quantity directly measured by this sampler as
\begin{equation}
 Y_{26}=\sum_{G\in\mathcal S_{26}}w_H(G)[z^{-1}]D_G(z).
 \label{gff:loc:Y}
\end{equation}
$Y_{26}$ is not an MS residue. Its conversion is indispensable.

\subsection{Conversion to the same MS coefficient as the six-loop assembly}
Write the primitive factors as
\begin{equation}
 I_\gamma=\frac{r_\gamma}{z}+f_\gamma+O(z),\qquad
 I_\rho=\frac{r_\rho}{z}+f_\rho+O(z).
 \label{gff:loc:primitiveLaurent}
\end{equation}
For a single proper primitive subdivergence the incomplete MS operation
is $R'_GI_G=I_G-(r_\gamma/z)I_\rho$. Taking only the simple pole gives
\begin{equation}
 \boxed{[z^{-1}]R'_GI_G=[z^{-1}]D_G
       -\frac hLr_\gamma f_\rho
       +\frac kLf_\gamma r_\rho.}
 \label{gff:loc:MSsingle}
\end{equation}
Higher poles of $R'_GI_G$ are not assumed to vanish. Equation
\eqref{gff:loc:MSsingle} is a Laurent-coefficient identity and does not confuse
this auxiliary kinematic subtraction with MS.

There are three primitive five-loop inventory classes. In zero-based
array notation they are 123, 125 and 126, with coupling weights
$2,3/2,12$, respectively. Let $f_i$ denote their individual unweighted
finite parts, evaluated in the massive normalization
\eqref{gff:loc:normalization}, and put
\begin{equation}
 a_{5,0}^{\rm p}=2f_{123}+\frac32f_{125}+12f_{126}.
 \label{gff:loc:a50}
\end{equation}
The finite parts of conformally equivalent completed graphs are not
identified: completion identities fix their residues, not their
mass-dependent finite terms. All three terms in \eqref{gff:loc:a50} were
integrated separately.

The one-loop bubble is exactly $I_1=2/z$, with zero finite part.
For the $h=3$ class the two primitive factors are isomorphic, so its
conversion terms in \eqref{gff:loc:MSsingle} cancel. Exact aggregation over
the other 25 classes gives
\begin{equation}
 \boxed{M_{26}:=\sum_{G\in\mathcal S_{26}}w_H[z^{-1}]R'_GI_G
              =Y_{26}+\frac52a_{5,0}^{\rm p}.}
 \label{gff:loc:MSaggregate}
\end{equation}
For example, the three finite-part coefficients on the right are
$5,15/4,30$, which are exactly $5/2$ times the weights in
\eqref{gff:loc:a50}. The direct contribution to the physical assembly is
\begin{equation}
 \boxed{\Theta_{26}=\frac{(3c_6-b_6)_{26}}{2187}
             =-\frac{2}{729}M_{26}.}
 \label{gff:loc:Theta}
\end{equation}
The last sign follows from the even-loop sign in the bare amplitudes:
$(3c_6-b_6)_{26}=-6M_{26}$.

That earlier estimator measured the same weighted simple-pole coefficient
of the unsubtracted amplitudes,
\[
 H_{26}^{\rm bare}=\sum_{G\in\mathcal S_{26}}w_H(G)[z^{-1}]I_G(z).
\]
For a comparison with it, let $r_3,f_3$ be the unweighted primitive
three-loop residue and finite part. A second exact aggregate identity is
\begin{equation}
 \boxed{H_{26}^{\rm bare}=Y_{26}-\frac{31}{2}a_{5,0}^{\rm p}-4r_3f_3,
        \qquad r_3=2Q_2.}
 \label{gff:loc:rawAggregate}
\end{equation}
Comparing $Y_{26}$ directly to $H_{26}^{\rm bare}$ would therefore give a spurious
normalization discrepancy.

\subsection{The two-subgraph classes and a common finite kernel}
There are 40 six-loop classes with exactly two proper logarithmic
subgraphs and no power-divergent subgraph: 23 have a nested pair and 17
have a vertex-disjoint pair. Every decomposition used here has two
one-loop primitive leaves and one four-loop primitive leaf. At five
loops there are 12 analogous classes (seven nested, five disjoint);
only five of them carry a nonzero mass weight. Their primitive leaves
are two bubbles and the three-loop tetrahedron.

For $G$ let $U=U_G$ and let $U_1,U_2$ be its leading polynomials when
respectively $\gamma_1,\gamma_2$ shrink. For a subgraph of loop number
$h_i$, $U_i$ consists exactly of the monomials of $U$ containing $h_i$
edges from $\gamma_i$. The common face is $W=U_{12}$. Set
\begin{equation}
 X=U_1-W,\qquad Y=U_2-W,\qquad
 Z=U-U_1-U_2+W.
 \label{gff:two:facePartition}
\end{equation}
Thus $U=W+X+Y+Z$, $U_1=W+X$, $U_2=W+Y$.
The four polynomials partition the spanning-tree monomials and have
nonnegative coefficients. For nested subgraphs put
\begin{equation}
 \sigma_1=\sum_{e\in\gamma_1}x_e,\quad
 \sigma_2=\sum_{e\in\gamma_2\setminus\gamma_1}x_e,\quad
 \sigma_3=\sum_{e\notin\gamma_2}x_e,
 \qquad S=\sigma_1+\sigma_2+\sigma_3.
 \label{gff:two:nestedScales}
\end{equation}
For disjoint subgraphs, the three sums run over $\gamma_1$, $\gamma_2$
and their complement.

The two auxiliary forest kernels, including the common projective scale
weight but not the common propagator measure, are
\begin{align}
 K_{\rm n}&=\frac1U-\frac{\sigma_2+\sigma_3}{S U_1}
           -\frac{\sigma_3}{S U_2}+\frac{\sigma_3}{S W},
 \label{gff:two:nestedKernel}\\
 K_{\rm d}&=\frac1U-\frac{\sigma_2+\sigma_3}{S U_1}
           -\frac{\sigma_1+\sigma_3}{S U_2}+\frac{\sigma_3}{S W}.
 \label{gff:two:disjointKernel}
\end{align}
These are auxiliary subtractions, not MS graph contributions; the conversion follows below.
For stable evaluation write either kernel as
\begin{equation}
 K_\star=\frac{N_\star}{S\,U\,U_1\,U_2\,W},
 \label{gff:two:commonDenominator}
\end{equation}
where the numerators are the following explicit finite polynomials:
\begin{align}
 N_{\rm n}={}&\sigma_1W(W+X)(W+Y)
 -\sigma_2(Y+Z)W(W+Y)\nonumber\\
 &+\sigma_3\left[XY(2W+X+Y+Z)-ZW^2\right],
 \label{gff:two:nestedNumerator}\\
 N_{\rm d}={}&-\sigma_1W(W+X)(X+Z)
 -\sigma_2W(W+Y)(Y+Z)\nonumber\\
 &+\sigma_3\left[XY(2W+X+Y+Z)-ZW^2\right].
 \label{gff:two:disjointNumerator}
\end{align}
Equations \eqref{gff:two:nestedNumerator}--\eqref{gff:two:disjointNumerator}
are exact rational identities. The code evaluates these polynomials,
not a floating-point subtraction of the four singular fractions.

\paragraph{Why the subtracted integrals converge.}
In a Hepp sector, choose the largest edge equal to one and write
$x_{\pi(j)}=\prod_{k=j}^{2L-1}t_k$. For a proper edge subset $s$ let
$h_P(s)$ be the smallest degree of the polynomial $P$ in the edges of $s$.
For a scale sum $\sigma_i$ this is one if every edge in its group belongs
to $s$, and zero otherwise. Let $h_N(s)$ be the minimum degree among the
nonzero products displayed in \eqref{gff:two:nestedNumerator} or
\eqref{gff:two:disjointNumerator}. No cancellation of leading terms is used
to increase this minimum. The residual exponent is
\begin{equation}
 \alpha_s=\frac{|s|}{2}+h_N(s)-h_U(s)-h_{U_1}(s)-h_{U_2}(s)-h_W(s).
 \label{gff:two:sectorRate}
\end{equation}
An exact integer check over all proper edge subsets gives
$\alpha_s\ge1/2$ for each of the 40 six-loop and 12 five-loop kernels:
there are $40(2^{12}-2)+12(2^{10}-2)=176024$ checks. This controls all
nested endpoint limits, not just separate limits of individual edges.
It establishes a common analytic neighborhood of $z=0$ for the angular
integral after removing the overall $\Gamma(Lz/2)$ pole.

For completeness, the sector weights are computed by
\begin{equation}
 D(\varnothing)=1,\qquad
 D(s)=\frac1{\alpha_s}\sum_{e\in s}D(s\setminus\{e\}),
 \qquad \alpha_{\mathrm{all}}:=1.
 \label{gff:two:sectorDP}
\end{equation}
A descending edge choice has probability proportional to
$D(s\setminus\{e\})$; each continuous coordinate has density
$\alpha_s t^{\alpha_s-1}$. Graphs are sampled proportionally to
$|w_G|D(\mathrm{all})$. After dividing numerator and denominator
polynomials by their exact leading Hepp monomials, the residual factor
has no endpoint pole. The overall residue prefactor is $2/(L\pi^L)$.
The numerical estimates below include the graph sign; positivity of the
primitive estimator is not asserted for the forest estimator.

\subsection{Exact conversion of the auxiliary residues to MS}
\label{gff:two:conversion}
The conversion needs finite parts, and at this order also the coefficient
linear in $z$ of a primitive graph. They cannot be replaced by period
residues. Use the unweighted amplitudes
\begin{equation}
 I_1=\frac2z,\qquad
 I_p=\frac{r_p}{z}+f_p+a_p z+O(z^2),\qquad p=3,4,
 \label{gff:two:primitiveJets}
\end{equation}
where
\begin{equation}
 r_3=2Q_2,\qquad r_4=P_4(2)/2,\qquad
 Q_2=\frac{\Gamma(1/4)^8}{24\pi^4}.
 \label{gff:two:residueReferences}
\end{equation}
For a graph with one logarithmic child $\gamma$ of loop number $h$, define
\begin{equation}
 D_G(z)=I_G(z)-\frac{L-h}{L}I_\gamma(z)I_{G/\gamma}(z)
       =\frac{d_G}{z}+f_G+O(z).
 \label{gff:two:singleD}
\end{equation}
For the two-loop auxiliary graph this gives
\begin{equation}
 D_2=\frac{A}{z}+f_2+O(z),\qquad
 f_2=-4.8368122825689194750929364789022677856768458003883\ldots .
 \label{gff:two:f2}
\end{equation}
The value in \eqref{gff:two:f2} was checked with the two-${}_3F_2(1/4)$ formula
\eqref{gff:gd:twoLoopHyper} and two Cauchy grids; it is a finite part in the
specified massive prescription, not a physical epsilon coefficient.

\paragraph{General nested and disjoint identities.}
For a nested pair, write $\gamma_1=G_a$, $\gamma_2=J$, $J/\gamma_1=G_b$,
$G/\gamma_2=G_c$, $G/\gamma_1=T$, with primitive loop numbers $h,k,m$ and
$L=h+k+m$. Beta integration of the relative scale variables gives the
exact all-$z$ identity
\begin{equation}
 D_G=I_G-\frac{k+m}{L}I_aI_T-\frac mL I_JI_c
                    +\frac mL I_aI_bI_c.
 \label{gff:two:nestedD}
\end{equation}
Meanwhile
\begin{equation}
 R'I_G=I_G-\frac{r_a}{z}I_T+Z_J I_c,\qquad
 Z_J=-\KP\left(I_J-\frac{r_a}{z}I_b\right).
 \label{gff:two:nestedMS}
\end{equation}
where $\KP$ keeps the negative powers of $z$ (the pole part), as in \eqref{gff:eq:MSrecursion}. For a disjoint pair with primitive children $G_a,G_b$ and primitive
cograph $G_c$, the analogous identities are
\begin{align}
 D_G={}&I_G-\frac{L-h}{L}I_aI_{G/G_a}
             -\frac{L-k}{L}I_bI_{G/G_b}
             +\frac mL I_aI_bI_c,
 \label{gff:two:disjointD}\\
 R'I_G={}&I_G-\frac{r_a}{z}I_{G/G_a}
              -\frac{r_b}{z}I_{G/G_b}
              +\frac{r_ar_b}{z^2}I_c.
 \label{gff:two:disjointMS}
\end{align}
The factors $m/L$ and $(L-h)/L$ arise from the beta mean of the remaining
scale fraction, not from fitted subtraction constants. Substituting
\eqref{gff:two:singleD} into \eqref{gff:two:nestedD}--\eqref{gff:two:disjointMS}
and extracting $z^{-1}$ proves each graph's conversion. All graph
identifications, including the inherited subgraph vertex labels, are
checked against the six-loop inventory before the rational weights are
collected.

Let $\mathcal S_{6,2}$ be the set of the 40 six-loop classes with exactly two
proper logarithmic subgraphs, let
$Y_{40}=\sum_{G\in\mathcal S_{6,2}}w_H(G)[z^{-1}]D_G$, and abbreviate
$f_{5,j}=f_{\mathrm{L5G}j}$, $d_{5,j}=d_{\mathrm{L5G}j}$. Define
\begin{equation}\begin{aligned}
 F_5={}&\frac{15}{4}f_{5,052}+\frac{15}{2}f_{5,079}
       +13f_{5,100}+26f_{5,114}
       +\frac{13}{4}f_{5,125}+\frac{13}{2}f_{5,129}.
\end{aligned}\label{gff:two:F5}\end{equation}
The exact aggregate result is
\begin{equation}\boxed{\begin{aligned}
 M_{40}&:=\sum_{G\in\mathcal S_{6,2}}w_H(G)[z^{-1}]R'I_G\\
       &=Y_{40}+F_5+9Af_4-15a_4-18f_2r_4,\\
 \Theta_{40}&=-\frac2{729}M_{40}.
\end{aligned}}\label{gff:two:conversion40}\end{equation}
Thus the primitive four-loop coefficient $a_4$, not just its finite part,
is essential. The same integrations give the entire one-subdivergence
five-loop mass sector:
\begin{equation}\boxed{
 c_5^{(1)}=\frac52d_{5,052}+5d_{5,079}+3f_4.}
 \label{gff:two:c5single}
\end{equation}
Only the first two of the six fifth-order classes have nonzero mass
weights, although all six finite parts enter \eqref{gff:two:F5}.

\paragraph{A further five-loop mass sector.}
The mass-carrying two-subgraph classes are exactly
$\mathrm{L5G053},\mathrm{L5G054},\mathrm{L5G056},\mathrm{L5G077},
\mathrm{L5G080}$, with respective $w_B=(1,1,1/2,1/2,1)$.
All five have nested subdivergences. Put
$Y_{5B}=\sum w_B(G)[z^{-1}]D_G$ over this set, and define $f_{4,j}$
using \eqref{gff:two:singleD} for the indicated fourth-order class.
The independent rational collection gives
\begin{equation}\boxed{\begin{aligned}
 c_5^{(2)}=5\left[
 Y_{5B}+\frac{2A}{5}f_3-a_3-\frac35 f_2r_3
 -\frac45f_{4,013}+\frac65f_{4,024}\right].
\end{aligned}}\label{gff:two:c5double}\end{equation}
These classes are disjoint from the one-subgraph and closed Gamma
classes. Their contribution to the sixth-order combination is
$-6Ac_5^{(2)}/2187$, not a sixth-order primitive period.

\subsection{Finite-part integrands and numerical results}
For any convergent auxiliary kernel $D_G$, differentiation of the common
propagator measure at $z=0$ gives the logarithmic multiplier
\begin{equation}
 \mathcal J_L=L\log2+\frac14\sum_{j=1}^{2L-1}j\log t_j
                         -\frac L2\log S.
 \label{gff:two:finiteMultiplier}
\end{equation}
Multiplying its residue integrand by \eqref{gff:two:finiteMultiplier} gives
$f_G$. For primitive jets through $a_p$, the code uses the full Taylor
coefficients of
$z\Gamma(pz/2)/\Gamma(1/2+z/4)^{2p}$ and the exponential of the logarithmic
measure, retaining the exact factorials. No new period identity for the
massive finite parts is assumed.

Each of the two main production runs uses three independent streams:
$64\,000\,000$ observations for $Y_{40}$, the six fifth-order auxiliary
jets, and the primitive four-loop jets, respectively. Each stream is
split into 128 batches. Runs A and B use distinct seeds. Primitive
residues are fixed-expectation controls for their finite and linear
coefficients; the control slopes are fitted only to excluded
one-million-sample pilots. The fourth-order slopes are
$(-8.0444695540,40.8894797679)$ for $(f_4,a_4)$.

\begin{center}\small
\begin{tabular}{@{}lrr@{}}\toprule
Quantity & Run A & Run B\\\midrule
$Y_{40}$ & $761.536908\pm42.870149$ & $790.704780\pm44.081215$\\
$\Theta_{40}$ & $220.727969\pm0.224779$ & $220.412578\pm0.227748$\\
$c_5^{(1)}$ & $1651.795516\pm1.874636$ & $1649.773937\pm1.603267$\\
$\Theta_{40}-6Ac_5^{(1)}/2187$
 & $208.163498\pm0.218856$ & $207.863484\pm0.223344$\\\bottomrule
\end{tabular}
\end{center}
The displayed errors are sampling standard errors, not certified
uncertainties. Like every partial sum of this section, $\Theta_{40}$ is
neither an estimate nor a bound for $C_6(2)$; graph-sector contributions are
defined only in the stated MS prescription. The main runs agree within one combined standard error
for the MS six-loop result. Equal weighting, without fitting to a target
value, gives
\begin{align}
 Y_{40}&=776.1208440\pm30.7449312,&
 F_5&=6208.4320415\pm48.4413942,\nonumber\\
 f_4&=-957.9657216\pm0.0502308,&
 a_4&=4958.8036029\pm0.6170590,\nonumber\\
 M_{40}&=-80397.8646408\pm58.3185778,&
 c_5^{(1)}&=1650.7847264\pm1.2333617.
 \label{gff:two:inputNumerics}
\end{align}
The auxiliary residue $Y_{40}$ and the MS residue $M_{40}$ differ
substantially; interchanging them would give an incorrect coefficient.
Keeping the within-stream covariance of finite parts and residues yields
\begin{equation}\boxed{
 \Theta_{40}-\frac{6A}{2187}c_5^{(1)}
 =208.01349115\ldots,\qquad
 \widehat\sigma=0.15634950\ldots .}
 \label{gff:two:jointNumerical}
\end{equation}
Without the primitive controls, the same raw observations give
$\Theta_{40}=220.5459148\pm0.1606925$ and the joint combination
$207.9871304\pm0.1571710$. Controls were not fitted to an assumed value
of $C_6$ or to equality of the two runs.

The additional five-loop calculation uses two independently seeded runs,
each with $32\,000\,000$ observations in each of its three streams
($Y_{5B}$, the two fourth-order auxiliary finite parts, and primitive
three-loop jets), again with 128 batches. It gives
\begin{align}
 c_5^{(2)}&=2431.5047824\pm1.8267989,\nonumber\\
 -\frac{6A}{2187}c_5^{(2)}&=-18.49537102\pm0.01389564.
 \label{gff:two:additionalMassNumber}
\end{align}
The two independent estimates of $c_5^{(2)}$ differ by $0.053$ combined
standard errors. The analysis retains the covariance of $f_3$ and $a_3$,
as well as that of every residue/finite-part pair from the same stream.

\paragraph{Checks of the two-subgraph kernels.}\label{gff:mgc:twoChecks}
These checks were made by OpenAI GPT-6 astra with the kernels above (research-archive package
\path{lri_C6_d2_two_forest_update/}; the result record is
\path{lri_C6_d2_complete_update/reproduction/lri_C6_d2_two_forest_update/verification/C6_d2_continuation_result.json}).
The rational numerator identities \eqref{gff:two:nestedNumerator} and \eqref{gff:two:disjointNumerator} were
checked symbolically. An independent scalar implementation compares the normalized floating-point
numerator with the original four fractions at 500 digits; over 312 configurations, including a
coordinate as small as $10^{-140}$, the largest normalized error is $2.49\times10^{-16}$. This is a
pointwise test of the kernel, not an integration error bound.

The third-order nested and disjoint instances of the kernel have the unweighted auxiliary residues
\[
 [z^{-1}]D_{\mathrm{L3G004}}=A^2,\qquad
 [z^{-1}]D_{\mathrm{L3G007}}=\frac{A^2+4f_2}{3},
\]
which follow from the Gamma counterterms of Section~\ref{gff:two:gammaSector} and the exact two-loop finite
part \eqref{gff:two:f2} through the conversion of Section~\ref{gff:two:conversion}. A sixteen-million-observation
control, with the graph weights inserted, deviates by $-0.185$ and $-0.180$ standard errors; a separate
four-million-observation run by $-0.192$ and $1.542$. For the finite-part extension a
sixteen-million-observation two-loop run gives $[z^{-1}]D_2=2.77218438\pm0.00048568$ and
$[z^0]D_2=-4.83492035\pm0.00210898$, at $-0.833$ and $0.897$ standard errors from $A$ and
\eqref{gff:two:f2}. The reference $r_4=P_4(2)/2$ was evaluated afresh from \eqref{gff:fp:P4two} at 65 digits,
and two Cauchy evaluations of $f_2$ differ by $2.0\times10^{-64}$; this is consistency of the
numerical special-function evaluation, not an interval certificate.


\section{\texorpdfstring{Analytic power-sector compensation at $d=2,3$}{Analytic power-sector compensation at d=2,3}}\label{gff:app:C6d2Complete}\label{gff:app:C6d3Complete}

\subsection{An exact Gamma-reducible sector}
\label{gff:two:gammaSector}
The momentum-space normalization at $d=2$ has $r(2,z)=1$. With one
nonzero momentum $p$, a convolution of $n$ parallel propagators with
indices $\alpha_i$ reduces to a single propagator of index
$\eta=\sum_i\alpha_i-(n-1)$ and multiplier
\begin{equation}
 \mathcal B_n(\boldsymbol\alpha)
 =\frac{\Gamma(\eta)}{\Gamma(1-\eta)}
   \prod_{i=1}^n\frac{\Gamma(1-\alpha_i)}{\Gamma(\alpha_i)},
 \qquad I_{\parallel}(p)=\mathcal B_n(\boldsymbol\alpha)(p^2)^{-\eta}.
 \label{gff:two:parallel}
\end{equation}
The integration convention in this formula is $d^2k/\pi$ per loop.
Successive use of the two-propagator convolution proves the displayed
$n$-propagator version by telescoping Gamma factors. A degree-two vertex
with zero momentum injection adds the two adjacent indices. Initially
every line has index $1/2+z/4$; the final line of an accepted $L$-loop
reduction has index $Lz/2$. Thus the bare amplitude is a finite product of
Gamma functions at affine half-integer arguments.

This rearrangement is deliberately restricted. All original degree-two
vertices must be among the two momentum terminals. Both terminals are
inherited under every contraction. An internal index reaching the
infrared boundary $\alpha=1$ at $z=0$, an identically singular Gamma
factor, a tadpole quotient, a scaleless quotient, or an unsolved cograph
causes that reduction to be rejected. Power-divergent mass subgraphs are
analytically continued in the existing critical-mass prescription; they
are not assigned an extra logarithmic vertex counterterm. These
restrictions distinguish the accepted rearrangements from an unrestricted
infrared rearrangement that would require additional $R^*$ counterterms
\cite{gff:Herzog,gff:BaldufPanzerTropicalLocal}.

As in \eqref{gff:two:nestedMS}, $\KP$ projects a Laurent series onto its negative powers of $z$.
The UV counterterm is computed recursively as
\begin{equation}
 Z_G=-\KP\left[I_G+
 \sum_{\substack{S\ne\varnothing\ \mathrm{proper}\ \mathrm{spinney}}}
 \left(\prod_{\gamma\in S}Z_\gamma\right)I_{G/S}\right].
 \label{gff:two:MSrecursion}
\end{equation}
Spinney components here are vertex-disjoint. Nested subdivergences are
already included recursively in $Z_\gamma$. Confusing this convention
with the edge-compatible scale-derivative forests of the finite kernels of
Section~\ref{gff:cn:finiteKernels} changes individual graph counterterms. Formula
\eqref{gff:two:MSrecursion} is the usual MS $R$-operation; the special graph
selection, Gamma reductions and exact arithmetic are the present
calculation \cite{gff:Herzog}.

If a graph is a one-vertex join of two or more logarithmic blocks, its
amplitude factorizes. Induction in \eqref{gff:two:MSrecursion}, or the
multiplicativity of the counterterm character for this join, gives
\begin{equation}
 Z_{G_1\vee\cdots\vee G_k}=\prod_{j=1}^kZ_{G_j},\qquad k\ge2.
 \label{gff:two:articulation}
\end{equation}
Every factor on the right has strictly negative Laurent degree, so the
product has no simple pole. This establishes the simple-pole zero even
when the individual finite amplitudes are not evaluated. Only joins
whose blocks are all logarithmic are used in this additional zero test.

\paragraph{Exact arithmetic and closed counts.}
Gamma recurrence, reflection and half-integer polygamma values reduce the
Laurent arithmetic to an exact rational polynomial ring in $\ell=\log2$
and zeta values \cite{gff:DLMF5,gff:DLMFPolygamma}. Euler's constant cancels
explicitly. No integer-relation fit is used. Requiring both the original
graph and every necessary cograph/counterterm to close gives
\begin{center}
\begin{tabular}{@{}lrrrrrr@{}}\toprule
Loop order & 1&2&3&4&5&6\\\midrule
Closed Gamma-sector classes &1&2&6&15&47&137\\\bottomrule
\end{tabular}
\end{center}
The bare two-terminal graph being reducible by itself is not sufficient:
unsupported cograph finite parts are never filled with zero.

Define $b_L^\Gamma,c_L^\Gamma$ by precisely this closed class set, with
all original rational graph weights. Two useful exact results are
\begin{equation}
 c_5^\Gamma=
 \frac{12692}{3}\ell^4-208\ell^2+\frac{1085}{4}\ell\zeta(3)
 -\frac{224}{3}\ell-\frac{73}{6},
 \label{gff:two:c5Gamma}
\end{equation}
\begin{equation}\boxed{\begin{aligned}
 2187\,\Xi_\Gamma
 &:=(3c_6^\Gamma-b_6^\Gamma)-24\ell c_5^\Gamma\\
 &=\frac{5518904}{75}\ell^5-\frac{216256}{45}\ell^3
   +\frac{5763}{20}\ell^2\zeta(3)-2784\ell^2
   +\frac{6368}{15}\ell-69\zeta(3)
   +\frac{83397}{640}\zeta(5).
\end{aligned}}\label{gff:two:GammaAnswer}\end{equation}
Numerically,
\begin{equation}
 \frac{3c_6^\Gamma-b_6^\Gamma}{2187}
 =12.17588503555361667474529\ldots,\qquad
 \Xi_\Gamma=4.274679502470384138347130\ldots .
 \label{gff:two:GammaNumbers}
\end{equation}
Neither expression includes the missing graph classes. The full exact
counterterm of every accepted graph, not just the weighted totals, is
recorded in \path{sp_exact.json} of
\path{lri_C6_d2_complete_update/reproduction/lri_C6_d2_two_forest_update/analytic/verification/}.

For example, the closed third-order sector gives
$b_3^\Gamma=336\ell^2-4$ and $c_3^\Gamma=52\ell^2-4/3$.
Restoring the retained wineglass residue $A^2/3$ with its original weights,
and the primitive tetrahedron for the beta function, recovers
\begin{equation}
 c_3=\frac72A^2-\frac43,\qquad
 b_3=\frac{87}{4}A^2-4+6Q_2.
 \label{gff:two:lowerGammaCheck}
\end{equation}
This check takes the wineglass residue and the tetrahedron as inputs; it is
not a new independent evaluation of either. The two-loop
checks $b_2=-6A,c_2=-A$ are entirely within the new Gamma calculation.

\paragraph{Checks of the Gamma sector.}\label{gff:mgc:gammaChecks}
These checks were made by OpenAI GPT-6 astra (code \path{sp_exact.py}, \path{sp_ms.py} and
\path{check_gamma_coverage.py} in, and records \path{sp_exact.json} and \path{sp_ms_90.json} under,
\path{lri_C6_d2_complete_update/reproduction/lri_C6_d2_two_forest_update/analytic/}). The 793 pole
coefficients of the exact table were compared with 90-digit Gamma and polygamma arithmetic; the largest
absolute difference is below $4.5\times10^{-78}$. There are 214 accepted terminal-routing instances
through six loops. Wherever alternative routings are admitted their complete pole parts agree, and
changing $\log p^2$ of the terminal momentum from zero to one leaves the pole parts unchanged; exact
arithmetic checks both cancellations before numerical specialization, and the 90-digit residuals are
below $3.8\times10^{-85}$ and $2.0\times10^{-85}$.


\subsection{An analytic advance special to two dimensions}
Write \(d^2x/\pi\) for the normalized complex-plane measure. Riesz
convolution gives
\begin{equation}
 \mathcal B(u,v)=
 \frac{\Gamma(1-u/2)\Gamma(1-v/2)\Gamma((u+v)/2-1)}
 {\Gamma(u/2)\Gamma(v/2)\Gamma(2-(u+v)/2)}.
 \label{gff:done:Riesz}
\end{equation}
It replaces two distance powers \(u,v\) joined at an integrated vertex by
the power \(u+v-2\). This formula is used meromorphically, with all
normalizing factors retained.

For the five-edge two-variable integral define
\begin{equation}
 \mathcal J(\boldsymbol p)=\frac1{\pi^2}\int_{\mathbb C^2}
 \frac{d^2x\,d^2y}{|x|^{p_1}|1-x|^{p_2}|y|^{p_3}|1-y|^{p_4}|x-y|^{p_5}}.
 \label{gff:done:kite}
\end{equation}
Put \(a=-p_1/2,b=-p_2/2,A'=-p_3/2,B'=-p_4/2,c=-p_5/2\), and
\begin{align}
 \mathcal O(a,b,A',B',c)&=
 B(a+1,c+1)B(a+A'+c+2,B'+1)\notag\\
 &\quad\times{}_3F_2\!\left(
 \begin{matrix}-b,a+1,a+A'+c+2\\a+c+2,a+A'+B'+c+3\end{matrix};1\right),
 \label{gff:done:ordered}\end{align}
\begin{align}
 \mathcal M(a,b,A',B',c)&=
 B(a+1,-a-b-c-1)B(B'+1,-A'-B'-c-1)\notag\\
 &\quad\times{}_3F_2\!\left(
 \begin{matrix}-c,a+1,B'+1\\-b-c,-A'-c\end{matrix};1\right).
 \label{gff:done:mixed}\end{align}
Here \(A',B'\) are temporary exponent names, not the RG constant
\(A=4\log2\). A split-contour identity gives
\begin{align}\boxed{\begin{aligned}
 \mathcal J(\boldsymbol p)=\frac1{\pi^2}\bigl[&
 \sin(\pi a)\sin(\pi B')\,
 \mathcal O(a,b,A',B',c)\mathcal M(a,b,A',B',c)\\
 &+\sin(\pi A')\sin(\pi b)\,
 \mathcal O(A',B',a,b,c)\mathcal M(A',B',a,b,c)\bigr].
 \end{aligned}}\label{gff:done:kiteFactorization}
\end{align}
Thus the needed nonsymmetric kite is evaluated by four ordinary
\({}_3F_2(1)\) values, not by a noisy four-real-dimensional integration.
This is a bilinear identity. For complex regulator values the products in
\eqref{gff:done:kiteFactorization} are not replaced by products involving complex
conjugation.

For clarity, the contour argument starts in a common convergence region.
Split the two antiholomorphic contours so that one lies left of zero and
the other right of one. The two possible orderings of the remaining
holomorphic coordinates give \(\mathcal O\) and its exchanged version.
The discontinuities of the two outer contours give the sine factors in
\eqref{gff:done:kiteFactorization}. In the mixed domain the changes
\(x=-u/(1-u)\), \(y=1/(1-v)\) give \(\mathcal M\) by two Euler beta
integrations. The one-variable normalization is fixed by
\(-\sin(\pi b)B(a+1,b+1)B(-a-b-1,b+1)\); applying the construction twice
fixes the positive sign and the \(\pi^{-2}\) in \eqref{gff:done:kiteFactorization}.
Analytic continuation then proves the meromorphic identity away from its
polar divisors. This is the two-integration specialization of the usual
KLT contour factorization \cite{gff:C6KLT}; the relation between two-dimensional
massless Feynman integrals and bilinear hypergeometric expressions is also
established in \cite{gff:C6DuhrPorkert}.

The symmetric cases reduce independently to the complex Selberg Gamma
product \cite{gff:NeretinComplex}. In particular, setting all five powers to one
recovers \(\mathcal J=3Q_2\). Three nonsingular symmetric tests, five
nontrivial completion transformations at a nonsymmetric point, and a
changed-tail hypergeometric calculation verify the normalization separately
from the graph recursion. These are consistency checks of the analytic
formula, not a fit for its coefficients.

\paragraph{Counterterm closure and infrared conditions.}
A graph is accepted by the coordinate reduction only when its raw graph,
all required subgraphs, and every inherited-root quotient have a declared
reduction. A root pair is rejected if it is contained in a proper induced
collision set with \(e\geq2(v-1)\). This prevents a spurious infrared
singularity being interpreted as a UV counterterm. An unanchored
one-attachment massless block is assigned its analytic scaleless value zero
only when its superficial degree at \(d=2\) is the nonzero odd number
\(1\) or \(-1\). No logarithmic scaleless factor is assigned zero by this
rule. Its logarithmic subtractions are still performed before the graph
counterterm is extracted.

Cauchy extraction of the regular function \(z^2\mathcal J(z)\) uses two
settings: 32 nodes at radius 0.02 with 45-digit function evaluation, and
48 nodes at radius 0.016 with 55-digit evaluation. The two runs compare
2699 pole coefficients; their maximum absolute discrepancy is
\(7.38\times10^{-30}\). The maximum discrepancy between distinct admissible
roots in the higher-precision run is \(7.97\times10^{-45}\). The Gamma-only
counterterms of Section~\ref{gff:two:gammaSector} provide additional independent
checks. No zero is
inserted because a nonzero result is numerically inconvenient.

The closed graph/subgraph/quotient reductions and the two contractions below cover the nonzero power-sector target. The other logarithmic graph contributions are evaluated by the finite-scheme calculation, without requiring separate numerical values for every individual MS finite part.

\subsection{The complete power sector is explicitly accounted for}
The only two nonzero-target power graphs not closed by the preceding
procedure are \(\mathrm{L6G106}\) and \(\mathrm{L6G117}\). Both are a
sunset insertion on the four-loop primitive \(\mathrm{L4G027}\). Put
\(p=1-z/2\). Integrating the chain with powers \(p,3p,p\) gives
\begin{equation}
 \mathcal S(p)=
 \frac{\Gamma(1-p/2)^2\Gamma(1-3p/2)\Gamma(5p/2-2)}
 {\Gamma(p/2)^2\Gamma(3p/2)\Gamma(3-5p/2)},
 \qquad \mathcal S(1)=-4.
 \label{gff:done:sunset}
\end{equation}
Its effective edge power is \(5p-4\). The remaining graph is primitive,
so the deformation of this edge does not affect the required simple
pole. Each graph has \(w_A=2,w_B=0\), giving \(16P_4(2)\) to \(b_6\).
Together they contribute \(32P_4(2)\) to \(b_6\) and zero to \(c_6\).
This contraction was checked by exact graph isomorphism and not inferred
from an expected numerical coefficient.

Let a superscript \(p\) denote the sum of graphs containing at least one
proper power-divergent subgraph. Define the lower-order polynomial of the constant-coefficient MS assembly, that is, the terms of \eqref{gff:c6:RGassembly} that contain only $A$, $b_3$, $c_3$, $b_4$ and $c_4$, with $A$ at the present dimension (Section~\ref{gff:mgc:c6assembly} shows that \eqref{gff:c6:RGassembly} holds at every dimension with this $A$):
\begin{align}
 \mathcal R(b_3,c_3,b_4,c_4)={}&2016A^5-448A^3b_3+672A^3c_3
 -84A^2b_4+168A^2c_4\notag\\
 &+\frac{49}{3}Ab_3^2-42Ab_3c_3
 +\frac73b_3b_4-4b_3c_4-3b_4c_3.
 \label{gff:done:lowerR}
\end{align}
Subtracting the constant-coefficient MS assemblies for the full and
power-free graph sets gives the exact bookkeeping identity
\begin{align}\boxed{\begin{aligned}
 2187\,\Delta_p={}&(3c_6-b_6)^p-12A b_5^p+30A c_5^p\\
 &+\mathcal R(b_3,c_3,b_4,c_4)
 -\mathcal R(b_3-b_3^p,c_3-c_3^p,b_4-b_4^p,c_4-c_4^p),\\
 C_6^{\rm full}(2)={}&C_6^{\rm power\mbox{-}free}(2)+\Delta_p.
 \end{aligned}}\label{gff:done:powerCorrection}
\end{align}
Only the complete combination \((3c_6-b_6)^p\), rather than its two
individual summands, is needed. Any remaining individual sixth-order
power graph with \(3w_B=w_A\) has zero weight in this combination and is
excluded by an exact rational identity.

The actual MS sums include
\begin{align}
 b_3^p&=-4,&c_3^p&=-\frac43,\notag\\
 b_4^p&=16A+24,&c_4^p&=\frac83A+8,\notag\\
 b_5^p&=-1540.38932947914290838\ldots,&
 c_5^p&=-251.621917341000466476\ldots .
 \label{gff:done:powerSums}
\end{align}
The evaluated higher-order part of \eqref{gff:done:powerCorrection}, divided
by 2187, is \(1.535501191157190785360393154334\ldots\); the lower-order
difference is its negative at the computed precision. The two Cauchy
settings give \(\Delta_p=-8.02\times10^{-34}\) and
\(\Delta_p=6.91\times10^{-47}\), respectively. These residuals are
numerical consistency checks, not a proof of an all-orders
power-subgraph cancellation theorem. The calculation retains the explicitly
evaluated correction \eqref{gff:done:powerCorrection}; it does not omit the
power graphs by assumption. Their residual numerical uncertainty is
negligible on the scale of the reported $C_6(2)$ uncertainty.

\subsection{Power subgraphs: a separate three-dimensional audit}
It would not be justified to infer cancellation of power graphs from the
two-dimensional result. The correction is instead evaluated in the MS
prescription before being added to the finite-scheme power-free result.
The coordinate convolution with distance exponents $a,b$ is
\begin{equation}
 B_d(a,b)=\Gamma(d/2)
 \frac{\Gamma((d-a)/2)\Gamma((d-b)/2)\Gamma((a+b-d)/2)}
 {\Gamma(a/2)\Gamma(b/2)\Gamma((2d-a-b)/2)}.
 \label{gff:d3six:convolution}
\end{equation}
The dimension-three implementation uses affine edge powers
$m(3-z)/2$, not the two-dimensional powers with a numerical prefactor
changed. It applies the Taylor reductions of the kite master \eqref{gff:gs:masterJ} at its critical patterns A, B and T only to their declared orders. Gamma factors, logarithmic
counterterms, contractions and safe nonzero-distance roots are retained
recursively. A requested coefficient beyond the available master Taylor
order is rejected, not assigned zero. Nonlogarithmic one-attachment
massless factors are removed only in the specified analytic critical-mass
prescription; this test is not used to discard a logarithmic scaleless
UV counterterm.

There are 1,5,31 relevant power-containing classes at loops 3,4,5. At six
loops all 121 power classes with $3w_B-w_A\ne0$ are included; the other
52 power classes have exactly zero weight in the sixth-order target.
A zero sixth-order target weight does not license dropping a lower-order
input. In particular the five-loop graph described next remains necessary.

\paragraph{Two sunset contractions.}
The graphs $\mathrm{L6G106}$ and $\mathrm{L6G117}$ are sunset insertions on
the four-loop primitive. The chain of three effective edges with
critical powers $p,3p,p$, $p=d/2$, has factor
\begin{equation}
 M_d=\frac{\Gamma(-q)\Gamma(2q)^2}{\Gamma(3q)},\qquad
 M_3=-\frac{16\pi}{15}.
 \label{gff:d3six:sunset}
\end{equation}
Each graph contributes $-4M_3P_4(3)$ to $b_6$ and zero to $c_6$.
Only the primitive residue is needed after the contraction; no unknown
edge-deformation finite term is inserted. The inherited graph
identification and its weights are checked in the audit.

\subsection{An exact cancellation avoids two unknown core integrals}
The nontrivial extra obstruction in dimension three is a four-loop
two-terminal core. Using the terminal labels $2,3$, its seven internal
edges are
\begin{equation}
 (01),(02),(03),(03),(12),(12),(13).
 \label{gff:d3six:coreEdges}
\end{equation}
Every edge has distance power $p=(3-z)/2$. Equivalently its two-integrated-
vertex kite powers are $(p,2p,2p,p,p)$. The required expansion has the form
\begin{equation}
 J_{\rm core}(z)=\frac8{z^2}+\frac{2A}{z}+u_0+u_1z+O(z^2),
 \qquad A=2+4\log2-\pi.
 \label{gff:d3six:coreJet}
\end{equation}
The pole coefficients follow from its proper logarithmic subgraphs.
The finite coefficients $u_0,u_1$ are not computed. Their cancellation from the physical target is
proved below, rather than assumed.

Let $C(z)$ be the product of the two external chain convolutions and
$F(z)=C(z)J_{\rm core}(z)$. For reference,
\begin{equation}
 C(0)=M_3,\qquad C'(0)=\frac74M_3(A+N),\qquad
 N=\psi_0(-q)+\psi_0(3q)-\psi_0(1)-\psi_0(2q)
 =\frac{62}{15}-\pi-4\log2.
 \label{gff:d3six:chainJet}
\end{equation}
Denote by $F_0,F_1$ the unknown-dependent regular coefficients. Known
pole-dependent terms are retained elsewhere in the audit. Write $\doteq$
for equality only of the dependence on $F_0,F_1$. If $r_G=[z^{-1}]R'I_G$,
then the three affected graphs satisfy
\begin{equation}\begin{aligned}
 r_{\mathrm{L5G086}}&\doteq\frac25F_0,\\
 r_{\mathrm{L6G332}}&\doteq\frac A3F_0+\frac2{15}F_1,\\
 r_{\mathrm{L6G333}}&\doteq\frac A3F_0-\frac2{15}F_1.
\end{aligned}\label{gff:d3six:coreResidues}\end{equation}
These expressions follow by applying
$B_d(d-kz,d/2-lz)=2/(kz)+A+O(z)$ to the attached convolutions.
For $\mathrm{L6G333}$ the external bubble counterterm contributes
$-4F_1/5$; omitting it would spoil the identity.

The five-loop weights are $(w_A,w_B)=(3/4,1/4)$ and both six-loop
weights are $(3/2,1/4)$. Since $b_L=(-1)^{L+1}Lw_Ar_G$ and
$c_L=(-1)^{L+1}Lw_Br_G$,
the unknown part of the physical sixth-order MS numerator is precisely
\begin{equation}\boxed{
 \frac92\bigl(r_{\mathrm{L6G332}}+r_{\mathrm{L6G333}}\bigr)
 -\frac{15}{2}A\,r_{\mathrm{L5G086}}\doteq0.
}\label{gff:d3six:unknownCancellation}\end{equation}
This is an exact rational-polynomial identity in $A,F_0,F_1$. It
eliminates only the two unknown core coefficients, not the other parts
of these graphs. Substituting $(u_0,u_1)=(0,0)$ and $(101,-37)$ in the
full graph recursion changes the three individual residues exactly as
predicted by \eqref{gff:d3six:coreResidues}, to better than $4\times10^{-73}$.
Thus those trial assignments are algebraic checks, not estimates of
$u_0,u_1$. The saved auxiliary individual $b_5,b_6$ sums obtained with a
trial assignment must not be reported as separately evaluated physical
MS coefficients.

\subsection{Assembled power correction and completeness}
The constant-MS lower-order polynomial is the same function $\mathcal R$ as in \eqref{gff:done:lowerR}; only $A$ changes, to its $d=3$ value $2+4\log2-\pi$ of \eqref{gff:d3six:coreJet}. It is repeated here for reference:
\begin{align}
 \mathcal R(b_3,c_3,b_4,c_4)={}&2016A^5-448A^3b_3+672A^3c_3
 -84A^2b_4+168A^2c_4\notag\\
 &+\frac{49}{3}Ab_3^2-42Ab_3c_3
 +\frac73b_3b_4-4b_3c_4-3b_4c_3.
 \label{gff:d3six:Rlower}
\end{align}
A superscript $p$ denotes the complete power-containing graph sum.
Subtracting the full and power-free MS assemblies gives
\begin{equation}\boxed{\begin{aligned}
 2187\Delta_p={}&(3c_6-b_6)^p-12Ab_5^p+30Ac_5^p\\
 &+\mathcal R(b_3,c_3,b_4,c_4)
 -\mathcal R(b_3-b_3^p,c_3-c_3^p,b_4-b_4^p,c_4-c_4^p),\\
 C_6^{\rm full}(3)={}&C_6^{\rm pf}(3)+\Delta_p.
\end{aligned}}\label{gff:d3six:powerCorrection}\end{equation}
For the three- and four-loop power pieces the graph and counterterm recursion of this section agrees with
\begin{equation}\begin{aligned}
 b_3^p&=M,&c_3^p&=M/3,\\
 b_4^p&=-\tfrac{11}{2}AM-\tfrac32MN,&
 c_4^p&=-\tfrac76AM-\tfrac12MN,
\end{aligned}\label{gff:d3six:lowerPower}\end{equation}
where $M=M_3$. Setting $R=\psi_1(2q)-\psi_1(1)$,
$U=\psi_2(q)$, $V=\psi_2(1)+\psi_2(2q)$, $Q=Q_3$, and $P=P_4(3)$,
exact assembly algebra gives the lower-order part of
\eqref{gff:d3six:powerCorrection},
$\delta\mathcal R=\mathcal R(b_3,c_3,b_4,c_4)-\mathcal R(b_3-b_3^p,c_3-c_3^p,b_4-b_4^p,c_4-c_4^p)$
with the complete lower-order coefficients \eqref{gff:gd:bc3}, \eqref{gff:gs:b4} and
\eqref{gff:gs:c4}, in closed form; it is $2187\,\Delta_{p,\rm lower}$ below:
\begin{equation}\begin{aligned}
 \delta\mathcal R=\frac M{24}\bigl(&2727A^3+603A^2N-8AM+72AQ+63AR\\
 &-216NQ+27NR-192P-318U+36V\bigr).
\end{aligned}\label{gff:d3six:deltaR}\end{equation}
This polynomial alone does not prove the higher graph sum. The latter
was separately computed by the graph and counterterm recursion.

At 90 working digits the two parts of \eqref{gff:d3six:powerCorrection},
including its factor $1/2187$, are
\begin{align}
 \Delta_{p,\rm higher}&=+0.0736516970441648694692229455611\ldots,\notag\\
 \Delta_{p,\rm lower}&=-0.0736516970441648694692229455611\ldots.
 \label{gff:d3six:powerNumbers}
\end{align}
Their sum is $-1.02\times10^{-89}$. A separate 75-digit run with different
core assignments gives $-1.94\times10^{-75}$. A third algebraic stress
test replaces the values of the masters $Q_3$, $P_4(3)$, $\KK_3$ and $Z_3$
by the test numbers $7$, $11$, $5$ and $13$, and the core assignments by
$(-17,29)$: the combined
residual is $1.52\times10^{-65}$ at 65 working digits. This last test is
not a physical evaluation of those masters; it checks that common
master-dependent terms cancel consistently in the target.

These are high-precision numerical checks of the \emph{complete}
power correction, together with an exact proof of the unknown-core
cancellation. They are not an all-orders theorem allowing individual
power graphs to be dropped. The tiny computed residual is negligible
relative to the sampling error, but it is not an interval-certified error
bound. Every nonzero-weight sixth-order power graph and every required
fifth-order power input has been covered. Two unrelated non-power
four-loop graphs rejected by the coordinate reduction are instead
included in the finite-kernel integration and are not a missing power
contribution.

\section{Finite-scheme sixth-order assembly}\label{gff:cn:finiteAssembly}

\subsection{A finite zero-momentum scheme avoids separate MS finite parts}
In this subsection only, \(\widehat b_L(z),\widehat c_L(z)\) denote the
\emph{epsilon-dependent finite-scheme} RG coefficients of the power-free
graph set. They are not the constant MS coefficients. Since the coupling and
the mass are defined by bare amplitudes at zero external momentum, this scheme
is called finite-MOM (finite momentum subtraction) below. Let
\begin{equation}
 \mathcal A(h,z)=1+\sum_{L=1}^6(-1)^L\mathcal A_L(z)h^L,
 \qquad
 \mathcal B(h,z)=1+\sum_{L=1}^6(-1)^L\mathcal B_L(z)h^L,
 \label{gff:done:bareSeries}
\end{equation}
where the graph weights \eqref{gff:c6:weights} and the bubble normalization are
those of Sections~\ref{gff:cn:finiteKernels} and~\ref{gff:app:C6d2Complete}. Defining \(g=h\mathcal A(h,z)\), differentiation at
fixed bare parameters gives
\begin{equation}
 \beta(g,z)=-z\,h\,\partial_h\bigl(h\mathcal A(h,z)\bigr),\qquad
 \gamma_r(g,z)=-z\,h\,\partial_h\log\mathcal B(h,z).
 \label{gff:done:finiteRG}
\end{equation}
The right sides are reexpanded in \(g\). They have the form
\begin{equation}
 \beta=-zg+\sum_{L=1}^6\widehat b_L(z)g^{L+1},\qquad
 \gamma_r=\sum_{L=1}^6\widehat c_L(z)g^L,
 \qquad \widehat b_1=3,\quad\widehat c_1=1.
 \label{gff:done:finiteCoefficients}
\end{equation}
The fixed-point exponent is unchanged by this finite coupling and
operator redefinition: the coupling is reparametrized, and the extra
operator-anomalous-dimension term is proportional to \(\beta\), hence
vanishes at the fixed point. This assertion is used only for the same
power-free amplitudes in the two prescriptions. The correction for the
omitted power-containing amplitudes is separately supplied by
\eqref{gff:done:powerCorrection}.

For a compatible forest \(F\) of proper logarithmic collision sets,
let its blocks be the contracted subgraphs and the remaining root block.
Write \(h_R\) for the loop number of that root block and
\(I_B\) for a bare block amplitude. The finite scale-derivative combination is
\begin{equation}
 D_G(z)=\sum_F(-1)^{|F|}\frac{h_R}{L}\prod_{B\in F\cup\{G\}}I_B(z).
 \label{gff:done:forestD}
\end{equation}
The notation on the product means the forest blocks after their children
have been contracted, not a repeated product of unreduced parent graphs.
Terms with \(h_R=0\) vanish. Complete rational polynomial comparison with
independent inversion of \eqref{gff:done:bareSeries} proves
\begin{equation}
 \widehat b_L(z)=(-1)^{L+1}Lz\sum_G w_A(G)D_G(z),\qquad
 \widehat c_L(z)=(-1)^{L+1}Lz\sum_G w_B(G)D_G(z).
 \label{gff:done:forestRG}
\end{equation}
This comparison was rerun at all six loop orders. The sixth-order
polynomials contain 532 beta and 383 mass monomials; both differences
vanish over the rationals. It is not a numerical agreement at selected
values of the bare integrals. The use of common parametric forest kernels
is consistent with the general parametric renormalization framework
\cite{gff:BrownKreimerLocal}.

\paragraph{The independent inversion.}\label{gff:mgc:forestReversion}
The comparison uses the Lagrange--B\"urmann inversion of $g=h\mathcal A(h,z)$ in \eqref{gff:done:finiteRG}.
With primes denoting $\partial_h$ at fixed $z$ and $\mathcal A,\mathcal B$ the series
\eqref{gff:done:bareSeries},
\begin{equation}\begin{aligned}
 \widehat b_L(z)&=-\frac{z}{L+1}\,[h^L]\bigl(2h\mathcal A'+h^2\mathcal A''\bigr)\mathcal A^{-L-1},\\
 \widehat c_L(z)&=-\frac{z}{L}\,[h^{L-1}]\,\frac{d}{dh}\!\left(\frac{h\mathcal B'}{\mathcal B}\right)\mathcal A^{-L}.
\end{aligned}\label{gff:mgc:reversion}\end{equation}
Indeed $[g^n]F(h(g))=\frac1n[h^{n-1}]F'(h)\mathcal A^{-n}$; for $F=-zh(h\mathcal A)'$ one has
$F'=-z(\mathcal A+3h\mathcal A'+h^2\mathcal A'')$, and $[h^L](\mathcal A+h\mathcal A')\mathcal A^{-L-1}=0$
for $L\geq1$. The same comparison of every monomial was made against the forest polynomials
\eqref{gff:done:forestRG}. For the six-loop inputs the 487 power-free graphs, with 66,198 forests, combine
into 395 nonzero rational kernels; 92 graph kernels vanish identically, and the largest retained kernel
has 52 distinct denominators. These scale-derivative coefficients carry explicit regulator dependence;
they are not the constant MS coefficients and are never inserted into an MS fixed-point assembly
\cite{gff:BaldufPanzerTropicalLocal}. The inversion formulas \eqref{gff:mgc:reversion}, the comparison and the
counts are those of the earlier two-dimensional calculation by OpenAI GPT-6 astra (package
\path{LRI_C6_D2_HPC/reference/}); the derivation of \eqref{gff:mgc:reversion} from $[g^n]F(h(g))$ was written
out by a Claude Code session on 16 October 2026, when this paragraph was merged into the note.


\subsection{A common, locally subtracted integration kernel}
For a forest, let \(U_F\) be the product of the Symanzik polynomials of its
blocks, embedded in the original edge variables, and let \(E_R(F)\) be the
root-block edge set. With \(S=\sum_e x_e\), define
\begin{equation}\boxed{
 K_G(x)=\sum_F(-1)^{|F|}
 \frac{\sum_{e\in E_R(F)}x_e}{S\,U_F(x)}.
}\label{gff:done:commonKernel}
\end{equation}
It is homogeneous of degree \(-L\). All terms are combined before
numerical evaluation. At \(d=2\), the normalization
\(r(2,z)=\Gamma(1+z/2)/\Gamma(1+z/2)\) is identically one, and
\begin{equation}
 D_G(z)=\frac{\Gamma(Lz/2)}{\Gamma(1/2+z/4)^{2L}}
 \int d\Omega\,\prod_e x_e^{-1/2+z/4}K_G(x).
 \label{gff:done:commonIntegral}
\end{equation}
To verify the factor in \eqref{gff:done:forestD}, integrate first over the
relative scales of the blocks. Their Dirichlet parameters are
\(h_Bz/2\), so insertion of the root scale fraction has mean
\((h_Rz/2)/(Lz/2)=h_R/L\). This proves the exact common-integral identity
for positive regulator in its initial convergence region, and then by
continuation. No fitted subtraction constant is involved.

The forest cancellations remove the proper logarithmic endpoint
singularities. The remaining overall pole is isolated by
\(\Gamma(Lz/2)\). Consequently the integral in
\eqref{gff:done:commonIntegral} and its required logarithmic moments are
finite at \(z=0\). The individual terms of \eqref{gff:done:commonKernel} need
not be separately integrable there and are never sampled as unrelated
singular integrals.

There are 1, 5, 16, 76 and 395 nonzero finite kernels at loop orders
2 through 6. For all 493 kernels, an independent exact-rational check
compares the co-tree-mask implementation with a product of Kirchhoff
cofactor determinants constructed directly from the forest blocks. Every
difference is zero. This validates the common-integral mapping separately
from the RG polynomial check.

\subsection{Retaining the regulator dependence through sixth order}
Write
\(\widehat b_L(z)=\sum_j b_{L,j}z^j\) and
\(\widehat c_L(z)=\sum_j c_{L,j}z^j\). Terms through \(j=6-L\) are
required. In a projective sector with one largest edge fixed at one,
put
\begin{equation}
 s_L(x)=\frac14\sum_e\log x_e-\frac L2\log S,
 \qquad
 \mathcal G_L(z)=
 \frac{\pi^L\Gamma(1+Lz/2)}{\Gamma(1/2+z/4)^{2L}}.
 \label{gff:done:momentWeight}
\end{equation}
The normalized Gamma factor has \(\mathcal G_L(0)=1\) and
\begin{equation}
 \log\mathcal G_L(z)=\sum_{j\geq1}a_{L,j}z^j,\quad
 a_{L,1}=L\log2,\quad
 a_{L,j}=\frac{(-1)^j\zeta(j)}{j}
 \left[\left(\frac L2\right)^j-
 \frac{2L(2^j-1)}{4^j}\right]\quad(j\geq2).
 \label{gff:done:GammaJets}
\end{equation}
Only \(j\leq4\) is used. If \(t=s_L+a_{L,1}\), the required polynomials
are
\begin{align}
 E_0&=1,&E_1&=t,&E_2&=\frac{t^2}{2}+a_{L,2},\notag\\
 E_3&=\frac{t^3}{6}+a_{L,2}t+a_{L,3},\qquad
 E_4&=\frac{t^4}{24}+\frac{a_{L,2}t^2}{2}+a_{L,3}t
       +a_{L,4}+\frac{a_{L,2}^2}{2}.
 \label{gff:done:momentPolynomials}
\end{align}
Integrating \(E_jK_G\) with the critical projective measure and multiplying
by \(2(-1)^{L+1}w_{A,B}/\pi^L\) gives \(b_{L,j},c_{L,j}\).
This computes Taylor moments of a finite integral, not a numerical
subtraction of Laurent poles.

The two-loop jets are evaluated by the convergent two-\({}_3F_2(1/4)\)
formula \eqref{gff:gd:twoLoopHyper} for \(S_d(z)\) at \(d=2\), rather than by Monte
Carlo. Explicitly,
\begin{equation}
 \widehat b_2(z)=-6z\left[S_2(z)-\frac2{z^2}\right],\qquad
 \widehat c_2(z)=\frac16\widehat b_2(z).
 \label{gff:done:twoLoopMOM}
\end{equation}
Fresh 65-digit evaluations with radii 0.05 and 0.08 and respectively
64 and 80 Cauchy nodes agree with the saved input through \(z^4\) to
better than \(6.8\times10^{-61}\). The five beta coefficients begin
\begin{equation}
 \begin{aligned}
 b_{2,0}&=-16.63553233344,& b_{2,1}&=29.02087369541,\\
 b_{2,2}&=-47.11795797162,& b_{2,3}&=74.14063002868,\\
 b_{2,4}&=-114.64528966290.&&
 \end{aligned}
 \label{gff:done:twoLoopNumbers}
\end{equation}

Let \(\widehat B(g,z)=\sum_{L\geq1}\widehat b_L(z)g^{L-1}\).
The physical fixed point satisfies \(\epsilon=g_*\widehat B(g_*,\epsilon)\).
Lagrange inversion with \(z\) held fixed, followed by \(z=\epsilon\), gives
\begin{equation}\boxed{
 C_n^{\rm power\mbox{-}free}(2)=
 -\frac12\delta_{n1}+
 \sum_{k=1}^n\frac1k
 [g^{k-1}z^{n-k}]
 \partial_g\gamma_r(g,z)\widehat B(g,z)^{-k},\qquad n\geq1.
}\label{gff:done:bivariateLagrange}
\end{equation}
The constant term is \(C_0=1\). Only finite polynomials enter this
extraction. Independent exact-rational fixed-point substitution and
bivariate Lagrange extraction agree in 192 checks, including 24 checks of
the complete constant-coefficient MS sixth-order assembly.

\textbf{The regulator jets cannot be discarded.} Substituting only
\(\widehat b_L(0),\widehat c_L(0)\) into the constant-coefficient MS assembly
(the analogue of \eqref{gff:c6:RGassembly}) would use a
different truncated RG problem. All required jets in
\eqref{gff:done:bivariateLagrange} are retained, including the two-loop
\(z^4\), three-loop \(z^3\), four-loop \(z^2\), and fifth-loop \(z^1\)
terms. At sixth order only the combination
\(3c_{6,0}-b_{6,0}\) is required, as in the MS assembly.

\subsection{Importance sampling and retained sector controls}
\label{gff:done:samplingControls}
The integration uses Hepp sectors and subset dynamic programming, in the
spirit of tropical Feynman quadrature \cite{gff:BorinskyD2Focus}. For each
proper nonempty edge subset \(S\), choose a positive proposal exponent
\(\alpha(S)\). Numerical leading-power diagnostics, performed on the
\emph{combined} kernel, guide this choice; they are not presented as
rigorous valuations. The normalization is computed without enumerating
all \((2L)!\) edge orders:
\begin{equation}
 Z(\varnothing)=1,\qquad
 Z(S)=\frac1{\alpha(S)}\sum_{e\in S}Z(S\setminus\{e\}),
 \qquad
 Z(E)=\sum_{e\in E}Z(E\setminus\{e\}).
 \label{gff:done:proposalDP}
\end{equation}
An edge is removed with probability proportional to
\(Z(S\setminus\{e\})\); the next sector ratio is sampled with density
\(\alpha(S)t^{\alpha(S)-1}\). The exact reciprocal proposal weight is
applied to every measurement. Changing the positive exponents changes
the variance, not the target integral. Diagram probabilities are fixed
from preceding pilot data with a nonzero uniform floor, and likewise
corrected by their reciprocal. A stated sample count is the number of
random graph--sector--point draws for an entire loop order, not that
number of samples for each individual graph.

Co-tree monomials are evaluated relative to their largest logarithm.
A cancellation monitor compares the signed forest value with the sum of
absolute term contributions. Evaluation starts in extended precision and
is repeated in binary128 when that ratio is below \(10^{-9}\); a
100-digit fallback is available below \(10^{-23}\) at the second stage.
No sample is clipped. The primary run, run C below, invokes binary128 21703 times
and never invokes the final fallback. These diagnostics are precision
controls, not interval arithmetic.

The 100-digit level of that fallback was later found too low. The following account comes from the
three-dimensional sixth-order bundle \path{lri_C6_d3_complete_update/}, written by OpenAI GPT-6 astra (model
stated by the author): a deliberately extreme stress test of the graph $\mathrm{L3G007}$ in that package exposed
insufficient relative accuracy, and the fallback was raised to 300 digits before the test was repeated. This
is the level used by the later sixth-order runs, including the 2026 productions
(Sections~\ref{gff:prec:d2production2026} and~\ref{gff:prec:d3production2026}). No production run of the
three-dimensional package called the emergency path, so none used the defective branch, and two same-seed
128,000-sample replays at loops three and six agreed batch for batch before and after the change.
None of the runs A--D described next reached that branch: the counters of the final fallback in their
records (\path{lri_C6_d2_complete_update/reproduction/mom_prod_*.json} of the research archive) are zero
at every loop order, so their values are unaffected.


Runs A and B, which are unrelated to the two runs of \eqref{gff:two:inputNumerics},
integrate all 395 nonzero sixth-order finite kernels. In
run C the 10 primitive, 26 one-subgraph and 40 two-subgraph kernels are
replaced by their already measured finite-scheme values, with their
uncertainties propagated. The correct replacement is
\begin{equation}\begin{aligned}
 K_{76}^{\rm MOM}
 &=C_{6,\rm primitive}-\frac2{729}(Y_{26}+Y_{40})\\
 &=528.2946383532544,\qquad
 \widehat\sigma(K_{76}^{\rm MOM})=0.1373774929551.
 \end{aligned}\label{gff:done:knownMOM76}
\end{equation}
It is \emph{not} the sum of the separately MS-converted
\(\Theta_{26},\Theta_{40}\). This distinction avoids double counting
finite-scheme conversion terms. Exact zero weights remove another 28
kernels from the sixth-order physical combination, so the remaining
sixth-order integration in C has 291 kernels. Its proposal exponents are
0.85 times those of A and B. A fourth run, D, used the same substitution as C
with the proposal exponents of A and B, a different seed and 16,000,000 draws per
loop order (A and B used 8,000,000 and C 32,000,000). The measurements behind
\eqref{gff:done:knownMOM76} were made before, and independently of, the loop-order
samples and are not used to calibrate an unknown \(C_6\). Runs A--D, combined by
inverse covariance with the weights 0.041, 0.039, 0.702 and 0.219, gave the first
complete two-dimensional estimate, superseded by the 2026
production (Section~\ref{gff:prec:d2production2026}); their records are in the
research archive \texttt{LRI\_complete\_research\_bundle.zip}, which is not part
of this repository, and the input \eqref{gff:done:knownMOM76} is also in
\path{LRI_C6_D2_HPC/reference/baseline_previous/baseline/lri_C6_d2_complete_update/reproduction/mom_known_sectors.json}.

\subsection{The changes required by dimension three}
The long-range regulator is varied at fixed dimension, with
$q=d/4$ and edge exponent $a=q+z/4$. The general-dimensional
renormalization setting follows \cite{gff:BGHS}; the sixth-order numerical
work and its evidence are described here and are not part of that
three-loop paper.

For a proper logarithmic forest $F$ of a power-free $L$-loop graph $G$,
let $U_F$ be the product of the first Symanzik polynomials of its contracted
blocks, $E_R(F)$ the edges in its root block, and $S=\sum_e x_e$.
The correct common kernel in dimension three is
\begin{equation}\boxed{
 K_G^{(3)}(x)=\sum_F(-1)^{|F|}
 \frac{\sum_{e\in E_R(F)}x_e}{S\,U_F(x)^{3/2}}.
}\label{gff:d3six:kernel}\end{equation}
The fractional power applies to each complete forest denominator; it
cannot be obtained by raising an already summed two-dimensional kernel
to $3/2$. Each term is evaluated on the same point before the sum is
sampled. This preserves local cancellation of logarithmic endpoint
singularities. General parametric forest renormalization provides the
framework \cite{gff:BrownKreimerLocal,gff:Herzog}; the retained exact forest-to-RG
polynomial identity remains dimension independent.

With $r(d,z)=\Gamma(d/2+z/2)/\Gamma(1+z/2)$ and $d=3$, the corresponding
scale-derivative amplitude is
\begin{equation}
 D_G(3,z)=\frac{r(3,z)^L\Gamma(Lz/2)}{\Gamma(3/4+z/4)^{2L}}
 \int d\Omega\,\prod_e x_e^{-1/4+z/4}\,K_G^{(3)}(x).
 \label{gff:d3six:integral}
\end{equation}
Relative block-scale integration has Dirichlet parameters $h_Bz/2$.
The inserted root fraction therefore has mean $h_R/L$, proving the same
forest coefficient as in the two-dimensional construction. The finite
RG coefficients, distinguished by hats from constant MS coefficients, are
\begin{equation}
 \widehat b_L(z)=(-1)^{L+1}Lz\sum_G w_A(G)D_G(3,z),\qquad
 \widehat c_L(z)=(-1)^{L+1}Lz\sum_G w_B(G)D_G(3,z).
 \label{gff:d3six:forestRG}
\end{equation}
The inherited inventory contains 660 sixth-order graph classes. There are
487 power-free classes and 173 classes with a proper power subgraph.
After exact forest cancellation, the power-free numerical inputs at
$L=2,3,4,5,6$ contain respectively $1,5,16,76,395$ nonzero kernels.
Their edge lists and rational weights were checked again against the
inventory. This is reuse and checking of the project's graph inventory,
not a second independent enumeration.

\paragraph{Dimension-correct analytic prefactors.}
Unlike at $d=2$, $r(3,z)$ is not identically one. In a largest-edge
projective sector define
\begin{align}
 s_L(x)&=\frac14\sum_e\log x_e-\frac L2\log S,\notag\\
 G_{L,d}(z)&=
 \frac{r(d,z)^L\Gamma(1+Lz/2)}{\Gamma(q+z/4)^{2L}}
 \frac{\Gamma(q)^{2L}}{\Gamma(2q)^L},\qquad q=\frac34.
 \label{gff:d3six:G}
\end{align}
Then $G_{L,d}(0)=1$ and $\log G_{L,d}(z)=\sum_{j\geq1}a_{L,j}z^j$, where
\begin{equation}\boxed{
 a_{L,j}=\frac{
 L2^{-j}\psi_{j-1}(2q)
 +\big[(L/2)^j-L2^{-j}\big]\psi_{j-1}(1)
 -2L4^{-j}\psi_{j-1}(q)}{j!}.
}\label{gff:d3six:GammaJets}\end{equation}
The overall critical normalization is
\begin{equation}
 n_L=2(-1)^{L+1}\frac{\Gamma(2q)^L}{\Gamma(q)^{2L}}.
 \label{gff:d3six:norm}
\end{equation}
Integrating the coefficients of $\exp(s_Lz)G_{L,d}(z)$ against the
critical projective kernel and multiplying by $n_Lw_{A,B}$ gives the
required jets. These coefficients are the finite Bell polynomials
\eqref{gff:done:momentPolynomials}. Their maximum order is four; no
pole-fitting operation is used in the Monte Carlo integration.
All 20 Gamma logarithmic jets and five normalizations were freshly
checked by independent 65-digit Cauchy evaluation of the analytic Gamma
product. The maximum absolute discrepancy is $2.27\times10^{-59}$.

\subsection{Regulator dependence and finite-order extraction}
The two-loop input is recomputed from the explicit general-dimensional
two-${}_3F_2(1/4)$ formula \eqref{gff:gd:twoLoopHyper} for $S_d(z)$, using
\begin{equation}
 \widehat b_2(z)=-6z\left(S_3(z)-\frac2{z^2}\right),\qquad
 \widehat c_2(z)=\frac16\widehat b_2(z).
 \label{gff:d3six:twoLoop}
\end{equation}
Its beta jets through $z^4$ are
\begin{equation}\begin{aligned}
 b_{2,0}&=-9.7859764118999280,& b_{2,1}&=13.1771369571188313,\\
 b_{2,2}&=-16.3759608428540408,&b_{2,3}&=19.5740514096297382,\\
 b_{2,4}&=-22.7913170046010398.&&
\end{aligned}\label{gff:d3six:twoLoopValues}\end{equation}
Two independent 60-digit evaluations use radii $0.05,0.07$ and respectively
48,64 Cauchy nodes. Their full outputs, rather than these rounded entries,
are kept in the package records listed in Section~\ref{gff:cn:reproduction}.

Let $\widehat B(g,z)=\sum_{L\geq1}\widehat b_L(z)g^{L-1}$, with
$\widehat b_1=3,\widehat c_1=1$, and
$\gamma_r(g,z)=\sum_L\widehat c_L(z)g^L$. The power-free physical
fixed point is defined by $\epsilon=g_*\widehat B(g_*,\epsilon)$.
Keeping $z$ independent during Lagrange inversion gives
\begin{equation}\boxed{
 C_n^{\rm pf}(3)=-\frac12\delta_{n1}+
 \sum_{k=1}^{n}\frac1k[g^{k-1}z^{n-k}]
 \partial_g\gamma_r(g,z)\widehat B(g,z)^{-k},\qquad n\geq1,
}\label{gff:d3six:lagrange}\end{equation}
with $C_0=3/2$. The extraction is from finite polynomials, not from an
unevaluated master. The needed orders are $z^{6-L}$ at each loop $L$.
Thus the fifth-loop mass and coupling finite-scheme jets are included
separately, while at six loops only $3c_{6,0}-b_{6,0}$ is needed.
Replacing every hatted function by its constant term would be incorrect.
An independent rational-polynomial implementation of implicit fixed-point
substitution and bivariate Lagrange inversion gives 256 exact checks,
including 32 checks of the constant-MS sixth-order assembly.

\section{Final numerical design, covariance and precision}\label{gff:app:C6Precision}\label{gff:app:C6Methods}

\paragraph{Current and superseded estimates.}
This section serves both dimensions. Only the 2026 productions of
Sections~\ref{gff:prec:d2production2026} and~\ref{gff:prec:d3production2026} are
current. Each has two campaigns, $a$ and $b$, with 64 replicates at every loop
order, and they give $C_6(2)=362.366\pm0.004$ and $C_6(3)=33.5754\pm0.0005$. The
other campaigns described in this Part belong to superseded estimates; they are
kept because the productions inherit their kernels, allocations and checks.
The superseded sixth-order method-refinement calculation at $d=2,3$, called
the earlier sixth-order calculation in this section, with its refinement campaigns,
previous estimators, method benchmarks (Section~\ref{gff:mgs:benchmarks}), control
experiment (Section~\ref{gff:meth:crossdimension}) and method assessment
(Section~\ref{gff:mgs:selection}), was run by OpenAI GPT-6 astra (model stated by
the author) and written up by it in a methods note delivered with the package;
that note is merged into this section. Its code and records are in
\path{LRI_C6_D2_HPC/reference/}; record paths given below for the earlier
calculation are relative to that folder. Three parts were added by a Claude
Code session when the methods note was merged, on 16 October 2026: formula
\eqref{gff:mgs:linearControl}, rederived from the stored covariance, where it
reproduces the reported ratios $0.985$ and $0.962$; the descriptions of the
benchmark variants, of the pilot counts and of the trained maps in
Section~\ref{gff:mgs:benchmarks}, taken from the package code
(\path{code/experiment.cpp}, \path{code/pilot.py}, \path{code/train_grid.py});
and the proof of \eqref{gff:mgs:rao}.
\begin{itemize}
\item Runs A--D of Section~\ref{gff:done:samplingControls} gave the first
complete two-dimensional estimate.
\item The previous estimators of Section~\ref{gff:meth:pooling} pooled, at each
dimension, two baseline campaigns A and B with 64 replicates
(Sections~\ref{gff:prec:stratification} and~\ref{gff:prec:baseline}) and two
refinement campaigns, also labeled A and B, with 32 replicates
(Section~\ref{gff:meth:refinement}). They gave $C_6(2)=362.3489\pm0.1526$ and
$C_6(3)=33.5682\pm0.0242$ (sampling standard errors).
\end{itemize}
Upper-case letters label superseded campaigns; lower-case $a$ and $b$ label the
campaigns of the 2026 productions, here and in Sections~\ref{gff:app:C7d2},
\ref{gff:app:C7d3}, \ref{gff:app:C8d2} and~\ref{gff:app:C8d3}.

\subsection{Integrated RG Taylor vectors}
For a power-free $L$-loop graph $G$ with $E=2L$ edges, the retained
finite forest kernel is
\begin{equation}
 K_G^{(d)}(x)=\sum_F(-1)^{|F|}
 \frac{\displaystyle\sum_{e\in E_R(F)}x_e}
      {S(x)\,U_F(x)^{d/2}},\qquad S(x)=\sum_{e=1}^{E}x_e.
 \label{gff:prec:kernel}
\end{equation}
The sum is over the compatible forests of \eqref{gff:done:forestD}; $U_F$ is the
product of the first Symanzik polynomials of the forest blocks.  Terms
with identical denominators are combined before floating-point
integration.  The rational graph weights, normalized edge measure,
critical-mass prescription, and signs are those of
Section~\ref{gff:cn:finiteAssembly}.  The finite-kernel selection is fixed by exact forest cancellation.
There are respectively 5, 16, 76 and 395 nonzero finite kernels at loops
3, 4, 5 and 6.  This is a count after exact algebraic cancellation, not
the count of all six-loop graph classes.

The dimension-dependent prefactor, including its regulator dependence, is
\begin{equation}
 \frac{\Gamma(d/2)^L}{\Gamma(d/4)^{2L}}\,G_{L,d}(z)=\Gamma(1+Lz/2)
  \left[\frac{\Gamma(d/2+z/2)}{\Gamma(1+z/2)}\right]^L
  \Gamma(d/4+z/4)^{-2L}.
 \label{gff:prec:prefactor}
\end{equation}
Here $G_{L,d}$ is the normalized prefactor \eqref{gff:d3six:G} with $q=d/4$, so $G_{L,d}(0)=1$; at $d=2$ it is $\mathcal G_L$ of \eqref{gff:done:momentWeight}. The overall sign and factor are $n_L=2(-1)^{L+1}\Gamma(d/2)^L/\Gamma(d/4)^{2L}$, as in \eqref{gff:d3six:norm}, times the
retained coupling or mass graph weight.  Put
\begin{equation}
 a_{L,j}=[z^j]\log G_{L,d}(z),\qquad
 s_G(x)=\frac14\sum_e\log x_e-\frac L2\log S(x).
 \label{gff:prec:logjet}
\end{equation}
At each integration point the Taylor coefficients needed through
$z^{6-L}$ are generated by the finite recurrence
\begin{equation}
 e_0(x)=1,\qquad
 e_n(x)=\frac1n\sum_{j=1}^{n}j
       \bigl(a_{L,j}+\delta_{j1}s_G(x)\bigr)e_{n-j}(x).
 \label{gff:prec:jetrecursion}
\end{equation}
The scalar prefactors and their logarithmic Taylor coefficients were
recomputed at 70 decimal working digits and checked against the values stored with the superseded estimator.  The separately evaluated two-loop data remain inputs.  Four
additional independently seeded two-loop sampling checks use both ordering
maps in both dimensions and a different, constant-exponent importance
density.  All five required Taylor coefficients agree with the exact two-loop
master \eqref{gff:gd:twoLoopHyper} within 1.46 estimated standard errors.  Their total 8,388,608
samples are controls only and do not enter the final estimate.

For every loop the stored vector contains all required beta-function
Taylor coefficients followed by all required mass-anomalous-dimension
Taylor coefficients:
\begin{equation}
 \boldsymbol v_L=(b_{L,0},\ldots,b_{L,6-L},
                   c_{L,0},\ldots,c_{L,6-L}).
 \label{gff:prec:jetvector}
\end{equation}
Their covariance is not discarded.  Setting $z=\epsilon$ only after
forming these Taylor series, the coefficient map $\mathcal F_6$, which returns
the coefficient $C_6$ of $\Delta_{\mathcal E}$ at dimension $d$, is defined by
\begin{align}
 0&=-\epsilon g_*+3g_*^2+
    \sum_{L=2}^{6}\widehat b_L(\epsilon)g_*^{L+1}
       +O(\epsilon^8),\\
 \Delta_{\mathcal E}
   &=\frac{d-\epsilon}{2}+g_*+
    \sum_{L=2}^{6}\widehat c_L(\epsilon)g_*^L
       +O(\epsilon^7).
 \label{gff:prec:assembly}
\end{align}
Using only the values of the finite-scheme RG functions at $z=0$ would
not reproduce this calculation.  An independent finite two-variable
Lagrange expansion checks the fixed-point iteration used by $\mathcal F_6$.

\subsection{Graph stratification and fixed pilot allocation}
\label{gff:prec:stratification}
The estimator integrates every graph as a separate stratum, using a
fixed sample count $n_G$ within each independent randomized replicate.
For a generic component of the integral,
\begin{equation}
 \widehat I_b=\sum_G w_G\frac1{n_G}
       \sum_{i=0}^{n_G-1}f_G(X_{bGi})/p_G(X_{bGi}).
 \label{gff:prec:stratified}
\end{equation}
Here $p_G$ is the normalized, unchanged importance density for graph $G$;
there is no inverse random-graph selection probability.  This removes
random fluctuations in the number of times a graph is visited without
changing the integral being estimated.

The counts are powers of two, at least four.  A separate pilot and two
validation calculations estimate the graphwise variance of the linearized
$C_6$ contribution.  These estimates determine the allocation, with a
uniform floor; the allocation is frozen before either production
campaign.  No sample of a campaign is reused as a training pilot for that
campaign, and no
allocation is selected by agreement with a known physical coefficient.
This target-aware allocation is a variance-reduction design, not an
identity for $C_6$ and not a claim of mathematically optimal allocation.

Each graph stratum uses a linearly scrambled Sobol net with a separate
random digital shift.  Powers-of-two net sizes preserve the intended
balance; the first point is retained and the sequence is not thinned.
The direction matrices are generated with the documented SciPy Sobol
construction, then checked against its public point generator
\cite{gff:PrecisionSobol}.  Every baseline campaign has 64 independent randomized
replicates at every loop order, as does each campaign of the 2026
productions; the refinement campaigns of Section~\ref{gff:meth:refinement} have
32.  The same linear scrambling is shared
within a replicate while different graph shifts are independently
seeded; covariance is measured from the \emph{total} replicate vector.

Campaign A uses 52-bit digital coordinates and the original sequence of
conditional draws for the edge ordering.  Campaign B uses 60-bit digital
coordinates and the equivalent one-coordinate ordering construction below.
The integer net coordinate $k$ is mapped to $(k+1/2)/2^b$ before taking
logarithms.  This avoids evaluating an endpoint singularity at zero.
Finite digital resolution is itself an approximation; comparing 52 and
60 bits is an empirical control, not a proof of a bound on its bias.
The production arithmetic begins at long-double precision in both
campaigns; the optional guarded-double experimental build is not used.

\subsection{An exact compressed map for the edge ordering}
Let $S$ be the set of edges not yet selected and let $D(S)$ be the
positive dynamic-programming weight used by the importance sampler.
The probability of choosing edge $j\in S$ next is
\begin{equation}
 p(j\mid S)=\frac{D(S\setminus\{j\})}
                  {\sum_{k\in S}D(S\setminus\{k\})}.
 \label{gff:prec:conditional}
\end{equation}
One may use a fresh uniform coordinate for every such draw.  Alternatively,
start with a single $u\in(0,1)$, choose the interval $[P_j,P_j+p_j)$
containing $u$, and replace it by
\begin{equation}
 u\longmapsto\frac{u-P_j}{p_j}.
 \label{gff:prec:compression}
\end{equation}
Conditional on selecting $j$, the new variable is again uniform on
$(0,1)$.  Induction proves that an entire ordering $\sigma$ has exactly
its original probability $\prod p(j\mid S)$ under this map.  The gap
coordinates are sampled separately, so the continuous density and the
importance weight do not change.

For $E=2L$, the original map uses $2E-1$ digital coordinates after graph
stratification: $E$ ordering draws and $E-1$ scale draws.  The compressed
map uses $E$ coordinates: one ordering coordinate and $E-1$ scales.
At six loops the counts are thus 23 and 12.  This reduces the dimension
of the sampling map, \emph{not} the physical Schwinger integral.  It is
not claimed to improve every integrand, and the two campaigns test its
empirical behavior rather than assuming an improvement.  The probability
identity also passed exhaustive exact rational checks for 870 edge
permutations in smaller examples.

\subsection{Baseline campaigns of the previous estimators}
\label{gff:prec:baseline}
{\footnotesize
\setlength{\tabcolsep}{3.2pt}
\begin{center}\begin{tabular}{ccrrrrrr}
\toprule
$d$ & Run & $L=3$ & $L=4$ & $L=5$ & $L=6$ & $C_6$ & SE\\\midrule
3 & A & 4,456,448 & 9,502,720 & 32,768,000 & 67,076,096 & 33.590772 & 0.057948\\
3 & B & 4,456,448 & 9,502,720 & 32,768,000 & 67,076,096 & 33.560344 & 0.053265\\
\bottomrule
\end{tabular}\end{center}
}
The two three-dimensional baseline campaigns use 227,606,528 kernel samples
over all four loop orders; pilot and validation samples were additional and
excluded. They entered only the previous three-dimensional estimator
\eqref{gff:meth:combined}, and the two corresponding two-dimensional campaigns,
with 115,969,024 samples each, only the previous two-dimensional estimator.
Both are superseded by the productions of
Sections~\ref{gff:prec:d2production2026} and~\ref{gff:prec:d3production2026}.

Counts in this table are sums over graph strata, not counts per graph.
They are the actual rounded power-of-two allocations, not nominal
budgets.  Both campaigns use the same frozen graph allocation and
importance exponents; seeds, net precision and ordering maps differ.
The two campaigns are pooled with the fixed weights $1/2,1/2$ at the
level of the RG Taylor vectors.  The weights were not optimized after
observing the production central values or their empirical variances.

Separate campaigns have independent sampled points.

\subsection{Replicate covariance and fixed-point propagation}
For campaign $t$, loop $L$ and $R_{L,t}$ independent randomized replicates,
let $v_{L,t,r}$ be the complete vector of coupling, mass and regulator
coefficients. Define
\begin{equation}\begin{aligned}
 \mu_{L,t}&=\frac1{R_{L,t}}\sum_r v_{L,t,r},\\
 \widehat\Sigma_{L,t}
 &=\frac1{R_{L,t}(R_{L,t}-1)}
 \sum_r(v_{L,t,r}-\mu_{L,t})(v_{L,t,r}-\mu_{L,t})^{\mathsf T}.
\end{aligned}\label{gff:prec:jetcov}\end{equation}
This is the estimated covariance of the mean. The unit of statistical
independence is a randomized replicate, not a point within one Sobol net
\cite{gff:PrecisionQMC}. Independent loops have independent randomizations;
within each loop all beta/mass and Taylor-order covariances are retained.
The fixed weights, shared-input terms and final Jacobian propagation are
given in \eqref{gff:meth:combined}--\eqref{gff:meth:covariance} below.

No known physical coefficient is used to calibrate the estimator and no
outlying replicate is discarded; the quoted central value is the direct
pooled estimator. For the previous four-campaign estimators the nonlinear
fixed-point map was checked by an independent finite Lagrange expansion,
and a delete-one-replicate calculation diagnosed the leading finite-sample
bias. These archived checks concern the inherited map and estimator, not
the replicate data of the 2026 productions. A replicate bootstrap checks
variance propagation. In the previous two-dimensional estimator it drew
each shared retained input once per complete bootstrap sample; modeling
that input by a normal draw does not prove normality or give an interval
certificate. The 2026 productions have no shared retained input.

\subsection{The quantity being integrated}
For each loop order $L$, the sampled vector contains the coupling and
mass Taylor coefficients through regulator order $6-L$:
\begin{equation}
 \boldsymbol\mu_L=(b_{L,0},\ldots,b_{L,6-L},
 c_{L,0},\ldots,c_{L,6-L}),\qquad L=3,4,5,6.
 \label{gff:meth:jets}
\end{equation}
The one-loop input is exact, and the two-loop input is the special-function
evaluation of Section~\ref{gff:cn:finiteAssembly} at each dimension. The physical
coefficient is evaluated as a finite polynomial
\begin{equation}
 C_6=\mathcal F_6(\boldsymbol\mu_3,\boldsymbol\mu_4,
 \boldsymbol\mu_5,\boldsymbol\mu_6)+\eta.
 \label{gff:meth:functional}
\end{equation}
Here $\mathcal F_6$ solves the regulator-dependent finite-MOM fixed-point
problem before setting $z=\epsilon$; $\eta$ accounts only for separately
retained numerical kernels. Only the previous two-dimensional estimator used
such kernels, so $\eta=0$ in both current estimators. In particular,
\begin{equation}
 \frac{\partial C_6}{\partial(b_{6,0},c_{6,0})}
 =\frac{(-1,3)}{2187}.
 \label{gff:meth:lastgradient}
\end{equation}
One must not replace the Taylor vectors by their values at $z=0$ and
then apply the constant-coefficient MS assembly. The field-theoretic derivation of these
finite vectors and their power-sector compensation is unchanged.

The finite kernel retains the pointwise forest sum \eqref{gff:prec:kernel}: every
numerator and denominator is evaluated on the same integration point before
accumulating the graph estimate. Tropical importance
sampling uses the known endpoint geometry of these parametric
integrals~\cite{gff:BorinskyD2Focus}. More samples cannot in general repair
a representation which postpones subtraction of divergent terms until
after separate numerical integrations.

\subsection{The two-dimensional production of September 2026}\label{gff:prec:d2production2026}
At $d=2$ the current estimate comes from a single new production with two
independent campaigns, $a$ and $b$. It integrates the finite forest kernels
\eqref{gff:prec:kernel} of all 5, 16, 76 and 395 nonzero classes at loops
3--6, generates the regulator jets by \eqref{gff:prec:jetrecursion}, and
applies the fixed-point map \eqref{gff:prec:assembly} with the retained
analytic power-sector compensation of Section~\ref{gff:app:C6d2Complete}. The
code is an optimized, restartable reimplementation of the same
calculation. Its kernel keeps the addition order of every polynomial sum
and the binary128 and 300-decimal fallbacks at the unchanged thresholds
$10^{-9}$ and $10^{-23}$. In its legacy summation mode, replays of all 492
kernels against the original executable agree to a normalized difference
below $9\times10^{-20}$, with identical promotion counts; production uses
compensated summation within each task.

Both campaigns have 64 independent randomized replicates at every loop
order and use the conditional-draw ordering map with $2E-1$ digital
coordinates. Campaign $a$ uses 52-bit and campaign $b$ 60-bit
coordinates, and only campaign $a$ uses, at $L=6$, the pilot-frozen
separable map \eqref{gff:meth:transport}. Every stratum receives the frozen
per-replicate count of the refinement allocation \eqref{gff:meth:allocation},
raised to at least four, multiplied by $2^{13}$; the floor affects only the
28 six-loop kernels whose $C_6$ projection vanishes exactly. Per replicate
this gives 1,107,296,256, 2,164,260,864 and 9,487,515,648 points at
$L=3,4,5$, and 38,076,809,216 points for the 385 nonprimitive six-loop
kernels. The ten primitive six-loop kernels, which the previous estimator
retained as $P_{10}$, are integrated in every replicate with 5,645,533,184
points per replicate, divided among them by a variance-based (Neyman) allocation
frozen from archived paired records of the earlier calculation (\path{paired_new_a_d2_L6.json},
\path{paired_new_b_d2_L6.json}, \path{paired_prod_a_d2_L6.json} and \path{paired_prod_b_d2_L6.json} in
\path{LRI_C6_D2_HPC/reference/records/}, read by \path{tools/primitive_allocation.py}). The two campaigns contain
7,229,621,141,504 kernel samples. Group seeds are derived by SHA-256 from a
fresh base seed, the campaign, the loop and the role. No production sample
was used to choose the allocation, and pilot, smoke and timing samples are
excluded. Every required kernel thus has expectation weight one in every
replicate, no numerical input is retained, and
\begin{equation}\boxed{
 \widehat C_6(2)=
 \mathcal F_6\left(\frac{\boldsymbol\mu^{a}+\boldsymbol\mu^{b}}2\right),
 \qquad\widehat\sigma_{\rm common}=0,
}\label{gff:prec:d2combined}\end{equation}
with the covariance \eqref{gff:meth:covariance} at $w_a=w_b=1/2$.

The production ran on SDSC Expanse as 428 independent 64-thread array
elements, 12,535 core-hours in total. Sample identities do not depend on
node, thread count or checkpoint boundaries: two blocks of 1024 six-loop
tasks run at 64 and at 16 threads gave identical rows, including 152,717
binary128 promotions. The campaigns give
\begin{center}\begin{tabular}{@{}lrr@{}}\toprule
Campaign & $C_6(2)$ & Sampling SE\\\midrule
$a$ & 362.36672022 & 0.00097561\\
$b$ & 362.36622331 & 0.00126747\\
Combined & 362.36647176 & 0.00079973\\\bottomrule
\end{tabular}\end{center}
They differ by 0.31 combined standard errors. The combined standard error
is about half the value 0.0015 projected from the earlier campaigns under
Monte Carlo scaling, in which the variance falls as $1/n$ and the standard
error as $1/\sqrt n$. This single comparison does not determine a
convergence rate; the replicate covariance, not the projection, is
reported. The new estimator lies 0.12 previous standard errors from the
previous estimator $362.3489138\pm0.1526445$.

The reported allowance is four sampling standard errors, 0.0032, rounded
up to a multiple of $0.001$, which gives $C_6(2)=362.366\pm0.004$. The
four-standard-error rule was written down before the production result was
available; the rounding step was refined afterwards from the $0.01$ first
proposed to $0.001$, which changes only the rounding of the allowance. An independent two-dimensional calculation through seventh
order, in which the mixed local-mass and logarithmic forests are integrated
explicitly instead of through the analytic compensation, gives as a
by-product $C_6(2)=362.3727\pm0.0086$ (sampling standard error), $+0.73$
combined standard errors from \eqref{gff:prec:d2combined}. It shares the graph
and renormalization framework and therefore does not test that framework.
The previous two-dimensional baseline and refinement campaigns, the
retained groups \eqref{gff:meth:retained} and the estimator built from them are
superseded; their records are in \path{LRI_C6_D2_HPC/reference/}
(Section~\ref{gff:cn:reproduction}).

\subsection{The three-dimensional production of September 2026}\label{gff:prec:d3production2026}
At $d=3$ the current estimate likewise comes from a single new production
with two independent campaigns, $a$ and $b$. It integrates the finite forest
kernels \eqref{gff:prec:kernel} of all 5, 16, 76 and 395 nonzero classes at loops
3--6, the ten primitive six-loop kernels included, generates the regulator
jets by \eqref{gff:prec:jetrecursion}, and applies the fixed-point map
\eqref{gff:prec:assembly} with the retained analytic power-sector compensation of
Section~\ref{gff:app:C6d3Complete}. The code is the three-dimensional counterpart
of the one in Section~\ref{gff:prec:d2production2026}, with the binary128 and
300-decimal fallbacks at the same thresholds $10^{-9}$ and $10^{-23}$. In its
legacy summation mode it replays the archived three-dimensional refinement
streams, 385,744,896 samples on their archived nets and seeds, with fallback
counts equal to the archive in every stream, and reproduces the
refinement-only estimate $33.56084035$ bit for bit; production uses
compensated summation within each task. Section~\ref{gff:done:samplingControls}
explains why the last fallback has 300 decimal digits.

Both campaigns have 64 independent randomized replicates at every loop
order and use the conditional-draw ordering map with $2E-1$ digital
coordinates and the baseline importance map. Campaign $a$ uses 52-bit and
campaign $b$ 60-bit coordinates. Every stratum receives the frozen
per-replicate count of the three-dimensional refinement allocation
\eqref{gff:meth:allocation}, multiplied by $2^{13}$. Per replicate this gives
1,241,513,984, 2,499,805,184, 9,743,368,192 and 35,891,576,832 points at
$L=3,4,5,6$, and the two campaigns contain 6,320,161,816,576 kernel samples.
Group seeds are derived by SHA-256 from a fresh base seed, the campaign, the
loop and the role. No production sample was used to choose the allocation,
and the samples of the thread-identity check below are excluded. Every
required kernel has expectation weight one in every replicate, no numerical
input is retained, and
\begin{equation}\boxed{
 \widehat C_6(3)=
 \mathcal F_6\left(\frac{\boldsymbol\mu^{a}+\boldsymbol\mu^{b}}2\right),
 \qquad\widehat\sigma_{\rm common}=0,
}\label{gff:prec:d3combined}\end{equation}
with the covariance \eqref{gff:meth:covariance} at $w_a=w_b=1/2$.

The production ran on SDSC Expanse on 29 September 2026 as 458 independent
64-thread array elements. Sample identities do not depend on node, thread
count or checkpoint boundaries: two blocks of 1024 six-loop tasks run at 64
and at 16 threads gave identical rows, including 141,994 binary128
promotions. The campaigns give
\begin{center}\begin{tabular}{@{}lrr@{}}\toprule
Campaign & $C_6(3)$ & Sampling SE\\\midrule
$a$ & 33.57562307 & 0.00015062\\
$b$ & 33.57523021 & 0.00018311\\
Combined & 33.57542664 & 0.00011855\\\bottomrule
\end{tabular}\end{center}
They differ by 1.66 combined standard errors. The combined standard error
is 0.54 of the value 0.00022 projected from the refinement campaigns by
plain Monte Carlo scaling, in which the standard error falls as $1/\sqrt n$;
the replicate covariance, not the projection, is reported. The estimator lies
0.30 standard errors of the superseded estimator \eqref{gff:meth:combined} from
its value $33.5681968\pm0.0241840$, and 0.52 refinement standard errors from
the refinement-only estimate.

The reported allowance is four sampling standard errors, 0.00047, rounded
up to a multiple of $0.0001$, which gives $C_6(3)=33.5754\pm0.0005$. The
estimator, the four-standard-error rule and the consistency checks were
written down and confirmed before any production result was available; the
rounding step was refined afterwards from the $0.001$ first registered to
$0.0001$, which changes only the rounding of the allowance. An independent
three-dimensional calculation through seventh order
(Section~\ref{gff:app:C7d3}), in which the power-containing sector is
integrated explicitly instead of through the analytic compensation, gives
as a by-product $C_6(3)=33.57387\pm0.00198$ (sampling standard error),
$-0.79$ combined standard errors from the estimator. It shares the graph
and renormalization framework and therefore does not test that framework.
The previous
three-dimensional baseline and refinement campaigns and the estimator
\eqref{gff:meth:combined} built from them are superseded. Their records are in
the reference tree that the two- and three-dimensional sixth-order packages
share and that is shipped once, in \path{LRI_C6_D2_HPC/reference/}
(\path{records/}, \path{baseline_previous/records/} and
\path{verification/method_precision_results.json}); the replay of
Section~\ref{gff:cn:reproduction} recomputes them from there.

\subsection{Result and interpretation of the errors}
The current estimates are
\begin{equation}\boxed{\begin{aligned}
 C_6(2)&=362.366\pm0.004,\\
 C_6(3)&=33.5754\pm0.0005.
\end{aligned}}\label{gff:meth:reported}\end{equation}
The allowances are four estimated sampling standard errors rounded
upward, to a multiple of $0.001$ at $d=2$ and of $0.0001$ at $d=3$. They are empirical
numerical allowances, not rigorous enclosures, guaranteed confidence
intervals, or truncation errors for the physical $\epsilon$ expansion.
The underlying estimator digits are recorded for reproducibility, not
as independently certified physical digits.
\begin{center}\begin{tabular}{@{}crr@{}}\toprule
$d$ & Estimator & Sampling standard error\\\midrule
2 & 362.36647176 & 0.00079973 \\
3 & 33.57542664 & 0.00011855 \\
\bottomrule\end{tabular}\end{center}
At $d=2$ and $d=3$ the estimates are the 2026 productions of
Sections~\ref{gff:prec:d2production2026} and~\ref{gff:prec:d3production2026} alone.
Fresh pilot and method-validation samples are excluded.
The lower-loop allocations were planned in part from the variances measured
in the superseded baseline campaigns (Section~\ref{gff:prec:baseline}); since the
allocation is fixed and the production samples are independently randomized,
this does not change the expectation of the estimate. No known $C_3,C_4,C_5$
is used as a fitted control or imposed as a constraint.

\subsection{A faster evaluator with the same integral}
Each monomial in a loop-$L$ forest polynomial has exactly $L$ distinct
edge factors. Associate a bit mask $m$ to a monomial and let $j$ be one
of its set bits. Precompute a directed acyclic product graph obeying
\begin{equation}
 p_0=1,\qquad p_m=x_jp_{m\setminus\{j\}}.
 \label{gff:meth:dag}
\end{equation}
All forest terms share these products. This replaces repeated
$\exp(\sum_{e\in m}\log x_e)$ evaluations by one exponential per edge
and a shared set of multiplications. For very small coordinates the
same common-monomial scaling as the original implementation is used, with
$h$ the largest of the monomial logarithms $\sum_{e\in m}\log x_e$ in the
forest sum, as in Section~\ref{gff:done:samplingControls}:
\begin{equation}
 p_m=\prod_{e\in m}\exp\left(\log x_e-\frac{h}{L}\right)
 =\exp\left(\sum_{e\in m}\log x_e-h\right).
 \label{gff:meth:scaled}
\end{equation}
The common factor is restored outside the forest sum. The direct
unscaled path is used only in a range checked against the exponent
range of extended precision. The original binary128 and 300-decimal
fallbacks for severe cancellation are retained. The numerical kernels
are not evaluated at lower precision to obtain the speedup.

An identical-point comparison used 1,095,680 six-loop evaluations for
each dimension, with identical seeds, digital matrices, allocations,
and importance distributions:
\begin{center}\begin{tabular}{@{}crrr@{}}\toprule
$d$ & Original evaluator (s) & Shared products (s) & Speedup\\\midrule
2 & 7.88237 & 2.77982 & 2.836\\
3 & 7.74461 & 2.69407 & 2.875\\\bottomrule
\end{tabular}\end{center}
These are elapsed times on one four-CPU machine, not
hardware-independent throughput claims. The maximum normalized
per-graph replicate difference was below $9.0\times10^{-13}$.
The improvement is in cost per evaluation, not in the variance of an
individual sample. Independently constructed Kirchhoff cofactor
integrands, built from cofactor determinants rather than from the sampler's list of
spanning-tree monomials, supplied 1,038 additional checks over both $d=2$ and $d=3$ and all
finite kernels, with reference values at 70 digits at ordinary points and at 650 digits at
extreme ones, including endpoint coordinates down to $10^{-130}$; the largest normalized discrepancy
was $6.88\times10^{-18}$. Neither comparison is an interval certificate
for the full integral.

\subsection{Sampling methods tested at matched budgets}\label{gff:mgs:benchmarks}
The baseline sampler is the one of Section~\ref{gff:prec:stratification}:
independently randomized Sobol nets inside each fixed graph stratum, with
power-of-two counts chosen separately for every stratum. The uncertainty is
always estimated between independent randomized nets, never from the dependent
points within one net \cite{gff:PrecisionSobol,gff:PrecisionQMC}. The benchmark target
is the six-loop influence \eqref{gff:meth:lastgradient}: with every six-loop finite
kernel sampled, let $y_{G,r}=(3c_{6,0}-b_{6,0})_{G,r}/2187$ be the contribution
of graph $G$ in replicate $r$. Its variance is estimated by
\begin{equation}
 \widehat V_{\rm bench}=\frac1R\sum_G\widehat V_G,\qquad
 \widehat V_G=\frac1{R-1}\sum_{r=1}^{R}\bigl(y_{G,r}-\bar y_G\bigr)^2,
 \label{gff:mgs:benchVariance}
\end{equation}
the sum of the graph-stratum replicate variances divided by the number $R$ of
replicates. It contains no covariance between strata; the production analysis
instead keeps the full covariance of the total replicate vectors,
\eqref{gff:prec:jetcov}.

Each pilot comparison used $R=16$ randomizations and 1,095,680 six-loop kernel
evaluations. The per-graph counts are powers of two, at least 64, the same at
both dimensions, and were fixed in advance from the graph probabilities of an
earlier, independent pilot, not from this benchmark
(\path{LRI_C6_D2_HPC/reference/code/pilot.py}). Antithetic sampling used half
as many points with two kernel evaluations each. The table gives
$\widehat V_{\rm bench}$ relative to the baseline at the same number of kernel
evaluations; smaller is better.
\begin{center}\begin{tabular}{@{}p{0.59\textwidth}rr@{}}\toprule
Method & $d=2$ & $d=3$\\\midrule
Randomized Sobol nets (baseline) & 1.000 & 1.000 \\
Independent pseudo-random Monte Carlo & 1.899 & 1.499 \\
Radial coordinates placed first & 1.105 & 1.126 \\
One-coordinate edge-order encoding & 1.083 & 1.055 \\
Endpoint rate $0.6\kappa$ & 1.667 & 2.301 \\
Endpoint rate $0.5\kappa$ & 4.194 & 6.194 \\
Endpoint rate $1.0\kappa$ & 4.452 & 1.797 \\
Radial antithetic pairs & 1.584 & 1.629 \\
\bottomrule\end{tabular}\end{center}
The variants are those of \path{code/experiment.cpp} in the same package. In
the baseline layout the $E$ ordering draws and the $E-1$ scale draws of the
conditional-draw map alternate in the $2E-1$ net coordinates. ``Radial
coordinates placed first'' puts the $E-1$ scale (radial) coordinates in the
first net dimensions and the ordering draws after them. The one-coordinate
encoding is the compressed map \eqref{gff:prec:compression}, with $E$ coordinates.
The antithetic variant evaluates each point together with its reflection
$v\mapsto1-v$ in the radial coordinates only, keeping the ordering coordinates,
and averages the pair. The endpoint rates are defined in the next paragraph.
The endpoint-rate and antithetic variants were run in the radial-first layout,
so their ratios include the layout effect of the third row; for pseudo-random
sampling the layout is immaterial.

The pilot differences are empirical and subject to finite-replicate
uncertainty; they are not theorems about the methods. Coordinate reordering,
and compression of all order choices into one uniform coordinate, gave no
convincing improvement. The compressed map was used only by baseline
campaign~B (Section~\ref{gff:prec:stratification}); both refinement campaigns, like
both campaigns of the 2026 productions, use the conditional-draw map.
Independent pseudo-random sampling was worse than randomized Sobol sampling.

\paragraph{Why a smoother endpoint map was not the winner.}
Write $\kappa(S)$ for the nominal endpoint power of the sector variable that
closes the edge subset $S$. The proposal exponents $\alpha(S)$ of the baseline
sampler, used as in \eqref{gff:done:proposalDP}, correspond to the nominal endpoint
rate $\alpha(S)=0.8\,\kappa(S)$. The rows ``endpoint rate'' multiply all of
them by $0.75$, $0.625$ and $1.25$, which gives
$\alpha=0.6\kappa$, $0.5\kappa$ and $1.0\kappa$; the dynamic-programming
normalization is recomputed with the changed exponents and the exact
reciprocal proposal weight is applied, so the target integral is unchanged. An
illustrative endpoint integral is
\begin{equation}
 \int_0^1t^{\kappa-1}f(t)\,dt
 =\frac1\alpha\int_0^1
 u^{\kappa/\alpha-1}f(u^{1/\alpha})\,du,\qquad t=u^{1/\alpha};
 \label{gff:mgs:power}
\end{equation}
sampling $t$ with the density $\alpha t^{\alpha-1}$ is this substitution with
uniform $u$. Reducing $\alpha$ makes the endpoint factor $u^{\kappa/\alpha-1}$
smoother, but it simultaneously changes the normalized sampling distribution in
every radial coordinate and can move variance into a narrow central region. The
variants $0.6\kappa$ and $0.5\kappa$ lost to the baseline in both dimensions,
and so did $1.0\kappa$. Smoothness alone is not a sufficient optimization
criterion for these finite, signed, high-dimensional kernels.


\paragraph{A pilot-frozen separable importance map.}
A separate training stream estimates the conditional second moment
in 32 bins for each radial coordinate. Its square root is smoothed,
clipped, and blended equally with a uniform density. Inverting this
positive piecewise-linear cumulative distribution defines $u=T(v)$.
The exact change of measure is
\begin{equation}
 \int_{[0,1]^s}f(u)\,du
 =\int_{[0,1]^s}f(T(v))
 \prod_{j=1}^sT_j'(v_j)\,dv.
 \label{gff:meth:transport}
\end{equation}
All Jacobian factors are included in the kernel. This is a small
VEGAS-inspired separable transport, not an implementation or benchmark
of the full VEGAS+ algorithm. VEGAS+ additionally uses adaptive
stratification and has documented advantages for some multi-peak and
diagonal structures~\cite{gff:MethodVegas}.

Fresh validation used 32 randomizations and 4,382,720 evaluations per
map. Relative variance per evaluation was
\begin{center}\begin{tabular}{@{}lrr@{}}\toprule
Frozen transport & $d=2$ & $d=3$\\\midrule
None & 1.0000 & 1.0000\\
Three trained axes & 0.9533 & 0.9871\\
All eleven trained axes & 0.8985 & 0.9905\\\bottomrule
\end{tabular}\end{center}
Thus the all-axis map supplied a modest approximately 10\% variance
reduction at $d=2$, but no material improvement at $d=3$. Among the refinement
campaigns only campaign A at $d=2,L=6$ uses this map; refinement campaign B
keeps the baseline map. In the 2026 production only campaign $a$ uses it
(Section~\ref{gff:prec:d2production2026}). Both three-dimensional campaigns use the baseline
map. No map is adapted using its production samples.
The mathematical map has endpoints $0,1$. Delivered decimal maps have
endpoint roundoff below $7.8\times10^{-16}$; a same-point replay with
endpoints snapped exactly to $0,1$ changed the six-loop mean by
$-2.95\times10^{-13}$. This is a floating-point consistency check,
not a bound on all digital-discretization effects.

The trained maps are built by \path{code/train_grid.py} of the same package.
At six loops each graph has $E-1=11$ radial coordinates. On every axis the
training histogram of the second moment is smoothed with the five-bin filter
$(1,2,3,2,1)/9$; its square root is normalized to mean one on that axis and
clipped to $[0.15,4]$, and the density is half this clipped score and half
uniform. ``Three trained axes'' transports, for each graph, the three radial
coordinates whose normalized score varies most across the 32 bins and leaves
the others uniform; ``all eleven'' transports every radial coordinate.

\paragraph{Conditional stratification and linear controls.}
Exact stratification of the first edge-order choice was also tested. The first
edge $j$ is chosen with the probability $p(j\mid E)$ of \eqref{gff:prec:conditional};
in the stratified variant each choice receives its own randomized net, with a
count proportional to $p(j\mid E)$, rounded to a power of two and at least
four, and its estimate enters with the weight $p(j\mid E)$. On fresh validation
streams with 32 randomizations, its variance times the number of kernel
evaluations, relative to the same product for the matched unstratified
validation run, was 1.043 at $d=2$ and 1.062 at $d=3$. It was not adopted.

Separately, let $\Sigma$ be the covariance of the estimates
$(\widehat C_3,\widehat C_4,\widehat C_5,\widehat C_6)$ of the baseline
campaigns, and $\Sigma_{\rm low}$ its $3\times3$ block for $C_3,C_4,C_5$. The
covariance-optimal linear control by the known $C_3,C_4,C_5$ would reduce the
standard error of $\widehat C_6$ by the factor
\begin{equation}
 \frac{\widehat\sigma_6^{\rm lin}}{\widehat\sigma_6}
 =\Bigl(1-\sum_{k,l=3}^{5}\Sigma_{6k}\bigl(\Sigma_{\rm low}^{-1}\bigr)_{kl}
 \Sigma_{l6}\big/\Sigma_{66}\Bigr)^{1/2},
 \label{gff:mgs:linearControl}
\end{equation}
which is 0.985 at $d=2$ and 0.962 at $d=3$. (Formula \eqref{gff:mgs:linearControl}
reproduces both ratios from the covariance stored for the baseline-campaign estimator, field \texttt{previous} of
\path{LRI_C6_D2_HPC/reference/verification/method_precision_results.json}.)
This is a diagnostic within that particular linear span, not a universal bound,
and those controls were neither fitted nor applied. They cannot directly remove
six-loop sampling variation, since $C_3,C_4,C_5$ contain no six-loop input.


\subsection{Refinement campaigns and cost-aware allocation}
\label{gff:meth:refinement}
Let $a_g/n_g$ model the variance of the target contribution from graph
$g$, and let $c_g$ denote the measured cost of one sample. Minimizing
this model at fixed cost gives
\begin{equation}
 \min_{n_g>0}\sum_g\frac{a_g}{n_g},\qquad
 \sum_gc_gn_g=T
 \quad\Longrightarrow\quad
 n_g=\lambda\sqrt{\frac{a_g}{c_g}}.
 \label{gff:meth:allocation}
\end{equation}
The derivative of $C_6$, rather than raw beta-function variance,
defines the pilot target. New six-loop allocations use the independent
method-validation streams, and lower-loop allocations use the variances
measured in the baseline campaigns of Section~\ref{gff:prec:baseline}; no sample
of a refinement campaign sets its own allocation. Each count is rounded to a power of two and frozen before
production integration. This is an MC-variance-model allocation used for
RQMC, not a claim that the RQMC variance is exactly proportional to
$1/n_g$ or that the allocation is globally optimal.

The baseline samples were used twice by the previous estimators: to plan the lower-loop
refinement allocations, and as two of the four pooled campaigns of
Section~\ref{gff:meth:pooling}. This does not change the expectation of the pooled
jet means. The allocation only sets the variance of the independently
randomized refinement samples, and the combination weights, $1/4$ for each
campaign and, at $d=2$, $1/2$ for the 66 nonprimitive kernels added
separately, were fixed before the refinement production
(\path{LRI_C6_D2_HPC/reference/inputs/frozen_production_plan.json}). They are
neither inverse-variance weights nor chosen from any result.


Each refinement campaign has 32 independently randomized replicates at
each loop order. The per-replicate point counts are
\begin{center}\begin{tabular}{@{}crrrr@{}}\toprule
$d$ & $L=3$ & $L=4$ & $L=5$ & $L=6$\\\midrule
2 & 135,168 & 264,192 & 1,158,144 & 4,647,936\\
3 & 151,552 & 305,152 & 1,189,376 & 4,381,184\\\bottomrule
\end{tabular}\end{center}
The refinement campaigns comprise 385,744,896 kernel evaluations at $d=3$
and 397,148,160 at $d=2$. Both are superseded; each row of the table is the
frozen allocation that the corresponding 2026 production multiplies by
$2^{13}$ (Sections~\ref{gff:prec:d2production2026}
and~\ref{gff:prec:d3production2026}). Of the refinement-stage runs, only these
two campaigns entered the previous estimators, together with the two baseline
campaigns. The 34,144,256 two-dimensional evaluations of the
dimension-coupled control experiment of Section~\ref{gff:meth:crossdimension} are
counted separately; that experiment entered neither the previous nor the
current estimators.
Pilot, method-validation, speed-test, and endpoint-replay samples are
not included. Campaigns A and B use different seeds and 52-bit and
60-bit digital coordinates, respectively. The source and input hashes
were frozen before production and rechecked afterward.

The three-dimensional refinement-campaign determinations, before
combination with the baseline campaigns, were
\begin{center}\begin{tabular}{@{}lrr@{}}\toprule
& $C_6(3)$ & Sampling SE\\\midrule
Refinement A & 33.59417512 & 0.040029 \\
Refinement B & 33.52750475 & 0.039502 \\
Refinement only & 33.56084035 & 0.028119 \\
\bottomrule\end{tabular}\end{center}
At $d=3$ the campaigns share no retained input, so the standard error of the
difference of two estimates is the quadrature sum of their standard errors.
Refinement A minus refinement B is $+1.19$ such standard errors, and the
refinement-only value minus the estimator of the two baseline campaigns alone,
$33.57555796\pm0.039354$, is $-0.30$ such standard errors (fields
\texttt{fresh\_campaign\_pull} and \texttt{fresh\_minus\_old\_pull} of the
record \path{verification/method_precision_results.json} of the earlier
calculation). Pooling all four campaigns in the previous estimator
\eqref{gff:meth:combined} lowered the three-dimensional sampling standard error
from 0.039354, for the baseline campaigns alone, to 0.024184, a factor 1.63,
which also reflects the larger number of production samples.
The 2026 code reproduces the refinement-only value bit for bit by replaying
the archived streams (Section~\ref{gff:prec:d3production2026}).

At $d=2$ the corresponding determinations were
\begin{center}\begin{tabular}{@{}lrr@{}}\toprule
& $C_6(2)$ & Sampling SE\\\midrule
Refinement A & 362.41502790 & 0.275726 \\
Refinement B & 362.46528830 & 0.273306 \\
Refinement only & 362.44015668 & 0.196310 \\
\bottomrule\end{tabular}\end{center}
The refinement campaigns integrated the 66 nonprimitive kernels of the groups
\eqref{gff:meth:retained} but not the ten primitive ones. Each value in this table
therefore adds $P_{10}$ once, and each standard error contains
$\widehat\sigma_P$ once. Because $P_{10}$ is common to the three rows, its
variance is removed from campaign-difference tests but not from the
uncertainty of each complete value. With it removed, campaigns A and B differ
by $-0.13$ standard errors of the difference,
$(\widehat\sigma_A^2+\widehat\sigma_B^2-2\widehat\sigma_P^2)^{1/2}$. The
refinement-only value lies $+0.62$ standard errors of the difference, again
with the $P_{10}$ variance removed, from the estimator of the baseline
campaigns alone, $362.25767208\pm0.226352$, which contains $P_{10}$ too (fields \texttt{fresh\_campaign\_pull} and
\texttt{fresh\_minus\_old\_pull} of
\path{LRI_C6_D2_HPC/reference/verification/method_precision_results.json}).
Pooling the four campaigns lowered the sampling standard error of the
two-dimensional estimator from 0.226352, for the baseline campaigns alone, to
0.152644, a factor 1.48. This gain also uses more production samples; it is
not attributed to the change of sampling algorithm alone. The speedup of the
shared-product evaluator \eqref{gff:meth:dag} lowers the cost per evaluation, a
different quantity from the variance per evaluation, and the two are not
combined.


\subsection{Campaign pooling, retained inputs and covariance}
\label{gff:meth:pooling}
The previous three-dimensional estimator pooled the four campaign jet means,
from two baseline and two refinement campaigns, with fixed weights $1/4$;
every vector is complete:
\begin{equation}\boxed{
 \widehat C_6(3)=
 \mathcal F_6\left(\frac{\boldsymbol\mu^{\rm old,A}
 +\boldsymbol\mu^{\rm old,B}+\boldsymbol\mu^{\rm new,A}
 +\boldsymbol\mu^{\rm new,B}}4\right).
}\label{gff:meth:combined}
\end{equation}
There is no separately retained numerical kernel at $d=3$.

The previous two-dimensional estimator pooled the same four campaign
generations, but replaced the 76 six-loop finite kernels of
Section~\ref{gff:cn:finiteKernels} by three independently measured groups. In
the finite-MOM normalization these were
\begin{equation}\begin{aligned}
 P_{10}&=718.7129769223133, &\widehat\sigma_P&=0.0414043282783,\\
 S_{26}&=-188.2890632768604, &\widehat\sigma_{26}&=0.1002179257784,\\
 S_{40}&=-2.129275292198507, &\widehat\sigma_{40}&=0.0843482336693.
\end{aligned}\label{gff:meth:retained}\end{equation}
These are not the separately MS-converted partial sectors
$\Theta_{26},\Theta_{40}$ of Section~\ref{gff:cn:finiteKernels}. Its baseline campaigns omitted all 76 kernels and its
refinement campaigns only the ten primitive ones, so it added
$P_{10}+\tfrac12S_{66}$, with $S_{66}=S_{26}+S_{40}$, and the common variance
\begin{equation}
 \widehat\sigma_{\rm common}^2
 =\widehat\sigma_P^2+
 \frac14(\widehat\sigma_{26}^2+\widehat\sigma_{40}^2).
 \label{gff:meth:common}
\end{equation}
The 2026 productions \eqref{gff:prec:d2combined} and \eqref{gff:prec:d3combined}
integrate every kernel in every replicate. Both current estimators therefore have
$\widehat\sigma_{\rm common}=0$, and no retained input is averaged away as
though independent.

\paragraph{The previous two-dimensional estimator.}
With the baseline campaigns labeled
old and the refinement campaigns new, it was
\begin{equation}
 \widehat C_6^{\rm prev}(2)=
 \mathcal F_6\left(\frac{\boldsymbol\mu^{\rm old,A}
 +\boldsymbol\mu^{\rm old,B}+\boldsymbol\mu^{\rm new,A}
 +\boldsymbol\mu^{\rm new,B}}4\right)
 +P_{10}+\frac12 S_{66}.
 \label{gff:mgs:previous2}
\end{equation}
The six-loop vectors of the two baseline campaigns omit all 76 kernels of
\eqref{gff:meth:retained}, those of the two refinement campaigns omit only the ten
primitive ones, and all lower-loop vectors are complete. Inside the average the
66 nonprimitive kernels therefore have total weight $\tfrac14+\tfrac14=\tfrac12$
and the ten primitive kernels weight zero, so adding $P_{10}+\tfrac12S_{66}$
gives every required contribution total expectation weight one. The groups
$P_{10}$, $S_{26}$, $S_{40}$ were measured independently, so the added term has
variance $\widehat\sigma_P^2+\tfrac14(\widehat\sigma_{26}^2+\widehat\sigma_{40}^2)$,
which is \eqref{gff:meth:common}, and it enters \eqref{gff:meth:covariance} once. The
baseline records also contain four-point diagnostic evaluations of the 76
omitted kernels; they are dropped and never enter as measurements. The
estimators of the refinement campaigns alone add $P_{10}$ once, with variance
$\widehat\sigma_P^2$, and that of the baseline campaigns alone adds
$P_{10}+S_{66}$. The last sum is the replacement used in run~C of
Section~\ref{gff:done:samplingControls}:
$P_{10}+S_{26}+S_{40}=528.2946383532544$ and
$(\widehat\sigma_P^2+\widehat\sigma_{26}^2+\widehat\sigma_{40}^2)^{1/2}=0.13737749$
agree with $K_{76}^{\rm MOM}$ and its standard error in
\eqref{gff:done:knownMOM76} to all printed digits.

The sampling standard error of \eqref{gff:mgs:previous2}, propagated with
\eqref{gff:meth:covariance}, was $0.152644$, made of $0.005635$, $0.037539$,
$0.080576$ and $0.096763$ from loops 3 to 6 and $0.077485$ from the common
term, combined in quadrature. That of the previous three-dimensional estimator
\eqref{gff:meth:combined} was $0.024184$, made of $0.001145$, $0.005894$,
$0.012458$ and $0.019839$ from loops 3 to 6, with no common term, combined in
quadrature (\path{verification/method_precision_results.json} of the earlier
calculation).


For each campaign and loop, let $\widehat\Sigma_{L,t}$ be the covariance
of its mean estimated from independent randomized replicates. With
$J_L=\partial\mathcal F_6/\partial\boldsymbol\mu_L$, the propagated
sampling variance is
\begin{equation}
 \widehat{\operatorname{Var}}(\widehat C_6)
 =\sum_{L,t}w_t^2J_L\widehat\Sigma_{L,t}J_L^{\mathsf T}
 +\widehat\sigma_{\rm common}^2.
 \label{gff:meth:covariance}
\end{equation}
All coupling/mass and Taylor-order covariance is retained. The error
budget, in units of $C_6$, is
\begin{center}\begin{tabular}{@{}crrrrrr@{}}\toprule
$d$ & $L=3$ & $L=4$ & $L=5$ & $L=6$ & Common & Total\\\midrule
2 & 0.000005 & 0.000113 & 0.000365 & 0.000702 & 0.000000 & 0.000800 \\
3 & 0.000001 & 0.000021 & 0.000058 & 0.000101 & 0.000000 & 0.000119 \\
\bottomrule\end{tabular}\end{center}
Both rows are the 2026 productions, and the common term is zero in both. Contributions in this table combine in quadrature, not by addition of
standard errors.

\subsection{A dimension-coupled primitive control experiment}\label{gff:meth:crossdimension}
This control experiment belongs to the earlier calculation. It concerns only
the previous three-dimensional estimator \eqref{gff:meth:combined}, and it
entered neither that estimator nor the 2026 production \eqref{gff:prec:d3combined}.

The Sobol matrices and graph shifts of a three-dimensional campaign can also
evaluate the ten primitive six-loop kernels at $d=2$, on the same randomized
net points; each dimension keeps its own importance exponents. This uses
common random numbers; it assumes no interpolation in $d$ and no relation
between the two physical coefficients. Let $\widehat P^{\rm pair}(2)$ be this
paired estimate, normalized as $P_{10}$ in \eqref{gff:meth:retained}, and
$\widehat P^{\rm ref}(2)$ an independent estimate of the same quantity with
standard error $\sigma_{\rm ref}$. For a fixed coefficient $\theta_{\rm cv}$,
\begin{equation}
 \widehat C_6^{\rm cv}(3)=\widehat C_6(3)
 -\theta_{\rm cv}\bigl(\widehat P^{\rm pair}(2)-\widehat P^{\rm ref}(2)\bigr).
 \label{gff:meth:crosscontrol}
\end{equation}
The difference in parentheses has expectation zero. The paired estimate is
correlated with the sampled three-dimensional primitive block; the reference
is independent of both and not exact, so its variance enters once:
\begin{equation}\begin{aligned}
 \operatorname{Var}\widehat C_6^{\rm cv}(3)
 &=\operatorname{Var}\widehat C_6(3)
 +\theta_{\rm cv}^2\bigl(\operatorname{Var}\widehat P^{\rm pair}(2)
 +\sigma_{\rm ref}^2\bigr)
 -2\theta_{\rm cv}\operatorname{Cov}\bigl(\widehat C_6(3),
 \widehat P^{\rm pair}(2)\bigr),\\
 \theta_{\rm cv}^{\rm opt}
 &=\frac{\operatorname{Cov}\bigl(\widehat C_6(3),\widehat P^{\rm pair}(2)\bigr)}
 {\operatorname{Var}\widehat P^{\rm pair}(2)+\sigma_{\rm ref}^2}.
\end{aligned}\label{gff:meth:noisycontrol}\end{equation}
The optimum assumes a known population covariance; it does not justify
fitting production coefficients without error control.

An independent paired pilot fitted one coefficient for each of the groups
$P_{10}$, $S_{26}$ and $S_{40}$ of \eqref{gff:meth:retained}, the primitive,
one-subgraph and two-subgraph blocks: the least-squares slope of the
three-dimensional block on its two-dimensional partner over the pilot
replicates, that is, the plain ratio of covariance to variance, without the
term $\sigma_{\rm ref}^2$ of \eqref{gff:meth:noisycontrol}.
On independent validation streams with 32 randomizations the residual-variance
ratios were 0.234, 0.743 and 0.597; these ratios also leave out the reference
term. The control doubles the kernel evaluations of a block, so the chosen
equal-cost rule kept a block only if twice its ratio was below one
(\path{inputs/cross_primitive_production_plan.json}). Only the primitive block
passed, with the frozen values
\begin{equation}
 \theta_{\rm cv}=0.09955611941934894,\qquad
 \widehat P^{\rm ref}(2)=P_{10},\qquad
 \sigma_{\rm ref}=\widehat\sigma_P.
 \label{gff:meth:crossbeta}
\end{equation}
No production sample was used to fit $\theta_{\rm cv}$ or to choose the block.
Evaluating the primitive block at $d=2$ on the stored net points of all four
campaigns of \eqref{gff:meth:combined} added 34,144,256 two-dimensional
evaluations and no three-dimensional sample, with identical seeds and digital
matrices verified (\path{verification/cross_production_manifest.json}). An
independent three-dimensional replay of 894,464 primitive evaluations, the
ten primitive graphs in the first two replicates of baseline campaign B and of
refinement campaign A, reproduced the stored graph-replicate values with
maximum difference zero (\path{verification/pairing_checks.json}). With the
campaign weights $1/4$ of \eqref{gff:meth:combined} the controlled result is
\begin{equation}
 \widehat C_6^{\rm cv}(3)=33.5681256836,\qquad
 \widehat\sigma_{\rm cv}=0.0243339897
 \label{gff:meth:cvresult}
\end{equation}
(\path{verification/controlled_precision_result.json}), against
$33.5681968\pm0.0241840$ for the plain estimator.

The validation variance of the primitive block fell by a factor of about 4.3,
but this block carries only about 5.8\% of the plain variance. In the
production its sampling error fell from 0.00584436 to 0.00325730; the
reference adds $\theta_{\rm cv}\widehat\sigma_P=0.00412205$, giving 0.00525370
in all (\path{verification/cross_control_budget.json}). In variance, the
control lowered the sampling part of the complete coefficient by
$9.7\times10^{-6}$ (its six-loop error fell from 0.019839 to 0.019593), while
the reference added $(\theta_{\rm cv}\widehat\sigma_P)^2=1.70\times10^{-5}$.
Because the reference term exceeds the gain, the controlled error is slightly
larger than the plain 0.02418404; this small difference is not evidence of a
general disadvantage of controls. With the reference term included, the
optimum of \eqref{gff:meth:noisycontrol} estimated after the production is
$0.0579$ (same record), a diagnostic that was not used. The experiment shows
that reducing the variance of a strongly correlated subintegral need not
reduce the uncertainty of the complete coefficient, and that an uncertain
auxiliary value cannot be treated as an exact control mean. The plain
estimator \eqref{gff:meth:combined} remained the previous three-dimensional
result; the controlled value served only as a consistency check and a method
benchmark.

\subsection{Final checks and evidence limits}
The same uncalibrated samples give the lower-order discrepancy checks
\begin{center}\begin{tabular}{@{}crrr@{}}\toprule
$d$ & $C_3$ pull & $C_4$ pull & $C_5$ pull\\\midrule
2 & -0.97 & -1.54 & -1.07 \\
3 & -1.69 & -1.09 & +0.92 \\
\bottomrule\end{tabular}\end{center}
Both rows are the 2026 productions. The reference special-function
values are checks only; their numerical values are not refitted or replaced
by the Monte Carlo estimates.
For the previous estimators, an archived independent finite two-variable
Lagrange expansion agreed with the fixed-point iteration to
$7\times10^{-13}$ at $d=2$ and $5\times10^{-14}$ at $d=3$, and the
delete-one-replicate bias estimates were below $3.4\times10^{-7}$ and
$6.5\times10^{-8}$ in $C_6$ units; these check the inherited map
$\mathcal F_6$, not the 2026 replicate data. Bootstrap resampling of the
independent replicate vectors gives sampling errors 0.00080124 at $d=2$
and 0.00011534 at $d=3$ (4096 resamples each).
The bootstrap is a diagnostic, not an interval certification.
The bootstrap means differ from the direct estimates by
$-1.3\times10^{-5}$ at $d=2$ and $+2.1\times10^{-8}$ at $d=3$.

For the previous two-dimensional estimator \eqref{gff:mgs:previous2}, the same
uncalibrated samples gave the pulls
$+0.09$, $+0.50$ and $+0.60$ for $C_3$, $C_4$ and $C_5$. Its replicate
bootstrap, which draws the common inputs once per resample according to
\eqref{gff:meth:common}, gave the sampling error 0.14987068 from 4096 resamples,
against the propagated 0.15264446.

For the previous three-dimensional estimator \eqref{gff:meth:combined} the pulls
were $+0.44$, $+0.80$ and $-0.21$ for $C_3$, $C_4$ and $C_5$, and its replicate
bootstrap gave the sampling error 0.02376720 from 4096 resamples, against the
propagated 0.02418404 (\path{verification/method_precision_results.json} and
\path{verification/extra_checks.json} of the earlier calculation).

\paragraph{Verification scope.}
The final numbers retain the analytic power-sector compensation and the
same graph/renormalization framework; narrower sampling errors do not bound
shared errors in that framework. The reference $C_3$--$C_5$ values were used
for comparisons only, not imposed as constraints. The reported results are
$C_6(2)=362.366\pm0.004$ and $C_6(3)=33.5754\pm0.0005$; all earlier,
less precise full-coefficient estimates and partial sums are superseded.

The method comparison of Sections~\ref{gff:mgs:benchmarks} and~\ref{gff:mgs:selection}
does not bound a shared error in graph generation, renormalization or analytic
continuation, and it did not rerun the analytic power-sector compensation.
Faster arithmetic and smaller sampling errors are neither an independent
field-theoretic validation nor rigorous error enclosures.

\subsection{Method selection and methods not adopted}\label{gff:mgs:selection}
\paragraph{Best supported choice for these integrands.}
The strongest tested practical combination is local finite forest subtraction
\eqref{gff:prec:kernel}, tropical importance sampling \eqref{gff:done:proposalDP},
graph-stratified randomized Sobol nets (Section~\ref{gff:prec:stratification}),
pilot-frozen cost-aware allocation \eqref{gff:meth:allocation}, and the shared
polynomial products \eqref{gff:meth:dag} with monitored high-precision fallback. A
separately trained transport \eqref{gff:meth:transport} is only modestly useful at
$d=2$ and unnecessary at $d=3$ in these tests. The runtime gain and the
variance gains are different quantities and are not conflated, and the total
precision gain of the refinement campaigns also comes from more production
samples (Section~\ref{gff:meth:refinement}).

\paragraph{Methods not benchmarked here.}
Rank-one shifted lattice rules with smooth inverse-CDF importance fits are an
established alternative for multi-loop integrals, including CPU and CUDA
implementations interfaced to pySecDec \cite{gff:mgs:BorowkaQMC}. That
implementation was studied but not benchmarked, so it is not declared better or
worse than the present Sobol method. Global tropical sampling across primitive
graphs has recently reached much higher primitive loop orders
\cite{gff:mgs:BorinskyGlobal}; those positive primitive examples do not by
themselves establish its performance for the signed, regulator-dependent
renormalized forest combinations of this calculation.

\paragraph{Conditional preintegration.}
A promising next structural improvement is conditional preintegration: a subset
of the coordinates is integrated exactly and only the rest, $u_{\rm s}$, is
sampled. If $Y$ is an unbiased single-point estimator, the law of total
variance,
\begin{equation}
 \operatorname{Var}(Y)
 =\mathbb E\bigl[\operatorname{Var}(Y\mid u_{\rm s})\bigr]
 +\operatorname{Var}\bigl(\mathbb E[Y\mid u_{\rm s}]\bigr),
\end{equation}
has a nonnegative first term, so
\begin{equation}
 \operatorname{Var}\bigl(\mathbb E[Y\mid u_{\rm s}]\bigr)\le\operatorname{Var}(Y).
 \label{gff:mgs:rao}
\end{equation}
This is useful computationally only when the cost of the exact inner
integration and any deterministic quadrature error are controlled. No integral
was replaced by an unchecked low-order quadrature, and no favorable variance
reduction is assumed. The first-edge stratification of
Section~\ref{gff:mgs:benchmarks}, which lost at equal cost, shows why a
conditionalization has to be benchmarked at equal cost.


\section{Full seventh order in one dimension: coordinate forests and binary-cluster integration}
\label{gff:app:C7complete}

\subsection{Result, normalization, and scope}
This section continues the numerical calculation to seventh order at
$d=1$, with the unchanged convention
\begin{equation}
 \epsilon=2s-1,\qquad
 \Delta_{\mathcal E}=\frac12-\frac\epsilon6+
 \sum_{n=2}^{7} C_n(1)\epsilon^n+O(\epsilon^8).
 \label{gff:c7f:convention}
\end{equation}
The full graph-based numerical estimate is
\begin{equation}\boxed{
 C_7(1)=-322550,\qquad
 \text{pre-registered empirical allowance: }\pm41.
}\label{gff:c7f:answer}\end{equation}
The unrounded estimator is $-322549.889$, with propagated sampling
standard error $10.088$; a bootstrap over the replicates of every stream
gives $10.111$. The estimator uses the radial quadrature through degree six.
The allowance is the output of a policy fixed before production: the
largest of four sampling errors, twice the last refinement plus four of
its errors, twice the degree six to seven check step plus four of its
errors, and the half-split displacement, rounded up to two significant
figures. It is not a rigorous enclosure, a probability guarantee, or an
estimate of the error of the truncated physical $\epsilon$ series.

The policy value is not a demonstrated total accuracy. It takes the
largest of its candidates rather than combining the estimated truncation
with the sampling error. The measured degree six to seven step is
$+10.40\pm0.43$, and a geometric tail beyond degree seven, extrapolated
with the last step ratio $0.600$, adds $+15.59$. Treating that tail as
uncertain by $100\%$ and neglecting the covariance between the estimate
and the check level, the truncation-aware estimate is
\begin{equation}
 C_7(1)\approx-322524\pm38,
 \label{gff:c7f:truncationAware}
\end{equation}
with the allowance taken as twice the combined error
$10.09\oplus0.43\oplus15.59=18.6$. This estimate is post hoc and depends on
the tail model, whose ratio at seventh order is larger than at fifth and
sixth order (Section~\ref{gff:c7f:checks}). The intervals of
\eqref{gff:c7f:answer} and \eqref{gff:c7f:truncationAware} overlap. Establishing
the total error at this level would require radial data at degree eight or
an independently validated tail model. The long decimal strings in the
machine-readable files specify the estimator, not certified physical
digits. In particular, no independent determination of the common graph
and renormalization framework is claimed.

Both logarithmically and power-subdivergent diagrams are included. No
seventh-order power correction is assigned zero, and no Pad\'e, ratio,
large-order, or interpolation estimate is used to fill a missing graph
class. The 44 seventh-order primitive classes are integrated directly in
a separate campaign, giving $-696179.472\pm0.867$. That campaign replaces
the primitive contribution $-696147.446\pm25.368$ of the first, smaller
seventh-order calculation (23 September 2026, $C_7(1)=-323405\pm5862$ sampling
standard error; record \path{GFF/results/C7_1/v1_46batches/archive_original.json}),
which it confirms within $1.26$ combined standard errors. The
primitive number alone is not the physical $C_7$.

\subsection{Complete inventory and exact algebraic coverage}
The eight-vertex enumeration was rebuilt from integer edge multiplicities.
There are 3986 vertex-isomorphism classes, representing 125894580 labeled
1PI graphs. Of these, 1091 contain an unmarked power-divergent two-point
subgraph. The seventh-order $3c_7-b_7$ target has nonzero weight for 845 of
these classes. The primitive sector has 44 classes.

For a graph with vertex automorphism group $\Aut(G)$, internal
multiplicities $m_{ij}$, and external half-edge counts
$e_i=4-\sum_jm_{ij}$, the class weights are
\begin{equation}
 w_A(G)=\frac{24}{|\Aut(G)|\prod_i e_i!\prod_{i<j}m_{ij}!},\qquad
 w_B(G)=\frac{w_A(G)}{12}\sum_i e_i(e_i-1).
 \label{gff:c7f:weights}
\end{equation}
They include the vertex orbit factor and use the full coordinate domain;
there is no second vertex-factorial division in the numerical integral.
All weights and all forest signs are generated exactly before numerical
conversion. The inventory reproduces the previously recorded class and
labeled counts.

The weighted coordinate forests were independently compared with formal
bare-series inversion, retaining every graph amplitude as an independent
symbol and canonically factoring articulation products. The comparison
is a coefficientwise rational-polynomial identity, not a test at a few
numerical inputs:
\begin{center}
\begin{tabular}{r rr r}
\toprule
Loop order & Coupling monomials & Mass monomials & Nonzero differences\\
\midrule
1&1&1&0\\2&2&2&0\\3&8&6&0\\4&29&25&0\\
5&134&103&0\\6&696&525&0\\7&4164&2935&0\\
\bottomrule
\end{tabular}
\end{center}
These counts are taken after articulation products are identified; without
that factorization the counts are larger.

An additional graph-by-graph exact reduction gives zero RG polynomials
for 691 of the seventh-order classes. Among the 845 power classes with
nonzero channel weight, 151 are such exact zeros; the remaining 694 are
integrated. For the fifth- and sixth-order power inputs, respectively
24 and 132 nonzero classes remain after the same algebraic zero test.
These exclusions are identities, not variance-based pruning. All other
nonzero classes have positive production allocation.

The new 44 primitive representatives were matched by completed-graph
isomorphism to all 14 periods used by the first calculation. Every
aggregated rational weight agrees, including the three lower-period
products. The final production integrates these 44 representatives
directly; the agreement of that integration with those periods is quoted in
Section~\ref{gff:c7f:production}.

\subsection{One common coordinate-space finite prescription}
Computing the power-free and power-containing parts in two different finite
prescriptions would require a conversion between them. Here both are
computed in one coordinate-space prescription. An edgewise
exponential infrared cutoff is used:
\begin{equation}
 C_\mu(x;z)=\frac{2^{z/2}}{\sqrt{\Gamma(1+z)}}
 |x|^{-(1-z)/2}e^{-\mu|x|}.
 \label{gff:c7f:propagator}
\end{equation}
The one-loop bubble is exactly
\begin{equation}
 \int_{\mathbb R}C_\mu(x;z)^2\,dx
 =\frac2z\mu^{-z}.
 \label{gff:c7f:bubble}
\end{equation}
The edgewise cutoff factorizes on products of subgraphs; replacing it by
an independently assigned cutoff on the span of each whole graph would
change finite parts and is not done here.

Let the dimensionless bare coupling be $x$, and write the weighted bare
amplitude series as
\begin{equation}
 \mathcal A(x,z)=1+\sum_{L\ge1}(-1)^L a_L(z)x^L,
 \qquad
 \mathcal B(x,z)=1+\sum_{L\ge1}(-1)^L e_L(z)x^L.
 \label{gff:c7f:bare}
\end{equation}
The finite coupling is $g=x\mathcal A(x,z)$. The two RG functions follow
from differentiating with respect to the common scale:
\begin{equation}
 \widehat\beta(g,z)+zg=-z x^2\partial_x\mathcal A(x,z),\qquad
 \widehat\gamma(g,z)=-z x\partial_x\log\mathcal B(x,z),
 \quad x=x(g,z).
 \label{gff:c7f:finiteRG}
\end{equation}
Thus the regulator-dependent functions, not just their $z=0$ values,
are required. The physical substitution is $z=\epsilon$ after the Taylor
jets and finite fixed point have been assembled.

For clarity, the independent algebraic check extracts the coupling and
mass coefficients directly from
\begin{align}
 \widehat b_L(z)&=-\frac{z}{L+1}[x^L]
 \left(\sum_{j\ge1}j(j+1)(-1)^j a_j(z)x^j\right)
 \mathcal A(x,z)^{-L-1},\\
 \widehat c_L(z)&=-\frac zL[x^{L-1}]
 \partial_x\!\left(\frac{x\partial_x\mathcal B(x,z)}
 {\mathcal B(x,z)}\right)\mathcal A(x,z)^{-L}.
 \label{gff:c7f:bareInversion}
\end{align}
The factor $z$ is carried symbolically before Laurent cancellation. The
code does not estimate these expressions by subtracting independently
sampled divergent amplitudes.

\subsection{Local full-graph forest integrands}
For every compatible proper logarithmic forest $F$, replace its blocks
by their quotient graphs and glue the quotients on the original labeled
vertices. Let $G_F$ be the resulting edge multigraph. Terms whose quotient
is a tadpole are set to zero in the same analytically continued critical
mass prescription as the preceding sections. Unmarked power subgraphs
are not dropped: their coordinate integrals are continued meromorphically
as described below.

For projective coordinates $u$, put
\begin{equation}
 H_F(u)=\sum_{e\in G_F}|u_{e_1}-u_{e_2}|,
 \qquad
 R_F(u)=\sum_{e\in\mathrm{root}(F)}|u_{e_1}-u_{e_2}|.
 \label{gff:c7f:HR}
\end{equation}
The common finite RG integral for an $L$-loop graph is generated by
\begin{equation}
 \mathcal D_G(z)=
 \frac{2^{Lz}\Gamma(1+Lz)}{\Gamma(1+z)^L}
 \int d\Omega\sum_F(-1)^{|F|}
 \frac{R_F(u)}{H_F(u)}H_F(u)^{-Lz}
 \prod_{e\in G_F}|u_{e_1}-u_{e_2}|^{-(1-z)/2}.
 \label{gff:c7f:coordinateKernel}
\end{equation}
Then $\widehat b_L=(-1)^{L+1}\sum_Gw_A(G)\mathcal D_G$ and
$\widehat c_L=(-1)^{L+1}\sum_Gw_B(G)\mathcal D_G$.
The factor $R_F/H_F$ is produced by the scale derivative on the root
block, followed by the common scale integral. Explicitly the latter is
$\int_0^\infty\lambda^{Lz}\exp(-\lambda H_F)d\lambda$.
The exact bare-polynomial check above independently verifies all forest
and root weights. This is a local representation of renormalization in
the spirit of scale-dependent forest formulas\cite{gff:BrownKreimerLocal}, but the
specific kernel and its combinatorial certificate here are part of the
present calculation.

All terms in a forest combination are evaluated on the same coordinates.
For a fixed forest term the Taylor variable appears through
\begin{equation}
 \exp\left\{z\left[
 \frac12\sum_e\log|u_{e_1}-u_{e_2}|-L\log H_F(u)
 \right]\right\}.
 \label{gff:c7f:logJet}
\end{equation}
The common Gamma factor in \eqref{gff:c7f:coordinateKernel} has logarithmic
Taylor coefficients
\begin{equation}
 \ell_1=L\log2,\qquad
 \ell_j=\frac{L^j-L}{j!}\psi_{j-1}(1),\quad j\ge2.
 \label{gff:c7f:gammaJet}
\end{equation}
These are included through the required order, with all products and
Laurent poles retained before the finite coefficients are extracted.

\subsection{Binary vertex clusters instead of fully ranked gaps}
Order the $L+1$ vertices along the line. A binary tree is obtained by
splitting at the largest gap, then repeating inside each of the two
subintervals. This partitions the positive-gap domain except for
measure-zero ties. The internal nodes are contiguous vertex clusters;
there are $L$ of them. Let the root scale be $\lambda$, and let
$t_C\in[0,1]$ be the ratio between each nonroot node's gap and its parent's
gap. Then
\begin{equation}
 \delta_C=\lambda\prod_{D\preceq C,\ D\ne\mathrm{root}}t_D,
 \qquad
 \left|\frac{\partial(\delta_1,\ldots,\delta_L)}
 {\partial(\lambda,(t_C))}\right|
 =\lambda^{L-1}\prod_{C\ne\mathrm{root}}t_C^{|C|-2}.
 \label{gff:c7f:treeJacobian}
\end{equation}
The coordinate of every vertex distance is formed as a sum of positive
gaps, never by subtracting two nearly equal vertex positions.

Unlike a full ranking of all gaps, this construction does not impose an
unnecessary relative ordering between gaps in disjoint clusters. It has
$(L+1)!\,\mathrm{Cat}_L$, with $\mathrm{Cat}_L=(2L)!/[L!(L+1)!]$,
ordered tree sectors rather than
$(L+1)!L!$ ranked-gap sectors, while keeping $L-1$ continuous projective
variables. At seven loops there are six continuous cluster ratios. The
sector sum is sampled, not exhaustively enumerated.

The tree probabilities are computed without enumerating all trees. For
a vertex subset $S$, let $W_{\{i\}}=1$ and
\begin{equation}
 W_S=\frac1{\alpha_S}
 \sum_{\substack{A\sqcup B=S,\ A,B\ne\varnothing\\\min S\in A}}2W_AW_B,
 \qquad \alpha_{\mathrm{root}}=1.
 \label{gff:c7f:treeDP}
\end{equation}
At a split, choose the unordered pair with probability proportional to
$W_AW_B$ and then choose its orientation uniformly. The normalization is
$Z=W_{\mathrm{root}}$. A particular ordered tree has probability
$1/(Z\prod_{C\ne\mathrm{root}}\alpha_C)$.

For the continuous power-free calculation, the proposal sees every
forest term, not only the original graph. With $e_F(C)$ the number of
edges internal to $C$ in $G_F$, choose
\begin{equation}
 p_C=|C|-1-\frac12\max_F e_F(C),\qquad
 \alpha_C=r\max(\tfrac12,p_C),\qquad
 t_C=u_C^{1/\alpha_C}.
 \label{gff:c7f:treeRates}
\end{equation}
Here $r=0.8$ or $0.7$ in the two independent campaigns. The combined tree
and radial importance weight is simply
\begin{equation}
 Z\prod_{C\ne\mathrm{root}}t_C^{|C|-1-\alpha_C}.
 \label{gff:c7f:treeWeight}
\end{equation}
Using only the original graph's edge count can miss a more singular
forest term and can give infinite sampling variance; this happens for
seven-loop forests. The present construction examines the
whole finite forest kernel.

The Jacobian is checked by exact rational determinants for 84 random
cluster geometries. As an independent normalization check, set
$\alpha_C=|C|-1$ and remove all edge factors. The dynamic program then
gives $Z=(L+1)!L$ for $L=1,\ldots,7$, exactly the full projective volume.
Explicit enumeration of the small trees confirms the proposal
probabilities with arbitrary rational subset rates.

\subsection{Power subgraphs: meromorphic radial integration and paired refinement}
In a power-containing graph, a cluster may have negative radial exponent.
It is not sampled as an ordinary convergent endpoint integral. For each
sampled tree, the regular factor is interpolated on a tensor
Chebyshev--Lobatto grid of degree $N$, and the singular powers are
integrated analytically. If $L_i(t)=\sum_k\ell_{ik}t^k$, the one-variable
moment is
\begin{equation}
 \int_0^1t^{a+bz}L_i(t)\,dt
 =\sum_k\frac{\ell_{ik}}{a+k+1+bz},
 \label{gff:c7f:meromorphicMoment}
\end{equation}
under meromorphic continuation. A zero denominator at $z=0$ supplies its
Laurent pole explicitly; a negative half-integer power is continued by
the same rational expression. No endpoint cutoff is fitted or sent to
zero numerically. Tensor products of these moments are formed before
the finite Taylor coefficient is selected.

For this discrete tree sampling only, the rates use
$\alpha_C=r\max(1/2,|p_C|)$. The tree probability is positive and the
inverse probability is $Z\prod_C\alpha_C$. It is a probability on a
finite set of trees, not a claim that the unregularized radial power
integral converges.

Let $Y_N$ denote the complete per-tree Laurent/Taylor vector at degree
$N$. The production estimator uses the exact telescoping organization
\begin{equation}
 \mathbb E[Y_6]=\mathbb E[Y_1]
 +\mathbb E[Y_3-Y_1]
 +\mathbb E[Y_5-Y_3]
 +\mathbb E[Y_6-Y_5].
 \label{gff:c7f:levels}
\end{equation}
Within every difference, both resolutions use the same graph, binary
tree, and random stream. Different levels use independent streams.
This is the multilevel variance-reduction principle\cite{gff:C7Giles}; it
reduces the noise of the refinement differences without extrapolating
from lower physical coefficients. The residual difference between
finite degree six and the meromorphic integral is assessed by the
successive paired refinements, including the last degree-five/six
comparison,
not assumed to be zero by definition.

Negative Laurent coefficients of the fully subtracted RG functions
vanish. Fixed linear combinations of these exact-zero quantities are
used as control variates. The control matrices were fitted to separate
coarse pilot samples, with a fixed regularization, and then frozen before
production; they were not chosen to match any physical reference. A separate first seventh-order
degree-one/three refinement run was used to freeze a cost-aware graph
allocation; that training stream is excluded from the final seventh-order
refinement mean. Every nonzero remaining stratum keeps positive coverage.
The same matrix is applied at every level in \eqref{gff:c7f:levels}.
Consequently any finite-grid error in a control is also refined by the
telescoping calculation. Neither $C_5$, $C_6$, nor a conjecture that the
seventh-order power correction vanishes is imposed as a constraint.

\subsection{Seventh-order assembly and covariance}
Write
\begin{equation}
 \widehat\beta(g,\epsilon)=-\epsilon g+
 \sum_{L=1}^7\widehat b_L(\epsilon)g^{L+1},\qquad
 \widehat\gamma(g,\epsilon)=
 \sum_{L=1}^7\widehat c_L(\epsilon)g^L,
 \label{gff:c7f:jets}
\end{equation}
with $\widehat b_1=3$, $\widehat c_1=1$, and retain each jet through
$\epsilon^{7-L}$. Solving the finite fixed point is equivalent to
\begin{equation}
 C_7=[\epsilon^7]\sum_{k=1}^7\frac{\epsilon^k}{k}[g^{k-1}]
 \partial_g\widehat\gamma(g,\epsilon)
 \left[3+\sum_{L=2}^7\widehat b_L(\epsilon)g^{L-1}\right]^{-k}.
 \label{gff:c7f:finiteLagrange}
\end{equation}
In a regulator-independent MS scheme this reduces to
$\frac17[g^6]\gamma_r'(g)B(g)^{-7}$, but that shorter expression cannot
be used on just the constant terms of the present finite-scheme jets.

For fixed one- through four-loop inputs, $C_5,C_6,C_7$ are affine in all
fifth-, sixth-, and seventh-loop jets. Two such high-loop insertions
would first contribute beyond order seven. Thus, if $v_L$ is the vector
of coupling and mass coefficients (with negative Laurent orders kept
for controls),
\begin{equation}
 \boldsymbol C=\boldsymbol K_{\le4}
 +\sum_{L=5}^7J_L(v_L^{\rm free}+v_L^{\rm power}).
 \label{gff:c7f:affine}
\end{equation}
Here $\boldsymbol C$ is the vector of the coefficients $C_n$,
$\boldsymbol K_{\le4}$ its value from the one- through four-loop inputs alone,
and $J_L$ the exact matrix of its dependence on $v_L$. The complete matrices
$J_L$ and the lower-loop vectors are in
\path{LRI_C7_HPC/reference/verification/full_affine_assembly.json}, and the
frozen control matrices in
\path{LRI_C7_HPC/reference/verification/power_control_training.json}. Thirty-two exact rational tests
verify both \eqref{gff:c7f:finiteLagrange} against an independently iterated
fixed point and the affine identity \eqref{gff:c7f:affine}.

All coupling/mass and Taylor-order covariances are retained. Each random
replicate is projected through the full matrix before computing its
covariance. The two power-free campaigns are averaged with equal weights
fixed independently of their observed values. The direct primitive
campaign is added only once, with variance $(0.8667)^2$. The two
independent coarse power campaigns are combined with preassigned
sample-count weights. Independent refinement-level covariance matrices
are then added. Independent production batches of one frozen design are
combined with equal weights, and their replicate covariances are
propagated. The finite high-loop assembly itself has no nonlinear
Monte Carlo bias, because it is affine.

\subsection{Production results and refinement}\label{gff:c7f:production}
The final production consists of 399 independent batches of one frozen
design, each with its own seeds. The design was fixed before production
from a separate pilot: the replicate counts and sample multipliers of every
campaign were chosen by minimum-variance allocation at the measured pilot
variances and costs, with every multiplier a power of two. Each batch contains the two
independent power-free campaigns at five, six and seven loops (128
randomized replicates each); the two independent coarse power campaigns
(128 and 64 replicates); the paired refinements $1\to3$, $3\to5$ and $5\to6$
of \eqref{gff:c7f:levels} (128, 128 and 64 replicates); and a paired $6\to7$
check level (64 replicates), which is never part of the estimate. The 44
primitive classes are integrated once, in a separate campaign with
$2^{23}$ points per graph and 64 replicates. It gives
$-696179.472\pm0.867$, which agrees with the period-based value
$-696147.446\pm25.368$ of the first calculation within $1.26$ combined
standard errors. The
production used about $2.5\times10^5$ core-hours.

The telescoping resolution sequence for the complete coefficient is
\begin{center}\begin{tabular}{rrr}
\toprule
Power-grid degree & Full $C_7$ estimate & Sampling standard error\\
\midrule
1 & $-331{,}309.98$ & 6.14\\
3 & $-322{,}614.21$ & 9.21\\
5 & $-322{,}567.22$ & 10.01\\
6 & $-322{,}549.89$ & 10.09\\
\bottomrule\end{tabular}\end{center}
The individual refinement corrections, in absolute units of $C_7$, are
\begin{center}\begin{tabular}{lrr}\toprule
Paired degrees & Correction & Sampling standard error\\\midrule
$1\to3$ & $+8{,}695.78$ & 6.86\\
$3\to5$ & $+46.99$ & 3.93\\
$5\to6$ & $+17.33$ & 1.21\\
$6\to7$ (check only) & $+10.40$ & 0.43\\
\bottomrule\end{tabular}\end{center}
The corrections decrease, but the ratio of successive steps grows, from
$0.37$ for $5\to6$ over $3\to5$ to $0.60$ for $6\to7$ over $5\to6$. No
limit has been replaced by a fitted extrapolation in
\eqref{gff:c7f:answer}. The allowance candidates are $4\times10.09=40.35$,
$2\times17.33+4\times1.21=39.50$, $2\times10.40+4\times0.43=22.49$ and the
half-split displacement $12.58$, which gives $\pm41$.

The scatter of the 399 batch estimates about their mean gives
$\chi^2=348.9$ for 398 degrees of freedom, so the replicate-based errors
describe the batch-to-batch variation. A bootstrap over the replicates of
every stream gives the standard error $10.11$, compared with the covariance
value $10.09$. Splitting the replicates of every stream into halves gives
$-322537.30$ and $-322562.47$. The
two independent power-free campaigns agree at each loop order to within
$0.42$ of their combined sampling errors. The largest absolute
standardized residual among the 24 monitored spurious Laurent poles is
$2.31$. With the zero-pole controls the estimate is $71.6\pm41.0$ above the
uncontrolled estimator, $-322621.5\pm41.1$ (the same convention as at eighth
order, Section~\ref{gff:c8f:checks}).

\subsection{Lower-order checks and the radial truncation}\label{gff:c7f:checks}
The same calculation gives the following uncalibrated lower-order values:
\begin{center}\begin{tabular}{rrrr}\toprule
Coefficient & Degree-six value & Sampling error & Reference\\\midrule
$C_5(1)$ & $-731.174185$ & $0.008355$ & $-731.139477$ (exact)\\
$C_6(1)$ & $14{,}332.435930$ & $0.305002$ & $14{,}331.687033$ (deterministic)\\
\bottomrule\end{tabular}\end{center}
The differences are $-4.15$ and $+2.46$ sampling errors. At this precision
the lower orders resolve the truncation of the radial quadrature at degree
six. Adding the measured degree six to seven step and a geometric tail with
each coefficient's own last step ratio gives
\begin{center}\begin{tabular}{lrrr}\toprule
 & Degree six & $+$ step $6\to7$ & $+$ geometric tail\\\midrule
$C_5(1)$, difference in errors & $-4.15$ & $-1.88$ & $-0.30$\\
$C_6(1)$, difference in errors & $+2.46$ & $+0.97$ & $-0.22$\\
$C_7(1)$, value & $-322549.9$ & $-322539.5$ & $-322523.9$\\
\bottomrule\end{tabular}\end{center}
The tail ratios are $0.408$ for $C_5$, $0.443$ for $C_6$ and $0.600$ for
$C_7$. The lower-order agreement is consistency with the tail model at
fifth and sixth order. It does not validate the seventh-order tail, which
converges more slowly, and it does not exclude other systematic errors.
This is the origin of the truncation-aware estimate
\eqref{gff:c7f:truncationAware}. Deterministic low-loop assembly also gives
$C_3=-3.71985632290349$ and $C_4=45.4248801371166$; their discrepancies
from the retained higher-precision values are respectively
$1.5\times10^{-12}$ and $3.4\times10^{-8}$.

The lower-order reference coefficients are used only as checks. They are
not fitted into the estimator; in particular, $C_6(1)$ is not an input to
the mean in \eqref{gff:c7f:affine}. The one- through four-loop inputs are
evaluated by deterministic coordinate quadrature. Changing the four-loop
degree from 16 to 20 moves the seventh-order assembly by less than $0.001$
with the higher-loop inputs held fixed; the preceding degree changes
decrease rapidly. This is a convergence observation, not an interval
bound.

The tree point evaluator was compared with an independent positive-gap
reconstruction at 650 decimal digits in 60 cases, including cluster ratios
exponentially smaller than the production minima. The maximum normalized
discrepancy was $1.18\times10^{-14}$; the extreme-endpoint subset had
maximum $3.45\times10^{-18}$. Continuous integration uses extended
precision, binary128 on strongly canceling samples, and a 100-digit
fallback for more severe cancellations. The 650-digit comparison tests the
values that the evaluator actually returns, including the chosen fallbacks
(extended precision in 51 of the 60 cases, binary128 in one and 100 digits in
eight), rather than only the uncancelled summands
(\path{LRI_C7_HPC/reference/verification/tree_point_checks.json}). The
meromorphic weights are generated at 100 working digits and stored in
extended precision. Independent monomial-moment tests of the stored weights
include the Laurent-pole and negative-power cases; the largest normalized
error in 1152 tests is $1.1\times10^{-11}$
(\path{LRI_C7_HPC/reference/verification/meromorphic_weight_checks.json}).
The production implementation evaluates the radial tensor quadrature in
bounded tiles with compensated accumulation. Replayed against binary128
on 112 degree-seven elements, its aggregate projected error in $C_7$ is
$2.8\times10^{-12}$, compared with $1.9\times10^{-11}$ for the original
dense implementation. The worst individual elements, which cancel strongly,
reach $4.2\times10^{-14}$ relative. These errors are many orders of
magnitude below the sampling error.

The two point streams are randomized differently. In the continuous
power-free integration every graph stratum receives a power-of-two number
of points of one linearly scrambled Sobol net, and the uncertainty is
estimated between independent randomized replicates, not between the
correlated points of one net \cite{gff:PrecisionSobol}. The export of the
direction matrices first reconstructs the sample values of the public SciPy
generator and checks them exactly. The power-sector quadrature instead draws
its finite tree sectors from independent 64-bit Mersenne Twister streams,
seeded per batch and per graph; the paired degrees of \eqref{gff:c7f:levels}
evaluate the same sampled trees, which is what makes the paired refinements
small in variance. The meromorphic weight files are native extended-precision
binaries; on a platform with a different extended-precision format they are
regenerated from the shipped code, not read. For the first, smaller
calculation, a direct fixed-point assembly of the raw Taylor coefficients,
followed by the recorded zero-pole control, agreed with the affine
assembly \eqref{gff:c7f:affine} to $1.8\times10^{-10}$.

\section{Full seventh order in two dimensions: mixed local-mass forests}\label{gff:app:C7d2}

\subsection{Result and scope}\label{gff:c7d2:result}
At $d=2$ the seventh-order coefficient uses the same physical regulator
and finite zero-momentum prescription as the sixth-order calculation of
Sections~\ref{gff:cn:finiteKernels}--\ref{gff:app:C6Precision}, with $\epsilon=2s-2$,
\begin{equation}
 \Delta_{\mathcal E}(2,\epsilon)=1-\frac\epsilon2+\gamma_r(g_*,\epsilon),\qquad
 \beta(g,z)=-zg+\sum_{L=1}^{7}\widehat b_L(z)g^{L+1},\qquad
 \gamma_r(g,z)=\sum_{L=1}^{7}\widehat c_L(z)g^L,
 \label{gff:c7d2:rg}
\end{equation}
$\widehat b_1=3$, $\widehat c_1=1$ and $\beta(g_*,\epsilon)=0$. Every loop
order keeps its Taylor polynomial through $z^{7-L}$ before $z=\epsilon$ is
set. Unlike the sixth-order calculation, the power-containing sector is
integrated explicitly instead of being taken from an analytic compensation.
Write $f_L=(b_L^{\rm pf},c_L^{\rm pf})$ and $p_L=(b_L^{\rm p},c_L^{\rm p})$
for the power-free and power-containing jets in the same prescription, and
$\mathcal F$ for the fixed-point series extraction. Then
\begin{equation}\boxed{
 C_7(2)=C_7^{\rm pf}(2)+\Delta_pC_7(2),\qquad
 \Delta_pC_7=\mathcal F_7(f+p)-\mathcal F_7(f).
}\label{gff:c7d2:difference}\end{equation}
The difference includes lower-loop nonlinear changes; it is not the
simple-pole sum of the seven-loop power graphs. The production of
Section~\ref{gff:c7d2:production} gives
\begin{center}\begin{tabular}{@{}lrr@{}}\toprule
Contribution & Estimate & Sampling standard error\\\midrule
Power-free coefficient $C_7^{\rm pf}(2)$ & $-3876.0700590$ & $0.5224502$\\
Power correction $\Delta_pC_7(2)$ & $-0.0560338$ & $0.0458637$\\
Complete coefficient $C_7(2)$ & $-3876.1260928$ & $0.5244497$\\\bottomrule
\end{tabular}\end{center}
The quoted uncertainty is the empirical reporting allowance
\begin{equation}
 E_{\rm emp}=\max\Bigl(4\widehat\sigma,\ \max_j\bigl|C_{7,j}^{\rm half}-C_7\bigr|\Bigr),
 \quad\text{rounded up},
 \label{gff:c7d2:allowance}
\end{equation}
where $\widehat\sigma$ is the sampling standard error and $C_{7,j}^{\rm half}$,
$j=1,2$, are the half-stream estimates: the estimator recomputed from the first
and from the second half of the replicates of every campaign, loop and sector.
The rule is that of an earlier, lower-precision integration of $C_7(2)$ with the
same graph set and prescription, called the earlier $C_7(2)$ calculation, or
the earlier calculation for short, in this section. The rule was applied
to the 2026 production, and its rounding step chosen, after the result was
available. Here $4\widehat\sigma=2.098$ exceeds both half-stream deviations,
$0.36$, and rounding up to a multiple of $0.1$ gives
\begin{equation}\boxed{C_7(2)=-3876.1\pm2.1,\qquad\epsilon=2s-2.}
\label{gff:c7d2:reported}\end{equation}
The earlier calculation rounded the same rule to a multiple of 10 at its
precision and reported $-3950\pm310$ (estimator $-3953.6336215$, sampling
standard error $76.7310362$); the new estimator lies 1.01 of its standard
errors from it. The allowance is empirical. It is not a probability
guarantee and does not bound an error shared by the graph and
renormalization framework; the additional estimator digits are given for
reproduction only.

The earlier calculation, its method and its frozen design were written by
OpenAI GPT-6 astra (model stated by the author); its code and records are in
\path{LRI_C7_D2_HPC/reference/} (Sections~\ref{gff:mg7:design}
and~\ref{gff:cn:reproduction}). When this method was merged into this section, on
16 October 2026, a Claude Code session wrote out or rechecked four parts: the
three-loop check \eqref{gff:c7d2:threeK}--\eqref{gff:c7d2:threeDerivative}, for which it
wrote out the graph weight, the forest cancellation \eqref{gff:mg7:bubblePair}, the
bound $\mathsf p_3/\mathsf p_2\le1/9$ and the series identity used for the
regulator derivative, and checked the quadrature and the Euler and Pfaff steps
numerically; the two-loop entries \eqref{gff:mg7:twoLoop}--\eqref{gff:mg7:b25} and the
polynomial \eqref{gff:mg7:MS7}, which it rechecked by a separate Cauchy evaluation
and an exact symbolic expansion; the primitive classes of
Section~\ref{gff:mg7:primitive}, for which it supplied the argument for $w_B=0$ and
the normalization \eqref{gff:mg7:primNorm} and compared the 40-digit values with the
closed forms; and the variance decomposition \eqref{gff:mg7:covDecomposition}, which
it evaluated with the 2026 record. It also checked that the group
configurations of the 2026 production hold the sample counts of
Section~\ref{gff:mg7:design}.

\subsection{Graph weights and complete coverage}
An $L$-loop graph has $L+1$ vertices and $2L$ individually labeled
lines. With edge multiplicities $m_{ij}$ and $e_i=4-\deg(i)$,
\begin{equation}
 w_A=\frac{24}{|\operatorname{Aut}G|\prod_{i<j}m_{ij}!\prod_ie_i!},\qquad
 s_G=\sum_ie_i(e_i-1),\qquad
 w_B=\frac{s_G}{12}w_A,
 \label{gff:c7d2:weights}
\end{equation}
so that $3w_B-w_A=(s_G-4)w_A/4$. This rational weight, not a sampled
value, determines the zero-channel exclusions at seven loops; at lower
loops the coupling and mass jets are kept separately even when the
final-order channel weight vanishes. A proper one-particle-irreducible
subgraph with $v$ vertices and $e$ internal lines is logarithmic if
$e=2v-2$ and an unmarked power-divergent two-point subgraph if $e=2v-1$.
Its loop number is $h=e-v+1$, so a power subgraph has $h=v$. All proper
edge subgraphs, including those not induced by their vertex set, were
enumerated independently; their mixed forest maps agree exactly with the
production maps for all 1,301 power-containing classes at loops three to
seven.
\begin{center}\begin{tabular}{@{}rrrrr@{}}\toprule
$L$ & All classes & Power-containing & Nonzero mixed kernels & Integrated power classes\\\midrule
3&8&1&1&1\\4&27&5&3&3\\5&129&31&20&20\\
6&660&173&107&107\\7&3986&1091&709&590\\\bottomrule
\end{tabular}\end{center}
At seven loops 459 classes have $3w_B-w_A=0$ and 639 further classes have
an exactly vanishing mixed forest map. The remaining 2,254 power-free
nonprimitive, 44 primitive and 590 power-containing classes are all
integrated in every replicate of the production below; there is no omitted
nonzero class. At loops three to six the integrated power-free and
power-containing classes number 5, 16, 76, 395 and 1, 3, 20, 107.

A graph is power-free if it has no unmarked power-divergent proper
two-point subgraph; nonprimitive graphs whose subdivergences are all
logarithmic are power-free. The power-free sector
partitions as
\begin{center}\begin{tabular}{@{}rrrrr@{}}\toprule
$L$ & Power-free classes & Nonzero finite kernels & Primitive among them & Zero kernels\\\midrule
2&2&1&0&1\\3&7&5&1&2\\4&22&16&1&6\\
5&98&76&3&22\\6&487&395&10&92\\7&2895&2441&44&454\\\bottomrule
\end{tabular}\end{center}
and at seven loops the 3,986 classes split by sector, kernel and channel
weight as
\begin{center}\begin{tabular}{@{}lrrr@{}}\toprule
Seven-loop classes & $3w_B-w_A\ne0$ & $3w_B-w_A=0$ & All\\\midrule
Power-free, nonzero kernel & 2298 & 143 & 2441\\
Power-free, zero kernel & 384 & 70 & 454\\
Power-containing, nonzero kernel & 590 & 119 & 709\\
Power-containing, zero kernel & 255 & 127 & 382\\\midrule
All & 3527 & 459 & 3986\\\bottomrule
\end{tabular}\end{center}
The first column contains the integrated classes, $2{,}298=2{,}254+44$
power-free and 590 power-containing, and the $639=384+255$ classes with a
vanishing mixed forest map; the 845 power-containing classes of nonzero
channel weight are $590+255$. The class-by-class assignment is in
\path{LRI_C7_D2_HPC/reference/verification/full_graph_coverage.json}. The
$262=143+119$ classes with a nonzero kernel and zero channel weight are not
integrated here; they first contribute at eighth order
(Section~\ref{gff:c8d2:changes}).

The production builders take their subgraphs from induced vertex sets. The
independent enumeration of all edge subgraphs also contains the
non-induced ones, for example a two-line logarithmic bubble inside a
triple-line power insertion, and the exact agreement of the forest maps for
all 1,301 power-containing classes shows that the forest terms of such
subgraphs cancel and were not dropped by assumption; the three-loop
example \eqref{gff:mg7:bubblePair} below shows the cancellation explicitly, and
Section~\ref{gff:c8d3:independent} proves it for every graph of the family.
After graph and line labels are matched, all 492 nonzero power-free kernels
at loops three to six are identical to the kernels of the sixth-order
calculation (Section~\ref{gff:cn:finiteAssembly}; record
\path{LRI_C7_D2_HPC/reference/LRI_C7_d2_progress/verification/legacy_mapping.json}).
A second construction from contracted graph Laplacians and exact rational
Kirchhoff determinants, which does not use the monomial selection,
reproduces 24 selected seven-loop power-free forest kernels exactly; the
analogous construction for the power-containing kernels is the check at 66
configurations quoted below.


\subsection{Local mass subtraction and the mixed forest kernel}
The Schwinger infrared factor does not perform critical mass tuning: an
unmarked self-energy insertion must have its local zero-momentum part
subtracted, as for the three-loop graph $I_3$ in
\cite[Appendix C.3.4]{gff:BGHS}. In two dimensions its superficial degree at
$z=0$ is one in momenta; the amplitude is even in the external momentum, so
the constant Taylor subtraction leaves a remainder of quadratic order. Put
\begin{equation}
 U_G(x)=\sum_T\prod_{e\notin T}x_e,\qquad S(x)=\sum_ex_e,\qquad
 d\Omega=\delta(1-S(x))\prod_edx_e,
\end{equation}
the sum running over spanning trees. A compatible mixed forest
$F=F_{\log}\cup F_{\rm p}$ consists of nested subgraphs, or edge-disjoint
subgraphs meeting in at most one vertex; nested blocks are taken after
contracting their children, and $U_F$ is the product of the block Symanzik
polynomials. Equivalently, a complement-of-tree monomial $m$ of $U_G$ is
kept in $U_F$ if and only if $|m\cap E_\gamma|=h_\gamma$ for every
$\gamma\in F$. With the logarithmic-root edge set
$E_R^{\log}(F)=E_G\setminus\bigcup_{\gamma\in F_{\log}}E_\gamma$, the
finite RG kernel is
\begin{equation}\boxed{
 K_G^{\rm full}(x)=\sum_F(-1)^{|F|}
 \frac{\sum_{e\in E_R^{\log}(F)}x_e}{S(x)\,U_F(x)}.
}\label{gff:c7d2:kernel}\end{equation}
Only logarithmic subgraphs are removed from the root numerator. A power
subgraph changes the denominator through its local Taylor subtraction but
does not remove its interaction order from the root amplitude; excluding
its lines from the numerator would implement a different scale derivative
and already changes the three-loop answer. To derive
\eqref{gff:c7d2:kernel}, apply $1-T^0_\gamma$, with $T^0_\gamma$ setting the
external momentum to zero, to every compatible local two-point insertion.
Gaussian integration replaces $U_G$ by $U_F$ for the selected power
subgraphs, while the vertex interaction order and the common Schwinger mass
factor are unchanged. The finite coupling and mass inversion of these
critically subtracted amplitudes then gives the same logarithmic root
derivative as in the power-free calculation, distributed over the
projective block scales as the root fraction
$\sum_{e\in E_R^{\log}}x_e/S$. Expanding both operations gives one sign
for each selected subgraph. This is a local implementation of parametric
forest renormalization \cite{gff:BrownKreimerLocal}. The scale exponents of the
blocks are derived in the next paragraph.
Identical selected denominators are grouped with integer numerator
coefficients and identically zero vectors are removed before any
floating-point arithmetic. An independent rational Kirchhoff-determinant
implementation agrees exactly at 66 configurations, including
endpoint-scaled points.

\paragraph{Scale exponents of the forest blocks.} This bookkeeping fixes the
numerator of \eqref{gff:c7d2:kernel}. Rescale the lines of a block $B$ of loop number $h_B$
by a common factor $\lambda$. The critical line measure
$\prod_ex_e^{-1/2+z/4}dx_e$ and the block Symanzik polynomial, of degree
$h_B$, give $\lambda^{|E_B|(1/2+z/4)-h_B}$. A logarithmic block has
$|E_B|=2h_B$ and exponent $h_Bz/2$, the Dirichlet parameter of its relative
scale in the proof of \eqref{gff:done:commonIntegral}; this produces the root
fraction and the sign of each logarithmic forest element. A power block
has $|E_B|=2h_B-1$ and exponent $-1/2+(2h_B-1)z/4$ before its Taylor
remainder. Its complement, with $2L-2h_B+1$ lines and loop number $L-h_B$,
has exponent $+1/2+(2L-2h_B+1)z/4$, the opposite engineering dimension, and
the two add to $Lz/2$, the exponent of the whole graph. A mass Taylor
operation therefore splits off no logarithmic scale: it carries no further
factor $h_B/L$ and introduces no second coupling-inversion root factor.
After the logarithmic inversion the root keeps its original lines,
including those of the power Taylor factors, and these form the edge set
$E_R^{\log}(F)$ of \eqref{gff:c7d2:kernel}. Nested blocks obey the same
argument after their children are contracted. In the monomial selection,
$m$ is the complement-of-tree mask of the individually labeled lines, so
parallel lines are separate variables.


In the normalization $r(2,z)=1$ the power-containing jets are
\begin{equation}\boxed{
 \begin{pmatrix}\widehat b_L^{\rm p}(z)\\ \widehat c_L^{\rm p}(z)\end{pmatrix}
 =\frac{2(-1)^{L+1}\Gamma(1+Lz/2)}{\Gamma(\tfrac12+z/4)^{2L}}
 \sum_{G\in{\rm p}}\begin{pmatrix}w_A(G)\\ w_B(G)\end{pmatrix}
 \int d\Omega\prod_ex_e^{-1/2+z/4}K_G^{\rm full}(x),
}\label{gff:c7d2:normalization}\end{equation}
with the same prefactor and finite coupling prescription as the
power-free data. Before local subtraction a connected power face has radial
exponent $e/2-h=-1/2$ at $z=0$; the Taylor remainder raises it to $+1/2$.
Logarithmic faces start at zero and their combined subtraction raises the
exponent; nested limits are handled by the compatible forest. Regulator
derivatives add finite powers of logarithms without removing these
margins. The importance maps use the forest denominator valuations and the
positive half-integer endpoint floor after subtraction.

Equation~\eqref{gff:c7d2:normalization} is \eqref{gff:done:forestRG} with $D_G$ of
\eqref{gff:done:commonIntegral} and $K_G$ replaced by $K_G^{\rm full}$, since
$Lz\,\Gamma(Lz/2)=2\Gamma(1+Lz/2)$: the prefactor is the logarithmic scale
derivative of $\Gamma(Lz/2)$, and at $z=0$ it equals $2(-1)^{L+1}\pi^{-L}$.
The power-free jets $f_L$ are given by the same formula with $K_G$ of
\eqref{gff:done:commonKernel}, so no regulator-independent MS coefficient
enters the input jets of \eqref{gff:c7d2:difference}.

\paragraph{Importance rates.} Let
$\nu_S(P)$ be the smallest degree of a polynomial $P$ in the variables of a
proper line subset $S$, and $N_F=\sum_{e\in E_R^{\log}(F)}x_e$. When the
lines of $S$ shrink by a common factor $\lambda$, the forest term $F$ of
\eqref{gff:c7d2:kernel}, times the line measure at $z=0$, scales as
$\lambda^{|S|/2+\nu_S(N_F)-\nu_S(U_F)}$. The builders store the rate
\begin{equation}
 \alpha_G(S)=\tfrac12\max\bigl(r_G(S),p_S\bigr),\qquad
 r_G(S)=\min_F\bigl[\,|S|+2\nu_S(N_F)-2\nu_S(U_F)\bigr],
 \label{gff:mg7:rates}
\end{equation}
with $p_S=1$ for odd and $p_S=2$ for even $|S|$. The termwise bound
$r_G(S)/2$ can be nonpositive where the forest sum cancels. Since the
kernel is rational and the measure contributes $|S|/2$, the exponent of the
summed kernel lies in $|S|/2+\mathbb Z$, and its positivity after
subtraction implies the floor $p_S/2$. The floor therefore sharpens a
nonpositive termwise bound without changing the integral;
Section~\ref{gff:c8d2:tails} uses the same rule at eighth order and shows that
the weight then stays bounded. At seven loops the floor strengthened the
rates of 81,632 even subsets in the power-free files (record
\path{LRI_C7_D2_HPC/reference/LRI_C7_d2_progress/verification/parity_rate_update.json});
at loops three to six the power-free rates are the frozen tables of the
sixth-order calculation, mapped exactly to the new labels. The campaign
factors $0.8$ and $0.9$ multiply these rates. All such choices change the
variance, not the integral, and no interval error is derived from these
power-counting arguments.

\paragraph{Routes not used.} Subtracting the local part of a self-energy
insertion is not the same as discarding the graph, and the absence of
wave-function renormalization at this order is not used to delete power
amplitudes: in the long-range prescription the local subtraction of a power
divergence leaves a finite two-point amplitude \cite{gff:BGHS}. Two other
routes to the power sector were tried. A coordinate-space routing audit
through the Gamma convolutions, complex Selberg specializations and
five-edge kite formulas of Section~\ref{gff:app:C6d2Complete}, with all
required logarithmic subgraphs and contraction quotients, closed
structurally for 740 of the 845 power-containing seven-loop classes of
nonzero channel weight and did not close for the other 105. It was an
audit of routes, not an evaluation, and no class was assigned zero. An
exploratory radial integration by parts did not implement the local
subtraction: it recovered a constant term, but its regulator derivatives
still produced divergent endpoint integrals. It was rejected, and none of
its samples enters a result. The kernel \eqref{gff:c7d2:kernel} instead
performs the local subtraction and the logarithmic forest combinations
before any sampling.


\subsection{An independent three-loop check}
The first power graph $G$ has three external-chain parameters with sum $R$
and internal parameters $y_1,y_2,y_3$. With $\mathsf P_2=y_1y_2+y_1y_3+y_2y_3$,
$\mathsf P_3=y_1y_2y_3$, $U_G=R\mathsf P_2+\mathsf P_3$ and $U_0=R\mathsf P_2$, its mixed kernel \eqref{gff:c7d2:kernel} is exactly
\begin{equation}
 K_G^{\rm full}=\frac1{U_G}-\frac1{U_0}=-\frac{\mathsf P_3}{R\mathsf P_2(R\mathsf P_2+\mathsf P_3)},
 \label{gff:c7d2:threeK}
\end{equation}
a rational remainder that can be integrated and differentiated at $z=0$.
Let $a=\tfrac12+z/4$, $k=3a-1$, and use simplex variables $u$ with
$\mathsf p_2=u_1u_2+u_1u_3+u_2u_3$ and $\mathsf p_3=u_1u_2u_3$. Integrating the external
parameters and the radial variable, and applying a Pfaff transformation,
gives
\begin{equation}
 \widehat b_3^{\rm p}(z)=
 -\frac{\Gamma(1+3z/2)\,\mathrm B(k,k)}{\Gamma(a)^3\Gamma(3a)}
 \int d\Omega_u\,\mathsf p_3^{z}\mathsf p_2^{-3a}\,
 {}_2F_1\!\left(\frac{3z}2,k;2k;1-\frac{\mathsf p_3}{\mathsf p_2}\right),
 \label{gff:c7d2:threeHyper}
\end{equation}
with $w_A=1/2$ included. At $z=0$ the hypergeometric function is one and
$\int d\Omega_u\,\mathsf p_2^{-3/2}=2\pi$, so $\widehat b_3^{\rm p}(0)=-4$ and
$\widehat c_3^{\rm p}=\widehat b_3^{\rm p}/3$, in agreement with the
subtracted $I_3$ normalization of \cite{gff:BGHS}. The first regulator
coefficient is
\begin{equation}
 [z]\widehat b_3^{\rm p}(z)=-\frac2\pi\int d\Omega_u\,\mathsf p_2^{-3/2}
 \left[-\frac32+\log \mathsf p_3-\frac34\log \mathsf p_2+
 3\log\frac2{1+\sqrt{\mathsf p_3/\mathsf p_2}}\right]=18.6743207750,
 \label{gff:c7d2:threeDerivative}
\end{equation}
from two independently tightened deterministic quadratures differing by
$4.6\times10^{-10}$; this reference uses neither the graph sampler nor its
monomial arrays. The production of Section~\ref{gff:c7d2:production} gives
$\widehat b_3^{\rm p}(0)=-3.99999979\pm0.00000042$ and
$[z]\widehat b_3^{\rm p}=18.6743200\pm0.0000021$, at $+0.50$ and $-0.36$
sampling standard errors from these values.

\paragraph{Derivation of \eqref{gff:c7d2:threeK}--\eqref{gff:c7d2:threeDerivative}.}
The polynomial
$U_G=R\mathsf P_2+\mathsf P_3$ fixes the topology of the single three-loop power-containing
class. Four vertices lie on a ring. One adjacent pair is joined by the
three lines $y_1,y_2,y_3$, the sunset $\gamma$, a power subgraph with
$v=2$, $e=3$ and $h=2$. The other three lines are in series, so their
parameters enter $U_G$ only through their sum $R$, and each of their two
inner vertices carries two external legs. Hence $|\Aut G|=2$, the triple
line contributes $3!$ and the legs $2!\,2!$, so $w_A=24/48=1/2$, and
$s_G=4$ gives $w_B=w_A/3$, which is the statement
$\widehat c_3^{\rm p}=\widehat b_3^{\rm p}/3$.

\emph{The kernel.} The only power subgraph is $\gamma$. The logarithmic
subgraphs are the three bubbles $\beta_{ij}=\{y_i,y_j\}$ inside $\gamma$,
which are not induced. Two bubbles share a line and are not compatible, so
the compatible forests are $\varnothing$, $\{\gamma\}$, $\{\beta_{ij}\}$
and $\{\beta_{ij},\gamma\}$. For $\varnothing$ and $\{\gamma\}$ the
numerator is $S$, since a power subgraph removes no line, and the terms
are $1/U_G$ and $-1/U_0$ with $U_0=U_\gamma U_{G/\gamma}=\mathsf P_2R$. For a bubble,
$U_{\beta_{ij}}=y_i+y_j$, and $G/\beta_{ij}$ is the one-vertex join of the
loop $y_k$ and the contracted chain, $\{i,j,k\}=\{1,2,3\}$. Since
$\gamma/\beta_{ij}$ is the loop $y_k$ and $G/\gamma$ the chain,
\begin{equation}
 U_{\{\beta_{ij}\}}=(y_i+y_j)\,y_kR=U_{\{\beta_{ij},\gamma\}},
 \label{gff:mg7:bubblePair}
\end{equation}
which are also the monomials of $U_G$ with exactly one of $y_i,y_j$. Both
forests have the numerator $R+y_k$ and they carry opposite signs, so they
cancel, and \eqref{gff:c7d2:threeK} remains.

\emph{The integral.} With $a=\tfrac12+z/4$ the line measure is
$\prod_ex_e^{a-1}dx_e$. Integrating the chain parameters at fixed $R$
gives $R^{3a-1}\Gamma(a)^3/\Gamma(3a)$. Put $R=r$ and $y_i=(1-r)u_i$ with
$\sum_iu_i=1$. Then $\mathsf P_2=(1-r)^2\mathsf p_2$, $\mathsf P_3=(1-r)^3\mathsf p_3$, the Jacobian is
$(1-r)^2$, $\prod_iy_i^{a-1}=(1-r)^{3a-3}\prod_iu_i^{a-1}$, and
$K_G^{\rm full}=-\mathsf p_3/[r(1-r)\mathsf p_2(r\mathsf p_2+(1-r)\mathsf p_3)]$. With $k=3a-1$ the $r$ integral is
Euler's,
\begin{equation}
 -\frac{\mathsf p_3}{\mathsf p_2}\int_0^1\frac{r^{k-1}(1-r)^{k-1}}{r\mathsf p_2+(1-r)\mathsf p_3}\,dr
 =-\frac{\mathrm B(k,k)}{\mathsf p_2}\,
 {}_2F_1\!\left(1,k;2k;1-\frac{\mathsf p_2}{\mathsf p_3}\right).
 \label{gff:mg7:threeEuler}
\end{equation}
Its argument is at most $-8$, because
$\mathsf p_3/\mathsf p_2=\bigl(\sum_iu_i^{-1}\bigr)^{-1}\le1/9$. The Pfaff transformation
${}_2F_1(1,k;2k;w)=(1-w)^{-k}{}_2F_1(2k-1,k;2k;w/(w-1))$, with
$1-w=\mathsf p_2/\mathsf p_3$, $w/(w-1)=1-\mathsf p_3/\mathsf p_2\in[8/9,1)$ and $2k-1=3z/2$, turns the
$u$ integrand into
$-\mathrm B(k,k)\,\mathsf p_2^{-1-k}\mathsf p_3^{k}\prod_iu_i^{a-1}\,
{}_2F_1(3z/2,k;2k;1-\mathsf p_3/\mathsf p_2)$. This equals
$-\mathrm B(k,k)\,\mathsf p_3^{z}\mathsf p_2^{-3a}\,{}_2F_1(3z/2,k;2k;1-\mathsf p_3/\mathsf p_2)$ because
$\prod_iu_i=\mathsf p_3$ and $4a-2=z$. With the prefactor
$2\Gamma(1+3z/2)/\Gamma(a)^6$ of \eqref{gff:c7d2:normalization} at $L=3$, the
weight $w_A=1/2$ and the chain factor $\Gamma(a)^3/\Gamma(3a)$, this is
\eqref{gff:c7d2:threeHyper}.

\emph{The two coefficients.} At $z=0$ one has $a=k=\tfrac12$,
$\mathrm B(\tfrac12,\tfrac12)=\pi$ and
$\Gamma(\tfrac12)^3\Gamma(\tfrac32)=\pi^2/2$, so the scalar prefactor is
$-2/\pi$, and $\int d\Omega_u\,\mathsf p_2^{-3/2}=2\pi$ gives
$\widehat b_3^{\rm p}(0)=-4$. The logarithmic derivative of the scalar
prefactor at $z=0$ is $\tfrac34[\psi(\tfrac12)-\psi(\tfrac32)]=-\tfrac32$,
and $\mathsf p_3^z\mathsf p_2^{-3a}$ contributes $\log \mathsf p_3-\tfrac34\log \mathsf p_2$. Since
$(3z/2)_n=O(z)$ for $n\ge1$, the $z$ dependence of $k$ enters the
hypergeometric function only at order $z^2$, and for $0\le\xi<1$
\[
 {}_2F_1\Bigl(\frac{3z}2,\frac12;1;\xi\Bigr)
 =1+\frac{3z}2\sum_{n\ge1}\frac{(\frac12)_n}{n!\,n}\,\xi^n+O(z^2)
 =1+3z\log\frac2{1+\sqrt{1-\xi}}+O(z^2),
\]
because $\sum_{n\ge1}(\frac12)_n\xi^n/(n!\,n)
=\int_0^\xi[(1-t)^{-1/2}-1]\,dt/t=2\log[2/(1+\sqrt{1-\xi})]$. At
$\xi=1-\mathsf p_3/\mathsf p_2$ the three contributions give \eqref{gff:c7d2:threeDerivative}.

\emph{Quadrature.} The integrand is symmetric in $u$. On the sector
$u_1\ge u_2\ge u_3$ put $u=(1,t,tv)/(1+t+tv)$ with $t,v\in(0,1)$, and
$t=\rho^2$. Then $d\Omega_u\,\mathsf p_2^{-3/2}=2(1+v+\rho^2v)^{-3/2}d\rho\,dv$, so
the six sectors give the weight $12(1+v+\rho^2v)^{-3/2}d\rho\,dv$, free of
the algebraic corner singularity, and the logarithms of
\eqref{gff:c7d2:threeDerivative} remain integrable. The adaptive error
estimates of the two deterministic quadratures are not interval
enclosures. In these variables a separate tanh-sinh quadrature at 20
digits reproduces $2\pi$ and $18.6743207750$, and \eqref{gff:mg7:threeEuler}
and the Pfaff step were confirmed numerically at a sample point.


\subsection{Regulator jets and physical assembly}
On a projective chart put $h(x)=\frac14\sum_e\log x_e-\frac L2\log S(x)$.
The scalar normalization logarithm has Taylor coefficients
\begin{equation}
 a_{L,j}=\frac{(L/2)^j\psi_{j-1}(1)-2L(1/4)^j\psi_{j-1}(\tfrac12)}{j!},
\end{equation}
and the integrand jets follow from $e_0=1$ and
\begin{equation}
 e_n(x)=\frac1n\sum_{j=1}^nj\bigl(a_{L,j}+\delta_{j1}h(x)\bigr)e_{n-j}(x),
 \qquad n\le7-L,
 \label{gff:c7d2:jets}
\end{equation}
so the maximum orders at loops two to seven are $5,4,3,2,1,0$. The
two-loop input is the unchanged exact-hypergeometric finite-scheme input
and the one-loop input is exact. With all jets present,
$g(\epsilon)=\frac13\bigl[\epsilon-\sum_{L=2}^7\widehat b_L(\epsilon)
g(\epsilon)^L\bigr]\bmod\epsilon^8$ is solved by eight iterations from
$g=0$ and substituted in \eqref{gff:c7d2:rg}. An independent bivariate
Lagrange extraction agrees exactly in 24 regulator-dependent rational
cases, and 24 further cases agree with the constant-coefficient MS
seventh-order polynomial; these test the assembly, not the sampled data.
No higher-order coefficient is extrapolated from $C_0,\ldots,C_6$.

By \eqref{gff:done:twoLoopMOM} and \eqref{gff:gd:twoLoopHyper} at $d=2$,
where $r(2,z)=1$ and $\Gamma(d/2)=1$,
\begin{equation}
 [z^j]\,\widehat b_2(z)=-6\,[z^{j+1}]\bigl(4F_1(2,z)-2W(2,z)F_2(2,z)\bigr),
 \qquad \widehat c_2=\tfrac16\widehat b_2,\qquad j=0,\ldots,5.
 \label{gff:mg7:twoLoop}
\end{equation}
Seventh order needs, beyond \eqref{gff:done:twoLoopNumbers},
\begin{equation}
 b_{2,5}=175.39598710648998744915099527723636\ldots .
 \label{gff:mg7:b25}
\end{equation}
Two Cauchy evaluations at 75 working digits, one with radius $0.05$ and 80
nodes and one with radius $0.08$ and 96 nodes, agree to
$1.75\times10^{-68}$ over the six entries
(\path{LRI_C7_D2_HPC/reference/LRI_C7_d2_progress/verification/two_loop_extension.json};
the stored input is \path{inputs/two_loop.json} of the same folder). This
is a numerical consistency check, not an interval bound. A separate
60-digit Cauchy evaluation, with radii $0.05$ and $0.08$ and 48 and 64
nodes, reproduces $b_{2,0},\ldots,b_{2,6}$, including the eighth-order
entry of Section~\ref{gff:c8d2:changes}.

For comparison only: in a regulator-independent MS scheme with
$b_2=-6A$, $c_2=-A$ and $A=A_2=4\log2$ \eqref{gff:gd:blocks}, the Lagrange
formula \eqref{gff:c6:Lagrange} at $n=7$ gives the seventh-order counterpart of
\eqref{gff:c6:RGassembly},
\begin{equation}
\begin{aligned}
 6561\,C_7={}&3c_7-b_7+12672A^6-3360A^4b_3+4320A^4c_3
 -672A^3b_4+1152A^3c_4\\
 &+192A^2b_3^2-432A^2b_3c_3-112A^2b_5+240A^2c_5
 +\frac{128}3Ab_3b_4-64Ab_3c_4\\
 &-48Ab_4c_3-14Ab_6+36Ac_6-\frac43b_3^3+4b_3^2c_3
 +\frac83b_3b_5-5b_3c_5\\
 &+\frac43b_4^2-4b_4c_4-3b_5c_3.
\end{aligned}\label{gff:mg7:MS7}
\end{equation}
It is the polynomial of the 24 constant-coefficient test cases above. An
exact symbolic expansion of \eqref{gff:c6:Lagrange} reproduces all 23
monomials of \eqref{gff:mg7:MS7} and, at $n=6$, \eqref{gff:c6:RGassembly}. The
polynomial must not be applied to the constant terms
$\widehat b_{L,0},\widehat c_{L,0}$ of the finite-scheme jets alone, which
define a different truncated problem; the finite-scheme coefficient is
extracted from the full jets, and independently by
\eqref{gff:c7f:finiteLagrange}, whose form does not depend on $d$. In both
schemes the seven-loop constants enter linearly, with
$\partial C_7/\partial(b_7,c_7)=(-1,3)/6561$, so that
$C_7(2)\supset(3\widehat c_{7,0}-\widehat b_{7,0})/6561$ and the classes
with $3w_B=w_A$ drop out at this order.


\subsection{Primitive classes: completed periods and three products}\label{gff:mg7:primitive}
A primitive class has $w_B=0$: a vertex with $e_i\ge2$ would
leave the other $L$ vertices with at least $2L-2$ lines, a logarithmic
subgraph, so each of four vertices carries one external leg and $s_G=0$.
For a primitive graph the forest sum \eqref{gff:done:forestD} has only its root
term, so by \eqref{gff:done:forestRG} and \eqref{gff:c6d2p:periodNorm}, with
$a_G=(2/7)\Pi_G$,
\begin{equation}
 C_{7,\rm prim}(2)=\frac{(3c_7-b_7)_{\rm prim}}{6561}
 =-\frac2{6561}\sum_{G\ \rm primitive}w_A(G)\,\Pi_G .
 \label{gff:mg7:primNorm}
\end{equation}
Completing each of the 44 classes to a four-regular graph on nine vertices
partitions them into 14 isomorphism classes, with periods
$\Pi^{(7)}_0,\ldots,\Pi^{(7)}_{13}$; the equality of periods within a class
is the completion invariance of Section~\ref{gff:cn:finiteKernels}. In the
zero-based order of the completion dictionary
\path{LRI_C7_D2_HPC/reference/LRI_C7_d2_progress/reference/primitive_completion.json},
the summed weights $\bar w_j=\sum w_A$ are
\[
 (\bar w_0,\ldots,\bar w_{13})
 =\Bigl(27,\tfrac{27}2,27,18,3,108,18,\tfrac{27}4,108,12,108,54,54,12\Bigr).
\]
Three completed graphs have three-vertex separators: $\Pi^{(7)}_6$ has one,
splitting it into two octahedra (four-loop factors); $\Pi^{(7)}_{12}$ has
one, splitting it into $K_5$ and the completed five-loop primitive; and
$\Pi^{(7)}_7$ has two, splitting it into three copies of $K_5$. The
factorization behind \eqref{gff:c6d2p:product}, applied once or twice, gives
\begin{equation}
 \Pi^{(7)}_6=P_4(2)^2,\qquad
 \Pi^{(7)}_7=P_3(2)^3=27Q_2^3,\qquad
 \Pi^{(7)}_{12}=P_3(2)P_5(2)=3Q_2P_5(2),
 \label{gff:mg7:products}
\end{equation}
and therefore
\begin{equation}\boxed{
 C_{7,\rm prim}(2)=-\frac2{6561}\Bigl[18P_4(2)^2+\frac{729}4Q_2^3
 +162\,Q_2P_5(2)+\sum_{j\notin\{6,7,12\}}\bar w_j\Pi^{(7)}_j\Bigr].
}\label{gff:mg7:primitive7}\end{equation}
The remaining sum has eleven periods. The representatives, the weights and
the map from the 44 classes to the 14 periods are in the completion
dictionary and in \path{verification/primitive_completion_checks.json} of
the same folder. At 40 working digits
\[\begin{aligned}
 Q_2&=12.7713811986469576577497636845240084\ldots,\\
 P_4(2)&=250.487315581867285531851756887911407\ldots,\\
 P_5(2)&=1548.08264023312061513243907394578934\ldots,
\end{aligned}\]
and the three product terms of \eqref{gff:mg7:primitive7} contribute
$-1436.35574062253405165231921281174843\ldots$ to $C_7(2)$. These values of
$Q_2$ and $P_4(2)$ agree with \eqref{gff:gc:Q2number} and with the fifty-digit
$P_4(2)$ of Section~\ref{gff:gc:d2closedForms}, and solving
\eqref{gff:gc:C5twoClosed} for $P_5(2)$ with the 36-digit $C_5(2)$ of
Section~\ref{gff:gdr:checks} reproduces $P_5(2)$ to 34 digits.

The earlier calculation used \eqref{gff:mg7:primitive7} as follows. Two
independently randomized integrations of all 44 classes, each with 16
replicates of 4,096 points per class, measured the three product blocks;
all six comparisons with the exact references lie within 1.88 sampling
standard errors. The measured blocks were then replaced by the references,
which gave the primitive contribution $-8657.4262196\pm1.2165092$ (the
equal-weight mean of the two integrations, with its sampling standard
error). It entered the earlier estimate once, as a retained input. The
2026 production integrates the 44 classes directly in every replicate and
does not use the reduction, so \eqref{gff:mg7:primitive7} remains an exact
identity available as a check. The product references have negligible
numerical error on the Monte Carlo scale; no interval certification is
asserted.


\subsection{The 2026 production}\label{gff:c7d2:production}
The production uses an optimized, restartable reimplementation of the same
calculation. Its kernel keeps the reference addition order and the
extended-precision first pass, with promotions to binary128 and to the
final tier at the unchanged thresholds $10^{-9}$ and $10^{-23}$. The final
tier uses 350-decimal GMP arithmetic instead of the original 300-decimal
tier; this raises the working precision and was fixed before the production.
On the production host both backends reproduce the original executable at
all 3,511 graph kernels with zero normalized difference and identical
promotion counts, and pass the independent high-precision point checks.
Results do not depend on the thread count: 5,776 seven-loop rows run at 64
and at 16 threads were identical, including 31,663 promotions to the GMP
tier.

The original executable of the earlier calculation was itself checked against independently ordered
high-precision monomial sums: 69 point comparisons for power-free and 174
for power-containing kernels, with maximum errors, normalized by one plus
the sum of the absolute forest terms, of $1.64\times10^{-32}$ and
$2.0\times10^{-32}$ in binary128 and of $1.85\times10^{-302}$ and
$1.3\times10^{-302}$ in the 300-decimal tier. It used compensated
accumulation and strict floating-point flags, without fast-math, and shared
monomial products in place of repeated exponentials, as in
\eqref{gff:meth:dag}, which changes the rounding order but neither the
precision nor the expression. These are pointwise arithmetic checks, not
bounds on an integral.


Two independent campaigns $a$ and $b$ have 64 randomized replicates for
every loop and sector, 52- and 60-bit digital coordinates, and importance
rates multiplied by $0.8$ and $0.9$. Every stratum receives its frozen
per-graph allocation from the earlier pilots, at least four points,
multiplied by $2^{12}$. The 44 primitive classes, which the earlier
calculation retained from a separate integration, are integrated with
$2^{26}$ points each per replicate. Counts are powers of two, the base seed
is fresh, and a cost model from an independent pilot affects only the
placement of tasks. The two campaigns contain 1,094,688,178,176 kernel
samples; they ran on the symmetry cluster on 14 AMD and 14 Intel nodes in
about 6.5 hours, and every merge and the assembly used the package's own
tools.

\paragraph{The frozen design.}\label{gff:mg7:design} The per-graph counts were fixed from pilot runs that never
entered a result. Each stratum receives a count proportional to the square
root of its projected contribution to the variance of $C_7$ per unit cost,
the rule \eqref{gff:meth:allocation} with $C_7$ in place of $C_6$, rounded
upward to a power of two and frozen before production, separately in the
power-free and in the power-containing sector. Per replicate the frozen
counts are
\begin{center}\begin{tabular}{@{}lrrrrrr@{}}\toprule
Sector & $L=3$ & $L=4$ & $L=5$ & $L=6$ & $L=7$ & Total\\\midrule
Power-free, nonprimitive & 10816 & 24704 & 97600 & 188480 & 251472 & 573072\\
Power-containing & 32768 & 163840 & 226304 & 130688 & 240384 & 793984\\\bottomrule
\end{tabular}\end{center}
together with 4,096 points for each primitive class (records
\path{LRI_C7_D2_HPC/reference/LRI_C7_d2_progress/verification/frozen_allocation.json}
and \path{LRI_C7_D2_HPC/reference/power_mass/inputs/production_plan.json}).
With the factor $2^{12}$ and the $2^{26}$ primitive points of the 2026
production, a replicate has $2^{12}(573{,}072+793{,}984)+44\cdot2^{26}$
points, and the two campaigns of 64 replicates contain the
1,094,688,178,176 samples quoted above; the group configurations
\path{GFF/data/C7_2/run1-20260930-m4096/groups/*/config.json} hold these
counts. The earlier calculation used the frozen counts themselves, with 16
replicates in each of two campaigns per sector, rate multipliers $0.8$ and
$0.9$, and 24,105,472 power-free and primitive and 25,407,488
power-containing kernel evaluations. Its power-free coefficient
$-3955.383753\pm76.541680$ (sampling standard error) includes the primitive
contribution of Section~\ref{gff:mg7:primitive} once; adding the power correction
$1.750131\pm5.424390$ (sampling standard error) gave its estimate. There $4\widehat\sigma=306.9$
exceeded both half-stream deviations, $89.3$, and rounding up to a multiple
of 10 gave the allowance 310 of Section~\ref{gff:c7d2:result}.


The campaign means are combined with fixed weights $1/2$. With $J_L$ and
$J_{f,L}$ the Jacobians of $\mathcal F$ at $f+p$ and at $f$, and
$V_{f,L},V_{p,L}$ the covariances of the campaign-mean input vectors, the
complete coefficient and the correction have covariances
\begin{equation}
 V_C=\sum_LJ_L(V_{f,L}+V_{p,L})J_L^{\mathsf T},\qquad
 V_\Delta=\sum_L(J_L-J_{f,L})V_{f,L}(J_L-J_{f,L})^{\mathsf T}
 +\sum_LJ_LV_{p,L}J_L^{\mathsf T}.
 \label{gff:c7d2:covariance}
\end{equation}
The earlier calculation added a retained primitive variance once; the 2026
production has no shared retained input. The two campaigns and the two
half streams give
\begin{center}\begin{tabular}{@{}lrr@{}}\toprule
Estimate & $C_7(2)$ & Sampling SE\\\midrule
Campaign $a$ & $-3875.3202$ & $0.7266$\\
Campaign $b$ & $-3876.9320$ & $0.7565$\\
First half stream & $-3876.4894$ & \\
Second half stream & $-3875.7627$ & \\
Combined & $-3876.1261$ & $0.5244$\\\bottomrule
\end{tabular}\end{center}
The campaigns differ by 1.54 combined standard errors. The error is
dominated by the power-free seven-loop graphs:
\begin{center}\begin{tabular}{@{}lrrrrr@{}}\toprule
Sector & $L=3$ & $L=4$ & $L=5$ & $L=6$ & $L=7$\\\midrule
Power-free & 0.0007 & 0.0047 & 0.0200 & 0.0590 & 0.5187\\
Power & 0.0000 & 0.0001 & 0.0013 & 0.0130 & 0.0440\\\bottomrule
\end{tabular}\end{center}
Entries are in units of $C_7$ and combine in quadrature.

\subsection{Checks}
The same uncalibrated samples give lower-order coefficients, compared with
the special-function values of Section~\ref{gff:sec:guide} and, for $C_6$, with
the two-dimensional production of Section~\ref{gff:prec:d2production2026}:
\begin{center}\begin{tabular}{@{}rrrr@{}}\toprule
Order & Estimate & Sampling SE & Pull\\\midrule
3 & $-0.8748513621$ & $1.25\times10^{-6}$ & $-1.08$\\
4 & $5.0488792961$ & $3.01\times10^{-5}$ & $+0.36$\\
5 & $-38.8344913344$ & $5.07\times10^{-4}$ & $-1.67$\\
6 & $362.3727456$ & $0.0085778$ & $+0.73$\\\bottomrule
\end{tabular}\end{center}
The reference values are checks only; no known coefficient is imposed and
no power correction is forced to vanish. The measured power corrections at
orders four to six are $(5.4\pm15.1)\times10^{-8}$,
$(1.10\pm1.85)\times10^{-5}$ and $(3.3\pm11.8)\times10^{-4}$; at order three
the difference is at the level of rounding. These are consistent with zero
but do not prove a cancellation at seventh order, where the correction is
$-0.056\pm0.046$. A bootstrap of 4,096 replicate resamples gives standard
error 0.52219, against 0.52445 from covariance propagation. The new value
lies 1.01 standard errors of the earlier calculation from its estimator,
and its power correction lies 0.33 of that calculation's standard errors
from its value $1.750131\pm5.424390$ (sampling standard error). Because the power sector is integrated
explicitly here, the sixth-order by-product also tests the analytic
power-sector compensation used at sixth order. These checks share the graph
and renormalization framework and do not validate it. The eighth-order
calculation of Section~\ref{gff:app:C8d2} gives $C_7(2)=-3875.55\pm0.29$ (sampling
standard error) as a by-product, $+0.97$ combined standard errors from
\eqref{gff:c7d2:reported}, and the independent code of
Section~\ref{gff:c8d2:independent} gives $-3875.47\pm0.35$, $+1.04$.

Besides \eqref{gff:c7d2:covariance}, the linearization
$\delta C_7=J\,\delta(f+p)=J_f\,\delta f+[(J-J_f)\,\delta f+J\,\delta p]$
gives the power-free variance and the cross-covariance of the power-free
coefficient with the correction,
\begin{equation}
 V_{\rm pf}=\sum_LJ_{f,L}V_{f,L}J_{f,L}^{\mathsf T},\qquad
 V_{f\Delta}=\sum_LJ_{f,L}V_{f,L}(J_L-J_{f,L})^{\mathsf T},\qquad
 V_C=V_{\rm pf}+V_\Delta+V_{f\Delta}+V_{f\Delta}^{\mathsf T},
 \label{gff:mg7:covDecomposition}
\end{equation}
where the last relation is an exact identity. The record
\path{GFF/results/C7_2/C7_result.json} gives $V_{f\Delta}=-5.09\times10^{-6}$
for $C_7$, a correlation of $-2\times10^{-4}$, and
$0.5224502^2+0.0458637^2-2\cdot5.09\times10^{-6}=0.5244497^2$ reproduces the
sampling standard error of Section~\ref{gff:c7d2:result}. In the earlier
calculation a retained primitive variance $\sigma_{\rm prim}^2$ entered
$V_C$ and $V_{\rm pf}$ once and canceled from $V_\Delta$, so neither the old
power-free error nor the common primitive error was averaged down by its
two power campaigns.


\section{Full seventh order in three dimensions}\label{gff:app:C7d3}

\subsection{Result and scope}\label{gff:c7d3:result}
At $d=3$, with $\epsilon=2s-3$, the seventh-order coefficient follows the
construction of Section~\ref{gff:app:C7d2} with the three-dimensional kernel
and normalization of Section~\ref{gff:cn:finiteAssembly}. With $\beta$ and
$\gamma_r$ as in \eqref{gff:c7d2:rg}, $\widehat b_1=3$ and $\widehat c_1=1$,
\begin{equation}
 \Delta_{\mathcal E}(3,\epsilon)=\frac32-\frac\epsilon2+\gamma_r(g_*,\epsilon),
 \label{gff:c7d3:rg}
\end{equation}
and every loop order keeps its Taylor polynomial through $z^{7-L}$ before
$z=\epsilon$ is set. As at $d=2$, and unlike the three-dimensional
sixth-order production of Section~\ref{gff:prec:d3production2026}, the
power-containing sector is integrated explicitly, and $C_7(3)$ is split as
in \eqref{gff:c7d2:difference}. The production of
Section~\ref{gff:c7d3:production} gives
\begin{center}\begin{tabular}{@{}lrr@{}}\toprule
Contribution & Estimate & Sampling standard error\\\midrule
Power-free coefficient $C_7^{\rm pf}(3)$ & $-220.3723574$ & $0.0277073$\\
Power correction $\Delta_pC_7(3)$ & $0.0178769$ & $0.0156478$\\
Complete coefficient $C_7(3)$ & $-220.3544805$ & $0.0318188$\\\bottomrule
\end{tabular}\end{center}
An earlier three-dimensional seventh-order calculation with the same worker,
whose code and records are kept in \path{LRI_C7_D3_HPC/reference/}
(Section~\ref{gff:cn:reproduction}), reported four sampling standard errors
rounded upward, with its two campaigns and its two half streams (the
estimates from the first and from the second half of the replicates of every
group) as diagnostics. The same rule, rounded up to a multiple of $0.01$, was
written down before the production result was available and confirmed
unchanged afterwards. Here $4\widehat\sigma=0.127$, which gives
\begin{equation}\boxed{C_7(3)=-220.35\pm0.13,\qquad\epsilon=2s-3.}
\label{gff:c7d3:reported}\end{equation}
Both half-stream estimates deviate from $C_7(3)$ by $0.034$, less than
$4\widehat\sigma$, so the two-dimensional rule \eqref{gff:c7d2:allowance} gives
the same allowance. The earlier $C_7(3)$ calculation reported $-228\pm16$
(estimator $-227.6181095$, sampling standard error $3.8273969$); the present
estimator lies 1.90 of its standard errors from it. The allowance is
empirical. It is not a probability guarantee and does not bound an error
shared by the graph and renormalization framework; the additional
estimator digits are given for reproduction only.

\subsection{Kernel, normalization and coverage}
The mixed forest kernel is \eqref{gff:c7d2:kernel} with every forest
denominator $U_F$ replaced by $U_F^{3/2}$, as in \eqref{gff:d3six:kernel}; the
fractional power is applied to each forest term before the terms are
summed. The edge measure is $\prod_ex_e^{-1/4+z/4}$ and the normalization
is that of \eqref{gff:d3six:integral} and \eqref{gff:d3six:G}--\eqref{gff:d3six:norm},
now with Gamma jets through order $7-L$. The integrand jets follow
\eqref{gff:c7d2:jets} with $a_{L,j}$ now taken from
\eqref{gff:d3six:GammaJets}; the logarithmic function $h(x)$ coincides with
$s_L(x)$ of \eqref{gff:d3six:G}. In three dimensions an unmarked self-energy
insertion has superficial degree $3/2$ in momenta at $z=0$. The amplitude is
even in the external momentum, so the constant Taylor subtraction again
leaves a remainder of quadratic order and the local mass subtraction of
Section~\ref{gff:app:C7d2} applies unchanged. Before subtraction a connected
power face has radial exponent $3e/4-3h/2=-3/4$ at $z=0$; the Taylor
remainder raises it to $+1/4$.

The graph classes, rational weights, zero-channel exclusions and forest
maps are those of Section~\ref{gff:app:C7d2}. The integrated classes at loops
three to seven number $5,16,76,395$ and $2{,}298$ ($2{,}254$ nonprimitive
and 44 primitive) in the power-free sector and $1,3,20,107,590$ in the
power-containing sector, $3{,}511$ kernels in all. For every one of them
the denominator polynomials and weights were checked to be unchanged from
the two-dimensional inventory; only the dimension-dependent exponents
differ. The two-loop input is \eqref{gff:d3six:twoLoop}, now through $z^5$, with
$b_{2,5}=26.0214190551330710$ in addition to \eqref{gff:d3six:twoLoopValues};
two 65-digit evaluations with radii $0.04,0.07$ and 64,80 Cauchy nodes
agree to $3.3\times10^{-57}$, and with the sixth-order input to
$1.3\times10^{-54}$. Thirty independent Cauchy evaluations of the Gamma
factors agree to $8.3\times10^{-54}$. The fixed point is solved as
in Section~\ref{gff:app:C7d2} and substituted in \eqref{gff:c7d3:rg}; 64 exact
rational cases test the assembly. The seven-loop constants enter only as
$C_7(3)\supset(3\widehat c_{7,0}-\widehat b_{7,0})/6561$.

The single three-loop power graph has the closed form
$\widehat b_3^{\rm p}(0)=-16\pi/15=-3.3510322$, recorded with the earlier
calculation, and $\widehat c_3^{\rm p}=\widehat b_3^{\rm p}/3$. The
production below gives $\widehat b_3^{\rm p}(0)=-3.3510653\pm0.0000262$,
$-1.26$ sampling standard errors from this value.

\subsection{The 2026 production}\label{gff:c7d3:production}
The production worker is byte-identical to that of the earlier
calculation. It keeps the extended-precision first pass,
with promotions to binary128 and to 350-decimal GMP arithmetic at the
unchanged thresholds $10^{-9}$ and $10^{-23}$. Before the campaign, the
worker replayed all 168,528 archived production rows on their archived
nets and seeds. Of these, 96\% reproduced bit for bit and the rest within
$1.0\times10^{-11}$ in normalized units, with identical promotion counts.
The rebuilt estimator of the earlier calculation was
$-227.61810947469922$, against its recorded value $-227.61810947469945$. An independent high-precision
reimplementation of the sampler and integrand reproduced 13 archived rows
across loops three to seven and both sectors. Results do not depend on the
thread count: 5,776 seven-loop rows run at 64 and at 16 threads were
identical, including 255,396 binary128 and 20,511 GMP promotions. The same
rows computed on AMD and on Intel processors differ by at most
$2.0\times10^{-11}$ in normalized units, with identical promotion counts;
this is about $10^{-7}$ of the relative statistical error.

Two independent campaigns $a$ and $b$ have 64 randomized replicates for
every loop and sector, 52- and 60-bit digital coordinates, and importance
rates multiplied by $0.8$ and $0.7$. Every stratum receives the frozen
per-graph allocation of the earlier calculation, at least 64 points,
multiplied by $2^{12}$. The exception is the 44 seven-loop primitive
classes. They are integrated with $2^{22}$ points each per replicate, a
flat allocation in place of their frozen counts, which changes only the
variance. Counts are powers of two and the base seed is fresh. The two
campaigns contain 867,784,720,384 kernel samples. They ran on the symmetry
cluster on 14 AMD and 14 Intel nodes and completed on 1 October 2026; every
merge and the assembly used the package's own tools. The covariances
follow \eqref{gff:c7d2:covariance}, and there is no shared retained input. The
two campaigns and the two half streams give
\begin{center}\begin{tabular}{@{}lrr@{}}\toprule
Estimate & $C_7(3)$ & Sampling SE\\\midrule
Campaign $a$ & $-220.3220$ & $0.0431$\\
Campaign $b$ & $-220.3870$ & $0.0468$\\
First half stream & $-220.3207$ & \\
Second half stream & $-220.3883$ & \\
Combined & $-220.3545$ & $0.0318$\\\bottomrule
\end{tabular}\end{center}
The campaigns differ by 1.02 combined standard errors. The error is
dominated by the seven-loop graphs of both sectors:
\begin{center}\begin{tabular}{@{}lrrrrr@{}}\toprule
Sector & $L=3$ & $L=4$ & $L=5$ & $L=6$ & $L=7$\\\midrule
Power-free & 0.0000 & 0.0006 & 0.0034 & 0.0119 & 0.0248\\
Power & 0.0000 & 0.0002 & 0.0014 & 0.0058 & 0.0145\\\bottomrule
\end{tabular}\end{center}
Entries are in units of $C_7$ and combine in quadrature.

\subsection{Checks}\label{gff:c7d3:checks}
The same uncalibrated samples give lower-order coefficients, compared with
the special-function values of Section~\ref{gff:sec:guide} and, for $C_6$, with
the three-dimensional production of Section~\ref{gff:prec:d3production2026},
with both uncertainties combined:
\begin{center}\begin{tabular}{@{}rrrr@{}}\toprule
Order & Estimate & Sampling SE & Pull\\\midrule
3 & $-0.3535806960$ & $3.59\times10^{-7}$ & $+2.30$\\
4 & $1.2280727249$ & $8.52\times10^{-6}$ & $+1.08$\\
5 & $-5.8503454757$ & $1.52\times10^{-4}$ & $+0.04$\\
6 & $33.5738672$ & $0.0019782$ & $-0.79$\\\bottomrule
\end{tabular}\end{center}
The reference values are checks only; no known coefficient is imposed and
no power correction is forced to vanish. A pull above three in magnitude
would have been reported without changing $C_7(3)$; none occurred. The
measured power corrections at orders four to six are
$(2.0\pm2.1)\times10^{-6}$, $(4.0\pm6.4)\times10^{-5}$ and
$(-3.7\pm10.2)\times10^{-4}$; at order three the difference is at the level
of rounding. These are consistent with zero but do not prove a cancellation
at seventh order, where the correction is $0.018\pm0.016$. A bootstrap of
4,096 replicate resamples gives standard error 0.03204, against 0.03182
from covariance propagation. The power correction lies 1.92 standard
errors of the earlier calculation from its value $-3.836223\pm2.007117$. Because the power
sector is integrated explicitly here, the sixth-order by-product also tests
the analytic power-sector compensation used for $C_6(3)$. These checks share
the graph and renormalization framework and do not validate it.

\section{Full eighth order in one dimension}\label{gff:app:C8complete}

\subsection{Result and scope}\label{gff:c8f:result}
The eighth-order coefficient uses the coordinate-forest construction of
Section~\ref{gff:app:C7complete}, extended by one order:
\begin{equation}\boxed{
 C_8(1)=8059767,\qquad \text{empirical allowance: }\pm17000,
}\label{gff:c8f:answer}\end{equation}
that is, $C_8(1)=(8.060\pm0.017)\times10^6$. The propagated sampling
standard error is $4210$; a bootstrap over the replicates of all
independent streams gives $4140$. The estimator is the mean of fourteen
independent production batches of one frozen design
(Section~\ref{gff:c8f:production}). Their individual values lie between
$8041935$ and $8089624$, each with standard error between $15090$ and
$16343$, and give $\chi^2=8.49$ for 13 degrees of freedom about the
combined value. The allowance rule was fixed before production and was
not changed: the largest of four sampling errors, twice the last refinement
plus four of its errors, twice the degree eight to nine check step plus four
of its errors, the half-split displacement and twice the last degree step of
the exact loop-five input, rounded up to two significant figures. The
four-error candidate, $4\times4209.58=16838.3$, is the largest. The number
of production batches was not fixed in advance. Under the stopping rule set
first, nine batches give $8061480\pm5300$ with allowance $\pm22000$; all
fourteen planned batches were completed after interim results of five to
eight batches had been seen (Section~\ref{gff:c8f:production}). As for the lower orders, the allowance is not a rigorous
enclosure and does not bound the omitted orders of the $\epsilon$ series.
The ratio $C_8(1)/C_7(1)=-25.0$ continues the sequence
$-12.2,-16.1,-19.6,-22.5$ of the preceding ratios.

\subsection{Changes relative to seventh order}\label{gff:c8f:changes}
The kernel, the binary cluster trees, the meromorphic radial quadrature,
the zero-pole controls and the affine assembly \eqref{gff:c7f:affine} are
unchanged. Each loop vector carries eight Laurent and Taylor slots per
channel instead of seven, and loops up to eight are integrated. Since two
insertions of loops five or higher first contribute beyond eighth order,
$C_2,\ldots,C_8$ are again affine in the loop-five to loop-eight jets. The
affine map shows that $C_8$ depends on loop eight only through the
combination $3c_8(0)-b_8(0)$: its two coefficients are in the ratio
$-3$ to rounding, so the classes with $3w_B=w_A$ contribute exactly zero.

The nine-vertex inventory has 26,540 classes. Of these, 7,231 contain a
power-divergent subgraph, 248 are primitive and power free, and 3,908 have an identically
vanishing forest polynomial. The weighted forests have 2,507,176
coordinate terms. The enumerator and the forest builder are the
seven-loop programs generalized to nine vertices with 128-bit packed
multiplicity keys. Regenerating loops one to seven reproduces the archived
inventories, forest arrays and exact-zero lists byte for byte. The
coefficientwise comparison with bare-series inversion holds as a rational
polynomial identity through eight loops. The exact class weights satisfy the
zero-dimensional sum rules \eqref{gff:eq:zeroDimSums} at every loop order through eight.
Every power-free class with an articulation vertex is an exact zero, so no
class is excluded beyond the exact zeros and the vanishing top-loop
channel.

Loops two to five are no longer sampled. Their vectors are deterministic
sums over all $(L+1)!\,\mathrm{Cat}_L/2$ reflection-reduced ordered trees,
each integrated with the full tensor quadrature: loops two to four at
degree 24 for the constant $K$, and loop five at degrees
$6,8,10,12,14$. The constant reproduces $C_2$ exactly and $C_3$ and $C_4$
to $1.3\times10^{-12}$ and $1.1\times10^{-9}$. Since $C_5$ depends only on
loops up to five, it is now a deterministic check: at the five degrees it
differs from the exact value by $-0.94$, $-0.13$, $-0.019$, $-0.0032$ and
$-5.8\times10^{-4}$. The loop-five contribution to $C_8$ at degree 14 is
$-1297392.1$, and its last degree step is $+2.8$.

Loops six to eight are sampled. The power-free sector uses the two
independent randomized Sobol campaigns of Section~\ref{gff:app:C7complete}.
The power sector is telescoped as in \eqref{gff:c7f:levels}, but to degree
eight:
\begin{equation}
 \mathbb E[Y_8]=\mathbb E[Y_1]+\mathbb E[Y_3-Y_1]+\mathbb E[Y_5-Y_3]
 +\mathbb E[Y_6-Y_5]+\mathbb E[Y_7-Y_6]+\mathbb E[Y_8-Y_7],
 \label{gff:c8f:levels}
\end{equation}
with a paired degree eight to nine level that is never added to the
estimate. With one tree per active eight-loop power graph, one replicate costs
about 5, 14, 32 and 67 core-hours at degrees six, seven, eight and nine. Degree eight is therefore affordable, whereas the seventh-order
estimate stopped at degree six, where its truncation dominated.

\subsection{Pilot, frozen design and production}\label{gff:c8f:production}
An independent pilot of all thirty sampled campaigns fixed the design and
trained the zero-pole controls on its coarse power campaigns; it never
enters the estimate. It also exposed a sampling hazard. The seventh-order
production gives $16372.8$ for the non-primitive contribution to the $C_7$
row. With 64 points per graph and 16 replicates, the two independent
seven-loop power-free campaigns of the pilot gave $-16793.9\pm9359.6$ and
$-13230.8\pm7406.2$ (sampling standard errors), that is $33166.7$ and
$29603.7$ below the production value, or $3.5$ and $4.0$ of their own
standard errors. A repetition of the first campaign with 4096 points per
graph and 32 replicates gives $16838.1\pm724.1$. That this value agrees to the unit with the
four-error candidate $16838.3$ of Section~\ref{gff:c8f:result} is a coincidence:
the two are unrelated quantities. The randomized estimator is unbiased, but at
small point counts its replicate distribution is strongly skewed. The
production therefore uses at least 4096 points per graph in every
power-free campaign. Per-graph sample counts follow a Neyman allocation
from the pilot, and the per-campaign multipliers are split into at most
128 replicates and a power-of-two sample multiplier.

The production consists of fourteen batches with fresh seeds, each of
498 work elements. The first four ran on SDSC Expanse in elements of 16
AMD EPYC 7742 cores and give $8066723\pm8015$ (sampling standard error).
Ten further batches of the same design were then planned together, at
first under a rule that started no new work after 48 hours and, once the
running work had finished, included the batches then complete. That rule
gives the first nine batches,
$8061480\pm5300$ with allowance $\pm22000$. After interim results of five
to eight batches had been seen, it was decided, on grounds of precision,
to complete all fourteen planned batches; from then on every planned batch
was to be included once complete, and all were. The design and the
allowance rule are unchanged. Seven batches ran on Intel Xeon Gold 6148
nodes of the symmetry cluster and seven on Expanse; each batch ran on one
processor type. The means of the two groups of seven differ by
$7037\pm8410$. All campaigns are verified
against the frozen design, including its counts, multipliers, numerical
input hashes and the loop-five input. No seed of the pilot or of the control
training is reused. A coverage guard requires every active class of loops six to
eight to be integrated with total weight one at every level, including
the check level. The calculation used about $9\times10^4$ physical core-hours
($1.4\times10^5$ as charged, counting hyperthreads on the symmetry cluster),
including development, the exact low-loop stage and the pilot.

The kernel was compared on identical tasks with the seventh-order
implementation. In power mode the seven common slots agree bitwise at loops
five to seven and degrees one to six; the randomized point evaluations
agree to $8\times10^{-17}$ relative. A binary128 replay of the contraction
gives $2.5\times10^{-17}$ of the row scale at loop eight.

\subsection{Composition and lower-order checks}\label{gff:c8f:checks}
The contributions to $C_8$ are
\begin{center}\begin{tabular}{@{}lrr@{}}\toprule
Piece & Contribution to $C_8$ & Sampling SE\\\midrule
Loops one to four, $K$ & $1531302.35$ & \\
Loop five, exact at degree 14 & $-1297392.1$ & \\
Power-free sector, loops six to eight & $7879086.8$ & $790.7$\\
Power sector, degree one & $316544.2$ & $1690.3$\\
Refinement $1\to3$ & $-364981.8$ & $2887.4$\\
Refinement $3\to5$ & $-4590.2$ & $2346.3$\\
Refinement $5\to6$ & $-26.5$ & $528.4$\\
Refinement $6\to7$ & $-202.2$ & $321.0$\\
Refinement $7\to8$ & $26.3$ & $117.8$\\
Check $8\to9$, not added & $-113.3$ & $68.1$\\\bottomrule
\end{tabular}\end{center}
The error is dominated by the eight-loop refinements $1\to3$ and
$3\to5$ and by the eight-loop coarse level. The last two refinements and
the check step are small and statistically consistent with zero. Adding
the check step would move the estimate by $-113$, less than one per cent
of the allowance; since it exceeds the last refinement in magnitude, no
geometric tail beyond degree nine is fitted. The radial truncation is
therefore negligible at the present precision, in contrast to seventh
order. With the frozen zero-pole controls the estimate is $11078$ above
the uncontrolled estimator, $8048689$, and its standard error falls from
$23730$ to $4210$. The controls are fixed combinations of 30 Laurent pole
coefficients that vanish exactly. Their standardized sample means are at
most $2.18$ in magnitude; the four above $2$ form two pairs of columns
that move together, and the squares sum to $39.8$, which is only
descriptive because the columns are correlated. Projected onto $C_8$, the
control correction, $+11078\pm24734$, is consistent with zero, but these
residuals alone do not exclude a nonzero pole expectation at the scale of
the allowance. Errors of the controls at finite degree are refined with the
levels and so enter the refinement terms of the allowance; a failure of
the exact vanishing itself would be an error of the framework, which the
allowance does not cover. The half-split
estimates are $8055451$ and $8064082$.

The same samples give the lower-order coefficients, uncalibrated:
\begin{center}\begin{tabular}{@{}rrrrr@{}}\toprule
Coefficient & Estimate & Sampling SE & Reference & Difference\\\midrule
$C_5(1)$ & $-731.140052$ & deterministic & $-731.139477$ (exact) & $-5.8\times10^{-4}$\\
$C_6(1)$ & $14332.8222$ & $2.52$ & $14331.68703$ (deterministic) & $+0.45$ SE\\
$C_7(1)$ & $-322495.6$ & $74.4$ & $-322549.9\pm10.1$ (Sec.~\ref{gff:app:C7complete}) & $+0.72$ SE\\\bottomrule
\end{tabular}\end{center}
The $C_7$ comparison combines both sampling errors. Against the
truncation-aware seventh-order value \eqref{gff:c7f:truncationAware} the
difference is $0.4$ combined standard errors. The references are checks only and are not
imposed. As at lower orders, these comparisons share the graph and
renormalization framework and do not validate it.

\needspace{12\baselineskip}
\section{Full eighth order in two dimensions}\label{gff:app:C8d2}

\subsection{Result and scope}\label{gff:c8d2:result}
At $d=2$, with $\epsilon=2s-2$, the eighth-order coefficient extends the
mixed local-mass forest calculation of Section~\ref{gff:app:C7d2} by one order.
The renormalization functions of \eqref{gff:c7d2:rg} are summed through eight
loops,
\begin{equation}
 \beta(g,z)=-zg+\sum_{L=1}^{8}\widehat b_L(z)g^{L+1},\qquad
 \gamma_r(g,z)=\sum_{L=1}^{8}\widehat c_L(z)g^L,
 \label{gff:c8d2:rg}
\end{equation}
with $\widehat b_1=3$, $\widehat c_1=1$ and
$\Delta_{\mathcal E}(2,\epsilon)=1-\epsilon/2+\gamma_r(g_*,\epsilon)$, and
every loop order keeps its Taylor polynomial through $z^{8-L}$ before
$z=\epsilon$ is set. Both sectors are integrated explicitly, and in the
notation of \eqref{gff:c7d2:difference}
\begin{equation}
 C_8(2)=C_8^{\rm pf}(2)+\Delta_pC_8(2),\qquad
 \Delta_pC_8=\mathcal F_8(f+p)-\mathcal F_8(f).
 \label{gff:c8d2:difference}
\end{equation}
The production of Section~\ref{gff:c8d2:production} gives
\begin{center}\begin{tabular}{@{}lrr@{}}\toprule
Contribution & Estimate & Sampling standard error\\\midrule
Power-free coefficient $C_8^{\rm pf}(2)$ & $45991.2684796$ & $7.6797436$\\
Power correction $\Delta_pC_8(2)$ & $-1.2658435$ & $2.9224044$\\
Complete coefficient $C_8(2)$ & $45990.0026361$ & $8.2168210$\\\bottomrule
\end{tabular}\end{center}
The estimate with its sampling standard error is $C_8(2)=45990.0\pm8.2$; the
production record quotes this, together with the weight bounds and tail
diagnostics of Section~\ref{gff:c8d2:tails}. As for the lower orders, the note
quotes instead the reporting allowance of the rule \eqref{gff:c7d2:allowance} of
Section~\ref{gff:app:C7d2}, with half streams defined there. Here
$4\widehat\sigma=32.87$ exceeds both half-stream deviations, $2.39$, and
rounding up to a multiple of $1$, two significant figures as for $C_7(2)$ and
$C_7(3)$, gives
\begin{equation}\boxed{C_8(2)=45990\pm33,\qquad\epsilon=2s-2.}
\label{gff:c8d2:reported}\end{equation}
The rule was chosen after the result was available. No allowance was
registered with the production decision, unlike the four-standard-error rules
of $C_6(2)$ and $C_6(3)$, which were written down before their results
(Sections~\ref{gff:prec:d2production2026} and~\ref{gff:prec:d3production2026}).
There is no earlier eighth-order calculation at $d=2$. A separately
developed code, run after this production, gives $46006.1\pm13.4$, $+1.03$
combined standard errors from it (Section~\ref{gff:c8d2:independent}). The
exact-path pilot
of Section~\ref{gff:c8d2:design}, with an independent seed, gave
$45620\pm214$, 1.73 combined standard errors from the production. The ratio
$C_8(2)/C_7(2)=-11.86$ continues the sequence $-5.77,-7.69,-9.33,-10.70$ of
the preceding ratios. The allowance is empirical. It is not a probability
guarantee and does not bound an error shared by the graph and
renormalization framework; the additional estimator digits are given for
reproduction only.

\subsection{Changes relative to seventh order}\label{gff:c8d2:changes}
The mixed forest kernel \eqref{gff:c7d2:kernel}, the power-sector normalization
\eqref{gff:c7d2:normalization}, the graph weights \eqref{gff:c7d2:weights}, the
importance sampler and the default precision policy of
Section~\ref{gff:c7d2:production} are unchanged. The native integrator is the
seven-loop one generalized to sixteen propagators and 31-dimensional digital
nets. It keeps an order-seven mode, in which it reproduces the seventh-order
production executables bit for bit (Section~\ref{gff:c8d2:design}).

\paragraph{Regulator jets and low-loop inputs.}
The integrand jets follow \eqref{gff:c7d2:jets} with $n\le8-L$, so the maximum
orders at loops two to eight are $6,5,4,3,2,1,0$. The Gamma table gains the
column $a_{L,6}$, used only by a two-loop control, and an eight-loop row
with normalization $-2/\pi^8$, whose logarithms are not used. The two-loop
input gains
\begin{equation}
 \widehat b_{2,6}=-266.52096796135938173225219261908980094902\ldots,
\end{equation}
from the same exact-hypergeometric numerator; its first six entries are
unchanged. Two Cauchy contours at 75 digits agree to $3.1\times10^{-66}$, and
a direct 130-digit Taylor series confirms the stored value to
$7.7\times10^{-68}$. Every printed Gamma jet is the correctly rounded value of
two independent Cauchy contours at 110 digits. The two-loop control graph,
integrated by the production code, reproduces the seven entries of
$\widehat b_2$ within 1.22 sampling standard errors.

\paragraph{Assembly.}
The physical assembly is a $d$-general eighth-order module shared with the
three-dimensional calculation, used byte for byte and pinned by hash in both.
Its floating-point inversion agrees with two exact rational inversions, the
bivariate Lagrange formula and the coefficientwise solution of
$\beta(g_*,\epsilon)=0$, to $7.0\times10^{-16}$ of the coefficient scale in 64
random eighth-order cases. With $z$-independent jets it reproduces the
constant-coefficient MS polynomials through eighth order; the eighth-order
polynomial, $19683\,C_8=3c_8-b_8+42Ac_7-16Ab_7+\cdots$ with 34 monomials after
$b_2=-6A$ and $c_2=-A$, was derived independently twice and agrees exactly. At
order seven the module reproduces the assembly of the seventh-order
production (Section~\ref{gff:c7d2:production}) to 2 ulp.
The inversion fixes how the new data enter:
\begin{itemize}
\item the top slot of every loop enters additively,
 $\partial C_8/\partial(\widehat b_{L,8-L},\widehat c_{L,8-L})=(-1,3)/3^{L+1}$
 for $L=3,\ldots,8$, without changing a lower $C_n$. The eight-loop constants
 therefore enter only as $C_8(2)\supset(3\widehat c_{8,0}-\widehat b_{8,0})/19683$,
 through the channel weight $3w_B-w_A$;
\item loop seven enters through both channels,
 $\partial C_8/\partial(\widehat b_{7,0},\widehat c_{7,0})=(-16A_2,42A_2)/19683$,
 with $A_2=-\widehat b_{2,0}/6=4\log2$ of \eqref{gff:gd:blocks}. This is not
 proportional to $(-1,3)$, so the seven-loop classes with $3w_B=w_A$, which
 Section~\ref{gff:app:C7d2} excludes, now contribute through
 $-2A_2\widehat b_{7,0}/19683$;
\item $C_2,\ldots,C_8$ are affine in the jets of loops five to eight taken
 together, since two insertions of loops five or higher first contribute
 beyond eighth order; loops two to four enter nonlinearly.
\end{itemize}
At the production means the three rules are hard gates: the top-slot and
loop-seven derivatives agree with their exact values to $6\times10^{-16}$
relative, and the response of $C_8$ to the jets of loops five to eight is
linear to rounding.

The eight-loop sums are complete for $C_8(2)$ but not individually. Since the
eight-loop classes with $3w_B=w_A$ are not integrated, the stored
$\widehat b_{8,0}$ and $\widehat c_{8,0}$ omit contributions with
$3\,\delta\widehat c_{8,0}=\delta\widehat b_{8,0}$, which cancel in
$3\widehat c_{8,0}-\widehat b_{8,0}$. They are therefore not the eight-loop
coefficients of $\beta$ and $\gamma_r$, and a quantity that depends on them
separately, such as the $\epsilon^8$ coefficient of $g_*$, is not an output of
this calculation. At loops three to seven no nonzero class is omitted, so those
jets are complete. For the restored seven-loop classes the top-slot combination
$3\widehat c_{7,1}-\widehat b_{7,1}$ vanishes by the weights, but both jets are
integrated, as for every other class. The independent code of
Section~\ref{gff:c8d2:independent} also leaves out the eight-loop classes with
$3w_B=w_A$, so the caveat applies equally to both codes.


\paragraph{Graph inventory and complete coverage.}
The nine-vertex inventory has the 26,540 classes of
Section~\ref{gff:c8f:changes}. Its $d=2$ per-graph input files, called payloads below (line masks, forest
terms, weights and importance rates), come from the seven-loop builders of
Section~\ref{gff:app:C7d2}, generalized to nine vertices with 128-bit class keys.
The only dimension-dependent part, the rate table, uses the parity floor of
the seven-loop power-free files: every proper edge subset has rate at least
$1/2$ if its size is odd and at least $1$ if it is even. Regenerating loops
three to seven reproduces every production file byte for byte, except the
rate lines of the power-free files at loops three to six, whose earlier rate
tables are kept unchanged. The same payloads, with the rate line rewritten,
serve the three-dimensional calculation.
\begin{center}\begin{tabular}{@{}rrrrrr@{}}\toprule
 & & & Nonzero & \multicolumn{2}{c@{}}{Integrated classes}\\
$L$ & All classes & Power-containing & mixed kernels & power-free & power\\\midrule
3&8&1&1&5&1\\4&27&5&3&16&3\\5&129&31&20&76&20\\
6&660&173&107&395&107\\7&3986&1091&709&2441&709\\
8&26540&7231&4891&16064&4336\\\bottomrule
\end{tabular}\end{center}
At eight loops 2,254 classes have $3w_B-w_A=0$ and 3,886 further classes have
an exactly vanishing mixed forest map. The remaining 15,816 power-free
nonprimitive, 248 primitive and 4,336 power-containing classes are integrated
in every replicate. At seven loops all 3,150 classes with a nonvanishing map
(2,397 nonprimitive, 44 primitive and 709 power-containing) are integrated,
262 more than in Section~\ref{gff:app:C7d2}. In all, 24,173 kernels are
integrated, about seven times the 3,511 of the seventh-order calculation, and
there is no omitted nonzero class. An eight-loop kernel has up to 4,032
complement-of-tree monomials and 280 forest terms. The exactly vanishing
kernels are exactly the classes with an articulation vertex (2,511
power-free and 2,340 power-containing at eight loops), and the exact class
weights satisfy the zero-dimensional sum rules \eqref{gff:eq:zeroDimSums} at
every loop order. An independent edge-subgraph enumeration rederives the
mixed forest map of every power-containing class (1,091 at seven and 7,231 at
eight loops) without a failure; exact rational Kirchhoff determinants agree
with the stored kernels at random rational points for 144 seven- and
eight-loop kernels; and an independent count of the compatible families of
logarithmic subgraphs (nested, or sharing at most one vertex; the empty family
included) gives 11,254,373 over all eight-loop classes, the total recorded with
the inventory of the one-dimensional calculation, and agrees with the builders'
count for each of the 19,309 power-free classes
(\path{LRI_C8_shared/inventory/inventory8_d2/gates/k_forests.json}).

\paragraph{Power sector and campaigns.}
The power-containing jets follow \eqref{gff:c7d2:normalization} through eight
loops. Their reporting follows a protocol registered before any eighth-order
data. The differences $\Delta_pC_n$ at orders three to seven are null rows,
that is, quantities expected to vanish: through sixth order the power sector
cancels in the exact coefficients (Section~\ref{gff:app:C6d2Complete}), and at
seventh order the measured difference is consistent with zero
(Section~\ref{gff:app:C7d2}). They are flagged when a pull exceeds three in
magnitude but never used as gates, and the eighth-order correction is reported
and never assumed to vanish. Two
independent campaigns $a$ and $b$ again use importance rates multiplied by
$0.8$ and $0.9$ and 52- and 60-bit digital coordinates, with group seeds
derived by hashing from a fresh base seed. The campaign means are combined
with fixed weights $1/2$, and the covariances follow \eqref{gff:c7d2:covariance}
through eight loops; there is no retained input.

\subsection{Pilot, design and numerics}\label{gff:c8d2:design}
An exact-path pilot integrated all 24,173 kernels in both campaigns with 8
replicates of 4,096 points per graph, on three nodes of the symmetry cluster for 227
core-hours. It gave $C_8(2)=45620\pm214$, with campaigns 0.64 standard errors
apart and lower-order pulls at most $2.11$ (at $C_6$). It fixed the design and
never enters the estimate; its seed is excluded from the production. With its
uniform allocation the power-free seven-loop graphs carried 71\% of the
variance at 2.4\% of the cost, and the power-free eight-loop graphs 18\% at
71\%.

Per-graph counts follow a Neyman allocation, $n_g\propto\sigma_g/\sqrt{c_g}$
with the pilot's per-point standard deviations $\sigma_g$ and measured costs
$c_g$, on a power-of-two lattice with at least 4,096 points per graph and 128
replicates in every campaign, sector and loop. After the pilot the author chose,
from a table of costs at four targets, a predicted standard error of
$2.5\times10^{-4}$ of $|C_8|\approx4.6\times10^4$, that is $11.50$, which
includes a safety factor $1.5$ on the raw prediction $7.67$. The design, fixed
before the production,
gives every power-free eight-loop graph at least 32,768 points per replicate
and contains 779,836,456,960 kernel samples in one batch.

\paragraph{Exact path.}
The default numerics are the precision policy of
Section~\ref{gff:c7d2:production}: a long-double first pass, binary128 below the
ratio $10^{-9}$ and 350-digit GMP arithmetic below $10^{-23}$. In order-seven
mode the default executables reproduce the seventh-order production
executables bit for bit, on both backends, on the two-loop graph and all 3,773
kernels of loops three to seven, including the 262 new seven-loop ones; the
order-eight slots common to order seven are identical. Each later kernel
speedup (parallel graph loading, and inline correctly rounded binary128
addition, subtraction and multiplication, tested against the compiler library
on $2^{32}$ random operand pairs per operation) had to reproduce the first
unoptimized eight-loop build bit for bit on a sample of 1,600 power-free and
400 power-containing eight-loop kernels, on both backends, so the default path
is no numerical change. The final executables passed these replays, the
450-digit point checks and the thread-identity tests again on the production
cluster.

\paragraph{The faster numerics profile.}
The production used an opt-in profile, chosen by the author after the pilot
and run only by an executable that passed an accuracy gate. Here $v$ is the
kernel value at a point and ${\rm err}$ the same forest sum with every
numerator term replaced by its absolute value, so that $|v|/{\rm err}$ is the
cancellation ratio of Section~\ref{gff:done:samplingControls}. A point that fails
the long-double test is first evaluated in double-double arithmetic, with a
106-bit significand, and accepted when $|v|\ge10^{-20}\,{\rm err}$; otherwise,
or outside the double-double range, it goes to the unchanged binary128 tier. A
point that reaches the multiprecision tier is first evaluated at 77 digits
(320 bits) and accepted when $|v|\ge10^{-63}\,{\rm err}$, and otherwise at the
unchanged 350 digits. The gate compared the production executable with the
exact path on the production hardware in four parts:
\begin{itemize}
\item[(i)] on 153,600 seven- and eight-loop points with 450-digit references,
 every accepted value is within the exact path's bound of $10^{-6}$ relative,
 with largest errors $5.2\times10^{-12}$ in double-double and
 $1.4\times10^{-11}$ in binary128; every 77-digit value equals the 350-digit
 one after conversion;
\item[(ii)] a paired campaign on identical streams, with four replicates of 256
 points for every kernel of loops three to eight in both campaigns (193,384
 tasks), gives
 $|\Delta C_8|+3\,{\rm SE}_\Delta=4.2\times10^{-16}$ against the bound $0.062$,
 one per cent of the standard error at the tightest target considered, with
 per-task normalized differences at most $5.6\times10^{-14}$.
\item[(iii)] the fractions of points promoted to each tier are unchanged;
\item[(iv)] the results are identical at 1, 8, 16 and 64 threads and in eight
 lanes of eight threads.
\end{itemize}
The profile runs about 2.5 times faster than the exact path at eight loops.
The production design keeps the exact pilot's variances and only re-costs the
allocation with a cost model measured under the profile, which (ii)
justifies; it costs 0.37 of the exact path at the same target.

\subsection{The 2026 production}\label{gff:c8d2:production}
The batch ran on the symmetry cluster on 6 and 7 October 2026, in lanes of
eight threads on AMD EPYC 7543 nodes of one runtime identity: eight nodes for
22.4 to 23.0 hours in a first round and six nodes for 4.4 to 5.6 hours after
the 24-hour limit. All 1,129 work indices, 6,192,256 tasks in 24,848
checkpoint blocks, ran once without a failure, in 28.7 hours from submission
to the final merge and 13,454 allocated core-hours, 0.93 of the design's
estimate. Every merge, the assembly and an independent recomputation of the
result record used the package's own tools. With development, the pilot and
the tail diagnostics the calculation used about $1.6\times10^4$ core-hours.
The two campaigns and the two half streams give
\begin{center}\begin{tabular}{@{}lrr@{}}\toprule
Estimate & $C_8(2)$ & Sampling SE\\\midrule
Campaign $a$ & $45993.58$ & $12.15$\\
Campaign $b$ & $45986.42$ & $11.06$\\
First half stream & $45987.61$ & \\
Second half stream & $45992.39$ & \\
Combined & $45990.00$ & $8.22$\\\bottomrule
\end{tabular}\end{center}
The campaigns differ by 0.44 combined standard errors. The error is dominated
by the power-free eight- and seven-loop graphs:
\begin{center}\begin{tabular}{@{}lrrrrrr@{}}\toprule
Sector & $L=3$ & $L=4$ & $L=5$ & $L=6$ & $L=7$ & $L=8$\\\midrule
Power-free & 0.01 & 0.07 & 0.22 & 0.65 & 3.67 & 6.71\\
Power & 0.00 & 0.01 & 0.11 & 0.43 & 1.02 & 2.70\\\bottomrule
\end{tabular}\end{center}
Entries are in units of $C_8$ and combine in quadrature. The realized standard
error is 1.07 times the raw design prediction, the size expected from
per-graph standard deviations estimated with seven degrees of freedom, and
well inside the safety factor; per stratum the realized variances at loops
seven and eight are 0.98 to 1.32 times the predictions.

\subsection{Weight bounds and tail diagnostics}\label{gff:c8d2:tails}
The package's heavy-tail rule, calibrated on synthetic families at $d=3$,
computes a fixed-$k$ ($k=16$) Hill estimate \cite{gff:Hill1975} $\widehat\xi$ of the
tail index of
the per-point weights of every graph and campaign. At production scale it
confirmed 4,832 eight-loop graphs (4,631 power-free and 201
power-containing) as possible infinite-variance cases, with two-campaign mean
$\widehat\xi>0.44$ at 32,768 or more points per replicate. Together they carry
27.6\% of the variance of $C_8(2)$, and no single graph in one campaign more
than 0.08\%. Since an infinite variance would invalidate the replicate standard
error, the weights were analyzed directly.

In an ordered sector the sampler draws successive scale gaps
$z_i=u_i^{1/(\rho_c\alpha_G(S_i))}$, with $u_i$ uniform, $S_i$ the set of lines
below the $i$th gap, $\alpha_G(S)$ the stored importance rate and
$\rho_c=0.8,0.9$ the campaign multiplier. The weight is
\begin{equation}
 t=Z_c\,K_G(x)\prod_iz_i^{|S_i|/2-\rho_c\alpha_G(S_i)},
 \label{gff:c8d2:weight}
\end{equation}
with $K_G$ the rational kernel \eqref{gff:c7d2:kernel} and $Z_c$ a constant. Let
$m_G(S)=|S|/2+\nu_S(K_G)$ be the exact leading exponent of the line measure
times $K_G$, after all forest cancellations, when the lines of a proper subset
$S$ are scaled to zero. Near that face
$t\sim z^{m_G(S)-\rho_c\alpha_G(S)}$, so the weight is bounded if
$m_G(S)\ge\rho_c\alpha_G(S)$ on every face and has a finite second moment if
$2m_G(S)>\rho_c\alpha_G(S)$. The builders set
$\alpha_G(S)=\max(r_G(S),p_S)/2$, where $r_G(S)/2$ is the smallest exponent of
the individual forest terms and $p_S$ is $1$ for odd and $2$ for even $|S|$.
Where $r_G(S)>0$ a sum cannot have a lower exponent than its terms, so
$m_G(S)\ge\alpha_G(S)$. Where $r_G(S)\le0$ the forest sum must cancel; since
$K_G$ is rational, $m_G(S)\in|S|/2+\mathbb Z$, and $m_G(S)>0$ already implies
$m_G(S)\ge\alpha_G(S)$. Once these cancellations are confirmed,
$m_G(S)\ge\alpha_G(S)>\rho_c\alpha_G(S)$ on every face, and the weight is
bounded at both rates. Nested limits are covered because every forest
denominator $U_F$ is the basis generating polynomial of a matroid, so that a
single monomial dominates all nested limits of a sector at once; this rests on
the builders' construction and was checked on random orderings.

The cancellations were verified at all 3,959,753 expanded subsets of the
24,173 kernels: every subset with $r_G(S)\le0$ of the kernels with the
builders' rate tables, and every subset of the 492 power-free kernels at loops
three to six, which keep their earlier tables. At each, $\sum_FN_F/U_F$ is
expanded in the scaling parameter in exact arithmetic modulo four primes
drawn uniformly from $[2^{61},2^{62})$, at two random points per prime and
graph. A nonzero residue proves a nonvanishing coefficient, so an error can
only overestimate a valuation. With explicit degree and height bounds and the
Schwartz--Zippel lemma, the union bound puts the probability of any wrong
valuation below $10^{-53}$, assuming independent uniform draws. Every stored
rate table is reproduced, every required cancellation holds, and the eight
trials agree everywhere. Nine archived controls behave as asserted: deleting
a forest term or doubling a numerator breaks the cancellation, and rate
tables multiplied by $3$ or $2.4$ give infinite variance with a finite mean at
both rates, or at rate $b$ only. Over all seven- and eight-loop kernels the
margins are $m_G-\rho_c\alpha_G\ge0.1$ and $0.05$ and
$2m_G-\rho_c\alpha_G\ge0.6$ and $0.55$ at rates $a$ and $b$. Seventeen
power-free kernels at loops five and six that keep their earlier rate tables
have unbounded weights of finite variance.

The estimated sampling standard error is $8.22$. Boundary analysis, supported
by this multi-prime randomized cancellation verification of all 24,173
integrated kernels, supports finite second moments at both production rates
and bounded weights for all seven- and eight-loop kernels; it is a randomized
verification, not a deterministic certificate in characteristic zero. The
4,832 eight-loop tail warnings are consistent with bounded, sharply peaked
weights. The confirmed power-free graphs have large canceling forest sums
(median 128 terms, against 16 for the cleared graphs), and the fixed-$k$ Hill
estimate falls with the number of points in every class, with log-log slopes
from $-0.34$ to $-0.17$; such finite-sample behavior cannot by itself prove
boundedness. Fresh nested sampling supports variance stability in the
selected main groups: over $2^{12}$ to $2^{17}$ points the slope of
$\log(N\,{\rm Var})$ is $-0.018$ $[-0.051,+0.030]$ for 16 confirmed
power-free and $-0.062$ $[-0.110,+0.060]$ for 8 confirmed power-containing
graphs, with 95\% intervals conditional on the selected panel. In the 380
power-containing eight-loop graphs with a subset that needs two orders of
cancellation (``double lift''; 0.058\% of the measured variance), a targeted
test of the 16 largest contributors (52\% of that class's measured variance)
shows rare-peak sampling effects, with per-point variance growing over
$2^{12}$ to $2^{19}$ points in the largest two. Inflating that class's
variance by factors 18.9, 45 or 335 gives standard errors 8.26, 8.32 or 8.97;
these variance-inflation calculations are sensitivity scenarios, not error
bounds. Every graph stays in the estimand.

\subsection{Checks}\label{gff:c8d2:checks}
The same uncalibrated samples give lower-order coefficients, compared with
the special-function values of Section~\ref{gff:sec:guide} and, for $C_6$ and
$C_7$, with the productions of Sections~\ref{gff:prec:d2production2026} and
\ref{gff:app:C7d2}, with both uncertainties combined:
\begin{center}\begin{tabular}{@{}rrrr@{}}\toprule
Order & Estimate & Sampling SE & Pull\\\midrule
3 & $-0.8748487258$ & $1.24\times10^{-6}$ & $+1.04$\\
4 & $5.0489759126$ & $6.80\times10^{-5}$ & $+1.58$\\
5 & $-38.8332791554$ & $9.37\times10^{-4}$ & $+0.39$\\
6 & $362.3625083$ & $0.0116244$ & $-0.34$\\
7 & $-3875.545222$ & $0.289774$ & $+0.97$\\\bottomrule
\end{tabular}\end{center}
The squared pulls sum to 4.78 for five checks. The reference values are checks
only; no known coefficient is imposed and no power correction is forced to
vanish. The power null rows at orders four to seven are
$(1.37\pm4.03)\times10^{-6}$, $(7.7\pm23.5)\times10^{-5}$,
$(2.28\pm4.86)\times10^{-3}$ and $(2.54\pm9.00)\times10^{-2}$, with pulls at
most $0.47$; at order three the difference vanishes by the weights. These are
consistent with zero but do not prove a cancellation at eighth order, where
the correction is $-1.27\pm2.92$. A bootstrap of 4,096 replicate resamples
gives standard error 8.09, against 8.22 from covariance propagation. Because
both sectors are integrated explicitly, the sixth-order by-product again tests
the analytic power-sector compensation used for $C_6(2)$, and the
seventh-order by-product, $-3875.55\pm0.29$, is an independent sampling of
the seventh-order integrands of Section~\ref{gff:app:C7d2}. These checks share
the graph and renormalization framework and do not validate it.

\subsection{An independent code}\label{gff:c8d2:independent}
A separately developed eighth-order code extends the seventh-order package of
Section~\ref{gff:app:C7d2} independently of the package used above. It generates
its own nine-vertex classes with 128-bit canonical keys and its own
eight-loop payloads, which are different files from those of
Section~\ref{gff:c8d2:changes}. It extends the kernel and the regulator jets to
eight loops on its own, derives the two-loop and Gamma inputs separately, and
solves for the eighth-order fixed point with its own assembly. It integrates the same 24,173 kernels:
16,064 power-free and 4,336 power-containing at eight loops, and 3,150 at
seven. It shares with the calculation above the mixed forest kernel
\eqref{gff:c7d2:kernel} and its precision policy, the importance sampler with
rates multiplied by $0.8$ and $0.9$, the digital nets and the seed hashing,
the finite renormalization scheme and, byte for byte, the payloads of loops
three to seven. It runs only the exact path of Section~\ref{gff:c8d2:design}
(long double, binary128 and 350-digit GMP), without the double-double and
77-digit tiers.

The independent code, its derivation of the eighth-order extension and the
construction checks under the next four headings were written by OpenAI
GPT-6 astra (model stated by the author), as were the caveat on the separate
eight-loop sums in Section~\ref{gff:c8d2:changes}, the reading of the bootstrap and
the reporting rule described below. They were delivered as the package
\path{LRI_C8_D2_HPC_optimized/}, together with an original package that holds
the graph payloads, the class generator and its check records (hosted, see
\path{GFF/data/C8_2_check/DATA.md}). A Claude Code session ported the code to
the clusters and ran the pilot and the production described afterwards. On
16 October 2026, when the delivered derivation was merged into this subsection, a
Claude Code session also
compared the shipped inputs of the two packages: the Gamma table of the
independent code agrees with that of Section~\ref{gff:c8d2:changes} to 65 digits,
the precision of that table, and its two-loop jets
$\widehat b_{2,0},\ldots,\widehat b_{2,6}$ agree with those of the calculation
above to 70 digits (largest relative difference $2.9\times10^{-70}$, at
$\widehat b_{2,6}$).

\paragraph{Inputs.}
The Gamma table \path{inputs/gamma_d2_8.txt} has one row for each
$L=2,\ldots,8$. A row holds the order $8-L$, the normalization
$2(-1)^{L+1}/\pi^L$ and the logarithms $a_{L,0},\ldots,a_{L,6}$ of
\eqref{gff:done:GammaJets}, with $a_{L,0}=0$; every row carries all seven
logarithms, whether or not its loop uses them.

For the two-loop input, at $d=2$ the parameters of \eqref{gff:gd:twoLoopHyperDefs}
are $a=\tfrac12+z/4$ and $b=\tfrac12+3z/4$. The jets are \eqref{gff:mg7:twoLoop} of
Section~\ref{gff:app:C7d2}, extended to $j=6$: the numerator
$z^2S_2(z)=4F_1(2,z)-2W(2,z)F_2(2,z)$ of \eqref{gff:gd:twoLoopHyper} is analytic at
$z=0$, where it equals 2, and the new entry $\widehat b_{2,6}$ needs its seventh
Taylor coefficient. At $d=2$ the upper parameter $2a=1+z/2$ of $F_2$ equals its
lower parameter $1+z/2$, and the series reduces to
\begin{equation}
 F_2(2,z)={}_2F_1\!\left(\begin{matrix}z,\ b\\ b+\tfrac12\end{matrix};\frac14\right).
 \label{gff:mg8:F2reduced}
\end{equation}
The code expands $F_1(2,z)$ and \eqref{gff:mg8:F2reduced}, both convergent at the
argument $1/4$, in finite Taylor arithmetic in $z$ at 120 decimal digits, where
183 terms of the series suffice, and it expands $\log W(2,z)$ analytically in
polygamma values. The result is
\[
 \widehat b_{2,6}=-266.52096796135938173225219261908980094902003629215748\ldots.
\]
Two evaluations of the unreduced representation on Cauchy contours, of
radius $0.05$ with 128 nodes and of radius $0.08$ with 160 nodes, differ from
the Taylor coefficients by at most $5.26\times10^{-112}$. The first six
coefficients are identical as binary64 numbers to the two-loop input of the
seventh-order calculation and differ from it by at most $4.06\times10^{-70}$ at high precision. The identity
$\widehat b_{2,0}=-24\log2$, that is $A_2=4\log2$ in \eqref{gff:gd:blocks}, holds
exactly and reproduces $C_2(2)=4\log2/9$ of \eqref{gff:gd:d2exact}. The 100-digit
values are in \path{inputs/two_loop.json}, written by
\path{tools/analytic_inputs.py}; its check record is
\path{verification/analytic_inputs.json} of the original package. These checks
concern the special-function input, not the sampled integrals.

\paragraph{Graphs and payloads.}
The code enumerates the one-particle-irreducible classes with its own
canonical keys of 128 bits. At eight loops the keys need up to 69 bits, and
5,068 of them exceed 64. Its class files of loops one to seven are byte for byte
those of the seventh-order calculation. At every loop through eight the exact
sums of $w_A$ and $w_B$ of \eqref{gff:c7d2:weights} over the classes equal the
zero-dimensional values, which the code computes separately from normal-ordered
Gaussian moments and a Legendre transform. For each of the 4,810 classes of
loops three to seven, in both sector inventories (9,620 rows), its forest
builder reproduces the weights, the monomial masks, the forest maps and the
exact-zero lists of the seventh-order files. The importance-rate arrays are left
out of this comparison, because the archived power-free rate arrays were
changed after they were generated; the archived rates of loops three to seven
are used as supplied, and the rates of the new eight-loop payloads follow the
rate formulas of the seventh-order builder. The rates affect only the
efficiency of the sampling, since the integral is defined by the complete
weighted kernel. At eight loops an edge-subgraph enumeration, which includes
the subgraphs not induced by their vertex sets, rederives the mixed forest maps
of all 7,231 power-containing classes, 10,451,732 forests in all, without a
failure. Thirty-six exact rational evaluations on 18 eight-loop kernels agree
with Kirchhoff determinants of the contracted forest blocks.

The census agrees with Section~\ref{gff:c8d2:changes}, with the exclusions taken in
the other order. Of the 19,309 power-free and 7,231 power-containing eight-loop
classes, 2,511 and 2,340 have an exactly vanishing kernel. Of the remaining
16,798 and 4,891 nonzero kernels, 734 and 555 have $3w_B=w_A$, and 16,064 and
4,336 are integrated (\path{inputs/coverage.json}). The 4,851 exact zeros and
1,289 channel exclusions are the $2{,}254+3{,}886=6{,}140$ excluded classes of
Section~\ref{gff:c8d2:changes}. At seven loops all 2,441 nonzero power-free
kernels, 44 of them primitive, and all 709 nonzero power-containing kernels are
integrated, including the 262 classes that the seventh-order production excluded
by $3w_B=w_A$ (the split table of Section~\ref{gff:app:C7d2}); they are already in the
seventh-order reference inventories, although its production payloads omit them. The exact-zero list
is computed from the mixed kernel \eqref{gff:c7d2:kernel} itself. The 3,908
classes with an identically vanishing coordinate-forest polynomial in the
one-dimensional calculation (Section~\ref{gff:c8f:changes}) belong to a different
kernel and are not used; the one-dimensional inventory serves only as a check
of class counts and topology.

\paragraph{Kernel arithmetic.}
Besides the 350-digit GMP final tier used in the production, the kernel has an
alternate final tier in Boost decimal arithmetic at 300 digits. In order-seven
mode, at eight points in each of two streams on all 3,773 kernels of loops
three to seven, including the 262 restored ones, all 60,368 point
contributions reproduce the arithmetic of the seventh-order kernel exactly on
both backends. At 37 points per backend, among them physical eight-loop
kernels and cancellation fixtures, the kernel values agree with an independent
450-digit evaluator to $4.55\times10^{-304}$ (decimal) and
$2.22\times10^{-354}$ (GMP), in units of the absolute-term scale ${\rm err}$ of
Section~\ref{gff:c8d2:design}. These are arithmetic comparisons, not interval
bounds on the integrals. Restart, chunking, complete coverage, immutable inputs
and idempotent merges were tested on native runs. The records are
\path{verification/kernel_checks.json} and \path{generator/reports/} of the
original package.

\paragraph{Assembly.}
The assembly evaluates the bivariate Lagrange formula
\eqref{gff:done:bivariateLagrange} through $n=8$, with
$\widehat B(g,z)=3+\sum_{L=2}^{8}\widehat b_L(z)g^{L-1}$ and $\gamma_r$ of
\eqref{gff:c8d2:rg}. Separately, it solves
$g=\frac13\bigl[\epsilon-\sum_{L=2}^8\widehat b_L(\epsilon)g^L\bigr]\bmod\epsilon^9$
by nine iterations from $g=0$ and substitutes the result in
$\Delta_{\mathcal E}=1-\epsilon/2+\gamma_r(g,\epsilon)$. In 24 exact rational
cases with regulator-dependent jets the two constructions agree exactly, and the
floating-point assembly used by the production agrees with them to a normalized
difference of $2.7\times10^{-15}$. The same test, \path{tests/assembly_formula.py}
(record \path{verification/assembly_formula.json} of the original package),
confirms $\partial C_8/\partial\widehat b_{8,0}=-1/19683$ and
$\partial C_8/\partial\widehat c_{8,0}=1/6561$, the contribution at eighth order
of the seventh-order classes with $3w_B=w_A$, the dependence of $C_8$ on the new
entry $\widehat b_{2,6}$, and the affine dependence on the jets of loops five to
eight, all as derived in Section~\ref{gff:c8d2:changes}.

\paragraph{Development runs (superseded).}
Before the port, OpenAI GPT-6 astra ran the delivered code in three
complete-coverage runs of its own, each with two campaigns $a$ and $b$ and small
point counts. A pilot gave $48626\pm2633$. A refinement with 8 replicates per
group and at least 32 points per eight-loop graph and replicate gave
$47068\pm1854$. A second refinement with 16 replicates and at least 128
power-free and 64 power-containing points per eight-loop graph and replicate
gave $45512\pm524$, quoted with the code as $45510\pm520$ (sampling standard
errors throughout). In the last run the two campaigns differed by 2.26 combined
standard errors, the power correction was $153\pm122$, and in 21 of the 25
graphs with the largest variances a single replicate carried at least 80\% of
the squared deviations in its campaign. This is the small-sample hazard also
met at one dimension (Section~\ref{gff:c8f:production}). The three values lie
$+0.99$, $+0.57$ and $-0.94$ combined standard errors from the production
below. They are superseded by it and enter neither its allocation, which was
computed from the pilot below alone, nor its estimate. The records of the
first two runs are in the original package; those of the third were not
delivered.

\paragraph{Port, pilot and production.}


The code was adapted to the Expanse cluster without a change to its
numerical files. Its optimized kernel reproduced the package's 52 frozen
fixtures bit for bit on the production hardware. It used a fresh base seed,
checked against the seeds of this section. A pilot with 32 replicates per
campaign gave $46178\pm417$ and fixed a Neyman allocation at 28,500 kernel
core-hours. The allocation has 64 replicates in every campaign, sector and
loop, at least 8,192 points per replicate for every power-free and 4,096 for
every power-containing eight-loop graph, and a predicted standard error of
$14.97$. The production ran on 9 and 10 October 2026 on eleven whole AMD EPYC
7742 nodes for 16.3 to 20.1 hours each. All 342 work indices ran once without
a failure, in 207 node-hours (26,526 core-hours), 96.5\% of it kernel time.
\begin{center}\begin{tabular}{@{}lrr@{}}\toprule
Estimate & $C_8(2)$ & Sampling SE\\\midrule
Campaign $a$ & $46012.55$ & $19.70$\\
Campaign $b$ & $45999.73$ & $18.23$\\
First half stream & $46008.16$ & \\
Second half stream & $46004.11$ & \\
Combined & $46006.14$ & $13.42$\\\bottomrule
\end{tabular}\end{center}
The campaigns differ by 0.48 combined standard errors, and a bootstrap of
4,096 replicate resamples gives standard error 13.38. The power-free part is
$46004.84\pm12.90$ and the explicitly integrated power correction
$+1.29\pm3.71$. The power-free eight- and seven-loop graphs carry 78\% and
14\% of the variance, and no single graph more than 0.2\%. At orders three to
seven the lower-order by-products have pulls $-0.01$, $+1.75$, $+0.57$,
$-1.18$ and $+1.04$ against the references of Section~\ref{gff:c8d2:checks}; the
squared pulls sum to 5.88 for five checks. The power null rows at orders four
to six have pulls of at most $0.76$.

The bootstrap resamples the 64 replicates of each group. For a mean of $R$
replicates the plug-in bootstrap variance is $(R-1)/R$ times the unbiased
estimate \eqref{gff:prec:jetcov}. To first order the bootstrap standard error is
therefore $\sqrt{(R-1)/R}$ times the propagated one: $0.992$ here, against the
observed $13.38/13.42=0.997$, and $0.996$ for the 128 replicates of the
calculation above, against $8.09/8.22=0.984$ (Section~\ref{gff:c8d2:checks}). The
finite number of resamples and the nonlinearity of $\mathcal F_8$ also enter
these ratios. The bootstrap mean differs from the estimator by $-0.11$, that is
$-0.008$ standard errors (\path{GFF/results/C8_2_check/prod1/C8_result.json}).
Neither a failure of the first-order propagation \eqref{gff:c7d2:covariance} nor a
finite-sample bias of the nonlinear map is therefore visible at this precision.
Resampling adds no integration points and cannot reveal regions that no
replicate sampled.

\paragraph{Comparison with the calculation above.}


The combined estimate differs from the estimator of \eqref{gff:c8d2:reported} by
$+16.1\pm15.7$, that is $+1.03$ combined standard errors, well inside the
allowance $\pm33$. The power-free parts differ by $0.90$ and the power
corrections by $0.54$ combined standard errors. The agreement tests the parts
that were developed independently: the eight-loop graph generation and
payloads, the eighth-order kernel and jet extension, the analytic inputs, the
assembly and the sampling. It does not test what the two codes share. An error
in the seventh-order kernel construction, the precision policy, the
renormalization scheme or the payloads of loops three to seven would be
common to both. The independent estimate is not combined with
\eqref{gff:c8d2:reported}, which is unchanged.

No reporting allowance is assigned to the independent estimate, and its
result record carries none. Under the reporting rule that came with the
independent code, an allowance and its rule are fixed only after fresh
campaigns, half streams, successive refinements, precision-tier replays and the
lower-order by-products have been compared, because a sampling standard error
alone does not establish final digits; no allowance was fixed for the run above.
The estimate therefore enters this note only through its sampling standard
error, as a check. The weight bounds of
Section~\ref{gff:c8d2:tails} were established for the rate tables of the payloads
of the calculation above. The independent eight-loop payloads carry their own
rate tables, built from the seventh-order rate formulas, and the bounds are
not claimed for them here.


\subsection{Records}\label{gff:c8d2:records}
The result record, the production, tail and pilot reports and the package are
listed in Section~\ref{gff:cn:reproduction}. The production's plan, design bundle
and the merged record of each of its 24 groups, from which the estimator and
every number of this section are recomputed, are in \path{GFF/data/C8_2/}. The
full campaign directories of the production
\path{prod1-2.5e-4-ddmp256-20261006-ba68f401e46e} and of the pilot
\path{pilot1-20261006-ba68f401e46e}, with every checkpoint, are archived as
listed in \path{GFF/data/C8_2/DATA.md}. The records of the independent code of
Section~\ref{gff:c8d2:independent} are in \path{GFF/data/C8_2_check/} and
\path{GFF/results/C8_2_check/}; its full campaign directory is hosted
(\path{GFF/data/C8_2_check/DATA.md}).

\needspace{12\baselineskip}
\section{Full eighth order in three dimensions}\label{gff:app:C8d3}

\subsection{Result and scope}\label{gff:c8d3:result}
At $d=3$, with $\epsilon=2s-3$, the eighth-order coefficient extends the
calculation of Section~\ref{gff:app:C7d3} by one order, as
Section~\ref{gff:app:C8d2} extends Section~\ref{gff:app:C7d2}. The renormalization
functions are summed through eight loops, $\Delta_{\mathcal E}(3,\epsilon)$ is
\eqref{gff:c7d3:rg}, and every loop order keeps its Taylor polynomial through
$z^{8-L}$ before $z=\epsilon$ is set. Both sectors are integrated explicitly
and $C_8(3)$ is split as in \eqref{gff:c8d2:difference}. The production of
Section~\ref{gff:c8d3:production} gives
\begin{center}\begin{tabular}{@{}lrr@{}}\toprule
Contribution & Estimate & Sampling standard error\\\midrule
Power-free coefficient $C_8^{\rm pf}(3)$ & $1601.9999290$ & $0.7279030$\\
Power correction $\Delta_pC_8(3)$ & $0.6052940$ & $0.6129303$\\
Complete coefficient $C_8(3)$ & $1602.6052230$ & $0.9516489$\\\bottomrule
\end{tabular}\end{center}
The quoted uncertainty is an allowance, not a standard error. It follows the
rule \eqref{gff:c7d2:allowance}, also used for $C_7(3)$ and $C_8(2)$: the larger
of four sampling standard errors and the deviations of the two half-stream
estimates (computed from the first and from the second half of the replicates
of every group), rounded up, here to a multiple of $0.1$, which keeps two
significant figures as for $C_7(3)$ and $C_8(2)$. The rules for $C_6(3)$ and
$C_7(3)$ were written down before their production results. This one was
chosen after the result was available and confirmed by the author; no
allowance was registered with the production decision. Here
$4\widehat\sigma=3.81$ exceeds both half-stream deviations, $0.046$, and the
rule gives
\begin{equation}\boxed{C_8(3)=1602.6\pm3.9,\qquad\epsilon=2s-3.}
\label{gff:c8d3:reported}\end{equation}
There is no earlier eighth-order calculation at $d=3$. The pilot of
Section~\ref{gff:c8d3:design}, with independent seeds, gave $1619\pm16$, 1.05
combined standard errors from the production. The ratio
$C_8(3)/C_7(3)=-7.27$ continues the sequence $-1.95,-3.47,-4.76,-5.74,-6.56$
of the preceding ratios. A separately developed code gives $1600.98\pm3.97$
(sampling standard error), $-0.40$ combined standard errors from the
estimator (Section~\ref{gff:c8d3:independent}). The allowance is empirical. It is
not a probability guarantee and does not bound an error shared by the graph
and renormalization framework, and its coverage is conditional on the finite
variance discussed in Section~\ref{gff:c8d3:tails}. The additional estimator
digits are given for reproduction only. The production record quotes the
sampling standard error alone, $C_8(3)=1602.6\pm1.0$, together with the tail
diagnostics of Section~\ref{gff:c8d3:tails}; the reported uncertainty is the
allowance $\pm3.9$ of \eqref{gff:c8d3:reported}.

\subsection{Changes relative to seventh order}\label{gff:c8d3:changes}
The three-dimensional kernel, edge measure and normalization of
Section~\ref{gff:app:C7d3} are unchanged: every forest denominator is raised to
$U_F^{3/2}$ before the forest terms are summed, and the Gamma jets
\eqref{gff:d3six:GammaJets} now run through order $8-L$. The graph inventory,
the mixed forest maps, the weights, the zero-channel rule and the shared
eighth-order assembly are those of Section~\ref{gff:c8d2:changes}, used with the
three-dimensional rate lines; the integrated classes at loops three to eight
number $5,16,76,395,2{,}441$ and $16{,}064$ in the power-free sector and
$1,3,20,107,709$ and $4{,}336$ in the power-containing sector, 24,173 kernels
in all. The worker is the three-dimensional seven-loop worker generalized to
sixteen propagators and 31-dimensional digital nets, with the extended-precision
first pass and promotions to binary128 and to 350-digit GMP arithmetic at the
unchanged thresholds $10^{-9}$ and $10^{-23}$. In its order-seven mode it
reproduces the unchanged seventh-order worker bit for bit on the two-loop
graph and all 3,511 kernels of loops three to seven, and it replays the archived
seventh-order production. Results do not depend on the thread count and
differ between processor types only in the last bits, at most
$2.0\times10^{-11}$ in normalized units between AMD and Intel processors;
every production job ran on one processor model, AMD EPYC 7543.

\paragraph{Low-loop inputs.}
The two-loop input gains
\begin{equation}
 \widehat b_{2,6}(3)=-29.258124480916806558671227221989158045345\ldots,
\end{equation}
from \eqref{gff:d3six:twoLoop}, now through $z^6$; its first six entries are those
of Section~\ref{gff:app:C7d3}. Two Cauchy contours at 65 digits
agree to $2.7\times10^{-56}$, and a 95-digit reference contour confirms the
stored entries to $1.7\times10^{-57}$. Every printed Gamma jet, 42 in all, is
the correctly rounded value of two independent Cauchy contours at 100 digits.
The single three-loop power graph has the closed form
$\widehat b_3^{\rm p}(0)=-16\pi/15$; the production gives
$-3.3509158\pm0.0000894$, $+1.30$ sampling standard errors from it.

\paragraph{Assembly gates.}
At the production means the top slots enter with
$\partial C_8/\partial(\widehat b_{L,8-L},\widehat c_{L,8-L})=(-1,3)/3^{L+1}$
to $4\times10^{-16}$ relative, so the eight-loop constants enter only as
$C_8(3)\supset(3\widehat c_{8,0}-\widehat b_{8,0})/19683$; loop seven enters
through $(-16A_3,42A_3)/19683$ with $A_3=-\widehat b_{2,0}/6=1.6309960686$, to
$8\times10^{-16}$; and the response of $C_8$ to the jets of loops five to eight
is linear to $3\times10^{-13}$. These three gates are hard; all pass.

\subsection{Pilot and design}\label{gff:c8d3:design}
A staged pilot fixed the design and never enters the estimate; its seeds are
excluded from the production. A screen at 4,096 points per replicate
(2,732 eight-loop graphs) selected the heavy-tail candidates. A full pilot
integrated all 24,173 kernels in both campaigns with 8 replicates of 4,096
points per graph, 16 replicates for the power-containing eight-loop graphs,
and gave $C_8(3)=1619\pm16$. A confirmation pilot integrated 710 selected
eight-loop graphs at 32,768 points per replicate, within a cap of a quarter of
the full pilot's eight-loop cost. The pilot used 422 core-hours.

Per-graph counts follow the Neyman allocation of Section~\ref{gff:c8d2:design}
on a power-of-two lattice with at least 4,096 points per graph and 128
replicates in every campaign, sector and loop. For the 710 confirmed graphs
the per-point standard deviation is the larger of the full-pilot and
confirmation estimates. After the pilot the author chose a target standard
error of 1, about $6\times10^{-4}$ of $C_8(3)$, and the cheapest of three
power-of-two designs, with a predicted standard error of $0.87$ ($1.31$ with
the safety factor $1.5$ of Section~\ref{gff:c8d2:design}) and an estimated cost
of 15,492 core-hours, to be topped up if the realized error exceeded 1. Two independent campaigns $a$ and $b$ use importance
rates multiplied by $0.8$ and $0.7$, as in Section~\ref{gff:app:C7d3}, and 52- and
60-bit digital coordinates, with group seeds derived from a fresh base seed.
The design contains 365,441,318,912 kernel samples.

\subsection{The 2026 production}\label{gff:c8d3:production}
The production ran on the symmetry cluster from 9 to 12 October 2026, in
lanes of eight threads on whole AMD EPYC 7543 nodes, at most eight at a time,
on ten nodes in all. All 1,941 work indices and 25,714 checkpoint blocks
completed. The lanes used 14,970 core-hours, 0.97 of the design's estimate,
and the merges 106; with the plan job and the rate probes of
Section~\ref{gff:c8d3:tails} the production used 15,102 core-hours. Every merge, the
assembly and an independent re-verification of the result record used the
package's own tools. The two campaigns and the two half streams give
\begin{center}\begin{tabular}{@{}lrr@{}}\toprule
Estimate & $C_8(3)$ & Sampling SE\\\midrule
Campaign $a$ & $1604.684$ & $1.239$\\
Campaign $b$ & $1600.527$ & $1.445$\\
First half stream & $1602.651$ & \\
Second half stream & $1602.559$ & \\
Combined & $1602.605$ & $0.952$\\\bottomrule
\end{tabular}\end{center}
The campaigns differ by 2.18 combined standard errors. Linearized with the
assembly Jacobian, the difference $4.16\pm1.90$ is spread over loops and
sectors, $+1.36$ and $+1.39$ from the power-free and power-containing
eight-loop graphs, $+0.90$ from the power-free seven-loop graphs and
$+0.19\pm0.10$ from the power-free five-loop graphs, with no term above two of
its own standard errors. Over the eight-loop graphs the per-graph campaign
differences have $\chi^2$ per graph $0.994$ (16,064 power-free) and $1.010$
(4,336 power-containing), the ten largest carry $+0.19$ and $-0.07$, and no
class of tail index carries the difference. No concentrated discrepancy is
therefore seen. This is consistent with a fluctuation, but it does not
exclude a small coherent effect of the importance rates spread over many
graphs, and it does not by itself establish that the per-graph standard
errors are calibrated. The error is dominated by the eight- and seven-loop
graphs of both sectors:
\begin{center}\begin{tabular}{@{}lrrrrrr@{}}\toprule
Sector & $L=3$ & $L=4$ & $L=5$ & $L=6$ & $L=7$ & $L=8$\\\midrule
Power-free & 0.001 & 0.010 & 0.048 & 0.160 & 0.362 & 0.609\\
Power & 0.000 & 0.003 & 0.018 & 0.079 & 0.220 & 0.566\\\bottomrule
\end{tabular}\end{center}
Entries are in units of $C_8$ and combine in quadrature. The realized standard
error is 1.09 times the design prediction.

\subsection{Tail diagnostics}\label{gff:c8d3:tails}
The fixed-$k$ ($k=16$) Hill estimate $\widehat\xi$ of Section~\ref{gff:c8d2:tails}
listed 10,034 eight-loop graphs with two-campaign mean $\widehat\xi>0.44$ at
production length. No graph carries more than 1.65\% of the variance of its
group, and none was flagged above 3\%. The listed graphs carry 62\% of the
power-free and 99.5\% of the power-containing eight-loop variance, so the
finiteness of their variance matters for the quoted error. The exact weight
bound of Section~\ref{gff:c8d2:tails} does not carry over, because the
three-dimensional kernel contains $U_F^{3/2}$. Two empirical tests replace it.

First, the importance rates were re-derived on 48 graphs: those with the
largest variance shares, representatives of the double-lift classes, two
controls and the 16 power-containing eight-loop graphs whose $\widehat\xi$
did not fall in the pilot. No asymptotic rate deficit was found. Five apparent deficits were
crossovers of the integrand from $t^{1.5}$ to $t^{1.0}$, with a zero of the
integrand in between, and the asymptotic exponent $1.0$ equals the table
rate. This is a probe at one base point, not a certificate.

Second, a true heavy tail keeps $\widehat\xi$ constant as the number $N$ of
points per replicate grows, while bounded, peaked weights make it fall. For
the 441 graphs confirmed in the pilot as possible infinite-variance cases
(246 power-free, 195 power-containing), $\widehat\xi$ fell from $N=4{,}096$
to $N=32{,}768$ as about $N^{-0.23}$ and $N^{-0.18}$. The production, with
independent seeds, mostly at intermediate lengths $N=8{,}192$ and
$16{,}384$, reproduces this dependence: the medians of $\widehat\xi$ lie at
0.996 to 1.031 of the falling trend interpolated between the pilot lengths,
13 to 38\% above the pilot's value at $N=32{,}768$, which a true heavy tail
would keep. The 13 graph-campaigns run beyond the pilot's range, at
$N=65{,}536$, continue to fall, to medians of 0.82 to 0.87 of their
$N=32{,}768$ values, at 0.94 to 1.02 of the extrapolated trend. At matched
sample length the campaign shift of $\widehat\xi$ in the power-free sector is
positive, $+0.03$ to $+0.05$, opposite in sign to the shift a rate-limited tail
would produce; in the power-containing sector it is inconclusive. The data
therefore behave as bounded, peaked weights would. This is empirical evidence, not a bound.

\subsection{Checks}\label{gff:c8d3:checks}
The same samples give lower-order coefficients, compared with the
special-function values of Section~\ref{gff:sec:guide} and with the productions of
Sections~\ref{gff:prec:d3production2026} and~\ref{gff:app:C7d3}, with both
uncertainties combined:
\begin{center}\begin{tabular}{@{}rrrr@{}}\toprule
Order & Estimate & Sampling SE & Pull\\\midrule
3 & $-0.3535809354$ & $1.38\times10^{-6}$ & $+0.43$\\
4 & $1.2280445577$ & $3.88\times10^{-5}$ & $-0.49$\\
5 & $-5.8499922821$ & $4.86\times10^{-4}$ & $+0.74$\\
6 & $33.5763740$ & $0.0050853$ & $+0.19$\\
7 & $-220.3847628$ & $0.0528950$ & $-0.49$\\\bottomrule
\end{tabular}\end{center}
The reference values are checks only; no known coefficient is imposed and no
power correction is forced to vanish. A pull above three in magnitude would
have been reported without changing $C_8(3)$; none occurred. The measured
power corrections at orders four to seven are
$(8.6\pm7.3)\times10^{-6}$, $(9.2\pm13.1)\times10^{-5}$,
$(3.3\pm19.5)\times10^{-4}$ and $0.009\pm0.024$; at order three the difference
is at the level of rounding. These are consistent with zero but do not prove a
cancellation at eighth order, where the correction is $0.61\pm0.61$. A
bootstrap of 4,096 replicate resamples gives standard error $0.940$, against
$0.952$ from covariance propagation. These checks share the graph and
renormalization framework and do not validate it.

\subsection{An independent code}\label{gff:c8d3:independent}
The eighth-order code of this subsection, \path{LRI_C8_D3_OPT/}, and its
derivation, given below, were written by OpenAI GPT-6 astra (model stated
by the author), which was given the seventh-order package of
Section~\ref{gff:app:C7d3} and an earlier version of this note; the cluster port
and the production were done within this project. The code extends that
package independently of the package used above. It
regenerates the graph inputs at every loop, restores the seven-loop classes
excluded at seventh order, extends the integrator and the regulator jets to
eight loops on its own, derives its own low-loop inputs, and solves for the
eighth-order fixed point with its own assembly; it integrates the same 24,173
kernels. It shares with the calculation above the three-dimensional mixed
forest kernel and the precision policy of the seventh-order integrator, the
importance sampler with rates multiplied by $0.8$ and $0.7$, the digital nets
and the finite renormalization scheme. An error in those shared parts would
not be detected.

\paragraph{Graph inputs.}
The code enumerates its own classes of the graph family of
Section~\ref{gff:app:C7d2}: connected one-particle-irreducible multigraphs
without self-lines, with $L+1$ vertices, $2L$ individually labeled lines and
at most four lines at each vertex. Classes are identified by a canonical form
from partition refinement, keyed by multiplicity strings, and their counts
agree at every loop with those of Section~\ref{gff:c8d2:changes}. An independent
nauty computation on two-color incidence graphs (vertices and labeled lines)
confirms all 26,540 eight-loop classes and their automorphism orders, and
reproduces the archived lower-loop inventories with their exact weights. At
eight loops the class weights sum to $\sum_Gw_A=49{,}985{,}915/768$ and
$\sum_Gw_B=33{,}994{,}099/6{,}912$, the zero-dimensional four-point and
mass-insertion values of \eqref{gff:eq:zeroDimSums} in class-weight form.

Its census of exact zeros and channel exclusions is that of
Sections~\ref{gff:c8d2:changes} and~\ref{gff:c8d2:independent}, which also explain
why its 4,851 eight-loop zeros differ from the 3,908 of
Section~\ref{gff:c8f:changes}; it integrates $26{,}540-4{,}851-1{,}289=20{,}400$
eight-loop kernels, all exclusions being exact. At seven loops the code
restores the 143 power-free and 119 power-containing nonzero classes with
$3w_B=w_A$ that Section~\ref{gff:app:C7d3} omits, giving 2,441 and 709 kernels.
Its builders reproduce all 3,511 payloads of the seventh-order calculation
(loops three to seven), that is, line masks, denominator lists, integer
numerators and importance rates, byte for byte after the header line,
whose weights are checked against the exact rational values. A coverage
record binds the census to the active inventory by its SHA-256 hash, and the
assembly refuses to run without it.

\paragraph{Non-induced divergent subgraphs.}
A subgraph is induced if it contains every line of $G$ between its vertices.
The code builds its forests from induced divergent subgraphs only, as do the
production builders whose maps the edge-subgraph enumerations of
Sections~\ref{gff:app:C7d2} and~\ref{gff:c8d2:changes} confirm. With parallel lines
a graph can also have proper divergent subgraphs that are not induced. The
forest terms that contain them cancel exactly, by the following argument. It
uses only the graph family and the structure of the forest sum
\eqref{gff:c7d2:kernel} with any fixed power of $U_F$ in its denominators, so it
holds in every dimension in which this sum is used, at $d=2$ and, with
$U_F^{3/2}$, at $d=3$.

Let $W$ be a proper nonempty subset of the vertices of $G$, $E(W)$ the set of
lines with both ends in $W$ and $\partial W$ the set of lines with exactly
one end in $W$. Since $G$ is connected and has no bridge, $|\partial W|\ge2$,
and the degree bound gives
\begin{equation}
 2|E(W)|+|\partial W|=\sum_{u\in W}\deg_G(u)\le4|W|,
 \qquad\text{hence}\qquad |E(W)|\le2|W|-1.
 \label{gff:c8d3:edgeBound}
\end{equation}
A proper divergent subgraph does not contain every vertex: on all $L+1$
vertices a logarithmic subgraph would need all $2L$ lines, and a power
subgraph one line more than $G$ has. A power subgraph on $v$ vertices has
$2v-1$ lines, the maximum allowed by \eqref{gff:c8d3:edgeBound}, so every proper
power subgraph is induced. A proper divergent subgraph $\gamma$ that is not
induced is therefore logarithmic, with $2v-2$ lines on its $v$ vertices, and
misses exactly one line $\ell_0$ of its induced closure $\bar\gamma$, the
subgraph of all $2v-1$ lines between the vertices of $\gamma$. Adding a line
to the one-particle-irreducible $\gamma$ keeps it one-particle irreducible,
so $\bar\gamma$ is a proper power subgraph, with loop number
$h_{\bar\gamma}=v$.

Let $m$ be a complement-of-tree monomial kept by $\gamma$, that is
$|m\cap E_\gamma|=h_\gamma=v-1$. The complementary spanning tree
then has $v-1$ lines in $\gamma$, which form a spanning tree of the vertices
of $\gamma$, so it cannot also contain $\ell_0$, which would close a cycle.
Hence $\ell_0\in m$ and $|m\cap E_{\bar\gamma}|=v=h_{\bar\gamma}$: the
constraint of $\bar\gamma$ follows from that of $\gamma$, and adding
$\bar\gamma$ to a forest that contains $\gamma$ does not change $U_F$.

Now take a compatible forest $F$ with at least one non-induced logarithmic
member, let $\gamma$ be an inclusion-maximal one, chosen by a fixed rule,
and toggle $\bar\gamma$: remove it from $F$ if it is a member and add it
otherwise. The result is again a compatible forest. By the definition of
compatibility before \eqref{gff:c7d2:kernel}, two members are either nested or
edge-disjoint and meeting in at most one vertex. Members inside $\gamma$
lie inside $\bar\gamma$. Members not nested with $\gamma$ meet its vertex set,
which is that of $\bar\gamma$, in at most one vertex, so they share no line
with $\bar\gamma$. Members containing $\gamma$ are induced, by the
maximality of $\gamma$ if they are logarithmic and by
\eqref{gff:c8d3:edgeBound} if they are power subgraphs, so they contain
$\bar\gamma$. The toggle does not change the set of non-induced logarithmic
members, so it selects the same $\gamma$ again and is an involution without
fixed points. It leaves $U_F$ unchanged, by the implied constraint, and the
root numerator $\sum_{e\in E_R^{\log}(F)}x_e$, because $\bar\gamma$ is a
power member and only logarithmic members remove lines from it; it reverses
the sign $(-1)^{|F|}$. The forest terms with a non-induced logarithmic member
therefore cancel in pairs, and the sum over forests of induced subgraphs
equals the full forest sum. $\square$

By the argument above, the closure of a non-induced divergent subgraph is a
proper power subgraph, so power-free graphs have no such subgraph. The code
also checks the cancellation by direct enumeration, which therefore needs
only the power-containing classes. Its all-edge audit visits
every proper subset of labeled lines of all 8,532 power-containing classes of
loops three to eight, 492,490,392 subsets, tests one-particle irreducibility,
builds the 11,123,924 unreduced mixed forests and reproduces every grouped
integer map, the zero maps and the maps with a vanishing top channel
included. At eight loops alone it covers the 7,231 power-containing classes,
with 40,679 non-induced logarithmic subgraphs and 10,451,732 forests.

\paragraph{Assembly and analytic inputs.}
The truncated fixed point is obtained by nine iterations of
$g\leftarrow\frac13\bigl[\epsilon-\sum_{L=2}^8\widehat b_L(\epsilon)g^L\bigr]
\bmod\epsilon^9$ from $g=0$, and in 64 exact rational cases,
regulator-dependent and constant-coefficient ones included, it agrees
exactly with the Lagrange formula \eqref{gff:d3six:lagrange} carried to eight
loops. The response to the jets of loops five to eight has a closed form.
Hold the jets of loops one to four fixed and put
\begin{equation}
 \Phi_{\le4}(g,z)=3g+\sum_{L=2}^{4}\widehat b_L(z)g^L,\qquad
 \gamma_{r,\le4}(g,z)=g+\sum_{L=2}^{4}\widehat c_L(z)g^L,\qquad
 \Phi_{\le4}\bigl(g_{\le4}(\epsilon),\epsilon\bigr)=\epsilon,
 \label{gff:c8d3:truncatedRG}
\end{equation}
so that $g_{\le4}(\epsilon)=\epsilon/3+O(\epsilon^2)$ is the fixed point of
the functions \eqref{gff:c8d2:rg} truncated after four loops. A change of the
jets of loops five to eight changes the exponent by
\begin{equation}
 \delta\Delta_{\mathcal E}(3,\epsilon)=\sum_{L=5}^{8}g_{\le4}^L
 \Bigl[\delta\widehat c_L(\epsilon)
 -\frac{\partial_g\gamma_{r,\le4}(g_{\le4},\epsilon)}
 {\partial_g\Phi_{\le4}(g_{\le4},\epsilon)}\,
 \delta\widehat b_L(\epsilon)\Bigr]\pmod{\epsilon^9},
 \label{gff:c8d3:variation}
\end{equation}
with $g_{\le4}=g_{\le4}(\epsilon)$. This first variation is exact through
order eight, because two insertions of loops five or higher first contribute
at order $5+5-1=9$. With $\widehat b_L(z)=\sum_j\widehat b_{L,j}z^j$ and
likewise for $\widehat c_L$, the Jacobian of $C_8$ is therefore, for
$L=5,\ldots,8$ and $0\le j\le8-L$, whatever the jets of loops five to eight,
\begin{equation}
 \frac{\partial C_8}{\partial\widehat c_{L,j}}=[\epsilon^{8-j}]\,g_{\le4}^L,
 \qquad
 \frac{\partial C_8}{\partial\widehat b_{L,j}}=-[\epsilon^{8-j}]\,
 \frac{\partial_g\gamma_{r,\le4}(g_{\le4},\epsilon)}
 {\partial_g\Phi_{\le4}(g_{\le4},\epsilon)}\,g_{\le4}^L.
 \label{gff:c8d3:jacobian}
\end{equation}
For $j=8-L$ it reduces to the top-slot rule $(-1,3)/3^{L+1}$ of
Section~\ref{gff:c8d3:changes}, so the eight-loop constants enter only as
$C_8(3)\supset(3\widehat c_{8,0}-\widehat b_{8,0})/19683$; for
$(L,j)=(7,0)$, with $A_3=-\widehat b_{2,0}/6$ (Section~\ref{gff:c8d3:changes})
and $\widehat c_{2,0}=\widehat b_{2,0}/6$, it gives
$(-16A_3,42A_3)/19683$. A seven-loop class with $3w_B=w_A$ has
$\delta\widehat b_{7,j}=3\delta\widehat c_{7,j}$. Its top slot cancels, but
through the mixing with the two-loop input its constant terms give
$\delta C_8=-2A_3\,\delta\widehat b_{7,0}/19683
=\widehat b_{2,0}\,\delta\widehat b_{7,0}/59049$, which is why the code
restores these classes.

The code derives its own two-loop and Gamma inputs. At 80 digits, Cauchy
contours of radii $0.04$ and $0.07$ with 80 and 96 nodes agree on all seven
two-loop entries to $3.5\times10^{-72}$; the six entries shared with
Section~\ref{gff:app:C7d3} agree with that calculation's input to
$4.5\times10^{-59}$; and a 112-node Cauchy extraction of the logarithm of the
Gamma product checks all 42 Gamma-jet coefficients (loops two to eight,
orders one to six) to $1.2\times10^{-73}$. These are high-precision
cross-checks, not interval bounds. By direct comparison of the files
\path{inputs/two_loop.json} and \path{inputs/gamma_d3_8.txt} of the two
packages, the two-loop entries and the Gamma table differ from those of the
calculation above by at most $1.7\times10^{-57}$, the precision of the
stored two-loop entries there (Section~\ref{gff:c8d3:changes}), and
$4.1\times10^{-58}$.

\paragraph{Worker.}
The production used the code's optimized worker, with the 350-digit GMP
tier. Its one change to the arithmetic is in the binary128 additions inside
the positive denominator polynomials: the 113-bit integer significands are
aligned with guard, round and sticky bits and rounded to nearest, ties to
even, so each sum is the correctly rounded one; zeros, subnormal and negative
operands, exceptional values and possible overflow go to the compiler's
addition. Signed forest sums, the addition order, the promotion thresholds
and the sampling are unchanged. A test of 12,051,840 additions agrees bit for
bit (NaN cases by outcome); on 179 graphs at 40,448 fixed points the
optimized and the unoptimized worker give identical coefficient vectors,
maximum contributions and promotion counts; 2,432 evaluations at shared
points reproduce the jets and promotion counts of the seventh-order worker of
Section~\ref{gff:app:C7d3} exactly; and 48 points of 12 eight-loop graphs agree with a 450-digit
reference within the tolerances of their precision tiers. Three builds on
Expanse gave one binary hash and passed all 14 build and test steps.

\paragraph{Pilot and allocation.}
The code was delivered with a pilot on its development host: both campaigns
with 8 replicates each and 32,768, 8,192, 2,048, 512, 64 and 16 points per
graph and replicate at loops three to eight, 21,342,208 evaluations in all.
It gave $C_8(3)=1514.6\pm161.7$ (sampling standard error), an exploratory
estimate. Its per-graph variances and timings fixed the allocation of the
production by the code's Neyman rule: the points per graph are proportional
to the square root of the ratio of per-point variance to per-point cost,
rounded up to a power of two, at least 4,096, and scaled to the budget. The
replay reproduces the allocation file byte for byte. The pilot is not part
of the estimate.

\paragraph{Production.}
The code was adapted to the Expanse cluster without a change to its
numerical files. With a fresh base seed, an allocation of about 2,000
core-hours of its development host, 16 replicates per campaign and at least
4,096 points per graph, it ran on 10 and 11 October 2026 on six whole AMD
EPYC 7742 nodes, the last 5 of its 1,994 jobs in a separate one-node job,
with $3.7\times10^{10}$ evaluations and about 9,800 service units in all. It gives
\begin{equation}
 C_8(3)=1600.98\pm3.97\quad\text{(sampling standard error)},
\end{equation}
$-0.40$ combined standard errors from the estimator. Its campaigns,
$1603.53\pm4.91$ and $1598.42\pm6.24$, differ by 0.64 combined standard
errors; its lower-order coefficients agree with the references within 1.73
standard errors (pulls $-1.41$, $-1.55$, $-0.10$, $+1.73$ and $-0.21$ at
orders three to seven; the squared pulls sum to 7.40 for these five checks),
and its sixth- and seventh-order by-products are
$33.599\pm0.014$ and $-220.40\pm0.19$. Its power-free part is
$1602.61\pm2.96$ and its power correction $-1.63\pm2.64$, with covariance
$1.9\times10^{-4}$; they differ from those of Section~\ref{gff:c8d3:result} by
$+0.20$ and $-0.82$ combined standard errors. The half streams give
$1599.79$ and $1602.17$, deviations of $-1.19$ and $+1.19$ from the
estimator, and a bootstrap of 4,096 replicate resamples gives standard error
$3.78$. The eight-loop graphs carry 42.1\% (power-free) and 31.5\%
(power-containing) of the variance, the seven-loop graphs 12.2\% and 12.6\%,
and the six-loop graphs 1.4\%. The power null rows at orders four to seven
are $(-4.66\pm2.22)\times10^{-5}$, $(-1.91\pm4.57)\times10^{-4}$,
$(-1.21\pm5.62)\times10^{-3}$ and $0.072\pm0.126$, with pulls $-2.10$,
$-0.42$, $-0.22$ and $+0.57$, whose squares sum to 4.95 for these four null
rows.

The quoted error is the covariance of the replicates propagated through the
fixed point. The code assigns no allowance, and its result record marks the
reportable precision as not established. There were no tail diagnostics like
those of Section~\ref{gff:c8d3:tails} and no convergence test beyond the
campaign and half-stream splits, and neither the delivered science code nor
this number had an independent review: the reviews of the port covered only
its cluster layer. The independent estimate is not combined with the
reported value.

\subsection{Records}\label{gff:c8d3:records}
The production result, the frozen plan and the merged replicate data that
reproduce the estimator are in \path{GFF/results/C8_3/}; the per-graph
replicate tables are hosted in \path{GFF/data/C8_3/}. The verification, the
job records and the analyses of Sections~\ref{gff:c8d3:production}
and~\ref{gff:c8d3:tails} are in
\path{LRI_C8_D3_HPC/benchmarks/symmetry_2026-10-c8/production/prod1/}.
The pilot report is in \path{pilot/evidence/} and the rate probes in
\path{production/margin_*} of the same \path{symmetry_2026-10-c8/} folder, and
the independent code's production is in \path{GFF/results/C8_3_check/}.
\path{LRI_C8_D3_HPC/public.mk} recomputes every number of
Sections~\ref{gff:c8d3:result}--\ref{gff:c8d3:checks} from these records. The
numbers of Section~\ref{gff:c8d3:independent} come from the independent package.
\path{LRI_C8_D3_OPT/public.mk} recomputes those of its production from
\path{GFF/results/C8_3_check/}, except the half streams and the bootstrap
error, which need the replicate matrices and are read from its result record.
The census, audit, input, worker and pilot checks are the package's delivered
evidence (\path{LRI_C8_D3_OPT/verification/}, \path{docs/} and
\path{README.md}), its builds are summarized in its \path{public/REVIEWS.md},
and the differences between the input files follow from the files
\path{inputs/two_loop.json} and \path{inputs/gamma_d3_8.txt} of the two
packages. The complete run tree of the production of
Section~\ref{gff:c8d3:production} is archived in two copies
(Section~\ref{gff:cn:reproduction}); the native outputs of the independent
production are not archived.

\section{Reproduction and data dependencies}\label{gff:cn:reproduction}\label{gff:sec:softwarefinal}
The code and records behind this note are in the folder \path{GFF/} of the
public repository of this work. Its folder \path{GFF/code/} holds the
packages named in the table below, each at the commit that produced the quoted
results; package paths in this note, in the text and in the table, are
relative to \path{GFF/code/}. The records from which the quoted numbers are
recomputed are in \path{GFF/data/}, and the published outputs in
\path{GFF/results/}; these paths are written in full. Every package has a
file \path{public.mk}: \texttt{make -f public.mk test} runs its tests, and
\texttt{make -f public.mk replay} recomputes from that package's shipped records
the numbers this note quotes for the results the package produced, and compares
them with the quoted digits. Numbers of earlier calculations that a package
keeps in \path{reference/}, checks made when the separate notes were merged into
this note, and construction checks whose records are hosted are not all
replayed; the rows below say which. Large raw
files (checkpoints, graph packs, payload archives) are hosted outside the
repository and fetched with \texttt{make -f public.mk fetch-data}; the
package's \path{DATA.md} says which target needs them, and its review
summary, \path{REVIEWS.md}, says which reviews it received. The present note is a
mathematical condensation, not a production integration campaign. A
package's finite parts and graph ordering must be used with that same
package's code; files with similar names from different schemes are not
interchangeable. Packages of the earlier research archive
\texttt{LRI\_complete\_research\_bundle.zip} that hold only superseded
estimates are not part of the repository. The separate \TeX{} notes that
came with several packages are merged into this note and are left out of the
repository: their derivations, methods, results and caveats are in the
sections cited in the table, and Section~\ref{gff:cn:provenance} says where each
block came from. Markdown records of derivation steps that ship inside a
package stay there as its records; for example, the delivered stage scripts
of \path{C6_research_note/} read or write such files under
\path{code/stages/} and \path{inputs/source_snapshots/}. No step of this note
cites them; the arguments are written out here, for $C_6(1)$ in
Appendix~\ref{gff:app:C6derivation}.

\paragraph{Provenance of the merged material.}\label{gff:cn:provenance}
The table lists the blocks that were merged into this note from separate
notes and records on 16 October 2026, with their source and author. None of
these sources names its author: their bylines read ``Research calculation
record'' or ``Calculation record'', or are empty. The entry ``OpenAI GPT-6
astra'' records the author's statement that the source was written by that
model (for the research bundle and the delivered packages, 13 October 2026;
for \path{C6_research_note/}, 16 October 2026); ``author not named'' marks a
source for which there is no such statement. Claude Code sessions wrote the
merged text itself and adapted it to the notation of this note.
{\small
\begin{longtable}{@{}>{\raggedright\arraybackslash}p{0.33\textwidth}>{\raggedright\arraybackslash}p{0.30\textwidth}>{\raggedright\arraybackslash}p{0.31\textwidth}@{}}
\toprule Block & Source & Author\\\midrule\endhead
Appendix~\ref{gff:app:C6derivation} (derivation of the $C_6(1)$ representation) & \path{C6_research_note/}: its note and input records & OpenAI GPT-6 astra\\
Section~\ref{gff:app:C7complete}: method, design and checks of the seventh order & \path{LRI_C7_HPC/reference/}: the two seventh-order notes & OpenAI GPT-6 astra\\
Sections~\ref{gff:sec:guide} (lattice model), \ref{gff:c1:notation} (the caveat on repeated parameters), \ref{gff:sec:graphs} (zero-dimensional benchmark), \ref{gff:app:KpsiProof} (checks of $\KK$), \ref{gff:sec:oldmasters} (decimal value of $\LL$), \ref{gff:sec:P5complete} (100-digit value of $P_5$), \ref{gff:sec:checks} & the earlier complete note in \path{LRI_C7_HPC/reference/} & OpenAI GPT-6 astra; the resummation of Section~\ref{gff:cnm:latticeComparison} recomputed at the merge by a Claude Code session\\
Sections~\ref{gff:ope:d1}, \ref{gff:ope:d1functional}; the $d=1$ checks of Section~\ref{gff:ope:checks} & \path{lri_ope_reproducibility/} and \cite{gff:GhoshSuperSum} & written at the merge by a Claude Code session\\
Section~\ref{gff:gs:masterTable}, values & the note's closed forms; \path{LRI_C5_d3/} records & a Claude Code session\\
Section~\ref{gff:gs:masterTable}, spectral route to $d=2$ & \path{lri_d3_reduction_update/}: its complete note & OpenAI GPT-6 astra\\
Section~\ref{gff:gc:primitive} (double residue series of $P_5(d)$) & \path{lri_d3_reduction_update/}: the general-dimensional addition & OpenAI GPT-6 astra; steps the source left implicit written out at the merge by a Claude Code session\\
\eqref{gff:fp:C5target} and the measure-pole bound & \path{lri_d3_reduction_update/}: the periodization addition & OpenAI GPT-6 astra\\
Section~\ref{gff:app:threeDimensionalReduction}: opening, count table, Section~\ref{gff:td:complexVariant}, values and dictionary of Section~\ref{gff:d3:values}, catalogue search & \path{lri_d3_reduction_update/}: the three-dimensional addition & OpenAI GPT-6 astra\\
Section~\ref{gff:rt:kummerDiagnostic}, rank-two checks, open problem & \path{lri_d3_matrix_update/}: its complete note & OpenAI GPT-6 astra\\
Section~\ref{gff:app:threeDimensionalReduction}: the link to $\mathfrak t_3(3;\frac32)$ & derived at the merge & a Claude Code session\\
Section~\ref{gff:gdr:halfShift} (half-shift identities at $d=3$) & \path{LRI_C5_d3/} & a Claude Code session (written out at the merge)\\
Section~\ref{gff:gdk:section}, including the growth proof and the continuation & \path{LRI_C5_general_d/K_general/} & a Claude Code session\\
\eqref{gff:gdk:Sconstant} & \path{LRI_C5_general_d/K_general/}; constant derived at the merge & a Claude Code session\\
Section~\ref{gff:done:samplingControls}, the 300-digit fallback & \path{lri_C6_d3_complete_update/} & OpenAI GPT-6 astra\\
Section~\ref{gff:meth:crossdimension}, $d=3$ rows of the method comparison & \path{LRI_C6_D2_HPC/reference/}: the numerical-methods note & OpenAI GPT-6 astra; the explanation of the failed control added at the merge by a Claude Code session\\
Section~\ref{gff:c8d3:independent} & \path{LRI_C8_D3_OPT/}: its derivation and calculation notes & OpenAI GPT-6 astra; cluster port and production by a Claude Code session\\
Part~V, general $d$, $d=3$ and the bootstrap comparison & \path{lri_ope_reproducibility/}: its note & author not named; supplied by the author on 16 October 2026; the checks of Section~\ref{gff:ope:checks} by Claude Code sessions\\
Section~\ref{gff:gc:d2closedForms}: numerical discovery of $\LL_2$, $\KK_2$ and the single-sum relations & \path{LRI_C5_d2_simplify/} & a Claude Code session\\
Sections~\ref{gff:gc:L2proof} (first and second proofs of $\LL_2$) and \ref{gff:gc:K2proof} (proof of $\KK_2$); in Section~\ref{gff:mgk:section} the sharper estimate, the per-spin identity, the second and third single-sum evaluations, the mixed-trace corollary and the Riemann bilinear proof & \path{LRI_C5_d2_proof/} & a Claude Code session with prover and reviewer subagents\\
Sections~\ref{gff:mgp:P4second}, \ref{gff:mgp:halfStrip}--\ref{gff:mgp:fourier}; the regularized spectral route of $\KK_2$ in Section~\ref{gff:mgk:section} & \path{LRI_C5_d2_simplify/proof/} & a Claude Code session with prover, verifier and synthesizer subagents\\
Section~\ref{gff:gdr:section}: the generic-dimensional reduction of $\LL_d$, $Z_d$ and its proofs & \path{LRI_C5_general_d/} & a Claude Code session; the explicit constants of \eqref{gff:mgf:gammaBounds}, \eqref{gff:mgf:multiplicityBound}, \eqref{gff:mgf:latticeSums} come from the growth proof of Section~\ref{gff:gdk:growthProof}\\
Sections~\ref{gff:cn:finiteKernels}--\ref{gff:cn:finiteAssembly}: the blocks merged from the earlier complete note & \path{LRI_C6_D2_HPC/reference/}: the complete note & OpenAI GPT-6 astra; one closed-form value and one derivation added at the merge by a Claude Code session\\
Checks merged from the same complete note into Sections~\ref{gff:mgc:higherPoles} (what the higher-pole identities cannot detect), \ref{gff:mgc:gdChecks}, \ref{gff:mgc:masterControls} (controls of the master evaluations), \ref{gff:mgc:gcChecks} (checks of the finite multivariable construction) and \ref{gff:mgc:fpChecks} & \path{LRI_C6_D2_HPC/reference/}: the complete note & OpenAI GPT-6 astra\\
Sections~\ref{gff:mgs:benchmarks}--\ref{gff:mgs:selection}: method blocks & \path{LRI_C6_D2_HPC/reference/}: the methods note & OpenAI GPT-6 astra; the variant descriptions read from the code, the linear-control formula and one proof added at the merge by a Claude Code session\\
Section~\ref{gff:app:C7d2}: merged blocks & \path{LRI_C7_D2_HPC/reference/}: the method and progress notes & OpenAI GPT-6 astra; the three-loop check, the two-loop entries, the primitive classes and the variance decomposition written out or rechecked at the merge by a Claude Code session\\
Section~\ref{gff:c8d2:independent}: merged blocks & \path{LRI_C8_D2_HPC_optimized/}: the calculation note and results fragment & OpenAI GPT-6 astra; the 65- and 70-digit comparison of the two packages' inputs made at the merge by a Claude Code session\\
Section~\ref{gff:ope:d2}, the $d=2$ rows of Section~\ref{gff:ope:bootstrap} and \eqref{gff:ope:b5MS} & written for this note from \path{lri_ope_reproducibility/} & a Claude Code session; checks in \path{ope_checks/d2/}\\
\bottomrule
\end{longtable}}

\begin{longtable}{@{}>{\raggedright\arraybackslash}p{0.30\textwidth}>{\raggedright\arraybackslash}p{0.65\textwidth}@{}}
\toprule Result or dependency & Package and principal record\\\midrule\endhead
One-dimensional analytic $C_5$ &
\path{lri_numerical_resolution/}: graph audit in
\path{verification/records/resolution_audit.json}, with the raw five-loop and
lower-loop Laurent records in \path{verification/records/}, the residue
certificate in \path{verification/inputs/residue_certificate.json} and the
kernels that produced the records in \path{verification/source/}. The packages
\path{lri_psi_complete_update/} and
\path{lri_fixed_hypergeometric_update/} supply the checks of the
polygamma evaluation of $\mathcal K$ and of the five-value $P_5$, with the
records \path{verification/numerical_checks.json} and
\path{verification/fixed_hypergeometric_checks.json}. The quoted values are
collected in \path{GFF/results/closed_forms_1/C5_1_quoted.json}, and
\texttt{make -f public.mk replay} in each package folder recomputes them from
the shipped files in under a minute on one core.\\
General $C_0$--$C_5$ &
\path{lri_general_special_functions_update/} (with the general-dimensional
$C_0$--$C_3$, graph-inventory, tetrahedron and two-loop files of
Section~\ref{gff:app:generalDimension}), \path{lri_general_closed_update/},
\path{lri_general_periodization_update/}; master evaluators, finite arrays,
residue and assembly checks. The quoted values and formulas, with their
records, are in \path{GFF/results/closed_forms_general/C0_C5/expected.json};
\texttt{make -f public.mk replay} in each package recomputes them.\\
Closed forms of $\mathcal L_2$, $\mathcal K_2$ and $C_5(2)$ &
\path{LRI_C5_d2_proof/} (proof routes with their machine checks; the computations behind the
quoted checks in \path{review_*/}; reviews summarized in \path{REVIEWS.md}) and
\path{LRI_C5_d2_simplify/} (discovery, $70$-digit single sums \path{confirm_dps70.json},
scripts and logs of the three further proofs in \path{proof/}, whose $\mathcal L_2$ parts are
written out in Sections~\ref{gff:mgp:halfStrip}--\ref{gff:mgp:fourier}; reviews summarized in \path{REVIEWS.md}); the independent spectral evaluation is
\path{LRI_C5_d3/wf_d2check/d2check_results.json}. Expected values:
\path{GFF/results/closed_forms_2/closed_forms_2.json}; \texttt{make -f public.mk replay} in
either package recomputes every quoted number.\\
Generic-dimensional reduction of $\mathcal L_d$, $Z_d$ and $\mathcal K_d$ &
\path{LRI_C5_general_d/}: the code and validation records of the general-$b$ periodization (Section~\ref{gff:gdr:periodization}) in \path{foundation/}; the specialization code and records in
\path{specialize/}; records of the checks in
\path{K_general/note_fragment/growth/checks/}, \path{K_general/note_fragment/topmode/results/},
\path{specialize/results/}, \path{L/results/} (with the code of Section~\ref{gff:mgl:section} in \path{L/}), \path{Z/results/} and the
end-to-end $C_5(1)$, $C_5(2)$, $C_5(3)$ in
\path{review/joint/numerics-end-to-end/e2e_report.json}; the review records
\path{review/VERIFICATION.md}, \path{review/lemmaG-proof/} and the check code of
\path{review/} (\path{lemmaG-numeric/}, \path{topmode-numeric/},
\path{joint/numerics-end-to-end/}), with all reviews of both sections, including
those of \path{K_general/note_fragment/}, summarized in \path{REVIEWS.md} (the
review transcripts are not included); independent spectral references
\path{references/reference_values.json}; $100$-digit $\mathcal K_d$ in
\path{K_general/position_space/K_d_values_mb_D100.json}. \texttt{make -f public.mk replay} compares
with these records every quoted number of Section~\ref{gff:gdk:section} (expected values
\path{public/expected_K_d.json}) and the numbers of Section~\ref{gff:gdr:section} listed in
\path{GFF/results/closed_forms_general/reduction_d/expected.json}, and reruns the fast checks among
them. It does not compare the numbers of Sections~\ref{gff:mgf:nodes}, \ref{gff:mgl:section} and
\ref{gff:mgz:proofs}, nor those of the items of Section~\ref{gff:gdr:checks} on the general-$b$
periodization, the review probes, the product form in detail and the two-node formula; these were
added on 16 October 2026 and are read from the records named there. \texttt{make -f public.mk
replay-full} recomputes part of those records (\path{L/results/val50.json},
\path{L/results/val70.json}, \path{Z/results/zj*.json}, \path{Z/results/zc*.json},
\path{Z/results/compare_refs.log}, \path{Z/results/check_formulas.log}) and compares them with the
shipped records, not with the text; the records in \path{foundation/results/},
\path{Z/results/offdiag60.json} and the other files of \path{L/results/} are recomputed by no target.
The $d=3$ inputs are in \path{LRI_C5_d3/} (\path{wf_L3/}, \path{wf_K3/K3_value.json}), with the
quoted values in \path{GFF/results/closed_forms_3/}.\\
Three-dimensional specialization &
\path{lri_d3_reduction_update/}: \path{code/general_d3/check_exact.py} (residue algebra,
polynomial identities and rational Horn coefficients; \path{verification/general_d3/exact_algebra.json},
with the eight real-node Horn arrays in \path{horn_arrays.json} of the same folder),
\path{code/general_d3/verify_d3.py} (formula-side values at 55 and 90 digits, the 75-digit
determinant, mode and complex-node checks and the literal hypergeometric checks;
\path{period_values.json}, \path{identities.json}, \path{literal_hypergeometric.json}) and
\path{code/general_d3/check_spectral.py} (fresh real-axis Gamma kernels; \path{fresh_spectral.json};
the residuals are its differences from \path{period_values.json}, also listed in the delivered record
\path{comparison_summary.json}). \path{lri_d3_matrix_update/}:
\path{code/rank_two_checks.py} (the 1107 exact checks and the degree-90 interior sums;
\path{verification/rank_two_checks.json}), \path{code/recompute_d3_periods.py} (the 100-digit values;
\path{verification/fresh_primitive_values.json}) and \path{code/scalar_catalog_search.py} (the
catalogue search of Section~\ref{gff:d3:values}; \path{verification/scalar_catalog_search.json}, rerun by
\texttt{make -f public.mk search}). Section~\ref{gff:app:threeDimensionalReduction} quotes these
records. The quoted values and their tolerances are in
\path{GFF/results/closed_forms_3/d3_specialization_quoted.json}.\\
$C_6(1)$ &
\path{GFF/results/C6_1/C6_final.json}, with the degree sequence in
\path{assembled.json}, the rounding estimates in \path{rounding.json}, the
stopping decisions in \path{decisions/}, the pre-registered plan
\path{PLAN.md} and the report \path{REPORT.md} of the same folder. The
results-review checks are \path{C7_crosscheck.json} and
\path{lower_amplitude_precision/} in that folder. The code is
\path{LRI_C6_1D_HPC/} at git commit \texttt{cef6fa7}; the production
ran at \texttt{8680864}. The merged level records and the engine outputs are
in \path{GFF/data/C6_1/}. The target \texttt{replay} of \path{public.mk} in
the code folder regenerates these records and checks every quoted number.
The superseded Monte Carlo estimator is not included in the public
repository.\\
Analytic representation of $C_6(1)$ &
\path{C6_research_note/}: the reassembly \path{outputs/reassembly.json}, the
interval certificate in
\path{outputs/rerun/complete_refined/} and the task runner
\path{code/reproduce_final.py}, which runs the 22 tasks of the final method
(Section~\ref{gff:c6d:numerics}). The derivation is
Appendix~\ref{gff:app:C6derivation}. The target \texttt{replay} of
\path{public.mk} in that folder recomputes the quoted numbers from these records
and compares them with \path{GFF/results/C6_1/analytic/quoted_numbers.json}.
The independent review is
\path{LRI_C6_closed_form_review/REVIEW.md}, with its check scripts and
their recorded outputs in \path{checks/} of the same folder.\\
Inherited $d=2$ finite kernels &
Superseded inputs of the previous two-dimensional estimator (the primitive
and locally subtracted kernels of Section~\ref{gff:cn:finiteKernels}), not
additional contributions to add again and not needed to reproduce any
current coefficient. The kernels, rates, two-loop input and fixed-point
map that the 2026 production uses are byte-identical copies in
\path{LRI_C6_D2_HPC/reference/}, bound by
\path{LRI_C6_D2_HPC/inputs/manifest.json}. The construction records
(\path{lri_C6_d2_primitive_update/}, \path{lri_C6_d2_subtraction_update/},
\path{lri_C6_d2_two_forest_update/}) are in the research archive
\texttt{LRI\_complete\_research\_bundle.zip}, which is not part of this
repository.\\
Power compensation and finite kernels &
\path{lri_C6_d2_complete_update/}: the part of the research-archive package
that the two-dimensional power-sector compensation and the audits of
Sections~\ref{gff:app:C6d2Complete} and~\ref{gff:cn:finiteAssembly} need, namely
the saved counterterms, the check scripts and the counterterm chain.
\texttt{make -f public.mk replay} reruns the checks and compares them with
the records in \path{GFF/results/C6_2/power_compensation/}, among them
\path{GFF/results/C6_2/power_compensation/verification/power_sector_audit.json};
\texttt{make -f public.mk regen} regenerates the counterterms from the
graph inventory. The three-dimensional compensation is
\path{lri_C6_d3_complete_update/}, with its full-coverage, analytic-extension,
power-sector and fixed-point audits; its quoted values are in
\path{GFF/results/C6_3/power_compensation_quoted.json}. The superseded first
full sampling run of $C_6(3)$ in that package is not included.\\
$C_6(2)$ &
\path{GFF/results/C6_2/C6_result.json}, with the production log, reporting
rule and checks in \path{GFF/results/C6_2/CAMPAIGN.md}. The frozen plan
\path{manifest.json}, the group configurations and the per-group
\path{merged.json} records are in \path{GFF/data/C6_2/run-20260928-m8192/},
and the complete run directory with the checkpoints is hosted outside the
repository and fetched by \texttt{make fetch-data}. The code is
\path{LRI_C6_D2_HPC/} at git commit \texttt{6aaf656};
\texttt{make -f public.mk replay} there recomputes every number quoted here
for $C_6(2)$ from the merged records, and \texttt{make -f public.mk remerge}
re-merges the checkpoints. The seventh-order cross-check is
\path{GFF/results/C7_2/C7_result.json}. The superseded estimator is
\path{LRI_C6_D2_HPC/reference/verification/method_precision_results.json}.\\

Superseded sixth-order estimators at $d=2,3$ and the method benchmarks &
\path{LRI_C6_D2_HPC/reference/}, the earlier method-refinement calculation
shared by the two sixth-order packages: the refinement, pilot, validation,
speed and endpoint-replay records in \path{records/} (the baseline campaigns in
\path{baseline_previous/records/}); the per-graph counts, binary Sobol direction
matrices, frozen transport tables \path{grid_d2_top3.txt},
\path{grid_d2_top11.txt} (and their $d=3$ counterparts) and the frozen plan
\path{frozen_production_plan.json} in \path{inputs/}; the source and input
hashes in \path{verification/production_source_hashes.json}. Without new
integration, \path{code/analyze.py} reassembles the previous estimators and
their campaign comparisons (\path{verification/method_precision_results.json};
the rounded reporting convention is stored separately in
\path{verification/reported_values.json}), \path{code/summarize_methods.py} the
benchmarks of Section~\ref{gff:mgs:benchmarks}
(\path{verification/method_benchmarks.json}) and \path{code/extra_checks.py} the
finite Lagrange, bootstrap and delete-one checks
(\path{verification/extra_checks.json}). These archived scripts rewrite their
output files in \path{verification/} in place, and on another NumPy build the
last bits can differ;
\path{LRI_C6_D2_HPC/tools/reassemble_reference.py} runs the archived
\path{code/analyze.py} for $d=2$ without writing and compares with the stored
result to relative $10^{-12}$. The integrator is rebuilt by
\path{code/build.sh} (C++17 with OpenMP, libquadmath and Boost.Multiprecision)
and replayed from the shipped direction matrices by \path{code/replay.py}.
\texttt{make -f public.mk replay} does not run these archived scripts; the
method numbers of Sections~\ref{gff:mgs:benchmarks}--\ref{gff:meth:pooling} that are
derived from stored values, such as the linear-control ratios and the
refinement pulls, are recomputed by no script.\\

$C_6(3)$ &
\path{GFF/results/C6_3/results_run-20260929-m8192/C6_result.json},
with the frozen plan \path{manifest.json}, the group configurations and the
per-group \path{merged.json} records in the same folder. The production log,
reporting rule and checks are in \path{GFF/results/C6_3/CAMPAIGN.md},
and the complete run directory is
\path{GFF/data/C6_3/run-20260929-m8192_full.tgz}, hosted outside the
repository and fetched by \texttt{make fetch-data}. The code is
\path{LRI_C6_D3_HPC/} at git commit \texttt{1c6dfb3}; its archived-stream
replay is
\path{benchmarks/expanse_2026-09/logs/validation_54510646/archived_stream_replay.json}
in that folder, and \texttt{make -f public.mk replay} there recomputes every
number quoted here for $C_6(3)$ from the merged records. The superseded estimator is
\path{reference/verification/method_precision_results.json} in the reference
tree that the package shares with \path{LRI_C6_D2_HPC/}.\\
Full $C_7(1)$ &
\path{GFF/results/C7_1/C7_final.json}, with the frozen design
\path{design_v1.json}, the batch inclusion list \path{inclusion.txt}, the
combined assembly \path{combined_assembly.json} and the report
\path{REPORT.md} in the same folder; the earlier 46-batch combination is in
\path{v1_46batches/} of that folder. The code is
\path{LRI_C7_HPC/}, exported from git commit \texttt{68dc6be},
whose code files equal those of commit \texttt{7320afc}, which ran the final
combination; the production workers ran commit \texttt{3213ce0}, which
differs from it only in the verification, assembly, combination and design
tools and their tests. Full forest arrays, the retained primitive input and
the earlier assembly are in \path{LRI_C7_HPC/reference/}. The target
\texttt{replay} of \path{public.mk} in the code folder recomputes the quoted
numbers from these records. The merged replicate records of the 399 batches
and of the primitive campaign are hosted data, described in
\path{GFF/data/C7_1/DATA.md}; from them the target \texttt{replay-full}
reruns the combination.\\
Full $C_7(2)$ &
\path{GFF/results/C7_2/C7_result.json}, with the frozen plan
\path{manifest.json} and the per-group \path{merged.json} records in
\path{GFF/data/C7_2/run1-20260930-m4096/}; \texttt{make -f public.mk replay}
in \path{LRI_C7_D2_HPC/} recomputes the quoted numbers from them for the 2026
production, together with the earlier estimator and power correction
(reassembled from \path{reference/} by \path{tools/reassemble_reference.py}), the
three-loop reference and the coverage tables; the other numbers of the earlier
calculation in Section~\ref{gff:app:C7d2} are records in \path{reference/},
reassembled by \path{reference/code/reproduce.py}, and are not replayed. The
production log and checks are in \path{GFF/results/C7_2/CAMPAIGN.md}; the
code is \path{LRI_C7_D2_HPC/} at git commit \texttt{fb64c47}. The code and records of the earlier
calculation (Section~\ref{gff:app:C7d2}), with its power-free and primitive campaigns in \path{reference/LRI_C7_d2_progress/} and its power-sector completion in \path{reference/power_mass/}, are in
\path{LRI_C7_D2_HPC/reference/}. The checkpoints and task tables of the run
are archived, not included.\\
Full $C_7(3)$ &
\path{GFF/results/C7_3/run1-20260930-m4096/C7_result.json},
with the frozen plan \path{manifest.json}, the group configurations and the
per-group \path{merged.json} records in the same folder. The production log,
reporting rule, errata and checks are in \path{GFF/results/C7_3/CAMPAIGN.md}, and
the complete run directory is the hosted
\path{GFF/data/C7_3/run1-20260930-m4096_full.tgz}. The production used
\path{LRI_C7_D3_HPC/}, the export of git commit \texttt{4de9412}; the
re-verification tool
\path{LRI_C7_D3_HPC/reverify_a6ce987/tools/report.py}, from commit
\texttt{a6ce987}, reproduces the result. The previous calculation and its
construction checks are in \path{LRI_C7_D3_HPC/reference/}, with its raw
records in the hosted \path{GFF/data/C7_3/reference_records.tar.gz}; the
archived-stream replay is in \path{GFF/results/C7_3/evidence/replay_report.json},
and the reviews are summarized in \path{GFF/results/C7_3/REVIEWS.md}.\\
Full $C_8(1)$ &
\path{GFF/results/C8_1/2026-10_ext1/C8_final_A3/C8_final.json}, the
fourteen-batch result, with the report \path{REPORT.md} one level up, the
nine-batch interim in \path{C8_A1_close/} and the records of both clusters in
\path{records_symmetry/} and \path{records_expanse/}; the extension rules are
\path{LRI_C8_HPC/docs/EXTENSION_A1.md}, \path{EXTENSION_A2.md} and
\path{EXTENSION_A3.md}, and the independent reviews are summarized in
\path{LRI_C8_HPC/REVIEWS.md}. The four-batch result
\path{GFF/results/C8_1/2026-10_expanse/C8_final.json} has the report
\path{REPORT.md}, the frozen design \path{frozen_design/}, the controlled affine
map \path{full_affine8.json}, the batch plans and the combined assembly in the
same folder. The merged records of all fourteen batches and of the
seven-loop test with 4096 points per graph are in \path{GFF/data/C8_1/};
\texttt{make -f public.mk replay} in \path{LRI_C8_HPC/} recombines them with
the production tools and checks every quoted number. The production code is
\path{LRI_C8_HPC/} at git commit \texttt{c0ffa63} (tools of the extension at
\texttt{15a110e} and \texttt{bd13a63}; the public copy is commit
\texttt{1282e37}, whose code equals \texttt{bd13a63}); its nine-vertex
inventory is \path{inputs8/} (the eight-loop coordinate file compressed in
\path{GFF/data/C8_1/inputs8/}), generated by \path{gen/}, and its kernel
regressions against the seventh-order binary are in \path{validation/}.\\
Full $C_8(2)$ &
\path{GFF/results/C8_2/C8_final.json}, with the tail diagnostics in
\path{GFF/results/C8_2/tailcheck/} and the production's replicate records
(plan, design bundle, merged groups) in \path{GFF/data/C8_2/}. The code is
\path{LRI_C8_D2_HPC/} at git commit \texttt{9df5b1c}, whose code equals that
of \texttt{ba68f40}, the export the production ran; its folder
\path{benchmarks/symmetry_2026-10-c8/} holds the result report
(\path{production/result/README.md}), the tail-check report
(\path{production/tailcheck/README.md}) and the pilot analysis
(\path{pilot/analysis/}), and \texttt{make -f public.mk replay} recomputes
every number of Sections~\ref{gff:c8d2:result}--\ref{gff:c8d2:checks} from the records. Its plan reviews and
code inspections are summarized in \path{LRI_C8_D2_HPC/REVIEWS.md}. The
shared nine-vertex inventory is \path{LRI_C8_shared/inventory/inventory8_d2/},
and the shared eighth-order assembly \path{LRI_C8_shared/assembly/} is
vendored byte for byte in \path{vendor/c8_assembly/}. The full campaign
directories are archived as listed in \path{GFF/data/C8_2/DATA.md}.\\
Independent check of $C_8(2)$ &
\path{GFF/results/C8_2_check/prod1/C8_result.json}, with the frozen plan, the
group configurations and merged records in \path{GFF/data/C8_2_check/prod1/},
the pilot in \path{GFF/data/C8_2_check/pilot1/} and the analysis in
\path{GFF/results/C8_2_check/prod1_analysis/}. The code is
\path{LRI_C8_D2_HPC_optimized/} at git commit \texttt{a791a1f};
\texttt{make -f public.mk replay} there recomputes the numbers of the
paragraphs on the port, pilot and production and on the comparison in
Section~\ref{gff:c8d2:independent}; the construction checks under Inputs, Graphs and
payloads, Kernel arithmetic and Assembly, and the development runs, rest on the
records of the hosted original package, and the comparison of the two
packages' inputs, the bootstrap ratios and the bootstrap mean are not
replayed. Its reviews are summarized in
\path{LRI_C8_D2_HPC_optimized/REVIEWS.md}. The full campaign directory and
the original package with its graph packs are hosted
(\path{GFF/data/C8_2_check/DATA.md}).\\
Full $C_8(3)$ &
\path{GFF/results/C8_3/C8_result.json}, with the frozen plan
\path{manifest.json}, the group configurations and the per-group
\path{merged.json} records that reproduce the estimator in the same folder;
the per-graph replicate tables \path{graph_replicates.csv} are hosted in
\path{GFF/data/C8_3/}. The production report \path{README.md}, the tail
diagnostics \path{tail_ntrend/} and the campaign split \path{ab_split/} are in
\path{LRI_C8_D3_HPC/benchmarks/symmetry_2026-10-c8/production/prod1/},
the rate probes in \path{production/margin_*}, the priced designs in
\path{production/price/} and the pilot report in \path{pilot/evidence/} of the
same \path{symmetry_2026-10-c8/} folder. The code is
\path{LRI_C8_D3_HPC/}: the production ran git commit
\texttt{4fbd1bae}, which is the code shipped, and the evidence is taken from
commit \texttt{6dd4a75a}; \path{public.mk} replays every quoted number. The
plan, its review log and the results review are not included; they are
summarized in \path{public/REVIEWS.md}. The shared nine-vertex inventory and
eighth-order assembly are those of the $C_8(2)$ row. The campaign directories
(13\,GB) are not included; they are archived in two copies.\\
Independent check of $C_8(3)$ &
\path{GFF/results/C8_3_check/runs/prod1/C8_result_54838171.json}, with the
frozen plan \path{manifest.json} and the per-graph statistics
\path{graph_statistics_54838171.json} in the same folder, the frozen
allocation \path{allocations/prod1_2000.json} and the production summary
\path{README.md} in \path{GFF/results/C8_3_check/}. The code is
\path{LRI_C8_D3_OPT/} at git commit \texttt{d5deeb9}; the target
\texttt{replay} of \path{public.mk} in that folder recomputes the estimator,
its standard error and the quoted checks from the means and covariances stored
in the result (\path{public/replay_check.py}). The port to the clusters, its
decisions and its reviews are summarized in \path{public/PROVENANCE.md} and
\path{public/REVIEWS.md} of the same folder; the logs, the admission ledger and
the lane summaries are not included. The graph payloads and the delivered pilot
are hosted data, and the native outputs of the campaign are not archived; both
are described in \path{GFF/data/C8_3_check/DATA.md}.\\
Bulk one-point coefficient (Part~V) &
\path{lri_ope_reproducibility/}, supplied by the author on 16 October 2026:
the engines in \path{code/}, the graph inventories and master inputs in
\path{data/}, the coefficient table \path{results/coefficients.json} and the
verification reports in \path{results/verification/}; the tasks of its
runner \path{run.py} are listed in Section~\ref{gff:ope:reproduce}.
\texttt{make -f public.mk test} checks the release hashes and runs the quick
task in a work folder, \texttt{make -f public.mk replay} also reruns the
checks of Section~\ref{gff:ope:checks}, which are in \path{ope_checks/} (folders
\path{d1/}, \path{d2/}, \path{d3/}, with their recorded outputs), and
\texttt{make -f public.mk full} reruns the remaining tasks (about ten minutes on
one core). The package's own note is merged into Part~V and is not part of
the repository.\\
\bottomrule
\end{longtable}
The mathematical definitions and intermediate steps above specify the
quantities being computed; raw numerical replicate files are evidence for
their estimators, not substitutes for missing master definitions. The
file \path{GFF/REPRODUCE.md} distinguishes the inexpensive replay from saved
records from graph regeneration and expensive integration. The numerical
allowances throughout this note are empirical and are not upgraded by
checksum or source-preservation checks.

\SourcePart{Bulk one-point data}

\section{The bulk one-point coefficient through sixth order}\label{gff:ope:section}

This part records a second observable built from the graphs and master
functions of Parts~I--III: the bulk one-point coefficient $a_{\phi^2}$ of the
composite $:\!\Phi^2\!:$ of the auxiliary bulk field, through $\epsilon^6$, at
fixed $0<d<4$. It is the horizontal coordinate of the numerical-bootstrap
bounds of \cite{gff:BLNV}. The calculation, its code and its checks were supplied
by the author on 16 October 2026 as a reproducibility package; this part
contains its derivations, results, conditions and caveats, so that the package
is needed only to rerun the computations. The package is
\path{lri_ope_reproducibility/}; package paths in this part are relative to
\path{GFF/code/}, as in Section~\ref{gff:cn:reproduction}.
Sections~\ref{gff:ope:status}--\ref{gff:ope:graphs} hold for general $d$, and
Sections~\ref{gff:ope:values}--\ref{gff:ope:d1functional} specialize to $d=1,2,3$.
Section~\ref{gff:ope:bootstrap} compares with the bootstrap,
and Section~\ref{gff:ope:checks} lists what was checked independently for this note.

\subsection{Status}\label{gff:ope:status}
The material has four logically different layers.
\begin{enumerate}
\item \emph{Definitions and the two lowest orders.} The observable, its
normalization and $a^{(2)}_{\phi^2}$, $a^{(3)}_{\phi^2}$ agree with the published
values of \cite{gff:BLNV} (Section~\ref{gff:ope:lowOrders}).
\item \emph{The calculation of $a^{(4)}_{\phi^2},a^{(5)}_{\phi^2},a^{(6)}_{\phi^2}$.}
It uses the graph set, local subtractions and master functions of this note.
Its checks are internal to that framework: the five- and six-vertex graphs are
reconstructed from shared master jets and checked with alternative regulator
directions, not integrated one by one. The four-vertex finite parts were also
evaluated at finite regulator.
\item \emph{The conversion to the thermal coefficients} $C_j(d)$
(Section~\ref{gff:ope:thermal}). It is exact algebra between two calculations that
share their inputs; it does not validate either of them independently.
\item \emph{Finite-coupling estimates} (Section~\ref{gff:ope:bootstrap}). The
Borel--Pad\'e values and the comparison with published bootstrap figures are
exploratory. No bootstrap problem was solved, the bootstrap data are readings of
published figures, and the spread between resummation prescriptions is not an
error bar.
\end{enumerate}
Residuals of the internal checks, however small, do not bound the omitted
orders of the series, the resummation bias, or an error shared by the
calculation and its inputs. The coefficients are numbers of a finite-order
calculation; $a^{(7)}_{\phi^2}$ and higher are not determined by the thermal
coefficients $C_6$, $C_7$, $C_8$ of Part~IV.

\subsection{The observable}\label{gff:ope:observable}
Write the long-range Ising model as the $d$-dimensional defect of a free scalar
$\Phi$ in a bulk of dimension $d+2-s=2+2\Delta_\phi$ and codimension $2-s$, the
local description of \cite{gff:BLNV,gff:GK}. In the notation of \cite{gff:BLNV} the
physical dimension is $p$, which is our $d$; the bulk dimension is its $d$ and
the codimension its $q$, whereas here $q=d/4$ throughout. The bulk propagator
and the defect interaction are
\begin{equation}
 \langle\Phi(X)\Phi(X')\rangle_0=|X-X'|^{-2\Delta_\phi},\qquad
 \phi(x)=\Phi(x,0),\qquad
 S_{\rm int}=\frac{\lambda_B}{4!}\int d^dx\,\phi(x)^4,
 \label{gff:ope:action}
\end{equation}
with $\Delta_\phi=(d-\epsilon)/4=q-\epsilon/4$ and $\epsilon=2s-d$, $q=d/4$, as in
\eqref{gff:gs:conventions}. The critical local mass is tuned away, and no
wave-function counterterm arises in the long-range regime, so $\Delta_\phi$ is
protected. For a bulk point at distance $|y|$ from the defect, define
\begin{equation}\boxed{
 a_{\phi^2}(d,\epsilon)=|y|^{2\Delta_\phi}\lim_{X'\to X}
 \Bigl[\langle\Phi(X)\Phi(X')\rangle-|X-X'|^{-2\Delta_\phi}\Bigr],\qquad
 \langle:\!\Phi^2\!:(x,y)\rangle=\frac{a_{\phi^2}(d,\epsilon)}{|y|^{2\Delta_\phi}}.
}\label{gff:ope:definition}\end{equation}
This is the coincident-limit convention of Appendix~B.4.2 of \cite{gff:BLNV}
(equation, figure and appendix numbers of \cite{gff:BLNV} refer to its arXiv
version~6). For the
unit-normalized composite $:\!\Phi^2\!:/\sqrt2$, divide by $\sqrt2$. After
tadpole and critical-mass subtraction there is no term of order $\lambda$, so
\begin{equation}
 a_{\phi^2}(d,\epsilon)=\sum_{n=2}^6a^{(n)}_{\phi^2}(d)\,\epsilon^n+O(\epsilon^7).
 \label{gff:ope:series}
\end{equation}
The superscript $(n)$ is an order of $\epsilon$; it is not an index of the
Laurent coefficients $a_{L,k}$ of Sections~\ref{gff:sec:graphs} and
\ref{gff:app:generalDimension}. Below, $a^{(n)}$ abbreviates $a^{(n)}_{\phi^2}(d)$.

Three quantities must be kept apart. The thermal dimension
$\Delta_{\Eop}=\sum_nC_n(d)\epsilon^n$ has the free value $d/2$; the squared defect
OPE coefficient $\lambda^2_{\phi\phi\Eop}$ has the free value $2$; and
$a_{\phi^2}$ has the free value $0$. Only the last is computed here. For the
second, only the one-loop relation is recorded:
\begin{equation}
 \lambda^2_{\phi\phi\Eop}=2-6C_2(d)\,\epsilon+O(\epsilon^2),
 \label{gff:ope:defectOPE}
\end{equation}
in agreement with the tree-level and one-loop terms $\sqrt2$ and
$-A\epsilon/(3\sqrt2)$ of $\lambda_{\phi\phi\Eop}$ in Eqs.~(3.24) and (3.33) of
\cite{gff:BLNV}. It is a separate low-order observable, and the higher-order
formulas of this part must not be used for it.
For the bootstrap, Eq.~(2.16) of \cite{gff:BLNV} allows at least
$-0.575408\le a_{\phi^2}\le0$ at $d=3$ and $-\frac78\le a_{\phi^2}\le0$ at $d=2$.
This range contains the generalized-free end, $a_{\phi^2}=0$, and the
short-range end of its Eq.~(2.15), where $\Delta_\phi$ is the short-range Ising
value, $0.5181489$ at $d=3$ and $\frac18$ at $d=2$; it is compatible with
bulk--defect crossing with nonnegative squared defect-channel coefficients.

\subsection{The result}\label{gff:ope:result}
With $A,R,U,V,X$ of \eqref{gff:gs:psiBasis}, the masters $Q_d$, $P_4(d)$,
$P_5(d)$, $\KK_d$, $\LL_d$, $Z_d$ of Section~\ref{gff:app:generalSpecial}, and the
reflected polygamma differences
\begin{equation}
 \mathsf e_j=\psi_j(-q)-\psi_j(q)=\psi_j(1-q)-\psi_j(q)+\frac{j!}{q^{j+1}},
 \qquad \mathsf e_0=\pi\cot(\pi q)+\frac1q,
 \label{gff:ope:ej}
\end{equation}
the calculation gives
\begin{equation}\boxed{
 a_{\phi^2}(d,\epsilon)=\epsilon^2\,
 \frac{\Gamma(2q)^3\,\Gamma(-q-\epsilon/4)}{54\,\Gamma(4q)\,\Gamma(q-\epsilon/4)}
 \Bigl[1+\mathsf H_1\epsilon+\mathsf H_2\epsilon^2+\mathsf H_3\epsilon^3
 +\mathsf H_4\epsilon^4\Bigr]+O(\epsilon^7),
}\label{gff:ope:factorized}\end{equation}
with
\begin{align}
 \mathsf H_1={}&\frac43A,\qquad
 \mathsf H_2=\frac{25}{36}A^2+\frac R6-\frac49Q_d,\notag\\
 \mathsf H_3={}&\frac{17}{108}A^3+\frac{AR}{12}-\frac49AQ_d+\frac4{27}P_4
 +\frac{71}{216}U-\frac V{36},\label{gff:ope:H123}
\end{align}
\begin{equation}\boxed{\begin{aligned}
 \mathsf H_4={}&-\frac{7A^4}{5184}-\frac{7A^2R}{288}-\frac{7A^2Q_d}{162}
 +\frac{14AP_4}{81}+\frac{301AU}{864}-\frac{7AV}{144}
 -\frac{28\KK_d}{81}-\frac{7\LL_d}{81}-\frac{16P_5}{81}\\
 &-\frac{103R^2}{5184}+\frac{7RQ_d}{54}-\frac{14Q_d^2}{81}
 -\frac{61X}{3888}+\frac{175Z_d}{162}.
\end{aligned}}\label{gff:ope:H4}\end{equation}
Here $P_4=P_4(d)$, $P_5=P_5(d)$, and $6Q_d$ is the corrected tetrahedron
$\mathcal T_d$ of \eqref{gff:gd:Qdef}. The leading coefficient is
\begin{equation}
 a^{(2)}=\frac{\Gamma(-q)\Gamma(2q)^3}{54\,\Gamma(q)\Gamma(4q)}
 =-\frac{\Gamma(1-q)\Gamma(2q)^3}{54\,\Gamma(1+q)\Gamma(4q)}<0 .
 \label{gff:ope:a2}
\end{equation}
Expanding the Gamma ratio,
$\Gamma(-q-\epsilon/4)\Gamma(q)/[\Gamma(q-\epsilon/4)\Gamma(-q)]
=\exp\sum_{k\ge1}(-\frac14)^k\mathsf e_{k-1}\epsilon^k/k!
=\sum_j\mathsf f_j\epsilon^j$, with
\begin{equation}\begin{aligned}
 \mathsf f_0&=1,\quad \mathsf f_1=-\frac{\mathsf e_0}4,\quad
 \mathsf f_2=\frac{\mathsf e_0^2+\mathsf e_1}{32},\quad
 \mathsf f_3=-\frac{\mathsf e_0^3+3\mathsf e_0\mathsf e_1+\mathsf e_2}{384},\\
 \mathsf f_4&=\frac{\mathsf e_0^4+6\mathsf e_0^2\mathsf e_1+3\mathsf e_1^2
 +4\mathsf e_0\mathsf e_2+\mathsf e_3}{6144},
\end{aligned}\label{gff:ope:fj}\end{equation}
the ordinary coefficients are the finite convolution
\begin{equation}\boxed{
 a^{(n)}=a^{(2)}\sum_{j=0}^{n-2}\mathsf H_j\,\mathsf f_{n-2-j},\qquad
 \mathsf H_0=1,\quad 2\le n\le6.
}\label{gff:ope:convolution}\end{equation}
For example $a^{(3)}/a^{(2)}=\frac43A-\frac{\mathsf e_0}{4}
=\frac43A-\frac\pi4\cot(\pi q)-\frac1d$.
Equations \eqref{gff:ope:ej}--\eqref{gff:ope:convolution} and the master values suffice
to evaluate the coefficients; no graph has to be rebuilt.

The form \eqref{gff:ope:factorized} is a statement about finite orders. Keeping the
Gamma ratio unexpanded generates terms beyond $\epsilon^6$; those terms are not
predicted. The organization itself is derived: all dependence on the
$\mathsf e_j$ comes from derivatives of the external-leg integral along one
direction (Section~\ref{gff:ope:lowOrders}), and the remaining coefficients
$\mathsf H_j$ contain only $A,R,U,V,X$ and the masters.

\subsection{Relation to the thermal coefficients}\label{gff:ope:thermal}
Eliminating $Q_d$, $P_4$ and $P_5$ with \eqref{gff:gs:C0C4}--\eqref{gff:gs:C5} gives
\begin{align}
 \mathsf H_2={}&6C_3+\frac{23}{36}A^2+\frac R9,\qquad
 \mathsf H_3=6C_4+5AC_3+\frac{13}{108}A^3+\frac{7AR}{108}+\frac U{12},
 \label{gff:ope:H23thermal}\\
 \mathsf H_4={}&6C_5+5AC_4+\frac{5A^2-R}{6}C_3+9C_3^2
 -\frac{A^4}{5184}+\frac{7A^2R}{2592}+\frac{235AU}{2592}-\frac{5AV}{432}
 -\frac{7R^2}{1728}-\frac{47X}{7776}\notag\\
 &+\frac{13Z_d-\LL_d-12\KK_d}{81}.\label{gff:ope:H4thermal}
\end{align}
The logarithm of the normalized series is shorter. Define
\begin{equation}
 \log\frac{a_{\phi^2}(d,\epsilon)}{a^{(2)}\epsilon^2}
 =\mathsf l_1\epsilon+\mathsf l_2\epsilon^2+\mathsf l_3\epsilon^3
 +\mathsf l_4\epsilon^4+O(\epsilon^5),
 \label{gff:ope:logdef}
\end{equation}
a formal Taylor series of a ratio that tends to one (no logarithm of a negative
number is taken). Then
\begin{equation}\boxed{\begin{aligned}
 \mathsf l_1={}&12C_2-\frac{\mathsf e_0}4,\qquad
 \mathsf l_2=6C_3-\frac{A^2}4+\frac R9+\frac{\mathsf e_1}{32},\\
 \mathsf l_3={}&6C_4-3AC_3+\frac{19A^3}{324}-\frac{AR}{12}+\frac U{12}
 -\frac{\mathsf e_2}{384},\\
 \mathsf l_4={}&6C_5-3AC_4+\Bigl(A^2-\frac56R\Bigr)C_3-9C_3^2
 -\frac{11A^4}{576}+\frac{37A^2R}{864}-\frac{53AU}{2592}-\frac{5AV}{432}
 -\frac{53R^2}{5184}\\
 &-\frac{47X}{7776}+\frac{13Z_d-\LL_d-12\KK_d}{81}+\frac{\mathsf e_3}{6144},
\end{aligned}}\label{gff:ope:logmap}\end{equation}
where $A=9C_2$. Exponentiation gives
\begin{equation}\begin{aligned}
 \frac{a^{(3)}}{a^{(2)}}&=\mathsf l_1,\qquad
 \frac{a^{(4)}}{a^{(2)}}=\mathsf l_2+\frac{\mathsf l_1^2}2,\qquad
 \frac{a^{(5)}}{a^{(2)}}=\mathsf l_3+\mathsf l_1\mathsf l_2+\frac{\mathsf l_1^3}6,\\
 \frac{a^{(6)}}{a^{(2)}}&=\mathsf l_4+\mathsf l_1\mathsf l_3+\frac{\mathsf l_2^2}2
 +\frac{\mathsf l_1^2\mathsf l_2}2+\frac{\mathsf l_1^4}{24}.
\end{aligned}\label{gff:ope:expmap}\end{equation}
Through $a^{(5)}$ the one-point coefficients are therefore fixed by
$C_2,\dots,C_4$, the polygamma values $A=9C_2$, $R$, $U$ and
$\mathsf e_0,\mathsf e_1,\mathsf e_2$, and the Gamma functions of
\eqref{gff:ope:a2}. At sixth order they are not fixed by
$C_2,\dots,C_5$: the combination $13Z_d-\LL_d-12\KK_d$ is independent of the
combination $-1536\KK_d-576\LL_d+7152Z_d$ that enters $C_5$ in \eqref{gff:gs:C5}, so
$a^{(6)}$ needs one more master value at a generic dimension. At $d=2$ the
closed forms of Section~\ref{gff:gc:d2closedForms} remove it (Section~\ref{gff:ope:d2}).
These relations eliminate shared integrals. They do not identify scaling
dimensions with OPE data, and they are checked by substitution, not by a new
loop calculation.

\subsection{Renormalization-group input}\label{gff:ope:RG}
The amplitudes of this part are massless position-space integrals, and
$\epsilon$ itself is their regulator. The coupling is
$g=\lambda/h_d=\pi^{d/2}\lambda/\Gamma(d/2)$, with $\lambda$ the quartic
coupling of \eqref{gff:ope:action}, so that every vertex integration carries the
normalized measure $h_d\,d^dx$ of Section~\ref{gff:gs:patterns}. It is renormalized by
minimal subtraction of the poles in $\epsilon$ (coordinate-space MS, with the
forest counterterms of Section~\ref{gff:ope:method}), with the critical-mass
prescription for local power subgraphs, and
\begin{equation}
 \beta(g)=-\epsilon g+3g^2-6Ag^3+b_3g^4+b_4g^5+b_5g^6+\cdots.
 \label{gff:ope:beta}
\end{equation}
This is not the massive momentum-space prescription of \eqref{gff:gd:RG} and
\eqref{gff:gd:normalization}, which has its own regulator $z$. The coefficients
$b_L$, $c_L$ of the two prescriptions nevertheless coincide where they were
compared: $b_3,c_3$ and $b_4,c_4$ of this prescription are those of
\eqref{gff:gd:bc3}, \eqref{gff:gs:b4} and \eqref{gff:gs:c4}, with $M$ and $N$ of
\eqref{gff:gs:MN}; its $3c_5-b_5$ is $D_5$ of \eqref{gff:gs:D5}; and its $b_5$ agrees
numerically with five-loop data of this note at $d=1,2,3$
(Section~\ref{gff:ope:checks}, items ``Exact algebra'' and ``Five-loop split'').
Finite parts are a different matter: finite parts of a massive-regulator
prescription and of the massless position-space prescription must not be
interchanged. The package's low-order matching relates a massive-regulator
coupling $g_{\rm alt}$ to $g$ through the regulator-dependent factor
$\mathsf F_{\rm sch}$; with both couplings minimally subtracted,
\begin{equation}\begin{aligned}
 \mathsf F_{\rm sch}(\epsilon)&=2^{-\epsilon}
 \frac{\Gamma(2q+\epsilon/2)}{\Gamma(2q)\Gamma(1+\epsilon/2)}
 \frac{\Gamma(q-\epsilon/4)^2}{\Gamma(q+\epsilon/4)^2}
 =1+\mathsf f^{\rm sch}_1\epsilon+\mathsf f^{\rm sch}_2\epsilon^2+\cdots,\\
 g_{\rm alt}&=\mathsf F_{\rm sch}\,g
 -\bigl[3\mathsf f^{\rm sch}_1+3\bigl(\mathsf f^{\rm sch}_2+(\mathsf f^{\rm sch}_1)^2\bigr)\epsilon+\cdots\bigr]g^2
 +\bigl[9(\mathsf f^{\rm sch}_1)^2+6A\,\mathsf f^{\rm sch}_1+O(\epsilon)\bigr]g^3+\cdots.
\end{aligned}\label{gff:ope:schemeF}\end{equation}
The package checks this matching at the fixed point through the order used
there (\path{lri_ope_reproducibility/code/a4/symbolic_check.py}). The
higher-order assembly below uses the coordinate-space prescription and its
forest counterterms throughout; it imports no unmatched finite amplitude.
The fixed point is
\begin{equation}\begin{aligned}
 g_*={}&\frac\epsilon3+\frac{2A}{9}\epsilon^2
 +\Bigl(\frac{8A^2}{27}-\frac{b_3}{81}\Bigr)\epsilon^3
 +\Bigl[\frac{40A^3}{81}-\frac{10Ab_3+b_4}{243}\Bigr]\epsilon^4\\
 &+\frac{672A^4-84A^2b_3-12Ab_4+b_3^2-b_5}{729}\epsilon^5+O(\epsilon^6).
\end{aligned}\label{gff:ope:gstar}\end{equation}
Since $a_{\phi^2}$ starts at $g^2$, $a^{(6)}$ needs $b_5$ but not $b_6$; it needs
$b_5$ itself, not only the thermal combination $D_5=3c_5-b_5$ of
\eqref{gff:gs:D5}. With $\widetilde N=\psi_1(-q)-\psi_1(3q)$, the five-loop forest
sums of the calculation are
\begin{equation}\begin{aligned}
 b_5={}&\frac{64341A^4}{128}+\frac{2145A^2M}{64}-\frac{17001A^2R}{64}
 +\frac{3159A^2Q_d}{8}+\frac{393AMN}{32}+63AP_4\\
 &+\frac{14907AU}{64}-\frac{603AV}{32}+87\KK_d+\frac{51}4\LL_d+\frac34M^2
 +\frac{81MN^2}{64}-\frac98MR+6MQ_d\\
 &+\frac{81M\widetilde N}{64}+24P_5+\frac{1839R^2}{128}
 -\frac{459RQ_d}{8}+63Q_d^2+\frac{565X}{128}-\frac{1311Z_d}{8},
\end{aligned}\label{gff:ope:b5}\end{equation}
\begin{equation}\begin{aligned}
 c_5={}&\frac{10549A^4}{128}+\frac{411A^2M}{64}-\frac{3561A^2R}{64}
 +\frac{267A^2Q_d}{8}+\frac{99AMN}{32}+3AP_4\\
 &+\frac{2531AU}{64}-\frac{75AV}{32}+21\KK_d+\frac54\LL_d+\frac14M^2
 +\frac{27MN^2}{64}-\frac18MR\\
 &+\frac{27M\widetilde N}{64}+\frac{463R^2}{128}-\frac{63RQ_d}{8}
 +\frac{415X}{384}-\frac{139Z_d}{8}.
\end{aligned}\label{gff:ope:c5}\end{equation}
Their combination $3c_5-b_5$ is exactly $D_5$ of \eqref{gff:gs:D5}; with the
Lagrange formula \eqref{gff:c6:Lagrange}, which holds at any $d$, the coefficients
\eqref{gff:gd:bc3}--\eqref{gff:gs:c4}, \eqref{gff:ope:b5} and \eqref{gff:ope:c5} give
$C_2,\dots,C_5$ of \eqref{gff:gs:C0C4}--\eqref{gff:gs:C5}. The split into $b_5$ and
$c_5$ is not visible in $D_5$: zero-weight vertex classes, for example, cancel
from $D_5$ but not from $b_5$. It was checked separately against this note's own data at $d=1,2,3$
(Section~\ref{gff:ope:checks}). The sensitivity is
$\partial a^{(6)}/\partial b_5=-\frac2{243}a^{(2)}$, about $2\times10^{-4}$ at
$d=3$.

For $a_{\phi^2}=\sum_{n\ge2}g^n\mathsf a_n(\epsilon)$, sixth order needs the
regular Taylor orders
\begin{center}
\begin{tabular}{@{}c*{5}{c}@{}}\toprule
target & $\mathsf a_2$ & $\mathsf a_3$ & $\mathsf a_4$ & $\mathsf a_5$ & $\mathsf a_6$\\\midrule
$a^{(4)}$ & $\epsilon^2$ & $\epsilon^1$ & $\epsilon^0$ & -- & --\\
$a^{(5)}$ & $\epsilon^3$ & $\epsilon^2$ & $\epsilon^1$ & $\epsilon^0$ & --\\
$a^{(6)}$ & $\epsilon^4$ & $\epsilon^3$ & $\epsilon^2$ & $\epsilon^1$ & $\epsilon^0$\\\bottomrule
\end{tabular}
\end{center}
together with $g_*$ through $\epsilon^{n-1}$. Extending the fixed point alone
does not give the next coefficient: each order needs new finite amplitudes.

\subsection{Graphs, kernels and finite parts}\label{gff:ope:graphs}

\subsubsection{External legs and the two lowest orders}\label{gff:ope:lowOrders}
Put the bulk point at unit distance from the defect. Both bulk legs end on
defect vertices $x_0,x_1$, and every two-point graph reduces to an internal
kernel $\propto|x_0-x_1|^{-2b}$ followed by the external-leg integral
\begin{equation}
 \mathsf J_d(a,b)=\int d^dx\,d^dy\,\frac{1}{(1+x^2)^a(1+y^2)^a|x-y|^{2b}}
 =\pi^d\frac{\Gamma(\frac d2-b)\Gamma(2a+b-d)\Gamma(a+b-\frac d2)^2}
 {\Gamma(a)^2\Gamma(\frac d2)\Gamma(2a+2b-d)} .
 \label{gff:ope:Jd}
\end{equation}
The closed form follows from Schwinger parameters (integrate the parameter of
$|x-y|$ first, then the two others as a Gamma and a Beta integral) in the
region of convergence; at the physical exponents it denotes the critical-mass
continuation. A convergent normalization check is $\mathsf J_1(1,\frac14)=\pi^2$.
For an $n$-vertex graph the $n-2$ internal integrations fix
$b=(2n-1)\Delta_\phi-(n-2)\frac d2$, and with $h_d$ of
Section~\ref{gff:app:generalSpecial}
\begin{equation}
 \mathsf L_n(\epsilon)=h_d^2\,\mathsf J_d\Bigl(\Delta_\phi,(2n-1)\Delta_\phi-\frac{(n-2)d}{2}\Bigr)
 =\frac{\Gamma(2q)\Gamma(-q+\frac{2n-1}4\epsilon)\Gamma(q-\frac{2n+1}4\epsilon)
 \Gamma(2q-\frac n2\epsilon)^2}{\Gamma(q-\frac\epsilon4)^2\,\Gamma(4q-n\epsilon)} .
 \label{gff:ope:Ln}
\end{equation}
All $\mathsf L_n(0)$ are equal, to $54a^{(2)}$. Along the direction
$(a,b)=(q+t,3q-t)$,
$\mathsf J_d=\pi^d\Gamma(2q)\Gamma(-q+t)/[\Gamma(4q)\Gamma(q+t)]$, so the
$k$th derivative of $\log\mathsf J_d$ along $\partial_a-\partial_b$ at
$(q,3q)$ is exactly $\mathsf e_{k-1}$ of \eqref{gff:ope:ej}. The Gamma ratio of
\eqref{gff:ope:factorized} is, up to a constant, $\mathsf J_d$ on this line at
$t=-\epsilon/4$; this shift is the origin of the factors $(-\frac14)^k$ in
\eqref{gff:ope:fj}. When the assembly is carried out with these
outer derivatives kept as independent symbols, they enter only through the
combination
$\mathsf H_4-\frac{\mathsf H_3}4\frac{\partial_-\mathsf J_d}{\mathsf J_d}
+\frac{\mathsf H_2}{32}\frac{\partial_-^2\mathsf J_d}{\mathsf J_d}
-\frac{\mathsf H_1}{384}\frac{\partial_-^3\mathsf J_d}{\mathsf J_d}
+\frac1{6144}\frac{\partial_-^4\mathsf J_d}{\mathsf J_d}$,
$\partial_-=\partial_a-\partial_b$ at $(a,b)=(q,3q)$, which is the origin of
\eqref{gff:ope:factorized}; no all-order statement about the interacting amplitude
is assumed.

The first graph is the two-vertex sunset with Wick weight $\frac16$. With the
normalized bubble
\begin{equation}
 \mathsf Y(\epsilon)=\Gamma(\tfrac d2)\frac{\Gamma(\frac\epsilon2)^2\Gamma(\frac d2-\epsilon)}
 {\Gamma(\frac{d-\epsilon}2)^2\Gamma(\epsilon)}
 =\frac4\epsilon+R\epsilon-\frac V2\epsilon^2
 +\Bigl(\frac{R^2}8+\frac{7X}{48}\Bigr)\epsilon^3+O(\epsilon^4),
 \label{gff:ope:bubble}
\end{equation}
whose constant term vanishes, the renormalized amplitudes of order $g^2$ and
$g^3$ (at $\mu|y|=1$) are
\begin{equation}
 \mathsf a_2=\frac{\mathsf L_2}{6},\qquad
 \mathsf a_3=\frac{\mathsf L_2}\epsilon-\frac14\mathsf Y\mathsf L_3 .
 \label{gff:ope:a23amp}
\end{equation}
The pole of $\mathsf a_3$ cancels because $\mathsf L_2(0)=\mathsf L_3(0)$. With
$g_*\simeq\epsilon/3$ the sunset gives \eqref{gff:ope:a2}; the sign is that of
$\Gamma(-q)<0$ and involves no positivity statement about a divergent integral.
The second arguments of $\mathsf L_2$ and $\mathsf L_3$ differ by
$\epsilon/2$, so the finite part of $\mathsf a_3$ is
$\frac12h_d^2\,\partial_b\mathsf J_d$ at $(a,b)=(q,3q)$.
Adding the $\epsilon$-dependence of $\mathsf L_2$, the finite part of
$\mathsf a_3$ and the $O(\epsilon^2)$ term of $g_*$ gives
$a^{(3)}/a^{(2)}=\frac43A+\frac14[\psi(q)-\psi(-q)]$. At $d=2$ and $d=3$,
\begin{equation}\begin{aligned}
 a_{\phi^2}(2,\epsilon)&=-\frac{\epsilon^2}{27}
 +\Bigl(\frac1{54}-\frac{8\log4}{81}\Bigr)\epsilon^3+O(\epsilon^4),\\
 a_{\phi^2}(3,\epsilon)&=-\frac{\pi^{5/2}\epsilon^2}{576\sqrt2\,\Gamma(\frac74)^2}
 +\frac{4\sqrt{2\pi}}{6075}\bigl(13\pi-28-64\log2\bigr)\Gamma\bigl(\tfrac94\bigr)^2\epsilon^3
 +O(\epsilon^4),
\end{aligned}\label{gff:ope:BLNVvalues}\end{equation}
which are the values of Eq.~(3.36) of \cite{gff:BLNV} (Section~\ref{gff:ope:checks}).
The same two coefficients also follow from the exact bulk--defect relation of
\cite{gff:BLNV} and the two-loop ratio of its defect-channel coefficients, without
any higher-order one-point input. Let
$b_0^{\Phi,-}/b_0^{\Phi,+}=\mathsf r_1\epsilon+\mathsf r_2\epsilon^2+\cdots$ be
the ratio of the bulk-to-defect coefficients of the two scalar defect modes of
$\Phi$, in the notation and normalization of \cite{gff:BLNV}. Its Eqs.~(3.25) and
(3.34) give
\begin{equation}
 \mathsf r_1=\frac{\Gamma(1-q)\Gamma(2q)}{3\sqrt6\,\Gamma(1+q)},\qquad
 \mathsf r_2=-\frac{\mathsf r_1}{12}\bigl[13\psi(q)+3\psi(-q)-8\psi(2q)-8\psi(1)\bigr],
 \label{gff:ope:r12}
\end{equation}
and expanding the exact relation, its Eq.~(2.11), through cubic order gives
\begin{equation}
 a^{(2)}=\frac{\Gamma(2q)\Gamma(q)}{\Gamma(4q)\Gamma(-q)}\,\mathsf r_1^2,\qquad
 a^{(3)}=\frac{\Gamma(2q)\Gamma(q)}{\Gamma(4q)\Gamma(-q)}
 \Bigl(2\mathsf r_1\mathsf r_2+\frac{\mathsf e_0}4\,\mathsf r_1^2\Bigr).
 \label{gff:ope:Rcheck}
\end{equation}
These are \eqref{gff:ope:a2} and the value of $a^{(3)}/a^{(2)}$ above. The
agreement is not independent of all low-order renormalization-group input:
the diagrammatic route of this subsection still uses the known two-loop
$\beta$ function, through $g_*$.

Read literally, two general formulas of \cite{gff:BLNV}, its Eqs.~(3.35) and
(3.47), disagree in places with these values (Section~\ref{gff:ope:d1}; the package
records the literal evaluations in
\path{lri_ope_reproducibility/results/verification/low_order_checks.json}).
These are apparent printed-formula inconsistencies, not conclusions about the
authors' numerical code. The conventions used here are fixed by the
operational definition \eqref{gff:ope:definition} and the explicit low-order
checks, not by a transcription of those formulas.

\subsubsection{Four vertices}
A tadpole-free two-point graph with $n$ quartic vertices is a symmetric
multiplicity matrix $m_{ij}\in\{0,1,2,3\}$ with zero diagonal, degrees $3$ at the
two roots, which carry the bulk legs, and $4$ elsewhere, and a connected
underlying graph. Up to root exchange and permutations of the $n-2$ internal
vertices, its Wick weight is
\begin{equation}
 w(m)=\frac{|{\rm orb}(m)|}{(n-2)!\prod_{i<j}m_{ij}!},
 \label{gff:ope:wick}
\end{equation}
and the sign $(-1)^n$ is kept separately. Reducible graphs are included. Graphs
with both bulk legs on one vertex give only local mass and tadpole terms. The
numbers of graphs for $n=2,\dots,6$ are $1,1,5,14,64$. With edges listed as
$(01,02,03,12,13,23)$ and squared-distance exponents, the four-vertex graphs are
\begin{center}\small
\begin{tabular}{@{}lll@{}}\toprule
edge word & weight & internal reduction\\\midrule
$(0,0,3,3,0,1)$ & $1/36$ & chain $\mathsf C(3\Delta_\phi,\Delta_\phi,3\Delta_\phi)$\\
$(0,1,2,2,1,1)$ & $1/4$ & crossed kite $\mathsf K_\times$\\
$(1,0,2,2,0,2)$ & $1/8$ & chain $\mathsf C(2\Delta_\phi,2\Delta_\phi,2\Delta_\phi)$\\
$(1,1,1,1,1,2)$ & $1/4$ & symmetric kite $\mathsf K_\parallel$\\
$(2,0,1,1,0,3)$ & $1/12$ & chain $\mathsf C(\Delta_\phi,3\Delta_\phi,\Delta_\phi)$\\\bottomrule
\end{tabular}
\end{center}
where the normalized three-link chain is
\begin{equation}
 \mathsf C(v_1,v_2,v_3)=\Gamma(\tfrac d2)^2\prod_{i=1}^3
 \frac{\Gamma(\frac d2-v_i)}{\Gamma(v_i)}\,
 \frac{\Gamma(v_1+v_2+v_3-d)}{\Gamma(\frac{3d}2-v_1-v_2-v_3)} .
 \label{gff:ope:chain}
\end{equation}
The two kites are, through the needed order,
\begin{equation}
 \mathsf K_\times=\frac8{\epsilon^2}+\frac{2A}\epsilon+4R+O(\epsilon),\qquad
 \mathsf K_\parallel=\frac4{\epsilon^2}+\frac{2A}\epsilon+\frac{A^2}2+\frac72R+O(\epsilon).
 \label{gff:ope:kites}
\end{equation}
The symmetric kite is the master \eqref{gff:gs:masterJ} at the absolute exponents
$(2\Delta_\phi,2\Delta_\phi,2\Delta_\phi,2\Delta_\phi,4\Delta_\phi)$, that is,
Pattern~A with the coordinates \eqref{gff:gs:Acoords} $t=\epsilon$,
$p_0=p_1=2\epsilon$ and $u=v=0$; its expansion
through $\epsilon^0$ is the Laurent expansion of $\mathcal P_A$ in
\eqref{gff:gs:Apole}, because the remainder \eqref{gff:gs:Ajet} is $O(\epsilon)$. The
crossed kite is not one of the patterns A, B, T: both of its internal vertices
have absolute exponent sum $2d-2\epsilon$. Its finite part comes from contour
deformation of a two-fold Mellin--Barnes representation, in which the crossed
residues and the single-contour pieces give the displayed terms and the
remaining double contour is $O(\epsilon)$. Keeping only the poles would give
the wrong $a^{(4)}$. Summing the weights gives the internal kernel
\begin{equation}
 \mathsf S_4=\frac9{2\epsilon^2}+\frac A\epsilon+\frac{A^2}8+3R+\frac M9
 +\epsilon\Bigl[\frac{M(A+N)}{18}+AR+\frac U8-\frac94V\Bigr]+O(\epsilon^2),
 \label{gff:ope:S4}
\end{equation}
and the renormalized amplitude
\begin{equation}
 \mathsf a_4=\Bigl(\frac9{2\epsilon^2}-\frac A\epsilon\Bigr)\mathsf L_2
 -\frac9{4\epsilon}\mathsf Y\mathsf L_3+\mathsf S_4\mathsf L_4 .
 \label{gff:ope:a4amp}
\end{equation}
Its finite limit is
\begin{equation}
 \mathsf a_4(0)=h_d^2\Bigl[\frac98\,\partial_b^2\mathsf J_d-A\,\partial_b\mathsf J_d
 +\Bigl(\frac{A^2}8+\frac34R+\frac M9\Bigr)\mathsf J_d\Bigr]_{(a,b)=(q,3q)} .
 \label{gff:ope:a4finite}
\end{equation}
With $\mathsf a_2$ through $\epsilon^2$, $\mathsf a_3$ through $\epsilon$ and
$g_*$ through $\epsilon^3$, the $M$ terms cancel and $\mathsf H_2$ of
\eqref{gff:ope:H123} follows. The full finite-$\epsilon$ Mellin--Barnes values of
both kites and the exact Gamma chains reproduce the assembled $a^{(4)}$ at
$d=1,2,3$ by extrapolation $\epsilon\to0$ to better than $10^{-13}$; at $d=2$ the
symmetric kite also has the exact complex-Selberg form
$2\gamma(a)^2\gamma(a+c)^2\gamma(2c)/[\gamma(2a+c)\gamma(2a+2c)\gamma(c)]$ with
$\gamma(z)=\Gamma(z)/\Gamma(1-z)$, $a=1-\Delta_\phi$, $c=-\Delta_\phi$.

\subsubsection{Five vertices}
The fourteen five-vertex graphs, with edges ordered
$(01,02,03,04,12,13,14,23,24,34)$ and roots $0,1$, have the finite internal
coefficients below (before the Wick weight). Graphs~1 and~13 are scaleless mass
terms. Graph~7 has no finite part but a pole: it is a marginal bubble times the
regular tetrahedron, $4Q_d/\epsilon+O(\epsilon)$, and the vanishing balanced first
derivative of the tetrahedron \eqref{gff:gs:Tjet} is what removes its finite part.
At $d=2$ the package also evaluates this graph as an exact three-variable
complex-Selberg integral at finite $\epsilon$; with the pole removed, its
extrapolation to $\epsilon=0$ is $1.3\times10^{-36}$, consistent with a zero
finite part (task \texttt{integrals}, record
\path{lri_ope_reproducibility/results/verification/a5_numeric/verification_results.json}).
\begin{center}\small
\renewcommand{\arraystretch}{1.25}
\needspace{24\baselineskip}
\begin{longtable}{@{}rllp{82mm}@{}}
\toprule
graph & edge word & weight & finite internal coefficient\\\midrule\endfirsthead
\toprule
graph & edge word & weight & finite internal coefficient\\\midrule\endhead
1 & \texttt{0003111300} & $1/36$ & $0$\\
2 & \texttt{0003120210} & $1/12$ & $-M(A+N)$\\
3 & \texttt{0012102300} & $1/24$ & $3M(A+N)$\\
4 & \texttt{0012111210} & $1/2$ & $5A^3/18+65AR/6+43U/36-163V/18$\\
5 & \texttt{0012120120} & $1/4$ & $-7A^3/18+77AR/6+7U/36-247V/18$\\
6 & \texttt{0012210111} & $1/4$ & $-A^3/3+11AR+11U/6-13V/3$\\
7 & \texttt{0111111111} & $1/6$ & $0$ (simple pole $4Q_d/\epsilon$)\\
8 & \texttt{1002011310} & $1/6$ & $2M(A+N)$\\
9 & \texttt{1002020220} & $1/16$ & $-40V$\\
10 & \texttt{1002110211} & $1/4$ & $2A^3/9+26AR/3+11U/9-166V/9$\\
11 & \texttt{1011011220} & $1/8$ & $8A^3/9+32AR/3+23U/9-160V/9$\\
12 & \texttt{1011101211} & $1/2$ & $3A^3/2+13AR/2+23U/12-17V/2$\\
13 & \texttt{2001001311} & $1/24$ & $0$\\
14 & \texttt{2001010221} & $1/8$ & $5M(A+N)$\\\bottomrule
\end{longtable}
\end{center}
The weighted sum gives
\begin{equation}\begin{aligned}
 \mathsf S_5={}&\frac{18}{\epsilon^3}+\frac{11A}{\epsilon^2}
 +\frac{b_3/9+3A^2/2+45R/2+4M/3}{\epsilon}\\
 &+\frac78A^3+M(A+N)+\frac{145}8AR+\frac{43}{16}U-\frac{181}8V+O(\epsilon),
\end{aligned}\label{gff:ope:S5}\end{equation}
whose finite term is new integral data; pole cancellation does not fix it. The
renormalized amplitude is
\begin{equation}
 \mathsf a_5=\Bigl[\frac{18}{\epsilon^3}-\frac{10A}{\epsilon^2}+\frac{b_3}{9\epsilon}\Bigr]\mathsf L_2
 +\Bigl[-\frac{27}{2\epsilon^2}+\frac{9A}{4\epsilon}\Bigr]\mathsf Y\mathsf L_3
 +\frac{12}\epsilon\mathsf S_4\mathsf L_4-\mathsf S_5\mathsf L_5 .
 \label{gff:ope:a5amp}
\end{equation}
The $O(\epsilon)$ term of $\mathsf S_4$ reaches the finite part through
$12/\epsilon$. After the lower Taylor terms and $g_*$ are inserted, all $M$ and
$MN$ terms cancel and $\mathsf H_3$ of \eqref{gff:ope:H123} remains.

\subsubsection{Six vertices}
The 64 six-vertex graphs give the kernel
\begin{equation}\begin{aligned}
 \mathsf S_6={}&\frac{135}{2\epsilon^4}+\frac{153A}{2\epsilon^3}
 +\frac{399A^2/8+23M/2+1053R/8+9Q_d}{\epsilon^2}\\
 &+\frac1\epsilon\Bigl[\frac{365}{16}A^3+\frac{283}{24}AM+\frac{2979}{16}AR
 +\frac92AQ_d+\frac{81}8MN+\frac12P_4+\frac{985}{32}U-\frac{2733}{16}V\Bigr]\\
 &+\mathsf w_6+O(\epsilon).
\end{aligned}\label{gff:ope:S6}\end{equation}
All four pole coefficients agree with the counterterm constraints, including
the parts multiplying derivatives of the external-leg integral; this checks the
subtractions and does not determine $\mathsf w_6$. The new data at sixth order
are $\mathsf w_4=[\epsilon^2]\mathsf S_4$, $\mathsf w_5=[\epsilon]\mathsf S_5$,
$\mathsf w_6=[\epsilon^0]\mathsf S_6$ and $b_5$:
\begin{align}
 \mathsf w_4={}&\frac{A^4}{384}+\frac{7A^2M}{288}+\frac{7A^2R}{64}+\frac{7AMN}{144}
 +\frac{11AU}{192}-\frac{47AV}{96}\notag\\
 &+\frac{\mathsf c_\times}4+\frac{7MN^2}{288}
 -\frac{M\widetilde N}{288}+\frac{103R^2}{128}+\frac{85X}{96},\label{gff:ope:w4}\\
 \mathsf w_5={}&\frac{113A^4}{576}+\frac{55A^2M}{96}+\frac{713A^2R}{96}
 +\frac{55AMN}{48}+\frac{619AU}{288}-\frac{2071AV}{144}+\frac92\mathsf c_\times
 +\frac{55MN^2}{96}+\frac{7MR}{12}\notag\\
 &+\frac{7M\widetilde N}{96}+\frac{643R^2}{64}+\frac32RQ_d+\frac{2179X}{192}
 -\frac{5Z_d}{12},\label{gff:ope:w5}\\
 \mathsf w_6={}&\frac{1139A^4}{128}+\frac{1697A^2M}{192}+\frac{9057A^2R}{64}
 +\frac98A^2Q_d+\frac{1513AMN}{96}+\frac{2549AU}{64}\notag\\
 &-\frac{6545AV}{32}+\frac{189}4\mathsf c_\times+5\KK_d+\frac14\LL_d+\frac{11M^2}{108}
 +\frac{459MN^2}{64}+\frac{89MR}8+\frac{123M\widetilde N}{64}\notag\\
 &+\frac{11345R^2}{128}+\frac{231}8RQ_d+\frac{40211X}{384}-\frac{89Z_d}8.
 \label{gff:ope:w6}
\end{align}
Here $\mathsf c_\times$ is the quadratic Taylor coefficient of the crossed-kite pattern
along its one undetermined direction. It is kept as a free symbol and cancels:
\begin{equation}
 10\cdot\frac14-\frac43\cdot\frac92+\frac2{27}\cdot\frac{189}4=0 ,
 \label{gff:ope:crossCancel}
\end{equation}
because the scalar assembly is
\begin{equation}\begin{aligned}
 \mathsf H_4={}&\frac{4811A^4}{1296}+\frac{23A^2M}{162}-\frac{839A^2R}{216}
 +\frac{253A^2Q_d}{81}-\frac{2AMN}{81}+\frac{56AP_4}{81}+\frac{521AU}{324}
 +\frac23AV-\frac{M^2}{729}\\
 &-\frac{MR}{18}+\frac{4MQ_d}{81}-\frac{161R^2}{144}-\frac{13RQ_d}{27}
 +\frac{28Q_d^2}{81}-\frac{35X}{24}
 +10\mathsf w_4-\frac43\mathsf w_5+\frac2{27}\mathsf w_6-\frac2{243}b_5 .
\end{aligned}\label{gff:ope:H4assembly}\end{equation}
Substituting \eqref{gff:ope:w4}--\eqref{gff:ope:w6} and \eqref{gff:ope:b5} removes
$M$, $N$, $\widetilde N$ and $\mathsf c_\times$ and gives \eqref{gff:ope:H4}. The masters
survive only through $\KK_d,\LL_d,Z_d,P_5$; no further master enters.

\subsubsection{Finite parts and local subtractions}\label{gff:ope:method}
Graphs that are not reduced by convolution \eqref{gff:gs:convolution} and
star--triangle are completed conformally as in \eqref{gff:gs:completion}. Their edge
exponents are written $\alpha_e=q\,r_e+\epsilon(v_e+\vartheta\,t_e)$ with
rational $r_e,v_e,t_e$; the auxiliary coordinate $\vartheta$ along the
direction $(t_e)$ separates the logarithmic collision poles. For a UV cluster
$\gamma$ the degree $(|\gamma|-1)\frac d2-\sum_{e\subset\gamma}\alpha_e$ is
$\epsilon$ times a linear function of $\vartheta$; IR poles are treated after
inversion. At each critical
hyperplane the residue factorizes into the normalized anchored subgraph period
and its contracted cograph, and the poles are subtracted before the analytic
remainder is reconstructed. With $m$ remaining integrations, sixth order needs
exponent jets through degree $4-m$. The pole-subtracted remainder is a
polynomial of that degree in the deformations; its restrictions to reference
hyperplanes with a vanishing completion edge reduce by star--triangle, and exact
rational linear algebra checks that these restrictions determine it. The jets
used are those of the patterns A, B and T, \eqref{gff:gs:Ajet}, \eqref{gff:gs:Bjet} and
\eqref{gff:gs:Tjet}, together with the wheel derivatives \eqref{gff:gs:wheelJets}; Gamma
jets are expanded in abstract polygammas, and no fit in the dimension is
made. The one direction not fixed by these restrictions is $\mathsf c_\times$.

The local counterterms are the recursive forest formula of
Section~\ref{gff:app:generalSpecial},
$Z_G=-\operatorname{PP}[I_G+\sum_{\mathcal F\neq\varnothing}
(\prod_{\gamma\in\mathcal F}Z_\gamma)I_{G/\mathcal F}]$ over disjoint proper
logarithmic subgraphs, with radial factor $2/(L\epsilon)$ for an $L$-loop
vertex graph. With the enumeration's sign convention,
\begin{equation}
 b_L=L(-1)^L\sum_Gw_4(G)\,[\epsilon^{-1}]Z_G,\qquad
 c_L=L(-1)^L\sum_Gw_2(G)\,[\epsilon^{-1}]Z_G ,
 \label{gff:ope:bLcL}
\end{equation}
over the 129 five-loop vertex classes of Section~\ref{gff:app:generalSpecial}, with
$w_4$ and $w_2$ the four-point and mass weights. The mass-weighted sum is
computed separately; $c_5$ is not obtained by fitting a thermal coefficient.
The reconstruction presumes the master identities and the analytic pole
structure of Parts~II--III. Alternative rational directions test its
consistency inside that framework, not its independence from those inputs.

\subsection{Values at \texorpdfstring{$d=1,2,3$}{d=1,2,3}}\label{gff:ope:values}
With the closed forms of Section~\ref{gff:c1:closedForms} at $d=1,2$ and the $d=3$
values of Sections~\ref{gff:d3:values} and \ref{gff:gdk:numerics} (with $\LL_3$ and $Z_3$
from the same evaluation as $C_5(3)$), \eqref{gff:ope:convolution} gives
\begin{center}\small
\begin{tabular}{@{}crrr@{}}\toprule
 & $d=1$ & $d=2$ & $d=3$\\\midrule
$a^{(2)}$ & $-0.139409946645025$ & $-0.0370370370370370$ & $-0.0254241874995486$\\
$a^{(3)}$ & $-0.850425350705174$ & $-0.118399443073569$ & $-0.0667823808153469$\\
$a^{(4)}$ & $1.67987388417139$ & $0.0717109289555818$ & $-0.0279145897482122$\\
$a^{(5)}$ & $-26.2720127156913$ & $-0.751719824443267$ & $-0.156637777571136$\\
$a^{(6)}$ & $450.850442331799$ & $6.09946795238583$ & $0.507382663942981$\\\midrule
$13Z_d-\LL_d-12\KK_d$ & $15277.4261926043$ & $913.335423400404$ & $163.599472783843$\\\bottomrule
\end{tabular}
\end{center}
The digits shown are exact roundings: the master values are known to at least
$40$ digits at all three dimensions. The supplied package evaluates the $d=3$
nonprimitive masters by double-precision spectral quadrature with analytic spin
tails; its $\KK_3$, $\LL_3$, $Z_3$ agree with this note's values to about 12
significant digits (largest difference $1.4\times10^{-11}$, in $\LL_3$, a
relative $2.2\times10^{-13}$), its $a^{(6)}(3)=0.507382663943009$ differs
from the table by $2.8\times10^{-14}$, and its $d=2$ value before the closed
forms of Section~\ref{gff:gc:d2closedForms} differed by $2.7\times10^{-13}$. The
table entries are coefficients, not estimates of $a_{\phi^2}$ at finite
$\epsilon$. At $d=1,2$ they alternate in sign from $a^{(3)}$ on, and at all
three dimensions they grow quickly, so a truncated series at moderate
$\epsilon$ is not reliable (Section~\ref{gff:ope:bootstrap}).

\subsection{Two dimensions}\label{gff:ope:d2}
At $d=2$, $q=\frac12$ and $A=4\log2$, $R=X=0$, $U=\psi_2(\frac12)=-14\zeta(3)$,
$V=2\psi_2(1)=-4\zeta(3)$, $\mathsf e_j=j!\,2^{j+1}$. The external-leg ratio is elementary,
$\Gamma(-\frac12-\frac\epsilon4)/\Gamma(\frac12-\frac\epsilon4)=-2/(1+\frac\epsilon2)$, so
\eqref{gff:ope:factorized} reads
\begin{equation}
 a_{\phi^2}(2,\epsilon)=-\frac{\epsilon^2}{27}\,
 \frac{1+\mathsf H_1\epsilon+\mathsf H_2\epsilon^2+\mathsf H_3\epsilon^3+\mathsf H_4\epsilon^4}{1+\epsilon/2}
 +O(\epsilon^7).
 \label{gff:ope:d2ratio}
\end{equation}
Every master in $\mathsf H_1,\dots,\mathsf H_4$ is known in closed form at $d=2$: $Q_2$ of
\eqref{gff:gc:Q2}; $P_4(2)$ and $P_5(2)$ of \eqref{gff:fp:P4two}, \eqref{gff:fp:P5two};
$\LL_2=-2\log2\,P_4(2)$ of \eqref{gff:gc:L2closed}, proved in Section~\ref{gff:gc:L2proof};
$\KK_2=32\Cat^2-28\zeta(3)\log2$ of \eqref{gff:gc:K2closed}, proved in Section~\ref{gff:gc:K2proof} for the
master-path definition \eqref{gff:gs:Kdefinition}, which is the definition used in \eqref{gff:ope:H4}; and
$Z_2=4\Cat Q_2$ of \eqref{gff:gs:Z2exact}, proved in Section~\ref{gff:gs:Z2section}. Given the supplied
calculation of $\mathsf H_1,\dots,\mathsf H_4$ (Section~\ref{gff:ope:status}), the coefficients are therefore
exact. With $Q_2$, $P_4=P_4(2)$ and $P_5=P_5(2)$,
\begin{equation}\begin{aligned}
 a^{(3)}(2)&=\frac1{54}-\frac{16\log2}{81},\qquad
 -27\,a^{(4)}(2)=\frac14-\frac{8\log2}3+\frac{100\log^22}9-\frac49Q_2,\\
 -27\,a^{(5)}(2)&=-\frac18+\frac{4\log2}3-\frac{50\log^22}9+\frac{272\log^32}{27}-\frac{485}{108}\zeta(3)\\
 &\quad+\Bigl(\frac29-\frac{16\log2}9\Bigr)Q_2+\frac4{27}P_4,\\
 -27\,a^{(6)}(2)&=\frac1{16}-\frac{2\log2}3+\frac{25\log^22}9-\frac{136\log^32}{27}-\frac{28\log^42}{81}\\
 &\quad+\frac{485}{216}\zeta(3)-\frac{2933}{324}\zeta(3)\log2-\frac{896}{81}\Cat^2\\
 &\quad+\Bigl(\frac{350}{81}\Cat-\frac{56\log^22}{81}+\frac{8\log2}9-\frac19\Bigr)Q_2-\frac{14}{81}Q_2^2
 +\Bigl(\frac{70\log2}{81}-\frac2{27}\Bigr)P_4-\frac{16}{81}P_5 ,
\end{aligned}\label{gff:ope:d2exact}\end{equation}
with $a^{(2)}(2)=-\frac1{27}$. So $a^{(6)}(2)$ needs no constant beyond those of $C_5(2)$ in
\eqref{gff:gc:C5twoClosed}: $\Gamma(\frac14)$, $\pi$, $\log2$, $\Cat$, $\zeta(3)$, $F_\pm$ and the five values
$\mathcal Y_r$. Its value is $a^{(6)}(2)=6.0994679523858305821\ldots$

\paragraph{The thermal map at $d=2$.} The extra combination of \eqref{gff:ope:logmap} is
\begin{equation}
 13Z_2-\LL_2-12\KK_2=52\Cat Q_2+2\log2\,P_4(2)-384\Cat^2+336\,\zeta(3)\log2
 =913.33542340040416069629\ldots,
 \label{gff:ope:d2extra}
\end{equation}
and \eqref{gff:gd:d2exact}, \eqref{gff:gc:d2coefficients} give, with $C_j$ standing for $C_j(2)$ in this paragraph,
\begin{equation}
 Q_2=2\log^22-\frac{27}{2}C_3,\qquad
 P_4(2)=\frac{81}{2}C_4+8\log^32-27\log2\,C_3+\frac{359}{16}\zeta(3).
 \label{gff:ope:d2eliminate}
\end{equation}
Inserting these into \eqref{gff:ope:logmap} removes \eqref{gff:ope:d2extra}:
\begin{equation}\begin{aligned}
 \mathsf l_1(2)&=\frac{16\log2}{3}-\frac12,\qquad \mathsf l_2(2)=6C_3-4\log^22+\frac18,\\
 \mathsf l_3(2)&=6C_4-12\log2\,C_3+\frac{304}{81}\log^32-\frac76\zeta(3)-\frac1{24},\\
 \mathsf l_4(2)&=6C_5-11\log2\,C_4+\frac{46\log^22-26\Cat}{3}C_3-9C_3^2-\frac{380}{81}\log^42\\
 &\quad+\frac{104}{81}\Cat\log^22-\frac{128}{27}\Cat^2+\frac{1303}{216}\zeta(3)\log2+\frac1{64}.
\end{aligned}\label{gff:ope:d2log}\end{equation}
At $d=2$ the one-point coefficients through $a^{(6)}$ thus follow from $C_3(2)$, $C_4(2)$,
$C_5(2)$ and the constants $\log2$, $\Cat$, $\zeta(3)$. Numerically $\mathsf l_1,\dots,\mathsf l_4$
are $3.19678496299$, $-7.04591213159$, $37.3758123361$ and $-277.339425824$.

\paragraph{Checks at $d=2$.} The following were checked with code of their own
(\path{ope_checks/d2/ope_d2_check.py}, record \path{ope_d2_check.json}), which takes
every $d=2$ input from this note. Exactly: the logarithm of \eqref{gff:ope:d2ratio}, with
$\mathsf H_1,\dots,\mathsf H_4$ of \eqref{gff:ope:H123}, \eqref{gff:ope:H4}, the closed forms above and
$C_3(2)$, $C_4(2)$, $C_5(2)$ from \eqref{gff:gc:d2coefficients}, gives \eqref{gff:ope:d2log}, and also
\eqref{gff:ope:logmap} at $d=2$; \eqref{gff:ope:d2exact} and the second form of \eqref{gff:ope:d2extra} follow
in the same way, and the two forms of $C_5(2)$, \eqref{gff:gc:d2coefficients} and
\eqref{gff:gc:C5twoClosed}, agree. All residuals are zero. Numerically, at 60 digits: the closed-form
masters reproduce the supplied package's $6Q_2$, $P_4(2)$, $P_5(2)$ and $Z_2$ to
$2\times10^{-58}$, and \eqref{gff:ope:d2log} with the numerical $C_3(2)$, $C_4(2)$, $C_5(2)$ reproduces
\eqref{gff:ope:convolution} to $10^{-59}$. These checks take $\mathsf H_4$ as given. They test the
two-dimensional specialization and the shared masters, not the six-vertex calculation.

The five-loop coefficient $b_5$ enters $a^{(6)}$ on its own, not only through $D_5=3c_5-b_5$ of
\eqref{gff:gs:D5} (Section~\ref{gff:ope:RG}). Our two-dimensional productions use the finite-MOM
coefficients $\widehat b_L(z)=\sum_jb_{L,j}z^j$ of \eqref{gff:done:finiteCoefficients}. The conversion of
Section~\ref{gff:ope:checks} (item ``Five-loop split'') to regulator-independent coefficients gives,
explicitly,
\begin{equation}
 b_5=b_{5,0}+\frac94b_{4,1}+\frac92b_{3,2}+\frac{27}4b_{2,3}-\frac12b_{2,1}b_{3,0}
 +\frac13b_{2,0}b_{3,1}+\frac12b_{2,0}b_{2,2}-\frac38b_{2,1}^2 ,
 \label{gff:ope:b5MS}
\end{equation}
where $b_{2,j}$ are the exact two-loop jets. With $D_5$ held at \eqref{gff:gs:D5}, the jets of the
seventh-order production (Section~\ref{gff:c7d2:production}) and of the eighth-order production
(Section~\ref{gff:app:C8d2}), which use independent random streams, give $b_5(2)=135708.65\pm0.27$ and
$135709.17\pm0.42$, together $135708.80\pm0.23$; \eqref{gff:ope:b5} gives $135708.425\ldots$. Without
that constraint the two productions give $135709.43\pm0.55$ and $135709.64\pm0.96$, because their
five-loop jets fluctuate along $D_5$ (by $-2.1$ and $-0.7$ standard errors). The same conversion
reproduces $b_4$ and $c_4$ of \eqref{gff:gs:b4}, \eqref{gff:gs:c4} within $2.0$ standard errors. Since
$a^{(6)}(2)$ changes by $3.05\times10^{-4}$ per unit of $b_5$, this checks the five-loop input of
$a^{(6)}(2)$ at the level of $10^{-4}$ in $a^{(6)}(2)$, against data of this note only
(\path{ope_checks/d2/ms_b5_from_jets.py}, records \path{ms_b5_c7d2.json},
\path{ms_b5_c8d2.json}).

\subsection{One dimension}\label{gff:ope:d1}
At $d=1$ every master in \eqref{gff:ope:H4} is a closed form proved in Parts~II--III:
$Q=\Gamma(\frac14)^4/\pi$, $P_4=13B_--3B_+$, $P_5$ in the five values
$\mathcal Y_r$ (Section~\ref{gff:c1:closedForms}), $\KK$ \eqref{gff:eq:KprovedFinal},
$\LL_1$ \eqref{gff:gdr:L1} and $Z_1=QH$ \eqref{gff:gs:d1Dictionary}. The coefficients
$a^{(2)}(1),\ldots,a^{(6)}(1)$ of Section~\ref{gff:ope:values} are therefore exact
expressions in the same fixed values as $C_5(1)$: $\Gamma(\frac14)$, $\pi$,
$\log2$, $\Cat$, $\zeta(3)$, $F_\pm$ and the $\mathcal Y_r$. They are
conditional only on the calculation itself (Section~\ref{gff:ope:status}). The
first two are
\begin{equation}
 a_{\phi^2}(1,\epsilon)=-\frac{2\pi^{3/2}\Gamma(\frac34)}{27\,\Gamma(\frac14)}\,\epsilon^2
 \Bigl[1+\Bigl(\frac{13\pi}{12}+\frac{16}{3}\log2-1\Bigr)\epsilon\Bigr]+O(\epsilon^4).
 \label{gff:ope:d1low}
\end{equation}
Both agree with the general formulas of \cite{gff:BLNV}: $a^{(2)}$ with its
quadratic coefficient for the Ising case, Eq.~(3.35), and $a^{(3)}$ with its
cubic coefficient for the long-range $O(N)$ model at $N=1$, Eq.~(3.47). Read
literally, its Ising cubic coefficient in Eq.~(3.35) and the sign of its
$O(N)$ quadratic coefficient in Eq.~(3.47) disagree with these values and with
its own Eq.~(3.36).

Two checks exist only at $d=1$. First, the package's reconstruction of all
six-vertex kernels and five-loop forest sums at 60 digits, with rational
regulator and reference directions other than those of the derivation and
$\mathsf c_\times=17$, runs at $d=1$. It differs from \eqref{gff:ope:H4} by
$7.6\times10^{-54}$ in $a^{(6)}(1)$, with residual poles below
$1.4\times10^{-52}$; it shares the master identities and the subtraction
framework (Section~\ref{gff:ope:status}). Second, the split of $D_5$ into $b_5$ and
$c_5$, which $a^{(6)}$ needs and \eqref{gff:gs:D5} does not fix, is compared at
$d=1$ with numerical graph integrations that do not use the master reduction
(Section~\ref{gff:ope:checks}). The residue certificate of
Appendix~\ref{gff:app:certificate} cannot serve for this: it omits the 28
five-loop classes of zero weight in $3c_5-b_5$, which is correct for $D_5$
but not for $b_5$ alone. Since $b_5$ enters \eqref{gff:ope:H4assembly} as
$-\frac2{243}b_5$, the coefficient $a^{(6)}(1)$ moves by
$-\frac2{243}a^{(2)}(1)=1.15\times10^{-3}$ per unit of $b_5$.

\subsection{A one-dimensional functional cross-check}\label{gff:ope:d1functional}
In one dimension the super sum rule of \cite{gff:GhoshSuperSum} gives a
determination of the low-order thermal coefficients that is independent of the
graphs. Its expansion parameter $\epsilon=1-4\Delta_\phi$ is ours at $d=1$.
Writing $\Delta_{\Eop}=2\Delta_\phi+\gamma^{(1)}\epsilon+\gamma^{(2)}\epsilon^2+\cdots$,
the functional $\widehat\beta_0=-\frac\epsilon4\widetilde\beta_0$ gives
$-\gamma^{(1)}+3(\gamma^{(1)})^2=0$, so $\gamma^{(1)}=\frac13=C_1+\frac12$ on
the nontrivial branch. At the next order, assuming that only $\Eop$
contributes, a high-precision numerical evaluation of the functional gives
$\gamma^{(2)}=\frac19(\pi+4\log2)=C_2(1)$. The cross-check therefore concerns
$C_1$ and $C_2(1)$; it does not test the one-point coefficients. At the next
order the sums over double-trace operators and over the twist-four family, and
the $\phi^4$ operator, contribute at $O(\epsilon^3)$ although each term starts
at $O(\epsilon^4)$, for example through
\begin{equation}
 \epsilon^4\sum_{n\ge1}n^{-1-\nu\epsilon}=\frac{\epsilon^3}{\nu}+\gamma_{\rm E}\epsilon^4+\cdots;
 \label{gff:ope:enhanced}
\end{equation}
their cancellation requires an exponent that \cite{gff:GhoshSuperSum} does not
determine, so \cite{gff:GhoshSuperSum} does not determine $C_3(1)$, and no
functional determination of it is made here. In that reference
$a_{\phi^2}$ denotes the squared defect OPE coefficient of $\phi^2$ in the
four-point function of $\phi$, not the bulk one-point coefficient of this part.

\subsection{Comparison with the numerical bootstrap}\label{gff:ope:bootstrap}
At fixed $d$ and $\Delta_\phi$, the bootstrap of \cite{gff:BLNV} constrains a region
in the plane $(a_{\phi^2},\Delta_{\Eop})$; $a_{\phi^2}$ is its horizontal
coordinate. Figure~11 of \cite{gff:BLNV} shows allowed regions at $d=3$ for four
external dimensions; their upper right corners (kinks) are the candidate
long-range Ising points. The comparison below inserts the series of this note
into one resummation prescription and places the result next to readings of
that figure. It does not rerun the bootstrap, it has no access to the
underlying numerical data, and it does not provide an excluding functional.
The package draws both comparisons, with its pixel traces of the published
boundaries and kinks, in \path{lri_ope_reproducibility/results/figures/}:
\path{3d_sixth_order_comparison.png} and \path{2d_sixth_order_comparison.png},
and, for the successive truncations and the two resummations at the weakest
coupling, \path{3d_weak_point_orders.png} and \path{2d_weak_point_orders.png}.
These plots use the older $C_6$ inputs; \path{results/bootstrap/current/}
holds the same plots with the current ones. The traces are approximate
readings, not bootstrap data, and the plots are not exclusion certificates.

\paragraph{Prescription.} For a truncated series $F(\epsilon)=\sum_kc_k\epsilon^k$
use the Borel transform $\sum_kc_kt^k/k!$, its $[L/M]$ Pad\'e approximant, and
the Laplace integral $\int_0^\infty e^{-t}[L/M](\epsilon t)\,dt$. Horizontally,
resum $a_{\phi^2}/\epsilon^2=a^{(2)}+\dots+a^{(6)}\epsilon^4$ with $[3/1]$ and
$[1/3]$, whose denominators have no positive real root at these points; $[2/2]$
has one at every point and is not used. Vertically, resum
$\Delta_{\Eop}$ through $C_6$ with $[4/2]$; the package also records $[3/2]$
through $C_5$ as an order diagnostic. Factoring out $\epsilon^2$ is part
of the prescription. No bootstrap point and no short-range endpoint is fitted,
and varying the Borel--Leroy parameter, the approximant or the factored power
would add uncertainty that has not been quantified.

\paragraph{Three dimensions.} The four external dimensions of Figure~11 are
equally spaced between $\Delta_\phi=\frac34$ and the short-range end
$\Delta_\phi=0.5181489$, the short-range Ising value used by \cite{gff:BLNV}
(Section~\ref{gff:ope:observable}), so $\epsilon=3-4\Delta_\phi$:
\begin{center}\small
\begin{tabular}{@{}rrrrrr@{}}\toprule
$\Delta_\phi$ & $\epsilon$ & $a_{\phi^2}$, $[3/1]$ & $a_{\phi^2}$, $[1/3]$ &
$\Delta_{\Eop}$, $[4/2]$ & kink of Fig.~11\\\midrule
$0.70363$ & $0.18548$ & $-0.00135602$ & $-0.00134382$ & $1.4738699$ & $(-0.00187,\,1.4700)$\\
$0.65726$ & $0.37096$ & $-0.00797138$ & $-0.00713938$ & $1.4541700$ & $(-0.00292,\,1.4500)$\\
$0.61089$ & $0.55644$ & $-0.0253361$ & $-0.0185657$ & $1.4383505$ & $(-0.0110,\,1.4318)$\\
$0.56452$ & $0.74192$ & $-0.0613651$ & $-0.0352226$ & $1.4248531$ & $(-0.0420,\,1.4154)$\\\bottomrule
\end{tabular}
\end{center}
The $[4/2]$ column uses $C_6(3)=33.5754$ of Section~\ref{gff:prec:d3production2026};
with the older input $33.57$ it is lower by $4.6\times10^{-8}$,
$1.2\times10^{-6}$, $7.6\times10^{-6}$ and $2.6\times10^{-5}$. The kinks are read
from the vector data of the published figure (\path{ope_checks/d3/check_fig11.py},
record \path{ope_checks/d3/check_fig11.json}; Section~\ref{gff:ope:checks}). The
figure file of the arXiv source, converted with \texttt{pdftocairo -svg}, draws
each region boundary as a polyline; its vertices are mapped to data
coordinates by the axis frame, and the kink is the vertex with the largest
turning angle in axis-normalized coordinates. Vertices that should lie on
round values, such as $a_{\phi^2}=-0.005$ or $\Delta_{\Eop}=1.45$, are
recovered to within $7\times10^{-7}$, so a kink is limited by the spacing of
the polyline vertices around it, $3\times10^{-4}$ to $4\times10^{-3}$ in
$a_{\phi^2}$, not by the reading. The supplied package used pixel readings
that agree to about one pixel.

At the weakest coupling, $\Delta_\phi=0.70363$, the drawn region at
$\Delta_{\Eop}=1.4739$ requires $a_{\phi^2}\lesssim-0.0032$, while both
resummations give $-0.0013$ to $-0.0014$: this resummed point lies outside the
drawn region. The truncated sums through $\epsilon^4$, $\epsilon^5$ and
$\epsilon^6$ are $-0.00133384$, $-0.00136823$ and $-0.00134757$, and the
cubic-order point that \cite{gff:BLNV} itself plots, $(-0.0013004,\,1.473078)$,
is also on the excluded side of its own boundary. The mismatch is therefore
not created by the higher orders; equally, it does not test them. At this
point the thermal terms $C_4\epsilon^4$, $C_5\epsilon^5$, $C_6\epsilon^6$ are
about $+0.00145$, $-0.00128$, $+0.00137$, so the vertical coordinate is not
known to six digits from apparent Pad\'e stability.

\paragraph{What the mismatch can and cannot mean.} An excluding functional must
be positive on every admissible exchanged dimension and spin. The bounds of
\cite{gff:BLNV} impose the long-range OPE relations on a grid of exchanged
dimensions (its App.~A.2: 3000 points, and up to 12000 in its refined
two-dimensional runs, where the kinks move up as the grid is refined). A
functional positive on the grid can be negative between grid points, which
would make a finite-grid boundary too restrictive. The two numerical
truncations have different logical effects: keeping fewer crossing
derivatives, with positivity still imposed on the continuum, only restricts
the available functionals and weakens the exclusions, so it cannot by itself
make a true solution inconsistent; sampling the exchanged dimension instead
can miss a negative interval. Near the generalized-free end the spectrum approaches
double-twist values at a rate set by $\epsilon^2$; the package estimates the
leading spin-2 shift as $-4\epsilon^2/[9d(d+2)]\approx-0.0010$ at the first
point, smaller than a grid spacing of order $0.007$ (this estimate was not
checked here). This is a possible mechanism, not a demonstrated one. The decisive
test would be to check the candidate functional between grid points, including
the smallest eigenvalue of each matrix sector and the large-dimension tail. If
positivity held and both coordinates had controlled errors, the discrepancy would
have to be resolved in the perturbation theory, the resummation or the spectral
assumptions. None of this has been done.

\paragraph{Two dimensions.} Figure~7 of \cite{gff:BLNV} shows, at $d=2$ and eight external
dimensions between $\Delta_\phi=0.13333$ and $0.425$, the kinks of runs with 3000, 6000 and 12000 grid
points and the authors' own two-sided Pad\'e estimates; it shows kinks, not allowed regions. With
$\epsilon=2-4\Delta_\phi$, the four weakest couplings give
\begin{center}\small
\begin{tabular}{@{}rrrrrrr@{}}\toprule
$\Delta_\phi$ & $\epsilon$ & $a_{\phi^2}$, $[3/1]$ & $a_{\phi^2}$, $[1/3]$ &
$\Delta_{\Eop}$, $[4/2]$ & kink of Fig.~7 & $\Delta_{\Eop}$, $[5/3]$ to $C_8$\\\midrule
$0.425$ & $0.3$ & $-0.00654824$ & $-0.00601150$ & $0.9678196$ & $(-0.0060,\,0.9600)$ & $0.9679055$\\
$0.35$ & $0.6$ & $-0.0414334$ & $-0.0277046$ & $0.9578452$ & $(-0.0430,\,0.9602)$ & $0.9590956$\\
$0.275$ & $0.9$ & $-0.134446$ & $-0.0621047$ & $0.9629664$ & $(-0.1420,\,0.9660)$ & $0.9680003$\\
$0.2$ & $1.2$ & $-0.325205$ & $-0.105209$ & $0.9803317$ & $(-0.3750,\,0.9878)$ & $0.9929583$\\\bottomrule
\end{tabular}
\end{center}
The $[4/2]$ column uses $C_6(2)=362.366$ of Section~\ref{gff:app:C6Precision}; with the older input
$362.3$ it is lower by $1.7\times10^{-6}$, $2.6\times10^{-5}$, $1.1\times10^{-4}$ and $2.7\times10^{-4}$,
and the allowance $\pm0.004$ of $C_6(2)$ moves the first row by $10^{-7}$. As at $d=3$, $[2/2]$ has a
positive Borel pole at every point. The kinks are those of the 12000-point runs. Each is the centre of
its cross marker, read from the marker's placement in the vector figure file of the arXiv source of
\cite{gff:BLNV} (\path{ope_checks/d2/check_fig7.py}, record \path{check_fig7.json}), which places it
to about $3\times10^{-9}$ in $a_{\phi^2}$ and $2\times10^{-10}$ in $\Delta_{\Eop}$;
the supplied package's pixel readings agree to $2\times10^{-4}$ in $a_{\phi^2}$, except at
$\Delta_\phi=0.425$, where its $-0.00676$ is $7.6\times10^{-4}$ off, and to $8\times10^{-5}$ in
$\Delta_{\Eop}$.

At the weakest coupling, $\Delta_\phi=0.425$, the kink at $a_{\phi^2}=-0.0060$ lies between the two
horizontal resummations, while the resummed $\Delta_{\Eop}=0.96782$ lies $0.0078$ above it. The
truncated sums $\sum_{n\le k}a^{(n)}(2)\,\epsilon^n$ at this point are $-0.00594925977$, $-0.00777593895$
and $-0.00332942681$ for $k=4,5,6$, so the last available term is not an error estimate. The kinks
of the same figure move up as the grid is refined, from $0.9575$ at 3000 points to $0.9592$ and
$0.9600$, and the authors' own Pad\'e estimates at this $\Delta_\phi$ are $0.9695$ with all their
perturbative data and $0.9723$ without the order $\epsilon^3$; \cite{gff:BLNV} itself notes that the
kinks of best agreement lie noticeably below its perturbative line. So the gap was present before
the orders computed here, as at $d=3$. At the three stronger couplings the resummed points lie
$0.0023$ to $0.0074$ below the kinks. The last column is a diagnostic only: it adds $C_7(2)$ of
Section~\ref{gff:app:C7d2} and $C_8(2)$ of Section~\ref{gff:app:C8d2} in $[5/3]$, which has no positive Borel
pole at the four points ($[4/3]$, $[3/4]$, $[4/4]$ and $[3/5]$ have one at $\epsilon=0.3$). It moves
$\Delta_{\Eop}$ at $\epsilon=0.3$ by $8.6\times10^{-5}$, and by at most $7.4\times10^{-6}$ more when
$C_7(2)$ and $C_8(2)$ vary over their allowances; at the three stronger couplings it moves it by
$1.3\times10^{-3}$, $5.0\times10^{-3}$ and $1.3\times10^{-2}$, comparable to the distances from the kinks.
Neither column is an error estimate (\path{ope_checks/d2/ope_d2_check.py}).

\subsection{Checks made for this note}\label{gff:ope:checks}
The following were done independently of the supplied code, with this note's
inputs; the scripts and their outputs are in \path{ope_checks/}.
\begin{enumerate}
\item \emph{Rerun.} All tasks of the package except the \TeX{} build (low-order
checks, symbolic reconstruction of the five- and six-vertex kernels and of the
129-class forest sums, the $d=1$ alternative-regulator reconstruction at 60
digits, the four-vertex finite-regulator checks, the $d=2,3$ spectral masters,
and the comparison plots) pass on a clean extraction on the symmetry cluster
(Intel Xeon Gold 6148, Python~3.13.3, the pinned library versions). Every
output equals the shipped one byte for byte, apart from recorded run times.
\item \emph{Exact algebra} (32 identities, own code): the package's
$C_2,\dots,C_5$ equal \eqref{gff:gs:C0C4}--\eqref{gff:gs:C5}; its $b_3,c_3,b_4,c_4$
equal \eqref{gff:gd:bc3}, \eqref{gff:gs:b4}, \eqref{gff:gs:c4}; $3c_5-b_5$ from
\eqref{gff:ope:b5}--\eqref{gff:ope:c5} equals \eqref{gff:gs:D5}, and Lagrange inversion with
these coefficients gives \eqref{gff:gs:C5}; \eqref{gff:ope:gstar} solves
\eqref{gff:ope:beta}; \eqref{gff:ope:H4assembly} with \eqref{gff:ope:w4}--\eqref{gff:ope:w6} and
\eqref{gff:ope:b5} gives \eqref{gff:ope:H4}; \eqref{gff:ope:fj}, \eqref{gff:ope:H23thermal},
\eqref{gff:ope:H4thermal}, \eqref{gff:ope:logmap} and \eqref{gff:ope:expmap} are exact.
\item \emph{Low orders.} $a^{(2)}$ and $a^{(3)}$ at $d=2,3$ equal Eq.~(3.36) of
\cite{gff:BLNV}, \eqref{gff:ope:BLNVvalues}, to 35 digits. At $d=1$ they equal the
quadratic coefficient of Eq.~(3.35) and the $N=1$ cubic coefficient of
Eq.~(3.47) of \cite{gff:BLNV}, and the ratios $R^{(1)}$, $R^{(2)}$ of its
Eq.~(3.34), to 80 digits; the cubic coefficient of its Eq.~(3.35) read
literally gives $-1.50734$ instead of $a^{(3)}(1)=-0.850425$, and the
quadratic coefficient of its Eq.~(3.47) has the opposite sign
(\path{ope_checks/d1/}, with its recorded run).
\item \emph{Amplitudes.} With \eqref{gff:ope:Ln}, \eqref{gff:ope:bubble},
\eqref{gff:ope:a23amp}, \eqref{gff:ope:a4amp}, \eqref{gff:ope:a5amp}, the kernels
\eqref{gff:ope:S4}, \eqref{gff:ope:S5} and $g_*$, the amplitudes $\mathsf a_3$,
$\mathsf a_4$, $\mathsf a_5$ are finite and the re-expanded
$\sum_ng_*^n\mathsf a_n$ reproduces $a^{(2)},\dots,a^{(5)}$ of
\eqref{gff:ope:convolution} to $10^{-31}$ at $d=3$ and at $d=2.6$. The bubble
expansion \eqref{gff:ope:bubble} and the closed form \eqref{gff:ope:Jd} (by direct
quadrature at $d=1$) were checked numerically.
\item \emph{Kernels.} The symmetric kite \eqref{gff:ope:kites} through $\epsilon^0$
follows from \eqref{gff:gs:Apole} at $d=1$, $2.6$ and $3$, including its $R$ term;
the four-vertex table with the exact chains \eqref{gff:ope:chain} gives the pole and
constant terms of \eqref{gff:ope:S4}; the five-vertex table sums to the finite part
of \eqref{gff:ope:S5}. The crossed kite \eqref{gff:ope:kites} enters the package
through a two-fold Mellin--Barnes representation, whose finite-$\epsilon$ values
reproduce the assembled $a^{(4)}$ (Section~\ref{gff:ope:graphs}). At $d=2$ this
note's kite formula \eqref{gff:done:kiteFactorization}, evaluated at
$\epsilon=0.02,0.04,0.06,0.1,0.15$ (the slowly convergent ${}_3F_2(1)$ after
Thomae's transformation), reproduces those finite-$\epsilon$ values of
$\mathsf K_\times$ to $2\times10^{-14}$ or better, and those of
$\mathsf K_\parallel$ to $10^{-16}$. At $d=1$ the full-line Selberg formula of
Section~\ref{gff:app:selberg} gives $\mathsf K_\parallel$ to all 15 digits of the
package's values; for $\mathsf K_\times$ at $d=1$ only a double-precision direct
quadrature was made, which agrees to $0.3\%$, within its own error (it misses
the exact $\mathsf K_\parallel$ by $0.7\%$). Since $R=0$ at $d=2$, the $4R$ term
of $\mathsf K_\times$ is tested only through the package's representation.
\item \emph{Five-loop split.} At $d=1$, $b_5$ and $c_5$ of
\eqref{gff:ope:b5}--\eqref{gff:ope:c5} agree to about $10^{-12}$ relative with the
two-channel five-loop graph sums of Section~\ref{gff:sec:checks} at quadrature
degree 18, renormalized with the degree-28 lower-loop inputs of the $C_6(1)$
calculation (records \path{amplitude5_maxgap_18.json} and
\path{amplitude5_span_18.json} in \path{lri_numerical_resolution/verification/records/},
assembled in \path{LRI_C6_1D_HPC/reference/assembly_lower6.json}). These are
numerical graph integrations that do not use the master reduction
(Section~\ref{gff:ope:d1}).
At $d=2$, the $M$-dependent parts of \eqref{gff:ope:b5} and \eqref{gff:ope:c5} equal
the minimal-subtraction power sums $b_5^p$ and $c_5^p$ of \eqref{gff:done:powerSums}
to all 21 printed digits, and their $b_3$, $b_4$ parts equal $b_3^p$, $b_4^p$
there. At $d=2$ and $d=3$ the complete $b_5$ was also extracted from the
finite-scheme Taylor jets of the seventh- and eighth-order productions. A
regular redefinition $g=g'+\sum_k\rho_k(z)g'^{k+1}$ with an operator
normalization brings the finite-scheme $\beta$ and $\gamma_r$ to
regulator-independent coefficients. Because of the $-zg'$ term this
redefinition is unique, so the result is the minimal-subtraction form used
here. It needs $b_2$ through $z^3$, $b_3$ through $z^2$, $b_4$ through $z^1$
and $b_5$ at $z^0$, and likewise for $c_L$. Holding $D_5$ at its proved value
\eqref{gff:gs:D5}, the $C_7(3)$ and $C_8(3)$ productions give
$b_5(3)=18822.676\pm0.097$ and $18822.42\pm0.28$, against $18822.627$ from
\eqref{gff:ope:b5}. The same conversion reproduces $b_3,c_3,b_4,c_4$ of
\eqref{gff:gd:bc3}--\eqref{gff:gs:c4} within the sampling errors. At $d=2$ it gives
$b_5(2)=135708.80\pm0.23$, against $135708.43$.
The individual $b_5$ sums of Section~\ref{gff:app:C6d3Complete} cannot serve
here, because they depend on trial core assignments.
\item \emph{Values.} The table of Section~\ref{gff:ope:values} was computed from
this note's master values, which reproduce $C_2,\dots,C_5$ of the
coefficient table of Section~\ref{gff:sec:guide} at $d=1,2,3$. The log map with
those $C_2,\dots,C_5$ and the master value of $13Z_d-\LL_d-12\KK_d$ reproduces it
to the precision of the displayed $C_j$.
\item \emph{Bootstrap rows.} An own Borel--Pad\'e evaluation reproduces the
four $d=3$ rows of Section~\ref{gff:ope:bootstrap} to all printed digits and finds
the positive Borel pole of $[2/2]$ at every point. The kinks and the
cubic-order point of Figure~11 were taken from the vector paths of the figure
file in the arXiv source of \cite{gff:BLNV}: the plotted point lies at
$(-0.0013004,1.473078)$, the package's recomputation of it is
$(-0.00130080,1.473065)$, and the package's four kink readings agree with the
vector kinks to within $9\times10^{-5}$ in $a_{\phi^2}$ and $2\times10^{-4}$ in
$\Delta_{\Eop}$.
\end{enumerate}

\subsection{Reproduction}\label{gff:ope:reproduce}
The package \path{lri_ope_reproducibility/} contains the low-order,
four-, five- and six-vertex and relation engines, the graph inventories
(\path{data/a6/graphs_n*.json}, \path{vertex_n*.json}), the master inputs, the
bootstrap figure readings, and stored certificates. Its runner stages each
engine in a temporary directory and does not overwrite inputs:
\begin{center}\small
\begin{tabular}{@{}lp{105mm}@{}}\toprule
task & content (run time on the cluster above)\\\midrule
\texttt{quick} & low-order diagrams, the thermal map, the coefficient table, and a regression against stored certificates (4~s)\\
\texttt{symbolic} & fresh graph inventories, five- and six-vertex symbolic kernels, forest sums, pole cancellation, $3c_5-b_5$, the compact result (2~min)\\
\texttt{numeric} & all sixth-order kernels and five-loop forest sums at 60 digits in $d=1$ with other rational regulator directions and $\mathsf c_\times=17$; agreement with \eqref{gff:ope:H4} to about $10^{-53}$, residual poles about $10^{-52}$ (5~min)\\
\texttt{integrals} & finite-regulator four-vertex kite and chain checks at $d=1,2,3$, the $d=2$ Selberg check, the five-vertex checks (1.5~min)\\
\texttt{masters} & $d=2,3$ spectral evaluation of $\KK_d,\LL_d,Z_d$ with varied quadrature (20~s)\\
\texttt{figures} & Borel--Pad\'e rows and comparison plots, with both $C_6$ inputs (30~s)\\\bottomrule
\end{tabular}
\end{center}
The command \texttt{python run.py}, followed by the task names and
\texttt{-{}-output-dir DIR}, runs them; each writes \path{verification/run_TASK.json}. The numerical
residuals of these tasks are consistency checks within the framework above,
not uncertainties of the coefficients.

\subsection{Open points}\label{gff:ope:open}
\begin{itemize}
\item An evaluation of selected five- and six-vertex finite parts by a method
that does not use the pattern jets, for example sector decomposition with all
local subtractions, would test the reconstruction beyond its alternative
directions. The $R$ term of the crossed-kite finite part \eqref{gff:ope:kites}
is checked only through the package's own Mellin--Barnes representation; an
evaluation at $d\ne2$ to many digits would close this.
\item The resummation has no controlled error: varying the Borel--Leroy
parameter, the approximants and the factored power, rejecting or treating
positive Borel poles, and propagating the input uncertainties remain to be done.
\item The bootstrap mismatch needs a continuum-positivity check of the
candidate functional, not more coefficients.
\item The apparent printed-formula inconsistencies in Eqs.~(3.35) and (3.47)
of \cite{gff:BLNV} (Sections~\ref{gff:ope:lowOrders} and \ref{gff:ope:d1}) should be
checked against the authors' ancillary calculation or clarified with them.
They should not be used to infer that the authors' numerical bootstrap code
shares them.
\item $a^{(7)}$ would need $\mathsf a_2,\dots,\mathsf a_6$ one order further, a
new finite $\mathsf a_7$ and $b_6$; neither $C_6$ nor the pattern of
\eqref{gff:ope:logmap} determines them, and new master derivatives may enter.
\end{itemize}

\appendix

\section{Complete retained residue certificate}\label{gff:app:certificate}
This appendix lists the residues that enter the one-dimensional assembly of Sections~\ref{gff:sec:shifts}--\ref{gff:sec:assembly}: the 26 non-wheel four-loop rows and 92 retained five-loop rows, as exact polynomials, so that the weighted sum can be inspected without the machine-readable certificate. The six wheel rows are in Section~\ref{gff:sec:wheelclassupdated}; the three primitive rows are determined by their periods; the 28 five-loop zero-weight classes contribute nothing to $3c_5-b_5$. Together these disjoint sets exhaust the stated graph list.

For each row, $i$ is the zero-based canonical class index and the word lists upper-triangular multiplicities in lexicographic vertex-pair order. The columns are \emph{RG weights}, not raw Wick weights:
\[
\omega_b=L(-1)^Lw_4/|\Aut G|,\qquad
\omega_c=L(-1)^Lw_2/|\Aut G|.
\]
A row with residue $R_j^{(L)}$ contributes $\omega_bR_j^{(L)}$ to $b_L$ and $\omega_cR_j^{(L)}$ to $c_L$. Equal residues are identified to keep the certificate legible. All expressions use the same $A,Q,\Cat,\mathcal K$ as the main result. Their transcription from the retained JSON is checked by exact symbolic substitution, including $\log2=(A-\pi)/4$. No numerical fit enters this table.
\subsection{4-loop rows and residue polynomials}
\begin{longtable}{@{}r l r r r@{}}\toprule Class & Edge word & $\omega_b$ & $\omega_c$ & Residue\\\midrule\endhead
0 & \texttt{0002011310} & $2$ & $\frac{2}{3}$ & $R_{0}^{(4)}$\\
1 & \texttt{0002020220} & $\frac{3}{4}$ & $\frac{1}{4}$ & $R_{0}^{(4)}$\\
2 & \texttt{0002021210} & $3$ & $\frac{1}{2}$ & $R_{0}^{(4)}$\\
3 & \texttt{0002110211} & $3$ & $1$ & $R_{0}^{(4)}$\\
4 & \texttt{0002111201} & $12$ & $2$ & $R_{0}^{(4)}$\\
5 & \texttt{0002112200} & $3$ & $\frac{1}{2}$ & $R_{0}^{(4)}$\\
6 & \texttt{0011011220} & $\frac{3}{2}$ & $\frac{1}{2}$ & $R_{1}^{(4)}$\\
7 & \texttt{0011012210} & $6$ & $1$ & $R_{2}^{(4)}$\\
8 & \texttt{0011101211} & $6$ & $2$ & $R_{2}^{(4)}$\\
9 & \texttt{0011102201} & $6$ & $1$ & $R_{3}^{(4)}$\\
10 & \texttt{0011102210} & $12$ & $2$ & $R_{4}^{(4)}$\\
11 & \texttt{0011102300} & $4$ & $\frac{2}{3}$ & $R_{5}^{(4)}$\\
12 & \texttt{0011111111} & $12$ & $2$ & $R_{6}^{(4)}$\\
13 & \texttt{0011111120} & $24$ & $4$ & $R_{7}^{(4)}$\\
14 & \texttt{0011201102} & $6$ & $1$ & $R_{8}^{(4)}$\\
15 & \texttt{0011201111} & $24$ & $4$ & $R_{9}^{(4)}$\\
16 & \texttt{0011201201} & $12$ & $2$ & $R_{3}^{(4)}$\\
17 & \texttt{0011202110} & $3$ & $1$ & $R_{10}^{(4)}$\\
18 & \texttt{0011202200} & $3$ & $\frac{1}{2}$ & $R_{11}^{(4)}$\\
19 & \texttt{0011211110} & $6$ & $1$ & $R_{12}^{(4)}$\\
20 & \texttt{0011301011} & $1$ & $\frac{1}{3}$ & $R_{0}^{(4)}$\\
21 & \texttt{0011301101} & $4$ & $\frac{2}{3}$ & $R_{13}^{(4)}$\\
22 & \texttt{0012102200} & $6$ & $0$ & $R_{11}^{(4)}$\\
23 & \texttt{0012111110} & $24$ & $0$ & $R_{6}^{(4)}$\\
24 & \texttt{0012111200} & $24$ & $0$ & $R_{4}^{(4)}$\\
25 & \texttt{0012210101} & $3$ & $0$ & $R_{0}^{(4)}$\\
\bottomrule\end{longtable}
\begin{equation*}\begin{aligned}R_{0}^{(4)}&=0.\end{aligned}\end{equation*}
\begin{equation*}\begin{aligned}R_{1}^{(4)}&=- \frac{2 A^{3}}{9} + \frac{\pi^{3}}{9} - 4 \zeta(3)\\ &- \frac{2 A \pi^{2}}{9}.\end{aligned}\end{equation*}
\begin{equation*}\begin{aligned}R_{2}^{(4)}&=- \frac{23 A^{3}}{36} + \frac{25 \pi^{3}}{36} + 19 \zeta(3)\\ &+\frac{A \pi^{2}}{36}.\end{aligned}\end{equation*}
\begin{equation*}\begin{aligned}R_{3}^{(4)}&=- \frac{23 A^{3}}{36} - \frac{11 \pi^{3}}{36} - 9 \zeta(3)\\ &+\frac{A \pi^{2}}{36}.\end{aligned}\end{equation*}
\begin{equation*}\begin{aligned}R_{4}^{(4)}&=- \frac{5 A^{3}}{18} + \frac{5 \pi^{3}}{36} + 3 \zeta(3)\\ &- \frac{A \pi^{2}}{9}.\end{aligned}\end{equation*}
\begin{equation*}\begin{aligned}R_{5}^{(4)}&=2 A \pi - 2 \pi^{2} - 4 \pi.\end{aligned}\end{equation*}
\begin{equation*}\begin{aligned}R_{6}^{(4)}&=- \frac{3 A Q}{2}.\end{aligned}\end{equation*}
\begin{equation*}\begin{aligned}R_{7}^{(4)}&=- \frac{4 A^{3}}{3} + \frac{5 \pi^{3}}{3} + 48 \zeta(3)\\ &+\frac{A \pi^{2}}{3}.\end{aligned}\end{equation*}
\begin{equation*}\begin{aligned}R_{8}^{(4)}&=- \frac{A^{3}}{9} + \frac{\pi^{3}}{18} + 2 \zeta(3)\\ &- \frac{A \pi^{2}}{9}.\end{aligned}\end{equation*}
\begin{equation*}\begin{aligned}R_{9}^{(4)}&=- \frac{4 A^{3}}{3} - \frac{\pi^{3}}{3} - 8 \zeta(3)\\ &+\frac{A \pi^{2}}{3}.\end{aligned}\end{equation*}
\begin{equation*}\begin{aligned}R_{10}^{(4)}&=4 \pi^{2} + 8 \pi.\end{aligned}\end{equation*}
\begin{equation*}\begin{aligned}R_{11}^{(4)}&=- \frac{A^{3}}{12} - \frac{\pi^{3}}{12} - \zeta(3)\\ &- \frac{A \pi^{2}}{12}.\end{aligned}\end{equation*}
\begin{equation*}\begin{aligned}R_{12}^{(4)}&=- \frac{5 A^{3}}{9} + \frac{5 \pi^{3}}{18} - 2 \zeta(3)\\ &- \frac{2 A \pi^{2}}{9}.\end{aligned}\end{equation*}
\begin{equation*}\begin{aligned}R_{13}^{(4)}&=2 A \pi + 2 \pi^{2} + 4 \pi.\end{aligned}\end{equation*}
\subsection{5-loop rows and residue polynomials}
\begin{longtable}{@{}r l r r r@{}}\toprule Class & Edge word & $\omega_b$ & $\omega_c$ & Residue\\\midrule\endhead
5 & \texttt{000020021201200} & $- \frac{15}{8}$ & $- \frac{5}{16}$ & $R_{0}^{(5)}$\\
6 & \texttt{000020021210110} & $- \frac{15}{2}$ & $- \frac{5}{4}$ & $R_{0}^{(5)}$\\
7 & \texttt{000020021301100} & $- \frac{5}{2}$ & $- \frac{5}{12}$ & $R_{0}^{(5)}$\\
8 & \texttt{000020022210100} & $- \frac{15}{8}$ & $- \frac{5}{16}$ & $R_{0}^{(5)}$\\
13 & \texttt{000020111111200} & $- \frac{15}{4}$ & $- \frac{5}{8}$ & $R_{0}^{(5)}$\\
14 & \texttt{000020111120110} & $-15$ & $- \frac{5}{2}$ & $R_{0}^{(5)}$\\
15 & \texttt{000020111121100} & $-15$ & $- \frac{5}{2}$ & $R_{0}^{(5)}$\\
16 & \texttt{000020111130010} & $- \frac{5}{2}$ & $- \frac{5}{12}$ & $R_{0}^{(5)}$\\
17 & \texttt{000020111220001} & $- \frac{15}{2}$ & $- \frac{5}{4}$ & $R_{0}^{(5)}$\\
18 & \texttt{000020112120100} & $- \frac{15}{2}$ & $- \frac{5}{4}$ & $R_{0}^{(5)}$\\
19 & \texttt{000020112220000} & $- \frac{15}{8}$ & $- \frac{5}{16}$ & $R_{0}^{(5)}$\\
20 & \texttt{000020120210011} & $- \frac{15}{4}$ & $- \frac{5}{8}$ & $R_{0}^{(5)}$\\
21 & \texttt{000020120211001} & $- \frac{15}{4}$ & $- \frac{5}{8}$ & $R_{0}^{(5)}$\\
22 & \texttt{000020120211010} & $- \frac{15}{4}$ & $- \frac{5}{8}$ & $R_{0}^{(5)}$\\
23 & \texttt{000021110110111} & $- \frac{15}{2}$ & $- \frac{5}{4}$ & $R_{0}^{(5)}$\\
27 & \texttt{000110012201200} & $- \frac{15}{2}$ & $- \frac{5}{4}$ & $R_{1}^{(5)}$\\
28 & \texttt{000110012210110} & $-15$ & $- \frac{5}{2}$ & $R_{2}^{(5)}$\\
29 & \texttt{000110012301100} & $-5$ & $- \frac{5}{6}$ & $R_{3}^{(5)}$\\
30 & \texttt{000110012310100} & $- \frac{5}{2}$ & $- \frac{5}{12}$ & $R_{0}^{(5)}$\\
31 & \texttt{000110013210100} & $- \frac{5}{2}$ & $- \frac{5}{12}$ & $R_{0}^{(5)}$\\
32 & \texttt{000110022201100} & $- \frac{15}{4}$ & $- \frac{5}{8}$ & $R_{4}^{(5)}$\\
37 & \texttt{000110102111200} & $-15$ & $- \frac{5}{2}$ & $R_{5}^{(5)}$\\
38 & \texttt{000110102120110} & $- \frac{15}{2}$ & $- \frac{5}{4}$ & $R_{6}^{(5)}$\\
39 & \texttt{000110102121100} & $-15$ & $- \frac{5}{2}$ & $R_{7}^{(5)}$\\
105 & \texttt{000120111310000} & $-5$ & $0$ & $R_{8}^{(5)}$\\
106 & \texttt{000120120201001} & $-15$ & $0$ & $R_{9}^{(5)}$\\
107 & \texttt{000120120202000} & $- \frac{15}{4}$ & $0$ & $R_{10}^{(5)}$\\
108 & \texttt{000120120211000} & $- \frac{15}{2}$ & $0$ & $R_{11}^{(5)}$\\
109 & \texttt{000120120301000} & $- \frac{5}{2}$ & $0$ & $R_{12}^{(5)}$\\
110 & \texttt{000120121210000} & $- \frac{15}{4}$ & $0$ & $R_{13}^{(5)}$\\
111 & \texttt{000121101101200} & $- \frac{15}{2}$ & $0$ & $R_{14}^{(5)}$\\
112 & \texttt{000121101110101} & $-30$ & $0$ & $R_{15}^{(5)}$\\
114 & \texttt{000121101110200} & $-30$ & $0$ & $R_{16}^{(5)}$\\
115 & \texttt{000121101200101} & $-15$ & $0$ & $R_{6}^{(5)}$\\
116 & \texttt{000121101210100} & $- \frac{15}{2}$ & $0$ & $R_{17}^{(5)}$\\
117 & \texttt{000121102110100} & $-15$ & $0$ & $R_{14}^{(5)}$\\
118 & \texttt{000121110110110} & $-15$ & $0$ & $R_{18}^{(5)}$\\
119 & \texttt{000121110200011} & $-15$ & $0$ & $R_{6}^{(5)}$\\
120 & \texttt{000121110200101} & $-15$ & $0$ & $R_{0}^{(5)}$\\
121 & \texttt{000121120200001} & $- \frac{15}{8}$ & $0$ & $R_{0}^{(5)}$\\
122 & \texttt{000131101110100} & $-5$ & $0$ & $R_{19}^{(5)}$\\
127 & \texttt{001111011101200} & $-30$ & $0$ & $R_{20}^{(5)}$\\
40 & \texttt{000110102130010} & $-5$ & $- \frac{5}{6}$ & $R_{21}^{(5)}$\\
41 & \texttt{000110102210101} & $-15$ & $- \frac{5}{2}$ & $R_{22}^{(5)}$\\
42 & \texttt{000110102211100} & $-15$ & $- \frac{5}{2}$ & $R_{23}^{(5)}$\\
43 & \texttt{000110102220001} & $- \frac{15}{2}$ & $- \frac{5}{4}$ & $R_{24}^{(5)}$\\
44 & \texttt{000110102220010} & $- \frac{15}{2}$ & $- \frac{5}{4}$ & $R_{9}^{(5)}$\\
45 & \texttt{000110102220100} & $- \frac{15}{2}$ & $- \frac{5}{4}$ & $R_{25}^{(5)}$\\
46 & \texttt{000110102310001} & $-5$ & $- \frac{5}{6}$ & $R_{3}^{(5)}$\\
48 & \texttt{000110103210100} & $-5$ & $- \frac{5}{6}$ & $R_{21}^{(5)}$\\
49 & \texttt{000110103220000} & $- \frac{5}{2}$ & $- \frac{5}{12}$ & $R_{12}^{(5)}$\\
52 & \texttt{000110111112100} & $-30$ & $-5$ & $R_{26}^{(5)}$\\
53 & \texttt{000110111201101} & $-30$ & $-5$ & $R_{27}^{(5)}$\\
54 & \texttt{000110111202100} & $-15$ & $- \frac{5}{2}$ & $R_{28}^{(5)}$\\
55 & \texttt{000110111211001} & $-15$ & $- \frac{5}{2}$ & $R_{29}^{(5)}$\\
56 & \texttt{000110111211010} & $-30$ & $-5$ & $R_{30}^{(5)}$\\
57 & \texttt{000110111301001} & $-10$ & $- \frac{5}{3}$ & $R_{31}^{(5)}$\\
59 & \texttt{000110112201100} & $-15$ & $- \frac{5}{2}$ & $R_{32}^{(5)}$\\
60 & \texttt{000110112210001} & $-15$ & $- \frac{5}{2}$ & $R_{33}^{(5)}$\\
61 & \texttt{000110112210010} & $-15$ & $- \frac{5}{2}$ & $R_{7}^{(5)}$\\
62 & \texttt{000110112210100} & $-15$ & $- \frac{5}{2}$ & $R_{34}^{(5)}$\\
63 & \texttt{000110112211000} & $-15$ & $- \frac{5}{2}$ & $R_{35}^{(5)}$\\
64 & \texttt{000110112220000} & $- \frac{15}{2}$ & $- \frac{5}{4}$ & $R_{11}^{(5)}$\\
66 & \texttt{000110112310000} & $-5$ & $- \frac{5}{6}$ & $R_{8}^{(5)}$\\
67 & \texttt{000110201210002} & $- \frac{15}{4}$ & $- \frac{5}{8}$ & $R_{36}^{(5)}$\\
68 & \texttt{000110201211001} & $-15$ & $- \frac{5}{2}$ & $R_{37}^{(5)}$\\
69 & \texttt{000110201220001} & $- \frac{15}{2}$ & $- \frac{5}{4}$ & $R_{24}^{(5)}$\\
71 & \texttt{000110202220000} & $- \frac{15}{8}$ & $- \frac{5}{16}$ & $R_{10}^{(5)}$\\
72 & \texttt{000110211211000} & $- \frac{15}{4}$ & $- \frac{5}{8}$ & $R_{38}^{(5)}$\\
76 & \texttt{000111101101201} & $-15$ & $- \frac{5}{2}$ & $R_{39}^{(5)}$\\
77 & \texttt{000111101102200} & $-15$ & $- \frac{5}{2}$ & $R_{16}^{(5)}$\\
79 & \texttt{000111101111110} & $-30$ & $-5$ & $R_{20}^{(5)}$\\
80 & \texttt{000111101111200} & $-30$ & $-5$ & $R_{40}^{(5)}$\\
81 & \texttt{000111101200102} & $-15$ & $- \frac{5}{2}$ & $R_{41}^{(5)}$\\
82 & \texttt{000111101201101} & $-30$ & $-5$ & $R_{42}^{(5)}$\\
83 & \texttt{000111101210101} & $-15$ & $- \frac{5}{2}$ & $R_{43}^{(5)}$\\
84 & \texttt{000111102200101} & $-15$ & $- \frac{5}{2}$ & $R_{22}^{(5)}$\\
86 & \texttt{000111102210100} & $- \frac{15}{2}$ & $- \frac{5}{4}$ & $R_{17}^{(5)}$\\
88 & \texttt{000111111200002} & $- \frac{15}{4}$ & $- \frac{5}{8}$ & $R_{0}^{(5)}$\\
89 & \texttt{000111111200011} & $-30$ & $-5$ & $R_{44}^{(5)}$\\
91 & \texttt{000111111201100} & $-15$ & $- \frac{5}{2}$ & $R_{45}^{(5)}$\\
92 & \texttt{000111111300001} & $- \frac{5}{2}$ & $- \frac{5}{12}$ & $R_{0}^{(5)}$\\
94 & \texttt{000111201210001} & $- \frac{15}{2}$ & $- \frac{5}{4}$ & $R_{46}^{(5)}$\\
95 & \texttt{000120012210100} & $- \frac{15}{4}$ & $0$ & $R_{0}^{(5)}$\\
96 & \texttt{000120021201100} & $- \frac{15}{2}$ & $0$ & $R_{47}^{(5)}$\\
97 & \texttt{000120102120100} & $-15$ & $0$ & $R_{9}^{(5)}$\\
98 & \texttt{000120102220000} & $- \frac{15}{8}$ & $0$ & $R_{10}^{(5)}$\\
100 & \texttt{000120111120010} & $-30$ & $0$ & $R_{48}^{(5)}$\\
101 & \texttt{000120111210001} & $-30$ & $0$ & $R_{34}^{(5)}$\\
102 & \texttt{000120111210010} & $-15$ & $0$ & $R_{48}^{(5)}$\\
103 & \texttt{000120111211000} & $-15$ & $0$ & $R_{49}^{(5)}$\\
104 & \texttt{000120111220000} & $-15$ & $0$ & $R_{11}^{(5)}$\\
\bottomrule\end{longtable}
\begin{equation*}\begin{aligned}R_{0}^{(5)}&=0.\end{aligned}\end{equation*}
\begin{equation*}\begin{aligned}R_{1}^{(5)}&=- \frac{19 A^{4}}{60} + 16 A \zeta(3) + \frac{19 \pi^{4}}{900}\\ &- \frac{A^{2} \pi^{2}}{30} + \frac{8 A \pi^{3}}{15}.\end{aligned}\end{equation*}
\begin{equation*}\begin{aligned}R_{2}^{(5)}&=- \frac{499 A^{4}}{720} + \frac{17 A^{2} \pi^{2}}{360} - \frac{3 \pi^{4}}{400}\\ &+\frac{191 A \pi^{3}}{180} + \frac{141 A \zeta(3)}{5}.\end{aligned}\end{equation*}
\begin{equation*}\begin{aligned}R_{3}^{(5)}&=- \frac{4 \pi^{3}}{3} - \frac{24 \pi^{2}}{5} - \frac{48 \pi}{5}\\ &+2 A^{2} \pi.\end{aligned}\end{equation*}
\begin{equation*}\begin{aligned}R_{4}^{(5)}&=- \frac{47 A^{4}}{720} - \frac{43 A^{2} \pi^{2}}{360} + \frac{83 \pi^{4}}{1200}\\ &+\frac{7 A \pi^{3}}{180} - \frac{3 A \zeta(3)}{5}.\end{aligned}\end{equation*}
\begin{equation*}\begin{aligned}R_{5}^{(5)}&=- \frac{499 A^{4}}{720} - \frac{4 \mathcal K}{5} - \frac{3 \pi^{4}}{400}\\ &+\frac{17 A^{2} \pi^{2}}{360} + \frac{191 A \pi^{3}}{180} + \frac{141 A \zeta(3)}{5}.\end{aligned}\end{equation*}
\begin{equation*}\begin{aligned}R_{6}^{(5)}&=- \frac{47 A^{4}}{180} - \frac{A^{2} \pi^{2}}{18} - \frac{11 \pi^{4}}{900}\\ &+\frac{2 A \pi^{3}}{9} + \frac{24 A \zeta(3)}{5}.\end{aligned}\end{equation*}
\begin{equation*}\begin{aligned}R_{7}^{(5)}&=- \frac{499 A^{4}}{720} + \frac{4 \mathcal K}{5} - \frac{3 \pi^{4}}{400}\\ &+\frac{17 A^{2} \pi^{2}}{360} + \frac{227 A \pi^{3}}{180} + \frac{169 A \zeta(3)}{5}.\end{aligned}\end{equation*}
\begin{equation*}\begin{aligned}R_{8}^{(5)}&=- \frac{4 \pi^{3}}{15} - \frac{12 \pi^{2}}{5} - \frac{24 \pi}{5}\\ &+\frac{3 A^{2} \pi}{5} + \frac{2 A \pi^{2}}{5} + \frac{4 A \pi}{5}.\end{aligned}\end{equation*}
\begin{equation*}\begin{aligned}R_{9}^{(5)}&=- \frac{A^{4}}{16} - A \zeta(3) + \frac{3 \pi^{4}}{400}\\ &- \frac{5 A^{2} \pi^{2}}{72} - \frac{A \pi^{3}}{12}.\end{aligned}\end{equation*}
\begin{equation*}\begin{aligned}R_{10}^{(5)}&=- \frac{A^{4}}{60} - \frac{4 A \zeta(3)}{5} - \frac{19 \pi^{4}}{900}\\ &- \frac{A^{2} \pi^{2}}{30} - \frac{A \pi^{3}}{15}.\end{aligned}\end{equation*}
\begin{equation*}\begin{aligned}R_{11}^{(5)}&=- \frac{11 A^{4}}{80} - \frac{47 A^{2} \pi^{2}}{360} - \frac{3 \pi^{4}}{400}\\ &+\frac{17 A \pi^{3}}{60} + \frac{33 A \zeta(3)}{5}.\end{aligned}\end{equation*}
\begin{equation*}\begin{aligned}R_{12}^{(5)}&=\frac{4 \pi^{3}}{3} + \frac{24 \pi^{2}}{5} + \frac{48 \pi}{5}\\ &+\frac{6 A^{2} \pi}{5} - \frac{12 A \pi^{2}}{5} - \frac{24 A \pi}{5}.\end{aligned}\end{equation*}
\begin{equation*}\begin{aligned}R_{13}^{(5)}&=- \frac{7 A^{4}}{180} - \frac{A^{2} \pi^{2}}{18} - \frac{11 \pi^{4}}{900}\\ &+\frac{2 A \pi^{3}}{45} + \frac{8 A \zeta(3)}{5}.\end{aligned}\end{equation*}
\begin{equation*}\begin{aligned}R_{14}^{(5)}&=- \frac{9 A^{2} Q}{10} - \frac{44 \Cat Q}{5} - \frac{5 Q \pi^{2}}{6}.\end{aligned}\end{equation*}
\begin{equation*}\begin{aligned}R_{15}^{(5)}&=- \frac{39 A^{2} Q}{20} - \frac{22 \Cat Q}{5} - \frac{7 Q \pi^{2}}{60}.\end{aligned}\end{equation*}
\begin{equation*}\begin{aligned}R_{16}^{(5)}&=- \frac{53 A^{4}}{180} + \frac{4 \mathcal K}{5} + \frac{11 \pi^{4}}{900}\\ &- \frac{A^{2} \pi^{2}}{6} + \frac{16 A \pi^{3}}{45} + 8 A \zeta(3).\end{aligned}\end{equation*}
\begin{equation*}\begin{aligned}R_{17}^{(5)}&=- \frac{7 A^{4}}{48} - A \zeta(3) + \frac{623 \pi^{4}}{3600}\\ &- \frac{A^{2} \pi^{2}}{8} - \frac{A \pi^{3}}{12}.\end{aligned}\end{equation*}
\begin{equation*}\begin{aligned}R_{18}^{(5)}&=- \frac{9 A^{2} Q}{20} + \frac{22 \Cat Q}{5} + \frac{7 Q \pi^{2}}{60}.\end{aligned}\end{equation*}
\begin{equation*}\begin{aligned}R_{19}^{(5)}&=\frac{24 Q \pi}{5}.\end{aligned}\end{equation*}
\begin{equation*}\begin{aligned}R_{20}^{(5)}&=- \frac{39 A^{2} Q}{20} + \frac{74 \Cat Q}{5} + \frac{89 Q \pi^{2}}{60}.\end{aligned}\end{equation*}
\begin{equation*}\begin{aligned}R_{21}^{(5)}&=\frac{4 \pi^{3}}{15} + \frac{12 \pi^{2}}{5} + \frac{24 \pi}{5}\\ &+\frac{13 A^{2} \pi}{5} - \frac{14 A \pi^{2}}{5} - \frac{28 A \pi}{5}.\end{aligned}\end{equation*}
\begin{equation*}\begin{aligned}R_{22}^{(5)}&=- \frac{499 A^{4}}{720} - 11 A \zeta(3) - \frac{3 \pi^{4}}{400}\\ &+\frac{17 A^{2} \pi^{2}}{360} - \frac{61 A \pi^{3}}{180}.\end{aligned}\end{equation*}
\begin{equation*}\begin{aligned}R_{23}^{(5)}&=- \frac{A^{4}}{3} - \frac{2 A^{2} \pi^{2}}{9} + \frac{2 A \pi^{3}}{5}\\ &+\frac{16 A \zeta(3)}{5}.\end{aligned}\end{equation*}
\begin{equation*}\begin{aligned}R_{24}^{(5)}&=- \frac{19 A^{4}}{60} - \frac{88 A \zeta(3)}{5} + \frac{19 \pi^{4}}{900}\\ &- \frac{A^{2} \pi^{2}}{30} - \frac{2 A \pi^{3}}{3}.\end{aligned}\end{equation*}
\begin{equation*}\begin{aligned}R_{25}^{(5)}&=- \frac{62 \pi^{3}}{15} - 16 \pi^{2} - 32 \pi\\ &- \frac{2 A^{2} \pi}{5} + 4 A \pi^{2} + 8 A \pi.\end{aligned}\end{equation*}
\begin{equation*}\begin{aligned}R_{26}^{(5)}&=- \frac{3 A^{2} Q}{2} + \frac{76 \Cat Q}{5} + \frac{41 Q \pi^{2}}{30}.\end{aligned}\end{equation*}
\begin{equation*}\begin{aligned}R_{27}^{(5)}&=- \frac{3 A^{2} Q}{2} - 4 \Cat Q - \frac{7 Q \pi^{2}}{30}.\end{aligned}\end{equation*}
\begin{equation*}\begin{aligned}R_{28}^{(5)}&=- \frac{287 A^{4}}{240} - \frac{8 \mathcal K}{5} + \frac{3 \pi^{4}}{400}\\ &+\frac{127 A^{2} \pi^{2}}{360} + \frac{37 A \pi^{3}}{20} + \frac{271 A \zeta(3)}{5}.\end{aligned}\end{equation*}
\begin{equation*}\begin{aligned}R_{29}^{(5)}&=- \frac{3 A^{2} Q}{2} - \frac{52 \Cat Q}{5} - \frac{23 Q \pi^{2}}{30}.\end{aligned}\end{equation*}
\begin{equation*}\begin{aligned}R_{30}^{(5)}&=- \frac{25 A^{4}}{12} - \frac{16 \mathcal K}{5} - \frac{11 \pi^{4}}{900}\\ &+\frac{5 A^{2} \pi^{2}}{6} + \frac{52 A \pi^{3}}{15} + \frac{512 A \zeta(3)}{5}.\end{aligned}\end{equation*}
\begin{equation*}\begin{aligned}R_{31}^{(5)}&=\frac{4 \pi^{3}}{15} + \frac{12 \pi^{2}}{5} + \frac{24 \pi}{5}\\ &+\frac{13 A^{2} \pi}{5} + \frac{14 A \pi^{2}}{5} + \frac{28 A \pi}{5}.\end{aligned}\end{equation*}
\begin{equation*}\begin{aligned}R_{32}^{(5)}&=- \frac{287 A^{4}}{240} + \frac{4 \mathcal K}{5} + \frac{3 \pi^{4}}{400}\\ &+\frac{127 A^{2} \pi^{2}}{360} + \frac{37 A \pi^{3}}{20} + \frac{271 A \zeta(3)}{5}.\end{aligned}\end{equation*}
\begin{equation*}\begin{aligned}R_{33}^{(5)}&=- \frac{47 A^{4}}{240} - \frac{11 A^{2} \pi^{2}}{120} + \frac{3 \pi^{4}}{400}\\ &+\frac{19 A \pi^{3}}{60} + \frac{43 A \zeta(3)}{5}.\end{aligned}\end{equation*}
\begin{equation*}\begin{aligned}R_{34}^{(5)}&=- \frac{53 A^{4}}{180} + 8 A \zeta(3) + \frac{11 \pi^{4}}{900}\\ &- \frac{A^{2} \pi^{2}}{6} + \frac{16 A \pi^{3}}{45}.\end{aligned}\end{equation*}
\begin{equation*}\begin{aligned}R_{35}^{(5)}&=- \frac{287 A^{4}}{240} - \frac{4 \mathcal K}{5} + \frac{3 \pi^{4}}{400}\\ &+\frac{127 A^{2} \pi^{2}}{360} + \frac{41 A \pi^{3}}{20} + \frac{299 A \zeta(3)}{5}.\end{aligned}\end{equation*}
\begin{equation*}\begin{aligned}R_{36}^{(5)}&=- \frac{19 A^{4}}{720} + A \zeta(3) - \frac{3 \pi^{4}}{400}\\ &- \frac{23 A^{2} \pi^{2}}{360} - \frac{A \pi^{3}}{180}.\end{aligned}\end{equation*}
\begin{equation*}\begin{aligned}R_{37}^{(5)}&=- \frac{287 A^{4}}{240} - \frac{177 A \zeta(3)}{5} + \frac{3 \pi^{4}}{400}\\ &+\frac{127 A^{2} \pi^{2}}{360} - \frac{27 A \pi^{3}}{20}.\end{aligned}\end{equation*}
\begin{equation*}\begin{aligned}R_{38}^{(5)}&=- \frac{73 A^{4}}{180} - 8 A \zeta(3) - \frac{389 \pi^{4}}{900}\\ &- \frac{7 A^{2} \pi^{2}}{18} + \frac{26 A \pi^{3}}{45}.\end{aligned}\end{equation*}
\begin{equation*}\begin{aligned}R_{39}^{(5)}&=- \frac{3 A^{2} Q}{2} + \frac{44 \Cat Q}{5} + \frac{5 Q \pi^{2}}{6}.\end{aligned}\end{equation*}
\begin{equation*}\begin{aligned}R_{40}^{(5)}&=- \frac{25 A^{4}}{12} + \frac{4 \mathcal K}{5} - \frac{11 \pi^{4}}{900}\\ &+\frac{5 A^{2} \pi^{2}}{6} + \frac{52 A \pi^{3}}{15} + \frac{512 A \zeta(3)}{5}.\end{aligned}\end{equation*}
\begin{equation*}\begin{aligned}R_{41}^{(5)}&=- \frac{13 A^{4}}{180} - \frac{A^{2} \pi^{2}}{6} + \frac{11 \pi^{4}}{900}\\ &+\frac{8 A \pi^{3}}{45} + \frac{24 A \zeta(3)}{5}.\end{aligned}\end{equation*}
\begin{equation*}\begin{aligned}R_{42}^{(5)}&=- \frac{25 A^{4}}{12} - \frac{4 \mathcal K}{5} - \frac{11 \pi^{4}}{900}\\ &+\frac{5 A^{2} \pi^{2}}{6} + \frac{28 A \pi^{3}}{15} + \frac{288 A \zeta(3)}{5}.\end{aligned}\end{equation*}
\begin{equation*}\begin{aligned}R_{43}^{(5)}&=- \frac{79 A^{4}}{80} - \frac{79 A \zeta(3)}{5} - \frac{623 \pi^{4}}{3600}\\ &- \frac{43 A^{2} \pi^{2}}{360} + \frac{A \pi^{3}}{60}.\end{aligned}\end{equation*}
\begin{equation*}\begin{aligned}R_{44}^{(5)}&=- \frac{25 A^{4}}{12} - 32 A \zeta(3) - \frac{11 \pi^{4}}{900}\\ &+\frac{5 A^{2} \pi^{2}}{6} - \frac{4 A \pi^{3}}{3}.\end{aligned}\end{equation*}
\begin{equation*}\begin{aligned}R_{45}^{(5)}&=- \frac{79 A^{4}}{80} - \frac{43 A^{2} \pi^{2}}{360} - \frac{623 \pi^{4}}{3600}\\ &+\frac{17 A \pi^{3}}{12} + \frac{117 A \zeta(3)}{5}.\end{aligned}\end{equation*}
\begin{equation*}\begin{aligned}R_{46}^{(5)}&=\frac{62 \pi^{3}}{15} + 16 \pi^{2} + 32 \pi\\ &+\frac{2 A^{2} \pi}{5} + 4 A \pi^{2} + 8 A \pi.\end{aligned}\end{equation*}
\begin{equation*}\begin{aligned}R_{47}^{(5)}&=- \frac{11 A^{4}}{80} + A \zeta(3) - \frac{3 \pi^{4}}{400}\\ &- \frac{47 A^{2} \pi^{2}}{360} + \frac{A \pi^{3}}{12}.\end{aligned}\end{equation*}
\begin{equation*}\begin{aligned}R_{48}^{(5)}&=- \frac{9 A^{2} Q}{10} + 4 \Cat Q + \frac{7 Q \pi^{2}}{30}.\end{aligned}\end{equation*}
\begin{equation*}\begin{aligned}R_{49}^{(5)}&=- \frac{9 A^{2} Q}{10} + \frac{52 \Cat Q}{5} + \frac{23 Q \pi^{2}}{30}.\end{aligned}\end{equation*}

\clearpage
\section{Derivation of the special-function representation of $C_6(1)$}\label{gff:app:C6derivation}
This appendix reproduces the derivation of the analytic calculation in the
package \path{C6_research_note/}, written by OpenAI GPT-6 astra, adapted to
the notation of this note. It starts from the reduced input
\eqref{gff:c6:reduced} of Section~\ref{gff:c6:analytic} and proves the conversions
stated there and in Section~\ref{gff:c1:c6analytic}: the ordinary-period sectors
(Section~\ref{gff:c6d:ordinary}), the wheel (Section~\ref{gff:c6d:wheel}), the
seventh-power sector (Section~\ref{gff:c6d:seventh}), the primitive
(Section~\ref{gff:c6d:primitive}) and the polylogarithmic sector
(Section~\ref{gff:c6d:polylog}). Sections~\ref{gff:c6d:numerics} and
\ref{gff:c6d:recurrence} give the numerical method and its error control.
Definitions already given in Sections~\ref{gff:c1:c6analytic} and
\ref{gff:c6:analytic} are used without restating them. The package's records of
these steps are the reduced setup
\path{C6_research_note/inputs/upstream_reduced_setup.md} (extracted from
\path{C6_research_note/inputs/source_snapshots/C6_central_binomial_wheel_expression.json}),
the derivation snapshots
\path{C6_research_note/inputs/source_snapshots/*_derivation.md} with their
expression and check files, and the stage folders
\path{C6_research_note/code/stages/} with their \path{DERIVATION.md} files,
inputs and code; the steps below are written out where they are used, so
these records are not cited again step by step, and the inputs that remain
are named in Section~\ref{gff:c6d:scope}.

\subsection{Scope, levels of statements and notation}\label{gff:c6d:scope}
\paragraph{Scope.}
The identity \eqref{gff:c6:reduced} is the supplied reduced analytic input of
the method; the graph-to-period reduction that produced it is derived neither
here nor in the package. The derivations below convert its sectors into the
formula \eqref{gff:c1:c6nine}. The model, the expansion parameter, the
conventions and the graph enumeration are those of this note
(Sections~\ref{gff:sec:guide}--\ref{gff:sec:graphs}), and Section~\ref{gff:c1:status}
gives the status of each statement.

\paragraph{Three levels of statements.}
The package keeps three kinds of statement apart, and so does this appendix.
\emph{Physical identification}: the graph-to-period identities and the
renormalization-group assembly behind \eqref{gff:c6:reduced} are inherited inputs;
the downstream checks do not establish them independently. \emph{Analytic
identity}: a finite algebraic, differential or contour derivation may prove
the equality of two representations, conditional on its stated input
identities. \emph{Numerical certification}: an outward-rounded interval plus
an analytic tail bound encloses a specified expression; agreement of
floating-point calculations alone is not a certificate. Input identities that
are supplied with the reduced expression, rather than derived from it, are
named as such where they are used. After the derivations below, the
conversion to the nine values uses three of them, the elliptic-moment
evaluations \eqref{gff:c6d:Iseries}--\eqref{gff:c6d:contourM2} on which the wheel's
periodization identities rest (Section~\ref{gff:c6d:wheelPeriod}), besides
published identities. The interval certificate uses more: it encloses an
earlier form of the reduced expression (Section~\ref{gff:c6d:interval}), which
reaches \eqref{gff:c6:reduced} only through further exact identities of the
package, among them the inherited wheel normalizations of $B_-+B_+$ and
$B_--B_+$. This appendix does not reproduce those identities; they are
inputs of the enclosure.

\paragraph{Symbols.}
The package writes $B_{\mathrm{pol}}$, $B_{\mathrm{low}}$, $B_6$, $B_8$ for the
sectors $S_{\mathrm{pol}}$, $S_{\mathrm{low}}$, $S_6$, $S_8$ of
\eqref{gff:c6:reduced}; $U_4,U_6,U_8$ and $H_-,H_0,H_+$ for the hypergeometric
values $\varpi_4,\varpi_6,\varpi_8$ and $\Phi_-,\Phi_0,\Phi_+$ of
\eqref{gff:c6:ordinarySectors}; $K_n$ for $\kappa_n$ of \eqref{gff:c6:residualSums},
whose sum $M_7$ it writes with a script letter; $h(x)$ for the Gauss function
$w(x)$ of
\eqref{gff:c6:primitivePairing}; $G$ for $\Cat$; $\ell$ for $\log2$; and
$\mathcal H_7^{(u)}$ for $\mathcal H_7$. It defines $P_b$ by $S_6=24P_b$, and
it writes $E_{\mathrm{real}}=E-\frac{476}{3}\pi^3(32\Cat+9\pi^2)$ for the
polynomial part of \eqref{gff:c1:c6nine}. To avoid clashes with symbols of this
note, the following local symbols are also renamed: the Barnes functional
(package $\mathcal B_c$) is $\mathfrak B_c$; the wheel sum that the package
writes as a script $S_7$ is $\Sigma_7$; the contour kernel of the primitive (package
$\mathcal K$) is $\widetilde{\mathcal N}$; the primitive Meijer kernels
(package $M_j$) are $M_j^{(3)}$; the Barnes kernel of the lower block (package
$K$) is $\widetilde M$; the wheel's parameter point and direction are
$\mathbf p_0$, $\mathbf v$; the coefficient matrices of the primitive are
$\mathsf M$, $\mathsf N$; and the Euler-weighted approximant (package $E_N$)
is $\widehat E_N$. Repetition is written $a^{\times k}$, as in
Section~\ref{gff:c1:c6analytic}.

\paragraph{Harmonic remainders.}
The constants $A=\pi+4\log2$, $Q=\Gamma(\frac14)^4/\pi$, $\Cat$ and the
sequences $a_n,h_n,\delta_n,y_n$ are those of Section~\ref{gff:c6:analytic}, with
$H_0^{(p)}=0$. Equivalently,
\begin{equation}\label{gff:c6d:deltaH}
 a_n=\frac{(1/2)_n}{n!},\qquad
 \delta_n=\frac{\pi^2}{3}-4H_{2n}^{(2)}+2H_n^{(2)}
 =\psi_1(n+\tfrac12)-\psi_1(n+1)>0
\end{equation}
for every $n\ge0$. For example $\delta_0=\pi^2/3$, $\delta_1=\pi^2/3-3$ and
$\delta_2=\pi^2/3-115/36$. The empty harmonic sum vanishes; the terms with
$n>0$ do not.

\paragraph{Polylogarithm and Meijer conventions.}
The multiple polylogarithm has the descending index order of
Section~\ref{gff:c6:analytic}. In the ordinary alphabet of harmonic
polylogarithms, $H_0(z)=\log z$ and $H_1(z)=-\log(1-z)$, with differential
letters $dz/z$ and $dz/(1-z)$. Thus $H_{011}(z)=\operatorname{Li}_{2,1}(z,1)$
and $H_{01101}(z)=\operatorname{Li}_{2,1,2}(z,1,1)$. Principal branches
continued from the stated convergence domains are used.

For an ordinary $G_{p,q}^{m,n}$ we use the inverse Mellin measure $ds/(2\pi\ii)$
and the argument factor $x^{-s}$. Its kernel is
\begin{equation}\label{gff:c6d:gkernel}
 \frac{\prod_{j=1}^{m}\Gamma(b_j+s)
       \prod_{j=1}^{n}\Gamma(1-a_j-s)}
      {\prod_{j=m+1}^{q}\Gamma(1-b_j-s)
       \prod_{j=n+1}^{p}\Gamma(a_j+s)}.
\end{equation}
The first $n$ upper and the first $m$ lower parameters are in numerator
groups. An identical numerical parameter in two different denominator groups
does not cancel. The repetition notation $a^{\times k}$ means a list of $k$
entries, not a power of a function \cite{gff:DLMF16,gff:C6dWolframMeijerG}.

For contour manipulations we use the functional
\begin{equation}\label{gff:c6d:Barnes}
 \mathfrak B_c[f]=\frac1{2\pi\ii}\int_{\operatorname{Re}s=c}f(s)\dd s,
\end{equation}
with the vertical contour directed upward. The kernels and contour majorants
are specified wherever this notation is used.

\paragraph{The restored Laurent coefficients.}
The correction of Section~\ref{gff:c6:analytic}, $R_{147,-3}=R_{148,-3}=\frac25$,
contributes to $r_6=2187\,C_6(1)$
\begin{align}\label{gff:c6d:repair}
\Delta r_6^{\mathrm{repair}}={}&
 \frac{1024}{25}A^5-\frac{1216}{15}A^3\pi^2
 +\frac{2672}{5}A^2\pi^3+\frac{60864}{5}A^2\zeta(3)\notag\\
&-\frac{2738}{75}A\pi^4+\frac{5963}{15}\pi^5
 -\frac{12046}{5}\pi^2\zeta(3)+\frac{711576}{5}\zeta(5).
\end{align}
Divided by $2187$ this is the contribution $+659.278983$ quoted in
Section~\ref{gff:c6:analytic}. It is already included in the polynomial $E$ of
\eqref{gff:c6:Epoly} and must not be added again. The package describes it as a
correction of the symbolic skeleton, not an adjustment of the independent
numerical estimate.

\subsection{Specification of the nine values}\label{gff:c6d:dictionary}
\subsubsection{The four ordinary Meijer values}\label{gff:c6d:ordinaryValues}
The values $\mathcal G_4,\mathcal G_6,\mathcal G_8,\mathcal G_7$ and their
parameters are defined in Section~\ref{gff:c1:c6analytic}. The package writes
$X_4=-C_d$ and $Y_4=-C_s$, so that $\rho_4=\sqrt{X_4^2+Y_4^2}$ and
$\alpha_4=\pi^{-1}\arctan(Y_4/X_4)$. Both $X_4,Y_4>0$ for the specified
constants (Section~\ref{gff:c6d:lower}); the real principal inverse tangent is
used. For the eighth-power block,
\begin{equation}\label{gff:c6d:alpha8}
 d_8=\frac1{2\pi}\arctan\frac{556}{1233-417\sqrt2},\qquad
 e_8=\frac1{2\pi}\operatorname{arcosh}\frac{1233-139\sqrt2}{\rho_8},\qquad
 \alpha_{8,\pm}=d_8\pm\ii e_8 .
\end{equation}
All displayed real square roots and real inverse functions take their
positive or principal values. The subscript $8$ denotes eighth-power kernel
content, not an eighth-order thermal coefficient. In the seventh block the
$\eta_{7,\pm}$ use real principal inverse tangents and $\beta_7$ uses the
principal complex inverse tangent; here $q_{7,2}>0$, $q_{7,1}>0$ and
$q_{7,0}<0$, as follows from $2/3<\log2<7/10$ and $3<\pi<22/7$. These are the
retained lower-, sixth-, eighth- and seventh-sector reductions.

\subsubsection{The bivariate convention: the kernel is part of the specification}\label{gff:c6d:Hconvention}
The bivariate $H$ convention is that of Andrade, Rathie and Farias
\cite{gff:C6dARF} (its bivariate Definition~B), as adopted by the package. All
three values have arguments $(1,1)$ and unit slopes. In a local variable $s$,
the factor is
\begin{equation}\label{gff:c6d:localH}
 \Theta(s)=
 \frac{\prod_{j=1}^{m}\Gamma(b_j+B_js)
       \prod_{j=1}^{n}\Gamma(1-a_j-A_js)}
      {\prod_{j=m+1}^{q}\Gamma(1-b_j-B_js)
       \prod_{j=n+1}^{p}\Gamma(a_j+A_js)}.
\end{equation}
For the shared group, the present cases have $m_0=0$. Its numerator factors
are $\Gamma(1-e_j-E_js-F_jt)$ for the first $n_0$ upper entries; its
denominator contains the remaining upper factors $\Gamma(e_j+E_js+F_jt)$ and
every lower factor $\Gamma(1-f_j-A_js-B_jt)$. The standard measure is
$ds\,dt/(2\pi\ii)^2$, with $x^{-s}y^{-t}$. Ordered pair and triple lists and
numerator counts fix each function without naming a residual integral as a
new function. Both the lists and the simplified kernels are given, to prevent
any ambiguity of convention; the kernels and contours are those displayed in
Section~\ref{gff:c1:c6analytic}.

\subsubsection{The wheel value}\label{gff:c6d:HWdef}
With $\alpha_W=\pi^{-1}\arctan(5/8)$ put
\begin{align}
 \boldsymbol a_W&=((1/2,1)^{\times4},(\alpha_W,1)),&
 \boldsymbol b_W&=((0,1)^{\times4},(\alpha_W,1)).
\end{align}
Then
\begin{equation}\label{gff:c6d:HW}
\mathcal H_W=H_{2,2;5,5;2,2}^{0,2;2,3;2,2}
 \left[\begin{matrix}
 (1/2;1,1)^{\times2}:\boldsymbol a_W;(1/2,1),(0,1)\\
 (-1/2;1,1)^{\times2}:\boldsymbol b_W;(0,1),(1/2,1)
 \end{matrix}\middle|1,1\right].
\end{equation}
Its kernel is $N_{\alpha_W}(s)P(t)/(1/2-s-t)^2$, the integrand of
Section~\ref{gff:c1:c6analytic}, where for any angle parameter $\alpha$
\begin{align}\label{gff:c6d:Nalpha0}
 N_\alpha(s)&=
 \frac{\Gamma(s)^2\Gamma(1/2-s)^3}
 {\Gamma(1-s)^2\Gamma(1/2+s)
  \Gamma(s+\alpha)\Gamma(1-s-\alpha)},\notag\\
 P(t)&=\Gamma(t)\Gamma(t+1/2)\Gamma(1/2-t)\Gamma(1-t),
\end{align}
so that the kernel $N(s)$ of Section~\ref{gff:c1:c6analytic} is $N_{\alpha_W}(s)$.
The upward lines $\operatorname{Re}s=1/4$,
$\operatorname{Re}t=1/8$ separate the poles and give absolute convergence. The
shared quotient is $\Gamma(1/2-s-t)^2/\Gamma(3/2-s-t)^2$. The normalization
$W_0=\sqrt{89}\,\mathcal H_W/\pi^2$, so that
$-393\pi^3W_0=-393\pi\sqrt{89}\,\mathcal H_W$, is derived in
Section~\ref{gff:c6d:wheelH} from the wheel identity \eqref{gff:c6d:wheelcore}, which
rests on the periodization identities \eqref{gff:c6d:wheelperiodinputs} proved in
Section~\ref{gff:c6d:wheelPeriod}.

\subsubsection{The unit-slope seventh-sector value}\label{gff:c6d:H7def}
With $\beta=\beta_7$, define the lists
\begin{align*}
\boldsymbol a_7&=((-1/4,1),(1/2,1)^{\times7},(\beta,1),(\bar\beta,1)),\\
\boldsymbol b_7&=((0,1)^{\times6},(1,1),(3/4,1),(\beta,1),(\bar\beta,1)),\\
\boldsymbol c_7&=((0,1),(1/2,1),(1/2,1),(1/4,1),(3/4,1)),\\
\boldsymbol d_7&=((0,1),(1/2,1),(-1/2,1),(1/4,1),(3/4,1)).
\end{align*}
Then
\begin{equation}\label{gff:c6d:H7}
 \mathcal H_7=H_{1,1;10,10;5,5}^{0,1;3,8;2,3}
 \left[\begin{matrix}(1/2;1,1):\boldsymbol a_7;\boldsymbol c_7\\
 (-1/2;1,1):\boldsymbol b_7;\boldsymbol d_7\end{matrix}\middle|1,1\right],
\end{equation}
with the simplified kernel
$\mathcal A_7(s)\cot(2\pi t)/[(1/2-t)(1/2-s-t)]$ and the contours
$\operatorname{Re}s=1/4$, $\operatorname{Re}t=1/16$ of
Section~\ref{gff:c1:c6analytic}. The second local gamma quotient equals
$\cot(2\pi t)/(1/2-t)$; its parameter lists are those in \eqref{gff:c6d:H7}. This
is a bivariate function, not an ordinary Meijer $G$ value. The normalization
$24\pi^3M_7=24\pi^{3/2}[\rho_7\mathcal G_7+196\pi\sqrt{185}\,\mathcal H_7]$ is
proved in Section~\ref{gff:c6d:seventhInner}.

\subsubsection{The primitive value}\label{gff:c6d:HPdef}
The unique root $\lambda_*\in(15,16)$ of $4\lambda^3-112\lambda^2+909\lambda-2250=0$
has the explicit form
\begin{equation}\label{gff:c6d:lambda}
 \lambda_*=\frac{28+\sqrt{409}\cos\!\left[
 \frac13\arccos\frac{7298}{409\sqrt{409}}\right]}3 .
\end{equation}
In the definitions of $e_*$, $\chi_*$, $B_\sigma$, $C_\sigma$ and
$\beta_\sigma$ of Section~\ref{gff:c1:c6analytic}, both real square-root radicands
are positive, and the complex inverse tangent is principal. These parameters
depend only on elementary and algebraic constants, not on a numerically fitted
primitive value. Define
\begin{align*}
 \boldsymbol a_\sigma&=((1/2,1)^{\times4},(\beta_\sigma,1),(\bar\beta_\sigma,1)),\\
 \boldsymbol b_\sigma&=((0,1)^{\times3},(\beta_\sigma,1),(\bar\beta_\sigma,1)).
\end{align*}
Then
\begin{equation}\label{gff:c6d:HP}
 \mathcal H_*=H_{0,1;6,5;6,5}^{0,0;3,3;3,3}
 \left[\begin{matrix}-:\boldsymbol a_+;\boldsymbol a_-\\
 (0;1,1):\boldsymbol b_+;\boldsymbol b_-\end{matrix}\middle|1,1\right].
\end{equation}
For each local group the first three upper and the first three lower
parameters are numerator entries. The kernel is
$\Psi_+(s)\Psi_-(t)/\Gamma(1-s-t)$ on the lines
$\operatorname{Re}s=\operatorname{Re}t=1/4$ (Section~\ref{gff:c1:c6analytic}). The
exact primitive normalizations are $P_a=80\pi^3\mathcal H_*$ and
$96P_a=7680\pi^3\mathcal H_*$. The primitive normalization is unchanged. The
global contour relation and the factorization producing this value are derived
in Section~\ref{gff:c6d:primitive} from the stated Gauss-product input.

\subsection{The ordinary-period sectors}\label{gff:c6d:ordinary}
\subsubsection{The input periods and their weights}
The three ordinary sectors of \eqref{gff:c6:reduced} have the hypergeometric input
forms \eqref{gff:c6:ordinarySectors}, with $S_6=24P_b$. They specify the analytic
quantities that are converted into $\mathcal G_4$, $\mathcal G_6$ and
$\mathcal G_8$; their graph identifications are part of the reduced input. The
formerly separated contribution $-144Q^3/\pi$ belongs to the complete $P_b$ and
is inside $S_6$ there (it comes from the $-1$ in $2\Phi_--1$); it is not an
additional term in $E$ and must not be appended to $E$. The parameter excesses are positive, so the ordinary series converge
at their unit arguments.

\subsubsection{A common gamma expansion}
The gamma recurrence and the integer and half-integer polygamma values give,
for $n\ge0$ and $\varepsilon\to0$,
\begin{equation}\label{gff:c6d:gammaLaurent}
 \frac{\Gamma(-n+\varepsilon)}{\Gamma(1/2-n+\varepsilon)}
 =\frac{a_n}{\sqrt\pi\,\varepsilon}
 \exp\left[2h_n\varepsilon+
 \frac{\delta_n-2\pi^2/3}{2}\varepsilon^2+O(\varepsilon^3)\right].
\end{equation}
For $n=0$, the first logarithmic coefficient is $2\log2$ and the second
derivative is $-\pi^2/3$. Moving to $-n$ adds $2(H_n-H_{2n})$ and
$-4H_{2n}^{(2)}+2H_n^{(2)}$ respectively. This gives \eqref{gff:c6d:gammaLaurent}
with exactly the remainders $h_n$ and $\delta_n$ of Section~\ref{gff:c6:analytic}
and \eqref{gff:c6d:deltaH}. The same expansion fixes the logarithmic
normalizations of the primitive and of the wheel residues.

\subsubsection{The lower block: two fourth-order kernels and sine addition}\label{gff:c6d:lower}
On the upward line $\operatorname{Re}s=1/4$, define
\begin{align}\label{gff:c6d:M12}
 M_1(s)&=\frac{\Gamma(s)\Gamma(1/2-s)^3}
 {\Gamma(1-s)^3\Gamma(1/2+s)},\notag\\
 M_2(s)&=\frac{\Gamma(s)^2\Gamma(1/2-s)^2}
 {\Gamma(1-s)^2\Gamma(1/2+s)^2}.
\end{align}
Gamma reflection writes these as $R_4(s)\tan^3(\pi s)$ and
$R_4(s)\tan^2(\pi s)$, where $R_4(s)=[\Gamma(s)/\Gamma(s+1/2)]^4$. Their
contour values are
\begin{equation}
 g_{13}=G^{1,3}_{4,4}\!\left(1\,\middle|\,\begin{matrix}(\frac12)^{\times4}\\
 0^{\times4}\end{matrix}\right),\qquad
 g_{22}=G^{2,2}_{4,4}\!\left(1\,\middle|\,\begin{matrix}(\frac12)^{\times4}\\
 0^{\times4}\end{matrix}\right).
\end{equation}
The residues of $M_1$ at $s=-n$ are $\pi a_n^4$; hence $g_{13}=\pi\varpi_4$.
To obtain the other normalization, the Bailey--Mishev Barnes and ${}_7F_6$
representations at $(1/2,1/2,1/2,1/2;1;1,1)$ \cite{gff:Mishev}, a published
identity, give
\begin{equation}
 \frac{\varpi_6}{2\pi}=\frac1{\pi^5}\mathfrak B_{1/4}[\widetilde M],\qquad
 \widetilde M(s)=\frac{\Gamma(s)^2\Gamma(1/2-s)^4}{\Gamma(1-s)^2}.
\end{equation}
Here and below $\mathfrak B_c$ is the upward Barnes functional
\eqref{gff:c6d:Barnes}. Set $M_0(s)=\Gamma(1/2-s)^4/\Gamma(1-s)^4$. Reflection
gives $\widetilde M=\pi^2(M_2+M_0)$. The contribution $\mathfrak B_{1/4}[M_0]$
is zero: under $s\mapsto1/2-s$ it is the inverse Mellin value of $R_4$ at $1$.
The beta integral identifies $\Gamma(s)/\Gamma(s+1/2)$ as the Laplace
transform of
\[
 b(t)=\frac1{\sqrt\pi\sqrt{1-e^{-t}}},\qquad t>0.
\]
Since $b(t)=t^{-1/2}(1+O(t))/\sqrt\pi$, its fourfold additive convolution is
$t(1+O(t))$ at $0$ and has value zero there. Inversion is justified by the
integrability of $R_4(1/4+\ii u)$. Therefore
\begin{equation}\label{gff:c6d:U46}
 \varpi_4=g_{13}/\pi,\qquad \varpi_6=2g_{22}/\pi^2.
\end{equation}
The gamma kernel $N_\alpha$ of \eqref{gff:c6d:Nalpha0} obeys
\begin{equation}\label{gff:c6d:Nalpha}
 N_\alpha=\cos(\pi\alpha)M_1+\sin(\pi\alpha)M_2.
\end{equation}
Using $X_4=-C_d$, $Y_4=-C_s$ and the angle $\alpha_4$,
\eqref{gff:c6d:U46} and \eqref{gff:c6d:Nalpha} yield
\[
 S_{\mathrm{low}}=4\pi^2(C_dg_{13}+C_sg_{22})
                  =-4\pi^2\rho_4\mathcal G_4.
\]
There is no contour boundary correction: the kernels are
$O((1+|\operatorname{Im}s|)^{-2})$ on the common line. The branch can be
checked without sampling the periods. The inequalities $\pi>3$ and
$\log2>1/2$ give $A>5$. Splitting Euler's integral for $\Gamma(1/4)$ at $t=1$,
and using $e^{-t}<1-t+t^2/2$ on $(0,1)$ and $t^{-3/4}<1$ on $(1,\infty)$,
\[
 \Gamma(1/4)
 <\int_0^1t^{-3/4}\Bigl(1-t+\frac{t^2}2\Bigr)\dd t+\int_1^\infty e^{-t}\dd t
 <4-\frac45+\frac29+\frac25<4 .
\]
Thus $Q<256/3$. These bounds imply
\[
 C_s<5120-\frac{105}{4}\,25-135-\frac{2469}{4}\,9=-\frac{2453}{2},\qquad
 C_d<8192-42\cdot25-270-1227\cdot9=-4171 .
\]
Hence $X_4,Y_4>0$, as used above.

\subsubsection{The sixth-power block}
The sixth-power sector $S_6=24P_b$ of \eqref{gff:c6:ordinarySectors}, with
$P_b=\frac{6Q^3}{\pi}(2\Phi_--1)-104\pi^5\Phi_0+\frac{2^{18}\pi^5}{Q^3}\Phi_+$,
has the absolutely convergent coefficient series
\begin{equation}\label{gff:c6d:Pbseries}
 P_b=24\pi^3\sum_{n\ge0}a_n^6
                  (312h_n^2+13\delta_n+88\pi h_n).
\end{equation}
It is derived from the three hypergeometric values by a finite Fourier
interpolation, not inferred from a numerical value. Define the proof-only
function
\begin{equation}\label{gff:c6d:f6}
 \mathsf f_6(t)=\pi^3\sum_{m\in\mathbb Z}
       \bigl[\Gamma(\tfrac34+m+t)\Gamma(\tfrac34-m-t)\bigr]^{-6}.
\end{equation}
It is entire, even and one-periodic. Gamma reflection and ratio asymptotics
\cite{gff:DLMF5,gff:DLMFStirling} give locally uniform convergence and, on a vertical
strip,
$\mathsf f_6(t)=O(e^{6\pi|\operatorname{Im}t|}(1+|\operatorname{Im}t|)^{-2})$.
Moving Fourier-coefficient contours vertically excludes the frequencies
$|k|\ge3$, including the boundary frequencies because of the factor
$(1+|\operatorname{Im}t|)^{-2}$. Hence $\mathsf f_6$ lies in the span of $1$,
$\cos2\pi t$ and $\cos4\pi t$. For every function $f$ in this
three-dimensional space,
\begin{equation}\label{gff:c6d:f6interp}
 384f(0)-104f(\tfrac14)+32f(\tfrac12)
 =312f(\tfrac14)-\frac{176}{\pi}f'(\tfrac14)+\frac{52}{\pi^2}f''(\tfrac14).
\end{equation}
Gamma reflection fixes the sample normalizations exactly:
\[
 \mathsf f_6(0)=\frac{Q^3}{64\pi^6}(2\Phi_--1),\qquad
 \mathsf f_6(\tfrac14)=\Phi_0,\qquad
 \mathsf f_6(\tfrac12)=\frac{8192}{Q^3}\Phi_+ .
\]
Consequently $P_b$ equals $\pi^5$ times the left side of \eqref{gff:c6d:f6interp}
with $f=\mathsf f_6$. At $t=\frac14$ each negative-index term of
\eqref{gff:c6d:f6} vanishes to order six. All derivatives needed here have order
at most two, so this half contributes zero. On the nonnegative half, polygamma
reflection and recurrence \cite{gff:DLMFPolygamma} yield
\[
 \mathsf f_6(\tfrac14)=\sum a_n^6,\qquad
 \mathsf f_6'(\tfrac14)=-12\sum a_n^6h_n,\qquad
 \mathsf f_6''(\tfrac14)=\sum a_n^6(144h_n^2+6\delta_n-6\pi^2).
\]
The constant $312$ term cancels the $-6\pi^2$ derivative term, and $\pi^5$
times the right side of \eqref{gff:c6d:f6interp} is \eqref{gff:c6d:Pbseries}. This is a
positive central-binomial series; it does not give separate finite evaluations
of $\Phi_-$, $\Phi_0$ or $\Phi_+$.

Introduce the proof kernels
\[
 M_j^{(6)}(s)=
 \left[\frac{\Gamma(s)}{\Gamma(s+1/2)}\right]^6
 \tan^{6-j}(\pi s),\qquad j=1,\ldots,5.
\]
Their reflection is $M_j^{(6)}(1/2-s)=M_{6-j}^{(6)}(s)$. For the ordinary values
\[
 X_{35}=G^{3,5}_{6,6}\!\left(1\,\middle|\,\begin{matrix}(\frac12)^{\times6}\\
 0^{\times6}\end{matrix}\right),\qquad
 X_{44}=G^{4,4}_{6,6}\!\left(1\,\middle|\,\begin{matrix}(\frac12)^{\times6}\\
 0^{\times6}\end{matrix}\right),
\]
the corresponding kernels satisfy
\begin{align}
 N_{35}&=\frac{\Gamma(s)^3\Gamma(1/2-s)^5}
                  {\Gamma(1-s)^3\Gamma(1/2+s)}
       =\pi^2(M_1^{(6)}+M_3^{(6)}),\notag\\
 N_{44}&=\frac{\Gamma(s)^4\Gamma(1/2-s)^4}
                  {\Gamma(1-s)^2\Gamma(1/2+s)^2}
       =\pi^2(M_2^{(6)}+M_4^{(6)}).
\end{align}
Expanding by \eqref{gff:c6d:gammaLaurent} and using the reflection of
$M_4^{(6)}$ gives
\begin{equation}
 X_{35}=3\pi^2\sum a_n^6(24h_n^2+\delta_n),\qquad
 X_{44}=24\pi^3\sum a_n^6h_n.
\end{equation}
Multiplication by $13$, $11$ and $8\pi$ now reproduces every term in
\eqref{gff:c6d:Pbseries}: $P_b=8\pi(13X_{35}+11X_{44})$. The kernel of
$\mathcal G_6$ is
\[
 N_6(s)=\frac{\Gamma(s)^4\Gamma(1/2-s)^5}
 {\Gamma(1-s)^2\Gamma(1/2+s)
   \Gamma(s+\alpha_6)\Gamma(1-s-\alpha_6)}.
\]
Sine addition gives $N_6=\cos(\pi\alpha_6)N_{35}+\sin(\pi\alpha_6)N_{44}$.
Because these two coefficients are $13/\sqrt{290}$ and $11/\sqrt{290}$,
respectively, this proves $S_6=192\pi\sqrt{290}\,\mathcal G_6$. All kernels
decay at least cubically on $\operatorname{Re}s=1/4$, and their residue series
converge absolutely at $1$, justifying these operations.

\subsubsection{The eighth-power block}
The eighth-power sector of \eqref{gff:c6:ordinarySectors} has the coefficient
series, with $y_n=n+1/4$,
\begin{align}\label{gff:c6d:B8series}
 S_8=8\pi^4\sum_{n\ge0}a_n^8\big[&
 (2466-834\sqrt2)\pi^2y_n-1112\pi(1-16y_nh_n)\notag\\
 &+139\sqrt2(256y_nh_n^2-32h_n+8y_n\delta_n)\big].
\end{align}
It is derived from $\varpi_8$, $Y_1$ and $Y_7$ as follows. Write
$\mathfrak Y=(17+12\sqrt2)Y_1-(17-12\sqrt2)Y_7$, so that
$S_8=\pi^6(4932-5004\sqrt2)\varpi_8+556\pi^2\mathfrak Y$, and note that
$Y_r=\sum_{n\ge0}(n+t)[\Gamma(n+t+\frac14)/\Gamma(n+t+\frac34)]^8$ with
$t=r/8$; in particular $Y_2=\frac{\pi^4}4\varpi_8$, so that
$S_8=\pi^2[(19728-20016\sqrt2)Y_2+556\mathfrak Y]$. Define the proof-only
function
\begin{equation}\label{gff:c6d:f8}
 \mathsf f_8(t)=\pi^4\sum_{m\in\mathbb Z}
 \frac{m+t}{[\Gamma(\frac34+m+t)\Gamma(\frac34-m-t)]^8}.
\end{equation}
It is odd and one-periodic. The same argument gives
$\mathsf f_8(t)=O(e^{8\pi|\operatorname{Im}t|}(1+|\operatorname{Im}t|)^{-2})$,
so the frequencies $|k|\ge4$ vanish, including $k=\pm4$. Thus $\mathsf f_8$
lies in the span of $\sin2\pi t$, $\sin4\pi t$ and $\sin6\pi t$, and
interpolation on these three sine modes gives the exact identity
\begin{equation}\label{gff:c6d:f8interp}
 \mathsf f_8(\tfrac18)=\frac{5\sqrt2}{8}\mathsf f_8(\tfrac14)
 -\frac1{4\pi}\mathsf f_8'(\tfrac14)
 +\frac{\sqrt2}{32\pi^2}\mathsf f_8''(\tfrac14).
\end{equation}
Gamma reflection, including the sign from the numerator of the negative half,
gives
\[
 Y_2=\pi^4\mathsf f_8(\tfrac14),\qquad \mathfrak Y=64\pi^4\mathsf f_8(\tfrac18).
\]
The factors $64$ are fixed by $\cos^8(\pi/8)=(17+12\sqrt2)/64$ and
$\sin^8(\pi/8)=(17-12\sqrt2)/64$. At the center the negative half now
vanishes to order eight. Put
\[
 u_n^{(8)}=1-16y_nh_n,\qquad
 v_n^{(8)}=256y_nh_n^2-32h_n+8y_n\delta_n-8\pi^2y_n.
\]
Then
\[
 \mathsf f_8(\tfrac14)=\sum a_n^8y_n,\qquad
 \mathsf f_8'(\tfrac14)=\sum a_n^8u_n^{(8)},\qquad
 \mathsf f_8''(\tfrac14)=\sum a_n^8v_n^{(8)}.
\]
Substituting these sums and \eqref{gff:c6d:f8interp} into the expression of $S_8$
through $Y_2$ and $\mathfrak Y$ gives \eqref{gff:c6d:B8series}. No eighth-shifted
gamma value or parameter derivative occurs on the right. This converts the
combined ordinary periods to explicit harmonic sums; it is not a finite
elementary evaluation.

To convert this coefficient polynomial, let $K_j$ be the contour value of
\[
 (1/4-s)\left[\frac{\Gamma(s)}{\Gamma(s+1/2)}\right]^8
                       \tan^{8-j}(\pi s).
\]
The factor $1/4-s=\Gamma(5/4-s)/\Gamma(1/4-s)$ supplies an ordinary contiguous
gamma pair. Equation~\eqref{gff:c6d:gammaLaurent} gives the first three residue
polynomials, each multiplied by $a_n^8$:
\begin{align}
 K_1:&\quad \pi^3y_n,\notag\\
 K_2:&\quad \pi^2(16y_nh_n-1),\notag\\
 K_3:&\quad \pi[y_n(128h_n^2+4\delta_n-\pi^2)-16h_n].
\end{align}
Set $a=1233-278\sqrt2$, $b=556$, $c=139\sqrt2$. Comparing the full polynomial
in \eqref{gff:c6d:B8series} gives
\begin{equation}
 S_8=16\pi^3(aK_1+bK_2+cK_3).
\end{equation}
The angle parameters $\alpha_{8,\pm}$ satisfy the exact identity
\begin{align}
 &a\sin^2(\pi s)+b\sin(\pi s)\cos(\pi s)+c\cos^2(\pi s)\notag\\
 &\hspace{12mm}=\rho_8\sin\pi(s+\alpha_{8,+})
                         \sin\pi(s+\alpha_{8,-}).
\end{align}
Indeed $\rho_8^2=(a-c)^2+b^2$ and $4ac-b^2=685548\sqrt2-618272>0$, fixing the
real $e_8$ and the conjugate angles \eqref{gff:c6d:alpha8}. Gamma reflection
converts the full kernel to
\begin{equation}
 \frac{\Gamma(s)^3\Gamma(1/2-s)^7\Gamma(5/4-s)}
 {\Gamma(1-s)^5\Gamma(1/2+s)\Gamma(1/4-s)
 \prod_{\sigma=\pm}\Gamma(s+\alpha_{8,\sigma})
                         \Gamma(1-s-\alpha_{8,\sigma})}.
\end{equation}
This is exactly the kernel of $\mathcal G_8$, proving
$S_8=16\pi^3\rho_8\mathcal G_8$. The common contour has cubic decay, so the
evaluation at $1$ needs no regularization.

\subsection{The wheel: one positive sequence and one bivariate value}\label{gff:c6d:wheel}
\subsubsection{Zero-index bookkeeping and direct convergence}
The legacy positive wheel and its zero-included version are
\begin{align}\label{gff:c6d:Wc}
 W_c&=\sum_{n=1}^{\infty}a_n^4
       [8+5\pi(4n+1)a_n^2]\delta_n,\notag\\
 W_0&=W_c+\frac{(8+5\pi)\pi^2}{3}.
\end{align}
The zero-included convention is used in \eqref{gff:c6:reduced} and
\eqref{gff:c1:c6nine}. The $n=0$ contribution to $-393\pi^3W_0$ is
$-1048\pi^5-655\pi^6$ and is already present there. The recurrence
\begin{equation}
 a_{n+1}=a_n\frac{2n+1}{2n+2},\qquad
 \delta_{n+1}=\delta_n-\frac4{(2n+1)^2}+\frac1{(n+1)^2}
\end{equation}
requires only one binomial sequence and one reciprocal-square accumulator.

For $x=n+1/4$, the integral
$\psi_1(z)=\int_0^\infty ue^{-zu}(1-e^{-u})^{-1}\dd u$ \cite{gff:DLMFPolygamma}
and \eqref{gff:c6d:deltaH} give
$\delta_n=\int_0^\infty ue^{-(n+1/2)u}(1+e^{-u/2})^{-1}\dd u$, that is,
\begin{equation}\label{gff:c6d:deltabound}
 \delta_n=\frac12\int_0^\infty u e^{-xu}\operatorname{sech}(u/4)\dd u,
 \qquad 0<\delta_n<\frac1{2x^2};
\end{equation}
the bound follows from $0<\operatorname{sech}(u/4)<1$ for $u>0$. Moreover
$a_n^2\le1/(3n+1)$ for every $n\ge0$: equality holds at $n=0$, and since
$a_{n+1}^2=a_n^2(2n+1)^2/(2n+2)^2$ and
$(2n+1)^2(3n+4)=(3n+1)(2n+2)^2-n$, the bound passes from $n$ to $n+1$. Hence
$a_n^2\le1/(3x)$ and $(4n+1)a_n^2\le4/3$. Together these give, for the literal
sum over
$1\le n<N$,
\begin{equation}\label{gff:c6d:wheeltail}
 0<W_c-W_{c,N}<\frac{8+20\pi/3}{54(N-3/4)^3}.
\end{equation}
The direct-tail estimate applies to the literal finite sum; the stronger
numerical construction of Section~\ref{gff:c6d:wheelcert} controls the corrected
tail.

\subsubsection{The global fourth/sixth-power identity}
Put
\begin{equation}
 \tau_n=12\zeta(3)-16H_{2n}^{(3)}+4H_n^{(3)},\quad
 D_4=\sum_{n\ge0}a_n^4\delta_n,\quad
 D_6=\sum_{n\ge0}(4n+1)a_n^6\delta_n.
\end{equation}
The key identity of the calculation is
\begin{equation}\label{gff:c6d:wheelcore}
 \sum_{n\ge0}a_n^4(8h_n\delta_n+\tau_n)=\pi^2D_6,
 \qquad W_0=8D_4+5\pi D_6.
\end{equation}
It is a global summation identity, not a termwise identification of fourth and
sixth binomial powers. The package derives it from a completed hypergeometric
symmetry \cite{gff:Mishev,gff:C6dBao} and the two periodization identities
\eqref{gff:c6d:wheelperiodinputs} below. The specialized symmetry calculation is
given here, and the two periodization identities are proved in
Section~\ref{gff:c6d:wheelPeriod}.

Write
\begin{align}
 q_n&=(4n+1)a_n^6,\notag\\
 A_n&=4H_{2n}-3H_n-2\log2+\frac1{4n+1},\notag\\
 R_n&=4H_{2n}^{(2)}-H_n^{(2)}+\frac2{(4n+1)^2}.
\end{align}
For parameters satisfying $e+f+g-a-b-c-d=1$, define the completed function
$L(a,b,c,d;e;f,g)$ by the following ${}_7F_6$ representation. Put
$\alpha=d+g-e$ and take $\operatorname{Re}(f-d)>0$:
\begin{align}\label{gff:c6d:completedL}
 L={}&\frac{\Gamma(1+\alpha)}
 {\pi\Gamma(g)\Gamma(1+g-e)\Gamma(f-d)
  \Gamma(1+a+d-e)\Gamma(1+b+d-e)\Gamma(1+c+d-e)}\notag\\
 &\quad\times{}_7F_6\!\left(\begin{matrix}
 \alpha,1+\alpha/2,g-a,g-b,g-c,d,1+d-e\\
 \alpha/2,1+a+d-e,1+b+d-e,1+c+d-e,1+g-e,g\end{matrix};1\right).
\end{align}
The input symmetry is invariance under
\begin{equation}
 \mathcal C(a,b,c,d;e;f,g)
 =(g-c,a,b,g-d;1+a+b-f;g,1+a+b-e).
\end{equation}
Let $\omega=e^{\ii\pi/4}$ and set
\begin{align}
 \mathbf p_0&=(1/2,1/2,1/2,1/2;1;1,1),\notag\\
 \mathbf v&=(\omega^2,\omega,1,1-\omega+\omega^2;
                  1+\omega^2;\omega^2-\omega^3,1+\omega^3).
\end{align}
Direct substitution gives $\mathcal C(\mathbf p_0+x\mathbf v)=\mathbf p_0+\omega x\mathbf v$.
Hence $L(\mathbf p_0+x\mathbf v)=L(\mathbf p_0+\omega x\mathbf v)$, and the
coefficient of $x^2$ vanishes. The parameter excess is one at $x=0$ and stays
positive in a sufficiently small closed neighborhood; the twice-differentiated
coefficients have a summable bound $n^{-2+\varepsilon}(1+\log^2 n)$ for some
$\varepsilon<1$.

Include the gamma prefactor in each summand $l_n(x)$ of
\eqref{gff:c6d:completedL}. At the base point, $l_n(0)=q_n/(2\pi)$. With
$\varrho=1-\omega+\omega^3=1-\sqrt2$, the sums of slopes and of squared slopes
needed for differentiation are
\begin{center}\small
\begin{tabular}{@{}lcc@{}}
\toprule
Gamma-factor group & Sum of slopes & Sum of squared slopes\\\midrule
Six half-parameter factors depending on $n$ & $4\varrho$ & $2\varrho^2$\\
Five unit-parameter denominator factors & $2\varrho$ & $0$\\
Half-parameter normalization factors & $2\varrho$ & $2\varrho^2$\\
\bottomrule
\end{tabular}
\end{center}
The very-well-poised linear factor has slope $\varrho$. The polygamma
recurrence then gives
\[
 (\log l_n)'(0)=2\varrho A_n,\qquad
 (\log l_n)''(0)=-2\varrho^2R_n.
\]
Therefore
\begin{equation}\label{gff:c6d:wheelvanish}
 l_n''(0)=\frac{q_n}{2\pi}4\varrho^2(A_n^2-R_n/2),\qquad
 \sum_{n\ge0}q_n(A_n^2-R_n/2)=0.
\end{equation}
Differentiating only the hypergeometric factor and omitting its gamma
normalization would not give this identity.

To apply \eqref{gff:c6d:wheelvanish}, define
\begin{align}
 L_n&=8H_{2n}-6H_n-4\log2+\frac2{4n+1}=2A_n,\notag\\
 M_n&=\frac{5\pi^2}{3}-16H_{2n}^{(2)}+6H_n^{(2)}
                           -\frac4{(4n+1)^2},\notag\\
 \Sigma_7&=\sum_{n\ge0}q_n(L_n^2+M_n),\qquad
 V_w=\sum_{n\ge0}q_n=\varpi_6.
\end{align}
The pointwise identity
\[
 L_n^2+M_n-(\pi^2+2\delta_n)=4(A_n^2-R_n/2)
\]
implies $\Sigma_7=\pi^2V_w+2D_6$. The two periodization identities specific to
the calculation are, with $B_s=2\pi^4V_w$ and
$\mathcal T=\sum a_n^4h_n(8h_n^2+\delta_n)$,
\begin{align}\label{gff:c6d:wheelperiodinputs}
 \Sigma_7&=\frac{3B_s}{2\pi^2}-\frac{32}{\pi^2}\mathcal T,\notag\\
 \sum_{n\ge0}a_n^4(128h_n^3+24h_n\delta_n+\tau_n)&=B_s/2.
\end{align}
They are proved in Section~\ref{gff:c6d:wheelPeriod} from a periodization of the
gamma ratio $\Gamma(z+\frac12)/\Gamma(z+1)$ and three elliptic-moment
evaluations that the package inherits; they do not follow from the symmetry
alone. Subtracting
$16\mathcal T$ from the second line gives
\begin{align*}
 \sum a_n^4(8h_n\delta_n+\tau_n)
 &=\frac{B_s}{2}-16\mathcal T
 =\frac{\pi^2\Sigma_7}{2}-\frac{B_s}{4}\\
 &=\frac{\pi^2}{2}(\pi^2V_w+2D_6)-\frac{2\pi^4V_w}{4}
 =\pi^2D_6.
\end{align*}
This proves \eqref{gff:c6d:wheelcore} from the symmetry and
\eqref{gff:c6d:wheelperiodinputs}. Every series used in this elimination is
absolutely convergent.

\subsubsection{Proof of the two periodization identities}\label{gff:c6d:wheelPeriod}
\paragraph{The quartic periodization.}
Put $\mathsf g(z)=\Gamma(z+\frac12)/\Gamma(z+1)$, so that
$\mathsf g(n)^4=\pi^2a_n^4$, and $\mathsf b(t)=\sum_{n\ge0}\mathsf g(n+t)^4$.
Half-integer recurrence and special values of the polygamma functions
\cite{gff:DLMFPolygamma} give
\begin{equation}\label{gff:c6d:loggDerivatives}
 (\log\mathsf g)'(n)=-2h_n,\qquad
 (\log\mathsf g)''(n)=\delta_n,\qquad
 (\log\mathsf g)'''(n)=-\tau_n .
\end{equation}
Gamma reflection \cite{gff:DLMF5} gives
$\mathsf g(-\frac12-z)=-\cot(\pi z)\,\mathsf g(z)$. Consequently
\[
 \mathfrak P(t)=\cos^4(\pi t)\,\mathsf b(t)+\sin^4(\pi t)\,\mathsf b(\tfrac12-t)
\]
is the cosine-multiplied bilateral sum
$\cos^4(\pi t)\sum_{n\in\mathbb Z}\mathsf g(n+t)^4$. It is entire and periodic
of period one: the cosine factor removes all fourth-order half-integer poles.
In a vertical strip, gamma-ratio asymptotics and reflection
\cite{gff:DLMF5,gff:DLMFStirling} give
$|\mathfrak P(\theta+\ii y)|=O(e^{4\pi|y|}/|y|)$; indeed the one-sided sum of
four gamma ratios is bounded by a constant times
$\sum_{n\ge0}|n+\theta+\ii y|^{-2}=O(1/|y|)$. Shifting the contour of a
Fourier coefficient vertically therefore excludes the frequencies $\pm2$ as
well as all larger ones. Reflection under $t\mapsto\frac12-t$, which leaves
$\mathfrak P$ invariant, removes the cosine mode, so $\mathfrak P$ is a
constant plus a multiple of $\sin(2\pi t)$. Since
$\sin^4(\pi t)\mathsf b(\frac12-t)$ starts at order four and
$\cos^4(\pi t)=1-2\pi^2t^2+O(t^4)$,
\[
 \mathfrak P'(0)=\mathsf b'(0),\qquad
 \mathfrak P'''(0)=\mathsf b'''(0)-12\pi^2\mathsf b'(0).
\]
The finite Fourier form gives $\mathfrak P'''(0)=-4\pi^2\mathfrak P'(0)$. Hence
\begin{equation}\label{gff:c6d:quarticCubic}
 \mathsf b'''(0)=8\pi^2\mathsf b'(0).
\end{equation}

\paragraph{Three inherited elliptic-moment evaluations.}
With $\mathsf k(t)={}_2F_1(\frac12,\frac12;1;t)=2K(t)/\pi$ in the parameter
convention, put
\[
 \mathfrak N(a,b)=\int_0^1t^{a-1}(1-t)^{b-1}\mathsf k(t)\mathsf k(1-t)\dd t,\qquad
 \mathfrak M(s)=\mathfrak N(s,1),\qquad
 \mathfrak I(s)=\mathfrak N(s,1-s),
\]
and $\mathfrak M_2(a,b)=\int_0^1t^{a-1}(1-t)^{b-1}\mathsf k(t)^2\dd t$. These
integrals and their fixed-order logarithmic derivatives are holomorphic in
their common positive-beta domain, and the endpoints have at most logarithmic
growth. The package takes three evaluations of these moments from the earlier
stages of the calculation and does not rederive them. First, the
very-well-poised Mellin series, for $0<\operatorname{Re}s<1$,
\begin{equation}\label{gff:c6d:Iseries}
 \mathfrak I(s)=\frac{\Gamma(1+s)\Gamma(s)^2\Gamma(1-s)}{\Gamma(s+\frac12)^2}\,
 {}_7F_6\!\left(\begin{matrix}s,s,s,s,1+\frac s2,\frac12,\frac12\\[1pt]
 \frac s2,1,1,1,s+\frac12,s+\frac12\end{matrix};1\right).
\end{equation}
Second, the mixed Mellin series
\begin{equation}\label{gff:c6d:Mseries}
 \mathfrak M(s)=\frac{\Gamma(s)^2}{\Gamma(s+\frac12)^2}\,
 {}_4F_3\!\left(\begin{matrix}\frac12,\frac12,s,s\\[1pt]
 1,s+\frac12,s+\frac12\end{matrix};1\right).
\end{equation}
Third, the upper-half-plane contour relation for the square moments, with
$c=2-a-b$,
\begin{equation}\label{gff:c6d:contourM2}
 0=\mathfrak M_2(a,b)-e^{\ii\pi a}\mathfrak M_2(a,c)
 -e^{-\ii\pi b}\bigl[\mathfrak M_2(c,b)-\mathfrak M_2(b,c)
 +2\ii\,\mathfrak N(c,b)\bigr],
\end{equation}
whose arcs vanish when
$\operatorname{Re}a,\operatorname{Re}b,\operatorname{Re}c>0$.

\paragraph{Consequences of the inputs.}
At $s=\frac12$ the $n$-th term of \eqref{gff:c6d:Iseries}, gamma prefactor
included, is $\frac{\pi^2}2q_n$, and its first and second logarithmic
derivatives in $s$ are $L_n$ and $M_n$. Hence
\begin{equation}\label{gff:c6d:Ihalf}
 \mathfrak I(\tfrac12)=\frac{\pi^2}2V_w=\frac{B_s}{4\pi^2},\qquad
 \Sigma_7=\frac2{\pi^2}\,\mathfrak I''(\tfrac12).
\end{equation}
Taking $(a,b,c)=(s,1,1-s)$ in \eqref{gff:c6d:contourM2} gives
$\mathfrak M_2(s,1-s)=2\mathfrak M(1-s)/\sin\pi s$. Taking
$(a,b,c)=(1,s,1-s)$, and its imaginary part, then gives
\begin{equation}\label{gff:c6d:IfromM}
 \mathfrak I(s)=\frac{\mathfrak M(s)-\mathfrak M(1-s)}{\cos\pi s}.
\end{equation}
The singularity at $s=\frac12$ is removable. All derivatives below include the
derivative of this denominator, with its sign. By \eqref{gff:c6d:Mseries},
$\mathfrak M(\frac12+t)=\pi^{-1}\sum_{n\ge0}\mathsf g(n)^2\mathsf g(n+t)^2$, and
the differentiated series converge locally normally.

\paragraph{The second identity.}
At its removable point \eqref{gff:c6d:IfromM} gives
$\mathfrak M'(\frac12)=-\frac\pi2\mathfrak I(\frac12)$, and termwise
differentiation gives $\mathfrak M'(\frac12)=\mathsf b'(0)/(2\pi)$. With
\eqref{gff:c6d:Ihalf},
\begin{equation}\label{gff:c6d:bprime}
 \mathsf b'(0)=-\frac{B_s}4,\qquad\text{equivalently}\qquad
 \sum_{n\ge0}a_n^4h_n=\frac{B_s}{32\pi^2}.
\end{equation}
By \eqref{gff:c6d:loggDerivatives},
\[
 \left.\frac{d^3}{dt^3}\mathsf g(n+t)^4\right|_{t=0}
 =-4\pi^2a_n^4(128h_n^3+24h_n\delta_n+\tau_n).
\]
Equations \eqref{gff:c6d:quarticCubic} and \eqref{gff:c6d:bprime} now give
$\sum_{n\ge0}a_n^4(128h_n^3+24h_n\delta_n+\tau_n)=-2\,\mathsf b'(0)=B_s/2$,
the second line of \eqref{gff:c6d:wheelperiodinputs}.

\paragraph{The first identity.}
Write $\langle f\rangle=\sum_{n\ge0}\mathsf g(n)^4f_n$ and
$\mathsf d_n=(\log\mathsf g)'(n)=-2h_n$. Expanding \eqref{gff:c6d:IfromM} about
$s=\frac12$ gives
$\mathfrak I''(\frac12)=-\frac2{3\pi}\bigl[\mathfrak M'''(\frac12)
+\pi^2\mathfrak M'(\frac12)\bigr]$. The third derivative of the series for
$\mathfrak M(\frac12+t)$ contains $\langle(\log\mathsf g)'''\rangle$, which
\eqref{gff:c6d:quarticCubic}, written as
$\langle64\mathsf d^3+48\mathsf d\delta+4(\log\mathsf g)'''\rangle
=8\pi^2\langle4\mathsf d\rangle$, eliminates. The result is
\[
 \mathfrak I''(\tfrac12)=\frac{16}{\pi^2}\langle\mathsf d^3\rangle
 +\frac8{\pi^2}\langle\mathsf d\delta\rangle-12\langle\mathsf d\rangle .
\]
With $\mathsf g(n)^4=\pi^2a_n^4$, $\mathsf d_n=-2h_n$, \eqref{gff:c6d:bprime} and
\eqref{gff:c6d:Ihalf}, this is the first line of \eqref{gff:c6d:wheelperiodinputs}:
\[
 \Sigma_7=\frac{3B_s}{2\pi^2}
 -\frac{32}{\pi^2}\sum_{n\ge0}a_n^4h_n(8h_n^2+\delta_n).
\]

\paragraph{What remains an input.}
The two periodization identities, and with them \eqref{gff:c6d:wheelcore} and the
wheel normalization of Section~\ref{gff:c6d:wheelH}, rest on the three inherited
evaluations \eqref{gff:c6d:Iseries}--\eqref{gff:c6d:contourM2}, besides gamma and
polygamma identities and the published symmetry used above. The quartic
periodization \eqref{gff:c6d:quarticCubic} is proved here. Both periodization
identities were also confirmed numerically (Section~\ref{gff:c6:analytic}); these
are floating-point checks, not part of the proof.

\subsubsection{From the common Barnes contour to one $H$ value}\label{gff:c6d:wheelH}
Use the kernels $M_1,M_2$ of \eqref{gff:c6d:M12}, and put
$d(s)=\psi_1(1/2-s)-\psi_1(1-s)$. At $s=-n$, $d(-n)=\delta_n$ and
$d'(-n)=\tau_n$. Equation~\eqref{gff:c6d:gammaLaurent} gives
\[
 \Res_{s=-n}(M_1d)=\pi a_n^4\delta_n,\qquad
 \Res_{s=-n}(M_2d)=a_n^4(8h_n\delta_n+\tau_n).
\]
Thus \eqref{gff:c6d:wheelcore} yields
\begin{equation}
 W_0=\frac1\pi\mathfrak B_{1/4}[(8M_1+5M_2)d]
     =\frac{\sqrt{89}}\pi\mathfrak B_{1/4}[N_{\alpha_W}d].
\end{equation}
The second equality uses \eqref{gff:c6d:Nalpha} and
$\cos(\pi\alpha_W)=8/\sqrt{89}$, $\sin(\pi\alpha_W)=5/\sqrt{89}$.

To remove the trigamma insertion, start with
\[
 d(s)=-4\int_0^1\frac{x^{-2s}\log x}{1+x}\dd x,
 \qquad \operatorname{Re}s<1/2.
\]
The Mellin representation
$(1+x)^{-1}=\mathfrak B_{1/4}[\Gamma(z)\Gamma(1-z)x^{-z}]$ and
$\int_0^1x^{a-1}(-\log x)\dd x=a^{-2}$ give
\[
 d(s)=4\mathfrak B_{1/4}\!
 \left[\frac{\Gamma(z)\Gamma(1-z)}{(1-2s-z)^2}\right].
\]
Substitute $z=2t$ and use gamma duplication,
\[
 \Gamma(2t)\Gamma(1-2t)
 =\frac{\Gamma(t)\Gamma(t+1/2)\Gamma(1/2-t)\Gamma(1-t)}{2\pi}
 =\frac{P(t)}{2\pi}.
\]
The Jacobian and the squared denominator then give
\begin{equation}\label{gff:c6d:wheelinner}
 \mathfrak B_{1/8}\!\left[\frac{P(t)}{(1/2-s-t)^2}\right]=\pi d(s).
\end{equation}
Substituting this identity proves $W_0=\sqrt{89}\,\mathcal H_W/\pi^2$, with
precisely the kernel of Section~\ref{gff:c6d:HWdef}. On the specified contours,
$|1/2-s-t|\ge1/8$, $P(t)=2\pi^2/\sin(2\pi t)$ decays exponentially, and
$N_{\alpha_W}(1/4+\ii u)=O((1+|u|)^{-2})$. These give an integrable product
majorant and justify the Mellin interchanges.

\subsection{The seventh-power sector}\label{gff:c6d:seventh}
\subsubsection{A complete central series, not just a logarithmic fragment}
The sum $M_7=\sum_{n\ge0}(-1)^na_n^7F_n$ and its coefficients are defined in
\eqref{gff:c6:residualSums}. The $n=0$ term is included. The definition includes
the complete ordinary and logarithmic seventh-sector contributions; no
separate ordinary correction is added to it.
For auditing the initial term,
\begin{align}
F_0={}&39886\log^32-10906\pi\log^22-22792\log^22
 -2188\pi^2\log2+8156\pi\log2\notag\\
 &-\frac{1679}{6}\pi^3+\frac{2330}{3}\pi^2-2156\zeta(3).
\end{align}

\subsubsection{Matching the complete coefficient with a Barnes kernel}
For the Barnes proof, put
\begin{equation}
 Y(s)=\tfrac14-s,\quad R(s)=\left[\frac{\Gamma(s)}{\Gamma(s+1/2)}\right]^7,
 \quad z(s)=\tan(\pi s),\quad
 \Psi(s)=\psi(1-2s)+\gamma_{\mathrm E}.
\end{equation}
Here $\gamma_{\mathrm E}$ is Euler's constant; it is used in the proof but is
not a parameter of the final functions. Define $p(z)=7z^2+13z+11$ and
$q(z)=q_{7,2}z^2+q_{7,1}z+q_{7,0}$. The common kernel is
\begin{equation}\label{gff:c6d:seventhBarnes}
 I(s)=Y(s)R(s)\sec(\pi s)z(s)^4[196p(z(s))\Psi(s)+q(z(s))].
\end{equation}
For $j=1,2,3$, the kernels $M_j^{(7)}(s)=Y(s)R(s)\sec(\pi s)z(s)^{7-j}$ have
the following residues at $s=-n$, with the common factor $(-1)^na_n^7$
omitted:
\begin{align}
 M_1^{(7)}:&\quad\pi^{5/2}y_n,\notag\\
 M_2^{(7)}:&\quad\pi^{3/2}(14y_nh_n-1),\notag\\
 M_3^{(7)}:&\quad\sqrt\pi\,[y_n(98h_n^2+7\delta_n/2-\pi^2/2)-14h_n].
\end{align}
They follow from \eqref{gff:c6d:gammaLaurent} and the tangent and secant
expansions. The digamma insertion satisfies
\begin{equation}
 \Psi(-n)=H_{2n},\quad
 \Psi'(-n)=2H_{2n}^{(2)}-\pi^2/3,\quad
 \Psi''(-n)=-8[\zeta(3)-H_{2n}^{(3)}].
\end{equation}
Taking the residues of the products with this Taylor expansion reproduces the
complete polynomial $F_n$:
\begin{equation}\label{gff:c6d:seventhresidue}
 \Res_{s=-n}I(s)=\sqrt\pi\,(-1)^na_n^7F_n.
\end{equation}
At negative half-integers, $R(s)$ has zeros of order seven. The tangent and
secant remove $8-j$ of these orders in $M_j^{(7)}$, leaving order $j-1\ge0$;
$\Psi$ is analytic there. There is consequently no additional left-pole
family. On the initial line $\operatorname{Re}s=1/4$ the kernels decay as
$O((1+|\operatorname{Im}s|)^{-5/2}e^{-\pi|\operatorname{Im}s|})$, with a
logarithmic factor for $\Psi$. On the closing lines
$\operatorname{Re}s=-N-1/4$, reflection gives an integrated bound
$O(N^{-5/2}\log(N+2))$. The horizontal sides vanish exponentially. Thus
closing to the left is justified and gives
\begin{equation}
 M_7=\pi^{-1/2}\mathfrak B_{1/4}[I].
\end{equation}

The parameters $\eta_{7,\pm}$ and $\beta_7,\bar\beta_7$ of
Section~\ref{gff:c1:c6analytic} factor both quadratics:
\begin{align}
 p(\tan\theta)&=\sqrt{185}\,
 \frac{\sin(\theta+\pi\beta_7)\sin(\theta+\pi\bar\beta_7)}{\cos^2\theta},\notag\\
 q(\tan\theta)&=\rho_7\,
 \frac{\sin(\theta+\pi\eta_{7,+})\sin(\theta+\pi\eta_{7,-})}{\cos^2\theta}.
\end{align}
For example the tangent parameters in the first line have sum $13/7$,
product $11/7$, and normalizing square $185/49$. Gamma reflection gives
\begin{align}
 YR\sec(\pi s)z^4p(z)&=\frac{\sqrt{185}}\pi N_p(s),\notag\\
 YR\sec(\pi s)z^4q(z)&=\frac{\rho_7}\pi N_7(s),\notag\\
 N_p(s)&=\frac{\Gamma(s)^3\Gamma(1/2-s)^7\Gamma(5/4-s)}
 {\Gamma(1-s)^4\Gamma(1/4-s)
  \prod_{b\in\{\beta_7,\bar\beta_7\}}\Gamma(s+b)\Gamma(1-s-b)}.
\end{align}
The kernel $N_7$ is the same expression with $\beta_7,\bar\beta_7$ replaced by
$\eta_{7,+},\eta_{7,-}$, and hence is the ordinary kernel of $\mathcal G_7$.
Also $\mathcal A_7(s)=(-s)N_p(s)$ by the gamma recurrence. These identities fix
both the ordinary contribution and the outer kernel of the bivariate
contribution.

\subsubsection{The inner digamma transform and the unit-slope value}\label{gff:c6d:seventhInner}
For $\operatorname{Re}s=1/4$, set
\begin{equation}
 J_\pm(s)=\mathfrak B_{1/8}\left[
 \frac{\Gamma(t)\Gamma(1-t)e^{\pm\ii\pi t}}
 {(1-t)(1-2s-t)}\right].
\end{equation}
Residues at $t=-k$ give $1/[(k+1)(k+1-2s)]$; their partial-fraction sum is
\begin{equation}
 J_-(s)=\frac{\psi(1-2s)+\gamma_{\mathrm E}}{-2s}
       =\frac{\Psi(s)}{-2s}.
\end{equation}
The phase and the alternating gamma residue cancel. Quadratic rational decay
in the undamped direction and exponential decay in the other direction justify
closing the contour. Gamma reflection gives
\[
 \Gamma(t)\Gamma(1-t)(e^{-\ii\pi t}-e^{\ii\pi t})=-2\pi\ii.
\]
Thus $J_--J_+$ is $-2\pi\ii$ times the contour value of
$[(1-t)(1-2s-t)]^{-1}$. Both poles of this rational function are on the right,
and its quadratic decay permits closing to the left. The difference is zero.
Averaging the two phases and substituting $t=2v$ now gives
\begin{equation}
 J_-(s)=\frac\pi2\mathfrak B_{1/16}\left[
 \frac{\cot(2\pi v)}{(1/2-v)(1/2-s-v)}\right].
\end{equation}
The factor is $\pi/2$: the measure contributes $2$, and the two denominators
contribute $4$. Therefore
\begin{equation}\label{gff:c6d:seventhinner}
 \mathfrak B_{1/16}\left[
 \frac{\cot(2\pi v)}{(1/2-v)(1/2-s-v)}\right]
 =-\frac{\Psi(s)}{\pi s}.
\end{equation}
Gamma reflection at quarter arguments gives the unit-slope quotient
\begin{equation}
 \cot(2\pi v)=
 \frac{\Gamma(v)\Gamma(v+1/2)\Gamma(1-v)\Gamma(1/2-v)}
 {\Gamma(v+1/4)\Gamma(v+3/4)\Gamma(1/4-v)\Gamma(3/4-v)}.
\end{equation}
There is no extra $\pi$ in this identity. Combining it with the two rational
factors yields exactly the lists in \eqref{gff:c6d:H7}. Inserting
$\mathcal A_7=(-s)N_p$ into \eqref{gff:c6d:seventhinner} gives
$\mathcal H_7=\mathfrak B_{1/4}[N_p\Psi]/\pi$. Consequently
\begin{align}\label{gff:c6d:seventhfinal}
 M_7
 &=\frac{\rho_7\mathcal G_7}{\pi^{3/2}}
       +\frac{196\sqrt{185}}{\sqrt\pi}\mathcal H_7,\notag\\
 24\pi^3M_7
 &=24\pi^{3/2}\big[\rho_7\mathcal G_7
                  +196\pi\sqrt{185}\,\mathcal H_7\big].
\end{align}
This proves the complete normalization of the seventh-power sector in
\eqref{gff:c1:c6nine}.

For $s=1/4+\ii u$ and $v=1/16+\ii w$, $|\cot(2\pi v)|\le1+\sqrt2$.
Cauchy--Schwarz bounds the integral of the absolute rational factors by
\[
 \int_{\RR}\frac{\dd w}
 {\sqrt{(7/16)^2+w^2}\sqrt{(3/16)^2+(u+w)^2}}
 \le\frac{16\pi}{\sqrt{21}}.
\]
The outer factor is $O((1+|u|)^{-3/2}e^{-\pi|u|})$. This gives a product
majorant for the bivariate integral and justifies the interchanges. The
logarithms of both final arguments $(1,1)$ are zero; conjugation of the
parameter pairs and of the imaginary contour coordinates also shows that the
specified value is real.

\subsection{The primitive: from the Gauss pairing to one bivariate value}\label{gff:c6d:primitive}
\subsubsection{The Gauss-product input}
The primitive input is \eqref{gff:c6:primitivePairing}, with
$v(x)={}_2F_1(\frac14,\frac14;1;x)$, $w(x)={}_2F_1(\frac14,\frac14;\frac12;x)$,
the products $p_0,p_1,p_2$ and the pairings $I_{ij}=I_{ji}$ defined there.
The zero-balanced Gauss factor has at most a logarithmic endpoint singularity,
so these pairings converge absolutely. The coefficient in $r_6$ is $96P_a$.

\subsubsection{The contour relation used in the final reduction}
On the open upper half-plane define
\[
 f(z)=v(z),\qquad
 g(z)=\frac{\sqrt Q}{\pi}w(1-z)-v(z),\qquad
 f_*=f(1-z),\quad g_*=g(1-z).
\]
The star means reflection, not complex conjugation. Principal Gauss branches
and separate principal square roots $\sqrt z,\sqrt{1-z}$ are used. The Gauss
cut-jump and connection formulas give the exterior boundary data
\begin{center}
\begin{tabular}{@{}ccccc@{}}
\toprule
Region & $f$ & $g$ & $f_*$ & $g_*$\\\midrule
$x<0$, $z=x+\ii0$ & $u$ & $r-\ii u$ & $a-\ii b$ & $a+\ii b$\\
$x>1$, $z=x+\ii0$ & $a+\ii b$ & $a-\ii b$ & $u$ & $r+\ii u$\\
\bottomrule
\end{tabular}
\end{center}
Here $u,r,a,b$ are real local values. For $x<0$, the argument $1-z$ approaches
the zero-balanced cut from below; this fixes the minus sign in $g=r-\ii u$. For
$x>1$, the connection formula
\[
 v(z)=\frac{\sqrt Q}{2\pi}w(1-z)
 -\frac2{\sqrt Q}\sqrt{1-z}\;
 {}_2F_1\!\left(\begin{matrix}3/4,3/4\\3/2\end{matrix};1-z\right)
\]
shows that $f$ and $g$ are conjugate on the upper bank. Reflection supplies the
remaining entries.

Use the holomorphic numerator
\begin{align}
 \mathcal N={}&g^2g_*^2+g^2f_*^2+f^2g_*^2-4fgf_*g_*-3f^2f_*^2\notag\\
 &+4\ii(fg f_*^2-f^2f_*g_*),\qquad
 \widetilde{\mathcal N}(z)=\frac{\mathcal N(z)}{\sqrt z\sqrt{1-z}}.
\end{align}
Substitution of either boundary row gives the polynomial identity
\[
 \mathcal N=2r(a-b)[(a+b)r-2(a-b)u]\in\RR.
\]
On $x<0$ and $x>1$, the reciprocal square-root factor is respectively
$-\ii/\sqrt{(-x)(1-x)}$ and $\ii/\sqrt{x(x-1)}$. Therefore
$\operatorname{Re}\widetilde{\mathcal N}(x+\ii0)=0$ outside $[0,1]$. At
infinity the Gauss connection expansion gives
$\widetilde{\mathcal N}(z)=O(|z|^{-2}(1+\log|z|)^4)$. Near $0$ and $1$ the
bounds are of inverse-square-root type with at most quadratic logarithms. The
large semicircle and the two endpoint indentations all vanish. Taking the real
part of Cauchy's theorem leaves
$\int_0^1\operatorname{Re}\widetilde{\mathcal N}(x)\dd x=0$.

On the interior interval, $p_0=f^2$, $p_1/\pi=f(f+g)$ and $p_2/\pi^2=(f+g)^2$.
Expanding the real numerator in these products, and using $I_{ij}=I_{ji}$,
gives
\begin{align}\label{gff:c6d:pairings}
 I_{22}-4I_{12}+4I_{02}-4I_{00}&=0,\notag\\
 P_a&=\pi^4[60(I_{12}-I_{02})+24I_{11}+56I_{00}].
\end{align}
Thus the logarithm-squared times logarithm-squared pairing $I_{22}$ of
\eqref{gff:c6:primitivePairing} is removed. The contour proof, not an equality of
individual power-series coefficients, establishes the relation.

\subsubsection{Identification of the three ordinary Meijer solutions}
All three $p_i$ solve the symmetric-square equation
\begin{equation}\label{gff:c6d:primitiveODE}
 [\theta^3-x(\theta+\tfrac12)^3]p=0,\qquad\theta=x\frac{d}{dx}.
\end{equation}
Define
\begin{equation}
 g_j(x)=G^{j,3-j}_{3,3}\!\left(x\,\middle|\,\begin{matrix}(\frac12)^{\times3}\\
 0^{\times3}\end{matrix}\right),\quad j=1,2,3.
\end{equation}
Their inverse-Mellin kernels are
\begin{equation}\label{gff:c6d:primitiveM}
 M_j^{(3)}(s)=\frac{\Gamma(s)^j\Gamma(1/2-s)^{3-j}}
 {\Gamma(1-s)^{3-j}\Gamma(1/2+s)^j}
 =\left[\frac{\Gamma(s)}{\Gamma(s+1/2)}\right]^3
                      \tan^{3-j}(\pi s).
\end{equation}
The standard Meijer differential equation is also \eqref{gff:c6d:primitiveODE}.
To determine the complete normalization, expand at $s=0$:
\[
 \left[\frac{\Gamma(s)}{\Gamma(s+1/2)}\right]^3
 =\frac{1+6s\log2+(18\log^22-\pi^2/2)s^2+O(s^3)}{\pi^{3/2}s^3}.
\]
Writing $L=-\log x$, the residues give
\begin{align}
 g_1(x)&=\sqrt\pi+O(x),\notag\\
 g_2(x)&=(L+6\log2)/\sqrt\pi+O(x|\log x|),\notag\\
 g_3(x)&=[(L+6\log2)^2-\pi^2]/(2\pi^{3/2})+O(x\log^2x).
\end{align}
On the Gauss side, $v(x)=1+O(x)$ and the zero-balanced connection formula is
$\sqrt Q\,w(1-x)=L+\pi+6\log2+O(x|\log x|)$. Matching all constant and
logarithmic coefficients gives
\begin{equation}\label{gff:c6d:gj}
 g_1=\sqrt\pi p_0,\qquad
 g_2=\frac{p_1-\pi p_0}{\sqrt\pi},\qquad
 g_3=\frac{p_2-2\pi p_1}{2\pi^{3/2}}.
\end{equation}
The Frobenius recurrence of \eqref{gff:c6d:primitiveODE} then fixes every higher
coefficient. Matching the differential equation without these boundary
polynomials would not suffice.

For the real numerical evaluator, introduce the elliptic coordinate
\begin{equation}\label{gff:c6d:ellipticCoordinate}
 m(x)=\frac{x}{2(1+\sqrt{1-x})},\qquad
 K_x=K(m(x)),\qquad
 K_x^c=K(1-m(x)),
\end{equation}
with $K(m)$ in the parameter convention of Section~\ref{gff:c1:notation}. The
argument is the elliptic parameter, not the modulus. The substitution
$x=4m(1-m)$ transforms the elliptic equation $m(1-m)Y''+(1-2m)Y'-Y/4=0$ into
$x(1-x)y''+(1-3x/2)y'-y/16=0$. Its regular normalization gives
$v(x)=2K_x/\pi$. At zero the complete logarithmic normalization is
$2(K_x+K_x^c)=-\log x+\pi+6\log2+O(x|\log x|)$, which matches the
zero-balanced connection formula above; the regular and logarithmic Frobenius
data then give $\sqrt Q\,w(1-x)=2(K_x+K_x^c)$. Therefore
\begin{equation}\label{gff:c6d:gelliptic}
 g_1=\frac4{\pi^{3/2}}K_x^2,\qquad
 g_2=\frac4{\pi^{3/2}}K_xK_x^c,\qquad
 g_3=\frac2{\pi^{3/2}}[(K_x^c)^2-K_x^2].
\end{equation}
These are finite kernel identities used to check the primitive numerically.
The package also contains a Wolfram implementation of them but claims no run
of it; its numerical values come from the equivalent Python implementations.

\subsubsection{Beta integration and the four intermediate pairings}
Write
\[
 \langle f,g\rangle=\int_0^1
           \frac{f(x)g(1-x)}{\sqrt{x(1-x)}}\dd x,
 \qquad H_{jk}=\langle g_j,g_k\rangle.
\]
Inserting the two kernels \eqref{gff:c6d:primitiveM} and applying Euler's beta
integral gives the bivariate kernel
\begin{equation}\label{gff:c6d:intermediateH}
 \frac{\Gamma(s)^j\Gamma(1/2-s)^{4-j}
       \Gamma(t)^k\Gamma(1/2-t)^{4-k}}
 {\Gamma(1-s)^{3-j}\Gamma(1/2+s)^j
  \Gamma(1-t)^{3-k}\Gamma(1/2+t)^k\Gamma(1-s-t)}.
\end{equation}
The added factor is $\Gamma(1/2-s)\Gamma(1/2-t)/\Gamma(1-s-t)$. On
$\operatorname{Re}s=\operatorname{Re}t=1/4$, each $M_j^{(3)}$ is
$O((1+|\operatorname{Im}s|)^{-3/2})$. Before integrating in $x$, the absolute
three-variable integral is bounded by $B(1/4,1/4)$ times the product of the
two integrable local majorants. This justifies the interchanges and the
absolute convergence of \eqref{gff:c6d:intermediateH}.

Inverting \eqref{gff:c6d:gj} gives
\[
 p_0=g_1/\sqrt\pi,\quad p_1=\sqrt\pi(g_1+g_2),\quad
 p_2=2\pi^{3/2}(g_1+g_2+g_3).
\]
Substituting these into \eqref{gff:c6d:pairings} gives
\begin{align*}
 P_a/\pi^3={}&120(H_{11}+2H_{12}+H_{13}+H_{22}+H_{23})\\
 &-120(H_{11}+H_{12}+H_{13})
   +24(H_{11}+2H_{12}+H_{22})+56H_{11}.
\end{align*}
Thus
\begin{equation}\label{gff:c6d:fourH}
 P_a=8\pi^3(10H_{11}+21H_{12}+18H_{22}+15H_{23}),\qquad
 H_{33}+2H_{13}-H_{22}=0.
\end{equation}
The second relation is the first line of \eqref{gff:c6d:pairings} in the $g_j$
basis. It is the relation used in the final compression.

\subsubsection{From the four pairings to the final single value}
The coefficient matrix and the zero-relation matrix are
\begin{equation}
 \mathsf M=\begin{pmatrix}10&21/2&0\\21/2&18&15/2\\0&15/2&0\end{pmatrix},\qquad
 \mathsf N=\begin{pmatrix}0&0&1\\0&-1&0\\1&0&1\end{pmatrix}.
\end{equation}
Adding $\lambda\mathsf N$ leaves the definite pairing unchanged. The cubic
$4\lambda^3-112\lambda^2+909\lambda-2250$ is $4\det(\mathsf M+\lambda\mathsf N)$.
For its root $\lambda_*$ and the parameters of $\mathcal H_*$,
$v_\sigma=(1,B_\sigma,C_\sigma)^T$ satisfies
\begin{equation}
 \mathsf M+\lambda_*\mathsf N=5(v_+v_-^T+v_-v_+^T).
\end{equation}
With $f_\sigma=g_1+B_\sigma g_2+C_\sigma g_3$, symmetry of the pairing gives
\[
 10H_{11}+21H_{12}+18H_{22}+15H_{23}=10\langle f_+,f_-\rangle.
\]
The kernel of $f_\sigma$ contains $P_\sigma(z)=z^2+B_\sigma z+C_\sigma$,
$z=\tan(\pi s)$. Put $D_\sigma=\sqrt{(1-C_\sigma)^2+B_\sigma^2}$. The angle
parameters $\beta_\sigma$ give
\[
 P_\sigma(\tan\theta)=D_\sigma
 \frac{\sin(\theta+\pi\beta_\sigma)
       \sin(\theta+\pi\bar\beta_\sigma)}{\cos^2\theta}.
\]
Gamma reflection consequently converts $f_\sigma/D_\sigma$ to one local Meijer
kernel. The normalization is exact:
\[
 P_+(\ii)P_-(\ii)=-\frac45-\frac35\ii,\qquad
 |P_\sigma(\ii)|=D_\sigma,\qquad D_+D_-=1.
\]
After beta integration the local kernels are precisely $\Psi_+,\Psi_-$ of
Section~\ref{gff:c1:c6analytic}, with the shared factor $1/\Gamma(1-s-t)$. It
follows that $\langle f_+,f_-\rangle=\mathcal H_*$, and hence
\begin{equation}\label{gff:c6d:PaH}
 P_a=80\pi^3\mathcal H_*,\qquad 96P_a=7680\pi^3\mathcal H_*.
\end{equation}
No term is extracted from an endpoint or transferred into the known polynomial
in this conversion.

\subsection{Conversion to the real log-sine and Clausen values}\label{gff:c6d:polylog}
The starting block is $S_{\mathrm{pol}}$ of Section~\ref{gff:c6:analytic}, with
the descending-index convention. For $l(\theta)=\log(2\sin(\theta/2))$ and
$0<\theta<2\pi$, the principal logarithm is
\begin{equation}
 \log(1-e^{\ii\theta})=l(\theta)+\frac\ii2(\theta-\pi).
\end{equation}
The derivative
\[
 \frac{d}{dz}\operatorname{Li}_{2,1}(z,1)=\frac{\log^2(1-z)}{2z}
\]
and $z=e^{\ii\theta}$ imply
\begin{equation}
 \frac{d}{d\theta}\operatorname{Im}\operatorname{Li}_{2,1}(e^{\ii\theta},1)
 =\frac12l(\theta)^2-\frac{(\theta-\pi)^2}{8}.
\end{equation}
The initial imaginary part at $\theta=0$ is zero. The logarithmic singularity
is integrable. Integrating to $\pi/2$ gives
\[
 D_3=-\frac12\operatorname{Ls}_3(\pi/2)
       -\int_0^{\pi/2}\frac{(\theta-\pi)^2}{8}\dd\theta
     =-\frac12L_3-\frac{7\pi^3}{192}.
\]
This is a differential identity. For weight five, the real coefficient
convention $\operatorname{Cl}_{2,1,2}(\sigma)=\sum_{n>m>k\ge1}\cos(n\sigma)/(n^2mk^2)$
of Section~\ref{gff:c1:c6analytic} gives directly $D_5=C_{212}$: at odd total
weight the multiple Clausen convention used here is the cosine component of
the descending multiple polylogarithm. This is a real-function convention, not
a new elementary evaluation of $D_5$. Absolute convergence follows from
\[
 \sum_{n>m>k\ge1}\frac1{n^2mk^2}
 \le\zeta(2)\sum_{n\ge1}\frac{H_{n-1}}{n^2}<\infty,
\]
so the cosine component may be taken term by term. The definitions of $L_3$
and $C_{212}$ fix established special functions \cite{gff:C6dBorweinStraub}; they
are not additional residual operations in the fixed-value formula
\eqref{gff:c1:c6nine}.

Substitute these two identities into the entire block, not just into its
individual special values:
\begin{align}
 S_{\mathrm{pol}}
 &=4352\big[(32\Cat+9\pi^2)(-L_3/2-7\pi^3/192)-64C_{212}\big]\notag\\
 &=-2176\big[(32\Cat+9\pi^2)L_3+128C_{212}\big]
       -\frac{476}{3}\pi^3(32\Cat+9\pi^2).
\end{align}
Adding the last polynomial to $E$ gives the polynomial
$E_{\mathrm{real}}=E-\frac{476}{3}\pi^3(32\Cat+9\pi^2)$ of \eqref{gff:c1:c6nine}.
The fully expanded $E$ has eighteen monomials in $A,Q,\Cat,\pi,\zeta(3),\zeta(5)$,
and only two of them change. Their new coefficients are
\[
 -\frac{544}{3}-\frac{476\cdot32}{3}=-\frac{15776}{3}\quad(\Cat\pi^3),\qquad
 \frac{37631}{128}-\frac{476\cdot9}{3}=-\frac{145153}{128}\quad(\pi^5).
\]
There is no new physical contribution, and the $Q$-dependent block is
unchanged. Together with the ordinary, wheel, seventh and primitive
conversions, this completes the derivation of \eqref{gff:c1:c6nine} from the
stated reduced input.

The formula \eqref{gff:c1:c6nine} has no explicit $\operatorname{Re}$ or
$\operatorname{Im}$ operation. Some Meijer $G$ and bivariate $H$ parameter
lists still contain complex-conjugate pairs. Removing the projection notation
does not remove those parameters, and the multiple Clausen name does not
constitute a reduction of arithmetic independence.

\subsection{Numerical evaluation and error control}\label{gff:c6d:numerics}
\subsubsection{Evaluation of the final expression}
A mathematically short special-function formula need not be the fastest way to
evaluate the constant. The preferred numerical routes are a rational
central-wheel tail correction, Euler-transformed seventh coefficients, a
geometric primitive recurrence, and algorithms for harmonic polylogarithms and
log-sine values for the real polylogarithmic sector. These compute the same
specified values using the proved conversion identities, not by subtracting
unknown parts from the target $C_6$.

\subsubsection{Wheel certificate: a finite correction, not an infinite asymptotic sum}\label{gff:c6d:wheelcert}
For $t\ge N$, put $x=t+1/4$, $y=x^{-2}$ and
\begin{equation}
 \Delta_r q(t)=q(t)-\left(\frac{t+1/2}{t+1}\right)^r q(t+1),\qquad
 d_j=\frac{E_{2j}(2j+1)}{2\,16^j}.
\end{equation}
The rational bases are
\begin{align}
 \phi_{-1}^{(4)}(t)&=\frac{(t-1/4)^3}{(t-1/2)^4},&
 \phi_j^{(4)}(t)&=\frac{t-1/4}{t^j(t-1/2)^{j+4}},\notag\\
 \phi_j^{(6)}(t)&=\frac{(t-1/4)^6}{t^j(t-1/2)^{j+6}}.
\end{align}
To make the matching and the bound explicit, define the Euler numbers by
$\operatorname{sech}v=\sum_{j\ge0}E_{2j}v^{2j}/(2j)!$ near zero. Write
\[
 q_M^{(4)}=\sum_{j=-1}^M\alpha_j\phi_j^{(4)},\qquad
 q_L^{(6)}=\sum_{j=0}^L\beta_j\phi_j^{(6)},\qquad L=M+1.
\]
Choose their rational coefficients by matching
\[
 \Delta_4q_M^{(4)}=\sum_{k=0}^{M+1}d_ky^{k+1}+O(y^{M+3}),\qquad
 \Delta_6q_L^{(6)}=4\sqrt y\sum_{k=0}^Ld_ky^k
                                    +O(\sqrt y\,y^{L+1}).
\]
The triangular leading multipliers are $3$ for the fourth-power index $-1$,
$2j+5$ for its indices $j\ge0$, and $2j+3$ in the sixth-power construction.
The initial coefficients are $\alpha_{-1}=1/6$, $\alpha_0=-77/480$ and
$\beta_0=2/3$, $\beta_1=-91/120$. The rational differences have the
denominator polynomials
\begin{align*}
 \mathcal D_4(y)&=(1-y/16)^M(1-9y/16)^{M+4},&
 \Delta_4q_M^{(4)}&=\mathcal P_4(y)/\mathcal D_4(y),\\
 \mathcal D_6(y)&=(1-y/16)^L(1-9y/16)^{L+6},&
 \Delta_6q_L^{(6)}&=\sqrt y\,\mathcal P_6(y)/\mathcal D_6(y).
\end{align*}
Let $\mathcal U_4(y)=\sum_{k=0}^{M+1}d_ky^{k+1}$ and
$\mathcal U_6(y)=4\sum_{k=0}^Ld_ky^k$. Exact polynomial division gives
\[
 \mathcal E_4(y)=\frac{\mathcal P_4-\mathcal D_4\mathcal U_4}{y^{M+3}},
 \qquad
 \mathcal E_6(y)=\frac{\mathcal P_6-\mathcal D_6\mathcal U_6}{y^{L+1}}.
\]
For $Y=(N+1/4)^{-2}$, the constants used below are
\begin{align*}
 K_4&=|d_{M+2}|+
       \frac{\sum_j|[y^j]\mathcal E_4(y)|Y^j}{\mathcal D_4(Y)},\\
 K_6&=4|d_{L+1}|+
       \frac{\sum_j|[y^j]\mathcal E_6(y)|Y^j}{\mathcal D_6(Y)}.
\end{align*}
Both denominators are positive for $0\le y\le Y$ when $N\ge2$. They imply, for
$t\ge N$ and $x=t+1/4$,
\[
 |\delta(t)-\Delta_4q_M^{(4)}(t)|\le K_4x^{-2M-6},\qquad
 |4x\delta(t)-\Delta_6q_L^{(6)}(t)|\le K_6x^{-2L-3}.
\]
Here $\delta(t)=\psi_1(t+1/2)-\psi_1(t+1)$ extends the integer remainder, and
\eqref{gff:c6d:deltabound} holds with $x=t+1/4$. The global Taylor estimate used is
\begin{equation}\label{gff:c6d:sechTaylor}
 \left|\operatorname{sech}v-
       \sum_{k=0}^{K}\frac{E_{2k}v^{2k}}{(2k)!}\right|
 \le\frac{|E_{2K+2}|v^{2K+2}}{(2K+2)!}\quad(v\in\RR).
\end{equation}
It follows from the positive product
$\cosh v=\prod_{k\ge1}\bigl(1+4v^2/((2k-1)^2\pi^2)\bigr)$ \cite{gff:C6dDLMFcosh}.
Since $(1+\mu\varsigma)^{-1}$ is the Laplace transform at $\varsigma$ of an
exponential variable with mean $\mu$,
$\operatorname{sech}\sqrt\varsigma=\mathbb E\,e^{-\varsigma\mathsf X}$, where
$\mathsf X$ is the sum of independent exponential variables with means
$4/((2k-1)^2\pi^2)$; these means are summable, with sum $\frac12$. Comparing
Taylor coefficients at $\varsigma=0$ gives
$\mathbb E\,\mathsf X^k=k!\,|E_{2k}|/(2k)!$. The finite remainder
$|e^{-u}-\sum_{k=0}^K(-u)^k/k!|\le u^{K+1}/(K+1)!$ for $u\ge0$, applied at
$u=v^2\mathsf X$ and averaged, gives \eqref{gff:c6d:sechTaylor}. Integrated against
\eqref{gff:c6d:deltabound}, it gives
$\delta(t)=\sum_{k=0}^{K}d_kx^{-2k-2}+\mathsf r_K$ with
$|\mathsf r_K|\le|d_{K+1}|x^{-2K-4}$, which supplies the next-term bounds
above. No infinite asymptotic expansion is summed at fixed $t$.

Rational triangular matching determines $q_M^{(4)}$ and $q_{M+1}^{(6)}$ so that
their finite differences match the truncated expansions of $\delta(t)$ and
$(4t+1)\delta(t)$. The constants $K_4,K_6$ above give the finite approximant
\begin{align}
 W_{N,M}={}&\sum_{n=1}^{N-1}a_n^4[8+5\pi(4n+1)a_n^2]\delta_n\notag\\
 &+8a_N^4q_M^{(4)}(N)+5\pi a_N^6q_{M+1}^{(6)}(N),\notag\\
 |W_c-W_{N,M}|&\le
 \frac{(8/9)K_4+(110/189)K_6}
 {(2M+7)(N-3/4)^{2M+7}}.\label{gff:c6d:wheelcertbound}
\end{align}
The added terms evaluate the two rational tails completely by telescoping,
since $a_{n+1}^r=a_n^r[(n+\frac12)/(n+1)]^r$. The upper boundaries vanish:
$q_M^{(4)}$ is $O(t^{-1})$ and $q_{M+1}^{(6)}$ is $O(1)$, whereas the
respective gamma weights $a_n^4$ and $a_n^6$ are $O(n^{-2})$ and $O(n^{-3})$.
Using $a_n^4\le1/(9x^2)$, $a_n^6\le1/(27x^3)$, $\pi<22/7$ and integral
comparison then gives \eqref{gff:c6d:wheelcertbound}.
Each approximant is a rational cubic in $\pi$, but its coefficients depend on
$N,M$ and the error is not identically zero. It does not prove that the exact
wheel lies in $\mathbb Q+\mathbb Q\pi+\mathbb Q\pi^2+\mathbb Q\pi^3$.

The remainder is controlled by a finite global Taylor inequality for
$\operatorname{sech}$, obtained from the positive product for $\cosh$, not by
assuming convergence of an infinite asymptotic expansion at fixed $t$. The
certificate independently bounds $\pi$ by Machin's formula and uses exact
rational arithmetic plus outward dyadic rounding.

\subsubsection{Euler-weighted alternating sums}
For alternating coefficients $b_n$, the evaluator uses the finite weights
\begin{equation}
 w_{N,n}=1-2^{-N}\sum_{k=0}^{n}\binom Nk,\qquad
 \widehat E_N=\sum_{n=0}^{N-1}(-1)^n w_{N,n}b_n.
\end{equation}
For the seventh-sector sum, $b_n=a_n^7F_n$, write $\widehat M_{7,N}$ for
$\widehat E_N$. The lemma below gives the tail bound
\begin{equation}\label{gff:c6d:M7tail}
 |M_7-\widehat M_{7,N}|\le2^{100-N},\qquad N\ge1.
\end{equation}
Finite arithmetic must still be enclosed; the dyadic interval code does so.
This bound applies to the Euler-weighted approximant, not to a raw partial
sum. The elementary remainder kernel $2^{-N}(1-x)^N/(1+x)$ and beta-moment
bounds underpin this family of evaluations; Section~\ref{gff:c6d:interval} applies
them to further alternating series.

\paragraph{Lemma (Euler tail of the seventh-sector sum).}
For $N\ge1$, $|M_7-\widehat M_{7,N}|\le2^{99-N}$; in particular
\eqref{gff:c6d:M7tail} holds with one bit of slack.

\emph{Proof.} Restrict $x$ and $s$ to the closed bidisk
$|x|,|s|\le\frac1{64}$, put $t=\frac14+s$, $\mathsf a=\frac12+s+x$,
$\mathsf d=\frac12-2x$, and define
\[
 \mathfrak Z(x,s)=\left[\frac{\Gamma(1-x)\cos\pi(s+x)}
 {\Gamma(\frac12+x)\cos\pi x}\right]^7
 \sum_{n\ge0}(-1)^n(n+t)
 \left[\frac{\Gamma(n+\mathsf a)}{\Gamma(n+\mathsf a+\mathsf d)}\right]^7 .
\]
By gamma reflection this is the nonnegative-index half of the finite-Fourier
family with which the package derives \eqref{gff:c6:residualSums} from its earlier
five-shift form. Its terms are $O(n^{-5/2+14\operatorname{Re}x})$, so the
series converges absolutely and uniformly on the bidisk, where $\mathfrak Z$ is
analytic. At $x=s=0$ the $n$-th term is $(-1)^na_n^7y_n$. Its local
logarithmic derivatives there follow from polygamma reflection, recurrence
and half-integer values \cite{gff:DLMFPolygamma}:
\[
 \partial_s:\ y_n^{-1}-14h_n,\qquad
 \partial_s^2:\ -y_n^{-2}-7\pi^2+7\delta_n,\qquad
 \partial_x:\ 14H_{2n},\qquad
 \partial_x\partial_s:\ R_n^{(7)},\qquad
 \partial_x\partial_s^2:\ S_n^{(7)},
\]
with $R_n^{(7)}=-7(\frac{\pi^2}3+4H_{2n}^{(2)})$ and
$S_n^{(7)}=-112[\zeta(3)-H_{2n}^{(3)}]$. With $u_n,v_n$ of
\eqref{gff:c6:residualSums}, the six derivatives at $(0,0)$ needed below are
\begin{align*}
 \mathfrak Z&=\sum(-1)^na_n^7y_n,&
 \mathfrak Z_x&=\sum(-1)^na_n^7\,14H_{2n}y_n,\\
 \mathfrak Z_s&=\sum(-1)^na_n^7u_n,&
 \mathfrak Z_{xs}&=\sum(-1)^na_n^7(14H_{2n}u_n+y_nR_n^{(7)}),\\
 \mathfrak Z_{ss}&=\sum(-1)^na_n^7v_n,&
 \mathfrak Z_{xss}&=\sum(-1)^na_n^7(14H_{2n}v_n+2u_nR_n^{(7)}+y_nS_n^{(7)}).
\end{align*}
The linear combination with the coefficients
\begin{center}\small
\begin{tabular}{@{}ll@{\qquad}ll@{}}
\toprule
Derivative & Coefficient & Derivative & Coefficient\\\midrule
$\mathfrak Z_x$ & $560\pi^2$ & $\mathfrak Z$ & $5920\pi^2\log2-3224\pi^3$\\
$\mathfrak Z_{xs}$ & $-182\pi$ & $\mathfrak Z_s$ & $1136\pi^2-1924\pi\log2$\\
$\mathfrak Z_{xss}$ & $77$ & $\mathfrak Z_{ss}$ & $814\log2-360\pi$\\
\bottomrule
\end{tabular}
\end{center}
reproduces $F_n$ of \eqref{gff:c6:residualSums} term by term (an identity of
polynomials in $\pi$, $\log2$, $y_n$, $h_n$, $\delta_n$, $H_{2n}$,
$H_{2n}^{(2)}$, $H_{2n}^{(3)}$ and $\zeta(3)$), so it equals $M_7$. Let
$\mathfrak Z_N$ be $\mathfrak Z$ with the $n$-sum replaced by its Euler
approximant with the weights $w_{N,n}$. The weights act termwise, so the same
combination of the derivatives of $\mathfrak Z_N$ at $(0,0)$ is
$\widehat M_{7,N}$.

Euler's beta integral \cite{gff:C6dDLMFBeta} gives the sevenfold beta-moment
representation
\[
 \mathfrak Z(x,s)=\mathfrak c(x,s)\int_{[0,1]^7}
 \Bigl[\sum_{n\ge0}(-1)^n(n+t)\xi^n\Bigr]
 \prod_{i=1}^7u_i^{\mathsf a-1}(1-u_i)^{\mathsf d-1}\dd u_i,
 \qquad \xi=u_1\cdots u_7,
\]
with
\[
 \mathfrak c(x,s)=\left[\frac{\Gamma(1-x)\cos\pi(s+x)}
 {\Gamma(\frac12+x)\cos\pi x\,\Gamma(\frac12-2x)}\right]^7 .
\]
The alternating kernel is
$\sum_{n\ge0}(-1)^n(n+t)\xi^n=t/(1+\xi)-\xi/(1+\xi)^2$, and its exact remainder
after $N$ Euler coefficients is
\[
 2^{-N}\left[\frac{t(1-\xi)^N}{1+\xi}
 -\frac{N\xi(1-\xi)^{N-1}}{1+\xi}
 -\frac{\xi(1-\xi)^N}{(1+\xi)^2}\right].
\]
Its absolute value is at most $(|t|+2)2^{-N}$, since
$N\xi(1-\xi)^{N-1}\le1$. On the bidisk
$\operatorname{Re}\mathsf a,\operatorname{Re}\mathsf d\ge\frac{15}{32}$ and
$B(\operatorname{Re}\mathsf a,\operatorname{Re}\mathsf d)<16$. Elementary
integral bounds give $|\Gamma(1-x)|<3$; reflection and the same integral
bounds give $|1/\Gamma(\frac12+x)|<2$ and $|1/\Gamma(\frac12-2x)|<2$. Also
$|\cos\pi(s+x)|<2$ and $|1/\cos\pi x|<2$. Thus
\[
 |\mathfrak c(x,s)|\,B(\operatorname{Re}\mathsf a,\operatorname{Re}\mathsf d)^7
 <(48\cdot16)^7<2^{70},
\]
and $|\mathfrak Z-\mathfrak Z_N|<2^{72-N}$ uniformly on the bidisk. Cauchy's
coefficient bound yields
\[
 |\partial_x^\mu\partial_s^\nu(\mathfrak Z-\mathfrak Z_N)(0,0)|
 \le\mu!\,\nu!\,64^{\mu+\nu}2^{72-N},\qquad \mu=0,1,\quad \nu=0,1,2.
\]
Using $\pi<4$ and $0<\log2<1$, the sum of the Cauchy-weighted absolute
coefficients of the table is less than
$2^{20}+2^{22}+2^{26}+2^{19}+2^{21}+2^{25}<2^{27}$. This proves the bound
$2^{99-N}$. The proof differentiates a bounded analytic remainder; it makes no
assumption that an infinite asymptotic expansion converges.

\subsubsection{The primitive's fast recurrence}
The primitive evaluator uses integer and half-integer Frobenius branches in a
nine-component system with logarithmic degree at most two.
Section~\ref{gff:c6d:recurrence} states both exact $9\times9$ matrices and all seed
polynomials. The recurrence is
\begin{align}\label{gff:c6d:fastRec}
 [(n+\lambda)I-A_0]z_{n,j}={}&C_0z_{n-1,j}
 +[(n-1+\lambda)I+A_0]z_{n-2,j}\notag\\
 &+(j+1)(z_{n-2,j+1}-z_{n,j+1}),
 \quad\lambda\in\{0,1/2\}.
\end{align}
The terminal log channel is zero, and $j$ is processed in descending order. The
first components are combined with elementary logarithmic moments at
$\rho=3-2\sqrt2$. The complete-interval evaluation uses an explicit
last-two-state norm majorant; its numerical values are recomputed for each run.
The evaluation weights and the last-two-state bound are specified in
Section~\ref{gff:c6d:recurrence}. They are implemented in the package's interval
code; no previous primitive value is used to initialize the recurrence.

\subsubsection{The complete interval certificate}\label{gff:c6d:interval}
\paragraph{What is enclosed.}
The interval of width below $1.183\times10^{-145}$ quoted in
Section~\ref{gff:c6:analytic} is built by the package's interval stage for the
form of the reduced expression that this stage receives,
\begin{equation}\label{gff:c6d:intervalForm}
\begin{aligned}
 2187\,C_6^{\mathrm{rep}}(1)={}&E_{\mathrm{central}}+835584\Lambda_{3,2}
 +391680\mathsf H_5+13056\sqrt2(32\Cat+3\pi^2)\mathsf H_3\\
 &-393\pi^3W_c+96P_a-784\pi^{3/2}\mathsf L_7 .
\end{aligned}
\end{equation}
Here $W_c$ is \eqref{gff:c6d:Wc} and $P_a$ is \eqref{gff:c6:primitivePairing};
$\Lambda_{3,2}=\operatorname{Re}\operatorname{Li}_{3,2}(\frac{1+\ii}2,1-\ii)$,
$\mathsf H_3=\sum_{n\ge0}a_n2^{-n}(2n+1)^{-3}$ and
$\mathsf H_5=\sum_{n\ge0}2^{-n}(n+1)^{-5}=2\operatorname{Li}_5(\frac12)$. The
sum $\mathsf L_7$ is a fixed combination of the five alternating series
$\mathsf L(t)$ below, and the polynomial $E_{\mathrm{central}}$ contains,
besides elementary constants, the ten ordinary series of the table below and
the five alternating series $\mathsf Z(t)$. The coefficients of
$E_{\mathrm{central}}$ and the weights of $\mathsf L_7$ are those of the
package's input expression for this stage. Exact identities established in
later stages of the package carry \eqref{gff:c6d:intervalForm} into
\eqref{gff:c6:reduced}. Among them are closed-form evaluations of two ordinary
values in $E_{\mathrm{central}}$, the identity that turns the complete
seventh-power contribution into $24\pi^3M_7$, the conversion of
$\Lambda_{3,2}$, $\mathsf H_3$ and $\mathsf H_5$ into the values $D_3$, $D_5$
of $S_{\mathrm{pol}}$, and the inherited wheel normalizations
$B_-+B_+=2\pi^4\varpi_6$, $B_--B_+=4\pi^3\varpi_4$ that write the lower block
through $\varpi_4$ and $\varpi_6$. These identities are not reproduced in this
appendix. Reassembling the later forms from the same intervals of the
unchanged constants, with $M_7$ enclosed directly through \eqref{gff:c6d:M7tail},
again gives a width below $1.183\times10^{-145}$; the values $D_3$ and $D_5$
receive no separate interval.

\paragraph{Alternating series of the earlier form.}
For the five shifts $t\in\{\frac18,\frac14,\frac38,\frac58,\frac78\}$ put
$\mathsf a_t=t+\frac14$, $\mathsf b_t=t+\frac34$,
$g_t=\Gamma(\mathsf a_t)/\Gamma(\mathsf b_t)$ and
\[
 \mathsf Z(t)=\sum_{n\ge0}(-1)^n(n+t)
 \left[\frac{\Gamma(n+\mathsf a_t)}{\Gamma(n+\mathsf b_t)}\right]^7,\qquad
 \mathsf L(t)=\sum_{n\ge0}(-1)^n(n+t)
 \left[\frac{\Gamma(n+\mathsf a_t)}{\Gamma(n+\mathsf b_t)}\right]^7
 H_{2n}(2\mathsf a_t),
\]
with $H_m(a)=\sum_{j=0}^{m-1}(j+a)^{-1}$. Their finite Euler approximants
$\widehat E_N$ use the weights $w_{N,n}$ above; the equivalent form
$w_{N,n}=\sum_{k=n}^{N-1}\binom kn2^{-k-1}$ follows from finite binomial
summation, and the first form avoids a forward-difference table. Euler's beta
integral \cite{gff:C6dDLMFBeta} writes
$[(\mathsf a_t)_n/(\mathsf b_t)_n]^7=\mathbb E\,\mathsf S^n$, where
$\mathsf S$ is the product of seven independent
$\mathrm{Beta}(\mathsf a_t,\frac12)$ variables. For fixed $0\le\xi\le1$ the
alternating generating function of $(n+t)\xi^n$ is
$t/(1+\xi)-\xi/(1+\xi)^2$, and its finite Euler remainder is, with
$u=(1-\xi)/2$,
\[
 \mathsf r_N(\xi)=\frac{tu^N}{1+\xi}-\frac{\xi u^N}{(1+\xi)^2}
 -\frac{N\xi(1-\xi)^{N-1}}{2^N(1+\xi)},\qquad
 |\mathsf r_N(\xi)|\le(t+2)2^{-N},
\]
since $N\xi(1-\xi)^{N-1}\le1$ for $N\ge1$, endpoints included. Integrating
against the positive beta measure proves
\begin{equation}\label{gff:c6d:eulerZ}
 |\mathsf Z(t)-\widehat E_N[\mathsf Z(t)]|\le(t+2)g_t^7\,2^{-N}.
\end{equation}
For $\mathsf L(t)$ deform each beta density to
$\mathrm{Beta}(\mathsf a_t+x,\frac12-2x)$. Its normalized $n$-th moment is
$(\mathsf a_t+x)_n/(\mathsf b_t-x)_n$, and for seven copies
\[
 \left.\frac{d}{dx}\left[\frac{(\mathsf a_t+x)_n}{(\mathsf b_t-x)_n}\right]^7
 \right|_{x=0}
 =14\left[\frac{(\mathsf a_t)_n}{(\mathsf b_t)_n}\right]^7H_{2n}(2\mathsf a_t),
\]
by the duplication $H_{2n}(2a)=[H_n(a)+H_n(a+\frac12)]/2$. The logarithmic
derivative of one density is $\varphi(\xi)-\mathbb E\,\varphi(\xi)$ with
$\varphi(\xi)=\log\xi-2\log(1-\xi)$; its mean absolute value is at most
$2\,\mathbb E|\varphi|$. With the factor $1/14$, the seven copies give a signed
measure whose total variation is at most
$\mathbb E|\varphi|\le\psi(\mathsf b_t)-\psi(\mathsf a_t)
+2[\psi(\mathsf b_t)-\psi(\frac12)]$. The polygamma series and an integral
comparison give $\psi_1(a)\le a^{-2}+a^{-1}$ \cite{gff:DLMFPolygamma}. Since
$\mathsf a_t\ge\frac38$, $\mathsf b_t\le\frac{13}8$ and
$\psi_1(\frac12)=\pi^2/2<5$,
\[
 \psi(\mathsf b_t)-\psi(\mathsf a_t)\le\frac{44}9,\qquad
 2[\psi(\mathsf b_t)-\psi(\tfrac12)]\le\frac{45}4,\qquad
 \frac{44}9+\frac{45}4<17 .
\]
Integrating $\mathsf r_N$ against this signed measure proves
\begin{equation}\label{gff:c6d:eulerL}
 |\mathsf L(t)-\widehat E_N[\mathsf L(t)]|\le17(t+2)g_t^7\,2^{-N}.
\end{equation}
The beta parameters remain positive near $x=0$ and logarithmic beta moments
are integrable, which justifies the differentiation and the finite Euler
operation. The derivative serves only the error bound; the evaluator contains
none.

\paragraph{Ordinary positive series.}
The ten positive unit-argument series have rational term ratios
$R(n)=t_{n+1}/t_n=\prod_i(n+a_i)/\prod_i(n+b_i)$, $t_0=1$, with all shifts
positive, equal numerator and denominator degrees, and
$\alpha=\sum_ib_i-\sum_ia_i>1$; gamma-product asymptotics give
$t_n=O(n^{-\alpha})$ \cite{gff:DLMFStirling,gff:DLMF16}. For a finite $M\ge1$ choose
$q_M(n)=\sum_{j=0}^Mc_jn^{1-j}$ with $q_M(n)-R(n)q_M(n+1)=1+O(n^{-M-1})$. This
is a finite triangular rational calculation: the leading multiplier for $c_j$
is $\alpha+j-1>0$. To turn it into a bound, put $y=1/n$,
$\mathcal A_R(y)=\prod_i(1+a_iy)$, $\mathcal B_R(y)=\prod_i(1+b_iy)$ and
$\mathcal C_R(y)=\sum_{j=0}^Mc_jy^j$. The exact polynomial
\begin{equation}\label{gff:c6d:telescopeE}
 \mathcal E(y)=y(1+y)^{M-1}\mathcal B_R(y)-(1+y)^{M-1}\mathcal B_R(y)\mathcal C_R(y)
 +\mathcal A_R(y)\sum_{j=0}^Mc_jy^j(1+y)^{M-j}
\end{equation}
has zero coefficients through degree $M+1$, and the exact rational residual is
\[
 \mathsf r(n)=1-q_M(n)+R(n)q_M(n+1)
 =\frac{\mathcal E(y)}{y(1+y)^{M-1}\mathcal B_R(y)} .
\]
Since the denominator factors other than $y$ are at least $1$ for $y\ge0$,
for $n\ge N$
\[
 |\mathsf r(n)|\le\varepsilon_{N,M}=\sum_{k\ge M+2}|[y^k]\mathcal E(y)|\,N^{1-k},
\]
a finite rational number, not an omitted asymptotic term. The identity
$t_nR(n)=t_{n+1}$ telescopes, and $t_nq_M(n)\to0$ because $\alpha>1$. For
$T_N=\sum_{n\ge N}t_n$ this gives $|T_N-t_Nq_M(N)|\le\varepsilon_{N,M}T_N$,
so that, whenever $\varepsilon_{N,M}<1$ and $q_M(N)>0$ (both checked exactly),
\begin{equation}\label{gff:c6d:telescopeBound}
 \frac{t_Nq_M(N)}{1+\varepsilon_{N,M}}\le T_N\le
 \frac{t_Nq_M(N)}{1-\varepsilon_{N,M}} .
\end{equation}
Adding the finite rational partial sum gives an exact rational interval for
each series. The parameter ratios, including the factorial denominator, are
\begin{center}\small
\begin{tabular}{@{}llll@{}}
\toprule
Series & Numerator shifts $a_i$ & Denominator shifts $b_i$ & $\alpha$\\\midrule
${}_4F_3$ factor of $B_-$ & $\frac14,\frac14,\frac12,\frac12$ & $\frac34,\frac34,1,1$ & $2$\\
${}_4F_3$ factor of $B_+$ & $\frac34,\frac34,\frac12,\frac12$ & $\frac54,\frac54,1,1$ & $2$\\
${}_5F_4$ value & $\frac14,(\frac12)^{\times4}$ & $\frac54,1^{\times4}$ & $3$\\
${}_6F_5$ value & $1,1,(\frac54)^{\times4}$ & $2,(\frac74)^{\times4},1$ & $3$\\
$\Phi_-$ & $(\frac14)^{\times6}$ & $(\frac34)^{\times6}$ & $3$\\
$\Phi_0$ & $(\frac12)^{\times6}$ & $1^{\times6}$ & $3$\\
$\Phi_+$ & $(\frac34)^{\times6}$ & $(\frac54)^{\times6}$ & $3$\\
$Y_r$, $r=1,2,7$, normalized & $t+1,(t+\frac14)^{\times8}$ & $t,(t+\frac34)^{\times8}$ & $3$\\
\bottomrule
\end{tabular}
\end{center}
Here $t=r/8$, the $Y_r$ are normalized by their $n=0$ term, which is restored
afterwards,
$B_-=\Gamma(\frac14)^4\,{}_4F_3(\frac14,\frac14,\frac12,\frac12;\frac34,\frac34,1;1)$,
$B_+=64\pi^4\Gamma(\frac14)^{-4}\,{}_4F_3(\frac34,\frac34,\frac12,\frac12;\frac54,\frac54,1;1)$,
and the ${}_5F_4$ and ${}_6F_5$ values are the two ordinary values of
$E_{\mathrm{central}}$ named above. Gamma prefactors and the linear
combinations built from these series are evaluated with outward-rounded
interval operations.

\paragraph{Elementary constants, polylogarithms and the wheel.}
All certificate arithmetic uses dyadic endpoints, integers divided by $2^P$.
Integer floors and ceilings enclose every multiplication and division, square
roots use integer square roots, and additions and subtractions of endpoint
integers are exact; no binary floating-point special-value call enters the
construction. The constant $\pi$ is enclosed by Machin's arctangent identity
with alternating rational tails. Logarithms use range reduction and a finite
$\operatorname{artanh}$ series with a geometric tail; exponentials use range
reduction and a factorial Taylor remainder. For a positive gamma argument $a$,
the argument is shifted to $a+K$, the finite Stirling expansion of
$\log\Gamma$ is applied, and the logarithm of the exact product
$a(a+1)\cdots(a+K-1)$ is subtracted; the real-positive remainder has the sign
of, and at most the magnitude of, its first omitted Bernoulli term
\cite{gff:DLMFStirling}. Exponentiating the resulting interval encloses $\Gamma$.
This uses a finite remainder theorem, not convergence of Stirling's formal
series. Catalan's constant, $\zeta(3)$ and $\zeta(5)$ are enclosed by
Euler-transformed alternating Hausdorff-moment sums with nonnegative tails at
most $2^{-N}$; the zeta values include the exact Dirichlet-eta prefactor, and
$\zeta(2)=\pi^2/6$. The trigonometric eighth-angle constants are algebraic
expressions in $\sqrt2$. The half-argument series $\mathsf H_3$ and
$\mathsf H_5$ are positive, with tails bounded by twice their first omitted
term. For $\Lambda_{3,2}$, the substitution $n=m+r$ in the strict double sum
and $\frac{1+\ii}2(1-\ii)=1$ give
\begin{gather*}
 \Lambda_{3,2}=\operatorname{Re}\sum_{r\ge1}z_\Lambda^r\Lambda_r,\qquad
 z_\Lambda=\frac{1+\ii}2,\\
 \Lambda_r=\frac{\zeta(3)-H_r^{(3)}}{r^2}+\frac{3\zeta(2)-2H_r^{(2)}}{r^3}
 -\frac{3H_r}{r^4}=\sum_{m\ge1}\frac1{m^2(m+r)^3}.
\end{gather*}
Since $\Lambda_r\le\zeta(2)/(r+1)^3$, $|z_\Lambda|<\frac34$ and
$\zeta(2)/(1-|z_\Lambda|)<8$, the tail after $N$ terms is less than
$8(\frac34)^{N+1}/(N+2)^3$; the real and imaginary parts of the powers of
$z_\Lambda$ are generated by rational interval recurrences. The wheel $W_c$ is
recomputed with the certificate of Section~\ref{gff:c6d:wheelcert} at $N=256$,
$M=40$ (orders $40$ and $41$ for the fourth- and sixth-power corrections),
using the new $\pi$ interval; earlier wheel endpoints are read only for an
overlap check. The primitive uses \eqref{gff:c6d:primitiveTail}.

\paragraph{Assembly.}
Each run evaluates the complete polynomial and all six further weighted
contributions of \eqref{gff:c6d:intervalForm} before division by $2187$; the
source repair is included once, in the polynomial. The three recorded
configurations are
\begin{center}\small
\begin{tabular}{@{}lrrrr@{}}
\toprule
Configuration & Arithmetic bits & Ordinary $N/M$ & Primitive $N$ & Full interval width\\\midrule
smoke & 640 & 128/40 & 140 & $<4.673\times10^{-67}$\\
release & 1024 & 256/80 & 240 & $<1.986\times10^{-136}$\\
refined & 1280 & 320/96 & 280 & $<1.183\times10^{-145}$\\
\bottomrule
\end{tabular}
\end{center}
The three intervals are nested, as checked with their exact rational
endpoints. All runs use the same wheel cutoff, so its finite tail bound
dominates the last run. In the refined run the weighted widths are bounded by
$1.1826\times10^{-145}$ (wheel), $6.1870\times10^{-166}$ (polynomial),
$2.1146\times10^{-205}$ (primitive), $2.7549\times10^{-244}$
($\Lambda_{3,2}$), $3.7886\times10^{-280}$ ($\mathsf H_3$),
$7.1367\times10^{-284}$ ($\mathsf H_5$) and $4.9568\times10^{-323}$
($\mathsf L_7$). These are outward-rounded upper bounds, not differences
between floating-point approximations. The direct enclosure of $M_7$ by
\eqref{gff:c6d:M7tail} initializes $\pi$ by Machin's formula, $\log2$ by a
rational $\operatorname{artanh}$ series and $\zeta(3)$ by an Euler-transformed
eta series, never starts from an earlier value of $M_7$, and gives nested
intervals of width below $1.245\times10^{-60}$, $6.110\times10^{-151}$ and
$3.000\times10^{-241}$ at $N=300$, $600$ and $900$.

\subsubsection{Recorded validation of the retained computation}
The execution records of the package contain real-function checks, interval
assemblies, residual enclosures, exact kernel transformations and separate
Barnes comparisons. These records were produced for the consolidation of the
package's sources. They are reruns of the supplied implementations, not
independent rederivations of every analytic input. The final-method profile of
the package (Section~\ref{gff:cn:reproduction}) selects the tasks needed for the
method described here.

The real check used 105 working decimal digits for its 90-digit request. The
formula \eqref{gff:c1:c6nine} was then freshly assembled from the rerun values and
the proved conversions. Its absolute discrepancy from the refined interval
midpoint was about $4.2620\times10^{-93}$. The baseline assembly from complex
polylogarithms and series differed by about $3.9697\times10^{-94}$. Using the
independently rerun 40-digit seventh Barnes block gave a discrepancy in the
full coefficient of about $4.1439\times10^{-37}$. These are floating-point
assembly checks, not new interval bounds for the bivariate Barnes integrals.

The refined full interval constructor of Section~\ref{gff:c6d:interval} was also
rerun. It uses 1280-bit outward interval arithmetic and returns a width below
$1.183\times10^{-145}$ for the retained analytic expression in the form
\eqref{gff:c6d:intervalForm}. The subsequent assembly in the real basis uses those
intervals through exact identities but is not itself advertised as a new
nine-function interval implementation.

\begin{center}\small
\begin{tabular}{@{}ll@{}}
\toprule
Quantity & Rerun numerical value (shortened)\\\midrule
$L_3$ & $-2.0335765060720546009120689697005182\ldots$\\
$C_{212}$ & $0.032383901800869146760958099191565833\ldots$\\
$W_0$ & $78.57031213654711716068274888671985244\ldots$\\
$M_7$ & $-15332.55419477457527697894653409303157\ldots$\\
$P_a$ & $334926.04524805183788322200124007778661\ldots$\\
$\mathcal H_W$ & $82.198372839948704474444779671977025\ldots$\\
$\mathcal H_7$ & $-12.961075412927996857090472232198888\ldots$\\
$\mathcal H_*$ & $135.02348601116092499084881978216882\ldots$\\
\bottomrule
\end{tabular}
\end{center}
The $\mathcal H$ entries are obtained through the proved beta, Barnes and
series equivalents; they are not three independent fully two-dimensional
numerical integrations. Their exact normalization maps are
$W_0=\sqrt{89}\,\mathcal H_W/\pi^2$, \eqref{gff:c6d:seventhfinal} and
\eqref{gff:c6d:PaH}.

The three residual contributions, in $C_6$ rather than $r_6$ units, are
\begin{align}
 -\frac{393\pi^3W_0}{2187}&=-437.775914461168923635033332127009\ldots,\notag\\
 \frac{24\pi^3M_7}{2187}&=-5217.069054363479839305067485280987\ldots,\notag\\
 \frac{96P_a}{2187}&=14701.829146690890003104395116162537\ldots.
\end{align}
Their sum is $9046.984177866241240164294298754541\ldots$. They are three
normalization-consistent contributions, not nine additional unknown constants.

\subsection{Primitive recurrence data}\label{gff:c6d:recurrence}
This subsection supplies the matrices, seeds, evaluation weights and the
a~posteriori bound used in \eqref{gff:c6d:fastRec}. They are the computational
input recorded in the supplied recurrence and interval implementation.
The global constant $A$ still means $\pi+4\log2$; the matrix $A_0$ below is a
different object.

\begin{equation}
 A_0=\begin{pmatrix}
0 & 2 & 0 & 2 & 0 & 0 & 0 & 0 & 0\\0 & \frac{1}{2} & 1 & 0 & 2 & 0 & 0 & 0 & 0\\0 & 0 & 1 & 0 & 0 & 2 & 0 & 0 & 0\\0 & 0 & 0 & 0 & 2 & 0 & 1 & 0 & 0\\0 & 0 & 0 & 0 & \frac{1}{2} & 1 & 0 & 1 & 0\\0 & 0 & 0 & 0 & 0 & 1 & 0 & 0 & 1\\0 & 0 & 0 & 0 & 0 & 0 & 0 & 2 & 0\\0 & 0 & 0 & 0 & 0 & 0 & 0 & \frac{1}{2} & 1\\0 & 0 & 0 & 0 & 0 & 0 & 0 & 0 & 1
\end{pmatrix}.
\end{equation}
\begin{equation}
 C_0=\begin{pmatrix}
-1 & -4 & 0 & -4 & 0 & 0 & 0 & 0 & 0\\\frac{1}{4} & 2 & -2 & 0 & -4 & 0 & 0 & 0 & 0\\0 & \frac{1}{2} & 5 & 0 & 0 & -4 & 0 & 0 & 0\\\frac{1}{4} & 0 & 0 & 1 & -4 & 0 & -2 & 0 & 0\\0 & \frac{1}{4} & 0 & \frac{1}{4} & 4 & -2 & 0 & -2 & 0\\0 & 0 & \frac{1}{4} & 0 & \frac{1}{2} & 7 & 0 & 0 & -2\\0 & 0 & 0 & \frac{1}{2} & 0 & 0 & 3 & -4 & 0\\0 & 0 & 0 & 0 & \frac{1}{2} & 0 & \frac{1}{4} & 6 & -2\\0 & 0 & 0 & 0 & 0 & \frac{1}{2} & 0 & \frac{1}{2} & 9
\end{pmatrix}.
\end{equation}

There are two branches $\lambda=0,1/2$. At each order $n$, the nine-component
vectors $z_{n,j}^{(\lambda)}$ have logarithmic channels $j=0,1,2$; set
$z_{n,3}^{(\lambda)}=0$. For $n\ge2$, compute the channels in descending order
using \eqref{gff:c6d:fastRec}. The matrix on its left is upper triangular, so each
solve is an exact back substitution.

The initial vectors are specified by finite polynomials. Let $L$ be a formal
variable, $x=L-A$, and define the three-vectors
\begin{align*}
 \mathbf e&=(1,0,0),\\
 \mathbf a&=(15x^2-12\pi x-\pi^2,\;15x-6\pi,\;15),\\
 \mathbf b&=\left(\frac{15(x^2+6x)-12\pi(x+3)-\pi^2}{8},\right.\\[-3pt]
 &\hspace{15mm}\left.\frac{15(x^2+8x+6)-12\pi(x+4)-\pi^2}{16},
                    \frac{15(x+7)-6\pi}{8}\right),\\
 \mathbf p&=(\pi+6x,3,0),\qquad
 \mathbf q=\left(\frac{\pi+6x+18}{8},\frac{\pi+6x+24}{16},\frac38\right).
\end{align*}
Here $\otimes$ is the Kronecker product with the first vector's index
outermost. Form the nine-component polynomials
\begin{align*}
 U_0&=Q\,\mathbf a\otimes\mathbf e,\\
 U_1&=Q[\mathbf b\otimes\mathbf e+
                \mathbf a\otimes(3/4,1/8,0)]
        -\frac{16\pi^4}{Q}\mathbf e\otimes(1,1/2,1/4),\\
 V_0&=8\pi^2\mathbf p\otimes(1,1/4,0),\\
 V_1&=8\pi^2[\mathbf q\otimes(1,1/4,0)
                  +\mathbf p\otimes(1,1/2,1/16)].
\end{align*}
With $[L^j]$ denoting coefficient extraction from these finite polynomials,
the seeds are
\begin{align*}
 z_{0,j}^{(0)}&=[L^j]U_0,&
 z_{1,j}^{(0)}&=[L^j](4U_1-3U_0-2\partial_LU_0),\\
 z_{0,j}^{(1/2)}&=2[L^j]V_0,&
 z_{1,j}^{(1/2)}&=[L^j](8V_1-8V_0-4\partial_LV_0).
\end{align*}
The half-integer branch has no degree-two logarithmic channel. The displayed
derivatives are derivatives of explicit finite polynomials used to initialize
the algorithm, not special-function parameter derivatives.

Put $\rho=3-2\sqrt2$, and write $u_{n,j}$ and $v_{n,j}$ for the first
components of $z_{n,j}^{(0)}$ and $z_{n,j}^{(1/2)}$. Define
\begin{equation}
 J_j(s)=\rho^s\sum_{r=0}^j\binom jr
          \frac{(-1)^r r!\log^{j-r}\rho}{s^{r+1}}.
\end{equation}
These are the elementary moments $\int_0^\rho x^{s-1}\log^j x\dd x$. The
finite primitive evaluator, retaining $n<N$, is
\begin{equation}\label{gff:c6d:primitiveFinite}
 P_{a,N}=4\sum_{n=0}^{N-1}\left[
 \sum_{j=0}^2u_{n,j}J_j(n+1/2)+
 \sum_{j=0}^1v_{n,j}J_j(n+1)\right].
\end{equation}
The recurrence represents $P_a$ as the convergent limit of these
approximants. No numerical primitive value is used as a seed.

For completeness, the error bound can be constructed from the final two
states. The matrix norms are $\|A_0\|_\infty=4$ and $\|C_0\|_\infty=10$. On all
logarithmic channels, the operator $(Dz)_j=(j+1)z_{j+1}$ has norm $2$.
Rearranging the recurrence and applying a Neumann inverse bound gives
\[
 \|z_n^{(\lambda)}\|_\infty\le
 \frac{10\|z_{n-1}^{(\lambda)}\|_\infty+
       (n+\lambda+5)\|z_{n-2}^{(\lambda)}\|_\infty}{n+\lambda-6}.
\]
Let $U_n^{(\lambda)}=\rho^n\|z_n^{(\lambda)}\|_\infty$. Since $\rho<1/5$, the
two scalar coefficients are bounded by $\mathfrak a_n=2/(n-6)$ and
$\mathfrak b_n=(n+5)/(25(n-6))$. For $n\ge48$,
\[
 \frac{\mathfrak a_n}{1/4}+\frac{\mathfrak b_n}{(1/4)^2}
       =\frac{16n+280}{25(n-6)}\le1.
\]
With $K_\lambda=\max\{U_{N-1}^{(\lambda)},U_{N-2}^{(\lambda)}/4\}$, induction
gives $\sum_{n\ge N}U_n^{(\lambda)}\le K_\lambda/3$. Also $|\log\rho|<2$ and
$j\le2$ imply, for $s\ge N\ge48$,
\[
 \sum_{j=0}^2|J_j(s)|/\rho^s
 \le 7/N+5/N^2+2/N^3<8/N.
\]
The resulting bound is
\begin{equation}\label{gff:c6d:primitiveTail}
 |P_a-P_{a,N}|\le\frac{32}{3N}
          (\sqrt\rho\,K_0+\rho K_{1/2}),\qquad N\ge48.
\end{equation}
The code uses outward upper norms of the computed states in this inequality
and encloses every finite operation in \eqref{gff:c6d:primitiveFinite}. Thus both
the arithmetic uncertainty and the remaining infinite tail are included.

\begin{thebibliography}{99}
\bibitem{gff:BrownKreimerLocal} F.~Brown and D.~Kreimer,
\emph{Angles, scales and parametric renormalization},
\href{https://arxiv.org/abs/1112.1180}{arXiv:1112.1180} (2011).

\bibitem{gff:BaldufPanzerTropicalLocal} P.-H.~Balduf and E.~Panzer,
\emph{Renormalized tropical field theory},
\href{https://arxiv.org/abs/2512.21091}{arXiv:2512.21091} (2025).

\bibitem{gff:BytevKniehlMoch} V.~Bytev, B.~Kniehl and S.~Moch,
\emph{Derivatives of Horn-type hypergeometric functions with respect to their parameters},
\href{https://arxiv.org/abs/1712.07579}{arXiv:1712.07579},
including the general Gamma-ratio definition and generalized Lauricella arrays.

\bibitem{gff:DLMFGaussUnity} NIST Digital Library of Mathematical Functions,
\S15.4(ii), Gauss's unity-argument summation, equation (15.4.20).
\url{https://dlmf.nist.gov/15.4.E20}.

\bibitem{gff:MikhailovVolchanskiy} S. V. Mikhailov and N. Volchanskiy,
\emph{Two-loop kite master integral for a correlator of two composite vertices},
\href{https://arxiv.org/abs/1812.02164}{arXiv:1812.02164},
especially eq.~(3.30) and Appendix~A for the generalized Lauricella notation.

\bibitem{gff:NeretinComplex} Yu.~A.~Neretin,
\emph{On the Dotsenko--Fateev complex twin of the Selberg integral and its extensions},
\href{https://arxiv.org/abs/2212.09112}{arXiv:2212.09112}, especially eq.~(1.13).

\bibitem{gff:BGHS} D.~Benedetti, R.~Gurau, S.~Harribey and K.~Suzuki,
\emph{Long-range multi-scalar models at three loops},
J. Phys. A \textbf{53} (2020) 445008,
\href{https://arxiv.org/abs/2007.04603}{arXiv:2007.04603}.

\bibitem{gff:BLNV} C.~Behan, E.~Lauria, M.~Nocchi and P.~van Vliet,
\emph{Analytic and numerical bootstrap for the long-range Ising model},
JHEP \textbf{03} (2024) 136, \href{https://arxiv.org/abs/2311.02742v6}{arXiv:2311.02742v6}.

\bibitem{gff:GK} S.~Giombi and H.~Khanchandani,
\emph{$O(N)$ models with boundary interactions and their long range generalizations},
JHEP \textbf{08} (2020) 010, \href{https://arxiv.org/abs/1912.08169}{arXiv:1912.08169}.

\bibitem{gff:GhoshSuperSum} K.~Ghosh, M.~F.~Paulos, N.~Suchel and Z.~Zheng,
\emph{Super sum rules for long-range models},
\href{https://arxiv.org/abs/2603.22395}{arXiv:2603.22395} (2026).

\bibitem{gff:cnm:UGA} K.~Uzelac, Z.~Glumac and A.~Ani\v ci\'c,
\emph{Critical behavior of the long-range Ising chain from the largest-cluster probability distribution},
Phys.\ Rev.\ E \textbf{63} (2001) 037101,
\href{https://arxiv.org/abs/cond-mat/0011306}{arXiv:cond-mat/0011306}.

\bibitem{gff:BGH} D.~Benedetti, R.~Gurau and S.~Harribey,
\emph{Addendum: Long-range multi-scalar models at three loops},
\href{https://arxiv.org/abs/2411.00805}{arXiv:2411.00805}.

\bibitem{gff:Selberg} P.~J.~Forrester and S.~O.~Warnaar,
\emph{The importance of the Selberg integral},
Bull. Amer. Math. Soc. \textbf{45} (2008) 489--534,
\href{https://arxiv.org/abs/0710.3981}{arXiv:0710.3981}.

\bibitem{gff:Mishev} I.~D.~Mishev, \emph{Coxeter group actions on Saalsch\"utzian ${}_4F_3(1)$ series and
very-well-poised ${}_7F_6(1)$ series},
\href{https://arxiv.org/abs/1008.1011}{arXiv:1008.1011}.

\bibitem{gff:C6dBao} D.~Bao, \emph{Coxeter symmetry in hypergeometric functions and elliptic integral moments},
\href{https://arxiv.org/abs/2609.34016}{arXiv:2609.34016} (2026); its Proposition~5.4 is used.

\bibitem{gff:C6dBorweinStraub} J.~M.~Borwein and A.~Straub, \emph{Special values of generalized log-sine integrals},
in \emph{Proceedings of ISSAC 2011} (ACM, 2011), pp.~43--50,
\href{https://doi.org/10.1145/1993886.1993899}{doi:10.1145/1993886.1993899},
\href{https://arxiv.org/abs/1103.4298}{arXiv:1103.4298}.

\bibitem{gff:C6dARF} J.~A.~A.~Andrade, P.~N.~Rathie and R.~B.~A.~Farias,
\emph{Exact Bayesian computation using H-functions}, Comput.\ Appl.\ Math.\ \textbf{37} (2018) 2277--2293,
\href{https://doi.org/10.1007/s40314-017-0451-z}{doi:10.1007/s40314-017-0451-z}; its bivariate Definition~B is used.

\bibitem{gff:C6dWolframMeijerG} Wolfram Language documentation, \emph{MeijerG},
\url{https://reference.wolfram.com/language/ref/MeijerG.html}.

\bibitem{gff:C6dDLMFBeta} NIST Digital Library of Mathematical Functions, \S5.12,
beta function and Euler's beta integral. \url{https://dlmf.nist.gov/5.12}.

\bibitem{gff:C6dDLMFcosh} NIST Digital Library of Mathematical Functions, \S4.36,
infinite products and partial fractions of the hyperbolic functions.
\url{https://dlmf.nist.gov/4.36}.

\bibitem{gff:Herzog} F.~Herzog, \emph{Zimmermann's forest formula, infrared divergences and the QCD beta function},
\href{https://arxiv.org/abs/1711.06121}{arXiv:1711.06121}.

\bibitem{gff:DLMF16} NIST Digital Library of Mathematical Functions,
\href{https://dlmf.nist.gov/16.2}{\S16.2 (convergence)},
\href{https://dlmf.nist.gov/16.4}{\S16.4 (unity-argument identities)},
\href{https://dlmf.nist.gov/16.5}{\S16.5 (Euler and Barnes integrals)},
\href{https://dlmf.nist.gov/16.17}{\S16.17 (Meijer $G$ definition)},
\href{https://dlmf.nist.gov/16.21}{\S16.21 (differential equation)}.

\bibitem{gff:DLMF15} NIST Digital Library of Mathematical Functions,
\href{https://dlmf.nist.gov/15.8}{\S15.8 (Gauss transformations and connection formulas)}.

\bibitem{gff:DIS} S.~E.~Derkachov, A.~P.~Isaev and L.~A.~Shumilov,
\emph{Ladder and zig-zag Feynman diagrams, operator formalism and conformal triangles},
\href{https://arxiv.org/abs/2302.11238}{arXiv:2302.11238}.

\bibitem{gff:Wan} J.~G.~Wan, \emph{Moments of products of elliptic integrals},
\href{https://arxiv.org/abs/1101.1132}{arXiv:1101.1132}.

\bibitem{gff:DLMF5} NIST Digital Library of Mathematical Functions, \S5.5, gamma reflection, recurrence and duplication. \url{https://dlmf.nist.gov/5.5}.

\bibitem{gff:DLMFStirling} NIST Digital Library of Mathematical Functions, \S5.11, gamma asymptotics. \url{https://dlmf.nist.gov/5.11}.

\bibitem{gff:DLMFHurwitz} NIST Digital Library of Mathematical Functions, \S25.11, Hurwitz zeta function. \url{https://dlmf.nist.gov/25.11}.

\bibitem{gff:FlajoletGourdonDumas} P.~Flajolet, X.~Gourdon and P.~Dumas, \emph{Mellin transforms and asymptotics: harmonic sums}, Theor.\ Comput.\ Sci.\ \textbf{144} (1995) 3--58.

\bibitem{gff:FlajoletOdlyzko} P.~Flajolet and A.~Odlyzko, \emph{Singularity analysis of generating functions}, SIAM J.\ Discrete Math.\ \textbf{3} (1990) 216--240.


\bibitem{gff:DLMFElliptic} NIST Digital Library of Mathematical Functions, \S19.8, complete elliptic integrals, AGM and Landen transformations. \url{https://dlmf.nist.gov/19.8}.

\bibitem{gff:DLMFPolygamma} NIST Digital Library of Mathematical Functions,
\S5.15, polygamma functions: series, reflection, and multiplication formulas.
\url{https://dlmf.nist.gov/5.15}.

\bibitem{gff:DLMFDigamma} NIST Digital Library of Mathematical Functions,
\S5.4(ii), special values of the psi (digamma) function.
\url{https://dlmf.nist.gov/5.4#ii}.

\bibitem{gff:DLMFEllipticConnection} NIST Digital Library of Mathematical Functions,
\S19.7, connection formulas for complete elliptic integrals.
\url{https://dlmf.nist.gov/19.7}.

\bibitem{gff:DLMFTheta} NIST Digital Library of Mathematical Functions,
\S20.5 (infinite products of theta functions) and \S22.2 (Jacobian elliptic functions as
quotients of theta functions). \url{https://dlmf.nist.gov/20.5}, \url{https://dlmf.nist.gov/22.2}.

\bibitem{gff:TitchmarshCarlson} E.~C.~Titchmarsh, \emph{The Theory of Functions}, 2nd ed.,
Oxford University Press (1939), \S5.81 (Carlson's theorem).

\bibitem{gff:WhittakerWatsonBinet} E.~T.~Whittaker and G.~N.~Watson, \emph{A Course of Modern
Analysis}, 4th ed., Cambridge University Press (1927), \S12.31 (Binet's first expression for
$\log\Gamma(z)$).


\bibitem{gff:MacdonaldHyperI} I.~G.~Macdonald,
\emph{Hypergeometric Functions I},
\href{https://arxiv.org/abs/1309.4568}{arXiv:1309.4568} (2013),
originally circulated in 1987--1988; partitional Pochhammer symbols,
Jack normalizations, and matrix-argument hypergeometric series.

\bibitem{gff:DLMFMatrixHyper} NIST Digital Library of Mathematical Functions,
\S35.8, generalized hypergeometric functions of real matrix argument,
especially equations (35.8.1) and (35.8.5).
\url{https://dlmf.nist.gov/35.8}.

\bibitem{gff:RichardsZhengMatrix} D.~Richards and Q.~Zheng,
\emph{A Reflection Formula for the Gaussian Hypergeometric Function of
Matrix Argument}, J.~Math.\ Anal.\ Appl.\ \textbf{531} (2024), 127862,
\href{https://arxiv.org/abs/2002.05248}{arXiv:2002.05248};
especially Sections~3 and~4, distinguishing real and Hermitian matrix arguments.

\bibitem{gff:BorinskyD2Focus} M.~Borinsky,
\emph{Tropical Monte Carlo quadrature for Feynman integrals},
Annales de l'Institut Henri Poincar\'e D \textbf{10} (2023), 635--685,
arXiv:2008.12310.

\bibitem{gff:SchnetzD2Focus} O.~Schnetz,
\emph{Numbers and functions in quantum field theory},
Phys. Rev. D \textbf{97} (2018), 085018, arXiv:1606.08598,
Sections 3.1--3.3 (weighted completion and product identities).
% END FOCUSED D2 PRIMITIVE UPDATE

\bibitem{gff:C6KLT} N.~E.~J.~Bjerrum-Bohr, P.~H.~Damgaard, T.~S\o ndergaard and P.~Vanhove,
\emph{The Momentum Kernel of Gauge and Gravity Theories},
JHEP \textbf{01} (2011), 001, arXiv:1010.3933.

\bibitem{gff:Hill1975} B.~M.~Hill,
\emph{A simple general approach to inference about the tail of a distribution},
Ann.\ Statist.\ \textbf{3} (1975), 1163--1174.

\bibitem{gff:C6DuhrPorkert} C.~Duhr and F.~Porkert,
\emph{Feynman integrals in two dimensions and single-valued hypergeometric functions},
JHEP \textbf{02} (2024), 179, arXiv:2309.12772.

\bibitem{gff:PrecisionSobol} SciPy developers, \emph{scipy.stats.qmc.Sobol}: LMS scrambling, digital shifts, powers-of-two balance and bit resolution. Documentation accessed 14 September 2026. \url{https://docs.scipy.org/doc/scipy/reference/generated/scipy.stats.qmc.Sobol.html}.

\bibitem{gff:PrecisionQMC} SciPy developers, \emph{Quasi-Monte Carlo submodule}: independent randomizations, approximate error estimation and limits of convergence claims. Documentation accessed 14 September 2026. \url{https://docs.scipy.org/doc/scipy/reference/stats.qmc.html}.

\bibitem{gff:Arb} F. Johansson, \emph{Arb: efficient arbitrary-precision midpoint-radius interval arithmetic},
IEEE Trans.\ Comput.\ \textbf{66} (2017) 1281--1292, \href{https://doi.org/10.1109/TC.2017.2690633}{doi:10.1109/TC.2017.2690633}.

\bibitem{gff:MethodVegas}G. P. Lepage,
\emph{Adaptive Multidimensional Integration: VEGAS Enhanced},
J. Comput. Phys. \textbf{439} (2021) 110386,
\href{https://arxiv.org/abs/2009.05112}{arXiv:2009.05112}.

\bibitem{gff:mgs:BorowkaQMC} S.~Borowka, G.~Heinrich, S.~Jahn, S.~P.~Jones,
M.~Kerner and J.~Schlenk,
\emph{A GPU compatible quasi-Monte Carlo integrator interfaced to pySecDec},
Comput.\ Phys.\ Commun.\ \textbf{240} (2019) 120--137,
\href{https://arxiv.org/abs/1811.11720}{arXiv:1811.11720}.

\bibitem{gff:mgs:BorinskyGlobal} M.~Borinsky,
\emph{Tropicalized quantum field theory and global tropical sampling},
Ann.\ Henri Poincar\'e (2026),
\href{https://doi.org/10.1007/s00023-026-01695-y}{doi:10.1007/s00023-026-01695-y},
\href{https://arxiv.org/abs/2508.14263}{arXiv:2508.14263}.


\bibitem{gff:C7Giles} M.~B.~Giles, \emph{Multilevel Monte Carlo methods},
arXiv:1304.5472 (2013), \url{https://arxiv.org/abs/1304.5472}.

\end{thebibliography}
